% !BIB TS-program = biber
\documentclass[10pt]{amsart}
\usepackage[normalem]{ulem}

%For spacing:
\newcommand{\myspace}{\vspace*{0.3cm}\newline}
%BibLaTeX:
\usepackage{filecontents}

\begin{filecontents*}{mrnumber.dbx}
\DeclareDatamodelFields[type=field,datatype=verbatim]{mrnumber}
\DeclareDatamodelEntryfields{mrnumber}
\end{filecontents*}

\usepackage[sorting=ynt, backend=biber, datamodel=mrnumber, style=numeric-comp, isbn=false, sortcites, maxbibnames=99, firstinits=true]{biblatex}

\assignrefcontextentries*[sorting=ynt]{*}

\addbibresource{ABCPNT.bib}
%Space between references:

%The note field is often printed in bibliography styles. So if you insist on using this field, you need to avoid printing it twice. This can be done by clearing note in the begentry bibliography macro and then restoring it in finentry. Read here for more info: https://tex.stackexchange.com/q/91304/48840

\usepackage{amsmath,amssymb,amsthm,geometry,bm}
%\usepackage{amsrefs}
\usepackage[normalem]{ulem}
\usepackage[pdftex]{graphicx}
\pdfcompresslevel=9
\usepackage[colorlinks=true, pdfstartview=FitH, linkcolor=blue,
citecolor=blue, urlcolor=blue]{hyperref}
\geometry{left=1in, right=1in, top=.75in, bottom=.75in}
\newfont{\footsc}{cmcsc10 at 8truept}
\newfont{\footbf}{cmbx10 at 8truept}
\newfont{\footrm}{cmr10 at 10truept}
%\usepackage{alchemy}
%\usepackage{showkeys}

\newcommand{\C}{{\mathbb C}}
\newcommand{\N}{{\mathbb N}}
\newcommand{\Q}{{\mathbb Q}}
\newcommand{\R}{{\mathbb R}}
\newcommand{\Z}{{\mathbb Z}}
\newcommand{\ep}{\varepsilon}

\newcommand{\cA}{{\mathcal A}}
\newcommand{\cB}{{\mathcal B}}
\newcommand{\cN}{{\mathcal N}}
\newcommand{\cR}{{\mathcal R}}
\newcommand{\cS}{{\mathcal S}}

\newcommand{\fA}{{\mathfrak A}}

\renewcommand{\mod}[1]{{\ifmmode\text{\rm\ (mod~$#1$)}\else\discretionary{}{}{\hbox{ }}\rm(mod~$#1$)\fi}}
\newcommand{\li}{\mathop{\rm li}}
\newcommand{\Li}{\mathop{\rm Li}}
\newcommand{\cmd}[1]{{\tt$\backslash$#1}}
\newcommand{\cM}{{\mathcal M}}
%\renewcommand{\P}{{\mathcal P}}
\newcommand{\cP}{{\mathcal P}}
\newcommand{\Deltadot}{\mathord{\mathring\Delta}}
\newcommand{\Edot}{\mathord{\mathring E}}

\numberwithin{equation}{section}%{subsection} %sets equation numbers <chapter>.<section>.<index>
%\numberwithin{thm}{section} %sets equation numbers <chapter>.<section>.<index>
\newtheorem{theorem}{Theorem}[subsection]
\newtheorem{lemma}[theorem]{Lemma}
\newtheorem{corollary}[theorem]{Corollary}
\newtheorem{conjecture}[theorem]{Conjecture}
\newtheorem{definition}[theorem]{Definition}
\newtheorem{remark}[theorem]{Remark}
\newtheorem{example}[theorem]{Example}
\newtheorem{problem}[theorem]{Problem}

\usepackage{verbatim}


\title[Annotated bibliography for comparative prime number theory]{An annotated bibliography for \\ comparative prime number theory}
\author[Martin, Yang, Bahrini, Bajpai, Benl\.{i}, Downey, Li, Liang, Parvardi, Simpson, White, and Yip]{Greg Martin, Pu Justin Scarfy Yang, Aram Bahrini, Prajeet Bajpai, \\ K\"ubra Benl\.{i}, Jenna Downey,  Yuan Yuan Li, Xiaoxuan Liang,\\ Amir Parvardi, Reginald Simpson, Ethan Patrick White, and Chi Hoi Yip \\ Last modified: \today}
\address{University of British Columbia\\ Department of Mathematics \\ Room 121\\ 1984 Mathematics Road\\Vancouver, BC Canada V6T 1Z2}
\email{gerg@math.ubc.ca}
\subjclass[2010]{11N13 (11Y35)}


\begin{document}
\nocite{*}

\begin{abstract}
The goal of this annotated bibliography is to record every publication on the topic of comparative prime number theory together with a summary of its results. We use a unified system of notation for the quantities being studied and for the hypotheses under which results are obtained. We encourage feedback on this manuscript (see the end of Section~1 for details).
\end{abstract}
\maketitle

\large
\section{Introduction}\label{intro section}

Comparative prime number theory is the study of number-theoretic quantities, such as functions that count primes with particular properties, and how they compare to one another. It certainly includes (but is not limited to) ``prime number races'', which examine inequalities between the counting functions of primes in arithmetic progressions to the same modulus; indeed, Chebyshev observing the apparent preponderance of primes of the form $4k+3$ over those of the form $4k+1$ was the historical beginning of comparative prime number theory. Studying inequalities between two functions can be rephrased as studying the sign of their difference, and so the methods of comparative prime number theory also extend to studying the sign (and changes of sign) of other number-theoretic quantities that are less directly related to prime-counting functions.

The phrase ``comparative prime number theory'' goes back at least as far as the title of a long sequence of papers of Knapowski and Tur\'an, beginning with~\cite{MR0146156}. That paper begins with a list of several questions that can be interpreted as an attempt to define the scope of the field, as does the first paper~\cite{MR0162771} in a sequel series by the same authors. Other surveys of these topics include papers by Kaczorowski~\cite{MR1645811} and by Ford and Konyagin~\cite{MR1985941}, as well as an expository introduction to the field by Granville and the first author~\cite{MR2202918}.

This being said, there is no ironclad definition of what is and is not comparative prime number theory. Most quantities in this field have ``explicit formulas'' that express them as sums of oscillatory functions indexed by the zeros of $L$-functions of some type (including the Riemann zeta-function). As such, suitably normalized versions of these quantities are expected to have limiting (logarithmic) distribution functions, which are measures that record the frequencies with which the normalized quantities take values in various intervals in the limit (``continuous histograms'' of their values). In our view, the existence of such a limiting distribution is the main criterion for deciding whether a topic does or does not belong to the field of comparative prime number theory.

The purpose of this annotated bibliography is to provide a single exhaustive resource that lists every publication in the field of comparative prime number theory, and provides a summary of the results of each publication included. Like any human endeavour, the fulfillment of that goal will be imperfect. More specifically, we have aimed for completeness for all publications through 2019 (in an alternate timeline without the COVID-19 pandemic, this work might have been completed in 2020), as well as an incomplete list of publications from 2020--2023.

The publications in comparative prime number theory over the 170 years of its existence have understandably used a wide variety of notations for the same objects. Another purpose of this work is to propose a unified system of notation for referring to the functions and quantities that are the main objects of study in comparative prime number theory, as well as uniform terminology for the assumptions on zeros of $L$-functions that arise repeatedly when trying to prove theorems about these quantities. In particular, in our summaries of each publication, we have translated the results into this modern unified notation whenever possible, rather than preserving the notation used by the authors. In this respect, this work is more of a scientific resource than a historical document, although of course we hope it has some utility in the latter role (and we have included authors' exact words on a few occasions, particularly when problems or conjectures were first proposed).

Section~\ref{notation section} is therefore a long section presenting this system of notation for elementary functions, prime counting functions and other summatory functions of number-theoretic quantities, their error terms (both normalized and unnormalized), weighted and averaged versions of these quantities, analogues of these quantities over number fields and function fields, functions that count the number of sign changes of these quantities, and (natural and logarithmic) limiting densities and limiting distribution functions. Section~\ref{other notation section}, beginning on page~\pageref{other notation section}, describes objects and theorems that frequently arise in this field, such as Dirichlet characters and $L$-functions, Landau's theorem, explicit formulas, the power-sum method, $k$-functions, and various hypotheses on the zeros of $L$-functions.
Section~\ref{questions section}, beginning on page~\pageref{questions section}, enumerates the types of questions that comparative number theory studies about the quantities from Section~\ref{notation section}.
The annotated bibliography proper begins on page~\pageref{first bib}.

The origin of this manuscript was a literature survey project by the first two authors in 2012; since then, the many other authors have contributed significantly and have greatly expanded the extent of this bibliography and the accompanying material.

Any document of this length is bound to contain imperfections. We will gratefully receive any feedback about the contents; moreover, we plan on updating this document next summer to reflect all of the input we have received. In that vein, we invite all readers to send us their comments aimed at improving that next version of this annotated bibliography. As all errors or omissions are unintentional and all differences of opinions are presented in good faith, we ask for constructive feedback communicated in a positive tone.

Comments and suggestions can be submitted by emailing the first author at \href{mailto:gerg@math.ubc.ca}{gerg@math.ubc.ca} (be careful with the spelling), with the subject line ``ABCPNT'' (the acronym for the title of this manuscript), by June 30, 2024.


\section{Notation related to number theory and real analysis}\label{notation section}

We use $\N$ to denote the set of positive integers, and similarly $\Z$, $\R$, and~$\C$ to denote the sets of integers, real numbers, and complex numbers, respectively. We reserve the letter $p$ to denote prime numbers, and sums and products such as $\sum_p$ and $\prod_{p\mid q}$ are restricted to prime values of~$p$.

We use the following standard conventions regarding magnitudes of complex-valued functions~$f$ and~$g$, real-valued functions~$h$, and nonnegative real-valued functions~$r$ and~$s$ (of a complex or real argument~$z$):
\begin{itemize}
    \item $f(z) \ll s(z)$ (due to Vinogradov) means that there exists a constant $C>0$ such that $|f(z)| \le Cs(z)$ for all values of~$z$ under consideration;
    \item $O(s(z))$ (due to Bachmann) represents an unspecified function $f(z)$ with the property that $f(z)\ll s(z)$;
    \item $r(z) \asymp s(z)$ (due to Hardy) means that both $r(z) \ll s(z)$ and $s(z) \ll r(z)$ are true;
    \item $f(z) \sim g(z)$ (also due to Hardy) means that $\lim f(z)/g(z)=1$, where the location of the limit is taken from context (often as $z\to\infty$ through real numbers);
    \item $f(z) = o(s(z))$ (due to Landau) means that $\lim f(z)/s(z)=0$;
    \item $f(z) = \Omega(s(z))$ (due to Hardy and Littlewood) is the negation of $f(z) = o(s(z))$, or equivalently the statement $\limsup |f(z)|/s(z)>0$;
    \item $h(z) = \Omega_+(s(z))$ and $h(z) = \Omega_-(s(z))$ (due in this form to Landau) mean, respectively, that $\limsup h(z)/s(z)>0$ and $\liminf h(z)/s(z)<0$, either of which implies $h(z) = \Omega(s(z))$;
    \item $h(z) = \Omega_\pm(s(z))$ means that both $h(z) = \Omega_+(s(z))$ and $h(z) = \Omega_-(s(z))$ are true.
\end{itemize}


\subsection{Elementary functions}

As is standard in number theory, we use $\phi(n)$ to denote the Euler totient function, which is the number of reduced residue classes modulo~$n$. We use $\omega(n)$ to denote the number of distinct prime factors of~$n$ and $\Omega(n)$ to denote the number of prime factors of $n$ counted with multiplicity. We let $\mu(n)$ and $\Lambda(n)$ denote the M\"obius and von Mangoldt functions, respectively:
\[
\mu(n) = \begin{cases}
(-1)^{\omega(n)}, &\text{if $n$ is squarefree}, \\ 0, &\text{otherwise;}
\end{cases}
\qquad
\Lambda(n) = \begin{cases}
\log p, &\text{if $n=p^r$ for some $r\in\N$}, \\ 0, &\text{otherwise.} 
\end{cases}
\]

We use $(a,q)$ as a shorthand for $\gcd(a,q)$.
For any $(a,q)=1$, we define
\[
c_q(a) = \#\{b\mod q\colon b^2\equiv a\mod q\}
\]
to be the number of ``square roots'' of $a$ modulo~$q$. For brevity we write $c_q=c_q(1)$, which is also the number of real Dirichlet characters\mod q, or equivalently the index $\bigl[ (\Z/q\Z)^\times : \bigl( (\Z/q\Z)^\times \bigr)^2 \bigr]$; it turns out that $c_q=2^{\omega(q)+\eta}$ where $\eta\in\{-1,0,1\}$ depends upon the power of $2$ dividing~$n$. For $(a,q)=1$, it is the case that $c_q(a)$ equals $c_q$ if $a$ is a square\mod q and $0$ otherwise. (Many sources define $c(q,a)$ to be $c_q(a)-1$, which is more convenient for some purposes and less convenient for others.)

We define two closely related logarithmic integrals
\begin{align*}
\li(x) &= \lim_{\ep\to0+} \bigg( \int_0^{1-\ep} \frac{dt}{\log t} + \int_{1+\ep}^x \frac{dt}{\log t} \bigg) = \sum_{k=1}^K \frac{(k-1)!x}{(\log x)^k} + O_K\bigg( \frac x{(\log x)^{K+1}} \bigg) \\
\Li(x) &= \int_2^x \frac{dt}{\log t} = \li(x) - \li(2) \approx \li(x) - 1.04516378.
\end{align*}


\subsection{Prime counting functions}

We use the standard notation for the prime counting functions
\[
\begin{split}
\pi(x) &= \# \{ p \le x \} = \sum_{p\le x} 1 \\
\Pi(x) &= \sum_{n\le x} \frac{\Lambda(n)}{\log n} = \sum_{p^k \le x} \frac1k = \sum_{k=1}^\infty \frac{\pi(x^{1/k})}k \\
\theta(x) &= \sum_{p \le x} \log p \\
%\Pi^*(x) &= \sum_{p^k \le x} 1 = \sum_{k=1}^\infty \pi(x^{1/k}) \\
\psi(x) &= \sum_{n\le x} \Lambda(n) = \sum_{p^k \le x} \log p = \sum_{p\le x} \bigg\lfloor \frac{\log x}{\log p} \bigg\rfloor \log p = \sum_{k=1}^\infty \frac{\theta(x^{1/k})}k.
\end{split}
\]
We may replace the cutoff variable $x$ with any set $S$ of real numbers, so that for example
\[
\psi(S) = \sum_{n\in S} \Lambda(n)
\quad\text{and}\quad
\Pi\big( (0,x] \big) = \Pi(x)
\quad\text{and}\quad
\theta\big( (x,y] \big) = \theta(y)-\theta(x).
\]

All of these functions have analogues for prime powers restricted to arithmetic progressions:
\begin{align*}
\pi(x;q,a) &= \# \{ p \le x \colon p\equiv a\mod q \} = \sum_{\substack{p\le x \\ p\equiv a\mod q}} 1 \\
\Pi(x;q,a) &= \sum_{\substack{n\le x \\ n\equiv a\mod q}} \frac{\Lambda(n)}{\log n} = \sum_{\substack{p^k\le x \\ p^k\equiv a\mod q}} \frac1k \\
\theta(x;q,a) &= \sum_{\substack{p\le x \\ p\equiv a\mod q}} \log p \\
\psi(x;q,a) &= \sum_{\substack{n\le x \\ n\equiv a\mod q}} \Lambda(n) = \sum_{\substack{p^k\le x \\ p^k\equiv a\mod q}} \log p.
\end{align*}
These counting functions are interesting only in the case $(a,q)=1$, a restriction that we will usually not state explicitly.
Here too we may replace the first argument with a set, so that for example $\pi(S;q,a) = \#\{p\in S\colon p\equiv a\mod q\}$.

When the third argument is a set rather than an integer, the function counts prime powers that are congruent modulo~$q$ to any element of that set; for example, $\theta(x;q,\{1,2\}) = \theta(x;q,1) + \theta(x;q,2)$. In this context, $\cN$ and $\cR$ always refer to the quadratic residues and nonresidues, respectively, among the reduced residues modulo~$q$, so that for example
\begin{align*}
\pi(x;q,\cR) &= \# \{ p \le x \colon p \text{ is a quadratic residue} \mod q \} \\
\pi(x;q,\cN) &= \# \{ p \le x \colon p \text{ is a quadratic nonresidue} \mod q \}.
\end{align*}
Note that $\cR$ contains $\phi(q)/c_q$ residue classes\mod q and $\cN$ contains the other $\phi(q)(1-1/{c_q})$ residue classes. We use $\cA = \cN\cup\cR$ to refer to the set of all reduced residue classes.

When these prime counting functions for arithmetic progressions appear with four arguments instead of three, the function is the difference of the counts for the two indicated arithmetic progressions; for example, $\psi(x;q,a,1) = \psi(x;q,a) - \psi(x;q,1)$.
(Warning: some authors use $\Delta$ for differences of this type, but we give a different meaning to $\Delta$ below in Section~\ref{errors section}.)
When these final two arguments are sets, we make the convention that the two counting functions being subtracted are individually normalized by the number of distinct reduced residue classes in each set; for example,
\begin{align}
\theta(x;7,\{1,2\},\{3,4,5\}) &= \tfrac12\theta(x;7,\{1,2\}) - \tfrac13\theta(x;7,\{3,4,5\}) \notag \\ \Pi(x;q,\cN,\cR) &= \frac1{\phi(q)-c_q} \Pi(x;q,\cN) - \frac1{c_q} \Pi(x;q,\cR). \label{N minus R}
\end{align}
(This convention is consistent with the four-argument notation when the last two arguments are single integers, although these is some dissonance between this convention and the three-argument notation when the last argument is a set, since that function is not normalized in this way.)
There is no need for the notation to admit the possibility of two different moduli, since such a difference can always be written using residue classes of the least common multiple of the moduli: for example,
\[
\pi(x;8,1) - \pi(x;5,2) = 4\pi\big( x;40, \{1,9,17,33\}, \{7,17,27,37\} \big) = 3\pi\big( x;40, \{1,9,33\}, \{7,27,37\} \big).
\]

The residue class~$1\mod q$ is special in some ways, and it is thus helpful to define the notation
\[
\pi(x;q,1,\max) = \pi(x;q,1) - \max_{\substack{a\in(\Z/q\Z)^\times \\ a\not\equiv1\mod q}} \pi(x;q,a),
\quad
\pi(x;q,1,\min) = \pi(x;q,1) - \min_{\substack{a\in(\Z/q\Z)^\times \\ a\not\equiv1\mod q}} \pi(x;q,a),
\]
and similarly for other prime counting functions.


\subsection{Prime ideal classes} \label{number fields sec}

For any number field~$K$ (finite extension of~$\Q$), we say that $\alpha\in K$ is totally positive if~$\alpha$ maps to a positive real number under all real embeddings of~$K$ in~$\C$.
We call ideals~$\mathfrak{a}$ and~$\mathfrak{b}$ of a number field~$K$ congruent modulo another ideal $\mathfrak{f}$ if both~$\mathfrak{a}$ and~$\mathfrak{b}$ are coprime to $\mathfrak{f}$ and there exist totally positive algebraic integers $\alpha$ and $\beta$ in $K$ with $\alpha\equiv\beta \equiv 1\mod {\mathfrak{f}}$ such that $\alpha\mathfrak{a}=\beta \mathfrak{b}$. The equivalence classes of ideals modulo~$\mathfrak{f}$ form a group under ideal multiplication, with the principal ideal class~$\mathfrak{K}_0$ as its identity element. For a character~$\chi$ of this group, we abuse notation slightly by defining $\chi(\mathfrak{a})$ on ideals~$\mathfrak{a}$ directly: if $\mathfrak{a}$ is coprime to $\mathfrak{f}$ then we set $\chi(\mathfrak{a})=\chi([\mathfrak{a}])$, where $[\mathfrak{a}]$ is the ideal class\mod {\mathfrak{f}} containing $\mathfrak{a}$, and if $\mathfrak{a}$ is not coprime to $\mathfrak{f}$ then we set $\chi(\mathfrak{a})=0$. We can now define the Hecke--Landau zeta-function $\zeta(s,\chi)$ to be the Dirichlet series
\[
\zeta(s,\chi)=\sum_{\mathfrak{a}} \frac{\chi(\mathfrak{a})}{\mathfrak{Na}^s},
\]
Finally, for an ideal class $\mathfrak{K}$, we define prime ideal counting functions such as
\[
\pi(x,\mathfrak{K})= \sum_{\substack{\mathfrak{Np}\leq x\\\mathfrak{p}\in\mathfrak{K}\\\mathfrak{p}\,\textbf{prime ideal}}} 1, \qquad
\psi(x,\mathfrak{K})= \sum_{\substack{\mathfrak{Np}^m\leq x\\\mathfrak{p}^m\in\mathfrak{K}\\\mathfrak{p}\,\textbf{prime ideal}}} \log \mathfrak{Np}.
\]


\subsection{Error terms for prime counting functions} \label{errors section}

These prime counting functions have well-known main terms, and it is useful to have a standard notation to refer to the error terms obtained by subtracting these main terms, as well as normalized versions of such error terms.
We use $\Delta$ to denote error terms for the standard prime counting functions:
\[
\Delta^\psi(x) = \psi(x) - x, \quad \Delta^\theta(x) = \theta(x) - x, \quad \Delta^\Pi(x) = \Pi(x) - \li(x), \quad \Delta^\pi(x) = \pi(x) - \li(x).
\]
(In this document's article summaries, we will use the above normalizations even when an article subtracts a slightly different main term: we do not distinguish here between $\li(x)$ and $\Li(x)$ and $\sum_{2\le n\le x} 1/\log n$, for example.)
We also use $E$ for normalized versions of these error terms:
\[
E^\psi(x) = \frac{\Delta^\psi(x)}{\sqrt x}, \quad E^\theta(x) = \frac{\Delta^\theta(x)}{\sqrt x}, \quad E^\Pi(x) = \frac{\Delta^\Pi(x)}{\sqrt x/\log x}, \quad E^\pi(x) = \frac{\Delta^\pi(x)}{\sqrt x/\log x}.
\]
While there is not a formula for starting with a general function $f$ and determining the correct denominator to use when defining $E^f$, the normalization factor chosen is the one for which the resulting~$E$ function is expected to have a limiting logarithmic distribution.

It's not uncommon to integrate these error terms: for $f\in\{\pi,\Pi,\theta,\psi\}$ we define $\fA^f_0(x) = \Delta^f(x)$ and, for $m\ge1$,
\[
\fA^f_m(x) = \int_0^x \fA^f_{m-1}(t)\,dt.
\]
(Again, we ignore the fact that some articles might use a different lower endpoint for such integrals.) This operation has a predictable effect on summatory functions and explicit formulas: for example, $\fA_m^\psi(x) = \sum_{n\le x} (\Lambda(n)-1)(x-n)^m/m!$ has an explicit formula containing terms of the form $x^{\rho+m}/\rho(\rho+1)\cdots(\rho+m)$. For repeated integration of the absolute error, we also define $\fA^f_{|0|}(x) = |\fA^f(x)|$ and, for $m\ge1$,
\[
\fA^f_{|m|}(x) = \int_0^x \fA^f_{|m-1|}(t)\,dt.
\]
There are similar logarithmic integration operators: we define $A^f_0(x) = \Delta^f(x)$ and, for $m\ge1$,
\[
A^f_m(x) = \int_0^x A^f_{m-1}(t)\,\frac{dt}t.
\]
This operation also predictably affects summatory functions and explicit formulas: for example, $A_m^\psi(x) = \sum_{n\le x} (\Lambda(n)-1)(\log\frac xn)^m$ has an explicit formula containing terms of the form $x^{\rho}/\rho^{m+1}$. We also use the notation $A^f_{|m|}(x)$ for repeated logarithmic integration of the absolute error.

When we count primes in arithmetic progressions, the error terms $\Delta$ include a factor of $\phi(q)$ for simplicity: for example,
\[
\Delta^\psi(x;q,a) = \phi(q)\psi(x;q,a) - x \quad\text{and}\quad \Delta^\pi(x;q,a) = \phi(q)\pi(x;q,a) - \li(x).
\]
The normalized error terms $E$ are then derived from these $\Delta$ as before: for example,
\[
E^\psi(x;q,a) = \frac{\Delta^\psi(x;q,a)}{\sqrt x} \quad\text{and}\quad E^\pi(x;q,a) = \frac{\Delta^\pi(x;q,a)}{\sqrt x/\log x}.
\]
It is convenient at times to use a prime counting function itself as the main term, and such error terms are denoted by the symbol~$\Deltadot$: for example,
\[
\Deltadot^\psi(x;q,a) = \phi(q)\psi(x;q,a) - \psi(x) \quad\text{and}\quad \Deltadot^\pi(x;q,a) = \phi(q)\pi(x;q,a) - \pi(x).
\]
(Typically this modification results in the same explicit formula with the principal character removed.)
The corresponding normalized error terms are denoted by~$\Edot$: for example,
\[
\Edot^\psi(x;q,a) = \frac{\Deltadot^\psi(x;q,a)}{\sqrt x} \quad\text{and}\quad \Edot^\pi(x;q,a) = \frac{\Deltadot^\pi(x;q,a)}{\sqrt x/\log x}.
\]
We extend our convention regarding counting functions in arithmetic progressions: for example,
\[
\Delta^\psi(x;q,a,b) = \Delta^\psi(x;q,a) - \Delta^\psi(x;q,b) \quad\text{and}\quad E^\pi(x;q,a,b) = E^\pi(x;q,a) - E^\pi(x;q,b).
\]
Note that functions of the first type are almost redundant, since (for example) $\Delta^\psi(x;q,a,b) = \phi(q) \psi(x;q,a,b)$ exactly. (And recall that some authors use $\Delta$ to mean this difference function without the factor $\phi(q)$.) However, there will be situations where each notation is useful to us; furthermore, this new use of~$\Delta$ already follows from existing notational conventions.

%Should we define
%$$
%E(x;q,\cN,\cR) =
%\frac{\log x}{\sqrt x} \big( \pi(x;q,\cN) - (c_q-1)\pi(x;q,\cR) \big) ?
%$$

It can also also convenient to define this notation for the function $\psi(x,\chi) = \sum_{n\le x} \Lambda(n) \chi(n)$, for any Dirichlet character~$\chi$ (see Section~\ref{Dirichlet section}), in the following way:
\[
\Delta^\psi(x,\chi) = \psi(x,\chi) - \begin{cases}
x, &\text{if } \chi=\chi_0, \\
0, &\text{if } \chi\ne\chi_0,
\end{cases}
\qquad
E^\psi(x,\chi) = \frac{\Delta^\psi(x,\chi)}{\sqrt x}.
\]

All of the functions in this section so far have been real-valued (except for the last paragraph where the functions are potentially complex-valued); in the context of primes in arithmetic progressions, it is often helpful to consider vector-valued functions. We use subscripts to indicate the modulus and residue classes---for example,
\[
\pi_{q;a_1,\dots,a_r}(x) = \bigl( \pi(x;q,a_1), \dots, \pi(x;q,a_r) \bigr)
\quad\text{and}\quad
\Edot^\psi_{q;a_1,\dots,a_r}(x) = \bigl( \Edot^\psi(x;q,a_1), \dots, \Edot^\psi(x;q,a_r) \bigr).
\]


\subsection{Weighted versions of prime counting functions}

It is common to vary these prime counting functions by attaching a weight to each term in the sum, changing for example $\sum_{n\le x} \Lambda(n)$ to $\sum_{n\le x} \Lambda(n) g(n)$. We use the following consistent notation for the most common of these variants.

As is standard, the subscript~$0$, as in the example $\psi_0(x) = \frac12\big( \psi(x-) + \psi(x+) \big)$, represents a modification of a function's value at a jump discontinuity to equal the average of the left- and right-hand limits.

The subscript $r$ represents weighting by a reciprocal factor (often resulting in a ``Mertens sum''); for example,
\[
\pi_r(x) = \sum_{p\le x} \frac 1p, \quad  \theta_r(x) = \sum_{p\le x} \frac{\log p}{p}, \quad\text{and}\quad \psi_r(x;q,a) = \sum_{\substack{n\le x \\ n\equiv a\mod q}} \frac{\Lambda(n)}n.
\]
If we wish to modify one of these Mertens sums at its jump discontinuities as above, we concatenate the two subscripts: for example, $\pi_{r0}(x) = \frac12\big( \pi_r(x-) + \pi_r(x+) \big)$. Indeed all of our previous notational variants can apply to these sums as well---for example,
\[
\Delta^{\pi_r}(x) = \pi_r(x) - (\log\log x + B) \quad\text{and}\quad E^{\pi_r}(x) = \sqrt x\log x \cdot \Delta^{\pi_r}(x)
\]
for the appropriate constant~$B$.

The subscript $e$ represents weighting by an exponentially decaying function of $x$ rather than cutting off abruptly at $x$; for example,
\[
\pi_e(x) = \sum_p e^{-p/x} \quad\text{and}\quad \psi_e(x;q,a) = \sum_{\substack{n\ge 1 \\ n\equiv a\mod q}} \Lambda(n)e^{-n/x}.
\]
In terms of their asymptotics, these exponentially weighted sums usually act like their abrupt-cutoff versions; for example, $\pi_e(x)$ has a similar size to $\pi(x)$. However, their oscillations are typically damped, often resulting in rather different properties when comparing two such functions to each other (such as the exponentially weighted version having a bias for one sign while the unweighted version exhibiting oscillations of sign).

The subscript $l$ represents weighting by a certain exponential factor with a squared logarithm, scaled by a second parameter $r$: for example,
\[
\pi_l(x,r) = \sum_p e^{-\frac1r(\log \frac px)^2} \quad\text{and}\quad \psi_l(x,r;q,a) = \sum_{\substack{n\ge 1 \\ n\equiv a\mod q}} \Lambda(n)e^{-\frac1r(\log \frac nx)^2}.
\]
In asymptotic terms, this weighting is similar to restricting the range of summation to approximately $[e^{-\sqrt r} x, e^{\sqrt r} x]$; again, the oscillatory nature of the weighted sum can be rather different.

When the weight function is a Dirichlet character $\chi$ (see Section~\ref{Dirichlet section}), we follow the tradition of putting $\chi$ as an extra argument rather than a subscript; for example,
\[
\theta(x,\chi) = \sum_{p\le x} \chi(p)\log p.
\]


\subsection{Summatory functions}

Certain summatory functions of multiplicative functions have been analyzed using the techniques of comparative prime number theory. Two notable examples are the sums of the M\"obius and Liouville functions, which are denoted by
\[
M(x) = \sum_{n\le x} \mu(n) \quad\text{and}\quad L(x) = \sum_{n\le x} (-1)^{\Omega(n)},
\]
respectively. (The Liouville function is typically denoted by $\lambda(n) = (-1)^{\Omega(n)}$, but we avoid that notation herein to free the symbol $\lambda$ for other uses.) Two conjectures that motivated substantial work in comparative prime number theory are the ``Mertens conjecture'', the assertion that $|M(x)| < \sqrt x$, and the ``P\'olya problem'', the assertion that $L(x) \le 0$. (The latter assertion is often mistakenly named ``P\'olya's conjecture'', but P\'olya only posed and studied the problem rather than making a definitive conjecture and indeed probably found it unlikely to be true.) Both assertions have been disproved (in~\cite{MR783538} and~\cite{MR0104638}, respectively), although research continues into the distribution of these two functions. The weak Mertens conjecture, namely the assertion that $M(x) \ll \sqrt x$, is still unresolved, although it was shown in~\cite{MR0006202} to be incompatible with the pair of conjectures RH and LI (see Section~\ref{hypotheses section}).

The notational conventions from the previous sections are used for weighted versions of these summary functions as well; for example,
\[
M(x;q,a) = \sum_{\substack{n\le x \\ n\equiv a\mod q}} \mu(n) \quad\text{and}\quad L_r(x) = \sum_{n\le x} \frac{(-1)^{\Omega(n)}}{n};
\]
the conjecture that the latter is always nonnegative (often attributed to Tur\'an, though again he only studied the problem rather than asserting a conjecture) was also disproved in~\cite{MR0104638}.
We also define the notation $\Delta^M(x) = M(x)$ and $\Delta^L(x) = L(x)$ and $\Delta^{L_r}(x) = L_r(x)$; while unprofitable on their own, these definitions allow us to employ the notation for repeated averaging described in Section~\ref{errors section}, as well as the notation $E^M(x) = M(x)/\sqrt x$ and $E^L(x) = L(x)/\sqrt x$ and $E^{L_r}(x) = L_r(x)\sqrt x$.

We also introduce some standard notation for $k$-free numbers, which are numbers not divisible by the $k$th power of any prime, so that squarefree numbers are the case $k=2$. Let $Q_k(x)$ denote the number of $k$-free integers up to~$x$, and define $\Delta^{Q_k}(x)=Q_k(x)-x/\zeta(k)$. For integers $k\ge2$, the generalized M\"obius function $\mu_k(n)$ is defined to be $\mu_k(n)=(-1)^{\Omega(n)}$ if $n$ is $k$-free and $\mu_k(n)=0$ otherwise. Note that these functions interpolate between $\mu_2(n)=\mu(n)$ and $\lim_{k\to\infty} \mu_k(n)=(-1)^{\Omega(n)}$. These quantities are related by the identity $Q_k(x)=\sum_{n \leq x}\mu_k^2(n)$.

Another summatory function studied using techniques that
overlap with those of comparitive prime number theory is
\[
  D(x) = \sum_{n \le x} \tau(n) ,
  \text{ where } \tau(n) = \#\{ d \colon d \mid n \} = \sum_{d \mid n} 1 .
\]
It was first proven by Dirichlet (see~\cite{MR1576550} for a discussion of the history) that 
\[
  D(x) = x\log x - (2C_0-1)x + O(\sqrt x),
\]
where $C_0$ is Euler's constant. The study of the error
term $\Delta^D(x) = D(x) - x\log x - (2C_0-1)x$ is intertwined with comparative prime number theory, and one early result by Hardy \cite{MR1576550} demonstrates that the techniques of comparative prime theory are often applicable to the study of this error term.


\subsection{Counting sign changes}

We use the letter $W$ generally to denote the function that counts the number of sign changes of another function on an interval.
To be pedantic, if $h$ is a function from $(1,\infty)$ to $\R$, then we define
\begin{multline*}
W(h;T) = \max \big\{ n\ge0\colon \text{there exist } 1 < t_0 < t_1 < \cdots < t_n < T \\
\text{ with } h(t_{j-1}) h(t_j) < 0 \text{ for all } 1\le j\le n \big\}.
\end{multline*}
(One could quibble over whether taking the value $0$ counts as a sign change regardless of its neighboring values; the results in this subject tend not to require this loophole.)
We can demand large oscillations to go along with our sign changes by adding a function as an additional argument:
\begin{multline*}
W\big( h;T; S(t) \big) = \max \big\{ n\ge0\colon \text{there exist } 1 < t_0 < t_1 < \cdots < t_n < T \\
\text{ with } h(t_{j-1}) h(t_j) < 0 \text{ for all } 1\le j\le n \text{ and } |h(t_j)| > S(t_j) \text{ for all } 0\le j\le n \big\}.
\end{multline*}

Given functions $f$ and $g$ from $(1,\infty)$ to $\R$, we further define $W(f,g;T) = W(f-g;T)$ to be the counting function of sign changes of the difference $f(x)-g(x)$. Certain special cases of this notation deserve a shorthand notation: we define $W^\pi(T) = W(\pi,\li;T)$ and $W^\Pi(T) = W(\Pi,\li;T)$, and also $W^\theta(T) = W(\theta,x;T)$ and $W^\psi(T) = W(\psi,x;T)$ where $x$ denotes the identity function. (As before, we do not distinguish in our summaries between $\li(x)$ and $\Li(x)$ and $\sum_{2\le n\le x} 1/\log n$ in this context.) The bare notation $W(T)$ is a further shorthand for~$W^\pi(T)$.

In addition, given a positive integer $q$ and distinct reduced residues $a$ and $b\mod q$, we define $W^\psi_{q;a,b}(T) = W \big( \psi(x;q,a),\psi(x;q,b);T \big)$, and similarly with $\psi$ replaced by $\theta$, $\Pi$, or $\pi$; we further shorten $W^\pi_{q;a,b}(T)$ to $W_{q;a,b}(T)$. We may add a function as an additional argument as above to indicate large oscillations, as in $W_{q;a,b}(T;S(t))$ for example; similarly, we may replace single residue classes with sets of residue classes, as in $W_{q;\cN,\cR}(T)$.


\subsection{Densities}

The natural density of a set $\cS$ of positive real numbers is
\[
\mathfrak d(\cS) = \lim_{x\to\infty} \frac{\mathop{\rm meas}\big( \{ 0<t\le x\colon x\in \cS \} \big) }x = \lim_{x\to\infty} \frac1x \int\limits_{\substack{0<t<x \\ x\in \cS}} dt,
\]
where ``meas'' denotes Lebesgue measure on~$\R$.
On the other hand, the logarithmic density of a set $\cS \subset (1,\infty)$ is
\[
\delta(\cS) = \lim_{x\to\infty} \frac1{\log x} \int\limits_{\substack{1<t<x \\ x\in \cS}} \frac{dt}t.
\]
An easy change of variables shows that the logarithmic density of $\cS$ equals the natural density of the set $\log \cS = \{ \log t\colon t\in \cS \}$. Moreover, a partial summation argument shows that if the natural density $\mathfrak d(\cS)$ exists, then the logarithmic density $\delta(\cS)$ also exists and has the same value. However, there are sets whose natural density does not exist but whose logarithmic density does exist; for example, the union (over $k\in\N$) of the intervals $[10^{2k-1},10^{2k})$ has logarithmic density equal to~$\frac12$ but does not have a natural density.
%(There are discrete analogues of natural and logarithmic density for subsets of the positive integers; this last counterexample is analogous to the set of all positive integers with an even number of digits, which does not have a natural density but whose logarithmic density equals~$\frac12$.)

We will use many variants of this logarithmic density notation. If $f_1,\ldots,f_r$ are functions from $(1,\infty)$ to~$\R$, then we define the shorthand notation
\[
\delta(f_1,f_2,\ldots,f_r) = \delta\big( \{ x>1\colon f_1(x) > f_2(x) > \cdots > f_r(x) \} \big).
\]
For example, $\delta(\li,\pi)$ is the logarithmic density of the set of real numbers $x>1$ for which $\li(x) > \pi(x)$. Certain special cases of this notation get even further shorthand. For example, let $q$ be a positive integer, and let $a_1,\ldots,a_r$ be distinct reduced residues\mod q. Then we define
\[
\delta_{q;a_1,\ldots,a_r} = \delta\big( \pi(x;q,a_1), \ldots, \pi(x;q,a_r) \big) = \delta\big( \{ x>1\colon \pi(x;q,a_1) > \cdots > \pi(x;q,a_r) \} \big).
\]
We also define
\[
\delta_{q;\cN,\cR} = \delta\big( \pi(x;q,\cN), \pi(x;q,\cR) \big) = \delta\big( \{ x>1\colon \pi(x;q,\cN) > \pi(x;q,\cR) \} \big)
\]
and similarly for $\delta_{q;\cR,\cN}$ (these definitions are sensible when $q$ has primitive roots).

Finally, we define the upper and lower logarithmic densities of $\cS$ (which always exist) as
\[
\overline\delta(\cS) = \limsup_{x\to\infty} \frac1{\log x} \int\limits_{\substack{1<t<x \\ x\in \cS}} \frac{dt}t, \quad
\underline\delta(\cS) = \liminf_{x\to\infty} \frac1{\log x} \int\limits_{\substack{1<t<x \\ x\in \cS}} \frac{dt}t,
\]
so that $\delta(\cS)$ exists if and only if $\overline\delta(\cS) = \underline\delta(\cS)$. This notation propagates through our shorthand notations as well; for instance, $\underline\delta_{q;\cN,\cR} = \underline\delta\big( \{ x>1\colon \pi(x;q,\cN) > \pi(x;q,\cR) \} \big)$.


\subsection{Limiting distributions and density functions}

Given a function $h\colon [0,\infty)\to\R$, the limiting (or asymptotic) cumulative distribution function of $h$ is the nondecreasing function
\[
\lim_{T\to\infty} \frac {\mathop{\rm meas}\{ t\in[0,T]\colon h(t) \le a \}}T = \lim_{T\to\infty} \biggl( \frac1T \int\limits_{\substack{0\le t\le T \\ h(t) \le a}} \,dt \biggr)
\]
if the limit exists (except at jump discontinuities, of which there are only a countable number). More common in comparative prime number theory is the limiting logarithmic cumulative distribution function, with the analogous definition
\[
\kappa^h(\alpha) = \lim_{U\to\infty} \frac1{\log U} \biggl( \int\limits_{\substack{1\le u\le U \\ h(u) \le \alpha}} \,\frac{du}u \biggr),
\]
which equivalently is the cumulative distribution function of $h(e^t)$. There is a corresponding limiting logarithmic density $\mu^h$, which is the measure satisfying
\[
\mu^h\bigl( (\alpha,\beta] \bigr) = \int_\alpha^\beta \kappa^h(x) \,dx
\]
for any real numbers $\alpha<\beta$. It has the property that for any bounded continuous function $f(x)$,
\[
\lim_{U\to\infty} \frac1{\log U} \biggl( \int_1^U f\bigl( h(u) \bigr) \,\frac{du}u \biggr) = \int_{\R} f(x)\,d\mu^h(x),
\]
and the continuity assumption can be omitted if $\mu^h$ is absolutely continuous with respect to Lebesgue measure. These logarithmic densities are probability measures and thus can be viewed as the densities of random variables. Vector-valued functions have analogous logarithmic cumulative distribution functions and logarithmic densities on~$\R^r$.


\section{Notation related to complex analysis} \label{other notation section}

As is usual in analytic number theory, we often use $s=\sigma+it$ to denote a complex variable and its real and imaginary parts; its argument will be denoted by $\arg(s)$, so that $s=|s|e^{i\arg s}$. If $\rho$ is a nontrivial zero of a Dirichlet (or other) $L$-function, including the Riemann zeta-function, we write $\rho=\beta+i\gamma$ to refer to its real and imaginary parts.

\subsection{Dirichlet characters and Dirichlet $L$-functions} \label{Dirichlet section}

As usual, a Dirichlet character with modulus $q$ is a completely multiplicative function on $\Z$ with period $q$ whose support is the set of integers coprime to~$q$. We call characters real, complex, quadratic, (im)primitive, and induced with their standard meanings; the conductor of a character $\chi$ is the modulus of the primitive character $\chi^*$ that induces it.

We use $\chi_0$ to denote the principal character (the modulus being understood from context). When $D\ne1$ is a fundamental discriminant, we let $\chi_D$ denote the associated quadratic character, which is a primitive character of conductor~$|D|$ that is even if $D>0$ and odd if $D<0$. When $q$ is prime, we use the shorthand $\chi_{\pm q}$ to mean $\chi_q$ if $q\equiv1\mod4$ and $\chi_{-q}$ if $q\equiv3\mod4$. On the other hand, by $\chi_1$ we mean a hypothetical quadratic character with an exceptional zero $\beta_1$.

Every Dirichlet character gives rise to a Dirichlet $L$-function $L(s,\chi) = \sum_{n=1}^\infty \chi(n)n^{-s}$. Like the Riemann zeta-function (which is the special case $q=1$ and $\chi=\chi_0$), Dirichlet $L$-functions have infinitely many nontrivial zeros $\rho=\beta+i\gamma$ in the critical strip $0<\beta<1$. These zeros are counted by the function
\[
N(T,\chi) = \#\{ \rho\colon L(\rho,\chi) = 0,\, 0<\beta<1,\, |\gamma|\le T\}.
\]
Note the slight dissonance with the traditional notation
\[
N(T) = \#\{ \rho\colon \zeta(\rho) = 0,\, 0<\beta<1,\, 0\le \gamma\le T\}
\]
which counts only nontrivial zeros of $\zeta(s)$ in the upper half-plane: this suffices for $\zeta(s)$ due to the Schwarz reflection principle, but Dirichlet $L$-functions do not all possess that symmetry.

Sums over zeros of Dirichlet $L$-functions (of the type that arise in explicit formulas, for example) often do not converge absolutely, and therefore we adopt the standing convention that sums over nontrivial zeros are limits of their symmetric truncations:
\[
\sum_\rho f(\rho) = \lim_{T\to\infty} \sum_{\substack{L(\rho,\chi)=0 \\ 0<\beta<1 \\ |\gamma|\le T}} f(\rho).
\]


\subsection{Landau's theorem}

For a real-valued function $A(x)$, define
\[
g(s)=\int_1^\infty \frac {A(x)}{x^s} \,dx 
\]
Typically there will be a real number $\sigma_0$ such that this integral converges when $\sigma>\sigma_0$ and diverges when $\sigma<\sigma_0$. 
Landau proved that if $A(x)$ is eventually positive or eventually negative, then $g(s)$ has a singularity at $s=\sigma_0$ (that is, $g(s)$ must have a rightmost singularity on the real axis).

The contrapositive of this theorem is a useful tool in comparative prime number theory:
Suppose that $g(s)$ has no singularities on the subray $\{ \sigma\in\R\colon \sigma>\sigma_1\}$ of the real axis (that is, $g(s)$ is analytic on a neighborhood of that ray), but $g(s)$ is not analytic in the half-plane $\{s\in\C\colon \sigma>\sigma_1\}$. Then $A(x)$ has arbitrarily large sign changes.


\subsection{Explicit formulas}\label{formulas section}

As mentioned earlier, one of the defining characteristics of comparative prime number theory is the presence of an ``explicit formula''. There is no precise definition of that term, but typically they contain a sum over the (nontrivial) zeros of some $L$-function. The prototypical example is the explicit formula
\[
\psi_0(x) = x - \sum_{\rho} \frac{x^\rho}\rho - \log2\pi - \frac12\log\bigg( 1 - \frac1{x^2} \bigg)
\]
for the Chebyshev function $\psi(x)$ modified at its jump discontinuities; the fact that this is an exact equality for all $x>1$ is one of the most beautiful statements in analytic number theory.

Explicit formulas for prime-counting functions yield explicit formulas for their normalized error terms; for example, assuming the generalized Riemann hypothesis,
\[
E^\theta(x;q,a,b) = c_q(b)-c_q(a) - \sum_{\chi\mod q} \big( \overline\chi(a) - \overline\chi(b) \big) \sum_{\substack{\gamma\in\R \\ L(1/2+i\gamma,\chi) = 0}} \frac{x^{i\gamma}}{\frac12+i\gamma} + O_q(x^{-1/6}).
\]
This formula is helpful for studying when $E^\theta(x;q,a,b) > 0$, or equivalently when $\theta(x;q,a) > \theta(x;q,b)$. Note that each summand in the inner sum oscillates around a circle of fixed radius (one that decreases as~$\gamma$ increases); while this inner sum is not literally bounded, it is bounded on average over~$x$ and possesses a limiting logarithmic distribution. Therefore $E^\theta(x;q,a,b)$ has a limiting logarithmic distribution with mean $c_q(b)-c_q(a)$, the sign of which depends on whether~$a$ and~$b$ are quadratic residues or nonresidues modulo~$q$.

\subsection{The power-sum method}

A great deal of early progress in comparative prime number theory, particularly the unconditional results, relied on the study of linear combinations of powers of complex numbers, namely sums of the shape
\[
s_{v}=\sum_{j=1}^{n} b_{j} z_{j}^{v}.
\]
Lower bounds for such sums were systematically developed by Tur\'an and S\'os. While there are many variants of these lower bounds that have been obtained, they can be grouped into two main categories.

The ``first main theorem'' is a type of result that applies when the $z_j$ are large. For example,
suppose that $z_1,\ldots,z_n$ are distinct complex numbers with $|z_n|\ge1$ for all~$n$. For any nonnegative integer $m$, there exists an integer $m+1\le v\le m+n$ such that
\begin{equation*}
|s_{v}| \ge \bigg(\frac{n}{A(m+n)}\bigg)^{n} |s_0|,
\end{equation*}
where $A$ is an absolute constant. 

The ``second main theorem'' is a type of result that applies when the $z_j$ are small. For example,
suppose that $z_1,\ldots,z_n$ are distinct complex numbers with $1\ge |z_1|\ge\cdots\ge|z_n|$. For any nonnegative integer $m$, there exists an integer $m+1\le v\le m+n$ such that
\begin{equation*}
|s_{v}| \ge \bigg(\frac{n}{B(m+n)}\bigg)^{n} \min _{1 \leq j \leq n} \bigg| \sum_{n=1}^{j} b_{n} \bigg|,
\end{equation*}
where $B$ is an absolute constant.

Instead of restricting the candidate exponents~$v$ to an interval of exactly $n$ consecutive integers, we may allow candidates from a longer range of exponents.
For example, in the ``second main theorem'' (so that $1\ge |z_1|\ge\cdots\ge|z_n|$), let $m\ge N\ge n$; then there exists an integer $m+1\le v\le m+N$ such that
\begin{equation*}
|s_{v}| \ge \bigg(\frac{N}{Bm}\bigg)^{N} \min_{1 \leq j \leq n} \bigg| \sum_{n=1}^{j} b_{n} \bigg|.
\end{equation*}
For the ``second main theorem'', one can also obtain better conclusions by adding an ``argument restriction'', that is, the assumption that each $|\arg z_j|\ge\ep$ for some fixed $\ep>0$. Stronger results can also be obtained by assuming that each $b_j$ is a nonnegative real number, and strengthened further by restricting to the special case $b_1=\cdots=b_n=1$.

Note that these results show that some $s_v$ is large in modulus but gives no information about its argument. Tur\'an (somewhat unhelpfully) calls these results ``two-sided'' theorems. There exist analogous results where the lower bound applies not just to $|s_v|$ but to $\Re s_v$ or $-\Re s_v$; Tur\'an calls such results ``one-sided'' theorems.


\subsection{$k$-functions}  \label{k-functions section}

A great deal of the work of Kaczorowski involves certain functions called $k$-functions, which are superficially similar to sums that appear in explicit formulas for $\psi(x,\chi)$. For $\Im z>0$, define
\begin{align*}
k(z, \chi) = \sum_{\gamma>0} e^{\rho z} \quad \text{and} \quad K(z, \chi) = \sum_{\gamma>0} \frac{e^{\rho z}}{\rho},
\end{align*}
where the sums are over zeros of $L(s,\chi)$ in the upper half-plane.

These functions can be regarded as having their domain equal to~$\cM$, the Riemann surface for $\log z$; every point on the surface can be uniquely written as $re^{ia}$ where $r>0$ and $a\in\R$. Let $z^c$ denote the natural extension of complex conjugation to~$\cM$, namely $(re^{ia})^c = re^{-ia}$; also let $z^*$ denote an extension of multiplication by $-1$ to~$\cM$, namely $(re^{ia})^* = re^{i(a - \pi)}$.

Certain functions appear frequently in connection to $k$-functions: define
\[
D(z, \chi) = - \sum_{\substack{\beta>0 \\ L(\beta, \chi)=0}} e^{\beta z} + \frac1{e^{2z}-1} \begin{cases}
e^{3z} + e^{2z}-1, &\text{if } \chi=\chi_0, \\
e^z, &\text{if } \chi\ne\chi_0 \text{ and } \chi(-1)=1, \\
e^{2z}, &\text{if } \chi(-1) = -1.
\end{cases}
\]
Further define
\begin{align*}
F(x,\chi) = \lim_{y \to 0^+} \bigg( K(x+iy, \chi) + \overline{K(x+iy, \overline{\chi})} \bigg)
\end{align*}
and
\[
R_1(x) = \frac12\log(1-e^{-2x}), \quad R_{-1}(x) = \frac12 \log \frac{e^x-1}{e^x+1}.
\]
Certain constants also appear frequently: define
\[
B(\chi) = \sum_{\substack{\beta>0 \\ L(\beta, \chi)=0}}  \frac{1}{\beta} - \frac{C_0}{2} -\frac{1}{2} \log \frac{\pi}{q} + F(0, \chi) - \begin{cases}
1, &\text{if } \chi=\chi_0, \\
0, &\text{if } \chi\ne\chi_0 \text{ and } \chi(-1)=1, \\
\log2, &\text{if } \chi(-1) = -1
\end{cases}
\]
(note that $B(\chi)$ is not the same as a constant of the same name related to the Hadamard product expansion of $L(s,\chi)$) and $C(\chi) = B(\chi) + C_0 + \log\frac{2\pi}q$.


\subsection{Hypotheses on zeros}\label{hypotheses section}

It is extremely difficult to obtain unconditional results in comparative prime number theory, particularly where limiting logarithmic distributions and densities are concerned. Certain assumptions on the zeros of Dirichlet $L$-functions therefore arise repeatedly in this subject. The most famous of these is the generalized Riemann hypothesis (GRH), sometimes called the Riemann--Piltz conjecture, which asserts that all nontrivial zeros of all Dirichlet $L$-functions have real part equal to~$\frac12$. We use $\sigma_0$-GRH to denote the weaker (but still currently inaccessible) assertion that $L(\sigma+it,\chi)$ does not vanish when $\sigma>\sigma_0$, so that $1$-GRH is trivial and $\frac12$-GRH is the same as the full GRH.

Given a nonempty set $X$ of Dirichlet $L$-functions (or, abusing notation slightly, Dirichlet characters), we let $\Theta(X)$ denote the supremum of the real parts of their zeros, that is, the smallest real number such that $\Theta(X)$-GRH holds. We use the abbreviation $\Theta(q)$ when~$X$ is the set of all Dirichlet characters modulo~$q$, as well as $\Theta(\chi)$ when~$X$ consists of the single Dirichlet character~$\chi$. The assertion that some Dirichlet $L$-function in~$X$ has a zero with real part exactly equal to $\Theta(X)$ is abbreviated SA for ``supremum attained'' (and sometimes referred to as ``Ingham's condition''). We may write $SA(X)$ to emphasize that we are considering a specific set of Dirichlet $L$-functions, but the set is often inferred from context (this remark applies similarly to the remainder of the notation in this section). We note that GRH implies SA but that $\Theta(X)=1$ is inconsistent with~SA.

Regarding the vertical distributions of the zeros, we use HC to denote the ``Haselgrove condition'' that no Dirichlet character (in the set under discussion) vanishes on the segment $0<\sigma<1$ of the real axis. Such a real zero would create a non-oscillatory term in relevant explicit formulas, one that could result in an unexpected source of bias. By continuity, HC implies that there exists a positive constant~$E_k$ such that these $L(s,\chi)$ are nonzero on the rectangle $\{ 0<\sigma<1,\, |t| \le E_k \}$; we write HC$(E_k)$ if we need to refer to this parameter.

The notation GRH$(H)$ (sometimes called the ``finite Riemann--Piltz'' conjecture) denotes the generalized Riemann hypothesis ``up to height~$H$'', namely the statement that if $\rho$ is a nontrivial zero of $L(s,\chi)$ with $|\gamma|\le H$, then $\beta=\frac12$. 
Note that HC$(E_k)$ implies GRH$(H)$ if $E_k \ge H$; on the other hand, GRH$(H)$ gives no constraint at all upon zeros on the critical line. We therefore use the notation GRH$(H,E_k)$ to denote the combination of GRH$(H)$ and HC$(E_k)$, the latter of which constrains only the zeros on the critical line when $E_k < H$.
Note also that GRH$(0)$ is almost the same as HC, except that GRH$(0)$ allows for the possibility of a zero at $s=\frac12$.

The arithmetic nature of the imaginary parts (ordinates) of zeros of $L(s,\chi)$ is also significant in comparative prime number theory. We write LI (sometimes called GSH for the ``grand simplicity hypothesis'') to denote the ``linear independence'' assertion that the multiset of nonnegative ordinates of zeros of the relevant Dirichlet $L$-functions is linearly independent over the rational numbers. In particular, LI implies that all zeros are simple and that $L(\frac12,\chi)\ne0$. We use LI$(\sigma)$ to denote the corresponding linear independence conjecture restricted to the zeros with real parts greater than or equal to~$\sigma$.

For the Riemann zeta-function, the Riemann hypothesis (RH) is the assertion that all nontrivial zeros of $\zeta(s)$ have real part equal to~$\frac12$. Almost all of the other notation above would be used in the same form when referring to $\zeta(s)$, although $\Theta(\{\zeta(s)\})$ is abbreviated simply to~$\Theta$. These same abbreviations are also used for analogous hypotheses on zeros of other $L$-functions that should be clear from context.


\section{Types of questions} \label{questions section}

Given two functions $f,g\colon (1,\infty)\to\R$ that are asymptotic to each other, such as $\pi(x)$ and $\li(x)$ or $\pi(x;4,1)$ and $\pi(x;4,3)$, the questions that comparative prime number theory tends to ask about the pair of functions are:
\begin{enumerate}
\item Are there arbitrarily large values of $x$ for which $f(x) > g(x)$, and arbitrarily large values of $x$ for which $g(x) < f(x)$? In other words, does the difference $f(x)-g(x)$ change signs infinitely often? (These are not quite mathematically identical because of the possibility of plentiful or carefully arranged ties $f(x)=g(x)$, so implicit in this question is asking whether such ties are rare.) The other alternative is that one of the functions exceeds the other for all sufficiently large~$x$.
\item How large and positive can the difference $f(x)-g(x)$ get? How large and negative can it get?
\item More generally, what is the distribution of values of $f(x)-g(x)$? Is it possible that some suitably normalized version of this difference, such as $(f(x)-g(x))/\sqrt x$, actually has a limiting distribution or a limiting logarithmic distribution?
\item How often does the difference $f(x)-g(x)$ change sign? How many sign changes are there in $(1,X)$ as a function of $X$? How close can we take $Y=Y(X)$ to $X$ to ensure that there is always a sign change in $[X,Y]$?
\item What is the natural density of the set of real numbers $x>1$ for which $f(x)>g(x)$? What is its logarithmic density~$\delta(f,g)$? (Typically we believe that the natural densities of such sets do not exist in prime number races, but that their logarithmic densities do exist.)
\item Given a family of races, such as $\pi(x;q,\cN)$ versus $\pi(x;q,\cR)$: how do answers to the above questions, such as $\delta_{q;\cN,\cR}$, depend upon the member of the family ($q$ in this case)? Do the distributions of the members of the family tend to some limit, such as a normal distribution?
\end{enumerate}
Some of the above questions have analogues for several functions $f_1,\ldots,f_r\colon(1,\infty)\to\R$ considered together:
\begin{enumerate}\setcounter{enumi}{6}
\item Are there arbitrarily large values of $x$ for which $f_1(x) > \cdots > f_r(x)$? Does this remain true no matter how we permute the~$f_j$?
\item More generally, what is the distribution of values of the vector $\big( f_1(x), \ldots, f_r(x) \big) \in \R^r$? Is it possible that some suitably normalized version of this difference actually has a limiting distribution or a limiting logarithmic distribution?
\item What is the natural density of the set of real numbers $x>1$ for which $f_1(x) > \cdots > f_r(x)$? What is its logarithmic density $\delta(f_1,\ldots,f_r)$? (As before, we believe that the natural densities of such sets do not exist in prime number races, but that their logarithmic densities do exist.)
\item Given a family of such $r$-way races, how do answers to the above questions depend upon the member of the family? Do the distributions of the members of the family tend to some limit, such as a multivariate normal distribution?
\end{enumerate}

The articles~\cite{MR0146156} and~\cite{MR0162771} by Knapowski and Tur\'an present organized schema for problems in comparative prime number theory, as do surveys of these topics by Kaczorowski~\cite{MR1645811} and by Ford and Konyagin~\cite{MR1985941}, although several of the questions listed above had not yet been investigated sufficiently deeply to make some of their lists.


\section*{Acknowledgments}

We gratefully thank Devang Agarwal, Michael Coons,  Alia Hamieh, Elchin Hasanalizade, Farid Jokar, and \LaTeX\ Stack Exchange user ``moewe'' for their contributions to this manuscript. Many authors' research was supported by the Natural Science and Engineering Research Council of Canada.

%%%%%%%%%%%%%%%%%

\section*{Chronological bibliography}

The annotated bibliography begins here, with all of the sources cited and summarized listed in chronological order; items in this chronological list are labeled by their number alone, such as~[123]. Following the annotated bibliography is a second list, in alphabetical order by author, of the same set of sources but without annotations; items in this alphabetical list have been given labels that are numbers following the letter ``A'' (for ``alphabetical''), such as~[A45], to distinguish them from the labels in the main list.

Our goal has been to describe the results using a single system of notation, both to avoid the need to define notation in individual annotations and to propose a unified notation for current and future practitioners of comparative prime number theory. Any notation in a summary that is not defined there can be found or deduced from the detailed material in Sections~2--3.

Readers who spot errors, omissions, or opportunities for improvement are invited to send in their feedback to be incorporated into the next version; please see the end of Section~1 for the specific way feedback can be reported.


\bigskip
\label{first bib}
   
\normalsize
\printbibliography[heading=none]

\section*{Alphabetic bibliography} \label{second bib}
\setlength{\bibitemsep}{-.2\baselineskip}
\togglefalse{showannotation}
\newrefcontext[sorting=nyt, labelprefix=A] 
\printbibliography[heading=none]

\end{document}
% $ biblatex auxiliary file $
% $ biblatex bbl format version 3.1 $
% Do not modify the above lines!
%
% This is an auxiliary file used by the 'biblatex' package.
% This file may safely be deleted. It will be recreated by
% biber as required.
%
\begingroup
\makeatletter
\@ifundefined{ver@biblatex.sty}
  {\@latex@error
     {Missing 'biblatex' package}
     {The bibliography requires the 'biblatex' package.}
      \aftergroup\endinput}
  {}
\endgroup


\refsection{0}
  \datalist[entry]{nyt/global/7FC56270E7A70FA81A5935B72EACBE29/global/global}
    \entry{MR3261965}{article}{}
      \name{author}{3}{}{%
        {{hash=3d7986c1f1aad61da9a4782349cfc695}{%
           family={Akbary},
           familyi={A\bibinitperiod},
           given={A.},
           giveni={A\bibinitperiod}}}%
        {{hash=585b691d1c27d2773434ed1ac9664075}{%
           family={Ng},
           familyi={N\bibinitperiod},
           given={N.},
           giveni={N\bibinitperiod}}}%
        {{hash=77b8d807a8642e05d0c1744c3120354e}{%
           family={Shahabi},
           familyi={S\bibinitperiod},
           given={M.},
           giveni={M\bibinitperiod}}}%
      }
      \strng{namehash}{87ed2f0893dc8a9968837b23ecc3e66b}
      \strng{fullhash}{87ed2f0893dc8a9968837b23ecc3e66b}
      \strng{bibnamehash}{87ed2f0893dc8a9968837b23ecc3e66b}
      \strng{authorbibnamehash}{87ed2f0893dc8a9968837b23ecc3e66b}
      \strng{authornamehash}{87ed2f0893dc8a9968837b23ecc3e66b}
      \strng{authorfullhash}{87ed2f0893dc8a9968837b23ecc3e66b}
      \field{sortinit}{A}
      \field{sortinithash}{a3dcedd53b04d1adfd5ac303ecd5e6fa}
      \field{labelprefix}{A}
      \field{labelnamesource}{author}
      \field{labeltitlesource}{title}
      \field{annotation}{The authors establish a sufficient condition for when general functions of the form \[\Re \sum_{\lambda_n\le X} r_n e^{i\lambda_n y} + E(y,X)\] (or vector-valued analogues) are $B^2$-almost-periodic and therefore possess limiting distributions. They apply this theorem, assuming the appropriate RH, to (logarithmically scaled versions of) functions such as: $\sum_{n\le x} \Lambda(n) a_n(\pi)$ for automorphic $L$-functions $L(s,\chi)$; the Mertens and Liouville functions $M(x)$ and $L(x)$, analogues of these for arithmetic progressions, and interpolations between these and $M_r(x)$ and $L_r(x)$; and (on the further assumption that Artin $L$-functions are entire) for prime counting functions in Chebotarev conjugacy classes. They provide further conclusions if the linear independence of the $\lambda_n$ is also assumed.\par This article cites~\cites{1938.Wintner,2000.Ng,MR0104638,MR598886,MR783538,MR1159354,MR1329368,MR1501802,MR1507933,MR1544509,MR1806291,MR2398787,MR2994374,MR3028758,MR3063909,MR3264673,MR3272280,MR2078705,MR584557}.}
      \field{issn}{0033-5606}
      \field{journaltitle}{Q. J. Math.}
      \field{number}{3}
      \field{title}{Limiting distributions of the classical error terms of prime number theory}
      \field{volume}{65}
      \field{year}{2014}
      \true{nocite}
      \field{pages}{743\bibrangedash 780}
      \range{pages}{38}
      \verb{doi}
      \verb 10.1093/qmath/hat059
      \endverb
      \verb{mrnumber}
      \verb 3261965
      \endverb
      \verb{urlraw}
      \verb https://doi.org/10.1093/qmath/hat059
      \endverb
      \verb{url}
      \verb https://doi.org/10.1093/qmath/hat059
      \endverb
    \endentry
    \entry{MR654520}{incollection}{}
      \name{author}{2}{}{%
        {{hash=44f716ae9c52c6e9e0952988d11b7582}{%
           family={Anderson},
           familyi={A\bibinitperiod},
           given={R.\bibnamedelimi J.},
           giveni={R\bibinitperiod\bibinitdelim J\bibinitperiod}}}%
        {{hash=32597e08070c5aa88e8fb07dd295112a}{%
           family={Stark},
           familyi={S\bibinitperiod},
           given={H.\bibnamedelimi M.},
           giveni={H\bibinitperiod\bibinitdelim M\bibinitperiod}}}%
      }
      \list{publisher}{1}{%
        {Springer, Berlin-New York}%
      }
      \strng{namehash}{b69969b5d1cab3b7bcc8edd374f284cb}
      \strng{fullhash}{b69969b5d1cab3b7bcc8edd374f284cb}
      \strng{bibnamehash}{b69969b5d1cab3b7bcc8edd374f284cb}
      \strng{authorbibnamehash}{b69969b5d1cab3b7bcc8edd374f284cb}
      \strng{authornamehash}{b69969b5d1cab3b7bcc8edd374f284cb}
      \strng{authorfullhash}{b69969b5d1cab3b7bcc8edd374f284cb}
      \field{sortinit}{A}
      \field{sortinithash}{a3dcedd53b04d1adfd5ac303ecd5e6fa}
      \field{labelprefix}{A}
      \field{labelnamesource}{author}
      \field{labeltitlesource}{title}
      \field{annotation}{The authors disprove several conjectures using a theorem of Landau and the locations of zeros of $\zeta(s)$. The authors also show that for conjectures of a certain type (such as von Sterneck's conjecture~\cites{1912.vonSterneck} $|M(x)|<\frac{1}{2}\sqrt{x}$ for $x>200$ and Pólya's problem~\cites{1919.Polya} $L(x)\leq 0$ for $x \geq 2$), a single counterexample implies infinitely many counterexamples. In particular, using the known counterexamples from Neubauer \cites{MR155787} that $E^M(x)>0.557$ for $x = 7.76 \times 10^9$, and from Lehman \cites{MR0120198} that $E^L(x)>0.023$ for $x=9.064 \times 10^8$, the authors conclude that $\limsup_{x \to \infty} E^M(x)>0.557$ and $\limsup_{x \to \infty} E^L(x)>0.023$.\par This article cites~\cites{MR0330069,MR0376566,MR0198145,MR0332685,MR0104638,MR0289452,MR0006202,MR0352026,MR0120198,MR549977,MR155787,MR280447,MR43841,MR51857}.}
      \field{bookpagination}{none}
      \field{booktitle}{Analytic number theory ({P}hiladelphia, {P}a., 1980)}
      \field{series}{Lecture Notes in Math.}
      \field{title}{Oscillation theorems}
      \field{volume}{899}
      \field{year}{1981}
      \true{nocite}
      \field{pages}{79\bibrangedash 106}
      \range{pages}{28}
      \verb{mrnumber}
      \verb 654520
      \endverb
    \endentry
    \entry{MR4369787}{article}{}
      \name{author}{1}{}{%
        {{hash=53e2b48286c0b37330d03fd888d20e26}{%
           family={Aymone},
           familyi={A\bibinitperiod},
           given={Marco},
           giveni={M\bibinitperiod}}}%
      }
      \strng{namehash}{53e2b48286c0b37330d03fd888d20e26}
      \strng{fullhash}{53e2b48286c0b37330d03fd888d20e26}
      \strng{bibnamehash}{53e2b48286c0b37330d03fd888d20e26}
      \strng{authorbibnamehash}{53e2b48286c0b37330d03fd888d20e26}
      \strng{authornamehash}{53e2b48286c0b37330d03fd888d20e26}
      \strng{authorfullhash}{53e2b48286c0b37330d03fd888d20e26}
      \field{sortinit}{A}
      \field{sortinithash}{a3dcedd53b04d1adfd5ac303ecd5e6fa}
      \field{labelprefix}{A}
      \field{labelnamesource}{author}
      \field{labeltitlesource}{title}
      \field{annotation}{The author observes that GRH for $L(s, \chi_{-4})$ would imply that $\sum_{p \leq x} \chi_{-4}(p)p^{-\sigma}$ would be eventually negative when $\sigma \to 1/2^+$. He investigates the connection between the sign changes of the weighted prime number races associated to a real nonprincipal Dirichlet character~$\chi$ and the zero-free region of the corresponding $L$-function $L(s, \chi)$, proving that there exists $0\leq\sigma<1$ such that $\sum_{p\leq x}\chi(p)p^{-\sigma}$ has only finitely many sign changes if and only if there exists $\ep>0$ such that $L(s,\chi)\neq 0$ in the half plane $\sigma>1-\ep$.\par This article cites~\cites{MR0882550, MR3773805, MR2202918}.}
      \field{issn}{1793-0421}
      \field{journaltitle}{Int. J. Number Theory}
      \field{number}{1}
      \field{title}{A note on prime number races and zero free regions for {$L$} functions}
      \field{volume}{18}
      \field{year}{2022}
      \true{nocite}
      \field{pages}{1\bibrangedash 8}
      \range{pages}{8}
      \verb{doi}
      \verb 10.1142/S1793042122500014
      \endverb
      \verb{mrnumber}
      \verb 4369787
      \endverb
      \verb{urlraw}
      \verb https://doi.org/10.1142/S1793042122500014
      \endverb
      \verb{url}
      \verb https://doi.org/10.1142/S1793042122500014
      \endverb
    \endentry
    \entry{MR4265352}{article}{}
      \name{author}{1}{}{%
        {{hash=49275b2a1e8fa5af0a0e7ef3941655a7}{%
           family={Bailleul},
           familyi={B\bibinitperiod},
           given={Alexandre},
           giveni={A\bibinitperiod}}}%
      }
      \strng{namehash}{49275b2a1e8fa5af0a0e7ef3941655a7}
      \strng{fullhash}{49275b2a1e8fa5af0a0e7ef3941655a7}
      \strng{bibnamehash}{49275b2a1e8fa5af0a0e7ef3941655a7}
      \strng{authorbibnamehash}{49275b2a1e8fa5af0a0e7ef3941655a7}
      \strng{authornamehash}{49275b2a1e8fa5af0a0e7ef3941655a7}
      \strng{authorfullhash}{49275b2a1e8fa5af0a0e7ef3941655a7}
      \field{extraname}{1}
      \field{sortinit}{B}
      \field{sortinithash}{8de16967003c7207dae369d874f1456e}
      \field{labelprefix}{A}
      \field{labelnamesource}{author}
      \field{labeltitlesource}{title}
      \field{annotation}{This article studies a generalization of Chebyshev's bias to the Chebotarev density theorem, which was first investigated by Ng~\cite{2000.Ng}. Namely, if $L/K$ is a Galois extension of number fields and $C$ is a conjugacy class of $\operatorname{Gal}(L/K)$, then $\pi_C(x)$ denotes the number of unramified prime ideals $\mathfrak{p}$ of $L$ with norm ${N(\mathfrak{p}) \leq x}$ and whose Artin symbol $(\frac{\mathfrak{p}}{L/K})$ equals $C$. For any two conjugacy classes $C_1$ and $C_2$ of $\operatorname{Gal}(L/K)$, a ``Chebotarev race'' measures the logarithmic density of values $x$ such that ${\pi_{C_1}(x)}/{\#C_1} > {\pi_{C_2}(x)}/{\#C_2}$. Fiorilli and Jouve~\cite{arXiv.FiorilliJouve} have shown that some Chebotarev races are extremely biased (in that the logarithmic density is close to 1) and the author continues this study in a new context.\par Using class field theory constructions and results on root numbers of quaternion extensions, the author analyzes Chebyshev's bias in certain families of number field extensions with Galois groups of $2$-power order, namely dihedral or generalized quaternion. Conditional on GRH and variants of LI, he performs an analysis in two different aspects: the ``horizontal aspect'' (fixed Galois group) and the ``vertical aspect'' (high-degree towers of $2$-power order). Each direction reveals an interesting connection between the central zeros of Artin $L$-functions and the bias of the Chebotarev race; this relationship is a key novelty of the article.\par This article cites~\cites{MR2474313,MR2811555,1853.Chebyshev,MR4120786,MR3803705,arXiv.FiorilliJouve,MR4400872,2000.Ng,MR3028758,MR3961320,MR3773805,MR1251926,MR1344249,MR0146156,MR0162771,MR2998146,MR3063909,MR4082248,MR1329368,MR0004255}.}
      \field{issn}{1937-0652}
      \field{journaltitle}{Algebra Number Theory}
      \field{number}{4}
      \field{title}{Chebyshev's bias in dihedral and generalized quaternion {G}alois groups}
      \field{volume}{15}
      \field{year}{2021}
      \true{nocite}
      \field{pages}{999\bibrangedash 1041}
      \range{pages}{43}
      \verb{doi}
      \verb 10.2140/ant.2021.15.999
      \endverb
      \verb{mrnumber}
      \verb 4265352
      \endverb
      \verb{urlraw}
      \verb https://doi.org/10.2140/ant.2021.15.999
      \endverb
      \verb{url}
      \verb https://doi.org/10.2140/ant.2021.15.999
      \endverb
    \endentry
    \entry{MR4447414}{article}{}
      \name{author}{1}{}{%
        {{hash=49275b2a1e8fa5af0a0e7ef3941655a7}{%
           family={Bailleul},
           familyi={B\bibinitperiod},
           given={Alexandre},
           giveni={A\bibinitperiod}}}%
      }
      \strng{namehash}{49275b2a1e8fa5af0a0e7ef3941655a7}
      \strng{fullhash}{49275b2a1e8fa5af0a0e7ef3941655a7}
      \strng{bibnamehash}{49275b2a1e8fa5af0a0e7ef3941655a7}
      \strng{authorbibnamehash}{49275b2a1e8fa5af0a0e7ef3941655a7}
      \strng{authornamehash}{49275b2a1e8fa5af0a0e7ef3941655a7}
      \strng{authorfullhash}{49275b2a1e8fa5af0a0e7ef3941655a7}
      \field{extraname}{2}
      \field{sortinit}{B}
      \field{sortinithash}{8de16967003c7207dae369d874f1456e}
      \field{labelprefix}{A}
      \field{labelnamesource}{author}
      \field{labeltitlesource}{title}
      \field{annotation}{The author proves general versions of Kronecker--Weyl theorems, both in the discrete and continuous settings and with no linear independence assumption, where the sets on which the equidistribution is guaranteed are explicitly constructed. The article consists of many applications of these explicit theorems, including sufficient conditions for lower bounds for the lower densities in certain Chebotarev races. The applications are given in a more general setting for random variables, where the densities are bounded by probabilities of certain explicitly given events. As a concrete application, in the last part of the article, the author studies densities for prime divisor races in geometric Galois extensions of function fields of one variable over finite fields, including the cases for which the LI hypothesis fails to hold.\par This article cites~\cites{MR3261965, MR2474313, MR3581815, MR2811555, MR4120786, MR4176773, MR4316675, arXiv.FiorilliJouve, MR1957110, MR4082248, MR4082248, 2000.Ng, MR1329368}.}
      \field{issn}{2522-0160}
      \field{journaltitle}{Res. Number Theory}
      \field{number}{3}
      \field{title}{Explicit {K}ronecker--{W}eyl theorems and applications to prime number races}
      \field{volume}{8}
      \field{year}{2022}
      \true{nocite}
      \field{pages}{Paper No. 43\bibrangessep 34}
      \range{pages}{-1}
      \verb{doi}
      \verb 10.1007/s40993-022-00349-2
      \endverb
      \verb{mrnumber}
      \verb 4447414
      \endverb
      \verb{urlraw}
      \verb https://doi.org/10.1007/s40993-022-00349-2
      \endverb
      \verb{url}
      \verb https://doi.org/10.1007/s40993-022-00349-2
      \endverb
    \endentry
    \entry{MR2013259}{article}{}
      \name{author}{2}{}{%
        {{hash=f8a84ad9ba564b4272e977c2c3bc0e70}{%
           family={Balanzario},
           familyi={B\bibinitperiod},
           given={E.\bibnamedelimi P.},
           giveni={E\bibinitperiod\bibinitdelim P\bibinitperiod}}}%
        {{hash=92c6e72ba89fceafe31367461d88d360}{%
           family={Hernández},
           familyi={H\bibinitperiod},
           given={S.},
           giveni={S\bibinitperiod}}}%
      }
      \strng{namehash}{63612cf56d8d86ce4a5d5b3a7102abf9}
      \strng{fullhash}{63612cf56d8d86ce4a5d5b3a7102abf9}
      \strng{bibnamehash}{63612cf56d8d86ce4a5d5b3a7102abf9}
      \strng{authorbibnamehash}{63612cf56d8d86ce4a5d5b3a7102abf9}
      \strng{authornamehash}{63612cf56d8d86ce4a5d5b3a7102abf9}
      \strng{authorfullhash}{63612cf56d8d86ce4a5d5b3a7102abf9}
      \field{sortinit}{B}
      \field{sortinithash}{8de16967003c7207dae369d874f1456e}
      \field{labelprefix}{A}
      \field{labelnamesource}{author}
      \field{labeltitlesource}{title}
      \field{annotation}{Fix distinct integers $a\ge0$ and $b\ge1$ and a real number $0<\theta<\frac1{2b}$, and define \[ h(x)=\sum_{n=1}^\infty \ell(n)e^{-n/x}-\frac{\Gamma(1/a)}{a\zeta(b/a)} x^{1/a}, \quad\text{where } \frac{\zeta(as)}{\zeta(bs)}=\sum_{n=1}^\infty\frac{\ell(n)}{n^s}. \] (When $a=0$, the second term in the definition of~$h$ is omitted.) Special cases include $\ell(n)=\mu(n)$ when $(a,b)=(0,1)$ and $\ell(n)=\lambda(n)$ when $(a,b)=(2,1)$. Assuming that the zeros of $\zeta(s)$, up to a height depending on $\theta$ and $b$, are simple and lie on the critical line, the authors show that $W(h;T;t^\theta) \gg \log T$. For example, in the special cases mentioned above, the authors verify the assumption up to height $10^5$ to deduce oscillations of size $t^{0.278}$. \par This article cites~\cites{MR777013,MR1193094}.}
      \field{issn}{0003-889X}
      \field{journaltitle}{Arch. Math. (Basel)}
      \field{number}{3}
      \field{title}{On the number of large oscillations of some arithmetical power series}
      \field{volume}{81}
      \field{year}{2003}
      \true{nocite}
      \field{pages}{285\bibrangedash 290}
      \range{pages}{6}
      \verb{doi}
      \verb 10.1007/s00013-003-4704-2
      \endverb
      \verb{mrnumber}
      \verb 2013259
      \endverb
    \endentry
    \entry{MR1109990}{article}{}
      \name{author}{1}{}{%
        {{hash=3a6f6cc642ae49cd63edd92858935cb9}{%
           family={Bartz},
           familyi={B\bibinitperiod},
           given={K.\bibnamedelimi M.},
           giveni={K\bibinitperiod\bibinitdelim M\bibinitperiod}}}%
      }
      \strng{namehash}{3a6f6cc642ae49cd63edd92858935cb9}
      \strng{fullhash}{3a6f6cc642ae49cd63edd92858935cb9}
      \strng{bibnamehash}{3a6f6cc642ae49cd63edd92858935cb9}
      \strng{authorbibnamehash}{3a6f6cc642ae49cd63edd92858935cb9}
      \strng{authornamehash}{3a6f6cc642ae49cd63edd92858935cb9}
      \strng{authorfullhash}{3a6f6cc642ae49cd63edd92858935cb9}
      \field{sortinit}{B}
      \field{sortinithash}{8de16967003c7207dae369d874f1456e}
      \field{labelprefix}{A}
      \field{labelnamesource}{author}
      \field{labeltitlesource}{title}
      \field{annotation}{Assuming RH and the simplicity of the zeros of $\zeta(s)$, Titchmarsh showed that \begin{align*} M_0(x)=\sum_\rho \frac{x^{\rho}}{\rho\zeta'(\rho)}-2-\sum_{n=1}^{\infty}\frac{(-1)^n(2\pi/x)^{2n}}{(2n)!n\zeta(2n+1)}. \end{align*} In Part I, the author investigates the function $m(z)=\sum_{\Im\rho>0}{e^{\rho z}}/{\zeta'(\rho)}$ (still assuming the simplicity of the zeros). He shows that $m(z)$ is a holomorphic function for $\Im z>0$ that has an analytic continuation to a meromorphic function on $\mathbb{C}$ satisfying the functional equation \begin{align*} m(z)+\overline{m(\overline{z})}=-2\sum_{n=1}^{\infty}\frac{\mu(n)}{n}\cos\bigg(\frac{2\pi}{n}e^{-z}\bigg). \end{align*} Using analytic properties of $m(z)$, he establishes Titchmarsh's formula for $M_0(x)$ without assuming RH. \par This article cites~\cites{1932.Ingham,MR0046485,MR783538,MR1083000}.}
      \field{issn}{0065-1036}
      \field{journaltitle}{Acta Arith.}
      \field{number}{4}
      \field{title}{On some complex explicit formulae connected with the {M}öbius function. {I}, {II}}
      \field{volume}{57}
      \field{year}{1991}
      \true{nocite}
      \field{pages}{283\bibrangedash 293\bibrangessep 295\bibrangedash 305}
      \range{pages}{22}
      \verb{doi}
      \verb 10.4064/aa-57-4-283-293
      \endverb
      \verb{mrnumber}
      \verb 1109990
      \endverb
      \verb{urlraw}
      \verb https://doi.org/10.4064/aa-57-4-283-293
      \endverb
      \verb{url}
      \verb https://doi.org/10.4064/aa-57-4-283-293
      \endverb
    \endentry
    \entry{MR0330069}{article}{}
      \name{author}{5}{}{%
        {{hash=7c2cfee9c4a819ab77cc793d9ce1bdab}{%
           family={Bateman},
           familyi={B\bibinitperiod},
           given={P.\bibnamedelimi T.},
           giveni={P\bibinitperiod\bibinitdelim T\bibinitperiod}}}%
        {{hash=6bdab84d1ccda38b53176b70d8f29951}{%
           family={Brown},
           familyi={B\bibinitperiod},
           given={J.\bibnamedelimi W.},
           giveni={J\bibinitperiod\bibinitdelim W\bibinitperiod}}}%
        {{hash=98ca3e62c2dc209695914011ea7ad03d}{%
           family={Hall},
           familyi={H\bibinitperiod},
           given={R.\bibnamedelimi S.},
           giveni={R\bibinitperiod\bibinitdelim S\bibinitperiod}}}%
        {{hash=82a56dc9f4059d331b665de076b6ccc8}{%
           family={Kloss},
           familyi={K\bibinitperiod},
           given={K.\bibnamedelimi E.},
           giveni={K\bibinitperiod\bibinitdelim E\bibinitperiod}}}%
        {{hash=af6770049aed53b1c6747cf202761a2c}{%
           family={Stemmler},
           familyi={S\bibinitperiod},
           given={Rosemarie\bibnamedelima M.},
           giveni={R\bibinitperiod\bibinitdelim M\bibinitperiod}}}%
      }
      \list{publisher}{1}{%
        {Academic Press, London}%
      }
      \strng{namehash}{818aad5926f1e301d22acb7b91a4fafa}
      \strng{fullhash}{e164d0e7c5218899980cab8c648e58c0}
      \strng{bibnamehash}{e164d0e7c5218899980cab8c648e58c0}
      \strng{authorbibnamehash}{e164d0e7c5218899980cab8c648e58c0}
      \strng{authornamehash}{818aad5926f1e301d22acb7b91a4fafa}
      \strng{authorfullhash}{e164d0e7c5218899980cab8c648e58c0}
      \field{sortinit}{B}
      \field{sortinithash}{8de16967003c7207dae369d874f1456e}
      \field{labelprefix}{A}
      \field{labelnamesource}{author}
      \field{labeltitlesource}{title}
      \field{annotation}{This article strengthens a result of Ingham~\cite{MR0006202} in a way that allows for computational exploration. Define a sum $\sum_{n=1}^N c_n\gamma_n$, where $\gamma_1, \gamma_2, \ldots$ are the positive ordinates of the nontrivial zeros of $\zeta(s)$, to be of type~(B) if each $c_n\in\{-2,-1,0,1,2\}$ with at most one $|c_n|$ equal to~$2$. The authors show that if at most finitely many sums of type (B) equal zero, then \begin{align*} \limsup_{x \to \infty} x^{-1/2}L(x)=\infty, &\quad \liminf_{x \to \infty} x^{-1/2}L(x)=-\infty, \\ \limsup_{x \to \infty} x^{-1/2}M(x)=\infty, &\quad \liminf_{x \to \infty} x^{-1/2}M(x)=-\infty. \end{align*} Tables of data are given for the smallest sums of type (B) for $N=1,\ldots,20$, as well as of type~(A) sums where only $c_n\in\{-1,0,1\}$ are allowed.\par This article cites~\cite{1932.Ingham, MR0006202}.}
      \field{booktitle}{Computers in number theory ({P}roc. {S}ci. {R}es. {C}ouncil {A}tlas {S}ympos. {N}o. 2, {O}xford, 1969)}
      \field{title}{Linear relations connecting the imaginary parts of the zeros of the zeta function}
      \field{year}{1971}
      \true{nocite}
      \field{pages}{11\bibrangedash 19}
      \range{pages}{9}
      \verb{mrnumber}
      \verb 0330069
      \endverb
    \endentry
    \entry{MR1816036}{article}{}
      \name{author}{4}{}{%
        {{hash=e542667d9161d12cece57b09b9c6c693}{%
           family={Bays},
           familyi={B\bibinitperiod},
           given={C.},
           giveni={C\bibinitperiod}}}%
        {{hash=e7e31df1803902b56a78f32b1c733598}{%
           family={Ford},
           familyi={F\bibinitperiod},
           given={K.},
           giveni={K\bibinitperiod}}}%
        {{hash=08e54e6381028d2260f342dd1c27f66e}{%
           family={Hudson},
           familyi={H\bibinitperiod},
           given={R.\bibnamedelimi H.},
           giveni={R\bibinitperiod\bibinitdelim H\bibinitperiod}}}%
        {{hash=9073d8f33230938a75c73529bcd5e1ef}{%
           family={Rubinstein},
           familyi={R\bibinitperiod},
           given={M.},
           giveni={M\bibinitperiod}}}%
      }
      \strng{namehash}{0543a8fb7e634d40caee3a99730ca685}
      \strng{fullhash}{4815739125bfd0a086513390034cdd23}
      \strng{bibnamehash}{4815739125bfd0a086513390034cdd23}
      \strng{authorbibnamehash}{4815739125bfd0a086513390034cdd23}
      \strng{authornamehash}{0543a8fb7e634d40caee3a99730ca685}
      \strng{authorfullhash}{4815739125bfd0a086513390034cdd23}
      \field{sortinit}{B}
      \field{sortinithash}{8de16967003c7207dae369d874f1456e}
      \field{labelprefix}{A}
      \field{labelnamesource}{author}
      \field{labeltitlesource}{title}
      \field{annotation}{This article is concerned with estimating the densities $\delta_{q,\cN,\cR}$, in particular by using the fact that the truncated explicit formula is an almost-periodic function whose largest ``quasi-period'' is dictated by its lowest zero. The authors present data supporting the fact that these densities depends strongly on the location of the first few zeros of $L(s,\chi_{\pm q})$ and the size of the first zero. Plots of $E(x;q,\cN,\cR)$ for all primes $q$ with $h(-q)=1$ are provided, and the difficulty of accurately estimating the densities from such data is pointed out. The authors note that the Chowla--Selberg formula implies that if $Q(\sqrt{-q} )$ is an imaginary quadratic field with class number~$1$ then $L(s,\chi_{-q})$ has a relatively low-lying zero; they also analyze the connection between low-lying zeros and class numbers $3$ and~$5$.\par This article cites~\cites{MR0476616, MR529694, MR1651741, 1853.Chebyshev, MR552467, MR1251926, MR1310738, MR0146156, MR0083001, 1914.Littlewood, MR1329368, MR0108470}.}
      \field{issn}{0022-314X}
      \field{journaltitle}{J. Number Theory}
      \field{number}{1}
      \field{title}{Zeros of {D}irichlet {$L$}-functions near the real axis and {C}hebyshev's bias}
      \field{volume}{87}
      \field{year}{2001}
      \true{nocite}
      \field{pages}{54\bibrangedash 76}
      \range{pages}{23}
      \verb{doi}
      \verb 10.1006/jnth.2000.2601
      \endverb
      \verb{mrnumber}
      \verb 1816036
      \endverb
      \verb{urlraw}
      \verb https://doi.org/10.1006/jnth.2000.2601
      \endverb
      \verb{url}
      \verb https://doi.org/10.1006/jnth.2000.2601
      \endverb
    \endentry
    \entry{MR0447090}{article}{}
      \name{author}{2}{}{%
        {{hash=e542667d9161d12cece57b09b9c6c693}{%
           family={Bays},
           familyi={B\bibinitperiod},
           given={C.},
           giveni={C\bibinitperiod}}}%
        {{hash=08e54e6381028d2260f342dd1c27f66e}{%
           family={Hudson},
           familyi={H\bibinitperiod},
           given={R.\bibnamedelimi H.},
           giveni={R\bibinitperiod\bibinitdelim H\bibinitperiod}}}%
      }
      \strng{namehash}{121290fcdac2ab7b7c2b8152b61c32dc}
      \strng{fullhash}{121290fcdac2ab7b7c2b8152b61c32dc}
      \strng{bibnamehash}{121290fcdac2ab7b7c2b8152b61c32dc}
      \strng{authorbibnamehash}{121290fcdac2ab7b7c2b8152b61c32dc}
      \strng{authornamehash}{121290fcdac2ab7b7c2b8152b61c32dc}
      \strng{authorfullhash}{121290fcdac2ab7b7c2b8152b61c32dc}
      \field{extraname}{1}
      \field{sortinit}{B}
      \field{sortinithash}{8de16967003c7207dae369d874f1456e}
      \field{labelprefix}{A}
      \field{labelnamesource}{author}
      \field{labeltitlesource}{title}
      \field{annotation}{The authors describe in detail a refinement of the segmented sieve of Eratosthenes, which they call the dual sieve, designed to lower the execution time. As an illustration, they record the number of primes in the eight reduced residue classes modulo $24$ (from which one can calculate the number of primes in residue classes modulo any divisor of~$24$) up to $10^{11},2 \times 10^{11},\ldots, 10^{12}$. From their table, one easily observes that $\pi(x;24,1)$ is consistently smaller than any other $\pi(x;24,a)$ by an amount that is very roughly $\frac12\pi(\sqrt x)$.}
      \field{journaltitle}{Nordisk Tidskr. Informationsbehandling (BIT)}
      \field{number}{2}
      \field{title}{The segmented sieve of {E}ratosthenes and primes in arithmetic progressions to {$10^{12}$}}
      \field{volume}{17}
      \field{year}{1977}
      \true{nocite}
      \field{pages}{121\bibrangedash 127}
      \range{pages}{7}
      \verb{doi}
      \verb 10.1007/bf01932283
      \endverb
      \verb{mrnumber}
      \verb 0447090
      \endverb
      \verb{urlraw}
      \verb https://doi.org/10.1007/bf01932283
      \endverb
      \verb{url}
      \verb https://doi.org/10.1007/bf01932283
      \endverb
    \endentry
    \entry{MR0476616}{article}{}
      \name{author}{2}{}{%
        {{hash=e542667d9161d12cece57b09b9c6c693}{%
           family={Bays},
           familyi={B\bibinitperiod},
           given={C.},
           giveni={C\bibinitperiod}}}%
        {{hash=08e54e6381028d2260f342dd1c27f66e}{%
           family={Hudson},
           familyi={H\bibinitperiod},
           given={R.\bibnamedelimi H.},
           giveni={R\bibinitperiod\bibinitdelim H\bibinitperiod}}}%
      }
      \strng{namehash}{121290fcdac2ab7b7c2b8152b61c32dc}
      \strng{fullhash}{121290fcdac2ab7b7c2b8152b61c32dc}
      \strng{bibnamehash}{121290fcdac2ab7b7c2b8152b61c32dc}
      \strng{authorbibnamehash}{121290fcdac2ab7b7c2b8152b61c32dc}
      \strng{authornamehash}{121290fcdac2ab7b7c2b8152b61c32dc}
      \strng{authorfullhash}{121290fcdac2ab7b7c2b8152b61c32dc}
      \field{extraname}{2}
      \field{sortinit}{B}
      \field{sortinithash}{8de16967003c7207dae369d874f1456e}
      \field{labelprefix}{A}
      \field{labelnamesource}{author}
      \field{labeltitlesource}{title}
      \field{annotation}{The authors determine that $x = 608,981,813,029$ is the smallest $x$ such that $\pi(x;3,2,1) = -1$. A faster version of a previous program of theirs (which had run up to $2.5 \times 10^{11}$) was used to find this sign change. The authors provide graphs of $\pi(x;3,2,1)$ near this first sign change; they highlight that $\pi(x;3,2,1)$ becomes negative at two separate regions near the sign change, before taking on values shortly after that are much more positive. The authors observe that neither $\pi(x;3,2,1)$ nor $\pi(x;4,3,1)$ becomes very negative near the occurrence of its first negative values; in attempts to determine a smaller Skewes number, consequently, they recommend evaluation of $\Delta^\pi(x)$ in regular intervals in order to not miss a ``shallow'' sign change.\par This article cites~\cites{MR529694,1853.Chebyshev,1932.Ingham,MR0162771,MR0083001,MR0202686,1914.Littlewood,MR0108470,MR0476615}.}
      \field{issn}{0025-5718}
      \field{journaltitle}{Math. Comp.}
      \field{number}{142}
      \field{title}{Details of the first region of integers {$x$} with {$\pi _{3,2}(x)<\pi _{3,1}(x)$}}
      \field{volume}{32}
      \field{year}{1978}
      \true{nocite}
      \field{pages}{571\bibrangedash 576}
      \range{pages}{6}
      \verb{doi}
      \verb 10.2307/2006165
      \endverb
      \verb{mrnumber}
      \verb 0476616
      \endverb
      \verb{urlraw}
      \verb https://doi.org/10.2307/2006165
      \endverb
      \verb{url}
      \verb https://doi.org/10.2307/2006165
      \endverb
    \endentry
    \entry{MR529694}{article}{}
      \name{author}{2}{}{%
        {{hash=e542667d9161d12cece57b09b9c6c693}{%
           family={Bays},
           familyi={B\bibinitperiod},
           given={C.},
           giveni={C\bibinitperiod}}}%
        {{hash=08e54e6381028d2260f342dd1c27f66e}{%
           family={Hudson},
           familyi={H\bibinitperiod},
           given={R.\bibnamedelimi H.},
           giveni={R\bibinitperiod\bibinitdelim H\bibinitperiod}}}%
      }
      \strng{namehash}{121290fcdac2ab7b7c2b8152b61c32dc}
      \strng{fullhash}{121290fcdac2ab7b7c2b8152b61c32dc}
      \strng{bibnamehash}{121290fcdac2ab7b7c2b8152b61c32dc}
      \strng{authorbibnamehash}{121290fcdac2ab7b7c2b8152b61c32dc}
      \strng{authornamehash}{121290fcdac2ab7b7c2b8152b61c32dc}
      \strng{authorfullhash}{121290fcdac2ab7b7c2b8152b61c32dc}
      \field{extraname}{3}
      \field{sortinit}{B}
      \field{sortinithash}{8de16967003c7207dae369d874f1456e}
      \field{labelprefix}{A}
      \field{labelnamesource}{author}
      \field{labeltitlesource}{title}
      \field{annotation}{The authors of this article determine by computation the locations where $\pi(x;4,3,1) < 0$, and where $\pi(x;8,a,1) < 0$ for any $a\in\{3,5,7\}$, for $x$ up to $10^{12}$. For $x < 10^9$, a check was made at every prime; for $10^9 \le x \le 10^{12}$, a check was made every $10^7$ integers, with additional checks in between if $\pi(x;q,a,b)$ was found to be near zero. They then organize these locations into ``axis-crossing regions'' (ACRs) $[m,n]$, where $\pi(m;q,a,1) = \pi(n;q,a,1) = -1$ and $\pi(x;q,a,1) \ge 0$ for all $x$ outside an ACR, with $m$ at least twice as large as the upper bound for the previous ACR. \par For $q=4$, they find six distinct ACRs under $10^{12}$. For $(q,a) = (8,5)$, they find two ACRs under $10^{12}$ and find no ACRs for $(q,a) = (8,3)$ or $(q,a) = (8,7)$. They compare their computations to earlier published results from Leech \cite{MR0083001}, Shanks \cite{MR0108470}, and an unpublished communication from Lehmer (dated October 29, 1975). While their results overlap with Leech and Shanks for $q=4$ for $x \le 3 \cdot 10^6$, they find that their new information contradicts a prior characterization of the ACRs as mostly consisting of sparse, tiny intervals. For example, one ACR below $x < 2 \cdot 10^{10}$ contains $5\cdot10^8$ integers where $\pi(x;4,3,1) < 0$; another ACR between $37\cdot10^9$ and $39\cdot10^9$ contains $1.2\cdot10^9$ integers with $\pi(x;8,5,1) < 0$. Consequently, they argue that for large $x$, such regions may be more typical than sign-changes being sparse, isolated points. \par This article cites~\cites{MR0476615,MR0447090,1853.Chebyshev,MR0146157,MR0146158,MR0083001,1914.Littlewood,MR0108470}.}
      \field{issn}{0161-1712}
      \field{journaltitle}{Internat. J. Math. Math. Sci.}
      \field{number}{1}
      \field{title}{Numerical and graphical description of all axis crossing regions for the moduli {$4$} and {$8$} which occur before {$10^{12}$}}
      \field{volume}{2}
      \field{year}{1979}
      \true{nocite}
      \field{pages}{111\bibrangedash 119}
      \range{pages}{9}
      \verb{doi}
      \verb 10.1155/S0161171279000119
      \endverb
      \verb{mrnumber}
      \verb 529694
      \endverb
      \verb{urlraw}
      \verb https://doi.org/10.1155/S0161171279000119
      \endverb
      \verb{url}
      \verb https://doi.org/10.1155/S0161171279000119
      \endverb
    \endentry
    \entry{MR1651741}{article}{}
      \name{author}{2}{}{%
        {{hash=e542667d9161d12cece57b09b9c6c693}{%
           family={Bays},
           familyi={B\bibinitperiod},
           given={C.},
           giveni={C\bibinitperiod}}}%
        {{hash=08e54e6381028d2260f342dd1c27f66e}{%
           family={Hudson},
           familyi={H\bibinitperiod},
           given={R.\bibnamedelimi H.},
           giveni={R\bibinitperiod\bibinitdelim H\bibinitperiod}}}%
      }
      \strng{namehash}{121290fcdac2ab7b7c2b8152b61c32dc}
      \strng{fullhash}{121290fcdac2ab7b7c2b8152b61c32dc}
      \strng{bibnamehash}{121290fcdac2ab7b7c2b8152b61c32dc}
      \strng{authorbibnamehash}{121290fcdac2ab7b7c2b8152b61c32dc}
      \strng{authornamehash}{121290fcdac2ab7b7c2b8152b61c32dc}
      \strng{authorfullhash}{121290fcdac2ab7b7c2b8152b61c32dc}
      \field{extraname}{4}
      \field{sortinit}{B}
      \field{sortinithash}{8de16967003c7207dae369d874f1456e}
      \field{labelprefix}{A}
      \field{labelnamesource}{author}
      \field{labeltitlesource}{title}
      \field{annotation}{The authors describe computations of both $\pi(x;4,3,1)$ for values of $x$ past $10^{12}$, as well as its estimate using a truncated explicit formula with 12,000 zeros of $L(s,\chi{-4})$ for values of $x$ up to $10^{1000}$. The estimate duplicates the true distribution with satisfying accuracy, rediscovers all known axis-crossing regions, and finds probable new axis-crossing regions. The method extends to other differences such as $\Delta^\pi(x)$ and $\pi(x;q,\cN,\cR)$.\par This article cites~\cites{MR0476616, MR529694, MR1651741, MR0447090, MR1310738, MR0083001, MR1329368, MR0108470}.}
      \field{issn}{0025-5718}
      \field{journaltitle}{Math. Comp.}
      \field{number}{230}
      \field{title}{Zeroes of {D}irichlet {$L$}-functions and irregularities in the distribution of primes}
      \field{volume}{69}
      \field{year}{2000}
      \true{nocite}
      \field{pages}{861\bibrangedash 866}
      \range{pages}{6}
      \verb{doi}
      \verb 10.1090/S0025-5718-99-01105-9
      \endverb
      \verb{mrnumber}
      \verb 1651741
      \endverb
      \verb{urlraw}
      \verb https://doi.org/10.1090/S0025-5718-99-01105-9
      \endverb
      \verb{url}
      \verb https://doi.org/10.1090/S0025-5718-99-01105-9
      \endverb
    \endentry
    \entry{MR0476615}{article}{}
      \name{author}{2}{}{%
        {{hash=dcdb9f9c423bbcbfa2a60d519ccd6915}{%
           family={Bays},
           familyi={B\bibinitperiod},
           given={Carter},
           giveni={C\bibinitperiod}}}%
        {{hash=9e0609c9931c1b620df97e8c58d07532}{%
           family={Hudson},
           familyi={H\bibinitperiod},
           given={Richard\bibnamedelima H.},
           giveni={R\bibinitperiod\bibinitdelim H\bibinitperiod}}}%
      }
      \strng{namehash}{0824391c51610edfe4bd3c404488dea5}
      \strng{fullhash}{0824391c51610edfe4bd3c404488dea5}
      \strng{bibnamehash}{0824391c51610edfe4bd3c404488dea5}
      \strng{authorbibnamehash}{0824391c51610edfe4bd3c404488dea5}
      \strng{authornamehash}{0824391c51610edfe4bd3c404488dea5}
      \strng{authorfullhash}{0824391c51610edfe4bd3c404488dea5}
      \field{extraname}{5}
      \field{sortinit}{B}
      \field{sortinithash}{8de16967003c7207dae369d874f1456e}
      \field{labelprefix}{A}
      \field{labelnamesource}{author}
      \field{labeltitlesource}{title}
      \field{annotation}{The authors describe the sixth ``axis crossing region'' (a term they define rigorously) for $\pi(x;4,3,1)$, in which there are $4.1 \times 10^8$ consecutive integers satisfying $\pi(n;4,3,1)<0$. The authors suggest from their findings of surprisingly large axis crossing regions that the number of integers number of integers $n\le x$ with $\pi(n;4,3,1)<0$ might be $\Omega(x/\log x)$, while still respecting the conjecture that this counting function is $o(x)$.\par This article cites~\cites{1853.Chebyshev,MR1555148,MR0460261,MR0184920,MR0146158,MR0156826,MR0162771,MR0174538,MR0068565,MR1544293A,MR0083001,1914.Littlewood,MR1576856,MR0108470}.}
      \field{issn}{0025-5718}
      \field{journaltitle}{Math. Comp.}
      \field{number}{141}
      \field{title}{On the fluctuations of {L}ittlewood for primes of the form {$4n \pm 1$}}
      \field{volume}{32}
      \field{year}{1978}
      \true{nocite}
      \field{pages}{281\bibrangedash 286}
      \range{pages}{6}
      \verb{doi}
      \verb 10.2307/2006277
      \endverb
      \verb{mrnumber}
      \verb 0476615
      \endverb
    \endentry
    \entry{MR0472726}{article}{}
      \name{author}{2}{}{%
        {{hash=dcdb9f9c423bbcbfa2a60d519ccd6915}{%
           family={Bays},
           familyi={B\bibinitperiod},
           given={Carter},
           giveni={C\bibinitperiod}}}%
        {{hash=9e0609c9931c1b620df97e8c58d07532}{%
           family={Hudson},
           familyi={H\bibinitperiod},
           given={Richard\bibnamedelima H.},
           giveni={R\bibinitperiod\bibinitdelim H\bibinitperiod}}}%
      }
      \strng{namehash}{0824391c51610edfe4bd3c404488dea5}
      \strng{fullhash}{0824391c51610edfe4bd3c404488dea5}
      \strng{bibnamehash}{0824391c51610edfe4bd3c404488dea5}
      \strng{authorbibnamehash}{0824391c51610edfe4bd3c404488dea5}
      \strng{authornamehash}{0824391c51610edfe4bd3c404488dea5}
      \strng{authorfullhash}{0824391c51610edfe4bd3c404488dea5}
      \field{extraname}{6}
      \field{sortinit}{B}
      \field{sortinithash}{8de16967003c7207dae369d874f1456e}
      \field{labelprefix}{A}
      \field{labelnamesource}{author}
      \field{labeltitlesource}{title}
      \field{annotation}{The authors describe their empirical observation that $\pi(x;24,13,1)<0$ for about a third of the integers between $0.978\times10^{12}$ and $1.094\times10^{12}$. They assert the expectation that the set of such real numbers~$x$ has density~$0$.\par This article cites~\cite{MR0476615, MR0447090, MR0476616, MR0460261, MR0146157, MR0162771, MR0174538, MR0068565, MR0108470}.}
      \field{issn}{0075-4102}
      \field{journaltitle}{J. Reine Angew. Math.}
      \field{title}{{The appearance of tens of billions of integers {$x$} with {$\pi_{24,13}(x)<\pi_{24,1}(x)$} in the vicinity of {$10^{12}$}}}
      \field{volume}{299/300}
      \field{year}{1978}
      \true{nocite}
      \field{pages}{234\bibrangedash 237}
      \range{pages}{4}
      \verb{doi}
      \verb 10.1515/crll.1978.299-300.234
      \endverb
      \verb{mrnumber}
      \verb 0472726
      \endverb
    \endentry
    \entry{MR1752093}{article}{}
      \name{author}{2}{}{%
        {{hash=dcdb9f9c423bbcbfa2a60d519ccd6915}{%
           family={Bays},
           familyi={B\bibinitperiod},
           given={Carter},
           giveni={C\bibinitperiod}}}%
        {{hash=9e0609c9931c1b620df97e8c58d07532}{%
           family={Hudson},
           familyi={H\bibinitperiod},
           given={Richard\bibnamedelima H.},
           giveni={R\bibinitperiod\bibinitdelim H\bibinitperiod}}}%
      }
      \strng{namehash}{0824391c51610edfe4bd3c404488dea5}
      \strng{fullhash}{0824391c51610edfe4bd3c404488dea5}
      \strng{bibnamehash}{0824391c51610edfe4bd3c404488dea5}
      \strng{authorbibnamehash}{0824391c51610edfe4bd3c404488dea5}
      \strng{authornamehash}{0824391c51610edfe4bd3c404488dea5}
      \strng{authorfullhash}{0824391c51610edfe4bd3c404488dea5}
      \field{extraname}{7}
      \field{sortinit}{B}
      \field{sortinithash}{8de16967003c7207dae369d874f1456e}
      \field{labelprefix}{A}
      \field{labelnamesource}{author}
      \field{labeltitlesource}{title}
      \field{annotation}{Using Lehman's theorem~\cite{MR0202686} together with the first $10^6$ zeros of $\zeta(s)$ (supplied by Odlyzko), the authors show that there exist many integers~$x$ in the range $[\exp(727.95209-.002),\exp(727.95209+.002)]$ for which $\pi(x)>\li(x)$. (This corresponds to the interval $[1.3954272\times10^{316}, 1.4010201\times10^{316})$, although the authors claim that the interval is $[1.398201\times10^{316}, 1.398244\times10^{316}]$.) They also report on computations approximating $\pi(x)-\li(x)$ for $x$ up to about $4\times10^{12370}$ (the significance being that the 20th region with many solutions to $\pi(x)>\li(x)$ is probably around this number) that support the conditional result of Rubinstein and Sarnak that $\delta(\li,\pi) \approx 0.99999974$.\par This article cites~\cites{MR1651741,MR1816036,1932.Ingham,MR0202686,1914.Littlewood,MR866118,MR1329368,MR1573970,MR0067145}.}
      \field{issn}{0025-5718}
      \field{journaltitle}{Math. Comp.}
      \field{number}{231}
      \field{title}{A new bound for the smallest {$x$} with {$\pi (x)>\li (x)$}}
      \field{volume}{69}
      \field{year}{2000}
      \true{nocite}
      \field{pages}{1285\bibrangedash 1296}
      \range{pages}{12}
      \verb{mrnumber}
      \verb 1752093
      \endverb
    \endentry
    \entry{MR675189}{article}{}
      \name{author}{1}{}{%
        {{hash=6bf42c5abf099333fae470886fff8616}{%
           family={Bentz},
           familyi={B\bibinitperiod},
           given={H.-J.},
           giveni={H\bibinithyphendelim J\bibinitperiod}}}%
      }
      \strng{namehash}{6bf42c5abf099333fae470886fff8616}
      \strng{fullhash}{6bf42c5abf099333fae470886fff8616}
      \strng{bibnamehash}{6bf42c5abf099333fae470886fff8616}
      \strng{authorbibnamehash}{6bf42c5abf099333fae470886fff8616}
      \strng{authornamehash}{6bf42c5abf099333fae470886fff8616}
      \strng{authorfullhash}{6bf42c5abf099333fae470886fff8616}
      \field{sortinit}{B}
      \field{sortinithash}{8de16967003c7207dae369d874f1456e}
      \field{labelprefix}{A}
      \field{labelnamesource}{author}
      \field{labeltitlesource}{title}
      \field{annotation}{For $0\le\alpha<\frac12$, the author shows unconditionally that \[ \sum_p \chi_{-4}(p) \frac{\log p}{p^\alpha} e^{-(\log x)^2/p} \sim -\frac{\sqrt{\pi x}}2 e^{x(1-2\alpha)^2/16}; \] when $\alpha=\frac12$, the right-hand side must be replaced by $\frac14\sqrt{\pi x}$. Both results remain valid if $\chi_{-4}$ is replaced by $\chi_{-3}$. These results can be interpreted as comparing (in a specific way) the residue class $1$ to the other reduced residue class modulo~$4$ or~$3$. Analogously, when $\alpha=\frac12$, the author establishes the same result when comparing $1\mod8$ to another reduced residue class\mod8; if two reduced residue classes\mod8 are compared, the resulting expression is bounded. The author asserts that the required hypotheses on zeros of relevant Dirichlet $L$-functions is that they do not vanish in the ``bowtie'' $\{s\colon \sigma>0,\, 0<|t|<|\sigma-\frac12|\}$. The author also presents some numerical data concerning the prime number race\mod 3.\par This article cites~\cites{MR595317, MR534235, MR552462, MR1555148, MR0460261, 1932.Ingham, MR0146156, MR0146157, MR0146158, MR0146159, MR0146160, MR0146161, MR0156826, MR0156827, MR0162771, MR0174538, MR0180539, MR0182616A, MR0182616B, MR0200250, MR0302585, MR1544293A, MR1544293B, MR0083001, MR0108470}.}
      \field{issn}{0022-314X}
      \field{journaltitle}{J. Number Theory}
      \field{number}{2}
      \field{title}{Discrepancies in the distribution of prime numbers}
      \field{volume}{15}
      \field{year}{1982}
      \true{nocite}
      \field{pages}{252\bibrangedash 274}
      \range{pages}{23}
      \verb{doi}
      \verb 10.1016/0022-314X(82)90030-0
      \endverb
      \verb{mrnumber}
      \verb 675189
      \endverb
      \verb{urlraw}
      \verb https://doi.org/10.1016/0022-314X(82)90030-0
      \endverb
      \verb{url}
      \verb https://doi.org/10.1016/0022-314X(82)90030-0
      \endverb
    \endentry
    \entry{MR595317}{article}{}
      \name{author}{2}{}{%
        {{hash=6bf42c5abf099333fae470886fff8616}{%
           family={Bentz},
           familyi={B\bibinitperiod},
           given={H.-J.},
           giveni={H\bibinithyphendelim J\bibinitperiod}}}%
        {{hash=c6d965df16e600c701d4bcc2888ae34c}{%
           family={Pintz},
           familyi={P\bibinitperiod},
           given={J.},
           giveni={J\bibinitperiod}}}%
      }
      \strng{namehash}{8ee030157e0efba486591e59a6b5401b}
      \strng{fullhash}{8ee030157e0efba486591e59a6b5401b}
      \strng{bibnamehash}{8ee030157e0efba486591e59a6b5401b}
      \strng{authorbibnamehash}{8ee030157e0efba486591e59a6b5401b}
      \strng{authornamehash}{8ee030157e0efba486591e59a6b5401b}
      \strng{authorfullhash}{8ee030157e0efba486591e59a6b5401b}
      \field{extraname}{1}
      \field{sortinit}{B}
      \field{sortinithash}{8de16967003c7207dae369d874f1456e}
      \field{labelprefix}{A}
      \field{labelnamesource}{author}
      \field{labeltitlesource}{title}
      \field{annotation}{The first section of this article offers a short, conventional history of prime number races, specifically citing Shanks's computational work and heuristics from~\cite{MR0108470} as motivation for its results. Let $\ell_1$ be a quadratic residue\mod q and $\ell_2$ a quadratic nonresidue\mod q. Suppose that Dirichlet $L$-functions\mod q satisfying the condition that all zeros $\beta+i\gamma$ satisfy the inequality $\beta^2-\gamma^2 \le \frac 14$ (a ``bowtie'' assumption). Then for $0 \le \alpha < 1/2$, \[ \sum_{p\equiv\ell_1\mod q} \frac{\log p}{p^\alpha} e^{-(\log p)^2/x} - \sum_{p\equiv\ell_2\mod q} \frac{\log p}{p^\alpha} e^{-(\log p)^2/x} \sim \frac{c_q}{\phi(q)} \sqrt{\pi x} \cdot e^{\frac x4 (\frac 12 - \alpha)^2}, \] and in particular tends to infinity. By computations of Spira, this result is unconditional for $q\le24$.\par This article cites~\cites{MR662791, MR534235, MR552462, 1853.Chebyshev, MR1555148, MR0146156, MR0146157, MR0146158, MR0146159, MR0146160, MR0146161, MR0156826, MR0156827, MR0162771, MR0174538, MR0180539, MR0182616A, MR0182616B, MR0200250, MR1544293A, MR1544293B, MR0083001, MR0108470}.}
      \field{issn}{0026-9255}
      \field{journaltitle}{Monatsh. Math.}
      \field{number}{2}
      \field{title}{Quadratic residues and the distribution of prime numbers}
      \field{volume}{90}
      \field{year}{1980}
      \true{nocite}
      \field{pages}{91\bibrangedash 100}
      \range{pages}{10}
      \verb{doi}
      \verb 10.1007/BF01303260
      \endverb
      \verb{mrnumber}
      \verb 595317
      \endverb
      \verb{urlraw}
      \verb https://doi.org/10.1007/BF01303260
      \endverb
      \verb{url}
      \verb https://doi.org/10.1007/BF01303260
      \endverb
    \endentry
    \entry{MR662791}{article}{}
      \name{author}{2}{}{%
        {{hash=6bf42c5abf099333fae470886fff8616}{%
           family={Bentz},
           familyi={B\bibinitperiod},
           given={H.-J.},
           giveni={H\bibinithyphendelim J\bibinitperiod}}}%
        {{hash=c6d965df16e600c701d4bcc2888ae34c}{%
           family={Pintz},
           familyi={P\bibinitperiod},
           given={J.},
           giveni={J\bibinitperiod}}}%
      }
      \list{language}{1}{%
        {German}%
      }
      \strng{namehash}{8ee030157e0efba486591e59a6b5401b}
      \strng{fullhash}{8ee030157e0efba486591e59a6b5401b}
      \strng{bibnamehash}{8ee030157e0efba486591e59a6b5401b}
      \strng{authorbibnamehash}{8ee030157e0efba486591e59a6b5401b}
      \strng{authornamehash}{8ee030157e0efba486591e59a6b5401b}
      \strng{authorfullhash}{8ee030157e0efba486591e59a6b5401b}
      \field{extraname}{2}
      \field{sortinit}{B}
      \field{sortinithash}{8de16967003c7207dae369d874f1456e}
      \field{labelprefix}{A}
      \field{labelnamesource}{author}
      \field{labeltitlesource}{title}
      \field{annotation}{The authors give a shorter proof that $\displaystyle\lim_{x\to\infty}\sum_p \chi_{-4}(p)\log p\cdot\exp\left(-\frac{\log^2 p}{x}\right)=-\infty$.\par This article cites~\cites{MR534235,MR552462,MR1555148,MR0174538,MR1544293A,MR1544293B,MR0108470,1853.Chebyshev}.}
      \field{issn}{0378-6218}
      \field{journaltitle}{Resultate Math.}
      \field{number}{1}
      \field{title}{Über das {T}schebyschef-{P}roblem}
      \field{volume}{5}
      \field{year}{1982}
      \true{nocite}
      \field{pages}{1\bibrangedash 5}
      \range{pages}{5}
      \verb{doi}
      \verb 10.1007/bf03323296
      \endverb
      \verb{mrnumber}
      \verb 662791
      \endverb
      \verb{urlraw}
      \verb https://doi.org/10.1007/bf03323296
      \endverb
      \verb{url}
      \verb https://doi.org/10.1007/bf03323296
      \endverb
    \endentry
    \entry{MR591612}{article}{}
      \name{author}{2}{}{%
        {{hash=9f64a229f3d656a4a04d7eda97b0e30a}{%
           family={Bentz},
           familyi={B\bibinitperiod},
           given={Hans-J.},
           giveni={H\bibinithyphendelim J\bibinitperiod}}}%
        {{hash=d192c3b49f427e2494b8aaea251e390b}{%
           family={Pintz},
           familyi={P\bibinitperiod},
           given={János},
           giveni={J\bibinitperiod}}}%
      }
      \list{language}{1}{%
        {German}%
      }
      \strng{namehash}{e6be682d0e7c51c0cf93fc1a6932936c}
      \strng{fullhash}{e6be682d0e7c51c0cf93fc1a6932936c}
      \strng{bibnamehash}{e6be682d0e7c51c0cf93fc1a6932936c}
      \strng{authorbibnamehash}{e6be682d0e7c51c0cf93fc1a6932936c}
      \strng{authornamehash}{e6be682d0e7c51c0cf93fc1a6932936c}
      \strng{authorfullhash}{e6be682d0e7c51c0cf93fc1a6932936c}
      \field{extraname}{3}
      \field{sortinit}{B}
      \field{sortinithash}{8de16967003c7207dae369d874f1456e}
      \field{labelprefix}{A}
      \field{labelnamesource}{author}
      \field{labeltitlesource}{title}
      \field{annotation}{Chebyshev conjectured that $\lim_{x \to \infty} \pi_e(x,\chi_{-4}) = -\infty$, which is equivalent~\cites{MR1555148,MR1544293A, MR1544293B} to GRH for $L(s,\chi_{-4})$. Knapowski and Turán showed that these assertions are further equivalent to $\lim_{x \to \infty} \theta_l(x,r,\chi_{-4}) = -\infty$. In this article, the authors establish a related implication: Let $\chi$ be a quadratic character\mod q. Suppose that all nonreal zeros $\rho= \beta + i\gamma$ of $L(s,\chi_1)$ satisfy $\beta^2-\gamma^2 < 1/4$, a property implied by GRH$(\frac{\sqrt3}2)$. Then for $0 \leq \alpha < 1/2$ they prove that \[ \lim_{x \to \infty} \sum_{p} \chi_1(p) \frac{\log p}{p^\alpha} \cdot \exp\biggl( - \frac{\log^2 p}{x} \biggr) = -\infty.\]\par This article cites~\cites{MR595317,MR552462,MR534235,MR1555148,MR0174538,MR0182616A,MR0087685,MR0108470,1853.Chebyshev}.}
      \field{issn}{0025-5874}
      \field{journaltitle}{Math. Z.}
      \field{number}{1}
      \field{title}{Über eine Verallgemeinerung des Tschebyschef-Problems}
      \field{volume}{174}
      \field{year}{1980}
      \true{nocite}
      \field{pages}{35\bibrangedash 41}
      \range{pages}{7}
      \verb{doi}
      \verb 10.1007/BF01215079
      \endverb
      \verb{mrnumber}
      \verb 591612
      \endverb
    \endentry
    \entry{MR534235}{article}{}
      \name{author}{1}{}{%
        {{hash=687d9a8bcde5764db23c9a6841909f8e}{%
           family={Besenfelder},
           familyi={B\bibinitperiod},
           given={H.-J.},
           giveni={H\bibinithyphendelim J\bibinitperiod}}}%
      }
      \list{language}{1}{%
        {German}%
      }
      \strng{namehash}{687d9a8bcde5764db23c9a6841909f8e}
      \strng{fullhash}{687d9a8bcde5764db23c9a6841909f8e}
      \strng{bibnamehash}{687d9a8bcde5764db23c9a6841909f8e}
      \strng{authorbibnamehash}{687d9a8bcde5764db23c9a6841909f8e}
      \strng{authornamehash}{687d9a8bcde5764db23c9a6841909f8e}
      \strng{authorfullhash}{687d9a8bcde5764db23c9a6841909f8e}
      \field{extraname}{1}
      \field{sortinit}{B}
      \field{sortinithash}{8de16967003c7207dae369d874f1456e}
      \field{labelprefix}{A}
      \field{labelnamesource}{author}
      \field{labeltitlesource}{title}
      \field{annotation}{Using an existing explicit formula for general Mellin-transform pairs, the author shows that \begin{equation*} \begin{split} 2 \sqrt{\pi y}\sum_{\substack{0<\sigma<1 \\ L(\sigma+i\gamma,\chi_{-4})=0}} e^{y(\sigma-1/2+i\gamma)^2 }&=\log \frac4\pi -2 \sum_{n=1}^\infty \frac{\Lambda(n)\chi_{-4}(n)}{\sqrt n}e^{-{(\log n)^2}/4y}\\ & \qquad{}-C_0+2\int_0^{\infty} \frac{e^{-x^2/4y + x/2}-1}{1-e^{2x}} \,dx. \end{split} \end{equation*} From this identity, he proves unconditionally that $$\lim_{x\to\infty} \sum_p \chi_{-4}(p) \frac{\log p }{\sqrt p} e^{-(\log^2{p})/x} = -\infty.$$ (Note: the author, Hans--Joachim Besenfelder, soon changed his last name to Bentz and began to publish under that name.)\par This article cites~\cites{MR1555148, MR0174538, MR1544293A, 1853.Chebyshev}.}
      \field{issn}{0075-4102}
      \field{journaltitle}{J. Reine Angew. Math.}
      \field{title}{Über eine {V}ermutung von {T}schebyschef. {I}}
      \field{volume}{307/308}
      \field{year}{1979}
      \true{nocite}
      \field{pages}{411\bibrangedash 417}
      \range{pages}{7}
      \verb{doi}
      \verb 10.1515/crll.1979.307-308.411
      \endverb
      \verb{mrnumber}
      \verb 534235
      \endverb
      \verb{urlraw}
      \verb https://doi.org/10.1515/crll.1979.307-308.411
      \endverb
      \verb{url}
      \verb https://doi.org/10.1515/crll.1979.307-308.411
      \endverb
    \endentry
    \entry{MR552462}{article}{}
      \name{author}{1}{}{%
        {{hash=687d9a8bcde5764db23c9a6841909f8e}{%
           family={Besenfelder},
           familyi={B\bibinitperiod},
           given={H.-J.},
           giveni={H\bibinithyphendelim J\bibinitperiod}}}%
      }
      \list{language}{1}{%
        {German}%
      }
      \strng{namehash}{687d9a8bcde5764db23c9a6841909f8e}
      \strng{fullhash}{687d9a8bcde5764db23c9a6841909f8e}
      \strng{bibnamehash}{687d9a8bcde5764db23c9a6841909f8e}
      \strng{authorbibnamehash}{687d9a8bcde5764db23c9a6841909f8e}
      \strng{authornamehash}{687d9a8bcde5764db23c9a6841909f8e}
      \strng{authorfullhash}{687d9a8bcde5764db23c9a6841909f8e}
      \field{extraname}{2}
      \field{sortinit}{B}
      \field{sortinithash}{8de16967003c7207dae369d874f1456e}
      \field{labelprefix}{A}
      \field{labelnamesource}{author}
      \field{labeltitlesource}{title}
      \field{annotation}{The author proves unconditionally that for all $0\le\alpha\le\frac12$, $$\lim_{x\to\infty} \sum_{p} \frac{\chi_{-4} (p) \log p}{p^{\alpha}} e^{-(\log^2 p)/x} = -\infty,$$ using an explicit formula similar to the one in his prior work~\cite{MR534235}.\par This article cites~\cites{MR1555148, MR0174538, MR1544293A, MR0300853, MR0300854,1853.Chebyshev, MR534235, MR0156826, MR1544293B}.}
      \field{issn}{0075-4102}
      \field{journaltitle}{J. Reine Angew. Math.}
      \field{title}{Über eine {V}ermutung von {T}schebyschef. {II}}
      \field{volume}{313}
      \field{year}{1980}
      \true{nocite}
      \field{pages}{52\bibrangedash 58}
      \range{pages}{7}
      \verb{doi}
      \verb 10.1515/crll.1980.313.52
      \endverb
      \verb{mrnumber}
      \verb 552462
      \endverb
      \verb{urlraw}
      \verb https://doi.org/10.1515/crll.1980.313.52
      \endverb
      \verb{url}
      \verb https://doi.org/10.1515/crll.1980.313.52
      \endverb
    \endentry
    \entry{MR3335903}{article}{}
      \name{author}{2}{}{%
        {{hash=251a94d49687997eac18f56ef1823028}{%
           family={Best},
           familyi={B\bibinitperiod},
           given={D.\bibnamedelimi G.},
           giveni={D\bibinitperiod\bibinitdelim G\bibinitperiod}}}%
        {{hash=f8faf70b3f8ac0edff4b9b5337e687bf}{%
           family={Trudgian},
           familyi={T\bibinitperiod},
           given={T.\bibnamedelimi S.},
           giveni={T\bibinitperiod\bibinitdelim S\bibinitperiod}}}%
      }
      \strng{namehash}{4ff3faf470a39e73d27ec82924abaeb8}
      \strng{fullhash}{4ff3faf470a39e73d27ec82924abaeb8}
      \strng{bibnamehash}{4ff3faf470a39e73d27ec82924abaeb8}
      \strng{authorbibnamehash}{4ff3faf470a39e73d27ec82924abaeb8}
      \strng{authornamehash}{4ff3faf470a39e73d27ec82924abaeb8}
      \strng{authorfullhash}{4ff3faf470a39e73d27ec82924abaeb8}
      \field{sortinit}{B}
      \field{sortinithash}{8de16967003c7207dae369d874f1456e}
      \field{labelprefix}{A}
      \field{labelnamesource}{author}
      \field{labeltitlesource}{title}
      \field{annotation}{Using the method of Grosswald~\cite{MR0332685}, the authors prove that \[ \liminf_{x \to \infty} E^M(x) < -1.6383 \quad \text{and} \quad \limsup_{x \to \infty} E^M(x) > 1.6383. \] The proof uses a refined variant of the LLL-algorithm and a kernel function used by Odlyzko and te Riele~\cite{MR783538}, as well as several techniques for speeding up the computation.\par In the appendix of the article, the authors list new lower bounds on $m$ for the ``$m$-dependence'' of the first $n$ positive ordinates of nontrivial zeros of $\zeta(s)$, for $n \leq 500$. This numerical evidence, combined with Grosswald's theorem, provides new proofs that $E^\pi(x)$, $L(x)$, and $L_r(x)$ change signs infinitely often.\par This article cites~\cite{MR429789,MR2282922,MR783538,MR0006202,MR0104638,MR0332685,MR0332684,MR0330069,MR654520}.}
      \field{issn}{0025-5718}
      \field{journaltitle}{Math. Comp.}
      \field{number}{294}
      \field{title}{Linear relations of zeroes of the zeta-function}
      \field{volume}{84}
      \field{year}{2015}
      \true{nocite}
      \field{pages}{2047\bibrangedash 2058}
      \range{pages}{12}
      \verb{doi}
      \verb 10.1090/S0025-5718-2014-02916-5
      \endverb
      \verb{mrnumber}
      \verb 3335903
      \endverb
      \verb{urlraw}
      \verb https://doi.org/10.1090/S0025-5718-2014-02916-5
      \endverb
      \verb{url}
      \verb https://doi.org/10.1090/S0025-5718-2014-02916-5
      \endverb
    \endentry
    \entry{MR3509711}{article}{}
      \name{author}{3}{}{%
        {{hash=922e48f9d159c924a55b336dfeb65261}{%
           family={Bhowmik},
           familyi={B\bibinitperiod},
           given={Gautami},
           giveni={G\bibinitperiod}}}%
        {{hash=d6ae4b9844776294c92ff23952e91b53}{%
           family={Ramaré},
           familyi={R\bibinitperiod},
           given={Olivier},
           giveni={O\bibinitperiod}}}%
        {{hash=dd239e78c668d9270b067bfae5530a1f}{%
           family={Schlage--Puchta},
           familyi={S\bibinithyphendelim \bibinithyphendelim P\bibinitperiod},
           given={Jan-Christoph},
           giveni={J\bibinithyphendelim C\bibinitperiod}}}%
      }
      \list{language}{2}{%
        {English, with English}%
        {French summaries}%
      }
      \strng{namehash}{83e0fb0de8538ce919bef0bbac9f375f}
      \strng{fullhash}{83e0fb0de8538ce919bef0bbac9f375f}
      \strng{bibnamehash}{83e0fb0de8538ce919bef0bbac9f375f}
      \strng{authorbibnamehash}{83e0fb0de8538ce919bef0bbac9f375f}
      \strng{authornamehash}{83e0fb0de8538ce919bef0bbac9f375f}
      \strng{authorfullhash}{83e0fb0de8538ce919bef0bbac9f375f}
      \field{sortinit}{B}
      \field{sortinithash}{8de16967003c7207dae369d874f1456e}
      \field{labelprefix}{A}
      \field{labelnamesource}{author}
      \field{labeltitlesource}{title}
      \field{annotation}{The authors establish an effective version of Landau's theorem for Dirichlet integrals, which they apply to the function $G_k(n) = \sum_{n_1 + \cdots n_k = n} \Lambda(n_1) \cdots \Lambda(n_k)$ counting the weighted number of representations of $n$ as the sum of $k$ primes. Further defining \[ \Delta_k(x) = \sum_{n\le x} G_k(n) - \frac{x^k}{k!} + k \sum_\rho \frac{x^{\rho+k-1}}{\rho(\rho+1)\cdots(\rho+k-1)}, \] they show under RH that $\Delta_2(n) = \Omega_\pm(x)$ and, for particular constants $c_k$, that $\Delta_k(n) - c_kx^{k-1} = \Omega_\pm(x^{k-1})$ for $k\ge3$.\par This article cites~\cites{MR0146156,MR2040112}.}
      \field{issn}{1246-7405}
      \field{journaltitle}{J. Théor. Nombres Bordeaux}
      \field{number}{2}
      \field{title}{Tauberian oscillation theorems and the distribution of {G}oldbach numbers}
      \field{volume}{28}
      \field{year}{2016}
      \true{nocite}
      \field{pages}{291\bibrangedash 299}
      \range{pages}{9}
      \verb{mrnumber}
      \verb 3509711
      \endverb
    \endentry
    \entry{MR2398787}{article}{}
      \name{author}{3}{}{%
        {{hash=86e7b84d8adeabc3f69abb52df304f6a}{%
           family={Borwein},
           familyi={B\bibinitperiod},
           given={Peter},
           giveni={P\bibinitperiod}}}%
        {{hash=a15eabe7e88a5dca9dbe96e2754e374a}{%
           family={Ferguson},
           familyi={F\bibinitperiod},
           given={Ron},
           giveni={R\bibinitperiod}}}%
        {{hash=454d3dedc781ec4a49112de10313b411}{%
           family={Mossinghoff},
           familyi={M\bibinitperiod},
           given={Michael\bibnamedelima J.},
           giveni={M\bibinitperiod\bibinitdelim J\bibinitperiod}}}%
      }
      \strng{namehash}{96894e6dfcaffbe33c893760d2cc6cc2}
      \strng{fullhash}{96894e6dfcaffbe33c893760d2cc6cc2}
      \strng{bibnamehash}{96894e6dfcaffbe33c893760d2cc6cc2}
      \strng{authorbibnamehash}{96894e6dfcaffbe33c893760d2cc6cc2}
      \strng{authornamehash}{96894e6dfcaffbe33c893760d2cc6cc2}
      \strng{authorfullhash}{96894e6dfcaffbe33c893760d2cc6cc2}
      \field{sortinit}{B}
      \field{sortinithash}{8de16967003c7207dae369d874f1456e}
      \field{labelprefix}{A}
      \field{labelnamesource}{author}
      \field{labeltitlesource}{title}
      \field{annotation}{The authors develop refined algorithms for computing the sums $L(x)$ and $L_r(x)$. They show that the smallest positive integer~$n$ for which $L_r(n)<0$ (that is, the first witness to the negative resolution of Turán's problem) is $72,185,376,951,205$; they also describe the many sign changes of $L_r(x)$ among the subsequent integers, and suggest that its next sign change occurs near $1.16\cdot10^{19}$. Similar results are obtained for $L(n)$, including a new region of sign changes starting at $351,100,332,278,253$; their calculation allows them to conclude that $L(n) > 0.061867\sqrt n$ infinitely often.\par This article \cites{MR27305,MR654520,MR0120198,MR0006202,MR0104638,MR0003644,MR584557}.}
      \field{issn}{0025-5718}
      \field{journaltitle}{Math. Comp.}
      \field{number}{263}
      \field{title}{Sign changes in sums of the {L}iouville function}
      \field{volume}{77}
      \field{year}{2008}
      \true{nocite}
      \field{pages}{1681\bibrangedash 1694}
      \range{pages}{14}
      \verb{doi}
      \verb 10.1090/S0025-5718-08-02036-X
      \endverb
      \verb{mrnumber}
      \verb 2398787
      \endverb
    \endentry
    \entry{MR0369287}{article}{}
      \name{author}{1}{}{%
        {{hash=a6d7fefac7d12a29485ec824a06cdbc6}{%
           family={Brent},
           familyi={B\bibinitperiod},
           given={Richard\bibnamedelima P.},
           giveni={R\bibinitperiod\bibinitdelim P\bibinitperiod}}}%
      }
      \strng{namehash}{a6d7fefac7d12a29485ec824a06cdbc6}
      \strng{fullhash}{a6d7fefac7d12a29485ec824a06cdbc6}
      \strng{bibnamehash}{a6d7fefac7d12a29485ec824a06cdbc6}
      \strng{authorbibnamehash}{a6d7fefac7d12a29485ec824a06cdbc6}
      \strng{authornamehash}{a6d7fefac7d12a29485ec824a06cdbc6}
      \strng{authorfullhash}{a6d7fefac7d12a29485ec824a06cdbc6}
      \field{sortinit}{B}
      \field{sortinithash}{8de16967003c7207dae369d874f1456e}
      \field{labelprefix}{A}
      \field{labelnamesource}{author}
      \field{labeltitlesource}{title}
      \field{annotation}{The author reports on extensive calculations of $\pi(x)$ and of the counting function for twin primes. He suggests from these computations that $E^\pi(x)$ has a limiting distribution with mean $-1$ and standard deviation approximately $0.21$ (which differs by less than $3\%$ from the true value under RH and LI). He verifies that $\pi(x)\le\li(x)$ for $x\le 8\times10^{10}$, and predicts from the observed distribution above that $\pi(x)\le\li(x)$ for~$x$ up to at least $10^{100}$. Shanks~\cite{MR0108470} conjectured that $\frac1x \int_0^x E^\pi(t)\, dt$ tends to~$-1$; the author suggests from his computations that this conjecture is false but that its logarithmic counterpart $\lim_{x\to\infty} \frac1{\log x} \int_1^x E^\pi(t)\, \frac{dt}t = -1$ should be true.\par The author also shows that the error term in the counting function of twin primes does not exceed $2.3x^{1/2}/\log x$ in absolute value for $x\leq 8\times 10^{10}$. He also gives the heuristic estimate $1.9021604\pm5\times10^{-7}$ for Brun's constant $\sum(\frac1q+\frac1{q+2})$ (where the sum is over twin primes~$q$), although no rigorous upper bound is given.\par This article cites~\cites{MR1555148,1932.Ingham, MR0202686,1914.Littlewood,MR0108470,MR0067145}.}
      \field{issn}{0025-5718}
      \field{journaltitle}{Math. Comp.}
      \field{note}{Collection of articles dedicated to Derrick Henry Lehmer on the occasion of his seventieth birthday}
      \field{title}{Irregularities in the distribution of primes and twin primes}
      \field{volume}{29}
      \field{year}{1975}
      \true{nocite}
      \field{pages}{43\bibrangedash 56}
      \range{pages}{14}
      \verb{doi}
      \verb 10.2307/2005460
      \endverb
      \verb{mrnumber}
      \verb 0369287
      \endverb
    \endentry
    \entry{MR3414306}{article}{}
      \name{author}{1}{}{%
        {{hash=b58171fe4e01fa0b55ff6f0dcd610186}{%
           family={Büthe},
           familyi={B\bibinitperiod},
           given={J.},
           giveni={J\bibinitperiod}}}%
      }
      \strng{namehash}{b58171fe4e01fa0b55ff6f0dcd610186}
      \strng{fullhash}{b58171fe4e01fa0b55ff6f0dcd610186}
      \strng{bibnamehash}{b58171fe4e01fa0b55ff6f0dcd610186}
      \strng{authorbibnamehash}{b58171fe4e01fa0b55ff6f0dcd610186}
      \strng{authornamehash}{b58171fe4e01fa0b55ff6f0dcd610186}
      \strng{authorfullhash}{b58171fe4e01fa0b55ff6f0dcd610186}
      \field{sortinit}{B}
      \field{sortinithash}{8de16967003c7207dae369d874f1456e}
      \field{labelprefix}{A}
      \field{labelnamesource}{author}
      \field{labeltitlesource}{title}
      \field{annotation}{Rosser and Schoenfeld \cite{MR0137689} observed that $\Delta^{\pi_r}(x) \ge 0$ for $1\le x\le 10^8$. The author proves that there exists an $x_0 \approx 1.91\times 10^{215}$ such that $\Delta^{\pi_r}(x) < 0$ for $x \in [x_0-6\times 10^{103},x_0]$. The proof is an adaptation of a method of Lehman~\cite{MR0202686} that improved upon Skewes's number, though with a different choice of kernel function.\par This article cites~\cites{MR0202686, MR729173, MR3356033, MR1573970, MR0137689}.}
      \field{issn}{0065-1036}
      \field{journaltitle}{Acta Arith.}
      \field{number}{2}
      \field{title}{On the first sign change in {M}ertens' theorem}
      \field{volume}{171}
      \field{year}{2015}
      \true{nocite}
      \field{pages}{183\bibrangedash 195}
      \range{pages}{13}
      \verb{doi}
      \verb 10.4064/aa171-2-5
      \endverb
      \verb{mrnumber}
      \verb 3414306
      \endverb
      \verb{urlraw}
      \verb https://doi.org/10.4064/aa171-2-5
      \endverb
      \verb{url}
      \verb https://doi.org/10.4064/aa171-2-5
      \endverb
    \endentry
    \entry{MR2811555}{article}{}
      \name{author}{2}{}{%
        {{hash=3417f8b5237d274ddd67adbd43553547}{%
           family={Cha},
           familyi={C\bibinitperiod},
           given={B},
           giveni={B\bibinitperiod}}}%
        {{hash=5fe0c1b9c49f664236b5bbd12a052411}{%
           family={Im},
           familyi={I\bibinitperiod},
           given={B.-H.},
           giveni={B\bibinithyphendelim H\bibinitperiod}}}%
      }
      \strng{namehash}{ea08ea1554b13d2b20e9dca09ca0fd3e}
      \strng{fullhash}{ea08ea1554b13d2b20e9dca09ca0fd3e}
      \strng{bibnamehash}{ea08ea1554b13d2b20e9dca09ca0fd3e}
      \strng{authorbibnamehash}{ea08ea1554b13d2b20e9dca09ca0fd3e}
      \strng{authornamehash}{ea08ea1554b13d2b20e9dca09ca0fd3e}
      \strng{authorfullhash}{ea08ea1554b13d2b20e9dca09ca0fd3e}
      \field{sortinit}{C}
      \field{sortinithash}{4c244ceae61406cdc0cc2ce1cb1ff703}
      \field{labelprefix}{A}
      \field{labelnamesource}{author}
      \field{labeltitlesource}{title}
      \field{annotation}{The authors study Chebyshev's bias in a finite Galois extension~$F$ of a global function field~$K$. Let $\pi_C(N)$ be the number of prime elements in~$K$ of degree~$N$ whose Frobenius lies in~$C$. Following the strategy of Rubinstein and Sarnak~\cite{MR1329368}, the authors study the limiting distribution of $\frac{X}{q^{X/2}}\sum_{N=1}^X \frac{|G|}{|C|} \big( \pi_c(N)-\pi(N) \big)$. The authors show that when $F/K$ is geometric and satisfies LI, there is a bias towards conjugacy class of the Galois group containing fewer square elements, as in the number field case. The authors give a complete description of the prime element race in the case of scalar field extensions of~$K$.\par This article cites~\cite{MR1329368,MR2202918,MR1806291,MR2600420,MR2474313}.}
      \field{issn}{0022-314X}
      \field{journaltitle}{J. Number Theory}
      \field{number}{10}
      \field{title}{{C}hebyshev's bias in {G}alois extensions of global function fields}
      \field{volume}{131}
      \field{year}{2011}
      \true{nocite}
      \field{pages}{1875\bibrangedash 1886}
      \range{pages}{12}
      \verb{doi}
      \verb 10.1016/j.jnt.2011.03.011
      \endverb
      \verb{mrnumber}
      \verb 2811555
      \endverb
    \endentry
    \entry{MR3581815}{article}{}
      \name{author}{3}{}{%
        {{hash=3417f8b5237d274ddd67adbd43553547}{%
           family={Cha},
           familyi={C\bibinitperiod},
           given={B.},
           giveni={B\bibinitperiod}}}%
        {{hash=889ecd9d3a804c5a7e9a5d8287766657}{%
           family={Fiorilli},
           familyi={F\bibinitperiod},
           given={D.},
           giveni={D\bibinitperiod}}}%
        {{hash=dfca5f0fbbbf59635f02988f421953ae}{%
           family={Jouve},
           familyi={J\bibinitperiod},
           given={F.},
           giveni={F\bibinitperiod}}}%
      }
      \strng{namehash}{5cfc0c94602fe7b3e73bd312fbce4b01}
      \strng{fullhash}{5cfc0c94602fe7b3e73bd312fbce4b01}
      \strng{bibnamehash}{5cfc0c94602fe7b3e73bd312fbce4b01}
      \strng{authorbibnamehash}{5cfc0c94602fe7b3e73bd312fbce4b01}
      \strng{authornamehash}{5cfc0c94602fe7b3e73bd312fbce4b01}
      \strng{authorfullhash}{5cfc0c94602fe7b3e73bd312fbce4b01}
      \field{sortinit}{C}
      \field{sortinithash}{4c244ceae61406cdc0cc2ce1cb1ff703}
      \field{labelprefix}{A}
      \field{labelnamesource}{author}
      \field{labeltitlesource}{title}
      \field{annotation}{In this article, the authors describe the prime number race for elliptic curves over the function field of a proper, smooth, and geometrically connected curve over a finite field. Let ${E/K}$ be a family of elliptic curves of unbounded conductor for which $L(E/K,T)$ satisfies LI (which can be confirmed in some cases) and such that $\text{rank}(E/K)=o(\sqrt{N_{E/K}})$ as $N_{E/K}\to\infty$. Then, the random variable ${\frac{\sqrt q-1}{\sqrt q}}X_E/\sqrt{N_{E/K}}$ converges in distribution to the standard Gaussian as $N_{E/K}\to\infty$. As a consequence, $\delta(E)$ (the proportion of those~$x$ for which more primes up to~$x$ have positive trace of Frobenius $a_p(E)$ than negative) tends to $\frac 12$.\par Moreover, the authors also study the behaviour of the function $T_E(X)= -{X}{q^{-X/2}}\sum_{\deg(v)\leq X} 2\cos{\theta_v}$ associated to the elliptic curves of Ulmer's family of elliptic curves over $\Bbb F_q[t]$. They discuss the cases of extreme bias and moderate bias for Ulmer's family. Moreover, through proving a central limit theorem, the authors shows that $\delta(E_f)-\frac 12 = \Omega(\frac1{\sqrt{d}})$ for quadratic twists~$E_f$.\par This article cites~\cites{MR3261965,MR2474313,MR3264673,MR3272280,MR3028758,MR1329368,2007.Sarnak}.}
      \field{issn}{0012-9593}
      \field{journaltitle}{Ann. Sci. Éc. Norm. Supér. (4)}
      \field{number}{5}
      \field{title}{Prime number races for elliptic curves over function fields}
      \field{volume}{49}
      \field{year}{2016}
      \true{nocite}
      \field{pages}{1239\bibrangedash 1277}
      \range{pages}{39}
      \verb{doi}
      \verb 10.24033/asens.2308
      \endverb
      \verb{mrnumber}
      \verb 3581815
      \endverb
      \verb{urlraw}
      \verb https://doi.org/10.24033/asens.2308
      \endverb
      \verb{url}
      \verb https://doi.org/10.24033/asens.2308
      \endverb
    \endentry
    \entry{MR2474313}{article}{}
      \name{author}{1}{}{%
        {{hash=7871a490afe3f3963421804f312b0ad7}{%
           family={Cha},
           familyi={C\bibinitperiod},
           given={Byungchul},
           giveni={B\bibinitperiod}}}%
      }
      \strng{namehash}{7871a490afe3f3963421804f312b0ad7}
      \strng{fullhash}{7871a490afe3f3963421804f312b0ad7}
      \strng{bibnamehash}{7871a490afe3f3963421804f312b0ad7}
      \strng{authorbibnamehash}{7871a490afe3f3963421804f312b0ad7}
      \strng{authornamehash}{7871a490afe3f3963421804f312b0ad7}
      \strng{authorfullhash}{7871a490afe3f3963421804f312b0ad7}
      \field{extraname}{1}
      \field{sortinit}{C}
      \field{sortinithash}{4c244ceae61406cdc0cc2ce1cb1ff703}
      \field{labelprefix}{A}
      \field{labelnamesource}{author}
      \field{labeltitlesource}{title}
      \field{annotation}{The author adapts the arguments of Rubinstein and Sarnak~\cite{MR1329368} to the function field setting, establishing (under the appropriate generalization of LI) the existence of limiting distributions and their symmetries, a bias towards quadratic nonresidues in two-way races, and central limit theorems. The author also provides some examples where LI is violated and the expected biases are overridden, as well as an example where LI can actually be verified.\par This article cites~\cites{1853.Chebyshev,MR2202918,MR1329368}.}
      \field{issn}{0010-437X}
      \field{journaltitle}{Compos. Math.}
      \field{number}{6}
      \field{title}{{C}hebyshev's bias in function fields}
      \field{volume}{144}
      \field{year}{2008}
      \true{nocite}
      \field{pages}{1351\bibrangedash 1374}
      \range{pages}{24}
      \verb{doi}
      \verb 10.1112/S0010437X08003631
      \endverb
      \verb{mrnumber}
      \verb 2474313
      \endverb
    \endentry
    \entry{MR3684399}{article}{}
      \name{author}{1}{}{%
        {{hash=7871a490afe3f3963421804f312b0ad7}{%
           family={Cha},
           familyi={C\bibinitperiod},
           given={Byungchul},
           giveni={B\bibinitperiod}}}%
      }
      \strng{namehash}{7871a490afe3f3963421804f312b0ad7}
      \strng{fullhash}{7871a490afe3f3963421804f312b0ad7}
      \strng{bibnamehash}{7871a490afe3f3963421804f312b0ad7}
      \strng{authorbibnamehash}{7871a490afe3f3963421804f312b0ad7}
      \strng{authornamehash}{7871a490afe3f3963421804f312b0ad7}
      \strng{authorfullhash}{7871a490afe3f3963421804f312b0ad7}
      \field{extraname}{2}
      \field{sortinit}{C}
      \field{sortinithash}{4c244ceae61406cdc0cc2ce1cb1ff703}
      \field{labelprefix}{A}
      \field{labelnamesource}{author}
      \field{labeltitlesource}{title}
      \field{annotation}{In this article, the author investigates an analogue of the Mertens sum for function fields. Let~$C$ be a nonsingular projective curve defined over a finite field $\mathbb{F}_q$ of characteristic $p>2$. The author defines the {M}öbius function $\mu_{C/\mathbb{F}_q}(D)$ of $C/\mathbb{F}_q$ for all effective divisors $D$ of $C$ to be $$ \mu_{C/\mathbb{F}_q}(D) = \begin{cases} 1, & \text{if $D=0$},\\ 0, & \text{if a prime divisor divides $D$ with order at least 2},\\ (-1)^t, & \text{if $D$ is a sum of $t$ distinct prime divisors}. \end{cases}$$ Let $c_{\mu}(N)= \sum_{\deg(D)=N} \mu_{C/\mathbb{F}_q}(D)$ and $M_{C/\mathbb{F}_q}(X)= \sum_{\deg D \le X} \mu_{C/\mathbb{F}_q}(D)$. The author justifies these new definitions by showing that the Dirichlet series $Z_\mu(u)$ associated with $\mu_{C/\mathbb{F}_q}(D)$ for a divisor~$D$ is closely related to the the zeta function $Z_{C/\mathbb{F}_q}(u)$ associated with the curve~$C$: $$ \frac{1}{Z_{C/\mathbb{F}_q}(u)}=Z_\mu(u)= \sum_{D \geq 0}\frac{\mu_{C/\mathbb{F}_q}(D)}{\mathcal{N}D^s} = \sum_{N=0}^\infty c_\mu(N)u^N, $$ where $u=q^{-s}$ and $\mathcal{N}D$ is the absolute norm of~$D$, so that $M_{C/\mathbb{F}_q}(X)=\sum_{N=0}^\infty c_\mu(N)$. By working with the coefficients, the author shows that $M_{C/\mathbb{F}_q}(X) \ll X^{2g-1}q^{X/2}$.\par Let $\gamma_1 = \sqrt{q}e^{i \theta_1} , \ldots, \gamma_{2g} = \sqrt{q}e^{i \theta_{2g}}$ be the inverse zeros of $Z_{C/\mathbb{F}_q}(u)$, where~$g$ is the genus of~$C$. The curve~$C$ is said to satisfy LI if the set $\{\theta_j \colon 0 \le \theta_j \le \pi \} \cup \{\pi\}$ is linearly independent over~$\mathbb{Q}$. The author shows that LI implies $$ D(C/\mathbb{F}_q) = \limsup_{X \to \infty} \frac{|M_{C/\mathbb{F}_q}(X)|}{q^{X/2}}=\sum_{j=1}^{2g} \left| \frac{\gamma_j}{Z_{C/\mathbb{F}_q}'(\gamma_j^{-1})}\frac{\gamma_j}{\gamma_j-1} \right|<\infty, $$ while $D(C/\mathbb{F}_q)=\infty$ if there is a zero with multiple order. The author also shows that a family of hyperelliptic curves satisfy LI and computes the geometric average of $D(C/\mathbb{F}_q)$ over this family using an equidistribution theorem due to Deligne.\par This article cites~\cites{MR2474313,MR2811555,MR3189983,MR0006202,1897.Mertens,MR2078705,MR783538,MR1329368}.}
      \field{issn}{0065-1036}
      \field{journaltitle}{Acta Arith.}
      \field{number}{4}
      \field{title}{The summatory function of the {M}öbius function in function fields}
      \field{volume}{179}
      \field{year}{2017}
      \true{nocite}
      \field{pages}{375\bibrangedash 395}
      \range{pages}{21}
      \verb{doi}
      \verb 10.4064/aa8590-1-2017
      \endverb
      \verb{mrnumber}
      \verb 3684399
      \endverb
      \verb{urlraw}
      \verb https://doi.org/10.4064/aa8590-1-2017
      \endverb
      \verb{url}
      \verb https://doi.org/10.4064/aa8590-1-2017
      \endverb
    \endentry
    \entry{MR2600420}{article}{}
      \name{author}{2}{}{%
        {{hash=7871a490afe3f3963421804f312b0ad7}{%
           family={Cha},
           familyi={C\bibinitperiod},
           given={Byungchul},
           giveni={B\bibinitperiod}}}%
        {{hash=1470ddce431ad26e652280417656b586}{%
           family={Kim},
           familyi={K\bibinitperiod},
           given={Seick},
           giveni={S\bibinitperiod}}}%
      }
      \strng{namehash}{6505b468496e6694d7b4710ddb7b703d}
      \strng{fullhash}{6505b468496e6694d7b4710ddb7b703d}
      \strng{bibnamehash}{6505b468496e6694d7b4710ddb7b703d}
      \strng{authorbibnamehash}{6505b468496e6694d7b4710ddb7b703d}
      \strng{authornamehash}{6505b468496e6694d7b4710ddb7b703d}
      \strng{authorfullhash}{6505b468496e6694d7b4710ddb7b703d}
      \field{sortinit}{C}
      \field{sortinithash}{4c244ceae61406cdc0cc2ce1cb1ff703}
      \field{labelprefix}{A}
      \field{labelnamesource}{author}
      \field{labeltitlesource}{title}
      \field{annotation}{The authors derive a formula, in the function field setting, for the logarithmic densities of races among the counting functions of irreducible polynomials of degree up to~$X$ in residue classes. The proofs follow~\cite{MR1329368} closely, although the authors introduce a smoothing function on the Fourier side to overcome the problem that the Fourier transforms decay slowly due to the finiteness of the number of zeros of the relevant $L$-functions.\par This article cites ~\cites{MR1329368,MR2474313,MR1806291,MR2202918}.}
      \field{issn}{0022-314X}
      \field{journaltitle}{J. Number Theory}
      \field{number}{4}
      \field{title}{Biases in the prime number race of function fields}
      \field{volume}{130}
      \field{year}{2010}
      \true{nocite}
      \field{pages}{1048\bibrangedash 1055}
      \range{pages}{8}
      \verb{doi}
      \verb 10.1016/j.jnt.2009.09.015
      \endverb
      \verb{mrnumber}
      \verb 2600420
      \endverb
    \endentry
    \entry{MR2652902}{article}{}
      \name{author}{2}{}{%
        {{hash=69f328db05bb9b15874d29cfcf79dfdb}{%
           family={Chao},
           familyi={C\bibinitperiod},
           given={Kuok\bibnamedelima Fai},
           giveni={K\bibinitperiod\bibinitdelim F\bibinitperiod}}}%
        {{hash=81a915ac987d508458c06ee0891cc3c2}{%
           family={Plymen},
           familyi={P\bibinitperiod},
           given={Roger},
           giveni={R\bibinitperiod}}}%
      }
      \strng{namehash}{0737652978438d918ff3cb160495ab49}
      \strng{fullhash}{0737652978438d918ff3cb160495ab49}
      \strng{bibnamehash}{0737652978438d918ff3cb160495ab49}
      \strng{authorbibnamehash}{0737652978438d918ff3cb160495ab49}
      \strng{authornamehash}{0737652978438d918ff3cb160495ab49}
      \strng{authorfullhash}{0737652978438d918ff3cb160495ab49}
      \field{sortinit}{C}
      \field{sortinithash}{4c244ceae61406cdc0cc2ce1cb1ff703}
      \field{labelprefix}{A}
      \field{labelnamesource}{author}
      \field{labeltitlesource}{title}
      \field{annotation}{The authors modify Lehman's theorem~\cite{MR0202686} by improving the bound on $\theta_1(x) = \frac{2\log x}3 \big( \frac{\pi(x)}{x\log x}-1)$ using an inequality of Panaitopol. Together with $2\times10^6$ zeros of $\zeta(s)$ (supplied by Odlyzko), they show that there exist at least $10^{154}$ consecutive integers $x$ in the range $[1.3978965\times 10^{316}, 1.398344\times 10^{316}]$ for which $\pi(x)>\li(x)$.\par This article cites~\cites{MR1752093,MR2420864,MR0202686,1914.Littlewood,MR866118,MR0137689,MR1573970,MR0067145}.}
      \field{issn}{1793-0421}
      \field{journaltitle}{Int. J. Number Theory}
      \field{number}{3}
      \field{title}{A new bound for the smallest {$x$} with {$\pi (x)>\li (x)$}}
      \field{volume}{6}
      \field{year}{2010}
      \true{nocite}
      \field{pages}{681\bibrangedash 690}
      \range{pages}{10}
      \verb{doi}
      \verb 10.1142/S1793042110003125
      \endverb
      \verb{mrnumber}
      \verb 2652902
      \endverb
    \endentry
    \entry{MR3306734}{article}{}
      \name{author}{3}{}{%
        {{hash=550b5a0b5fd42a70ac3723323eef8887}{%
           family={Chaubey},
           familyi={C\bibinitperiod},
           given={Sneha},
           giveni={S\bibinitperiod}}}%
        {{hash=259966712b379a8c7cdc8ece86803924}{%
           family={Lanius},
           familyi={L\bibinitperiod},
           given={Melinda},
           giveni={M\bibinitperiod}}}%
        {{hash=db9b605a4c37f759795e5a1699044f16}{%
           family={Zaharescu},
           familyi={Z\bibinitperiod},
           given={Alexandru},
           giveni={A\bibinitperiod}}}%
      }
      \strng{namehash}{f307e9735e20f51a6fae94d9f56d696c}
      \strng{fullhash}{f307e9735e20f51a6fae94d9f56d696c}
      \strng{bibnamehash}{f307e9735e20f51a6fae94d9f56d696c}
      \strng{authorbibnamehash}{f307e9735e20f51a6fae94d9f56d696c}
      \strng{authornamehash}{f307e9735e20f51a6fae94d9f56d696c}
      \strng{authorfullhash}{f307e9735e20f51a6fae94d9f56d696c}
      \field{sortinit}{C}
      \field{sortinithash}{4c244ceae61406cdc0cc2ce1cb1ff703}
      \field{labelprefix}{A}
      \field{labelnamesource}{author}
      \field{labeltitlesource}{title}
      \field{annotation}{Atanassov defined the ``irrational factor'' of the number $n=p_1^{r_1}\cdots p_k^{r_k}$ to be $I(n) = p_1^{1/r_1}\cdots p_k^{1/r_k}$. The authors establish an asymptotic formula for the summatory function of $I(n)$ over a reduced residue class. They then use a Landau-type argument to show that the race between the summatory functions of $I(n)$ over integers that are $1\mod 3$ and integers that are $2\mod 3$ has infinitely many lead changes; they remark that a theorem from~\cite{MR2506435} shows that the number of sign changes in $[1,T]$ is $\gg\log T$ when~$T$ is sufficiently large.\par This article cites~\cites{1914.Littlewood,MR0146156,MR858395,MR1511360,MR2506435}.}
      \field{issn}{0253-4142}
      \field{journaltitle}{Proc. Indian Acad. Sci. Math. Sci.}
      \field{number}{4}
      \field{title}{Irrational factor races}
      \field{volume}{124}
      \field{year}{2014}
      \true{nocite}
      \field{pages}{471\bibrangedash 479}
      \range{pages}{9}
      \verb{doi}
      \verb 10.1007/s12044-014-0198-z
      \endverb
      \verb{mrnumber}
      \verb 3306734
      \endverb
    \endentry
    \entry{1853.Chebyshev}{article}{}
      \name{author}{1}{}{%
        {{hash=0fa981e305ea35073dc938b976cda0e1}{%
           family={Chebyshev},
           familyi={C\bibinitperiod},
           given={P.},
           giveni={P\bibinitperiod}}}%
      }
      \list{language}{1}{%
        {French}%
      }
      \strng{namehash}{0fa981e305ea35073dc938b976cda0e1}
      \strng{fullhash}{0fa981e305ea35073dc938b976cda0e1}
      \strng{bibnamehash}{0fa981e305ea35073dc938b976cda0e1}
      \strng{authorbibnamehash}{0fa981e305ea35073dc938b976cda0e1}
      \strng{authornamehash}{0fa981e305ea35073dc938b976cda0e1}
      \strng{authorfullhash}{0fa981e305ea35073dc938b976cda0e1}
      \field{sortinit}{C}
      \field{sortinithash}{4c244ceae61406cdc0cc2ce1cb1ff703}
      \field{labelprefix}{A}
      \field{labelnamesource}{author}
      \field{labeltitlesource}{title}
      \field{annotation}{In this letter, the author remarks for the first time that there appear to be more primes of the form $4n+3$ than $4n+1$, in that their counting functions ``differ notably in their second terms'' (original French: {\it ``different notablement entre elles par leurs seconds termes''}). Several assertions are made (without proof): that $E^\pi(x;4,3,1)$ takes values arbitrarily close to $1$; that $\pi_e(x,\chi_{-4})\to-\infty$ as $x\to\infty$; and that if $f(x)$ is a decreasing function with $\lim_{x\to\infty} x^{1/2} f(x) \ne 0$, then the series $\sum_p \chi_{-4}(p) f(p)$ diverges.}
      \field{journaltitle}{Bull. de la Classe phys. math. de l'Acad. Imp. des Sciences St. Petersburg}
      \field{title}{{Lettre de M. le professeur Tchébychev a M. Fuss, sur un nouveau théorème rélatif aux nombres premiers contenus dans la formes $4n+1$ et $4n+3$}}
      \field{volume}{11}
      \field{year}{1853}
      \true{nocite}
      \field{pages}{208}
      \range{pages}{1}
    \endentry
    \entry{MR602236}{article}{}
      \name{author}{1}{}{%
        {{hash=97466ee083d3402c436ef85e5d255b89}{%
           family={Chen},
           familyi={C\bibinitperiod},
           given={W.\bibnamedelimi W.\bibnamedelimi L.},
           giveni={W\bibinitperiod\bibinitdelim W\bibinitperiod\bibinitdelim L\bibinitperiod}}}%
      }
      \strng{namehash}{97466ee083d3402c436ef85e5d255b89}
      \strng{fullhash}{97466ee083d3402c436ef85e5d255b89}
      \strng{bibnamehash}{97466ee083d3402c436ef85e5d255b89}
      \strng{authorbibnamehash}{97466ee083d3402c436ef85e5d255b89}
      \strng{authornamehash}{97466ee083d3402c436ef85e5d255b89}
      \strng{authorfullhash}{97466ee083d3402c436ef85e5d255b89}
      \field{sortinit}{C}
      \field{sortinithash}{4c244ceae61406cdc0cc2ce1cb1ff703}
      \field{labelprefix}{A}
      \field{labelnamesource}{author}
      \field{labeltitlesource}{title}
      \field{annotation}{It was conjectured in~\cites{MR0108470} that $\sum_{n\le x} E^\pi(n) \sim x$ and $\sum_{n\le x} E^\pi(n;4,3,1) \sim x$; in this article the author disproves these conjectures. Defining the function $R(x)= \sum_{n \leq x} E^{\psi}(n)$ he shows that $R(x)=\Omega_\pm(x^{{1/2}+\Theta-\ep})$ for every $\ep>0$, which can be improved to $R(x)=\Omega_{\pm}(x^{1/2+\Theta})$ under SA. The same results hold for $P(x)-x$ in place of $R(x)$, where $P(x)= - \sum_{n \leq x} E^{\pi}(n)$. In these theorems, the author notes that $\psi(n)$ and $\pi(n)$ can also be replaced by $\psi(n;4,1,3)$ and $\pi(n;4,1,3)$.\par This article cites~\cites{MR0108470}.}
      \field{issn}{0024-6107}
      \field{journaltitle}{J. London Math. Soc. (2)}
      \field{number}{1}
      \field{title}{On the error term of the prime number theorem and the difference between the number of primes in the residue classes modulo $4$}
      \field{volume}{23}
      \field{year}{1981}
      \true{nocite}
      \field{pages}{24\bibrangedash 40}
      \range{pages}{17}
      \verb{doi}
      \verb 10.1112/jlms/s2-23.1.24
      \endverb
      \verb{mrnumber}
      \verb 602236
      \endverb
    \endentry
    \entry{MR0233781}{article}{}
      \name{author}{2}{}{%
        {{hash=e0b33296fa8f31deb59123e86ff78d6e}{%
           family={Cohen},
           familyi={C\bibinitperiod},
           given={A.\bibnamedelimi M.},
           giveni={A\bibinitperiod\bibinitdelim M\bibinitperiod}}}%
        {{hash=ab9ca14356d28b56bf08d921be43f3b6}{%
           family={Mayhew},
           familyi={M\bibinitperiod},
           given={M.\bibnamedelimi J.\bibnamedelimi E.},
           giveni={M\bibinitperiod\bibinitdelim J\bibinitperiod\bibinitdelim E\bibinitperiod}}}%
      }
      \strng{namehash}{d09ba402a722f40a101c1ad5fa0eb9d0}
      \strng{fullhash}{d09ba402a722f40a101c1ad5fa0eb9d0}
      \strng{bibnamehash}{d09ba402a722f40a101c1ad5fa0eb9d0}
      \strng{authorbibnamehash}{d09ba402a722f40a101c1ad5fa0eb9d0}
      \strng{authornamehash}{d09ba402a722f40a101c1ad5fa0eb9d0}
      \strng{authorfullhash}{d09ba402a722f40a101c1ad5fa0eb9d0}
      \field{sortinit}{C}
      \field{sortinithash}{4c244ceae61406cdc0cc2ce1cb1ff703}
      \field{labelprefix}{A}
      \field{labelnamesource}{author}
      \field{labeltitlesource}{title}
      \field{annotation}{Using an unpublished manuscript of Turing as a starting point, and using computations of zeros of $\zeta(s)$ by Haselgrove, the authors show that $\pi(x)-\li(x)>0$ for some $x\le10^{10^{529.7}}$.\par This article cites~\cites{MR0202686}.}
      \field{issn}{0024-6115}
      \field{journaltitle}{Proc. London Math. Soc. (3)}
      \field{title}{On the difference {$\pi (x)-\li (x)$}}
      \field{volume}{18}
      \field{year}{1968}
      \true{nocite}
      \field{pages}{691\bibrangedash 713}
      \range{pages}{23}
      \verb{doi}
      \verb 10.1112/plms/s3-18.4.691
      \endverb
      \verb{mrnumber}
      \verb 0233781
      \endverb
    \endentry
    \entry{MR1544509}{article}{}
      \name{author}{1}{}{%
        {{hash=8815262f8ce833f8804c25a5afc1996b}{%
           family={Cramér},
           familyi={C\bibinitperiod},
           given={Harald},
           giveni={H\bibinitperiod}}}%
      }
      \list{language}{1}{%
        {German}%
      }
      \strng{namehash}{8815262f8ce833f8804c25a5afc1996b}
      \strng{fullhash}{8815262f8ce833f8804c25a5afc1996b}
      \strng{bibnamehash}{8815262f8ce833f8804c25a5afc1996b}
      \strng{authorbibnamehash}{8815262f8ce833f8804c25a5afc1996b}
      \strng{authornamehash}{8815262f8ce833f8804c25a5afc1996b}
      \strng{authorfullhash}{8815262f8ce833f8804c25a5afc1996b}
      \field{sortinit}{C}
      \field{sortinithash}{4c244ceae61406cdc0cc2ce1cb1ff703}
      \field{labelprefix}{A}
      \field{labelnamesource}{author}
      \field{labeltitlesource}{title}
      \field{annotation}{Under RH, the author shows that \[ \lim_{x \to \infty} \frac{1}{\log x} \int_2^x \left( \frac{\Delta^\psi(t)}{t} \right)^2 \, dt = \sum_\rho \left| \frac{n_\rho}{\rho} \right|^2, \] where $n_\rho$ is the multiplicity of the zero $\rho$. It follows that $\fA_1^\psi(x) \ll x^{3/2}$ (and the same for $\fA_1^\theta(x)$) and thus that $\fA_1^\pi(x) \ll x^{3/2}/\log x$ (and the same for $\fA_1^\Pi(x)$).\par \par This article cites~\cites{MR1555148,1909.Landau}.}
      \field{issn}{0025-5874}
      \field{journaltitle}{Math. Z.}
      \field{number}{1}
      \field{title}{Ein Mittelwertsatz in der Primzahltheorie}
      \field{volume}{12}
      \field{year}{1922}
      \true{nocite}
      \field{pages}{147\bibrangedash 153}
      \range{pages}{7}
      \verb{doi}
      \verb 10.1007/BF01482072
      \endverb
      \verb{mrnumber}
      \verb 1544509
      \endverb
    \endentry
    \entry{MR0352012}{article}{}
      \name{author}{2}{}{%
        {{hash=d4045cfd1e1488ca61ed5c9c26bdd0a6}{%
           family={Dancs},
           familyi={D\bibinitperiod},
           given={S.},
           giveni={S\bibinitperiod}}}%
        {{hash=e68e101772b88cb1d48e688164bece92}{%
           family={Turán},
           familyi={T\bibinitperiod},
           given={P.},
           giveni={P\bibinitperiod}}}%
      }
      \strng{namehash}{e1ec24cc5bfe1ad55de291dec425eaea}
      \strng{fullhash}{e1ec24cc5bfe1ad55de291dec425eaea}
      \strng{bibnamehash}{e1ec24cc5bfe1ad55de291dec425eaea}
      \strng{authorbibnamehash}{e1ec24cc5bfe1ad55de291dec425eaea}
      \strng{authornamehash}{e1ec24cc5bfe1ad55de291dec425eaea}
      \strng{authorfullhash}{e1ec24cc5bfe1ad55de291dec425eaea}
      \field{sortinit}{D}
      \field{sortinithash}{c438b3d5d027251ba63f5ed538d98af5}
      \field{labelprefix}{A}
      \field{labelnamesource}{author}
      \field{labeltitlesource}{title}
      \field{annotation}{Spurred by unpublished work of Knapowski, the authors establish the following more flexible version of the second main theorem of the power-sum method. In addition to the usual condition $1=|z_1|\geq |z_2| \geq \cdots \geq |z_n|$, let $m \geq 0$ be an integer, choose $\delta_1$ and $\delta_2$ with $\frac{m}{m+n} \leq \delta_2 \leq \delta_1 \leq 1$, and choose indices $\ell_1$ and $\ell_2$ such that $|z_{\ell_1}| \geq \delta_1$ and $|z_{\ell_2}| \leq \delta_2$ (or $\ell_2=n$ if no such $z_{\ell_2}$ exists). Then \[ \max_{m+1 \leq v \leq m+n} \bigg| \sum_{j=1}^n b_j z_j^v \bigg| \geq 2\bigg( \frac{\delta_1-\delta_2}{8e} \bigg)^n \min_{\ell_1 \leq j \leq \ell_2} |b_1+ \cdots +b_j|. \]\par This article cites~\cites{MR0060548}.}
      \field{issn}{0524-9007}
      \field{journaltitle}{Ann. Univ. Sci. Budapest. Eötvös Sect. Math.}
      \field{title}{Investigations in the powersum theory. {I}}
      \field{volume}{16}
      \field{year}{1973}
      \true{nocite}
      \field{pages}{47\bibrangedash 52 (1974)}
      \range{pages}{-1}
      \verb{mrnumber}
      \verb 0352012
      \endverb
    \endentry
    \entry{MR2047102}{article}{}
      \name{author}{3}{}{%
        {{hash=6f53fb227b01894d40689a82f6e1c775}{%
           family={Deléglise},
           familyi={D\bibinitperiod},
           given={Marc},
           giveni={M\bibinitperiod}}}%
        {{hash=34a11ef93294af142e84c4243d02208b}{%
           family={Dusart},
           familyi={D\bibinitperiod},
           given={Pierre},
           giveni={P\bibinitperiod}}}%
        {{hash=75094d519df9a079290d93e149ccf600}{%
           family={Roblot},
           familyi={R\bibinitperiod},
           given={Xavier-François},
           giveni={X\bibinithyphendelim F\bibinitperiod}}}%
      }
      \strng{namehash}{1237cd8bf0b8b34be0815571948c8580}
      \strng{fullhash}{1237cd8bf0b8b34be0815571948c8580}
      \strng{bibnamehash}{1237cd8bf0b8b34be0815571948c8580}
      \strng{authorbibnamehash}{1237cd8bf0b8b34be0815571948c8580}
      \strng{authornamehash}{1237cd8bf0b8b34be0815571948c8580}
      \strng{authorfullhash}{1237cd8bf0b8b34be0815571948c8580}
      \field{sortinit}{D}
      \field{sortinithash}{c438b3d5d027251ba63f5ed538d98af5}
      \field{labelprefix}{A}
      \field{labelnamesource}{author}
      \field{labeltitlesource}{title}
      \field{annotation}{Extending a well-known method of Meissel for computing $\pi(x)$, the authors give an algorithm for computing $\pi(x;q,a)$ in time $O(x^{2/3}/\log^2 x)$. They use this algorithm is used to compute $\pi(x;4,1)$ and $\pi(x;4,3)$ for $x = 10^j,2 \cdot 10^j,\ldots,9 \cdot 10^j$ for each $10 \leq j \leq 19$. In particular, a new region where $\pi(x;4,3)<\pi(x;4,1)$ is discovered near $x = 10^{18}$. \par This article cites~\cites{MR0476615,MR1816036,MR1806291,MR0184920,MR0083001,1914.Littlewood,MR1329368}.}
      \field{issn}{0025-5718}
      \field{journaltitle}{Math. Comp.}
      \field{number}{247}
      \field{title}{Counting primes in residue classes}
      \field{volume}{73}
      \field{year}{2004}
      \true{nocite}
      \field{pages}{1565\bibrangedash 1575}
      \range{pages}{11}
      \verb{doi}
      \verb 10.1090/S0025-5718-04-01649-7
      \endverb
      \verb{mrnumber}
      \verb 2047102
      \endverb
    \endentry
    \entry{MR4120786}{article}{}
      \name{author}{1}{}{%
        {{hash=c9734df31e44a7c0ee5c006cb13a5e86}{%
           family={Devin},
           familyi={D\bibinitperiod},
           given={Lucile},
           giveni={L\bibinitperiod}}}%
      }
      \strng{namehash}{c9734df31e44a7c0ee5c006cb13a5e86}
      \strng{fullhash}{c9734df31e44a7c0ee5c006cb13a5e86}
      \strng{bibnamehash}{c9734df31e44a7c0ee5c006cb13a5e86}
      \strng{authorbibnamehash}{c9734df31e44a7c0ee5c006cb13a5e86}
      \strng{authornamehash}{c9734df31e44a7c0ee5c006cb13a5e86}
      \strng{authorfullhash}{c9734df31e44a7c0ee5c006cb13a5e86}
      \field{extraname}{1}
      \field{sortinit}{D}
      \field{sortinithash}{c438b3d5d027251ba63f5ed538d98af5}
      \field{labelprefix}{A}
      \field{labelnamesource}{author}
      \field{labeltitlesource}{title}
      \field{annotation}{In this article, the author extends the study of prime number races to counting functions with explicit formulas associated to general $L$-functions from the Selberg class, including Dirichlet $L$-functions and Hasse--Weil $L$-functions. She unconditionally proves the existence of a limiting logarithmic distribution for these counting functions, and establishes properties of the limiting distribution (such as absolute continuity) under much weaker conditions than GRH and LI.\par This article cites~\cites{MR1329368,MR2202918,MR3261965, MR1909220}.}
      \field{issn}{0305-0041}
      \field{journaltitle}{Math. Proc. Cambridge Philos. Soc.}
      \field{number}{1}
      \field{title}{Chebyshev's bias for analytic {L}-functions}
      \field{volume}{169}
      \field{year}{2020}
      \true{nocite}
      \field{pages}{103\bibrangedash 140}
      \range{pages}{38}
      \verb{doi}
      \verb 10.1017/s0305004119000100
      \endverb
      \verb{mrnumber}
      \verb 4120786
      \endverb
      \verb{urlraw}
      \verb https://doi.org/10.1017/s0305004119000100
      \endverb
      \verb{url}
      \verb https://doi.org/10.1017/s0305004119000100
      \endverb
    \endentry
    \entry{MR4176773}{article}{}
      \name{author}{1}{}{%
        {{hash=c9734df31e44a7c0ee5c006cb13a5e86}{%
           family={Devin},
           familyi={D\bibinitperiod},
           given={Lucile},
           giveni={L\bibinitperiod}}}%
      }
      \strng{namehash}{c9734df31e44a7c0ee5c006cb13a5e86}
      \strng{fullhash}{c9734df31e44a7c0ee5c006cb13a5e86}
      \strng{bibnamehash}{c9734df31e44a7c0ee5c006cb13a5e86}
      \strng{authorbibnamehash}{c9734df31e44a7c0ee5c006cb13a5e86}
      \strng{authornamehash}{c9734df31e44a7c0ee5c006cb13a5e86}
      \strng{authorfullhash}{c9734df31e44a7c0ee5c006cb13a5e86}
      \field{extraname}{2}
      \field{sortinit}{D}
      \field{sortinithash}{c438b3d5d027251ba63f5ed538d98af5}
      \field{labelprefix}{A}
      \field{labelnamesource}{author}
      \field{labeltitlesource}{title}
      \field{annotation}{In this article, the author extends her earlier the weakening of the LI assumption to the setting of prime number races with many contestants. Assuming GRH but only weakened versions of the LI assumption, she shows the existence of logarithmic densities associated to $r$-way prime number races. Her methods apply equally well to the number field and function field settings (in the latter case, GRH is no longer a hypothesis). An erratum appeared in the same journal in 2021 (MR\href{http://www.ams.org/mathscinet-getitem?mr=4352666}{\nolinkurl{4352666}}).\par This article cites~\cites{MR1329368,MR3993767,MR4082248,MR1957110,MR2202918,MR1985941,MR3803705,MR4120786,MR2600420,MR2811555,MR3581815,MR2474313,MR3261965}.}
      \field{issn}{0008-4395}
      \field{journaltitle}{Canad. Math. Bull.}
      \field{number}{4}
      \field{title}{Limiting properties of the distribution of primes in an arbitrarily large number of residue classes}
      \field{volume}{63}
      \field{year}{2020}
      \true{nocite}
      \field{pages}{837\bibrangedash 849}
      \range{pages}{13}
      \verb{doi}
      \verb 10.4153/s0008439520000089
      \endverb
      \verb{mrnumber}
      \verb 4176773
      \endverb
      \verb{urlraw}
      \verb https://doi.org/10.4153/s0008439520000089
      \endverb
      \verb{url}
      \verb https://doi.org/10.4153/s0008439520000089
      \endverb
    \endentry
    \entry{arXiv.Devin}{online}{}
      \name{author}{1}{}{%
        {{hash=c9734df31e44a7c0ee5c006cb13a5e86}{%
           family={Devin},
           familyi={D\bibinitperiod},
           given={Lucile},
           giveni={L\bibinitperiod}}}%
      }
      \strng{namehash}{c9734df31e44a7c0ee5c006cb13a5e86}
      \strng{fullhash}{c9734df31e44a7c0ee5c006cb13a5e86}
      \strng{bibnamehash}{c9734df31e44a7c0ee5c006cb13a5e86}
      \strng{authorbibnamehash}{c9734df31e44a7c0ee5c006cb13a5e86}
      \strng{authornamehash}{c9734df31e44a7c0ee5c006cb13a5e86}
      \strng{authorfullhash}{c9734df31e44a7c0ee5c006cb13a5e86}
      \field{extraname}{3}
      \field{sortinit}{D}
      \field{sortinithash}{c438b3d5d027251ba63f5ed538d98af5}
      \field{labelprefix}{A}
      \field{labelnamesource}{author}
      \field{labeltitlesource}{title}
      \field{annotation}{Every $p\equiv1\mod4$ can be written uniquely as $p=a^2+4b^2$ with positive integers $a,b$. The author formulates two conjectures concerning counting functions related to the distribution of these numbers~$a$ and $2b$. First, she conjectures that the function $\sum_{p\le x,\, p\equiv1\mod4} \mathop{\rm sign}(a-2b)$ is negative for a set of logarithmic density strictly between~$\frac12$ and~$1$, and thus there is a bias towards the even square being larger than the odd square in such representations. Second, she conjectures that $\sum_{p\le x,\, p\equiv1\mod4} \chi_{-4}(a)$ is eventually always positive, so that there is a complete bias towards the positive odd number in such representations being $1\mod4$ rather than $3\mod4$. Both conjectures rely on her extension of existing results to cover counting functions that are governed by a sum of infinitely many $L$-functions (such as Hecke $L$-functions).\par This article cites~\cites{MR3261965,MR3581815,MR4120786,MR3264673,MR3272280,MR1957110,MR1329368,2007.Sarnak,MR0004255}, as well as an earlier draft of this annotated bibliography.}
      \field{title}{Discrepancies in the distribution of Gaussian primes}
      \field{year}{2021}
      \true{nocite}
      \verb{urlraw}
      \verb https://arxiv.org/abs/2105.02492
      \endverb
      \verb{url}
      \verb https://arxiv.org/abs/2105.02492
      \endverb
    \endentry
    \entry{MR4316675}{article}{}
      \name{author}{2}{}{%
        {{hash=c9734df31e44a7c0ee5c006cb13a5e86}{%
           family={Devin},
           familyi={D\bibinitperiod},
           given={Lucile},
           giveni={L\bibinitperiod}}}%
        {{hash=0e36a3a1f7c6c0678ce294e777ceb2f9}{%
           family={Meng},
           familyi={M\bibinitperiod},
           given={Xianchang},
           giveni={X\bibinitperiod}}}%
      }
      \strng{namehash}{452055ee754af6a220856a0aadedfe82}
      \strng{fullhash}{452055ee754af6a220856a0aadedfe82}
      \strng{bibnamehash}{452055ee754af6a220856a0aadedfe82}
      \strng{authorbibnamehash}{452055ee754af6a220856a0aadedfe82}
      \strng{authornamehash}{452055ee754af6a220856a0aadedfe82}
      \strng{authorfullhash}{452055ee754af6a220856a0aadedfe82}
      \field{sortinit}{D}
      \field{sortinithash}{c438b3d5d027251ba63f5ed538d98af5}
      \field{labelprefix}{A}
      \field{labelnamesource}{author}
      \field{labeltitlesource}{title}
      \field{annotation}{The authors study the counting function of polynomials in $\mathbb F_q[t]$ in sets of invertible residue classes modulo a fixed polynomial that have $k$ irreducible factors (these can be counted either with or without multiplicity). They prove asymptotic formulas for the normalized differences of two such counting functions in terms of the zeros of the relevant $L$-functions; these asymptotic formulas are valid for $k$ almost as large as $\sqrt{\log X}$ where $X$ bounds the degree of the polynomials being counted. In the important special case where the sets of residue classes are the quadratic residues and nonresidues, the formulas simplify and the sign of the bias analyzed: assuming LI, the bias is towards quadratic residues when we count irreducible factors without multiplicity, while the bias depends on the parity of~$k$ when we count them with multiplicity. Finally, the authors use the fact that some $L$-functions vanish at the critical point $q^{-1/2}$ to show that there exist completely biased and unbiased races, contrary to what is expected in the number field setting.\par This article cites~\cites{MR3261965,MR3581815, MR2474313, MR3272280, MR1985941, MR3028758, MR2778652, MR3803705, MR3778223, MR4082248, MR2034131, MR1329368,MR4176773}.}
      \field{issn}{0001-8708}
      \field{journaltitle}{Adv. Math.}
      \field{title}{Chebyshev's bias for products of irreducible polynomials}
      \field{volume}{392}
      \field{year}{2021}
      \true{nocite}
      \field{pages}{Paper No. 108040\bibrangessep 45}
      \range{pages}{-1}
      \verb{doi}
      \verb 10.1016/j.aim.2021.108040
      \endverb
      \verb{mrnumber}
      \verb 4316675
      \endverb
      \verb{urlraw}
      \verb https://doi.org/10.1016/j.aim.2021.108040
      \endverb
      \verb{url}
      \verb https://doi.org/10.1016/j.aim.2021.108040
      \endverb
    \endentry
    \entry{MR0332684}{article}{}
      \name{author}{1}{}{%
        {{hash=5b3dabbd7db4c81dcf1e7cffedcb1b4d}{%
           family={Diamond},
           familyi={D\bibinitperiod},
           given={H.\bibnamedelimi G.},
           giveni={H\bibinitperiod\bibinitdelim G\bibinitperiod}}}%
      }
      \list{publisher}{1}{%
        {Springer, Berlin}%
      }
      \strng{namehash}{5b3dabbd7db4c81dcf1e7cffedcb1b4d}
      \strng{fullhash}{5b3dabbd7db4c81dcf1e7cffedcb1b4d}
      \strng{bibnamehash}{5b3dabbd7db4c81dcf1e7cffedcb1b4d}
      \strng{authorbibnamehash}{5b3dabbd7db4c81dcf1e7cffedcb1b4d}
      \strng{authornamehash}{5b3dabbd7db4c81dcf1e7cffedcb1b4d}
      \strng{authorfullhash}{5b3dabbd7db4c81dcf1e7cffedcb1b4d}
      \field{extraname}{1}
      \field{sortinit}{D}
      \field{sortinithash}{c438b3d5d027251ba63f5ed538d98af5}
      \field{labelprefix}{A}
      \field{labelnamesource}{author}
      \field{labeltitlesource}{title}
      \field{annotation}{The author presents two variants of oscillation theorems analogous to those of Ingham in~\cite{MR0006202}. Let $F(s)=\int_0^\infty e^{-su}f(u) \, du$ denote the Laplace transform of the measurable function $f\colon [0,\infty)\to\R$. We suppose that the integral defining $F(s)$ converges for $\Re(s)>0$, and that $F(s)$ can be continued as a meromorphic function to a neighborhood of the imaginary axis; suppose further that all the poles of $F(s)$ on the imaginary axis are simple. Let $T$ be the set of positive real numbers $t$ such that $it$ is a pole of~$F(s)$, and let $a_t$ be the residue of $F(s)$ at $s=it$; furthermore, let $a_0$ be the residue (possibly $0$) of $F(s)$ at $s=0$.\par The author defines a subset $W\subset T$ to be ``weakly independent of order $N$'' if the only way to find integers $|n_t|\leq N$ $(t\in W)$ such that $\sum_{t\in W} n_t t \in T$ is to choose one $n_t$ equal to~$1$ and the rest equal to~$0$. Given such a weakly independent subset $W\subset T$ of order $N$, the author proves that \begin{align*} \lim_{x\to\infty}\mathop{\rm ess\ sup}_{u\geq x} f(u) &\geq a_0+\frac {2N}{N+1}\sum_{j\in J}|a_j| \\ \lim_{x\to\infty}\mathop{\rm ess\ inf}_{u\geq x} f(u) &\leq a_0-\frac {2N}{N+1}\sum_{j\in J}|a_j| \end{align*} (where these essential supremum and infemum denote the supremum/infemum when we may ignore a set of inputs of measure~$0$); equivalently, if $ \frac {2N}{N+1}\sum_{j\in J}|a_j|>|a_0| $ then $f(x)$ has arbitrarily large sign changes. (The author gives a slight strengthening of this theorem as well.)\par This article cites~\cites{MR0006202, MR0104638, MR0198145, MR0202685, MR0330069,MR0120198,MR280447}.}
      \field{booktitle}{The theory of arithmetic functions ({P}roc. {C}onf., {W}estern {M}ichigan {U}niv., {K}alamazoo, {M}ich., 1971)}
      \field{title}{Two oscillation theorems}
      \field{year}{1972}
      \true{nocite}
      \field{pages}{113\bibrangedash 118. Lecture Notes in Math.\bibrangessep Vol. 251}
      \range{pages}{-1}
      \verb{mrnumber}
      \verb 0332684
      \endverb
    \endentry
    \entry{MR2605532}{article}{}
      \name{author}{2}{}{%
        {{hash=5b3dabbd7db4c81dcf1e7cffedcb1b4d}{%
           family={Diamond},
           familyi={D\bibinitperiod},
           given={H.\bibnamedelimi G.},
           giveni={H\bibinitperiod\bibinitdelim G\bibinitperiod}}}%
        {{hash=c6d965df16e600c701d4bcc2888ae34c}{%
           family={Pintz},
           familyi={P\bibinitperiod},
           given={J.},
           giveni={J\bibinitperiod}}}%
      }
      \strng{namehash}{d3c019510c17f40c3aa1bd8a2407a489}
      \strng{fullhash}{d3c019510c17f40c3aa1bd8a2407a489}
      \strng{bibnamehash}{d3c019510c17f40c3aa1bd8a2407a489}
      \strng{authorbibnamehash}{d3c019510c17f40c3aa1bd8a2407a489}
      \strng{authornamehash}{d3c019510c17f40c3aa1bd8a2407a489}
      \strng{authorfullhash}{d3c019510c17f40c3aa1bd8a2407a489}
      \field{sortinit}{D}
      \field{sortinithash}{c438b3d5d027251ba63f5ed538d98af5}
      \field{labelprefix}{A}
      \field{labelnamesource}{author}
      \field{labeltitlesource}{title}
      \field{annotation}{In this article, the authors show that $-\sum_{p\leq x}\log(1-\frac{1}{p})-\log\log x-C_0=\Omega_\pm(1/(\sqrt{x}\log x))$, which implies that $\prod_{p\leq x}(1-\frac{1}{p})^{-1} - e^{{C_0}}\log x$ changes sign infinitely often. If RH is false, the result is a consequence of Landau's theorem. If RH is true, the authors first show that $$ \int_1^x\frac{d\Pi(t)}{t}-\int_1^x\frac{1-t^{-1}}{t\log t}\,dt=\Omega_\pm\biggl( \frac{\log\log\log x}{\sqrt{x}\log x} \biggr) $$ and then use Littlewood's~\cites{1914.Littlewood} result on the sign changes of $\Delta^\pi(x)$ and Cramér's bound $\Delta^{E^{\psi}}_1(x)\ll x$. The authors include a second proof in the RH case using a variant of the Wiener--Ikehara method due to Ingham~\cites{MR0006202}.\par This article cites~\cites{MR0376566,MR0006202,1914.Littlewood,MR0137689,MR43841,MR51857}.}
      \field{issn}{1246-7405}
      \field{journaltitle}{J. Théor. Nombres Bordeaux}
      \field{number}{3}
      \field{title}{Oscillation of {M}ertens' product formula}
      \field{volume}{21}
      \field{year}{2009}
      \true{nocite}
      \field{pages}{523\bibrangedash 533}
      \range{pages}{11}
      \verb{mrnumber}
      \verb 2605532
      \endverb
      \verb{urlraw}
      \verb http://jtnb.cedram.org/item?id=JTNB_2009__21_3_523_0
      \endverb
      \verb{url}
      \verb http://jtnb.cedram.org/item?id=JTNB_2009__21_3_523_0
      \endverb
    \endentry
    \entry{MR0376566}{article}{}
      \name{author}{1}{}{%
        {{hash=4ceaa3cbf82f92b37edd652b07fd8881}{%
           family={Diamond},
           familyi={D\bibinitperiod},
           given={Harold\bibnamedelima G.},
           giveni={H\bibinitperiod\bibinitdelim G\bibinitperiod}}}%
      }
      \strng{namehash}{4ceaa3cbf82f92b37edd652b07fd8881}
      \strng{fullhash}{4ceaa3cbf82f92b37edd652b07fd8881}
      \strng{bibnamehash}{4ceaa3cbf82f92b37edd652b07fd8881}
      \strng{authorbibnamehash}{4ceaa3cbf82f92b37edd652b07fd8881}
      \strng{authornamehash}{4ceaa3cbf82f92b37edd652b07fd8881}
      \strng{authorfullhash}{4ceaa3cbf82f92b37edd652b07fd8881}
      \field{extraname}{2}
      \field{sortinit}{D}
      \field{sortinithash}{c438b3d5d027251ba63f5ed538d98af5}
      \field{labelprefix}{A}
      \field{labelnamesource}{author}
      \field{labeltitlesource}{title}
      \field{annotation}{The author gives a new proof that $\pi(x)-\li(x)$ changes sign infinitely often, without using an explicit formula for prime-counting functions but instead establishing a Tauberian theorem of Wiener--Ikehara type. The author believes that the arguments could be extended to gain the extra factor of $\log\log\log x$ in the oscillations found by Littlewood.\par This article cites~\cites{MR0137689, MR0006202}.}
      \field{issn}{0013-8584}
      \field{journaltitle}{Enseignement Math. (2)}
      \field{number}{1}
      \field{title}{Changes of sign of {$\pi (x)-\li (x)$}}
      \field{volume}{21}
      \field{year}{1975}
      \true{nocite}
      \field{pages}{1\bibrangedash 14}
      \range{pages}{14}
      \verb{mrnumber}
      \verb 0376566
      \endverb
    \endentry
    \entry{MR3521338}{article}{}
      \name{author}{3}{}{%
        {{hash=d7aee8a672f37558eb9119074764e7f5}{%
           family={Dummit},
           familyi={D\bibinitperiod},
           given={D.},
           giveni={D\bibinitperiod}}}%
        {{hash=7c95bde388d46e453494c7b2fbd1b088}{%
           family={Granville},
           familyi={G\bibinitperiod},
           given={A.},
           giveni={A\bibinitperiod}}}%
        {{hash=a699d3ec8d37bdcc6b5f75e108b4b9d3}{%
           family={Kisilevsky},
           familyi={K\bibinitperiod},
           given={B.},
           giveni={B\bibinitperiod}}}%
      }
      \strng{namehash}{e1122a9ecacc1c1ca72d1f02ad518714}
      \strng{fullhash}{e1122a9ecacc1c1ca72d1f02ad518714}
      \strng{bibnamehash}{e1122a9ecacc1c1ca72d1f02ad518714}
      \strng{authorbibnamehash}{e1122a9ecacc1c1ca72d1f02ad518714}
      \strng{authornamehash}{e1122a9ecacc1c1ca72d1f02ad518714}
      \strng{authorfullhash}{e1122a9ecacc1c1ca72d1f02ad518714}
      \field{sortinit}{D}
      \field{sortinithash}{c438b3d5d027251ba63f5ed538d98af5}
      \field{labelprefix}{A}
      \field{labelnamesource}{author}
      \field{labeltitlesource}{title}
      \field{annotation}{This article establishes permanent, relatively large biases in races involving products of two primes. The authors show that if $\chi$ is a quadratic character\mod d, then \[ \frac{\#\{ pq\le x\colon \chi(p)=\chi(q)=-1 \}}{\#\{ pq\le x\colon (pq,d)=1 \}} = 1 - \bigg( \sum_p \frac{\chi(p)}p + o(1) \bigg) \frac1{\log\log x}, \] and the same statement with both minus signs changed to plus signs. In particular, the race between integers $pq$ with $\chi(p)=\chi(q)=-1$ and those with $\chi(p)=\chi(q)=1$ has a bias in favor of the sign of (and proportional to) $\sum_p {\chi(p)}/p$.\par For example, $\sum_p {\chi_5(p)}/p \approx -1.008$, and correspondingly integers $pq$ with both $p$ and $q$ quadratic nonresidues\mod 5 are 41.6\% more numerous up to $10^7$ than random chance would suggest. The authors conjecture that there exists $d\le x$ such that the right-hand side is as large as $1+\log\log\log x/{\log\log x}$. They also note that the same proof gives a bias for the ratio \[ \sum_{\substack{p\le x\\\chi(p)=1}}\frac1p\; \bigg/\sum_{\substack{p\le x\\\chi(p)=-1}}\frac1p = 1 + \bigg( 2\sum_p \frac{\chi(p)}p + o(1) \bigg) \frac1{\log\log x}, \] giving a permanent bias involving only primes (in contrast to the unweighted race between $\pi(x;d,R)$ and $\pi(x;d,N)$); the proof also generalizes to products of $k$ primes with prescribed quadratic character values (for fixed $k$). They remark on the possibility of counting $pq\le x$ where $p$ and $q$ are restricted to prescribed but arbitrary residue classes.\par This article cites~\cites{MR2778652}.}
      \field{issn}{0025-5793}
      \field{journaltitle}{Mathematika}
      \field{number}{2}
      \field{title}{Big biases amongst products of two primes}
      \field{volume}{62}
      \field{year}{2016}
      \true{nocite}
      \field{pages}{502\bibrangedash 507}
      \range{pages}{6}
      \verb{doi}
      \verb 10.1112/S0025579315000339
      \endverb
      \verb{mrnumber}
      \verb 3521338
      \endverb
      \verb{urlraw}
      \verb https://doi.org/10.1112/S0025579315000339
      \endverb
      \verb{url}
      \verb https://doi.org/10.1112/S0025579315000339
      \endverb
    \endentry
    \entry{MR814687}{book}{}
      \name{author}{2}{}{%
        {{hash=5ddf9a712233e9bcfe9ad3ab3585dcf4}{%
           family={Ellison},
           familyi={E\bibinitperiod},
           given={William},
           giveni={W\bibinitperiod}}}%
        {{hash=641b110292dfa35db45ceb8c7790ba2d}{%
           family={Ellison},
           familyi={E\bibinitperiod},
           given={Fern},
           giveni={F\bibinitperiod}}}%
      }
      \list{publisher}{1}{%
        {John Wiley \& Sons, Inc., New York; Hermann, Paris}%
      }
      \strng{namehash}{2f130d3d62e0e56f43edeac5f8911bca}
      \strng{fullhash}{2f130d3d62e0e56f43edeac5f8911bca}
      \strng{bibnamehash}{2f130d3d62e0e56f43edeac5f8911bca}
      \strng{authorbibnamehash}{2f130d3d62e0e56f43edeac5f8911bca}
      \strng{authornamehash}{2f130d3d62e0e56f43edeac5f8911bca}
      \strng{authorfullhash}{2f130d3d62e0e56f43edeac5f8911bca}
      \field{sortinit}{E}
      \field{sortinithash}{c554bd1a0b76ea92b9f105fe36d9c7b0}
      \field{labelprefix}{A}
      \field{labelnamesource}{author}
      \field{labeltitlesource}{title}
      \field{annotation}{This book is an English translation of~\cite{MR0417077}.}
      \field{isbn}{0-471-82653-7}
      \field{series}{A Wiley-Interscience Publication}
      \field{title}{Prime numbers}
      \field{year}{1985}
      \true{nocite}
      \field{pages}{xii+417}
      \range{pages}{-1}
      \verb{mrnumber}
      \verb 814687
      \endverb
    \endentry
    \entry{MR0417077}{book}{}
      \name{author}{1}{}{%
        {{hash=2f35569443e88f879c6356cd12f56c9c}{%
           family={Ellison},
           familyi={E\bibinitperiod},
           given={William\bibnamedelima John},
           giveni={W\bibinitperiod\bibinitdelim J\bibinitperiod}}}%
      }
      \list{language}{1}{%
        {French}%
      }
      \list{publisher}{1}{%
        {Hermann, Paris}%
      }
      \strng{namehash}{2f35569443e88f879c6356cd12f56c9c}
      \strng{fullhash}{2f35569443e88f879c6356cd12f56c9c}
      \strng{bibnamehash}{2f35569443e88f879c6356cd12f56c9c}
      \strng{authorbibnamehash}{2f35569443e88f879c6356cd12f56c9c}
      \strng{authornamehash}{2f35569443e88f879c6356cd12f56c9c}
      \strng{authorfullhash}{2f35569443e88f879c6356cd12f56c9c}
      \field{sortinit}{E}
      \field{sortinithash}{c554bd1a0b76ea92b9f105fe36d9c7b0}
      \field{labelprefix}{A}
      \field{labelnamesource}{author}
      \field{labeltitlesource}{title}
      \field{annotation}{Chapter~6 of the book is on irregularities in the distribution of prime counting functions. Complete proofs of classical oscillation results on $\Delta^{\psi}(x)$ and $\Delta^{\Pi}(x)$ are presented in Section~6.2. The Mertens conjecture and Pólya problem and their connections with Ingham's method~\cite{MR0006202} are discussed in Section~6.3. In Section~6.4, the authors discuss the disproof a conjecture of Shanks~\cite{MR0108470}, which is a precise version of Chebyshev's bias on the prime races between $1 \mod 4$ primes and $3 \mod 4$ primes. In the notes section of the same chapter, the authors discuss the difference between $\Delta^{\pi}(x)$ and its analogue in the setting of Gaussian primes, as well as sign changes of arithmetic functions.\par This chapter cites~\cites{1897.Mertens, MR155787,MR0006202,MR280447,MR0120198, MR0104638, MR0108470, 1853.Chebyshev, MR0083001, MR1555148, MR1511360, 1914.Littlewood, MR1573970, MR0067145, MR0137689, MR0233781, MR0202686, MR0257003,MR0202685,MR0133308,MR0146156,MR0146157,MR0146158,MR0146159,MR0146160,MR0146161}.}
      \field{note}{En collaboration avec Michel Mendès France; Publications de l'Institut de Mathématique de l'Université de Nancago, No. IX; Actualités Scientifiques et Industrielles, No. 1366}
      \field{title}{Les nombres premiers}
      \field{year}{1975}
      \true{nocite}
      \field{pages}{xiv+442}
      \range{pages}{-1}
      \verb{mrnumber}
      \verb 0417077
      \endverb
    \endentry
    \entry{MR1502996}{article}{}
      \name{author}{2}{}{%
        {{hash=1bdd6511a49b2ac9d9318811c581b790}{%
           family={Evelyn},
           familyi={E\bibinitperiod},
           given={C.\bibnamedelimi J.\bibnamedelimi A.},
           giveni={C\bibinitperiod\bibinitdelim J\bibinitperiod\bibinitdelim A\bibinitperiod}}}%
        {{hash=f88769f16846813993d3053353ecc570}{%
           family={Linfoot},
           familyi={L\bibinitperiod},
           given={E.\bibnamedelimi H.},
           giveni={E\bibinitperiod\bibinitdelim H\bibinitperiod}}}%
      }
      \strng{namehash}{34737a2e61b311a58a63b988e137345a}
      \strng{fullhash}{34737a2e61b311a58a63b988e137345a}
      \strng{bibnamehash}{34737a2e61b311a58a63b988e137345a}
      \strng{authorbibnamehash}{34737a2e61b311a58a63b988e137345a}
      \strng{authornamehash}{34737a2e61b311a58a63b988e137345a}
      \strng{authorfullhash}{34737a2e61b311a58a63b988e137345a}
      \field{sortinit}{E}
      \field{sortinithash}{c554bd1a0b76ea92b9f105fe36d9c7b0}
      \field{labelprefix}{A}
      \field{labelnamesource}{author}
      \field{labeltitlesource}{title}
      \field{annotation}{Using Landau's theorem, the authors show that $\Delta^{Q_k}(x)=\Omega_{\pm}(x^{1/2k})$.\par This article cites~\cites{1909.Landau}.}
      \field{issn}{0003-486X,1939-8980}
      \field{journaltitle}{Ann. of Math. (2)}
      \field{number}{2}
      \field{title}{On a problem in the additive theory of numbers}
      \field{volume}{32}
      \field{year}{1931}
      \true{nocite}
      \field{pages}{261\bibrangedash 270}
      \range{pages}{10}
      \verb{doi}
      \verb 10.2307/1968190
      \endverb
      \verb{mrnumber}
      \verb 1502996
      \endverb
      \verb{urlraw}
      \verb https://doi.org/10.2307/1968190
      \endverb
      \verb{url}
      \verb https://doi.org/10.2307/1968190
      \endverb
    \endentry
    \entry{MR43841}{article}{}
      \name{author}{1}{}{%
        {{hash=695e0d70149c6690cca3917ef5b74c9e}{%
           family={Fawaz},
           familyi={F\bibinitperiod},
           given={A.\bibnamedelimi Y.},
           giveni={A\bibinitperiod\bibinitdelim Y\bibinitperiod}}}%
      }
      \strng{namehash}{695e0d70149c6690cca3917ef5b74c9e}
      \strng{fullhash}{695e0d70149c6690cca3917ef5b74c9e}
      \strng{bibnamehash}{695e0d70149c6690cca3917ef5b74c9e}
      \strng{authorbibnamehash}{695e0d70149c6690cca3917ef5b74c9e}
      \strng{authornamehash}{695e0d70149c6690cca3917ef5b74c9e}
      \strng{authorfullhash}{695e0d70149c6690cca3917ef5b74c9e}
      \field{extraname}{1}
      \field{sortinit}{F}
      \field{sortinithash}{fb0c0faa89eb6abae8213bf60e6799ea}
      \field{labelprefix}{A}
      \field{labelnamesource}{author}
      \field{labeltitlesource}{title}
      \field{annotation}{Assuming RH and the simplicity of the zeros of $\zeta(s)$, the author records the identity \begin{align*} L_0(x) = \frac{\sqrt{x}}{\zeta(\frac{1}{2})}+\sum_{\rho}\frac{\zeta(2\rho)}{\rho\zeta'(\rho)}x^{\rho}+\frac{1}{2\pi i}\int_{1/4-i\infty}^{1/4+i\infty}\frac{\zeta(2s)}{s\zeta(s)}x^{s} \, ds. \end{align*} By a more delicate argument, he establishes the more explicit series representation \[ \frac{1}{2\pi i}\int_{1/4-i\infty}^{1/4+i\infty}\frac{\zeta(2s)}{s\zeta(s)}x^{s} \, ds = 2\sum_{n=1}^\infty \frac{q(n)\lambda(n)}n \big( C(\sqrt{nx}) + S(\sqrt{nx}) \big), \] where $q(n)$ is the largest integer whose square divides~$n$, and $C(t)$ and $S(t)$ are the normalized Fresnel integrals, which both tend to $\frac12$ as $t\to\infty$. With these identities, the author indicates a possible approach to resolving Pólya's problem in the negative. \par This article cites~\cites{1909.Landau,1919.Polya,1936.Ingham,MR0006202}.}
      \field{issn}{0024-6115}
      \field{journaltitle}{Proc. London Math. Soc. (3)}
      \field{title}{The explicit formula for {$L_0(x)$}}
      \field{volume}{1}
      \field{year}{1951}
      \true{nocite}
      \field{pages}{86\bibrangedash 103}
      \range{pages}{18}
      \verb{doi}
      \verb 10.1112/plms/s3-1.1.86
      \endverb
      \verb{mrnumber}
      \verb 43841
      \endverb
      \verb{urlraw}
      \verb https://doi-org.ezproxy.library.ubc.ca/10.1112/plms/s3-1.1.86
      \endverb
      \verb{url}
      \verb https://doi-org.ezproxy.library.ubc.ca/10.1112/plms/s3-1.1.86
      \endverb
    \endentry
    \entry{MR51857}{article}{}
      \name{author}{1}{}{%
        {{hash=695e0d70149c6690cca3917ef5b74c9e}{%
           family={Fawaz},
           familyi={F\bibinitperiod},
           given={A.\bibnamedelimi Y.},
           giveni={A\bibinitperiod\bibinitdelim Y\bibinitperiod}}}%
      }
      \strng{namehash}{695e0d70149c6690cca3917ef5b74c9e}
      \strng{fullhash}{695e0d70149c6690cca3917ef5b74c9e}
      \strng{bibnamehash}{695e0d70149c6690cca3917ef5b74c9e}
      \strng{authorbibnamehash}{695e0d70149c6690cca3917ef5b74c9e}
      \strng{authornamehash}{695e0d70149c6690cca3917ef5b74c9e}
      \strng{authorfullhash}{695e0d70149c6690cca3917ef5b74c9e}
      \field{extraname}{2}
      \field{sortinit}{F}
      \field{sortinithash}{fb0c0faa89eb6abae8213bf60e6799ea}
      \field{labelprefix}{A}
      \field{labelnamesource}{author}
      \field{labeltitlesource}{title}
      \field{annotation}{This article is a continuation of the author's previous work~\cite{MR43841}. Defining \[ I(x)=\frac{1}{2\pi i}\int_{a-i\infty}^{a+i\infty}\frac{\zeta(2s)}{s\zeta(s)}x^{s} \, ds \] for any $a \in (0,1/2)$, the author proves that \[ \liminf_{x \to \infty} E^L(x) = -\limsup_{x \to 0^+} \frac{I(x)}{\sqrt x} \quad\text{and}\quad \limsup_{x \to \infty} E^L(x) = -\liminf_{x \to 0^+} \frac{I(x)}{\sqrt x}, \] which suggests an alternative approach to studying $L(x)$. \par This article cites~\cite{MR43841,MR0003644,MR0006202,1932.Ingham,1919.Polya,MR0046485,1936.Ingham}.}
      \field{issn}{0033-5606}
      \field{journaltitle}{Quart. J. Math. Oxford Ser. (2)}
      \field{title}{On an unsolved problem in the analytic theory of numbers}
      \field{volume}{3}
      \field{year}{1952}
      \true{nocite}
      \field{pages}{282\bibrangedash 295}
      \range{pages}{14}
      \verb{doi}
      \verb 10.1093/qmath/3.1.282
      \endverb
      \verb{mrnumber}
      \verb 51857
      \endverb
      \verb{urlraw}
      \verb https://doi.org/10.1093/qmath/3.1.282
      \endverb
      \verb{url}
      \verb https://doi.org/10.1093/qmath/3.1.282
      \endverb
    \endentry
    \entry{MR1806291}{article}{}
      \name{author}{2}{}{%
        {{hash=0644385a54ab456bd85da1f460b152fd}{%
           family={Feuerverger},
           familyi={F\bibinitperiod},
           given={A.},
           giveni={A\bibinitperiod}}}%
        {{hash=feda7d3cba2a0ad4e772307fca655937}{%
           family={Martin},
           familyi={M\bibinitperiod},
           given={G.},
           giveni={G\bibinitperiod}}}%
      }
      \strng{namehash}{ad494ec99168c86b41ad9bf70473e1a4}
      \strng{fullhash}{ad494ec99168c86b41ad9bf70473e1a4}
      \strng{bibnamehash}{ad494ec99168c86b41ad9bf70473e1a4}
      \strng{authorbibnamehash}{ad494ec99168c86b41ad9bf70473e1a4}
      \strng{authornamehash}{ad494ec99168c86b41ad9bf70473e1a4}
      \strng{authorfullhash}{ad494ec99168c86b41ad9bf70473e1a4}
      \field{sortinit}{F}
      \field{sortinithash}{fb0c0faa89eb6abae8213bf60e6799ea}
      \field{labelprefix}{A}
      \field{labelnamesource}{author}
      \field{labeltitlesource}{title}
      \field{annotation}{Assuming GRH and LI, the authors establish a general formula for $\delta_{q;a_1,\ldots,a_r}$ as a linear combination of principal values of integers of dimensions up to $r-1$. As an example, they numerically compute all possible race densities when $q\in\{5,8,12\}$, all densities for two- and three-way races when $q\in\{7,9\}$, and all two-way race densities when $q=11$. They also list some families of symmetries among these densities, as well as some families of provable inequalities for three-way race densities. The authors conjecture that ``two-ties'' configurations are possible in the following sense: for any $1\le i,j,k,\ell\le r$, there exist arbitrarily large $x$ such that $\pi(x;q,a_i) = \pi(x;q,a_j)$ and $\pi(x;q,a_k) = \pi(x;q,a_\ell)$ (even if the relative sizes of all $r$ contestants are prescribed); they further conjecture that analogous ``three-ties'' configurations occur only finitely often. They also make some speculative conjectures about ``bias factors'' and ``order equivalences'' to introduce the idea of comparing densities to one another; they also describe the question of the possible asymptotic sizes of $\delta_{q;a_1,\ldots,a_r}$ if $r$ is a function of~$q$. Finally, they remark that their techniques also resolve a question of Knapowski and Turán~\cite{MR0146156} about simultaneous solutions to $\Delta^\pi(x;q,a_j) > 0$.\par This article cites~\cites{MR1367576,MR0146156,1914.Littlewood,MR1329368,MR0108470}.}
      \field{issn}{1058-6458}
      \field{journaltitle}{Experiment. Math.}
      \field{number}{4}
      \field{title}{Biases in the {S}hanks-{R}ényi prime number race}
      \field{volume}{9}
      \field{year}{2000}
      \true{nocite}
      \field{pages}{535\bibrangedash 570}
      \range{pages}{36}
      \verb{mrnumber}
      \verb 1806291
      \endverb
      \verb{urlraw}
      \verb http://projecteuclid.org/euclid.em/1045759521
      \endverb
      \verb{url}
      \verb http://projecteuclid.org/euclid.em/1045759521
      \endverb
    \endentry
    \entry{MR3103752}{book}{}
      \name{author}{1}{}{%
        {{hash=889ecd9d3a804c5a7e9a5d8287766657}{%
           family={Fiorilli},
           familyi={F\bibinitperiod},
           given={D.},
           giveni={D\bibinitperiod}}}%
      }
      \list{language}{1}{%
        {French}%
      }
      \list{publisher}{1}{%
        {ProQuest LLC, Ann Arbor, MI}%
      }
      \strng{namehash}{889ecd9d3a804c5a7e9a5d8287766657}
      \strng{fullhash}{889ecd9d3a804c5a7e9a5d8287766657}
      \strng{bibnamehash}{889ecd9d3a804c5a7e9a5d8287766657}
      \strng{authorbibnamehash}{889ecd9d3a804c5a7e9a5d8287766657}
      \strng{authornamehash}{889ecd9d3a804c5a7e9a5d8287766657}
      \strng{authorfullhash}{889ecd9d3a804c5a7e9a5d8287766657}
      \field{extraname}{1}
      \field{sortinit}{F}
      \field{sortinithash}{fb0c0faa89eb6abae8213bf60e6799ea}
      \field{labelprefix}{A}
      \field{labelnamesource}{author}
      \field{labeltitlesource}{title}
      \field{annotation}{Chapter 2 of this thesis is the article~\cite{MR3028758}. A result in the introduction explains why the method of~\cite{MR777013} is not capable of producing more sign changes in the $\pi$-vs.-$\li$ race: the author shows that if \[ \psi_k(x) = \frac1{k!} \sum_{n\le x} \Lambda(n) \big( \log \tfrac xn \big)^k = x - \sum_\rho \frac{x^\rho}{\rho^{k+1}} + \sum_{j=0}^k \frac{a_{k-j}}{j!} (\log x)^j \] is the $k$-fold logarithmic average of $\psi(x)$, then $W(\psi_k;T) = \frac{\gamma_1}{2\pi} \log T + O_k(1)$ for $k\ge5$, where $\gamma_1 \approx 14.135$ is the smallest ordinate of a nontrivial zero of $\zeta(s)$. The phenomenon is that the repeated averaging washes out all of the small-scale sign changes expected for $\psi(x)$ itself.}
      \field{isbn}{978-0494-89015-8}
      \field{note}{Thesis (Ph.D.)--Université de Montréal}
      \field{pagetotal}{261}
      \field{title}{Irrégularités dans la distribution des nombres premiers et des suites plus générales dans les progressions arithmétiques}
      \field{year}{2011}
      \true{nocite}
      \verb{mrnumber}
      \verb 3103752
      \endverb
    \endentry
    \entry{MR3264673}{article}{}
      \name{author}{1}{}{%
        {{hash=889ecd9d3a804c5a7e9a5d8287766657}{%
           family={Fiorilli},
           familyi={F\bibinitperiod},
           given={D.},
           giveni={D\bibinitperiod}}}%
      }
      \strng{namehash}{889ecd9d3a804c5a7e9a5d8287766657}
      \strng{fullhash}{889ecd9d3a804c5a7e9a5d8287766657}
      \strng{bibnamehash}{889ecd9d3a804c5a7e9a5d8287766657}
      \strng{authorbibnamehash}{889ecd9d3a804c5a7e9a5d8287766657}
      \strng{authornamehash}{889ecd9d3a804c5a7e9a5d8287766657}
      \strng{authorfullhash}{889ecd9d3a804c5a7e9a5d8287766657}
      \field{extraname}{2}
      \field{sortinit}{F}
      \field{sortinithash}{fb0c0faa89eb6abae8213bf60e6799ea}
      \field{labelprefix}{A}
      \field{labelnamesource}{author}
      \field{labeltitlesource}{title}
      \field{annotation}{In this article, the author extends the results from~\cite{2007.Sarnak} on ``elliptic curve prime number races'', the races between primes for which the number of points on a given elliptic curve over $\mathbb{F}_p$ is greater or less than $p+1$. Assuming GRH and LI for $L$-functions of elliptic curves, the author establishes an equivalence between two phenomena. The first phenomenon is unbounded analytic ranks for elliptic curves---more precisely, the conjecture that $\limsup_{N_E\to\infty}r_{\text{an}}(E)/\sqrt{\log(N_E)}=\infty$ where $N_E$ is the conductor of~$E$. The second phenomenon is the existence of arbitrarily biased elliptic curve prime races, that is, races for which the logarithmic density of the set $\{t:S(t)=-\sum_{p\le t}a_pp^{-1/2} \geq 0\}$ is arbitrarily close to~$1$. When showing that the first phenomenon implies the second, LI can be weakened to the hypothesis that the nonreal zeros of $L$-functions of any elliptic curve have multiplicities bounded by a universal constant. When showing that the second phenomenon implies the first, only LI is required, and in fact the author shows that the existence of arbitrarily biased races under LI would imply GRH for an infinite family of $L(s,E)$.\par This article cites \cite{2007.Sarnak, MR3261965, MR3028758, MR2202918, 1914.Littlewood, 2000.Ng, MR1761276, MR1329368}.}
      \field{issn}{1073-7928}
      \field{journaltitle}{Int. Math. Res. Not. IMRN}
      \field{number}{18}
      \field{title}{Elliptic curves of unbounded rank and {C}hebyshev's bias}
      \field{year}{2014}
      \true{nocite}
      \field{pages}{4997\bibrangedash 5024}
      \range{pages}{28}
      \verb{doi}
      \verb 10.1093/imrn/rnt103
      \endverb
      \verb{mrnumber}
      \verb 3264673
      \endverb
      \verb{urlraw}
      \verb https://doi.org/10.1093/imrn/rnt103
      \endverb
      \verb{url}
      \verb https://doi.org/10.1093/imrn/rnt103
      \endverb
    \endentry
    \entry{MR3272280}{article}{}
      \name{author}{1}{}{%
        {{hash=889ecd9d3a804c5a7e9a5d8287766657}{%
           family={Fiorilli},
           familyi={F\bibinitperiod},
           given={D.},
           giveni={D\bibinitperiod}}}%
      }
      \strng{namehash}{889ecd9d3a804c5a7e9a5d8287766657}
      \strng{fullhash}{889ecd9d3a804c5a7e9a5d8287766657}
      \strng{bibnamehash}{889ecd9d3a804c5a7e9a5d8287766657}
      \strng{authorbibnamehash}{889ecd9d3a804c5a7e9a5d8287766657}
      \strng{authornamehash}{889ecd9d3a804c5a7e9a5d8287766657}
      \strng{authorfullhash}{889ecd9d3a804c5a7e9a5d8287766657}
      \field{extraname}{3}
      \field{sortinit}{F}
      \field{sortinithash}{fb0c0faa89eb6abae8213bf60e6799ea}
      \field{labelprefix}{A}
      \field{labelnamesource}{author}
      \field{labeltitlesource}{title}
      \field{annotation}{It is known, assuming GRH and LI, that $\delta(p;N,R)$ tends to $\frac12$ as the prime modulus $p$ tends to infinity. In contrast, the author shows that the analogous density $\delta(q;N,R)$ takes a set of values that is dense in $(\frac12,1)$, under the same two assumptions (although LI can be weakened to a mild bound on the multiplicities of zeros of $L(s,\chi)$), with large biases corresponding to highly composite moduli in a quantitative sense (including a conjecture for the asymptotic size of the analogues of Skewes's number for these races). The author proves similar results for fairly general linear combinations of reduced residue classes.\par This article cites~\cites{MR3261965, MR1752093, MR1816036, MR1806291, MR3028758, MR2202918, MR1555148, MR1344249, MR2943671, MR1544293A, MR1544293B, MR2631002, MR634679, arXiv.Myerscough, 2000.Ng, MR2078705, MR1329368, MR1761276, MR1573970}.}
      \field{issn}{1937-0652}
      \field{journaltitle}{Algebra Number Theory}
      \field{number}{7}
      \field{title}{Highly biased prime number races}
      \field{volume}{8}
      \field{year}{2014}
      \true{nocite}
      \field{pages}{1733\bibrangedash 1767}
      \range{pages}{35}
      \verb{doi}
      \verb 10.2140/ant.2014.8.1733
      \endverb
      \verb{mrnumber}
      \verb 3272280
      \endverb
      \verb{urlraw}
      \verb https://doi.org/10.2140/ant.2014.8.1733
      \endverb
      \verb{url}
      \verb https://doi.org/10.2140/ant.2014.8.1733
      \endverb
    \endentry
    \entry{MR3356760}{article}{}
      \name{author}{1}{}{%
        {{hash=889ecd9d3a804c5a7e9a5d8287766657}{%
           family={Fiorilli},
           familyi={F\bibinitperiod},
           given={D.},
           giveni={D\bibinitperiod}}}%
      }
      \strng{namehash}{889ecd9d3a804c5a7e9a5d8287766657}
      \strng{fullhash}{889ecd9d3a804c5a7e9a5d8287766657}
      \strng{bibnamehash}{889ecd9d3a804c5a7e9a5d8287766657}
      \strng{authorbibnamehash}{889ecd9d3a804c5a7e9a5d8287766657}
      \strng{authornamehash}{889ecd9d3a804c5a7e9a5d8287766657}
      \strng{authorfullhash}{889ecd9d3a804c5a7e9a5d8287766657}
      \field{extraname}{4}
      \field{sortinit}{F}
      \field{sortinithash}{fb0c0faa89eb6abae8213bf60e6799ea}
      \field{labelprefix}{A}
      \field{labelnamesource}{author}
      \field{labeltitlesource}{title}
      \field{annotation}{The author studies the ``variance'' of primes in arithmetic progressions \[ V(x;q) = \sum_{(a,q)=1} \bigg| \psi(x;q,a) - \frac{\psi(x;\chi_0)}{\phi(q)} \bigg|^2 \] that was first examined by Hooley and is known to be asymptotic to $x\log q$ in certain ranges of~$q$. The author conjectures that $V(x;q) \sim x\log q$ uniformly for $(\log\log x)^{1+\delta} \le q \le x$ for any fixed $\delta>0$, a much wider range than previously anticipated. The conjecture arises from the computation, assuming GRH and LI, of the limiting logarithmic distribution of $V(x;q)$ for a given~$q$ and studying its large deviations.\par This article cites~\cites{2000.Ng,MR0004255,MR0506080,MR634679,MR1329368,MR1806291,MR1956261,MR2078705,MR2631002,MR2943671,MR2998146,MR3028758,MR3063909,MR3264673,MR3272280}.}
      \field{issn}{1073-7928}
      \field{journaltitle}{Int. Math. Res. Not. IMRN}
      \field{number}{12}
      \field{title}{The distribution of the variance of primes in arithmetic progressions}
      \field{year}{2015}
      \true{nocite}
      \field{pages}{4421\bibrangedash 4448}
      \range{pages}{28}
      \verb{doi}
      \verb 10.1093/imrn/rnu074
      \endverb
      \verb{mrnumber}
      \verb 3356760
      \endverb
    \endentry
    \entry{arXiv.FiorilliJouve}{online}{}
      \name{author}{2}{}{%
        {{hash=889ecd9d3a804c5a7e9a5d8287766657}{%
           family={Fiorilli},
           familyi={F\bibinitperiod},
           given={D.},
           giveni={D\bibinitperiod}}}%
        {{hash=dfca5f0fbbbf59635f02988f421953ae}{%
           family={Jouve},
           familyi={J\bibinitperiod},
           given={F.},
           giveni={F\bibinitperiod}}}%
      }
      \strng{namehash}{bdf4ae134537f1f2a294e43fbbf433aa}
      \strng{fullhash}{bdf4ae134537f1f2a294e43fbbf433aa}
      \strng{bibnamehash}{bdf4ae134537f1f2a294e43fbbf433aa}
      \strng{authorbibnamehash}{bdf4ae134537f1f2a294e43fbbf433aa}
      \strng{authornamehash}{bdf4ae134537f1f2a294e43fbbf433aa}
      \strng{authorfullhash}{bdf4ae134537f1f2a294e43fbbf433aa}
      \field{extraname}{1}
      \field{sortinit}{F}
      \field{sortinithash}{fb0c0faa89eb6abae8213bf60e6799ea}
      \field{labelprefix}{A}
      \field{labelnamesource}{author}
      \field{labeltitlesource}{title}
      \field{annotation}{For a class function $t$ on $G=\mathop{\rm Gal}(L/K)$, the authors consider the Frobenius counting function \[ \pi(x;L/K,t) = \sum_{\substack{\mathfrak{p}\subset \mathcal{O}_K \\ {\rm N}\mathfrak{p}\leq x}} t(\text{Frob}_\mathfrak{p}) \] and study the ``normalized Chebotarev error term'' \[ E(x;L/K,t) = x^{-\Theta}\log x\cdot (\pi(x;L/K,t) - \widehat{t}(1)\text{Li}(x)), \] where $\Theta$ is the supremum of real parts of zeros of certain Artin $L$-functions of $L/\mathbb{Q}$ relevant to~$t$. They show that this error admits a limiting logarithmic distribution, and they compute its mean, along with unconditional and conditional bounds on its variance. They also prove improved bounds on the error term of the Chebotarev density theorem assuming GRH and Artin's conjecture. They further consider the logarithmic densities $\delta(L/K;t)$ of the set of values of~$x$ for which $E(x;L/K,t)$ is positive, giving criteria for when the density is close to 1 or $\frac{1}{2}$, indicating a bias or lack thereof. Finally, these ideas are applied to the class function $t_{C_1,C_2} = \frac{|G|}{|C_1|}1_{C_1}-\frac{|G|}{|C_2|}1_{C_2}$ for special families of extensions ($S_n$, dihedral, radical, abelian, and Hilbert class fields of quadratic extensions), giving asymptotic bounds for the densities in each case.\par This article cites~\cites{MR3261965,MR2474313,MR2811555,MR3581815,MR4120786,MR4316675,MR3521338,MR3272280,MR3264673,MR3028758,MR2778652,MR2202918,MR1344249,MR1310738,MR0146156,MR0146157,MR0146158,MR2943671,MR3624386,1914.Littlewood,MR4082248,MR3803705,MR2631002,2000.Ng,MR1761276,MR1329368,2007.Sarnak,MR0004255,MR4265352,MR4194167} and this bibliography.}
      \field{note}{J. Inst. Math. Jussieu, to appear}
      \field{title}{Distributions of Frobenius elements in families of Galois extensions}
      \field{year}{2023}
      \true{nocite}
      \verb{urlraw}
      \verb https://arxiv.org/abs/2001.05428
      \endverb
      \verb{url}
      \verb https://arxiv.org/abs/2001.05428
      \endverb
    \endentry
    \entry{MR3028758}{article}{}
      \name{author}{2}{}{%
        {{hash=889ecd9d3a804c5a7e9a5d8287766657}{%
           family={Fiorilli},
           familyi={F\bibinitperiod},
           given={D.},
           giveni={D\bibinitperiod}}}%
        {{hash=feda7d3cba2a0ad4e772307fca655937}{%
           family={Martin},
           familyi={M\bibinitperiod},
           given={G.},
           giveni={G\bibinitperiod}}}%
      }
      \strng{namehash}{8f8b9fef368ad32f0c071fbab617f456}
      \strng{fullhash}{8f8b9fef368ad32f0c071fbab617f456}
      \strng{bibnamehash}{8f8b9fef368ad32f0c071fbab617f456}
      \strng{authorbibnamehash}{8f8b9fef368ad32f0c071fbab617f456}
      \strng{authornamehash}{8f8b9fef368ad32f0c071fbab617f456}
      \strng{authorfullhash}{8f8b9fef368ad32f0c071fbab617f456}
      \field{extraname}{1}
      \field{sortinit}{F}
      \field{sortinithash}{fb0c0faa89eb6abae8213bf60e6799ea}
      \field{labelprefix}{A}
      \field{labelnamesource}{author}
      \field{labeltitlesource}{title}
      \field{annotation}{Assuming GRH and LI, the authors establish an asymptotic series for $\delta_{q;a,b}$ whose error term can be taken to be any negative power of~$q$. The first asymptotic formula given by this series, in the case where~$a$ is a quadratic nonresidue and~$b$ is a quadratic residue\mod q, is \[ \delta_{q;a,b} = \frac12 + \frac{\rho(q)}{\sqrt{2\pi V(q;a,b)}} + O_\ep(q^{-3/2+\ep}), \] where $V(q;a,b) = \sum_{\chi\mod q} |\chi(b)-\chi(a)|^2\sum_\rho 1/(\frac14+\gamma^2)$ and $\rho(q)$ is the number of square roots of~$1\mod q$. The authors give a closed form for $V(q;a,b)$ in terms of arithmetic properties of its arguments, which allow them to compare various values of $\delta_{q;a,b}$; for example, for fixed $a\ne-1$, they show (always assuming GRH and LI) that when~$q$ is sufficiently large, $\delta_{q;-1,1} < \delta_{q;a,1}$ whenever~$-1$ and~$a$ are both quadratic nonresidues\mod q. Explicit bounds for the error terms in this formula allow the authors to calculate the list of all~$117$ two-way prime number races (up to symmetries) for which $\delta_{q;a,b} > 0.9$, the largest of which is $\delta_{24,5,1} \approx 0.999988$.\par This article cites~\cites{MR0108470,MR0137689,MR0506080,MR1329368,MR1806291,MR1816036}.}
      \field{issn}{0075-4102}
      \field{journaltitle}{J. Reine Angew. Math.}
      \field{title}{Inequities in the {S}hanks-{R}ényi prime number race: an asymptotic formula for the densities}
      \field{volume}{676}
      \field{year}{2013}
      \true{nocite}
      \field{pages}{121\bibrangedash 212}
      \range{pages}{92}
      \verb{mrnumber}
      \verb 3028758
      \endverb
    \endentry
    \entry{MR4400872}{article}{}
      \name{author}{2}{}{%
        {{hash=c511fbe6d97a5da2025e6199f40ab277}{%
           family={Fiorilli},
           familyi={F\bibinitperiod},
           given={Daniel},
           giveni={D\bibinitperiod}}}%
        {{hash=8e65f0b511fecc96916507c9bdad4520}{%
           family={Jouve},
           familyi={J\bibinitperiod},
           given={Florent},
           giveni={F\bibinitperiod}}}%
      }
      \strng{namehash}{927d699f5bb8bac56bfe890ba7e7af7f}
      \strng{fullhash}{927d699f5bb8bac56bfe890ba7e7af7f}
      \strng{bibnamehash}{927d699f5bb8bac56bfe890ba7e7af7f}
      \strng{authorbibnamehash}{927d699f5bb8bac56bfe890ba7e7af7f}
      \strng{authornamehash}{927d699f5bb8bac56bfe890ba7e7af7f}
      \strng{authorfullhash}{927d699f5bb8bac56bfe890ba7e7af7f}
      \field{extraname}{2}
      \field{sortinit}{F}
      \field{sortinithash}{fb0c0faa89eb6abae8213bf60e6799ea}
      \field{labelprefix}{A}
      \field{labelnamesource}{author}
      \field{labeltitlesource}{title}
      \field{annotation}{If $L/K$ is a Galois extension of number fields and~$C$ is a conjugacy class in $\operatorname{Gal}(L/K)$, then $\pi(x; C, L/K)$ denotes the number of unramified prime ideals $\mathfrak{p}$ of $L$ with norm ${N(\mathfrak{p}) \leq x}$ whose Frobenius conjugacy class $\text{Frob}_{\mathfrak{p}}$ equals~$C$. For two conjugacy classes $C_1, \, C_2$, let $\delta_{L/K; C_1, C_2}$ be the logarithmic density of the set $\bigl\{x\geq 1\colon \frac1{\#C_1}\pi(x; C_1, L/K) > \frac1{\#C_2}\pi(x; C_2, L/K) \bigr\}$. The authors unconditionally show that there are infinitely many Galois extensions $L/K$ and conjugacy classes $C_1, \, C_2$ such that $\delta_{L/K; C_1, C_2}=1$.\par This article cites~\cites{1853.Chebyshev, MR1344249, 2000.Ng,arXiv.FiorilliJouve,MR1329368}.}
      \field{issn}{2429-7100}
      \field{journaltitle}{J. Éc. polytech. Math.}
      \field{title}{Unconditional {C}hebyshev biases in number fields}
      \field{volume}{9}
      \field{year}{2022}
      \true{nocite}
      \field{pages}{671\bibrangedash 679}
      \range{pages}{9}
      \verb{doi}
      \verb 10.5802/jep.19
      \endverb
      \verb{mrnumber}
      \verb 4400872
      \endverb
      \verb{urlraw}
      \verb https://doi.org/10.5802/jep.19
      \endverb
      \verb{url}
      \verb https://doi.org/10.5802/jep.19
      \endverb
    \endentry
    \entry{arXiv.FM}{online}{}
      \name{author}{2}{}{%
        {{hash=c511fbe6d97a5da2025e6199f40ab277}{%
           family={{Fiorilli}},
           familyi={F\bibinitperiod},
           given={Daniel},
           giveni={D\bibinitperiod}}}%
        {{hash=c06cf923cb7824f2455f1750eec3b292}{%
           family={{Martin}},
           familyi={M\bibinitperiod},
           given={Greg},
           giveni={G\bibinitperiod}}}%
      }
      \strng{namehash}{887572e6be3fbbeb636b6dec60af9d05}
      \strng{fullhash}{887572e6be3fbbeb636b6dec60af9d05}
      \strng{bibnamehash}{887572e6be3fbbeb636b6dec60af9d05}
      \strng{authorbibnamehash}{887572e6be3fbbeb636b6dec60af9d05}
      \strng{authornamehash}{887572e6be3fbbeb636b6dec60af9d05}
      \strng{authorfullhash}{887572e6be3fbbeb636b6dec60af9d05}
      \field{extraname}{2}
      \field{sortinit}{F}
      \field{sortinithash}{fb0c0faa89eb6abae8213bf60e6799ea}
      \field{labelprefix}{A}
      \field{labelnamesource}{author}
      \field{labeltitlesource}{title}
      \field{annotation}{Define $G(x;q) = \sum_{(a,q)=1} \bigl( \Delta^\theta(x;q,a)/\phi(q) \bigr)^2$. When~$q$ is fixed, results of the form $E(x,\chi) = \Omega(\log\log\log x)$ imply that $G(x;q) \not\ll x\log q$; such oscillation results were proved by Littlewood~\cite{1914.Littlewood} for certain characters, and extended by G.~Davidoff (unpublished) to all real nonprincipal characters. Hooley conjectured that $G(x;q) \ll x\log q$ as soon as~$q$ tends to infinity with~$x$. In this article, the authors disprove this conjecture. While the basic approach using Diophantine approximation and the explicit formula is the same as in earlier results, the need to have estimates that are uniform in~$q$ requires significant technical attention.\par This article cites~\cites{MR3356760, MR3028758, MR0506080, MR858395, 1914.Littlewood,MR820251, MR1329368, MR3871592}.}
      \field{note}{J. Eur. Math. Soc. (JEMS), to appear}
      \field{title}{A disproof of Hooley's conjecture}
      \field{year}{2023}
      \true{nocite}
      \verb{urlraw}
      \verb https://arxiv.org/abs/2008.05837
      \endverb
      \verb{url}
      \verb https://arxiv.org/abs/2008.05837
      \endverb
    \endentry
    \entry{MR1985941}{incollection}{}
      \name{author}{2}{}{%
        {{hash=e7e31df1803902b56a78f32b1c733598}{%
           family={Ford},
           familyi={F\bibinitperiod},
           given={K.},
           giveni={K\bibinitperiod}}}%
        {{hash=c48086b9b194fded979a790687062dbb}{%
           family={Konyagin},
           familyi={K\bibinitperiod},
           given={S.},
           giveni={S\bibinitperiod}}}%
      }
      \list{publisher}{1}{%
        {Mosk. Gos. Univ. im. Lomonosova, Mekh.-Mat. Fak., Moscow}%
      }
      \strng{namehash}{3cc0b23765677fa7230ff0cbba627702}
      \strng{fullhash}{3cc0b23765677fa7230ff0cbba627702}
      \strng{bibnamehash}{3cc0b23765677fa7230ff0cbba627702}
      \strng{authorbibnamehash}{3cc0b23765677fa7230ff0cbba627702}
      \strng{authornamehash}{3cc0b23765677fa7230ff0cbba627702}
      \strng{authorfullhash}{3cc0b23765677fa7230ff0cbba627702}
      \field{extraname}{1}
      \field{sortinit}{F}
      \field{sortinithash}{fb0c0faa89eb6abae8213bf60e6799ea}
      \field{labelprefix}{A}
      \field{labelnamesource}{author}
      \field{labeltitlesource}{title}
      \field{annotation}{This article describes nine families of problems that are central to the study of comparative prime number theory. The first eight are taken from or inspired by the problems listed by Knapowski and Turán in~\cite{MR0146156}; the ninth problem, entitled ``Union-problems'', examines the distribution of $\frac1{\#\cA} \pi(x;q,\cA) - \frac1{\#\cB} \pi(x;q,\cB)$. Throughout the rest of the article, the authors provide an overview of what is already known about the first seven problems.\par This article cites~\cites{1853.Chebyshev, MR1511360, 1914.Littlewood, MR1555148, MR1544293A, MR1544293B, MR1573970, 1938.Wintner, MR0004255, MR0006202, MR0067145, MR0083001, MR0108470, MR0146156, MR0146157, MR0146158, MR0146159, MR0146160, MR0146161, MR0156826, MR0156827, MR0176967, MR0162771, MR0174538, MR0180539, MR0182616A, MR0182616B, MR0202686, MR0200250, MR0202685, MR0330069, MR0332684, MR0332685, MR0289452, MR0302585, MR0447090, MR0506080, MR0472726, MR0476615, MR529694, MR552467, MR634679, MR866118, MR1083000, MR1105608, MR1251926, MR1329368, MR1310738, MR1367576, MR1651741, MR1806291, MR1816036, MR1862054, MR1909220, MR1956261, MR1957110}.}
      \field{booktitle}{I{V} {I}nternational {C}onference ``{M}odern {P}roblems of {N}umber {T}heory and its {A}pplications'': {C}urrent {P}roblems, {P}art {II} ({R}ussian) ({T}ula, 2001)}
      \field{title}{Chebyshev's conjecture and the prime number race}
      \field{year}{2002}
      \true{nocite}
      \field{pages}{67\bibrangedash 91}
      \range{pages}{25}
      \verb{mrnumber}
      \verb 1985941
      \endverb
    \endentry
    \entry{MR1909220}{article}{}
      \name{author}{2}{}{%
        {{hash=e7e31df1803902b56a78f32b1c733598}{%
           family={Ford},
           familyi={F\bibinitperiod},
           given={K.},
           giveni={K\bibinitperiod}}}%
        {{hash=c48086b9b194fded979a790687062dbb}{%
           family={Konyagin},
           familyi={K\bibinitperiod},
           given={S.},
           giveni={S\bibinitperiod}}}%
      }
      \strng{namehash}{3cc0b23765677fa7230ff0cbba627702}
      \strng{fullhash}{3cc0b23765677fa7230ff0cbba627702}
      \strng{bibnamehash}{3cc0b23765677fa7230ff0cbba627702}
      \strng{authorbibnamehash}{3cc0b23765677fa7230ff0cbba627702}
      \strng{authornamehash}{3cc0b23765677fa7230ff0cbba627702}
      \strng{authorfullhash}{3cc0b23765677fa7230ff0cbba627702}
      \field{extraname}{2}
      \field{sortinit}{F}
      \field{sortinithash}{fb0c0faa89eb6abae8213bf60e6799ea}
      \field{labelprefix}{A}
      \field{labelnamesource}{author}
      \field{labeltitlesource}{title}
      \field{annotation}{This article introduces the term ``barrier'' for a hypothetical configuration of zeros of Dirichlet $L$-functions that causes prime number races not to be exhaustive. The authors show (through several complementary constructions) that for every three-way prime number races, there is a finite barrier that prevents at least one of the six possible orderings from occurring for large~$x$; moreover, the zeros in these barriers can be arbitrarily close to the critical line and arbitrarily far from the real axis. While most of their constructions of barriers involve zeros with linearly dependent ordinates, the authors construct in the final section a barrier with linearly independent ordinates.\par This article cites~\cites{MR0476616,1853.Chebyshev,MR1251926,MR1310738,MR1367576,MR0146156,MR0146157,MR0146158,MR0146159,MR0146160,MR0146161,MR0156826,MR0156827,MR0162771,MR0174538,MR0180539,MR0182616A,MR0182616B,MR0200250,1914.Littlewood,MR1329368,MR0108470}.}
      \field{issn}{0012-7094}
      \field{journaltitle}{Duke Math. J.}
      \field{number}{2}
      \field{title}{The prime number race and zeros of {$L$}-functions off the critical line}
      \field{volume}{113}
      \field{year}{2002}
      \true{nocite}
      \field{pages}{313\bibrangedash 330}
      \range{pages}{18}
      \verb{doi}
      \verb 10.1215/S0012-7094-02-11324-6
      \endverb
      \verb{mrnumber}
      \verb 1909220
      \endverb
      \verb{urlraw}
      \verb https://doi.org/10.1215/S0012-7094-02-11324-6
      \endverb
      \verb{url}
      \verb https://doi.org/10.1215/S0012-7094-02-11324-6
      \endverb
    \endentry
    \entry{MR2075622}{article}{}
      \name{author}{2}{}{%
        {{hash=e7e31df1803902b56a78f32b1c733598}{%
           family={Ford},
           familyi={F\bibinitperiod},
           given={K.},
           giveni={K\bibinitperiod}}}%
        {{hash=c48086b9b194fded979a790687062dbb}{%
           family={Konyagin},
           familyi={K\bibinitperiod},
           given={S.},
           giveni={S\bibinitperiod}}}%
      }
      \list{publisher}{1}{%
        {Univ. Bonn, Bonn}%
      }
      \strng{namehash}{3cc0b23765677fa7230ff0cbba627702}
      \strng{fullhash}{3cc0b23765677fa7230ff0cbba627702}
      \strng{bibnamehash}{3cc0b23765677fa7230ff0cbba627702}
      \strng{authorbibnamehash}{3cc0b23765677fa7230ff0cbba627702}
      \strng{authornamehash}{3cc0b23765677fa7230ff0cbba627702}
      \strng{authorfullhash}{3cc0b23765677fa7230ff0cbba627702}
      \field{extraname}{3}
      \field{sortinit}{F}
      \field{sortinithash}{fb0c0faa89eb6abae8213bf60e6799ea}
      \field{labelprefix}{A}
      \field{labelnamesource}{author}
      \field{labeltitlesource}{title}
      \field{annotation}{The authors continue to explore ``barriers'' for various statements in comparative prime number theory. In this article they make the distinction between the barrier itself, which is a multiset of complex numbers, and the consequences to orderings of prime-counting functions (or functions that are sufficiently close to such) that would follow from the $L(s,\chi)$ having their rightmost zeros precisely at the elements of the barrier. After establishing some lemmas on values of trigonometric polynomials, the authors prove several results concerning whether or not bounded (or finite) barriers exist for various races of the form $\pi(x;q,1)$ against many other $\pi(x;q,a)$. They show that if every pair from a set of $r$ functions is exhaustive, then the total number of $r$-way orderings is at least $r(r-1)/2+1$; they investigate whether barriers can exist that limit the number of orderings to this minimal value (``extremal barriers''), and construct barriers that do force at most $r(r-1)$ orderings.\par This article cites~\cites{MR1909220,MR1985941,MR1083000,MR1251926,MR0146156,MR0146157,MR0146158,MR0146159,MR0146160,MR0146161,MR0156826,MR0156827,MR0162771,MR0174538,MR0180539,MR0182616A,MR0182616B,MR0200250,1914.Littlewood}.}
      \field{booktitle}{Proceedings of the {S}ession in {A}nalytic {N}umber {T}heory and {D}iophantine {E}quations}
      \field{series}{Bonner Math. Schriften}
      \field{title}{The prime number race and zeros of {$L$}-functions off the critical line. {II}}
      \field{volume}{360}
      \field{year}{2003}
      \true{nocite}
      \field{pages}{40}
      \range{pages}{1}
      \verb{mrnumber}
      \verb 2075622
      \endverb
    \endentry
    \entry{MR3151605}{article}{}
      \name{author}{3}{}{%
        {{hash=e7e31df1803902b56a78f32b1c733598}{%
           family={Ford},
           familyi={F\bibinitperiod},
           given={K.},
           giveni={K\bibinitperiod}}}%
        {{hash=a01ad5bbada0f5c69f81f1149c07aefc}{%
           family={Lamzouri},
           familyi={L\bibinitperiod},
           given={Y.},
           giveni={Y\bibinitperiod}}}%
        {{hash=c48086b9b194fded979a790687062dbb}{%
           family={Konyagin},
           familyi={K\bibinitperiod},
           given={S.},
           giveni={S\bibinitperiod}}}%
      }
      \strng{namehash}{63af0df95eb4cedba38dcfa76456f7cf}
      \strng{fullhash}{63af0df95eb4cedba38dcfa76456f7cf}
      \strng{bibnamehash}{63af0df95eb4cedba38dcfa76456f7cf}
      \strng{authorbibnamehash}{63af0df95eb4cedba38dcfa76456f7cf}
      \strng{authornamehash}{63af0df95eb4cedba38dcfa76456f7cf}
      \strng{authorfullhash}{63af0df95eb4cedba38dcfa76456f7cf}
      \field{sortinit}{F}
      \field{sortinithash}{fb0c0faa89eb6abae8213bf60e6799ea}
      \field{labelprefix}{A}
      \field{labelnamesource}{author}
      \field{labeltitlesource}{title}
      \field{annotation}{The authors construct ``barriers'' for two-way prime number races, including $\pi(x) - \li(x)$. They show that a certain hypothetical configuration of zeros of $L(s,\chi)$ would imply that $\{x \geq 2 \colon \pi(x;q,b)>\pi(x;q,a)\}$ has asymptotic density~$0$. Furthermore, they show that a certain hypothetical configuration of the zeros of $\zeta(s)$ would imply that $\{x \geq 2 \colon \pi(x) > \li(x)\}$ has asymptotic density~$0$, while a related configuration would result in the same set having asymptotic density~$1$.\par This article cites~\cites{MR1909220,MR1985941,MR2075622,MR2202918,MR1251926,MR1256149,MR0146156,MR0146157,MR0146158,MR0146159,MR0146160,MR0146161,MR0156826,MR0156827,MR0162771,MR0174538,MR0180539,MR0182616A,MR0182616B,MR0200250,MR0302585,1914.Littlewood,MR1329368,MR2753165}.}
      \field{issn}{0033-5606}
      \field{journaltitle}{Q. J. Math.}
      \field{number}{4}
      \field{title}{The prime number race and zeros of {D}irichlet {$L$}-functions off the critical line: {P}art {III}}
      \field{volume}{64}
      \field{year}{2013}
      \true{nocite}
      \field{pages}{1091\bibrangedash 1098}
      \range{pages}{8}
      \verb{doi}
      \verb 10.1093/qmath/has021
      \endverb
      \verb{mrnumber}
      \verb 3151605
      \endverb
      \verb{urlraw}
      \verb https://doi.org/10.1093/qmath/has021
      \endverb
      \verb{url}
      \verb https://doi.org/10.1093/qmath/has021
      \endverb
    \endentry
    \entry{MR2778652}{article}{}
      \name{author}{2}{}{%
        {{hash=e7e31df1803902b56a78f32b1c733598}{%
           family={Ford},
           familyi={F\bibinitperiod},
           given={K.},
           giveni={K\bibinitperiod}}}%
        {{hash=ee6277cd0f78c9744eaedab7ee5add23}{%
           family={Sneed},
           familyi={S\bibinitperiod},
           given={J.},
           giveni={J\bibinitperiod}}}%
      }
      \strng{namehash}{06ab6a04f799f3b33364b611b68406f5}
      \strng{fullhash}{06ab6a04f799f3b33364b611b68406f5}
      \strng{bibnamehash}{06ab6a04f799f3b33364b611b68406f5}
      \strng{authorbibnamehash}{06ab6a04f799f3b33364b611b68406f5}
      \strng{authornamehash}{06ab6a04f799f3b33364b611b68406f5}
      \strng{authorfullhash}{06ab6a04f799f3b33364b611b68406f5}
      \field{sortinit}{F}
      \field{sortinithash}{fb0c0faa89eb6abae8213bf60e6799ea}
      \field{labelprefix}{A}
      \field{labelnamesource}{author}
      \field{labeltitlesource}{title}
      \field{annotation}{The authors examine Chebyshev's bias for integers which are the product of two primes. Let $\pi_2(x;q,a)$ denote the number of integers $n$ up to $x$ such that $n\equiv a\mod q$ and $\Omega(n)=2$.\par Assume GRH for Dirichlet $L$-functions modulo $q$ and also that the zeros of $L(s,\chi)$ are simple for each nonprincipal character $\chi\mod q$. If $f(x_1,\ldots,x_r)$ is the logarithmic density function of $\big( E(x;q,a_1,b_1),\ldots,E(x;q,a_r,b_r) \big),$ the authors show that the logarithmic density function of $\big( E_2(x;q,a_1,b_1),\ldots,E_2(x;q,a_r,b_r) \big)$ is \[f\left(\frac{c_q(b_1)-c_q(a_1)}2-x_1,\ldots,\frac{c_q(b_r)-c_q(a_r)}2-x_r\right).\] Consequently, assuming both GRH and LI for modulus $q$, the authors show that $\delta_2\big( \pi_2(x;q,a),\pi_2(x;q,b) \big)$ exists, and equals $\frac12$ if $a$ and $b$ are both quadratic residues or both quadratic nonresidues\mod q. Otherwise, if $a$ is a quadratic nonresidue and $b$ is a quadratic residue, then \[1-\delta_{q;a_1,a_2}<\delta_2\big( \pi_2(x;q,a),\pi_2(x;q,b) \big)<\frac 12.\]\par This article cites~\cites{1853.Chebyshev, 1914.Littlewood, MR0083001, MR0146156, MR1329368, MR1985941, MR2202918, MR2753165, MR3028758}.}
      \field{issn}{1058-6458}
      \field{journaltitle}{Experiment. Math.}
      \field{number}{4}
      \field{title}{Chebyshev's bias for products of two primes}
      \field{volume}{19}
      \field{year}{2010}
      \true{nocite}
      \field{pages}{385\bibrangedash 398}
      \range{pages}{14}
      \verb{doi}
      \verb 10.1080/10586458.2010.10390630
      \endverb
      \verb{mrnumber}
      \verb 2778652
      \endverb
      \verb{urlraw}
      \verb https://doi.org/10.1080/10586458.2010.10390630
      \endverb
      \verb{url}
      \verb https://doi.org/10.1080/10586458.2010.10390630
      \endverb
    \endentry
    \entry{MR3961320}{article}{}
      \name{author}{3}{}{%
        {{hash=bb758f27d76cefc7bd5244190e91843b}{%
           family={Ford},
           familyi={F\bibinitperiod},
           given={Kevin},
           giveni={K\bibinitperiod}}}%
        {{hash=d6e30aa748c8fb70505ca61ee684d871}{%
           family={Harper},
           familyi={H\bibinitperiod},
           given={Adam\bibnamedelima J.},
           giveni={A\bibinitperiod\bibinitdelim J\bibinitperiod}}}%
        {{hash=64363cb2afef9e33f967130adc38684b}{%
           family={Lamzouri},
           familyi={L\bibinitperiod},
           given={Youness},
           giveni={Y\bibinitperiod}}}%
      }
      \strng{namehash}{edca729ace4d98a40acfda652484ce64}
      \strng{fullhash}{edca729ace4d98a40acfda652484ce64}
      \strng{bibnamehash}{edca729ace4d98a40acfda652484ce64}
      \strng{authorbibnamehash}{edca729ace4d98a40acfda652484ce64}
      \strng{authornamehash}{edca729ace4d98a40acfda652484ce64}
      \strng{authorfullhash}{edca729ace4d98a40acfda652484ce64}
      \field{sortinit}{F}
      \field{sortinithash}{fb0c0faa89eb6abae8213bf60e6799ea}
      \field{labelprefix}{A}
      \field{labelnamesource}{author}
      \field{labeltitlesource}{title}
      \field{annotation}{The authors show that prime race densities $\delta_{q;a_1,\ldots,a_r}$ can have size significantly different from $\frac1{r!}$ when~$r$ is sufficiently large in terms of~$q$. More precisely, they show that there exists $\eta>0$ and $a_1,\ldots,a_r\mod q$ such that \[ \delta_{q;a_1,\ldots,a_r} \le \exp\bigg( - \frac{\eta\min\{n,\phi(q)^{1/50}\}}{\log q} \bigg) \frac1{r!}, \] and an analogous lower bound for large values without the negative sign on the right-hand side. They derive this result from a similar result concerning the density of real numbers~$x$ such that $\pi(x;q,a_1) > \cdots > \pi(x;q,a_k)$ are the~$k$ largest of the~$r$ prime counting functions while simultaneously $\pi(x;q,a_{k+1}) < \cdots < \pi(x;q,a_{2k})$ are the~$k$ smallest. The authors exploit the fact that certain $E(x;q,a_j)$ are known to have large correlations of size $\Omega_-(1/\log q)$; they also develop a comparison theorem to multivariate normal distributions with a relative rather than absolute error.\par This article cites~\cite{MR1985941,MR0506080,MR1329368,MR1806291,MR2075622,MR2202918,MR2998146,MR3028758,MR3063909,MR3272280,MR3773805}, as well as an earlier draft of this annotated bibliography.}
      \field{issn}{0025-5831}
      \field{journaltitle}{Math. Ann.}
      \field{number}{1-2}
      \field{title}{Extreme biases in prime number races with many contestants}
      \field{volume}{374}
      \field{year}{2019}
      \true{nocite}
      \field{pages}{517\bibrangedash 551}
      \range{pages}{35}
      \verb{doi}
      \verb 10.1007/s00208-019-01810-x
      \endverb
      \verb{mrnumber}
      \verb 3961320
      \endverb
      \verb{urlraw}
      \verb https://doi.org/10.1007/s00208-019-01810-x
      \endverb
      \verb{url}
      \verb https://doi.org/10.1007/s00208-019-01810-x
      \endverb
    \endentry
    \entry{MR1862054}{article}{}
      \name{author}{2}{}{%
        {{hash=bb758f27d76cefc7bd5244190e91843b}{%
           family={Ford},
           familyi={F\bibinitperiod},
           given={Kevin},
           giveni={K\bibinitperiod}}}%
        {{hash=9e0609c9931c1b620df97e8c58d07532}{%
           family={Hudson},
           familyi={H\bibinitperiod},
           given={Richard\bibnamedelima H.},
           giveni={R\bibinitperiod\bibinitdelim H\bibinitperiod}}}%
      }
      \strng{namehash}{3bfafe737308f5813c05132ae0dd9deb}
      \strng{fullhash}{3bfafe737308f5813c05132ae0dd9deb}
      \strng{bibnamehash}{3bfafe737308f5813c05132ae0dd9deb}
      \strng{authorbibnamehash}{3bfafe737308f5813c05132ae0dd9deb}
      \strng{authornamehash}{3bfafe737308f5813c05132ae0dd9deb}
      \strng{authorfullhash}{3bfafe737308f5813c05132ae0dd9deb}
      \field{sortinit}{F}
      \field{sortinithash}{fb0c0faa89eb6abae8213bf60e6799ea}
      \field{labelprefix}{A}
      \field{labelnamesource}{author}
      \field{labeltitlesource}{title}
      \field{annotation}{The authors generalize Lehman's method~\cite{MR0202686} to arithmetic progressions, enabling them to give new results on the location of negative values of $\pi(x;q,b,1)$ for moduli $q \mid 24$. They show that each $\pi(x;8,b,1)$ is negative for some $x < 5 \cdot 10^{19}$, each $\pi(x;12,b,1)$ is negative for some $x <10^{84}$, and each $\pi(x;24,b,1)$ is negative for some $x<10^{253}$. Further, if GRH$(630,000)$ holds for $L(s,\chi_{-4})$, then $\pi(x;4,1,3) > \sqrt{x}/\log x$ for some~$x \approx 8\cdot10^{34,171}$. \par This article cites \cites{MR0202686, MR0447090, MR0476616, MR0472726, MR1752093,MR1816036,MR1251926,MR1256149,MR1367576,MR0146159,MR0068565,MR0083001,MR0137689,MR1329368,MR0108470,MR1573970, MR0067145}.}
      \field{issn}{0065-1036}
      \field{journaltitle}{Acta Arith.}
      \field{number}{4}
      \field{title}{Sign changes in {$\pi _{q,a}(x)-\pi _{q,b}(x)$}}
      \field{volume}{100}
      \field{year}{2001}
      \true{nocite}
      \field{pages}{297\bibrangedash 314}
      \range{pages}{18}
      \verb{doi}
      \verb 10.4064/aa100-4-1
      \endverb
      \verb{mrnumber}
      \verb 1862054
      \endverb
    \endentry
    \entry{MR974088}{article}{}
      \name{author}{1}{}{%
        {{hash=ed673dd9786336c3366d68e7f72549d4}{%
           family={Fujii},
           familyi={F\bibinitperiod},
           given={A.},
           giveni={A\bibinitperiod}}}%
      }
      \strng{namehash}{ed673dd9786336c3366d68e7f72549d4}
      \strng{fullhash}{ed673dd9786336c3366d68e7f72549d4}
      \strng{bibnamehash}{ed673dd9786336c3366d68e7f72549d4}
      \strng{authorbibnamehash}{ed673dd9786336c3366d68e7f72549d4}
      \strng{authornamehash}{ed673dd9786336c3366d68e7f72549d4}
      \strng{authorfullhash}{ed673dd9786336c3366d68e7f72549d4}
      \field{sortinit}{F}
      \field{sortinithash}{fb0c0faa89eb6abae8213bf60e6799ea}
      \field{labelprefix}{A}
      \field{labelnamesource}{author}
      \field{labeltitlesource}{title}
      \field{annotation}{The author shows if $0<\alpha\le4$, then the statement $\lim_{x\to\infty}\sum_p \chi_{-4}(p)e^{-(p/x)^\alpha}= - \infty$ is equivalent to GRH for $L(s,\chi_{-4})$, generalizing earlier results for $\alpha=1$. Further, the author defines $\xi(x,k)$ by $\Gamma(s)^k=\int_0^{\infty}x^{s-1}\xi(x,k)\, dx$ and shows the statement $\lim_{x\to\infty}\sum_p \chi_{-4}(p) \log{p}\cdot\xi(p/x,2)=-\infty$ is also equivalent to GRH for $L(s,\chi_{-4})$. The author conjectures that the first equivalence holds for any $\alpha>0$ and the second (without the factor $\log p$) for any integer $k\ge1$. He also remarks that the method can be used for other weight functions involving exponentials and Bessel functions, including Knapowski--Turán's function $e^{-\log^2(p/x)}$. \par This article cites~\cites{MR595317,MR1544293A,MR1544293B}.}
      \field{issn}{0386-2194}
      \field{journaltitle}{Proc. Japan Acad. Ser. A Math. Sci.}
      \field{number}{7}
      \field{title}{Some generalizations of {C}hebyshev's conjecture}
      \field{volume}{64}
      \field{year}{1988}
      \true{nocite}
      \field{pages}{260\bibrangedash 263}
      \range{pages}{4}
      \verb{mrnumber}
      \verb 974088
      \endverb
      \verb{urlraw}
      \verb http://projecteuclid.org/euclid.pja/1195513180
      \endverb
      \verb{url}
      \verb http://projecteuclid.org/euclid.pja/1195513180
      \endverb
    \endentry
    \entry{MR598886}{article}{}
      \name{author}{1}{}{%
        {{hash=e38330a7a234159955b75bdfae8cd3cd}{%
           family={Gallagher},
           familyi={G\bibinitperiod},
           given={P.\bibnamedelimi X.},
           giveni={P\bibinitperiod\bibinitdelim X\bibinitperiod}}}%
      }
      \strng{namehash}{e38330a7a234159955b75bdfae8cd3cd}
      \strng{fullhash}{e38330a7a234159955b75bdfae8cd3cd}
      \strng{bibnamehash}{e38330a7a234159955b75bdfae8cd3cd}
      \strng{authorbibnamehash}{e38330a7a234159955b75bdfae8cd3cd}
      \strng{authornamehash}{e38330a7a234159955b75bdfae8cd3cd}
      \strng{authorfullhash}{e38330a7a234159955b75bdfae8cd3cd}
      \field{sortinit}{G}
      \field{sortinithash}{62eb2aa29549e4fdbd3cb154ec5711cb}
      \field{labelprefix}{A}
      \field{labelnamesource}{author}
      \field{labeltitlesource}{title}
      \field{annotation}{Under RH, the author proves that $E^\psi(x) \ll (\log \log x)^2$ except on a set of finite measure, and that for any function $f(x)$ tending to infinity, $E^\psi(x) \ll f(x)$ except on a set of density~$0$; the proofs are similar to showing that $E^\psi(x)$ has a limiting logarithmic distribution. The method also provides short proofs of Cramér's conditional estimates~\cite{MR1544509} $\int_1^X E^\psi(x)^2\, dx \ll X$ and $\int_1^X E^\psi(x)^2 \frac{dx}{x} \sim C \log X$ (for an explicit constant~$C$), as well as of Selberg's conditional result on the normal density of primes in short intervals.\par This article cites~\cites{MR1544509}.}
      \field{issn}{0065-1036}
      \field{journaltitle}{Acta Arith.}
      \field{title}{Some consequences of the {R}iemann hypothesis}
      \field{volume}{37}
      \field{year}{1980}
      \true{nocite}
      \field{pages}{339\bibrangedash 343}
      \range{pages}{5}
      \verb{doi}
      \verb 10.4064/aa-37-1-339-343
      \endverb
      \verb{mrnumber}
      \verb 598886
      \endverb
      \verb{urlraw}
      \verb https://doi.org/10.4064/aa-37-1-339-343
      \endverb
      \verb{url}
      \verb https://doi.org/10.4064/aa-37-1-339-343
      \endverb
    \endentry
    \entry{MR240062}{article}{}
      \name{author}{2}{}{%
        {{hash=b86a87765e61c9d06430e684291fa9f9}{%
           family={Good},
           familyi={G\bibinitperiod},
           given={I.\bibnamedelimi J.},
           giveni={I\bibinitperiod\bibinitdelim J\bibinitperiod}}}%
        {{hash=6e79ce655dce4f2195319beea947f637}{%
           family={Churchhouse},
           familyi={C\bibinitperiod},
           given={R.\bibnamedelimi F.},
           giveni={R\bibinitperiod\bibinitdelim F\bibinitperiod}}}%
      }
      \strng{namehash}{8abf904253f009c4ff49d49aab43c597}
      \strng{fullhash}{8abf904253f009c4ff49d49aab43c597}
      \strng{bibnamehash}{8abf904253f009c4ff49d49aab43c597}
      \strng{authorbibnamehash}{8abf904253f009c4ff49d49aab43c597}
      \strng{authornamehash}{8abf904253f009c4ff49d49aab43c597}
      \strng{authorfullhash}{8abf904253f009c4ff49d49aab43c597}
      \field{sortinit}{G}
      \field{sortinithash}{62eb2aa29549e4fdbd3cb154ec5711cb}
      \field{labelprefix}{A}
      \field{labelnamesource}{author}
      \field{labeltitlesource}{title}
      \field{annotation}{The authors describe a random model for partial sums of $\mu(n)$ over short intervals and conjecture that when $h$ is large, the limiting distribution of $M(x)-M(x-h)$ is normal with mean $0$ and variance $6h/\pi^2$. Based on the law of the iterated logarithm, they also conjecture that $\lim \sup_{x \to \infty} {M(x)}/{\sqrt{x \log \log x}}=\sqrt{12}/\pi$.\par This article cites~\cites{MR0104638,1932.Ingham,1914.Littlewood,MR155787,MR0046485}.}
      \field{issn}{0025-5718}
      \field{journaltitle}{Math. Comp.}
      \field{title}{The {R}iemann hypothesis and pseudorandom features of the {M}öbius sequence}
      \field{volume}{22}
      \field{year}{1968}
      \true{nocite}
      \field{pages}{857\bibrangedash 861}
      \range{pages}{5}
      \verb{doi}
      \verb 10.2307/2004584
      \endverb
      \verb{mrnumber}
      \verb 240062
      \endverb
      \verb{urlraw}
      \verb https://doi.org/10.2307/2004584
      \endverb
      \verb{url}
      \verb https://doi.org/10.2307/2004584
      \endverb
    \endentry
    \entry{MR2202918}{article}{}
      \name{author}{2}{}{%
        {{hash=7c95bde388d46e453494c7b2fbd1b088}{%
           family={Granville},
           familyi={G\bibinitperiod},
           given={A.},
           giveni={A\bibinitperiod}}}%
        {{hash=feda7d3cba2a0ad4e772307fca655937}{%
           family={Martin},
           familyi={M\bibinitperiod},
           given={G.},
           giveni={G\bibinitperiod}}}%
      }
      \strng{namehash}{c4ece346b079ec803e063600350bbf92}
      \strng{fullhash}{c4ece346b079ec803e063600350bbf92}
      \strng{bibnamehash}{c4ece346b079ec803e063600350bbf92}
      \strng{authorbibnamehash}{c4ece346b079ec803e063600350bbf92}
      \strng{authornamehash}{c4ece346b079ec803e063600350bbf92}
      \strng{authorfullhash}{c4ece346b079ec803e063600350bbf92}
      \field{sortinit}{G}
      \field{sortinithash}{62eb2aa29549e4fdbd3cb154ec5711cb}
      \field{labelprefix}{A}
      \field{labelnamesource}{author}
      \field{labeltitlesource}{title}
      \field{annotation}{The authors present an accessible survey of prime number races, explicit formulas for $\pi(x)$ and $\pi(x;q,a)$, the biases caused by squares of primes, and limiting distributions and densities. The final section describes unpublished research of G.~Davidoff (in connection with an REU group), including the theorem that $W_{q;\cN,\cR}(T)$ is unbounded for any prime modulus~$q$, assuming only that $L(s,\chi_{\pm q})$ has no real zero in the interval $[\Theta(\chi_{\pm q}),1]$.\par This article cites~\cites{1914.Littlewood,MR0004255,MR0108470,MR0146156,MR0146157,MR0146158,MR0146159,MR0146160,MR0146161,MR0156826,MR0156827,MR0289452,MR0476616,MR552467,MR1329368,MR1344249,MR1367576,2007.Sarnak,2000.Ng,MR1752093,MR1816036,MR1806291,MR1956261,MR1909220, MR2075622,MR0882550}.}
      \field{issn}{0002-9890}
      \field{journaltitle}{Amer. Math. Monthly}
      \field{number}{1}
      \field{title}{Prime number races}
      \field{volume}{113}
      \field{year}{2006}
      \true{nocite}
      \field{pages}{1\bibrangedash 33}
      \range{pages}{33}
      \verb{doi}
      \verb 10.2307/27641834
      \endverb
      \verb{mrnumber}
      \verb 2202918
      \endverb
      \verb{urlraw}
      \verb https://doi.org/10.2307/27641834
      \endverb
      \verb{url}
      \verb https://doi.org/10.2307/27641834
      \endverb
    \endentry
    \entry{MR0198145}{article}{}
      \name{author}{1}{}{%
        {{hash=f2bf37798cd76913bab45e98f2a4d625}{%
           family={Grosswald},
           familyi={G\bibinitperiod},
           given={Emil},
           giveni={E\bibinitperiod}}}%
      }
      \strng{namehash}{f2bf37798cd76913bab45e98f2a4d625}
      \strng{fullhash}{f2bf37798cd76913bab45e98f2a4d625}
      \strng{bibnamehash}{f2bf37798cd76913bab45e98f2a4d625}
      \strng{authorbibnamehash}{f2bf37798cd76913bab45e98f2a4d625}
      \strng{authornamehash}{f2bf37798cd76913bab45e98f2a4d625}
      \strng{authorfullhash}{f2bf37798cd76913bab45e98f2a4d625}
      \field{extraname}{1}
      \field{sortinit}{G}
      \field{sortinithash}{62eb2aa29549e4fdbd3cb154ec5711cb}
      \field{labelprefix}{A}
      \field{labelnamesource}{author}
      \field{labeltitlesource}{title}
      \field{annotation}{The author generalizes theorems of Landau~\cite{MR1511360} and Pólya~\cite{1930.Polya} on the relationship between sign changes of a real-valued function and singularities of the corresponding Mellin transform; the main results allow for logarithmic-type singularities as well as poles. The article asserts that its Theorem~6 follows from its Theorem~2, but it is remarked in~\cite{MR858395} that this implication is invalid.\par This article cites~\cites{1930.Polya,1914.Littlewood, MR1511360, MR1511206}.}
      \field{issn}{0021-2172}
      \field{journaltitle}{Israel J. Math.}
      \field{title}{On some generalizations of theorems by Landau and Pólya}
      \field{volume}{3}
      \field{year}{1965}
      \true{nocite}
      \field{pages}{211\bibrangedash 220}
      \range{pages}{10}
      \verb{doi}
      \verb 10.1007/BF03008399
      \endverb
      \verb{mrnumber}
      \verb 0198145
      \endverb
    \endentry
    \entry{MR0202685}{article}{}
      \name{author}{1}{}{%
        {{hash=f2bf37798cd76913bab45e98f2a4d625}{%
           family={Grosswald},
           familyi={G\bibinitperiod},
           given={Emil},
           giveni={E\bibinitperiod}}}%
      }
      \strng{namehash}{f2bf37798cd76913bab45e98f2a4d625}
      \strng{fullhash}{f2bf37798cd76913bab45e98f2a4d625}
      \strng{bibnamehash}{f2bf37798cd76913bab45e98f2a4d625}
      \strng{authorbibnamehash}{f2bf37798cd76913bab45e98f2a4d625}
      \strng{authornamehash}{f2bf37798cd76913bab45e98f2a4d625}
      \strng{authorfullhash}{f2bf37798cd76913bab45e98f2a4d625}
      \field{extraname}{2}
      \field{sortinit}{G}
      \field{sortinithash}{62eb2aa29549e4fdbd3cb154ec5711cb}
      \field{labelprefix}{A}
      \field{labelnamesource}{author}
      \field{labeltitlesource}{title}
      \field{annotation}{In this article, the author reproduces known oscillation theorems for $\Delta^\psi(x)$, $\Delta^\Pi(x)$, and $\Delta^\pi(x)$ with proofs based on extended versions of Landau's theorem~\cite{MR1511360}. The results on the number of sign changes of $\psi(q;a,b)$ can be compared with ones in~\cite{MR0146156, MR0146157, MR0146158, MR0146159, MR0146161}. Moreover, under a suitable assumption weaker than HC on the triples $(q,a,b)$ the author proves that the difference $\pi(x;q,a,b)$ changes sign infinitely often for $q=43, 47, 163$, noting that the cases when $q=4$ was proved by Hardy and Littlewood~\cite{MR1555148} and the cases when $q=3, 5, 8$ were proved by Knapowski and Turán~\cite{MR0146159}. It has been noted~\cite[page 2]{MR858395} that some of the implications used in this article to derive some of the theorems from the others may not be valid.\par This article cites~\cites{MR0198145,MR0104638,MR0146156,MR0146157,MR0146158,MR0146159,MR0146160,MR0146161,MR0156826,MR0156827,MR0162771,MR0177963,MR1511206,MR0060548,1936.Ingham,MR1511360,1914.Littlewood,1930.Polya,1853.Chebyshev}.}
      \field{issn}{0002-9947}
      \field{journaltitle}{Trans. Amer. Math. Soc.}
      \field{title}{Oscillation theorems of arithmetical functions}
      \field{volume}{126}
      \field{year}{1967}
      \true{nocite}
      \field{pages}{1\bibrangedash 28}
      \range{pages}{28}
      \verb{doi}
      \verb 10.2307/1994409
      \endverb
      \verb{mrnumber}
      \verb 0202685
      \endverb
    \endentry
    \entry{MR0332685}{incollection}{}
      \name{author}{1}{}{%
        {{hash=f2bf37798cd76913bab45e98f2a4d625}{%
           family={Grosswald},
           familyi={G\bibinitperiod},
           given={Emil},
           giveni={E\bibinitperiod}}}%
      }
      \list{publisher}{1}{%
        {Springer, Berlin}%
      }
      \strng{namehash}{f2bf37798cd76913bab45e98f2a4d625}
      \strng{fullhash}{f2bf37798cd76913bab45e98f2a4d625}
      \strng{bibnamehash}{f2bf37798cd76913bab45e98f2a4d625}
      \strng{authorbibnamehash}{f2bf37798cd76913bab45e98f2a4d625}
      \strng{authornamehash}{f2bf37798cd76913bab45e98f2a4d625}
      \strng{authorfullhash}{f2bf37798cd76913bab45e98f2a4d625}
      \field{extraname}{3}
      \field{sortinit}{G}
      \field{sortinithash}{62eb2aa29549e4fdbd3cb154ec5711cb}
      \field{labelprefix}{A}
      \field{labelnamesource}{author}
      \field{labeltitlesource}{title}
      \field{annotation}{This article improves results due to Grosswald~\cites{MR0202685} that relied on an extended Landau's theorem and a theorem of Diamond~\cites{MR0332684}. Say that a finite set of real numbers is $k$-independent if there are no nontrivial vanishing linear combinations of that set with integer coefficients bounded by~$k$ in absolute value. The author shows that the verification of the $5$-independence of the (ordinates of the) first $30$ zeros of $\zeta$ is sufficient to show that $\pi(x)-\li(x)$ changes sign infinitely often. Similarly, the verification of the $13$-independence of the first $75$ zeros is sufficient to disprove the Mertens conjecture, and the verification of the $16$-independence of the first $13$ zeros is sufficient to prove that $L(x)$ and $L_r(x)$ change sign infinitely often. Moreover, for integers $q,q'$ with $\phi(q)=\phi(q')$, the author investigates the sign changes of $\pi(x;q,a)-\pi(x;q',a')$. In particular, the author shows that $\pi(x;p,a,a')$ changes sign infinitely often for all primes $p\le19$ and all appropriate $a,a'$.\par This article cites~\cites{MR0330069,MR0332684,MR0202685,MR1555148,MR0104638,MR0006202,MR0146156,MR0289452,1914.Littlewood,MR280447}.}
      \field{bookpagination}{none}
      \field{booktitle}{The theory of arithmetic functions ({P}roc. {C}onf., {W}estern {M}ichigan {U}niv., {K}alamazoo, {M}ich., 1971)}
      \field{note}{Lecture Notes in Math., Vol. 251}
      \field{title}{Oscillation theorems}
      \field{year}{1972}
      \true{nocite}
      \field{pages}{141\bibrangedash 168}
      \range{pages}{28}
      \verb{mrnumber}
      \verb 0332685
      \endverb
    \endentry
    \entry{MR0003644}{article}{}
      \name{author}{1}{}{%
        {{hash=4f3424aa2d374362fe10ed8b28fd159f}{%
           family={Gupta},
           familyi={G\bibinitperiod},
           given={Hansraj},
           giveni={H\bibinitperiod}}}%
      }
      \strng{namehash}{4f3424aa2d374362fe10ed8b28fd159f}
      \strng{fullhash}{4f3424aa2d374362fe10ed8b28fd159f}
      \strng{bibnamehash}{4f3424aa2d374362fe10ed8b28fd159f}
      \strng{authorbibnamehash}{4f3424aa2d374362fe10ed8b28fd159f}
      \strng{authornamehash}{4f3424aa2d374362fe10ed8b28fd159f}
      \strng{authorfullhash}{4f3424aa2d374362fe10ed8b28fd159f}
      \field{sortinit}{G}
      \field{sortinithash}{62eb2aa29549e4fdbd3cb154ec5711cb}
      \field{labelprefix}{A}
      \field{labelnamesource}{author}
      \field{labeltitlesource}{title}
      \field{annotation}{Regarding Pólya's problem, the author verifies that $L(n) \leq 0$ for $2 \leq n \leq 20,000$. The author also conjectures that $L(x)\ll\sqrt{x}$ based on the computation.\par This article cites~\cite{1919.Polya}.}
      \field{journaltitle}{Proc. Indian Acad. Sci., Sect. A.}
      \field{title}{On a table of values of {$L(n)$}}
      \field{volume}{12}
      \field{year}{1940}
      \true{nocite}
      \field{pages}{407\bibrangedash 409}
      \range{pages}{3}
      \verb{mrnumber}
      \verb 0003644
      \endverb
    \endentry
    \entry{MR1576550}{article}{}
      \name{author}{1}{}{%
        {{hash=9d761c354d534c74e69f17afcc78fea8}{%
           family={Hardy},
           familyi={H\bibinitperiod},
           given={G.\bibnamedelimi H.},
           giveni={G\bibinitperiod\bibinitdelim H\bibinitperiod}}}%
      }
      \strng{namehash}{9d761c354d534c74e69f17afcc78fea8}
      \strng{fullhash}{9d761c354d534c74e69f17afcc78fea8}
      \strng{bibnamehash}{9d761c354d534c74e69f17afcc78fea8}
      \strng{authorbibnamehash}{9d761c354d534c74e69f17afcc78fea8}
      \strng{authornamehash}{9d761c354d534c74e69f17afcc78fea8}
      \strng{authorfullhash}{9d761c354d534c74e69f17afcc78fea8}
      \field{sortinit}{H}
      \field{sortinithash}{6db6145dae8dc9e1271a8d556090b50a}
      \field{labelprefix}{A}
      \field{labelnamesource}{author}
      \field{labeltitlesource}{title}
      \field{annotation}{This article begins with a short background on the history of the Dirichlet divisor problem and states the best bounds for $\Delta^D(x)$ available at the time, namely Voronoi's result $\Delta^D(x) \ll x^{1/3}\log x$. The author then proves that $\Delta^D(x) = \Omega_\pm(x^{1/4})$. He draws a direct comparison to Schmidt's proof~\cite{MR1511206} that $\Delta^\Pi(x) = \Omega_\pm(\sqrt x / \log x)$, and to a simplified version of Schmidt's proof by Landau~\cite{1909.Landau}; he states that his proof for $\Delta^D$ does not ``differ in principle'' from Landau's proof for $\Delta^\Pi$, and for the sake of comparison provides a simplified form of the proof that $\Delta^\psi(x) = \Omega_\pm(\sqrt x)$ based on Landau's methods. \par The article then discusses various generalizations of the problem. The first generalization discussed, due to Piltz, concerns the same problem for $D_k(x) = \sum_{n \le x} \tau_k(n)$, where $\tau_k(n)$ denotes the number of ways~$n$ can be decomposed into~$k$ ordered factors: Hardy asserts that the methods in this article can show that $\Delta^{D_k}(x) = \Omega_\pm(x^{(k-1)/(2k)})$. The second generalization concerns $R(x) = \sum_{n \le x} r(n)$, where $r(n)$ counts the number of ways~$n$ can be represented as the sum of two squares: Hardy likewise asserts that $\Delta^R(x) = \Omega_\pm(x^{1/4})$. Finally, the article concludes with a discussion of an explicit formula for $D(x)$, first found by Voronoi, and some alternative proof methods and their comparative advantages and disadvantages. \par This article cites~\cite{1914.Littlewood, 1909.Landau}.}
      \field{issn}{0024-6115}
      \field{journaltitle}{Proc. London Math. Soc. (2)}
      \field{title}{On {D}irichlet's {D}ivisor {P}roblem}
      \field{volume}{15}
      \field{year}{1916}
      \true{nocite}
      \field{pages}{1\bibrangedash 25}
      \range{pages}{25}
      \verb{doi}
      \verb 10.1112/plms/s2-15.1.1
      \endverb
      \verb{mrnumber}
      \verb 1576550
      \endverb
      \verb{urlraw}
      \verb https://doi-org.ezproxy.library.ubc.ca/10.1112/plms/s2-15.1.1
      \endverb
      \verb{url}
      \verb https://doi-org.ezproxy.library.ubc.ca/10.1112/plms/s2-15.1.1
      \endverb
    \endentry
    \entry{1915.HardyLittlewood}{article}{}
      \name{author}{2}{}{%
        {{hash=9d761c354d534c74e69f17afcc78fea8}{%
           family={Hardy},
           familyi={H\bibinitperiod},
           given={G.\bibnamedelimi H.},
           giveni={G\bibinitperiod\bibinitdelim H\bibinitperiod}}}%
        {{hash=56dd48151c55dce1b8a073ecf909a238}{%
           family={Littlewood},
           familyi={L\bibinitperiod},
           given={J.\bibnamedelimi E.},
           giveni={J\bibinitperiod\bibinitdelim E\bibinitperiod}}}%
      }
      \strng{namehash}{ce5cb2e2cdf2e499e451b481213bdb97}
      \strng{fullhash}{ce5cb2e2cdf2e499e451b481213bdb97}
      \strng{bibnamehash}{ce5cb2e2cdf2e499e451b481213bdb97}
      \strng{authorbibnamehash}{ce5cb2e2cdf2e499e451b481213bdb97}
      \strng{authornamehash}{ce5cb2e2cdf2e499e451b481213bdb97}
      \strng{authorfullhash}{ce5cb2e2cdf2e499e451b481213bdb97}
      \field{extraname}{1}
      \field{sortinit}{H}
      \field{sortinithash}{6db6145dae8dc9e1271a8d556090b50a}
      \field{labelprefix}{A}
      \field{labelnamesource}{author}
      \field{labeltitlesource}{title}
      \field{annotation}{A letter of Tchebychef asserts that $\lim_{x\to\infty} \pi_e(x) = -\infty$. The authors verify the conjecture assuming GRH for $L(s,\chi_{-4})$. Using a formula of Cahen and Mellin and an application of Cauchy's theorem, the authors show that $-\pi_e(x)\gg \sqrt{x/\log x}$.}
      \field{journaltitle}{Proc. London Math. Soc. (2)}
      \field{title}{On an assertion of {T}chebychef}
      \field{volume}{14}
      \field{year}{1915}
      \true{nocite}
      \field{pages}{xv\bibrangedash xvi}
      \range{pages}{2}
    \endentry
    \entry{MR1555148}{article}{}
      \name{author}{2}{}{%
        {{hash=9d761c354d534c74e69f17afcc78fea8}{%
           family={Hardy},
           familyi={H\bibinitperiod},
           given={G.\bibnamedelimi H.},
           giveni={G\bibinitperiod\bibinitdelim H\bibinitperiod}}}%
        {{hash=56dd48151c55dce1b8a073ecf909a238}{%
           family={Littlewood},
           familyi={L\bibinitperiod},
           given={J.\bibnamedelimi E.},
           giveni={J\bibinitperiod\bibinitdelim E\bibinitperiod}}}%
      }
      \strng{namehash}{ce5cb2e2cdf2e499e451b481213bdb97}
      \strng{fullhash}{ce5cb2e2cdf2e499e451b481213bdb97}
      \strng{bibnamehash}{ce5cb2e2cdf2e499e451b481213bdb97}
      \strng{authorbibnamehash}{ce5cb2e2cdf2e499e451b481213bdb97}
      \strng{authornamehash}{ce5cb2e2cdf2e499e451b481213bdb97}
      \strng{authorfullhash}{ce5cb2e2cdf2e499e451b481213bdb97}
      \field{extraname}{2}
      \field{sortinit}{H}
      \field{sortinithash}{6db6145dae8dc9e1271a8d556090b50a}
      \field{labelprefix}{A}
      \field{labelnamesource}{author}
      \field{labeltitlesource}{title}
      \field{annotation}{This article contains full proofs of several results that had been announced by (at least one of) the authors in the few years prior.\par In Section 2.2, the authors obtain an explicit formula for the exponentially weighted sum $\psi_e(x) - 1/(e^{1/x}-1) = \sum_{n=1}^\infty (\Lambda(n)-1) e^{-n/x}$; furthermore, assuming RH, they show that this expression is both $\ll \sqrt x$ and $\Omega_\pm(\sqrt x)$. From the latter they deduce that $\Delta^\psi(x) = \Omega_\pm(\sqrt x)$. In Section 2.3, they consider the function $\sum_{p\ge3} (-1)^{(p+1)/2} e^{-p/x} = -\pi_e(x,\chi_{-4})$. Assuming GRH for $L(s,\chi_{-4})$, they prove that $\pi_e(x,\chi_{-4})\to-\infty$ as $x\to\infty$, which is one way of justifying Chebyshev's observation that there are more primes congruent to $3\mod 4$ than to $1\mod 4$.\par In Section 5, the authors provide a full proof of ``Littlewood's theorem'' (announced in~\cite{1914.Littlewood}) on irregularities in the distribution of primes: they prove that \[ \Delta^\pi(x) = \Omega_\pm\bigg( \frac{\sqrt x}{\log x} \log\log\log x \bigg), \] which in particular refutes the conjecture that $\pi(x)<\li(x)$ for all $x>1$.\par Their proof, which begins with the assumption of RH thanks to prior work of Landau~\cite[Sections 201--3]{1909.Landau}, uses homogeneous Diophantine approximation for the imaginary parts of the zeros of $\zeta(s)$. They assert that it can be shown in a similar way that $\psi(x,\chi_{-4}) = \Omega_\pm(\sqrt x\log\log\log x)$ and \[ \pi(x;4,3,1) = \Omega_\pm\bigg( \frac{\sqrt x}{\log x} \log\log\log x \bigg); \] these results are actually dissonant with Chebyshev's observations. (As a side remark, this is also the article in which appears the proof of the asymptotic formula for the second moment of $\zeta(s)$ on the critical line.)\par This article cites~\cites{MR1511206,1909.Landau,1914.Littlewood,1915.HardyLittlewood}.}
      \field{journaltitle}{Acta Math.}
      \field{number}{1}
      \field{title}{{Contributions to the theory of the Riemann zeta-function and the theory of the distribution of primes}}
      \field{volume}{41}
      \field{year}{1916}
      \true{nocite}
      \field{pages}{119\bibrangedash 196}
      \range{pages}{78}
      \verb{mrnumber}
      \verb 1555148
      \endverb
    \endentry
    \entry{MR3773805}{article}{}
      \name{author}{2}{}{%
        {{hash=d6e30aa748c8fb70505ca61ee684d871}{%
           family={Harper},
           familyi={H\bibinitperiod},
           given={Adam\bibnamedelima J.},
           giveni={A\bibinitperiod\bibinitdelim J\bibinitperiod}}}%
        {{hash=64363cb2afef9e33f967130adc38684b}{%
           family={Lamzouri},
           familyi={L\bibinitperiod},
           given={Youness},
           giveni={Y\bibinitperiod}}}%
      }
      \strng{namehash}{d2f0ae4433a2a9de0fd6046d871e9ead}
      \strng{fullhash}{d2f0ae4433a2a9de0fd6046d871e9ead}
      \strng{bibnamehash}{d2f0ae4433a2a9de0fd6046d871e9ead}
      \strng{authorbibnamehash}{d2f0ae4433a2a9de0fd6046d871e9ead}
      \strng{authornamehash}{d2f0ae4433a2a9de0fd6046d871e9ead}
      \strng{authorfullhash}{d2f0ae4433a2a9de0fd6046d871e9ead}
      \field{sortinit}{H}
      \field{sortinithash}{6db6145dae8dc9e1271a8d556090b50a}
      \field{labelprefix}{A}
      \field{labelnamesource}{author}
      \field{labeltitlesource}{title}
      \field{annotation}{The authors investigate, assuming GRH and LI, the asymptotic behavior of $\delta_{q;a_1,\ldots,a_n}$ when the number of competitors~$n$ grows as a function of the modulus~$q$. They show that if $n\le \log q/(\log\log q)^4$ then $\delta_{q;a_1,\ldots,a_n} \sim \frac{1}{n!}$, resolving an unpublished conjecture by Ford and Lamzouri and strengthening results of Rubinstein and Sarnak~\cite{MR1329368} and Lamzouri~\cite{MR3063909}. They prove that this is not necessarily true for larger~$n$, as predicted by Feuerverger and Martin~\cite{MR1806291}: when $\phi(q)^{\ep} \leq n \leq \phi(q)$ one has $\liminf_{q\to\infty} n!\delta_{q;a_1,\ldots,a_n} < 1$.\par They further discuss the first~$k$ leaders in a prime number race: for each integer $1\leq k\leq n$, they define $\delta_k(q;a_1, \ldots, a_n)$ to be the logarithmic density of the set of real numbers $x\geq 2$ such that $$ \pi(x;q, a_1)>\pi(x;q, a_2) > \cdots> \pi(x;q, a_k)> \max_{k+1\leq j\leq n}\pi(x;q, a_j). $$ They show that if $2 \leq n \leq \phi(q)^{1 / 32}$ then $\delta_1(q;a_1, \ldots, a_n) \sim \frac1n$, and that the analogous result $\delta_k(q;a_1, \ldots, a_n) \sim \frac{(n-k)!}{n!}$ holds for $k(\log k)^{10}\le(\log q)/\log n$ but not for $n\ge \phi(q)^\ep$.\par In addition to using the circle method to control the average size of correlations in prime number races, the authors develop sophisticated probabilistic tools including an exchangeable pairs version of Stein's method and variants of ``normal comparison'' lemmas of Slepian.\par This article cites~\cites{MR1806291, MR3272280, MR3028758, MR1909220, MR3151605, MR2202918, MR1251926, MR1256149, MR1367576, MR0146156, MR0146157, MR0146158, MR0146159, MR0146160, MR0156826, MR0156827, MR2998146, MR3063909, 1914.Littlewood, MR4082248, MR1297543, MR1329368}, as well as an earlier draft of this annotated bibliography.}
      \field{issn}{0178-8051}
      \field{journaltitle}{Probab. Theory Related Fields}
      \field{number}{3-4}
      \field{title}{Orderings of weakly correlated random variables, and prime number races with many contestants}
      \field{volume}{170}
      \field{year}{2018}
      \true{nocite}
      \field{pages}{961\bibrangedash 1010}
      \range{pages}{50}
      \verb{doi}
      \verb 10.1007/s00440-017-0800-2
      \endverb
      \verb{mrnumber}
      \verb 3773805
      \endverb
      \verb{urlraw}
      \verb https://doi.org/10.1007/s00440-017-0800-2
      \endverb
      \verb{url}
      \verb https://doi.org/10.1007/s00440-017-0800-2
      \endverb
    \endentry
    \entry{MR0104638}{article}{}
      \name{author}{1}{}{%
        {{hash=4cabd3deed7eb4da21672a1718558bbe}{%
           family={Haselgrove},
           familyi={H\bibinitperiod},
           given={C.\bibnamedelimi B.},
           giveni={C\bibinitperiod\bibinitdelim B\bibinitperiod}}}%
      }
      \strng{namehash}{4cabd3deed7eb4da21672a1718558bbe}
      \strng{fullhash}{4cabd3deed7eb4da21672a1718558bbe}
      \strng{bibnamehash}{4cabd3deed7eb4da21672a1718558bbe}
      \strng{authorbibnamehash}{4cabd3deed7eb4da21672a1718558bbe}
      \strng{authornamehash}{4cabd3deed7eb4da21672a1718558bbe}
      \strng{authorfullhash}{4cabd3deed7eb4da21672a1718558bbe}
      \field{sortinit}{H}
      \field{sortinithash}{6db6145dae8dc9e1271a8d556090b50a}
      \field{labelprefix}{A}
      \field{labelnamesource}{author}
      \field{labeltitlesource}{title}
      \field{annotation}{Following Ingham's method \cite{MR0006202}, the author resolves Pólya's problem in the negative by showing that a truncated version of the explicit formula for $E^L(x)$, using zeros of $\zeta(s)$ up to height $1,000$, is positive at $x=e^{831.847}$; while this does not rigorously establish that $E^L(e^{831.847})$ is itself positive, the truncated value is a lower bound for $\limsup_{x\to\infty} E^L(x)$. Similarly, the author resolves Turán's problem in the negative by showing that the analogous truncation of the explicit formula for $E^{L_r}(x)$ is negative at $x=e^{853.853}$. The computations were done on an EDSAC~I and a Mark~I. \par This article cites~\cites{MR27305,MR0006202,1897.Mertens,1919.Polya}.}
      \field{issn}{0025-5793}
      \field{journaltitle}{Mathematika}
      \field{title}{A disproof of a conjecture of Pólya}
      \field{volume}{5}
      \field{year}{1958}
      \true{nocite}
      \field{pages}{141\bibrangedash 145}
      \range{pages}{5}
      \verb{doi}
      \verb 10.1112/S0025579300001480
      \endverb
      \verb{mrnumber}
      \verb 0104638
      \endverb
    \endentry
    \entry{MR1159354}{article}{}
      \name{author}{1}{}{%
        {{hash=5895b52bfd9795ff19a630ff28e5fe17}{%
           family={Heath-Brown},
           familyi={H\bibinithyphendelim B\bibinitperiod},
           given={D.\bibnamedelimi R.},
           giveni={D\bibinitperiod\bibinitdelim R\bibinitperiod}}}%
      }
      \strng{namehash}{5895b52bfd9795ff19a630ff28e5fe17}
      \strng{fullhash}{5895b52bfd9795ff19a630ff28e5fe17}
      \strng{bibnamehash}{5895b52bfd9795ff19a630ff28e5fe17}
      \strng{authorbibnamehash}{5895b52bfd9795ff19a630ff28e5fe17}
      \strng{authornamehash}{5895b52bfd9795ff19a630ff28e5fe17}
      \strng{authorfullhash}{5895b52bfd9795ff19a630ff28e5fe17}
      \field{sortinit}{H}
      \field{sortinithash}{6db6145dae8dc9e1271a8d556090b50a}
      \field{labelprefix}{A}
      \field{labelnamesource}{author}
      \field{labeltitlesource}{title}
      \field{annotation}{Let $a_1(x),a_2(x),\ldots$ be continuous real-valued functions of period $1$, and let $\gamma_1, \gamma_2,\ldots$ be nonzero constants such that $\lim_{N\to\infty}\limsup_{T\to\infty} \frac{1}{T}\int_0^T \min\{1,|F(x)-\sum_{n\leq N}a_n(\gamma_nx)|\} \,dx=0$; the author shows that the mean value $\frac1T \int_0^T p(F(t))\,dt$ converges to a limit as $T\to\infty$ for every continuous, integrable, piecewise differentiable function $p$ for which $\hat p$ is also integrable. Under further hypotheses, including the linear independence of the $\gamma_j$ over $\Q$, the author shows that $F(x)$ has a limiting distribution function. When such a limiting distribution function exists, the author establishes a necessary condition for certain normalized moments of $F(x)$ to converge in the limit. The author applies these theorems when $F(x)$ is the error term in: the classical divisor problem, the circle problem, the Piltz divisor problem for $\tau_3$, or the second moment for $\zeta(s)$ on the critical line. An important lemma is the evaluation, for continuous functions $b_1(t),\ldots,b_k(t)$ of period $1$ from $\R$ to $\C$, of $\lim_{T\to\infty} \frac{1}{T}\int_0^T b_1(\gamma_1 t)\ldots b_k(\gamma_k t)\,dt$.\par This article cites~\cites{MR1544509,MR0882550}.}
      \field{issn}{0065-1036}
      \field{journaltitle}{Acta Arith.}
      \field{number}{4}
      \field{title}{The distribution and moments of the error term in the {D}irichlet divisor problem}
      \field{volume}{60}
      \field{year}{1992}
      \true{nocite}
      \field{pages}{389\bibrangedash 415}
      \range{pages}{27}
      \verb{doi}
      \verb 10.4064/aa-60-4-389-415
      \endverb
      \verb{mrnumber}
      \verb 1159354
      \endverb
      \verb{urlraw}
      \verb https://doi.org/10.4064/aa-60-4-389-415
      \endverb
      \verb{url}
      \verb https://doi.org/10.4064/aa-60-4-389-415
      \endverb
    \endentry
    \entry{MR0506080}{article}{}
      \name{author}{1}{}{%
        {{hash=a8dc47e3fff14028c7ddc24efe554469}{%
           family={Hooley},
           familyi={H\bibinitperiod},
           given={C.},
           giveni={C\bibinitperiod}}}%
      }
      \strng{namehash}{a8dc47e3fff14028c7ddc24efe554469}
      \strng{fullhash}{a8dc47e3fff14028c7ddc24efe554469}
      \strng{bibnamehash}{a8dc47e3fff14028c7ddc24efe554469}
      \strng{authorbibnamehash}{a8dc47e3fff14028c7ddc24efe554469}
      \strng{authornamehash}{a8dc47e3fff14028c7ddc24efe554469}
      \strng{authorfullhash}{a8dc47e3fff14028c7ddc24efe554469}
      \field{sortinit}{H}
      \field{sortinithash}{6db6145dae8dc9e1271a8d556090b50a}
      \field{labelprefix}{A}
      \field{labelnamesource}{author}
      \field{labeltitlesource}{title}
      \field{annotation}{Assuming GRH and LI, the author shows that $E^\psi(x;q,a)/\sqrt{\phi(q)\log q}$ has a limiting logarithmic distribution function (depending only on~$q$ and not~$a$), as well as the central limit theorem that these functions tend to the standard normal distribution as $q\to\infty$. The proof proceeds by writing the characteristic function of this distribution as a product of Bessel functions. The same result holds for $E^\theta(x;q,a)/\sqrt{\phi(q)\log q}$, except that the mean of the limiting logarithmic distribution depends on whether~$a$ is a quadratic residue\mod q.\par This article cites~\cites{MR0004255}.}
      \field{issn}{0024-6107}
      \field{journaltitle}{J. London Math. Soc. (2)}
      \field{number}{1}
      \field{title}{On the {B}arban-{D}avenport-{H}alberstam theorem. {VII}}
      \field{volume}{16}
      \field{year}{1977}
      \true{nocite}
      \field{pages}{1\bibrangedash 8}
      \range{pages}{8}
      \verb{doi}
      \verb 10.1112/jlms/s2-16.1.1
      \endverb
      \verb{mrnumber}
      \verb 0506080
      \endverb
    \endentry
    \entry{MR3692499}{article}{}
      \name{author}{1}{}{%
        {{hash=16aef6e7e9117011b23f5df0b0f8e46b}{%
           family={Hough},
           familyi={H\bibinitperiod},
           given={Patrick},
           giveni={P\bibinitperiod}}}%
      }
      \strng{namehash}{16aef6e7e9117011b23f5df0b0f8e46b}
      \strng{fullhash}{16aef6e7e9117011b23f5df0b0f8e46b}
      \strng{bibnamehash}{16aef6e7e9117011b23f5df0b0f8e46b}
      \strng{authorbibnamehash}{16aef6e7e9117011b23f5df0b0f8e46b}
      \strng{authornamehash}{16aef6e7e9117011b23f5df0b0f8e46b}
      \strng{authorfullhash}{16aef6e7e9117011b23f5df0b0f8e46b}
      \field{sortinit}{H}
      \field{sortinithash}{6db6145dae8dc9e1271a8d556090b50a}
      \field{labelprefix}{A}
      \field{labelnamesource}{author}
      \field{labeltitlesource}{title}
      \field{annotation}{The author establishes a stronger version of a conjecture from~\cite{MR3521338} on the biases between products of two primes, by showing that sufficiently large $x$, there are at least $\exp\big((C-\ep)\sqrt{\log x}\big)$ integers $d \leq \exp(C\sqrt{\log x})$ such that \[ \frac{\#\{pq \leq x \colon \chi_d(p) = \chi_d(q) = -1 \} } {\frac{1}{4}\# \{pq \leq x \colon (pq,d) = 1\} } \geq 1 + \frac{\log \log \log x + O(1)}{\log \log x}, \] and similarly with ``${}\ge1+\cdots$'' replaced by ``${}\le1-\cdots$'' and/or with $-1$ replaced by $1$ on the left-hand side. The author also shows that GRH implies that these oscillations are best possible. The proof uses a result of Granville and Soundararajan on extremal values of $L(1,\chi)$.\par This article cites~\cite{MR3028758,MR2202918,MR3521338}.}
      \field{issn}{2363-9555}
      \field{journaltitle}{Res. Number Theory}
      \field{title}{A lower bound for biases amongst products of two primes}
      \field{volume}{3}
      \field{year}{2017}
      \true{nocite}
      \field{pages}{Art. 19\bibrangessep 11}
      \range{pages}{-1}
      \verb{doi}
      \verb 10.1007/s40993-017-0083-9
      \endverb
      \verb{mrnumber}
      \verb 3692499
      \endverb
    \endentry
    \entry{MR552467}{article}{}
      \name{author}{1}{}{%
        {{hash=9e0609c9931c1b620df97e8c58d07532}{%
           family={Hudson},
           familyi={H\bibinitperiod},
           given={Richard\bibnamedelima H.},
           giveni={R\bibinitperiod\bibinitdelim H\bibinitperiod}}}%
      }
      \strng{namehash}{9e0609c9931c1b620df97e8c58d07532}
      \strng{fullhash}{9e0609c9931c1b620df97e8c58d07532}
      \strng{bibnamehash}{9e0609c9931c1b620df97e8c58d07532}
      \strng{authorbibnamehash}{9e0609c9931c1b620df97e8c58d07532}
      \strng{authornamehash}{9e0609c9931c1b620df97e8c58d07532}
      \strng{authorfullhash}{9e0609c9931c1b620df97e8c58d07532}
      \field{extraname}{1}
      \field{sortinit}{H}
      \field{sortinithash}{6db6145dae8dc9e1271a8d556090b50a}
      \field{labelprefix}{A}
      \field{labelnamesource}{author}
      \field{labeltitlesource}{title}
      \field{annotation}{In this article, the author outlines a combinatorial principle that seeks to explain various effects and biases in comparative prime number theory. He highlights Riemann's original explicit formula $\pi(x)\sim \li(x)-\tfrac 12 \li (x^{1/2})-\tfrac 13 \li (x^{1/3})+\cdots$ and connects it to Chebyshev's observation, which can be seen as approximating $\pi(x;4,3,1)$ by half the number of prime squares. Arguing from a generalisation of an exact formula of Meissel, the author deduces, in the example of primes\mod 4, that an ``excess'' in the number of integers of the form $pq$, where $p$ and $q$ are prime, in the class $1\mod 4$ must result in a corresponding ``deficiency'' in the number of primes of exactly this magnitude, that is, half the number of prime squares. A combinatorial observation gives a reason for such an excess: in counting integers that are the product of two primes from a set, products of distinct primes are counted twice (as $pq$ and $qp$), while the prime squares are not. The author then provides similar arguments for why cubic and higher order effects should exist. Along with describing the combinatorial principle in generality, he provides details of some numerical investigations into these effects.\par This article cites~\cites{1853.Chebyshev, 1914.Littlewood, MR0108470, MR0162771, MR0174538, MR0289452}.}
      \field{issn}{0075-4102}
      \field{journaltitle}{J. Reine Angew. Math.}
      \field{title}{A common combinatorial principle underlies {R}iemann's formula, the {C}hebyshev phenomenon, and other subtle effects in comparative prime number theory. {I}}
      \field{volume}{313}
      \field{year}{1980}
      \true{nocite}
      \field{pages}{133\bibrangedash 150}
      \range{pages}{18}
      \verb{doi}
      \verb 10.1515/crll.1980.313.133
      \endverb
      \verb{mrnumber}
      \verb 552467
      \endverb
      \verb{urlraw}
      \verb https://doi.org/10.1515/crll.1980.313.133
      \endverb
      \verb{url}
      \verb https://doi.org/10.1515/crll.1980.313.133
      \endverb
    \endentry
    \entry{MR777286}{article}{}
      \name{author}{1}{}{%
        {{hash=9e0609c9931c1b620df97e8c58d07532}{%
           family={Hudson},
           familyi={H\bibinitperiod},
           given={Richard\bibnamedelima H.},
           giveni={R\bibinitperiod\bibinitdelim H\bibinitperiod}}}%
      }
      \strng{namehash}{9e0609c9931c1b620df97e8c58d07532}
      \strng{fullhash}{9e0609c9931c1b620df97e8c58d07532}
      \strng{bibnamehash}{9e0609c9931c1b620df97e8c58d07532}
      \strng{authorbibnamehash}{9e0609c9931c1b620df97e8c58d07532}
      \strng{authornamehash}{9e0609c9931c1b620df97e8c58d07532}
      \strng{authorfullhash}{9e0609c9931c1b620df97e8c58d07532}
      \field{extraname}{2}
      \field{sortinit}{H}
      \field{sortinithash}{6db6145dae8dc9e1271a8d556090b50a}
      \field{labelprefix}{A}
      \field{labelnamesource}{author}
      \field{labeltitlesource}{title}
      \field{annotation}{The author presents data on $E^\pi(x;6,5,1)$ for $x\le2.5\cdot10^{11}$. He proves (acknowledging A.~Schinzel's help) that $\lim_{x\to\infty} E^\pi(x;6,5,1) \ne 1$, and gives a heuristic argument that $\lim_{x\to\infty} E^{A^\pi_m}(x;6,5,1) \ne 1$ for every $m\in\N$. On the other hand, he shows that $\lim_{m\to\infty} \limsup_{x\to\infty} E^{A^\pi_m}(x;6,5,1) = 1$ and similarly for $\liminf_{x\to\infty}$.\par This article cites~\cites{MR0162771, MR0174538, MR0177963, MR0180539, MR0182616A, MR0182616B,MR0417077}.}
      \field{issn}{0025-5718}
      \field{journaltitle}{Math. Comp.}
      \field{number}{170}
      \field{title}{Averaging effects on irregularities in the distribution of primes in arithmetic progressions}
      \field{volume}{44}
      \field{year}{1985}
      \true{nocite}
      \field{pages}{561\bibrangedash 571}
      \range{pages}{11}
      \verb{doi}
      \verb 10.2307/2007974
      \endverb
      \verb{mrnumber}
      \verb 777286
      \endverb
    \endentry
    \entry{MR0460261}{article}{}
      \name{author}{2}{}{%
        {{hash=9e0609c9931c1b620df97e8c58d07532}{%
           family={Hudson},
           familyi={H\bibinitperiod},
           given={Richard\bibnamedelima H.},
           giveni={R\bibinitperiod\bibinitdelim H\bibinitperiod}}}%
        {{hash=dcdb9f9c423bbcbfa2a60d519ccd6915}{%
           family={Bays},
           familyi={B\bibinitperiod},
           given={Carter},
           giveni={C\bibinitperiod}}}%
      }
      \strng{namehash}{7660af3606957b9dca1512f72f50250e}
      \strng{fullhash}{7660af3606957b9dca1512f72f50250e}
      \strng{bibnamehash}{7660af3606957b9dca1512f72f50250e}
      \strng{authorbibnamehash}{7660af3606957b9dca1512f72f50250e}
      \strng{authornamehash}{7660af3606957b9dca1512f72f50250e}
      \strng{authorfullhash}{7660af3606957b9dca1512f72f50250e}
      \field{sortinit}{H}
      \field{sortinithash}{6db6145dae8dc9e1271a8d556090b50a}
      \field{labelprefix}{A}
      \field{labelnamesource}{author}
      \field{labeltitlesource}{title}
      \field{annotation}{The goal of this article is to find an elementary explanation (not involving complex analysis) for biases in prime number races. They argue from the Meissel-type formula \begin{align*} \pi(x;q,a) = \Phi(x,x^{1/3};q,a) +& \pi(x^{1/3};q,a) - \sum_{\substack{x^{1/3}<p\le x^{1/2} \\ p\nmid q}} \pi\bigg( \frac xp;q,ap^{-1} \bigg) \\ +& \sum_{\substack{x^{1/3}<p\le x^{1/2} \\ p\nmid q}} \pi(p;q,ap^{-1}) - \sum_{\substack{b\mod q \\ b^2\equiv a\mod q}} \pi\big( (x^{1/3},x^{1/2}] ; q,b\big), \end{align*} where $\Phi(x,y;q,a)$ denotes the number of integers in $(1,x]$ congruent to $a\mod q$ that are free of prime factors up to~$y$. They identify the final term as a persistent bias against quadratic residues~$a$, and support their interpretation with numerical evidence.\par The authors refine the above formula to reduce the occurrences of $x^{1/3}$ to $x^{1/4}$, thus introducing contributions from products of three primes; while observable in the numerical data, these contributions seem to be less significant for large~$x$. From their analysis they make some heuristic predictions, in particular that \[ \fA_1^\pi(x;q,\cN) - (c_q-1) \fA_1^\pi(x;q,\cR) \sim \frac 12 \sum_{n=1}^x \pi(n^{1/2}). \]\par This article cites~\cite{MR0447090,MR0476615,MR1555148,MR0146156,MR0146157,MR0146158,MR0146159,MR0146160,MR0146161,MR0156826,MR0156827,MR0162771,MR0174538,MR0180539,MR0182616A,MR0182616B,MR0200250,MR0177963,MR0068565,MR1544293A,MR0083001,1853.Chebyshev,1891.Phragmen,MR0108470,MR0302590,MR0289452,MR0184920}.}
      \field{issn}{0075-4102}
      \field{journaltitle}{J. Reine Angew. Math.}
      \field{title}{The mean behavior of primes in arithmetic progressions}
      \field{volume}{296}
      \field{year}{1977}
      \true{nocite}
      \field{pages}{80\bibrangedash 99}
      \range{pages}{20}
      \verb{doi}
      \verb 10.1515/crll.1977.296.80
      \endverb
      \verb{mrnumber}
      \verb 0460261
      \endverb
    \endentry
    \entry{MR2994374}{article}{}
      \name{author}{1}{}{%
        {{hash=48017d687b31d05c591036588e04dd66}{%
           family={Humphries},
           familyi={H\bibinitperiod},
           given={Peter},
           giveni={P\bibinitperiod}}}%
      }
      \strng{namehash}{48017d687b31d05c591036588e04dd66}
      \strng{fullhash}{48017d687b31d05c591036588e04dd66}
      \strng{bibnamehash}{48017d687b31d05c591036588e04dd66}
      \strng{authorbibnamehash}{48017d687b31d05c591036588e04dd66}
      \strng{authornamehash}{48017d687b31d05c591036588e04dd66}
      \strng{authorfullhash}{48017d687b31d05c591036588e04dd66}
      \field{extraname}{1}
      \field{sortinit}{H}
      \field{sortinithash}{6db6145dae8dc9e1271a8d556090b50a}
      \field{labelprefix}{A}
      \field{labelnamesource}{author}
      \field{labeltitlesource}{title}
      \field{annotation}{The author gives an excellent review of the history and known results concerning the weighted Liouville summatory function $L_\alpha(x) = \sum_{n\le x} \lambda(n)/n^\alpha$. Define $\delta_\alpha$ to be the logarithmic density of the set of positive real numbers~$x$ such that $L_\alpha(x) \le 0$. Assuming RH and LI and the estimate $\sum_{0<\gamma<T} |\zeta'(\rho)|^{-2} \ll T$, the author shows when $0\le\alpha<\frac12$ that $\frac12\le\delta_\alpha<1$, so that in particular, $L_\alpha(x)$ changes sign infinitely often. The author additionally shows under the same hypotheses that $\lim_{\alpha\to1/2-} \delta_\alpha = 1$. Moreover, the author proves (without needing the hypothesis LI) that $\delta_{1/2} = 1$, lending some credence to the conjecture that $L_{1/2}(x) \le 0$ for $x\ge17$.\par This article cites~\cites{1919.Polya,2000.Ng,MR4082248,MR0006202,MR0104638,MR0330069,MR634679,MR654520,MR1329368,MR1806291,MR2078705,MR2398787,MR3062002,MR3063909,MR98723,MR584557,MR43841}.}
      \field{issn}{0022-314X}
      \field{journaltitle}{J. Number Theory}
      \field{number}{2}
      \field{title}{The distribution of weighted sums of the {L}iouville function and {P}ólya's conjecture}
      \field{volume}{133}
      \field{year}{2013}
      \true{nocite}
      \field{pages}{545\bibrangedash 582}
      \range{pages}{38}
      \verb{doi}
      \verb 10.1016/j.jnt.2012.08.011
      \endverb
      \verb{mrnumber}
      \verb 2994374
      \endverb
    \endentry
    \entry{MR3163001}{article}{}
      \name{author}{1}{}{%
        {{hash=48017d687b31d05c591036588e04dd66}{%
           family={Humphries},
           familyi={H\bibinitperiod},
           given={Peter},
           giveni={P\bibinitperiod}}}%
      }
      \strng{namehash}{48017d687b31d05c591036588e04dd66}
      \strng{fullhash}{48017d687b31d05c591036588e04dd66}
      \strng{bibnamehash}{48017d687b31d05c591036588e04dd66}
      \strng{authorbibnamehash}{48017d687b31d05c591036588e04dd66}
      \strng{authornamehash}{48017d687b31d05c591036588e04dd66}
      \strng{authorfullhash}{48017d687b31d05c591036588e04dd66}
      \field{extraname}{2}
      \field{sortinit}{H}
      \field{sortinithash}{6db6145dae8dc9e1271a8d556090b50a}
      \field{labelprefix}{A}
      \field{labelnamesource}{author}
      \field{labeltitlesource}{title}
      \field{annotation}{The author gives a necessary and sufficient condition for an elliptic curve over a finite field to satisfy the function field analogue of the Mertens conjecture $\limsup_{X \to \infty} |M_{E/\mathbb{F}_q}(X)|/q^{X/2}\leq 1$ (as studied by the author in~\cite{MR3189983}), in terms of the order of the finite field and the trace of the Frobenius endomorphism acting on the curve. Moreover, the author shows if the Mertens conjecture holds for a given elliptic curve, then in fact $\limsup_{X \to \infty} |M_{E/\mathbb{F}_q}(X)|/q^{X/2} = 1$.\par This article cites~\cites{MR0330069,MR3335903,MR3684399,MR0006202,MR2282922,1897.Mertens,MR783538}.}
      \field{issn}{0004-9727}
      \field{journaltitle}{Bull. Aust. Math. Soc.}
      \field{number}{1}
      \field{title}{On the {M}ertens conjecture for elliptic curves over finite fields}
      \field{volume}{89}
      \field{year}{2014}
      \true{nocite}
      \field{pages}{19\bibrangedash 32}
      \range{pages}{14}
      \verb{doi}
      \verb 10.1017/S0004972712001116
      \endverb
      \verb{mrnumber}
      \verb 3163001
      \endverb
      \verb{urlraw}
      \verb https://doi.org/10.1017/S0004972712001116
      \endverb
      \verb{url}
      \verb https://doi.org/10.1017/S0004972712001116
      \endverb
    \endentry
    \entry{MR3189983}{article}{}
      \name{author}{1}{}{%
        {{hash=48017d687b31d05c591036588e04dd66}{%
           family={Humphries},
           familyi={H\bibinitperiod},
           given={Peter},
           giveni={P\bibinitperiod}}}%
      }
      \strng{namehash}{48017d687b31d05c591036588e04dd66}
      \strng{fullhash}{48017d687b31d05c591036588e04dd66}
      \strng{bibnamehash}{48017d687b31d05c591036588e04dd66}
      \strng{authorbibnamehash}{48017d687b31d05c591036588e04dd66}
      \strng{authornamehash}{48017d687b31d05c591036588e04dd66}
      \strng{authorfullhash}{48017d687b31d05c591036588e04dd66}
      \field{extraname}{3}
      \field{sortinit}{H}
      \field{sortinithash}{6db6145dae8dc9e1271a8d556090b50a}
      \field{labelprefix}{A}
      \field{labelnamesource}{author}
      \field{labeltitlesource}{title}
      \field{annotation}{Cha~\cite{MR3684399} introduced the Mertens function $M_{C/\mathbb{F}_q}(X)$ of a smooth projective curve defined over a finite field, and showed under LI that $E^{M_{C/\mathbb{F}_q}}(X)$ is bounded where $E^{M_{C/\mathbb{F}_q}}(X)=M_{C/\mathbb{F}_q}(X)/q^{X/2}$. Thus, a natural analogue of the Mertens conjecture for function fields would be $$\limsup_{X \to \infty} |E^{M_{C/\mathbb{F}_q}}(X)| \leq 1.$$ For fixed~$q$ and~$g$, let $\mathcal{H}_{2g+1,q^n}$ denote the set of hyperelliptic curves $y^2=f(x)$ over $\mathbb{F}_{q^n}$ arising from squarefree monic polynomials in $\mathbb{F}_{q^n}[x]$ of degree $2g+1$. The author shows that as $n \to \infty$, almost all curves in $\mathcal{H}_{2g+1,q^n}$ satisfy $\limsup_{X \to \infty} |E^{M_{C/\mathbb{F}_{q^n}}}(X)| > 1$, while for any $\beta>1$, a positive proportion of curves in $\mathcal{H}_{2g+1,q^n}$ satisfy $\limsup_{X \to \infty} |E^{M_{C/\mathbb{F}_{q^n}}}(X)| \leq \beta$.\par This article cites~\cites{MR0330069,MR3163001,MR3335903,MR3684399,MR0006202,MR2282922,1897.Mertens,MR2078705,MR783538,1913.vonSterneck}.}
      \field{issn}{1793-0421}
      \field{journaltitle}{Int. J. Number Theory}
      \field{number}{2}
      \field{title}{On the {M}ertens conjecture for function fields}
      \field{volume}{10}
      \field{year}{2014}
      \true{nocite}
      \field{pages}{341\bibrangedash 361}
      \range{pages}{21}
      \verb{doi}
      \verb 10.1142/S1793042113500978
      \endverb
      \verb{mrnumber}
      \verb 3189983
      \endverb
      \verb{urlraw}
      \verb https://doi.org/10.1142/S1793042113500978
      \endverb
      \verb{url}
      \verb https://doi.org/10.1142/S1793042113500978
      \endverb
    \endentry
    \entry{MR3996180}{article}{}
      \name{author}{3}{}{%
        {{hash=48017d687b31d05c591036588e04dd66}{%
           family={Humphries},
           familyi={H\bibinitperiod},
           given={Peter},
           giveni={P\bibinitperiod}}}%
        {{hash=d346d7afb6639c668a44a4e3ce8a8b98}{%
           family={Shekatkar},
           familyi={S\bibinitperiod},
           given={Snehal\bibnamedelima M.},
           giveni={S\bibinitperiod\bibinitdelim M\bibinitperiod}}}%
        {{hash=1a3f18d66237c8283c6ec23cd9a8c12d}{%
           family={Wong},
           familyi={W\bibinitperiod},
           given={Tian\bibnamedelima An},
           giveni={T\bibinitperiod\bibinitdelim A\bibinitperiod}}}%
      }
      \strng{namehash}{0212718ca31c4c1791497f3ee79e82ac}
      \strng{fullhash}{0212718ca31c4c1791497f3ee79e82ac}
      \strng{bibnamehash}{0212718ca31c4c1791497f3ee79e82ac}
      \strng{authorbibnamehash}{0212718ca31c4c1791497f3ee79e82ac}
      \strng{authornamehash}{0212718ca31c4c1791497f3ee79e82ac}
      \strng{authorfullhash}{0212718ca31c4c1791497f3ee79e82ac}
      \field{sortinit}{H}
      \field{sortinithash}{6db6145dae8dc9e1271a8d556090b50a}
      \field{labelprefix}{A}
      \field{labelnamesource}{author}
      \field{labeltitlesource}{title}
      \field{annotation}{Given a set~$S$ of integers, define $\Omega_S(n) = \#\{p\mid n\colon p\in S\}$. The authors study $\Sigma_S(x) = \sum_{n\le x} (-1)^{\Omega_S(n)}$ when~$S$ is a union of residue classes modulo~$q$. If~$S$ is the complete set of reduced residue classes\mod q then $\Sigma_S(x)$ is similar to $L(x)$; if~$q$ is prime and~$S$ is the set of quadratic nonresidues\mod q, then $\Sigma_S(x)$ is similar to $\sum_{n\le x} (\frac nq)$ except that $(-1)^{\Omega_S(n)}$ takes nonzero values at multiples of~$q$ as well.\par The authors show that the general behaviour of $\Sigma_S(x)$ is more predictable than these special cases: if~$S$ is the union of~$r$ reduced residue classes\mod q with $r\notin \{ \phi(q), \frac{\phi(q)}2 \}$, then $\Sigma_S(x) \sim b_0 {x}/{(\log x)^{2-2r/\varphi(q)}}$ for some explicit constant~$b_0$ that has the same sign as $\frac{\phi(q)}2-r$. If $r=\frac{\phi(q)}2$, then still $\Sigma_S(x)=o(x)$ provided that~$S$ is not a set of the form $\{n\colon \chi(n)=-1\}$ for some quadratic Dirichlet character $\chi\mod q$. Finally, suppose that $S = \{n\colon \chi(n)\ne-1\}$ for some quadratic $\chi\mod q$. Assuming GRH, LI, and $\sum_{0<\gamma\leq T} {|L'(\rho,\chi_q )|^-2} \ll T^{\theta}$ for some $1<\theta<3-\sqrt{3}$, the authors show that the logarithmic density~$\delta$ of $\{x>0\colon \Sigma_S(x)\ge 0 \}$ exists and satisfies $1/2\leq \delta<1$.\par This article cites~\cites{1919.Polya, MR3261965, MR3335903, MR2398787, MR0104638, MR2994374, MR0006202, MR3604916, MR3778223}.}
      \field{issn}{1246-7405}
      \field{journaltitle}{J. Théor. Nombres Bordeaux}
      \field{number}{1}
      \field{title}{Biases in prime factorizations and {L}iouville functions for arithmetic progressions}
      \field{volume}{31}
      \field{year}{2019}
      \true{nocite}
      \field{pages}{1\bibrangedash 25}
      \range{pages}{25}
      \verb{mrnumber}
      \verb 3996180
      \endverb
      \verb{urlraw}
      \verb http://jtnb.cedram.org/item?id=JTNB_2019__31_1_1_0
      \endverb
      \verb{url}
      \verb http://jtnb.cedram.org/item?id=JTNB_2019__31_1_1_0
      \endverb
    \endentry
    \entry{MR3739227}{article}{}
      \name{author}{1}{}{%
        {{hash=28c386582a9310ce9c19f006559d9e7a}{%
           family={Hurst},
           familyi={H\bibinitperiod},
           given={Greg},
           giveni={G\bibinitperiod}}}%
      }
      \strng{namehash}{28c386582a9310ce9c19f006559d9e7a}
      \strng{fullhash}{28c386582a9310ce9c19f006559d9e7a}
      \strng{bibnamehash}{28c386582a9310ce9c19f006559d9e7a}
      \strng{authorbibnamehash}{28c386582a9310ce9c19f006559d9e7a}
      \strng{authornamehash}{28c386582a9310ce9c19f006559d9e7a}
      \strng{authorfullhash}{28c386582a9310ce9c19f006559d9e7a}
      \field{sortinit}{H}
      \field{sortinithash}{6db6145dae8dc9e1271a8d556090b50a}
      \field{labelprefix}{A}
      \field{labelnamesource}{author}
      \field{labeltitlesource}{title}
      \field{annotation}{The authors extend Odlyzko and te Riele's disproof of Mertens conjecture~\cite{MR783538} by obtaining stronger bounds using more modern algorithms and computers. In particular the authors prove that \[ \liminf_{x \to \infty} E^M(x) < -1.837625 \quad \text{and} \quad \limsup_{x \to \infty} E^M(x) > 1.826054. \] The authors compute $M(x)$ for $x\le10^{16}$ and for $x=2^{54},\ldots,2^{73}$. Their algorithm operates in logspace, lowering the amount of storage required during computations. The authors note that cache misses were a significant source of increased computing time for larger values of $M(x)$; they suggest that a future algorithm could further reduce cache size by storing Möbius values in two bits, rather than bytes.\par This article cites~\cite{MR783538,MR0006202,MR0202685,MR2078705,MR2118272,MR0046485,MR3335903}.}
      \field{issn}{0025-5718}
      \field{journaltitle}{Math. Comp.}
      \field{number}{310}
      \field{title}{Computations of the Mertens function and improved bounds on the Mertens conjecture}
      \field{volume}{87}
      \field{year}{2018}
      \true{nocite}
      \field{pages}{1013\bibrangedash 1028}
      \range{pages}{16}
      \verb{doi}
      \verb 10.1090/mcom/3275
      \endverb
      \verb{mrnumber}
      \verb 3739227
      \endverb
    \endentry
    \entry{1932.Ingham}{book}{}
      \name{author}{1}{}{%
        {{hash=a9ad153a68cf8125996aecb67d0f1f8c}{%
           family={Ingham},
           familyi={I\bibinitperiod},
           given={A.\bibnamedelimi E.},
           giveni={A\bibinitperiod\bibinitdelim E\bibinitperiod}}}%
      }
      \list{publisher}{1}{%
        {Cambridge Tracts in Mathematics and Mathematical Physics. 30. London: Cambridge University Press}%
      }
      \strng{namehash}{a9ad153a68cf8125996aecb67d0f1f8c}
      \strng{fullhash}{a9ad153a68cf8125996aecb67d0f1f8c}
      \strng{bibnamehash}{a9ad153a68cf8125996aecb67d0f1f8c}
      \strng{authorbibnamehash}{a9ad153a68cf8125996aecb67d0f1f8c}
      \strng{authornamehash}{a9ad153a68cf8125996aecb67d0f1f8c}
      \strng{authorfullhash}{a9ad153a68cf8125996aecb67d0f1f8c}
      \field{extraname}{1}
      \field{sortinit}{I}
      \field{sortinithash}{9417e9a1288a9371e2691d999083ed39}
      \field{labelprefix}{A}
      \field{labelnamesource}{author}
      \field{labeltitlesource}{title}
      \field{annotation}{In Chapter III of this book (Further Theory of $\zeta(s)$.~Applications), the author shows how information on the zero-free region of $\zeta(s)$ leads to improved estimates for $\Delta^\psi(x)$. Let $\eta(t)$ be a real-valued decreasing function defined on $t \geq 0$ satisfying $0 \leq \eta(t) \leq \frac{1}{2}$ and $1/\eta(t)=O(\log t)$, and suppose that $\eta'(t)$ is continuous and $\eta'(t)=o(1)$. If $\zeta(s)$ has no zeros in the region $\{\sigma>1-\eta(|t|)\}$, then for any $\alpha \in (0,1)$ we have $\Delta^\psi(x) \ll x \exp\bigl(-\frac{1}{2}\alpha \omega(x)\bigr)$, where $\omega(x) = \min_{t\ge1} \bigl( \eta(t) \log x+\log t \bigr)$.\par Chapter V (Irregularities of Distribution) gathers known results on comparative prime number theory from the literature. The author also outlines the proof of \[ \pi(x;4,3,1) = \Omega_\pm\bigg( \frac{\sqrt x}{\log x} \log\log\log x \bigg), \] which was asserted by Hardy and Littlewood~\cite{MR1555148}.\par This book cites~\cites{1909.Landau,MR1544509,MR1555148,MR1511360,MR1544293A,MR1544293B,1914.Littlewood,MR1574052,1891.Phragmen,1930.Polya}.}
      \field{pagetotal}{114}
      \field{title}{{The distribution of prime numbers}}
      \field{year}{1932}
      \true{nocite}
    \endentry
    \entry{1936.Ingham}{article}{}
      \name{author}{1}{}{%
        {{hash=a9ad153a68cf8125996aecb67d0f1f8c}{%
           family={Ingham},
           familyi={I\bibinitperiod},
           given={A.\bibnamedelimi E.},
           giveni={A\bibinitperiod\bibinitdelim E\bibinitperiod}}}%
      }
      \strng{namehash}{a9ad153a68cf8125996aecb67d0f1f8c}
      \strng{fullhash}{a9ad153a68cf8125996aecb67d0f1f8c}
      \strng{bibnamehash}{a9ad153a68cf8125996aecb67d0f1f8c}
      \strng{authorbibnamehash}{a9ad153a68cf8125996aecb67d0f1f8c}
      \strng{authornamehash}{a9ad153a68cf8125996aecb67d0f1f8c}
      \strng{authorfullhash}{a9ad153a68cf8125996aecb67d0f1f8c}
      \field{extraname}{2}
      \field{sortinit}{I}
      \field{sortinithash}{9417e9a1288a9371e2691d999083ed39}
      \field{labelprefix}{A}
      \field{labelnamesource}{author}
      \field{labeltitlesource}{title}
      \field{annotation}{This article gives another proof of an explicit version of Littlewood's theorem~\cite{MR1555148} and establishes the following stronger result: Assuming SA for $\zeta(s)$, there exists an absolute constant $A>1$ such that, for all $x>1$, the interval $(x,Ax)$ contains a sign change of $\pi(x)-\Li(x)$. The author highlights his use of Fejér kernels in the proof, in contrast with the Poisson kernel used by Skewes~\cite{MR1573970}.\par This article cites~\cites{MR1555148, 1930.Polya, MR1573970}.}
      \field{journaltitle}{Acta Arith.}
      \field{title}{{A note on the distribution of primes}}
      \field{volume}{1}
      \field{year}{1936}
      \true{nocite}
      \field{pages}{201\bibrangedash 211}
      \range{pages}{11}
    \endentry
    \entry{MR0006202}{article}{}
      \name{author}{1}{}{%
        {{hash=a9ad153a68cf8125996aecb67d0f1f8c}{%
           family={Ingham},
           familyi={I\bibinitperiod},
           given={A.\bibnamedelimi E.},
           giveni={A\bibinitperiod\bibinitdelim E\bibinitperiod}}}%
      }
      \strng{namehash}{a9ad153a68cf8125996aecb67d0f1f8c}
      \strng{fullhash}{a9ad153a68cf8125996aecb67d0f1f8c}
      \strng{bibnamehash}{a9ad153a68cf8125996aecb67d0f1f8c}
      \strng{authorbibnamehash}{a9ad153a68cf8125996aecb67d0f1f8c}
      \strng{authornamehash}{a9ad153a68cf8125996aecb67d0f1f8c}
      \strng{authorfullhash}{a9ad153a68cf8125996aecb67d0f1f8c}
      \field{extraname}{3}
      \field{sortinit}{I}
      \field{sortinithash}{9417e9a1288a9371e2691d999083ed39}
      \field{labelprefix}{A}
      \field{labelnamesource}{author}
      \field{labeltitlesource}{title}
      \field{annotation}{In this article, Ingham probes conjectured bounds for the summatory functions $M(x)$ and $L(x)$. He proves that the truth, for sufficiently large $x$, of any one of the inequalities $M(x) < Kx^{1/2}$, $M(x) >- Kx^{1/2}$, $L(x) < Kx^{1/2}$, or $L(x) >- Kx^{1/2}$ (where $K$ is a constant) would imply not only RH and the simplicity of the zeros of $\zeta(s)$ (as was ``well known''), but also the falsity of LI. Of this last assumption, the author writes: ``It would be easy to relax this hypothesis a little, but there seems no obvious way of replacing it by anything essentially easier to verify.'' Indeed, he shows that if there are only finitely many rational linear relations among the positive imaginary parts of these zeros, then $E^M(x)$ and $E^L(x)$ would be unbounded both above and below, contrary to existing conjectures.\par The method of the proof is similar to Littlewood's disproof of the conjecture $\pi(x) < \li(x)$ in~\cite{MR1555148}, including a reliance on trigonometric polynomials involving the zeros of $\zeta(s)$, except that Dirichlet's theorem on homogeneous Diophantine approximation is replaced by Kronecker's theorem on inhomogeneous Diophantine approximation. For the proof, the author establishes two main results, one concerning Laplace transforms of real trigonometric polynomials, and the other establishing the divergence (assuming RH) of the two residue series $\sum_{\gamma>0} 1/\rho\zeta'(\rho)$ and $\sum_{\gamma>0} \zeta(2\rho)/\rho\zeta'(\rho)$.\par This article cites~\cites{1897.Mertens, 1897.vonSterneck, 1901.vonSterneck, 1905.Stieltjes, 1913.vonSterneck, 1919.Polya,MR0003644}.}
      \field{issn}{0002-9327}
      \field{journaltitle}{Amer. J. Math.}
      \field{title}{On two conjectures in the theory of numbers}
      \field{volume}{64}
      \field{year}{1942}
      \true{nocite}
      \field{pages}{313\bibrangedash 319}
      \range{pages}{7}
      \verb{doi}
      \verb 10.2307/2371685
      \endverb
      \verb{mrnumber}
      \verb 0006202
      \endverb
      \verb{urlraw}
      \verb https://doi.org/10.2307/2371685
      \endverb
      \verb{url}
      \verb https://doi.org/10.2307/2371685
      \endverb
    \endentry
    \entry{MR0184920}{book}{}
      \name{author}{1}{}{%
        {{hash=a9ad153a68cf8125996aecb67d0f1f8c}{%
           family={Ingham},
           familyi={I\bibinitperiod},
           given={A.\bibnamedelimi E.},
           giveni={A\bibinitperiod\bibinitdelim E\bibinitperiod}}}%
      }
      \list{publisher}{1}{%
        {Stechert-Hafner, Inc., New York}%
      }
      \strng{namehash}{a9ad153a68cf8125996aecb67d0f1f8c}
      \strng{fullhash}{a9ad153a68cf8125996aecb67d0f1f8c}
      \strng{bibnamehash}{a9ad153a68cf8125996aecb67d0f1f8c}
      \strng{authorbibnamehash}{a9ad153a68cf8125996aecb67d0f1f8c}
      \strng{authornamehash}{a9ad153a68cf8125996aecb67d0f1f8c}
      \strng{authorfullhash}{a9ad153a68cf8125996aecb67d0f1f8c}
      \field{extraname}{4}
      \field{sortinit}{I}
      \field{sortinithash}{9417e9a1288a9371e2691d999083ed39}
      \field{labelprefix}{A}
      \field{labelnamesource}{author}
      \field{labeltitlesource}{title}
      \field{annotation}{This is a reprint of \cite{1932.Ingham}.}
      \field{series}{Cambridge Tracts in Mathematics and Mathematical Physics, No. 30}
      \field{title}{The distribution of prime numbers}
      \field{year}{1964}
      \true{nocite}
      \field{pages}{v+114}
      \range{pages}{-1}
      \verb{mrnumber}
      \verb 0184920
      \endverb
    \endentry
    \entry{MR1501802}{article}{}
      \name{author}{2}{}{%
        {{hash=5acf2ef9928231591e18fa4f528b72bd}{%
           family={Jessen},
           familyi={J\bibinitperiod},
           given={B.},
           giveni={B\bibinitperiod}}}%
        {{hash=0ecb7b4bc82132c8d3551c023cbe496a}{%
           family={Wintner},
           familyi={W\bibinitperiod},
           given={A.},
           giveni={A\bibinitperiod}}}%
      }
      \strng{namehash}{8f1d2457d99bfffa1d82626b836e82d6}
      \strng{fullhash}{8f1d2457d99bfffa1d82626b836e82d6}
      \strng{bibnamehash}{8f1d2457d99bfffa1d82626b836e82d6}
      \strng{authorbibnamehash}{8f1d2457d99bfffa1d82626b836e82d6}
      \strng{authornamehash}{8f1d2457d99bfffa1d82626b836e82d6}
      \strng{authorfullhash}{8f1d2457d99bfffa1d82626b836e82d6}
      \field{sortinit}{J}
      \field{sortinithash}{c45040a764d616897e7f5b30174d7b92}
      \field{labelprefix}{A}
      \field{labelnamesource}{author}
      \field{labeltitlesource}{title}
      \field{annotation}{This article contains a systematic study of infinite convolutions of distribution functions (densities) on $\R^k$, and consequently of infinite sums of independent random variables, via their Fourier transforms; they give criteria for when these convolutions converge and when the result is singular or absolutely continuous. They analyze distributions of functions on convex closed curves in~$\R^2$. The results are applied to the limiting distribution functions of almost-periodic functions of multiple types. As an application, the authors derive results on the limiting distributions of $\log\zeta(\sigma+it)$ and $\zeta(\sigma+it)$ for $\frac12<\sigma\le1$ (the case $\sigma>1$ being simpler).}
      \field{issn}{0002-9947}
      \field{journaltitle}{Trans. Amer. Math. Soc.}
      \field{number}{1}
      \field{title}{Distribution functions and the Riemann zeta function}
      \field{volume}{38}
      \field{year}{1935}
      \true{nocite}
      \field{pages}{48\bibrangedash 88}
      \range{pages}{41}
      \verb{doi}
      \verb 10.2307/1989728
      \endverb
      \verb{mrnumber}
      \verb 1501802
      \endverb
    \endentry
    \entry{MR429789}{article}{}
      \name{author}{2}{}{%
        {{hash=b603cd4e91f630615d9bcde63b8a7ffa}{%
           family={Jurkat},
           familyi={J\bibinitperiod},
           given={W.},
           giveni={W\bibinitperiod}}}%
        {{hash=c33ce895e089b8d5217e8f2e68f087ec}{%
           family={Peyerimhoff},
           familyi={P\bibinitperiod},
           given={A.},
           giveni={A\bibinitperiod}}}%
      }
      \strng{namehash}{a2de3531375e01cede375ae506c851d2}
      \strng{fullhash}{a2de3531375e01cede375ae506c851d2}
      \strng{bibnamehash}{a2de3531375e01cede375ae506c851d2}
      \strng{authorbibnamehash}{a2de3531375e01cede375ae506c851d2}
      \strng{authornamehash}{a2de3531375e01cede375ae506c851d2}
      \strng{authorfullhash}{a2de3531375e01cede375ae506c851d2}
      \field{sortinit}{J}
      \field{sortinithash}{c45040a764d616897e7f5b30174d7b92}
      \field{labelprefix}{A}
      \field{labelnamesource}{author}
      \field{labeltitlesource}{title}
      \field{annotation}{Similar to the first author's work~\cite{MR0352026}, the authors reduce the problem of improving the bounds on $\limsup E^M(x)$ and $\liminf E^M(x)$ to finding a reasonably good solution to an inhomogenous Diophantine approximation problem. Their new constructive algorithm leads to the following improved bounds related to the Mertens conjecture: \[ \liminf_{x \to \infty} E^M(x) < -0.638 \quad \text{and} \quad \limsup_{x \to \infty} E^M(x) > 0.779. \] The authors remark that the smallest counterexample to the Mertens conjecture is likely to be at least $\exp(4.16 \times 10^{14})$. \par This article cites~\cite{MR0330069,MR0332684,1897.Mertens,1897.vonSterneck,1901.vonSterneck,1912.vonSterneck,1913.vonSterneck,MR0046485,MR155787,MR0352026,MR0006202,MR228456}.}
      \field{issn}{0075-4102}
      \field{journaltitle}{J. Reine Angew. Math.}
      \field{title}{A constructive approach to {K}ronecker approximations and its application to the {M}ertens conjecture}
      \field{volume}{286(287)}
      \field{year}{1976}
      \true{nocite}
      \field{pages}{322\bibrangedash 340}
      \range{pages}{19}
      \verb{doi}
      \verb 10.1515/crll.1976.286-287.322
      \endverb
      \verb{mrnumber}
      \verb 429789
      \endverb
      \verb{urlraw}
      \verb https://doi.org/10.1515/crll.1976.286-287.322
      \endverb
      \verb{url}
      \verb https://doi.org/10.1515/crll.1976.286-287.322
      \endverb
    \endentry
    \entry{MR0352026}{article}{}
      \name{author}{1}{}{%
        {{hash=b9201f1e3293ed411cec93a708d9b9b4}{%
           family={Jurkat},
           familyi={J\bibinitperiod},
           given={W.\bibnamedelimi B.},
           giveni={W\bibinitperiod\bibinitdelim B\bibinitperiod}}}%
      }
      \list{publisher}{1}{%
        {Amer. Math. Soc., Providence, R.I.}%
      }
      \strng{namehash}{b9201f1e3293ed411cec93a708d9b9b4}
      \strng{fullhash}{b9201f1e3293ed411cec93a708d9b9b4}
      \strng{bibnamehash}{b9201f1e3293ed411cec93a708d9b9b4}
      \strng{authorbibnamehash}{b9201f1e3293ed411cec93a708d9b9b4}
      \strng{authornamehash}{b9201f1e3293ed411cec93a708d9b9b4}
      \strng{authorfullhash}{b9201f1e3293ed411cec93a708d9b9b4}
      \field{sortinit}{J}
      \field{sortinithash}{c45040a764d616897e7f5b30174d7b92}
      \field{labelprefix}{A}
      \field{labelnamesource}{author}
      \field{labeltitlesource}{title}
      \field{annotation}{Define $C=2+\sum_{n=1}^\infty {(-1)^n(2\pi x)^{2n}}/{n(2n)!\zeta(2n+1)} \approx -.505$. The author proves unconditionally that $\limsup_{x\to\infty} E^M(x) > 1+C$ and $\liminf_{x\to\infty} E^M(x) < C$; in particular, the latter inequality disproves the conjecture of von Sterneck~\cite{1901.vonSterneck} that $|M(x)| \le \frac12\sqrt x$. The main tool is a general oscillation result for almost periodic functions in the distributional sense (combined with Landau's theorem to reduce to the RH case). The author's general theorems recover $\Omega$-results proved by Hardy and Littlewood for $\Delta^\psi(x)$ and $\Delta^D(x)$.\par This article cites~\cites{MR0330069,MR0332684,MR1576550,1932.Ingham,1936.Ingham,MR0006202,1914.Littlewood,MR155787,MR1573970,1901.vonSterneck,MR0046485}.}
      \field{booktitle}{Analytic number theory ({P}roc. {S}ympos. {P}ure {M}ath., {V}ol. {XXIV}, {S}t. {L}ouis {U}niv., {S}t. {L}ouis, {M}o., 1972)}
      \field{title}{On the {M}ertens conjecture and related general $\Omega$-theorems}
      \field{year}{1973}
      \true{nocite}
      \field{pages}{147\bibrangedash 158}
      \range{pages}{12}
      \verb{mrnumber}
      \verb 0352026
      \endverb
    \endentry
    \entry{MR777013}{article}{}
      \name{author}{1}{}{%
        {{hash=d1e4ad71dae945a6977ecf6f9ca65c30}{%
           family={Kaczorowski},
           familyi={K\bibinitperiod},
           given={J.},
           giveni={J\bibinitperiod}}}%
      }
      \strng{namehash}{d1e4ad71dae945a6977ecf6f9ca65c30}
      \strng{fullhash}{d1e4ad71dae945a6977ecf6f9ca65c30}
      \strng{bibnamehash}{d1e4ad71dae945a6977ecf6f9ca65c30}
      \strng{authorbibnamehash}{d1e4ad71dae945a6977ecf6f9ca65c30}
      \strng{authornamehash}{d1e4ad71dae945a6977ecf6f9ca65c30}
      \strng{authorfullhash}{d1e4ad71dae945a6977ecf6f9ca65c30}
      \field{extraname}{1}
      \field{sortinit}{K}
      \field{sortinithash}{d3edc18d54b9438a72c24c925bfb38f4}
      \field{labelprefix}{A}
      \field{labelnamesource}{author}
      \field{labeltitlesource}{title}
      \field{annotation}{This article establishes a lower bound on the growth of the number of sign changes of $\Delta^\psi(x)$ and $\Delta^\Pi(x)$. Specifically, the author proves that $W^\psi(T) \ge \frac{\gamma_1}{4\pi} \log T$ (and the same for $W^\Pi(x)$) when $T$ is sufficiently large (effectively), where $\gamma_1 \approx 14.1347$. The key technique is to bound $W^\psi(T)$ below by $W\big( A_n^\psi;T \big)$, the number of sign changes of repeated logarithmic integrals of $\Delta^\psi(x)$; using the fact that the second-lowest nontrivial zero of $\zeta(s)$ has imaginary part exceeding $15$, the author derives an explicit formula for $W\big( A_n^\psi;T \big)$ when $n\asymp \log T$ is suitably chosen.\par This article cites~\cites{MR1555148, 1936.Ingham, 1932.Ingham, MR0133309, MR0133308, MR0439771, MR0439772, MR0202686, MR0439774, 1914.Littlewood, MR0439773, MR607089, MR630377, 1930.Polya, MR1511206, MR1573970, MR0067145}.}
      \field{issn}{0065-1036}
      \field{journaltitle}{Acta Arith.}
      \field{number}{4}
      \field{title}{On sign-changes in the remainder-term of the prime-number formula. {I}}
      \field{volume}{44}
      \field{year}{1984}
      \true{nocite}
      \field{pages}{365\bibrangedash 377}
      \range{pages}{13}
      \verb{doi}
      \verb 10.4064/aa-44-4-365-377
      \endverb
      \verb{mrnumber}
      \verb 777013
      \endverb
      \verb{urlraw}
      \verb https://doi.org/10.4064/aa-44-4-365-377
      \endverb
      \verb{url}
      \verb https://doi.org/10.4064/aa-44-4-365-377
      \endverb
    \endentry
    \entry{MR791085}{article}{}
      \name{author}{1}{}{%
        {{hash=d1e4ad71dae945a6977ecf6f9ca65c30}{%
           family={Kaczorowski},
           familyi={K\bibinitperiod},
           given={J.},
           giveni={J\bibinitperiod}}}%
      }
      \strng{namehash}{d1e4ad71dae945a6977ecf6f9ca65c30}
      \strng{fullhash}{d1e4ad71dae945a6977ecf6f9ca65c30}
      \strng{bibnamehash}{d1e4ad71dae945a6977ecf6f9ca65c30}
      \strng{authorbibnamehash}{d1e4ad71dae945a6977ecf6f9ca65c30}
      \strng{authornamehash}{d1e4ad71dae945a6977ecf6f9ca65c30}
      \strng{authorfullhash}{d1e4ad71dae945a6977ecf6f9ca65c30}
      \field{extraname}{2}
      \field{sortinit}{K}
      \field{sortinithash}{d3edc18d54b9438a72c24c925bfb38f4}
      \field{labelprefix}{A}
      \field{labelnamesource}{author}
      \field{labeltitlesource}{title}
      \field{annotation}{The author proves unconditionally that $W^\pi(T) \gg \log T$, though with an ineffective constant, and the same for $W^\theta(T)$. He also proves unconditionally that $\liminf_{T\to\infty} {W^\psi(T)}/{\log T} \ge {\gamma(\Theta)}/{\pi}$, where ${\gamma(\Theta)}$ is the smallest $\gamma>0$ such that $\zeta(\Theta+i\gamma)=0$ (or ${\gamma(\Theta)}=\infty$ if $\Theta$ is not attained); this improves a result of Pólya~\cite{1930.Polya}, which had $\limsup$ in place of $\liminf$. He remarks that if RH is false, then the proof of this latter result can be extended to $W^\theta(T)$ and (with a bit more difficulty) to $W^\Pi(T)$ and $W^\pi(T)$. As in a previous article, the proofs of both theorems make use of the iterated averages $A^f_n(x)$.\par This article cites~\cites{MR1555148, 1936.Ingham, MR777013, MR0133309, MR0133308, MR0439771, MR0439772, MR0439774, 1914.Littlewood, MR0439773, 1930.Polya}.}
      \field{issn}{0065-1036}
      \field{journaltitle}{Acta Arith.}
      \field{number}{1}
      \field{title}{On sign-changes in the remainder-term of the prime-number formula. {II}}
      \field{volume}{45}
      \field{year}{1985}
      \true{nocite}
      \field{pages}{65\bibrangedash 74}
      \range{pages}{10}
      \verb{doi}
      \verb 10.4064/aa-45-1-65-74
      \endverb
      \verb{mrnumber}
      \verb 791085
      \endverb
      \verb{urlraw}
      \verb https://doi.org/10.4064/aa-45-1-65-74
      \endverb
      \verb{url}
      \verb https://doi.org/10.4064/aa-45-1-65-74
      \endverb
    \endentry
    \entry{MR927376}{article}{}
      \name{author}{1}{}{%
        {{hash=d1e4ad71dae945a6977ecf6f9ca65c30}{%
           family={Kaczorowski},
           familyi={K\bibinitperiod},
           given={J.},
           giveni={J\bibinitperiod}}}%
      }
      \strng{namehash}{d1e4ad71dae945a6977ecf6f9ca65c30}
      \strng{fullhash}{d1e4ad71dae945a6977ecf6f9ca65c30}
      \strng{bibnamehash}{d1e4ad71dae945a6977ecf6f9ca65c30}
      \strng{authorbibnamehash}{d1e4ad71dae945a6977ecf6f9ca65c30}
      \strng{authornamehash}{d1e4ad71dae945a6977ecf6f9ca65c30}
      \strng{authorfullhash}{d1e4ad71dae945a6977ecf6f9ca65c30}
      \field{extraname}{3}
      \field{sortinit}{K}
      \field{sortinithash}{d3edc18d54b9438a72c24c925bfb38f4}
      \field{labelprefix}{A}
      \field{labelnamesource}{author}
      \field{labeltitlesource}{title}
      \field{annotation}{This article examines $W(\Delta^\psi_e;T)$, the number of sign changes of $\Delta^\psi_e(x) = \sum_{n=1}^\infty (\Lambda(n) - 1)e^{-n/x} $ in the interval $[0,T]$. The author first proves that SA (which he calls ``Ingham's condition'') is equivalent to the assertion that $W\big( \Delta^\psi_e;[T,eT] \big) \ll 1$ uniformly for $T>0$, which is further equivalent to each of $\limsup_{T\to\infty} W( \Delta^\psi_e;T ) < \infty$ and $\liminf_{T\to\infty} W( \Delta^\psi_e;T ) < \infty$. Assuming SA and LI$(\Theta)$, the author proves that $W(\Delta^\psi_e;T) \sim \kappa \log T$ as $T \to \infty$, for a constant $\kappa$ (given by an explicit integral) depending on the zeros of $\zeta(s)$ on the line $\sigma = \Theta$. Finally, assuming RH, the author proves that $W(\Delta^\psi_e;T) = \frac{\gamma_1}{\pi} \log T + O(1)$, where $\gamma_1 \approx 14.13$, and indeed that these sign changes are extremely regularly spaced and correspond to oscillations that are $\gg \sqrt x$. These results support the author's conjecture that $W(\Delta^\psi_e;T) \sim c\log T$ as $T \to \infty$. \par This article cites~\cites{MR1555148,1936.Ingham,MR777013,MR791085,MR0131410}.}
      \field{issn}{0065-1036}
      \field{journaltitle}{Acta Arith.}
      \field{number}{4}
      \field{title}{On sign-changes in the remainder-term of the prime-number formula. {III}}
      \field{volume}{48}
      \field{year}{1987}
      \true{nocite}
      \field{pages}{347\bibrangedash 371}
      \range{pages}{25}
      \verb{doi}
      \verb 10.4064/aa-48-4-347-371
      \endverb
      \verb{mrnumber}
      \verb 927376
      \endverb
      \verb{urlraw}
      \verb https://doi.org/10.4064/aa-48-4-347-371
      \endverb
      \verb{url}
      \verb https://doi.org/10.4064/aa-48-4-347-371
      \endverb
    \endentry
    \entry{MR945273}{article}{}
      \name{author}{1}{}{%
        {{hash=d1e4ad71dae945a6977ecf6f9ca65c30}{%
           family={Kaczorowski},
           familyi={K\bibinitperiod},
           given={J.},
           giveni={J\bibinitperiod}}}%
      }
      \strng{namehash}{d1e4ad71dae945a6977ecf6f9ca65c30}
      \strng{fullhash}{d1e4ad71dae945a6977ecf6f9ca65c30}
      \strng{bibnamehash}{d1e4ad71dae945a6977ecf6f9ca65c30}
      \strng{authorbibnamehash}{d1e4ad71dae945a6977ecf6f9ca65c30}
      \strng{authornamehash}{d1e4ad71dae945a6977ecf6f9ca65c30}
      \strng{authorfullhash}{d1e4ad71dae945a6977ecf6f9ca65c30}
      \field{extraname}{4}
      \field{sortinit}{K}
      \field{sortinithash}{d3edc18d54b9438a72c24c925bfb38f4}
      \field{labelprefix}{A}
      \field{labelnamesource}{author}
      \field{labeltitlesource}{title}
      \field{annotation}{The author proves, when $\Theta>\frac12$, that for any $\ep>0$ we have $\max_{T \le x \le (1+\ep)T} |\Delta_e^\psi(x)| \gg_\ep T^{\Theta-\ep}$. In light of the author's previous results~\cite{MR927376} that assumed RH, it follows that unconditionally (but ineffectively), $\max_{T \le x \le (1+\ep)T} |\Delta_e^\psi(x)| \gg_\ep \sqrt T$. The author also deduces that $W_e^\psi(T) = o(\log^2 T)$, and sketches a construction (of a ``barrier'') showing that this result cannot be improved without further information on the zeros of~$\zeta(s)$.\par This article cites~\cites{MR0125822,MR0131410,MR777013,MR791085,MR927376, MR0060548}.}
      \field{issn}{0065-1036}
      \field{journaltitle}{Acta Arith.}
      \field{number}{1}
      \field{title}{On sign-changes in the remainder-term of the prime-number formula. {IV}}
      \field{volume}{50}
      \field{year}{1988}
      \true{nocite}
      \field{pages}{15\bibrangedash 21}
      \range{pages}{7}
      \verb{doi}
      \verb 10.4064/aa-50-1-15-21
      \endverb
      \verb{mrnumber}
      \verb 945273
      \endverb
      \verb{urlraw}
      \verb https://doi.org/10.4064/aa-50-1-15-21
      \endverb
      \verb{url}
      \verb https://doi.org/10.4064/aa-50-1-15-21
      \endverb
    \endentry
    \entry{MR1083000}{article}{}
      \name{author}{1}{}{%
        {{hash=d1e4ad71dae945a6977ecf6f9ca65c30}{%
           family={Kaczorowski},
           familyi={K\bibinitperiod},
           given={J.},
           giveni={J\bibinitperiod}}}%
      }
      \strng{namehash}{d1e4ad71dae945a6977ecf6f9ca65c30}
      \strng{fullhash}{d1e4ad71dae945a6977ecf6f9ca65c30}
      \strng{bibnamehash}{d1e4ad71dae945a6977ecf6f9ca65c30}
      \strng{authorbibnamehash}{d1e4ad71dae945a6977ecf6f9ca65c30}
      \strng{authornamehash}{d1e4ad71dae945a6977ecf6f9ca65c30}
      \strng{authorfullhash}{d1e4ad71dae945a6977ecf6f9ca65c30}
      \field{extraname}{5}
      \field{sortinit}{K}
      \field{sortinithash}{d3edc18d54b9438a72c24c925bfb38f4}
      \field{labelprefix}{A}
      \field{labelnamesource}{author}
      \field{labeltitlesource}{title}
      \field{annotation}{This is the first in a series of articles on the ``$k$-functions'' $k(z,\chi)$ and $K(z,\chi)$ and certain limiting values $F(x,\chi)$ of the latter (see Section~\ref{k-functions section} of this annotated bibliography for definitions). In Section~$3$, the author proves that $k(z,\chi)$ can be analytically continued to a meromorphic function on the Riemann surface~$\cM$ for $\log z$, and indeed that \begin{align*} k(z, \chi) - \frac{1}{2\pi i} \frac{e^z}{e^z-1} \log z \end{align*} is meromorphic and single-valued on $\mathbb C$. Indeed, the author finds all of the singularities of $k(z,\chi)$ on~$\cM$ (all simple poles) and their residues. He also establishes the functional equations \[ k(z,\chi) + e^z k(z^*, \overline{\chi}) = D(z,\chi), \qquad k(z,\chi) + \overline{k(z^c, \overline{\chi})} = e^zD(-z,\chi). \] In Section~$4$, the author establishes explicit formulas for $\psi_0(x,\chi)$ and $\psi_{0r}(x,\chi)$, stated in the forms \begin{align*} \qquad F(x, \chi) +\sum_{\substack{\beta>0 \\ L(\beta, \chi)=0}} \frac{e^{\beta x}}{\beta} = -\psi_0(e^x, \chi) - R_{\chi(-1)}(x) + B(\chi) + \begin{cases} e^x, &\text{if } \chi=\chi_0, \\ -x, &\text{if } \chi\ne\chi_0 \text{ and } \chi(-1)=1, \\ 0, &\text{if } \chi(-1) = -1 \end{cases} \end{align*} for $x>0$, and \begin{align*} \qquad \,\, F(x, \chi) + \sum_{\substack{\beta>0 \\ L(\beta, \chi)=0}} \frac{e^{\beta x}}{\beta} = \psi_{0r}(e^{|x|}, \chi) + R_{-\chi(-1)}(|x|) + C(\chi) + \begin{cases} e^x, &\text{if } \chi=\chi_0, \\ x, &\text{if } \chi\ne\chi_0 \text{ and } \chi(-1)=1, \\ 0, &\text{if } \chi(-1) = -1 \end{cases} \end{align*} for $x<0$. The author then shows that the left-hand side is equal to the series $\sum_{\rho} {e^{\rho x}}/{\rho}$ as in the classical explicit formulas for the right-hand sides.\par This article cites~\cites{1932.Ingham}.}
      \field{issn}{0065-1036}
      \field{journaltitle}{Acta Arith.}
      \field{number}{3}
      \field{title}{The {$k$}-functions in multiplicative number theory. {I}. {O}n complex explicit formulae}
      \field{volume}{56}
      \field{year}{1990}
      \true{nocite}
      \field{pages}{195\bibrangedash 211}
      \range{pages}{17}
      \verb{doi}
      \verb 10.4064/aa-56-3-195-211
      \endverb
      \verb{mrnumber}
      \verb 1083000
      \endverb
      \verb{urlraw}
      \verb https://doi.org/10.4064/aa-56-3-195-211
      \endverb
      \verb{url}
      \verb https://doi.org/10.4064/aa-56-3-195-211
      \endverb
    \endentry
    \entry{MR1083001}{article}{}
      \name{author}{1}{}{%
        {{hash=d1e4ad71dae945a6977ecf6f9ca65c30}{%
           family={Kaczorowski},
           familyi={K\bibinitperiod},
           given={J.},
           giveni={J\bibinitperiod}}}%
      }
      \strng{namehash}{d1e4ad71dae945a6977ecf6f9ca65c30}
      \strng{fullhash}{d1e4ad71dae945a6977ecf6f9ca65c30}
      \strng{bibnamehash}{d1e4ad71dae945a6977ecf6f9ca65c30}
      \strng{authorbibnamehash}{d1e4ad71dae945a6977ecf6f9ca65c30}
      \strng{authornamehash}{d1e4ad71dae945a6977ecf6f9ca65c30}
      \strng{authorfullhash}{d1e4ad71dae945a6977ecf6f9ca65c30}
      \field{extraname}{6}
      \field{sortinit}{K}
      \field{sortinithash}{d3edc18d54b9438a72c24c925bfb38f4}
      \field{labelprefix}{A}
      \field{labelnamesource}{author}
      \field{labeltitlesource}{title}
      \field{annotation}{Let $0 < \gamma_1 \leq \gamma_2 \leq \cdots$ denote the imaginary parts of nontrivial zeros of $L(s, \chi)$ in the upper half plane. In this article Kaczorowski defines a positive Toeplitz matrix $\boldsymbol{A}=(a_{nk})$ by $a_{nk} = e^{-\gamma_k} \gamma_k^n \left(\sum_{h=1}^\infty e^{-\gamma_h}\gamma_h^n\right)^{-1}$ (for $n, k \geq 1$), and proves that for any nonzero real $x$, the sequence $(x\gamma_n)_{n=1}^\infty$ is $\boldsymbol{A}$-uniformly distributed\mod1; the known result that the $x\gamma_n$ are uniformly distributed\mod1 (in the sense of Weyl) follows as a corollary.\par This article cites~\cites{MR1083000}.}
      \field{issn}{0065-1036}
      \field{journaltitle}{Acta Arith.}
      \field{number}{3}
      \field{title}{The {$k$}-functions in multiplicative number theory. {II}. {U}niform distribution of zeta zeros}
      \field{volume}{56}
      \field{year}{1990}
      \true{nocite}
      \field{pages}{213\bibrangedash 224}
      \range{pages}{12}
      \verb{doi}
      \verb 10.4064/aa-56-3-213-224
      \endverb
      \verb{mrnumber}
      \verb 1083001
      \endverb
      \verb{urlraw}
      \verb https://doi.org/10.4064/aa-56-3-213-224
      \endverb
      \verb{url}
      \verb https://doi.org/10.4064/aa-56-3-213-224
      \endverb
    \endentry
    \entry{MR1105605}{article}{}
      \name{author}{1}{}{%
        {{hash=d1e4ad71dae945a6977ecf6f9ca65c30}{%
           family={Kaczorowski},
           familyi={K\bibinitperiod},
           given={J.},
           giveni={J\bibinitperiod}}}%
      }
      \strng{namehash}{d1e4ad71dae945a6977ecf6f9ca65c30}
      \strng{fullhash}{d1e4ad71dae945a6977ecf6f9ca65c30}
      \strng{bibnamehash}{d1e4ad71dae945a6977ecf6f9ca65c30}
      \strng{authorbibnamehash}{d1e4ad71dae945a6977ecf6f9ca65c30}
      \strng{authornamehash}{d1e4ad71dae945a6977ecf6f9ca65c30}
      \strng{authorfullhash}{d1e4ad71dae945a6977ecf6f9ca65c30}
      \field{extraname}{7}
      \field{sortinit}{K}
      \field{sortinithash}{d3edc18d54b9438a72c24c925bfb38f4}
      \field{labelprefix}{A}
      \field{labelnamesource}{author}
      \field{labeltitlesource}{title}
      \field{annotation}{Continuing the previous article in this series, the author defines an ``$\boldsymbol{A}$-discrepancy'' \begin{align*} D_n^*(x) = \sup_{0\leq t \leq 1} \bigg| \bigg( \sum_{\substack{k\ge1\\ \{x\gamma_k\} < t}} e^{-\gamma_k} \gamma_k^n \bigg) \bigg/ \bigg( \sum_{k=1}^\infty e^{-\gamma_k} \gamma_k^n \bigg) - t\bigg|. \end{align*} He shows that $D_n^*(x) \ll (\log \log n/\log n)^{2/3}$ for every real number $x \ne 0$. Under a certain $\boldsymbol{A}$-variant of zero-density theorems for Dirichlet $L$-functions, he proves that $D_n^*(x) \ll 1/\log n$ and conjectures that $D_n^*(x) \sim \alpha(x)/\log n$ for some constant~$\alpha(x)$.\par This article cites~\cites{MR1083001}.}
      \field{issn}{0065-1036}
      \field{journaltitle}{Acta Arith.}
      \field{number}{3}
      \field{title}{The {$k$}-functions in multiplicative number theory. {III}. {U}niform distribution of zeta zeros; discrepancy}
      \field{volume}{57}
      \field{year}{1991}
      \true{nocite}
      \field{pages}{199\bibrangedash 210}
      \range{pages}{12}
      \verb{doi}
      \verb 10.4064/aa-57-3-199-210
      \endverb
      \verb{mrnumber}
      \verb 1105605
      \endverb
      \verb{urlraw}
      \verb https://doi.org/10.4064/aa-57-3-199-210
      \endverb
      \verb{url}
      \verb https://doi.org/10.4064/aa-57-3-199-210
      \endverb
    \endentry
    \entry{MR1105608}{article}{}
      \name{author}{1}{}{%
        {{hash=d1e4ad71dae945a6977ecf6f9ca65c30}{%
           family={Kaczorowski},
           familyi={K\bibinitperiod},
           given={J.},
           giveni={J\bibinitperiod}}}%
      }
      \strng{namehash}{d1e4ad71dae945a6977ecf6f9ca65c30}
      \strng{fullhash}{d1e4ad71dae945a6977ecf6f9ca65c30}
      \strng{bibnamehash}{d1e4ad71dae945a6977ecf6f9ca65c30}
      \strng{authorbibnamehash}{d1e4ad71dae945a6977ecf6f9ca65c30}
      \strng{authornamehash}{d1e4ad71dae945a6977ecf6f9ca65c30}
      \strng{authorfullhash}{d1e4ad71dae945a6977ecf6f9ca65c30}
      \field{extraname}{8}
      \field{sortinit}{K}
      \field{sortinithash}{d3edc18d54b9438a72c24c925bfb38f4}
      \field{labelprefix}{A}
      \field{labelnamesource}{author}
      \field{labeltitlesource}{title}
      \field{annotation}{Using $k$-functions, the author proves an oscillation theorem for a specific class of almost periodic functions, generalizing results of Ingham. He then applies the theorem to prime counting functions: for instance, he shows that for every positive constant $A>0$, there exists a number $c>1$ such that for all $x>1$, we have $\sup E^\psi(x) \ge A$ and $\inf E^\psi(x) \le -A$ on the interval $(x, cx)$. The author observes that Ingham's refutation of the weak Mertens conjecture under RH and LI can be adapted to give the analogous results for $E^\pi(x;q,a)$ and $E^\pi(x;q,a,b)$; he reproves these results, assuming GRH but replacing LI by a condition on the intersection of the appropriate infinite subtorus and the diagonal subtorus. Finally, the author establishes an interesting connection between the range of $E^\psi(x;q,a)$ and the range of $x^{1/2}\Delta_r^\psi(x;q,a^{-1})$.\par This article cites~\cite{MR654520, MR0330069, MR0332684, MR0104638, 1936.Ingham, MR0006202, MR0352026, MR1083000, MR783538, MR0289452}.}
      \field{issn}{0065-1036}
      \field{journaltitle}{Acta Arith.}
      \field{number}{3}
      \field{title}{The {$k$}-functions in multiplicative number theory. {IV}. {O}n a method of {A}. {E}. {I}ngham}
      \field{volume}{57}
      \field{year}{1991}
      \true{nocite}
      \field{pages}{231\bibrangedash 244}
      \range{pages}{14}
      \verb{doi}
      \verb 10.4064/aa-57-3-231-244
      \endverb
      \verb{mrnumber}
      \verb 1105608
      \endverb
      \verb{urlraw}
      \verb https://doi.org/10.4064/aa-57-3-231-244
      \endverb
      \verb{url}
      \verb https://doi.org/10.4064/aa-57-3-231-244
      \endverb
    \endentry
    \entry{MR1133236}{article}{}
      \name{author}{1}{}{%
        {{hash=d1e4ad71dae945a6977ecf6f9ca65c30}{%
           family={Kaczorowski},
           familyi={K\bibinitperiod},
           given={J.},
           giveni={J\bibinitperiod}}}%
      }
      \strng{namehash}{d1e4ad71dae945a6977ecf6f9ca65c30}
      \strng{fullhash}{d1e4ad71dae945a6977ecf6f9ca65c30}
      \strng{bibnamehash}{d1e4ad71dae945a6977ecf6f9ca65c30}
      \strng{authorbibnamehash}{d1e4ad71dae945a6977ecf6f9ca65c30}
      \strng{authornamehash}{d1e4ad71dae945a6977ecf6f9ca65c30}
      \strng{authorfullhash}{d1e4ad71dae945a6977ecf6f9ca65c30}
      \field{extraname}{9}
      \field{sortinit}{K}
      \field{sortinithash}{d3edc18d54b9438a72c24c925bfb38f4}
      \field{labelprefix}{A}
      \field{labelnamesource}{author}
      \field{labeltitlesource}{title}
      \field{annotation}{Assuming GRH and HC, the author shows that \begin{align*} \liminf_{T \to \infty} \frac{W_{q, a}^{\psi} (T)}{\log T} \geq \frac{\gamma_0}{\pi}+ \kappa, \end{align*} where $\gamma$ is the lowest ``uncancelled'' zero of Dirichlet $L$-functions modulo~$q$---more precisely, the minimal $\gamma>0$ such that $\frac12+i \gamma$ is a pole of $\sum_{\chi \mod q} \overline\chi(a) \frac{L'}L(s, \chi)$---and where $\kappa=\kappa(q,a)$ is a nonnegative number, defined as the density of zeros of a certain linear combination of $K$-functions. When $q=a=1$, this result implies that \begin{align*} \liminf_{T \to \infty} \frac{W^{\psi} (T)}{\log T} \geq \frac{\gamma_1}{\pi}+10^{-250}, \end{align*} with $\gamma_1=14.135\ldots$; this result breaks the barrier $\frac{\gamma_1}{\pi}$ that had been achieved in prior results.\par This article cites~\cites{MR1083000, MR1105608, MR777013, MR791085, MR927376, MR858395, MR891057}.}
      \field{issn}{0065-1036}
      \field{journaltitle}{Acta Arith.}
      \field{number}{1}
      \field{title}{The {$k$}-functions in multiplicative number theory. {V}. {C}hanges of sign of some arithmetical error terms}
      \field{volume}{59}
      \field{year}{1991}
      \true{nocite}
      \field{pages}{37\bibrangedash 58}
      \range{pages}{22}
      \verb{doi}
      \verb 10.4064/aa-59-1-37-58
      \endverb
      \verb{mrnumber}
      \verb 1133236
      \endverb
      \verb{urlraw}
      \verb https://doi.org/10.4064/aa-59-1-37-58
      \endverb
      \verb{url}
      \verb https://doi.org/10.4064/aa-59-1-37-58
      \endverb
    \endentry
    \entry{MR1251926}{article}{}
      \name{author}{1}{}{%
        {{hash=d1e4ad71dae945a6977ecf6f9ca65c30}{%
           family={Kaczorowski},
           familyi={K\bibinitperiod},
           given={J.},
           giveni={J\bibinitperiod}}}%
      }
      \strng{namehash}{d1e4ad71dae945a6977ecf6f9ca65c30}
      \strng{fullhash}{d1e4ad71dae945a6977ecf6f9ca65c30}
      \strng{bibnamehash}{d1e4ad71dae945a6977ecf6f9ca65c30}
      \strng{authorbibnamehash}{d1e4ad71dae945a6977ecf6f9ca65c30}
      \strng{authornamehash}{d1e4ad71dae945a6977ecf6f9ca65c30}
      \strng{authorfullhash}{d1e4ad71dae945a6977ecf6f9ca65c30}
      \field{extraname}{10}
      \field{sortinit}{K}
      \field{sortinithash}{d3edc18d54b9438a72c24c925bfb38f4}
      \field{labelprefix}{A}
      \field{labelnamesource}{author}
      \field{labeltitlesource}{title}
      \field{annotation}{The author shows, assuming GRH, that there are infinitely many integers $m$ such that $\pi(m;q,1)$ is larger than any other $\pi(m;q,a)$, as well as infinitely many integers such that $\pi(m;q,1)$ is the smallest. Indeed, the sets of such integers have positive lower density, even if a stronger inequality is demanded: the author proves that for any positive real number~$u$, there exist constants $c>1$ and $b>0$ such that for every $T\geq 1$, \begin{align*} \#\big\{T\leq m\leq c T\colon E^\pi(m;q,1) - \max_{a\not\equiv 1\mod q} E^\pi(m;q,a) \ge u \big\}\geq b T, \end{align*} and similarly for $E^\psi$; and the same result holds with${}\ge u$ replaced by${}\le-u$. The author uses his $k$-functions that appeared in prior work, as well as an examination of the boundary values of Dirichlet series.\par The author formulates the following ``Strong Race Hypothesis'': for every permutation of the set $\{a_1,a_2,\ldots,a_{\phi(q)}\}$ of reduced residue classes modulo~$q$, the set of integers $m$ such that \[ \pi(m;q,a_1)<\pi(m;q,a_2)<\cdots<\pi(m;q,a_{\phi(q)}) \] has positive lower density. \par This article cites~\cites{MR0108470, MR0146156, MR1083000, MR1105608}.}
      \field{issn}{0033-5606}
      \field{journaltitle}{Quart. J. Math. Oxford Ser. (2)}
      \field{number}{176}
      \field{title}{A contribution to the {S}hanks-{R}ényi race problem}
      \field{volume}{44}
      \field{year}{1993}
      \true{nocite}
      \field{pages}{451\bibrangedash 458}
      \range{pages}{8}
      \verb{doi}
      \verb 10.1093/qmath/44.4.451
      \endverb
      \verb{mrnumber}
      \verb 1251926
      \endverb
      \verb{urlraw}
      \verb https://doi.org/10.1093/qmath/44.4.451
      \endverb
      \verb{url}
      \verb https://doi.org/10.1093/qmath/44.4.451
      \endverb
    \endentry
    \entry{MR1367576}{article}{}
      \name{author}{1}{}{%
        {{hash=d1e4ad71dae945a6977ecf6f9ca65c30}{%
           family={Kaczorowski},
           familyi={K\bibinitperiod},
           given={J.},
           giveni={J\bibinitperiod}}}%
      }
      \strng{namehash}{d1e4ad71dae945a6977ecf6f9ca65c30}
      \strng{fullhash}{d1e4ad71dae945a6977ecf6f9ca65c30}
      \strng{bibnamehash}{d1e4ad71dae945a6977ecf6f9ca65c30}
      \strng{authorbibnamehash}{d1e4ad71dae945a6977ecf6f9ca65c30}
      \strng{authornamehash}{d1e4ad71dae945a6977ecf6f9ca65c30}
      \strng{authorfullhash}{d1e4ad71dae945a6977ecf6f9ca65c30}
      \field{extraname}{11}
      \field{sortinit}{K}
      \field{sortinithash}{d3edc18d54b9438a72c24c925bfb38f4}
      \field{labelprefix}{A}
      \field{labelnamesource}{author}
      \field{labeltitlesource}{title}
      \field{annotation}{The author's earlier result~\cite{MR1251926} implies that there exists a permutation $\{\sigma_j\}$ of the reduced residue classes\mod q that begins with $1$ such that the set of $x$ for which the $\pi(x;q,\sigma_j)$ are in the given order has positive lower density. In this article, he provides a method for computing explicit permutations with this property. As an application, he gives permutations for each prime modulus $\le 29$ that satisfy these conditions and therefore provably occur with positive lower density; for example, modulo $13$ he provides the permutation $(1,7,8,9,2,6,12,10,11,5,3,4)$ with this property. Similar results apply to permutations with $1$ in last place, and to the functions $\psi(x;q,\sigma_j)$.\par This article cites~\cites{MR0146156, MR1083000, MR1105608,MR1251926, MR1256149, MR1310738}.}
      \field{issn}{0065-1036}
      \field{journaltitle}{Acta Arith.}
      \field{number}{1}
      \field{title}{On the {S}hanks-{R}ényi race problem}
      \field{volume}{74}
      \field{year}{1996}
      \true{nocite}
      \field{pages}{31\bibrangedash 46}
      \range{pages}{16}
      \verb{doi}
      \verb 10.4064/aa-74-1-31-46
      \endverb
      \verb{mrnumber}
      \verb 1367576
      \endverb
      \verb{urlraw}
      \verb https://doi.org/10.4064/aa-74-1-31-46
      \endverb
      \verb{url}
      \verb https://doi.org/10.4064/aa-74-1-31-46
      \endverb
    \endentry
    \entry{MR858395}{article}{}
      \name{author}{2}{}{%
        {{hash=d1e4ad71dae945a6977ecf6f9ca65c30}{%
           family={Kaczorowski},
           familyi={K\bibinitperiod},
           given={J.},
           giveni={J\bibinitperiod}}}%
        {{hash=c6d965df16e600c701d4bcc2888ae34c}{%
           family={Pintz},
           familyi={P\bibinitperiod},
           given={J.},
           giveni={J\bibinitperiod}}}%
      }
      \strng{namehash}{2529f5112eea5e5d50df838061479c33}
      \strng{fullhash}{2529f5112eea5e5d50df838061479c33}
      \strng{bibnamehash}{2529f5112eea5e5d50df838061479c33}
      \strng{authorbibnamehash}{2529f5112eea5e5d50df838061479c33}
      \strng{authornamehash}{2529f5112eea5e5d50df838061479c33}
      \strng{authorfullhash}{2529f5112eea5e5d50df838061479c33}
      \field{extraname}{1}
      \field{sortinit}{K}
      \field{sortinithash}{d3edc18d54b9438a72c24c925bfb38f4}
      \field{labelprefix}{A}
      \field{labelnamesource}{author}
      \field{labeltitlesource}{title}
      \field{annotation}{In this article, the authors improve upon and extend results of Landau~\cite{MR1511360}, Pólya~\cite{1930.Polya}, and Grosswald~\cite{MR0198145}. Given a Dirichlet intergal $F(s) = \int_{x_0}^\infty f(x)x^{-s-1}\,dx$ converging on a right half-plane $\{\sigma>\theta\}$, with a continuation to a larger half-plane except for perhaps countably many poles or logarithmic singularities (in a precise sense), the authors show that $\liminf_{T\to\infty} \frac{W(f,T)}{\log T}\ge \frac\gamma\pi$ where $\gamma = \inf\{|t| \colon F(s)\text{ is not regular at }\theta+it\}$. This result implies $\gg \log T$ sign changes of the functions $M(x)$ and $\Delta^{Q_k}(x)$, as well as of $\psi(x;q,\ell_1,\ell_2)$ assuming~HC. For a slightly more restricted class of functions, the authors prove an effective version of a similar result, again guaranteeing $\gg \log T$ sign changes (essentially using a single singularity of $F(s)$) but now with effective constants.\par This article cites~\cites{MR1511360, 1930.Polya, MR0439774, MR0198145, MR0202685, MR0146160, MR0146161, MR688630, MR0218318, MR777013, MR791085}.}
      \field{issn}{0236-5294}
      \field{journaltitle}{Acta Math. Hungar.}
      \field{number}{1-2}
      \field{title}{Oscillatory properties of arithmetical functions. {I}}
      \field{volume}{48}
      \field{year}{1986}
      \true{nocite}
      \field{pages}{173\bibrangedash 185}
      \range{pages}{13}
      \verb{doi}
      \verb 10.1007/BF01949062
      \endverb
      \verb{mrnumber}
      \verb 858395
      \endverb
      \verb{urlraw}
      \verb https://doi.org/10.1007/BF01949062
      \endverb
      \verb{url}
      \verb https://doi.org/10.1007/BF01949062
      \endverb
    \endentry
    \entry{MR891057}{article}{}
      \name{author}{2}{}{%
        {{hash=d1e4ad71dae945a6977ecf6f9ca65c30}{%
           family={Kaczorowski},
           familyi={K\bibinitperiod},
           given={J.},
           giveni={J\bibinitperiod}}}%
        {{hash=c6d965df16e600c701d4bcc2888ae34c}{%
           family={Pintz},
           familyi={P\bibinitperiod},
           given={J.},
           giveni={J\bibinitperiod}}}%
      }
      \strng{namehash}{2529f5112eea5e5d50df838061479c33}
      \strng{fullhash}{2529f5112eea5e5d50df838061479c33}
      \strng{bibnamehash}{2529f5112eea5e5d50df838061479c33}
      \strng{authorbibnamehash}{2529f5112eea5e5d50df838061479c33}
      \strng{authornamehash}{2529f5112eea5e5d50df838061479c33}
      \strng{authorfullhash}{2529f5112eea5e5d50df838061479c33}
      \field{extraname}{2}
      \field{sortinit}{K}
      \field{sortinithash}{d3edc18d54b9438a72c24c925bfb38f4}
      \field{labelprefix}{A}
      \field{labelnamesource}{author}
      \field{labeltitlesource}{title}
      \field{annotation}{The authors extend their earlier results to obtain $\gg \log T$ sign changes for functions such as $\Delta^\Pi(x;q,a)$, $\Delta^\pi(x;q,a)$ where $a$ is a quadratic nonresidue, $\Delta^\Pi(x;q,a,b)$ where $a\not\equiv b\mod q$, and so on. They also similarly obtain sign changes (in relatively short intervals) for the error term in the asymptotic formula for the counting function of irreducible elements in the ring of integers $\mathcal{O}_K$ of a number field~$K$, assuming the Dedekind zeta function of the Hilbert class field of $K$ does not vanish on the interval $[1/2,1)$ and has at least one simple zero in the half-plane $\sigma>1/2$.\par This article cites~\cites{1936.Ingham, MR777013, MR791085, MR858395}.}
      \field{issn}{0236-5294}
      \field{journaltitle}{Acta Math. Hungar.}
      \field{number}{3-4}
      \field{title}{Oscillatory properties of arithmetical functions. {II}}
      \field{volume}{49}
      \field{year}{1987}
      \true{nocite}
      \field{pages}{441\bibrangedash 453}
      \range{pages}{13}
      \verb{doi}
      \verb 10.1007/BF01951008
      \endverb
      \verb{mrnumber}
      \verb 891057
      \endverb
      \verb{urlraw}
      \verb https://doi.org/10.1007/BF01951008
      \endverb
      \verb{url}
      \verb https://doi.org/10.1007/BF01951008
      \endverb
    \endentry
    \entry{MR0980524}{article}{}
      \name{author}{2}{}{%
        {{hash=d1e4ad71dae945a6977ecf6f9ca65c30}{%
           family={Kaczorowski},
           familyi={K\bibinitperiod},
           given={J.},
           giveni={J\bibinitperiod}}}%
        {{hash=36befda65f61ea63f097e5518d3ec577}{%
           family={Staś},
           familyi={S\bibinitperiod},
           given={W.},
           giveni={W\bibinitperiod}}}%
      }
      \strng{namehash}{d5a8f4983f5607d22545d2b33d58100f}
      \strng{fullhash}{d5a8f4983f5607d22545d2b33d58100f}
      \strng{bibnamehash}{d5a8f4983f5607d22545d2b33d58100f}
      \strng{authorbibnamehash}{d5a8f4983f5607d22545d2b33d58100f}
      \strng{authornamehash}{d5a8f4983f5607d22545d2b33d58100f}
      \strng{authorfullhash}{d5a8f4983f5607d22545d2b33d58100f}
      \field{extraname}{1}
      \field{sortinit}{K}
      \field{sortinithash}{d3edc18d54b9438a72c24c925bfb38f4}
      \field{labelprefix}{A}
      \field{labelnamesource}{author}
      \field{labeltitlesource}{title}
      \field{annotation}{The author shows that if $K$ is a number field such that HC holds for $\zeta_K(s)$, then for any $0<\ep <1$ one has $W^{\psi_K}(T) \geq (1-\ep) \frac{\gamma_K}{\pi} \log T$ when $T$ is sufficiently large in terms of $K$ and $\ep$. The assumption of no real zeros can be removed if the contribution from the real zeros is subtracted from $E^{\psi_K}(x)$. \par This article cites~\cite{1936.Ingham,MR777013,MR817334,MR0132728}.}
      \field{issn}{0010-1354}
      \field{journaltitle}{Colloq. Math.}
      \field{number}{1}
      \field{title}{On the number of sign changes in the remainder-term of the prime-ideal theorem}
      \field{volume}{56}
      \field{year}{1988}
      \true{nocite}
      \field{pages}{185\bibrangedash 197}
      \range{pages}{13}
      \verb{doi}
      \verb 10.4064/cm-56-1-185-197
      \endverb
      \verb{mrnumber}
      \verb 980524
      \endverb
      \verb{urlraw}
      \verb https://doi.org/10.4064/cm-56-1-185-197
      \endverb
      \verb{url}
      \verb https://doi.org/10.4064/cm-56-1-185-197
      \endverb
    \endentry
    \entry{MR1256149}{article}{}
      \name{author}{1}{}{%
        {{hash=6fa84c8cc1e455847a05b9946a1bd205}{%
           family={Kaczorowski},
           familyi={K\bibinitperiod},
           given={Jerzy},
           giveni={J\bibinitperiod}}}%
      }
      \strng{namehash}{6fa84c8cc1e455847a05b9946a1bd205}
      \strng{fullhash}{6fa84c8cc1e455847a05b9946a1bd205}
      \strng{bibnamehash}{6fa84c8cc1e455847a05b9946a1bd205}
      \strng{authorbibnamehash}{6fa84c8cc1e455847a05b9946a1bd205}
      \strng{authornamehash}{6fa84c8cc1e455847a05b9946a1bd205}
      \strng{authorfullhash}{6fa84c8cc1e455847a05b9946a1bd205}
      \field{extraname}{12}
      \field{sortinit}{K}
      \field{sortinithash}{d3edc18d54b9438a72c24c925bfb38f4}
      \field{labelprefix}{A}
      \field{labelnamesource}{author}
      \field{labeltitlesource}{title}
      \field{annotation}{In this article, the author uses $k$-functions to prove several results. Assuming RH, he shows that the set $\{ m\in\N\colon E^\pi(x) > u \}$ has positive lower density for any $u\in\R$, and similarly with $E^\pi(x)$ replaced by $-E^\pi(x)$, $E^\psi(x)$, $-E^\psi(x)$; assuming GRH, $E^\pi(x)$ may be further replaced by $E^\pi(x;q,1,\max)$ and $-E^\pi(x;q,1,\min)$ and their counterparts with~$\psi$. Unconditionally he gives analogous weaker density statements such as $\# \{ m\le M\colon E^\pi(x) > u \} = \Omega(M^{1-\ep})$. All these results are ineffective, but the author provides effective counterparts such as $\liminf_{N\to\infty} \sum_{k=1}^{N} 2^{-k} \# \{2^k \leq m < 2^{k+1}\colon E^\pi(x) > u \} > 0$.\par This article cites~\cite{MR0476616, MR0476615, 1936.Ingham, MR1083000, MR1105608, MR1251926, MR1310738, MR1256149, MR0146156, MR1511360, 1914.Littlewood, MR0087685, MR0108470}.}
      \field{issn}{0075-4102}
      \field{journaltitle}{J. Reine Angew. Math.}
      \field{title}{Results on the distribution of primes}
      \field{volume}{446}
      \field{year}{1994}
      \true{nocite}
      \field{pages}{89\bibrangedash 113}
      \range{pages}{25}
      \verb{doi}
      \verb 10.1515/crll.1994.446.89
      \endverb
      \verb{mrnumber}
      \verb 1256149
      \endverb
    \endentry
    \entry{MR1344249}{article}{}
      \name{author}{1}{}{%
        {{hash=6fa84c8cc1e455847a05b9946a1bd205}{%
           family={Kaczorowski},
           familyi={K\bibinitperiod},
           given={Jerzy},
           giveni={J\bibinitperiod}}}%
      }
      \strng{namehash}{6fa84c8cc1e455847a05b9946a1bd205}
      \strng{fullhash}{6fa84c8cc1e455847a05b9946a1bd205}
      \strng{bibnamehash}{6fa84c8cc1e455847a05b9946a1bd205}
      \strng{authorbibnamehash}{6fa84c8cc1e455847a05b9946a1bd205}
      \strng{authornamehash}{6fa84c8cc1e455847a05b9946a1bd205}
      \strng{authorfullhash}{6fa84c8cc1e455847a05b9946a1bd205}
      \field{extraname}{13}
      \field{sortinit}{K}
      \field{sortinithash}{d3edc18d54b9438a72c24c925bfb38f4}
      \field{labelprefix}{A}
      \field{labelnamesource}{author}
      \field{labeltitlesource}{title}
      \field{annotation}{The author examines upper and lower natural densities in the prime number race modulo~$4$. Assuming GRH for $L(s, \chi_{-4})$, he proves that $\mathfrak d^*(4;1,3) \geq 0.04054045$ and $\mathfrak d^*(4;3,1) \geq 0.99998936$, and that $\mathfrak d_*(4; 1,3) <0.0000106$ and $\mathfrak d_*(4; 3,1) < 0.9594595$.\par This article cites~\cite{MR0476615, 1936.Ingham, MR1083000, MR1105608, MR1251926, MR1310738, MR1256149, MR0146156, MR1511360, 1914.Littlewood, MR0087685, MR0108470}.}
      \field{issn}{0174-4747}
      \field{journaltitle}{Analysis}
      \field{number}{2}
      \field{title}{On the distribution of primes (mod $4$)}
      \field{volume}{15}
      \field{year}{1995}
      \true{nocite}
      \field{pages}{159\bibrangedash 171}
      \range{pages}{13}
      \verb{doi}
      \verb 10.1524/anly.1995.15.2.159
      \endverb
      \verb{mrnumber}
      \verb 1344249
      \endverb
    \endentry
    \entry{MR1310738}{article}{}
      \name{author}{1}{}{%
        {{hash=6fa84c8cc1e455847a05b9946a1bd205}{%
           family={Kaczorowski},
           familyi={K\bibinitperiod},
           given={Jerzy},
           giveni={J\bibinitperiod}}}%
      }
      \strng{namehash}{6fa84c8cc1e455847a05b9946a1bd205}
      \strng{fullhash}{6fa84c8cc1e455847a05b9946a1bd205}
      \strng{bibnamehash}{6fa84c8cc1e455847a05b9946a1bd205}
      \strng{authorbibnamehash}{6fa84c8cc1e455847a05b9946a1bd205}
      \strng{authornamehash}{6fa84c8cc1e455847a05b9946a1bd205}
      \strng{authorfullhash}{6fa84c8cc1e455847a05b9946a1bd205}
      \field{extraname}{14}
      \field{sortinit}{K}
      \field{sortinithash}{d3edc18d54b9438a72c24c925bfb38f4}
      \field{labelprefix}{A}
      \field{labelnamesource}{author}
      \field{labeltitlesource}{title}
      \field{annotation}{Assuming GRH, the author shows that for any permutation $(a_1,a_2,a_3,a_4)$ of $(1,2,3,4)$, there exist constants $b>0$ and $c_0>1$ such that \begin{align*} \#\big\{T\le x\le c_0T \colon \psi(x;5,a_1)>\psi(x;5,a_2)>&\psi(x;5,a_3)>\psi(x;5,a_4), \\&\min_{1\le j\le 3}\big( \psi(x;5,a_j,a_{j+1}) \big) \geq b\sqrt x \big\} \gg T. \end{align*} The proof is another application of the author's theory of $k$-functions and involves explicit calculations using the Dirichlet $L$-functions\mod 5 and exponential sums corresponding to each permutation.\par This article cites~\cites{MR1083000, MR1251926, MR0146156}.}
      \field{issn}{0022-314X}
      \field{journaltitle}{J. Number Theory}
      \field{number}{1}
      \field{title}{On the {S}hanks-{R}ényi race problem mod {$5$}}
      \field{volume}{50}
      \field{year}{1995}
      \true{nocite}
      \field{pages}{106\bibrangedash 118}
      \range{pages}{13}
      \verb{doi}
      \verb 10.1006/jnth.1995.1006
      \endverb
      \verb{mrnumber}
      \verb 1310738
      \endverb
      \verb{urlraw}
      \verb https://doi.org/10.1006/jnth.1995.1006
      \endverb
      \verb{url}
      \verb https://doi.org/10.1006/jnth.1995.1006
      \endverb
    \endentry
    \entry{MR1645811}{incollection}{}
      \name{author}{1}{}{%
        {{hash=6fa84c8cc1e455847a05b9946a1bd205}{%
           family={Kaczorowski},
           familyi={K\bibinitperiod},
           given={Jerzy},
           giveni={J\bibinitperiod}}}%
      }
      \list{publisher}{1}{%
        {Birkhäuser, Basel}%
      }
      \strng{namehash}{6fa84c8cc1e455847a05b9946a1bd205}
      \strng{fullhash}{6fa84c8cc1e455847a05b9946a1bd205}
      \strng{bibnamehash}{6fa84c8cc1e455847a05b9946a1bd205}
      \strng{authorbibnamehash}{6fa84c8cc1e455847a05b9946a1bd205}
      \strng{authornamehash}{6fa84c8cc1e455847a05b9946a1bd205}
      \strng{authorfullhash}{6fa84c8cc1e455847a05b9946a1bd205}
      \field{extraname}{15}
      \field{sortinit}{K}
      \field{sortinithash}{d3edc18d54b9438a72c24c925bfb38f4}
      \field{labelprefix}{A}
      \field{labelnamesource}{author}
      \field{labeltitlesource}{title}
      \field{annotation}{From the Math Review by M.~Jutila: ``Various aspects and problems as well as the history of comparative prime number theory are surveyed, including recent important work by the author [using $K$-functions]. The central topic is the comparative study of the frequency of primes in different arithmetic progressions. The article ends with a list of open problems and an extensive bibliography with 30 references.''\par This article cites~\cites{MR0476615,MR1159354,MR0506080,MR777013,MR791085,MR927376,MR945273,MR1083000,MR1105608,MR1133236,MR1251926,MR1310738,MR1256149,MR1344249,MR858395,MR891057,MR0146156,MR0162771,MR1544293A,MR1544293B,MR783538,MR634679,MR1329368,MR0137689,MR0004255}.}
      \field{bookpagination}{none}
      \field{booktitle}{European {C}ongress of {M}athematics, {V}ol. {I} ({B}udapest, 1996)}
      \field{series}{Progr. Math.}
      \field{title}{Boundary values of {D}irichlet series and the distribution of primes}
      \field{volume}{168}
      \field{year}{1998}
      \true{nocite}
      \field{pages}{237\bibrangedash 254}
      \range{pages}{18}
      \verb{mrnumber}
      \verb 1645811
      \endverb
    \endentry
    \entry{MR2340242}{article}{}
      \name{author}{1}{}{%
        {{hash=6fa84c8cc1e455847a05b9946a1bd205}{%
           family={Kaczorowski},
           familyi={K\bibinitperiod},
           given={Jerzy},
           giveni={J\bibinitperiod}}}%
      }
      \strng{namehash}{6fa84c8cc1e455847a05b9946a1bd205}
      \strng{fullhash}{6fa84c8cc1e455847a05b9946a1bd205}
      \strng{bibnamehash}{6fa84c8cc1e455847a05b9946a1bd205}
      \strng{authorbibnamehash}{6fa84c8cc1e455847a05b9946a1bd205}
      \strng{authornamehash}{6fa84c8cc1e455847a05b9946a1bd205}
      \strng{authorfullhash}{6fa84c8cc1e455847a05b9946a1bd205}
      \field{extraname}{16}
      \field{sortinit}{K}
      \field{sortinithash}{d3edc18d54b9438a72c24c925bfb38f4}
      \field{labelprefix}{A}
      \field{labelnamesource}{author}
      \field{labeltitlesource}{title}
      \field{annotation}{The author shows that \begin{align*} \sum_{n=1}^\infty \mu(n) \left(\cos \left(\frac{x}{n}\right) - 1\right) = \Omega_{\pm} (x^{1/2} \log \log \log x), \end{align*} which comes tantalizingly close to the assertion that $E^M(x) = \Omega_\pm(\log\log\log x)$. The proof uses an explicit formula for a relative of $k(z)$, namely \begin{align*} m(z) = \frac{1}{2\pi i} \int_{\mathcal C} \frac{e^{zs}}{\zeta(s)} \, ds \end{align*} where $\mathcal C$ is the boundary of a half strip enclosing the part of the critical strip lying in the upper half plane. \par This article cites~\cite{MR1109990, MR0006202, MR1105608, MR2329690, MR2118272, MR2078705, MR783538, MR891440, MR0882550}.}
      \field{issn}{0024-6107}
      \field{journaltitle}{J. Lond. Math. Soc. (2)}
      \field{number}{2}
      \field{title}{Results on the {M}öbius function}
      \field{volume}{75}
      \field{year}{2007}
      \true{nocite}
      \field{pages}{509\bibrangedash 521}
      \range{pages}{13}
      \verb{doi}
      \verb 10.1112/jlms/jdm006
      \endverb
      \verb{mrnumber}
      \verb 2340242
      \endverb
      \verb{urlraw}
      \verb https://doi.org/10.1112/jlms/jdm006
      \endverb
      \verb{url}
      \verb https://doi.org/10.1112/jlms/jdm006
      \endverb
    \endentry
    \entry{MR2470105}{article}{}
      \name{author}{1}{}{%
        {{hash=6fa84c8cc1e455847a05b9946a1bd205}{%
           family={Kaczorowski},
           familyi={K\bibinitperiod},
           given={Jerzy},
           giveni={J\bibinitperiod}}}%
      }
      \strng{namehash}{6fa84c8cc1e455847a05b9946a1bd205}
      \strng{fullhash}{6fa84c8cc1e455847a05b9946a1bd205}
      \strng{bibnamehash}{6fa84c8cc1e455847a05b9946a1bd205}
      \strng{authorbibnamehash}{6fa84c8cc1e455847a05b9946a1bd205}
      \strng{authornamehash}{6fa84c8cc1e455847a05b9946a1bd205}
      \strng{authorfullhash}{6fa84c8cc1e455847a05b9946a1bd205}
      \field{extraname}{17}
      \field{sortinit}{K}
      \field{sortinithash}{d3edc18d54b9438a72c24c925bfb38f4}
      \field{labelprefix}{A}
      \field{labelnamesource}{author}
      \field{labeltitlesource}{title}
      \field{annotation}{Given a number field~$K$, the author studies the oscillations of the error term $E_K(x)$ for the counting function of irreducible elements of~$K$ (up to units). Let $\zeta_{K_H}(s)$ denote the Dedekind zeta function of the Hilbert class field of~$K$, and let $\rho$ denote a nontrivial zero of $\zeta_{K_H}$. The author introduces a complicated integer-valued quantity $m^*(\rho,K)$ related to the multiplicity of $\rho$; under the assumption that some $m^*(\rho,K)$ is nonzero, he proves that $E_K(x) = \Omega_\pm\big({\sqrt{x}(\log\log x)^{D(K)-1}/\log{x}}\big)$, where $D(K)$ is the Davenport constant of the class group of~$K$. As a result, when $K$ with class number~$2$, we unconditionally have $E_k(x) = \Omega_\pm\big({\sqrt{x}\log\log{x}/\log{x}}\big)$; when the class number is an odd prime~$p$, we have $E_k(x)=\Omega_\pm({\sqrt{x}(\log\log{x})^{p-1}/\log{x}})$ provided that $\zeta_{K_H}(s)$ has at least one nontrivial zero of multiplicity not divisible by~$p$. All of these oscillations have logarithmic frequency. Assuming RH for $\zeta_{K_H}(s)$, the author sketches a proof of $H_K(x)\ll \sqrt x(\log{x})^{D(K)+1}$.\par This article cites~\cites{MR791085,MR891057,MR1957110,MR2329690,MR2506435,1914.Littlewood,MR0882550}.}
      \field{issn}{0026-9255}
      \field{journaltitle}{Monatsh. Math.}
      \field{number}{1}
      \field{title}{On the distribution of irreducible algebraic integers}
      \field{volume}{156}
      \field{year}{2009}
      \true{nocite}
      \field{pages}{47\bibrangedash 71}
      \range{pages}{25}
      \verb{doi}
      \verb 10.1007/s00605-008-0559-8
      \endverb
      \verb{mrnumber}
      \verb 2470105
      \endverb
    \endentry
    \entry{MR2730494}{article}{}
      \name{author}{1}{}{%
        {{hash=6fa84c8cc1e455847a05b9946a1bd205}{%
           family={Kaczorowski},
           familyi={K\bibinitperiod},
           given={Jerzy},
           giveni={J\bibinitperiod}}}%
      }
      \strng{namehash}{6fa84c8cc1e455847a05b9946a1bd205}
      \strng{fullhash}{6fa84c8cc1e455847a05b9946a1bd205}
      \strng{bibnamehash}{6fa84c8cc1e455847a05b9946a1bd205}
      \strng{authorbibnamehash}{6fa84c8cc1e455847a05b9946a1bd205}
      \strng{authornamehash}{6fa84c8cc1e455847a05b9946a1bd205}
      \strng{authorfullhash}{6fa84c8cc1e455847a05b9946a1bd205}
      \field{extraname}{18}
      \field{sortinit}{K}
      \field{sortinithash}{d3edc18d54b9438a72c24c925bfb38f4}
      \field{labelprefix}{A}
      \field{labelnamesource}{author}
      \field{labeltitlesource}{title}
      \field{annotation}{The author continues his examination of~$E_K(x)$ from~\cite{MR2470105}, under the assumption that no entire function from the Selberg class vanishes at $s=1$ (which follows from Selberg's orthonormality conjecture). He shows that this assumption implies that at least one of the quantities $m^*(\rho,K)$ from his previous work is nonzero, and therefore that this nonvanishing assumption implies that $E_K(x) = \Omega_\pm\big({\sqrt{x}(\log\log x)^{D(K)-1}/\log{x}}\big)$, where $D(K)$ is the Davenport constant of the class group of~$K$. He also shows unconditionally that $E_K(x)=\Omega_\pm({\sqrt{x}/(\log{x})^{O_K(1)}})$. All these oscillations have logarithmic frequency.\par This article cites~\cites{MR777013,MR791085,MR2470105,MR891057,MR2506435}.}
      \field{issn}{0025-584X}
      \field{journaltitle}{Math. Nachr.}
      \field{number}{9}
      \field{title}{{$\Omega$}-estimates related to irreducible algebraic integers}
      \field{volume}{283}
      \field{year}{2010}
      \true{nocite}
      \field{pages}{1291\bibrangedash 1303}
      \range{pages}{13}
      \verb{doi}
      \verb 10.1002/mana.200710158
      \endverb
      \verb{mrnumber}
      \verb 2730494
      \endverb
    \endentry
    \entry{MR1957110}{article}{}
      \name{author}{2}{}{%
        {{hash=6fa84c8cc1e455847a05b9946a1bd205}{%
           family={Kaczorowski},
           familyi={K\bibinitperiod},
           given={Jerzy},
           giveni={J\bibinitperiod}}}%
        {{hash=d6ae4b9844776294c92ff23952e91b53}{%
           family={Ramaré},
           familyi={R\bibinitperiod},
           given={Olivier},
           giveni={O\bibinitperiod}}}%
      }
      \strng{namehash}{ea4952aa08d573beb06051bb8ebb2f86}
      \strng{fullhash}{ea4952aa08d573beb06051bb8ebb2f86}
      \strng{bibnamehash}{ea4952aa08d573beb06051bb8ebb2f86}
      \strng{authorbibnamehash}{ea4952aa08d573beb06051bb8ebb2f86}
      \strng{authornamehash}{ea4952aa08d573beb06051bb8ebb2f86}
      \strng{authorfullhash}{ea4952aa08d573beb06051bb8ebb2f86}
      \field{sortinit}{K}
      \field{sortinithash}{d3edc18d54b9438a72c24c925bfb38f4}
      \field{labelprefix}{A}
      \field{labelnamesource}{author}
      \field{labeltitlesource}{title}
      \field{annotation}{This article investigates, assuming GRH, the distribution of values of a large class of functions of arithmetic significance (related to the Selberg class) using boundary values of $k$-functions. They establish an explicit formula for the appropriate functions as well as almost periodicity in the $L^2$ sense of Stepanov and the existence of a limiting logarithmic density.\par This article cites~\cites{MR1083000,MR1105608,MR1344249,MR1645811,MR1329368}.}
      \field{issn}{0065-1036}
      \field{journaltitle}{Acta Arith.}
      \field{number}{3}
      \field{title}{Almost periodicity of some error terms in prime number theory}
      \field{volume}{106}
      \field{year}{2003}
      \true{nocite}
      \field{pages}{277\bibrangedash 297}
      \range{pages}{21}
      \verb{doi}
      \verb 10.4064/aa106-3-6
      \endverb
      \verb{mrnumber}
      \verb 1957110
      \endverb
    \endentry
    \entry{MR1042465}{article}{}
      \name{author}{2}{}{%
        {{hash=6fa84c8cc1e455847a05b9946a1bd205}{%
           family={Kaczorowski},
           familyi={K\bibinitperiod},
           given={Jerzy},
           giveni={J\bibinitperiod}}}%
        {{hash=b39a346ae2b44071e01b12c1003e1246}{%
           family={Staś},
           familyi={S\bibinitperiod},
           given={Włodzimierz},
           giveni={W\bibinitperiod}}}%
      }
      \strng{namehash}{e170c0eabd81bb1165f3c0582dd5b3a2}
      \strng{fullhash}{e170c0eabd81bb1165f3c0582dd5b3a2}
      \strng{bibnamehash}{e170c0eabd81bb1165f3c0582dd5b3a2}
      \strng{authorbibnamehash}{e170c0eabd81bb1165f3c0582dd5b3a2}
      \strng{authornamehash}{e170c0eabd81bb1165f3c0582dd5b3a2}
      \strng{authorfullhash}{e170c0eabd81bb1165f3c0582dd5b3a2}
      \field{extraname}{2}
      \field{sortinit}{K}
      \field{sortinithash}{d3edc18d54b9438a72c24c925bfb38f4}
      \field{labelprefix}{A}
      \field{labelnamesource}{author}
      \field{labeltitlesource}{title}
      \field{annotation}{This article is identical to~\cite{MR0980524}.}
      \field{issn}{0137-9747}
      \field{journaltitle}{Discuss. Math.}
      \field{title}{On the number of sign-changes in the remainder-term of the prime-ideal theorem}
      \field{volume}{9}
      \field{year}{1988}
      \true{nocite}
      \field{pages}{83\bibrangedash 102 (1989)}
      \range{pages}{-1}
      \verb{mrnumber}
      \verb 1042465
      \endverb
    \endentry
    \entry{MR2329690}{article}{}
      \name{author}{2}{}{%
        {{hash=6fa84c8cc1e455847a05b9946a1bd205}{%
           family={Kaczorowski},
           familyi={K\bibinitperiod},
           given={Jerzy},
           giveni={J\bibinitperiod}}}%
        {{hash=87807ce1c4509fc3bd17938ccd4fe3e6}{%
           family={Wiertelak},
           familyi={W\bibinitperiod},
           given={Kazimierz},
           giveni={K\bibinitperiod}}}%
      }
      \strng{namehash}{5e2b3cb8731c2c2fbd035cbb54743005}
      \strng{fullhash}{5e2b3cb8731c2c2fbd035cbb54743005}
      \strng{bibnamehash}{5e2b3cb8731c2c2fbd035cbb54743005}
      \strng{authorbibnamehash}{5e2b3cb8731c2c2fbd035cbb54743005}
      \strng{authornamehash}{5e2b3cb8731c2c2fbd035cbb54743005}
      \strng{authorfullhash}{5e2b3cb8731c2c2fbd035cbb54743005}
      \field{extraname}{1}
      \field{sortinit}{K}
      \field{sortinithash}{d3edc18d54b9438a72c24c925bfb38f4}
      \field{labelprefix}{A}
      \field{labelnamesource}{author}
      \field{labeltitlesource}{title}
      \field{annotation}{The authors consider a general class of functions of the form $F(z)=\sum_{n=1}^\infty a_ne^{i\omega_nz}$ on the upper-half plane, whose boundary values $P(x) = \lim_{y\to x^+} \Re F(x+iy)$ exist for sufficiently large $x$. If certain growth conditions are satisfied, they establish oscillation results for $P(x)$. Applying their result to $K(z)e^{-z/2}$ recovers Littlewood's theorem $\Delta^\psi(x)=\Omega_\pm(x^\frac{1}{2}\log\log\log{x})$, and similar consequences hold for $\Delta^\psi(x;q,1)$ and $\psi(x;q,1,a)$. As another application, the authors establish that the error term in the asymptotic formula for $\sum_{n\leq x}2^{\omega(n)}$ is $\Omega_\pm\big(x^{1/4} (\log\log{x})^{1/2} / (\log\log\log{x})^{3/2} \big)$. (As an aside, in this variant of the divisor problem, obtaining even $\ll x^{1/2-\delta}$ for the error term requires~RH.)\par This article cites~\cites{MR1083000,MR1105608,MR1645811,1914.Littlewood}.}
      \field{issn}{0305-0041}
      \field{journaltitle}{Math. Proc. Cambridge Philos. Soc.}
      \field{number}{3}
      \field{title}{{$\Omega$}-estimates for a class of arithmetic error terms}
      \field{volume}{142}
      \field{year}{2007}
      \true{nocite}
      \field{pages}{385\bibrangedash 394}
      \range{pages}{10}
      \verb{doi}
      \verb 10.1017/S0305004107000035
      \endverb
      \verb{mrnumber}
      \verb 2329690
      \endverb
    \endentry
    \entry{MR2506435}{article}{}
      \name{author}{2}{}{%
        {{hash=6fa84c8cc1e455847a05b9946a1bd205}{%
           family={Kaczorowski},
           familyi={K\bibinitperiod},
           given={Jerzy},
           giveni={J\bibinitperiod}}}%
        {{hash=87807ce1c4509fc3bd17938ccd4fe3e6}{%
           family={Wiertelak},
           familyi={W\bibinitperiod},
           given={Kazimierz},
           giveni={K\bibinitperiod}}}%
      }
      \strng{namehash}{5e2b3cb8731c2c2fbd035cbb54743005}
      \strng{fullhash}{5e2b3cb8731c2c2fbd035cbb54743005}
      \strng{bibnamehash}{5e2b3cb8731c2c2fbd035cbb54743005}
      \strng{authorbibnamehash}{5e2b3cb8731c2c2fbd035cbb54743005}
      \strng{authornamehash}{5e2b3cb8731c2c2fbd035cbb54743005}
      \strng{authorfullhash}{5e2b3cb8731c2c2fbd035cbb54743005}
      \field{extraname}{2}
      \field{sortinit}{K}
      \field{sortinithash}{d3edc18d54b9438a72c24c925bfb38f4}
      \field{labelprefix}{A}
      \field{labelnamesource}{author}
      \field{labeltitlesource}{title}
      \field{annotation}{The authors use a combination of methods from~\cites{MR777013,MR791085,MR858395,MR2329690} to prove that $$ W\bigl( E^\psi(T);T;\log\log H(T) \bigr)\gg (\log{T})/H(T) $$ for sufficiently large~$T$, where $H(T)$ is a function with $1\ll H(T)<\log T$. Letting $H(T)=\exp((\log\log{T})^c)$ with a small $c>0$, this theorem recovers the classical result of Littlewood that $\Delta^\psi(x)=\Omega_\pm(\sqrt{x}\log\log\log{x})$. They show the analogous bound for $E^\pi(T;q,a,1)$ assuming HC for Dirichlet $L$-functions. The authors also investigate the summatory function $D_2(x)=\sum_{n\leq x}2^{\omega(n)}$. Assuming that all the simple zeros $\rho=1/2+i\gamma$ of $\zeta(s)$ on the critical line satisfy ${\zeta'(\rho)}\gg|\gamma|^{-O(1)}$, they show that $$ W\bigl( D_2(T);T;T^{1/4}(\log H(T))^{1/2}(\log\log H(T))^{-3/2} \bigr) \gg (\log T)/H(T) $$ for any function $1\ll H(T)\leq \log T$.\par This article cites~\cites{MR777013,MR791085,MR927376,MR945273,MR858395,MR1083000,MR1105608,MR2329690,MR2271616,1914.Littlewood,MR0882550}.}
      \field{issn}{0002-9947}
      \field{journaltitle}{Trans. Amer. Math. Soc.}
      \field{number}{9}
      \field{title}{Oscillations of a given size of some arithmetic error terms}
      \field{volume}{361}
      \field{year}{2009}
      \true{nocite}
      \field{pages}{5023\bibrangedash 5039}
      \range{pages}{17}
      \verb{doi}
      \verb 10.1090/S0002-9947-09-04803-X
      \endverb
      \verb{mrnumber}
      \verb 2506435
      \endverb
    \endentry
    \entry{MR2684490}{article}{}
      \name{author}{2}{}{%
        {{hash=6fa84c8cc1e455847a05b9946a1bd205}{%
           family={Kaczorowski},
           familyi={K\bibinitperiod},
           given={Jerzy},
           giveni={J\bibinitperiod}}}%
        {{hash=87807ce1c4509fc3bd17938ccd4fe3e6}{%
           family={Wiertelak},
           familyi={W\bibinitperiod},
           given={Kazimierz},
           giveni={K\bibinitperiod}}}%
      }
      \strng{namehash}{5e2b3cb8731c2c2fbd035cbb54743005}
      \strng{fullhash}{5e2b3cb8731c2c2fbd035cbb54743005}
      \strng{bibnamehash}{5e2b3cb8731c2c2fbd035cbb54743005}
      \strng{authorbibnamehash}{5e2b3cb8731c2c2fbd035cbb54743005}
      \strng{authornamehash}{5e2b3cb8731c2c2fbd035cbb54743005}
      \strng{authorfullhash}{5e2b3cb8731c2c2fbd035cbb54743005}
      \field{extraname}{3}
      \field{sortinit}{K}
      \field{sortinithash}{d3edc18d54b9438a72c24c925bfb38f4}
      \field{labelprefix}{A}
      \field{labelnamesource}{author}
      \field{labeltitlesource}{title}
      \field{annotation}{In this article, the authors study the oscillations of $f(x)=\sum_{n\leq x}\phi(n)-\frac{3}{\pi^2}x^2$, the best known results being Walfisz's estimate $f(x)\ll x(\log{x})^{2/3}(\log\log{x})^{4/3}$ and Montgomery's result $f(x)=\Omega_\pm(x\sqrt{\log\log{x}})$. The authors show that $f(x) = f^{AR}(x)+f^{AN}(x)$, where the arithmetic part $f^{AR}(x) = \sum_{n=1}^\infty\frac{\mu(n)}{n}\{\frac{x}{n}\}$ and the analytic part $f^{AN}(x) = -\int_0^x f^{AR}(t)\frac{dt}{t} = \frac{1}{2}\sum_{n=1}^\infty\mu(n)\{\frac{x}{n}\}^2+\frac{1}{2}$. They establish Montgomery-type oscillations for $f^{AR}(x)$ and similar functions, including examples arising from coefficients of newforms. The analytic part $f^{AN}(x)$ has an explicit formula (thus is $o(x)$ unconditionally in particular) and is interesting in its own right, and the authors establish Littlewood-type oscillations for it. \par This article cites~\cites{MR777013,MR791085,MR2470105,MR2730494,MR858395,MR2329690,MR2506435,MR2759800}.}
      \field{issn}{0022-314X}
      \field{journaltitle}{J. Number Theory}
      \field{number}{12}
      \field{title}{Oscillations of the remainder term related to the {E}uler totient function}
      \field{volume}{130}
      \field{year}{2010}
      \true{nocite}
      \field{pages}{2683\bibrangedash 2700}
      \range{pages}{18}
      \verb{doi}
      \verb 10.1016/j.jnt.2010.06.010
      \endverb
      \verb{mrnumber}
      \verb 2684490
      \endverb
    \endentry
    \entry{MR2759800}{article}{}
      \name{author}{2}{}{%
        {{hash=6fa84c8cc1e455847a05b9946a1bd205}{%
           family={Kaczorowski},
           familyi={K\bibinitperiod},
           given={Jerzy},
           giveni={J\bibinitperiod}}}%
        {{hash=87807ce1c4509fc3bd17938ccd4fe3e6}{%
           family={Wiertelak},
           familyi={W\bibinitperiod},
           given={Kazimierz},
           giveni={K\bibinitperiod}}}%
      }
      \strng{namehash}{5e2b3cb8731c2c2fbd035cbb54743005}
      \strng{fullhash}{5e2b3cb8731c2c2fbd035cbb54743005}
      \strng{bibnamehash}{5e2b3cb8731c2c2fbd035cbb54743005}
      \strng{authorbibnamehash}{5e2b3cb8731c2c2fbd035cbb54743005}
      \strng{authornamehash}{5e2b3cb8731c2c2fbd035cbb54743005}
      \strng{authorfullhash}{5e2b3cb8731c2c2fbd035cbb54743005}
      \field{extraname}{4}
      \field{sortinit}{K}
      \field{sortinithash}{d3edc18d54b9438a72c24c925bfb38f4}
      \field{labelprefix}{A}
      \field{labelnamesource}{author}
      \field{labeltitlesource}{title}
      \field{annotation}{The authors observe that the order of magnitude of arithmetic error terms is often uninfluenced by smoothing; for example, $\Delta^\psi(x) = \Omega_\pm (x\log\log\log x)$ unconditionally and $\Delta^\psi(x) \ll x(\log x)^2$ on RH, while $A_1^\psi(x)=\Omega_\pm(\sqrt{x})$ unconditionally and $A_1^\psi(x)=O(\sqrt{x})$ on RH. If we define $f(x)=\sum_{n\leq x}\phi(n)-\frac{3}{\pi^2}x^2$, then Walfisz showed that $f(x)\ll x(\log{x})^{2/3}(\log\log{x})^{4/3}$ and Montgomery showed that $f(x)=\Omega_\pm(x\sqrt{\log\log{x}})$. In contrast, the authors show that $A_1^f(x)=\Omega_\pm(\sqrt{x}\log\log\log{x})$ unconditionally and $A_1^f(x) \le \sqrt{x}\exp\big( O(\log{x}/\log\log{x}) \big)$ on RH.\par This article cites~\cites{MR777013, MR791085,MR2470105,MR2730494,MR858395,MR1083000,MR1105608,MR1105608,MR2329690,MR2506435,1914.Littlewood}.}
      \field{issn}{0025-584X}
      \field{journaltitle}{Math. Nachr.}
      \field{number}{11}
      \field{title}{Smoothing arithmetic error terms: the case of the {E}uler {$\phi$} function}
      \field{volume}{283}
      \field{year}{2010}
      \true{nocite}
      \field{pages}{1637\bibrangedash 1645}
      \range{pages}{9}
      \verb{doi}
      \verb 10.1002/mana.200810048
      \endverb
      \verb{mrnumber}
      \verb 2759800
      \endverb
    \endentry
    \entry{MR2256805}{article}{}
      \name{author}{1}{}{%
        {{hash=1ad9cfd8df25881b7891382dee75005c}{%
           family={Karatsuba},
           familyi={K\bibinitperiod},
           given={A.\bibnamedelimi A.},
           giveni={A\bibinitperiod\bibinitdelim A\bibinitperiod}}}%
      }
      \list{language}{1}{%
        {Russian}%
      }
      \strng{namehash}{1ad9cfd8df25881b7891382dee75005c}
      \strng{fullhash}{1ad9cfd8df25881b7891382dee75005c}
      \strng{bibnamehash}{1ad9cfd8df25881b7891382dee75005c}
      \strng{authorbibnamehash}{1ad9cfd8df25881b7891382dee75005c}
      \strng{authornamehash}{1ad9cfd8df25881b7891382dee75005c}
      \strng{authorfullhash}{1ad9cfd8df25881b7891382dee75005c}
      \field{extraname}{1}
      \field{sortinit}{K}
      \field{sortinithash}{d3edc18d54b9438a72c24c925bfb38f4}
      \field{labelprefix}{A}
      \field{labelnamesource}{author}
      \field{labeltitlesource}{title}
      \field{annotation}{Let $R_1(x)=\pi(x)-\li(x)+\frac{1}{2}\li(\sqrt{x})$. The author follows Kaczorowski's method in~\cite{MR777013, MR791085} to show that for $W(R_1,T) \gg \log T$ and $W(\int_0^x R_1(t)\,dt,T) \gg \log T$. \par This article cites~\cites{MR2169423,MR777013, MR791085}.}
      \field{issn}{0869-5652}
      \field{journaltitle}{Dokl. Akad. Nauk}
      \field{number}{4}
      \field{title}{Behavior of the function {$R_1(x)$} and of its mean value}
      \field{volume}{404}
      \field{year}{2005}
      \true{nocite}
      \field{pages}{439\bibrangedash 442}
      \range{pages}{4}
      \verb{mrnumber}
      \verb 2256805
      \endverb
    \endentry
    \entry{MR2169423}{article}{}
      \name{author}{1}{}{%
        {{hash=1ad9cfd8df25881b7891382dee75005c}{%
           family={Karatsuba},
           familyi={K\bibinitperiod},
           given={A.\bibnamedelimi A.},
           giveni={A\bibinitperiod\bibinitdelim A\bibinitperiod}}}%
      }
      \list{language}{1}{%
        {Russian}%
      }
      \strng{namehash}{1ad9cfd8df25881b7891382dee75005c}
      \strng{fullhash}{1ad9cfd8df25881b7891382dee75005c}
      \strng{bibnamehash}{1ad9cfd8df25881b7891382dee75005c}
      \strng{authorbibnamehash}{1ad9cfd8df25881b7891382dee75005c}
      \strng{authornamehash}{1ad9cfd8df25881b7891382dee75005c}
      \strng{authorfullhash}{1ad9cfd8df25881b7891382dee75005c}
      \field{extraname}{2}
      \field{sortinit}{K}
      \field{sortinithash}{d3edc18d54b9438a72c24c925bfb38f4}
      \field{labelprefix}{A}
      \field{labelnamesource}{author}
      \field{labeltitlesource}{title}
      \field{annotation}{Let $R_1(x)=\pi(x)-\li(x)+\frac{1}{2}\li(\sqrt{x})$, and recursively define $$ R_{j+1}(x) = \int_2^{x} R_j(t)\,dt = -\frac{x^{j+\frac{1}{2}}}{\log x} \sum_{\gamma} \frac{x^{\frac{1}{2}+i \gamma}}{(\frac{1}{2}+i\gamma) \cdots (j+\frac{1}{2}+i \gamma)} +O\bigg(\frac{x^{j+\frac{1}{2}}}{\log^2 x}\bigg) $$ for each $j \geq 1$. Ingham~\cite{1932.Ingham} conjectured that~$\pi(x)$ is better approximated by $\li(x)-\frac{1}{2}\li(\sqrt{x})$ than by $\li(x)$. Under RH, the author validates this conjecture by showing that each of the functions $R_j(x)$ changes signs infinitely often, in contrast to $\fA^\pi_1(x) < -\frac35 x^{3/2}/\log x + O(x^{3/2}/\log^2 x)$. \par This article cites~\cites{1932.Ingham}.}
      \field{issn}{2226-8383}
      \field{journaltitle}{Chebyshevskii Sb.}
      \field{number}{4(12)}
      \field{title}{On the approximation of {$\pi(x)$}}
      \field{volume}{5}
      \field{year}{2005}
      \true{nocite}
      \field{pages}{5\bibrangedash 20}
      \range{pages}{16}
      \verb{mrnumber}
      \verb 2169423
      \endverb
    \endentry
    \entry{MR2262605}{article}{}
      \name{author}{1}{}{%
        {{hash=1ad9cfd8df25881b7891382dee75005c}{%
           family={Karatsuba},
           familyi={K\bibinitperiod},
           given={A.\bibnamedelimi A.},
           giveni={A\bibinitperiod\bibinitdelim A\bibinitperiod}}}%
      }
      \list{language}{1}{%
        {Russian}%
      }
      \strng{namehash}{1ad9cfd8df25881b7891382dee75005c}
      \strng{fullhash}{1ad9cfd8df25881b7891382dee75005c}
      \strng{bibnamehash}{1ad9cfd8df25881b7891382dee75005c}
      \strng{authorbibnamehash}{1ad9cfd8df25881b7891382dee75005c}
      \strng{authornamehash}{1ad9cfd8df25881b7891382dee75005c}
      \strng{authorfullhash}{1ad9cfd8df25881b7891382dee75005c}
      \field{extraname}{3}
      \field{sortinit}{K}
      \field{sortinithash}{d3edc18d54b9438a72c24c925bfb38f4}
      \field{labelprefix}{A}
      \field{labelnamesource}{author}
      \field{labeltitlesource}{title}
      \field{annotation}{Let $R_1(x)=\pi(x)-\li(x)+\frac{1}{2}\li(\sqrt{x})$ and $R_{j+1}(x)=\int_2^{x} R_j(t)\,dt$ for $j\ge1$. The author generalizes his earlier article~\cite{MR2256805} to show that for $W(R_j,T) \gg \log T$ for all positive integers $j$. \par This article cites~\cites{MR0973808A, MR0973808B, MR777013, MR791085, MR2169423}.}
      \field{isbn}{5-87954-388-9}
      \field{issn}{2226-8383}
      \field{journaltitle}{Chebyshevskii Sb.}
      \field{number}{2(14)}
      \field{title}{On the number of sign changes of the function {$R_1(x)$} and its mean values}
      \field{volume}{6}
      \field{year}{2005}
      \true{nocite}
      \field{pages}{163\bibrangedash 183}
      \range{pages}{21}
      \verb{mrnumber}
      \verb 2262605
      \endverb
    \endentry
    \entry{MR0176967}{article}{}
      \name{author}{1}{}{%
        {{hash=cc661844b1342c55df563c91418ea523}{%
           family={Kátai},
           familyi={K\bibinitperiod},
           given={I.},
           giveni={I\bibinitperiod}}}%
      }
      \list{language}{1}{%
        {German}%
      }
      \strng{namehash}{cc661844b1342c55df563c91418ea523}
      \strng{fullhash}{cc661844b1342c55df563c91418ea523}
      \strng{bibnamehash}{cc661844b1342c55df563c91418ea523}
      \strng{authorbibnamehash}{cc661844b1342c55df563c91418ea523}
      \strng{authornamehash}{cc661844b1342c55df563c91418ea523}
      \strng{authorfullhash}{cc661844b1342c55df563c91418ea523}
      \field{extraname}{1}
      \field{sortinit}{K}
      \field{sortinithash}{d3edc18d54b9438a72c24c925bfb38f4}
      \field{labelprefix}{A}
      \field{labelnamesource}{author}
      \field{labeltitlesource}{title}
      \field{annotation}{Let $\ell_1$ and $\ell_2$ be distinct reduced residues\mod k. Assuming HC, the author proves that \[ \limsup_{x \to \infty} \frac{\psi(x,k,\ell_1,\ell_2)}{\sqrt{x}}>0\] (and hence the corresponding statement for $\liminf$). When $\ell_1$ and $\ell_2$ are either both quadratic residues or both quadratic nonresidues, it follows that \[ \limsup_{x \to \infty} \frac{\pi(x,k,\ell_1,\ell_2)}{{\sqrt{x}}/{\log x}} >0,\] (and the corresponding statement for $\liminf$). The proof uses an idea of Littlewood, namely to estimate the iterated integrals $\fA_n(x) = \int_2^x \fA_{n-1}(u) \ du$ where $\fA_0(x) = \psi(x;k,\ell_1,\ell_2) +O(\log x)$ is the explicit sum over zeros of Dirichlet $L$-functions\mod k.\par Assuming GRH, the author can make the above statements quantitative and localized to intervals of the form $(x_0,ax_0)$, thus obtaining the lower bounds $W^\psi_{k;\ell_1,\ell_2}(T) \gg \log x$ and (under the same assumption on $\ell_1$ and $\ell_2$) the same estimate for $W^\psi_{k;\ell_1,\ell_2}(T) \gg \log x$.\par This article cites~\cites{MR1511360, MR0146156,MR0146157,MR0146158,MR0146159,MR0146160,MR0146161,MR0156826,MR0156827, MR1574052}.}
      \field{issn}{0524-9007}
      \field{journaltitle}{Ann. Univ. Sci. Budapest. Eötvös Sect. Math.}
      \field{title}{Eine {B}emerkung zur ``{C}omparative prime-number theory {I}-{VIII}'' von {S}. {K}napowski und {P}. {T}urán}
      \field{volume}{7}
      \field{year}{1964}
      \true{nocite}
      \field{pages}{33\bibrangedash 40}
      \range{pages}{8}
      \verb{mrnumber}
      \verb 0176967
      \endverb
    \endentry
    \entry{MR0207665}{article}{}
      \name{author}{1}{}{%
        {{hash=cc661844b1342c55df563c91418ea523}{%
           family={Kátai},
           familyi={K\bibinitperiod},
           given={I.},
           giveni={I\bibinitperiod}}}%
      }
      \list{language}{1}{%
        {Russian}%
      }
      \strng{namehash}{cc661844b1342c55df563c91418ea523}
      \strng{fullhash}{cc661844b1342c55df563c91418ea523}
      \strng{bibnamehash}{cc661844b1342c55df563c91418ea523}
      \strng{authorbibnamehash}{cc661844b1342c55df563c91418ea523}
      \strng{authornamehash}{cc661844b1342c55df563c91418ea523}
      \strng{authorfullhash}{cc661844b1342c55df563c91418ea523}
      \field{extraname}{2}
      \field{sortinit}{K}
      \field{sortinithash}{d3edc18d54b9438a72c24c925bfb38f4}
      \field{labelprefix}{A}
      \field{labelnamesource}{author}
      \field{labeltitlesource}{title}
      \field{annotation}{The author establishes various theorems on bounds for various arithmetic functions using the ideas of Rodosskiĭ. These results include bounds related to $M(x)$ and $\psi(x;q,a,b)$ and $\psi_r(x;q,a)$, as well as estimates related to prime races modulo~$8$.\par This article cites~\cites{MR0087685,MR0146159,MR0146160,MR0146161,MR0156826,MR0156827,MR0177957}.}
      \field{issn}{0001-5954}
      \field{journaltitle}{Acta Math. Acad. Sci. Hungar}
      \field{title}{Comparative theory of prime numbers}
      \field{volume}{18}
      \field{year}{1967}
      \true{nocite}
      \field{pages}{133\bibrangedash 149}
      \range{pages}{17}
      \verb{doi}
      \verb 10.1007/BF02020967
      \endverb
      \verb{mrnumber}
      \verb 0207665
      \endverb
    \endentry
    \entry{MR0218318}{article}{}
      \name{author}{1}{}{%
        {{hash=cc661844b1342c55df563c91418ea523}{%
           family={Kátai},
           familyi={K\bibinitperiod},
           given={I.},
           giveni={I\bibinitperiod}}}%
      }
      \strng{namehash}{cc661844b1342c55df563c91418ea523}
      \strng{fullhash}{cc661844b1342c55df563c91418ea523}
      \strng{bibnamehash}{cc661844b1342c55df563c91418ea523}
      \strng{authorbibnamehash}{cc661844b1342c55df563c91418ea523}
      \strng{authornamehash}{cc661844b1342c55df563c91418ea523}
      \strng{authorfullhash}{cc661844b1342c55df563c91418ea523}
      \field{extraname}{3}
      \field{sortinit}{K}
      \field{sortinithash}{d3edc18d54b9438a72c24c925bfb38f4}
      \field{labelprefix}{A}
      \field{labelnamesource}{author}
      \field{labeltitlesource}{title}
      \field{annotation}{The author establishes an oscillation theorem of Landau type for Dirichlet integrals with a nonreal pole of arbitrary multiplicity, with the additional feature that the oscillations can be localized to explicit intervals of the form $[T,T^K]$. From this, he deduces many unconditional number-theoretical results. For example, all of the following oscillations can be found in all sufficiently large intervals of the form $[T,T^{7+4\sqrt3}]$: \begin{itemize} \item $M(x) = \Omega_\pm(\sqrt x)$, and the same for $M(x,\chi_{-4})$ and $M(x;4,1)$ and $M(x;4,3)$ and $M_e(x)$ \item $M_r(x) = \Omega_\pm(1/\sqrt x)$, and the same for $\sum_{n=1}^\infty \frac{\mu(n)}n e^{-(x/n)^2}$ \item for any $k\ge2$, $\Delta^{Q_k}(x)=\Omega_\pm(x^{1/2k})$, and similarly for the sum of $e^{-n/x}$ over all $k$-free numbers \item $\psi(x;8,\ell_1,\ell_2) = \Omega_\pm(x^{1/2})$ and, if $\ell_1,\ell_2\not\equiv1\mod8$, then $\pi(x;8,\ell_1,\ell_2) = \Omega_\pm(x^{1/2})/\log x$ \end{itemize} Furthermore, by considering separately the cases where RH or GRH is true or false, the author finds all the following oscillations in intervals of the form $[T,T^{1+\ep}]$ for any fixed $\ep>0$: \begin{itemize} \item $M(x) = \Omega_\pm(x^{\Theta-\ep})$, and similarly for the other functions in the previous list \item under HC, $\psi(x;k,\ell_1,\ell_2) = \Omega_\pm(x^{\Theta(k)-\ep})$, and the same for $\psi_e(x;k,\ell_1,\ell_2)$ (indeed, a hypothesis slightly weaker than HC is required, in that real zeros of different $L$-functions could cancel each other out) \end{itemize} These last theorems imply qualitative improvements on the number of sign changes of their respective functions: for example, $W(M,T)/\log\log T\to\infty$ and, under HC, $W^\psi_{k,\ell_1,\ell_2}(T)/\log\log T\to\infty$. \par This article cites~\cites{MR1511360, MR1555148, MR0146156, MR0146157, MR0146158, MR0146159, MR0146160, MR0146161, MR0156826, MR0156827}.}
      \field{issn}{0001-5954}
      \field{journaltitle}{Acta Math. Acad. Sci. Hungar.}
      \field{title}{On investigations in the comparative prime number theory}
      \field{volume}{18}
      \field{year}{1967}
      \true{nocite}
      \field{pages}{379\bibrangedash 391}
      \range{pages}{13}
      \verb{doi}
      \verb 10.1007/BF02280297
      \endverb
      \verb{mrnumber}
      \verb 0218318
      \endverb
      \verb{urlraw}
      \verb https://doi.org/10.1007/BF02280297
      \endverb
      \verb{url}
      \verb https://doi.org/10.1007/BF02280297
      \endverb
    \endentry
    \entry{MR0219496}{article}{}
      \name{author}{1}{}{%
        {{hash=cc661844b1342c55df563c91418ea523}{%
           family={Kátai},
           familyi={K\bibinitperiod},
           given={I.},
           giveni={I\bibinitperiod}}}%
      }
      \strng{namehash}{cc661844b1342c55df563c91418ea523}
      \strng{fullhash}{cc661844b1342c55df563c91418ea523}
      \strng{bibnamehash}{cc661844b1342c55df563c91418ea523}
      \strng{authorbibnamehash}{cc661844b1342c55df563c91418ea523}
      \strng{authornamehash}{cc661844b1342c55df563c91418ea523}
      \strng{authorfullhash}{cc661844b1342c55df563c91418ea523}
      \field{extraname}{4}
      \field{sortinit}{K}
      \field{sortinithash}{d3edc18d54b9438a72c24c925bfb38f4}
      \field{labelprefix}{A}
      \field{labelnamesource}{author}
      \field{labeltitlesource}{title}
      \field{annotation}{Assuming RH$(H)$, the author shows that $\max_{T^\kappa\le x\le T} M(x) \gg \sqrt x$ for~$T$ sufficiently large in terms of~$H$ and for an explicit $\kappa=\kappa(H)$ that increases to $1$ as $H\to\infty$ (and similarly with $M(x)$ replaced by $-M(x)$). All constants are effectively computable; for example, the value $\kappa=0.36$ follows from Rosser and Schoenfeld's calculations of zeros of $\zeta(s)$. The author proves analogous results for several weighted versions of $M(x)$, such as \[ S(x)=\sum_{n=1}^\infty \frac{\mu(n)}{n}e^{-(x/n)^2} = \sum_{k=1}^\infty \frac{(-x^2)^k}{k!\zeta(2k+1)}, \] and also establishes lower bounds of magnitude $T^{\kappa/2}$ for various averages of $M(x)$ and its variants. \par This article cites~\cites{MR0218318,MR0231801,MR0096630,MR0155802,MR0142520,MR1555148,MR0177957},}
      \field{issn}{0065-1036}
      \field{journaltitle}{Acta Arith.}
      \field{sortyear}{1967}
      \field{title}{On oscillations of number-theoretic functions}
      \field{volume}{13}
      \field{year}{1967/1968}
      \true{nocite}
      \field{pages}{107\bibrangedash 122}
      \range{pages}{16}
      \verb{doi}
      \verb 10.4064/aa-13-1-107-122
      \endverb
      \verb{mrnumber}
      \verb 0219496
      \endverb
    \endentry
    \entry{MR0233782}{article}{}
      \name{author}{1}{}{%
        {{hash=cc661844b1342c55df563c91418ea523}{%
           family={Kátai},
           familyi={K\bibinitperiod},
           given={I.},
           giveni={I\bibinitperiod}}}%
      }
      \strng{namehash}{cc661844b1342c55df563c91418ea523}
      \strng{fullhash}{cc661844b1342c55df563c91418ea523}
      \strng{bibnamehash}{cc661844b1342c55df563c91418ea523}
      \strng{authorbibnamehash}{cc661844b1342c55df563c91418ea523}
      \strng{authornamehash}{cc661844b1342c55df563c91418ea523}
      \strng{authorfullhash}{cc661844b1342c55df563c91418ea523}
      \field{extraname}{5}
      \field{sortinit}{K}
      \field{sortinithash}{d3edc18d54b9438a72c24c925bfb38f4}
      \field{labelprefix}{A}
      \field{labelnamesource}{author}
      \field{labeltitlesource}{title}
      \field{annotation}{For the moduli $q\in\{3,4,5,6,7,8,9,10,11,12,19,24\}$, and on every sufficiently large interval of the form $x\in[T,T^{7+4\sqrt3}]$, the author shows that $\Deltadot^\pi(x;q,a) = \Omega_\pm(\sqrt x/\log x)$ for all quadratic nonresidues~$a$, and also that $\pi(x;q,a,b) = \Omega_\pm(\sqrt x/\log x)$ when $a$ and $b$ are both quadratic residues or both quadratic nonresidues; the same results hold for the exponentially weighted functions $\Deltadot^\pi_e(x;q,a)$ and $\pi_e(x;q,a,b)$.\par This article cites~\cites{1914.Littlewood,MR0067145,MR0133309,MR0202686,MR0156827,MR0177963,MR0087685}.}
      \field{issn}{0001-6969}
      \field{journaltitle}{Acta Sci. Math. (Szeged)}
      \field{title}{On oscillation of the number of primes in an arithmetical progression.}
      \field{volume}{29}
      \field{year}{1968}
      \true{nocite}
      \field{pages}{271\bibrangedash 282}
      \range{pages}{12}
      \verb{mrnumber}
      \verb 0233782
      \endverb
    \endentry
    \entry{MR0231801}{article}{}
      \name{author}{1}{}{%
        {{hash=fa8135c285260ff6fe48a19f384aa417}{%
           family={Kátai},
           familyi={K\bibinitperiod},
           given={Imre},
           giveni={I\bibinitperiod}}}%
      }
      \list{language}{1}{%
        {Hungarian}%
      }
      \strng{namehash}{fa8135c285260ff6fe48a19f384aa417}
      \strng{fullhash}{fa8135c285260ff6fe48a19f384aa417}
      \strng{bibnamehash}{fa8135c285260ff6fe48a19f384aa417}
      \strng{authorbibnamehash}{fa8135c285260ff6fe48a19f384aa417}
      \strng{authornamehash}{fa8135c285260ff6fe48a19f384aa417}
      \strng{authorfullhash}{fa8135c285260ff6fe48a19f384aa417}
      \field{extraname}{6}
      \field{sortinit}{K}
      \field{sortinithash}{d3edc18d54b9438a72c24c925bfb38f4}
      \field{labelprefix}{A}
      \field{labelnamesource}{author}
      \field{labeltitlesource}{title}
      \field{annotation}{Assuming RH, the author proves that $M(x)=\Omega_{\pm}(Tx^{1/2}e^{-c(\log\log{x})^2})$ for some $c>0$. Analogously, assuming RH$(B)$ for any given $B>10^5$, the author proves that $M(x)=\Omega_{\pm}(x^{1/2-\delta}e^{-c(\log\log{x})^2})$, where $\delta\ll(\log\log B)/\log B$.\par This article cites~\cites{MR0046485,MR0096630,MR0133287}.}
      \field{journaltitle}{Magyar Tud. Akad. Mat. Fiz. Oszt. Közl.}
      \field{title}{The {$\Omega$}-estimation of the arithmetic mean of the Möbius function}
      \field{volume}{15}
      \field{year}{1965}
      \true{nocite}
      \field{pages}{15\bibrangedash 18}
      \range{pages}{4}
      \verb{mrnumber}
      \verb 0231801
      \endverb
    \endentry
    \entry{MR0241374}{article}{}
      \name{author}{1}{}{%
        {{hash=fa8135c285260ff6fe48a19f384aa417}{%
           family={Kátai},
           familyi={K\bibinitperiod},
           given={Imre},
           giveni={I\bibinitperiod}}}%
      }
      \list{language}{1}{%
        {Hungarian, with English summary}%
      }
      \strng{namehash}{fa8135c285260ff6fe48a19f384aa417}
      \strng{fullhash}{fa8135c285260ff6fe48a19f384aa417}
      \strng{bibnamehash}{fa8135c285260ff6fe48a19f384aa417}
      \strng{authorbibnamehash}{fa8135c285260ff6fe48a19f384aa417}
      \strng{authornamehash}{fa8135c285260ff6fe48a19f384aa417}
      \strng{authorfullhash}{fa8135c285260ff6fe48a19f384aa417}
      \field{extraname}{7}
      \field{sortinit}{K}
      \field{sortinithash}{d3edc18d54b9438a72c24c925bfb38f4}
      \field{labelprefix}{A}
      \field{labelnamesource}{author}
      \field{labeltitlesource}{title}
      \field{annotation}{The author summarizes many oscillation results from his earlier work, most unconditional and some assuming HC, for functions including $M(x)$, $M_r(x)$, $M_e(x)$, $\psi(x,k,\ell_1,\ell_2)$ and $\pi(x,k,\ell_1,\ell_2)$, $\psi_e(x,k,\ell_1,\ell_2)$, and $\sum_{n=1}^\infty \frac{\mu(n)}n e^{-(x/n)^2}$, as well as functions involving the distribution of $k$-free integers. An example that typifies the strengths of the results is $$ \max_{T<x<T^{7+4\sqrt3}} M(x) = \Omega_\pm(\sqrt x) \quad\text{and}\quad \max_{T<x<T^{1+\ep}} M(x) = \Omega_\pm(x^{\Theta-\ep}). $$\par This article cites~\cites{MR1555148,MR1511360,MR1544293B,1914.Littlewood,MR1574052,MR1575079,MR0146156,MR0146157,MR0146158,MR0146159,MR0146160,MR0146161,MR0156826,MR0156827,MR0096630,MR0142520,MR0155802,MR0172856,MR0087685,MR0177957,MR0046485,MR0067145}.}
      \field{journaltitle}{Magyar Tud. Akad. Mat. Fiz. Oszt. Közl.}
      \field{title}{Omega-type investigations in prime number theory}
      \field{volume}{16}
      \field{year}{1966}
      \true{nocite}
      \field{pages}{369\bibrangedash 396}
      \range{pages}{28}
      \verb{mrnumber}
      \verb 0241374
      \endverb
    \endentry
    \entry{MR3414302}{article}{}
      \name{author}{2}{}{%
        {{hash=6b393826a3dd69e5d16ac8319b2612dd}{%
           family={Kisilevsky},
           familyi={K\bibinitperiod},
           given={H.},
           giveni={H\bibinitperiod}}}%
        {{hash=8aabc8fd166d5aa179c76a42ffba2cc5}{%
           family={Rubinstein},
           familyi={R\bibinitperiod},
           given={M.\bibnamedelimi O.},
           giveni={M\bibinitperiod\bibinitdelim O\bibinitperiod}}}%
      }
      \strng{namehash}{f6f50421343a63c3dbf092af65a8d0d5}
      \strng{fullhash}{f6f50421343a63c3dbf092af65a8d0d5}
      \strng{bibnamehash}{f6f50421343a63c3dbf092af65a8d0d5}
      \strng{authorbibnamehash}{f6f50421343a63c3dbf092af65a8d0d5}
      \strng{authornamehash}{f6f50421343a63c3dbf092af65a8d0d5}
      \strng{authorfullhash}{f6f50421343a63c3dbf092af65a8d0d5}
      \field{sortinit}{K}
      \field{sortinithash}{d3edc18d54b9438a72c24c925bfb38f4}
      \field{labelprefix}{A}
      \field{labelnamesource}{author}
      \field{labeltitlesource}{title}
      \field{annotation}{Let $P_{\rm odd}= \{ 2, 5, 11, 17, 23, 31, \ldots \}$ be the set consisting of every other prime (the odd-index primes). In this article, the authors show that $P_{\rm odd}$ cannot be written as a finite union of sets of the form $\{p\colon p\equiv a\mod q\}$, even up to finitely many exceptions. More generally, call a set of $P$ prime ideals of a number field $K$ a Chebotarev set if there are finitely many finite Galois extensions $L_i/K$ and conjugacy classes $C_i$ such that the symmetric difference of the sets $P$ and $\bigcup_{i=1} \pi_{C_i}(L_i/K)$ is finite. Using explicit formulas, the authors show that if~$P$ is a Chebotarev set of density $\beta\in\mathbb{Q}$ with $0<\beta<1$, then $\Delta_{P}(x)=P(x)-\beta\pi_K(x)=\Omega(x^{1/2}/\log x)$. In particular, $P_{\rm odd}$ is not a Chebatorev set (hence not a finite union of sets of primes in residue classes) since the counting function of $P_{\rm odd}$ is $\frac12\pi(x) + O(1)$.\par This article cites~\cites{1932.Ingham,MR1329368}.}
      \field{issn}{0065-1036}
      \field{journaltitle}{Acta Arith.}
      \field{number}{2}
      \field{title}{Chebotarev sets}
      \field{volume}{171}
      \field{year}{2015}
      \true{nocite}
      \field{pages}{97\bibrangedash 124}
      \range{pages}{28}
      \verb{doi}
      \verb 10.4064/aa171-2-1
      \endverb
      \verb{mrnumber}
      \verb 3414302
      \endverb
      \verb{urlraw}
      \verb https://doi.org/10.4064/aa171-2-1
      \endverb
      \verb{url}
      \verb https://doi.org/10.4064/aa171-2-1
      \endverb
    \endentry
    \entry{MR0096622}{article}{}
      \name{author}{1}{}{%
        {{hash=d33db82fc1e83dc064a23d93e5a21891}{%
           family={Knapowski},
           familyi={K\bibinitperiod},
           given={S.},
           giveni={S\bibinitperiod}}}%
      }
      \strng{namehash}{d33db82fc1e83dc064a23d93e5a21891}
      \strng{fullhash}{d33db82fc1e83dc064a23d93e5a21891}
      \strng{bibnamehash}{d33db82fc1e83dc064a23d93e5a21891}
      \strng{authorbibnamehash}{d33db82fc1e83dc064a23d93e5a21891}
      \strng{authornamehash}{d33db82fc1e83dc064a23d93e5a21891}
      \strng{authorfullhash}{d33db82fc1e83dc064a23d93e5a21891}
      \field{extraname}{1}
      \field{sortinit}{K}
      \field{sortinithash}{d3edc18d54b9438a72c24c925bfb38f4}
      \field{labelprefix}{A}
      \field{labelnamesource}{author}
      \field{labeltitlesource}{title}
      \field{annotation}{This article begins by recalling that $\Delta^\pi(x;q,1) \ll_\ep x^{\theta+\ep}$ implies $\Delta^\pi(x;q,a) \ll_\ep x^{\theta+\ep}$ for all $(a,q)=1$. The author establishes an explicit inequality relating the two, namely that \[ \max_{x\le T} |\Delta^\pi(x;q,a)| \le T^{\delta(T)} \exp\bigg( \frac{(1+q^{-1})\log T}{\sqrt{\log\log T}} \bigg) \big( \max_{x\le T} |\Delta^\pi(x;q,1)| + \phi(q)\sqrt T \big). \] Here $\delta(T) = \ep(\sqrt T)-\ep\big(\exp(\sqrt{\log\log T})\big)$, where $\ep(H)$ is the largest real part of the zeros of $\zeta(s)$ up to height~$H$; in particular, $\delta(T)$ is eventually identically~$0$ assuming~SA, and $\delta(T)\to0$ unconditionally. The author derives this result from the analogous result for $\Delta^\psi(x;q,a)$, which does not contain the term $\phi(q)\sqrt T$; he remarks upon a similar result for $\Delta^\Pi(x;q,a)$ assuming~SA. \par This article cites~\cites{1909.Landau,MR0060548}.}
      \field{issn}{0065-1036}
      \field{journaltitle}{Acta Arith.}
      \field{title}{On prime numbers in an arithmetical progression}
      \field{volume}{4}
      \field{year}{1958}
      \true{nocite}
      \field{pages}{57\bibrangedash 70}
      \range{pages}{14}
      \verb{doi}
      \verb 10.4064/aa-4-1-57-70
      \endverb
      \verb{mrnumber}
      \verb 0096622
      \endverb
      \verb{urlraw}
      \verb https://doi.org/10.4064/aa-4-1-57-7
      \endverb
      \verb{url}
      \verb https://doi.org/10.4064/aa-4-1-57-7
      \endverb
    \endentry
    \entry{MR0096630}{article}{}
      \name{author}{1}{}{%
        {{hash=d33db82fc1e83dc064a23d93e5a21891}{%
           family={Knapowski},
           familyi={K\bibinitperiod},
           given={S.},
           giveni={S\bibinitperiod}}}%
      }
      \strng{namehash}{d33db82fc1e83dc064a23d93e5a21891}
      \strng{fullhash}{d33db82fc1e83dc064a23d93e5a21891}
      \strng{bibnamehash}{d33db82fc1e83dc064a23d93e5a21891}
      \strng{authorbibnamehash}{d33db82fc1e83dc064a23d93e5a21891}
      \strng{authornamehash}{d33db82fc1e83dc064a23d93e5a21891}
      \strng{authorfullhash}{d33db82fc1e83dc064a23d93e5a21891}
      \field{extraname}{2}
      \field{sortinit}{K}
      \field{sortinithash}{d3edc18d54b9438a72c24c925bfb38f4}
      \field{labelprefix}{A}
      \field{labelnamesource}{author}
      \field{labeltitlesource}{title}
      \field{annotation}{The author proves that the assumption $\int_1^T ({M(x)}/x)^2 \, dx \ll \log T$, which is stronger than the Mertens conjecture and is known to imply both RH and the simplicity of the zeros of~$\zeta(s)$, also implies that $M(x) = \Omega\big( x^{1/2}\exp(-{\log x}/{\sqrt{\log\log x}}) \big)$.\par This article cites~\cite{MR0060548}.}
      \field{issn}{0065-1036}
      \field{journaltitle}{Acta Arith.}
      \field{title}{On the Möbius function}
      \field{volume}{4}
      \field{year}{1958}
      \true{nocite}
      \field{pages}{209\bibrangedash 216}
      \range{pages}{8}
      \verb{doi}
      \verb 10.4064/aa-4-3-209-216
      \endverb
      \verb{mrnumber}
      \verb 0096630
      \endverb
    \endentry
    \entry{MR0111722}{article}{}
      \name{author}{1}{}{%
        {{hash=d33db82fc1e83dc064a23d93e5a21891}{%
           family={Knapowski},
           familyi={K\bibinitperiod},
           given={S.},
           giveni={S\bibinitperiod}}}%
      }
      \list{language}{1}{%
        {English, with Russian summary}%
      }
      \strng{namehash}{d33db82fc1e83dc064a23d93e5a21891}
      \strng{fullhash}{d33db82fc1e83dc064a23d93e5a21891}
      \strng{bibnamehash}{d33db82fc1e83dc064a23d93e5a21891}
      \strng{authorbibnamehash}{d33db82fc1e83dc064a23d93e5a21891}
      \strng{authornamehash}{d33db82fc1e83dc064a23d93e5a21891}
      \strng{authorfullhash}{d33db82fc1e83dc064a23d93e5a21891}
      \field{extraname}{3}
      \field{sortinit}{K}
      \field{sortinithash}{d3edc18d54b9438a72c24c925bfb38f4}
      \field{labelprefix}{A}
      \field{labelnamesource}{author}
      \field{labeltitlesource}{title}
      \field{annotation}{The author continues his work in~\cite{MR0096630} by considering lower bounds for the logarithmic averages $A^\psi_{|1|}(x)$ and $A^M_{|1|}(x)$. The author shows that if $\rho_0=\beta_0+i\gamma_0$ is a zero of $\zeta(s)$, then \[ A^\psi_{|1|}(x) > x^{\beta_0}\exp\bigg( -14\frac{\log{x}}{\sqrt{\log\log{x}}} \bigg) \] when $x$ is sufficiently large in terms of~$|\rho_0|$. He also shows that the conjectured upper bound $A^M_{|1|}(x) \ll \sqrt x$ actually implies the lower bound \[ A^M_{|1|}(x) > \sqrt x \exp\bigg( -\frac{\log x}{\sqrt{\log\log x}} \bigg) \] when $x$ is sufficiently large.\par This article cites~\cites{MR0096630, MR0068565, MR0077579, MR0043121, MR0060548}.}
      \field{issn}{0001-5954}
      \field{journaltitle}{Acta Math. Acad. Sci. Hungar.}
      \field{title}{On the mean values of certain functions in prime number theory}
      \field{volume}{10}
      \field{year}{1959}
      \true{nocite}
      \field{pages}{375\bibrangedash 390. (unbound insert)}
      \range{pages}{-1}
      \verb{mrnumber}
      \verb 0111722
      \endverb
    \endentry
    \entry{MR0125822}{article}{}
      \name{author}{1}{}{%
        {{hash=d33db82fc1e83dc064a23d93e5a21891}{%
           family={Knapowski},
           familyi={K\bibinitperiod},
           given={S.},
           giveni={S\bibinitperiod}}}%
      }
      \strng{namehash}{d33db82fc1e83dc064a23d93e5a21891}
      \strng{fullhash}{d33db82fc1e83dc064a23d93e5a21891}
      \strng{bibnamehash}{d33db82fc1e83dc064a23d93e5a21891}
      \strng{authorbibnamehash}{d33db82fc1e83dc064a23d93e5a21891}
      \strng{authornamehash}{d33db82fc1e83dc064a23d93e5a21891}
      \strng{authorfullhash}{d33db82fc1e83dc064a23d93e5a21891}
      \field{extraname}{4}
      \field{sortinit}{K}
      \field{sortinithash}{d3edc18d54b9438a72c24c925bfb38f4}
      \field{labelprefix}{A}
      \field{labelnamesource}{author}
      \field{labeltitlesource}{title}
      \field{annotation}{Assuming GRH$(q^7)$, the author shows that \[ \int_{T\exp(-(\log T)^{3/4})}^T \frac{|{\Delta^\psi(x;q,a)}|}{x} \,dx>{\phi(q)}T^{1/2}\exp\bigg(-\frac{2\log T}{\log\log T}\bigg) \] with $T$ is sufficiently large in terms of~$q$, which in particular implies an $\Omega$-result for $\Delta^\psi(x;q,a)$. Unconditionally the right-hand side can be replaced by $T^{1/4}$. \par This article cites~\cites{MR0096622, MR0087685,MR0060548}.}
      \field{issn}{0065-1036}
      \field{journaltitle}{Acta Arith.}
      \field{sortyear}{1960}
      \field{title}{Contributions to the theory of the distribution of prime numbers in arithmetical progressions. I}
      \field{volume}{6}
      \field{year}{1960/1961}
      \true{nocite}
      \field{pages}{415\bibrangedash 434}
      \range{pages}{20}
      \verb{doi}
      \verb 10.4064/aa-6-4-415-434
      \endverb
      \verb{mrnumber}
      \verb 0125822
      \endverb
    \endentry
    \entry{MR0142520}{article}{}
      \name{author}{1}{}{%
        {{hash=d33db82fc1e83dc064a23d93e5a21891}{%
           family={Knapowski},
           familyi={K\bibinitperiod},
           given={S.},
           giveni={S\bibinitperiod}}}%
      }
      \strng{namehash}{d33db82fc1e83dc064a23d93e5a21891}
      \strng{fullhash}{d33db82fc1e83dc064a23d93e5a21891}
      \strng{bibnamehash}{d33db82fc1e83dc064a23d93e5a21891}
      \strng{authorbibnamehash}{d33db82fc1e83dc064a23d93e5a21891}
      \strng{authornamehash}{d33db82fc1e83dc064a23d93e5a21891}
      \strng{authorfullhash}{d33db82fc1e83dc064a23d93e5a21891}
      \field{extraname}{5}
      \field{sortinit}{K}
      \field{sortinithash}{d3edc18d54b9438a72c24c925bfb38f4}
      \field{labelprefix}{A}
      \field{labelnamesource}{author}
      \field{labeltitlesource}{title}
      \field{annotation}{Assuming GRH$(q^7)$, the author shows that \begin{align*} \int_{T\exp(-(\log T)^{3/4})}^T \frac{|\psi(x;q,a_1,a_2)|}{x} \,dx>T^{1/2}\exp\bigg(-\frac{2\log T}{\log\log T}\bigg) \end{align*} when $T$ is sufficiently large in terms of~$q$, and the analogous result with $\psi$ replaced by~$\Pi$. (Unconditionally the right-hand side can be replaced by $T^{1/4}$.) If both $a_1$ and $a_2$ are nonsquares\mod q, then the same result holds for~$\pi$ as well. \par This article cites~\cites{MR0125822,MR0087685}.}
      \field{issn}{0065-1036}
      \field{journaltitle}{Acta Arith}
      \field{sortyear}{1961}
      \field{title}{Contributions to the theory of the distributon of prime numbers in arithmetical progressions. II}
      \field{volume}{7}
      \field{year}{1961/1962}
      \true{nocite}
      \field{pages}{325\bibrangedash 335}
      \range{pages}{11}
      \verb{doi}
      \verb 10.4064/aa-7-4-325-335
      \endverb
      \verb{mrnumber}
      \verb 0142520
      \endverb
    \endentry
    \entry{MR0133287}{article}{}
      \name{author}{1}{}{%
        {{hash=d33db82fc1e83dc064a23d93e5a21891}{%
           family={Knapowski},
           familyi={K\bibinitperiod},
           given={S.},
           giveni={S\bibinitperiod}}}%
      }
      \strng{namehash}{d33db82fc1e83dc064a23d93e5a21891}
      \strng{fullhash}{d33db82fc1e83dc064a23d93e5a21891}
      \strng{bibnamehash}{d33db82fc1e83dc064a23d93e5a21891}
      \strng{authorbibnamehash}{d33db82fc1e83dc064a23d93e5a21891}
      \strng{authornamehash}{d33db82fc1e83dc064a23d93e5a21891}
      \strng{authorfullhash}{d33db82fc1e83dc064a23d93e5a21891}
      \field{extraname}{6}
      \field{sortinit}{K}
      \field{sortinithash}{d3edc18d54b9438a72c24c925bfb38f4}
      \field{labelprefix}{A}
      \field{labelnamesource}{author}
      \field{labeltitlesource}{title}
      \field{annotation}{Supposing that $A^M_{|1|}(T)<aT^{1/2}$ for $T\ge1$, the author exhibits a constant $H(a)$ such that \[ \int_{T/H(a)}^T\frac{|M(x)|}{x}\,dx>\frac{T^{1/2}}{H(a)}. \] for $T\ge H(a)$, refining a theorem from~\cite{MR0111722} by following an idea due to Ingham in a letter to Turán. The author also states another theorem which is proved later in~\cite{MR0142500}.\par This article cites~\cites{MR0111722,MR0096630,MR0142500}.}
      \field{issn}{0065-1036}
      \field{journaltitle}{Acta Arith.}
      \field{title}{Mean-value estimations for the {M}öbius function. {I}}
      \field{volume}{7}
      \field{year}{1961}
      \true{nocite}
      \field{pages}{121\bibrangedash 130}
      \range{pages}{10}
      \verb{doi}
      \verb 10.4064/aa-7-2-121-130
      \endverb
      \verb{mrnumber}
      \verb 0133287
      \endverb
      \verb{urlraw}
      \verb https://doi.org/10.4064/aa-7-2-121-130
      \endverb
      \verb{url}
      \verb https://doi.org/10.4064/aa-7-2-121-130
      \endverb
    \endentry
    \entry{MR0142500}{article}{}
      \name{author}{1}{}{%
        {{hash=d33db82fc1e83dc064a23d93e5a21891}{%
           family={Knapowski},
           familyi={K\bibinitperiod},
           given={S.},
           giveni={S\bibinitperiod}}}%
      }
      \strng{namehash}{d33db82fc1e83dc064a23d93e5a21891}
      \strng{fullhash}{d33db82fc1e83dc064a23d93e5a21891}
      \strng{bibnamehash}{d33db82fc1e83dc064a23d93e5a21891}
      \strng{authorbibnamehash}{d33db82fc1e83dc064a23d93e5a21891}
      \strng{authornamehash}{d33db82fc1e83dc064a23d93e5a21891}
      \strng{authorfullhash}{d33db82fc1e83dc064a23d93e5a21891}
      \field{extraname}{7}
      \field{sortinit}{K}
      \field{sortinithash}{d3edc18d54b9438a72c24c925bfb38f4}
      \field{labelprefix}{A}
      \field{labelnamesource}{author}
      \field{labeltitlesource}{title}
      \field{annotation}{This article presents a proof of the theorem which was announced in~\cite{MR0133287}: When $T$ is sufficiently large, RH($\omega$) implies that for $T \leq \exp(\omega^{10})$. \[ \int_{T\exp(-100\frac{\log T}{\log\log T}\log\log\log T)}^T\frac{|M(x)|}{x}\,dx>T^{1/2}\exp\biggl(-12\frac{\log T}{\log\log T}\log\log\log T\biggr). \] \par This article cites~\cites{MR0133287,MR0125822,MR0060548}.}
      \field{issn}{0065-1036}
      \field{journaltitle}{Acta Arith.}
      \field{title}{Mean-value estimations for the {M}öbius function. {II}}
      \field{volume}{7}
      \field{year}{1961}
      \true{nocite}
      \field{pages}{337\bibrangedash 343}
      \range{pages}{7}
      \verb{doi}
      \verb 10.4064/aa-7-4-337-343
      \endverb
      \verb{mrnumber}
      \verb 0142500
      \endverb
      \verb{urlraw}
      \verb https://doi.org/10.4064/aa-7-4-337-343
      \endverb
      \verb{url}
      \verb https://doi.org/10.4064/aa-7-4-337-343
      \endverb
    \endentry
    \entry{MR0133309}{article}{}
      \name{author}{1}{}{%
        {{hash=d33db82fc1e83dc064a23d93e5a21891}{%
           family={Knapowski},
           familyi={K\bibinitperiod},
           given={S.},
           giveni={S\bibinitperiod}}}%
      }
      \strng{namehash}{d33db82fc1e83dc064a23d93e5a21891}
      \strng{fullhash}{d33db82fc1e83dc064a23d93e5a21891}
      \strng{bibnamehash}{d33db82fc1e83dc064a23d93e5a21891}
      \strng{authorbibnamehash}{d33db82fc1e83dc064a23d93e5a21891}
      \strng{authornamehash}{d33db82fc1e83dc064a23d93e5a21891}
      \strng{authorfullhash}{d33db82fc1e83dc064a23d93e5a21891}
      \field{extraname}{8}
      \field{sortinit}{K}
      \field{sortinithash}{d3edc18d54b9438a72c24c925bfb38f4}
      \field{labelprefix}{A}
      \field{labelnamesource}{author}
      \field{labeltitlesource}{title}
      \field{annotation}{The author shows that if $\rho_0 = \beta_0 + i \gamma_0$ is any zero of $\zeta(s)$, then for $T$ sufficiently large in terms of~$\gamma_0$, \begin{align*} \Delta^\psi(t) &= \Omega_\pm \bigg( t^{\beta_0} \exp\bigg( \frac{-15\log t}{\sqrt{\log \log t}} \bigg) \bigg), \end{align*} and the same for $\Delta^\Pi(t)$ (which implies an $\Omega_-$-result, though not an $\Omega_+$-result, for $\Delta^\pi(t)$). A slightly more precise version of this result implies that $\liminf_{T\to\infty} W^\psi(T) /\log\log T \ge 1/\log 2$.\par This article cites~\cites{1936.Ingham, MR0043121, MR0067145}.}
      \field{issn}{0024-6107}
      \field{journaltitle}{J. London Math. Soc.}
      \field{title}{On sign-changes in the remainder-term in the prime-number formula}
      \field{volume}{36}
      \field{year}{1961}
      \true{nocite}
      \field{pages}{451\bibrangedash 460}
      \range{pages}{10}
      \verb{doi}
      \verb 10.1112/jlms/s1-36.1.451
      \endverb
      \verb{mrnumber}
      \verb 0133309
      \endverb
      \verb{urlraw}
      \verb https://doi.org/10.1112/jlms/s1-36.1.451
      \endverb
      \verb{url}
      \verb https://doi.org/10.1112/jlms/s1-36.1.451
      \endverb
    \endentry
    \entry{MR0133308}{article}{}
      \name{author}{1}{}{%
        {{hash=d33db82fc1e83dc064a23d93e5a21891}{%
           family={Knapowski},
           familyi={K\bibinitperiod},
           given={S.},
           giveni={S\bibinitperiod}}}%
      }
      \strng{namehash}{d33db82fc1e83dc064a23d93e5a21891}
      \strng{fullhash}{d33db82fc1e83dc064a23d93e5a21891}
      \strng{bibnamehash}{d33db82fc1e83dc064a23d93e5a21891}
      \strng{authorbibnamehash}{d33db82fc1e83dc064a23d93e5a21891}
      \strng{authornamehash}{d33db82fc1e83dc064a23d93e5a21891}
      \strng{authorfullhash}{d33db82fc1e83dc064a23d93e5a21891}
      \field{extraname}{9}
      \field{sortinit}{K}
      \field{sortinithash}{d3edc18d54b9438a72c24c925bfb38f4}
      \field{labelprefix}{A}
      \field{labelnamesource}{author}
      \field{labeltitlesource}{title}
      \field{annotation}{This article is concerned with explicit lower bounds for the number $W(T)$ of sign changes of the function $\li(x)-\pi(x)$. Previously, Ingham~\cite{1936.Ingham} had shown, assuming SA, that $ \liminf_{T \to \infty} {W(T)}/{\log T} >0, $ while the author~\cite{MR0133309} had proved unconditionally that $ \liminf_{T \to \infty} {W(T)}/{\log \log T} >0. $ Skewes~\cite{MR0067145} famously found that $W\big(e^{e^{e^{e^{7.705}}}}\big) \ge 1$. \par In this article, the author shows unconditionally that $W(T) \ge e^{-35} \log \log \log \log T$ for $T \ge e^{e^{e^{e^{35}}}}$ (the author did not try to optimize these constants). Similar to~\cite{MR0067145}, the proof is divided into two cases, first when RH is ``nearly true'' and then the contrary case. More specifically, setting $X = \sqrt[7]{\log \log T}$, the author defines a hypothesis (C) (the ``nearly true'' case) as follows: \emph{Every zero $\rho = \beta + i\gamma$ for which $|\gamma| \le X^3$ satisfies $\beta - \frac{1}{2} \le 2/(3X^3 \log X)$}. The author then proves the inequality above first assuming (C) and then again assuming its negation (NC).\par This article cites~\cites{1914.Littlewood, 1932.Ingham, MR0067145, MR0133309}.}
      \field{issn}{0065-1036}
      \field{journaltitle}{Acta Arith.}
      \field{sortyear}{1961}
      \field{title}{On sign-changes of the difference {$\pi (x)-{\rm li}\,x$}}
      \field{volume}{7}
      \field{year}{1961/1962}
      \true{nocite}
      \field{pages}{107\bibrangedash 119}
      \range{pages}{13}
      \verb{doi}
      \verb 10.4064/aa-7-2-107-119
      \endverb
      \verb{mrnumber}
      \verb 0133308
      \endverb
      \verb{urlraw}
      \verb https://doi.org/10.4064/aa-7-2-107-119
      \endverb
      \verb{url}
      \verb https://doi.org/10.4064/aa-7-2-107-119
      \endverb
    \endentry
    \entry{MR0142521}{article}{}
      \name{author}{1}{}{%
        {{hash=d33db82fc1e83dc064a23d93e5a21891}{%
           family={Knapowski},
           familyi={K\bibinitperiod},
           given={S.},
           giveni={S\bibinitperiod}}}%
      }
      \strng{namehash}{d33db82fc1e83dc064a23d93e5a21891}
      \strng{fullhash}{d33db82fc1e83dc064a23d93e5a21891}
      \strng{bibnamehash}{d33db82fc1e83dc064a23d93e5a21891}
      \strng{authorbibnamehash}{d33db82fc1e83dc064a23d93e5a21891}
      \strng{authornamehash}{d33db82fc1e83dc064a23d93e5a21891}
      \strng{authorfullhash}{d33db82fc1e83dc064a23d93e5a21891}
      \field{extraname}{10}
      \field{sortinit}{K}
      \field{sortinithash}{d3edc18d54b9438a72c24c925bfb38f4}
      \field{labelprefix}{A}
      \field{labelnamesource}{author}
      \field{labeltitlesource}{title}
      \field{annotation}{Assuming GRH$(q^7)$, the author shows that \begin{align*} \int_{T\exp(-(\log T)^{3/4})}^T \frac{|\pi(x;q,a_1,a_2)|}{x} \,dx>T^{1/2}\exp\bigg(-\frac{7\log T}{\log\log T}\bigg); \end{align*} unlike in his earlier work, one no longer needs to assume that $a_1$ and $a_2$ are nonsquares. \par This article cites~\cites{MR0125822,MR0142520,MR0087685,MR0060548}.}
      \field{issn}{0065-1036}
      \field{journaltitle}{Acta Arith}
      \field{sortyear}{1962}
      \field{title}{Contributions to the theory of the distributon of prime numbers in arithmetical progressions. III}
      \field{volume}{8}
      \field{year}{1962/1963}
      \true{nocite}
      \field{pages}{97\bibrangedash 105}
      \range{pages}{9}
      \verb{doi}
      \verb 10.4064/aa-8-1-97-105
      \endverb
      \verb{mrnumber}
      \verb 0142521
      \endverb
    \endentry
    \entry{MR0155802}{article}{}
      \name{author}{1}{}{%
        {{hash=d33db82fc1e83dc064a23d93e5a21891}{%
           family={Knapowski},
           familyi={K\bibinitperiod},
           given={S.},
           giveni={S\bibinitperiod}}}%
      }
      \strng{namehash}{d33db82fc1e83dc064a23d93e5a21891}
      \strng{fullhash}{d33db82fc1e83dc064a23d93e5a21891}
      \strng{bibnamehash}{d33db82fc1e83dc064a23d93e5a21891}
      \strng{authorbibnamehash}{d33db82fc1e83dc064a23d93e5a21891}
      \strng{authornamehash}{d33db82fc1e83dc064a23d93e5a21891}
      \strng{authorfullhash}{d33db82fc1e83dc064a23d93e5a21891}
      \field{extraname}{11}
      \field{sortinit}{K}
      \field{sortinithash}{d3edc18d54b9438a72c24c925bfb38f4}
      \field{labelprefix}{A}
      \field{labelnamesource}{author}
      \field{labeltitlesource}{title}
      \field{annotation}{Assuming RH$(H)$ (but with no assumptions on the simplicity of zeros of $\zeta(s)$), the author shows that \[ \max_{1\le x\le T} M(x) > \sqrt T \exp\bigg( -\frac{15\log T\log\log\log T}{\log\log T} \bigg) \] for $T\le e^{H^{10}}$, and similarly for $-M(x)$; in the same range he deduces that $W(M,T) \gg \log T$. \par This article cites~\cites{MR0096630,MR0111722,MR0133309,MR0146158,MR0132728,MR0133287}.}
      \field{issn}{0065-1036}
      \field{journaltitle}{Acta Arith.}
      \field{sortyear}{1962}
      \field{title}{On oscillations of certain means formed from the Möbius series. I}
      \field{volume}{8}
      \field{year}{1962/1963}
      \true{nocite}
      \field{pages}{311\bibrangedash 320}
      \range{pages}{10}
      \verb{doi}
      \verb 10.4064/aa-8-3-311-320
      \endverb
      \verb{mrnumber}
      \verb 0155802
      \endverb
    \endentry
    \entry{MR0172856}{article}{}
      \name{author}{1}{}{%
        {{hash=d33db82fc1e83dc064a23d93e5a21891}{%
           family={Knapowski},
           familyi={K\bibinitperiod},
           given={S.},
           giveni={S\bibinitperiod}}}%
      }
      \strng{namehash}{d33db82fc1e83dc064a23d93e5a21891}
      \strng{fullhash}{d33db82fc1e83dc064a23d93e5a21891}
      \strng{bibnamehash}{d33db82fc1e83dc064a23d93e5a21891}
      \strng{authorbibnamehash}{d33db82fc1e83dc064a23d93e5a21891}
      \strng{authornamehash}{d33db82fc1e83dc064a23d93e5a21891}
      \strng{authorfullhash}{d33db82fc1e83dc064a23d93e5a21891}
      \field{extraname}{12}
      \field{sortinit}{K}
      \field{sortinithash}{d3edc18d54b9438a72c24c925bfb38f4}
      \field{labelprefix}{A}
      \field{labelnamesource}{author}
      \field{labeltitlesource}{title}
      \field{annotation}{The author investigates oscillations of the Mertens sum $M(x)$ over intervals. Assuming RH, the author shows that for all~$T>1$, \[ \max_{I\subset[Te^{-6(\log T)^{5/6}},Te^{6(\log T)^{5/6}}]} M(I) > T^{1/2}e^{-(\log T)^{3/4}} \] and the symmetric result for the minimum; indeed, only the assumption RH$(H)$ is needed to establish these oscillations for $T\le e^{H^6}$. The author also establishes oscillations of the same size for the absolute logarithmic average $A^M_{|1|}\big( [Te^{-6(\log T)^{5/6}},Te^{6(\log T)^{5/6}}] \big)$. \par This article cites~\cites{MR0155802,MR0146158,MR0174538}.}
      \field{issn}{0065-1036}
      \field{journaltitle}{Acta Arith.}
      \field{title}{On oscillations of certain means formed from the {M}öbius series. {II}}
      \field{volume}{10}
      \field{year}{1964}
      \true{nocite}
      \field{pages}{377\bibrangedash 386}
      \range{pages}{10}
      \verb{doi}
      \verb 10.4064/aa-10-4-377-386
      \endverb
      \verb{mrnumber}
      \verb 0172856
      \endverb
      \verb{urlraw}
      \verb https://doi.org/10.4064/aa-10-4-377-386
      \endverb
      \verb{url}
      \verb https://doi.org/10.4064/aa-10-4-377-386
      \endverb
    \endentry
    \entry{MR0131410}{article}{}
      \name{author}{2}{}{%
        {{hash=d33db82fc1e83dc064a23d93e5a21891}{%
           family={Knapowski},
           familyi={K\bibinitperiod},
           given={S.},
           giveni={S\bibinitperiod}}}%
        {{hash=36befda65f61ea63f097e5518d3ec577}{%
           family={Staś},
           familyi={S\bibinitperiod},
           given={W.},
           giveni={W\bibinitperiod}}}%
      }
      \strng{namehash}{3d6f57664844e9e6d022d0b7cd9c3631}
      \strng{fullhash}{3d6f57664844e9e6d022d0b7cd9c3631}
      \strng{bibnamehash}{3d6f57664844e9e6d022d0b7cd9c3631}
      \strng{authorbibnamehash}{3d6f57664844e9e6d022d0b7cd9c3631}
      \strng{authornamehash}{3d6f57664844e9e6d022d0b7cd9c3631}
      \strng{authorfullhash}{3d6f57664844e9e6d022d0b7cd9c3631}
      \field{sortinit}{K}
      \field{sortinithash}{d3edc18d54b9438a72c24c925bfb38f4}
      \field{labelprefix}{A}
      \field{labelnamesource}{author}
      \field{labeltitlesource}{title}
      \field{annotation}{The authors examine the function $\Delta^{\psi_e}(x)=\sum_{n=1}^{\infty}(\Lambda(n)-1)e^{-n/x}$. They prove, unconditionally and effectively, that $$ \max_{1\le y\le x} |\Delta^{\psi_e}(y)| > x^{1/2} \exp\bigg( - \frac{4\log x\log \log \log x}{\log \log x} \bigg) $$ when $x$ is sufficiently large. The main tool of the proof is a version of the second main theorem of the power-sum method. \par This article cites~\cites{MR1555148,MR0125822,MR0060548}.}
      \field{issn}{0065-1036}
      \field{journaltitle}{Acta Arith.}
      \field{sortyear}{1961}
      \field{title}{A note on a theorem of Hardy and Littlewood}
      \field{volume}{7}
      \field{year}{1961/1962}
      \true{nocite}
      \field{pages}{161\bibrangedash 166}
      \range{pages}{6}
      \verb{doi}
      \verb 10.4064/aa-7-2-161-166
      \endverb
      \verb{mrnumber}
      \verb 0131410
      \endverb
    \endentry
    \entry{MR0146156}{article}{}
      \name{author}{2}{}{%
        {{hash=d33db82fc1e83dc064a23d93e5a21891}{%
           family={Knapowski},
           familyi={K\bibinitperiod},
           given={S.},
           giveni={S\bibinitperiod}}}%
        {{hash=e68e101772b88cb1d48e688164bece92}{%
           family={Turán},
           familyi={T\bibinitperiod},
           given={P.},
           giveni={P\bibinitperiod}}}%
      }
      \strng{namehash}{b6f197e0531d2bbfab6a59c048e3b28d}
      \strng{fullhash}{b6f197e0531d2bbfab6a59c048e3b28d}
      \strng{bibnamehash}{b6f197e0531d2bbfab6a59c048e3b28d}
      \strng{authorbibnamehash}{b6f197e0531d2bbfab6a59c048e3b28d}
      \strng{authornamehash}{b6f197e0531d2bbfab6a59c048e3b28d}
      \strng{authorfullhash}{b6f197e0531d2bbfab6a59c048e3b28d}
      \field{extraname}{1}
      \field{sortinit}{K}
      \field{sortinithash}{d3edc18d54b9438a72c24c925bfb38f4}
      \field{labelprefix}{A}
      \field{labelnamesource}{author}
      \field{labeltitlesource}{title}
      \field{annotation}{The authors start by introducing ten problems of interest in ``comparative prime-number theory'' to the modulus~$k$, the first seven concerning the sign changes and extreme values of $\pi(x;k,\ell_1,\ell_2)$ and the natural density of the solutions to $\pi(x;k,\ell_1,\ell_2)>0$. The eighth problem, which the authors call the ``race-problem of Shanks--Rényi'' (which is perhaps the first time Rényi's name was linked to comparative prime number theory) is whether there are arbitrarily large solutions $x$ to $\pi(x;k,\ell_1) < \cdots < \pi(x;k,\ell_{\phi(k)})$; the last two problems concern the simultaneous inequalities $\pi(x;k,\ell_j) > \frac1{\phi(k)} \li(x)$.\par The authors allude to variants of these ten problems generated by replacing $\pi(x;k,\ell)$ by $\pi_e(x;k,\ell)$ (and, where needed, $\Li(x)$ by $\int_2^{\infty} \frac{e^{-t/x}}{\log{t}} \,dt$), and further vary these problems by replacing $\pi$ with $\psi$ or~$\Pi$. They are aware that such problems could be further varied (``\ldots{} the analogous problems concerning the distribution of primes in binary quadratic forms with fixed discriminant or of the prime ideals of a fixed field in various idealclasses'').\par In Section~4, the authors discuss how some of the problems involving $\psi(x;k,\ell)$ can be conditionally solved using Landau's theorem (and hence unconditionally for moduli dividing $24$). In Section~5, the authors discuss the results they have so far concerning the problems involving~$\pi(x;k,\ell)$, and in Section~8, they briefly discuss what is known about prime number races modulo~$4$. Throughout the rest of this article, the authors introduce the results that will be proved in the next seven articles of the series.\par This article cites~\cites{1853.Chebyshev, 1891.Phragmen, MR1555148, 1930.Polya, 1936.Ingham, MR0067145, MR0108470}.}
      \field{issn}{0001-5954}
      \field{journaltitle}{Acta Math. Acad. Sci. Hungar.}
      \field{title}{Comparative prime-number theory. {I}. {I}ntroduction}
      \field{volume}{13}
      \field{year}{1962}
      \true{nocite}
      \field{pages}{299\bibrangedash 314}
      \range{pages}{16}
      \verb{doi}
      \verb 10.1007/BF02020796
      \endverb
      \verb{mrnumber}
      \verb 0146156
      \endverb
      \verb{urlraw}
      \verb https://doi.org/10.1007/BF02020796
      \endverb
      \verb{url}
      \verb https://doi.org/10.1007/BF02020796
      \endverb
    \endentry
    \entry{MR0146157}{article}{}
      \name{author}{2}{}{%
        {{hash=d33db82fc1e83dc064a23d93e5a21891}{%
           family={Knapowski},
           familyi={K\bibinitperiod},
           given={S.},
           giveni={S\bibinitperiod}}}%
        {{hash=e68e101772b88cb1d48e688164bece92}{%
           family={Turán},
           familyi={T\bibinitperiod},
           given={P.},
           giveni={P\bibinitperiod}}}%
      }
      \strng{namehash}{b6f197e0531d2bbfab6a59c048e3b28d}
      \strng{fullhash}{b6f197e0531d2bbfab6a59c048e3b28d}
      \strng{bibnamehash}{b6f197e0531d2bbfab6a59c048e3b28d}
      \strng{authorbibnamehash}{b6f197e0531d2bbfab6a59c048e3b28d}
      \strng{authornamehash}{b6f197e0531d2bbfab6a59c048e3b28d}
      \strng{authorfullhash}{b6f197e0531d2bbfab6a59c048e3b28d}
      \field{extraname}{2}
      \field{sortinit}{K}
      \field{sortinithash}{d3edc18d54b9438a72c24c925bfb38f4}
      \field{labelprefix}{A}
      \field{labelnamesource}{author}
      \field{labeltitlesource}{title}
      \field{annotation}{This article focuses on the race between the residue class $1\mod k$ and other residue classes $\ell\not\equiv 1 \mod k$, for a fixed modulus~$k$ for which HC holds.\par Fix a character $\chi\mod k$ such that $\chi(\ell)\ne1$, and let $\rho_0=\beta_0+i\gamma_0$ be a zero of $L(s,\chi)$. The authors prove that for $T$ large enough, \begin{align*} \max_{T^{1/3}\le x\le T}\psi(x;k,1,\ell) &> T^{\beta_0}\exp \left(-41\frac{\log T \log\log\log T}{\log \log T}\right)\\ \max_{T^{1/3}\le x\le T}\Pi(x;k,1,\ell) &> T^{\beta_0}\exp \left(-41\frac{\log T \log\log\log T}{\log \log T}\right) \end{align*} and symmetric results for the minimum. As one might expect, the oscillations obtained for $\pi(x;k,1,\ell)$ are worse, but the authors do prove \[ \max_{\exp (\log_3^{1/130}T)\le x\le T}\left(\frac{\log x}{\sqrt x}\right)\pi(x;k,1,\ell)> \frac{1}{100} \log\log\log\log\log T \] and the symmetric result for the minimum. Each of these results yields lower bounds on the corresponding number of sign changes, although for $\pi(x;k,1,\ell)$ the bound is very low---improving this bound is addressed directly in~\cite{MR0146158}. The authors' methods can also compare primes congruent to $1\mod k$ to the average number of primes in other residue classes\mod k: \[ \max_{\exp (\log_3^{1/130}T)\le x\le T}\left(\frac{\log x}{\sqrt x}\right)\bigg(\pi(x;k,1)-\frac1{\phi(k)-1}\sum_{\substack{(\ell,k)=1\\\ell\neq 1}}\pi(x;k,\ell)\bigg)> \frac{1}{100} \log_5 T \] and the symmetric result for the minimum.\par The proofs follow from the application of power-sum method for bounding exponential sums. Siegel's theorem on the existence of zeros in certain rectangles, coupled with the verification of HC for certain moduli up to~$24$, give for these moduli the first unconditional results about the size of the fluctuations of the above functions and the number of their sign changes.\par This article cites~\cites{1914.Littlewood, MR1555148, 1936.Ingham, MR0067145, MR0133309}.}
      \field{issn}{0001-5954}
      \field{journaltitle}{Acta Math. Acad. Sci. Hungar.}
      \field{title}{Comparative prime-number theory. {II}. {C}omparison of the progressions {$\equiv 1$} {${\rm mod}\ k$} and {$\equiv l$} {${\rm mod}\ k,\,l\not\equiv 1$} {${\rm mod}\ k$}}
      \field{volume}{13}
      \field{year}{1962}
      \true{nocite}
      \field{pages}{315\bibrangedash 342}
      \range{pages}{28}
      \verb{doi}
      \verb 10.1007/BF02020797
      \endverb
      \verb{mrnumber}
      \verb 0146157
      \endverb
      \verb{urlraw}
      \verb https://doi.org/10.1007/BF02020797
      \endverb
      \verb{url}
      \verb https://doi.org/10.1007/BF02020797
      \endverb
    \endentry
    \entry{MR0146158}{article}{}
      \name{author}{2}{}{%
        {{hash=d33db82fc1e83dc064a23d93e5a21891}{%
           family={Knapowski},
           familyi={K\bibinitperiod},
           given={S.},
           giveni={S\bibinitperiod}}}%
        {{hash=e68e101772b88cb1d48e688164bece92}{%
           family={Turán},
           familyi={T\bibinitperiod},
           given={P.},
           giveni={P\bibinitperiod}}}%
      }
      \strng{namehash}{b6f197e0531d2bbfab6a59c048e3b28d}
      \strng{fullhash}{b6f197e0531d2bbfab6a59c048e3b28d}
      \strng{bibnamehash}{b6f197e0531d2bbfab6a59c048e3b28d}
      \strng{authorbibnamehash}{b6f197e0531d2bbfab6a59c048e3b28d}
      \strng{authornamehash}{b6f197e0531d2bbfab6a59c048e3b28d}
      \strng{authorfullhash}{b6f197e0531d2bbfab6a59c048e3b28d}
      \field{extraname}{3}
      \field{sortinit}{K}
      \field{sortinithash}{d3edc18d54b9438a72c24c925bfb38f4}
      \field{labelprefix}{A}
      \field{labelnamesource}{author}
      \field{labeltitlesource}{title}
      \field{annotation}{The authors continue their comparison of $\pi(x;k,\ell)$ and $\pi(x;k,1)$, where $\ell\not\equiv1\mod k$ and $(\ell,k)=1$, assuming HC for the modulus~$k$. They show that $W_{k;\ell,1}(T) > k^{-c}\log\log\log\log T$ for sufficiently large~$T$, where $c$ is an absolute effective constant. The proof technique involves Dirichlet's box principle and bounds obtained on $\Pi(x;k,\ell,1)$ in~\cite{MR0146157}.\par When $\ell$ is a quadratic residue\mod k, they also show that if $\rho_0=\beta_0+i\gamma_0$ is a zero of $L(s,\chi)$ for some character $\chi$ such that $\chi(\ell)\neq 1$, then for sufficiently large~$T$, \[ \max_{T^{1/3}\leq x \leq T} \pi(x;k,\ell,1) >T^{\beta_0} \exp\left(-42 \frac{\log T\log \log \log T}{\log \log T}\right),\] and a similar statement holds for the minimum; consequently, for such $k$ and $\ell$, the inequality $W_{k;\ell,1}(T) > \frac1{\log 3} \log\log T + O(1)$ holds for sufficiently large~$T$. The proof involves power-sum method for bounds on exponential sums.\par The authors remark that both theorems hold as well for $W\big(\Delta^\pi(x;k,1);T\big)$ and $W\big(\Deltadot^\pi(x;k,1);T\big)$.\par This article cites~\cites{MR1555148,1936.Ingham,MR0146156,MR0146157,MR1573970}.}
      \field{issn}{0001-5954}
      \field{journaltitle}{Acta Math. Acad. Sci. Hungar.}
      \field{title}{Comparative prime-number theory. {III}. {C}ontinuation of the study of comparison of the progressions {$\equiv 1$} {${\rm mod}\ k$} and {$\equiv l$} {${\rm mod}\ k$}}
      \field{volume}{13}
      \field{year}{1962}
      \true{nocite}
      \field{pages}{343\bibrangedash 364}
      \range{pages}{22}
      \verb{doi}
      \verb 10.1007/BF02020798
      \endverb
      \verb{mrnumber}
      \verb 0146158
      \endverb
      \verb{urlraw}
      \verb https://doi.org/10.1007/BF02020798
      \endverb
      \verb{url}
      \verb https://doi.org/10.1007/BF02020798
      \endverb
    \endentry
    \entry{MR0146159}{article}{}
      \name{author}{2}{}{%
        {{hash=d33db82fc1e83dc064a23d93e5a21891}{%
           family={Knapowski},
           familyi={K\bibinitperiod},
           given={S.},
           giveni={S\bibinitperiod}}}%
        {{hash=e68e101772b88cb1d48e688164bece92}{%
           family={Turán},
           familyi={T\bibinitperiod},
           given={P.},
           giveni={P\bibinitperiod}}}%
      }
      \strng{namehash}{b6f197e0531d2bbfab6a59c048e3b28d}
      \strng{fullhash}{b6f197e0531d2bbfab6a59c048e3b28d}
      \strng{bibnamehash}{b6f197e0531d2bbfab6a59c048e3b28d}
      \strng{authorbibnamehash}{b6f197e0531d2bbfab6a59c048e3b28d}
      \strng{authornamehash}{b6f197e0531d2bbfab6a59c048e3b28d}
      \strng{authorfullhash}{b6f197e0531d2bbfab6a59c048e3b28d}
      \field{extraname}{4}
      \field{sortinit}{K}
      \field{sortinithash}{d3edc18d54b9438a72c24c925bfb38f4}
      \field{labelprefix}{A}
      \field{labelnamesource}{author}
      \field{labeltitlesource}{title}
      \field{annotation}{The authors apply techniques from earlier in this series of articles to the modulus $k=8$, when $\ell_1,\ell_2 \in \{3,5,7\}$ are distinct quadratic nonresidues. They show that \[ \max_{T^{1/3}\leq x \leq T} \pi(x;8,\ell_1,\ell_2) > \sqrt T \exp\left(-23 \frac{\log T\log \log \log T}{\log \log T}\right),\] and similarly for $\Pi(x; 8, \ell_1,\ell_2)$ and $\psi(x; 8, \ell_1, \ell_2)$. Since $\ell_1$ and $\ell_2$ can be interchanged, this result implies the inequality $W_{8;\ell_1,\ell_2}(T) > \frac1{\log 3} \log\log T + O(1)$, and similarly for $W^\Pi_{8;\ell_1,\ell_2}(T)$ and $W^\psi_{8;\ell_1,\ell_2}(T)$.\par In the third section, the authors remark on the modulus $k=5$. The case $(\ell_1, \ell_2)=(\ell,1)$ has already been handled earlier in this series; the authors mention that the case $(\ell_1, \ell_2)=(2,3)$, where both are quadratic nonresidues, can be handled in a similar way to the $k=8$ cases treated in this article. The remaining cases have $\ell_1=4$, a quadratic residue not equal to $1\mod5$, and $\ell_2\in\{2,3\}$, and these cases yield ``an unpleasant (or pleasant?)~surprise'': the authors cannot establish sign changes for $\pi(x;5, 4, 2)$ and $\pi(x;5, 4, 3)$ even assuming GRH (although the methods do work for the $\Pi$ and $\psi$ versions, a situation the authors discuss further in the later articles of this series).\par This article cites~\cites{MR0146156,MR0146157, MR0146158}.}
      \field{issn}{0001-5954}
      \field{journaltitle}{Acta Math. Acad. Sci. Hungar.}
      \field{title}{Comparative prime-number theory. {IV}. {P}aradigma to the general case, {$k=8$} and {$5$}}
      \field{volume}{14}
      \field{year}{1963}
      \true{nocite}
      \field{pages}{31\bibrangedash 42}
      \range{pages}{12}
      \verb{doi}
      \verb 10.1007/BF01901928
      \endverb
      \verb{mrnumber}
      \verb 0146159
      \endverb
      \verb{urlraw}
      \verb https://doi.org/10.1007/BF01901928
      \endverb
      \verb{url}
      \verb https://doi.org/10.1007/BF01901928
      \endverb
    \endentry
    \entry{MR0146160}{article}{}
      \name{author}{2}{}{%
        {{hash=d33db82fc1e83dc064a23d93e5a21891}{%
           family={Knapowski},
           familyi={K\bibinitperiod},
           given={S.},
           giveni={S\bibinitperiod}}}%
        {{hash=e68e101772b88cb1d48e688164bece92}{%
           family={Turán},
           familyi={T\bibinitperiod},
           given={P.},
           giveni={P\bibinitperiod}}}%
      }
      \strng{namehash}{b6f197e0531d2bbfab6a59c048e3b28d}
      \strng{fullhash}{b6f197e0531d2bbfab6a59c048e3b28d}
      \strng{bibnamehash}{b6f197e0531d2bbfab6a59c048e3b28d}
      \strng{authorbibnamehash}{b6f197e0531d2bbfab6a59c048e3b28d}
      \strng{authornamehash}{b6f197e0531d2bbfab6a59c048e3b28d}
      \strng{authorfullhash}{b6f197e0531d2bbfab6a59c048e3b28d}
      \field{extraname}{5}
      \field{sortinit}{K}
      \field{sortinithash}{d3edc18d54b9438a72c24c925bfb38f4}
      \field{labelprefix}{A}
      \field{labelnamesource}{author}
      \field{labeltitlesource}{title}
      \field{annotation}{Under the assumption of a ``finite Riemann--Piltz conjecture'' GRH$(H,E_k)$, the authors establish the following result for any distinct reduced residues $\ell_1,\ell_2\mod k$ and for $T$ sufficiently large (explicitly quantified in the article): \[ \max_{T^{1/3} \leq x \leq T } \Pi(x;k,\ell_1,\ell_2) > \sqrt{T}\exp \bigg( {-44\frac{\log{T}\log\log\log{T}}{\log\log{T}}} \bigg),\] and the same with $\Pi$ replaced by~$\psi$. A central element to their proof is the estimation of the integral \begin{align*} J(T) &=-\frac{1}{2\pi i}\int_2 \left(\frac{e^{y_1s}}{s}\right)^v\frac{(\omega_0L_1^{v_0})^s}{s^{v_0+1}}\cdot \frac{1}{\phi(k)} \bigg\{\sum_{\substack{\chi\mod k \\ \chi(\ell_1)\neq \chi(\ell_2)}}\left(\overline{\chi}(\ell_1)-\overline{\chi}(\ell_2)\right)\frac{L'}{L}(s,\chi)\bigg\}\ ds \\ &= \frac{1}{(v+v_0)!}\int_1^{Y_1} \Pi(x;k,\ell_1,\ell_2) \frac d{dx} \left( \bigg( \log\frac{Y_1}{x} \bigg)^{v+v_0} \log x \right) \, dx. \end{align*} Their approach involves an application of power-sum method similar to what appears in the previous articles of the series. The authors note that their main theorem gives similar bounds on $\Delta^\Pi(x;k,\ell)/\phi(k)$ and $\Delta^\psi(x;k,\ell)/\phi(k)$. In particular, this result implies the lower bounds $W^\Pi_{k;\ell_1,\ell_2}(T) > \frac1{\log 3}\log\log T + O(1)$ and $W^\psi_{k;\ell_1,\ell_2}(T) > \frac1{\log 3}\log\log T + O(1)$.\par This article cites~\cites{MR0125822,MR0146156,MR0146157,MR0146158,MR0146159}.}
      \field{issn}{0001-5954}
      \field{journaltitle}{Acta Math. Acad. Sci. Hungar.}
      \field{title}{Comparative prime-number theory. {V}. {S}ome theorems concerning the general case}
      \field{volume}{14}
      \field{year}{1963}
      \true{nocite}
      \field{pages}{43\bibrangedash 63}
      \range{pages}{21}
      \verb{doi}
      \verb 10.1007/BF01901929
      \endverb
      \verb{mrnumber}
      \verb 0146160
      \endverb
      \verb{urlraw}
      \verb https://doi.org/10.1007/BF01901929
      \endverb
      \verb{url}
      \verb https://doi.org/10.1007/BF01901929
      \endverb
    \endentry
    \entry{MR0146161}{article}{}
      \name{author}{2}{}{%
        {{hash=d33db82fc1e83dc064a23d93e5a21891}{%
           family={Knapowski},
           familyi={K\bibinitperiod},
           given={S.},
           giveni={S\bibinitperiod}}}%
        {{hash=e68e101772b88cb1d48e688164bece92}{%
           family={Turán},
           familyi={T\bibinitperiod},
           given={P.},
           giveni={P\bibinitperiod}}}%
      }
      \strng{namehash}{b6f197e0531d2bbfab6a59c048e3b28d}
      \strng{fullhash}{b6f197e0531d2bbfab6a59c048e3b28d}
      \strng{bibnamehash}{b6f197e0531d2bbfab6a59c048e3b28d}
      \strng{authorbibnamehash}{b6f197e0531d2bbfab6a59c048e3b28d}
      \strng{authornamehash}{b6f197e0531d2bbfab6a59c048e3b28d}
      \strng{authorfullhash}{b6f197e0531d2bbfab6a59c048e3b28d}
      \field{extraname}{6}
      \field{sortinit}{K}
      \field{sortinithash}{d3edc18d54b9438a72c24c925bfb38f4}
      \field{labelprefix}{A}
      \field{labelnamesource}{author}
      \field{labeltitlesource}{title}
      \field{annotation}{Under the assumption of a ``finite Riemann--Piltz conjecture'' GRH$(H,E_k)$, the authors establish the following result in the case that $\ell_1$ and $\ell_2$ are either both quadratic residues or both quadratic nonresidues$\mod k$: when $T$ is sufficiently large in terms of $k$ (the authors give an explicit lower bound), the inequalities \begin{align*} \max_{T^{1/3}\leq x \leq T} \pi(x,k,\ell_1,\ell_2) &> \sqrt{T}\exp \bigg( {-44\frac{\log{T}\log\log\log{T}}{\log\log{T}}} \bigg) \\ \min_{T^{1/3}\leq x \leq T} \pi(x,k,\ell_1,\ell_2) &< -\sqrt{T}\exp \bigg( {-44\frac{\log{T}\log\log\log{T}}{\log\log{T}}} \bigg) \end{align*} hold. As usual, this result implies the lower bound $W_{k;\ell_1,\ell_2}(T) > \frac1{\log 3}\log\log T + O(1)$. The authors rely on multiple lemmas from their use of power-sum method in previous articles of this series. \par This article cites~\cites{MR0146157, MR0146158, MR0146160}.}
      \field{issn}{0001-5954}
      \field{journaltitle}{Acta Math. Acad. Sci. Hungar.}
      \field{title}{Comparative prime-number theory. {VI}. {C}ontinuation of the general case}
      \field{volume}{14}
      \field{year}{1963}
      \true{nocite}
      \field{pages}{65\bibrangedash 78}
      \range{pages}{14}
      \verb{doi}
      \verb 10.1007/BF01901930
      \endverb
      \verb{mrnumber}
      \verb 0146161
      \endverb
      \verb{urlraw}
      \verb https://doi.org/10.1007/BF01901930
      \endverb
      \verb{url}
      \verb https://doi.org/10.1007/BF01901930
      \endverb
    \endentry
    \entry{MR0156826}{article}{}
      \name{author}{2}{}{%
        {{hash=d33db82fc1e83dc064a23d93e5a21891}{%
           family={Knapowski},
           familyi={K\bibinitperiod},
           given={S.},
           giveni={S\bibinitperiod}}}%
        {{hash=e68e101772b88cb1d48e688164bece92}{%
           family={Turán},
           familyi={T\bibinitperiod},
           given={P.},
           giveni={P\bibinitperiod}}}%
      }
      \strng{namehash}{b6f197e0531d2bbfab6a59c048e3b28d}
      \strng{fullhash}{b6f197e0531d2bbfab6a59c048e3b28d}
      \strng{bibnamehash}{b6f197e0531d2bbfab6a59c048e3b28d}
      \strng{authorbibnamehash}{b6f197e0531d2bbfab6a59c048e3b28d}
      \strng{authornamehash}{b6f197e0531d2bbfab6a59c048e3b28d}
      \strng{authorfullhash}{b6f197e0531d2bbfab6a59c048e3b28d}
      \field{extraname}{7}
      \field{sortinit}{K}
      \field{sortinithash}{d3edc18d54b9438a72c24c925bfb38f4}
      \field{labelprefix}{A}
      \field{labelnamesource}{author}
      \field{labeltitlesource}{title}
      \field{annotation}{This article gives a general conditional proof that $\psi(x; k, \ell_1,\ell_2)$ changes sign infinitely often. The authors show, assuming HC$(E_k)$ for the modulus $k$ for some constant $0<A_k\le1$, that there exists a positive constant $c$ such that $\psi(x; k, \ell_1,\ell_2)$ changes sign in every interval of the form $\omega \leq x \leq \exp(2\sqrt{\omega})$ as long as $$ \omega \geq \max \big\{ e^{k^c},e^{2/{E(k)^3}} \big\}. $$ This result immediately implies results for the first sign change of $\psi(x; k, \ell_1,\ell_2)$ and for its number of sign changes.\par This article cites~\cites{MR0146156, MR0146157, MR0146158, MR0146160, MR0146161}.}
      \field{issn}{0001-5954}
      \field{journaltitle}{Acta Math. Acad. Sci. Hungar}
      \field{title}{Comparative prime-number theory. {VII}. {T}he problem of sign-changes in the general case}
      \field{volume}{14}
      \field{year}{1963}
      \true{nocite}
      \field{pages}{241\bibrangedash 250}
      \range{pages}{10}
      \verb{doi}
      \verb 10.1007/BF01895712
      \endverb
      \verb{mrnumber}
      \verb 0156826
      \endverb
      \verb{urlraw}
      \verb https://doi.org/10.1007/BF01895712
      \endverb
      \verb{url}
      \verb https://doi.org/10.1007/BF01895712
      \endverb
    \endentry
    \entry{MR0156827}{article}{}
      \name{author}{2}{}{%
        {{hash=d33db82fc1e83dc064a23d93e5a21891}{%
           family={Knapowski},
           familyi={K\bibinitperiod},
           given={S.},
           giveni={S\bibinitperiod}}}%
        {{hash=e68e101772b88cb1d48e688164bece92}{%
           family={Turán},
           familyi={T\bibinitperiod},
           given={P.},
           giveni={P\bibinitperiod}}}%
      }
      \strng{namehash}{b6f197e0531d2bbfab6a59c048e3b28d}
      \strng{fullhash}{b6f197e0531d2bbfab6a59c048e3b28d}
      \strng{bibnamehash}{b6f197e0531d2bbfab6a59c048e3b28d}
      \strng{authorbibnamehash}{b6f197e0531d2bbfab6a59c048e3b28d}
      \strng{authornamehash}{b6f197e0531d2bbfab6a59c048e3b28d}
      \strng{authorfullhash}{b6f197e0531d2bbfab6a59c048e3b28d}
      \field{extraname}{8}
      \field{sortinit}{K}
      \field{sortinithash}{d3edc18d54b9438a72c24c925bfb38f4}
      \field{labelprefix}{A}
      \field{labelnamesource}{author}
      \field{labeltitlesource}{title}
      \field{annotation}{Hardy--Littlewood and Landau had already shown that the assertion $\lim_{x\to\infty} \pi_e(x;4,1,3) = -\infty$ is equivalent to GRH for $L(s,\chi_{-4})$. In this article the authors obtain an analogous equivalence concerning the races between $1$ and a nonsquare\mod 8: slightly modifying the arguments for the\mod4 case, they show that the assertion $\lim_{x\to\infty} \theta_e(x;8,1,\ell) = -\infty$ for all $\ell\not\equiv1\mod8$ is equivalent to GRH for the three nonprincipal Dirichlet $L$-functions\mod8, and the same for the assertion $\lim_{x\to\infty} \pi_e(x;8,1,\ell) = -\infty$.\par They further show that the race between two nonsquares switches infinitely often---more precisely, for $\ell_1\not\equiv \ell_2\not\equiv 1\mod8$, they unconditionally show that when $T$ is large enough, \[ \max_{T^{1/3}\le x\le T}\theta_e(x;8,\ell_1,\ell_2)>\sqrt T \exp\left(-22\frac{\log T\log\log\log T}{\log\log T}\right). \] They indicate that this result is ``deeper'', and in particular that they cannot yet replace $\theta_e$ with $\pi_e$ in this result.\par The proofs rely on the power-sum method, as well as some explicit numerical data for the low-lying zeros of the $L$-functions\mod 8. \par This article cites~\cites{MR1555148, MR1544293A,MR1544293B}.}
      \field{issn}{0001-5954}
      \field{journaltitle}{Acta Math. Acad. Sci. Hungar}
      \field{title}{Comparative prime-number theory. {VIII}. {C}hebyshev's problem for {$k=8$}}
      \field{volume}{14}
      \field{year}{1963}
      \true{nocite}
      \field{pages}{251\bibrangedash 268}
      \range{pages}{18}
      \verb{doi}
      \verb 10.1007/BF01895713
      \endverb
      \verb{mrnumber}
      \verb 0156827
      \endverb
      \verb{urlraw}
      \verb https://doi.org/10.1007/BF01895713
      \endverb
      \verb{url}
      \verb https://doi.org/10.1007/BF01895713
      \endverb
    \endentry
    \entry{MR0162771}{article}{}
      \name{author}{2}{}{%
        {{hash=d33db82fc1e83dc064a23d93e5a21891}{%
           family={Knapowski},
           familyi={K\bibinitperiod},
           given={S.},
           giveni={S\bibinitperiod}}}%
        {{hash=e68e101772b88cb1d48e688164bece92}{%
           family={Turán},
           familyi={T\bibinitperiod},
           given={P.},
           giveni={P\bibinitperiod}}}%
      }
      \strng{namehash}{b6f197e0531d2bbfab6a59c048e3b28d}
      \strng{fullhash}{b6f197e0531d2bbfab6a59c048e3b28d}
      \strng{bibnamehash}{b6f197e0531d2bbfab6a59c048e3b28d}
      \strng{authorbibnamehash}{b6f197e0531d2bbfab6a59c048e3b28d}
      \strng{authornamehash}{b6f197e0531d2bbfab6a59c048e3b28d}
      \strng{authorfullhash}{b6f197e0531d2bbfab6a59c048e3b28d}
      \field{extraname}{9}
      \field{sortinit}{K}
      \field{sortinithash}{d3edc18d54b9438a72c24c925bfb38f4}
      \field{labelprefix}{A}
      \field{labelnamesource}{author}
      \field{labeltitlesource}{title}
      \field{annotation}{The first two sections offer a short summary of comparative prime number theory up to 1964. The authors classify the subject into 48 separate problems over 12 categories (and additional variants on these), which is of interest to those interested in the history of the field. They then move on to prove results about ``strongly localized accumulation problems''.\par Most generally, assuming HC for the modulus~$k$, they show that if $T$ is sufficiently large in terms of $k$, then for any $(\ell,k) = 1$ with $\ell\not\equiv1\mod k$, \[ \max_I \psi\big( I; k, \ell, 1 \big) > \sqrt{T}e^{-\log^{11/12} T} \quad\text{and}\quad \min_I \psi\big( I; k, \ell, 1 \big) < -\sqrt{T}e^{-\log^{11/12} T}, \] where the maximum and minimum are taken over all subintervals $I$ of $\big[ Te^{-\log^{11/12} T} , T \big]$. The central argument involves the evaluation of the integral \[ \frac 1{2\pi i} \int e^{As}\bigg(\frac{e^{Bs}-e^{-Bs}}{2Bs}\bigg)^r \bigg( \frac1{\phi(k)} \sum_{\chi\mod k} \overline{\chi}(\ell) \frac{L'}L(s,\chi) \bigg) \, ds \] for positive constants $A$ and~$B$.\par This article cites~\cites{1853.Chebyshev, MR1555148, MR1544293A, MR1544293B, MR0060548, MR0133309, MR0146156, MR0146157, MR0146158, MR0146159, MR0146160, MR0146161, MR0156826, MR0156827}.}
      \field{issn}{0065-1036}
      \field{journaltitle}{Acta Arith.}
      \field{title}{Further developments in the comparative prime-number theory. {I}}
      \field{volume}{9}
      \field{year}{1964}
      \true{nocite}
      \field{pages}{23\bibrangedash 40}
      \range{pages}{18}
      \verb{doi}
      \verb 10.4064/aa-9-1-23-40
      \endverb
      \verb{mrnumber}
      \verb 0162771
      \endverb
      \verb{urlraw}
      \verb https://doi.org/10.4064/aa-9-1-23-40
      \endverb
      \verb{url}
      \verb https://doi.org/10.4064/aa-9-1-23-40
      \endverb
    \endentry
    \entry{MR0174538}{article}{}
      \name{author}{2}{}{%
        {{hash=d33db82fc1e83dc064a23d93e5a21891}{%
           family={Knapowski},
           familyi={K\bibinitperiod},
           given={S.},
           giveni={S\bibinitperiod}}}%
        {{hash=e68e101772b88cb1d48e688164bece92}{%
           family={Turán},
           familyi={T\bibinitperiod},
           given={P.},
           giveni={P\bibinitperiod}}}%
      }
      \strng{namehash}{b6f197e0531d2bbfab6a59c048e3b28d}
      \strng{fullhash}{b6f197e0531d2bbfab6a59c048e3b28d}
      \strng{bibnamehash}{b6f197e0531d2bbfab6a59c048e3b28d}
      \strng{authorbibnamehash}{b6f197e0531d2bbfab6a59c048e3b28d}
      \strng{authornamehash}{b6f197e0531d2bbfab6a59c048e3b28d}
      \strng{authorfullhash}{b6f197e0531d2bbfab6a59c048e3b28d}
      \field{extraname}{10}
      \field{sortinit}{K}
      \field{sortinithash}{d3edc18d54b9438a72c24c925bfb38f4}
      \field{labelprefix}{A}
      \field{labelnamesource}{author}
      \field{labeltitlesource}{title}
      \field{annotation}{The authors establish several theorems, all assuming HC$(E_k)$ with $E_k \ll \sqrt{\log k}/k$; most of their results concern the function $\theta_l(x,r;k,\ell_1,\ell_2)$ where $r=r(x,k)$ satisfies $\frac{\log k}{E_k} \ll r \le \log x$. Their most general result (Theorem VI) is that for any quadratic nonresidue $\ell_1\mod k$ and quadratic residue $\ell_2\mod k$, if $L(s,\chi)$ satisfies GRH for all characters $\chi\mod k$ such that $\chi(\ell_1) \ne \chi(\ell_2)$, then $\theta_l(x,r;k,\ell_1,\ell_2) \gg \sqrt{x}$ for $x$ sufficiently large. \par They also establish the following result. Let $\ell$ be a quadratic nonresidue\mod k, and suppose that there exists a character $\chi\mod k$ with $\chi(\ell)\ne1$ such that $L(s,\chi)$ has a zero $\rho_0=\beta_0+i\gamma_0$ with $\beta_0>\frac12$. Then for all $T$ (sufficiently large in terms of $k$ and $\rho_0$), there exist subintervals $I_\pm$ of $\big[ Te^{-5\log^{20/21}T} , Te^{5\log^{20/21}T} \big]$ such that $\pi(I_+;k,\ell,1) > T^{\beta_0} \exp\big( -(2+\gamma_0^2)(\log T)^{5/7} \big) $ and $ \pi(I_-;k,\ell,1) < - T^{\beta_0} \exp\big( -(2+\gamma_0^2)(\log T)^{5/7} \big)$. They deduce this theorem from an analogous theorem involving $\theta_l(x,r;k,\ell,1)$, which serves as a sort of inverse to Theorem VI. \par Together, these results imply, given a quadratic nonresidue $\ell\mod k$, that $\lim_{x\to\infty} \theta_l(x,r;k,\ell,1) = +\infty$ holds if and only if $L(s,\chi)$ satisfies GRH for all characters $\chi \mod k$ with $\chi(\ell)\ne1$. It follows that $\lim_{x\to\infty} \theta_l(x,r;k,\ell,1) = +\infty$ holds for all quadratic nonresidues\mod k if and only if $L(s,\chi)$ satisfies GRH for every nonprincipal character $\chi \mod k$. \par The authors also make some remarks about races between residue classes to different moduli, showing for example how the race between $\pi(x;3,1)$ and $\pi(x;4,1)$ reduces to a race between residue classes modulo~$12$, to which their results apply. \par This article cites~\cites{1853.Chebyshev, MR1555148, MR1544293A, MR0060548, MR0133309, MR0146157, MR0146159, MR0156827}.}
      \field{issn}{0065-1036}
      \field{journaltitle}{Acta Arith.}
      \field{title}{Further developments in the comparative prime-number theory. {II}. {A} modification of {C}hebyshev's assertion}
      \field{volume}{10}
      \field{year}{1964}
      \true{nocite}
      \field{pages}{293\bibrangedash 313}
      \range{pages}{21}
      \verb{doi}
      \verb 10.4064/aa-10-3-293-313
      \endverb
      \verb{mrnumber}
      \verb 0174538
      \endverb
      \verb{urlraw}
      \verb https://doi.org/10.4064/aa-10-3-293-313
      \endverb
      \verb{url}
      \verb https://doi.org/10.4064/aa-10-3-293-313
      \endverb
    \endentry
    \entry{MR0180539}{article}{}
      \name{author}{2}{}{%
        {{hash=d33db82fc1e83dc064a23d93e5a21891}{%
           family={Knapowski},
           familyi={K\bibinitperiod},
           given={S.},
           giveni={S\bibinitperiod}}}%
        {{hash=e68e101772b88cb1d48e688164bece92}{%
           family={Turán},
           familyi={T\bibinitperiod},
           given={P.},
           giveni={P\bibinitperiod}}}%
      }
      \strng{namehash}{b6f197e0531d2bbfab6a59c048e3b28d}
      \strng{fullhash}{b6f197e0531d2bbfab6a59c048e3b28d}
      \strng{bibnamehash}{b6f197e0531d2bbfab6a59c048e3b28d}
      \strng{authorbibnamehash}{b6f197e0531d2bbfab6a59c048e3b28d}
      \strng{authornamehash}{b6f197e0531d2bbfab6a59c048e3b28d}
      \strng{authorfullhash}{b6f197e0531d2bbfab6a59c048e3b28d}
      \field{extraname}{11}
      \field{sortinit}{K}
      \field{sortinithash}{d3edc18d54b9438a72c24c925bfb38f4}
      \field{labelprefix}{A}
      \field{labelnamesource}{author}
      \field{labeltitlesource}{title}
      \field{annotation}{This article concerns the weighted function $\theta_l(x,r;k,\ell,1)$, for a modulus $k$ satisfying HC and a quadratic residue $\ell\not\equiv1\mod k$. Let $\beta_0$ be the real part of any zero of an $L(s,\chi)$ where $\chi \mod k$ a character such that $\chi(\ell) \ne 1$. The authors exhibit extreme values of $\theta_l(x,r;k,\ell,1)$ where $x$ is near $T$ and $r$ is near $(\log T)^{2/3}$; more precisely, there exists a positive constant $c$ such that for $T$ sufficiently large, there exist $x_1,x_2 \in (Te^{-(\log T)^{5/6}}, Te^{(\log T)^{11/15}})$ such that for suitable $r_1,r_2 \in [(2\log T)^{2/3}, (2\log T)^{2/3} + (2\log T)^{2/5}]$, \[ \theta_l(x_1,r_1;k,\ell,1) > T^{\beta_0} e^{-c(\log T)^{5/6}} \quad\text{and}\quad \theta_l(x_2,r_2;k,\ell,1) < -T^{\beta_0} e^{-c(\log T)^{5/6}}. \] The authors then state that using the methods of their prior article~\cite{MR0174538}, it follows that for $T$ sufficiently large, there exist closed subintervals $I,J \subseteq [Te^{-(\log T)^{6/7}}, Te^{(\log T)^{6/7}}]$ such that one has the ``strongly localized accumulations'' \[ \pi(I;k,\ell,1) > \sqrt{T}e^{-c(\log T)^{5/6}} \quad\text{and}\quad \pi(J;k,\ell,1) < -\sqrt{T}e^{-c(\log T)^{5/6}}. \] \par This article cites~\cites{1853.Chebyshev, MR1555148, MR0156827, MR0174538, MR0146157, MR1544293A, MR1544293B, MR0060548}.}
      \field{issn}{0065-1036}
      \field{journaltitle}{Acta Arith.}
      \field{title}{Further developments in the comparative prime-number theory. {III}}
      \field{volume}{11}
      \field{year}{1965}
      \true{nocite}
      \field{pages}{115\bibrangedash 127}
      \range{pages}{13}
      \verb{doi}
      \verb 10.4064/aa-11-1-115-127
      \endverb
      \verb{mrnumber}
      \verb 0180539
      \endverb
      \verb{urlraw}
      \verb https://doi.org/10.4064/aa-11-1-115-127
      \endverb
      \verb{url}
      \verb https://doi.org/10.4064/aa-11-1-115-127
      \endverb
    \endentry
    \entry{MR0182616A}{article}{}
      \name{author}{2}{}{%
        {{hash=d33db82fc1e83dc064a23d93e5a21891}{%
           family={Knapowski},
           familyi={K\bibinitperiod},
           given={S.},
           giveni={S\bibinitperiod}}}%
        {{hash=e68e101772b88cb1d48e688164bece92}{%
           family={Turán},
           familyi={T\bibinitperiod},
           given={P.},
           giveni={P\bibinitperiod}}}%
      }
      \strng{namehash}{b6f197e0531d2bbfab6a59c048e3b28d}
      \strng{fullhash}{b6f197e0531d2bbfab6a59c048e3b28d}
      \strng{bibnamehash}{b6f197e0531d2bbfab6a59c048e3b28d}
      \strng{authorbibnamehash}{b6f197e0531d2bbfab6a59c048e3b28d}
      \strng{authornamehash}{b6f197e0531d2bbfab6a59c048e3b28d}
      \strng{authorfullhash}{b6f197e0531d2bbfab6a59c048e3b28d}
      \field{extraname}{12}
      \field{sortinit}{K}
      \field{sortinithash}{d3edc18d54b9438a72c24c925bfb38f4}
      \field{labelprefix}{A}
      \field{labelnamesource}{author}
      \field{labeltitlesource}{title}
      \field{annotation}{Let $\ell_1$ and $\ell_2$ be quadratic nonresidues modulo a sufficiently large modulus~$k$. Let $\eta$ be sufficiently small in terms of~$k$, and suppose that the Dirichlet $L$-functions\mod k satisfy GRH$(2/\sqrt\eta,E_k)$ for suitable $E_k$. Then, when $T$ is sufficiently large in terms of $k$ and~$\eta$, there exist $x_+$ and $x_-$ in the interval $[T^{1-\sqrt{\eta}} , Te^{\log^{3/4}{T}}]$, and $\eta_1$ and $\eta_2$ in the interval $[2\eta \log{T} , 2\eta \log{T}+\sqrt{\log{T}}]$, such that $\theta_l(x_+,v_+;k,\ell_1,\ell_2) > T^{1/2 - 4\sqrt{\eta}}$ and $\theta_l(x_-,v_-;k,\ell_1,\ell_2) < - T^{1/2 - 4\sqrt{\eta}}$. Furthermore, under the same assumptions, there exist subintervals $I_+$, $I_-$ of $[T^{1-4\sqrt{\eta}},T^{1+4\sqrt{\eta}}]$ such that $\pi(I_+; k, \ell_1, \ell_2) >T^{1/2 - 5\sqrt{\eta}}$ and $\pi(I_-; k, \ell_1, \ell_2) < -T^{1/2 - 5\sqrt{\eta}}$. \par This article cites~\cites{MR0060548, MR0146158, MR0146160, MR0174538}.}
      \field{issn}{0065-1036}
      \field{journaltitle}{Acta Arith. 11 (1965), 147-161; ibid.}
      \field{title}{Further developments in the comparative prime-number theory. {IV}}
      \field{volume}{11}
      \field{year}{1965}
      \true{nocite}
      \field{pages}{147\bibrangedash 161}
      \range{pages}{15}
      \verb{doi}
      \verb 10.4064/aa-11-2-193-202
      \endverb
      \verb{mrnumber}
      \verb 0182616
      \endverb
      \verb{urlraw}
      \verb https://doi.org/10.4064/aa-11-2-193-202
      \endverb
      \verb{url}
      \verb https://doi.org/10.4064/aa-11-2-193-202
      \endverb
    \endentry
    \entry{MR0182616B}{article}{}
      \name{author}{2}{}{%
        {{hash=d33db82fc1e83dc064a23d93e5a21891}{%
           family={Knapowski},
           familyi={K\bibinitperiod},
           given={S.},
           giveni={S\bibinitperiod}}}%
        {{hash=e68e101772b88cb1d48e688164bece92}{%
           family={Turán},
           familyi={T\bibinitperiod},
           given={P.},
           giveni={P\bibinitperiod}}}%
      }
      \strng{namehash}{b6f197e0531d2bbfab6a59c048e3b28d}
      \strng{fullhash}{b6f197e0531d2bbfab6a59c048e3b28d}
      \strng{bibnamehash}{b6f197e0531d2bbfab6a59c048e3b28d}
      \strng{authorbibnamehash}{b6f197e0531d2bbfab6a59c048e3b28d}
      \strng{authornamehash}{b6f197e0531d2bbfab6a59c048e3b28d}
      \strng{authorfullhash}{b6f197e0531d2bbfab6a59c048e3b28d}
      \field{extraname}{13}
      \field{sortinit}{K}
      \field{sortinithash}{d3edc18d54b9438a72c24c925bfb38f4}
      \field{labelprefix}{A}
      \field{labelnamesource}{author}
      \field{labeltitlesource}{title}
      \field{annotation}{This short article is distinct among the second series by Knapowski and Turán, in that, rather than making use of ``one-sided theorems'' of the power-sum method, it uses a different, ``two-sided'' theorem to obtain its results: If $m > 0$ and $z_1,\ldots,z_n \in \C$ nondecreasing in absolute value with $|z_1| = 1$, then for any $b_1,\ldots,b_n \in \C$, there exists an integer $\nu$ such that $m \le \nu \le m+n$ and \[ \bigg| \sum_{j=1}^n b_jz_j \bigg| \ge \frac 1{2n} \left( \frac n{8e(m+n)} \right)^n \min_{1\le k\le n} \bigg| \sum_{j=1}^k b_j \bigg|. \] In addition to the use of the ``two-sided'' theorem above, the authors use a modified idea attributed to Kreisel involving a sequence of integrals. \par The main result of this article is a single theorem, for residues $\ell\not\equiv1\mod k$ for sufficiently large moduli~$k$, under the assumption that there exists $0 < \delta < \frac 1{10}$ such that no function $L(s,\chi)$ with $\chi(\ell) \ne 1$ vanishes in the closed disk $|s-1| \le \frac 12 + 4\delta$. (This assumption is stronger than HC$(2\sqrt\delta)$ but weaker than HC$(\frac 12 + 4\delta)$.) For any sufficiently large $T$, the interval $I = [T,e^{(\log T)^2(\log\log T)^3}]$ contains $x_1,x_2$ such that \[ \psi(x_1;k,1,\ell) \ge x_1^{1/2-4\delta} \quad \text{and} \quad \psi(x_2;k,1,\ell) \le -x_1^{1/2-4\delta}. \] The authors compare their result to~\cite[Theorem 1.1]{MR0146157}, which uses more conventional methods and yields a more localized sign change. \par This article cites~\cites{MR0133308, MR0146157, MR0132728}.}
      \field{issn}{0065-1036}
      \field{journaltitle}{Acta Arith. 11 (1965), 147-161; ibid.}
      \field{title}{Further developments in the comparative prime-number theory. {V}}
      \field{volume}{11}
      \field{year}{1965}
      \true{nocite}
      \field{pages}{193\bibrangedash 202}
      \range{pages}{10}
      \verb{doi}
      \verb 10.4064/aa-11-2-193-202
      \endverb
      \verb{mrnumber}
      \verb 0182616
      \endverb
      \verb{urlraw}
      \verb https://doi.org/10.4064/aa-11-2-193-202
      \endverb
      \verb{url}
      \verb https://doi.org/10.4064/aa-11-2-193-202
      \endverb
    \endentry
    \entry{MR0177963}{article}{}
      \name{author}{2}{}{%
        {{hash=d33db82fc1e83dc064a23d93e5a21891}{%
           family={Knapowski},
           familyi={K\bibinitperiod},
           given={S.},
           giveni={S\bibinitperiod}}}%
        {{hash=e68e101772b88cb1d48e688164bece92}{%
           family={Turán},
           familyi={T\bibinitperiod},
           given={P.},
           giveni={P\bibinitperiod}}}%
      }
      \strng{namehash}{b6f197e0531d2bbfab6a59c048e3b28d}
      \strng{fullhash}{b6f197e0531d2bbfab6a59c048e3b28d}
      \strng{bibnamehash}{b6f197e0531d2bbfab6a59c048e3b28d}
      \strng{authorbibnamehash}{b6f197e0531d2bbfab6a59c048e3b28d}
      \strng{authornamehash}{b6f197e0531d2bbfab6a59c048e3b28d}
      \strng{authorfullhash}{b6f197e0531d2bbfab6a59c048e3b28d}
      \field{extraname}{14}
      \field{sortinit}{K}
      \field{sortinithash}{d3edc18d54b9438a72c24c925bfb38f4}
      \field{labelprefix}{A}
      \field{labelnamesource}{author}
      \field{labeltitlesource}{title}
      \field{annotation}{The authors begin by remarking on some variants of the result of Hardy--Littlewood--Landau~\cite{MR1544293A, MR1544293B, MR1555148} that Chebyshev's assertion, namely that $\lim_{x\to\infty} \pi_e(x;4,1,3) = -\infty$, is equivalent to GRH for $L(s,\chi_{-4})$. The same methods would show the ``Abelian preponderance-relations'' that the limit $\lim_{x\to\infty} \pi_e(x;3,1,2) = -\infty$ if and only if GRH holds for $L(s,\chi_{-3})$, while $\lim_{x\to\infty} \pi_e(x;8,1,\ell) = -\infty$ for all $\ell\in\{3,5,7\}$ if and only if GRH holds for all nonprincipal Dirichlet $L$-functions\mod 8, and (``mutatis mutandis'') $\lim_{x\to\infty} \pi_e(x;12,1,\ell) = -\infty$ for all $\ell\in\{5,7,11\}$ if and only if GRH holds for all nonprincipal Dirichlet $L$-functions\mod{12}. All of these results, they point out, hold with $\pi_e$ replaced by~$\theta_e$. \par For the modulus $k=8$, in the case where $\ell_1,\ell_2\in\{3,5,7\}$ are distinct quadratic nonresidues, the authors had shown ~\cite{MR0156827} that \begin{align*} \max_{T^{1/3} \leq x \leq T}\theta_e(x; 8, \ell_1, \ell_2)> \sqrt{T} \exp\left(-22 \frac{\log T \log\log\log T}{\log \log T}\right); \end{align*} however, they point out that the method failed to yield the analogous result for the ``properly Čebyšev'' function~$\pi_e$. In this article, the authors do establish analogous large oscillations (without identifying the signs of those oscillations) in the form \begin{align*} \max_{T^{1/3} \leq x \leq T} \left|\pi_e(x; 8, \ell_1, \ell_2)\right| \geq \sqrt{T} \exp\left(-23 \frac{\log T \log\log\log T}{\log \log T}\right) \end{align*} for any distinct reduced residues $\ell_1,\ell_2\mod8$, as well as the analogous statement for $\pi_e(x;4,1,3)$. The additional technical tool is a result (then unpublished) of Szegö that derives estimates for $\sum_{j=1}^n b_j e^{-jy} \log j$ from estimates for $\sum_{j=1}^n b_j e^{-jy}$. \par This article cites~\cites{1853.Chebyshev, MR1544293A, MR1544293B, MR1555148, MR0156827}.}
      \field{issn}{0021-7670}
      \field{journaltitle}{J. Analyse Math.}
      \field{title}{On an assertion of Čebyšev}
      \field{volume}{14}
      \field{year}{1965}
      \true{nocite}
      \field{pages}{267\bibrangedash 274}
      \range{pages}{8}
      \verb{doi}
      \verb 10.1007/BF02806393
      \endverb
      \verb{mrnumber}
      \verb 0177963
      \endverb
      \verb{urlraw}
      \verb https://doi.org/10.1007/BF02806393
      \endverb
      \verb{url}
      \verb https://doi.org/10.1007/BF02806393
      \endverb
    \endentry
    \entry{MR0200250}{article}{}
      \name{author}{2}{}{%
        {{hash=d33db82fc1e83dc064a23d93e5a21891}{%
           family={Knapowski},
           familyi={K\bibinitperiod},
           given={S.},
           giveni={S\bibinitperiod}}}%
        {{hash=e68e101772b88cb1d48e688164bece92}{%
           family={Turán},
           familyi={T\bibinitperiod},
           given={P.},
           giveni={P\bibinitperiod}}}%
      }
      \strng{namehash}{b6f197e0531d2bbfab6a59c048e3b28d}
      \strng{fullhash}{b6f197e0531d2bbfab6a59c048e3b28d}
      \strng{bibnamehash}{b6f197e0531d2bbfab6a59c048e3b28d}
      \strng{authorbibnamehash}{b6f197e0531d2bbfab6a59c048e3b28d}
      \strng{authornamehash}{b6f197e0531d2bbfab6a59c048e3b28d}
      \strng{authorfullhash}{b6f197e0531d2bbfab6a59c048e3b28d}
      \field{extraname}{15}
      \field{sortinit}{K}
      \field{sortinithash}{d3edc18d54b9438a72c24c925bfb38f4}
      \field{labelprefix}{A}
      \field{labelnamesource}{author}
      \field{labeltitlesource}{title}
      \field{annotation}{The authors consider a ``modified Abelian means'' race between two quadratic residues $\ell_1,\ell_2\mod k$, under the assumption $GRH(\frac3{\sqrt\eta},E_k)$ for suitable constants $\eta$ and~$E_k$. Their main result is that there exist $x\in [T^{1-\sqrt\eta},T\log T]$ and $\nu \sim 2\eta\log T$ such that \[ \theta_l(x,\nu;k,\ell_1,\ell_2) > T^{1/2-2\sqrt{\eta}} \] (and thus the symmetric result for a large negative value). A corollary is the existence of an interval $I \subset [T^{1-4\sqrt{\eta}},T^{1+4\sqrt{\eta}}]$ such that \[ \pi(I;k,\ell_1,\ell_2) > T^{1/2 - 3 \sqrt{\eta}}. \] The authors obtain similar bounds for two quadratic nonresidues in~\cite{MR0182616A}, but emphasize that they have not been able to extend the results to races where $\ell_1 \not\equiv 1 \mod{k}$ is a quadratic residue and $\ell_2$ is a quadratic nonresidue. They employ the power-sum method for exponential sums.\par This article cites~\cites{MR0146158,MR0174538,MR0180539,MR0182616A}.}
      \field{issn}{0065-1036}
      \field{journaltitle}{Acta Arith.}
      \field{title}{Further developments in the comparative prime-number theory. {VI}. {A}ccumulation theorems for residue-classes representing quadratic residues {${\rm mod}\,k$}}
      \field{volume}{12}
      \field{year}{1966}
      \true{nocite}
      \field{pages}{85\bibrangedash 96}
      \range{pages}{12}
      \verb{doi}
      \verb 10.4064/aa-12-1-85-96
      \endverb
      \verb{mrnumber}
      \verb 0200250
      \endverb
      \verb{urlraw}
      \verb https://doi.org/10.4064/aa-12-1-85-96
      \endverb
      \verb{url}
      \verb https://doi.org/10.4064/aa-12-1-85-96
      \endverb
    \endentry
    \entry{MR0272729}{article}{}
      \name{author}{2}{}{%
        {{hash=d33db82fc1e83dc064a23d93e5a21891}{%
           family={Knapowski},
           familyi={K\bibinitperiod},
           given={S.},
           giveni={S\bibinitperiod}}}%
        {{hash=e68e101772b88cb1d48e688164bece92}{%
           family={Turán},
           familyi={T\bibinitperiod},
           given={P.},
           giveni={P\bibinitperiod}}}%
      }
      \list{language}{1}{%
        {German}%
      }
      \list{publisher}{1}{%
        {Plenum, New York}%
      }
      \strng{namehash}{b6f197e0531d2bbfab6a59c048e3b28d}
      \strng{fullhash}{b6f197e0531d2bbfab6a59c048e3b28d}
      \strng{bibnamehash}{b6f197e0531d2bbfab6a59c048e3b28d}
      \strng{authorbibnamehash}{b6f197e0531d2bbfab6a59c048e3b28d}
      \strng{authornamehash}{b6f197e0531d2bbfab6a59c048e3b28d}
      \strng{authorfullhash}{b6f197e0531d2bbfab6a59c048e3b28d}
      \field{extraname}{16}
      \field{sortinit}{K}
      \field{sortinithash}{d3edc18d54b9438a72c24c925bfb38f4}
      \field{labelprefix}{A}
      \field{labelnamesource}{author}
      \field{labeltitlesource}{title}
      \field{annotation}{The authors prove that for all sufficiently large~$T$, there exist positive reals $x<y\le T$ such that $\theta\bigl((x,y];4,3,1\bigr)>\sqrt{y}$. The authors also prove an analogous bound for $\theta\bigl((x,y];4,1,3\bigr)$.\par This article cites~\cites{1853.Chebyshev,MR1555148,MR1544293A,MR1544293B,1914.Littlewood,MR0162771,MR1511360,MR0174538,MR0132728,MR0146157}}
      \field{booktitle}{Number Theory and Analysis (Papers in Honor of Edmund Landau)}
      \field{title}{Über einige Fragen der vergleichenden Primzahltheorie}
      \field{year}{1969}
      \true{nocite}
      \field{pages}{157\bibrangedash 171}
      \range{pages}{15}
      \verb{mrnumber}
      \verb 0272729
      \endverb
    \endentry
    \entry{MR0302585}{article}{}
      \name{author}{2}{}{%
        {{hash=d33db82fc1e83dc064a23d93e5a21891}{%
           family={Knapowski},
           familyi={K\bibinitperiod},
           given={S.},
           giveni={S\bibinitperiod}}}%
        {{hash=e68e101772b88cb1d48e688164bece92}{%
           family={Turán},
           familyi={T\bibinitperiod},
           given={P.},
           giveni={P\bibinitperiod}}}%
      }
      \strng{namehash}{b6f197e0531d2bbfab6a59c048e3b28d}
      \strng{fullhash}{b6f197e0531d2bbfab6a59c048e3b28d}
      \strng{bibnamehash}{b6f197e0531d2bbfab6a59c048e3b28d}
      \strng{authorbibnamehash}{b6f197e0531d2bbfab6a59c048e3b28d}
      \strng{authornamehash}{b6f197e0531d2bbfab6a59c048e3b28d}
      \strng{authorfullhash}{b6f197e0531d2bbfab6a59c048e3b28d}
      \field{extraname}{17}
      \field{sortinit}{K}
      \field{sortinithash}{d3edc18d54b9438a72c24c925bfb38f4}
      \field{labelprefix}{A}
      \field{labelnamesource}{author}
      \field{labeltitlesource}{title}
      \field{annotation}{The authors show that for large enough $T$, there exist numbers $U_1, U_2, U_3, U_4$ with \begin{align*} \log \log \log T &\leq U_2 \exp\left(-\log^{15/16} U_2\right) \leq U_1 < U_2 \leq T\\ \log \log \log T &\leq U_4 \exp\left(-\log^{15/16} U_4\right) \leq U_3 < U_4 \leq T \end{align*} such that $\theta([U_1, U_2]; 4, 1, 3) > \sqrt{U_2}$ and $\theta([U_3, U_4]; 4, 1, 3)< - \sqrt{U_4}$. In particular, there exist consecutive primes $p_n$ and $p_{n+1}$, both congruent to $1\mod4$, satisfying $\log \log \log T \leq p_n < p_{n+1} \leq T$. \par This article cites~\cites{MR0146157, MR0174538, MR0060548}.}
      \field{issn}{0065-1036}
      \field{journaltitle}{Acta Arith.}
      \field{title}{Further developments in the comparative prime number theory. {VII}}
      \field{volume}{21}
      \field{year}{1972}
      \true{nocite}
      \field{pages}{193\bibrangedash 201}
      \range{pages}{9}
      \verb{doi}
      \verb 10.4064/aa-21-1-193-201
      \endverb
      \verb{mrnumber}
      \verb 0302585
      \endverb
      \verb{urlraw}
      \verb https://doi.org/10.4064/aa-21-1-193-201
      \endverb
      \verb{url}
      \verb https://doi.org/10.4064/aa-21-1-193-201
      \endverb
    \endentry
    \entry{MR0439771}{inproceedings}{}
      \name{author}{2}{}{%
        {{hash=d33db82fc1e83dc064a23d93e5a21891}{%
           family={Knapowski},
           familyi={K\bibinitperiod},
           given={S.},
           giveni={S\bibinitperiod}}}%
        {{hash=e68e101772b88cb1d48e688164bece92}{%
           family={Turán},
           familyi={T\bibinitperiod},
           given={P.},
           giveni={P\bibinitperiod}}}%
      }
      \list{publisher}{1}{%
        {North-Holland, Amsterdam}%
      }
      \strng{namehash}{b6f197e0531d2bbfab6a59c048e3b28d}
      \strng{fullhash}{b6f197e0531d2bbfab6a59c048e3b28d}
      \strng{bibnamehash}{b6f197e0531d2bbfab6a59c048e3b28d}
      \strng{authorbibnamehash}{b6f197e0531d2bbfab6a59c048e3b28d}
      \strng{authornamehash}{b6f197e0531d2bbfab6a59c048e3b28d}
      \strng{authorfullhash}{b6f197e0531d2bbfab6a59c048e3b28d}
      \field{extraname}{18}
      \field{sortinit}{K}
      \field{sortinithash}{d3edc18d54b9438a72c24c925bfb38f4}
      \field{labelprefix}{A}
      \field{labelnamesource}{author}
      \field{labeltitlesource}{title}
      \field{annotation}{The authors prove that $W^\pi(T) = \Omega((\log T)^{1/4}(\log\log T)^{-4})$ (a slightly weaker result is claimed). In light of Ingham's stronger result~\cite{1936.Ingham} assuming SA, it suffices to assume that RH is false. Given any zero $\rho=\beta+i\gamma$ of $\zeta(s)$ with $\beta>\frac12$, the authors show that for $x \in [Y,Y\exp((\log Y)^{3/4} (\log\log Y)^4)]$ one has $\Delta^\Pi(x) = \Omega_\pm(Y^\beta \exp(-\sqrt{\log Y}))$ when $Y$ is sufficiently large in terms of $\rho$. The authors make special mention of their application of a ``two-sided'' power-sum theorem to get a strong ``one-sided'' result.\par This article cites~\cite{MR1511206, 1914.Littlewood, MR1555148, MR1573970, MR1575079, MR0067145, MR0202686, MR0137689, 1930.Polya, 1936.Ingham, MR0133309, MR0182616B}.}
      \field{booktitle}{Topics in number theory ({P}roc. {C}olloq., {D}ebrecen, 1974)}
      \field{title}{On the sign changes of {$(\pi (x)-{\rm li}\ x)$}. {I}}
      \field{year}{1976}
      \true{nocite}
      \field{pages}{153\bibrangedash 169. Colloq. Math. Soc. János Bolyai\bibrangessep Vol. 13}
      \range{pages}{-1}
      \verb{mrnumber}
      \verb 0439771
      \endverb
    \endentry
    \entry{MR0439772}{article}{}
      \name{author}{2}{}{%
        {{hash=d33db82fc1e83dc064a23d93e5a21891}{%
           family={Knapowski},
           familyi={K\bibinitperiod},
           given={S.},
           giveni={S\bibinitperiod}}}%
        {{hash=e68e101772b88cb1d48e688164bece92}{%
           family={Turán},
           familyi={T\bibinitperiod},
           given={P.},
           giveni={P\bibinitperiod}}}%
      }
      \strng{namehash}{b6f197e0531d2bbfab6a59c048e3b28d}
      \strng{fullhash}{b6f197e0531d2bbfab6a59c048e3b28d}
      \strng{bibnamehash}{b6f197e0531d2bbfab6a59c048e3b28d}
      \strng{authorbibnamehash}{b6f197e0531d2bbfab6a59c048e3b28d}
      \strng{authornamehash}{b6f197e0531d2bbfab6a59c048e3b28d}
      \strng{authorfullhash}{b6f197e0531d2bbfab6a59c048e3b28d}
      \field{extraname}{19}
      \field{sortinit}{K}
      \field{sortinithash}{d3edc18d54b9438a72c24c925bfb38f4}
      \field{labelprefix}{A}
      \field{labelnamesource}{author}
      \field{labeltitlesource}{title}
      \field{annotation}{Following ideas of Littlewood, Ingham, and Skewes, the authors show unconditionally that $W(Y) \gg \log\log\log Y$ for sufficiently large $Y$, where the implied constants are effective.\par The proof itself is divided into two cases. First, supposing the existence of an RH-violating zero $\beta+i\gamma$ of $\zeta(s)$ such that $\beta \ge \frac 12 + 2\log^{-1/5} Y$ and $0 < \gamma \le \log^{1/5} Y$, the authors establish the much stronger lower bound \[ V_1(Y) > \frac 12 \bigg(\frac{\log Y}{2\log^{5/6} Y} \bigg)^{1/5} > \frac 14 \log^{1/30} Y. \] The second case, where there is no such zero, is more technical and relies as usual upon Dirichlet's box principle.\par This article cites~\cites{1914.Littlewood, MR1555148, MR0067145, 1930.Polya, 1936.Ingham, MR0439771, MR0133308, 1932.Ingham}.}
      \field{journaltitle}{Monatsh. Math.}
      \field{number}{2}
      \field{title}{On the sign changes of {$(\pi (x)-{\rm li}\ x)$}. {II}}
      \field{volume}{82}
      \field{year}{1976}
      \true{nocite}
      \field{pages}{163\bibrangedash 175}
      \range{pages}{13}
      \verb{doi}
      \verb 10.1007/BF01305997
      \endverb
      \verb{mrnumber}
      \verb 0439772
      \endverb
      \verb{urlraw}
      \verb https://doi.org/10.1007/BF01305997
      \endverb
      \verb{url}
      \verb https://doi.org/10.1007/BF01305997
      \endverb
    \endentry
    \entry{MR0466043}{incollection}{}
      \name{author}{2}{}{%
        {{hash=d33db82fc1e83dc064a23d93e5a21891}{%
           family={Knapowski},
           familyi={K\bibinitperiod},
           given={S.},
           giveni={S\bibinitperiod}}}%
        {{hash=e68e101772b88cb1d48e688164bece92}{%
           family={Turán},
           familyi={T\bibinitperiod},
           given={P.},
           giveni={P\bibinitperiod}}}%
      }
      \list{publisher}{1}{%
        {Academic Press, New York}%
      }
      \strng{namehash}{b6f197e0531d2bbfab6a59c048e3b28d}
      \strng{fullhash}{b6f197e0531d2bbfab6a59c048e3b28d}
      \strng{bibnamehash}{b6f197e0531d2bbfab6a59c048e3b28d}
      \strng{authorbibnamehash}{b6f197e0531d2bbfab6a59c048e3b28d}
      \strng{authornamehash}{b6f197e0531d2bbfab6a59c048e3b28d}
      \strng{authorfullhash}{b6f197e0531d2bbfab6a59c048e3b28d}
      \field{extraname}{20}
      \field{sortinit}{K}
      \field{sortinithash}{d3edc18d54b9438a72c24c925bfb38f4}
      \field{labelprefix}{A}
      \field{labelnamesource}{author}
      \field{labeltitlesource}{title}
      \field{annotation}{The authors show, with $p_\nu$ denoting the $\nu$th prime, that when $T$ is sufficiently large, \[ \sum_{\substack{p_\nu\le T\\p_\nu\equiv p_{\nu+1}\equiv 1\mod 4}} 1 > \log^B T \] for some effective constant~$B$. The fact that the left-hand side tends to $\infty$ follows from Littlewood's oscillation theorem for primes\mod4, but no quantitative rate of growth had been established. The proof uses their result from~\cite{MR0302585}, as well as Ingham's application of Fejér kernels as in~\cite{1936.Ingham}. The authors note the open problem about the existence of infinitely many triples of consecutive primes congruent to~$1\mod4$. They also guess that the four possibilities for the pair of congruence classes $(p_\nu, p_{\nu+1}) \mod 4$ are not equally likely.\par This article cites \cite{MR0174538, 1932.Ingham, MR0302585}.}
      \field{booktitle}{Number theory and algebra}
      \field{title}{On prime numbers {$\equiv 1$} resp. {$3\mod4$}}
      \field{year}{1977}
      \true{nocite}
      \field{pages}{157\bibrangedash 165}
      \range{pages}{9}
      \verb{mrnumber}
      \verb 0466043
      \endverb
    \endentry
    \entry{MR820241}{incollection}{}
      \name{author}{2}{}{%
        {{hash=c5f23f7a8c75fe8b8720ab22fc314bda}{%
           family={Kolesnik},
           familyi={K\bibinitperiod},
           given={G.},
           giveni={G\bibinitperiod}}}%
        {{hash=b79ccf3603e8ce64e49a0d4684120097}{%
           family={Straus},
           familyi={S\bibinitperiod},
           given={E.\bibnamedelimi G.},
           giveni={E\bibinitperiod\bibinitdelim G\bibinitperiod}}}%
      }
      \list{publisher}{1}{%
        {Birkhäuser, Basel}%
      }
      \strng{namehash}{db6a61d464b7f706be1353a113c7a3d2}
      \strng{fullhash}{db6a61d464b7f706be1353a113c7a3d2}
      \strng{bibnamehash}{db6a61d464b7f706be1353a113c7a3d2}
      \strng{authorbibnamehash}{db6a61d464b7f706be1353a113c7a3d2}
      \strng{authornamehash}{db6a61d464b7f706be1353a113c7a3d2}
      \strng{authorfullhash}{db6a61d464b7f706be1353a113c7a3d2}
      \field{sortinit}{K}
      \field{sortinithash}{d3edc18d54b9438a72c24c925bfb38f4}
      \field{labelprefix}{A}
      \field{labelnamesource}{author}
      \field{labeltitlesource}{title}
      \field{annotation}{The authors examine the second main theorem of the power-sum method, obtaining the constant $B=4e$ which is best possible by~\cites{MR0159791}. When $m\leq n$ they also give the larger lower bound \[ |s_v| \geq \frac{n!(2m+n)!}{2^n(2m+2n)! \sqrt{2m+2n+1}} \min_{1 \leq j \leq n}|b_1+ \cdots + b_j|. \]\par This article cites~\cites{MR0077579,MR0159791}.}
      \field{bookpagination}{none}
      \field{booktitle}{Studies in pure mathematics}
      \field{title}{On the sum of powers of complex numbers}
      \field{year}{1983}
      \true{nocite}
      \field{pages}{427\bibrangedash 442}
      \range{pages}{16}
      \verb{mrnumber}
      \verb 820241
      \endverb
    \endentry
    \entry{MR2420864}{article}{}
      \name{author}{1}{}{%
        {{hash=9491eb14d70047b89ad523bb740e7125}{%
           family={Kotnik},
           familyi={K\bibinitperiod},
           given={Tadej},
           giveni={T\bibinitperiod}}}%
      }
      \strng{namehash}{9491eb14d70047b89ad523bb740e7125}
      \strng{fullhash}{9491eb14d70047b89ad523bb740e7125}
      \strng{bibnamehash}{9491eb14d70047b89ad523bb740e7125}
      \strng{authorbibnamehash}{9491eb14d70047b89ad523bb740e7125}
      \strng{authornamehash}{9491eb14d70047b89ad523bb740e7125}
      \strng{authorfullhash}{9491eb14d70047b89ad523bb740e7125}
      \field{sortinit}{K}
      \field{sortinithash}{d3edc18d54b9438a72c24c925bfb38f4}
      \field{labelprefix}{A}
      \field{labelnamesource}{author}
      \field{labeltitlesource}{title}
      \field{annotation}{The author shows that $\pi(x) <\li(x)$ for $2 \leq x \leq 10^{14}$, improving unpublished work of Odlyzko showing that $\pi(x) <\li(x)$ for $x \leq 1.59 \times 10^{13}$. The rigorous computation used an Eratosthenes sieve for $\pi(x)$ and Ramanujan's formula $\li(x)=\sqrt{x} \sum_{n=1}^{\infty} a_n \log ^n x + \log \log x + C_0$, truncated at $n=75$ and linearly interpolated when~$x>10^{10}$. Based on the data set, the author conjectures that $|\Delta^\pi(x)|<\sqrt{x}$ and ${-\frac{2}{5}x^{3/2} <\fA^\pi_1(x)<0}$ for $x>2$. The program for computation and storage of the data was written in Delphi~6.0 and run on a PC with a 2.4~GHz Intel Pentium~4 processor and 512~MB of RAM.\par This article cites~\cites{MR1752093,MR0369287,1932.Ingham,MR0202686,1914.Littlewood,MR0137689,MR0067145,MR866118}.}
      \field{issn}{1019-7168}
      \field{journaltitle}{Adv. Comput. Math.}
      \field{number}{1}
      \field{title}{The prime-counting function and its analytic approximations: {$\pi (x)$} and its approximations}
      \field{volume}{29}
      \field{year}{2008}
      \true{nocite}
      \field{pages}{55\bibrangedash 70}
      \range{pages}{16}
      \verb{doi}
      \verb 10.1007/s10444-007-9039-2
      \endverb
      \verb{mrnumber}
      \verb 2420864
      \endverb
    \endentry
    \entry{MR2118272}{article}{}
      \name{author}{2}{}{%
        {{hash=9491eb14d70047b89ad523bb740e7125}{%
           family={Kotnik},
           familyi={K\bibinitperiod},
           given={Tadej},
           giveni={T\bibinitperiod}}}%
        {{hash=1ed18e6c700f1143da76930e91f0e194}{%
           family={Lune},
           familyi={L\bibinitperiod},
           given={Jan},
           giveni={J\bibinitperiod},
           prefix={van\bibnamedelima de},
           prefixi={v\bibinitperiod\bibinitdelim d\bibinitperiod}}}%
      }
      \strng{namehash}{01c804ddeb08d741c890deb05472a70b}
      \strng{fullhash}{01c804ddeb08d741c890deb05472a70b}
      \strng{bibnamehash}{01c804ddeb08d741c890deb05472a70b}
      \strng{authorbibnamehash}{01c804ddeb08d741c890deb05472a70b}
      \strng{authornamehash}{01c804ddeb08d741c890deb05472a70b}
      \strng{authorfullhash}{01c804ddeb08d741c890deb05472a70b}
      \field{sortinit}{K}
      \field{sortinithash}{d3edc18d54b9438a72c24c925bfb38f4}
      \field{labelprefix}{A}
      \field{labelnamesource}{author}
      \field{labeltitlesource}{title}
      \field{annotation}{The authors use the first million zeros of $\zeta(s)$ to approximate $E^M(x)$ for $x \leq 10^{10^{10}}$. Based on the data, the authors conjecture that $E^M(x) = \Omega_{\pm} (\sqrt{\log \log \log x} )$, with the data suggesting approximately the constant $\frac12$ for each sign. In particular, the authors propose that earlier conjectures~\cite{MR280447,MR240062,MR2078705} on the extreme values of $E^M(x)$ do not hold. \par This article cites~\cites{1914.Littlewood,1897.Mertens,MR2078705,MR783538,MR891440,MR280447,MR240062}.}
      \field{issn}{1058-6458}
      \field{journaltitle}{Experiment. Math.}
      \field{number}{4}
      \field{title}{On the order of the {M}ertens function}
      \field{volume}{13}
      \field{year}{2004}
      \true{nocite}
      \field{pages}{473\bibrangedash 481}
      \range{pages}{9}
      \verb{mrnumber}
      \verb 2118272
      \endverb
    \endentry
    \entry{MR2282922}{incollection}{}
      \name{author}{2}{}{%
        {{hash=9491eb14d70047b89ad523bb740e7125}{%
           family={Kotnik},
           familyi={K\bibinitperiod},
           given={Tadej},
           giveni={T\bibinitperiod}}}%
        {{hash=c83d046dcdb536f5163e46eb8fe07827}{%
           family={Riele},
           familyi={R\bibinitperiod},
           given={Herman},
           giveni={H\bibinitperiod},
           prefix={te},
           prefixi={t\bibinitperiod}}}%
      }
      \list{publisher}{1}{%
        {Springer, Berlin}%
      }
      \strng{namehash}{ed0be9de3d43cca36b74eb71c7163f7a}
      \strng{fullhash}{ed0be9de3d43cca36b74eb71c7163f7a}
      \strng{bibnamehash}{ed0be9de3d43cca36b74eb71c7163f7a}
      \strng{authorbibnamehash}{ed0be9de3d43cca36b74eb71c7163f7a}
      \strng{authornamehash}{ed0be9de3d43cca36b74eb71c7163f7a}
      \strng{authorfullhash}{ed0be9de3d43cca36b74eb71c7163f7a}
      \field{sortinit}{K}
      \field{sortinithash}{d3edc18d54b9438a72c24c925bfb38f4}
      \field{labelprefix}{A}
      \field{labelnamesource}{author}
      \field{labeltitlesource}{title}
      \field{annotation}{Based on the work of Odlyzko and te Riele~\cite{MR783538}, a refined version of the LLL algorithm, and improved precision for zeros of $\zeta$ function, the authors show that \[ \liminf_{x \to \infty} E^M(x) < -1.229 \quad \text{and} \quad \limsup_{x \to \infty} E^M(x) > 1.218. \] Using a similar idea, the authors show that a counterexample to the Mertens conjecture occurs before $\exp(1.59 \times 10^{40})$, improving Pintz's result~\cite{MR891440}. The authors also provide new numerical evidence to support the conjecture $M(x) = \Omega_{\pm} (\sqrt{x\log \log \log x} )$ made in~\cite{MR2118272}. \par This article cites~\cite{MR0046485,MR891440,MR2078705,MR783538,1897.Mertens,MR2118272,MR429789,MR0330069,MR0352026,MR0006202,MR549977,MR280447,MR240062}.}
      \field{booktitle}{Algorithmic number theory}
      \field{series}{Lecture Notes in Comput. Sci.}
      \field{title}{The {M}ertens conjecture revisited}
      \field{volume}{4076}
      \field{year}{2006}
      \true{nocite}
      \field{pages}{156\bibrangedash 167}
      \range{pages}{12}
      \verb{mrnumber}
      \verb 2282922
      \endverb
    \endentry
    \entry{MR2987108}{article}{}
      \name{author}{2}{}{%
        {{hash=3028d3b5860b77006bb46a3a93b688b4}{%
           family={Kunik},
           familyi={K\bibinitperiod},
           given={Matthias},
           giveni={M\bibinitperiod}}}%
        {{hash=cdaf4b1735b29361604e44f0ef594303}{%
           family={Lucht},
           familyi={L\bibinitperiod},
           given={Lutz\bibnamedelima G.},
           giveni={L\bibinitperiod\bibinitdelim G\bibinitperiod}}}%
      }
      \strng{namehash}{0ff5403c60f18b96de7464513572a2ac}
      \strng{fullhash}{0ff5403c60f18b96de7464513572a2ac}
      \strng{bibnamehash}{0ff5403c60f18b96de7464513572a2ac}
      \strng{authorbibnamehash}{0ff5403c60f18b96de7464513572a2ac}
      \strng{authornamehash}{0ff5403c60f18b96de7464513572a2ac}
      \strng{authorfullhash}{0ff5403c60f18b96de7464513572a2ac}
      \field{sortinit}{K}
      \field{sortinithash}{d3edc18d54b9438a72c24c925bfb38f4}
      \field{labelprefix}{A}
      \field{labelnamesource}{author}
      \field{labeltitlesource}{title}
      \field{annotation}{In this article, the authors investigate the series $F(w)= \sum_{n=1}^{\infty} \frac{\Lambda(n)}{n}w^n$ for $|w|\le1$. The authors use a more convenient direct method compared to that of Hardy and Littlewood in~\cite{MR1555148} to derive the explicit formulas for the series $\sum \Lambda(n)e^{-nz} = 1/z + (\cosh(z)-1)\log z + T(z) - \sum_ \rho \Gamma(\rho)z^{-\rho}$ and for $F(e^{-z})$; in particular the authors obtain a closed form for the entire function $T(z)$. The main result reveals logarithmic singularities of $F(e^{2\pi it})$ at the reduced rational numbers $t=\frac{a}{q}$ with squarefree denominator~$q \in \N$.\par This article cites~\cites{MR1555148}.}
      \field{isbn}{978-83-232-2442-6}
      \field{issn}{0208-6573}
      \field{journaltitle}{Funct. Approx. Comment. Math.}
      \field{number}{part 1}
      \field{part}{part 1}
      \field{title}{Power series with the von {M}angoldt function}
      \field{volume}{47}
      \field{year}{2012}
      \true{nocite}
      \field{pages}{15\bibrangedash 33}
      \range{pages}{19}
      \verb{doi}
      \verb 10.7169/facm/2012.47.1.2
      \endverb
      \verb{mrnumber}
      \verb 2987108
      \endverb
    \endentry
    \entry{MR2943671}{article}{}
      \name{author}{1}{}{%
        {{hash=a01ad5bbada0f5c69f81f1149c07aefc}{%
           family={Lamzouri},
           familyi={L\bibinitperiod},
           given={Y.},
           giveni={Y\bibinitperiod}}}%
      }
      \strng{namehash}{a01ad5bbada0f5c69f81f1149c07aefc}
      \strng{fullhash}{a01ad5bbada0f5c69f81f1149c07aefc}
      \strng{bibnamehash}{a01ad5bbada0f5c69f81f1149c07aefc}
      \strng{authorbibnamehash}{a01ad5bbada0f5c69f81f1149c07aefc}
      \strng{authornamehash}{a01ad5bbada0f5c69f81f1149c07aefc}
      \strng{authorfullhash}{a01ad5bbada0f5c69f81f1149c07aefc}
      \field{extraname}{1}
      \field{sortinit}{L}
      \field{sortinithash}{dad3efd0836470093a7b4a7bb756eb8c}
      \field{labelprefix}{A}
      \field{labelnamesource}{author}
      \field{labeltitlesource}{title}
      \field{annotation}{The author refines tail estimates for the distribution $\mu_{q;a_1,\ldots,a_r}$ given by Rubinstein and Sarnak~\cite{MR1329368}. From the Math Review by D.~R.~Heath-Brown, with slight changes in notation: ``Define \[ \sigma(q) = \biggl( 2 \sum_{\substack{\chi\mod q \\ \chi\ne\chi_0}} \sum_{\gamma>0} \frac1{1/4+\gamma^2} \biggr)^{1/2}, \] with $\gamma$ running over ordinates of zeros of $L(s,\chi)$. \ldots{} It is proved here that if $0<\lambda\le\sqrt{\log\log q}$ then \[ \mu_{q;a_1,\ldots,a_r} \bigl( \|t\|^2 > \lambda\sigma(q) \bigr) = (2\pi)^{-r/2} \int_{\|t\|_2>\lambda} \exp\bigl( -\tfrac12\|t\|_2 \bigr) \,dt + O_r\bigl( (\log q)^{-2} \bigr). \] \ldots{} For large~$V$ it is shown that \[ \exp\biggl( -c_1 \frac{V^2}{\phi(q)\log q} \biggr) \ll_q \mu_{q;a_1,\ldots,a_r} \bigl( \|t\|^2 > V \bigr) \ll_q \exp\biggl( -c_2 \frac{V^2}{\phi(q)\log q} \biggr) \] for $(\phi(q)\log q)^{1/2} \ll V \le A\phi(q)\log q$, with constants $c_1,c_2$ depending only on~$r$ and~$A$. For even larger~$V$ the paper gives further results, showing that there are transitions in the behaviour, occuring around $V=\phi(q)\log q$ and $V=\phi(q)\log^2 q$.''\par This article cites~\cites{MR3028758,MR0146156,MR0146157,MR0146158,MR0146157,MR0146160,MR0146161,MR0156826,MR0156827,MR1251926,MR1367576,MR2631002,MR634679,MR1329368}.}
      \field{issn}{0305-0041}
      \field{journaltitle}{Math. Proc. Cambridge Philos. Soc.}
      \field{number}{1}
      \field{title}{Large deviations of the limiting distribution in the {S}hanks--{R}ényi prime number race}
      \field{volume}{153}
      \field{year}{2012}
      \true{nocite}
      \field{pages}{147\bibrangedash 166}
      \range{pages}{20}
      \verb{doi}
      \verb 10.1017/S030500411200014X
      \endverb
      \verb{mrnumber}
      \verb 2943671
      \endverb
      \verb{urlraw}
      \verb https://doi.org/10.1017/S030500411200014X
      \endverb
      \verb{url}
      \verb https://doi.org/10.1017/S030500411200014X
      \endverb
    \endentry
    \entry{MR2998146}{article}{}
      \name{author}{1}{}{%
        {{hash=a01ad5bbada0f5c69f81f1149c07aefc}{%
           family={Lamzouri},
           familyi={L\bibinitperiod},
           given={Y.},
           giveni={Y\bibinitperiod}}}%
      }
      \strng{namehash}{a01ad5bbada0f5c69f81f1149c07aefc}
      \strng{fullhash}{a01ad5bbada0f5c69f81f1149c07aefc}
      \strng{bibnamehash}{a01ad5bbada0f5c69f81f1149c07aefc}
      \strng{authorbibnamehash}{a01ad5bbada0f5c69f81f1149c07aefc}
      \strng{authornamehash}{a01ad5bbada0f5c69f81f1149c07aefc}
      \strng{authorfullhash}{a01ad5bbada0f5c69f81f1149c07aefc}
      \field{extraname}{2}
      \field{sortinit}{L}
      \field{sortinithash}{dad3efd0836470093a7b4a7bb756eb8c}
      \field{labelprefix}{A}
      \field{labelnamesource}{author}
      \field{labeltitlesource}{title}
      \field{annotation}{In this article, the author examines the order of magnitude of the density $\delta_{q;a_1,\ldots,a_r}$ when the number of contestants~$r$ tends to infinity as $q\to\infty$. Assuming GRH and LI, he shows that $\delta_{q;a_1,\ldots,a_r}=\frac{1}{r!}(1+O(\frac{r^2}{\log{q}}))$ when $r\leq \sqrt{\log{q}}$, while $\delta_{q;a_1,\ldots,a_r}=\exp(-r\log{r}+r+O(\log{r}+\frac{r^2}{\log{q}}))$ when $\sqrt{\log{q}}\ll r\leq (1-\ep)\log{q}/\log\log{q}$. For larger~$r$, he provides the upper bound $\delta_{q;a_1,\ldots,a_r} \ll_{\ep} q^{-1+2\ep}$ for $(1-\ep)\log{q}/\log\log{q}\leq r\leq\phi(q)$.\par This article cites~\cites{MR0460261,MR1806291,MR3028758,MR1909220,MR2075622,MR1251926,MR1256149,MR1367576,MR0146156,MR0146157,MR0146158,MR0146157,MR0146160,MR0146161,MR0156826,MR0156827,MR2943671,1914.Littlewood,MR1956261,MR1329368}.}
      \field{issn}{1073-2780}
      \field{journaltitle}{Math. Res. Lett.}
      \field{number}{3}
      \field{title}{The {S}hanks-{R}ényi prime number race with many contestants}
      \field{volume}{19}
      \field{year}{2012}
      \true{nocite}
      \field{pages}{649\bibrangedash 666}
      \range{pages}{18}
      \verb{doi}
      \verb 10.4310/MRL.2012.v19.n3.a11
      \endverb
      \verb{mrnumber}
      \verb 2998146
      \endverb
      \verb{urlraw}
      \verb https://doi.org/10.4310/MRL.2012.v19.n3.a11
      \endverb
      \verb{url}
      \verb https://doi.org/10.4310/MRL.2012.v19.n3.a11
      \endverb
    \endentry
    \entry{MR3063909}{article}{}
      \name{author}{1}{}{%
        {{hash=a01ad5bbada0f5c69f81f1149c07aefc}{%
           family={Lamzouri},
           familyi={L\bibinitperiod},
           given={Y.},
           giveni={Y\bibinitperiod}}}%
      }
      \strng{namehash}{a01ad5bbada0f5c69f81f1149c07aefc}
      \strng{fullhash}{a01ad5bbada0f5c69f81f1149c07aefc}
      \strng{bibnamehash}{a01ad5bbada0f5c69f81f1149c07aefc}
      \strng{authorbibnamehash}{a01ad5bbada0f5c69f81f1149c07aefc}
      \strng{authornamehash}{a01ad5bbada0f5c69f81f1149c07aefc}
      \strng{authorfullhash}{a01ad5bbada0f5c69f81f1149c07aefc}
      \field{extraname}{3}
      \field{sortinit}{L}
      \field{sortinithash}{dad3efd0836470093a7b4a7bb756eb8c}
      \field{labelprefix}{A}
      \field{labelnamesource}{author}
      \field{labeltitlesource}{title}
      \field{annotation}{Assuming GRH and LI, the author describes phenomena that occur for multi-way prime number races that are not present in two-way races. When $r\ge3$, the author shows that $\bigl| \delta_{q;a_1,\ldots,a_r} - \frac1{r!} \bigr| = \Omega(1/\log q)$, in contrast to the $r=2$ case~\cite{MR3028758} where $\bigl| \delta_{q;a_1,a_2} - \frac12 \bigr| \ll q^{-1/2+o(1)}$. The author also shows that when~$q$ is sufficiently large in terms of~$r$, there are always~$r$ squares $a_1,\ldots,a_r$ modulo~$q$ for which $\delta_{q;a_1,\ldots,a_r} \ne \frac1{r!}$ (and similarly~$r$ nonsquares). The method of proof is to study the approximation of the characteristic function $\hat{\mu}_{q;a_1,\ldots,a_r}$ by that of a multivariate Gaussian, paying particular attention to secondary main terms that arise.\par This article cites~\cites{MR0447090,MR1806291,MR3028758,MR1909220,MR2202918,MR0146156,MR0146157,MR0146158,MR0146159,MR0146160,MR0146161,MR0156826,MR0156827,MR1251926,MR1256149,MR1367576,MR2998146,MR1329368}.}
      \field{issn}{0025-5831}
      \field{journaltitle}{Math. Ann.}
      \field{number}{3}
      \field{title}{Prime number races with three or more competitors}
      \field{volume}{356}
      \field{year}{2013}
      \true{nocite}
      \field{pages}{1117\bibrangedash 1162}
      \range{pages}{46}
      \verb{doi}
      \verb 10.1007/s00208-012-0874-1
      \endverb
      \verb{mrnumber}
      \verb 3063909
      \endverb
      \verb{urlraw}
      \verb https://doi.org/10.1007/s00208-012-0874-1
      \endverb
      \verb{url}
      \verb https://doi.org/10.1007/s00208-012-0874-1
      \endverb
    \endentry
    \entry{MR3455269}{article}{}
      \name{author}{1}{}{%
        {{hash=a01ad5bbada0f5c69f81f1149c07aefc}{%
           family={Lamzouri},
           familyi={L\bibinitperiod},
           given={Y.},
           giveni={Y\bibinitperiod}}}%
      }
      \strng{namehash}{a01ad5bbada0f5c69f81f1149c07aefc}
      \strng{fullhash}{a01ad5bbada0f5c69f81f1149c07aefc}
      \strng{bibnamehash}{a01ad5bbada0f5c69f81f1149c07aefc}
      \strng{authorbibnamehash}{a01ad5bbada0f5c69f81f1149c07aefc}
      \strng{authornamehash}{a01ad5bbada0f5c69f81f1149c07aefc}
      \strng{authorfullhash}{a01ad5bbada0f5c69f81f1149c07aefc}
      \field{extraname}{4}
      \field{sortinit}{L}
      \field{sortinithash}{dad3efd0836470093a7b4a7bb756eb8c}
      \field{labelprefix}{A}
      \field{labelnamesource}{author}
      \field{labeltitlesource}{title}
      \field{annotation}{Let $\mathcal{M}$ be the set of real numbers $x>1$ such that $\prod_{p \leq x } ( 1- 1/{p})^{-1} > e^{C_0} \log x$. Assuming RH, the author proves that both $\mathcal{M}$ and its complement have positive lower logarithmic density. Further assuming LI, the author shows that the logarithmic density of $\mathcal{M}$ equals $\delta(\li,\pi) \approx 0.99999974$. The author also conjectures that \[ \limsup_{x\to\infty} \bigg( \bigg( \prod_{p\le x} \bigg( 1-\frac1p \bigg)^{-1} - e^{C_0}\log x \bigg) \bigg/ \frac{(\log\log\log x)^2}{\sqrt x} \bigg) = \frac{e^{C_0}}{2\pi} \] and the symmetric result for the $\liminf$. \par This article cites~\cites{MR2605532,MR3151605,MR2202918,1914.Littlewood,MR2631002,MR634679,MR0137689,MR1329368}, as well as an earlier draft of this annotated bibliography.}
      \field{issn}{1793-0421}
      \field{journaltitle}{Int. J. Number Theory}
      \field{number}{1}
      \field{title}{A bias in {M}ertens' product formula}
      \field{volume}{12}
      \field{year}{2016}
      \true{nocite}
      \field{pages}{97\bibrangedash 109}
      \range{pages}{13}
      \verb{doi}
      \verb 10.1142/S1793042116500068
      \endverb
      \verb{mrnumber}
      \verb 3455269
      \endverb
      \verb{urlraw}
      \verb https://doi.org/10.1142/S1793042116500068
      \endverb
      \verb{url}
      \verb https://doi.org/10.1142/S1793042116500068
      \endverb
    \endentry
    \entry{MR4012359}{article}{}
      \name{author}{2}{}{%
        {{hash=64363cb2afef9e33f967130adc38684b}{%
           family={Lamzouri},
           familyi={L\bibinitperiod},
           given={Youness},
           giveni={Y\bibinitperiod}}}%
        {{hash=5509e2d8aa4558a4fd44867b61f0c363}{%
           family={Martin},
           familyi={M\bibinitperiod},
           given={Bruno},
           giveni={B\bibinitperiod}}}%
      }
      \strng{namehash}{52cf16fd2d29d30832e4c5dcecfe7ec4}
      \strng{fullhash}{52cf16fd2d29d30832e4c5dcecfe7ec4}
      \strng{bibnamehash}{52cf16fd2d29d30832e4c5dcecfe7ec4}
      \strng{authorbibnamehash}{52cf16fd2d29d30832e4c5dcecfe7ec4}
      \strng{authornamehash}{52cf16fd2d29d30832e4c5dcecfe7ec4}
      \strng{authorfullhash}{52cf16fd2d29d30832e4c5dcecfe7ec4}
      \field{sortinit}{L}
      \field{sortinithash}{dad3efd0836470093a7b4a7bb756eb8c}
      \field{labelprefix}{A}
      \field{labelnamesource}{author}
      \field{labeltitlesource}{title}
      \field{annotation}{Let~$A$ and~$B$ be disjoint sets of reduced residue classes modulo~$q$ with $\#A=\#B$. If~$A$ and~$B$ contain the same number of quadratic residues, the authors prove that~HC implies $\pi(x;q,A,B)=\Omega_{\pm} (\sqrt{x}/\log x)$. Under GRH and LI, the authors show that $\delta_{q;A,B}=\frac12$ if~$A$ and~$B$ contain the same number of quadratic residues while $\delta_{q;A,B}<\frac12$ if~$A$ contains more quadratic residues than~$B$.\par Using these results, the authors study the oscillation and the limiting logarithmic distribution of the function $S_k(x)=\sum_{p \leq x} (-1)^{s_k(p)}$, where $s_k(p)$ is the sum of the last $k$ binary digits of $p$ (note $S_2(x)=\pi(x;4,3,1)-1$ recovers Chebyshev's bias \cite{1853.Chebyshev}). The first result above implies that $S_k(x)$ changes sign infinitely often for all $k\ge2$, while the second result above implies that $\delta(\{x\colon S_k(x)>0\}) = \frac12$ for $k\ge4$. (In both cases the assumptions refer to Dirichlet $L$-functions modulo~$2^k$.) The authors further compute that $\delta(\{x\colon S_3(x)>0\})=\delta(\{x\colon \pi(x;8,3)+\pi(x;8,5)>\pi(x;8,1)+\pi(x;8,7)\})\approx 0.9822$ using the approach of Rubinstein--Sarnak~\cite{MR1329368}.\par This article cites \cite{MR1985941,MR3151605,MR1555148,MR1251926,MR0233782,MR0146157,MR1329368}.}
      \field{issn}{0208-6573}
      \field{journaltitle}{Funct. Approx. Comment. Math.}
      \field{number}{1}
      \field{title}{On the race between primes with an odd versus an even sum of the last {$k$} binary digits}
      \field{volume}{61}
      \field{year}{2019}
      \true{nocite}
      \field{pages}{7\bibrangedash 25}
      \range{pages}{19}
      \verb{doi}
      \verb 10.7169/facm/1687
      \endverb
      \verb{mrnumber}
      \verb 4012359
      \endverb
      \verb{urlraw}
      \verb https://doi.org/10.7169/facm/1687
      \endverb
      \verb{url}
      \verb https://doi.org/10.7169/facm/1687
      \endverb
    \endentry
    \entry{MR1511360}{article}{}
      \name{author}{1}{}{%
        {{hash=45dbea41eb1e7f07fcfe10d098ce7670}{%
           family={Landau},
           familyi={L\bibinitperiod},
           given={E.},
           giveni={E\bibinitperiod}}}%
      }
      \list{language}{1}{%
        {German}%
      }
      \strng{namehash}{45dbea41eb1e7f07fcfe10d098ce7670}
      \strng{fullhash}{45dbea41eb1e7f07fcfe10d098ce7670}
      \strng{bibnamehash}{45dbea41eb1e7f07fcfe10d098ce7670}
      \strng{authorbibnamehash}{45dbea41eb1e7f07fcfe10d098ce7670}
      \strng{authornamehash}{45dbea41eb1e7f07fcfe10d098ce7670}
      \strng{authorfullhash}{45dbea41eb1e7f07fcfe10d098ce7670}
      \field{extraname}{1}
      \field{sortinit}{L}
      \field{sortinithash}{dad3efd0836470093a7b4a7bb756eb8c}
      \field{labelprefix}{A}
      \field{labelnamesource}{author}
      \field{labeltitlesource}{title}
      \field{annotation}{Phragmén~\cite{1891.Phragmen} proved Chebyshev's claim that~$1$ is a limit point of the function $E^\pi(x;4,3,1)$, using the theory of functions of complex arguments. The author of this article gives a simpler proof based on the theory of Dirichlet series.\par This article cites~\cite{MR1511206}.}
      \field{issn}{0025-5831}
      \field{journaltitle}{Math. Ann.}
      \field{number}{4}
      \field{title}{Über einen {S}atz von {T}schebyschef}
      \field{volume}{61}
      \field{year}{1906}
      \true{nocite}
      \field{pages}{527\bibrangedash 550}
      \range{pages}{24}
      \verb{doi}
      \verb 10.1007/BF01449495
      \endverb
      \verb{mrnumber}
      \verb 1511360
      \endverb
      \verb{urlraw}
      \verb https://doi.org/10.1007/BF01449495
      \endverb
      \verb{url}
      \verb https://doi.org/10.1007/BF01449495
      \endverb
    \endentry
    \entry{1909.Landau}{book}{}
      \name{author}{1}{}{%
        {{hash=45dbea41eb1e7f07fcfe10d098ce7670}{%
           family={Landau},
           familyi={L\bibinitperiod},
           given={E.},
           giveni={E\bibinitperiod}}}%
      }
      \list{language}{1}{%
        {German}%
      }
      \list{publisher}{1}{%
        {Leipzig und Berlin, B. G. Teubner}%
      }
      \strng{namehash}{45dbea41eb1e7f07fcfe10d098ce7670}
      \strng{fullhash}{45dbea41eb1e7f07fcfe10d098ce7670}
      \strng{bibnamehash}{45dbea41eb1e7f07fcfe10d098ce7670}
      \strng{authorbibnamehash}{45dbea41eb1e7f07fcfe10d098ce7670}
      \strng{authornamehash}{45dbea41eb1e7f07fcfe10d098ce7670}
      \strng{authorfullhash}{45dbea41eb1e7f07fcfe10d098ce7670}
      \field{extraname}{2}
      \field{sortinit}{L}
      \field{sortinithash}{dad3efd0836470093a7b4a7bb756eb8c}
      \field{labelprefix}{A}
      \field{labelnamesource}{author}
      \field{labeltitlesource}{title}
      \field{annotation}{In Sections~199 and~200, the author generalizes his proof from~\cite{MR1511360}, showing that $(c_q(b)-c_q(a))/\phi(q)$ is a limit point of the function $E^\pi(x;q,a,b)$. He also shows that if $c>\sqrt{1/4+\gamma_1^2}$, then both $E^\Pi(x) > \frac1{29}$ and $E^\Pi(x) < -\frac1{29}$ occur infinitely often, as well as the same result with~$\Pi(x)$ in the definition of~$E^\Pi$ replaced by $\pi(x)+\frac12\Li(\sqrt x)$.}
      \field{title}{Handbuch der Lehre von der Verteilung der Primzahlen. 2 Bände}
      \field{year}{1909}
      \true{nocite}
      \field{pages}{xviii+pp. 1\bibrangedash 564\bibrangessep ix+pp. 565\bibrangedash 961}
      \range{pages}{-1}
    \endentry
    \entry{MR1544293A}{article}{}
      \name{author}{1}{}{%
        {{hash=45dbea41eb1e7f07fcfe10d098ce7670}{%
           family={Landau},
           familyi={L\bibinitperiod},
           given={E.},
           giveni={E\bibinitperiod}}}%
      }
      \list{language}{1}{%
        {German}%
      }
      \strng{namehash}{45dbea41eb1e7f07fcfe10d098ce7670}
      \strng{fullhash}{45dbea41eb1e7f07fcfe10d098ce7670}
      \strng{bibnamehash}{45dbea41eb1e7f07fcfe10d098ce7670}
      \strng{authorbibnamehash}{45dbea41eb1e7f07fcfe10d098ce7670}
      \strng{authornamehash}{45dbea41eb1e7f07fcfe10d098ce7670}
      \strng{authorfullhash}{45dbea41eb1e7f07fcfe10d098ce7670}
      \field{extraname}{3}
      \field{sortinit}{L}
      \field{sortinithash}{dad3efd0836470093a7b4a7bb756eb8c}
      \field{labelprefix}{A}
      \field{labelnamesource}{author}
      \field{labeltitlesource}{title}
      \field{annotation}{Chebyshev asserted that $\pi_e(x,\chi_{-4})\to-\infty$ as $x\to\infty$; the author shows that this assertion implies GRH for $L(s,\chi_{-4})$. Assuming this GRH, he shows that $\pi(x;4,3,1) \ll \sqrt x\log x$. Assuming the original assertion of Chebyshev, he proves that $\sum_p \chi_{-4}(p)f(p)$ converges whenever $f(p)$ is strictly decreasing and satisfies $f(p) \ll x^{-1/2-\delta}$ for some $\delta>0$.\par The author also examines the series \[ \frac{L'(s,\chi_D)}{L(s,\chi_D)}=-\sum_{p} \frac{\chi_D(p) \log p}{p^s-\chi_D(p)}, \] the identity originally valid for $\Re s>1$; he shows that this series diverges at $s=\frac12$, disproving a conjecture of Lerch. Finally, the author shows that $\sum_{n=1}^\infty \mu(n)/\sqrt n$ diverges, disproving a conjecture of Stieltjes, and also shows that $M(x) = \Omega(\sqrt x)$.\par This article cites~\cite{1853.Chebyshev}.}
      \field{issn}{0025-5874}
      \field{journaltitle}{Math. Z.}
      \field{number}{2-3}
      \field{title}{Über einige ältere {V}ermutungen und {B}ehauptungen in der {P}rimzahltheorie}
      \field{volume}{1}
      \field{year}{1918}
      \true{nocite}
      \field{pages}{1\bibrangedash 24}
      \range{pages}{24}
      \verb{doi}
      \verb 10.1007/BF01203613
      \endverb
      \verb{mrnumber}
      \verb 1544293
      \endverb
      \verb{urlraw}
      \verb https://doi.org/10.1007/BF01203613
      \endverb
      \verb{url}
      \verb https://doi.org/10.1007/BF01203613
      \endverb
    \endentry
    \entry{MR1544293B}{article}{}
      \name{author}{1}{}{%
        {{hash=45dbea41eb1e7f07fcfe10d098ce7670}{%
           family={Landau},
           familyi={L\bibinitperiod},
           given={E.},
           giveni={E\bibinitperiod}}}%
      }
      \list{language}{1}{%
        {German}%
      }
      \strng{namehash}{45dbea41eb1e7f07fcfe10d098ce7670}
      \strng{fullhash}{45dbea41eb1e7f07fcfe10d098ce7670}
      \strng{bibnamehash}{45dbea41eb1e7f07fcfe10d098ce7670}
      \strng{authorbibnamehash}{45dbea41eb1e7f07fcfe10d098ce7670}
      \strng{authornamehash}{45dbea41eb1e7f07fcfe10d098ce7670}
      \strng{authorfullhash}{45dbea41eb1e7f07fcfe10d098ce7670}
      \field{extraname}{4}
      \field{sortinit}{L}
      \field{sortinithash}{dad3efd0836470093a7b4a7bb756eb8c}
      \field{labelprefix}{A}
      \field{labelnamesource}{author}
      \field{labeltitlesource}{title}
      \field{annotation}{The author gives a simplified proof of the result of Hardy and Littlewood~\cite[Section 2.3]{MR1555148} that GRH for $L(s,\chi_{-4})$ implies Chebyshev's assertion $\pi_e(x,\chi_{-4})\to-\infty$ as $x\to\infty$.}
      \field{issn}{0025-5874}
      \field{journaltitle}{Math. Z.}
      \field{number}{2-3}
      \field{title}{Über einige ältere {V}ermutungen und {B}ehauptungen in der {P}rimzahltheorie}
      \field{volume}{1}
      \field{year}{1918}
      \true{nocite}
      \field{pages}{213\bibrangedash 219}
      \range{pages}{7}
      \verb{doi}
      \verb 10.1007/BF01203613
      \endverb
      \verb{mrnumber}
      \verb 1544293
      \endverb
      \verb{urlraw}
      \verb https://doi.org/10.1007/BF01203613
      \endverb
      \verb{url}
      \verb https://doi.org/10.1007/BF01203613
      \endverb
    \endentry
    \entry{MR0068565}{book}{}
      \name{author}{1}{}{%
        {{hash=45dbea41eb1e7f07fcfe10d098ce7670}{%
           family={Landau},
           familyi={L\bibinitperiod},
           given={E.},
           giveni={E\bibinitperiod}}}%
      }
      \list{language}{1}{%
        {German}%
      }
      \list{publisher}{1}{%
        {Chelsea Publishing Co., New York}%
      }
      \strng{namehash}{45dbea41eb1e7f07fcfe10d098ce7670}
      \strng{fullhash}{45dbea41eb1e7f07fcfe10d098ce7670}
      \strng{bibnamehash}{45dbea41eb1e7f07fcfe10d098ce7670}
      \strng{authorbibnamehash}{45dbea41eb1e7f07fcfe10d098ce7670}
      \strng{authornamehash}{45dbea41eb1e7f07fcfe10d098ce7670}
      \strng{authorfullhash}{45dbea41eb1e7f07fcfe10d098ce7670}
      \field{extraname}{5}
      \field{sortinit}{L}
      \field{sortinithash}{dad3efd0836470093a7b4a7bb756eb8c}
      \field{labelprefix}{A}
      \field{labelnamesource}{author}
      \field{labeltitlesource}{title}
      \field{annotation}{This is the second edition of~\cite{1909.Landau}.}
      \field{note}{2d ed; With an appendix by Paul T. Bateman}
      \field{title}{Handbuch der Lehre von der Verteilung der Primzahlen. 2 Bände}
      \field{year}{1953}
      \true{nocite}
      \field{pages}{xviii+pp. 1\bibrangedash 564\bibrangessep ix+pp. 565\bibrangedash 1001}
      \range{pages}{-1}
      \verb{mrnumber}
      \verb 0068565
      \endverb
    \endentry
    \entry{MR1916279}{article}{}
      \name{author}{1}{}{%
        {{hash=5cf7064e6cd6ca7d20eeb2d3e5e04811}{%
           family={Lau},
           familyi={L\bibinitperiod},
           given={Yuk-Kam},
           giveni={Y\bibinithyphendelim K\bibinitperiod}}}%
      }
      \strng{namehash}{5cf7064e6cd6ca7d20eeb2d3e5e04811}
      \strng{fullhash}{5cf7064e6cd6ca7d20eeb2d3e5e04811}
      \strng{bibnamehash}{5cf7064e6cd6ca7d20eeb2d3e5e04811}
      \strng{authorbibnamehash}{5cf7064e6cd6ca7d20eeb2d3e5e04811}
      \strng{authornamehash}{5cf7064e6cd6ca7d20eeb2d3e5e04811}
      \strng{authorfullhash}{5cf7064e6cd6ca7d20eeb2d3e5e04811}
      \field{sortinit}{L}
      \field{sortinithash}{dad3efd0836470093a7b4a7bb756eb8c}
      \field{labelprefix}{A}
      \field{labelnamesource}{author}
      \field{labeltitlesource}{title}
      \field{annotation}{This article investigates the limiting distribution of almost periodic functions. The author proves the existence of the limiting distribution of a class of functions which are bounded and can be approximated by periodic functions in $L_1$-norm. By using the quantitative version of the continuity theorem, the author is able to investigate the rate of convergence of some cases. Compared to Heath-Brown's work~\cites{MR1159354} on distribution of the error term in the Dirichlet divisor problem, the result here is more general and requires weaker hypotheses.\par This article cites ~\cites{MR1159354,MR1329368}.}
      \field{issn}{0022-314X}
      \field{journaltitle}{J. Number Theory}
      \field{number}{2}
      \field{title}{On the existence of limiting distributions of some number-theoretic error terms}
      \field{volume}{94}
      \field{year}{2002}
      \true{nocite}
      \field{pages}{359\bibrangedash 374}
      \range{pages}{16}
      \verb{doi}
      \verb 10.1006/jnth.2001.2734
      \endverb
      \verb{mrnumber}
      \verb 1916279
      \endverb
    \endentry
    \entry{arXiv.Lay}{online}{}
      \name{author}{1}{}{%
        {{hash=00f8cf5ecf192c1d41eced997c115cb8}{%
           family={Lay},
           familyi={L\bibinitperiod},
           given={J.},
           giveni={J\bibinitperiod}}}%
      }
      \strng{namehash}{00f8cf5ecf192c1d41eced997c115cb8}
      \strng{fullhash}{00f8cf5ecf192c1d41eced997c115cb8}
      \strng{bibnamehash}{00f8cf5ecf192c1d41eced997c115cb8}
      \strng{authorbibnamehash}{00f8cf5ecf192c1d41eced997c115cb8}
      \strng{authornamehash}{00f8cf5ecf192c1d41eced997c115cb8}
      \strng{authorfullhash}{00f8cf5ecf192c1d41eced997c115cb8}
      \field{sortinit}{L}
      \field{sortinithash}{dad3efd0836470093a7b4a7bb756eb8c}
      \field{labelprefix}{A}
      \field{labelnamesource}{author}
      \field{labeltitlesource}{title}
      \field{annotation}{The author shows that the functions $E^\theta_r(x)$ and $E^\pi_r(x)$ change sign infinitely often. Under RH, following a similar proof used by Diamond and Pintz~\cite{MR2605532}, the author shows $E^\theta_r(x) = \Omega_{\pm} \big( \log \log \log x \big)$ and the same for $E^\pi_r(x)$. Similar to the work of Lamzouri~\cite{MR3455269}, assuming RH and LI, the author proves that both the logarithmic densities of $\{x>1\colon E^\theta_r(x)>0\}$ and $\{x>1\colon E^\pi_r(x)>0\}$ equal $\delta(\li,\pi) \approx 0.99999974$.\par This article cites~\cites{MR3261965,MR2605532,1932.Ingham,MR3455269,MR0137689,MR1329368,MR0004255,MR3455269}.}
      \field{title}{Sign changes in Mertens' first and second theorems}
      \field{year}{2015}
      \true{nocite}
      \verb{urlraw}
      \verb https://arxiv.org/abs/1505.03589
      \endverb
      \verb{url}
      \verb https://arxiv.org/abs/1505.03589
      \endverb
    \endentry
    \entry{MR2037293}{article}{}
      \name{author}{1}{}{%
        {{hash=13724bf2435ff7e0fb4fe7a4020ba08c}{%
           family={Leboeuf},
           familyi={L\bibinitperiod},
           given={P.},
           giveni={P\bibinitperiod}}}%
      }
      \strng{namehash}{13724bf2435ff7e0fb4fe7a4020ba08c}
      \strng{fullhash}{13724bf2435ff7e0fb4fe7a4020ba08c}
      \strng{bibnamehash}{13724bf2435ff7e0fb4fe7a4020ba08c}
      \strng{authorbibnamehash}{13724bf2435ff7e0fb4fe7a4020ba08c}
      \strng{authornamehash}{13724bf2435ff7e0fb4fe7a4020ba08c}
      \strng{authorfullhash}{13724bf2435ff7e0fb4fe7a4020ba08c}
      \field{sortinit}{L}
      \field{sortinithash}{dad3efd0836470093a7b4a7bb756eb8c}
      \field{labelprefix}{A}
      \field{labelnamesource}{author}
      \field{labeltitlesource}{title}
      \field{annotation}{In Section~$4$, the author heuristically computes higher moments of $E^\pi(x)$ from the explicit formula, determining that the third moment should be asymptotic to $-\frac3{\sqrt x} \sum_{\gamma>0} (\frac14+\gamma^2)^2$; while this vanishes in the limit as $x\to\infty$, it demonstrates an asymmetry for finite~$x$. The author also calculates the limiting fourth moment and notes that it differs from that of the Gaussian approximation. In addition, the author shows that there are persistent correlations between values of $E^\pi(x)$ at arguments separated by up to~$\frac x{\gamma_1}$, where $\gamma_1\approx 14.135$ is the smallest ordinate of a nontrivial of $\zeta(s)$. These observations are supported by numerical calculations. \par This article cites~\cites{MR1329368,MR1651741,MR0202686}.}
      \field{issn}{1424-0637}
      \field{journaltitle}{Ann. Henri Poincaré}
      \field{number}{suppl. 2}
      \field{title}{Prime correlations and fluctuations}
      \field{volume}{4}
      \field{year}{2003}
      \true{nocite}
      \field{pages}{S727\bibrangedash S752}
      \range{pages}{-1}
      \verb{doi}
      \verb 10.1007/s00023-003-0958-2
      \endverb
      \verb{mrnumber}
      \verb 2037293
      \endverb
    \endentry
    \entry{MR0083001}{article}{}
      \name{author}{1}{}{%
        {{hash=0d7fbc6c344349f0368ab70890c25508}{%
           family={Leech},
           familyi={L\bibinitperiod},
           given={J.},
           giveni={J\bibinitperiod}}}%
      }
      \strng{namehash}{0d7fbc6c344349f0368ab70890c25508}
      \strng{fullhash}{0d7fbc6c344349f0368ab70890c25508}
      \strng{bibnamehash}{0d7fbc6c344349f0368ab70890c25508}
      \strng{authorbibnamehash}{0d7fbc6c344349f0368ab70890c25508}
      \strng{authornamehash}{0d7fbc6c344349f0368ab70890c25508}
      \strng{authorfullhash}{0d7fbc6c344349f0368ab70890c25508}
      \field{sortinit}{L}
      \field{sortinithash}{dad3efd0836470093a7b4a7bb756eb8c}
      \field{labelprefix}{A}
      \field{labelnamesource}{author}
      \field{labeltitlesource}{title}
      \field{annotation}{The author uses the EDSAC at Cambridge to compute $\pi(x;4,1)$ and $\pi(x;4,3)$ for $x$ up to $3\times10^6$. He discovers that $\pi(x;4,1)>\pi(x;4,3)$ at $x = 26,861$, for which $\pi(x;4,1) = 1,473$ and $\pi(x;4,3) = 1,472$. The other values of $x$ above $26,863$ for which $\pi(x;4,1)>\pi(x;4,3)$ are between $616,000$ and $634,000$; the greatest difference found is at $x = 623,681$, for which $\pi(x;4,1) = 25,444$ and $\pi(x;4,3) = 25,436$.\par The author notes that $\pi_i(x) = 2\pi(x;4,1) + \pi(\sqrt x;4,3) + 1$, the number of Gaussian primes with norm at most $x$ (up to associates, and for $x\ge2$), is consequently large near this latter range as well; the most extreme value found is at $x = 617,537$, for which $\pi_i(x)= 50,509 \approx \li(x)+19.5$. \par When examining the explicit formula for $\pi(x;4,3,1)$ at $x = 620,000$, the author found that the first $20$ pairs of zeros of $L(s,\chi_{-4})$, whose imaginary parts ranged from $\pm 6.020948$ to $\pm 49.723129$, included $16$ pairs that give negative contributions to the explicit formula, while subsequent zeros gave more or less random contributions.\par This article cites~\cites{1932.Ingham,MR1573970}.}
      \field{issn}{0024-6107}
      \field{journaltitle}{J. London Math. Soc.}
      \field{title}{Note on the distribution of prime numbers}
      \field{volume}{32}
      \field{year}{1957}
      \true{nocite}
      \field{pages}{56\bibrangedash 58}
      \range{pages}{3}
      \verb{doi}
      \verb 10.1112/jlms/s1-32.1.56
      \endverb
      \verb{mrnumber}
      \verb 0083001
      \endverb
      \verb{urlraw}
      \verb https://doi.org/10.1112/jlms/s1-32.1.56
      \endverb
      \verb{url}
      \verb https://doi.org/10.1112/jlms/s1-32.1.56
      \endverb
    \endentry
    \entry{MR0202686}{article}{}
      \name{author}{1}{}{%
        {{hash=d67ac0f50fddf9c6d09a97d86037e761}{%
           family={Lehman},
           familyi={L\bibinitperiod},
           given={R.\bibnamedelimi S.},
           giveni={R\bibinitperiod\bibinitdelim S\bibinitperiod}}}%
      }
      \strng{namehash}{d67ac0f50fddf9c6d09a97d86037e761}
      \strng{fullhash}{d67ac0f50fddf9c6d09a97d86037e761}
      \strng{bibnamehash}{d67ac0f50fddf9c6d09a97d86037e761}
      \strng{authorbibnamehash}{d67ac0f50fddf9c6d09a97d86037e761}
      \strng{authornamehash}{d67ac0f50fddf9c6d09a97d86037e761}
      \strng{authorfullhash}{d67ac0f50fddf9c6d09a97d86037e761}
      \field{extraname}{1}
      \field{sortinit}{L}
      \field{sortinithash}{dad3efd0836470093a7b4a7bb756eb8c}
      \field{labelprefix}{A}
      \field{labelnamesource}{author}
      \field{labeltitlesource}{title}
      \field{annotation}{The author shows that $\pi(x)>\li(x)$ for some $x<1.65\times10^{1,165}$; more precisely, he shows that there is an interval of length $10^{500}$ between $1.53\times10^{1,165}$ and $1.65\times10^{1,165}$ on which $\pi(x)>\li(x)$. The computations, involving the zeros of $\zeta(s)$ up to height $12,000$, were performed on an IBM 7090 at Berkeley. \par This article cites~\cites{MR0067145,MR0133308}.}
      \field{issn}{0065-1036}
      \field{journaltitle}{Acta Arith.}
      \field{title}{On the difference {$\pi (x)-\li (x)$}}
      \field{volume}{11}
      \field{year}{1966}
      \true{nocite}
      \field{pages}{397\bibrangedash 410}
      \range{pages}{14}
      \verb{doi}
      \verb 10.4064/aa-11-4-397-410
      \endverb
      \verb{mrnumber}
      \verb 0202686
      \endverb
    \endentry
    \entry{MR0120198}{article}{}
      \name{author}{1}{}{%
        {{hash=019b07c7177a5ca71ef4716a49411b64}{%
           family={Lehman},
           familyi={L\bibinitperiod},
           given={R.\bibnamedelimi Sherman},
           giveni={R\bibinitperiod\bibinitdelim S\bibinitperiod}}}%
      }
      \strng{namehash}{019b07c7177a5ca71ef4716a49411b64}
      \strng{fullhash}{019b07c7177a5ca71ef4716a49411b64}
      \strng{bibnamehash}{019b07c7177a5ca71ef4716a49411b64}
      \strng{authorbibnamehash}{019b07c7177a5ca71ef4716a49411b64}
      \strng{authornamehash}{019b07c7177a5ca71ef4716a49411b64}
      \strng{authorfullhash}{019b07c7177a5ca71ef4716a49411b64}
      \field{extraname}{2}
      \field{sortinit}{L}
      \field{sortinithash}{dad3efd0836470093a7b4a7bb756eb8c}
      \field{labelprefix}{A}
      \field{labelnamesource}{author}
      \field{labeltitlesource}{title}
      \field{annotation}{The author describes computations leading to an explicit counterexample to Pólya's problem, showing that $L(906,180,359)=1$ is a small value for which the inequality fails. In order to speed up the computation for $L(x)$, which was performed on an IBM 704 at Berkeley, the author uses (a slightly more complicated version of) the recursive formula $$ L(x)=\sum_{m \leq x/m} \mu(m) \bigg \lfloor \sqrt{\frac{x}{m}}\bigg \rfloor -\sum_{k <v} \lambda(k)\bigg( \bigg \lfloor\frac{x}{km}\bigg \rfloor-\bigg \lfloor\frac{x}{mv}\bigg \rfloor \bigg) -\sum_{x/w<l \leq x/v} L \bigg(\frac{x}{l}\bigg) \sum_{\substack{m \mid l \\ m \leq x/w}} \mu(m). $$ The author uses a conditional truncated explicit formula for a weighted variant of $L(e^u)$ to identify promising candidates for positive values of $L(x)$.\par This article cites~\cite{1912.vonSterneck,1919.Polya,MR0006202,1932.Ingham,MR0104638,MR43841,MR51857}.}
      \field{issn}{0025-5718}
      \field{journaltitle}{Math. Comp.}
      \field{title}{On {L}iouville's function}
      \field{volume}{14}
      \field{year}{1960}
      \true{nocite}
      \field{pages}{311\bibrangedash 320}
      \range{pages}{10}
      \verb{doi}
      \verb 10.2307/2003890
      \endverb
      \verb{mrnumber}
      \verb 0120198
      \endverb
      \verb{urlraw}
      \verb https://doi.org/10.2307/2003890
      \endverb
      \verb{url}
      \verb https://doi.org/10.2307/2003890
      \endverb
    \endentry
    \entry{MR115965}{article}{}
      \name{author}{2}{}{%
        {{hash=63196bddebd8e713df6bac51c837c623}{%
           family={Lehmer},
           familyi={L\bibinitperiod},
           given={D.\bibnamedelimi H.},
           giveni={D\bibinitperiod\bibinitdelim H\bibinitperiod}}}%
        {{hash=74d133cbc34e02461a902602a69cfb30}{%
           family={Selberg},
           familyi={S\bibinitperiod},
           given={S.},
           giveni={S\bibinitperiod}}}%
      }
      \strng{namehash}{9aa6cea78280f1fd434f36a0a19ee629}
      \strng{fullhash}{9aa6cea78280f1fd434f36a0a19ee629}
      \strng{bibnamehash}{9aa6cea78280f1fd434f36a0a19ee629}
      \strng{authorbibnamehash}{9aa6cea78280f1fd434f36a0a19ee629}
      \strng{authornamehash}{9aa6cea78280f1fd434f36a0a19ee629}
      \strng{authorfullhash}{9aa6cea78280f1fd434f36a0a19ee629}
      \field{sortinit}{L}
      \field{sortinithash}{dad3efd0836470093a7b4a7bb756eb8c}
      \field{labelprefix}{A}
      \field{labelnamesource}{author}
      \field{labeltitlesource}{title}
      \field{annotation}{Using Landau's Theorem, the authors show that $\fA_1^{M_r}(x)-K$ changes signs infinitely often for any constant~$K$. Numerical calculations show that the first $56$ sign changes of $\fA_1^{M_r}(x)-2$ are nearly in a geometric progression, and a heuristic explanation of this phenomenon is derived from the explicit formula $\fA_1^{M_r}(x)-2=-\sum {x^\rho}/{\rho(1-\rho) \zeta'(\rho)}$.}
      \field{issn}{0065-1036}
      \field{journaltitle}{Acta Arith.}
      \field{title}{A sum involving the function of {M}öbius}
      \field{volume}{6}
      \field{year}{1960}
      \true{nocite}
      \field{pages}{111\bibrangedash 114}
      \range{pages}{4}
      \verb{doi}
      \verb 10.4064/aa-6-1-111-114
      \endverb
      \verb{mrnumber}
      \verb 115965
      \endverb
      \verb{urlraw}
      \verb https://doi.org/10.4064/aa-6-1-111-114
      \endverb
      \verb{url}
      \verb https://doi.org/10.4064/aa-6-1-111-114
      \endverb
    \endentry
    \entry{MR3624386}{article}{}
      \name{author}{2}{}{%
        {{hash=5cb6f78b503191fd0bfbfcba676ac134}{%
           family={Lemke\bibnamedelima Oliver},
           familyi={L\bibinitperiod\bibinitdelim O\bibinitperiod},
           given={R.\bibnamedelimi J.},
           giveni={R\bibinitperiod\bibinitdelim J\bibinitperiod}}}%
        {{hash=549bdccc730fda930fd587b6ae3caaba}{%
           family={Soundararajan},
           familyi={S\bibinitperiod},
           given={K.},
           giveni={K\bibinitperiod}}}%
      }
      \strng{namehash}{ae4a066eca6b7243e3cef24fd0e898c7}
      \strng{fullhash}{ae4a066eca6b7243e3cef24fd0e898c7}
      \strng{bibnamehash}{ae4a066eca6b7243e3cef24fd0e898c7}
      \strng{authorbibnamehash}{ae4a066eca6b7243e3cef24fd0e898c7}
      \strng{authornamehash}{ae4a066eca6b7243e3cef24fd0e898c7}
      \strng{authorfullhash}{ae4a066eca6b7243e3cef24fd0e898c7}
      \field{extraname}{1}
      \field{sortinit}{L}
      \field{sortinithash}{dad3efd0836470093a7b4a7bb756eb8c}
      \field{labelprefix}{A}
      \field{labelnamesource}{author}
      \field{labeltitlesource}{title}
      \field{annotation}{Given a tuple $\textbf{a} = (a_1,\ldots,a_r)$ of reduced residues modulo~$q$, let \[ \pi(x;q,\textbf{a}) = \# \{ p_n \leq x \colon p_{n+i-1} \equiv a_i \mod{q} \text{ for each } 1 \leq i \leq r\}\] count the occurrences of the pattern of residues defined by~$\textbf{a}$. The authors observe (among other things) that repeated residues appear less frequently than changing ones; for example, for $x_0 = p_{10^8}$, we have $\pi(x_0;10,(1,3)) = 7,429,438$ but $\pi(x_0;10,(1,1)) = 4,623,042$. The authors conjecture that all patterns occur equally often in the limit, but that lower order terms create predictable biases; in contrast to prime number races and their infinity of sign changes, some of these inequalities should always hold, such as $\pi(x;3,(1,-1)) > \pi(x;3,(1,1))$ for $x \geq 5$. The authors provide numerical evidence for their conjectures, as well as heuristic justification related to the Hardy--Littlewood prime $k$-tuples conjecture.\par This article cites~\cites{MR1329368,MR0466043,MR2202918,MR3521338}.}
      \field{issn}{1091-6490}
      \field{journaltitle}{Proc. Natl. Acad. Sci. USA}
      \field{number}{31}
      \field{title}{Unexpected biases in the distribution of consecutive primes}
      \field{volume}{113}
      \field{year}{2016}
      \true{nocite}
      \field{pages}{E4446\bibrangedash E4454}
      \range{pages}{-1}
      \verb{doi}
      \verb 10.1073/pnas.1605366113
      \endverb
      \verb{mrnumber}
      \verb 3624386
      \endverb
      \verb{urlraw}
      \verb https://doi.org/10.1073/pnas.1605366113
      \endverb
      \verb{url}
      \verb https://doi.org/10.1073/pnas.1605366113
      \endverb
    \endentry
    \entry{MR4043824}{article}{}
      \name{author}{2}{}{%
        {{hash=f5fefda26be806a5e5f1affbb4e80fe9}{%
           family={Lemke\bibnamedelima Oliver},
           familyi={L\bibinitperiod\bibinitdelim O\bibinitperiod},
           given={Robert\bibnamedelima J.},
           giveni={R\bibinitperiod\bibinitdelim J\bibinitperiod}}}%
        {{hash=db8c7134ab46023e3103bbf8edcf2ca5}{%
           family={Soundararajan},
           familyi={S\bibinitperiod},
           given={Kannan},
           giveni={K\bibinitperiod}}}%
      }
      \strng{namehash}{b4ea39c0bfa1f7e30e1685d948198229}
      \strng{fullhash}{b4ea39c0bfa1f7e30e1685d948198229}
      \strng{bibnamehash}{b4ea39c0bfa1f7e30e1685d948198229}
      \strng{authorbibnamehash}{b4ea39c0bfa1f7e30e1685d948198229}
      \strng{authornamehash}{b4ea39c0bfa1f7e30e1685d948198229}
      \strng{authorfullhash}{b4ea39c0bfa1f7e30e1685d948198229}
      \field{extraname}{2}
      \field{sortinit}{L}
      \field{sortinithash}{dad3efd0836470093a7b4a7bb756eb8c}
      \field{labelprefix}{A}
      \field{labelnamesource}{author}
      \field{labeltitlesource}{title}
      \field{annotation}{The authors continue their study of biases in the patterns of consecutive primes\mod q. Their conjectures in~\cite{MR3624386} contained a tertiary main term with complicated constants $c_2(q,{\bf a})$ whose distribution was not easily understood. In this article, they connect these constants with both the discrete Fourier transform of Dedekind sums and the error term $R(x) = \sum_{n\le x} \phi(n) - \frac{3}{\pi^2}x^2$. They show that $q^{-1}c_2(q, (a,b))$ has a continuous limiting distribution as $q\to\infty$ (connected to that of the related Dedekind sums), and that $x^{-1}R(x)$ has a limiting distribution with doubly exponential decay. \par This article cites~\cites{MR3624386, MR1297543}.}
      \field{issn}{0305-0041}
      \field{journaltitle}{Math. Proc. Cambridge Philos. Soc.}
      \field{number}{1}
      \field{title}{The distribution of consecutive prime biases and sums of sawtooth random variables}
      \field{volume}{168}
      \field{year}{2020}
      \true{nocite}
      \field{pages}{149\bibrangedash 169}
      \range{pages}{21}
      \verb{doi}
      \verb 10.1017/s0305004118000592
      \endverb
      \verb{mrnumber}
      \verb 4043824
      \endverb
      \verb{urlraw}
      \verb https://doi-org.ezproxy.library.ubc.ca/10.1017/s0305004118000592
      \endverb
      \verb{url}
      \verb https://doi-org.ezproxy.library.ubc.ca/10.1017/s0305004118000592
      \endverb
    \endentry
    \entry{MR0439774}{inproceedings}{}
      \name{author}{1}{}{%
        {{hash=66650e0cb2c39da0339e84e2bc736eab}{%
           family={Levinson},
           familyi={L\bibinitperiod},
           given={N.},
           giveni={N\bibinitperiod}}}%
      }
      \list{publisher}{1}{%
        {North-Holland, Amsterdam}%
      }
      \strng{namehash}{66650e0cb2c39da0339e84e2bc736eab}
      \strng{fullhash}{66650e0cb2c39da0339e84e2bc736eab}
      \strng{bibnamehash}{66650e0cb2c39da0339e84e2bc736eab}
      \strng{authorbibnamehash}{66650e0cb2c39da0339e84e2bc736eab}
      \strng{authornamehash}{66650e0cb2c39da0339e84e2bc736eab}
      \strng{authorfullhash}{66650e0cb2c39da0339e84e2bc736eab}
      \field{sortinit}{L}
      \field{sortinithash}{dad3efd0836470093a7b4a7bb756eb8c}
      \field{labelprefix}{A}
      \field{labelnamesource}{author}
      \field{labeltitlesource}{title}
      \field{annotation}{Assuming that SA is false, the author shows that $\limsup_{T \to \infty} W^\pi(T)/\log T = \infty$. The author adapts a Landau-type argument of Pólya~\cite{1930.Polya,MR0263757} to treat functions with logarithmic singularities. \par This article cites~\cite{1936.Ingham,MR0133309,MR0133308,MR0439771,1930.Polya,MR0263757}.}
      \field{booktitle}{Topics in number theory ({P}roc. {C}olloq., {D}ebrecen, 1974)}
      \field{title}{On the number of sign changes of {$\pi (x)-\li x$}}
      \field{year}{1976}
      \true{nocite}
      \field{pages}{171\bibrangedash 177. Colloq. Math. Soc. János Bolyai\bibrangessep Vol. 13}
      \range{pages}{-1}
      \verb{mrnumber}
      \verb 0439774
      \endverb
    \endentry
    \entry{MR3993767}{article}{}
      \name{author}{3}{}{%
        {{hash=d6c5b18d6dc761cc36407ac5d5441ece}{%
           family={Lichtman},
           familyi={L\bibinitperiod},
           given={J.\bibnamedelimi D.},
           giveni={J\bibinitperiod\bibinitdelim D\bibinitperiod}}}%
        {{hash=feda7d3cba2a0ad4e772307fca655937}{%
           family={Martin},
           familyi={M\bibinitperiod},
           given={G.},
           giveni={G\bibinitperiod}}}%
        {{hash=a658c262e68541245018e2e8dd4064a7}{%
           family={Pomerance},
           familyi={P\bibinitperiod},
           given={C.},
           giveni={C\bibinitperiod}}}%
      }
      \strng{namehash}{c132fc8c2248d719686bb8b2ea4c74aa}
      \strng{fullhash}{c132fc8c2248d719686bb8b2ea4c74aa}
      \strng{bibnamehash}{c132fc8c2248d719686bb8b2ea4c74aa}
      \strng{authorbibnamehash}{c132fc8c2248d719686bb8b2ea4c74aa}
      \strng{authornamehash}{c132fc8c2248d719686bb8b2ea4c74aa}
      \strng{authorfullhash}{c132fc8c2248d719686bb8b2ea4c74aa}
      \field{sortinit}{L}
      \field{sortinithash}{dad3efd0836470093a7b4a7bb756eb8c}
      \field{labelprefix}{A}
      \field{labelnamesource}{author}
      \field{labeltitlesource}{title}
      \field{annotation}{Assuming GRH and LI, the authors show that the set of primes~$q$ for which $\pi(q)>\li(q)$ has a well-defined logarithmic relative density in the set of all primes, whose value equals $\delta(\pi,\li)$ (approximately $2.6\cdot10^{-7}$). The same is true of the set of primes~$q$ such that $e^{C_0} \log q > \prod_{p\le q} (1-1/p)^{-1}$, and the set of primes~$q$ such that $1/\log q > \sum_{p\ge q} 1/p\log p$. The methods apply to discrete subsets of~$\R$ that are well-distributed in almost all short intervals (such as the integers themselves). \par This article cites~\cite{MR3261965,MR1752093,MR4120786,MR3028758,MR3455269,1914.Littlewood,MR1329368}.}
      \field{journaltitle}{Proc. Amer. Math. Soc.}
      \field{number}{9}
      \field{title}{Primes in prime number races}
      \field{volume}{147}
      \field{year}{2019}
      \true{nocite}
      \field{pages}{3743\bibrangedash 3757}
      \range{pages}{15}
    \endentry
    \entry{MR4379402}{article}{}
      \name{author}{2}{}{%
        {{hash=ab91abc046ab17798c39d00cb11edfba}{%
           family={Lin},
           familyi={L\bibinitperiod},
           given={Jiawei},
           giveni={J\bibinitperiod}}}%
        {{hash=c06cf923cb7824f2455f1750eec3b292}{%
           family={Martin},
           familyi={M\bibinitperiod},
           given={Greg},
           giveni={G\bibinitperiod}}}%
      }
      \strng{namehash}{2a2204ec026337a0144e11ff56499eee}
      \strng{fullhash}{2a2204ec026337a0144e11ff56499eee}
      \strng{bibnamehash}{2a2204ec026337a0144e11ff56499eee}
      \strng{authorbibnamehash}{2a2204ec026337a0144e11ff56499eee}
      \strng{authornamehash}{2a2204ec026337a0144e11ff56499eee}
      \strng{authorfullhash}{2a2204ec026337a0144e11ff56499eee}
      \field{sortinit}{L}
      \field{sortinithash}{dad3efd0836470093a7b4a7bb756eb8c}
      \field{labelprefix}{A}
      \field{labelnamesource}{author}
      \field{labeltitlesource}{title}
      \field{annotation}{Assuming GRH and LI, the authors establish an asymptotic formula for $\delta_{q;a_1,a_2,a_3}$ in the case when $a_1^2\equiv a_2^2\equiv a_3^2\mod q$, with a power savings in $q$ in the error term; the main term, which is the arctangent of a specified algebraic function of the quantities $b(\chi) = \sum_{L(\rho,\chi)=0} 1/|\rho|^2$, arises as an orthant probability for a multivariate normal random variable. The congruence hypothesis allows for an atypical normalization of the $E^\pi(x;q,a_j)$ for which any given $\chi\mod q$ appears in at most one of the three such expressions, allowing one to model the error terms by independent random variables. The authors propose that to minimize their error terms, asymptotic formulas for $\delta_{q;a_1,\ldots,a_r}$ always be phrased in terms of orthant probabilities.\par This article cites~\cites{MR3028758,MR3773805,MR2998146,MR3063909,MR1956261,MR1329368}.}
      \field{issn}{0008-414X}
      \field{journaltitle}{Canad. J. Math.}
      \field{number}{1}
      \field{title}{Densities in certain three-way prime number races}
      \field{volume}{74}
      \field{year}{2022}
      \true{nocite}
      \field{pages}{232\bibrangedash 265}
      \range{pages}{34}
      \verb{doi}
      \verb 10.4153/S0008414X20000747
      \endverb
      \verb{mrnumber}
      \verb 4379402
      \endverb
      \verb{urlraw}
      \verb https://doi.org/10.4153/S0008414X20000747
      \endverb
      \verb{url}
      \verb https://doi.org/10.4153/S0008414X20000747
      \endverb
    \endentry
    \entry{1914.Littlewood}{article}{}
      \name{author}{1}{}{%
        {{hash=56dd48151c55dce1b8a073ecf909a238}{%
           family={Littlewood},
           familyi={L\bibinitperiod},
           given={J.\bibnamedelimi E.},
           giveni={J\bibinitperiod\bibinitdelim E\bibinitperiod}}}%
      }
      \list{language}{1}{%
        {French}%
      }
      \strng{namehash}{56dd48151c55dce1b8a073ecf909a238}
      \strng{fullhash}{56dd48151c55dce1b8a073ecf909a238}
      \strng{bibnamehash}{56dd48151c55dce1b8a073ecf909a238}
      \strng{authorbibnamehash}{56dd48151c55dce1b8a073ecf909a238}
      \strng{authornamehash}{56dd48151c55dce1b8a073ecf909a238}
      \strng{authorfullhash}{56dd48151c55dce1b8a073ecf909a238}
      \field{extraname}{1}
      \field{sortinit}{L}
      \field{sortinithash}{dad3efd0836470093a7b4a7bb756eb8c}
      \field{labelprefix}{A}
      \field{labelnamesource}{author}
      \field{labeltitlesource}{title}
      \field{annotation}{Assuming RH, the author shows that \begin{align*} \Delta^\psi(x) &= \Omega_\pm\big( \sqrt{x}\log\log\log{x} \big) \\ \Delta^\pi(x) &= \Omega_\pm\bigg( \frac{\sqrt{x}\log\log\log{x}}{\log{x}} \bigg). \end{align*} It follows that the conjectured inequality $\pi(x)<\Li(x)$ (``{\em présumée par divers auteurs pour des raisons empiriques}'') cannot hold for all values of~$x$.\par This article cites~\cite{1909.Landau}.}
      \field{journaltitle}{Comptes Rendus de l'Acad. Sci. Paris}
      \field{title}{{Sur la distribution des nombres premiers}}
      \field{volume}{158}
      \field{year}{1914}
      \true{nocite}
      \field{pages}{1869\bibrangedash 1872}
      \range{pages}{4}
    \endentry
    \entry{MR1574052}{article}{}
      \name{author}{1}{}{%
        {{hash=56dd48151c55dce1b8a073ecf909a238}{%
           family={Littlewood},
           familyi={L\bibinitperiod},
           given={J.\bibnamedelimi E.},
           giveni={J\bibinitperiod\bibinitdelim E\bibinitperiod}}}%
      }
      \strng{namehash}{56dd48151c55dce1b8a073ecf909a238}
      \strng{fullhash}{56dd48151c55dce1b8a073ecf909a238}
      \strng{bibnamehash}{56dd48151c55dce1b8a073ecf909a238}
      \strng{authorbibnamehash}{56dd48151c55dce1b8a073ecf909a238}
      \strng{authornamehash}{56dd48151c55dce1b8a073ecf909a238}
      \strng{authorfullhash}{56dd48151c55dce1b8a073ecf909a238}
      \field{extraname}{2}
      \field{sortinit}{L}
      \field{sortinithash}{dad3efd0836470093a7b4a7bb756eb8c}
      \field{labelprefix}{A}
      \field{labelnamesource}{author}
      \field{labeltitlesource}{title}
      \field{annotation}{The author shows that $\Delta^{\psi}(x)=\Omega_{\pm}(\sqrt{x})$. The main idea is to examine the recursive sequence of integrals $\fA^{\psi}_n(x)$; since $x^{-1}\fA^{\psi}_n(x)$ is an average of $\fA^{\psi}_{n-1}(x)$, it suffices to prove that $\fA^{\psi}_n(x)=\Omega_{\pm}(x^{n+\frac{1}{2}})$ when $n$ is sufficiently large.\par This article cites~\cites{1909.Landau}.}
      \field{journaltitle}{J. London Math. Soc.}
      \field{number}{1}
      \field{title}{Mathematical Notes: 3; on a Theorem Concerning the Distribution of Prime Numbers}
      \field{volume}{2}
      \field{year}{1927}
      \true{nocite}
      \field{pages}{41\bibrangedash 45}
      \range{pages}{5}
      \verb{doi}
      \verb 10.1112/jlms/s1-2.1.41
      \endverb
      \verb{mrnumber}
      \verb 1574052
      \endverb
    \endentry
    \entry{MR1575079}{article}{}
      \name{author}{1}{}{%
        {{hash=56dd48151c55dce1b8a073ecf909a238}{%
           family={Littlewood},
           familyi={L\bibinitperiod},
           given={J.\bibnamedelimi E.},
           giveni={J\bibinitperiod\bibinitdelim E\bibinitperiod}}}%
      }
      \strng{namehash}{56dd48151c55dce1b8a073ecf909a238}
      \strng{fullhash}{56dd48151c55dce1b8a073ecf909a238}
      \strng{bibnamehash}{56dd48151c55dce1b8a073ecf909a238}
      \strng{authorbibnamehash}{56dd48151c55dce1b8a073ecf909a238}
      \strng{authornamehash}{56dd48151c55dce1b8a073ecf909a238}
      \strng{authorfullhash}{56dd48151c55dce1b8a073ecf909a238}
      \field{extraname}{3}
      \field{sortinit}{L}
      \field{sortinithash}{dad3efd0836470093a7b4a7bb756eb8c}
      \field{labelprefix}{A}
      \field{labelnamesource}{author}
      \field{labeltitlesource}{title}
      \field{annotation}{This short note is mathematically concerned with the maximum values of trigonometric polynomials. The author reveals that his motivation was to establish an explicit upper bound for the first sign change of $\Delta^\pi(x)$ without assuming RH. Ultimately, the author claims to have found alternative means for finding such an upper bound.\par This note cites~\cites{MR1573970}.}
      \field{journaltitle}{J. London Math. Soc.}
      \field{number}{3}
      \field{title}{Mathematical Notes (12): An Inequality for a Sum of Cosines}
      \field{volume}{12}
      \field{year}{1937}
      \true{nocite}
      \field{pages}{217\bibrangedash 221}
      \range{pages}{5}
      \verb{doi}
      \verb 10.1112/jlms/s1-12.2.217
      \endverb
      \verb{mrnumber}
      \verb 1575079
      \endverb
    \endentry
    \entry{MR0159791}{article}{}
      \name{author}{1}{}{%
        {{hash=044834fd4db0262d94552d21e42c5f4d}{%
           family={Makai},
           familyi={M\bibinitperiod},
           given={E.},
           giveni={E\bibinitperiod}}}%
      }
      \strng{namehash}{044834fd4db0262d94552d21e42c5f4d}
      \strng{fullhash}{044834fd4db0262d94552d21e42c5f4d}
      \strng{bibnamehash}{044834fd4db0262d94552d21e42c5f4d}
      \strng{authorbibnamehash}{044834fd4db0262d94552d21e42c5f4d}
      \strng{authornamehash}{044834fd4db0262d94552d21e42c5f4d}
      \strng{authorfullhash}{044834fd4db0262d94552d21e42c5f4d}
      \field{sortinit}{M}
      \field{sortinithash}{2e5c2f51f7fa2d957f3206819bf86dc3}
      \field{labelprefix}{A}
      \field{labelnamesource}{author}
      \field{labeltitlesource}{title}
      \field{annotation}{The author shows that the constant~$B$ in the second main theorem of the power-sum method cannot be taken smaller than~$4e$. \par This article cites~\cites{MR0077579}.}
      \field{issn}{0001-5954}
      \field{journaltitle}{Acta Math. Acad. Sci. Hungar.}
      \field{title}{On a minimum problem. II}
      \field{volume}{15}
      \field{year}{1964}
      \true{nocite}
      \field{pages}{63\bibrangedash 66}
      \range{pages}{4}
      \verb{doi}
      \verb 10.1007/BF01897022
      \endverb
      \verb{mrnumber}
      \verb 0159791
      \endverb
    \endentry
    \entry{MR1956261}{inproceedings}{}
      \name{author}{1}{}{%
        {{hash=feda7d3cba2a0ad4e772307fca655937}{%
           family={Martin},
           familyi={M\bibinitperiod},
           given={G.},
           giveni={G\bibinitperiod}}}%
      }
      \list{publisher}{1}{%
        {A K Peters, Natick, MA}%
      }
      \strng{namehash}{feda7d3cba2a0ad4e772307fca655937}
      \strng{fullhash}{feda7d3cba2a0ad4e772307fca655937}
      \strng{bibnamehash}{feda7d3cba2a0ad4e772307fca655937}
      \strng{authorbibnamehash}{feda7d3cba2a0ad4e772307fca655937}
      \strng{authornamehash}{feda7d3cba2a0ad4e772307fca655937}
      \strng{authorfullhash}{feda7d3cba2a0ad4e772307fca655937}
      \field{sortinit}{M}
      \field{sortinithash}{2e5c2f51f7fa2d957f3206819bf86dc3}
      \field{labelprefix}{A}
      \field{labelnamesource}{author}
      \field{labeltitlesource}{title}
      \field{annotation}{The author begins with the known values (assuming GRH and LI), for $q=8$ and $q=12$, of $\delta_{q;a,1}$ where $a\not\equiv1\mod q$ and of $\delta_{q;a,b,c}$ where $\{a,b,c\}$ is a permutation of the three nonidentity elements of $(\Z/q\Z)^\times$. He investigates how one could have predicted the relative sizes of these densities using the values and conductors of the nonprincipal Dirichlet characters\mod q, by arguing by analogy with independent random variables with variances of different sizes.\par The author also comments on the equality of a family of variant definitions of the logarithmic density of a set of positive real numbers, as well as some conjectures on the rarity of ties $\pi(x;q,a)=\pi(x;q,b)$.\par This article cites~\cites{MR0083001,MR0476616,MR529694,MR552467,MR1329368,MR1806291}.}
      \field{booktitle}{Number theory for the millennium, {II} ({U}rbana, {IL}, 2000)}
      \field{title}{Asymmetries in the {S}hanks-{R}ényi prime number race}
      \field{year}{2002}
      \true{nocite}
      \field{pages}{403\bibrangedash 415}
      \range{pages}{13}
      \verb{mrnumber}
      \verb 1956261
      \endverb
    \endentry
    \entry{MR4591762}{article}{}
      \name{author}{3}{}{%
        {{hash=c06cf923cb7824f2455f1750eec3b292}{%
           family={Martin},
           familyi={M\bibinitperiod},
           given={Greg},
           giveni={G\bibinitperiod}}}%
        {{hash=95a7b4ff9d58f1b7a8d6dc0d7ba8604f}{%
           family={Mossinghoff},
           familyi={M\bibinitperiod},
           given={Michael},
           giveni={M\bibinitperiod}}}%
        {{hash=427a3ba3d2c76d68bc8c1bfc231edce0}{%
           family={Trudgian},
           familyi={T\bibinitperiod},
           given={Timothy},
           giveni={T\bibinitperiod}}}%
      }
      \strng{namehash}{8e4514da6e580a9113eaccb5513e18fc}
      \strng{fullhash}{8e4514da6e580a9113eaccb5513e18fc}
      \strng{bibnamehash}{8e4514da6e580a9113eaccb5513e18fc}
      \strng{authorbibnamehash}{8e4514da6e580a9113eaccb5513e18fc}
      \strng{authornamehash}{8e4514da6e580a9113eaccb5513e18fc}
      \strng{authorfullhash}{8e4514da6e580a9113eaccb5513e18fc}
      \field{sortinit}{M}
      \field{sortinithash}{2e5c2f51f7fa2d957f3206819bf86dc3}
      \field{labelprefix}{A}
      \field{labelnamesource}{author}
      \field{labeltitlesource}{title}
      \field{annotation}{The authors investigate comparative number theoretic results for a family of arithmetic functions called ``fake~$\mu$'s'', which are multiplicative functions $\digamma(n)$ taking values in $\{-1,0,1\}$; fake~$\mu$'s include many commonly studied functions such as $\mu(n)$, $\mu^2(n)$, $(-1)^{\omega(n)}$, and $\mu_k(n)$. Generalizing a technique of Tanaka~\cite{MR0605090}, the authors show that if $f(n)$ is a fake~$\mu$ with $f(p)=-1$ and $f(p^2)=1$ for all primes~$p$, then its summatory function $F(x)$ satisfes $$ \liminf _{x \rightarrow \infty} \frac{F(x)}{\sqrt{x}}<b \quad \text { and } \quad \limsup _{x \rightarrow \infty} \frac{F(x)}{\sqrt{x}}>b, $$ where~$b$ is twice the residue at~$\frac12$ of the Dirichlet series corresponding to~$f(n)$. The authors also determine the minimum and maximum of the above constant $b$ and show that both extreme values can be achieved by particular fake~$\mu$'s.\par This article cites~\cites{MR2398787, MR1806291, MR0104638, MR2994374, MR3996180,MR3739227, MR0006202,MR0120198,MR3604916, 1897.Mertens, MR3062002,MR3700043,MR2078705,MR783538,1919.Polya,MR1329368,MR584557,MR0882550,MR0605090, MR4194167,MR4132130}.}
      \field{issn}{0002-9939}
      \field{journaltitle}{Proc. Amer. Math. Soc.}
      \field{number}{8}
      \field{title}{Fake mu's}
      \field{volume}{151}
      \field{year}{2023}
      \true{nocite}
      \field{pages}{3229\bibrangedash 3244}
      \range{pages}{16}
      \verb{doi}
      \verb 10.1090/proc/16186
      \endverb
      \verb{mrnumber}
      \verb 4591762
      \endverb
      \verb{urlraw}
      \verb https://doi.org/10.1090/proc/16186
      \endverb
      \verb{url}
      \verb https://doi.org/10.1090/proc/16186
      \endverb
    \endentry
    \entry{MR4082248}{article}{}
      \name{author}{2}{}{%
        {{hash=c06cf923cb7824f2455f1750eec3b292}{%
           family={Martin},
           familyi={M\bibinitperiod},
           given={Greg},
           giveni={G\bibinitperiod}}}%
        {{hash=3d2daa810172929af247bf36116d8aaa}{%
           family={Ng},
           familyi={N\bibinitperiod},
           given={Nathan},
           giveni={N\bibinitperiod}}}%
      }
      \strng{namehash}{11060c47f44aeca22f1c306f4f674b92}
      \strng{fullhash}{11060c47f44aeca22f1c306f4f674b92}
      \strng{bibnamehash}{11060c47f44aeca22f1c306f4f674b92}
      \strng{authorbibnamehash}{11060c47f44aeca22f1c306f4f674b92}
      \strng{authornamehash}{11060c47f44aeca22f1c306f4f674b92}
      \strng{authorfullhash}{11060c47f44aeca22f1c306f4f674b92}
      \field{sortinit}{M}
      \field{sortinithash}{2e5c2f51f7fa2d957f3206819bf86dc3}
      \field{labelprefix}{A}
      \field{labelnamesource}{author}
      \field{labeltitlesource}{title}
      \field{annotation}{The authors investigate, under GRH, how to weaken the LI assumption in earlier results. Say that a zero $\frac12+i\gamma$ of some Dirichlet $L$-function\mod q is ``self-sufficient'' if it cannot be written as a rational linear combination of other such zeros (so that LI is true if and only if all such zeros are self-sufficient). The authors show that a certain finite number of self-sufficient zeros is enough for the logarithmic limiting distributions of $E^\pi(x;q,a_1,\ldots,a_r)$ to exist. If the sum of reciprocals of ordinates of self-sufficient zeros is sufficiently large in terms of~$q$, the authors show that all logarithmic densities $\delta_{q;a_1,\ldots,a_r}$ exist and are positive; if that reciprocal sum diverges, then the distribution of $E^\pi(x;q,a_1,\ldots,a_r)$ assigns mass to every open set in $\R^r$.\par This article cites~\cites{MR3261965,MR4120786,MR0332684,MR3028758,MR3151605,MR2075622,MR2202918,MR0332685,MR2994374,MR0006202,MR1105608,MR1251926,MR0176967,MR1329368,MR2753165,MR0289452,1938.Wintner}.}
      \field{issn}{0002-9947}
      \field{journaltitle}{Trans. Amer. Math. Soc.}
      \field{number}{5}
      \field{title}{Inclusive prime number races}
      \field{volume}{373}
      \field{year}{2020}
      \true{nocite}
      \field{pages}{3561\bibrangedash 3607}
      \range{pages}{47}
      \verb{doi}
      \verb 10.1090/tran/7996
      \endverb
      \verb{mrnumber}
      \verb 4082248
      \endverb
      \verb{urlraw}
      \verb https://doi.org/10.1090/tran/7996
      \endverb
      \verb{url}
      \verb https://doi.org/10.1090/tran/7996
      \endverb
    \endentry
    \entry{MR3604916}{article}{}
      \name{author}{1}{}{%
        {{hash=5d7cfdf14768e19ed402154ca9e7f128}{%
           family={Meng},
           familyi={M\bibinitperiod},
           given={X.},
           giveni={X\bibinitperiod}}}%
      }
      \strng{namehash}{5d7cfdf14768e19ed402154ca9e7f128}
      \strng{fullhash}{5d7cfdf14768e19ed402154ca9e7f128}
      \strng{bibnamehash}{5d7cfdf14768e19ed402154ca9e7f128}
      \strng{authorbibnamehash}{5d7cfdf14768e19ed402154ca9e7f128}
      \strng{authornamehash}{5d7cfdf14768e19ed402154ca9e7f128}
      \strng{authorfullhash}{5d7cfdf14768e19ed402154ca9e7f128}
      \field{extraname}{1}
      \field{sortinit}{M}
      \field{sortinithash}{2e5c2f51f7fa2d957f3206819bf86dc3}
      \field{labelprefix}{A}
      \field{labelnamesource}{author}
      \field{labeltitlesource}{title}
      \field{annotation}{Assuming RH, this article connects the normalized error term $E^{Q_k}(x) = (Q_k(x) - x/\zeta(k))/x^{1/2k}$ for the distribution of $k$-free numbers with the sum $J_{-1}(T) = \sum _{0<\gamma \leq T} |\zeta'(\rho)|^{-2}$ over nontrivial zeros of $\zeta(s)$. The author first shows that $J_{-1}(T) \ll_{\ep} T^{1+\ep}$ holding for all $\ep>0$ is equivalent to $\int^{X}_1 E^{Q_k}(x)^2 \,\frac{dx}{x} \ll_{k}\log X$ holding for all $k \geq 2$. If in fact $J_{-1}(T)\ll _{\ep}T^{1+\ep}$ for all $\ep>0$, the author proves that for each $k \geq 2$, \[ \lim_{X\to\infty} \frac1{\log X} \int^{X} _1 E^{Q_k}(x)^2 \frac{dx}{x} = \sum_{\gamma>0}\frac{2|\zeta(\rho/k)|^2}{|\rho \zeta'(\rho)|^2} \] and $E^{Q_k}(x)$ has a limiting logarithmic distribution. The author establishes analogous results for $M(x)$. Assuming (RH still and) the assumption $J_{-1}(T)\ll T^{2-\delta}$ for any fixed $\delta>0$, he establishes a weaker version of Mertens conjecture, namely $\int_2^X (M(x)/x)^2 \,dx \ll \log X$, and proves that $M(x) \ll x^{1/2}(\log \log x)^{3/2}$ except on a set of finite logarithmic measure. He also shows that \[ \lim_{X\to\infty} \frac1{\log X} \int^{X}_1 E^M(x)^2 \,\frac{dx}{x} = \sum_{\gamma>0}\frac{2}{|\rho \zeta'(\rho)|^2} \] and $E^M(x)$ has a limiting logarithmic distribution.\par This article cites \cites{MR3261965,MR0882550,MR2078705,MR3261965,MR1502996}.}
      \field{issn}{0305-0041}
      \field{journaltitle}{Math. Proc. Cambridge Philos. Soc.}
      \field{number}{2}
      \field{title}{The distribution of {$k$}-free numbers and the derivative of the {R}iemann zeta-function}
      \field{volume}{162}
      \field{year}{2017}
      \true{nocite}
      \field{pages}{293\bibrangedash 317}
      \range{pages}{25}
      \verb{doi}
      \verb 10.1017/S0305004116000554
      \endverb
      \verb{mrnumber}
      \verb 3604916
      \endverb
      \verb{urlraw}
      \verb https://doi.org/10.1017/S0305004116000554
      \endverb
      \verb{url}
      \verb https://doi.org/10.1017/S0305004116000554
      \endverb
    \endentry
    \entry{MR3803705}{article}{}
      \name{author}{1}{}{%
        {{hash=5d7cfdf14768e19ed402154ca9e7f128}{%
           family={Meng},
           familyi={M\bibinitperiod},
           given={X.},
           giveni={X\bibinitperiod}}}%
      }
      \strng{namehash}{5d7cfdf14768e19ed402154ca9e7f128}
      \strng{fullhash}{5d7cfdf14768e19ed402154ca9e7f128}
      \strng{bibnamehash}{5d7cfdf14768e19ed402154ca9e7f128}
      \strng{authorbibnamehash}{5d7cfdf14768e19ed402154ca9e7f128}
      \strng{authornamehash}{5d7cfdf14768e19ed402154ca9e7f128}
      \strng{authorfullhash}{5d7cfdf14768e19ed402154ca9e7f128}
      \field{extraname}{2}
      \field{sortinit}{M}
      \field{sortinithash}{2e5c2f51f7fa2d957f3206819bf86dc3}
      \field{labelprefix}{A}
      \field{labelnamesource}{author}
      \field{labeltitlesource}{title}
      \field{annotation}{Define $\pi_k(x;q,a) = \#\{n \leq x \colon n \equiv a \mod{q},\, \omega(n) = k \}$, and define $N_k(x;q,a)$ to be the \mbox{analogous} function with $\omega(n)$ replaced by $\Omega(n)$; define $\pi_k(x;q,a,b)=\pi_k(x;q,a)-\pi_k(x;q,b)$ and similarly for~$N_k(x;q,a,b)$. Assuming GRH and the simplicity of the zeros of $L(s,\chi)$, the author establishes asymptotic formulas relating $\pi_k(x;q,a,b)$ and $N_k(x;q,a,b)$ to $N_1(x;q,a,b) = \pi(x;q,a,b)$. Further assuming LI, the author establishes the existence of the densities $\delta(\pi_k(x;q,a,b))$ and $\delta(N_k(x;q,a,b))$ (including computations of several such), characterizes when these densities are less than or greater than $\frac12$, and demonstrates how the distances from such densities to $\frac12$ decrease as functions of~$k$. \par This article cites~\cites{1853.Chebyshev,MR3028758,MR1985941,MR2778652,MR2202918,MR552467,MR0146156,MR3063909,MR0068565,MR0083001,1914.Littlewood,MR1329368,MR0882550}.}
      \field{issn}{1937-0652}
      \field{journaltitle}{Algebra Number Theory}
      \field{number}{2}
      \field{title}{Chebyshev's bias for products of $k$ primes}
      \field{volume}{12}
      \field{year}{2018}
      \true{nocite}
      \field{pages}{305\bibrangedash 341}
      \range{pages}{37}
      \verb{doi}
      \verb 10.2140/ant.2018.12.305
      \endverb
      \verb{mrnumber}
      \verb 3803705
      \endverb
    \endentry
    \entry{MR3778223}{article}{}
      \name{author}{1}{}{%
        {{hash=5d7cfdf14768e19ed402154ca9e7f128}{%
           family={Meng},
           familyi={M\bibinitperiod},
           given={X.},
           giveni={X\bibinitperiod}}}%
      }
      \strng{namehash}{5d7cfdf14768e19ed402154ca9e7f128}
      \strng{fullhash}{5d7cfdf14768e19ed402154ca9e7f128}
      \strng{bibnamehash}{5d7cfdf14768e19ed402154ca9e7f128}
      \strng{authorbibnamehash}{5d7cfdf14768e19ed402154ca9e7f128}
      \strng{authornamehash}{5d7cfdf14768e19ed402154ca9e7f128}
      \strng{authorfullhash}{5d7cfdf14768e19ed402154ca9e7f128}
      \field{extraname}{3}
      \field{sortinit}{M}
      \field{sortinithash}{2e5c2f51f7fa2d957f3206819bf86dc3}
      \field{labelprefix}{A}
      \field{labelnamesource}{author}
      \field{labeltitlesource}{title}
      \field{annotation}{Given $k \geq 2$ and a multiset $\{a_1,\ldots,a_k\}$ of reduced residues modulo~$q$, the author determines an asymptotic formula for the number of squarefree integers up to~$x$ that can be written as $p_1p_2\cdots p_k$ with each $p_j \equiv a_j \mod q$, with an explicit second-order term of relative size $1/\log\log x$ compared to the main term. In particular, certain such multisets are biased over others, and those biases are eventually permanent. The constants in this second-order term are closely related to the constants in Mertens sum for primes in arithmetic progressions, which are heavily influenced by the least prime appearing. The author's results are given both for fixed $k$ and with $k$ growing like a constant multiple of $\log\log x$. \par This article cites~\cites{MR3521338,MR2778652,MR3803705,MR2034131}.}
      \field{issn}{0025-5793}
      \field{journaltitle}{Mathematika}
      \field{number}{1}
      \field{title}{Large bias for integers with prime factors in arithmetic progressions}
      \field{volume}{64}
      \field{year}{2018}
      \true{nocite}
      \field{pages}{237\bibrangedash 252}
      \range{pages}{16}
      \verb{mrnumber}
      \verb 3778223
      \endverb
    \endentry
    \entry{MR4077187}{article}{}
      \name{author}{1}{}{%
        {{hash=0e36a3a1f7c6c0678ce294e777ceb2f9}{%
           family={Meng},
           familyi={M\bibinitperiod},
           given={Xianchang},
           giveni={X\bibinitperiod}}}%
      }
      \strng{namehash}{0e36a3a1f7c6c0678ce294e777ceb2f9}
      \strng{fullhash}{0e36a3a1f7c6c0678ce294e777ceb2f9}
      \strng{bibnamehash}{0e36a3a1f7c6c0678ce294e777ceb2f9}
      \strng{authorbibnamehash}{0e36a3a1f7c6c0678ce294e777ceb2f9}
      \strng{authornamehash}{0e36a3a1f7c6c0678ce294e777ceb2f9}
      \strng{authorfullhash}{0e36a3a1f7c6c0678ce294e777ceb2f9}
      \field{extraname}{4}
      \field{sortinit}{M}
      \field{sortinithash}{2e5c2f51f7fa2d957f3206819bf86dc3}
      \field{labelprefix}{A}
      \field{labelnamesource}{author}
      \field{labeltitlesource}{title}
      \field{annotation}{Assuming GRH and LI, the author proves that both $\sum_{n\le x} \omega(n)\chi_{-4}(n) < 0$ and $\sum_{n\le x} \Omega(n)\chi_{-4}(n) > 0$ on sets of logarithmic density~$1$. This result gives a conditional confirmation of a conjecture of G.~Martin.\par This article cites~\cites{MR3261965,1853.Chebyshev, MR3604916, MR3803705, MR3778223, MR1329368}.}
      \field{issn}{0033-5606}
      \field{journaltitle}{Q. J. Math.}
      \field{number}{1}
      \field{title}{Number of prime factors over arithmetic progressions}
      \field{volume}{71}
      \field{year}{2020}
      \true{nocite}
      \field{pages}{97\bibrangedash 121}
      \range{pages}{25}
      \verb{doi}
      \verb 10.1093/qmathj/haz040
      \endverb
      \verb{mrnumber}
      \verb 4077187
      \endverb
      \verb{urlraw}
      \verb https://doi.org/10.1093/qmathj/haz040
      \endverb
      \verb{url}
      \verb https://doi.org/10.1093/qmathj/haz040
      \endverb
    \endentry
    \entry{1897.Mertens}{article}{}
      \name{author}{1}{}{%
        {{hash=a2b41763522272edb01791c27c2476d9}{%
           family={Mertens},
           familyi={M\bibinitperiod},
           given={F.},
           giveni={F\bibinitperiod}}}%
      }
      \list{language}{1}{%
        {German}%
      }
      \strng{namehash}{a2b41763522272edb01791c27c2476d9}
      \strng{fullhash}{a2b41763522272edb01791c27c2476d9}
      \strng{bibnamehash}{a2b41763522272edb01791c27c2476d9}
      \strng{authorbibnamehash}{a2b41763522272edb01791c27c2476d9}
      \strng{authornamehash}{a2b41763522272edb01791c27c2476d9}
      \strng{authorfullhash}{a2b41763522272edb01791c27c2476d9}
      \field{sortinit}{M}
      \field{sortinithash}{2e5c2f51f7fa2d957f3206819bf86dc3}
      \field{labelprefix}{A}
      \field{labelnamesource}{author}
      \field{labeltitlesource}{title}
      \field{annotation}{The author publishes a table of the values of $\mu(n)$ and $M(n)$ for $n\le10,000$, using the formula $\sum_{d=1}^{\lfloor\sqrt n\rfloor} \big( \mu(d) \lfloor \frac nd \rfloor + M(\frac nd) \big) = \lfloor\sqrt n\rfloor M(\lfloor\sqrt n\rfloor) + 1$ to check the output. Based on the table, the author conjectures that $|M(x)|< \sqrt{x}$ for all $x>1$. Using the formula $\psi(x)=-\sum_{n \leq x} \mu(n) \lfloor \frac xn \rfloor \log n$, the author shows that $M(x)\ll\sqrt x$ would imply $\Delta^\psi(x) \ll x^{3/4} \log x$ and $\Delta^\Pi(x) \ll x^{3/4}$ as well as RH, and that $M(x)\ll\sqrt x/(\log x)^{1+\delta}$ for some $\delta>0$ would imply that $\zeta(s)$ has no nontrivial zeros.}
      \field{journaltitle}{Sitzungsberichte Akad. Wien}
      \field{title}{Über eine zahlentheoretische Funktion}
      \field{volume}{106}
      \field{year}{1897}
      \true{nocite}
      \field{pages}{761\bibrangedash 830}
      \range{pages}{70}
    \endentry
    \entry{MR2631002}{book}{}
      \name{author}{1}{}{%
        {{hash=fbe9313e8e03d825339429599a21e678}{%
           family={Monach},
           familyi={M\bibinitperiod},
           given={William\bibnamedelima Reynolds},
           giveni={W\bibinitperiod\bibinitdelim R\bibinitperiod}}}%
      }
      \list{publisher}{1}{%
        {ProQuest LLC, Ann Arbor, MI}%
      }
      \strng{namehash}{fbe9313e8e03d825339429599a21e678}
      \strng{fullhash}{fbe9313e8e03d825339429599a21e678}
      \strng{bibnamehash}{fbe9313e8e03d825339429599a21e678}
      \strng{authorbibnamehash}{fbe9313e8e03d825339429599a21e678}
      \strng{authornamehash}{fbe9313e8e03d825339429599a21e678}
      \strng{authorfullhash}{fbe9313e8e03d825339429599a21e678}
      \field{sortinit}{M}
      \field{sortinithash}{2e5c2f51f7fa2d957f3206819bf86dc3}
      \field{labelprefix}{A}
      \field{labelnamesource}{author}
      \field{labeltitlesource}{title}
      \field{annotation}{In Chapter 2 of this thesis, assuming RH and LI, the author shows that the limiting logarithmic distribution of $E^\psi$ is equal to the distribution function of the random variable $\sum_{\gamma>0} 2\sin(2\pi\theta_\gamma)/|\rho|$, where the~$\theta_\gamma$ are independent random variables uniformly distributed in $[0,1)$. The remainder of the chapter describes methods of computing bounds on this distribution function on the interval $[-1,1]$ (a table of values is included, as are the programs used to generate said table) and determining an asymptotic formula for its derivative.\par This article cites~\cite{MR0104638,MR0137689,MR1501802}.}
      \field{note}{Thesis (Ph.D.)--University of Michigan)}
      \field{pagetotal}{180}
      \field{title}{Numerical Investigation of Several Problems in Number Theory}
      \field{year}{1980}
      \true{nocite}
      \verb{mrnumber}
      \verb 2631002
      \endverb
      \verb{urlraw}
      \verb http://gateway.proquest.com/openurl?url_ver=Z39.88-2004&rft_val_fmt=info:ofi/fmt:kev:mtx:dissertation&res_dat=xri:pqdiss&rft_dat=xri:pqdiss:8106192
      \endverb
      \verb{url}
      \verb http://gateway.proquest.com/openurl?url_ver=Z39.88-2004&rft_val_fmt=info:ofi/fmt:kev:mtx:dissertation&res_dat=xri:pqdiss&rft_dat=xri:pqdiss:8106192
      \endverb
    \endentry
    \entry{MR634679}{inproceedings}{}
      \name{author}{1}{}{%
        {{hash=98296eedbd8e7fdb593797362d5ac2c7}{%
           family={Montgomery},
           familyi={M\bibinitperiod},
           given={H.\bibnamedelimi L.},
           giveni={H\bibinitperiod\bibinitdelim L\bibinitperiod}}}%
      }
      \list{publisher}{1}{%
        {Queen's Univ., Kingston, Ont.}%
      }
      \strng{namehash}{98296eedbd8e7fdb593797362d5ac2c7}
      \strng{fullhash}{98296eedbd8e7fdb593797362d5ac2c7}
      \strng{bibnamehash}{98296eedbd8e7fdb593797362d5ac2c7}
      \strng{authorbibnamehash}{98296eedbd8e7fdb593797362d5ac2c7}
      \strng{authornamehash}{98296eedbd8e7fdb593797362d5ac2c7}
      \strng{authorfullhash}{98296eedbd8e7fdb593797362d5ac2c7}
      \field{extraname}{1}
      \field{sortinit}{M}
      \field{sortinithash}{2e5c2f51f7fa2d957f3206819bf86dc3}
      \field{labelprefix}{A}
      \field{labelnamesource}{author}
      \field{labeltitlesource}{title}
      \field{annotation}{Section~3 of this article examines random variables of the form $X=\sum_{k=1}^\infty r_k \sin(2\pi\theta_k)$ for $\{r_k\}$ a decreasing $\ell^2$ sequence, where the $\theta_k$ are independently uniformly distributed on~$\R/\Z$. The author establishes, for any integer $K\ge1$, the bounds \begin{align*} P\bigg(X\geq 2\sum_{k=1}^K r_k\bigg) &\leq \exp\left(-\frac 34 \bigg(\sum_{k=1}^K r_k\bigg)^2\bigg(\sum_{k=K+1}^\infty r_k^2\bigg)^{-1}\right) \\ P\bigg(X\geq \frac 12 \sum_{k=1}^K r_k\bigg) &\geq \frac 1{2^{40}}\exp\left(-100 \bigg(\sum_{k=1}^K r_k\bigg)^2\bigg(\sum_{k=K+1}^\infty r_k^2\bigg)^{-1}\right); \end{align*} in addition, if $\delta$ is sufficiently small and $\sum_{k\colon r_k>\delta}(r_k-\delta)\geq V$, then \[ P(X\geq V)\geq \frac 12 \exp\bigg(-\frac 12 \sum_{k\colon r_k>\delta} \log\frac{\pi^2 r_k}{2\delta}\bigg). \] These results can be applied to the limiting logarithmic distribution function of $E^\psi(x)$, which (assuming RH and LI) is the same as the distribution of the random variable $Y=\sum_{\gamma>0}\frac 2{|\rho|}\sin(2\pi\theta_\rho)$. In particular, the second result implies that there exist constants $0<c_1<c_2$ such that \[ \exp\big( -c_2\sqrt{v}e^{\sqrt{2\pi v}} \big) \le P(Y > v) \le \exp\big( -c_1\sqrt{v}e^{\sqrt{2\pi v}} \big), \] which suggests the conjecture $\displaystyle\limsup_{x\to\infty}\frac{E^\psi(x)}{(\log\log\log{x})^2}=\frac 1{2\pi}$ and $\displaystyle\liminf_{x\to\infty}\frac{E^\psi(x)}{(\log\log\log{x})^2}=-\frac 1{2\pi}$.}
      \field{booktitle}{Proceedings of the {Q}ueen's {N}umber {T}heory {C}onference, 1979 ({K}ingston, {O}nt., 1979)}
      \field{series}{Queen's Papers in Pure and Appl. Math.}
      \field{title}{The zeta function and prime numbers}
      \field{volume}{54}
      \field{year}{1980}
      \true{nocite}
      \field{pages}{1\bibrangedash 31}
      \range{pages}{31}
      \verb{mrnumber}
      \verb 634679
      \endverb
    \endentry
    \entry{MR1297543}{book}{}
      \name{author}{1}{}{%
        {{hash=d6fc16fe2b971a8009f6a09376b27581}{%
           family={Montgomery},
           familyi={M\bibinitperiod},
           given={Hugh\bibnamedelima L.},
           giveni={H\bibinitperiod\bibinitdelim L\bibinitperiod}}}%
      }
      \list{publisher}{1}{%
        {Published for the Conference Board of the Mathematical Sciences, Washington, DC; by the American Mathematical Society, Providence, RI}%
      }
      \strng{namehash}{d6fc16fe2b971a8009f6a09376b27581}
      \strng{fullhash}{d6fc16fe2b971a8009f6a09376b27581}
      \strng{bibnamehash}{d6fc16fe2b971a8009f6a09376b27581}
      \strng{authorbibnamehash}{d6fc16fe2b971a8009f6a09376b27581}
      \strng{authornamehash}{d6fc16fe2b971a8009f6a09376b27581}
      \strng{authorfullhash}{d6fc16fe2b971a8009f6a09376b27581}
      \field{extraname}{2}
      \field{sortinit}{M}
      \field{sortinithash}{2e5c2f51f7fa2d957f3206819bf86dc3}
      \field{labelprefix}{A}
      \field{labelnamesource}{author}
      \field{labeltitlesource}{title}
      \field{annotation}{In Chapter~$5$ the author gives proofs of the first and second main theorems of the power-sum method using duality, and improves the first main theorem for longer ranges of~$v$ to $\max_{1\leq v\leq N^2} |s_v|\gg {|s_0|}/{N}$. The author also deduces Fabry's gap theorem from the first main theorem, and gives improved lower bounds for the specific cases where all the coefficients~$b_j$ are all nonnegative or all equal to~$1$.\par This chapter cites~\cites{MR0133309,MR0133308,MR1575079,MR820241,MR749389}.}
      \field{isbn}{0-8218-0737-4}
      \field{series}{CBMS Regional Conference Series in Mathematics}
      \field{title}{Ten lectures on the interface between analytic number theory and harmonic analysis}
      \field{volume}{84}
      \field{year}{1994}
      \true{nocite}
      \field{pages}{xiv+220}
      \range{pages}{-1}
      \verb{doi}
      \verb 10.1090/cbms/084
      \endverb
      \verb{mrnumber}
      \verb 1297543
      \endverb
    \endentry
    \entry{MR2271616}{article}{}
      \name{author}{2}{}{%
        {{hash=d6fc16fe2b971a8009f6a09376b27581}{%
           family={Montgomery},
           familyi={M\bibinitperiod},
           given={Hugh\bibnamedelima L.},
           giveni={H\bibinitperiod\bibinitdelim L\bibinitperiod}}}%
        {{hash=855d3da694fbeeb0998eb3736ff8e5d1}{%
           family={Vorhauer},
           familyi={V\bibinitperiod},
           given={Ulrike\bibnamedelimb M.\bibnamedelimi A.},
           giveni={U\bibinitperiod\bibinitdelim M\bibinitperiod\bibinitdelim A\bibinitperiod}}}%
      }
      \strng{namehash}{a14ec9a880532023786b3f0b7cf9fee9}
      \strng{fullhash}{a14ec9a880532023786b3f0b7cf9fee9}
      \strng{bibnamehash}{a14ec9a880532023786b3f0b7cf9fee9}
      \strng{authorbibnamehash}{a14ec9a880532023786b3f0b7cf9fee9}
      \strng{authornamehash}{a14ec9a880532023786b3f0b7cf9fee9}
      \strng{authorfullhash}{a14ec9a880532023786b3f0b7cf9fee9}
      \field{sortinit}{M}
      \field{sortinithash}{2e5c2f51f7fa2d957f3206819bf86dc3}
      \field{labelprefix}{A}
      \field{labelnamesource}{author}
      \field{labeltitlesource}{title}
      \field{annotation}{The authors show that $\Delta^\psi(x)$ takes values of both signs in every interval $[x,19x]$ for $x\geq 1$ (where the~$19$ is best possible), and takes values of both signs in every interval $[x,2.02x]$ when $x$ is sufficiently large.\par This article cites~\cites{1936.Ingham, MR777013, MR1133236, 1914.Littlewood, MR1574052,1930.Polya}.}
      \field{issn}{0208-6573}
      \field{journaltitle}{Funct. Approx. Comment. Math.}
      \field{title}{Changes of sign of the error term in the prime number theorem}
      \field{volume}{35}
      \field{year}{2006}
      \true{nocite}
      \field{pages}{235\bibrangedash 247}
      \range{pages}{13}
      \verb{doi}
      \verb 10.7169/facm/1229442626
      \endverb
      \verb{mrnumber}
      \verb 2271616
      \endverb
    \endentry
    \entry{MR2034131}{article}{}
      \name{author}{1}{}{%
        {{hash=40a927b0e919f747a0218929b31aef12}{%
           family={Moree},
           familyi={M\bibinitperiod},
           given={Pieter},
           giveni={P\bibinitperiod}}}%
      }
      \strng{namehash}{40a927b0e919f747a0218929b31aef12}
      \strng{fullhash}{40a927b0e919f747a0218929b31aef12}
      \strng{bibnamehash}{40a927b0e919f747a0218929b31aef12}
      \strng{authorbibnamehash}{40a927b0e919f747a0218929b31aef12}
      \strng{authornamehash}{40a927b0e919f747a0218929b31aef12}
      \strng{authorfullhash}{40a927b0e919f747a0218929b31aef12}
      \field{sortinit}{M}
      \field{sortinithash}{2e5c2f51f7fa2d957f3206819bf86dc3}
      \field{labelprefix}{A}
      \field{labelnamesource}{author}
      \field{labeltitlesource}{title}
      \field{annotation}{Define $N(x;q,a)$ to be the number of integers up to $x$ all of whose prime factors are congruent to $a\mod q$. The author shows that $\min \{N(x;3,2), N(x;4,3)\} \geq \max\{N(x;3,1), N(x;4,1)\}$ for all $x\ge1$, using the fact that $N(x;q,a)$ is the summatory function of a multiplicative function whose values on primes alone has summatory function $\pi(x;q,a)$; consequently, biases in the distribution of small primes in residue classes modulo~$q$ can be magnified into complete biases among the $N(x;q,a)$. Indeed, Wirsing's method is already enough to establish such inequalities for sufficiently large~$x$; the author develops effective versions of Wirsing's method to conclude the same for all $x\ge1$. \par This article cites~\cites{MR0476616,1853.Chebyshev,MR1862054,MR1645811,MR0162771,1914.Littlewood,MR1756780,MR0137689,MR1329368}.}
      \field{issn}{0025-5718}
      \field{journaltitle}{Math. Comp.}
      \field{number}{245}
      \field{title}{{C}hebyshev's bias for composite numbers with restricted prime divisors}
      \field{volume}{73}
      \field{year}{2004}
      \true{nocite}
      \field{pages}{425\bibrangedash 449}
      \range{pages}{25}
      \verb{doi}
      \verb 10.1090/S0025-5718-03-01536-9
      \endverb
      \verb{mrnumber}
      \verb 2034131
      \endverb
    \endentry
    \entry{MR4194167}{article}{}
      \name{author}{3}{}{%
        {{hash=454d3dedc781ec4a49112de10313b411}{%
           family={Mossinghoff},
           familyi={M\bibinitperiod},
           given={Michael\bibnamedelima J.},
           giveni={M\bibinitperiod\bibinitdelim J\bibinitperiod}}}%
        {{hash=ca1fa3eb66152a77dcfaa9c5daf01ef9}{%
           family={Oliveira\bibnamedelimb e\bibnamedelima Silva},
           familyi={O\bibinitperiod\bibinitdelim e\bibinitperiod\bibinitdelim S\bibinitperiod},
           given={Tomás},
           giveni={T\bibinitperiod}}}%
        {{hash=66490333fa52c597ed64f7abf27c7710}{%
           family={Trudgian},
           familyi={T\bibinitperiod},
           given={Timothy\bibnamedelima S.},
           giveni={T\bibinitperiod\bibinitdelim S\bibinitperiod}}}%
      }
      \strng{namehash}{7efc60c309fdd97a9edaf0af2c43c3d9}
      \strng{fullhash}{7efc60c309fdd97a9edaf0af2c43c3d9}
      \strng{bibnamehash}{7efc60c309fdd97a9edaf0af2c43c3d9}
      \strng{authorbibnamehash}{7efc60c309fdd97a9edaf0af2c43c3d9}
      \strng{authornamehash}{7efc60c309fdd97a9edaf0af2c43c3d9}
      \strng{authorfullhash}{7efc60c309fdd97a9edaf0af2c43c3d9}
      \field{sortinit}{M}
      \field{sortinithash}{2e5c2f51f7fa2d957f3206819bf86dc3}
      \field{labelprefix}{A}
      \field{labelnamesource}{author}
      \field{labeltitlesource}{title}
      \field{annotation}{For each $k \geq 2$, set $$ C_k=\min\biggl\{ \limsup_{x \to \infty} \frac{\Delta^{Q_k}(x)}{x^{1/2k}}, \biggl| \liminf_{x \to \infty} \frac{\Delta^{Q_k}(x)}{x^{1/2k}} \biggr| \biggr\}. $$ The authors show that $C_k>3$ for $2 \leq k \leq 5$ and $C_k>2$ for $6\leq k \leq 10$, and that $C_k>0.74969$ for all sufficiently large $k$, improving previously known lower bounds on $C_k$ significantly. The proof is based on a variant of Ingham's method developed by Anderson and Stark~\cites{MR654520}. To achieve that, the authors establish weak linear dependence relations among a subset of ordinates of $\zeta$ via the LLL algorithm. The massive computation required approximately~$5$ core-years.\par This article cites~\cites{MR654520, MR3335903,MR0006202,MR429789,MR3604916,MR3062002,MR3700043,MR4132130,MR2078705,MR0199161,MR0882550,MR1502996,MR0724708}.}
      \field{issn}{0025-5718,1088-6842}
      \field{journaltitle}{Math. Comp.}
      \field{number}{328}
      \field{title}{The distribution of {$k$}-free numbers}
      \field{volume}{90}
      \field{year}{2021}
      \true{nocite}
      \field{pages}{907\bibrangedash 929}
      \range{pages}{23}
      \verb{doi}
      \verb 10.1090/mcom/3581
      \endverb
      \verb{mrnumber}
      \verb 4194167
      \endverb
      \verb{urlraw}
      \verb https://doi.org/10.1090/mcom/3581
      \endverb
      \verb{url}
      \verb https://doi.org/10.1090/mcom/3581
      \endverb
    \endentry
    \entry{MR3062002}{article}{}
      \name{author}{2}{}{%
        {{hash=454d3dedc781ec4a49112de10313b411}{%
           family={Mossinghoff},
           familyi={M\bibinitperiod},
           given={Michael\bibnamedelima J.},
           giveni={M\bibinitperiod\bibinitdelim J\bibinitperiod}}}%
        {{hash=66490333fa52c597ed64f7abf27c7710}{%
           family={Trudgian},
           familyi={T\bibinitperiod},
           given={Timothy\bibnamedelima S.},
           giveni={T\bibinitperiod\bibinitdelim S\bibinitperiod}}}%
      }
      \strng{namehash}{901cb3f7453290119b926d8dd538a562}
      \strng{fullhash}{901cb3f7453290119b926d8dd538a562}
      \strng{bibnamehash}{901cb3f7453290119b926d8dd538a562}
      \strng{authorbibnamehash}{901cb3f7453290119b926d8dd538a562}
      \strng{authornamehash}{901cb3f7453290119b926d8dd538a562}
      \strng{authorfullhash}{901cb3f7453290119b926d8dd538a562}
      \field{extraname}{1}
      \field{sortinit}{M}
      \field{sortinithash}{2e5c2f51f7fa2d957f3206819bf86dc3}
      \field{labelprefix}{A}
      \field{labelnamesource}{author}
      \field{labeltitlesource}{title}
      \field{annotation}{For $\alpha\in[0,1]$, define $L_\alpha(x) = \sum_{n\le x} (-1)^{\Omega(n)}/n^\alpha$. The authors generalize classical results of Pólya and Turán by proving that RH is equivalent to the estimate $L_\alpha(x) \ll_{\alpha,\ep} x^{1/2-\alpha+\ep}$ for all $\ep>0$, the equivalence being for any (hence all) $\alpha\in[0,1]$. Modifying the function slightly by defining \[ {\mathcal L}_\alpha(x) = L_\alpha(x) - \begin{cases} 0, &\text{if }0\le\alpha<\frac12, \\ \log x/2\zeta(\frac12), &\text{if }\alpha=\frac12, \\ \zeta(2\alpha)/\zeta(\alpha), &\text{if }\frac12<\alpha\le1, \end{cases} \] they also show that the assertion, for any $c\in\R$, that ${\mathcal L}_\alpha(x) - cx^{1/2-\alpha}$ has constant sign implies that all the zeros of $\zeta(s)$ are simple but is incompatible with~LI, even with finitely many exceptions. (The assertion that $L_{1/2}(x) - c$ itself has constant sign is shown to imply that all zeros of $\zeta(s)$ have multiplicity at most~$2$, but is not known to be inconsistent with~LI.) The authors state the problem of showing that ${\mathcal L}_\alpha(x)$ has infinitely many sign changes, which is known only for $\alpha=0$ and $\alpha=1$~\cite{MR0104638}, and also of showing that $L_{1/2}(x)\le0$ for~$x$ sufficiently large, perhaps $x\ge17$.\par This article cites~\cites{MR27305,MR0006202,MR0104638,MR1329368,MR2398787,MR0120198,MR0882550,MR584557}.}
      \field{issn}{1446-7887}
      \field{journaltitle}{J. Aust. Math. Soc.}
      \field{number}{1--2}
      \field{title}{Between the problems of {P}ólya and {T}urán}
      \field{volume}{93}
      \field{year}{2012}
      \true{nocite}
      \field{pages}{157\bibrangedash 171}
      \range{pages}{15}
      \verb{doi}
      \verb 10.1017/S1446788712000201
      \endverb
      \verb{mrnumber}
      \verb 3062002
      \endverb
    \endentry
    \entry{MR3700043}{incollection}{}
      \name{author}{2}{}{%
        {{hash=454d3dedc781ec4a49112de10313b411}{%
           family={Mossinghoff},
           familyi={M\bibinitperiod},
           given={Michael\bibnamedelima J.},
           giveni={M\bibinitperiod\bibinitdelim J\bibinitperiod}}}%
        {{hash=66490333fa52c597ed64f7abf27c7710}{%
           family={Trudgian},
           familyi={T\bibinitperiod},
           given={Timothy\bibnamedelima S.},
           giveni={T\bibinitperiod\bibinitdelim S\bibinitperiod}}}%
      }
      \list{publisher}{1}{%
        {Springer, Cham}%
      }
      \strng{namehash}{901cb3f7453290119b926d8dd538a562}
      \strng{fullhash}{901cb3f7453290119b926d8dd538a562}
      \strng{bibnamehash}{901cb3f7453290119b926d8dd538a562}
      \strng{authorbibnamehash}{901cb3f7453290119b926d8dd538a562}
      \strng{authornamehash}{901cb3f7453290119b926d8dd538a562}
      \strng{authorfullhash}{901cb3f7453290119b926d8dd538a562}
      \field{extraname}{2}
      \field{sortinit}{M}
      \field{sortinithash}{2e5c2f51f7fa2d957f3206819bf86dc3}
      \field{labelprefix}{A}
      \field{labelnamesource}{author}
      \field{labeltitlesource}{title}
      \field{annotation}{The authors show that \begin{align*} \liminf_{x \to \infty} E^L(x) < -2.3723 \quad &\text{and} \quad \limsup_{x \to \infty} E^L(x) > 1.0028 \\ \liminf_{x \to \infty} E^{L_r}(x) < -1.0028 \quad &\text{and} \quad \limsup_{x \to \infty} E^{L_r}(x) > 2.3723. \end{align*} They further define $L_\alpha(x)=\sum_{n \leq x} \lambda(n)n^{-\alpha}$ for all $\alpha \in (0,1)$, and show that $L_\alpha(x)$ changes signs infinitely often for each $\alpha \in (0, 0.29714\ldots)$, while $L_\alpha(x)-\zeta(2\alpha)/\zeta(\alpha)$ changes signs infinitely often for each $\alpha \in (0.70285,\ldots,1)$. The theoretical and computational reasoning employed is similar to~\cite{MR3335903}. \par This article cites~\cite{MR27305,MR584557,1905.Stieltjes,MR2398787,MR783538,MR0120198,MR429789,MR0006202,MR2994374,MR0104638,MR0332685,MR0202685,MR0332684,MR3335903,MR0330069,MR654520,MR3261965}.}
      \field{booktitle}{Exploring the {R}iemann zeta function}
      \field{title}{The {L}iouville function and the {R}iemann hypothesis}
      \field{year}{2017}
      \true{nocite}
      \field{pages}{201\bibrangedash 221}
      \range{pages}{21}
      \verb{mrnumber}
      \verb 3700043
      \endverb
    \endentry
    \entry{MR4132130}{incollection}{}
      \name{author}{2}{}{%
        {{hash=454d3dedc781ec4a49112de10313b411}{%
           family={Mossinghoff},
           familyi={M\bibinitperiod},
           given={Michael\bibnamedelima J.},
           giveni={M\bibinitperiod\bibinitdelim J\bibinitperiod}}}%
        {{hash=66490333fa52c597ed64f7abf27c7710}{%
           family={Trudgian},
           familyi={T\bibinitperiod},
           given={Timothy\bibnamedelima S.},
           giveni={T\bibinitperiod\bibinitdelim S\bibinitperiod}}}%
      }
      \list{publisher}{1}{%
        {Amer. Math. Soc., [Providence], RI}%
      }
      \strng{namehash}{901cb3f7453290119b926d8dd538a562}
      \strng{fullhash}{901cb3f7453290119b926d8dd538a562}
      \strng{bibnamehash}{901cb3f7453290119b926d8dd538a562}
      \strng{authorbibnamehash}{901cb3f7453290119b926d8dd538a562}
      \strng{authornamehash}{901cb3f7453290119b926d8dd538a562}
      \strng{authorfullhash}{901cb3f7453290119b926d8dd538a562}
      \field{extraname}{3}
      \field{sortinit}{M}
      \field{sortinithash}{2e5c2f51f7fa2d957f3206819bf86dc3}
      \field{labelprefix}{A}
      \field{labelnamesource}{author}
      \field{labeltitlesource}{title}
      \field{annotation}{Let $H(x)=\sum_{n\leq x}(-1)^{\omega(n)}$. The authors prove that $H(x)>1.7\sqrt{x}$ and $H(x)<-1.7\sqrt{x}$ infinitely often, complementing existing results on oscillations for $L(x)$. The proof is very similar to that of the follow-up paper~\cite{MR4194167}.\par This article cites~\cites{MR654520,MR0330069, MR3335903,MR2398787,MR0332684,MR0202685,MR0332685,MR0104638,MR2994374,MR3739227,MR0006202,MR429789,MR0120198,1897.Mertens,MR3062002,MR3700043,MR2078705,MR783538,1919.Polya,MR1329368,MR0882550}.}
      \field{booktitle}{75 years of mathematics of computation}
      \field{isbn}{978-1-4704-5163-9}
      \field{series}{Contemp. Math.}
      \field{title}{A tale of two omegas}
      \field{volume}{754}
      \field{year}{2020}
      \true{nocite}
      \field{pages}{343\bibrangedash 364}
      \range{pages}{22}
      \verb{mrnumber}
      \verb 4132130
      \endverb
    \endentry
    \entry{arXiv.Myerscough}{online}{}
      \name{author}{1}{}{%
        {{hash=3a7e707950739f23c95b5c3b3aee3641}{%
           family={Myerscough},
           familyi={M\bibinitperiod},
           given={C.},
           giveni={C\bibinitperiod}}}%
      }
      \strng{namehash}{3a7e707950739f23c95b5c3b3aee3641}
      \strng{fullhash}{3a7e707950739f23c95b5c3b3aee3641}
      \strng{bibnamehash}{3a7e707950739f23c95b5c3b3aee3641}
      \strng{authorbibnamehash}{3a7e707950739f23c95b5c3b3aee3641}
      \strng{authornamehash}{3a7e707950739f23c95b5c3b3aee3641}
      \strng{authorfullhash}{3a7e707950739f23c95b5c3b3aee3641}
      \field{sortinit}{M}
      \field{sortinithash}{2e5c2f51f7fa2d957f3206819bf86dc3}
      \field{labelprefix}{A}
      \field{labelnamesource}{author}
      \field{labeltitlesource}{title}
      \field{annotation}{The author studies the family $P_\mu(t)$ of density functions of the sum of random variables \begin{align*} \sum_\chi|\chi(a)-\chi(b)|\sum_{\gamma>\mu}\frac{2\Re Z_\gamma}{\sqrt{1/4+\gamma^2}} \end{align*} where $Z_\gamma$ are independent random variables, uniformly distributed on the unit circle, indexed by the ordinates of zeros of $L(s,\chi)$; the motivation is that $P_0(t)$ has the same distribution as the prime number race measure $\mu_{q;a,b}$. The author discovers that when $t$ is small, the behaviour of $P_\mu(t)$ approaches the normal distribution over a wider and wider range of $t/\sigma_\mu$ as $\mu$ increases, while when $t$ is large, $P_\mu(t)$ decays significantly faster than a normal distribution.\par To explicitly evaluate $P_\mu(t)$, the author compares the methods of the steepest descent (to third order), numerical convolution methods, and the Rubinstein--Sarnak method, and notes that there is an agreement to within $0.001\%$ for $1\leq t\leq 1.2$ (and similar observations for other ranges of~$t$). The author points out that Rubinstein--Sarnak is the best way to obtain results which are much larger than the absolute accuracy of computation, while the steepest descent method is desirable for extreme deviations, and the convolution method is valuable for some intermediate ranges.\par This article cites~\cites{MR1651741, MR1752093, MR1816036, MR1806291, MR3272280, MR3028758, 1932.Ingham, MR2943671, MR2037293, MR0202686, 1914.Littlewood, MR1956261, MR2631002, MR1329368}}
      \field{title}{Application of an accurate remainder term in the calculation of residue class distributions}
      \field{year}{2013}
      \true{nocite}
      \verb{urlraw}
      \verb https://arxiv.org/abs/1301.1434
      \endverb
      \verb{url}
      \verb https://arxiv.org/abs/1301.1434
      \endverb
    \endentry
    \entry{MR1756780}{book}{}
      \name{author}{1}{}{%
        {{hash=8c77965dc9a09815d7992a73cfd267ae}{%
           family={Narkiewicz},
           familyi={N\bibinitperiod},
           given={Władysław},
           giveni={W\bibinitperiod}}}%
      }
      \list{publisher}{1}{%
        {Springer-Verlag, Berlin}%
      }
      \strng{namehash}{8c77965dc9a09815d7992a73cfd267ae}
      \strng{fullhash}{8c77965dc9a09815d7992a73cfd267ae}
      \strng{bibnamehash}{8c77965dc9a09815d7992a73cfd267ae}
      \strng{authorbibnamehash}{8c77965dc9a09815d7992a73cfd267ae}
      \strng{authornamehash}{8c77965dc9a09815d7992a73cfd267ae}
      \strng{authorfullhash}{8c77965dc9a09815d7992a73cfd267ae}
      \field{sortinit}{N}
      \field{sortinithash}{98cf339a479c0454fe09153a08675a15}
      \field{labelprefix}{A}
      \field{labelnamesource}{author}
      \field{labeltitlesource}{title}
      \field{annotation}{Section~6.6 is devoted to the sign of the difference $\pi(x)-\li(x)$. The author sketches the proof of Littlewood's theorem and gives a detailed survey of results about $W(T)$ and $W^{\psi}(T)$.\par This article cites~\cites{1914.Littlewood, 1930.Polya, 1936.Ingham,MR1511206, MR1573970, MR0067145, MR0133309, MR0133308, MR0202686, MR0233782, MR0439771, MR0439772, MR0439774, MR607089, MR0447151, MR866118}.}
      \field{isbn}{3-540-66289-8}
      \field{series}{Springer Monographs in Mathematics}
      \field{title}{The development of prime number theory}
      \field{year}{2000}
      \true{nocite}
      \field{pages}{xii+448}
      \range{pages}{-1}
      \verb{doi}
      \verb 10.1007/978-3-662-13157-2
      \endverb
      \verb{mrnumber}
      \verb 1756780
      \endverb
      \verb{urlraw}
      \verb https://doi.org/10.1007/978-3-662-13157-2
      \endverb
      \verb{url}
      \verb https://doi.org/10.1007/978-3-662-13157-2
      \endverb
    \endentry
    \entry{MR155787}{article}{}
      \name{author}{1}{}{%
        {{hash=36809eb1c6a9cf47ae1c9468614c32b9}{%
           family={Neubauer},
           familyi={N\bibinitperiod},
           given={Gerhard},
           giveni={G\bibinitperiod}}}%
      }
      \list{language}{1}{%
        {German}%
      }
      \strng{namehash}{36809eb1c6a9cf47ae1c9468614c32b9}
      \strng{fullhash}{36809eb1c6a9cf47ae1c9468614c32b9}
      \strng{bibnamehash}{36809eb1c6a9cf47ae1c9468614c32b9}
      \strng{authorbibnamehash}{36809eb1c6a9cf47ae1c9468614c32b9}
      \strng{authornamehash}{36809eb1c6a9cf47ae1c9468614c32b9}
      \strng{authorfullhash}{36809eb1c6a9cf47ae1c9468614c32b9}
      \field{sortinit}{N}
      \field{sortinithash}{98cf339a479c0454fe09153a08675a15}
      \field{labelprefix}{A}
      \field{labelnamesource}{author}
      \field{labeltitlesource}{title}
      \field{annotation}{The author conducts an empirical study on the function $M(n)$, disproving two conjectures. The first conjecture, by von Sterneck~\cite{1913.vonSterneck} asserts that $|M(n)| \le \frac12\sqrt n$ for $n>200$, but the author finds that $n=7,760,000,000$ is a counterexample. The second conjecture, by Miller, asserts that $\sum_{n \leq x} M(n)<0$ for $x \geq 3$, but the author shows that it first fails at $x=21,067$. The author also points out some errors in the values of $M(n)$ listed in~\cite{1901.vonSterneck,1912.vonSterneck,1913.vonSterneck}. \par This article cites~\cite{1897.vonSterneck,1901.vonSterneck,1912.vonSterneck,1913.vonSterneck}.}
      \field{issn}{0029-599X}
      \field{journaltitle}{Numer. Math.}
      \field{title}{Eine empirische {U}ntersuchung zur {M}ertensschen {F}unktion}
      \field{volume}{5}
      \field{year}{1963}
      \true{nocite}
      \field{pages}{1\bibrangedash 13}
      \range{pages}{13}
      \verb{doi}
      \verb 10.1007/BF01385874
      \endverb
      \verb{mrnumber}
      \verb 155787
      \endverb
      \verb{urlraw}
      \verb https://doi-org.ezproxy.library.ubc.ca/10.1007/BF01385874
      \endverb
      \verb{url}
      \verb https://doi-org.ezproxy.library.ubc.ca/10.1007/BF01385874
      \endverb
    \endentry
    \entry{2000.Ng}{online}{}
      \name{author}{1}{}{%
        {{hash=585b691d1c27d2773434ed1ac9664075}{%
           family={Ng},
           familyi={N\bibinitperiod},
           given={N.},
           giveni={N\bibinitperiod}}}%
      }
      \strng{namehash}{585b691d1c27d2773434ed1ac9664075}
      \strng{fullhash}{585b691d1c27d2773434ed1ac9664075}
      \strng{bibnamehash}{585b691d1c27d2773434ed1ac9664075}
      \strng{authorbibnamehash}{585b691d1c27d2773434ed1ac9664075}
      \strng{authornamehash}{585b691d1c27d2773434ed1ac9664075}
      \strng{authorfullhash}{585b691d1c27d2773434ed1ac9664075}
      \field{extraname}{1}
      \field{sortinit}{N}
      \field{sortinithash}{98cf339a479c0454fe09153a08675a15}
      \field{labelprefix}{A}
      \field{labelnamesource}{author}
      \field{labeltitlesource}{title}
      \field{annotation}{In Section~2 of this thesis, the author gives examples of Artin $L$-functions for which LI is false. In Section~5.1, the author derives (under GRH) an explicit formula for $\frac{\log{x}}{\sqrt{x}}\left(\frac{|G|}{|C|}\pi_C(x)-\pi_K(x)\right)$ as $x\to\infty$, where $L/K$ is a normal extension of number fields, $G=\operatorname{Gal}(L/K)$ and $C$ is a conjugacy class of~$G$, and $\pi_K(x)=\#\{\mathfrak p\subset\mathcal{O}_K \colon \text{N}\mathfrak p\leq x\}$ and $\pi_C(x)=\#\{ \mathfrak p\subset\mathcal{O}_K \colon \text{N}\mathfrak p\leq x,\, \sigma_{\mathfrak p}=C\}$ where $\sigma_{\mathfrak p}$ is the Frobenius. Using the explicit formula, the author obtains analogues of Chebychev's bias in Galois groups. In Section~6.1, an explicit formula is given for $\frac{\log{x}}{\sqrt{x}}\pi(x,\mathfrak{a_1},\mathfrak{a_2})$, where $\mathfrak{a_1}$ and $\mathfrak{a_2}$ are ideal classes of a number field $K$. This formula is used to prove analogues of Chebychev's bias for primes in ideal classes. In Chapter 7, the author conditionally establishes the existence of a limiting distribution for $E^M(x)$ and conjectures that its maximal order is $(\log\log\log x)^{5/4}$.\par This article cites~ \cites{MR0330069, MR1816036, MR1544509, MR1806291, MR598886, MR0104638, MR1159354, MR0006202, MR634679, MR783538, MR1329368, MR0882550}}
      \field{note}{Thesis (Ph.D.)--University of British Columbia}
      \field{pagetotal}{161}
      \field{title}{{Limiting processes and Zeros of Artin L-Functions}}
      \field{year}{2000}
      \true{nocite}
      \verb{urlraw}
      \verb http://www.cs.uleth.ca/~nathanng/RESEARCH/phd.thesis.pdf
      \endverb
      \verb{url}
      \verb http://www.cs.uleth.ca/~nathanng/RESEARCH/phd.thesis.pdf
      \endverb
    \endentry
    \entry{MR2078705}{article}{}
      \name{author}{1}{}{%
        {{hash=585b691d1c27d2773434ed1ac9664075}{%
           family={Ng},
           familyi={N\bibinitperiod},
           given={N.},
           giveni={N\bibinitperiod}}}%
      }
      \strng{namehash}{585b691d1c27d2773434ed1ac9664075}
      \strng{fullhash}{585b691d1c27d2773434ed1ac9664075}
      \strng{bibnamehash}{585b691d1c27d2773434ed1ac9664075}
      \strng{authorbibnamehash}{585b691d1c27d2773434ed1ac9664075}
      \strng{authornamehash}{585b691d1c27d2773434ed1ac9664075}
      \strng{authorfullhash}{585b691d1c27d2773434ed1ac9664075}
      \field{extraname}{2}
      \field{sortinit}{N}
      \field{sortinithash}{98cf339a479c0454fe09153a08675a15}
      \field{labelprefix}{A}
      \field{labelnamesource}{author}
      \field{labeltitlesource}{title}
      \field{annotation}{This article contains several results on the Mertens sum $M(x)$, all conditional on both RH and the estimate $\sum_{0<\gamma\le T} 1/{|\zeta'(\rho)|^2}\ll T$. The author shows that $M(x)\ll x^{1/2}(\log{x})^{3/2}$, and furthermore that $M(x)\ll x^{1/2}(\log\log{x})^{3/2}$ except on a set of finite logarithmic measure. He also proves that \[ \int_1^X \bigg( \frac{M(x)}x \bigg)^2\,dx \sim \log X \cdot \sum_{\gamma>0}\frac 2{|\rho\zeta'(\rho)|^2}, \] which in particular implies the ``weak Mertens conjecture'' $\int_2^Y ({M(x)}/x)^2 \,dx\ll \log{Y}$. Under the additional assumption of LI, he shows that $E^M(x)$ has a limiting logarithmic distribution whose Fourier transform can be written down explicitly. Partly building on unpublished work by Gonek, the author conjectures that $E^M(x)=\Omega_{\pm} \big((\log \log \log x)^{5/4}\big).$ \par This article cites~\cites{MR1544509, 1932.Ingham, MR0006202, MR0104638, MR0330069, MR634679, MR783538, MR1159354, MR1329368, 2000.Ng}.}
      \field{issn}{0024-6115}
      \field{journaltitle}{Proc. London Math. Soc. (3)}
      \field{number}{2}
      \field{title}{The distribution of the summatory function of the {M}öbius function}
      \field{volume}{89}
      \field{year}{2004}
      \true{nocite}
      \field{pages}{361\bibrangedash 389}
      \range{pages}{29}
      \verb{doi}
      \verb 10.1112/S0024611504014741
      \endverb
      \verb{mrnumber}
      \verb 2078705
      \endverb
      \verb{urlraw}
      \verb https://doi.org/10.1112/S0024611504014741
      \endverb
      \verb{url}
      \verb https://doi.org/10.1112/S0024611504014741
      \endverb
    \endentry
    \entry{MR783538}{article}{}
      \name{author}{2}{}{%
        {{hash=5763e72244b318a02691f3fc37f706a2}{%
           family={Odlyzko},
           familyi={O\bibinitperiod},
           given={A.\bibnamedelimi M.},
           giveni={A\bibinitperiod\bibinitdelim M\bibinitperiod}}}%
        {{hash=f6fbcde367a0ed5df3da06b96b6d8eff}{%
           family={Riele},
           familyi={R\bibinitperiod},
           given={H.\bibnamedelimi J.\bibnamedelimi J.},
           giveni={H\bibinitperiod\bibinitdelim J\bibinitperiod\bibinitdelim J\bibinitperiod},
           prefix={te},
           prefixi={t\bibinitperiod}}}%
      }
      \strng{namehash}{d88121d28b2f42829b064ec49a086281}
      \strng{fullhash}{d88121d28b2f42829b064ec49a086281}
      \strng{bibnamehash}{d88121d28b2f42829b064ec49a086281}
      \strng{authorbibnamehash}{d88121d28b2f42829b064ec49a086281}
      \strng{authornamehash}{d88121d28b2f42829b064ec49a086281}
      \strng{authorfullhash}{d88121d28b2f42829b064ec49a086281}
      \field{sortinit}{O}
      \field{sortinithash}{ff8d4eeb5101e3cf3809959b3592d942}
      \field{labelprefix}{A}
      \field{labelnamesource}{author}
      \field{labeltitlesource}{title}
      \field{annotation}{The authors disprove the Mertens conjecture by showing that \[ \liminf_{x \to \infty} E^M(x) <-1.009 \quad \text{and} \quad \limsup_{x \to \infty} E^M(x) > 1.06. \] The method, based upon work of Ingham, is to find values of $h_K(y) = \sum_{\rho}k(\gamma) \frac{e^{i\gamma y}}{\rho \zeta'(\rho)}$ that are large in absolute value, using the kernel $k(t) = g(t/T)$ with $T = 2515.286\ldots$ the height of the $2000$th zero of~$\zeta(s)$, where \[ g(t) = \begin{cases} (1-|t|)\cos(\pi t) + \pi^{-1}\sin(\pi |t|), & |t| \leq 1 \\ 0 , & |t| \geq 1 . \end{cases} \] The key development is the algorithm due to Lenstra, Lenstra, and Lovász for finding short vectors in lattices, which reduces the computation time needed to find an appropriate inhomogeneous Diophantine approximation. The authors begin with a summary of past work on the conjecture and conclude with remarks on future work towards finding explicit counterexamples to Mertens conjecture, discussing limitations of their method.\par This article cites~\cites{MR0046485,MR654520,MR0330069,MR0104638,MR0006202,MR0352026,MR736725,1897.vonSterneck,1901.vonSterneck,1912.vonSterneck,1913.vonSterneck,MR0120198,MR429789,MR549977,MR280447,MR240062}.}
      \field{issn}{0075-4102}
      \field{journaltitle}{J. Reine Angew. Math.}
      \field{title}{Disproof of the Mertens conjecture}
      \field{volume}{357}
      \field{year}{1985}
      \true{nocite}
      \field{pages}{138\bibrangedash 160}
      \range{pages}{23}
      \verb{doi}
      \verb 10.1515/crll.1985.357.138
      \endverb
      \verb{mrnumber}
      \verb 783538
      \endverb
    \endentry
    \entry{MR3206041}{article}{}
      \name{author}{1}{}{%
        {{hash=9f2b575b642afbcf502733e2015df915}{%
           family={Petrushov},
           familyi={P\bibinitperiod},
           given={O.\bibnamedelimi A.},
           giveni={O\bibinitperiod\bibinitdelim A\bibinitperiod}}}%
      }
      \strng{namehash}{9f2b575b642afbcf502733e2015df915}
      \strng{fullhash}{9f2b575b642afbcf502733e2015df915}
      \strng{bibnamehash}{9f2b575b642afbcf502733e2015df915}
      \strng{authorbibnamehash}{9f2b575b642afbcf502733e2015df915}
      \strng{authornamehash}{9f2b575b642afbcf502733e2015df915}
      \strng{authorfullhash}{9f2b575b642afbcf502733e2015df915}
      \field{sortinit}{P}
      \field{sortinithash}{bb5b15f2db90f7aef79bb9e83defefcb}
      \field{labelprefix}{A}
      \field{labelnamesource}{author}
      \field{labeltitlesource}{title}
      \field{annotation}{The author establishes results on the limit supremum and infimum of a real continuous function $\nu(t)$ under assumptions on the behaviour of the Laplace transform $\int_0^\infty e^{-st} \,d\nu(t)$ near its singular point. As a corollary, the author obtains results on the limiting behaviour of the function $M_a(x) = \sum_{n < x} \mu(n) n^{-a}$ for $0\leq a \le \frac12$ in terms of functions related to $1/\zeta(s)$.\par This article cites~\cite{MR783538}.}
      \field{issn}{0001-4346}
      \field{journaltitle}{Math. Notes}
      \field{number}{5--6}
      \field{title}{Asymptotic estimates of functions based on the behavior of their {L}aplace transforms near singular points}
      \field{volume}{93}
      \field{year}{2013}
      \true{nocite}
      \field{pages}{906\bibrangedash 916}
      \range{pages}{11}
      \verb{doi}
      \verb 10.1134/S0001434613050283
      \endverb
      \verb{mrnumber}
      \verb 3206041
      \endverb
    \endentry
    \entry{1891.Phragmen}{article}{}
      \name{author}{1}{}{%
        {{hash=24a636b82c755af93afc8c09bfe2de5f}{%
           family={Phragmén},
           familyi={P\bibinitperiod},
           given={P.},
           giveni={P\bibinitperiod}}}%
      }
      \list{language}{1}{%
        {French}%
      }
      \strng{namehash}{24a636b82c755af93afc8c09bfe2de5f}
      \strng{fullhash}{24a636b82c755af93afc8c09bfe2de5f}
      \strng{bibnamehash}{24a636b82c755af93afc8c09bfe2de5f}
      \strng{authorbibnamehash}{24a636b82c755af93afc8c09bfe2de5f}
      \strng{authornamehash}{24a636b82c755af93afc8c09bfe2de5f}
      \strng{authorfullhash}{24a636b82c755af93afc8c09bfe2de5f}
      \field{sortinit}{P}
      \field{sortinithash}{bb5b15f2db90f7aef79bb9e83defefcb}
      \field{labelprefix}{A}
      \field{labelnamesource}{author}
      \field{labeltitlesource}{title}
      \field{annotation}{The author establishes a general proposition (similar to the eventual ``Landau's theorem'') capable of establishing that particular functions change sign for arbitrarily large arguments. From this proposition, the author shows that $\Pi(x) - (\li(x)-\log2)$ changes sign infinitely often, as do the differences $\psi(x) - (x-\log\frac{\pi}2)$ and $\Pi^*_r(x) - (\log\log x+C_0)$ and $\psi_r(x) - (\log x - C_0)$. Also, each of the differences $\Pi_r(x;4,1) - \big( \frac12\Li(x) - \frac12\log \frac{\log x}{\log 2} - \log 2 \big)$ and $\Pi_r(x;4,3) - \big( \frac12\Li(x) - \frac12\log \frac{\log x}{\log 2} \big)$ and $\Pi_r(x;4,1,3) + \log 2$ changes signs infinitely often. Finally, the author establishes the assertion of Chebyshev that~$1$ is a limit point of the function $E^\pi(x;4,3,1)$. \par This article cites~\cites{1853.Chebyshev}.}
      \field{journaltitle}{Öfversigt af Kongl. Vetenskaps--Akademiens Föhandlingar.}
      \field{title}{{Sur le logarithme intégral et la fonction $f(x)$ de Riemann}}
      \field{volume}{48}
      \field{year}{1891}
      \true{nocite}
      \field{pages}{599\bibrangedash 616}
      \range{pages}{18}
    \endentry
    \entry{1884.Piltz}{article}{}
      \name{author}{1}{}{%
        {{hash=d60e0007aa708cf28f344f23b440f29d}{%
           family={Piltz},
           familyi={P\bibinitperiod},
           given={A.},
           giveni={A\bibinitperiod}}}%
      }
      \list{language}{1}{%
        {German}%
      }
      \strng{namehash}{d60e0007aa708cf28f344f23b440f29d}
      \strng{fullhash}{d60e0007aa708cf28f344f23b440f29d}
      \strng{bibnamehash}{d60e0007aa708cf28f344f23b440f29d}
      \strng{authorbibnamehash}{d60e0007aa708cf28f344f23b440f29d}
      \strng{authornamehash}{d60e0007aa708cf28f344f23b440f29d}
      \strng{authorfullhash}{d60e0007aa708cf28f344f23b440f29d}
      \field{sortinit}{P}
      \field{sortinithash}{bb5b15f2db90f7aef79bb9e83defefcb}
      \field{labelprefix}{A}
      \field{labelnamesource}{author}
      \field{labeltitlesource}{title}
      \field{annotation}{Several authors cite this habilitation thesis as the first appearance of GRH, the generalized Riemann hypothesis for Dirichlet $L$-functions (``Riemann--Piltz conjecture''). On page 25 the author writes, ``{\it Wie in der Einleitung erwähnt wurde, hatte Riemann die Vermuthung, dass diese Verschwindungstellen was die Funktion $\zeta(s)$ betrifft, sich sämmtlich durch die Form $1/2+\alpha i$ wo $\alpha$ reell ist, bringen lassen. Dieser Satz gilt nicht nur für die Funktion $\zeta(s)$ sondern auch für die $\zeta(s,\nu)$.}''}
      \field{journaltitle}{Habilitationsschrift, Friedrich--Schiller--Universität Jena}
      \field{title}{Über die Häufigkeit der Primzahlen in arithmetischen Progressionen und über verwandte Gesetze}
      \field{year}{1884}
      \true{nocite}
    \endentry
    \entry{MR0439773}{article}{}
      \name{author}{1}{}{%
        {{hash=c6d965df16e600c701d4bcc2888ae34c}{%
           family={Pintz},
           familyi={P\bibinitperiod},
           given={J.},
           giveni={J\bibinitperiod}}}%
      }
      \list{language}{1}{%
        {German}%
      }
      \strng{namehash}{c6d965df16e600c701d4bcc2888ae34c}
      \strng{fullhash}{c6d965df16e600c701d4bcc2888ae34c}
      \strng{bibnamehash}{c6d965df16e600c701d4bcc2888ae34c}
      \strng{authorbibnamehash}{c6d965df16e600c701d4bcc2888ae34c}
      \strng{authornamehash}{c6d965df16e600c701d4bcc2888ae34c}
      \strng{authorfullhash}{c6d965df16e600c701d4bcc2888ae34c}
      \field{extraname}{1}
      \field{sortinit}{P}
      \field{sortinithash}{bb5b15f2db90f7aef79bb9e83defefcb}
      \field{labelprefix}{A}
      \field{labelnamesource}{author}
      \field{labeltitlesource}{title}
      \field{annotation}{The author shows that there exists $c>0$ such that $W(T) \gg (\log\log T)^{c}$ when $T$ is sufficiently large. Indeed, this is a special case of a more general result that establishes many large oscillations of $\Delta(x)$: let $D$ be sufficiently large and set $\mu=D/\log\log\log\log T$. Then there are at least $\exp\big( (\log\log\log T)^{1-\mu} \big)$ sign changes of $\Delta(x)$ up to $T$, with oscillations as large as \[ \Delta(x) > \bigg( \frac12-\frac{3\log D}{D} \bigg)\mu \cdot \frac{\sqrt x\log\log\log x}{\log x} \] (and the negative analogue), which, when $D$ is so large that $\mu\gg1$, provides oscillations as large as those established by Littlewood.\par This article cites~\cites{1936.Ingham, MR0439772, MR0067145, 1914.Littlewood}.}
      \field{journaltitle}{Monatsh. Math.}
      \field{number}{3}
      \field{title}{Bemerkungen zur {A}rbeit: ``{O}n the sign changes of {$(\pi (x)-{\rm li}\ x)$}. {II}'' ({M}onatsh. {M}ath. {\bf 82} (1976), no. 2, 163--175) von {S}. {K}napowski und {P}. {T}urán}
      \field{volume}{82}
      \field{year}{1976}
      \true{nocite}
      \field{pages}{199\bibrangedash 206}
      \range{pages}{8}
      \verb{doi}
      \verb 10.1007/BF01526326
      \endverb
      \verb{mrnumber}
      \verb 0439773
      \endverb
      \verb{urlraw}
      \verb https://doi.org/10.1007/BF01526326
      \endverb
      \verb{url}
      \verb https://doi.org/10.1007/BF01526326
      \endverb
    \endentry
    \entry{MR607089}{article}{}
      \name{author}{1}{}{%
        {{hash=c6d965df16e600c701d4bcc2888ae34c}{%
           family={Pintz},
           familyi={P\bibinitperiod},
           given={J.},
           giveni={J\bibinitperiod}}}%
      }
      \strng{namehash}{c6d965df16e600c701d4bcc2888ae34c}
      \strng{fullhash}{c6d965df16e600c701d4bcc2888ae34c}
      \strng{bibnamehash}{c6d965df16e600c701d4bcc2888ae34c}
      \strng{authorbibnamehash}{c6d965df16e600c701d4bcc2888ae34c}
      \strng{authornamehash}{c6d965df16e600c701d4bcc2888ae34c}
      \strng{authorfullhash}{c6d965df16e600c701d4bcc2888ae34c}
      \field{extraname}{2}
      \field{sortinit}{P}
      \field{sortinithash}{bb5b15f2db90f7aef79bb9e83defefcb}
      \field{labelprefix}{A}
      \field{labelnamesource}{author}
      \field{labeltitlesource}{title}
      \field{annotation}{This article establishes new results on the number of sign changes of $\pi(x)-\li(x)$. In particular it proves the effective result $ W(T)\gg \sqrt{\log T}/\log\log T $ and corresponding results for $W^\Pi(T)$, $W^\theta(T)$, and $W^\psi(T)$ with the same lower bound. Moreover, it establishes that there is necessarily a sign change in the interval $[T,T\exp(63\sqrt{\log T}\log\log T)]$ for $T$ large enough in each of these cases, although the lower bound on such $T$ is effectively computable only in the $\Pi(x)$ and $\psi(x)$ versions.\par Under RH, Ingham's result from 1936 gives in fact a stronger localization theorem. The author here uses power-sum method, and in particular a result of Sós--Turán, to achieve a result under the assumption that RH fails. Ingham's idea to use Fejér kernels is also applied to prove the effective lower bound on $W(T)$ in the absence of an effective localization result.\par This article cites~\cites{1914.Littlewood, 1930.Polya, 1936.Ingham, MR0067145, MR0133309, MR0133308, MR0439771, MR0439772, MR0439773}.}
      \field{issn}{0081-6906}
      \field{journaltitle}{Studia Sci. Math. Hungar.}
      \field{number}{3-4}
      \field{title}{On the remainder term of the prime number formula. {III}. {S}ign changes of {$\pi (x)-{\rm li}x$}}
      \field{volume}{12}
      \field{year}{1977}
      \true{nocite}
      \field{pages}{345\bibrangedash 369 (1980)}
      \range{pages}{-1}
      \verb{mrnumber}
      \verb 607089
      \endverb
    \endentry
    \entry{MR0447151}{incollection}{}
      \name{author}{1}{}{%
        {{hash=c6d965df16e600c701d4bcc2888ae34c}{%
           family={Pintz},
           familyi={P\bibinitperiod},
           given={J.},
           giveni={J\bibinitperiod}}}%
      }
      \list{publisher}{1}{%
        {Soc. Math. France, Paris}%
      }
      \strng{namehash}{c6d965df16e600c701d4bcc2888ae34c}
      \strng{fullhash}{c6d965df16e600c701d4bcc2888ae34c}
      \strng{bibnamehash}{c6d965df16e600c701d4bcc2888ae34c}
      \strng{authorbibnamehash}{c6d965df16e600c701d4bcc2888ae34c}
      \strng{authornamehash}{c6d965df16e600c701d4bcc2888ae34c}
      \strng{authorfullhash}{c6d965df16e600c701d4bcc2888ae34c}
      \field{extraname}{3}
      \field{sortinit}{P}
      \field{sortinithash}{bb5b15f2db90f7aef79bb9e83defefcb}
      \field{labelprefix}{A}
      \field{labelnamesource}{author}
      \field{labeltitlesource}{title}
      \field{annotation}{The author begins with a thorough summary of results on sign changes of $\Delta^\pi(x)$ and related problems; he then announces, without proof, several new results of this type. He claims that for $T$ sufficiently large, unconditionally, $W(T) \gg \sqrt{\log T}/\log\log T$, and that there exists $c>0$ such that every interval of the form $[T^c,T]$ contains a sign change of $\Delta^\pi(x)$; ineffectively one can narrow these intervals to the form $[Te^{-\sqrt{\log T}\log\log T},T]$. Even if one restricts to ``big sign changes'', where $\Delta^\pi(x) = \Omega_\pm(\sqrt x\log\log\log x / \log x)$, the author asserts that the number of such sign changes up to $T$ is $\gg \sqrt{\log T} e^{-\sqrt{\log\log T}}$ effectively and $\gg \sqrt{\log T}/(\log\log T)^2$ ineffectively; these sign changes can be localized as well, and the latter inequality even holds for large sign changes of the average of $\Delta^\pi(x)$ over intervals of length $x/\log\log x$. The author further asserts that analogous theorems can be proved for the other prime counting functions, as well as for $\pi(x;4,1,3)$ and some other class of prime races.\par This article cites~\cites{1914.Littlewood, 1936.Ingham, MR0067145, MR0146156, MR0146157, MR0146158, MR0146159, MR0146160, MR0146161, MR0156826, MR0156827, MR0439771, MR0439772, MR0439773}.}
      \field{booktitle}{Journées {A}rithmétiques de {C}aen ({U}niv. {C}aen, {C}aen, 1976)}
      \field{title}{On the sign changes of {$\pi (x)-{\rm li}(x)$}}
      \field{year}{1977}
      \true{nocite}
      \field{pages}{255--265. Astérisque No. 41\bibrangedash 42}
      \range{pages}{-1}
      \verb{mrnumber}
      \verb 0447151
      \endverb
    \endentry
    \entry{MR630377}{article}{}
      \name{author}{1}{}{%
        {{hash=c6d965df16e600c701d4bcc2888ae34c}{%
           family={Pintz},
           familyi={P\bibinitperiod},
           given={J.},
           giveni={J\bibinitperiod}}}%
      }
      \strng{namehash}{c6d965df16e600c701d4bcc2888ae34c}
      \strng{fullhash}{c6d965df16e600c701d4bcc2888ae34c}
      \strng{bibnamehash}{c6d965df16e600c701d4bcc2888ae34c}
      \strng{authorbibnamehash}{c6d965df16e600c701d4bcc2888ae34c}
      \strng{authornamehash}{c6d965df16e600c701d4bcc2888ae34c}
      \strng{authorfullhash}{c6d965df16e600c701d4bcc2888ae34c}
      \field{extraname}{4}
      \field{sortinit}{P}
      \field{sortinithash}{bb5b15f2db90f7aef79bb9e83defefcb}
      \field{labelprefix}{A}
      \field{labelnamesource}{author}
      \field{labeltitlesource}{title}
      \field{annotation}{This article establishes lower bounds for the number of sign changes for the error terms of classical prime-counting functions. The main theorem of this article states that for $f=\pi$, $\Pi$, $\theta$, or $\psi$, there exists an absolute constant $Y(f)$ such that for $Y>Y(f)$, \[ W^f(Y)>\frac 1{10^{11}}\frac{\log{Y}}{(\log\log{Y})^3}. \] Interestingly, $Y(\Pi)$ and $Y(\psi)$ are effectively computable in the author's proof, whereas $Y(\pi)$ and $Y(\theta)$ are ineffective constants. \par This article cites~\cites{1914.Littlewood, 1936.Ingham, MR0067145, MR0133309, MR0133308, MR0439771, MR0439772, MR0439773, MR607089}.}
      \field{issn}{0081-6906}
      \field{journaltitle}{Studia Sci. Math. Hungar.}
      \field{number}{1-2}
      \field{title}{On the remainder term of the prime number formula. {IV}. {S}ign changes of {$\pi (x)-{\li}x$}}
      \field{volume}{13}
      \field{year}{1978}
      \true{nocite}
      \field{pages}{29\bibrangedash 42 (1981)}
      \range{pages}{-1}
      \verb{mrnumber}
      \verb 630377
      \endverb
    \endentry
    \entry{MR585891}{article}{}
      \name{author}{1}{}{%
        {{hash=c6d965df16e600c701d4bcc2888ae34c}{%
           family={Pintz},
           familyi={P\bibinitperiod},
           given={J.},
           giveni={J\bibinitperiod}}}%
      }
      \strng{namehash}{c6d965df16e600c701d4bcc2888ae34c}
      \strng{fullhash}{c6d965df16e600c701d4bcc2888ae34c}
      \strng{bibnamehash}{c6d965df16e600c701d4bcc2888ae34c}
      \strng{authorbibnamehash}{c6d965df16e600c701d4bcc2888ae34c}
      \strng{authornamehash}{c6d965df16e600c701d4bcc2888ae34c}
      \strng{authorfullhash}{c6d965df16e600c701d4bcc2888ae34c}
      \field{extraname}{5}
      \field{sortinit}{P}
      \field{sortinithash}{bb5b15f2db90f7aef79bb9e83defefcb}
      \field{labelprefix}{A}
      \field{labelnamesource}{author}
      \field{labeltitlesource}{title}
      \field{annotation}{This article contains explicit oscillation results under the assumption that RH is false. Suppose that $\rho_0 = \beta_0+i\gamma_0$ is a nontrivial zero of~$\zeta(s)$. Let $0<\ep\le0.02$, and set $A=40000 \ep^{-2}\log\gamma_0$. Then for $H$ sufficiently large in terms of $\rho_0$, there exist $x_+$ and $x_-$ in the interval $[H,H^A]$ such that \[ \Delta^\pi(x_+) > (1-\ep)\frac{x_+^{\beta_0}}{|\rho_0|\log x_+} \quad\text{and}\quad \Delta^\pi(x_-) < -(1-\ep)\frac{x_-^{\beta_0}}{|\rho_0|\log x_-}, \] and the same for $\Delta^\Pi(x)$; similarly, the result holds without the factor $\log x$ in the denominator for $\Delta^\psi(x)$ and~$\Delta^\theta(x)$. In all these theorems, if in addition $\beta_0>\frac12+\ep$ and $\gamma_0$ is sufficiently large in terms of $\ep$, then by replacing the factor $(1-\ep)/|\rho_0|$ by the smaller $1/\gamma_0^{1+\ep}$, the localization can be improved to the interval $[H,H^{1+\ep}]$. Consequently, $W^f(T)/\log\log T$ tends to infinity for each of the four functions $f\in\{\pi,\Pi,\theta,\psi\}$ (the case where RH is true having been handled by Ingham~\cite{1936.Ingham}).\par This article cites~\cites{MR0146156, MR0146157, MR0146158, MR0146159, MR0146160, MR0146161, MR0156826, MR0156827, MR0162771, MR0174538, MR0180539, MR0182616A, MR0182616B, 1914.Littlewood, 1936.Ingham, MR0439771, MR634679}.}
      \field{issn}{0065-1036}
      \field{journaltitle}{Acta Arith.}
      \field{number}{4}
      \field{title}{On the remainder term of the prime number formula. {I}. {O}n a problem of {L}ittlewood}
      \field{volume}{36}
      \field{year}{1980}
      \true{nocite}
      \field{pages}{341\bibrangedash 365}
      \range{pages}{25}
      \verb{doi}
      \verb 10.4064/aa-36-4-341-365
      \endverb
      \verb{mrnumber}
      \verb 585891
      \endverb
      \verb{urlraw}
      \verb https://doi.org/10.4064/aa-36-4-341-365
      \endverb
      \verb{url}
      \verb https://doi.org/10.4064/aa-36-4-341-365
      \endverb
    \endentry
    \entry{MR598876}{article}{}
      \name{author}{1}{}{%
        {{hash=c6d965df16e600c701d4bcc2888ae34c}{%
           family={Pintz},
           familyi={P\bibinitperiod},
           given={J.},
           giveni={J\bibinitperiod}}}%
      }
      \strng{namehash}{c6d965df16e600c701d4bcc2888ae34c}
      \strng{fullhash}{c6d965df16e600c701d4bcc2888ae34c}
      \strng{bibnamehash}{c6d965df16e600c701d4bcc2888ae34c}
      \strng{authorbibnamehash}{c6d965df16e600c701d4bcc2888ae34c}
      \strng{authornamehash}{c6d965df16e600c701d4bcc2888ae34c}
      \strng{authorfullhash}{c6d965df16e600c701d4bcc2888ae34c}
      \field{extraname}{6}
      \field{sortinit}{P}
      \field{sortinithash}{bb5b15f2db90f7aef79bb9e83defefcb}
      \field{labelprefix}{A}
      \field{labelnamesource}{author}
      \field{labeltitlesource}{title}
      \field{annotation}{This article investigates the connection between the zero free region of $\zeta(s)$ and the size of the remainder term in the prime number theorem. Let $\eta\colon [1,\infty)\to(0,\frac12]$ be a continuous, decreasing function, and suppose that $\zeta(s)$ does not vanish when $\sigma>1-\eta(|t|)$. If we define $\omega(x) = \min_{t\ge1} (\eta(t)\log x + \log t)$, then for any $0<\ep<1$, \[ \Delta^\psi(x) \ll_\ep x/e^{(1-\ep)\omega(x)}, \] and the same is true for $\Delta^\theta(x)$ and $\Delta^\Pi(x)$ and~$\Delta^\pi(x)$. This is an improvement of a result of Ingham~\cite[Theorem 22]{1932.Ingham}, which had a factor of $\frac12$ in the exponent (and additional conditions upon~$\eta$). In particular, when combined with a 1960/61 theorem of Staś, this result provides a nearly lossless relationship between zero-free regions for $\zeta(s)$ and error terms in the prime number theorem. It follows that an $\Omega$-theorem for any of the four error terms given above actually implies $\Omega_\pm$-theorems, of the same order of magnitude up to an $\ep$ in the exponent, for all four error terms, an implication that seems extremely difficult to prove directly.\par This article cites~\cites{MR0146156, MR0146157, MR0146158, MR0146159, MR0146160, MR0146161, MR0156826, MR0156827, MR0162771, MR0174538, MR0180539, MR0182616A, MR0182616B, 1914.Littlewood, 1932.Ingham, MR0439771, MR634679}.}
      \field{issn}{0065-1036}
      \field{journaltitle}{Acta Arith.}
      \field{title}{On the remainder term of the prime number formula. {II}. {O}n a theorem of {I}ngham}
      \field{volume}{37}
      \field{year}{1980}
      \true{nocite}
      \field{pages}{209\bibrangedash 220}
      \range{pages}{12}
      \verb{doi}
      \verb 10.4064/aa-37-1-209-220
      \endverb
      \verb{mrnumber}
      \verb 598876
      \endverb
      \verb{urlraw}
      \verb https://doi.org/10.4064/aa-37-1-209-220
      \endverb
      \verb{url}
      \verb https://doi.org/10.4064/aa-37-1-209-220
      \endverb
    \endentry
    \entry{MR681441}{article}{}
      \name{author}{1}{}{%
        {{hash=c6d965df16e600c701d4bcc2888ae34c}{%
           family={Pintz},
           familyi={P\bibinitperiod},
           given={J.},
           giveni={J\bibinitperiod}}}%
      }
      \strng{namehash}{c6d965df16e600c701d4bcc2888ae34c}
      \strng{fullhash}{c6d965df16e600c701d4bcc2888ae34c}
      \strng{bibnamehash}{c6d965df16e600c701d4bcc2888ae34c}
      \strng{authorbibnamehash}{c6d965df16e600c701d4bcc2888ae34c}
      \strng{authornamehash}{c6d965df16e600c701d4bcc2888ae34c}
      \strng{authorfullhash}{c6d965df16e600c701d4bcc2888ae34c}
      \field{extraname}{7}
      \field{sortinit}{P}
      \field{sortinithash}{bb5b15f2db90f7aef79bb9e83defefcb}
      \field{labelprefix}{A}
      \field{labelnamesource}{author}
      \field{labeltitlesource}{title}
      \field{annotation}{For any of the functions $f\in\{\pi,\Pi,\theta,\psi\}$, the author establishes lower bounds for the integrated absolute error term $\fA_{|1|}^f(x)$. The main theorem of this article states that if $\beta_0+i\gamma_0$ is a zero of the Riemann zeta function, then $\fA_{|1|}^f(Y)/Y \ge Y^{\beta_0}e^{-2\sqrt{\log{Y}}(\log\log{Y})^2}$ when $Y$ is sufficiently large in terms of~$\gamma_0$. The author sketches a modification of the proof that yields the stronger lower bound $\fA_{|1|}^f(Y)/Y \ge Y^{\beta_0}e^{-18(\log{Y})^{1/3}(\log\log{Y})^{4/3}}$.\par This article cites~\cites{MR0125822, MR0043121, MR585891, MR598876, MR607089, MR630377, MR681442,MR1544509,MR1575079,MR0111722,MR0417077}.}
      \field{issn}{0081-6906}
      \field{journaltitle}{Studia Sci. Math. Hungar.}
      \field{number}{1-3}
      \field{title}{On the remainder term of the prime number formula. {V}. {E}ffective mean value theorems}
      \field{volume}{15}
      \field{year}{1980}
      \true{nocite}
      \field{pages}{215\bibrangedash 223}
      \range{pages}{9}
      \verb{mrnumber}
      \verb 681441
      \endverb
    \endentry
    \entry{MR681442}{article}{}
      \name{author}{1}{}{%
        {{hash=c6d965df16e600c701d4bcc2888ae34c}{%
           family={Pintz},
           familyi={P\bibinitperiod},
           given={J.},
           giveni={J\bibinitperiod}}}%
      }
      \strng{namehash}{c6d965df16e600c701d4bcc2888ae34c}
      \strng{fullhash}{c6d965df16e600c701d4bcc2888ae34c}
      \strng{bibnamehash}{c6d965df16e600c701d4bcc2888ae34c}
      \strng{authorbibnamehash}{c6d965df16e600c701d4bcc2888ae34c}
      \strng{authornamehash}{c6d965df16e600c701d4bcc2888ae34c}
      \strng{authorfullhash}{c6d965df16e600c701d4bcc2888ae34c}
      \field{extraname}{8}
      \field{sortinit}{P}
      \field{sortinithash}{bb5b15f2db90f7aef79bb9e83defefcb}
      \field{labelprefix}{A}
      \field{labelnamesource}{author}
      \field{labeltitlesource}{title}
      \field{annotation}{This article concerns the absolute averages $\fA_{|1|}$ of various standard error terms for prime counting functions. When $Y$ is sufficiently large (ineffectively), the author proves that \begin{align*} \fA_{|1|}^\pi(Y) &>0.62 \frac{Y^{3/2}}{\log Y}, & \fA_{|1|}^\Pi(Y) &>9 \cdot 10^{-5} \frac{Y^{3/2}}{\log Y}, \\ \fA_{|1|}^\theta(Y) &>0.62{Y^{3/2}}, & \fA_{|1|}^\psi(Y) &> 10^{-4} {Y^{3/2}}. \end{align*} Thanks to work of Cramér~\cite{MR1544509}, if RH is true then these bounds are best possible up to the leading constants (and even those constants are not too far off). Under RH, the author can improve some of these constants and also better localize the implied large values of the error terms; indeed, the lower bounds for the $\fA_{|1|}$ are derived from the existence of large oscillations of the error terms, rather than the other way around. \par This article cites~\cites{MR1544509,1936.Ingham,MR681441}.}
      \field{issn}{0081-6906}
      \field{journaltitle}{Studia Sci. Math. Hungar.}
      \field{number}{1-3}
      \field{title}{On the remainder term of the prime number formula. {VI}. {I}neffective mean value theorems}
      \field{volume}{15}
      \field{year}{1980}
      \true{nocite}
      \field{pages}{225\bibrangedash 230}
      \range{pages}{6}
      \verb{mrnumber}
      \verb 681442
      \endverb
    \endentry
    \entry{MR688630}{article}{}
      \name{author}{1}{}{%
        {{hash=c6d965df16e600c701d4bcc2888ae34c}{%
           family={Pintz},
           familyi={P\bibinitperiod},
           given={J.},
           giveni={J\bibinitperiod}}}%
      }
      \strng{namehash}{c6d965df16e600c701d4bcc2888ae34c}
      \strng{fullhash}{c6d965df16e600c701d4bcc2888ae34c}
      \strng{bibnamehash}{c6d965df16e600c701d4bcc2888ae34c}
      \strng{authorbibnamehash}{c6d965df16e600c701d4bcc2888ae34c}
      \strng{authornamehash}{c6d965df16e600c701d4bcc2888ae34c}
      \strng{authorfullhash}{c6d965df16e600c701d4bcc2888ae34c}
      \field{extraname}{9}
      \field{sortinit}{P}
      \field{sortinithash}{bb5b15f2db90f7aef79bb9e83defefcb}
      \field{labelprefix}{A}
      \field{labelnamesource}{author}
      \field{labeltitlesource}{title}
      \field{annotation}{The author shows that $M(x)$ changes sign in the interval $[Y \exp(-3 \log_2^{3/2} Y),Y]$ for all sufficiently large~$Y$, which improves the main result in~\cite{MR660714}. The proof uses upper and lower bounds for an integral of the form $\int_v^w |M(x)|x^{-1-\theta} \,dx$, where $\theta$ is the maximal real part of the zeros of $\zeta(s)$ for $|t| \lesssim \log Y$.\par This article cites~\cites{MR0207665,MR0219496,MR0155802,1909.Landau,MR630377,MR660714,MR678996,MR0087685}.}
      \field{issn}{0081-6906}
      \field{journaltitle}{Studia Sci. Math. Hungar.}
      \field{number}{4}
      \field{title}{Oscillatory properties of {$M(x)=\sum _{n\leq x}\mu (n)$}. {II}}
      \field{volume}{15}
      \field{year}{1980}
      \true{nocite}
      \field{pages}{491\bibrangedash 496}
      \range{pages}{6}
      \verb{mrnumber}
      \verb 688630
      \endverb
    \endentry
    \entry{MR660714}{article}{}
      \name{author}{1}{}{%
        {{hash=c6d965df16e600c701d4bcc2888ae34c}{%
           family={Pintz},
           familyi={P\bibinitperiod},
           given={J.},
           giveni={J\bibinitperiod}}}%
      }
      \strng{namehash}{c6d965df16e600c701d4bcc2888ae34c}
      \strng{fullhash}{c6d965df16e600c701d4bcc2888ae34c}
      \strng{bibnamehash}{c6d965df16e600c701d4bcc2888ae34c}
      \strng{authorbibnamehash}{c6d965df16e600c701d4bcc2888ae34c}
      \strng{authornamehash}{c6d965df16e600c701d4bcc2888ae34c}
      \strng{authorfullhash}{c6d965df16e600c701d4bcc2888ae34c}
      \field{extraname}{10}
      \field{sortinit}{P}
      \field{sortinithash}{bb5b15f2db90f7aef79bb9e83defefcb}
      \field{labelprefix}{A}
      \field{labelnamesource}{author}
      \field{labeltitlesource}{title}
      \field{annotation}{Assuming RH$(T)$, the author shows that $M(x)$ changes sign in the interval $[Y^{1-1/(T-2)}, Y^{1+1/(T-2)}]$ for every sufficiently large~$Y$. Consequently, using Brent's verification of $RH$ for the first $7 \cdot 10^7$ zeros of $\zeta(s)$, the author concludes that $M(x)$ changes sign in every interval of the form $[Y^{1-10^{-7}},Y]$ for all sufficiently large~$Y$, which improves the main result in ~\cite{MR0219496}.\par This article cites~\cites{MR0207665,MR0219496,MR0155802,MR0172856}.}
      \field{issn}{0174-4747}
      \field{journaltitle}{Analysis}
      \field{number}{3}
      \field{title}{On the sign changes of {$M(x)=\sum _{n\leq x}\mu (n)$}}
      \field{volume}{1}
      \field{year}{1981}
      \true{nocite}
      \field{pages}{191\bibrangedash 195}
      \range{pages}{5}
      \verb{doi}
      \verb 10.1524/anly.1981.1.3.191
      \endverb
      \verb{mrnumber}
      \verb 660714
      \endverb
      \verb{urlraw}
      \verb https://doi.org/10.1524/anly.1981.1.3.191
      \endverb
      \verb{url}
      \verb https://doi.org/10.1524/anly.1981.1.3.191
      \endverb
    \endentry
    \entry{MR678996}{article}{}
      \name{author}{1}{}{%
        {{hash=c6d965df16e600c701d4bcc2888ae34c}{%
           family={Pintz},
           familyi={P\bibinitperiod},
           given={J.},
           giveni={J\bibinitperiod}}}%
      }
      \strng{namehash}{c6d965df16e600c701d4bcc2888ae34c}
      \strng{fullhash}{c6d965df16e600c701d4bcc2888ae34c}
      \strng{bibnamehash}{c6d965df16e600c701d4bcc2888ae34c}
      \strng{authorbibnamehash}{c6d965df16e600c701d4bcc2888ae34c}
      \strng{authornamehash}{c6d965df16e600c701d4bcc2888ae34c}
      \strng{authorfullhash}{c6d965df16e600c701d4bcc2888ae34c}
      \field{extraname}{11}
      \field{sortinit}{P}
      \field{sortinithash}{bb5b15f2db90f7aef79bb9e83defefcb}
      \field{labelprefix}{A}
      \field{labelnamesource}{author}
      \field{labeltitlesource}{title}
      \field{annotation}{This article is concerned with oscillations in the Mertens sum. The natural difficulty of this problem comes from the fact that the explicit formula for $M(x)$ contains terms of the form $x^{\rho}/\rho\zeta'(\rho)$, which are more difficult to handle than the terms $x^{\rho}/\rho$ appearing in the explicit formula for $\Delta^\psi(x)$. The author proves that if $\rho_0=\beta_0+i\gamma_0$ is a zero of $\zeta(s)$, then for $Y> e^{|\gamma_0|+4}$, \begin{align*} \frac{1}{Y} \int_{Y/(100 \log Y)}^Y |M(x)|\, dx > \frac{1}{6|\rho_0|^3}Y^{\beta_0} \quad \text{and} \quad \max_{x \leq Y} |M(x)| \geq \frac{1}{6|\rho_0|^3}Y^{\beta_0}. \end{align*} Consequently, using the first zeta zero $\frac12 + i\gamma_1$ with $\gamma_1\approx14.13$, \begin{align*} \max_{Y/(100 \log Y) \le x \le Y} |M(x)| \geq \frac{1}{17000}\sqrt{Y} \end{align*} for $Y\ge2$; the constant $1/17000$ can be improved but not enough to disprove the Mertens conjecture.\par This article cites~\cites{MR0218318, MR0043121}.}
      \field{issn}{0065-1036}
      \field{journaltitle}{Acta Arith.}
      \field{number}{1}
      \field{title}{Oscillatory properties of {$M(x)=\sum _{n\leq x}\mu (n)$}. {I}}
      \field{volume}{42}
      \field{year}{1982}
      \true{nocite}
      \field{pages}{49\bibrangedash 55}
      \range{pages}{7}
      \verb{doi}
      \verb 10.4064/aa-42-1-49-55
      \endverb
      \verb{mrnumber}
      \verb 678996
      \endverb
      \verb{urlraw}
      \verb https://doi.org/10.4064/aa-42-1-49-55
      \endverb
      \verb{url}
      \verb https://doi.org/10.4064/aa-42-1-49-55
      \endverb
    \endentry
    \entry{MR0724708}{article}{}
      \name{author}{1}{}{%
        {{hash=c6d965df16e600c701d4bcc2888ae34c}{%
           family={Pintz},
           familyi={P\bibinitperiod},
           given={J.},
           giveni={J\bibinitperiod}}}%
      }
      \strng{namehash}{c6d965df16e600c701d4bcc2888ae34c}
      \strng{fullhash}{c6d965df16e600c701d4bcc2888ae34c}
      \strng{bibnamehash}{c6d965df16e600c701d4bcc2888ae34c}
      \strng{authorbibnamehash}{c6d965df16e600c701d4bcc2888ae34c}
      \strng{authornamehash}{c6d965df16e600c701d4bcc2888ae34c}
      \strng{authorfullhash}{c6d965df16e600c701d4bcc2888ae34c}
      \field{extraname}{12}
      \field{sortinit}{P}
      \field{sortinithash}{bb5b15f2db90f7aef79bb9e83defefcb}
      \field{labelprefix}{A}
      \field{labelnamesource}{author}
      \field{labeltitlesource}{title}
      \field{annotation}{The author shows $\fA_{|1|}^{Q_k}(x) \gg Y^{1/2k}$ for each $k \geq 2$ using an unexpectedly simple argument, with effective estimates whose dependence on~$k$ is explicit. These results significantly strengthen work of Evelyn--Linfoot~\cites{MR1502996} and Kátai~\cite{MR0219496}.\par This article cites~\cites{MR1502996, MR0219496}.}
      \field{issn}{0024-6107,1469-7750}
      \field{journaltitle}{J. London Math. Soc. (2)}
      \field{number}{3}
      \field{title}{On the distribution of square-free numbers}
      \field{volume}{28}
      \field{year}{1983}
      \true{nocite}
      \field{pages}{401\bibrangedash 405}
      \range{pages}{5}
      \verb{doi}
      \verb 10.1112/jlms/s2-28.3.401
      \endverb
      \verb{mrnumber}
      \verb 724708
      \endverb
      \verb{urlraw}
      \verb https://doi.org/10.1112/jlms/s2-28.3.401
      \endverb
      \verb{url}
      \verb https://doi.org/10.1112/jlms/s2-28.3.401
      \endverb
    \endentry
    \entry{MR820251}{incollection}{}
      \name{author}{1}{}{%
        {{hash=c6d965df16e600c701d4bcc2888ae34c}{%
           family={Pintz},
           familyi={P\bibinitperiod},
           given={J.},
           giveni={J\bibinitperiod}}}%
      }
      \list{publisher}{1}{%
        {Birkhäuser, Basel}%
      }
      \strng{namehash}{c6d965df16e600c701d4bcc2888ae34c}
      \strng{fullhash}{c6d965df16e600c701d4bcc2888ae34c}
      \strng{bibnamehash}{c6d965df16e600c701d4bcc2888ae34c}
      \strng{authorbibnamehash}{c6d965df16e600c701d4bcc2888ae34c}
      \strng{authornamehash}{c6d965df16e600c701d4bcc2888ae34c}
      \strng{authorfullhash}{c6d965df16e600c701d4bcc2888ae34c}
      \field{extraname}{13}
      \field{sortinit}{P}
      \field{sortinithash}{bb5b15f2db90f7aef79bb9e83defefcb}
      \field{labelprefix}{A}
      \field{labelnamesource}{author}
      \field{labeltitlesource}{title}
      \field{annotation}{The author establishes two theorems that improve and simplify prior work of Turán and of Ingham. The first theorem states that if $\zeta(\beta_0 + i \gamma_0)=0$, then for $T$ sufficiently large in terms of $\gamma_0$, there exists an $x \in [T^{1/4}, T]$ for which $|\Delta^\psi(x)| \gg_{\gamma_0} x^{\beta_1}$ (with an explicit dependence on~$\gamma_0$). The second theorem assumes that $\zeta(s)\ne0$ in a region of the shape $\sigma \ge 1-\eta(t)$ where $\eta$ is continuous and decreasing, and, defining $\omega(x) = \min_{t\ge0} (\eta(t)\log x+\log t)$, concludes that $\Delta^\psi(x) = \Omega(x/e^{54\omega(x)})$. The main tool is the power-sum estimate of Sós and Turán. Near-optimal improvements (by the author) of these two results appeared slightly earlier~\cite{MR585891, MR598876}.\par This article cites~\cites{1891.Phragmen, MR585891, MR598876,1932.Ingham,MR0043121,MR0049219,MR0077579}.}
      \field{booktitle}{Studies in pure mathematics}
      \field{title}{Oscillatory properties of the remainder term of the prime number formula}
      \field{year}{1983}
      \true{nocite}
      \field{pages}{551\bibrangedash 560}
      \range{pages}{10}
      \verb{mrnumber}
      \verb 820251
      \endverb
    \endentry
    \entry{MR781183}{incollection}{}
      \name{author}{1}{}{%
        {{hash=c6d965df16e600c701d4bcc2888ae34c}{%
           family={Pintz},
           familyi={P\bibinitperiod},
           given={J.},
           giveni={J\bibinitperiod}}}%
      }
      \list{publisher}{1}{%
        {North-Holland, Amsterdam}%
      }
      \strng{namehash}{c6d965df16e600c701d4bcc2888ae34c}
      \strng{fullhash}{c6d965df16e600c701d4bcc2888ae34c}
      \strng{bibnamehash}{c6d965df16e600c701d4bcc2888ae34c}
      \strng{authorbibnamehash}{c6d965df16e600c701d4bcc2888ae34c}
      \strng{authornamehash}{c6d965df16e600c701d4bcc2888ae34c}
      \strng{authorfullhash}{c6d965df16e600c701d4bcc2888ae34c}
      \field{extraname}{14}
      \field{sortinit}{P}
      \field{sortinithash}{bb5b15f2db90f7aef79bb9e83defefcb}
      \field{labelprefix}{A}
      \field{labelnamesource}{author}
      \field{labeltitlesource}{title}
      \field{annotation}{The author investigates sign changes of $M(x)$ using lower bounds for $\fA^M_{|1|}(Y)$. He shows that if $\zeta(\rho_0)=0$, then $\fA^M_{|1|}(Y) > Y^{1+\beta_0}/6|\rho_0|^3$ for $Y>e^{|\gamma_0|+4}$, an effective improvement of an ineffective inequality of Kátai that provides the first effective disproof for the Mertens conjecture when RH is false. The author further concludes that $M(x)$ changes sign in every interval of the form $\bigl[ Y \exp(-3 \log _2 ^{3/2} Y), Y \bigr]$ when~$Y$ is sufficiently large.\par This article cites~\cites{MR0006202, MR0231801, MR0241374, MR0096630, MR0133287, MR0142500, 1909.Landau, 1897.Mertens, 1905.Stieltjes}.}
      \field{bookpagination}{none}
      \field{booktitle}{Topics in classical number theory, {V}ol. {I}, {II} ({B}udapest, 1981)}
      \field{series}{Colloq. Math. Soc. János Bolyai}
      \field{title}{On the partial sums of the {M}öbius function}
      \field{volume}{34}
      \field{year}{1984}
      \true{nocite}
      \field{pages}{1229\bibrangedash 1250}
      \range{pages}{22}
      \verb{mrnumber}
      \verb 781183
      \endverb
    \endentry
    \entry{MR756094}{incollection}{}
      \name{author}{1}{}{%
        {{hash=c6d965df16e600c701d4bcc2888ae34c}{%
           family={Pintz},
           familyi={P\bibinitperiod},
           given={J.},
           giveni={J\bibinitperiod}}}%
      }
      \list{publisher}{1}{%
        {Springer, Berlin}%
      }
      \strng{namehash}{c6d965df16e600c701d4bcc2888ae34c}
      \strng{fullhash}{c6d965df16e600c701d4bcc2888ae34c}
      \strng{bibnamehash}{c6d965df16e600c701d4bcc2888ae34c}
      \strng{authorbibnamehash}{c6d965df16e600c701d4bcc2888ae34c}
      \strng{authornamehash}{c6d965df16e600c701d4bcc2888ae34c}
      \strng{authorfullhash}{c6d965df16e600c701d4bcc2888ae34c}
      \field{extraname}{15}
      \field{sortinit}{P}
      \field{sortinithash}{bb5b15f2db90f7aef79bb9e83defefcb}
      \field{labelprefix}{A}
      \field{labelnamesource}{author}
      \field{labeltitlesource}{title}
      \field{annotation}{This article is primarily a summary of the results to be proved in the series \cite{MR585891,MR598876,MR688630,MR681441,MR681442, MR660714} by the author. The main functions of interest are $S(x) = \max_{0 \leq u \leq x} |\Delta^\psi(u)|$ and $\fA^\psi_{|1|}(x) = \int_0^x |\Delta^\psi(u)| \, du$. The following theorem is proved: Define $\omega(x) =\log \frac{x}{Z(x)}$, where $Z(x) = \max_{\rho} \frac{x^\beta}{|\gamma|}$. Then \[ \log \frac{x}{S(x)} \sim \log \frac{x^2}{\fA^\psi_{|1|}(x)} \sim \omega(x).\] In particular, this implies that $S(x)$ and $\frac1x\fA^\psi_{|1|}(x)$ are close in value, that is, the mean and maximum of $|\Delta^\psi(u)|$ are close. The proof uses a zero-density theorem of Carlson (for the upper bounds) and the power-sum method (for the lower bounds).\par This article cites~\cites{1932.Ingham,MR777013,MR0133309,MR0133308,MR0439771,MR0439772,1914.Littlewood,MR585891,MR598876,MR688630,MR681441,MR681442, MR660714,MR0067145,MR0111722,MR0146153}.}
      \field{booktitle}{Number theory, {N}oordwijkerhout 1983 ({N}oordwijkerhout, 1983)}
      \field{series}{Lecture Notes in Math.}
      \field{title}{On the remainder term of the prime number formula and the zeros of {R}iemann's zeta-function}
      \field{volume}{1068}
      \field{year}{1984}
      \true{nocite}
      \field{pages}{186\bibrangedash 197}
      \range{pages}{12}
      \verb{mrnumber}
      \verb 756094
      \endverb
    \endentry
    \entry{MR736725}{article}{}
      \name{author}{1}{}{%
        {{hash=c6d965df16e600c701d4bcc2888ae34c}{%
           family={Pintz},
           familyi={P\bibinitperiod},
           given={J.},
           giveni={J\bibinitperiod}}}%
      }
      \strng{namehash}{c6d965df16e600c701d4bcc2888ae34c}
      \strng{fullhash}{c6d965df16e600c701d4bcc2888ae34c}
      \strng{bibnamehash}{c6d965df16e600c701d4bcc2888ae34c}
      \strng{authorbibnamehash}{c6d965df16e600c701d4bcc2888ae34c}
      \strng{authornamehash}{c6d965df16e600c701d4bcc2888ae34c}
      \strng{authorfullhash}{c6d965df16e600c701d4bcc2888ae34c}
      \field{extraname}{16}
      \field{sortinit}{P}
      \field{sortinithash}{bb5b15f2db90f7aef79bb9e83defefcb}
      \field{labelprefix}{A}
      \field{labelnamesource}{author}
      \field{labeltitlesource}{title}
      \field{annotation}{By refining the proof method in his previous work~\cite{MR688630}, the author proves that if $\rho_0 = \beta_0 + i\gamma_0$ is a zero of $\zeta(s)$, then when $Y$ is sufficiently large in terms of~$\gamma_0$, \[ \max_{x\in [Y\exp( -5(\log\log Y)^{5/2} ),Y]} \frac{M(x)}{x^{\beta_0}} > \frac1{48|\rho_0|^3}, \quad \min_{x\in [Y\exp( -5(\log\log Y)^{5/2} ),Y]} \frac{M(x)}{x^{\beta_0}} < -\frac1{48|\rho_0|^3}. \] \par This article cites~\cites{MR0207665,MR630377,MR678996,MR688630,MR0155802,MR0231801,MR0219496}.}
      \field{issn}{0065-1036}
      \field{journaltitle}{Acta Arith.}
      \field{number}{2}
      \field{title}{Oscillatory properties of {$M(x)=\sum _{n\leq x}\mu (n)$}. {III}}
      \field{volume}{43}
      \field{year}{1984}
      \true{nocite}
      \field{pages}{105\bibrangedash 113}
      \range{pages}{9}
      \verb{doi}
      \verb 10.4064/aa-43-2-105-113
      \endverb
      \verb{mrnumber}
      \verb 736725
      \endverb
      \verb{urlraw}
      \verb https://doi.org/10.4064/aa-43-2-105-113
      \endverb
      \verb{url}
      \verb https://doi.org/10.4064/aa-43-2-105-113
      \endverb
    \endentry
    \entry{MR891440}{article}{}
      \name{author}{1}{}{%
        {{hash=c6d965df16e600c701d4bcc2888ae34c}{%
           family={Pintz},
           familyi={P\bibinitperiod},
           given={J.},
           giveni={J\bibinitperiod}}}%
      }
      \strng{namehash}{c6d965df16e600c701d4bcc2888ae34c}
      \strng{fullhash}{c6d965df16e600c701d4bcc2888ae34c}
      \strng{bibnamehash}{c6d965df16e600c701d4bcc2888ae34c}
      \strng{authorbibnamehash}{c6d965df16e600c701d4bcc2888ae34c}
      \strng{authornamehash}{c6d965df16e600c701d4bcc2888ae34c}
      \strng{authorfullhash}{c6d965df16e600c701d4bcc2888ae34c}
      \field{extraname}{17}
      \field{sortinit}{P}
      \field{sortinithash}{bb5b15f2db90f7aef79bb9e83defefcb}
      \field{labelprefix}{A}
      \field{labelnamesource}{author}
      \field{labeltitlesource}{title}
      \field{annotation}{The disproof of the Mertens conjecture by Odlyzko and te Riele~\cites{MR783538} does not provide any effective counterexample. In this article, the author shows that there exists $x<\exp(3.21 \cdot 10^{64})$, such that $|M(x)| > \sqrt{x}$.\par This article cites~\cites{MR783538, 1897.Mertens, MR0006202,MR429789}.}
      \field{issn}{0303-1179}
      \field{journaltitle}{Astérisque}
      \field{number}{147-148}
      \field{title}{An effective disproof of the Mertens conjecture}
      \field{year}{1987}
      \true{nocite}
      \field{pages}{325\bibrangedash 333\bibrangessep 346}
      \range{pages}{10}
      \verb{mrnumber}
      \verb 891440
      \endverb
    \endentry
    \entry{MR802311}{article}{}
      \name{author}{2}{}{%
        {{hash=c6d965df16e600c701d4bcc2888ae34c}{%
           family={Pintz},
           familyi={P\bibinitperiod},
           given={J.},
           giveni={J\bibinitperiod}}}%
        {{hash=c420715a04171fd2f8ed0fee9cd57cc0}{%
           family={Salerno},
           familyi={S\bibinitperiod},
           given={S.},
           giveni={S\bibinitperiod}}}%
      }
      \strng{namehash}{587c0ae365c9a78befc74bd5d1ffd220}
      \strng{fullhash}{587c0ae365c9a78befc74bd5d1ffd220}
      \strng{bibnamehash}{587c0ae365c9a78befc74bd5d1ffd220}
      \strng{authorbibnamehash}{587c0ae365c9a78befc74bd5d1ffd220}
      \strng{authornamehash}{587c0ae365c9a78befc74bd5d1ffd220}
      \strng{authorfullhash}{587c0ae365c9a78befc74bd5d1ffd220}
      \field{extraname}{1}
      \field{sortinit}{P}
      \field{sortinithash}{bb5b15f2db90f7aef79bb9e83defefcb}
      \field{labelprefix}{A}
      \field{labelnamesource}{author}
      \field{labeltitlesource}{title}
      \field{annotation}{The authors elaborate on their work in~\cite{MR756697} to handle prime number races where a bias is present. Again assuming a finite Riemann--Piltz conjecture, they show that when $Y$ is sufficiently large, \begin{align*} \frac{1}{Y}\int_{Y^{1-7/\lambda}}^{Y} |\pi(x;q,\ell_1,\ell_2)| \, dx \ge \sqrt Y \exp\left(-\frac{9 \log Y}{\lambda} - c_3 q \lambda (\log\log Y)^2\right), \end{align*} (and the same for $\theta$ in place of $\pi$) for any $\lambda$ satisfying \begin{align*} \frac{\sqrt{\log Y}}{\sqrt q \log \log Y} < \lambda < \frac{c_2 \log Y}{q \left(\log \log Y\right)^2}. \end{align*} They first deal with the case when both $\ell_1$ and $\ell_2$ are quadratic residues\mod q, using an explicit formula involving zeros of both $L(s,\chi)$ and $L(2s,\chi)$. In the remaining case when $\ell_1$ is a residue and $\ell_2$ is a nonresidue, there is an additional term corresponding to the pole of $L(2s, \chi_0)$ at $s=\frac12$.\par This article cites~\cites{MR756697,MR0142521,MR0142520,MR0125822}.}
      \field{issn}{0003-889X}
      \field{journaltitle}{Arch. Math. (Basel)}
      \field{number}{4}
      \field{title}{Irregularities in the distribution of primes in arithmetic progressions. {II}}
      \field{volume}{43}
      \field{year}{1984}
      \true{nocite}
      \field{pages}{351\bibrangedash 357}
      \range{pages}{7}
      \verb{doi}
      \verb 10.1007/BF01196659
      \endverb
      \verb{mrnumber}
      \verb 802311
      \endverb
      \verb{urlraw}
      \verb https://doi.org/10.1007/BF01196659
      \endverb
      \verb{url}
      \verb https://doi.org/10.1007/BF01196659
      \endverb
    \endentry
    \entry{MR764945}{article}{}
      \name{author}{2}{}{%
        {{hash=c6d965df16e600c701d4bcc2888ae34c}{%
           family={Pintz},
           familyi={P\bibinitperiod},
           given={J.},
           giveni={J\bibinitperiod}}}%
        {{hash=c420715a04171fd2f8ed0fee9cd57cc0}{%
           family={Salerno},
           familyi={S\bibinitperiod},
           given={S.},
           giveni={S\bibinitperiod}}}%
      }
      \strng{namehash}{587c0ae365c9a78befc74bd5d1ffd220}
      \strng{fullhash}{587c0ae365c9a78befc74bd5d1ffd220}
      \strng{bibnamehash}{587c0ae365c9a78befc74bd5d1ffd220}
      \strng{authorbibnamehash}{587c0ae365c9a78befc74bd5d1ffd220}
      \strng{authornamehash}{587c0ae365c9a78befc74bd5d1ffd220}
      \strng{authorfullhash}{587c0ae365c9a78befc74bd5d1ffd220}
      \field{extraname}{2}
      \field{sortinit}{P}
      \field{sortinithash}{bb5b15f2db90f7aef79bb9e83defefcb}
      \field{labelprefix}{A}
      \field{labelnamesource}{author}
      \field{labeltitlesource}{title}
      \field{annotation}{In this article, the authors obtain new estimates on $\psi(x;q,\ell_1,\ell_2)$ for arbitrary residues $\ell_1,\ell_2$. Assuming GRH$(cq^2 \log^6 q,E(q))$, the authors prove that for $Y$ sufficiently large, there exists \[ x \in \left[ Y\exp\left( - \frac{cq}{\sqrt{E(q)}}(\log Y)^{1/2} (\log \log Y)^{3/2} \right), Y \right]\] such that \[ \psi(x;q,\ell_1,\ell_2) > \sqrt{Y}\exp\left( - \frac{cq}{\sqrt{E(q)}}(\log Y)^{1/2} (\log \log Y)^{3/2} \right). \] This is an improvement over the work of Knapowski and Turán both in the localization and in the lower bound. The authors improve the power-sum bounds used by Knapowski and Turán to prove their results.\par This article cites~\cites{1909.Landau,MR0198145,MR0146156,MR0146157,MR0146158,MR0146159,MR0146160,MR0146161,MR0156826,MR0156827}.}
      \field{issn}{0391-173X}
      \field{journaltitle}{Ann. Scuola Norm. Sup. Pisa Cl. Sci. (4)}
      \field{number}{2}
      \field{title}{On the comparative theory of primes}
      \field{volume}{11}
      \field{year}{1984}
      \true{nocite}
      \field{pages}{245\bibrangedash 260}
      \range{pages}{16}
      \verb{mrnumber}
      \verb 764945
      \endverb
      \verb{urlraw}
      \verb http://www.numdam.org/item?id=ASNSP_1984_4_11_2_245_0
      \endverb
      \verb{url}
      \verb http://www.numdam.org/item?id=ASNSP_1984_4_11_2_245_0
      \endverb
    \endentry
    \entry{MR819064}{article}{}
      \name{author}{2}{}{%
        {{hash=c6d965df16e600c701d4bcc2888ae34c}{%
           family={Pintz},
           familyi={P\bibinitperiod},
           given={J.},
           giveni={J\bibinitperiod}}}%
        {{hash=c420715a04171fd2f8ed0fee9cd57cc0}{%
           family={Salerno},
           familyi={S\bibinitperiod},
           given={S.},
           giveni={S\bibinitperiod}}}%
      }
      \strng{namehash}{587c0ae365c9a78befc74bd5d1ffd220}
      \strng{fullhash}{587c0ae365c9a78befc74bd5d1ffd220}
      \strng{bibnamehash}{587c0ae365c9a78befc74bd5d1ffd220}
      \strng{authorbibnamehash}{587c0ae365c9a78befc74bd5d1ffd220}
      \strng{authornamehash}{587c0ae365c9a78befc74bd5d1ffd220}
      \strng{authorfullhash}{587c0ae365c9a78befc74bd5d1ffd220}
      \field{extraname}{3}
      \field{sortinit}{P}
      \field{sortinithash}{bb5b15f2db90f7aef79bb9e83defefcb}
      \field{labelprefix}{A}
      \field{labelnamesource}{author}
      \field{labeltitlesource}{title}
      \field{annotation}{Assuming HC$(E_q)$ and GRH$(D)$ with $D\geq D_0=c_0q^2\log^6q$ and $\lambda\geq 20D_0$, the authors show that when~$T$ is sufficiently large, there exists some~$k$ and some~$x$ near~$T$ such that \[ \psi_l(x,4k;q,a,b) > \sqrt{x}\exp\biggl( -\frac{LD^2}{\lambda^2}-\frac{cq^2}{E_q}\lambda\log^3L \biggr). \] The main tool in the proof is a ``one-sided'' version of the second main theorem of the power-sum method. When both~$a$ and~$b$ are quadratic residues or nonresidues, the same result holds with $\psi_l$ replaced by $\theta_l$.\par This article cites~\cites{1853.Chebyshev,1909.Landau,MR0162771,MR0174538,MR0180539,MR0182616A,MR0087685}.}
      \field{issn}{0236-5294}
      \field{journaltitle}{Acta Math. Hungar.}
      \field{number}{1-2}
      \field{title}{Accumulation theorems for primes in arithmetic progressions}
      \field{volume}{46}
      \field{year}{1985}
      \true{nocite}
      \field{pages}{151\bibrangedash 172}
      \range{pages}{22}
      \verb{doi}
      \verb 10.1007/BF01961016
      \endverb
      \verb{mrnumber}
      \verb 819064
      \endverb
    \endentry
    \entry{MR845900}{article}{}
      \name{author}{2}{}{%
        {{hash=c6d965df16e600c701d4bcc2888ae34c}{%
           family={Pintz},
           familyi={P\bibinitperiod},
           given={J.},
           giveni={J\bibinitperiod}}}%
        {{hash=c420715a04171fd2f8ed0fee9cd57cc0}{%
           family={Salerno},
           familyi={S\bibinitperiod},
           given={S.},
           giveni={S\bibinitperiod}}}%
      }
      \strng{namehash}{587c0ae365c9a78befc74bd5d1ffd220}
      \strng{fullhash}{587c0ae365c9a78befc74bd5d1ffd220}
      \strng{bibnamehash}{587c0ae365c9a78befc74bd5d1ffd220}
      \strng{authorbibnamehash}{587c0ae365c9a78befc74bd5d1ffd220}
      \strng{authornamehash}{587c0ae365c9a78befc74bd5d1ffd220}
      \strng{authorfullhash}{587c0ae365c9a78befc74bd5d1ffd220}
      \field{extraname}{4}
      \field{sortinit}{P}
      \field{sortinithash}{bb5b15f2db90f7aef79bb9e83defefcb}
      \field{labelprefix}{A}
      \field{labelnamesource}{author}
      \field{labeltitlesource}{title}
      \field{annotation}{In this article, the author establishes the effective results $\max_{Y^{7/8} \leq x \leq Y} E^\psi(x;q,a_1,a_2) > e^{-q^{c_1}}$ and $\max_{B_2(Y) \leq x \leq Y} E^\psi(x;q,a_1,a_2) \gg (\log Y)^{-q^{c_2}}$, where $B_2(Y) \gg Y(\log Y)^{-q^{c_3}}$, as well as the same statements where~$\psi$ is replaced by~$\theta$, $\Pi$, or~$\pi$. Furthermore, the author shows that $W^\psi(Y) \gg_q (\log Y)/\log\log Y$ with explicit constants, and similarly with~$\psi$ replaced by~$\Pi$ or, if~$a_1$ and~$a_2$ are both quadratic residues or both quadratic nonresidues modulo~$q$, with~$\theta$ or~$\pi$.\par This article cites~\cites{MR0146156, MR0146157, MR0146158, MR0146159, MR0146160,MR0146161, MR0156826, MR0156827, MR756697, MR764945}.}
      \field{issn}{0022-314X}
      \field{journaltitle}{J. Number Theory}
      \field{number}{2}
      \field{title}{Some consequences of the general Riemann hypothesis in the comparative theory of primes}
      \field{volume}{23}
      \field{year}{1986}
      \true{nocite}
      \field{pages}{183\bibrangedash 194}
      \range{pages}{12}
      \verb{doi}
      \verb 10.1016/0022-314X(86)90088-0
      \endverb
      \verb{mrnumber}
      \verb 845900
      \endverb
    \endentry
    \entry{MR756697}{article}{}
      \name{author}{2}{}{%
        {{hash=d192c3b49f427e2494b8aaea251e390b}{%
           family={Pintz},
           familyi={P\bibinitperiod},
           given={János},
           giveni={J\bibinitperiod}}}%
        {{hash=d91f53901946e1db0f20eee95cc2a03e}{%
           family={Salerno},
           familyi={S\bibinitperiod},
           given={Saverio},
           giveni={S\bibinitperiod}}}%
      }
      \strng{namehash}{13a7448d38e929fce6e0619ccab09b13}
      \strng{fullhash}{13a7448d38e929fce6e0619ccab09b13}
      \strng{bibnamehash}{13a7448d38e929fce6e0619ccab09b13}
      \strng{authorbibnamehash}{13a7448d38e929fce6e0619ccab09b13}
      \strng{authornamehash}{13a7448d38e929fce6e0619ccab09b13}
      \strng{authorfullhash}{13a7448d38e929fce6e0619ccab09b13}
      \field{extraname}{5}
      \field{sortinit}{P}
      \field{sortinithash}{bb5b15f2db90f7aef79bb9e83defefcb}
      \field{labelprefix}{A}
      \field{labelnamesource}{author}
      \field{labeltitlesource}{title}
      \field{annotation}{Assuming a finite Riemann--Piltz conjecture, the authors show that when $Y$ is sufficiently large, \begin{align*} \int_{Y^{1-7/\lambda}}^{Y} \psi(x;q,\ell_1,\ell_2) \,\frac{dx}x \gg \sqrt Y \exp\left(-\frac{2 \log Y}{\lambda} - c_3 q \lambda \log^2 Y\right) \end{align*} (and the same for $\Pi$ in place of $\psi$) for any $\lambda$ satisfying \begin{align*} \frac{\sqrt{\log Y}}{\sqrt q \log \log Y} < \lambda < \frac{c_2 \log Y}{q (\log \log Y)^2}. \end{align*} Essentially any such choice of $\lambda$ improves upon analogous results of Knapowski~\cite{MR0125822, MR0142520, MR0142521}. The proof also works for $\psi$ replaced by $\theta$ or $\pi$, but only if $\ell_1$ and $\ell_2$ are both quadratic nonresidues\mod q.\par This article cites~\cites{MR0182616A,MR819064,MR0142521,MR0142520,MR0125822}.}
      \field{issn}{0003-889X}
      \field{journaltitle}{Arch. Math. (Basel)}
      \field{number}{5}
      \field{title}{Irregularities in the distribution of primes in arithmetic progressions. {I}}
      \field{volume}{42}
      \field{year}{1984}
      \true{nocite}
      \field{pages}{439\bibrangedash 447}
      \range{pages}{9}
      \verb{doi}
      \verb 10.1007/BF01190694
      \endverb
      \verb{mrnumber}
      \verb 756697
      \endverb
      \verb{urlraw}
      \verb https://doi.org/10.1007/BF01190694
      \endverb
      \verb{url}
      \verb https://doi.org/10.1007/BF01190694
      \endverb
    \endentry
    \entry{MR3454375}{article}{}
      \name{author}{2}{}{%
        {{hash=b7137f17a7ecf01647eb7f9483185c94}{%
           family={Platt},
           familyi={P\bibinitperiod},
           given={D.\bibnamedelimi J.},
           giveni={D\bibinitperiod\bibinitdelim J\bibinitperiod}}}%
        {{hash=f8faf70b3f8ac0edff4b9b5337e687bf}{%
           family={Trudgian},
           familyi={T\bibinitperiod},
           given={T.\bibnamedelimi S.},
           giveni={T\bibinitperiod\bibinitdelim S\bibinitperiod}}}%
      }
      \strng{namehash}{c14c29110c40a8b34a8d4a8b0fedcbce}
      \strng{fullhash}{c14c29110c40a8b34a8d4a8b0fedcbce}
      \strng{bibnamehash}{c14c29110c40a8b34a8d4a8b0fedcbce}
      \strng{authorbibnamehash}{c14c29110c40a8b34a8d4a8b0fedcbce}
      \strng{authornamehash}{c14c29110c40a8b34a8d4a8b0fedcbce}
      \strng{authorfullhash}{c14c29110c40a8b34a8d4a8b0fedcbce}
      \field{extraname}{1}
      \field{sortinit}{P}
      \field{sortinithash}{bb5b15f2db90f7aef79bb9e83defefcb}
      \field{labelprefix}{A}
      \field{labelnamesource}{author}
      \field{labeltitlesource}{title}
      \field{annotation}{The authors compute that $\Delta^\theta(x) < 0$ for $0 \le x \le 1.39 \times 10^{17}$. By partial summation, this implies that $\Delta^\pi(x) < 0$ for $2 < x \le 1.39 \times 10^{17}$. The authors also prove that there exists $x \approx 1.3971623 \times 10^{316}$ for which $\Delta^\theta(x) > 0$.\par This article cites~\cites{MR1752093, 1932.Ingham, MR2420864, MR0202686, 1914.Littlewood, MR2271616, MR3356033, MR0067145}.}
      \field{issn}{0025-5718}
      \field{journaltitle}{Math. Comp.}
      \field{number}{299}
      \field{title}{On the first sign change of {$\theta(x)-x$}}
      \field{volume}{85}
      \field{year}{2016}
      \true{nocite}
      \field{pages}{1539\bibrangedash 1547}
      \range{pages}{9}
      \verb{doi}
      \verb 10.1090/mcom/3021
      \endverb
      \verb{mrnumber}
      \verb 3454375
      \endverb
      \verb{urlraw}
      \verb https://doi.org/10.1090/mcom/3021
      \endverb
      \verb{url}
      \verb https://doi.org/10.1090/mcom/3021
      \endverb
    \endentry
    \entry{MR3925757}{article}{}
      \name{author}{2}{}{%
        {{hash=7876864d71873e82aad2393b8cf18fe4}{%
           family={Platt},
           familyi={P\bibinitperiod},
           given={Dave},
           giveni={D\bibinitperiod}}}%
        {{hash=af8f95636cfb8ee32db62432b50fcea2}{%
           family={Trudgian},
           familyi={T\bibinitperiod},
           given={Tim},
           giveni={T\bibinitperiod}}}%
      }
      \strng{namehash}{d0b290fda16e446719a77447f29d0e95}
      \strng{fullhash}{d0b290fda16e446719a77447f29d0e95}
      \strng{bibnamehash}{d0b290fda16e446719a77447f29d0e95}
      \strng{authorbibnamehash}{d0b290fda16e446719a77447f29d0e95}
      \strng{authornamehash}{d0b290fda16e446719a77447f29d0e95}
      \strng{authorfullhash}{d0b290fda16e446719a77447f29d0e95}
      \field{extraname}{2}
      \field{sortinit}{P}
      \field{sortinithash}{bb5b15f2db90f7aef79bb9e83defefcb}
      \field{labelprefix}{A}
      \field{labelnamesource}{author}
      \field{labeltitlesource}{title}
      \field{annotation}{The authors show that GRH for $L(s,\chi_{-4})$ is equivalent, for any $0<\alpha<20.40442$, to \[ \lim_{x\to\infty} \sum_p \chi_{-4}(p) e^{-(x/n)^\alpha} = -\infty; \] they show also that the given method cannot be extended to $\alpha<20.40443$.\par This article cites~\cite{MR1985941,MR974088,MR1555148,MR1544293A,MR0202686,MR1329368}.}
      \field{issn}{1793-0421}
      \field{journaltitle}{Int. J. Number Theory}
      \field{number}{3}
      \field{title}{Fujii's development on {C}hebyshev's conjecture}
      \field{volume}{15}
      \field{year}{2019}
      \true{nocite}
      \field{pages}{639\bibrangedash 644}
      \range{pages}{6}
      \verb{doi}
      \verb 10.1142/S1793042119500337
      \endverb
      \verb{mrnumber}
      \verb 3925757
      \endverb
      \verb{urlraw}
      \verb https://doi.org/10.1142/S1793042119500337
      \endverb
      \verb{url}
      \verb https://doi.org/10.1142/S1793042119500337
      \endverb
    \endentry
    \entry{MR4249594}{book}{}
      \name{author}{1}{}{%
        {{hash=81a915ac987d508458c06ee0891cc3c2}{%
           family={Plymen},
           familyi={P\bibinitperiod},
           given={Roger},
           giveni={R\bibinitperiod}}}%
      }
      \list{publisher}{1}{%
        {American Mathematical Society, Providence, RI}%
      }
      \strng{namehash}{81a915ac987d508458c06ee0891cc3c2}
      \strng{fullhash}{81a915ac987d508458c06ee0891cc3c2}
      \strng{bibnamehash}{81a915ac987d508458c06ee0891cc3c2}
      \strng{authorbibnamehash}{81a915ac987d508458c06ee0891cc3c2}
      \strng{authornamehash}{81a915ac987d508458c06ee0891cc3c2}
      \strng{authorfullhash}{81a915ac987d508458c06ee0891cc3c2}
      \field{sortinit}{P}
      \field{sortinithash}{bb5b15f2db90f7aef79bb9e83defefcb}
      \field{labelprefix}{A}
      \field{labelnamesource}{author}
      \field{labeltitlesource}{title}
      \field{annotation}{Chapter~6 provides a comprehensive overview of the methods used to study the oscillations of $\pi(x)-\li(x)$. Chapter 7 outlines known results on the logarithmic density of the race between $\pi(x)$ and $\li(x)$, as well as upper bounds for the Skewes number. \par This book cites~\cites{MR1752093, MR2652902,1932.Ingham, MR0202686, MR3993767, 1914.Littlewood, MR866118, MR1329368, MR2684372, MR1573970, MR0067145,MR2813366,MR0882550}.}
      \field{isbn}{978-1-4704-6257-4}
      \field{series}{Student Mathematical Library}
      \field{title}{The {G}reat {P}rime {N}umber {R}ace}
      \field{volume}{92}
      \field{year}{2020}
      \true{nocite}
      \field{pages}{138}
      \range{pages}{1}
      \verb{mrnumber}
      \verb 4249594
      \endverb
    \endentry
    \entry{1919.Polya}{article}{}
      \name{author}{1}{}{%
        {{hash=85d5bb5658a4749b33270d8f0e8be228}{%
           family={Pólya},
           familyi={P\bibinitperiod},
           given={G.},
           giveni={G\bibinitperiod}}}%
      }
      \list{language}{1}{%
        {German}%
      }
      \strng{namehash}{85d5bb5658a4749b33270d8f0e8be228}
      \strng{fullhash}{85d5bb5658a4749b33270d8f0e8be228}
      \strng{bibnamehash}{85d5bb5658a4749b33270d8f0e8be228}
      \strng{authorbibnamehash}{85d5bb5658a4749b33270d8f0e8be228}
      \strng{authornamehash}{85d5bb5658a4749b33270d8f0e8be228}
      \strng{authorfullhash}{85d5bb5658a4749b33270d8f0e8be228}
      \field{extraname}{1}
      \field{sortinit}{P}
      \field{sortinithash}{bb5b15f2db90f7aef79bb9e83defefcb}
      \field{labelprefix}{A}
      \field{labelnamesource}{author}
      \field{labeltitlesource}{title}
      \field{annotation}{In Section~III, the author empirically observes that $L(n)\le 0$ for $2\le n\le1500$, with equality only for $n\in\{2,4,6,10,16,26,40,96,586\}$; he explains some of these equalities as resulting from imaginary quadratic fields $\mathbb Q(\sqrt{-p})$ with class number~$1$. The author also notes that RH would follow from the assertion that $L(n)\le 0$ for all large~$n$. However, he does not formulate this as a conjecture, saying only: ``{\it Ich teile diese Beobachtung mit, um evtl.~weitere numerische Untersuchung zu veranlassen. Der Bewies von [$L(n)\le 0$], sogar nur für hinreichend grosses $n$, würde den Beweis der Riemannschen Vermutung nach sich ziehen\ldots.}''.\par This article cites~\cite{1909.Landau}.}
      \field{journaltitle}{Jahresbericht der deutschen Math.--Vereinigung}
      \field{title}{{Verschiedene Bemerkungen zur Zahlentheorie}}
      \field{volume}{28}
      \field{year}{1919}
      \true{nocite}
      \field{pages}{31\bibrangedash 40}
      \range{pages}{10}
    \endentry
    \entry{1930.Polya}{article}{}
      \name{author}{1}{}{%
        {{hash=85d5bb5658a4749b33270d8f0e8be228}{%
           family={Pólya},
           familyi={P\bibinitperiod},
           given={G.},
           giveni={G\bibinitperiod}}}%
      }
      \list{language}{1}{%
        {German}%
      }
      \strng{namehash}{85d5bb5658a4749b33270d8f0e8be228}
      \strng{fullhash}{85d5bb5658a4749b33270d8f0e8be228}
      \strng{bibnamehash}{85d5bb5658a4749b33270d8f0e8be228}
      \strng{authorbibnamehash}{85d5bb5658a4749b33270d8f0e8be228}
      \strng{authornamehash}{85d5bb5658a4749b33270d8f0e8be228}
      \strng{authorfullhash}{85d5bb5658a4749b33270d8f0e8be228}
      \field{extraname}{2}
      \field{sortinit}{P}
      \field{sortinithash}{bb5b15f2db90f7aef79bb9e83defefcb}
      \field{labelprefix}{A}
      \field{labelnamesource}{author}
      \field{labeltitlesource}{title}
      \field{annotation}{The author proves that $\limsup_{T\to\infty} W^\psi(T)/\log T \ge \gamma_1/\pi$, where $\gamma_1 \approx 14.135$ is the smallest ordinate of a nontrivial zero of $\zeta(s)$. Indeed, $\gamma_1$ can be replaced by the minimum positive ordinate of the zeros of $\zeta(s)$ having maximal real part, or~$\infty$ if no such zero exists. It was known for a long while that this article contained an error that could be mended; the corrected version finally appeared as~\cite{MR0263757}.\par This article cites~\cite{1891.Phragmen,MR1511206}.}
      \field{journaltitle}{Gött. Nachr.}
      \field{title}{{Über das Vorzeichen des Restgliedes im Primzahltheorie}}
      \field{year}{1930}
      \true{nocite}
      \field{pages}{19\bibrangedash 27}
      \range{pages}{9}
    \endentry
    \entry{MR1576856}{article}{}
      \name{author}{1}{}{%
        {{hash=85d5bb5658a4749b33270d8f0e8be228}{%
           family={Pólya},
           familyi={P\bibinitperiod},
           given={G.},
           giveni={G\bibinitperiod}}}%
      }
      \strng{namehash}{85d5bb5658a4749b33270d8f0e8be228}
      \strng{fullhash}{85d5bb5658a4749b33270d8f0e8be228}
      \strng{bibnamehash}{85d5bb5658a4749b33270d8f0e8be228}
      \strng{authorbibnamehash}{85d5bb5658a4749b33270d8f0e8be228}
      \strng{authornamehash}{85d5bb5658a4749b33270d8f0e8be228}
      \strng{authorfullhash}{85d5bb5658a4749b33270d8f0e8be228}
      \field{extraname}{3}
      \field{sortinit}{P}
      \field{sortinithash}{bb5b15f2db90f7aef79bb9e83defefcb}
      \field{labelprefix}{A}
      \field{labelnamesource}{author}
      \field{labeltitlesource}{title}
      \field{annotation}{Let $\omega(u)$ be a real-valued function and set $\Phi(s) = \int_1^\infty \omega(u)u^{-s} \, du$. If the number of sign changes $W^\omega(x)$ is $O(\log x)$, then a rightmost singularity~$\sigma+it$ of $\Phi(s)$ (that is, a singularity with $\sigma$ maximal) exists and satisfies $0\le \gamma \le \limsup_{x\to\infty} \pi W^\omega(x)/\log x$. This result generalizes Landau's theorem, which is the case $W^\omega(x)\ll1$.}
      \field{issn}{0024-6115}
      \field{journaltitle}{Proc. London Math. Soc. (2)}
      \field{number}{2}
      \field{title}{On polar singularities of power series and of Dirichlet series}
      \field{volume}{33}
      \field{year}{1931}
      \true{nocite}
      \field{pages}{85\bibrangedash 101}
      \range{pages}{17}
      \verb{doi}
      \verb 10.1112/plms/s2-33.1.85
      \endverb
      \verb{mrnumber}
      \verb 1576856
      \endverb
    \endentry
    \entry{MR0263757}{article}{}
      \name{author}{1}{}{%
        {{hash=85d5bb5658a4749b33270d8f0e8be228}{%
           family={Pólya},
           familyi={P\bibinitperiod},
           given={G.},
           giveni={G\bibinitperiod}}}%
      }
      \list{language}{1}{%
        {German}%
      }
      \list{publisher}{1}{%
        {Plenum, New York}%
      }
      \strng{namehash}{85d5bb5658a4749b33270d8f0e8be228}
      \strng{fullhash}{85d5bb5658a4749b33270d8f0e8be228}
      \strng{bibnamehash}{85d5bb5658a4749b33270d8f0e8be228}
      \strng{authorbibnamehash}{85d5bb5658a4749b33270d8f0e8be228}
      \strng{authornamehash}{85d5bb5658a4749b33270d8f0e8be228}
      \strng{authorfullhash}{85d5bb5658a4749b33270d8f0e8be228}
      \field{extraname}{4}
      \field{sortinit}{P}
      \field{sortinithash}{bb5b15f2db90f7aef79bb9e83defefcb}
      \field{labelprefix}{A}
      \field{labelnamesource}{author}
      \field{labeltitlesource}{title}
      \field{annotation}{This is a corrected version of the author's earlier article~\cite{1930.Polya}.}
      \field{booktitle}{Number Theory and Analysis (Papers in Honor of Edmund Landau)}
      \field{title}{Über das Vorzeichen des Restgliedes im Primzahlsatz}
      \field{year}{1969}
      \true{nocite}
      \field{pages}{233\bibrangedash 244}
      \range{pages}{12}
      \verb{mrnumber}
      \verb 0263757
      \endverb
    \endentry
    \entry{MR0087685}{book}{}
      \name{author}{1}{}{%
        {{hash=fa022d8efe9c2814c0c7cd5fe921a192}{%
           family={Prachar},
           familyi={P\bibinitperiod},
           given={Karl},
           giveni={K\bibinitperiod}}}%
      }
      \list{language}{1}{%
        {German}%
      }
      \list{publisher}{1}{%
        {Springer-Verlag, Berlin-Göttingen-Heidelberg}%
      }
      \strng{namehash}{fa022d8efe9c2814c0c7cd5fe921a192}
      \strng{fullhash}{fa022d8efe9c2814c0c7cd5fe921a192}
      \strng{bibnamehash}{fa022d8efe9c2814c0c7cd5fe921a192}
      \strng{authorbibnamehash}{fa022d8efe9c2814c0c7cd5fe921a192}
      \strng{authornamehash}{fa022d8efe9c2814c0c7cd5fe921a192}
      \strng{authorfullhash}{fa022d8efe9c2814c0c7cd5fe921a192}
      \field{sortinit}{P}
      \field{sortinithash}{bb5b15f2db90f7aef79bb9e83defefcb}
      \field{labelprefix}{A}
      \field{labelnamesource}{author}
      \field{labeltitlesource}{title}
      \field{annotation}{The contents of this book are divided into the following ten main chapters, plus an Appendix that provides additional background on topics such as partial summation and the $\Gamma$-function. The chapters are: I. Elementary results; II. Sieve methods; III. The prime number theorem; IV. Primes in an arithmetical progression; V. Different applications; VI. The Goldbach problem; VII. Function theoretic properties of the $L$-functions. Explicit formulae and their applications; VIII. Trigonometric sums; IX. Theorems on the density of zeros of the $L$-functions and their application in prime number theory; X. The smallest prime in an arithmetical progression.}
      \field{title}{Primzahlverteilung}
      \field{year}{1957}
      \true{nocite}
      \field{pages}{x+415 pp.}
      \range{pages}{-1}
      \verb{mrnumber}
      \verb 0087685
      \endverb
    \endentry
    \entry{MR1761276}{article}{}
      \name{author}{1}{}{%
        {{hash=2e461074b5178ec74d42e5ab1268d913}{%
           family={Puchta},
           familyi={P\bibinitperiod},
           given={J.-C.},
           giveni={J\bibinithyphendelim C\bibinitperiod}}}%
      }
      \strng{namehash}{2e461074b5178ec74d42e5ab1268d913}
      \strng{fullhash}{2e461074b5178ec74d42e5ab1268d913}
      \strng{bibnamehash}{2e461074b5178ec74d42e5ab1268d913}
      \strng{authorbibnamehash}{2e461074b5178ec74d42e5ab1268d913}
      \strng{authornamehash}{2e461074b5178ec74d42e5ab1268d913}
      \strng{authorfullhash}{2e461074b5178ec74d42e5ab1268d913}
      \field{sortinit}{P}
      \field{sortinithash}{bb5b15f2db90f7aef79bb9e83defefcb}
      \field{labelprefix}{A}
      \field{labelnamesource}{author}
      \field{labeltitlesource}{title}
      \field{annotation}{Assuming GRH, the author shows that $E^\psi(x;q,1,\min) = \Omega_+(\log\log\log x)$ and $E^\psi(x;q,1,\max) = \Omega_-(\log\log\log x)$ where all constants are effective. In particular, the $q=1$ case gives an effective version of the classical result of Littlewood~\cites{1914.Littlewood}. One tool is an evaluation of the moments of $E^{\psi_\chi}(u)$ for Dirichlet characters~$\chi$: the author shows that $\frac{1}{y}\int_0^y E^{\psi_\chi}(u)^k \, \frac{du}u \to (-1)^k \sum_{\gamma_1 + \cdots + \gamma_k=0} {1/\rho_1 \cdots \rho_k}$, where the sum runs over all nontrivial zeros of $L(s,\chi)$.\par This article cites \cites{1914.Littlewood, MR0004255,MR1083000,MR1083001,MR1105605,MR1105608, MR1310738,MR1251926}.}
      \field{issn}{0236-5294}
      \field{journaltitle}{Acta Math. Hungar.}
      \field{number}{3}
      \field{title}{On large oscillations of the remainder of the prime number theorems}
      \field{volume}{87}
      \field{year}{2000}
      \true{nocite}
      \field{pages}{213\bibrangedash 227}
      \range{pages}{15}
      \verb{mrnumber}
      \verb 1761276
      \endverb
    \endentry
    \entry{MR2110393}{article}{}
      \name{author}{1}{}{%
        {{hash=2a44d6571c930fba6c9dbf36ab288ad2}{%
           family={Radziejewski},
           familyi={R\bibinitperiod},
           given={Maciej},
           giveni={M\bibinitperiod}}}%
      }
      \strng{namehash}{2a44d6571c930fba6c9dbf36ab288ad2}
      \strng{fullhash}{2a44d6571c930fba6c9dbf36ab288ad2}
      \strng{bibnamehash}{2a44d6571c930fba6c9dbf36ab288ad2}
      \strng{authorbibnamehash}{2a44d6571c930fba6c9dbf36ab288ad2}
      \strng{authornamehash}{2a44d6571c930fba6c9dbf36ab288ad2}
      \strng{authorfullhash}{2a44d6571c930fba6c9dbf36ab288ad2}
      \field{extraname}{1}
      \field{sortinit}{R}
      \field{sortinithash}{b9c68a358aea118dfa887b6e902414a7}
      \field{labelprefix}{A}
      \field{labelnamesource}{author}
      \field{labeltitlesource}{title}
      \field{annotation}{Given any number field~$K$ and any subgroup~$\Gamma$ of the narrow class group of~$K$, there is a notion of irreducibility of ideals in~$\Gamma$; the corresponding factorization into irreducibles can be nonunique, and indeed the lengths of such factorizations can also be nonunique. The author examines the counting functions of ideals in~$\Gamma$ whose set of factorization lengths has cardinality lying in a prescribed interval, or (alternatively) contains one of a prescribed set of positive integers. The author obtains oscillation results for the corresponding error terms of the form $\Omega(x^{1/2-\ep})$, as well as lower bounds of size $\log X$ for the number of sign changes up to~$X$. \par This article cites~\cites{MR891057,MR2123959}.}
      \field{issn}{0065-1036}
      \field{journaltitle}{Acta Arith.}
      \field{number}{2}
      \field{title}{On the distribution of algebraic numbers with prescribed factorization properties}
      \field{volume}{116}
      \field{year}{2005}
      \true{nocite}
      \field{pages}{153\bibrangedash 171}
      \range{pages}{19}
      \verb{doi}
      \verb 10.4064/aa116-2-4
      \endverb
      \verb{mrnumber}
      \verb 2110393
      \endverb
    \endentry
    \entry{MR2123959}{article}{}
      \name{author}{1}{}{%
        {{hash=2a44d6571c930fba6c9dbf36ab288ad2}{%
           family={Radziejewski},
           familyi={R\bibinitperiod},
           given={Maciej},
           giveni={M\bibinitperiod}}}%
      }
      \strng{namehash}{2a44d6571c930fba6c9dbf36ab288ad2}
      \strng{fullhash}{2a44d6571c930fba6c9dbf36ab288ad2}
      \strng{bibnamehash}{2a44d6571c930fba6c9dbf36ab288ad2}
      \strng{authorbibnamehash}{2a44d6571c930fba6c9dbf36ab288ad2}
      \strng{authornamehash}{2a44d6571c930fba6c9dbf36ab288ad2}
      \strng{authorfullhash}{2a44d6571c930fba6c9dbf36ab288ad2}
      \field{extraname}{2}
      \field{sortinit}{R}
      \field{sortinithash}{b9c68a358aea118dfa887b6e902414a7}
      \field{labelprefix}{A}
      \field{labelnamesource}{author}
      \field{labeltitlesource}{title}
      \field{annotation}{Kaczorowski and Pintz~\cite{MR891057} considered functions of a real variable whose Mellin transforms had singularities of the form $(s-\rho)^w P(\log(s-\rho))$ for some $w\in\C$. The author extends this class of functions to those with singularities that are linear combinations of this type, showing that such functions (suitably normalized) have oscillations of the form $\Omega(x^{1/2-\ep})$, as well as lower bounds of size $\log X$ for the number of sign changes up to~$X$. The motivation was to address error terms of counting functions corresponding to certain ideal factorization problems in number fields~\cite{MR2110393}, and the author provides an application to counting the number of ideals in a subgroup of the narrow class group all of whose restricted factorizations have the same length.\par This article cites~\cites{MR0198145,MR1555148,1936.Ingham,MR1256149,MR0219496,MR1511360,MR0439774,1930.Polya,MR891057,MR858395}.}
      \field{issn}{0026-9255}
      \field{journaltitle}{Monatsh. Math.}
      \field{number}{2}
      \field{title}{Oscillations of error terms associated with certain arithmetical functions}
      \field{volume}{144}
      \field{year}{2005}
      \true{nocite}
      \field{pages}{113\bibrangedash 130}
      \range{pages}{18}
      \verb{doi}
      \verb 10.1007/s00605-003-0147-x
      \endverb
      \verb{mrnumber}
      \verb 2123959
      \endverb
    \endentry
    \entry{MR3179660}{article}{}
      \name{author}{1}{}{%
        {{hash=2a44d6571c930fba6c9dbf36ab288ad2}{%
           family={Radziejewski},
           familyi={R\bibinitperiod},
           given={Maciej},
           giveni={M\bibinitperiod}}}%
      }
      \strng{namehash}{2a44d6571c930fba6c9dbf36ab288ad2}
      \strng{fullhash}{2a44d6571c930fba6c9dbf36ab288ad2}
      \strng{bibnamehash}{2a44d6571c930fba6c9dbf36ab288ad2}
      \strng{authorbibnamehash}{2a44d6571c930fba6c9dbf36ab288ad2}
      \strng{authornamehash}{2a44d6571c930fba6c9dbf36ab288ad2}
      \strng{authorfullhash}{2a44d6571c930fba6c9dbf36ab288ad2}
      \field{extraname}{3}
      \field{sortinit}{R}
      \field{sortinithash}{b9c68a358aea118dfa887b6e902414a7}
      \field{labelprefix}{A}
      \field{labelnamesource}{author}
      \field{labeltitlesource}{title}
      \field{annotation}{The author summarizes many existing Landau-type oscillation theorems, and establishes a new oscillation theorem for functions whose Mellin transforms satisfy certain growth conditions (weaker than usual) and have certain singularities (more general than usual). As an application, he shows that the counting function of numbers that can be written as the sum of two squares has oscillations of size $x^{1/2}(\log x)^{-3/2-\ep}$; he generalizes this result to $\sum_{n\le x} (\frac14r(n))^z$, where $r(n)$ is the number of representations of~$n$ as the sum of two squares, for generic~$z$ (but not $z=1$).\par This article cites~\cites{1891.Phragmen,MR0146156,MR0146159,MR0156827,MR858395,MR891057,MR1511360,MR2110393,MR2123959,MR2470105,MR2730494}.}
      \field{issn}{0033-5606}
      \field{journaltitle}{Q. J. Math.}
      \field{number}{1}
      \field{title}{Oscillatory properties of real functions with weakly bounded Mellin transform}
      \field{volume}{65}
      \field{year}{2014}
      \true{nocite}
      \field{pages}{249\bibrangedash 266}
      \range{pages}{18}
      \verb{doi}
      \verb 10.1093/qmath/has036
      \endverb
      \verb{mrnumber}
      \verb 3179660
      \endverb
    \endentry
    \entry{MR549977}{article}{}
      \name{author}{1}{}{%
        {{hash=f6fbcde367a0ed5df3da06b96b6d8eff}{%
           family={Riele},
           familyi={R\bibinitperiod},
           given={H.\bibnamedelimi J.\bibnamedelimi J.},
           giveni={H\bibinitperiod\bibinitdelim J\bibinitperiod\bibinitdelim J\bibinitperiod},
           prefix={te},
           prefixi={t\bibinitperiod}}}%
      }
      \strng{namehash}{f6fbcde367a0ed5df3da06b96b6d8eff}
      \strng{fullhash}{f6fbcde367a0ed5df3da06b96b6d8eff}
      \strng{bibnamehash}{f6fbcde367a0ed5df3da06b96b6d8eff}
      \strng{authorbibnamehash}{f6fbcde367a0ed5df3da06b96b6d8eff}
      \strng{authornamehash}{f6fbcde367a0ed5df3da06b96b6d8eff}
      \strng{authorfullhash}{f6fbcde367a0ed5df3da06b96b6d8eff}
      \field{extraname}{1}
      \field{sortinit}{R}
      \field{sortinithash}{b9c68a358aea118dfa887b6e902414a7}
      \field{labelprefix}{A}
      \field{labelnamesource}{author}
      \field{labeltitlesource}{title}
      \field{annotation}{The author introduces some modifications of the method of Jurkat and Peyerimhoff~\cite{MR429789} and proves that \[ \liminf_{x \to \infty} E^M(x) < -0.843 \quad \text{and} \quad \limsup_{x \to \infty} E^M(x) > 0.860. \] The calculation took several hundred CPU-hours on a CDC Cyber 73/173 system. \par This article cites~\cite{1897.Mertens,MR0330069,MR228456,MR429789}.}
      \field{issn}{0075-4102}
      \field{journaltitle}{J. Reine Angew. Math.}
      \field{title}{Computations concerning the conjecture of {M}ertens}
      \field{volume}{311(312)}
      \field{year}{1979}
      \true{nocite}
      \field{pages}{356\bibrangedash 360}
      \range{pages}{5}
      \verb{doi}
      \verb 10.1515/crll.1979.311-312.356
      \endverb
      \verb{mrnumber}
      \verb 549977
      \endverb
      \verb{urlraw}
      \verb https://doi.org/10.1515/crll.1979.311-312.356
      \endverb
      \verb{url}
      \verb https://doi.org/10.1515/crll.1979.311-312.356
      \endverb
    \endentry
    \entry{MR866118}{article}{}
      \name{author}{1}{}{%
        {{hash=492b43ecdc1e86c9cf9402cb28b077f2}{%
           family={Riele},
           familyi={R\bibinitperiod},
           given={Herman\bibnamedelimb J.\bibnamedelimi J.},
           giveni={H\bibinitperiod\bibinitdelim J\bibinitperiod\bibinitdelim J\bibinitperiod},
           prefix={te},
           prefixi={t\bibinitperiod}}}%
      }
      \strng{namehash}{492b43ecdc1e86c9cf9402cb28b077f2}
      \strng{fullhash}{492b43ecdc1e86c9cf9402cb28b077f2}
      \strng{bibnamehash}{492b43ecdc1e86c9cf9402cb28b077f2}
      \strng{authorbibnamehash}{492b43ecdc1e86c9cf9402cb28b077f2}
      \strng{authornamehash}{492b43ecdc1e86c9cf9402cb28b077f2}
      \strng{authorfullhash}{492b43ecdc1e86c9cf9402cb28b077f2}
      \field{extraname}{2}
      \field{sortinit}{R}
      \field{sortinithash}{b9c68a358aea118dfa887b6e902414a7}
      \field{labelprefix}{A}
      \field{labelnamesource}{author}
      \field{labeltitlesource}{title}
      \field{annotation}{The author shows that $\pi(x) > \li(x)$ for some $6.62 \times 10^{370} \leq x \leq 6.69 \times 10^{370}$, thereby improving the previous best estimate, $1.65 \times 10^{1165}$, for Skewes's number found by Lehman~\cite{MR0202686}. Using an explicit formula for $E^\pi(e^u)$ averaged by a Gaussian kernel, Lehman had found three candidates for $x$ near which $\pi(x) > \li(x)$, namely $e^{727.952}$, $e^{853.853}$, and $e^{2682.977}$. Lehman showed that $e^{2682.977}$ produced an actual example; using the zeros of $\zeta(s)$ up to height $5 \times 10^4$, found on a CYBER 205 supercomputer located at the Academic Computer Centre Amsterdam, the author shows that $e^{853.853}$ produces an actual example. The author speculates that zeros up to height $4 \times 10^5$ would be required to determine whether there is an actual example around $e^{727.952}$. \par This article cites~\cites{MR0202686,MR0067145,1914.Littlewood}.}
      \field{issn}{0025-5718}
      \field{journaltitle}{Math. Comp.}
      \field{number}{177}
      \field{title}{On the sign of the difference {$\pi(x)-\li(x)$}}
      \field{volume}{48}
      \field{year}{1987}
      \true{nocite}
      \field{pages}{323\bibrangedash 328}
      \range{pages}{6}
      \verb{doi}
      \verb 10.2307/2007893
      \endverb
      \verb{mrnumber}
      \verb 866118
      \endverb
      \verb{urlraw}
      \verb https://doi.org/10.2307/2007893
      \endverb
      \verb{url}
      \verb https://doi.org/10.2307/2007893
      \endverb
    \endentry
    \entry{MR729173}{incollection}{}
      \name{author}{1}{}{%
        {{hash=4440c3774c4e033603367ee77d77249f}{%
           family={Robin},
           familyi={R\bibinitperiod},
           given={G.},
           giveni={G\bibinitperiod}}}%
      }
      \list{language}{1}{%
        {French}%
      }
      \list{publisher}{1}{%
        {Birkhäuser Boston, Boston, MA}%
      }
      \strng{namehash}{4440c3774c4e033603367ee77d77249f}
      \strng{fullhash}{4440c3774c4e033603367ee77d77249f}
      \strng{bibnamehash}{4440c3774c4e033603367ee77d77249f}
      \strng{authorbibnamehash}{4440c3774c4e033603367ee77d77249f}
      \strng{authornamehash}{4440c3774c4e033603367ee77d77249f}
      \strng{authorfullhash}{4440c3774c4e033603367ee77d77249f}
      \field{extraname}{1}
      \field{sortinit}{R}
      \field{sortinithash}{b9c68a358aea118dfa887b6e902414a7}
      \field{labelprefix}{A}
      \field{labelnamesource}{author}
      \field{labeltitlesource}{title}
      \field{annotation}{The author investigates the function $\frac{\sigma(n)}n$ and proves that RH is true if and only if $\frac{\sigma(n)}n <e^{C_0}\log\log{n}$ when~$n$ is sufficiently large. If~$(C_k)$ is the sequence of colossally abundant numbers, the author further shows that $(\sigma(C_k)/C_k\log\log C_k)$ has infinitely many local extrema. Finally, the author shows that if RH is false, then both $\Delta^{\pi_r}(x)$ and $\prod_{p\leq x} (1-\frac1p)-1/e^{C_0}\log{x}$ change signs infinitely often.\par This article cites~\cites{MR0417077,MR0202685,MR0137689}}
      \field{bookpagination}{none}
      \field{booktitle}{Seminar on number theory, {P}aris 1981--82 ({P}aris, 1981/1982)}
      \field{series}{Progr. Math.}
      \field{title}{Sur l'ordre maximum de la fonction somme des diviseurs}
      \field{volume}{38}
      \field{year}{1983}
      \true{nocite}
      \field{pages}{233\bibrangedash 244}
      \range{pages}{12}
      \verb{mrnumber}
      \verb 729173
      \endverb
    \endentry
    \entry{MR928842}{article}{}
      \name{author}{1}{}{%
        {{hash=28b00c8013a39aa054af5aea7f8c5286}{%
           family={Robin},
           familyi={R\bibinitperiod},
           given={Guy},
           giveni={G\bibinitperiod}}}%
      }
      \list{language}{1}{%
        {French}%
      }
      \strng{namehash}{28b00c8013a39aa054af5aea7f8c5286}
      \strng{fullhash}{28b00c8013a39aa054af5aea7f8c5286}
      \strng{bibnamehash}{28b00c8013a39aa054af5aea7f8c5286}
      \strng{authorbibnamehash}{28b00c8013a39aa054af5aea7f8c5286}
      \strng{authornamehash}{28b00c8013a39aa054af5aea7f8c5286}
      \strng{authorfullhash}{28b00c8013a39aa054af5aea7f8c5286}
      \field{extraname}{2}
      \field{sortinit}{R}
      \field{sortinithash}{b9c68a358aea118dfa887b6e902414a7}
      \field{labelprefix}{A}
      \field{labelnamesource}{author}
      \field{labeltitlesource}{title}
      \field{annotation}{This article examines, assuming HC, the asymptotic behaviour of the weighted average $ \cP(x)=\sum_{n\leq x} \mathring\Delta(n;k,\ell) n^{-\alpha}\log^\beta{n} $ where $\alpha$ and $\beta$ are fixed real numbers. If GRH is false, then $\cP(x) \ll 1+x^{1-\alpha+\Theta_k}\log^{\beta-1}{x}$ and $\cP(x)=\Omega_{\pm}(x^{1-\alpha+\Theta_k-\ep})$; under the additional assumption of SA, we have $\cP(x)=\Omega_\pm (x^{1-\alpha+\Theta_k}\log^{\beta-1}{x})$ (which is thus best possible for $\alpha < 1+\Theta_k$).\par If GRH is true, the behaviour depends more significantly upon~$\alpha$ and~$\beta$. When $\alpha>\frac32$, we have $\cP(x)\ll1$. When $\alpha=\frac32$, we have $\cP(x)\ll1$ if $\beta<0$, and $\cP(x)=(1-c_k(\ell))\log\log{x}+O(1)$ if $\beta=0$, and $\cP(x)=(1-c_k(\ell))(\log x)^\beta/\beta+O((\log{x})^{\beta-1}\log\log{x})$ if $\beta>0$. Finally, when $\alpha<\frac32$, we have $\cP(x) \ll x^{3/2-\alpha}\log^{\beta-1}{x}$. This theorem disproves Shanks's conjecture $\sum_{n\le x}\pi(n;4,3,1) n^{1/2}/{\pi(n)}\sim x$, as well as corresponding conjectures for other moduli. Moreover, it shows that Brent's conjecture $\sum_{n\le x}\pi(n;4,3,1) /n^{1/2}{\pi(n)}\sim \log x$ is equivalent to GRH.\par Again assuming GRH and $\alpha<\frac32$, the author asserts that for certain moduli including $3$, $4$, $5$, $6$, $7$, $8$, $9$, $10$, $12$, there exists a constant $\alpha_{k,\ell}$ such that for $\alpha>\alpha_{k,\ell}$, when $x$ is sufficiently large then $\cP(x)<0$ if $\ell$ is a quadratic residue and $\cP(x)>0$ if $\ell$ is a quadratic nonresidue. (It seems that this result actually holds for all moduli $k\ge3$.) On the other hand, for some moduli including $23, 43, 67, 163$, there exists a constant $\alpha'_{k,\ell}$ such that for $\alpha<\alpha'_{k,\ell}$, we have $\cP(x)=\Omega_\pm(x^{3/2-\alpha}\log^{\beta-1}{x})$.\par This article cites~\cites{MR1555148, MR1544293A, MR1544293B, MR0083001, MR0108470, MR0146156, MR0146157, MR0146158, MR0146159, MR0146160, MR0146161, MR0156826, MR0156827, MR0162771, MR0174538, MR0180539, MR0182616A, MR0182616B, MR0200250, MR0289452, MR0302585, MR0369287, MR0417077, MR595317, MR654520, MR602236, MR675189}.}
      \field{issn}{0240-2955}
      \field{journaltitle}{Ann. Fac. Sci. Toulouse Math. (5)}
      \field{number}{2}
      \field{title}{Irrégularités dans la distribution des nombres premiers dans les progressions arithmétiques}
      \field{volume}{8}
      \field{year}{1986}
      \true{nocite}
      \field{pages}{159\bibrangedash 173}
      \range{pages}{15}
      \verb{mrnumber}
      \verb 928842
      \endverb
      \verb{urlraw}
      \verb http://www.numdam.org/item?id=AFST_1986-1987_5_8_2_159_0
      \endverb
      \verb{url}
      \verb http://www.numdam.org/item?id=AFST_1986-1987_5_8_2_159_0
      \endverb
    \endentry
    \entry{MR0137689}{article}{}
      \name{author}{2}{}{%
        {{hash=14b76cd756023555a4b6632dbee72243}{%
           family={Rosser},
           familyi={R\bibinitperiod},
           given={J.\bibnamedelimi Barkley},
           giveni={J\bibinitperiod\bibinitdelim B\bibinitperiod}}}%
        {{hash=197eb6ee8cad0f814cbfe07e59001ef3}{%
           family={Schoenfeld},
           familyi={S\bibinitperiod},
           given={Lowell},
           giveni={L\bibinitperiod}}}%
      }
      \strng{namehash}{748c2386a8972928ef3f5dc56357fa0e}
      \strng{fullhash}{748c2386a8972928ef3f5dc56357fa0e}
      \strng{bibnamehash}{748c2386a8972928ef3f5dc56357fa0e}
      \strng{authorbibnamehash}{748c2386a8972928ef3f5dc56357fa0e}
      \strng{authornamehash}{748c2386a8972928ef3f5dc56357fa0e}
      \strng{authorfullhash}{748c2386a8972928ef3f5dc56357fa0e}
      \field{sortinit}{R}
      \field{sortinithash}{b9c68a358aea118dfa887b6e902414a7}
      \field{labelprefix}{A}
      \field{labelnamesource}{author}
      \field{labeltitlesource}{title}
      \field{annotation}{Sprinkled among the many explicit inequalities for prime counting functions are a few comments directly related to comparative prime number theory, particularly on pages 72--73. The authors point out that it is unknown whether $n/\phi(n) \ge e^{C_0}\log\log n$ infinitely often, where $C_0$ is Euler's constant. They also speculate as to the existence of counterexamples to the inequalities \begin{align*} \log\log x + B &< \sum_{p\le x} \frac1p < \log\log x + B + \frac2{\sqrt x\log x} \\ \log x + E &< \sum_{p\le x} \frac{\log p}p < \log x + E + \frac{2.06123}{\sqrt x} \\ e^{C_0}\log x &< \prod_{p\le x} \frac p{p-1} < e^{C_0} \bigg( \log x + \frac2{\sqrt x} \bigg), \end{align*} where the constant $B$ is the unique real number with the property that the differences among the expressions on the first line tends to~$0$, and similarly with $E$ and the second line.}
      \field{issn}{0019-2082}
      \field{journaltitle}{Illinois J. Math.}
      \field{title}{Approximate formulas for some functions of prime numbers}
      \field{volume}{6}
      \field{year}{1962}
      \true{nocite}
      \field{pages}{64\bibrangedash 94}
      \range{pages}{31}
      \verb{mrnumber}
      \verb 0137689
      \endverb
    \endentry
    \entry{MR1329368}{article}{}
      \name{author}{2}{}{%
        {{hash=9073d8f33230938a75c73529bcd5e1ef}{%
           family={Rubinstein},
           familyi={R\bibinitperiod},
           given={M.},
           giveni={M\bibinitperiod}}}%
        {{hash=ea78b04077f0596f70f50f0559b2a1f2}{%
           family={Sarnak},
           familyi={S\bibinitperiod},
           given={P.},
           giveni={P\bibinitperiod}}}%
      }
      \strng{namehash}{78bb8618e127f18952100dffae455403}
      \strng{fullhash}{78bb8618e127f18952100dffae455403}
      \strng{bibnamehash}{78bb8618e127f18952100dffae455403}
      \strng{authorbibnamehash}{78bb8618e127f18952100dffae455403}
      \strng{authornamehash}{78bb8618e127f18952100dffae455403}
      \strng{authorfullhash}{78bb8618e127f18952100dffae455403}
      \field{sortinit}{R}
      \field{sortinithash}{b9c68a358aea118dfa887b6e902414a7}
      \field{labelprefix}{A}
      \field{labelnamesource}{author}
      \field{labeltitlesource}{title}
      \field{annotation}{This is the article that really placed in central roles the logarithmic limiting distributions and logarithmic densities of prime number races.\par Assuming GRH: the authors show that \begin{align*} E_{q;a_1,\ldots,a_r}(x) = \frac{\log x}{\sqrt x} \big( \phi(q) \pi(x;q,a_1) - \pi(x), \ldots, \phi(q) \pi(x;q,a_r) - \pi(x) \big) \end{align*} has a limiting logarithmic distribution $\mu_{q;a_1,\ldots,a_r}$ on $\mathbb R^r$. They give an exponential upper bound for the ``tail'' of $\mu_{q;a_1,\ldots,a_r}$ (that is, the mass assigned to the exterior of a large ball), as well as a doubly exponential lower bound for the portion of that tail lying in certain specific orthants. They note the analogous results for the race between $\pi(x)$ and $\Li(x)$, as well as for the race between $\pi(x;q,N)$ and $\pi(x;q,R)$; in these two-way races, it follows that $\underline\delta(\pi,\Li)$, $\underline\delta(\Li,\pi)$, $\underline\delta_{q;R,N}$, and $\underline\delta_{q;N,R}$ are strictly positive.\par Assuming GRH and LI: they give the formula for the Fourier transform of $\mu_{q;a_1,\ldots,a_r}$. From it they deduce that the densities $\delta_{q;a_1,\ldots,a_r}$ exist and are strictly positive. They characterize the races (all with $r\le 3$) for which $\mu_{q;a_1,\ldots ,a_r}$ is symmetric under all permutations of the coordinates. They show that $\delta_{q;a_1,\ldots ,a_r}$ tends to $1/r!$ as $q$ tends to infinity, and establish a central limit theorem for $E_{q; N,R}(x)/\sqrt{\log q}$. They also compute $\delta(\Li,\pi)$, and $\delta_{q;N,R}$ for $q\in\{3,4,5,7,11,13\}$, to several decimal places.\par This article cites~\cites{1914.Littlewood,MR1573970,MR0004255,MR0083001,MR0108470,MR0146156,MR0476616,MR866118,1938.Wintner,1932.Ingham,MR0506080,MR1159354}.}
      \field{issn}{1058-6458}
      \field{journaltitle}{Experiment. Math.}
      \field{number}{3}
      \field{title}{Chebyshev's bias}
      \field{volume}{3}
      \field{year}{1994}
      \true{nocite}
      \field{pages}{173\bibrangedash 197}
      \range{pages}{25}
      \verb{mrnumber}
      \verb 1329368
      \endverb
      \verb{urlraw}
      \verb http://projecteuclid.org/euclid.em/1048515870
      \endverb
      \verb{url}
      \verb http://projecteuclid.org/euclid.em/1048515870
      \endverb
    \endentry
    \entry{MR1908877}{article}{}
      \name{author}{1}{}{%
        {{hash=b8f6e7b5ec0a55f3a2fbfd17aab7dfa9}{%
           family={Ruzsa},
           familyi={R\bibinitperiod},
           given={Imre\bibnamedelima Z.},
           giveni={I\bibinitperiod\bibinitdelim Z\bibinitperiod}}}%
      }
      \strng{namehash}{b8f6e7b5ec0a55f3a2fbfd17aab7dfa9}
      \strng{fullhash}{b8f6e7b5ec0a55f3a2fbfd17aab7dfa9}
      \strng{bibnamehash}{b8f6e7b5ec0a55f3a2fbfd17aab7dfa9}
      \strng{authorbibnamehash}{b8f6e7b5ec0a55f3a2fbfd17aab7dfa9}
      \strng{authornamehash}{b8f6e7b5ec0a55f3a2fbfd17aab7dfa9}
      \strng{authorfullhash}{b8f6e7b5ec0a55f3a2fbfd17aab7dfa9}
      \field{sortinit}{R}
      \field{sortinithash}{b9c68a358aea118dfa887b6e902414a7}
      \field{labelprefix}{A}
      \field{labelnamesource}{author}
      \field{labeltitlesource}{title}
      \field{annotation}{The author proves that the number of pairs of consecutive primes up to~$x$ that are both congruent to $1\mod4$ is $\gg x\log\log x/\log^2x$, improving a result of Shiu. A generalization holds where the single residue class $1\mod4$ is replaced by an arbitrary set of reduced residue classes modulo~$q$ of size $\phi(q)/2$. The proof uses Maier's method. \par This article cites \cites{MR0466043,MR1297543}.}
      \field{issn}{0019-3577}
      \field{journaltitle}{Indag. Math. (N.S.)}
      \field{number}{4}
      \field{title}{Consecutive primes modulo 4}
      \field{volume}{12}
      \field{year}{2001}
      \true{nocite}
      \field{pages}{489\bibrangedash 503}
      \range{pages}{15}
      \verb{doi}
      \verb 10.1016/S0019-3577(01)80038-0
      \endverb
      \verb{mrnumber}
      \verb 1908877
      \endverb
    \endentry
    \entry{MR0207632}{article}{}
      \name{author}{1}{}{%
        {{hash=56c8a95f3c694f8332b16213f065326f}{%
           family={Ryan},
           familyi={R\bibinitperiod},
           given={J.\bibnamedelimi T.},
           giveni={J\bibinitperiod\bibinitdelim T\bibinitperiod}}}%
      }
      \strng{namehash}{56c8a95f3c694f8332b16213f065326f}
      \strng{fullhash}{56c8a95f3c694f8332b16213f065326f}
      \strng{bibnamehash}{56c8a95f3c694f8332b16213f065326f}
      \strng{authorbibnamehash}{56c8a95f3c694f8332b16213f065326f}
      \strng{authornamehash}{56c8a95f3c694f8332b16213f065326f}
      \strng{authorfullhash}{56c8a95f3c694f8332b16213f065326f}
      \field{sortinit}{R}
      \field{sortinithash}{b9c68a358aea118dfa887b6e902414a7}
      \field{labelprefix}{A}
      \field{labelnamesource}{author}
      \field{labeltitlesource}{title}
      \field{annotation}{Let $\pi_a(x;m,b) = \#\{ n\le x\colon n\equiv b\mod m,\, \Omega(n)=a\}$. Via combinatorial reasoning, the author conjectures that $\pi_a(x;m,b)/\pi_a(x;1,1)$ is asymptotically \begin{align*} \frac 1{\phi(m)}-(-1)^{a+1} \bigg( 1-\frac1{c_m} \bigg) x^{-1/2} &\quad\text{if $b$ is a quadratic residue,} \\ \frac 1{\phi(m)}+\frac{(-1)^{a+1}}{c_m}x^{-1/2} &\quad\text{if $b$ is a quadratic nonresidue.} \end{align*} The author provides some numerical evidence supporting the conjecture.\par This article cites~\cite{1891.Phragmen, MR1511360, MR0068565, MR0108470}.}
      \field{issn}{0002-9890}
      \field{journaltitle}{Amer. Math. Monthly}
      \field{number}{1}
      \field{title}{One more ``many-more'' assertion}
      \field{volume}{74}
      \field{year}{1967}
      \true{nocite}
      \field{pages}{19\bibrangedash 24}
      \range{pages}{6}
      \verb{doi}
      \verb 10.2307/2314046
      \endverb
      \verb{mrnumber}
      \verb 0207632
      \endverb
    \endentry
    \entry{MR280447}{article}{}
      \name{author}{1}{}{%
        {{hash=7bb09541026c884b0f336e16312810b6}{%
           family={Saffari},
           familyi={S\bibinitperiod},
           given={Bahman},
           giveni={B\bibinitperiod}}}%
      }
      \list{language}{1}{%
        {French}%
      }
      \strng{namehash}{7bb09541026c884b0f336e16312810b6}
      \strng{fullhash}{7bb09541026c884b0f336e16312810b6}
      \strng{bibnamehash}{7bb09541026c884b0f336e16312810b6}
      \strng{authorbibnamehash}{7bb09541026c884b0f336e16312810b6}
      \strng{authornamehash}{7bb09541026c884b0f336e16312810b6}
      \strng{authorfullhash}{7bb09541026c884b0f336e16312810b6}
      \field{sortinit}{S}
      \field{sortinithash}{c319cff79d99c853d775f88277d4e45f}
      \field{labelprefix}{A}
      \field{labelnamesource}{author}
      \field{labeltitlesource}{title}
      \field{annotation}{The author investigates the connection between Mertens conjecture and a finite version of LI, using a method similar to that of Ingham~\cite{MR0006202}. Let $\gamma_1<\gamma_2<\cdots$ denote the positive ordinates of the nontrivial zeros of $\zeta(s)$. Let $P(h)$ be the statement that there are no nontrivial linear relations $\sum_{k=1}^{h} a_{k} \gamma_{k} \ne 0$ with $\sum_{k=1}^h |a_h| \le h$, a finite and computationally verifiable assertion. The author shows that $P(28000)$ would imply $\lim \sup _{x \rightarrow \infty}|M(x)| / \sqrt{x}>1.179$, which disproves Mertens conjecture. \par This article cites~\cite{MR0330069,1897.Mertens,1897.vonSterneck,1901.vonSterneck,1913.vonSterneck,MR155787,MR0006202,1905.Stieltjes}.}
      \field{issn}{0151-0509}
      \field{journaltitle}{C. R. Acad. Sci. Paris Sér. A-B}
      \field{title}{Sur la fausseté de la conjecture de {M}ertens. ({W}ith discussion.)}
      \field{volume}{271}
      \field{year}{1970}
      \true{nocite}
      \field{pages}{A1097\bibrangedash A1101}
      \range{pages}{-1}
      \verb{mrnumber}
      \verb 280447
      \endverb
    \endentry
    \entry{MR1193094}{article}{}
      \name{author}{1}{}{%
        {{hash=c5aa7061ae0ee03e7795baf500605b3a}{%
           family={Sankaranarayanan},
           familyi={S\bibinitperiod},
           given={A.},
           giveni={A\bibinitperiod}}}%
      }
      \strng{namehash}{c5aa7061ae0ee03e7795baf500605b3a}
      \strng{fullhash}{c5aa7061ae0ee03e7795baf500605b3a}
      \strng{bibnamehash}{c5aa7061ae0ee03e7795baf500605b3a}
      \strng{authorbibnamehash}{c5aa7061ae0ee03e7795baf500605b3a}
      \strng{authornamehash}{c5aa7061ae0ee03e7795baf500605b3a}
      \strng{authorfullhash}{c5aa7061ae0ee03e7795baf500605b3a}
      \field{sortinit}{S}
      \field{sortinithash}{c319cff79d99c853d775f88277d4e45f}
      \field{labelprefix}{A}
      \field{labelnamesource}{author}
      \field{labeltitlesource}{title}
      \field{annotation}{Using a method of Kaczorowski~\cite{MR777013}, the author shows that the number of sign changes of $\Delta^{Q_2}(x)$ for $2\le x\le T$ is $\gg \log T$ with an effective constant. The author claims the analogous result for $\Delta^{Q_k}$ (with the implicit constant depending on~$k$) for all $2 \le k\le 10^8$. \par This article cites~\cites{MR945273,MR891057,MR1029966,MR0973808A,MR0973808B, MR777013}.}
      \field{issn}{0003-889X}
      \field{journaltitle}{Arch. Math. (Basel)}
      \field{number}{1}
      \field{title}{On the sign changes in the remainder term of an asymptotic formula for the number of squarefree numbers}
      \field{volume}{60}
      \field{year}{1993}
      \true{nocite}
      \field{pages}{51\bibrangedash 57}
      \range{pages}{7}
      \verb{doi}
      \verb 10.1007/BF01194239
      \endverb
      \verb{mrnumber}
      \verb 1193094
      \endverb
    \endentry
    \entry{MR2684372}{article}{}
      \name{author}{2}{}{%
        {{hash=99298fc0cd61b25871f9b3ad564fa78b}{%
           family={Saouter},
           familyi={S\bibinitperiod},
           given={Yannick},
           giveni={Y\bibinitperiod}}}%
        {{hash=447f22d7de07f4cbd918f6b2abb97288}{%
           family={Demichel},
           familyi={D\bibinitperiod},
           given={Patrick},
           giveni={P\bibinitperiod}}}%
      }
      \strng{namehash}{74715dbf8807c6d6381a8af29072c8b3}
      \strng{fullhash}{74715dbf8807c6d6381a8af29072c8b3}
      \strng{bibnamehash}{74715dbf8807c6d6381a8af29072c8b3}
      \strng{authorbibnamehash}{74715dbf8807c6d6381a8af29072c8b3}
      \strng{authornamehash}{74715dbf8807c6d6381a8af29072c8b3}
      \strng{authorfullhash}{74715dbf8807c6d6381a8af29072c8b3}
      \field{sortinit}{S}
      \field{sortinithash}{c319cff79d99c853d775f88277d4e45f}
      \field{labelprefix}{A}
      \field{labelnamesource}{author}
      \field{labeltitlesource}{title}
      \field{annotation}{The authors improve the error term for the function $I(\omega, \eta)$ in Lehman's theorem~\cite{MR0202686}, which allows them to conclude that there are more than $6.09 \times 10^{150}$ successive integers in the vicinity of $\exp(727.951335792) \approx 1.397166707819\times10^{316}$ such that $\pi(x)>\li(x)$. They show that $\pi(x)-\li(x)>9.1472 \times 10^{149}$ for some $x \in [1.39715131\times10^{316}, 1.39718211\times10^{316}]$; assuming RH, they show that $\pi(x)-\li(x)>1.7503\times 10^{148}$ for some~$x$ in the smaller interval $[1.3971619476\times10^{316},1.3971714624\times10^{316}]$.\par This article cites \cite{MR0137689, MR0202686, MR866118, MR552467, MR1752093, MR2652902}.}
      \field{issn}{0025-5718}
      \field{journaltitle}{Math. Comp.}
      \field{number}{272}
      \field{title}{A sharp region where {$\pi (x)-\li (x)$} is positive}
      \field{volume}{79}
      \field{year}{2010}
      \true{nocite}
      \field{pages}{2395\bibrangedash 2405}
      \range{pages}{11}
      \verb{doi}
      \verb 10.1090/S0025-5718-10-02351-3
      \endverb
      \verb{mrnumber}
      \verb 2684372
      \endverb
    \endentry
    \entry{MR3120597}{article}{}
      \name{author}{2}{}{%
        {{hash=99298fc0cd61b25871f9b3ad564fa78b}{%
           family={Saouter},
           familyi={S\bibinitperiod},
           given={Yannick},
           giveni={Y\bibinitperiod}}}%
        {{hash=c83d046dcdb536f5163e46eb8fe07827}{%
           family={Riele},
           familyi={R\bibinitperiod},
           given={Herman},
           giveni={H\bibinitperiod},
           prefix={te},
           prefixi={t\bibinitperiod}}}%
      }
      \strng{namehash}{c854407c034da79afefc129f9cbec9ec}
      \strng{fullhash}{c854407c034da79afefc129f9cbec9ec}
      \strng{bibnamehash}{c854407c034da79afefc129f9cbec9ec}
      \strng{authorbibnamehash}{c854407c034da79afefc129f9cbec9ec}
      \strng{authornamehash}{c854407c034da79afefc129f9cbec9ec}
      \strng{authorfullhash}{c854407c034da79afefc129f9cbec9ec}
      \field{sortinit}{S}
      \field{sortinithash}{c319cff79d99c853d775f88277d4e45f}
      \field{labelprefix}{A}
      \field{labelnamesource}{author}
      \field{labeltitlesource}{title}
      \field{annotation}{The authors refine Pintz's effective disproof~\cite{MR891440} of the Mertens conjecture and show that $|M(x)|>1.0088 \sqrt{x}$ for some $x<\exp(1.004 \times 10^{33})$. The authors also discuss possibilities for obtaining smaller counterexamples. \par This article cites~\cite{1897.Mertens,MR2282922,MR783538,MR891440,MR2684372,MR2340242}.}
      \field{issn}{0025-5718}
      \field{journaltitle}{Math. Comp.}
      \field{number}{285}
      \field{title}{Improved results on the {M}ertens conjecture}
      \field{volume}{83}
      \field{year}{2014}
      \true{nocite}
      \field{pages}{421\bibrangedash 433}
      \range{pages}{13}
      \verb{doi}
      \verb 10.1090/S0025-5718-2013-02716-0
      \endverb
      \verb{mrnumber}
      \verb 3120597
      \endverb
      \verb{urlraw}
      \verb https://doi.org/10.1090/S0025-5718-2013-02716-0
      \endverb
      \verb{url}
      \verb https://doi.org/10.1090/S0025-5718-2013-02716-0
      \endverb
    \endentry
    \entry{MR3356033}{article}{}
      \name{author}{3}{}{%
        {{hash=99298fc0cd61b25871f9b3ad564fa78b}{%
           family={Saouter},
           familyi={S\bibinitperiod},
           given={Yannick},
           giveni={Y\bibinitperiod}}}%
        {{hash=427a3ba3d2c76d68bc8c1bfc231edce0}{%
           family={Trudgian},
           familyi={T\bibinitperiod},
           given={Timothy},
           giveni={T\bibinitperiod}}}%
        {{hash=447f22d7de07f4cbd918f6b2abb97288}{%
           family={Demichel},
           familyi={D\bibinitperiod},
           given={Patrick},
           giveni={P\bibinitperiod}}}%
      }
      \strng{namehash}{fcb20b42ca212ff9c12bc49fe9926ddf}
      \strng{fullhash}{fcb20b42ca212ff9c12bc49fe9926ddf}
      \strng{bibnamehash}{fcb20b42ca212ff9c12bc49fe9926ddf}
      \strng{authorbibnamehash}{fcb20b42ca212ff9c12bc49fe9926ddf}
      \strng{authornamehash}{fcb20b42ca212ff9c12bc49fe9926ddf}
      \strng{authorfullhash}{fcb20b42ca212ff9c12bc49fe9926ddf}
      \field{sortinit}{S}
      \field{sortinithash}{c319cff79d99c853d775f88277d4e45f}
      \field{labelprefix}{A}
      \field{labelnamesource}{author}
      \field{labeltitlesource}{title}
      \field{annotation}{The authors incorporate both theoretical and computational improvements to show that there are more than $7.17\times10^{152}$ consecutive integers in the interval $[1.397165243588\times10^{316}, 1.397167149324\times10^{316}]$ for which $\pi(x)>\li(x)$.\par This article cites~\cites{MR2420864,MR0202686,1914.Littlewood,MR2684372}.}
      \field{issn}{0025-5718}
      \field{journaltitle}{Math. Comp.}
      \field{number}{295}
      \field{title}{A still sharper region where {$\pi (x)-\li (x)$} is positive}
      \field{volume}{84}
      \field{year}{2015}
      \true{nocite}
      \field{pages}{2433\bibrangedash 2446}
      \range{pages}{14}
      \verb{doi}
      \verb 10.1090/S0025-5718-2015-02930-5
      \endverb
      \verb{mrnumber}
      \verb 3356033
      \endverb
    \endentry
    \entry{2007.Sarnak}{online}{}
      \name{author}{1}{}{%
        {{hash=ea78b04077f0596f70f50f0559b2a1f2}{%
           family={Sarnak},
           familyi={S\bibinitperiod},
           given={P.},
           giveni={P\bibinitperiod}}}%
      }
      \strng{namehash}{ea78b04077f0596f70f50f0559b2a1f2}
      \strng{fullhash}{ea78b04077f0596f70f50f0559b2a1f2}
      \strng{bibnamehash}{ea78b04077f0596f70f50f0559b2a1f2}
      \strng{authorbibnamehash}{ea78b04077f0596f70f50f0559b2a1f2}
      \strng{authornamehash}{ea78b04077f0596f70f50f0559b2a1f2}
      \strng{authorfullhash}{ea78b04077f0596f70f50f0559b2a1f2}
      \field{sortinit}{S}
      \field{sortinithash}{c319cff79d99c853d775f88277d4e45f}
      \field{labelprefix}{A}
      \field{labelnamesource}{author}
      \field{labeltitlesource}{title}
      \field{annotation}{Let $\lambda(p)=\tau(p)/p^{11/2}$ denote the normalized $p$th Fourier coefficient of the Ramanujan function $\Delta(z)$, a holomorphic cusp form of weight~$12$. Under various assumptions (functional equation, RH, LI) for the symmetric power $L$-functions associated to $\Delta(z)$, the author proves the existence of a limiting logarithmic distribution for the cumulative sum $x^{-1/2}\log x \cdot \sum_{p\le x}\lambda(p)$. This distribution has mean~$1$ but infinite variance, so that even though there is a bias towards being positive, the set of such $x$ has logarithmic density $\frac12$. The author provides a similar analysis for $x^{-1/2}\log x \cdot \sum_{p\le x}a_E(p)/\sqrt p$, where $a_E(p)$ are the coefficients of the weight-$2$ cusp form attached to an elliptic curve $E/\mathbb{Q}$; here the mean is $1-2r(E)$ where $r(E)$ is the rank of~$E$. Here the variance is (conjecturally) finite, and so the logarithmic density of the set of $x$ for which $\sum_{p\le x}a_E(p)/\sqrt p > 0$ is strictly between~$\frac12$ and~$1$ when $r(E)=0$, but strictly between~$0$ and~$\frac12$ when $r(E)\ge1$. The author makes analogous remarks about the symmetric powers of these elliptic curve $L$-functions, where the finiteness of the variance corresponds to whether $E$ has complex multiplication.\par This letter cites~\cites{MR1329368}.}
      \field{title}{Letter to {B}arry {M}azur on `{C}hebyshev's bias' for {$\tau (p)$}}
      \field{year}{2007}
      \true{nocite}
      \verb{urlraw}
      \verb http://web.math.princeton.edu/sarnak/MazurLtrMay08.PDF
      \endverb
      \verb{url}
      \verb http://web.math.princeton.edu/sarnak/MazurLtrMay08.PDF
      \endverb
    \endentry
    \entry{MR2040112}{article}{}
      \name{author}{1}{}{%
        {{hash=b772e92036ff60f0205b68346a3f61db}{%
           family={Schlage--Puchta},
           familyi={S\bibinithyphendelim \bibinithyphendelim P\bibinitperiod},
           given={J.-C.},
           giveni={J\bibinithyphendelim C\bibinitperiod}}}%
      }
      \strng{namehash}{b772e92036ff60f0205b68346a3f61db}
      \strng{fullhash}{b772e92036ff60f0205b68346a3f61db}
      \strng{bibnamehash}{b772e92036ff60f0205b68346a3f61db}
      \strng{authorbibnamehash}{b772e92036ff60f0205b68346a3f61db}
      \strng{authornamehash}{b772e92036ff60f0205b68346a3f61db}
      \strng{authorfullhash}{b772e92036ff60f0205b68346a3f61db}
      \field{extraname}{1}
      \field{sortinit}{S}
      \field{sortinithash}{c319cff79d99c853d775f88277d4e45f}
      \field{labelprefix}{A}
      \field{labelnamesource}{author}
      \field{labeltitlesource}{title}
      \field{annotation}{This article examines the race between $\pi(x;q,1)$ and $\pi(x;q,{\ne}1) = \max_{a\not\equiv1\mod q} \pi(x;q,a)$. Set $C = \exp\big(\max\{q,e^{1260}\}^{170}+e^{18c_q}\big)$. Assuming GRH, the author proves that $W(\pi(x;q,1), \pi(x;q,{\ne}1); T) > (\log T)/C-1$, and in particular that there exists $x<e^C$ such that $\pi(x;q,1) > \pi(x;q,{\ne}1)$. The analogous results are proved for the race between $\pi(x)$ and $\li(x)$, with $C=\exp(e^{16.7})$.\par This article cites~\cites{1936.Ingham, MR1251926, MR1310738, MR0202686, MR607089, MR630377, MR1761276, MR866118, MR1573970}.}
      \field{issn}{0236-5294}
      \field{journaltitle}{Acta Math. Hungar.}
      \field{number}{4}
      \field{title}{Sign changes of {$\pi(x,q,1)-\pi(x,q,a)$}}
      \field{volume}{102}
      \field{year}{2004}
      \true{nocite}
      \field{pages}{305\bibrangedash 320}
      \range{pages}{16}
      \verb{doi}
      \verb 10.1023/B:AMHU.0000024681.23784.d1
      \endverb
      \verb{mrnumber}
      \verb 2040112
      \endverb
      \verb{urlraw}
      \verb https://doi.org/10.1023/B:AMHU.0000024681.23784.d1
      \endverb
      \verb{url}
      \verb https://doi.org/10.1023/B:AMHU.0000024681.23784.d1
      \endverb
    \endentry
    \entry{MR3871592}{article}{}
      \name{author}{1}{}{%
        {{hash=b772e92036ff60f0205b68346a3f61db}{%
           family={Schlage--Puchta},
           familyi={S\bibinithyphendelim \bibinithyphendelim P\bibinitperiod},
           given={J.-C.},
           giveni={J\bibinithyphendelim C\bibinitperiod}}}%
      }
      \strng{namehash}{b772e92036ff60f0205b68346a3f61db}
      \strng{fullhash}{b772e92036ff60f0205b68346a3f61db}
      \strng{bibnamehash}{b772e92036ff60f0205b68346a3f61db}
      \strng{authorbibnamehash}{b772e92036ff60f0205b68346a3f61db}
      \strng{authornamehash}{b772e92036ff60f0205b68346a3f61db}
      \strng{authorfullhash}{b772e92036ff60f0205b68346a3f61db}
      \field{extraname}{2}
      \field{sortinit}{S}
      \field{sortinithash}{c319cff79d99c853d775f88277d4e45f}
      \field{labelprefix}{A}
      \field{labelnamesource}{author}
      \field{labeltitlesource}{title}
      \field{annotation}{Using the power-sum method and Pintz's technique of kernel functions, the author improves Pintz's result ~\cites{MR820251} on localized oscillations of $\Delta^\psi$. The main result is that if $0<\ep<\frac1e$, and $\rho_0=\sigma_0+i \gamma_0$ is a zero of $\zeta(s)$ with $\sigma_0 \geq \frac{1}{2}+\ep$ and $\gamma_0 > 8.3 ^{1/\ep}$, then for each~$T$ sufficiently large (depending explicitly on~$\gamma_0$ and~$\ep$), there exists $x \in [T, T^{1+\ep}]$ such that $|\Delta^\psi (x)| \gg {x^{\sigma_0}}/{\gamma_0^{1+\ep}}$.\par This article cites~\cites{MR820241,MR1575079, MR820251,MR749389,MR0043121}.}
      \field{issn}{0236-5294}
      \field{journaltitle}{Acta Math. Hungar.}
      \field{number}{2}
      \field{title}{Oscillations of the error term in the prime number theorem}
      \field{volume}{156}
      \field{year}{2018}
      \true{nocite}
      \field{pages}{303\bibrangedash 308}
      \range{pages}{6}
      \verb{doi}
      \verb 10.1007/s10474-018-0884-x
      \endverb
      \verb{mrnumber}
      \verb 3871592
      \endverb
      \verb{urlraw}
      \verb https://doi.org/10.1007/s10474-018-0884-x
      \endverb
      \verb{url}
      \verb https://doi.org/10.1007/s10474-018-0884-x
      \endverb
    \endentry
    \entry{MR1511206}{article}{}
      \name{author}{1}{}{%
        {{hash=2c1103066f165b2530a1d823ff4f4b10}{%
           family={Schmidt},
           familyi={S\bibinitperiod},
           given={Erhard},
           giveni={E\bibinitperiod}}}%
      }
      \list{language}{1}{%
        {German}%
      }
      \strng{namehash}{2c1103066f165b2530a1d823ff4f4b10}
      \strng{fullhash}{2c1103066f165b2530a1d823ff4f4b10}
      \strng{bibnamehash}{2c1103066f165b2530a1d823ff4f4b10}
      \strng{authorbibnamehash}{2c1103066f165b2530a1d823ff4f4b10}
      \strng{authornamehash}{2c1103066f165b2530a1d823ff4f4b10}
      \strng{authorfullhash}{2c1103066f165b2530a1d823ff4f4b10}
      \field{sortinit}{S}
      \field{sortinithash}{c319cff79d99c853d775f88277d4e45f}
      \field{labelprefix}{A}
      \field{labelnamesource}{author}
      \field{labeltitlesource}{title}
      \field{annotation}{The author shows that $\Delta^\Pi(x) = \Omega_\pm(x^{\Theta-\ep})$ for every~$\ep>0$, and that both $E^\Pi(x) > \frac1{29}$ and $E^\Pi(x) < -\frac1{29}$ occur infinitely often. He also shows the same results with~$\Pi(x)$ in the definitions of~$\Delta^\Pi$ and~$E^\Pi$ replaced by $\pi(x)+\frac12\Li(\sqrt x)$.}
      \field{issn}{0025-5831}
      \field{journaltitle}{Math. Ann.}
      \field{number}{2}
      \field{title}{Über die {A}nzahl der {P}rimzahlen unter gegebener {G}renze}
      \field{volume}{57}
      \field{year}{1903}
      \true{nocite}
      \field{pages}{195\bibrangedash 204}
      \range{pages}{10}
      \verb{doi}
      \verb 10.1007/BF01444344
      \endverb
      \verb{mrnumber}
      \verb 1511206
      \endverb
      \verb{urlraw}
      \verb https://doi.org/10.1007/BF01444344
      \endverb
      \verb{url}
      \verb https://doi.org/10.1007/BF01444344
      \endverb
    \endentry
    \entry{MR4449835}{article}{}
      \name{author}{1}{}{%
        {{hash=b1b4d664f6bf440515ac846f7b16d530}{%
           family={Sedrati},
           familyi={S\bibinitperiod},
           given={Youssef},
           giveni={Y\bibinitperiod}}}%
      }
      \strng{namehash}{b1b4d664f6bf440515ac846f7b16d530}
      \strng{fullhash}{b1b4d664f6bf440515ac846f7b16d530}
      \strng{bibnamehash}{b1b4d664f6bf440515ac846f7b16d530}
      \strng{authorbibnamehash}{b1b4d664f6bf440515ac846f7b16d530}
      \strng{authornamehash}{b1b4d664f6bf440515ac846f7b16d530}
      \strng{authorfullhash}{b1b4d664f6bf440515ac846f7b16d530}
      \field{sortinit}{S}
      \field{sortinithash}{c319cff79d99c853d775f88277d4e45f}
      \field{labelprefix}{A}
      \field{labelnamesource}{author}
      \field{labeltitlesource}{title}
      \field{annotation}{The author extends results for classical prime number races~\cites{MR1806291,MR3063909} to the function field setting. Given a prime power~$q$ and a monic polynomial $m\in \mathbb{F}_q[T]$, the function $\pi_q(a,m, N)$ counts the number of irreducible polynomials $P\in\mathbb{F}_q[T]$ of degree~$N$ such that $P \equiv a \mod m$. Let $\delta_{m;a_1,\ldots, a_r}$ be the density of integers $X$ such that $\sum_{N=1}^X \pi_q(a_1,m,N)>\cdots > \sum_{N=1}^X \pi_q(a_r,m,N)$. Assuming LI, the author establishes an asymptotic formula for $\delta_{m;a_1,\ldots, a_r}$ as the degree of~$m$ tends to infinity. The author further shows that races with three or more competitors behave differently than two-way races in the sense that $\delta_{m;a_1,a_2} - \frac12 \ll q^{(-1/2+o(1))\deg m}$ while $|\delta_{m;a_1,\ldots, a_r} - \frac1{r!}| = \Omega(1/\deg m)$ when $r\ge3$. When $a_1, \ldots, a_r$ have bounded degree, the author provides a simple criteria for $\delta_{m;a_1,\ldots, a_r}$ to exhibit an extreme bias, and proves (still assuming LI) that when $r\ge3$, there are always unbiased races involving only quadratic residues or only quadratic residues when $\deg m$ is large enough. The author concludes by giving a few examples of races where LI is actually false.\par This article cites~\cites{MR4265352,MR2474313,MR3581815,MR2811555,MR2600420,MR4176773,MR3803705,MR1985941,MR1806291,MR3028758,MR2202918,MR3063909,1914.Littlewood,MR1329368}.}
      \field{issn}{0025-5793}
      \field{journaltitle}{Mathematika}
      \field{number}{3}
      \field{title}{Inequities in the {S}hanks-{R}enyi prime number race over function fields}
      \field{volume}{68}
      \field{year}{2022}
      \true{nocite}
      \field{pages}{840\bibrangedash 895}
      \range{pages}{56}
      \verb{doi}
      \verb 10.1112/mtk.12150
      \endverb
      \verb{mrnumber}
      \verb 4449835
      \endverb
      \verb{urlraw}
      \verb https://doi.org/10.1112/mtk.12150
      \endverb
      \verb{url}
      \verb https://doi.org/10.1112/mtk.12150
      \endverb
    \endentry
    \entry{MR0108470}{article}{}
      \name{author}{1}{}{%
        {{hash=22b786017fa0c36179547815bccb7685}{%
           family={Shanks},
           familyi={S\bibinitperiod},
           given={D.},
           giveni={D\bibinitperiod}}}%
      }
      \strng{namehash}{22b786017fa0c36179547815bccb7685}
      \strng{fullhash}{22b786017fa0c36179547815bccb7685}
      \strng{bibnamehash}{22b786017fa0c36179547815bccb7685}
      \strng{authorbibnamehash}{22b786017fa0c36179547815bccb7685}
      \strng{authornamehash}{22b786017fa0c36179547815bccb7685}
      \strng{authorfullhash}{22b786017fa0c36179547815bccb7685}
      \field{sortinit}{S}
      \field{sortinithash}{c319cff79d99c853d775f88277d4e45f}
      \field{labelprefix}{A}
      \field{labelnamesource}{author}
      \field{labeltitlesource}{title}
      \field{annotation}{In this article, the author investigates Chebyshev's assertion that there are more primes of the form $4m-1$ than of the form $4m+1$. Define $\tau(n)=\pi(n;4,3,1)\sqrt{n}/\pi(n)$ (which is asymptotically equivalent to $\pi(n;4,3,1)\log n/\sqrt n$ but is easier to manipulate numerically). Upon computing $\pi(n;4,3,1)$ for values of $n$ up to 3 million, he analyzes the values $\tau(1000k)$ for $1\le k\le2000$, noting that their histogram is ``roughly normal with a mean of (nearly)~$1$''. The author conjectures that \[ \lim_{x\to\infty} \frac1x \sum_{n\le x} \tau(n) = \lim_{x\to\infty} \frac1x \sum_{n\le x} \frac{\pi(x;4,3,1)\sqrt{n}}{\pi(n)} = 1, \] and notes that weaker versions of the conjecture---namely, that the above limit holds under GRH for $L(s,\chi_{-4})$, or that the above limit either equals $1$ or fails to exist---are also open. (He can show, under GRH, that the mean value inside the limit is positive and bounded away from $0$ for sufficiently large~$x$.) \par Next, the author discusses the distribution of primes in the residue classes modulo $8$, $10$, and $12$. Both from examining the collected data and from combinatorial reasoning involving the multiplicative groups of those moduli, he concludes that the quadratic residues are the ones with a smaller number of primes (on average). For the specific modulus $4$, he outlines an argument, based on combinatorial reasoning with the quantities $\#\{ n\le x\colon \Omega(n)=a,\, n\equiv\pm1\mod4\}$, that shows that the mean value of $\tau(n)$ should be~$1$. Indeed, he remarks that the generalization of this mean value to the integers with $a$ prime factors (counted with multiplicity) predicts that it is the residue class $(-1)^a\mod4$ that should have more such integers; in other words, the bias switches according to the parity of the number of prime factors. A related remark is that there is a bias towards integers for which $(-1)^{\Omega(n)}\chi_{-4}(n)$ equals~$1$ over those for which it equals~$-1$. \par This article ends with a discussion of how similar arguments to those laid out in this article could be used to analyze the relationship between $\pi(x)$ and $\li(x)$. \par This article cites~\cites{1853.Chebyshev, 1891.Phragmen, MR1511360, MR1555148, MR1573970, MR0068565, MR0083001}.}
      \field{issn}{0891-6837}
      \field{journaltitle}{Math. Tables Aids Comput.}
      \field{title}{Quadratic residues and the distribution of primes}
      \field{volume}{13}
      \field{year}{1959}
      \true{nocite}
      \field{pages}{272\bibrangedash 284}
      \range{pages}{13}
      \verb{mrnumber}
      \verb 0108470
      \endverb
    \endentry
    \entry{MR0302590}{article}{}
      \name{author}{2}{}{%
        {{hash=41686ca8e3d6aab7a821e927fb151ebd}{%
           family={Shanks},
           familyi={S\bibinitperiod},
           given={Daniel},
           giveni={D\bibinitperiod}}}%
        {{hash=f6052b9f2b8df2f5c47462bdead6fa53}{%
           family={Lal},
           familyi={L\bibinitperiod},
           given={Mohan},
           giveni={M\bibinitperiod}}}%
      }
      \strng{namehash}{2b4db4c1ac7b9014f0cffbd9d5021787}
      \strng{fullhash}{2b4db4c1ac7b9014f0cffbd9d5021787}
      \strng{bibnamehash}{2b4db4c1ac7b9014f0cffbd9d5021787}
      \strng{authorbibnamehash}{2b4db4c1ac7b9014f0cffbd9d5021787}
      \strng{authornamehash}{2b4db4c1ac7b9014f0cffbd9d5021787}
      \strng{authorfullhash}{2b4db4c1ac7b9014f0cffbd9d5021787}
      \field{sortinit}{S}
      \field{sortinithash}{c319cff79d99c853d775f88277d4e45f}
      \field{labelprefix}{A}
      \field{labelnamesource}{author}
      \field{labeltitlesource}{title}
      \field{annotation}{Set $\alpha_a(p)=3$ if~$a$ is a cubic residue modulo~$p$ and $\alpha_a(p)=0$ otherwise, and define $k_a = \prod_{p\nmid a} (p-\alpha_a(p))/(p-1)$, a conditionally convergent constant relevant to the conjectured asymptotic formula for the number of prime values of $n^3+a$. The authors describe their computations of~$k_2$ and~$k_3$ accurate to six decimal places. Further heuristics invoke the race between those primes for which~$a$ is a cubic residue and those for which it is not, which is a Chebyshev density theorem race for the extension $\Q(\sqrt[3]a,e^{2\pi i/3})/\Q$.}
      \field{issn}{0025-5718}
      \field{journaltitle}{Math. Comp.}
      \field{title}{Bateman's constants reconsidered and the distribution of cubic residues}
      \field{volume}{26}
      \field{year}{1972}
      \true{nocite}
      \field{pages}{265\bibrangedash 285}
      \range{pages}{21}
      \verb{doi}
      \verb 10.2307/2004737
      \endverb
      \verb{mrnumber}
      \verb 0302590
      \endverb
      \verb{urlraw}
      \verb https://doi.org/10.2307/2004737
      \endverb
      \verb{url}
      \verb https://doi.org/10.2307/2004737
      \endverb
    \endentry
    \entry{web.Shchebetov}{online}{}
      \name{author}{1}{}{%
        {{hash=a61d2f8619a2d8b3611da9d61d4a659a}{%
           family={Shchebetov},
           familyi={S\bibinitperiod},
           given={A.},
           giveni={A\bibinitperiod}}}%
      }
      \strng{namehash}{a61d2f8619a2d8b3611da9d61d4a659a}
      \strng{fullhash}{a61d2f8619a2d8b3611da9d61d4a659a}
      \strng{bibnamehash}{a61d2f8619a2d8b3611da9d61d4a659a}
      \strng{authorbibnamehash}{a61d2f8619a2d8b3611da9d61d4a659a}
      \strng{authornamehash}{a61d2f8619a2d8b3611da9d61d4a659a}
      \strng{authorfullhash}{a61d2f8619a2d8b3611da9d61d4a659a}
      \field{sortinit}{S}
      \field{sortinithash}{c319cff79d99c853d775f88277d4e45f}
      \field{labelprefix}{A}
      \field{labelnamesource}{author}
      \field{labeltitlesource}{title}
      \field{annotation}{This website contains the downloadable application ``Chebyshev's bias visualizer'', which permits graphical exploration of prime number races with many parameters customizable by the user. The author has explored several races for larger values of~$x$ than previously and, in particular, discovered the first region where $\pi(x;12,5,1)<0$ (starting at $25,726,067,172,577$) and the first region where $\pi(x;12,7,1)<0$ (starting at $27,489,101,529,529$).}
      \field{title}{Chebyshev's bias visualizer}
      \field{year}{2021}
      \true{nocite}
      \verb{urlraw}
      \verb http://math101.guru/en/downloads-2/repository/
      \endverb
      \verb{url}
      \verb http://math101.guru/en/downloads-2/repository/
      \endverb
    \endentry
    \entry{MR1573970}{article}{}
      \name{author}{1}{}{%
        {{hash=ed5a88e0fd4d9ca35b105db5f61affef}{%
           family={Skewes},
           familyi={S\bibinitperiod},
           given={S.},
           giveni={S\bibinitperiod}}}%
      }
      \strng{namehash}{ed5a88e0fd4d9ca35b105db5f61affef}
      \strng{fullhash}{ed5a88e0fd4d9ca35b105db5f61affef}
      \strng{bibnamehash}{ed5a88e0fd4d9ca35b105db5f61affef}
      \strng{authorbibnamehash}{ed5a88e0fd4d9ca35b105db5f61affef}
      \strng{authornamehash}{ed5a88e0fd4d9ca35b105db5f61affef}
      \strng{authorfullhash}{ed5a88e0fd4d9ca35b105db5f61affef}
      \field{extraname}{1}
      \field{sortinit}{S}
      \field{sortinithash}{c319cff79d99c853d775f88277d4e45f}
      \field{labelprefix}{A}
      \field{labelnamesource}{author}
      \field{labeltitlesource}{title}
      \field{annotation}{This article shows, assuming RH, that $\pi(x)>\li(x)$ for some $x<10^{10^{10^{34}}}$. Littlewood~\cites{1914.Littlewood,MR1555148} proved the existence of such an $x$ by considering the function $F(z) = \sum_{\gamma>0} {e^{-\gamma (\xi+i\eta)}}/{\gamma}$ for $0\le\xi\le1$ and $\eta\ge1$, which is relevant since the explicit formula yields $-2\Im F(i\log x) = E^\psi(x) +O(1)$. Using the Dirichlet box principle, Littlewood showed that $\Im F(\xi+i\eta)$ has large values (on the order of $\log \log \eta$) of either prescribed sign, with $\xi$ tending to $0$, and then used a modified form of the Phragmén--Lindelöf principle to show that an equally large value of $-\Im F(i\eta)$ must be attained.\par The Phragmén--Lindelöf principle, that the maximum of an analytic function defined on a semi-infinite strip (with suitable growth conditions) must occur on the boundary of the strip, is only an existence result; the author strengthens the result to give quantitative bounds on when an approximation to an interior value is attained on the boundary of the strip. In this way he is able to make Littlewood's result explicit (although the details are not included), which was not clearly possible beforehand.\par (Estimates of the smallest $x$ such that $\pi(x)>\li(x)$ have since been called ``Skewes numbers''. For comparison, the best estimates today are about $1.4\times 10^{316}$.)\par This article cites~\cites{1932.Ingham}.}
      \field{journaltitle}{J. London Math. Soc.}
      \field{number}{4}
      \field{title}{On the {D}ifference {$\pi(x)-{\rm li}\,(x)$} ({I})}
      \field{volume}{8}
      \field{year}{1933}
      \true{nocite}
      \field{pages}{277\bibrangedash 283}
      \range{pages}{7}
      \verb{doi}
      \verb 10.1112/jlms/s1-8.4.277
      \endverb
      \verb{mrnumber}
      \verb 1573970
      \endverb
      \verb{urlraw}
      \verb https://doi.org/10.1112/jlms/s1-8.4.277
      \endverb
      \verb{url}
      \verb https://doi.org/10.1112/jlms/s1-8.4.277
      \endverb
    \endentry
    \entry{MR0067145}{article}{}
      \name{author}{1}{}{%
        {{hash=ed5a88e0fd4d9ca35b105db5f61affef}{%
           family={Skewes},
           familyi={S\bibinitperiod},
           given={S.},
           giveni={S\bibinitperiod}}}%
      }
      \strng{namehash}{ed5a88e0fd4d9ca35b105db5f61affef}
      \strng{fullhash}{ed5a88e0fd4d9ca35b105db5f61affef}
      \strng{bibnamehash}{ed5a88e0fd4d9ca35b105db5f61affef}
      \strng{authorbibnamehash}{ed5a88e0fd4d9ca35b105db5f61affef}
      \strng{authornamehash}{ed5a88e0fd4d9ca35b105db5f61affef}
      \strng{authorfullhash}{ed5a88e0fd4d9ca35b105db5f61affef}
      \field{extraname}{2}
      \field{sortinit}{S}
      \field{sortinithash}{c319cff79d99c853d775f88277d4e45f}
      \field{labelprefix}{A}
      \field{labelnamesource}{author}
      \field{labeltitlesource}{title}
      \field{annotation}{This article provides an unconditional explicit estimate for a sign change of the difference $\Delta^\pi(x)$: if we define $X_1 = {e^{e^{e^{7.703}}}}$ and $X_2 = e^{4X_1^{30}} < e^{e^{e^{e^{7.705}}}} < 10^{10^{10^{10^3}}}$, then the author shows that there exists some $x<X_2$ such that $\pi(x)>\li(x)$. The author divides the proof into two cases, first when RH is ``nearly true'' and then the contrary case. More specifically, he defines a hypothesis (H) (the ``nearly true'' case) as follows: \emph{Every zero $\rho = \beta + i\gamma$ for which $|\gamma| < X_1^3$ satisfies $\beta - \frac{1}{2} \le X_1^{-3}\log^{-2}X_1$}. \par For the case where (H) holds, the author modifies Ingham's technique from~\cite{1936.Ingham}, which assumed RH but improved the estimation of $\Delta^{\psi_0}(x)$ by showing that zeros with $\gamma$ large relative to $x$ do not contribute meaningfully to the sum. Ultimately the author's argument boils down to estimation of the sum $\sum_{0<\gamma<500} \frac{\sin \gamma\omega}{\gamma}\left(1-\frac{\gamma}{500}\right)$; Dirichlet's box principle is used again, in conjunction with estimates of the values of the $269$ zeros of $\zeta(s)$ with $0<\gamma<500$. \par For the contrary case, which the author calls (NH), he remarks that it no longer suffices to work first with $\psi(x)$ and then pass to $\pi(x)$ with standard partial summation techniques. Instead, he works directly from the explicit formula for $\Delta^{\Pi_0}(x)$, introducing a smoothing factor to amplify the contribution from the hypothesized (H)-violating zero. Throughout, the author uses explicit estimates for sums over nontrivial zeros of $\zeta(s)$, such as $|N(T+h)-N(T)| < \frac{1}{2\pi} (h+1.77)\log T+ 8.7$ for $7.1<h<\frac T2$. \par This article cites~\cites{1914.Littlewood,1932.Ingham,MR1573970,1936.Ingham}.}
      \field{issn}{0024-6115}
      \field{journaltitle}{Proc. London Math. Soc. (3)}
      \field{title}{On the difference {$\pi(x)-{\rm li}\,x$}. {II}}
      \field{volume}{5}
      \field{year}{1955}
      \true{nocite}
      \field{pages}{48\bibrangedash 70}
      \range{pages}{23}
      \verb{doi}
      \verb 10.1112/plms/s3-5.1.48
      \endverb
      \verb{mrnumber}
      \verb 0067145
      \endverb
      \verb{urlraw}
      \verb https://doi.org/10.1112/plms/s3-5.1.48
      \endverb
      \verb{url}
      \verb https://doi.org/10.1112/plms/s3-5.1.48
      \endverb
    \endentry
    \entry{MR2753165}{book}{}
      \name{author}{1}{}{%
        {{hash=9670e498e8fb9912d411c47aed724ede}{%
           family={Sneed},
           familyi={S\bibinitperiod},
           given={J.\bibnamedelimi P.},
           giveni={J\bibinitperiod\bibinitdelim P\bibinitperiod}}}%
      }
      \list{publisher}{1}{%
        {ProQuest LLC, Ann Arbor, MI}%
      }
      \strng{namehash}{9670e498e8fb9912d411c47aed724ede}
      \strng{fullhash}{9670e498e8fb9912d411c47aed724ede}
      \strng{bibnamehash}{9670e498e8fb9912d411c47aed724ede}
      \strng{authorbibnamehash}{9670e498e8fb9912d411c47aed724ede}
      \strng{authornamehash}{9670e498e8fb9912d411c47aed724ede}
      \strng{authorfullhash}{9670e498e8fb9912d411c47aed724ede}
      \field{sortinit}{S}
      \field{sortinithash}{c319cff79d99c853d775f88277d4e45f}
      \field{labelprefix}{A}
      \field{labelnamesource}{author}
      \field{labeltitlesource}{title}
      \field{annotation}{In Chapter~2, the author verifies HC for moduli $q\le100$ and shows that all (two-way) prime number races for such moduli have infinitely many lead changes. In Chapter~3, the author explores semiprime races, which concern the functions $\pi_2(x;q,\ell)=\#\{p_1p_2\leq x\colon p_1p_2\equiv\ell\mod q\}$ and $E_2(x;q,\ell_1,\ell_2)=\big( \pi_2(x;q,\ell_1)-\pi_2(x;q,\ell_2) \big)\log x/\sqrt x\log\log x$. Assuming GRH and LI, he proves that $E_2(x;4,3,1)$ has mean $-\frac12$ (negative, unlike $E(x;4,3,1)$) and takes negative values with logarithmic density $\approx 0.894280$. \par This article cites~\cites{MR0447090,1853.Chebyshev,MR0332684,MR1862054,MR1985941,MR0202685,MR2202918,MR0006202,MR0176967,MR1511360,MR0146157,MR0083001,1914.Littlewood,MR1329368,MR0289452,MR0004255,MR429789}.}
      \field{isbn}{978-1124-04248-0}
      \field{note}{Thesis (Ph.D.)--University of Illinois at Urbana-Champaign}
      \field{pagetotal}{83}
      \field{title}{Prime and quasi-prime number races}
      \field{year}{2009}
      \true{nocite}
      \verb{mrnumber}
      \verb 2753165
      \endverb
      \verb{urlraw}
      \verb http://gateway.proquest.com/openurl?url_ver=Z39.88-2004&rft_val_fmt=info:ofi/fmt:kev:mtx:dissertation&res_dat=xri:pqdiss&rft_dat=xri:pqdiss:3411454
      \endverb
      \verb{url}
      \verb http://gateway.proquest.com/openurl?url_ver=Z39.88-2004&rft_val_fmt=info:ofi/fmt:kev:mtx:dissertation&res_dat=xri:pqdiss&rft_dat=xri:pqdiss:3411454
      \endverb
    \endentry
    \entry{MR0077579}{article}{}
      \name{author}{2}{}{%
        {{hash=7220d8e752844cc34f659367559952a5}{%
           family={Sós},
           familyi={S\bibinitperiod},
           given={Vera\bibnamedelima T.},
           giveni={V\bibinitperiod\bibinitdelim T\bibinitperiod}}}%
        {{hash=e68e101772b88cb1d48e688164bece92}{%
           family={Turán},
           familyi={T\bibinitperiod},
           given={P.},
           giveni={P\bibinitperiod}}}%
      }
      \list{language}{1}{%
        {English, with Russian summary}%
      }
      \strng{namehash}{94e15cc05318c39c603ade60525f4392}
      \strng{fullhash}{94e15cc05318c39c603ade60525f4392}
      \strng{bibnamehash}{94e15cc05318c39c603ade60525f4392}
      \strng{authorbibnamehash}{94e15cc05318c39c603ade60525f4392}
      \strng{authornamehash}{94e15cc05318c39c603ade60525f4392}
      \strng{authorfullhash}{94e15cc05318c39c603ade60525f4392}
      \field{sortinit}{S}
      \field{sortinithash}{c319cff79d99c853d775f88277d4e45f}
      \field{labelprefix}{A}
      \field{labelnamesource}{author}
      \field{labeltitlesource}{title}
      \field{annotation}{The authors give bounds for the valid constants $A$ and $B$ appearing in the first and second main theorems of the power-sum method. $A=2e$ is known to be valid in the first main theorem, and the authors show that $A=\frac4\pi$ is too small. They also show that $B=2e^{1+4/e}$ is valid in the second main theorem (improving the previous $B=24e^2$) while $B=1.321$ is not; they mention that the improvement in the upper bound leads to improved constants in a zero-density theorem for $\zeta(s)$ and in the maximal gap between primes. The authors also refer to a ``third main theorem'', in which the normalized lower bound for the power-sum is independent of the~$b_j$ rather than of the~$z_j$. \par This article cites~\cites{MR0060548}.}
      \field{issn}{0001-5954}
      \field{journaltitle}{Acta Math. Acad. Sci. Hungar.}
      \field{title}{On some new theorems in the theory of Diophantine approximations}
      \field{volume}{6}
      \field{year}{1955}
      \true{nocite}
      \field{pages}{241\bibrangedash 255}
      \range{pages}{15}
      \verb{doi}
      \verb 10.1007/BF02024389
      \endverb
      \verb{mrnumber}
      \verb 0077579
      \endverb
    \endentry
    \entry{MR228456}{article}{}
      \name{author}{1}{}{%
        {{hash=e653ad19c187960606202bfb90c8f746}{%
           family={Spira},
           familyi={S\bibinitperiod},
           given={Robert},
           giveni={R\bibinitperiod}}}%
      }
      \strng{namehash}{e653ad19c187960606202bfb90c8f746}
      \strng{fullhash}{e653ad19c187960606202bfb90c8f746}
      \strng{bibnamehash}{e653ad19c187960606202bfb90c8f746}
      \strng{authorbibnamehash}{e653ad19c187960606202bfb90c8f746}
      \strng{authornamehash}{e653ad19c187960606202bfb90c8f746}
      \strng{authorfullhash}{e653ad19c187960606202bfb90c8f746}
      \field{sortinit}{S}
      \field{sortinithash}{c319cff79d99c853d775f88277d4e45f}
      \field{labelprefix}{A}
      \field{labelnamesource}{author}
      \field{labeltitlesource}{title}
      \field{annotation}{Haselgrove~\cite{MR0104638} resolved {P}ólya's problem in the negative; in this article, the author provides numerical evidence that confirms Haselgrove's computations. The author does a similar computation for the Mertens function, showing that \[ \liminf_{x \to \infty} E^M(x) < -0.6027 \quad \text{and} \quad \limsup_{x \to \infty} E^M(x) > 0.5355. \] The calculations were performed on an IBM 7040 at the University of Tennessee Computing Center. \par This article cites~\cites{MR27305,MR0104638,MR0137689,MR0120198}.}
      \field{issn}{0025-5718}
      \field{journaltitle}{Math. Comp.}
      \field{title}{Zeros of sections of the zeta function. {II}}
      \field{volume}{22}
      \field{year}{1968}
      \true{nocite}
      \field{pages}{163\bibrangedash 173}
      \range{pages}{11}
      \verb{doi}
      \verb 10.2307/2004774
      \endverb
      \verb{mrnumber}
      \verb 228456
      \endverb
      \verb{urlraw}
      \verb https://doi.org/10.2307/2004774
      \endverb
      \verb{url}
      \verb https://doi.org/10.2307/2004774
      \endverb
    \endentry
    \entry{MR0300853}{article}{}
      \field{sortinit}{S}
      \field{sortinithash}{c319cff79d99c853d775f88277d4e45f}
      \field{labelprefix}{A}
      \field{labeltitlesource}{title}
      \field{annotation}{This article is a short biography of Stanisław Knapowski's life.}
      \field{issn}{0010-1354}
      \field{journaltitle}{Colloq. Math.}
      \field{title}{Stanisław {K}napowski (19. V. 1931--28. IX. 1967)}
      \field{volume}{23}
      \field{year}{1971}
      \true{nocite}
      \field{pages}{309\bibrangedash 310}
      \range{pages}{2}
      \verb{mrnumber}
      \verb 0300853
      \endverb
    \endentry
    \entry{MR0199161}{article}{}
      \name{author}{1}{}{%
        {{hash=32597e08070c5aa88e8fb07dd295112a}{%
           family={Stark},
           familyi={S\bibinitperiod},
           given={H.\bibnamedelimi M.},
           giveni={H\bibinitperiod\bibinitdelim M\bibinitperiod}}}%
      }
      \strng{namehash}{32597e08070c5aa88e8fb07dd295112a}
      \strng{fullhash}{32597e08070c5aa88e8fb07dd295112a}
      \strng{bibnamehash}{32597e08070c5aa88e8fb07dd295112a}
      \strng{authorbibnamehash}{32597e08070c5aa88e8fb07dd295112a}
      \strng{authornamehash}{32597e08070c5aa88e8fb07dd295112a}
      \strng{authorfullhash}{32597e08070c5aa88e8fb07dd295112a}
      \field{extraname}{1}
      \field{sortinit}{S}
      \field{sortinithash}{c319cff79d99c853d775f88277d4e45f}
      \field{labelprefix}{A}
      \field{labelnamesource}{author}
      \field{labeltitlesource}{title}
      \field{annotation}{Evelyn and Linfoot~\cite{MR1502996} showed that $\Delta^{Q_k}(x)=\Omega_{\pm}(x^{1/2k})$. Following Ingham's method~\cites{MR0006202}, the author makes their result effective by showing that \[ \liminf_{x \to \infty} \Delta^{Q_k}(x) /x^{1/2k} < -C_k \quad \text{and} \quad \limsup_{x \to \infty} \Delta^{Q_k}(x) /x^{1/2k} > C_k; \] here $C_k=2(1-\frac{\gamma_1}{\gamma_2}) \bigl| \zeta(\frac{\rho_1}k)/\rho_1\zeta'(\rho_1) \bigr|$, where $\rho_j=\frac12+i\gamma_j$ ($j=1,2$) are the first two zeros of $\zeta(s)$ above the real axis.\par This article cites~\cites{MR0006202,MR1502996}.}
      \field{issn}{0002-9939,1088-6826}
      \field{journaltitle}{Proc. Amer. Math. Soc.}
      \field{title}{On the asymptotic density of the {$k$}-free integers}
      \field{volume}{17}
      \field{year}{1966}
      \true{nocite}
      \field{pages}{1211\bibrangedash 1214}
      \range{pages}{4}
      \verb{doi}
      \verb 10.2307/2036123
      \endverb
      \verb{mrnumber}
      \verb 199161
      \endverb
      \verb{urlraw}
      \verb https://doi.org/10.2307/2036123
      \endverb
      \verb{url}
      \verb https://doi.org/10.2307/2036123
      \endverb
    \endentry
    \entry{MR0289452}{article}{}
      \name{author}{1}{}{%
        {{hash=32597e08070c5aa88e8fb07dd295112a}{%
           family={Stark},
           familyi={S\bibinitperiod},
           given={H.\bibnamedelimi M.},
           giveni={H\bibinitperiod\bibinitdelim M\bibinitperiod}}}%
      }
      \strng{namehash}{32597e08070c5aa88e8fb07dd295112a}
      \strng{fullhash}{32597e08070c5aa88e8fb07dd295112a}
      \strng{bibnamehash}{32597e08070c5aa88e8fb07dd295112a}
      \strng{authorbibnamehash}{32597e08070c5aa88e8fb07dd295112a}
      \strng{authornamehash}{32597e08070c5aa88e8fb07dd295112a}
      \strng{authorfullhash}{32597e08070c5aa88e8fb07dd295112a}
      \field{extraname}{2}
      \field{sortinit}{S}
      \field{sortinithash}{c319cff79d99c853d775f88277d4e45f}
      \field{labelprefix}{A}
      \field{labelnamesource}{author}
      \field{labeltitlesource}{title}
      \field{annotation}{This article gives a Tauberian theorem of Landau type (with a proof sketch and a reference to an unpublished manuscript of the author) and uses it to show that \begin{align*} \limsup_{x\to\infty} \bigl( & E^\pi(x;k,a) - E^\pi(x;q,b) \bigr) \ge c_q(b) - c_k(a) \\ &{} + \sum_{\substack{\chi\mod q \\ \chi\ne\chi_0}} \sum_{|\gamma|<T} \frac{\overline\chi(b)}\rho \bigg( 1 - \frac\gamma T \bigg) e^{(\rho-1/2)u} - \sum_{\substack{\chi\mod k \\ \chi\ne\chi_0}} \sum_{|\gamma|<T} \frac{\overline\chi(a)}\rho \bigg( 1 - \frac\gamma T \bigg) e^{(\rho-1/2)u} \end{align*} for any $T>0$ and any $u\in\R$, under the assumption of GRH$(0)$ for nonprincipal $L$-functions modulo~$k$ and~$q$. In particular, the $\limsup$ is positive if either $a$ is a nonsquare\mod k or $b$ is a square\mod q (these are the cases for which $c_q(b) - c_k(a) \ge 0$). Under the same assumption, the author further proves \begin{align*} \limsup_{x\to\infty} \bigl( E^\pi(x;k,a) & {}- E^\pi(x;q,b) \bigr) \ge c_q(b) - c_k(a) \\ &{} + \sum_{\substack{\chi\mod q \\ \chi\ne\chi_0}} \sum_{|\gamma|<T} \frac{\overline\chi(b)}\rho e^{(\rho-1/2)u} - \sum_{\substack{\chi\mod k \\ \chi\ne\chi_0}} \sum_{|\gamma|<T} \frac{\overline\chi(a)}\rho e^{(\rho-1/2)u} \end{align*} for any $T>0$ and any $u\ne0$. Finally, when $u<0$ the author obtains an exact formula for this right-hand side as $T\to\infty$, a version that implies that the $\limsup$ is infinite if $a\equiv1\mod k$ or $b\equiv1\mod q$ but not both. The author also uses this exact formula and some shrewd explicit computation to show for the first time that $\pi(x;5,4,2) = \Omega_+(\sqrt x/\log x)$.\par This article cites~\cites{MR0006202, MR0146156, MR0146157, MR0146158, MR0146159, MR0146160, MR0146161, MR0156826, MR0156827}.}
      \field{issn}{0065-1036}
      \field{journaltitle}{Acta Arith.}
      \field{title}{A problem in comparative prime number theory}
      \field{volume}{18}
      \field{year}{1971}
      \true{nocite}
      \field{pages}{311\bibrangedash 320}
      \range{pages}{10}
      \verb{doi}
      \verb 10.4064/aa-18-1-311-320
      \endverb
      \verb{mrnumber}
      \verb 0289452
      \endverb
      \verb{urlraw}
      \verb https://doi.org/10.4064/aa-18-1-311-320
      \endverb
      \verb{url}
      \verb https://doi.org/10.4064/aa-18-1-311-320
      \endverb
    \endentry
    \entry{MR0146153}{article}{}
      \name{author}{1}{}{%
        {{hash=36befda65f61ea63f097e5518d3ec577}{%
           family={Staś},
           familyi={S\bibinitperiod},
           given={W.},
           giveni={W\bibinitperiod}}}%
      }
      \list{language}{1}{%
        {German}%
      }
      \strng{namehash}{36befda65f61ea63f097e5518d3ec577}
      \strng{fullhash}{36befda65f61ea63f097e5518d3ec577}
      \strng{bibnamehash}{36befda65f61ea63f097e5518d3ec577}
      \strng{authorbibnamehash}{36befda65f61ea63f097e5518d3ec577}
      \strng{authornamehash}{36befda65f61ea63f097e5518d3ec577}
      \strng{authorfullhash}{36befda65f61ea63f097e5518d3ec577}
      \field{extraname}{1}
      \field{sortinit}{S}
      \field{sortinithash}{c319cff79d99c853d775f88277d4e45f}
      \field{labelprefix}{A}
      \field{labelnamesource}{author}
      \field{labeltitlesource}{title}
      \field{annotation}{Ingham \cite{1932.Ingham} showed that any precise information on the zero-free region of $\zeta(s)$ would lead to a correspondingly precise estimate on $\Delta^\psi(x)$. The author uses the power-sum method to obtain a partial converse of this theorem. As an illustration, the author shows that $\Delta^\psi(x) \ll x/ (\log x)^{1/10}$ would imply that $\zeta(s) \neq 0$ in the region $\sigma > 1 - \frac{1}{400}(\log t)t^{-20}$ when~$t$ is sufficiently large.\par This article cites~\cites{1932.Ingham,MR0060548}.}
      \field{issn}{0065-1036}
      \field{journaltitle}{Acta Arith.}
      \field{sortyear}{1960}
      \field{title}{Über die Umkehrung eines Satzes von Ingham}
      \field{volume}{6}
      \field{year}{1960/1961}
      \true{nocite}
      \field{pages}{435\bibrangedash 446}
      \range{pages}{12}
      \verb{doi}
      \verb 10.4064/aa-6-4-435-446
      \endverb
      \verb{mrnumber}
      \verb 0146153
      \endverb
    \endentry
    \entry{MR0177957}{article}{}
      \name{author}{1}{}{%
        {{hash=36befda65f61ea63f097e5518d3ec577}{%
           family={Staś},
           familyi={S\bibinitperiod},
           given={W.},
           giveni={W\bibinitperiod}}}%
      }
      \strng{namehash}{36befda65f61ea63f097e5518d3ec577}
      \strng{fullhash}{36befda65f61ea63f097e5518d3ec577}
      \strng{bibnamehash}{36befda65f61ea63f097e5518d3ec577}
      \strng{authorbibnamehash}{36befda65f61ea63f097e5518d3ec577}
      \strng{authornamehash}{36befda65f61ea63f097e5518d3ec577}
      \strng{authorfullhash}{36befda65f61ea63f097e5518d3ec577}
      \field{extraname}{2}
      \field{sortinit}{S}
      \field{sortinithash}{c319cff79d99c853d775f88277d4e45f}
      \field{labelprefix}{A}
      \field{labelnamesource}{author}
      \field{labeltitlesource}{title}
      \field{annotation}{Assuming RH, the author shows that \[ \max_{T^{1-\ep}\le x\le T} \bigg| \sum_{n=1}^\infty \frac{\mu(n)}n e^{-(x/n)^2} \bigg| \ge T^{-1/2-\ep} \] when~$T$ is sufficiently large in terms of~$\ep$. \par This article cites~\cite{MR0174538}.}
      \field{issn}{0065-1036}
      \field{journaltitle}{Acta Arith.}
      \field{sortyear}{1964}
      \field{title}{Some remarks on a series of Ramanujan}
      \field{volume}{10}
      \field{year}{1964/1965}
      \true{nocite}
      \field{pages}{359\bibrangedash 368}
      \range{pages}{10}
      \verb{doi}
      \verb 10.4064/aa-10-4-359-368
      \endverb
      \verb{mrnumber}
      \verb 0177957
      \endverb
    \endentry
    \entry{MR0429797}{article}{}
      \name{author}{2}{}{%
        {{hash=36befda65f61ea63f097e5518d3ec577}{%
           family={Staś},
           familyi={S\bibinitperiod},
           given={W.},
           giveni={W\bibinitperiod}}}%
        {{hash=b983eb20838b58a1908783452401fc09}{%
           family={Wiertelak},
           familyi={W\bibinitperiod},
           given={K.},
           giveni={K\bibinitperiod}}}%
      }
      \strng{namehash}{9af3bba0a834e4a5fc9ac880f2af2cce}
      \strng{fullhash}{9af3bba0a834e4a5fc9ac880f2af2cce}
      \strng{bibnamehash}{9af3bba0a834e4a5fc9ac880f2af2cce}
      \strng{authorbibnamehash}{9af3bba0a834e4a5fc9ac880f2af2cce}
      \strng{authornamehash}{9af3bba0a834e4a5fc9ac880f2af2cce}
      \strng{authorfullhash}{9af3bba0a834e4a5fc9ac880f2af2cce}
      \field{sortinit}{S}
      \field{sortinithash}{c319cff79d99c853d775f88277d4e45f}
      \field{labelprefix}{A}
      \field{labelnamesource}{author}
      \field{labeltitlesource}{title}
      \field{annotation}{Let $\mathfrak{K}_1,\mathfrak{K}_2$ be ideal classes\mod{\mathfrak{f}} in a number field. In this article, the authors use the second main theorem of the power-sum method to bound $\psi(x,\mathfrak{K}_1)-\psi(x,\mathfrak{K}_2)$ given a zero-free region for Hecke--Landau zeta functions of relevant Hecke characters~$\chi$, and establish a converse as well: Let~$\gamma_1$ be the supremum of numbers~$\gamma$ for which $\psi(x,\mathfrak{K}_1)-\psi(x,\mathfrak{K}_2) \ll xe^{-a(\log x)^\gamma}$ for some positive constant~$a$, and let~$\gamma_2$ be the infimum of numbers~$\gamma$ for which $\prod_{\chi(\mathfrak{K}_1)\ne\chi(\mathfrak{K}_2)}\zeta(s,\chi)$ does not vanish in the region $\sigma > 1-{b}/{(\log |t|)^\gamma}$ for some positive constant~$b$. The authors prove that $\gamma_1 = 1/(1+\gamma_2)$.\par This article cites~\cites{1932.Ingham,MR0087685,MR0146153,MR0060548}.}
      \field{issn}{0065-1036}
      \field{journaltitle}{Acta Arith.}
      \field{number}{2}
      \field{title}{Further applications of Turán's methods to the distribution of prime ideals in ideal classes\mod f}
      \field{volume}{31}
      \field{year}{1976}
      \true{nocite}
      \field{pages}{153\bibrangedash 165}
      \range{pages}{13}
      \verb{doi}
      \verb 10.4064/aa-31-2-153-165
      \endverb
      \verb{mrnumber}
      \verb 0429797
      \endverb
    \endentry
    \entry{MR817334}{article}{}
      \name{author}{1}{}{%
        {{hash=b39a346ae2b44071e01b12c1003e1246}{%
           family={Staś},
           familyi={S\bibinitperiod},
           given={Włodzimierz},
           giveni={W\bibinitperiod}}}%
      }
      \strng{namehash}{b39a346ae2b44071e01b12c1003e1246}
      \strng{fullhash}{b39a346ae2b44071e01b12c1003e1246}
      \strng{bibnamehash}{b39a346ae2b44071e01b12c1003e1246}
      \strng{authorbibnamehash}{b39a346ae2b44071e01b12c1003e1246}
      \strng{authornamehash}{b39a346ae2b44071e01b12c1003e1246}
      \strng{authorfullhash}{b39a346ae2b44071e01b12c1003e1246}
      \field{extraname}{3}
      \field{sortinit}{S}
      \field{sortinithash}{c319cff79d99c853d775f88277d4e45f}
      \field{labelprefix}{A}
      \field{labelnamesource}{author}
      \field{labeltitlesource}{title}
      \field{annotation}{In this article, the author proves an analogue of a result of Knapowski~\cite{MR0133309} for the error term of the prime ideal theorem. More precisely he shows that for a number field~$K$ whose Dedekind zeta function $\zeta_K(s)$ satisfies HC, and an arbitrary zero $\rho_0 = \beta_0 + i\gamma_0$ of $\zeta_K$, \[ \Delta^\psi(T,K) = \Omega_\pm\biggl( T^{\beta_0}\exp\biggl( -15\frac{\log T}{\sqrt{\log \log T}}\biggr)\biggr). \] The author also gives a lower bound on the number of sign changes of $\Delta^\psi(T,K)$.\par This article cites~\cites{MR0133309,MR0174538, MR585891,MR0060548,MR0132728}.}
      \field{issn}{0208-6573}
      \field{journaltitle}{Funct. Approx. Comment. Math.}
      \field{title}{On sign-changes in the remainder term of the prime ideal formula}
      \field{volume}{13}
      \field{year}{1982}
      \true{nocite}
      \field{pages}{159\bibrangedash 166}
      \range{pages}{8}
      \verb{mrnumber}
      \verb 817334
      \endverb
    \endentry
    \entry{MR1440155}{article}{}
      \name{author}{2}{}{%
        {{hash=a5c0de227f384c5812b85df50f0ac6e8}{%
           family={Stechkin},
           familyi={S\bibinitperiod},
           given={S.\bibnamedelimi B.},
           giveni={S\bibinitperiod\bibinitdelim B\bibinitperiod}}}%
        {{hash=548c78bd85b3dd0d14bb46a59f54e09e}{%
           family={Popov},
           familyi={P\bibinitperiod},
           given={A.\bibnamedelimi Yu.},
           giveni={A\bibinitperiod\bibinitdelim Y\bibinitperiod}}}%
      }
      \strng{namehash}{78455df493ce58458371e5d01ed8136b}
      \strng{fullhash}{78455df493ce58458371e5d01ed8136b}
      \strng{bibnamehash}{78455df493ce58458371e5d01ed8136b}
      \strng{authorbibnamehash}{78455df493ce58458371e5d01ed8136b}
      \strng{authornamehash}{78455df493ce58458371e5d01ed8136b}
      \strng{authorfullhash}{78455df493ce58458371e5d01ed8136b}
      \field{sortinit}{S}
      \field{sortinithash}{c319cff79d99c853d775f88277d4e45f}
      \field{labelprefix}{A}
      \field{labelnamesource}{author}
      \field{labeltitlesource}{title}
      \field{annotation}{The authors prove that when~$x$ sufficiently large, \[ \int_x^{2x} |\Delta^\psi(u)| \,du \ge \frac{x^{3/2}}{200} \quad\text{and}\quad \int_x^{2x} |\Delta^\pi(u)| \,du \ge \frac{x^{3/2}}{\log x}, \] and also that there exists a constant $A>1$ such that \[ \int_x^{Ax} \max\{\Delta^\psi(u)\} \,du \ge x^{3/2} \quad\text{and}\quad \int_x^{Ax} \min\{\Delta^\psi(u)\} \,du \le -x^{3/2}. \] Assuming RH, these latter inequalities hold with $A=212$. If RH is false, however, then the constant~$A$ (as well as the constant implied by ``sufficiently large'') is ineffective.\par This article cites~\cites{MR1511206, MR1544509}.}
      \field{issn}{0042-1316}
      \field{journaltitle}{Uspekhi Mat. Nauk}
      \field{number}{6(312)}
      \field{title}{Asymptotic distribution of prime numbers in the mean}
      \field{volume}{51}
      \field{year}{1996}
      \true{nocite}
      \field{pages}{21\bibrangedash 88}
      \range{pages}{68}
      \verb{doi}
      \verb 10.1070/RM1996v051n06ABEH003000
      \endverb
      \verb{mrnumber}
      \verb 1440155
      \endverb
      \verb{urlraw}
      \verb https://doi.org/10.1070/RM1996v051n06ABEH003000
      \endverb
      \verb{url}
      \verb https://doi.org/10.1070/RM1996v051n06ABEH003000
      \endverb
    \endentry
    \entry{MR0257003}{article}{}
      \name{author}{1}{}{%
        {{hash=56a0c36192f225fe8031e07a1c09cdd1}{%
           family={Steinig},
           familyi={S\bibinitperiod},
           given={J.},
           giveni={J\bibinitperiod}}}%
      }
      \strng{namehash}{56a0c36192f225fe8031e07a1c09cdd1}
      \strng{fullhash}{56a0c36192f225fe8031e07a1c09cdd1}
      \strng{bibnamehash}{56a0c36192f225fe8031e07a1c09cdd1}
      \strng{authorbibnamehash}{56a0c36192f225fe8031e07a1c09cdd1}
      \strng{authornamehash}{56a0c36192f225fe8031e07a1c09cdd1}
      \strng{authorfullhash}{56a0c36192f225fe8031e07a1c09cdd1}
      \field{sortinit}{S}
      \field{sortinithash}{c319cff79d99c853d775f88277d4e45f}
      \field{labelprefix}{A}
      \field{labelnamesource}{author}
      \field{labeltitlesource}{title}
      \field{annotation}{The author fills a gap in the proof of a refinement of Landau's theorem by Pólya and applies it to summatory functions of arithmetic functions whose associated Dirichlet series have functional equations, obtaining lower bounds for the number of sign changes of the real and imaginary parts of the corresponding error terms. Particular applications include the divisor function $d(n)$, Ramanujan's function $\tau(n)$, and the number of representations of~$n$ as the sum of~$k$ squares (whose associated Dirichlet series is the Epstein zeta-function $\zeta_k(s)$).\par This article cites~\cites{MR1511360}.}
      \field{issn}{0010-2571}
      \field{journaltitle}{Comment. Math. Helv.}
      \field{title}{The changes of sign of certain arithmetical error-terms}
      \field{volume}{44}
      \field{year}{1969}
      \true{nocite}
      \field{pages}{385\bibrangedash 400}
      \range{pages}{16}
      \verb{doi}
      \verb 10.1007/BF02564539
      \endverb
      \verb{mrnumber}
      \verb 0257003
      \endverb
    \endentry
    \entry{1897.vonSterneck}{article}{}
      \name{author}{1}{}{%
        {{hash=f4689d7596dd95e869b4ea317492a6f8}{%
           family={Sterneck},
           familyi={S\bibinitperiod},
           given={R.\bibnamedelimi D.},
           giveni={R\bibinitperiod\bibinitdelim D\bibinitperiod},
           prefix={von},
           prefixi={v\bibinitperiod}}}%
      }
      \list{language}{1}{%
        {German}%
      }
      \strng{namehash}{f4689d7596dd95e869b4ea317492a6f8}
      \strng{fullhash}{f4689d7596dd95e869b4ea317492a6f8}
      \strng{bibnamehash}{f4689d7596dd95e869b4ea317492a6f8}
      \strng{authorbibnamehash}{f4689d7596dd95e869b4ea317492a6f8}
      \strng{authornamehash}{f4689d7596dd95e869b4ea317492a6f8}
      \strng{authorfullhash}{f4689d7596dd95e869b4ea317492a6f8}
      \field{extraname}{1}
      \field{sortinit}{S}
      \field{sortinithash}{c319cff79d99c853d775f88277d4e45f}
      \field{labelprefix}{A}
      \field{labelnamesource}{author}
      \field{labeltitlesource}{title}
      \field{annotation}{The author provides a table of values for $M(x)$. From the table, the author observes that $|M(n)|< \frac{1}{2} \sqrt{n}$ for $201 \leq n \leq 150,000$, and conjectures that $|M(x)|<\sqrt{x}$ for all $x>1$. The author also compares $M(x)$ with a random walk at the squarefree numbers, commenting that the corresponding expectation of the absolute value of the random walk, namely $\sqrt{\frac{12}{\pi^2}}\sqrt{x} \approx 1.10266\sqrt{x}$, exceeds even the maximal values of $|M(x)|$ from the table. \par This article cites~\cites{1897.Mertens}.}
      \field{journaltitle}{Sitzungsberichte Akad. Wiss. Wien IIa}
      \field{title}{Empirische Untersuchung über den Verlauf der zahlentheoretischen Funktion $\sigma(n) = \sum_{x=1}^{x=n}\mu(x)$ im Intervalle von $0$ bis $150000$}
      \field{volume}{106}
      \field{year}{1897}
      \true{nocite}
      \field{pages}{835\bibrangedash 1024}
      \range{pages}{190}
    \endentry
    \entry{1901.vonSterneck}{article}{}
      \name{author}{1}{}{%
        {{hash=f4689d7596dd95e869b4ea317492a6f8}{%
           family={Sterneck},
           familyi={S\bibinitperiod},
           given={R.\bibnamedelimi D.},
           giveni={R\bibinitperiod\bibinitdelim D\bibinitperiod},
           prefix={von},
           prefixi={v\bibinitperiod}}}%
      }
      \list{language}{1}{%
        {German}%
      }
      \strng{namehash}{f4689d7596dd95e869b4ea317492a6f8}
      \strng{fullhash}{f4689d7596dd95e869b4ea317492a6f8}
      \strng{bibnamehash}{f4689d7596dd95e869b4ea317492a6f8}
      \strng{authorbibnamehash}{f4689d7596dd95e869b4ea317492a6f8}
      \strng{authornamehash}{f4689d7596dd95e869b4ea317492a6f8}
      \strng{authorfullhash}{f4689d7596dd95e869b4ea317492a6f8}
      \field{extraname}{2}
      \field{sortinit}{S}
      \field{sortinithash}{c319cff79d99c853d775f88277d4e45f}
      \field{labelprefix}{A}
      \field{labelnamesource}{author}
      \field{labeltitlesource}{title}
      \field{annotation}{The author verifies that $|M(n)| < \frac12\sqrt n$ for $201 \le n\le500,000$. \par This article cites~\cites{1897.vonSterneck}.}
      \field{journaltitle}{SBer. Kais. Akad. Wissensch. Wien IIa}
      \field{title}{Empirische Untersuchung über den Verlauf der zahlentheoretischen Funktion $\sigma(n) = \sum_{x=1}^{x=n}\mu(x)$ im Intervalle von $150000$ bis $500000$}
      \field{volume}{110}
      \field{year}{1901}
      \true{nocite}
      \field{pages}{1053\bibrangedash 1102}
      \range{pages}{50}
    \endentry
    \entry{1912.vonSterneck}{article}{}
      \name{author}{1}{}{%
        {{hash=f4689d7596dd95e869b4ea317492a6f8}{%
           family={Sterneck},
           familyi={S\bibinitperiod},
           given={R.\bibnamedelimi D.},
           giveni={R\bibinitperiod\bibinitdelim D\bibinitperiod},
           prefix={von},
           prefixi={v\bibinitperiod}}}%
      }
      \list{language}{1}{%
        {German}%
      }
      \strng{namehash}{f4689d7596dd95e869b4ea317492a6f8}
      \strng{fullhash}{f4689d7596dd95e869b4ea317492a6f8}
      \strng{bibnamehash}{f4689d7596dd95e869b4ea317492a6f8}
      \strng{authorbibnamehash}{f4689d7596dd95e869b4ea317492a6f8}
      \strng{authornamehash}{f4689d7596dd95e869b4ea317492a6f8}
      \strng{authorfullhash}{f4689d7596dd95e869b4ea317492a6f8}
      \field{extraname}{3}
      \field{sortinit}{S}
      \field{sortinithash}{c319cff79d99c853d775f88277d4e45f}
      \field{labelprefix}{A}
      \field{labelnamesource}{author}
      \field{labeltitlesource}{title}
      \field{annotation}{The author calculates $16$ selected values of $M(n)$ with $600,000 \le n \le 5,000,000$ and verifies they all satisfy $|M(n)| \le \frac12\sqrt n$. The computation uses the following formula, which is a refined version of a formula of Mertens~\cite{1897.Mertens}: for $j=1,2,3,4$, if $n$ exceeds the product of the first~$j$ primes, then $ \sum_{d=1}^{\lfloor \sqrt{n} \rfloor} \mu(d) \omega_j(n/d) +\sum_{d'} M(n/d) = \omega_j\bigl( \lfloor \sqrt{n} \rfloor \bigr) M\bigl( \lfloor \sqrt{n} \rfloor \bigr), $ where $\omega_j(x)$ is the number of positive integers $\leq x$ that are not divisible by any of the first~$j$ primes, and $d'$ runs over all such numbers $\leq g$. The author claims that the bound $|M(n)| \le \frac12\sqrt n$ for all $n>200$ represents an unproved but extremely probable number theoretic law, and thus RH could be also regarded as correct with a high degree of probability (original German: ``{\it $|M(n)| \le \frac12\sqrt n$ zwar ein unbewiesenes, aber ausserordentlich wahrscheinliches zahlentheoretisches Gesetz darstellt, und somit auch die Riemann'sche Vermutung mit einem hohen Grad von Wahrscheinlichkeit als richtig angesehen werden kann}''). \par This article cites~\cites{1897.Mertens,1897.vonSterneck,1901.vonSterneck}.}
      \field{journaltitle}{SBer. Kais. Akad. Wissensch. Wien IIa}
      \field{title}{Die zahlentheoretische Funktion $\sigma(n)$ bis zur Grenze $5000000$}
      \field{volume}{121}
      \field{year}{1912}
      \true{nocite}
      \field{pages}{1083\bibrangedash 1096}
      \range{pages}{14}
    \endentry
    \entry{1913.vonSterneck}{article}{}
      \name{author}{1}{}{%
        {{hash=f4689d7596dd95e869b4ea317492a6f8}{%
           family={Sterneck},
           familyi={S\bibinitperiod},
           given={R.\bibnamedelimi D.},
           giveni={R\bibinitperiod\bibinitdelim D\bibinitperiod},
           prefix={von},
           prefixi={v\bibinitperiod}}}%
      }
      \list{language}{1}{%
        {German}%
      }
      \list{publisher}{1}{%
        {Cambridge University Press}%
      }
      \strng{namehash}{f4689d7596dd95e869b4ea317492a6f8}
      \strng{fullhash}{f4689d7596dd95e869b4ea317492a6f8}
      \strng{bibnamehash}{f4689d7596dd95e869b4ea317492a6f8}
      \strng{authorbibnamehash}{f4689d7596dd95e869b4ea317492a6f8}
      \strng{authornamehash}{f4689d7596dd95e869b4ea317492a6f8}
      \strng{authorfullhash}{f4689d7596dd95e869b4ea317492a6f8}
      \field{extraname}{4}
      \field{sortinit}{S}
      \field{sortinithash}{c319cff79d99c853d775f88277d4e45f}
      \field{labelprefix}{A}
      \field{labelnamesource}{author}
      \field{labeltitlesource}{title}
      \field{annotation}{This is a summary of the previous three articles of the author. \par This article cites~\cites{1897.Mertens,1897.vonSterneck, 1901.vonSterneck,1912.vonSterneck}.}
      \field{journaltitle}{Proc. 5th International Congress of Mathematicians}
      \field{title}{Neue empirische Daten über die zahlentheoretische Funktion $\sigma(n)$}
      \field{volume}{1}
      \field{year}{1913}
      \true{nocite}
      \field{pages}{341\bibrangedash 343}
      \range{pages}{3}
    \endentry
    \entry{1905.Stieltjes}{book}{}
      \name{author}{1}{}{%
        {{hash=39453b92c259f9e27936958ace288ccd}{%
           family={Stieltjes},
           familyi={S\bibinitperiod},
           given={T.\bibnamedelimi J.},
           giveni={T\bibinitperiod\bibinitdelim J\bibinitperiod}}}%
      }
      \list{language}{1}{%
        {French}%
      }
      \list{publisher}{1}{%
        {Gauthier--Villars, Imprimeur--Libraire, Paris}%
      }
      \strng{namehash}{39453b92c259f9e27936958ace288ccd}
      \strng{fullhash}{39453b92c259f9e27936958ace288ccd}
      \strng{bibnamehash}{39453b92c259f9e27936958ace288ccd}
      \strng{authorbibnamehash}{39453b92c259f9e27936958ace288ccd}
      \strng{authornamehash}{39453b92c259f9e27936958ace288ccd}
      \strng{authorfullhash}{39453b92c259f9e27936958ace288ccd}
      \field{sortinit}{S}
      \field{sortinithash}{c319cff79d99c853d775f88277d4e45f}
      \field{labelprefix}{A}
      \field{labelnamesource}{author}
      \field{labeltitlesource}{title}
      \field{annotation}{In Lettre~79 (starting on page~160), the author states the conjectural bound $\big| \sum_{n\le x} \mu(n) \big| \le \sqrt x$ for the Mertens function.}
      \field{title}{Correspondance d'Hermite et de Stieltjes}
      \field{year}{1905}
      \true{nocite}
      \field{pages}{xxi+pp. 1\bibrangedash 477}
      \range{pages}{-1}
    \endentry
    \entry{MR2813366}{article}{}
      \name{author}{2}{}{%
        {{hash=ba700e1ffe0ab3bd31564865c0880f36}{%
           family={Stoll},
           familyi={S\bibinitperiod},
           given={Douglas\bibnamedelima A.},
           giveni={D\bibinitperiod\bibinitdelim A\bibinitperiod}}}%
        {{hash=447f22d7de07f4cbd918f6b2abb97288}{%
           family={Demichel},
           familyi={D\bibinitperiod},
           given={Patrick},
           giveni={P\bibinitperiod}}}%
      }
      \strng{namehash}{307314f703dd0ce2792f1f0923e08ca4}
      \strng{fullhash}{307314f703dd0ce2792f1f0923e08ca4}
      \strng{bibnamehash}{307314f703dd0ce2792f1f0923e08ca4}
      \strng{authorbibnamehash}{307314f703dd0ce2792f1f0923e08ca4}
      \strng{authornamehash}{307314f703dd0ce2792f1f0923e08ca4}
      \strng{authorfullhash}{307314f703dd0ce2792f1f0923e08ca4}
      \field{sortinit}{S}
      \field{sortinithash}{c319cff79d99c853d775f88277d4e45f}
      \field{labelprefix}{A}
      \field{labelnamesource}{author}
      \field{labeltitlesource}{title}
      \field{annotation}{The authors analyze $E^\pi(x)$ for $x<10^{10^{13}}$ using the first $2\times10^{11}$ nontrivial zeros of $\zeta(s)$. Based on numerical computation, the authors suggest that there may exist an~$x$ near $1.397162914 \times 10^{316}$ such that $\pi(x)>\li(x)$. The authors also conjecture that $|E^\pi(x)|<\frac1e(\log \log \log x+e+1)$.\par This article cites~\cites{MR1752093,MR0369287,MR2652902,1932.Ingham,MR2420864,MR0202686,1914.Littlewood,MR0137689,MR1329368,MR2684372,MR1573970,MR0067145}.}
      \field{issn}{0025-5718}
      \field{journaltitle}{Math. Comp.}
      \field{number}{276}
      \field{title}{The impact of {$\zeta(s)$} complex zeros on {$\pi(x)$} for {$x<10^{10^{13}}$}}
      \field{volume}{80}
      \field{year}{2011}
      \true{nocite}
      \field{pages}{2381\bibrangedash 2394}
      \range{pages}{14}
      \verb{doi}
      \verb 10.1090/S0025-5718-2011-02477-4
      \endverb
      \verb{mrnumber}
      \verb 2813366
      \endverb
    \endentry
    \entry{MR0973808A}{article}{}
      \name{author}{1}{}{%
        {{hash=83e5472a007b4ceae548c8cea09a150a}{%
           family={Szydło},
           familyi={S\bibinitperiod},
           given={B.},
           giveni={B\bibinitperiod}}}%
      }
      \list{language}{1}{%
        {German}%
      }
      \strng{namehash}{83e5472a007b4ceae548c8cea09a150a}
      \strng{fullhash}{83e5472a007b4ceae548c8cea09a150a}
      \strng{bibnamehash}{83e5472a007b4ceae548c8cea09a150a}
      \strng{authorbibnamehash}{83e5472a007b4ceae548c8cea09a150a}
      \strng{authornamehash}{83e5472a007b4ceae548c8cea09a150a}
      \strng{authorfullhash}{83e5472a007b4ceae548c8cea09a150a}
      \field{extraname}{1}
      \field{sortinit}{S}
      \field{sortinithash}{c319cff79d99c853d775f88277d4e45f}
      \field{labelprefix}{A}
      \field{labelnamesource}{author}
      \field{labeltitlesource}{title}
      \field{annotation}{This article concerns effective results of the type $W(f;X) \ge c\log X$ for $X\ge X_0$, using a modification of Kaczorowski's method with a more general averaging operator $x^{-d-k}\int_0^{x}f(\xi) \xi^{d+k-1} \,d\xi$.\par This article cites~\cite{1914.Littlewood, MR777013, MR791085, MR858395, MR0973808B}.}
      \field{journaltitle}{Math. Ann.}
      \field{title}{{Über Vorzeichenwechsel einiger arithmetischer Funktionen. I}}
      \field{volume}{283}
      \field{year}{1989}
      \true{nocite}
      \field{pages}{139\bibrangedash 149}
      \range{pages}{11}
      \verb{mrnumber}
      \verb 0973808
      \endverb
    \endentry
    \entry{MR0973808B}{article}{}
      \name{author}{1}{}{%
        {{hash=83e5472a007b4ceae548c8cea09a150a}{%
           family={Szydło},
           familyi={S\bibinitperiod},
           given={B.},
           giveni={B\bibinitperiod}}}%
      }
      \list{language}{1}{%
        {German}%
      }
      \strng{namehash}{83e5472a007b4ceae548c8cea09a150a}
      \strng{fullhash}{83e5472a007b4ceae548c8cea09a150a}
      \strng{bibnamehash}{83e5472a007b4ceae548c8cea09a150a}
      \strng{authorbibnamehash}{83e5472a007b4ceae548c8cea09a150a}
      \strng{authornamehash}{83e5472a007b4ceae548c8cea09a150a}
      \strng{authorfullhash}{83e5472a007b4ceae548c8cea09a150a}
      \field{extraname}{2}
      \field{sortinit}{S}
      \field{sortinithash}{c319cff79d99c853d775f88277d4e45f}
      \field{labelprefix}{A}
      \field{labelnamesource}{author}
      \field{labeltitlesource}{title}
      \field{annotation}{The author uses his results from~\cite{MR0973808A} to give explicit lower bounds for the number of sign changes of $\Delta^\psi(X)$. He shows unconditionally that $W^\psi(X)\ge 0.013\log X$ when $X\ge 10^{2250}$, the proof requiring minimal computation, as well as the similar results $W^\psi(X)\ge 0.994 \frac{\gamma_1}{\pi} \log X$ for $X\ge \exp (198,594)$ and $W^\psi(X)\ge 0.99999997 \frac{\gamma_1}{\pi} \log X$ for $X\ge \exp (9 \times 10^{14})$. Assuming RH$(H)$ for some $H\ge501.5$, he further shows that $W^\psi(X)\ge (1-\frac{3}{H})\frac{\gamma_1}{\pi}\log X$ when $X\ge \exp(0.09\max\{4400,H\})$.\par This article cites~\cite{MR777013, MR791085, MR858395, MR0973808A}.}
      \field{journaltitle}{Math. Ann.}
      \field{title}{{Über Vorzeichenwechsel einiger arithmetischer Funktionen. II}}
      \field{volume}{283}
      \field{year}{1989}
      \true{nocite}
      \field{pages}{151\bibrangedash 163}
      \range{pages}{13}
      \verb{mrnumber}
      \verb 0973808
      \endverb
    \endentry
    \entry{MR1029966}{article}{}
      \name{author}{1}{}{%
        {{hash=83e5472a007b4ceae548c8cea09a150a}{%
           family={Szydło},
           familyi={S\bibinitperiod},
           given={B.},
           giveni={B\bibinitperiod}}}%
      }
      \list{language}{1}{%
        {German}%
      }
      \strng{namehash}{83e5472a007b4ceae548c8cea09a150a}
      \strng{fullhash}{83e5472a007b4ceae548c8cea09a150a}
      \strng{bibnamehash}{83e5472a007b4ceae548c8cea09a150a}
      \strng{authorbibnamehash}{83e5472a007b4ceae548c8cea09a150a}
      \strng{authornamehash}{83e5472a007b4ceae548c8cea09a150a}
      \strng{authorfullhash}{83e5472a007b4ceae548c8cea09a150a}
      \field{extraname}{3}
      \field{sortinit}{S}
      \field{sortinithash}{c319cff79d99c853d775f88277d4e45f}
      \field{labelprefix}{A}
      \field{labelnamesource}{author}
      \field{labeltitlesource}{title}
      \field{annotation}{The author generalizes his results~\cites{MR0973808A, MR0973808B} on the number of sign changes in $W^\psi(x)$ to algebraic number fields. Let $K$ be an algebraic number field of degree $n$ and $d$ the absolute value of its discriminant. In this article, he extends previous work~\cite{MR0980524} by giving effective lower bounds for the function $W^{\psi_K}(X)$ which counts the sign changes of $\sum_{{\mathfrak Np}^r \leq x}\log {\mathfrak Np}-x$. Assume GRH$(H,0)$ for the Dedekind zeta function~$\zeta_K$. When~$H$ and~$X$ are sufficiently large (in a precise way given with effective constants), the author shows that $W^{\psi_K}(X)\ge (1-\frac{10}{H})\frac{\gamma_K}{\pi}\log X$, where $\gamma_K$ is the positive imaginary part of the nontrivial zero of~$\zeta_K$ closest to the real axis.\par This article cites~\cite{MR777013, MR791085, MR858395, MR891057, MR0973808A, MR0973808B,MR0980524}.}
      \field{journaltitle}{Monatsh. Math.}
      \field{title}{{Über Vorzeichenwechsel einiger arithmetischer Funktionen. III}}
      \field{volume}{108}
      \field{year}{1989}
      \true{nocite}
      \field{pages}{325\bibrangedash 336}
      \range{pages}{12}
      \verb{mrnumber}
      \verb 1029966
      \endverb
    \endentry
    \entry{MR584557}{article}{}
      \name{author}{1}{}{%
        {{hash=31d751abf6c11534c2c4933716b5a5c6}{%
           family={Tanaka},
           familyi={T\bibinitperiod},
           given={Minoru},
           giveni={M\bibinitperiod}}}%
      }
      \strng{namehash}{31d751abf6c11534c2c4933716b5a5c6}
      \strng{fullhash}{31d751abf6c11534c2c4933716b5a5c6}
      \strng{bibnamehash}{31d751abf6c11534c2c4933716b5a5c6}
      \strng{authorbibnamehash}{31d751abf6c11534c2c4933716b5a5c6}
      \strng{authornamehash}{31d751abf6c11534c2c4933716b5a5c6}
      \strng{authorfullhash}{31d751abf6c11534c2c4933716b5a5c6}
      \field{extraname}{1}
      \field{sortinit}{T}
      \field{sortinithash}{51f9faf24c60c62ca764a77f78cf5666}
      \field{labelprefix}{A}
      \field{labelnamesource}{author}
      \field{labeltitlesource}{title}
      \field{annotation}{This article reports on the sign changes of $L(x)$ for $x\leq 10^9$. In particular, the author shows that $906,150,256$ is the smallest integer $n \geq 2$ such that $L(n)>0$.\par This article cites~\cites{MR0006202,MR0120198}.}
      \field{issn}{0387-3870}
      \field{journaltitle}{Tokyo J. Math.}
      \field{number}{1}
      \field{title}{A numerical investigation on cumulative sum of the {L}iouville function}
      \field{volume}{3}
      \field{year}{1980}
      \true{nocite}
      \field{pages}{187\bibrangedash 189}
      \range{pages}{3}
      \verb{doi}
      \verb 10.3836/tjm/1270216093
      \endverb
      \verb{mrnumber}
      \verb 584557
      \endverb
      \verb{urlraw}
      \verb https://doi.org/10.3836/tjm/1270216093
      \endverb
      \verb{url}
      \verb https://doi.org/10.3836/tjm/1270216093
      \endverb
    \endentry
    \entry{MR0605090}{article}{}
      \name{author}{1}{}{%
        {{hash=31d751abf6c11534c2c4933716b5a5c6}{%
           family={Tanaka},
           familyi={T\bibinitperiod},
           given={Minoru},
           giveni={M\bibinitperiod}}}%
      }
      \strng{namehash}{31d751abf6c11534c2c4933716b5a5c6}
      \strng{fullhash}{31d751abf6c11534c2c4933716b5a5c6}
      \strng{bibnamehash}{31d751abf6c11534c2c4933716b5a5c6}
      \strng{authorbibnamehash}{31d751abf6c11534c2c4933716b5a5c6}
      \strng{authornamehash}{31d751abf6c11534c2c4933716b5a5c6}
      \strng{authorfullhash}{31d751abf6c11534c2c4933716b5a5c6}
      \field{extraname}{2}
      \field{sortinit}{T}
      \field{sortinithash}{51f9faf24c60c62ca764a77f78cf5666}
      \field{labelprefix}{A}
      \field{labelnamesource}{author}
      \field{labeltitlesource}{title}
      \field{annotation}{Define $B_k=\zeta\left(\frac{k}2\right) / \zeta\left(\frac{1}{2}\right) \zeta(k)$ if $k\ge3$ is odd and $B_k=1 / \zeta\left(\frac{1}{2}\right) \zeta\left(\frac{k}2\right)$ if $k\ge4$ is even. The author shows that $\sum_{n \leq x}\mu_k(x)-B_k \sqrt{x}=\Omega_{\pm}(\sqrt{x})$ for each $k\ge3$, generalizing classical oscillation results for $M(x)$ (which is the case $k=2$) and $L(x)$ (which is the case $k\to\infty$).\par This article cites~\cites{MR0068565,MR0006202}.}
      \field{issn}{0387-3870}
      \field{journaltitle}{Tokyo J. Math.}
      \field{number}{2}
      \field{title}{On the {M}öbius and allied functions}
      \field{volume}{3}
      \field{year}{1980}
      \true{nocite}
      \field{pages}{215\bibrangedash 218}
      \range{pages}{4}
      \verb{doi}
      \verb 10.3836/tjm/1270472994
      \endverb
      \verb{mrnumber}
      \verb 605090
      \endverb
      \verb{urlraw}
      \verb https://doi.org/10.3836/tjm/1270472994
      \endverb
      \verb{url}
      \verb https://doi.org/10.3836/tjm/1270472994
      \endverb
    \endentry
    \entry{MR0017310}{article}{}
      \name{author}{1}{}{%
        {{hash=a8001a783b2aab7a2865fcdf627bcf26}{%
           family={Tietze},
           familyi={T\bibinitperiod},
           given={Heinrich},
           giveni={H\bibinitperiod}}}%
      }
      \list{language}{1}{%
        {German}%
      }
      \strng{namehash}{a8001a783b2aab7a2865fcdf627bcf26}
      \strng{fullhash}{a8001a783b2aab7a2865fcdf627bcf26}
      \strng{bibnamehash}{a8001a783b2aab7a2865fcdf627bcf26}
      \strng{authorbibnamehash}{a8001a783b2aab7a2865fcdf627bcf26}
      \strng{authornamehash}{a8001a783b2aab7a2865fcdf627bcf26}
      \strng{authorfullhash}{a8001a783b2aab7a2865fcdf627bcf26}
      \field{sortinit}{T}
      \field{sortinithash}{51f9faf24c60c62ca764a77f78cf5666}
      \field{labelprefix}{A}
      \field{labelnamesource}{author}
      \field{labeltitlesource}{title}
      \field{annotation}{The author provides tables of $\phi(q)\pi(x;q,\cA,\Gamma)$ for various moduli~$q$ (with a particular focus on~$q=262$, the smallest even modulus for which $3,5,7,11,13$ are all quadratic residues) and various sets~$\Gamma$ of residue classes, often subgroups of~$\cA$.}
      \field{journaltitle}{Abh. Bayer. Akad. Wiss. Math.-Nat. Abt. (N.F.)}
      \field{number}{55}
      \field{title}{Einige Tabellen zur Verteilung der Primzahlen auf Untergruppen der teilerfremden Restklassen nach gegebenem Modul}
      \field{volume}{1944}
      \field{year}{1944}
      \true{nocite}
      \field{pages}{31}
      \range{pages}{1}
      \verb{mrnumber}
      \verb 0017310
      \endverb
    \endentry
    \entry{MR0046485}{book}{}
      \name{author}{1}{}{%
        {{hash=baf219d417d1941af57ccbb309deb548}{%
           family={Titchmarsh},
           familyi={T\bibinitperiod},
           given={E.\bibnamedelimi C.},
           giveni={E\bibinitperiod\bibinitdelim C\bibinitperiod}}}%
      }
      \list{publisher}{1}{%
        {Oxford, at the Clarendon Press}%
      }
      \strng{namehash}{baf219d417d1941af57ccbb309deb548}
      \strng{fullhash}{baf219d417d1941af57ccbb309deb548}
      \strng{bibnamehash}{baf219d417d1941af57ccbb309deb548}
      \strng{authorbibnamehash}{baf219d417d1941af57ccbb309deb548}
      \strng{authornamehash}{baf219d417d1941af57ccbb309deb548}
      \strng{authorfullhash}{baf219d417d1941af57ccbb309deb548}
      \field{extraname}{1}
      \field{sortinit}{T}
      \field{sortinithash}{51f9faf24c60c62ca764a77f78cf5666}
      \field{labelprefix}{A}
      \field{labelnamesource}{author}
      \field{labeltitlesource}{title}
      \field{annotation}{This is the first edition of~\cite{MR0882550}.}
      \field{title}{The {T}heory of the {R}iemann {Z}eta-{F}unction}
      \field{year}{1951}
      \true{nocite}
      \field{pages}{vi+346}
      \range{pages}{-1}
      \verb{mrnumber}
      \verb 0046485
      \endverb
    \endentry
    \entry{MR0882550}{book}{}
      \name{author}{1}{}{%
        {{hash=baf219d417d1941af57ccbb309deb548}{%
           family={Titchmarsh},
           familyi={T\bibinitperiod},
           given={E.\bibnamedelimi C.},
           giveni={E\bibinitperiod\bibinitdelim C\bibinitperiod}}}%
      }
      \list{publisher}{1}{%
        {The Clarendon Press, Oxford University Press, New York}%
      }
      \strng{namehash}{baf219d417d1941af57ccbb309deb548}
      \strng{fullhash}{baf219d417d1941af57ccbb309deb548}
      \strng{bibnamehash}{baf219d417d1941af57ccbb309deb548}
      \strng{authorbibnamehash}{baf219d417d1941af57ccbb309deb548}
      \strng{authornamehash}{baf219d417d1941af57ccbb309deb548}
      \strng{authorfullhash}{baf219d417d1941af57ccbb309deb548}
      \field{extraname}{2}
      \field{sortinit}{T}
      \field{sortinithash}{51f9faf24c60c62ca764a77f78cf5666}
      \field{labelprefix}{A}
      \field{labelnamesource}{author}
      \field{labeltitlesource}{title}
      \field{annotation}{This book is the second edition of the classical and comprehensive treatise on the theory of the Riemann zeta-function, now with a large number of chapter-end notes by Heath--Brown. Chapter~XII contains the asymptotic formula for the error term in the divisor problem and relevant upper bounds and oscillation results. Chapter~XIV contains an explicit formula for $M(x)$ and its upper bounds and oscillations, as well as a connection between its mean-square average and the simplicity of zeros of $\zeta(s)$; in the chapter-end notes, it is noted that the Mertens conjecture has been disproved and Turán's problem resolved in the negative. \par This article cites~\cite{MR0006202, MR0104638, MR429789, MR783538}.}
      \field{edition}{Second}
      \field{isbn}{0-19-853369-1}
      \field{note}{Edited and with a preface by D. R. Heath-Brown}
      \field{title}{The theory of the {R}iemann zeta-function}
      \field{year}{1986}
      \true{nocite}
      \field{pages}{x+412}
      \range{pages}{-1}
      \verb{mrnumber}
      \verb 0882550
      \endverb
    \endentry
    \entry{MR0049219}{article}{}
      \name{author}{1}{}{%
        {{hash=e68e101772b88cb1d48e688164bece92}{%
           family={Turán},
           familyi={T\bibinitperiod},
           given={P.},
           giveni={P\bibinitperiod}}}%
      }
      \list{language}{1}{%
        {English, with Russian summary}%
      }
      \strng{namehash}{e68e101772b88cb1d48e688164bece92}
      \strng{fullhash}{e68e101772b88cb1d48e688164bece92}
      \strng{bibnamehash}{e68e101772b88cb1d48e688164bece92}
      \strng{authorbibnamehash}{e68e101772b88cb1d48e688164bece92}
      \strng{authornamehash}{e68e101772b88cb1d48e688164bece92}
      \strng{authorfullhash}{e68e101772b88cb1d48e688164bece92}
      \field{extraname}{1}
      \field{sortinit}{T}
      \field{sortinithash}{51f9faf24c60c62ca764a77f78cf5666}
      \field{labelprefix}{A}
      \field{labelnamesource}{author}
      \field{labeltitlesource}{title}
      \field{annotation}{The author shows that if $\Delta^\pi(x) \ll x\exp\bigl( -a(\log x)^{1/(1+\beta)} \bigr)$ for some constants $a>0$ and $0<\beta<1$, then there exists a constant $b>0$ such that $\zeta(s)$ does not vanish in the domain $\sigma>1-{b}/{\log^\beta|t|}$; this is the converse of a special case of a theorem of Ingham. The main tool of the proof is a version of the ``second main theorem'' of the power-sum method (the statement in the article seems to contain the typo $\max|z_j|\ge1$ instead of $\max|z_j|\le1$) for longer ranges of the exponent.\par This article cites~\cites{1932.Ingham}.}
      \field{issn}{0001-5954}
      \field{journaltitle}{Acta Math. Acad. Sci. Hungar.}
      \field{title}{On the remainder-term of the prime-number formula. {II}}
      \field{volume}{1}
      \field{year}{1950}
      \true{nocite}
      \field{pages}{155\bibrangedash 166}
      \range{pages}{12}
      \verb{doi}
      \verb 10.1007/BF02021308
      \endverb
      \verb{mrnumber}
      \verb 0049219
      \endverb
    \endentry
    \entry{MR115977}{article}{}
      \name{author}{1}{}{%
        {{hash=e68e101772b88cb1d48e688164bece92}{%
           family={Turán},
           familyi={T\bibinitperiod},
           given={P.},
           giveni={P\bibinitperiod}}}%
      }
      \strng{namehash}{e68e101772b88cb1d48e688164bece92}
      \strng{fullhash}{e68e101772b88cb1d48e688164bece92}
      \strng{bibnamehash}{e68e101772b88cb1d48e688164bece92}
      \strng{authorbibnamehash}{e68e101772b88cb1d48e688164bece92}
      \strng{authornamehash}{e68e101772b88cb1d48e688164bece92}
      \strng{authorfullhash}{e68e101772b88cb1d48e688164bece92}
      \field{extraname}{2}
      \field{sortinit}{T}
      \field{sortinithash}{51f9faf24c60c62ca764a77f78cf5666}
      \field{labelprefix}{A}
      \field{labelnamesource}{author}
      \field{labeltitlesource}{title}
      \field{annotation}{The main topic of this article is the connection between RH and zeros of partial sums of the series for $\zeta(s)$. Herein, however, the author mentions that the assertion that $L_r(x)$ is nonnegative (eventually or even always) has been called ``Turán's conjecture'', but that he never made such a claim even implicitly. The exact quote, in which ``${}^1$'' refers to~\cite{MR27305} and ``(3.2)'' refers to the inequality $L_r(x) \ge 0$, is: ``{\it Merkwürdigerweise bezeichnetendie an~${}^1$ anschlie{ß}enden Arbeiten schlechthin (3.2) als ,,Tur~insche Vermutung``, sogar für $n\ge1$ behauptet. (3.2) kommt in~${}^1$ nirgends vor, für $n\ge1$ nicht einmal implizite behauptet.}'' \par This article cites~\cites{MR27305,MR0104638,MR1511360,MR0046485}.}
      \field{issn}{0001-5954}
      \field{journaltitle}{Acta Math. Acad. Sci. Hungar.}
      \field{title}{Nachtrag zu meiner {A}bhandlung ``{O}n some approximative {D}irichlet polynomials in the theory of zeta-function of {R}iemann''}
      \field{volume}{10}
      \field{year}{1959}
      \true{nocite}
      \field{pages}{277\bibrangedash 298 (unbound insert)}
      \range{pages}{-1}
      \verb{doi}
      \verb 10.1007/BF02024493
      \endverb
      \verb{mrnumber}
      \verb 115977
      \endverb
      \verb{urlraw}
      \verb https://doi-org.ezproxy.library.ubc.ca/10.1007/BF02024493
      \endverb
      \verb{url}
      \verb https://doi-org.ezproxy.library.ubc.ca/10.1007/BF02024493
      \endverb
    \endentry
    \entry{MR0132728}{article}{}
      \name{author}{1}{}{%
        {{hash=e68e101772b88cb1d48e688164bece92}{%
           family={Turán},
           familyi={T\bibinitperiod},
           given={P.},
           giveni={P\bibinitperiod}}}%
      }
      \list{language}{1}{%
        {English, with Russian summary}%
      }
      \strng{namehash}{e68e101772b88cb1d48e688164bece92}
      \strng{fullhash}{e68e101772b88cb1d48e688164bece92}
      \strng{bibnamehash}{e68e101772b88cb1d48e688164bece92}
      \strng{authorbibnamehash}{e68e101772b88cb1d48e688164bece92}
      \strng{authornamehash}{e68e101772b88cb1d48e688164bece92}
      \strng{authorfullhash}{e68e101772b88cb1d48e688164bece92}
      \field{extraname}{3}
      \field{sortinit}{T}
      \field{sortinithash}{51f9faf24c60c62ca764a77f78cf5666}
      \field{labelprefix}{A}
      \field{labelnamesource}{author}
      \field{labeltitlesource}{title}
      \field{annotation}{The author establishes a ``one-sided'' version of the second main theorem of the power-sum method with an argument restriction. Let $m$ be a nonnegative integer and choose $0 < \alpha \leq \frac{\pi}2 $, and suppose that $z_1,\ldots,z_n$ are complex numbers satisfying $1=|z_1| \geq |z_2| \geq \cdots \geq |z_n|$ and $|\arg(z_j)| > \alpha$, and that $b_1,\ldots,b_n$ are complex numbers with $\min_{1\le\mu\le n} \Re \sum_{j=1}^{\mu} b_j > 0$. Then there exists an integer $m+1 \le v \le m+n(3+\frac\pi\alpha)$ such that $$ \Re \sum_{j=1}^{n} b_jz_j^v \ge \frac1{2n+1}\bigg( \frac{n}{24e^3(m+n(3+\frac\pi\alpha))} \bigg)^{2n} \min_{1\le\mu\le n} \Re \sum_{j=1}^{\mu} b_j, $$ with the analogous result for $-\Re \sum_{j=1}^{n} b_jz_j^v$. \par This article cites~\cites{MR0060548}.}
      \field{issn}{0001-5954}
      \field{journaltitle}{Acta Math. Acad. Sci. Hungar.}
      \field{title}{On some further one-sided theorems of new type in the theory of Diophantine approximations}
      \field{volume}{12}
      \field{year}{1961}
      \true{nocite}
      \field{pages}{455\bibrangedash 468}
      \range{pages}{14}
      \verb{doi}
      \verb 10.1007/BF02023928
      \endverb
      \verb{mrnumber}
      \verb 0132728
      \endverb
    \endentry
    \entry{MR0229595}{inproceedings}{}
      \name{author}{1}{}{%
        {{hash=e68e101772b88cb1d48e688164bece92}{%
           family={Turán},
           familyi={T\bibinitperiod},
           given={P.},
           giveni={P\bibinitperiod}}}%
      }
      \list{publisher}{1}{%
        {Izdat. ``Nauka'', Leningrad}%
      }
      \strng{namehash}{e68e101772b88cb1d48e688164bece92}
      \strng{fullhash}{e68e101772b88cb1d48e688164bece92}
      \strng{bibnamehash}{e68e101772b88cb1d48e688164bece92}
      \strng{authorbibnamehash}{e68e101772b88cb1d48e688164bece92}
      \strng{authornamehash}{e68e101772b88cb1d48e688164bece92}
      \strng{authorfullhash}{e68e101772b88cb1d48e688164bece92}
      \field{extraname}{4}
      \field{sortinit}{T}
      \field{sortinithash}{51f9faf24c60c62ca764a77f78cf5666}
      \field{labelprefix}{A}
      \field{labelnamesource}{author}
      \field{labeltitlesource}{title}
      \field{annotation}{The author states some results whose detailed proofs later appeared in the series~\cite{MR0146156} with Knapowski. First, the interval $[T^{1/3},T]$ contains a sign change of $\psi(x;4,3,1)$ when $T$ is large, while $E^\pi(x;4,3,1) = \Omega_\pm(\log\log\log\log\log x)$. Analogous results hold for the moduli $q=3$ and $q=6$. For $q=8$, when $\ell=3,5,7$ the analogues hold for $\pi(x;8,\ell,1)$, while when $\ell_1,\ell_2\neq1$ are distinct then $\pi(x;8,\ell_1,\ell_2)=\Omega_\pm(\sqrt{x}e^{-{\log{x}}/{\sqrt{\log\log{x}}}})$ for values $x$ in all large intervals $[T^{1/3},T]$, and similarly for $\pi_e(x;8,\ell_1,\ell_2)$. When $\ell\in\{3,5,7\}$, he shows that $\lim_{x\to\infty} \pi_e(x;8,\ell,1) = \infty$ if and only GRH$(\chi)$ holds for all $\chi\mod8$ with $\chi(\ell)\neq1$. The analogues hold for $q=5$ (with the understanding that the distinction between $1$ and $3,5,7\mod8$ becomes the distinction between quadratic residues and nonresidues\mod5), with the exception that the $\Omega$-result for $\pi_e(x;5,4,1)$ requires that the complex $L(s,\chi)$ vanishes in the region $\sigma>1-(\log\log t)^{-1/3}$.}
      \field{booktitle}{Proc. {F}ourth {A}ll-{U}nion {M}ath. {C}ongr ({L}eningrad, 1961) ({R}ussian), {V}ol. {II}}
      \field{title}{On a comparative theory of primes}
      \field{year}{1964}
      \true{nocite}
      \field{pages}{137\bibrangedash 142}
      \range{pages}{6}
      \verb{mrnumber}
      \verb 0229595
      \endverb
    \endentry
    \entry{MR27305}{article}{}
      \name{author}{1}{}{%
        {{hash=f4ee93efcace35c40df58fa6a8eb26ba}{%
           family={Turán},
           familyi={T\bibinitperiod},
           given={Paul},
           giveni={P\bibinitperiod}}}%
      }
      \strng{namehash}{f4ee93efcace35c40df58fa6a8eb26ba}
      \strng{fullhash}{f4ee93efcace35c40df58fa6a8eb26ba}
      \strng{bibnamehash}{f4ee93efcace35c40df58fa6a8eb26ba}
      \strng{authorbibnamehash}{f4ee93efcace35c40df58fa6a8eb26ba}
      \strng{authornamehash}{f4ee93efcace35c40df58fa6a8eb26ba}
      \strng{authorfullhash}{f4ee93efcace35c40df58fa6a8eb26ba}
      \field{extraname}{5}
      \field{sortinit}{T}
      \field{sortinithash}{51f9faf24c60c62ca764a77f78cf5666}
      \field{labelprefix}{A}
      \field{labelnamesource}{author}
      \field{labeltitlesource}{title}
      \field{annotation}{The author shows that if there exists a positive constant $K$ such that $L_r(x)>-K/\sqrt{x}$ for sufficiently large $x$, then RH holds. The author also reports that $L_r(x)\geq 0$ for $x \leq 1000$. While people often refer to the statement that $L_r(x)$ is never negative as ``Turán's conjecture'', that claim is never made in this article. \par This article cites~\cite{1919.Polya,1930.Polya,MR0003644}.}
      \field{issn}{0023-3323}
      \field{journaltitle}{Danske Vid. Selsk. Mat.-Fys. Medd.}
      \field{number}{17}
      \field{title}{On some approximative {D}irichlet-polynomials in the theory of the zeta-function of {R}iemann}
      \field{volume}{24}
      \field{year}{1948}
      \true{nocite}
      \field{pages}{36}
      \range{pages}{1}
      \verb{mrnumber}
      \verb 27305
      \endverb
    \endentry
    \entry{MR0043121}{article}{}
      \name{author}{1}{}{%
        {{hash=f4ee93efcace35c40df58fa6a8eb26ba}{%
           family={Turán},
           familyi={T\bibinitperiod},
           given={Paul},
           giveni={P\bibinitperiod}}}%
      }
      \list{language}{1}{%
        {English, with Russian summary}%
      }
      \strng{namehash}{f4ee93efcace35c40df58fa6a8eb26ba}
      \strng{fullhash}{f4ee93efcace35c40df58fa6a8eb26ba}
      \strng{bibnamehash}{f4ee93efcace35c40df58fa6a8eb26ba}
      \strng{authorbibnamehash}{f4ee93efcace35c40df58fa6a8eb26ba}
      \strng{authornamehash}{f4ee93efcace35c40df58fa6a8eb26ba}
      \strng{authorfullhash}{f4ee93efcace35c40df58fa6a8eb26ba}
      \field{extraname}{6}
      \field{sortinit}{T}
      \field{sortinithash}{51f9faf24c60c62ca764a77f78cf5666}
      \field{labelprefix}{A}
      \field{labelnamesource}{author}
      \field{labeltitlesource}{title}
      \field{annotation}{This article investigates lower bounds for $|\Delta^\psi(x)|$ and $|\Delta^\pi(x)|$, specifically to answer Littlewood's call~\cite{MR1575079} to replace Landau's theorem, which depends upon the smallest half-plane containing all the zeros of $\zeta(s)$, by an estimation that depends only upon one single zero $\rho_0=\beta_0+i\gamma_0$ of $\zeta(s)$. The author shows, using a version of his ``first main theorem'', that when $T$ is sufficiently large in terms of $\rho_0$, $$ \max_{1\le x\le T}|\Delta^\psi(x)| > \frac{T^{\beta_0}}{{|\rho_0|}^{(10\log T)/\log \log T}}\exp \left(-\frac{c\log T \log\log\log T}{\log \log T}\right), $$ where $c$ is a positive constant depending on~$\rho_0$. \par This article cites~\cites{1914.Littlewood, MR1555148, 1930.Polya, 1936.Ingham, MR0067145, MR0133309, MR1511206, MR1575079}.}
      \field{issn}{0001-5954}
      \field{journaltitle}{Acta Math. Acad. Sci. Hungar.}
      \field{title}{On the remainder-term of the prime-number formula. {I}}
      \field{volume}{1}
      \field{year}{1950}
      \true{nocite}
      \field{pages}{48\bibrangedash 63}
      \range{pages}{16}
      \verb{doi}
      \verb 10.1007/BF02022552
      \endverb
      \verb{mrnumber}
      \verb 0043121
      \endverb
    \endentry
    \entry{MR0060548}{book}{}
      \name{author}{1}{}{%
        {{hash=f4ee93efcace35c40df58fa6a8eb26ba}{%
           family={Turán},
           familyi={T\bibinitperiod},
           given={Paul},
           giveni={P\bibinitperiod}}}%
      }
      \list{language}{1}{%
        {German}%
      }
      \list{publisher}{1}{%
        {Akadémiai Kiadó, Budapest}%
      }
      \strng{namehash}{f4ee93efcace35c40df58fa6a8eb26ba}
      \strng{fullhash}{f4ee93efcace35c40df58fa6a8eb26ba}
      \strng{bibnamehash}{f4ee93efcace35c40df58fa6a8eb26ba}
      \strng{authorbibnamehash}{f4ee93efcace35c40df58fa6a8eb26ba}
      \strng{authornamehash}{f4ee93efcace35c40df58fa6a8eb26ba}
      \strng{authorfullhash}{f4ee93efcace35c40df58fa6a8eb26ba}
      \field{extraname}{7}
      \field{sortinit}{T}
      \field{sortinithash}{51f9faf24c60c62ca764a77f78cf5666}
      \field{labelprefix}{A}
      \field{labelnamesource}{author}
      \field{labeltitlesource}{title}
      \field{annotation}{This book contains a coherent exposition of power-sum method which is also referred as ``Turàn's method''. Many applications of the method are given, including a result that derives zero-free regions from bounds for $E^\psi(x)$ and improved zero-density results.}
      \field{pagetotal}{196}
      \field{title}{Eine neue Methode in der Analysis und deren Anwendungen}
      \field{year}{1953}
      \true{nocite}
      \verb{mrnumber}
      \verb 0060548
      \endverb
    \endentry
    \entry{MR0300854}{article}{}
      \name{author}{1}{}{%
        {{hash=f4ee93efcace35c40df58fa6a8eb26ba}{%
           family={Turán},
           familyi={T\bibinitperiod},
           given={Paul},
           giveni={P\bibinitperiod}}}%
      }
      \strng{namehash}{f4ee93efcace35c40df58fa6a8eb26ba}
      \strng{fullhash}{f4ee93efcace35c40df58fa6a8eb26ba}
      \strng{bibnamehash}{f4ee93efcace35c40df58fa6a8eb26ba}
      \strng{authorbibnamehash}{f4ee93efcace35c40df58fa6a8eb26ba}
      \strng{authornamehash}{f4ee93efcace35c40df58fa6a8eb26ba}
      \strng{authorfullhash}{f4ee93efcace35c40df58fa6a8eb26ba}
      \field{extraname}{8}
      \field{sortinit}{T}
      \field{sortinithash}{51f9faf24c60c62ca764a77f78cf5666}
      \field{labelprefix}{A}
      \field{labelnamesource}{author}
      \field{labeltitlesource}{title}
      \field{annotation}{The author recalls theorems from Knapowski's work after he died in a car accident at age~$36$.\par This article cites~\cites{MR0096622, MR0142520, MR0142521, MR0146156, MR0146157, MR0146158, MR0133287}.}
      \field{issn}{0010-1354}
      \field{journaltitle}{Colloq. Math.}
      \field{title}{Commemoration on Stanisław Knapowski}
      \field{volume}{23}
      \field{year}{1971}
      \true{nocite}
      \field{pages}{310\bibrangedash 318}
      \range{pages}{9}
      \verb{doi}
      \verb 10.4064/cm-23-2-309-321
      \endverb
      \verb{mrnumber}
      \verb 0300854
      \endverb
    \endentry
    \entry{MR749389}{book}{}
      \name{author}{1}{}{%
        {{hash=f4ee93efcace35c40df58fa6a8eb26ba}{%
           family={Turán},
           familyi={T\bibinitperiod},
           given={Paul},
           giveni={P\bibinitperiod}}}%
      }
      \list{publisher}{1}{%
        {John Wiley \& Sons, Inc., New York}%
      }
      \strng{namehash}{f4ee93efcace35c40df58fa6a8eb26ba}
      \strng{fullhash}{f4ee93efcace35c40df58fa6a8eb26ba}
      \strng{bibnamehash}{f4ee93efcace35c40df58fa6a8eb26ba}
      \strng{authorbibnamehash}{f4ee93efcace35c40df58fa6a8eb26ba}
      \strng{authornamehash}{f4ee93efcace35c40df58fa6a8eb26ba}
      \strng{authorfullhash}{f4ee93efcace35c40df58fa6a8eb26ba}
      \field{extraname}{9}
      \field{sortinit}{T}
      \field{sortinithash}{51f9faf24c60c62ca764a77f78cf5666}
      \field{labelprefix}{A}
      \field{labelnamesource}{author}
      \field{labeltitlesource}{title}
      \field{annotation}{This is a comprehensive work on the power-sum method, completed with the assistance of G.~Halász and J.~Pintz and including a foreword by V.~T.~Sós. The first part of the book introduces the method, while the second and larger part discusses a wide range of applications. Chapters 48--56 are directly related to comparative prime number theory.}
      \field{isbn}{0-471-89255-6}
      \field{series}{Pure and Applied Mathematics (New York)}
      \field{title}{On a new method of analysis and its applications}
      \field{year}{1984}
      \true{nocite}
      \field{pages}{xvi+584}
      \range{pages}{-1}
      \verb{mrnumber}
      \verb 749389
      \endverb
    \endentry
    \entry{MR1507933}{article}{}
      \name{author}{1}{}{%
        {{hash=0ecb7b4bc82132c8d3551c023cbe496a}{%
           family={Wintner},
           familyi={W\bibinitperiod},
           given={A.},
           giveni={A\bibinitperiod}}}%
      }
      \strng{namehash}{0ecb7b4bc82132c8d3551c023cbe496a}
      \strng{fullhash}{0ecb7b4bc82132c8d3551c023cbe496a}
      \strng{bibnamehash}{0ecb7b4bc82132c8d3551c023cbe496a}
      \strng{authorbibnamehash}{0ecb7b4bc82132c8d3551c023cbe496a}
      \strng{authornamehash}{0ecb7b4bc82132c8d3551c023cbe496a}
      \strng{authorfullhash}{0ecb7b4bc82132c8d3551c023cbe496a}
      \field{extraname}{1}
      \field{sortinit}{W}
      \field{sortinithash}{1af34bd8c148ffb32de1494636b49713}
      \field{labelprefix}{A}
      \field{labelnamesource}{author}
      \field{labeltitlesource}{title}
      \field{annotation}{Assuming RH, the author investigates the tail $\sum_{|\rho|>T} x^{i\gamma}/\rho$ of the sum appearing in the explicit formula for $\psi(x)$, showing that its mean-square average tends to~$0$ as $T\to\infty$. This implies that the full sum converges (and represents, in the Besikovitch sense, the function of which it is the ``Fourier series'') and has a logarithmic distribution function. The author also mentions the potential relevance of LI to whether this distribution function can be constant on intervals.\par This article cites~\cites{MR1501802}.}
      \field{issn}{0002-9327}
      \field{journaltitle}{Amer. J. Math.}
      \field{number}{3}
      \field{title}{On the Asymptotic Distribution of the Remainder Term of the Prime-Number Theorem}
      \field{volume}{57}
      \field{year}{1935}
      \true{nocite}
      \field{pages}{534\bibrangedash 538}
      \range{pages}{5}
      \verb{doi}
      \verb 10.2307/2371183
      \endverb
      \verb{mrnumber}
      \verb 1507933
      \endverb
    \endentry
    \entry{1938.Wintner}{article}{}
      \name{author}{1}{}{%
        {{hash=0ecb7b4bc82132c8d3551c023cbe496a}{%
           family={Wintner},
           familyi={W\bibinitperiod},
           given={A.},
           giveni={A\bibinitperiod}}}%
      }
      \strng{namehash}{0ecb7b4bc82132c8d3551c023cbe496a}
      \strng{fullhash}{0ecb7b4bc82132c8d3551c023cbe496a}
      \strng{bibnamehash}{0ecb7b4bc82132c8d3551c023cbe496a}
      \strng{authorbibnamehash}{0ecb7b4bc82132c8d3551c023cbe496a}
      \strng{authornamehash}{0ecb7b4bc82132c8d3551c023cbe496a}
      \strng{authorfullhash}{0ecb7b4bc82132c8d3551c023cbe496a}
      \field{extraname}{2}
      \field{sortinit}{W}
      \field{sortinithash}{1af34bd8c148ffb32de1494636b49713}
      \field{labelprefix}{A}
      \field{labelnamesource}{author}
      \field{labeltitlesource}{title}
      \field{annotation}{The chapter titles are: 1.~Distribution functions; 2.~Integrals and convolutions; 3.~Moments of distribution functions; 4.~Fourier transforms; 5.~Infinite convolutions; 6.~Smoothness criteria for distribution functions; 7.~Convergence criteria for infinite convolutions; 8.~Asymptotic distributions; 9.~The Riemann zeta-function; 10.~Poisson convolutions; 11.~Convolutions and the theory of probability; 12.~Bernoulli convolutions; 13.~Almost periodic functions with linearly independent exponents; 14.~Symmetric distribution functions in $k$ dimensions; 15.~Two-dimensional convolutions and the Riemann zeta-function; 16.~The addition of convex curves.\par These notes cite~\cites{MR1501802,MR1507933}.}
      \field{journaltitle}{Lecture notes distributed by the Institute for Advanced Study (Princeton)}
      \field{title}{Asymptotic distributions and infinite convolutions}
      \field{year}{1938}
      \true{nocite}
    \endentry
    \entry{MR0004255}{article}{}
      \name{author}{1}{}{%
        {{hash=0ecb7b4bc82132c8d3551c023cbe496a}{%
           family={Wintner},
           familyi={W\bibinitperiod},
           given={A.},
           giveni={A\bibinitperiod}}}%
      }
      \strng{namehash}{0ecb7b4bc82132c8d3551c023cbe496a}
      \strng{fullhash}{0ecb7b4bc82132c8d3551c023cbe496a}
      \strng{bibnamehash}{0ecb7b4bc82132c8d3551c023cbe496a}
      \strng{authorbibnamehash}{0ecb7b4bc82132c8d3551c023cbe496a}
      \strng{authornamehash}{0ecb7b4bc82132c8d3551c023cbe496a}
      \strng{authorfullhash}{0ecb7b4bc82132c8d3551c023cbe496a}
      \field{extraname}{3}
      \field{sortinit}{W}
      \field{sortinithash}{1af34bd8c148ffb32de1494636b49713}
      \field{labelprefix}{A}
      \field{labelnamesource}{author}
      \field{labeltitlesource}{title}
      \field{annotation}{This article investigates the normalized remainder term $E^\Psi(x)$ by establishing the existence of its limiting logarithmic cumulative distribution function. Assuming RH, the author proves that the spectrum of this limiting distribution is unbounded both above and below, which implies the same for $E^\Psi(x)$. He also gives an estimate for all of the moments of this limiting distribution, in an effort to determine whether these moments uniquely determine the given distribution.\par The techniques involved in the proof come primarily from the application of the theory of almost-periodic functions (of both the uniform and Besicovitch varieties) to the sum over the nontrivial zeros in the fundamental explicit formula.\par This article cites~\cites{1932.Ingham,MR1501802,MR1507933}.}
      \field{issn}{0002-9327}
      \field{journaltitle}{Amer. J. Math.}
      \field{title}{On the distribution function of the remainder term of the prime number theorem}
      \field{volume}{63}
      \field{year}{1941}
      \true{nocite}
      \field{pages}{233\bibrangedash 248}
      \range{pages}{16}
      \verb{doi}
      \verb 10.2307/2371519
      \endverb
      \verb{mrnumber}
      \verb 0004255
      \endverb
      \verb{urlraw}
      \verb https://doi.org/10.2307/2371519
      \endverb
      \verb{url}
      \verb https://doi.org/10.2307/2371519
      \endverb
    \endentry
    \entry{MR29414}{article}{}
      \name{author}{1}{}{%
        {{hash=86b5f5639455f1050b0d4e2301179477}{%
           family={Wintner},
           familyi={W\bibinitperiod},
           given={Aurel},
           giveni={A\bibinitperiod}}}%
      }
      \strng{namehash}{86b5f5639455f1050b0d4e2301179477}
      \strng{fullhash}{86b5f5639455f1050b0d4e2301179477}
      \strng{bibnamehash}{86b5f5639455f1050b0d4e2301179477}
      \strng{authorbibnamehash}{86b5f5639455f1050b0d4e2301179477}
      \strng{authornamehash}{86b5f5639455f1050b0d4e2301179477}
      \strng{authorfullhash}{86b5f5639455f1050b0d4e2301179477}
      \field{extraname}{4}
      \field{sortinit}{W}
      \field{sortinithash}{1af34bd8c148ffb32de1494636b49713}
      \field{labelprefix}{A}
      \field{labelnamesource}{author}
      \field{labeltitlesource}{title}
      \field{annotation}{The author shows that $M(x)\ll\sqrt{x}$ is equivalent to $\sum_{n \leq x} \frac{\mu(n)}{\sqrt{n}}\ll1$; the forward direction follows from partial summation, while the converse uses a Tauberian theorem due to Riesz.}
      \field{issn}{0034-7760}
      \field{journaltitle}{Rev. Ci. (Lima)}
      \field{title}{A note on {M}ertens' hypothesis}
      \field{volume}{50}
      \field{year}{1948}
      \true{nocite}
      \field{pages}{181\bibrangedash 184}
      \range{pages}{4}
      \verb{mrnumber}
      \verb 29414
      \endverb
    \endentry
    \entry{MR98723}{article}{}
      \name{author}{1}{}{%
        {{hash=86b5f5639455f1050b0d4e2301179477}{%
           family={Wintner},
           familyi={W\bibinitperiod},
           given={Aurel},
           giveni={A\bibinitperiod}}}%
      }
      \strng{namehash}{86b5f5639455f1050b0d4e2301179477}
      \strng{fullhash}{86b5f5639455f1050b0d4e2301179477}
      \strng{bibnamehash}{86b5f5639455f1050b0d4e2301179477}
      \strng{authorbibnamehash}{86b5f5639455f1050b0d4e2301179477}
      \strng{authornamehash}{86b5f5639455f1050b0d4e2301179477}
      \strng{authorfullhash}{86b5f5639455f1050b0d4e2301179477}
      \field{extraname}{5}
      \field{sortinit}{W}
      \field{sortinithash}{1af34bd8c148ffb32de1494636b49713}
      \field{labelprefix}{A}
      \field{labelnamesource}{author}
      \field{labeltitlesource}{title}
      \field{annotation}{The author considers an analogue of the weak Mertens conjecture for the summatory function of the Liouville function. The author shows that $L(x)\ll\sqrt{x}$ is equivalent to ${\sum_{n \leq x} {\lambda(n)}/{\sqrt{n}}=c\log x+O(1)}$, where $c=1/\zeta(1/2)<0$. The proof follows the similar lines as of the author's earlier work~\cite{MR29414}.\par This article cites~\cite{MR0006202,MR29414}.}
      \field{issn}{0002-9327}
      \field{journaltitle}{Amer. J. Math.}
      \field{title}{On the {$\lambda $}-variant of {M}ertens' {$\mu $}-hypothesis}
      \field{volume}{80}
      \field{year}{1958}
      \true{nocite}
      \field{pages}{639\bibrangedash 642}
      \range{pages}{4}
      \verb{mrnumber}
      \verb 98723
      \endverb
    \endentry
  \enddatalist
  \datalist[entry]{ynt/global//global/global}
    \entry{1853.Chebyshev}{article}{}
      \name{author}{1}{}{%
        {{hash=0fa981e305ea35073dc938b976cda0e1}{%
           family={Chebyshev},
           familyi={C\bibinitperiod},
           given={P.},
           giveni={P\bibinitperiod}}}%
      }
      \list{language}{1}{%
        {French}%
      }
      \strng{namehash}{0fa981e305ea35073dc938b976cda0e1}
      \strng{fullhash}{0fa981e305ea35073dc938b976cda0e1}
      \strng{bibnamehash}{0fa981e305ea35073dc938b976cda0e1}
      \strng{authorbibnamehash}{0fa981e305ea35073dc938b976cda0e1}
      \strng{authornamehash}{0fa981e305ea35073dc938b976cda0e1}
      \strng{authorfullhash}{0fa981e305ea35073dc938b976cda0e1}
      \field{sortinit}{1}
      \field{sortinithash}{50c6687d7fc80f50136d75228e3c59ba}
      \field{labelnamesource}{author}
      \field{labeltitlesource}{title}
      \field{annotation}{In this letter, the author remarks for the first time that there appear to be more primes of the form $4n+3$ than $4n+1$, in that their counting functions ``differ notably in their second terms'' (original French: {\it ``different notablement entre elles par leurs seconds termes''}). Several assertions are made (without proof): that $E^\pi(x;4,3,1)$ takes values arbitrarily close to $1$; that $\pi_e(x,\chi_{-4})\to-\infty$ as $x\to\infty$; and that if $f(x)$ is a decreasing function with $\lim_{x\to\infty} x^{1/2} f(x) \ne 0$, then the series $\sum_p \chi_{-4}(p) f(p)$ diverges.}
      \field{journaltitle}{Bull. de la Classe phys. math. de l'Acad. Imp. des Sciences St. Petersburg}
      \field{title}{{Lettre de M. le professeur Tchébychev a M. Fuss, sur un nouveau théorème rélatif aux nombres premiers contenus dans la formes $4n+1$ et $4n+3$}}
      \field{volume}{11}
      \field{year}{1853}
      \true{nocite}
      \field{pages}{208}
      \range{pages}{1}
    \endentry
    \entry{1884.Piltz}{article}{}
      \name{author}{1}{}{%
        {{hash=d60e0007aa708cf28f344f23b440f29d}{%
           family={Piltz},
           familyi={P\bibinitperiod},
           given={A.},
           giveni={A\bibinitperiod}}}%
      }
      \list{language}{1}{%
        {German}%
      }
      \strng{namehash}{d60e0007aa708cf28f344f23b440f29d}
      \strng{fullhash}{d60e0007aa708cf28f344f23b440f29d}
      \strng{bibnamehash}{d60e0007aa708cf28f344f23b440f29d}
      \strng{authorbibnamehash}{d60e0007aa708cf28f344f23b440f29d}
      \strng{authornamehash}{d60e0007aa708cf28f344f23b440f29d}
      \strng{authorfullhash}{d60e0007aa708cf28f344f23b440f29d}
      \field{sortinit}{1}
      \field{sortinithash}{50c6687d7fc80f50136d75228e3c59ba}
      \field{labelnamesource}{author}
      \field{labeltitlesource}{title}
      \field{annotation}{Several authors cite this habilitation thesis as the first appearance of GRH, the generalized Riemann hypothesis for Dirichlet $L$-functions (``Riemann--Piltz conjecture''). On page 25 the author writes, ``{\it Wie in der Einleitung erwähnt wurde, hatte Riemann die Vermuthung, dass diese Verschwindungstellen was die Funktion $\zeta(s)$ betrifft, sich sämmtlich durch die Form $1/2+\alpha i$ wo $\alpha$ reell ist, bringen lassen. Dieser Satz gilt nicht nur für die Funktion $\zeta(s)$ sondern auch für die $\zeta(s,\nu)$.}''}
      \field{journaltitle}{Habilitationsschrift, Friedrich--Schiller--Universität Jena}
      \field{title}{Über die Häufigkeit der Primzahlen in arithmetischen Progressionen und über verwandte Gesetze}
      \field{year}{1884}
      \true{nocite}
    \endentry
    \entry{1891.Phragmen}{article}{}
      \name{author}{1}{}{%
        {{hash=24a636b82c755af93afc8c09bfe2de5f}{%
           family={Phragmén},
           familyi={P\bibinitperiod},
           given={P.},
           giveni={P\bibinitperiod}}}%
      }
      \list{language}{1}{%
        {French}%
      }
      \strng{namehash}{24a636b82c755af93afc8c09bfe2de5f}
      \strng{fullhash}{24a636b82c755af93afc8c09bfe2de5f}
      \strng{bibnamehash}{24a636b82c755af93afc8c09bfe2de5f}
      \strng{authorbibnamehash}{24a636b82c755af93afc8c09bfe2de5f}
      \strng{authornamehash}{24a636b82c755af93afc8c09bfe2de5f}
      \strng{authorfullhash}{24a636b82c755af93afc8c09bfe2de5f}
      \field{sortinit}{1}
      \field{sortinithash}{50c6687d7fc80f50136d75228e3c59ba}
      \field{labelnamesource}{author}
      \field{labeltitlesource}{title}
      \field{annotation}{The author establishes a general proposition (similar to the eventual ``Landau's theorem'') capable of establishing that particular functions change sign for arbitrarily large arguments. From this proposition, the author shows that $\Pi(x) - (\li(x)-\log2)$ changes sign infinitely often, as do the differences $\psi(x) - (x-\log\frac{\pi}2)$ and $\Pi^*_r(x) - (\log\log x+C_0)$ and $\psi_r(x) - (\log x - C_0)$. Also, each of the differences $\Pi_r(x;4,1) - \big( \frac12\Li(x) - \frac12\log \frac{\log x}{\log 2} - \log 2 \big)$ and $\Pi_r(x;4,3) - \big( \frac12\Li(x) - \frac12\log \frac{\log x}{\log 2} \big)$ and $\Pi_r(x;4,1,3) + \log 2$ changes signs infinitely often. Finally, the author establishes the assertion of Chebyshev that~$1$ is a limit point of the function $E^\pi(x;4,3,1)$. \par This article cites~\cites{1853.Chebyshev}.}
      \field{journaltitle}{Öfversigt af Kongl. Vetenskaps--Akademiens Föhandlingar.}
      \field{title}{{Sur le logarithme intégral et la fonction $f(x)$ de Riemann}}
      \field{volume}{48}
      \field{year}{1891}
      \true{nocite}
      \field{pages}{599\bibrangedash 616}
      \range{pages}{18}
    \endentry
    \entry{1897.Mertens}{article}{}
      \name{author}{1}{}{%
        {{hash=a2b41763522272edb01791c27c2476d9}{%
           family={Mertens},
           familyi={M\bibinitperiod},
           given={F.},
           giveni={F\bibinitperiod}}}%
      }
      \list{language}{1}{%
        {German}%
      }
      \strng{namehash}{a2b41763522272edb01791c27c2476d9}
      \strng{fullhash}{a2b41763522272edb01791c27c2476d9}
      \strng{bibnamehash}{a2b41763522272edb01791c27c2476d9}
      \strng{authorbibnamehash}{a2b41763522272edb01791c27c2476d9}
      \strng{authornamehash}{a2b41763522272edb01791c27c2476d9}
      \strng{authorfullhash}{a2b41763522272edb01791c27c2476d9}
      \field{sortinit}{1}
      \field{sortinithash}{50c6687d7fc80f50136d75228e3c59ba}
      \field{labelnamesource}{author}
      \field{labeltitlesource}{title}
      \field{annotation}{The author publishes a table of the values of $\mu(n)$ and $M(n)$ for $n\le10,000$, using the formula $\sum_{d=1}^{\lfloor\sqrt n\rfloor} \big( \mu(d) \lfloor \frac nd \rfloor + M(\frac nd) \big) = \lfloor\sqrt n\rfloor M(\lfloor\sqrt n\rfloor) + 1$ to check the output. Based on the table, the author conjectures that $|M(x)|< \sqrt{x}$ for all $x>1$. Using the formula $\psi(x)=-\sum_{n \leq x} \mu(n) \lfloor \frac xn \rfloor \log n$, the author shows that $M(x)\ll\sqrt x$ would imply $\Delta^\psi(x) \ll x^{3/4} \log x$ and $\Delta^\Pi(x) \ll x^{3/4}$ as well as RH, and that $M(x)\ll\sqrt x/(\log x)^{1+\delta}$ for some $\delta>0$ would imply that $\zeta(s)$ has no nontrivial zeros.}
      \field{journaltitle}{Sitzungsberichte Akad. Wien}
      \field{title}{Über eine zahlentheoretische Funktion}
      \field{volume}{106}
      \field{year}{1897}
      \true{nocite}
      \field{pages}{761\bibrangedash 830}
      \range{pages}{70}
    \endentry
    \entry{1897.vonSterneck}{article}{}
      \name{author}{1}{}{%
        {{hash=f4689d7596dd95e869b4ea317492a6f8}{%
           family={Sterneck},
           familyi={S\bibinitperiod},
           given={R.\bibnamedelimi D.},
           giveni={R\bibinitperiod\bibinitdelim D\bibinitperiod},
           prefix={von},
           prefixi={v\bibinitperiod}}}%
      }
      \list{language}{1}{%
        {German}%
      }
      \strng{namehash}{f4689d7596dd95e869b4ea317492a6f8}
      \strng{fullhash}{f4689d7596dd95e869b4ea317492a6f8}
      \strng{bibnamehash}{f4689d7596dd95e869b4ea317492a6f8}
      \strng{authorbibnamehash}{f4689d7596dd95e869b4ea317492a6f8}
      \strng{authornamehash}{f4689d7596dd95e869b4ea317492a6f8}
      \strng{authorfullhash}{f4689d7596dd95e869b4ea317492a6f8}
      \field{extraname}{1}
      \field{sortinit}{1}
      \field{sortinithash}{50c6687d7fc80f50136d75228e3c59ba}
      \field{labelnamesource}{author}
      \field{labeltitlesource}{title}
      \field{annotation}{The author provides a table of values for $M(x)$. From the table, the author observes that $|M(n)|< \frac{1}{2} \sqrt{n}$ for $201 \leq n \leq 150,000$, and conjectures that $|M(x)|<\sqrt{x}$ for all $x>1$. The author also compares $M(x)$ with a random walk at the squarefree numbers, commenting that the corresponding expectation of the absolute value of the random walk, namely $\sqrt{\frac{12}{\pi^2}}\sqrt{x} \approx 1.10266\sqrt{x}$, exceeds even the maximal values of $|M(x)|$ from the table. \par This article cites~\cites{1897.Mertens}.}
      \field{journaltitle}{Sitzungsberichte Akad. Wiss. Wien IIa}
      \field{title}{Empirische Untersuchung über den Verlauf der zahlentheoretischen Funktion $\sigma(n) = \sum_{x=1}^{x=n}\mu(x)$ im Intervalle von $0$ bis $150000$}
      \field{volume}{106}
      \field{year}{1897}
      \true{nocite}
      \field{pages}{835\bibrangedash 1024}
      \range{pages}{190}
    \endentry
    \entry{1901.vonSterneck}{article}{}
      \name{author}{1}{}{%
        {{hash=f4689d7596dd95e869b4ea317492a6f8}{%
           family={Sterneck},
           familyi={S\bibinitperiod},
           given={R.\bibnamedelimi D.},
           giveni={R\bibinitperiod\bibinitdelim D\bibinitperiod},
           prefix={von},
           prefixi={v\bibinitperiod}}}%
      }
      \list{language}{1}{%
        {German}%
      }
      \strng{namehash}{f4689d7596dd95e869b4ea317492a6f8}
      \strng{fullhash}{f4689d7596dd95e869b4ea317492a6f8}
      \strng{bibnamehash}{f4689d7596dd95e869b4ea317492a6f8}
      \strng{authorbibnamehash}{f4689d7596dd95e869b4ea317492a6f8}
      \strng{authornamehash}{f4689d7596dd95e869b4ea317492a6f8}
      \strng{authorfullhash}{f4689d7596dd95e869b4ea317492a6f8}
      \field{extraname}{2}
      \field{sortinit}{1}
      \field{sortinithash}{50c6687d7fc80f50136d75228e3c59ba}
      \field{labelnamesource}{author}
      \field{labeltitlesource}{title}
      \field{annotation}{The author verifies that $|M(n)| < \frac12\sqrt n$ for $201 \le n\le500,000$. \par This article cites~\cites{1897.vonSterneck}.}
      \field{journaltitle}{SBer. Kais. Akad. Wissensch. Wien IIa}
      \field{title}{Empirische Untersuchung über den Verlauf der zahlentheoretischen Funktion $\sigma(n) = \sum_{x=1}^{x=n}\mu(x)$ im Intervalle von $150000$ bis $500000$}
      \field{volume}{110}
      \field{year}{1901}
      \true{nocite}
      \field{pages}{1053\bibrangedash 1102}
      \range{pages}{50}
    \endentry
    \entry{MR1511206}{article}{}
      \name{author}{1}{}{%
        {{hash=2c1103066f165b2530a1d823ff4f4b10}{%
           family={Schmidt},
           familyi={S\bibinitperiod},
           given={Erhard},
           giveni={E\bibinitperiod}}}%
      }
      \list{language}{1}{%
        {German}%
      }
      \strng{namehash}{2c1103066f165b2530a1d823ff4f4b10}
      \strng{fullhash}{2c1103066f165b2530a1d823ff4f4b10}
      \strng{bibnamehash}{2c1103066f165b2530a1d823ff4f4b10}
      \strng{authorbibnamehash}{2c1103066f165b2530a1d823ff4f4b10}
      \strng{authornamehash}{2c1103066f165b2530a1d823ff4f4b10}
      \strng{authorfullhash}{2c1103066f165b2530a1d823ff4f4b10}
      \field{sortinit}{1}
      \field{sortinithash}{50c6687d7fc80f50136d75228e3c59ba}
      \field{labelnamesource}{author}
      \field{labeltitlesource}{title}
      \field{annotation}{The author shows that $\Delta^\Pi(x) = \Omega_\pm(x^{\Theta-\ep})$ for every~$\ep>0$, and that both $E^\Pi(x) > \frac1{29}$ and $E^\Pi(x) < -\frac1{29}$ occur infinitely often. He also shows the same results with~$\Pi(x)$ in the definitions of~$\Delta^\Pi$ and~$E^\Pi$ replaced by $\pi(x)+\frac12\Li(\sqrt x)$.}
      \field{issn}{0025-5831}
      \field{journaltitle}{Math. Ann.}
      \field{number}{2}
      \field{title}{Über die {A}nzahl der {P}rimzahlen unter gegebener {G}renze}
      \field{volume}{57}
      \field{year}{1903}
      \true{nocite}
      \field{pages}{195\bibrangedash 204}
      \range{pages}{10}
      \verb{doi}
      \verb 10.1007/BF01444344
      \endverb
      \verb{mrnumber}
      \verb 1511206
      \endverb
      \verb{urlraw}
      \verb https://doi.org/10.1007/BF01444344
      \endverb
      \verb{url}
      \verb https://doi.org/10.1007/BF01444344
      \endverb
    \endentry
    \entry{1905.Stieltjes}{book}{}
      \name{author}{1}{}{%
        {{hash=39453b92c259f9e27936958ace288ccd}{%
           family={Stieltjes},
           familyi={S\bibinitperiod},
           given={T.\bibnamedelimi J.},
           giveni={T\bibinitperiod\bibinitdelim J\bibinitperiod}}}%
      }
      \list{language}{1}{%
        {French}%
      }
      \list{publisher}{1}{%
        {Gauthier--Villars, Imprimeur--Libraire, Paris}%
      }
      \strng{namehash}{39453b92c259f9e27936958ace288ccd}
      \strng{fullhash}{39453b92c259f9e27936958ace288ccd}
      \strng{bibnamehash}{39453b92c259f9e27936958ace288ccd}
      \strng{authorbibnamehash}{39453b92c259f9e27936958ace288ccd}
      \strng{authornamehash}{39453b92c259f9e27936958ace288ccd}
      \strng{authorfullhash}{39453b92c259f9e27936958ace288ccd}
      \field{sortinit}{1}
      \field{sortinithash}{50c6687d7fc80f50136d75228e3c59ba}
      \field{labelnamesource}{author}
      \field{labeltitlesource}{title}
      \field{annotation}{In Lettre~79 (starting on page~160), the author states the conjectural bound $\big| \sum_{n\le x} \mu(n) \big| \le \sqrt x$ for the Mertens function.}
      \field{title}{Correspondance d'Hermite et de Stieltjes}
      \field{year}{1905}
      \true{nocite}
      \field{pages}{xxi+pp. 1\bibrangedash 477}
      \range{pages}{-1}
    \endentry
    \entry{MR1511360}{article}{}
      \name{author}{1}{}{%
        {{hash=45dbea41eb1e7f07fcfe10d098ce7670}{%
           family={Landau},
           familyi={L\bibinitperiod},
           given={E.},
           giveni={E\bibinitperiod}}}%
      }
      \list{language}{1}{%
        {German}%
      }
      \strng{namehash}{45dbea41eb1e7f07fcfe10d098ce7670}
      \strng{fullhash}{45dbea41eb1e7f07fcfe10d098ce7670}
      \strng{bibnamehash}{45dbea41eb1e7f07fcfe10d098ce7670}
      \strng{authorbibnamehash}{45dbea41eb1e7f07fcfe10d098ce7670}
      \strng{authornamehash}{45dbea41eb1e7f07fcfe10d098ce7670}
      \strng{authorfullhash}{45dbea41eb1e7f07fcfe10d098ce7670}
      \field{extraname}{1}
      \field{sortinit}{1}
      \field{sortinithash}{50c6687d7fc80f50136d75228e3c59ba}
      \field{labelnamesource}{author}
      \field{labeltitlesource}{title}
      \field{annotation}{Phragmén~\cite{1891.Phragmen} proved Chebyshev's claim that~$1$ is a limit point of the function $E^\pi(x;4,3,1)$, using the theory of functions of complex arguments. The author of this article gives a simpler proof based on the theory of Dirichlet series.\par This article cites~\cite{MR1511206}.}
      \field{issn}{0025-5831}
      \field{journaltitle}{Math. Ann.}
      \field{number}{4}
      \field{title}{Über einen {S}atz von {T}schebyschef}
      \field{volume}{61}
      \field{year}{1906}
      \true{nocite}
      \field{pages}{527\bibrangedash 550}
      \range{pages}{24}
      \verb{doi}
      \verb 10.1007/BF01449495
      \endverb
      \verb{mrnumber}
      \verb 1511360
      \endverb
      \verb{urlraw}
      \verb https://doi.org/10.1007/BF01449495
      \endverb
      \verb{url}
      \verb https://doi.org/10.1007/BF01449495
      \endverb
    \endentry
    \entry{1909.Landau}{book}{}
      \name{author}{1}{}{%
        {{hash=45dbea41eb1e7f07fcfe10d098ce7670}{%
           family={Landau},
           familyi={L\bibinitperiod},
           given={E.},
           giveni={E\bibinitperiod}}}%
      }
      \list{language}{1}{%
        {German}%
      }
      \list{publisher}{1}{%
        {Leipzig und Berlin, B. G. Teubner}%
      }
      \strng{namehash}{45dbea41eb1e7f07fcfe10d098ce7670}
      \strng{fullhash}{45dbea41eb1e7f07fcfe10d098ce7670}
      \strng{bibnamehash}{45dbea41eb1e7f07fcfe10d098ce7670}
      \strng{authorbibnamehash}{45dbea41eb1e7f07fcfe10d098ce7670}
      \strng{authornamehash}{45dbea41eb1e7f07fcfe10d098ce7670}
      \strng{authorfullhash}{45dbea41eb1e7f07fcfe10d098ce7670}
      \field{extraname}{2}
      \field{sortinit}{1}
      \field{sortinithash}{50c6687d7fc80f50136d75228e3c59ba}
      \field{labelnamesource}{author}
      \field{labeltitlesource}{title}
      \field{annotation}{In Sections~199 and~200, the author generalizes his proof from~\cite{MR1511360}, showing that $(c_q(b)-c_q(a))/\phi(q)$ is a limit point of the function $E^\pi(x;q,a,b)$. He also shows that if $c>\sqrt{1/4+\gamma_1^2}$, then both $E^\Pi(x) > \frac1{29}$ and $E^\Pi(x) < -\frac1{29}$ occur infinitely often, as well as the same result with~$\Pi(x)$ in the definition of~$E^\Pi$ replaced by $\pi(x)+\frac12\Li(\sqrt x)$.}
      \field{title}{Handbuch der Lehre von der Verteilung der Primzahlen. 2 Bände}
      \field{year}{1909}
      \true{nocite}
      \field{pages}{xviii+pp. 1\bibrangedash 564\bibrangessep ix+pp. 565\bibrangedash 961}
      \range{pages}{-1}
    \endentry
    \entry{1912.vonSterneck}{article}{}
      \name{author}{1}{}{%
        {{hash=f4689d7596dd95e869b4ea317492a6f8}{%
           family={Sterneck},
           familyi={S\bibinitperiod},
           given={R.\bibnamedelimi D.},
           giveni={R\bibinitperiod\bibinitdelim D\bibinitperiod},
           prefix={von},
           prefixi={v\bibinitperiod}}}%
      }
      \list{language}{1}{%
        {German}%
      }
      \strng{namehash}{f4689d7596dd95e869b4ea317492a6f8}
      \strng{fullhash}{f4689d7596dd95e869b4ea317492a6f8}
      \strng{bibnamehash}{f4689d7596dd95e869b4ea317492a6f8}
      \strng{authorbibnamehash}{f4689d7596dd95e869b4ea317492a6f8}
      \strng{authornamehash}{f4689d7596dd95e869b4ea317492a6f8}
      \strng{authorfullhash}{f4689d7596dd95e869b4ea317492a6f8}
      \field{extraname}{3}
      \field{sortinit}{1}
      \field{sortinithash}{50c6687d7fc80f50136d75228e3c59ba}
      \field{labelnamesource}{author}
      \field{labeltitlesource}{title}
      \field{annotation}{The author calculates $16$ selected values of $M(n)$ with $600,000 \le n \le 5,000,000$ and verifies they all satisfy $|M(n)| \le \frac12\sqrt n$. The computation uses the following formula, which is a refined version of a formula of Mertens~\cite{1897.Mertens}: for $j=1,2,3,4$, if $n$ exceeds the product of the first~$j$ primes, then $ \sum_{d=1}^{\lfloor \sqrt{n} \rfloor} \mu(d) \omega_j(n/d) +\sum_{d'} M(n/d) = \omega_j\bigl( \lfloor \sqrt{n} \rfloor \bigr) M\bigl( \lfloor \sqrt{n} \rfloor \bigr), $ where $\omega_j(x)$ is the number of positive integers $\leq x$ that are not divisible by any of the first~$j$ primes, and $d'$ runs over all such numbers $\leq g$. The author claims that the bound $|M(n)| \le \frac12\sqrt n$ for all $n>200$ represents an unproved but extremely probable number theoretic law, and thus RH could be also regarded as correct with a high degree of probability (original German: ``{\it $|M(n)| \le \frac12\sqrt n$ zwar ein unbewiesenes, aber ausserordentlich wahrscheinliches zahlentheoretisches Gesetz darstellt, und somit auch die Riemann'sche Vermutung mit einem hohen Grad von Wahrscheinlichkeit als richtig angesehen werden kann}''). \par This article cites~\cites{1897.Mertens,1897.vonSterneck,1901.vonSterneck}.}
      \field{journaltitle}{SBer. Kais. Akad. Wissensch. Wien IIa}
      \field{title}{Die zahlentheoretische Funktion $\sigma(n)$ bis zur Grenze $5000000$}
      \field{volume}{121}
      \field{year}{1912}
      \true{nocite}
      \field{pages}{1083\bibrangedash 1096}
      \range{pages}{14}
    \endentry
    \entry{1913.vonSterneck}{article}{}
      \name{author}{1}{}{%
        {{hash=f4689d7596dd95e869b4ea317492a6f8}{%
           family={Sterneck},
           familyi={S\bibinitperiod},
           given={R.\bibnamedelimi D.},
           giveni={R\bibinitperiod\bibinitdelim D\bibinitperiod},
           prefix={von},
           prefixi={v\bibinitperiod}}}%
      }
      \list{language}{1}{%
        {German}%
      }
      \list{publisher}{1}{%
        {Cambridge University Press}%
      }
      \strng{namehash}{f4689d7596dd95e869b4ea317492a6f8}
      \strng{fullhash}{f4689d7596dd95e869b4ea317492a6f8}
      \strng{bibnamehash}{f4689d7596dd95e869b4ea317492a6f8}
      \strng{authorbibnamehash}{f4689d7596dd95e869b4ea317492a6f8}
      \strng{authornamehash}{f4689d7596dd95e869b4ea317492a6f8}
      \strng{authorfullhash}{f4689d7596dd95e869b4ea317492a6f8}
      \field{extraname}{4}
      \field{sortinit}{1}
      \field{sortinithash}{50c6687d7fc80f50136d75228e3c59ba}
      \field{labelnamesource}{author}
      \field{labeltitlesource}{title}
      \field{annotation}{This is a summary of the previous three articles of the author. \par This article cites~\cites{1897.Mertens,1897.vonSterneck, 1901.vonSterneck,1912.vonSterneck}.}
      \field{journaltitle}{Proc. 5th International Congress of Mathematicians}
      \field{title}{Neue empirische Daten über die zahlentheoretische Funktion $\sigma(n)$}
      \field{volume}{1}
      \field{year}{1913}
      \true{nocite}
      \field{pages}{341\bibrangedash 343}
      \range{pages}{3}
    \endentry
    \entry{1914.Littlewood}{article}{}
      \name{author}{1}{}{%
        {{hash=56dd48151c55dce1b8a073ecf909a238}{%
           family={Littlewood},
           familyi={L\bibinitperiod},
           given={J.\bibnamedelimi E.},
           giveni={J\bibinitperiod\bibinitdelim E\bibinitperiod}}}%
      }
      \list{language}{1}{%
        {French}%
      }
      \strng{namehash}{56dd48151c55dce1b8a073ecf909a238}
      \strng{fullhash}{56dd48151c55dce1b8a073ecf909a238}
      \strng{bibnamehash}{56dd48151c55dce1b8a073ecf909a238}
      \strng{authorbibnamehash}{56dd48151c55dce1b8a073ecf909a238}
      \strng{authornamehash}{56dd48151c55dce1b8a073ecf909a238}
      \strng{authorfullhash}{56dd48151c55dce1b8a073ecf909a238}
      \field{extraname}{1}
      \field{sortinit}{1}
      \field{sortinithash}{50c6687d7fc80f50136d75228e3c59ba}
      \field{labelnamesource}{author}
      \field{labeltitlesource}{title}
      \field{annotation}{Assuming RH, the author shows that \begin{align*} \Delta^\psi(x) &= \Omega_\pm\big( \sqrt{x}\log\log\log{x} \big) \\ \Delta^\pi(x) &= \Omega_\pm\bigg( \frac{\sqrt{x}\log\log\log{x}}{\log{x}} \bigg). \end{align*} It follows that the conjectured inequality $\pi(x)<\Li(x)$ (``{\em présumée par divers auteurs pour des raisons empiriques}'') cannot hold for all values of~$x$.\par This article cites~\cite{1909.Landau}.}
      \field{journaltitle}{Comptes Rendus de l'Acad. Sci. Paris}
      \field{title}{{Sur la distribution des nombres premiers}}
      \field{volume}{158}
      \field{year}{1914}
      \true{nocite}
      \field{pages}{1869\bibrangedash 1872}
      \range{pages}{4}
    \endentry
    \entry{1915.HardyLittlewood}{article}{}
      \name{author}{2}{}{%
        {{hash=9d761c354d534c74e69f17afcc78fea8}{%
           family={Hardy},
           familyi={H\bibinitperiod},
           given={G.\bibnamedelimi H.},
           giveni={G\bibinitperiod\bibinitdelim H\bibinitperiod}}}%
        {{hash=56dd48151c55dce1b8a073ecf909a238}{%
           family={Littlewood},
           familyi={L\bibinitperiod},
           given={J.\bibnamedelimi E.},
           giveni={J\bibinitperiod\bibinitdelim E\bibinitperiod}}}%
      }
      \strng{namehash}{ce5cb2e2cdf2e499e451b481213bdb97}
      \strng{fullhash}{ce5cb2e2cdf2e499e451b481213bdb97}
      \strng{bibnamehash}{ce5cb2e2cdf2e499e451b481213bdb97}
      \strng{authorbibnamehash}{ce5cb2e2cdf2e499e451b481213bdb97}
      \strng{authornamehash}{ce5cb2e2cdf2e499e451b481213bdb97}
      \strng{authorfullhash}{ce5cb2e2cdf2e499e451b481213bdb97}
      \field{extraname}{1}
      \field{sortinit}{1}
      \field{sortinithash}{50c6687d7fc80f50136d75228e3c59ba}
      \field{labelnamesource}{author}
      \field{labeltitlesource}{title}
      \field{annotation}{A letter of Tchebychef asserts that $\lim_{x\to\infty} \pi_e(x) = -\infty$. The authors verify the conjecture assuming GRH for $L(s,\chi_{-4})$. Using a formula of Cahen and Mellin and an application of Cauchy's theorem, the authors show that $-\pi_e(x)\gg \sqrt{x/\log x}$.}
      \field{journaltitle}{Proc. London Math. Soc. (2)}
      \field{title}{On an assertion of {T}chebychef}
      \field{volume}{14}
      \field{year}{1915}
      \true{nocite}
      \field{pages}{xv\bibrangedash xvi}
      \range{pages}{2}
    \endentry
    \entry{MR1576550}{article}{}
      \name{author}{1}{}{%
        {{hash=9d761c354d534c74e69f17afcc78fea8}{%
           family={Hardy},
           familyi={H\bibinitperiod},
           given={G.\bibnamedelimi H.},
           giveni={G\bibinitperiod\bibinitdelim H\bibinitperiod}}}%
      }
      \strng{namehash}{9d761c354d534c74e69f17afcc78fea8}
      \strng{fullhash}{9d761c354d534c74e69f17afcc78fea8}
      \strng{bibnamehash}{9d761c354d534c74e69f17afcc78fea8}
      \strng{authorbibnamehash}{9d761c354d534c74e69f17afcc78fea8}
      \strng{authornamehash}{9d761c354d534c74e69f17afcc78fea8}
      \strng{authorfullhash}{9d761c354d534c74e69f17afcc78fea8}
      \field{sortinit}{1}
      \field{sortinithash}{50c6687d7fc80f50136d75228e3c59ba}
      \field{labelnamesource}{author}
      \field{labeltitlesource}{title}
      \field{annotation}{This article begins with a short background on the history of the Dirichlet divisor problem and states the best bounds for $\Delta^D(x)$ available at the time, namely Voronoi's result $\Delta^D(x) \ll x^{1/3}\log x$. The author then proves that $\Delta^D(x) = \Omega_\pm(x^{1/4})$. He draws a direct comparison to Schmidt's proof~\cite{MR1511206} that $\Delta^\Pi(x) = \Omega_\pm(\sqrt x / \log x)$, and to a simplified version of Schmidt's proof by Landau~\cite{1909.Landau}; he states that his proof for $\Delta^D$ does not ``differ in principle'' from Landau's proof for $\Delta^\Pi$, and for the sake of comparison provides a simplified form of the proof that $\Delta^\psi(x) = \Omega_\pm(\sqrt x)$ based on Landau's methods. \par The article then discusses various generalizations of the problem. The first generalization discussed, due to Piltz, concerns the same problem for $D_k(x) = \sum_{n \le x} \tau_k(n)$, where $\tau_k(n)$ denotes the number of ways~$n$ can be decomposed into~$k$ ordered factors: Hardy asserts that the methods in this article can show that $\Delta^{D_k}(x) = \Omega_\pm(x^{(k-1)/(2k)})$. The second generalization concerns $R(x) = \sum_{n \le x} r(n)$, where $r(n)$ counts the number of ways~$n$ can be represented as the sum of two squares: Hardy likewise asserts that $\Delta^R(x) = \Omega_\pm(x^{1/4})$. Finally, the article concludes with a discussion of an explicit formula for $D(x)$, first found by Voronoi, and some alternative proof methods and their comparative advantages and disadvantages. \par This article cites~\cite{1914.Littlewood, 1909.Landau}.}
      \field{issn}{0024-6115}
      \field{journaltitle}{Proc. London Math. Soc. (2)}
      \field{title}{On {D}irichlet's {D}ivisor {P}roblem}
      \field{volume}{15}
      \field{year}{1916}
      \true{nocite}
      \field{pages}{1\bibrangedash 25}
      \range{pages}{25}
      \verb{doi}
      \verb 10.1112/plms/s2-15.1.1
      \endverb
      \verb{mrnumber}
      \verb 1576550
      \endverb
      \verb{urlraw}
      \verb https://doi-org.ezproxy.library.ubc.ca/10.1112/plms/s2-15.1.1
      \endverb
      \verb{url}
      \verb https://doi-org.ezproxy.library.ubc.ca/10.1112/plms/s2-15.1.1
      \endverb
    \endentry
    \entry{MR1555148}{article}{}
      \name{author}{2}{}{%
        {{hash=9d761c354d534c74e69f17afcc78fea8}{%
           family={Hardy},
           familyi={H\bibinitperiod},
           given={G.\bibnamedelimi H.},
           giveni={G\bibinitperiod\bibinitdelim H\bibinitperiod}}}%
        {{hash=56dd48151c55dce1b8a073ecf909a238}{%
           family={Littlewood},
           familyi={L\bibinitperiod},
           given={J.\bibnamedelimi E.},
           giveni={J\bibinitperiod\bibinitdelim E\bibinitperiod}}}%
      }
      \strng{namehash}{ce5cb2e2cdf2e499e451b481213bdb97}
      \strng{fullhash}{ce5cb2e2cdf2e499e451b481213bdb97}
      \strng{bibnamehash}{ce5cb2e2cdf2e499e451b481213bdb97}
      \strng{authorbibnamehash}{ce5cb2e2cdf2e499e451b481213bdb97}
      \strng{authornamehash}{ce5cb2e2cdf2e499e451b481213bdb97}
      \strng{authorfullhash}{ce5cb2e2cdf2e499e451b481213bdb97}
      \field{extraname}{2}
      \field{sortinit}{1}
      \field{sortinithash}{50c6687d7fc80f50136d75228e3c59ba}
      \field{labelnamesource}{author}
      \field{labeltitlesource}{title}
      \field{annotation}{This article contains full proofs of several results that had been announced by (at least one of) the authors in the few years prior.\par In Section 2.2, the authors obtain an explicit formula for the exponentially weighted sum $\psi_e(x) - 1/(e^{1/x}-1) = \sum_{n=1}^\infty (\Lambda(n)-1) e^{-n/x}$; furthermore, assuming RH, they show that this expression is both $\ll \sqrt x$ and $\Omega_\pm(\sqrt x)$. From the latter they deduce that $\Delta^\psi(x) = \Omega_\pm(\sqrt x)$. In Section 2.3, they consider the function $\sum_{p\ge3} (-1)^{(p+1)/2} e^{-p/x} = -\pi_e(x,\chi_{-4})$. Assuming GRH for $L(s,\chi_{-4})$, they prove that $\pi_e(x,\chi_{-4})\to-\infty$ as $x\to\infty$, which is one way of justifying Chebyshev's observation that there are more primes congruent to $3\mod 4$ than to $1\mod 4$.\par In Section 5, the authors provide a full proof of ``Littlewood's theorem'' (announced in~\cite{1914.Littlewood}) on irregularities in the distribution of primes: they prove that \[ \Delta^\pi(x) = \Omega_\pm\bigg( \frac{\sqrt x}{\log x} \log\log\log x \bigg), \] which in particular refutes the conjecture that $\pi(x)<\li(x)$ for all $x>1$.\par Their proof, which begins with the assumption of RH thanks to prior work of Landau~\cite[Sections 201--3]{1909.Landau}, uses homogeneous Diophantine approximation for the imaginary parts of the zeros of $\zeta(s)$. They assert that it can be shown in a similar way that $\psi(x,\chi_{-4}) = \Omega_\pm(\sqrt x\log\log\log x)$ and \[ \pi(x;4,3,1) = \Omega_\pm\bigg( \frac{\sqrt x}{\log x} \log\log\log x \bigg); \] these results are actually dissonant with Chebyshev's observations. (As a side remark, this is also the article in which appears the proof of the asymptotic formula for the second moment of $\zeta(s)$ on the critical line.)\par This article cites~\cites{MR1511206,1909.Landau,1914.Littlewood,1915.HardyLittlewood}.}
      \field{journaltitle}{Acta Math.}
      \field{number}{1}
      \field{title}{{Contributions to the theory of the Riemann zeta-function and the theory of the distribution of primes}}
      \field{volume}{41}
      \field{year}{1916}
      \true{nocite}
      \field{pages}{119\bibrangedash 196}
      \range{pages}{78}
      \verb{mrnumber}
      \verb 1555148
      \endverb
    \endentry
    \entry{MR1544293A}{article}{}
      \name{author}{1}{}{%
        {{hash=45dbea41eb1e7f07fcfe10d098ce7670}{%
           family={Landau},
           familyi={L\bibinitperiod},
           given={E.},
           giveni={E\bibinitperiod}}}%
      }
      \list{language}{1}{%
        {German}%
      }
      \strng{namehash}{45dbea41eb1e7f07fcfe10d098ce7670}
      \strng{fullhash}{45dbea41eb1e7f07fcfe10d098ce7670}
      \strng{bibnamehash}{45dbea41eb1e7f07fcfe10d098ce7670}
      \strng{authorbibnamehash}{45dbea41eb1e7f07fcfe10d098ce7670}
      \strng{authornamehash}{45dbea41eb1e7f07fcfe10d098ce7670}
      \strng{authorfullhash}{45dbea41eb1e7f07fcfe10d098ce7670}
      \field{extraname}{3}
      \field{sortinit}{1}
      \field{sortinithash}{50c6687d7fc80f50136d75228e3c59ba}
      \field{labelnamesource}{author}
      \field{labeltitlesource}{title}
      \field{annotation}{Chebyshev asserted that $\pi_e(x,\chi_{-4})\to-\infty$ as $x\to\infty$; the author shows that this assertion implies GRH for $L(s,\chi_{-4})$. Assuming this GRH, he shows that $\pi(x;4,3,1) \ll \sqrt x\log x$. Assuming the original assertion of Chebyshev, he proves that $\sum_p \chi_{-4}(p)f(p)$ converges whenever $f(p)$ is strictly decreasing and satisfies $f(p) \ll x^{-1/2-\delta}$ for some $\delta>0$.\par The author also examines the series \[ \frac{L'(s,\chi_D)}{L(s,\chi_D)}=-\sum_{p} \frac{\chi_D(p) \log p}{p^s-\chi_D(p)}, \] the identity originally valid for $\Re s>1$; he shows that this series diverges at $s=\frac12$, disproving a conjecture of Lerch. Finally, the author shows that $\sum_{n=1}^\infty \mu(n)/\sqrt n$ diverges, disproving a conjecture of Stieltjes, and also shows that $M(x) = \Omega(\sqrt x)$.\par This article cites~\cite{1853.Chebyshev}.}
      \field{issn}{0025-5874}
      \field{journaltitle}{Math. Z.}
      \field{number}{2-3}
      \field{title}{Über einige ältere {V}ermutungen und {B}ehauptungen in der {P}rimzahltheorie}
      \field{volume}{1}
      \field{year}{1918}
      \true{nocite}
      \field{pages}{1\bibrangedash 24}
      \range{pages}{24}
      \verb{doi}
      \verb 10.1007/BF01203613
      \endverb
      \verb{mrnumber}
      \verb 1544293
      \endverb
      \verb{urlraw}
      \verb https://doi.org/10.1007/BF01203613
      \endverb
      \verb{url}
      \verb https://doi.org/10.1007/BF01203613
      \endverb
    \endentry
    \entry{MR1544293B}{article}{}
      \name{author}{1}{}{%
        {{hash=45dbea41eb1e7f07fcfe10d098ce7670}{%
           family={Landau},
           familyi={L\bibinitperiod},
           given={E.},
           giveni={E\bibinitperiod}}}%
      }
      \list{language}{1}{%
        {German}%
      }
      \strng{namehash}{45dbea41eb1e7f07fcfe10d098ce7670}
      \strng{fullhash}{45dbea41eb1e7f07fcfe10d098ce7670}
      \strng{bibnamehash}{45dbea41eb1e7f07fcfe10d098ce7670}
      \strng{authorbibnamehash}{45dbea41eb1e7f07fcfe10d098ce7670}
      \strng{authornamehash}{45dbea41eb1e7f07fcfe10d098ce7670}
      \strng{authorfullhash}{45dbea41eb1e7f07fcfe10d098ce7670}
      \field{extraname}{4}
      \field{sortinit}{1}
      \field{sortinithash}{50c6687d7fc80f50136d75228e3c59ba}
      \field{labelnamesource}{author}
      \field{labeltitlesource}{title}
      \field{annotation}{The author gives a simplified proof of the result of Hardy and Littlewood~\cite[Section 2.3]{MR1555148} that GRH for $L(s,\chi_{-4})$ implies Chebyshev's assertion $\pi_e(x,\chi_{-4})\to-\infty$ as $x\to\infty$.}
      \field{issn}{0025-5874}
      \field{journaltitle}{Math. Z.}
      \field{number}{2-3}
      \field{title}{Über einige ältere {V}ermutungen und {B}ehauptungen in der {P}rimzahltheorie}
      \field{volume}{1}
      \field{year}{1918}
      \true{nocite}
      \field{pages}{213\bibrangedash 219}
      \range{pages}{7}
      \verb{doi}
      \verb 10.1007/BF01203613
      \endverb
      \verb{mrnumber}
      \verb 1544293
      \endverb
      \verb{urlraw}
      \verb https://doi.org/10.1007/BF01203613
      \endverb
      \verb{url}
      \verb https://doi.org/10.1007/BF01203613
      \endverb
    \endentry
    \entry{1919.Polya}{article}{}
      \name{author}{1}{}{%
        {{hash=85d5bb5658a4749b33270d8f0e8be228}{%
           family={Pólya},
           familyi={P\bibinitperiod},
           given={G.},
           giveni={G\bibinitperiod}}}%
      }
      \list{language}{1}{%
        {German}%
      }
      \strng{namehash}{85d5bb5658a4749b33270d8f0e8be228}
      \strng{fullhash}{85d5bb5658a4749b33270d8f0e8be228}
      \strng{bibnamehash}{85d5bb5658a4749b33270d8f0e8be228}
      \strng{authorbibnamehash}{85d5bb5658a4749b33270d8f0e8be228}
      \strng{authornamehash}{85d5bb5658a4749b33270d8f0e8be228}
      \strng{authorfullhash}{85d5bb5658a4749b33270d8f0e8be228}
      \field{extraname}{1}
      \field{sortinit}{1}
      \field{sortinithash}{50c6687d7fc80f50136d75228e3c59ba}
      \field{labelnamesource}{author}
      \field{labeltitlesource}{title}
      \field{annotation}{In Section~III, the author empirically observes that $L(n)\le 0$ for $2\le n\le1500$, with equality only for $n\in\{2,4,6,10,16,26,40,96,586\}$; he explains some of these equalities as resulting from imaginary quadratic fields $\mathbb Q(\sqrt{-p})$ with class number~$1$. The author also notes that RH would follow from the assertion that $L(n)\le 0$ for all large~$n$. However, he does not formulate this as a conjecture, saying only: ``{\it Ich teile diese Beobachtung mit, um evtl.~weitere numerische Untersuchung zu veranlassen. Der Bewies von [$L(n)\le 0$], sogar nur für hinreichend grosses $n$, würde den Beweis der Riemannschen Vermutung nach sich ziehen\ldots.}''.\par This article cites~\cite{1909.Landau}.}
      \field{journaltitle}{Jahresbericht der deutschen Math.--Vereinigung}
      \field{title}{{Verschiedene Bemerkungen zur Zahlentheorie}}
      \field{volume}{28}
      \field{year}{1919}
      \true{nocite}
      \field{pages}{31\bibrangedash 40}
      \range{pages}{10}
    \endentry
    \entry{MR1544509}{article}{}
      \name{author}{1}{}{%
        {{hash=8815262f8ce833f8804c25a5afc1996b}{%
           family={Cramér},
           familyi={C\bibinitperiod},
           given={Harald},
           giveni={H\bibinitperiod}}}%
      }
      \list{language}{1}{%
        {German}%
      }
      \strng{namehash}{8815262f8ce833f8804c25a5afc1996b}
      \strng{fullhash}{8815262f8ce833f8804c25a5afc1996b}
      \strng{bibnamehash}{8815262f8ce833f8804c25a5afc1996b}
      \strng{authorbibnamehash}{8815262f8ce833f8804c25a5afc1996b}
      \strng{authornamehash}{8815262f8ce833f8804c25a5afc1996b}
      \strng{authorfullhash}{8815262f8ce833f8804c25a5afc1996b}
      \field{sortinit}{1}
      \field{sortinithash}{50c6687d7fc80f50136d75228e3c59ba}
      \field{labelnamesource}{author}
      \field{labeltitlesource}{title}
      \field{annotation}{Under RH, the author shows that \[ \lim_{x \to \infty} \frac{1}{\log x} \int_2^x \left( \frac{\Delta^\psi(t)}{t} \right)^2 \, dt = \sum_\rho \left| \frac{n_\rho}{\rho} \right|^2, \] where $n_\rho$ is the multiplicity of the zero $\rho$. It follows that $\fA_1^\psi(x) \ll x^{3/2}$ (and the same for $\fA_1^\theta(x)$) and thus that $\fA_1^\pi(x) \ll x^{3/2}/\log x$ (and the same for $\fA_1^\Pi(x)$).\par \par This article cites~\cites{MR1555148,1909.Landau}.}
      \field{issn}{0025-5874}
      \field{journaltitle}{Math. Z.}
      \field{number}{1}
      \field{title}{Ein Mittelwertsatz in der Primzahltheorie}
      \field{volume}{12}
      \field{year}{1922}
      \true{nocite}
      \field{pages}{147\bibrangedash 153}
      \range{pages}{7}
      \verb{doi}
      \verb 10.1007/BF01482072
      \endverb
      \verb{mrnumber}
      \verb 1544509
      \endverb
    \endentry
    \entry{MR1574052}{article}{}
      \name{author}{1}{}{%
        {{hash=56dd48151c55dce1b8a073ecf909a238}{%
           family={Littlewood},
           familyi={L\bibinitperiod},
           given={J.\bibnamedelimi E.},
           giveni={J\bibinitperiod\bibinitdelim E\bibinitperiod}}}%
      }
      \strng{namehash}{56dd48151c55dce1b8a073ecf909a238}
      \strng{fullhash}{56dd48151c55dce1b8a073ecf909a238}
      \strng{bibnamehash}{56dd48151c55dce1b8a073ecf909a238}
      \strng{authorbibnamehash}{56dd48151c55dce1b8a073ecf909a238}
      \strng{authornamehash}{56dd48151c55dce1b8a073ecf909a238}
      \strng{authorfullhash}{56dd48151c55dce1b8a073ecf909a238}
      \field{extraname}{2}
      \field{sortinit}{1}
      \field{sortinithash}{50c6687d7fc80f50136d75228e3c59ba}
      \field{labelnamesource}{author}
      \field{labeltitlesource}{title}
      \field{annotation}{The author shows that $\Delta^{\psi}(x)=\Omega_{\pm}(\sqrt{x})$. The main idea is to examine the recursive sequence of integrals $\fA^{\psi}_n(x)$; since $x^{-1}\fA^{\psi}_n(x)$ is an average of $\fA^{\psi}_{n-1}(x)$, it suffices to prove that $\fA^{\psi}_n(x)=\Omega_{\pm}(x^{n+\frac{1}{2}})$ when $n$ is sufficiently large.\par This article cites~\cites{1909.Landau}.}
      \field{journaltitle}{J. London Math. Soc.}
      \field{number}{1}
      \field{title}{Mathematical Notes: 3; on a Theorem Concerning the Distribution of Prime Numbers}
      \field{volume}{2}
      \field{year}{1927}
      \true{nocite}
      \field{pages}{41\bibrangedash 45}
      \range{pages}{5}
      \verb{doi}
      \verb 10.1112/jlms/s1-2.1.41
      \endverb
      \verb{mrnumber}
      \verb 1574052
      \endverb
    \endentry
    \entry{1930.Polya}{article}{}
      \name{author}{1}{}{%
        {{hash=85d5bb5658a4749b33270d8f0e8be228}{%
           family={Pólya},
           familyi={P\bibinitperiod},
           given={G.},
           giveni={G\bibinitperiod}}}%
      }
      \list{language}{1}{%
        {German}%
      }
      \strng{namehash}{85d5bb5658a4749b33270d8f0e8be228}
      \strng{fullhash}{85d5bb5658a4749b33270d8f0e8be228}
      \strng{bibnamehash}{85d5bb5658a4749b33270d8f0e8be228}
      \strng{authorbibnamehash}{85d5bb5658a4749b33270d8f0e8be228}
      \strng{authornamehash}{85d5bb5658a4749b33270d8f0e8be228}
      \strng{authorfullhash}{85d5bb5658a4749b33270d8f0e8be228}
      \field{extraname}{2}
      \field{sortinit}{1}
      \field{sortinithash}{50c6687d7fc80f50136d75228e3c59ba}
      \field{labelnamesource}{author}
      \field{labeltitlesource}{title}
      \field{annotation}{The author proves that $\limsup_{T\to\infty} W^\psi(T)/\log T \ge \gamma_1/\pi$, where $\gamma_1 \approx 14.135$ is the smallest ordinate of a nontrivial zero of $\zeta(s)$. Indeed, $\gamma_1$ can be replaced by the minimum positive ordinate of the zeros of $\zeta(s)$ having maximal real part, or~$\infty$ if no such zero exists. It was known for a long while that this article contained an error that could be mended; the corrected version finally appeared as~\cite{MR0263757}.\par This article cites~\cite{1891.Phragmen,MR1511206}.}
      \field{journaltitle}{Gött. Nachr.}
      \field{title}{{Über das Vorzeichen des Restgliedes im Primzahltheorie}}
      \field{year}{1930}
      \true{nocite}
      \field{pages}{19\bibrangedash 27}
      \range{pages}{9}
    \endentry
    \entry{MR1502996}{article}{}
      \name{author}{2}{}{%
        {{hash=1bdd6511a49b2ac9d9318811c581b790}{%
           family={Evelyn},
           familyi={E\bibinitperiod},
           given={C.\bibnamedelimi J.\bibnamedelimi A.},
           giveni={C\bibinitperiod\bibinitdelim J\bibinitperiod\bibinitdelim A\bibinitperiod}}}%
        {{hash=f88769f16846813993d3053353ecc570}{%
           family={Linfoot},
           familyi={L\bibinitperiod},
           given={E.\bibnamedelimi H.},
           giveni={E\bibinitperiod\bibinitdelim H\bibinitperiod}}}%
      }
      \strng{namehash}{34737a2e61b311a58a63b988e137345a}
      \strng{fullhash}{34737a2e61b311a58a63b988e137345a}
      \strng{bibnamehash}{34737a2e61b311a58a63b988e137345a}
      \strng{authorbibnamehash}{34737a2e61b311a58a63b988e137345a}
      \strng{authornamehash}{34737a2e61b311a58a63b988e137345a}
      \strng{authorfullhash}{34737a2e61b311a58a63b988e137345a}
      \field{sortinit}{1}
      \field{sortinithash}{50c6687d7fc80f50136d75228e3c59ba}
      \field{labelnamesource}{author}
      \field{labeltitlesource}{title}
      \field{annotation}{Using Landau's theorem, the authors show that $\Delta^{Q_k}(x)=\Omega_{\pm}(x^{1/2k})$.\par This article cites~\cites{1909.Landau}.}
      \field{issn}{0003-486X,1939-8980}
      \field{journaltitle}{Ann. of Math. (2)}
      \field{number}{2}
      \field{title}{On a problem in the additive theory of numbers}
      \field{volume}{32}
      \field{year}{1931}
      \true{nocite}
      \field{pages}{261\bibrangedash 270}
      \range{pages}{10}
      \verb{doi}
      \verb 10.2307/1968190
      \endverb
      \verb{mrnumber}
      \verb 1502996
      \endverb
      \verb{urlraw}
      \verb https://doi.org/10.2307/1968190
      \endverb
      \verb{url}
      \verb https://doi.org/10.2307/1968190
      \endverb
    \endentry
    \entry{MR1576856}{article}{}
      \name{author}{1}{}{%
        {{hash=85d5bb5658a4749b33270d8f0e8be228}{%
           family={Pólya},
           familyi={P\bibinitperiod},
           given={G.},
           giveni={G\bibinitperiod}}}%
      }
      \strng{namehash}{85d5bb5658a4749b33270d8f0e8be228}
      \strng{fullhash}{85d5bb5658a4749b33270d8f0e8be228}
      \strng{bibnamehash}{85d5bb5658a4749b33270d8f0e8be228}
      \strng{authorbibnamehash}{85d5bb5658a4749b33270d8f0e8be228}
      \strng{authornamehash}{85d5bb5658a4749b33270d8f0e8be228}
      \strng{authorfullhash}{85d5bb5658a4749b33270d8f0e8be228}
      \field{extraname}{3}
      \field{sortinit}{1}
      \field{sortinithash}{50c6687d7fc80f50136d75228e3c59ba}
      \field{labelnamesource}{author}
      \field{labeltitlesource}{title}
      \field{annotation}{Let $\omega(u)$ be a real-valued function and set $\Phi(s) = \int_1^\infty \omega(u)u^{-s} \, du$. If the number of sign changes $W^\omega(x)$ is $O(\log x)$, then a rightmost singularity~$\sigma+it$ of $\Phi(s)$ (that is, a singularity with $\sigma$ maximal) exists and satisfies $0\le \gamma \le \limsup_{x\to\infty} \pi W^\omega(x)/\log x$. This result generalizes Landau's theorem, which is the case $W^\omega(x)\ll1$.}
      \field{issn}{0024-6115}
      \field{journaltitle}{Proc. London Math. Soc. (2)}
      \field{number}{2}
      \field{title}{On polar singularities of power series and of Dirichlet series}
      \field{volume}{33}
      \field{year}{1931}
      \true{nocite}
      \field{pages}{85\bibrangedash 101}
      \range{pages}{17}
      \verb{doi}
      \verb 10.1112/plms/s2-33.1.85
      \endverb
      \verb{mrnumber}
      \verb 1576856
      \endverb
    \endentry
    \entry{1932.Ingham}{book}{}
      \name{author}{1}{}{%
        {{hash=a9ad153a68cf8125996aecb67d0f1f8c}{%
           family={Ingham},
           familyi={I\bibinitperiod},
           given={A.\bibnamedelimi E.},
           giveni={A\bibinitperiod\bibinitdelim E\bibinitperiod}}}%
      }
      \list{publisher}{1}{%
        {Cambridge Tracts in Mathematics and Mathematical Physics. 30. London: Cambridge University Press}%
      }
      \strng{namehash}{a9ad153a68cf8125996aecb67d0f1f8c}
      \strng{fullhash}{a9ad153a68cf8125996aecb67d0f1f8c}
      \strng{bibnamehash}{a9ad153a68cf8125996aecb67d0f1f8c}
      \strng{authorbibnamehash}{a9ad153a68cf8125996aecb67d0f1f8c}
      \strng{authornamehash}{a9ad153a68cf8125996aecb67d0f1f8c}
      \strng{authorfullhash}{a9ad153a68cf8125996aecb67d0f1f8c}
      \field{extraname}{1}
      \field{sortinit}{1}
      \field{sortinithash}{50c6687d7fc80f50136d75228e3c59ba}
      \field{labelnamesource}{author}
      \field{labeltitlesource}{title}
      \field{annotation}{In Chapter III of this book (Further Theory of $\zeta(s)$.~Applications), the author shows how information on the zero-free region of $\zeta(s)$ leads to improved estimates for $\Delta^\psi(x)$. Let $\eta(t)$ be a real-valued decreasing function defined on $t \geq 0$ satisfying $0 \leq \eta(t) \leq \frac{1}{2}$ and $1/\eta(t)=O(\log t)$, and suppose that $\eta'(t)$ is continuous and $\eta'(t)=o(1)$. If $\zeta(s)$ has no zeros in the region $\{\sigma>1-\eta(|t|)\}$, then for any $\alpha \in (0,1)$ we have $\Delta^\psi(x) \ll x \exp\bigl(-\frac{1}{2}\alpha \omega(x)\bigr)$, where $\omega(x) = \min_{t\ge1} \bigl( \eta(t) \log x+\log t \bigr)$.\par Chapter V (Irregularities of Distribution) gathers known results on comparative prime number theory from the literature. The author also outlines the proof of \[ \pi(x;4,3,1) = \Omega_\pm\bigg( \frac{\sqrt x}{\log x} \log\log\log x \bigg), \] which was asserted by Hardy and Littlewood~\cite{MR1555148}.\par This book cites~\cites{1909.Landau,MR1544509,MR1555148,MR1511360,MR1544293A,MR1544293B,1914.Littlewood,MR1574052,1891.Phragmen,1930.Polya}.}
      \field{pagetotal}{114}
      \field{title}{{The distribution of prime numbers}}
      \field{year}{1932}
      \true{nocite}
    \endentry
    \entry{MR1573970}{article}{}
      \name{author}{1}{}{%
        {{hash=ed5a88e0fd4d9ca35b105db5f61affef}{%
           family={Skewes},
           familyi={S\bibinitperiod},
           given={S.},
           giveni={S\bibinitperiod}}}%
      }
      \strng{namehash}{ed5a88e0fd4d9ca35b105db5f61affef}
      \strng{fullhash}{ed5a88e0fd4d9ca35b105db5f61affef}
      \strng{bibnamehash}{ed5a88e0fd4d9ca35b105db5f61affef}
      \strng{authorbibnamehash}{ed5a88e0fd4d9ca35b105db5f61affef}
      \strng{authornamehash}{ed5a88e0fd4d9ca35b105db5f61affef}
      \strng{authorfullhash}{ed5a88e0fd4d9ca35b105db5f61affef}
      \field{extraname}{1}
      \field{sortinit}{1}
      \field{sortinithash}{50c6687d7fc80f50136d75228e3c59ba}
      \field{labelnamesource}{author}
      \field{labeltitlesource}{title}
      \field{annotation}{This article shows, assuming RH, that $\pi(x)>\li(x)$ for some $x<10^{10^{10^{34}}}$. Littlewood~\cites{1914.Littlewood,MR1555148} proved the existence of such an $x$ by considering the function $F(z) = \sum_{\gamma>0} {e^{-\gamma (\xi+i\eta)}}/{\gamma}$ for $0\le\xi\le1$ and $\eta\ge1$, which is relevant since the explicit formula yields $-2\Im F(i\log x) = E^\psi(x) +O(1)$. Using the Dirichlet box principle, Littlewood showed that $\Im F(\xi+i\eta)$ has large values (on the order of $\log \log \eta$) of either prescribed sign, with $\xi$ tending to $0$, and then used a modified form of the Phragmén--Lindelöf principle to show that an equally large value of $-\Im F(i\eta)$ must be attained.\par The Phragmén--Lindelöf principle, that the maximum of an analytic function defined on a semi-infinite strip (with suitable growth conditions) must occur on the boundary of the strip, is only an existence result; the author strengthens the result to give quantitative bounds on when an approximation to an interior value is attained on the boundary of the strip. In this way he is able to make Littlewood's result explicit (although the details are not included), which was not clearly possible beforehand.\par (Estimates of the smallest $x$ such that $\pi(x)>\li(x)$ have since been called ``Skewes numbers''. For comparison, the best estimates today are about $1.4\times 10^{316}$.)\par This article cites~\cites{1932.Ingham}.}
      \field{journaltitle}{J. London Math. Soc.}
      \field{number}{4}
      \field{title}{On the {D}ifference {$\pi(x)-{\rm li}\,(x)$} ({I})}
      \field{volume}{8}
      \field{year}{1933}
      \true{nocite}
      \field{pages}{277\bibrangedash 283}
      \range{pages}{7}
      \verb{doi}
      \verb 10.1112/jlms/s1-8.4.277
      \endverb
      \verb{mrnumber}
      \verb 1573970
      \endverb
      \verb{urlraw}
      \verb https://doi.org/10.1112/jlms/s1-8.4.277
      \endverb
      \verb{url}
      \verb https://doi.org/10.1112/jlms/s1-8.4.277
      \endverb
    \endentry
    \entry{MR1501802}{article}{}
      \name{author}{2}{}{%
        {{hash=5acf2ef9928231591e18fa4f528b72bd}{%
           family={Jessen},
           familyi={J\bibinitperiod},
           given={B.},
           giveni={B\bibinitperiod}}}%
        {{hash=0ecb7b4bc82132c8d3551c023cbe496a}{%
           family={Wintner},
           familyi={W\bibinitperiod},
           given={A.},
           giveni={A\bibinitperiod}}}%
      }
      \strng{namehash}{8f1d2457d99bfffa1d82626b836e82d6}
      \strng{fullhash}{8f1d2457d99bfffa1d82626b836e82d6}
      \strng{bibnamehash}{8f1d2457d99bfffa1d82626b836e82d6}
      \strng{authorbibnamehash}{8f1d2457d99bfffa1d82626b836e82d6}
      \strng{authornamehash}{8f1d2457d99bfffa1d82626b836e82d6}
      \strng{authorfullhash}{8f1d2457d99bfffa1d82626b836e82d6}
      \field{sortinit}{1}
      \field{sortinithash}{50c6687d7fc80f50136d75228e3c59ba}
      \field{labelnamesource}{author}
      \field{labeltitlesource}{title}
      \field{annotation}{This article contains a systematic study of infinite convolutions of distribution functions (densities) on $\R^k$, and consequently of infinite sums of independent random variables, via their Fourier transforms; they give criteria for when these convolutions converge and when the result is singular or absolutely continuous. They analyze distributions of functions on convex closed curves in~$\R^2$. The results are applied to the limiting distribution functions of almost-periodic functions of multiple types. As an application, the authors derive results on the limiting distributions of $\log\zeta(\sigma+it)$ and $\zeta(\sigma+it)$ for $\frac12<\sigma\le1$ (the case $\sigma>1$ being simpler).}
      \field{issn}{0002-9947}
      \field{journaltitle}{Trans. Amer. Math. Soc.}
      \field{number}{1}
      \field{title}{Distribution functions and the Riemann zeta function}
      \field{volume}{38}
      \field{year}{1935}
      \true{nocite}
      \field{pages}{48\bibrangedash 88}
      \range{pages}{41}
      \verb{doi}
      \verb 10.2307/1989728
      \endverb
      \verb{mrnumber}
      \verb 1501802
      \endverb
    \endentry
    \entry{MR1507933}{article}{}
      \name{author}{1}{}{%
        {{hash=0ecb7b4bc82132c8d3551c023cbe496a}{%
           family={Wintner},
           familyi={W\bibinitperiod},
           given={A.},
           giveni={A\bibinitperiod}}}%
      }
      \strng{namehash}{0ecb7b4bc82132c8d3551c023cbe496a}
      \strng{fullhash}{0ecb7b4bc82132c8d3551c023cbe496a}
      \strng{bibnamehash}{0ecb7b4bc82132c8d3551c023cbe496a}
      \strng{authorbibnamehash}{0ecb7b4bc82132c8d3551c023cbe496a}
      \strng{authornamehash}{0ecb7b4bc82132c8d3551c023cbe496a}
      \strng{authorfullhash}{0ecb7b4bc82132c8d3551c023cbe496a}
      \field{extraname}{1}
      \field{sortinit}{1}
      \field{sortinithash}{50c6687d7fc80f50136d75228e3c59ba}
      \field{labelnamesource}{author}
      \field{labeltitlesource}{title}
      \field{annotation}{Assuming RH, the author investigates the tail $\sum_{|\rho|>T} x^{i\gamma}/\rho$ of the sum appearing in the explicit formula for $\psi(x)$, showing that its mean-square average tends to~$0$ as $T\to\infty$. This implies that the full sum converges (and represents, in the Besikovitch sense, the function of which it is the ``Fourier series'') and has a logarithmic distribution function. The author also mentions the potential relevance of LI to whether this distribution function can be constant on intervals.\par This article cites~\cites{MR1501802}.}
      \field{issn}{0002-9327}
      \field{journaltitle}{Amer. J. Math.}
      \field{number}{3}
      \field{title}{On the Asymptotic Distribution of the Remainder Term of the Prime-Number Theorem}
      \field{volume}{57}
      \field{year}{1935}
      \true{nocite}
      \field{pages}{534\bibrangedash 538}
      \range{pages}{5}
      \verb{doi}
      \verb 10.2307/2371183
      \endverb
      \verb{mrnumber}
      \verb 1507933
      \endverb
    \endentry
    \entry{1936.Ingham}{article}{}
      \name{author}{1}{}{%
        {{hash=a9ad153a68cf8125996aecb67d0f1f8c}{%
           family={Ingham},
           familyi={I\bibinitperiod},
           given={A.\bibnamedelimi E.},
           giveni={A\bibinitperiod\bibinitdelim E\bibinitperiod}}}%
      }
      \strng{namehash}{a9ad153a68cf8125996aecb67d0f1f8c}
      \strng{fullhash}{a9ad153a68cf8125996aecb67d0f1f8c}
      \strng{bibnamehash}{a9ad153a68cf8125996aecb67d0f1f8c}
      \strng{authorbibnamehash}{a9ad153a68cf8125996aecb67d0f1f8c}
      \strng{authornamehash}{a9ad153a68cf8125996aecb67d0f1f8c}
      \strng{authorfullhash}{a9ad153a68cf8125996aecb67d0f1f8c}
      \field{extraname}{2}
      \field{sortinit}{1}
      \field{sortinithash}{50c6687d7fc80f50136d75228e3c59ba}
      \field{labelnamesource}{author}
      \field{labeltitlesource}{title}
      \field{annotation}{This article gives another proof of an explicit version of Littlewood's theorem~\cite{MR1555148} and establishes the following stronger result: Assuming SA for $\zeta(s)$, there exists an absolute constant $A>1$ such that, for all $x>1$, the interval $(x,Ax)$ contains a sign change of $\pi(x)-\Li(x)$. The author highlights his use of Fejér kernels in the proof, in contrast with the Poisson kernel used by Skewes~\cite{MR1573970}.\par This article cites~\cites{MR1555148, 1930.Polya, MR1573970}.}
      \field{journaltitle}{Acta Arith.}
      \field{title}{{A note on the distribution of primes}}
      \field{volume}{1}
      \field{year}{1936}
      \true{nocite}
      \field{pages}{201\bibrangedash 211}
      \range{pages}{11}
    \endentry
    \entry{MR1575079}{article}{}
      \name{author}{1}{}{%
        {{hash=56dd48151c55dce1b8a073ecf909a238}{%
           family={Littlewood},
           familyi={L\bibinitperiod},
           given={J.\bibnamedelimi E.},
           giveni={J\bibinitperiod\bibinitdelim E\bibinitperiod}}}%
      }
      \strng{namehash}{56dd48151c55dce1b8a073ecf909a238}
      \strng{fullhash}{56dd48151c55dce1b8a073ecf909a238}
      \strng{bibnamehash}{56dd48151c55dce1b8a073ecf909a238}
      \strng{authorbibnamehash}{56dd48151c55dce1b8a073ecf909a238}
      \strng{authornamehash}{56dd48151c55dce1b8a073ecf909a238}
      \strng{authorfullhash}{56dd48151c55dce1b8a073ecf909a238}
      \field{extraname}{3}
      \field{sortinit}{1}
      \field{sortinithash}{50c6687d7fc80f50136d75228e3c59ba}
      \field{labelnamesource}{author}
      \field{labeltitlesource}{title}
      \field{annotation}{This short note is mathematically concerned with the maximum values of trigonometric polynomials. The author reveals that his motivation was to establish an explicit upper bound for the first sign change of $\Delta^\pi(x)$ without assuming RH. Ultimately, the author claims to have found alternative means for finding such an upper bound.\par This note cites~\cites{MR1573970}.}
      \field{journaltitle}{J. London Math. Soc.}
      \field{number}{3}
      \field{title}{Mathematical Notes (12): An Inequality for a Sum of Cosines}
      \field{volume}{12}
      \field{year}{1937}
      \true{nocite}
      \field{pages}{217\bibrangedash 221}
      \range{pages}{5}
      \verb{doi}
      \verb 10.1112/jlms/s1-12.2.217
      \endverb
      \verb{mrnumber}
      \verb 1575079
      \endverb
    \endentry
    \entry{1938.Wintner}{article}{}
      \name{author}{1}{}{%
        {{hash=0ecb7b4bc82132c8d3551c023cbe496a}{%
           family={Wintner},
           familyi={W\bibinitperiod},
           given={A.},
           giveni={A\bibinitperiod}}}%
      }
      \strng{namehash}{0ecb7b4bc82132c8d3551c023cbe496a}
      \strng{fullhash}{0ecb7b4bc82132c8d3551c023cbe496a}
      \strng{bibnamehash}{0ecb7b4bc82132c8d3551c023cbe496a}
      \strng{authorbibnamehash}{0ecb7b4bc82132c8d3551c023cbe496a}
      \strng{authornamehash}{0ecb7b4bc82132c8d3551c023cbe496a}
      \strng{authorfullhash}{0ecb7b4bc82132c8d3551c023cbe496a}
      \field{extraname}{2}
      \field{sortinit}{1}
      \field{sortinithash}{50c6687d7fc80f50136d75228e3c59ba}
      \field{labelnamesource}{author}
      \field{labeltitlesource}{title}
      \field{annotation}{The chapter titles are: 1.~Distribution functions; 2.~Integrals and convolutions; 3.~Moments of distribution functions; 4.~Fourier transforms; 5.~Infinite convolutions; 6.~Smoothness criteria for distribution functions; 7.~Convergence criteria for infinite convolutions; 8.~Asymptotic distributions; 9.~The Riemann zeta-function; 10.~Poisson convolutions; 11.~Convolutions and the theory of probability; 12.~Bernoulli convolutions; 13.~Almost periodic functions with linearly independent exponents; 14.~Symmetric distribution functions in $k$ dimensions; 15.~Two-dimensional convolutions and the Riemann zeta-function; 16.~The addition of convex curves.\par These notes cite~\cites{MR1501802,MR1507933}.}
      \field{journaltitle}{Lecture notes distributed by the Institute for Advanced Study (Princeton)}
      \field{title}{Asymptotic distributions and infinite convolutions}
      \field{year}{1938}
      \true{nocite}
    \endentry
    \entry{MR0003644}{article}{}
      \name{author}{1}{}{%
        {{hash=4f3424aa2d374362fe10ed8b28fd159f}{%
           family={Gupta},
           familyi={G\bibinitperiod},
           given={Hansraj},
           giveni={H\bibinitperiod}}}%
      }
      \strng{namehash}{4f3424aa2d374362fe10ed8b28fd159f}
      \strng{fullhash}{4f3424aa2d374362fe10ed8b28fd159f}
      \strng{bibnamehash}{4f3424aa2d374362fe10ed8b28fd159f}
      \strng{authorbibnamehash}{4f3424aa2d374362fe10ed8b28fd159f}
      \strng{authornamehash}{4f3424aa2d374362fe10ed8b28fd159f}
      \strng{authorfullhash}{4f3424aa2d374362fe10ed8b28fd159f}
      \field{sortinit}{1}
      \field{sortinithash}{50c6687d7fc80f50136d75228e3c59ba}
      \field{labelnamesource}{author}
      \field{labeltitlesource}{title}
      \field{annotation}{Regarding Pólya's problem, the author verifies that $L(n) \leq 0$ for $2 \leq n \leq 20,000$. The author also conjectures that $L(x)\ll\sqrt{x}$ based on the computation.\par This article cites~\cite{1919.Polya}.}
      \field{journaltitle}{Proc. Indian Acad. Sci., Sect. A.}
      \field{title}{On a table of values of {$L(n)$}}
      \field{volume}{12}
      \field{year}{1940}
      \true{nocite}
      \field{pages}{407\bibrangedash 409}
      \range{pages}{3}
      \verb{mrnumber}
      \verb 0003644
      \endverb
    \endentry
    \entry{MR0004255}{article}{}
      \name{author}{1}{}{%
        {{hash=0ecb7b4bc82132c8d3551c023cbe496a}{%
           family={Wintner},
           familyi={W\bibinitperiod},
           given={A.},
           giveni={A\bibinitperiod}}}%
      }
      \strng{namehash}{0ecb7b4bc82132c8d3551c023cbe496a}
      \strng{fullhash}{0ecb7b4bc82132c8d3551c023cbe496a}
      \strng{bibnamehash}{0ecb7b4bc82132c8d3551c023cbe496a}
      \strng{authorbibnamehash}{0ecb7b4bc82132c8d3551c023cbe496a}
      \strng{authornamehash}{0ecb7b4bc82132c8d3551c023cbe496a}
      \strng{authorfullhash}{0ecb7b4bc82132c8d3551c023cbe496a}
      \field{extraname}{3}
      \field{sortinit}{1}
      \field{sortinithash}{50c6687d7fc80f50136d75228e3c59ba}
      \field{labelnamesource}{author}
      \field{labeltitlesource}{title}
      \field{annotation}{This article investigates the normalized remainder term $E^\Psi(x)$ by establishing the existence of its limiting logarithmic cumulative distribution function. Assuming RH, the author proves that the spectrum of this limiting distribution is unbounded both above and below, which implies the same for $E^\Psi(x)$. He also gives an estimate for all of the moments of this limiting distribution, in an effort to determine whether these moments uniquely determine the given distribution.\par The techniques involved in the proof come primarily from the application of the theory of almost-periodic functions (of both the uniform and Besicovitch varieties) to the sum over the nontrivial zeros in the fundamental explicit formula.\par This article cites~\cites{1932.Ingham,MR1501802,MR1507933}.}
      \field{issn}{0002-9327}
      \field{journaltitle}{Amer. J. Math.}
      \field{title}{On the distribution function of the remainder term of the prime number theorem}
      \field{volume}{63}
      \field{year}{1941}
      \true{nocite}
      \field{pages}{233\bibrangedash 248}
      \range{pages}{16}
      \verb{doi}
      \verb 10.2307/2371519
      \endverb
      \verb{mrnumber}
      \verb 0004255
      \endverb
      \verb{urlraw}
      \verb https://doi.org/10.2307/2371519
      \endverb
      \verb{url}
      \verb https://doi.org/10.2307/2371519
      \endverb
    \endentry
    \entry{MR0006202}{article}{}
      \name{author}{1}{}{%
        {{hash=a9ad153a68cf8125996aecb67d0f1f8c}{%
           family={Ingham},
           familyi={I\bibinitperiod},
           given={A.\bibnamedelimi E.},
           giveni={A\bibinitperiod\bibinitdelim E\bibinitperiod}}}%
      }
      \strng{namehash}{a9ad153a68cf8125996aecb67d0f1f8c}
      \strng{fullhash}{a9ad153a68cf8125996aecb67d0f1f8c}
      \strng{bibnamehash}{a9ad153a68cf8125996aecb67d0f1f8c}
      \strng{authorbibnamehash}{a9ad153a68cf8125996aecb67d0f1f8c}
      \strng{authornamehash}{a9ad153a68cf8125996aecb67d0f1f8c}
      \strng{authorfullhash}{a9ad153a68cf8125996aecb67d0f1f8c}
      \field{extraname}{3}
      \field{sortinit}{1}
      \field{sortinithash}{50c6687d7fc80f50136d75228e3c59ba}
      \field{labelnamesource}{author}
      \field{labeltitlesource}{title}
      \field{annotation}{In this article, Ingham probes conjectured bounds for the summatory functions $M(x)$ and $L(x)$. He proves that the truth, for sufficiently large $x$, of any one of the inequalities $M(x) < Kx^{1/2}$, $M(x) >- Kx^{1/2}$, $L(x) < Kx^{1/2}$, or $L(x) >- Kx^{1/2}$ (where $K$ is a constant) would imply not only RH and the simplicity of the zeros of $\zeta(s)$ (as was ``well known''), but also the falsity of LI. Of this last assumption, the author writes: ``It would be easy to relax this hypothesis a little, but there seems no obvious way of replacing it by anything essentially easier to verify.'' Indeed, he shows that if there are only finitely many rational linear relations among the positive imaginary parts of these zeros, then $E^M(x)$ and $E^L(x)$ would be unbounded both above and below, contrary to existing conjectures.\par The method of the proof is similar to Littlewood's disproof of the conjecture $\pi(x) < \li(x)$ in~\cite{MR1555148}, including a reliance on trigonometric polynomials involving the zeros of $\zeta(s)$, except that Dirichlet's theorem on homogeneous Diophantine approximation is replaced by Kronecker's theorem on inhomogeneous Diophantine approximation. For the proof, the author establishes two main results, one concerning Laplace transforms of real trigonometric polynomials, and the other establishing the divergence (assuming RH) of the two residue series $\sum_{\gamma>0} 1/\rho\zeta'(\rho)$ and $\sum_{\gamma>0} \zeta(2\rho)/\rho\zeta'(\rho)$.\par This article cites~\cites{1897.Mertens, 1897.vonSterneck, 1901.vonSterneck, 1905.Stieltjes, 1913.vonSterneck, 1919.Polya,MR0003644}.}
      \field{issn}{0002-9327}
      \field{journaltitle}{Amer. J. Math.}
      \field{title}{On two conjectures in the theory of numbers}
      \field{volume}{64}
      \field{year}{1942}
      \true{nocite}
      \field{pages}{313\bibrangedash 319}
      \range{pages}{7}
      \verb{doi}
      \verb 10.2307/2371685
      \endverb
      \verb{mrnumber}
      \verb 0006202
      \endverb
      \verb{urlraw}
      \verb https://doi.org/10.2307/2371685
      \endverb
      \verb{url}
      \verb https://doi.org/10.2307/2371685
      \endverb
    \endentry
    \entry{MR0017310}{article}{}
      \name{author}{1}{}{%
        {{hash=a8001a783b2aab7a2865fcdf627bcf26}{%
           family={Tietze},
           familyi={T\bibinitperiod},
           given={Heinrich},
           giveni={H\bibinitperiod}}}%
      }
      \list{language}{1}{%
        {German}%
      }
      \strng{namehash}{a8001a783b2aab7a2865fcdf627bcf26}
      \strng{fullhash}{a8001a783b2aab7a2865fcdf627bcf26}
      \strng{bibnamehash}{a8001a783b2aab7a2865fcdf627bcf26}
      \strng{authorbibnamehash}{a8001a783b2aab7a2865fcdf627bcf26}
      \strng{authornamehash}{a8001a783b2aab7a2865fcdf627bcf26}
      \strng{authorfullhash}{a8001a783b2aab7a2865fcdf627bcf26}
      \field{sortinit}{1}
      \field{sortinithash}{50c6687d7fc80f50136d75228e3c59ba}
      \field{labelnamesource}{author}
      \field{labeltitlesource}{title}
      \field{annotation}{The author provides tables of $\phi(q)\pi(x;q,\cA,\Gamma)$ for various moduli~$q$ (with a particular focus on~$q=262$, the smallest even modulus for which $3,5,7,11,13$ are all quadratic residues) and various sets~$\Gamma$ of residue classes, often subgroups of~$\cA$.}
      \field{journaltitle}{Abh. Bayer. Akad. Wiss. Math.-Nat. Abt. (N.F.)}
      \field{number}{55}
      \field{title}{Einige Tabellen zur Verteilung der Primzahlen auf Untergruppen der teilerfremden Restklassen nach gegebenem Modul}
      \field{volume}{1944}
      \field{year}{1944}
      \true{nocite}
      \field{pages}{31}
      \range{pages}{1}
      \verb{mrnumber}
      \verb 0017310
      \endverb
    \endentry
    \entry{MR27305}{article}{}
      \name{author}{1}{}{%
        {{hash=f4ee93efcace35c40df58fa6a8eb26ba}{%
           family={Turán},
           familyi={T\bibinitperiod},
           given={Paul},
           giveni={P\bibinitperiod}}}%
      }
      \strng{namehash}{f4ee93efcace35c40df58fa6a8eb26ba}
      \strng{fullhash}{f4ee93efcace35c40df58fa6a8eb26ba}
      \strng{bibnamehash}{f4ee93efcace35c40df58fa6a8eb26ba}
      \strng{authorbibnamehash}{f4ee93efcace35c40df58fa6a8eb26ba}
      \strng{authornamehash}{f4ee93efcace35c40df58fa6a8eb26ba}
      \strng{authorfullhash}{f4ee93efcace35c40df58fa6a8eb26ba}
      \field{extraname}{1}
      \field{sortinit}{1}
      \field{sortinithash}{50c6687d7fc80f50136d75228e3c59ba}
      \field{labelnamesource}{author}
      \field{labeltitlesource}{title}
      \field{annotation}{The author shows that if there exists a positive constant $K$ such that $L_r(x)>-K/\sqrt{x}$ for sufficiently large $x$, then RH holds. The author also reports that $L_r(x)\geq 0$ for $x \leq 1000$. While people often refer to the statement that $L_r(x)$ is never negative as ``Turán's conjecture'', that claim is never made in this article. \par This article cites~\cite{1919.Polya,1930.Polya,MR0003644}.}
      \field{issn}{0023-3323}
      \field{journaltitle}{Danske Vid. Selsk. Mat.-Fys. Medd.}
      \field{number}{17}
      \field{title}{On some approximative {D}irichlet-polynomials in the theory of the zeta-function of {R}iemann}
      \field{volume}{24}
      \field{year}{1948}
      \true{nocite}
      \field{pages}{36}
      \range{pages}{1}
      \verb{mrnumber}
      \verb 27305
      \endverb
    \endentry
    \entry{MR29414}{article}{}
      \name{author}{1}{}{%
        {{hash=86b5f5639455f1050b0d4e2301179477}{%
           family={Wintner},
           familyi={W\bibinitperiod},
           given={Aurel},
           giveni={A\bibinitperiod}}}%
      }
      \strng{namehash}{86b5f5639455f1050b0d4e2301179477}
      \strng{fullhash}{86b5f5639455f1050b0d4e2301179477}
      \strng{bibnamehash}{86b5f5639455f1050b0d4e2301179477}
      \strng{authorbibnamehash}{86b5f5639455f1050b0d4e2301179477}
      \strng{authornamehash}{86b5f5639455f1050b0d4e2301179477}
      \strng{authorfullhash}{86b5f5639455f1050b0d4e2301179477}
      \field{extraname}{4}
      \field{sortinit}{1}
      \field{sortinithash}{50c6687d7fc80f50136d75228e3c59ba}
      \field{labelnamesource}{author}
      \field{labeltitlesource}{title}
      \field{annotation}{The author shows that $M(x)\ll\sqrt{x}$ is equivalent to $\sum_{n \leq x} \frac{\mu(n)}{\sqrt{n}}\ll1$; the forward direction follows from partial summation, while the converse uses a Tauberian theorem due to Riesz.}
      \field{issn}{0034-7760}
      \field{journaltitle}{Rev. Ci. (Lima)}
      \field{title}{A note on {M}ertens' hypothesis}
      \field{volume}{50}
      \field{year}{1948}
      \true{nocite}
      \field{pages}{181\bibrangedash 184}
      \range{pages}{4}
      \verb{mrnumber}
      \verb 29414
      \endverb
    \endentry
    \entry{MR0049219}{article}{}
      \name{author}{1}{}{%
        {{hash=e68e101772b88cb1d48e688164bece92}{%
           family={Turán},
           familyi={T\bibinitperiod},
           given={P.},
           giveni={P\bibinitperiod}}}%
      }
      \list{language}{1}{%
        {English, with Russian summary}%
      }
      \strng{namehash}{e68e101772b88cb1d48e688164bece92}
      \strng{fullhash}{e68e101772b88cb1d48e688164bece92}
      \strng{bibnamehash}{e68e101772b88cb1d48e688164bece92}
      \strng{authorbibnamehash}{e68e101772b88cb1d48e688164bece92}
      \strng{authornamehash}{e68e101772b88cb1d48e688164bece92}
      \strng{authorfullhash}{e68e101772b88cb1d48e688164bece92}
      \field{extraname}{2}
      \field{sortinit}{1}
      \field{sortinithash}{50c6687d7fc80f50136d75228e3c59ba}
      \field{labelnamesource}{author}
      \field{labeltitlesource}{title}
      \field{annotation}{The author shows that if $\Delta^\pi(x) \ll x\exp\bigl( -a(\log x)^{1/(1+\beta)} \bigr)$ for some constants $a>0$ and $0<\beta<1$, then there exists a constant $b>0$ such that $\zeta(s)$ does not vanish in the domain $\sigma>1-{b}/{\log^\beta|t|}$; this is the converse of a special case of a theorem of Ingham. The main tool of the proof is a version of the ``second main theorem'' of the power-sum method (the statement in the article seems to contain the typo $\max|z_j|\ge1$ instead of $\max|z_j|\le1$) for longer ranges of the exponent.\par This article cites~\cites{1932.Ingham}.}
      \field{issn}{0001-5954}
      \field{journaltitle}{Acta Math. Acad. Sci. Hungar.}
      \field{title}{On the remainder-term of the prime-number formula. {II}}
      \field{volume}{1}
      \field{year}{1950}
      \true{nocite}
      \field{pages}{155\bibrangedash 166}
      \range{pages}{12}
      \verb{doi}
      \verb 10.1007/BF02021308
      \endverb
      \verb{mrnumber}
      \verb 0049219
      \endverb
    \endentry
    \entry{MR0043121}{article}{}
      \name{author}{1}{}{%
        {{hash=f4ee93efcace35c40df58fa6a8eb26ba}{%
           family={Turán},
           familyi={T\bibinitperiod},
           given={Paul},
           giveni={P\bibinitperiod}}}%
      }
      \list{language}{1}{%
        {English, with Russian summary}%
      }
      \strng{namehash}{f4ee93efcace35c40df58fa6a8eb26ba}
      \strng{fullhash}{f4ee93efcace35c40df58fa6a8eb26ba}
      \strng{bibnamehash}{f4ee93efcace35c40df58fa6a8eb26ba}
      \strng{authorbibnamehash}{f4ee93efcace35c40df58fa6a8eb26ba}
      \strng{authornamehash}{f4ee93efcace35c40df58fa6a8eb26ba}
      \strng{authorfullhash}{f4ee93efcace35c40df58fa6a8eb26ba}
      \field{extraname}{3}
      \field{sortinit}{1}
      \field{sortinithash}{50c6687d7fc80f50136d75228e3c59ba}
      \field{labelnamesource}{author}
      \field{labeltitlesource}{title}
      \field{annotation}{This article investigates lower bounds for $|\Delta^\psi(x)|$ and $|\Delta^\pi(x)|$, specifically to answer Littlewood's call~\cite{MR1575079} to replace Landau's theorem, which depends upon the smallest half-plane containing all the zeros of $\zeta(s)$, by an estimation that depends only upon one single zero $\rho_0=\beta_0+i\gamma_0$ of $\zeta(s)$. The author shows, using a version of his ``first main theorem'', that when $T$ is sufficiently large in terms of $\rho_0$, $$ \max_{1\le x\le T}|\Delta^\psi(x)| > \frac{T^{\beta_0}}{{|\rho_0|}^{(10\log T)/\log \log T}}\exp \left(-\frac{c\log T \log\log\log T}{\log \log T}\right), $$ where $c$ is a positive constant depending on~$\rho_0$. \par This article cites~\cites{1914.Littlewood, MR1555148, 1930.Polya, 1936.Ingham, MR0067145, MR0133309, MR1511206, MR1575079}.}
      \field{issn}{0001-5954}
      \field{journaltitle}{Acta Math. Acad. Sci. Hungar.}
      \field{title}{On the remainder-term of the prime-number formula. {I}}
      \field{volume}{1}
      \field{year}{1950}
      \true{nocite}
      \field{pages}{48\bibrangedash 63}
      \range{pages}{16}
      \verb{doi}
      \verb 10.1007/BF02022552
      \endverb
      \verb{mrnumber}
      \verb 0043121
      \endverb
    \endentry
    \entry{MR43841}{article}{}
      \name{author}{1}{}{%
        {{hash=695e0d70149c6690cca3917ef5b74c9e}{%
           family={Fawaz},
           familyi={F\bibinitperiod},
           given={A.\bibnamedelimi Y.},
           giveni={A\bibinitperiod\bibinitdelim Y\bibinitperiod}}}%
      }
      \strng{namehash}{695e0d70149c6690cca3917ef5b74c9e}
      \strng{fullhash}{695e0d70149c6690cca3917ef5b74c9e}
      \strng{bibnamehash}{695e0d70149c6690cca3917ef5b74c9e}
      \strng{authorbibnamehash}{695e0d70149c6690cca3917ef5b74c9e}
      \strng{authornamehash}{695e0d70149c6690cca3917ef5b74c9e}
      \strng{authorfullhash}{695e0d70149c6690cca3917ef5b74c9e}
      \field{extraname}{1}
      \field{sortinit}{1}
      \field{sortinithash}{50c6687d7fc80f50136d75228e3c59ba}
      \field{labelnamesource}{author}
      \field{labeltitlesource}{title}
      \field{annotation}{Assuming RH and the simplicity of the zeros of $\zeta(s)$, the author records the identity \begin{align*} L_0(x) = \frac{\sqrt{x}}{\zeta(\frac{1}{2})}+\sum_{\rho}\frac{\zeta(2\rho)}{\rho\zeta'(\rho)}x^{\rho}+\frac{1}{2\pi i}\int_{1/4-i\infty}^{1/4+i\infty}\frac{\zeta(2s)}{s\zeta(s)}x^{s} \, ds. \end{align*} By a more delicate argument, he establishes the more explicit series representation \[ \frac{1}{2\pi i}\int_{1/4-i\infty}^{1/4+i\infty}\frac{\zeta(2s)}{s\zeta(s)}x^{s} \, ds = 2\sum_{n=1}^\infty \frac{q(n)\lambda(n)}n \big( C(\sqrt{nx}) + S(\sqrt{nx}) \big), \] where $q(n)$ is the largest integer whose square divides~$n$, and $C(t)$ and $S(t)$ are the normalized Fresnel integrals, which both tend to $\frac12$ as $t\to\infty$. With these identities, the author indicates a possible approach to resolving Pólya's problem in the negative. \par This article cites~\cites{1909.Landau,1919.Polya,1936.Ingham,MR0006202}.}
      \field{issn}{0024-6115}
      \field{journaltitle}{Proc. London Math. Soc. (3)}
      \field{title}{The explicit formula for {$L_0(x)$}}
      \field{volume}{1}
      \field{year}{1951}
      \true{nocite}
      \field{pages}{86\bibrangedash 103}
      \range{pages}{18}
      \verb{doi}
      \verb 10.1112/plms/s3-1.1.86
      \endverb
      \verb{mrnumber}
      \verb 43841
      \endverb
      \verb{urlraw}
      \verb https://doi-org.ezproxy.library.ubc.ca/10.1112/plms/s3-1.1.86
      \endverb
      \verb{url}
      \verb https://doi-org.ezproxy.library.ubc.ca/10.1112/plms/s3-1.1.86
      \endverb
    \endentry
    \entry{MR0046485}{book}{}
      \name{author}{1}{}{%
        {{hash=baf219d417d1941af57ccbb309deb548}{%
           family={Titchmarsh},
           familyi={T\bibinitperiod},
           given={E.\bibnamedelimi C.},
           giveni={E\bibinitperiod\bibinitdelim C\bibinitperiod}}}%
      }
      \list{publisher}{1}{%
        {Oxford, at the Clarendon Press}%
      }
      \strng{namehash}{baf219d417d1941af57ccbb309deb548}
      \strng{fullhash}{baf219d417d1941af57ccbb309deb548}
      \strng{bibnamehash}{baf219d417d1941af57ccbb309deb548}
      \strng{authorbibnamehash}{baf219d417d1941af57ccbb309deb548}
      \strng{authornamehash}{baf219d417d1941af57ccbb309deb548}
      \strng{authorfullhash}{baf219d417d1941af57ccbb309deb548}
      \field{extraname}{1}
      \field{sortinit}{1}
      \field{sortinithash}{50c6687d7fc80f50136d75228e3c59ba}
      \field{labelnamesource}{author}
      \field{labeltitlesource}{title}
      \field{annotation}{This is the first edition of~\cite{MR0882550}.}
      \field{title}{The {T}heory of the {R}iemann {Z}eta-{F}unction}
      \field{year}{1951}
      \true{nocite}
      \field{pages}{vi+346}
      \range{pages}{-1}
      \verb{mrnumber}
      \verb 0046485
      \endverb
    \endentry
    \entry{MR51857}{article}{}
      \name{author}{1}{}{%
        {{hash=695e0d70149c6690cca3917ef5b74c9e}{%
           family={Fawaz},
           familyi={F\bibinitperiod},
           given={A.\bibnamedelimi Y.},
           giveni={A\bibinitperiod\bibinitdelim Y\bibinitperiod}}}%
      }
      \strng{namehash}{695e0d70149c6690cca3917ef5b74c9e}
      \strng{fullhash}{695e0d70149c6690cca3917ef5b74c9e}
      \strng{bibnamehash}{695e0d70149c6690cca3917ef5b74c9e}
      \strng{authorbibnamehash}{695e0d70149c6690cca3917ef5b74c9e}
      \strng{authornamehash}{695e0d70149c6690cca3917ef5b74c9e}
      \strng{authorfullhash}{695e0d70149c6690cca3917ef5b74c9e}
      \field{extraname}{2}
      \field{sortinit}{1}
      \field{sortinithash}{50c6687d7fc80f50136d75228e3c59ba}
      \field{labelnamesource}{author}
      \field{labeltitlesource}{title}
      \field{annotation}{This article is a continuation of the author's previous work~\cite{MR43841}. Defining \[ I(x)=\frac{1}{2\pi i}\int_{a-i\infty}^{a+i\infty}\frac{\zeta(2s)}{s\zeta(s)}x^{s} \, ds \] for any $a \in (0,1/2)$, the author proves that \[ \liminf_{x \to \infty} E^L(x) = -\limsup_{x \to 0^+} \frac{I(x)}{\sqrt x} \quad\text{and}\quad \limsup_{x \to \infty} E^L(x) = -\liminf_{x \to 0^+} \frac{I(x)}{\sqrt x}, \] which suggests an alternative approach to studying $L(x)$. \par This article cites~\cite{MR43841,MR0003644,MR0006202,1932.Ingham,1919.Polya,MR0046485,1936.Ingham}.}
      \field{issn}{0033-5606}
      \field{journaltitle}{Quart. J. Math. Oxford Ser. (2)}
      \field{title}{On an unsolved problem in the analytic theory of numbers}
      \field{volume}{3}
      \field{year}{1952}
      \true{nocite}
      \field{pages}{282\bibrangedash 295}
      \range{pages}{14}
      \verb{doi}
      \verb 10.1093/qmath/3.1.282
      \endverb
      \verb{mrnumber}
      \verb 51857
      \endverb
      \verb{urlraw}
      \verb https://doi.org/10.1093/qmath/3.1.282
      \endverb
      \verb{url}
      \verb https://doi.org/10.1093/qmath/3.1.282
      \endverb
    \endentry
    \entry{MR0068565}{book}{}
      \name{author}{1}{}{%
        {{hash=45dbea41eb1e7f07fcfe10d098ce7670}{%
           family={Landau},
           familyi={L\bibinitperiod},
           given={E.},
           giveni={E\bibinitperiod}}}%
      }
      \list{language}{1}{%
        {German}%
      }
      \list{publisher}{1}{%
        {Chelsea Publishing Co., New York}%
      }
      \strng{namehash}{45dbea41eb1e7f07fcfe10d098ce7670}
      \strng{fullhash}{45dbea41eb1e7f07fcfe10d098ce7670}
      \strng{bibnamehash}{45dbea41eb1e7f07fcfe10d098ce7670}
      \strng{authorbibnamehash}{45dbea41eb1e7f07fcfe10d098ce7670}
      \strng{authornamehash}{45dbea41eb1e7f07fcfe10d098ce7670}
      \strng{authorfullhash}{45dbea41eb1e7f07fcfe10d098ce7670}
      \field{extraname}{5}
      \field{sortinit}{1}
      \field{sortinithash}{50c6687d7fc80f50136d75228e3c59ba}
      \field{labelnamesource}{author}
      \field{labeltitlesource}{title}
      \field{annotation}{This is the second edition of~\cite{1909.Landau}.}
      \field{note}{2d ed; With an appendix by Paul T. Bateman}
      \field{title}{Handbuch der Lehre von der Verteilung der Primzahlen. 2 Bände}
      \field{year}{1953}
      \true{nocite}
      \field{pages}{xviii+pp. 1\bibrangedash 564\bibrangessep ix+pp. 565\bibrangedash 1001}
      \range{pages}{-1}
      \verb{mrnumber}
      \verb 0068565
      \endverb
    \endentry
    \entry{MR0060548}{book}{}
      \name{author}{1}{}{%
        {{hash=f4ee93efcace35c40df58fa6a8eb26ba}{%
           family={Turán},
           familyi={T\bibinitperiod},
           given={Paul},
           giveni={P\bibinitperiod}}}%
      }
      \list{language}{1}{%
        {German}%
      }
      \list{publisher}{1}{%
        {Akadémiai Kiadó, Budapest}%
      }
      \strng{namehash}{f4ee93efcace35c40df58fa6a8eb26ba}
      \strng{fullhash}{f4ee93efcace35c40df58fa6a8eb26ba}
      \strng{bibnamehash}{f4ee93efcace35c40df58fa6a8eb26ba}
      \strng{authorbibnamehash}{f4ee93efcace35c40df58fa6a8eb26ba}
      \strng{authornamehash}{f4ee93efcace35c40df58fa6a8eb26ba}
      \strng{authorfullhash}{f4ee93efcace35c40df58fa6a8eb26ba}
      \field{extraname}{4}
      \field{sortinit}{1}
      \field{sortinithash}{50c6687d7fc80f50136d75228e3c59ba}
      \field{labelnamesource}{author}
      \field{labeltitlesource}{title}
      \field{annotation}{This book contains a coherent exposition of power-sum method which is also referred as ``Turàn's method''. Many applications of the method are given, including a result that derives zero-free regions from bounds for $E^\psi(x)$ and improved zero-density results.}
      \field{pagetotal}{196}
      \field{title}{Eine neue Methode in der Analysis und deren Anwendungen}
      \field{year}{1953}
      \true{nocite}
      \verb{mrnumber}
      \verb 0060548
      \endverb
    \endentry
    \entry{MR0067145}{article}{}
      \name{author}{1}{}{%
        {{hash=ed5a88e0fd4d9ca35b105db5f61affef}{%
           family={Skewes},
           familyi={S\bibinitperiod},
           given={S.},
           giveni={S\bibinitperiod}}}%
      }
      \strng{namehash}{ed5a88e0fd4d9ca35b105db5f61affef}
      \strng{fullhash}{ed5a88e0fd4d9ca35b105db5f61affef}
      \strng{bibnamehash}{ed5a88e0fd4d9ca35b105db5f61affef}
      \strng{authorbibnamehash}{ed5a88e0fd4d9ca35b105db5f61affef}
      \strng{authornamehash}{ed5a88e0fd4d9ca35b105db5f61affef}
      \strng{authorfullhash}{ed5a88e0fd4d9ca35b105db5f61affef}
      \field{extraname}{2}
      \field{sortinit}{1}
      \field{sortinithash}{50c6687d7fc80f50136d75228e3c59ba}
      \field{labelnamesource}{author}
      \field{labeltitlesource}{title}
      \field{annotation}{This article provides an unconditional explicit estimate for a sign change of the difference $\Delta^\pi(x)$: if we define $X_1 = {e^{e^{e^{7.703}}}}$ and $X_2 = e^{4X_1^{30}} < e^{e^{e^{e^{7.705}}}} < 10^{10^{10^{10^3}}}$, then the author shows that there exists some $x<X_2$ such that $\pi(x)>\li(x)$. The author divides the proof into two cases, first when RH is ``nearly true'' and then the contrary case. More specifically, he defines a hypothesis (H) (the ``nearly true'' case) as follows: \emph{Every zero $\rho = \beta + i\gamma$ for which $|\gamma| < X_1^3$ satisfies $\beta - \frac{1}{2} \le X_1^{-3}\log^{-2}X_1$}. \par For the case where (H) holds, the author modifies Ingham's technique from~\cite{1936.Ingham}, which assumed RH but improved the estimation of $\Delta^{\psi_0}(x)$ by showing that zeros with $\gamma$ large relative to $x$ do not contribute meaningfully to the sum. Ultimately the author's argument boils down to estimation of the sum $\sum_{0<\gamma<500} \frac{\sin \gamma\omega}{\gamma}\left(1-\frac{\gamma}{500}\right)$; Dirichlet's box principle is used again, in conjunction with estimates of the values of the $269$ zeros of $\zeta(s)$ with $0<\gamma<500$. \par For the contrary case, which the author calls (NH), he remarks that it no longer suffices to work first with $\psi(x)$ and then pass to $\pi(x)$ with standard partial summation techniques. Instead, he works directly from the explicit formula for $\Delta^{\Pi_0}(x)$, introducing a smoothing factor to amplify the contribution from the hypothesized (H)-violating zero. Throughout, the author uses explicit estimates for sums over nontrivial zeros of $\zeta(s)$, such as $|N(T+h)-N(T)| < \frac{1}{2\pi} (h+1.77)\log T+ 8.7$ for $7.1<h<\frac T2$. \par This article cites~\cites{1914.Littlewood,1932.Ingham,MR1573970,1936.Ingham}.}
      \field{issn}{0024-6115}
      \field{journaltitle}{Proc. London Math. Soc. (3)}
      \field{title}{On the difference {$\pi(x)-{\rm li}\,x$}. {II}}
      \field{volume}{5}
      \field{year}{1955}
      \true{nocite}
      \field{pages}{48\bibrangedash 70}
      \range{pages}{23}
      \verb{doi}
      \verb 10.1112/plms/s3-5.1.48
      \endverb
      \verb{mrnumber}
      \verb 0067145
      \endverb
      \verb{urlraw}
      \verb https://doi.org/10.1112/plms/s3-5.1.48
      \endverb
      \verb{url}
      \verb https://doi.org/10.1112/plms/s3-5.1.48
      \endverb
    \endentry
    \entry{MR0077579}{article}{}
      \name{author}{2}{}{%
        {{hash=7220d8e752844cc34f659367559952a5}{%
           family={Sós},
           familyi={S\bibinitperiod},
           given={Vera\bibnamedelima T.},
           giveni={V\bibinitperiod\bibinitdelim T\bibinitperiod}}}%
        {{hash=e68e101772b88cb1d48e688164bece92}{%
           family={Turán},
           familyi={T\bibinitperiod},
           given={P.},
           giveni={P\bibinitperiod}}}%
      }
      \list{language}{1}{%
        {English, with Russian summary}%
      }
      \strng{namehash}{94e15cc05318c39c603ade60525f4392}
      \strng{fullhash}{94e15cc05318c39c603ade60525f4392}
      \strng{bibnamehash}{94e15cc05318c39c603ade60525f4392}
      \strng{authorbibnamehash}{94e15cc05318c39c603ade60525f4392}
      \strng{authornamehash}{94e15cc05318c39c603ade60525f4392}
      \strng{authorfullhash}{94e15cc05318c39c603ade60525f4392}
      \field{sortinit}{1}
      \field{sortinithash}{50c6687d7fc80f50136d75228e3c59ba}
      \field{labelnamesource}{author}
      \field{labeltitlesource}{title}
      \field{annotation}{The authors give bounds for the valid constants $A$ and $B$ appearing in the first and second main theorems of the power-sum method. $A=2e$ is known to be valid in the first main theorem, and the authors show that $A=\frac4\pi$ is too small. They also show that $B=2e^{1+4/e}$ is valid in the second main theorem (improving the previous $B=24e^2$) while $B=1.321$ is not; they mention that the improvement in the upper bound leads to improved constants in a zero-density theorem for $\zeta(s)$ and in the maximal gap between primes. The authors also refer to a ``third main theorem'', in which the normalized lower bound for the power-sum is independent of the~$b_j$ rather than of the~$z_j$. \par This article cites~\cites{MR0060548}.}
      \field{issn}{0001-5954}
      \field{journaltitle}{Acta Math. Acad. Sci. Hungar.}
      \field{title}{On some new theorems in the theory of Diophantine approximations}
      \field{volume}{6}
      \field{year}{1955}
      \true{nocite}
      \field{pages}{241\bibrangedash 255}
      \range{pages}{15}
      \verb{doi}
      \verb 10.1007/BF02024389
      \endverb
      \verb{mrnumber}
      \verb 0077579
      \endverb
    \endentry
    \entry{MR0083001}{article}{}
      \name{author}{1}{}{%
        {{hash=0d7fbc6c344349f0368ab70890c25508}{%
           family={Leech},
           familyi={L\bibinitperiod},
           given={J.},
           giveni={J\bibinitperiod}}}%
      }
      \strng{namehash}{0d7fbc6c344349f0368ab70890c25508}
      \strng{fullhash}{0d7fbc6c344349f0368ab70890c25508}
      \strng{bibnamehash}{0d7fbc6c344349f0368ab70890c25508}
      \strng{authorbibnamehash}{0d7fbc6c344349f0368ab70890c25508}
      \strng{authornamehash}{0d7fbc6c344349f0368ab70890c25508}
      \strng{authorfullhash}{0d7fbc6c344349f0368ab70890c25508}
      \field{sortinit}{1}
      \field{sortinithash}{50c6687d7fc80f50136d75228e3c59ba}
      \field{labelnamesource}{author}
      \field{labeltitlesource}{title}
      \field{annotation}{The author uses the EDSAC at Cambridge to compute $\pi(x;4,1)$ and $\pi(x;4,3)$ for $x$ up to $3\times10^6$. He discovers that $\pi(x;4,1)>\pi(x;4,3)$ at $x = 26,861$, for which $\pi(x;4,1) = 1,473$ and $\pi(x;4,3) = 1,472$. The other values of $x$ above $26,863$ for which $\pi(x;4,1)>\pi(x;4,3)$ are between $616,000$ and $634,000$; the greatest difference found is at $x = 623,681$, for which $\pi(x;4,1) = 25,444$ and $\pi(x;4,3) = 25,436$.\par The author notes that $\pi_i(x) = 2\pi(x;4,1) + \pi(\sqrt x;4,3) + 1$, the number of Gaussian primes with norm at most $x$ (up to associates, and for $x\ge2$), is consequently large near this latter range as well; the most extreme value found is at $x = 617,537$, for which $\pi_i(x)= 50,509 \approx \li(x)+19.5$. \par When examining the explicit formula for $\pi(x;4,3,1)$ at $x = 620,000$, the author found that the first $20$ pairs of zeros of $L(s,\chi_{-4})$, whose imaginary parts ranged from $\pm 6.020948$ to $\pm 49.723129$, included $16$ pairs that give negative contributions to the explicit formula, while subsequent zeros gave more or less random contributions.\par This article cites~\cites{1932.Ingham,MR1573970}.}
      \field{issn}{0024-6107}
      \field{journaltitle}{J. London Math. Soc.}
      \field{title}{Note on the distribution of prime numbers}
      \field{volume}{32}
      \field{year}{1957}
      \true{nocite}
      \field{pages}{56\bibrangedash 58}
      \range{pages}{3}
      \verb{doi}
      \verb 10.1112/jlms/s1-32.1.56
      \endverb
      \verb{mrnumber}
      \verb 0083001
      \endverb
      \verb{urlraw}
      \verb https://doi.org/10.1112/jlms/s1-32.1.56
      \endverb
      \verb{url}
      \verb https://doi.org/10.1112/jlms/s1-32.1.56
      \endverb
    \endentry
    \entry{MR0087685}{book}{}
      \name{author}{1}{}{%
        {{hash=fa022d8efe9c2814c0c7cd5fe921a192}{%
           family={Prachar},
           familyi={P\bibinitperiod},
           given={Karl},
           giveni={K\bibinitperiod}}}%
      }
      \list{language}{1}{%
        {German}%
      }
      \list{publisher}{1}{%
        {Springer-Verlag, Berlin-Göttingen-Heidelberg}%
      }
      \strng{namehash}{fa022d8efe9c2814c0c7cd5fe921a192}
      \strng{fullhash}{fa022d8efe9c2814c0c7cd5fe921a192}
      \strng{bibnamehash}{fa022d8efe9c2814c0c7cd5fe921a192}
      \strng{authorbibnamehash}{fa022d8efe9c2814c0c7cd5fe921a192}
      \strng{authornamehash}{fa022d8efe9c2814c0c7cd5fe921a192}
      \strng{authorfullhash}{fa022d8efe9c2814c0c7cd5fe921a192}
      \field{sortinit}{1}
      \field{sortinithash}{50c6687d7fc80f50136d75228e3c59ba}
      \field{labelnamesource}{author}
      \field{labeltitlesource}{title}
      \field{annotation}{The contents of this book are divided into the following ten main chapters, plus an Appendix that provides additional background on topics such as partial summation and the $\Gamma$-function. The chapters are: I. Elementary results; II. Sieve methods; III. The prime number theorem; IV. Primes in an arithmetical progression; V. Different applications; VI. The Goldbach problem; VII. Function theoretic properties of the $L$-functions. Explicit formulae and their applications; VIII. Trigonometric sums; IX. Theorems on the density of zeros of the $L$-functions and their application in prime number theory; X. The smallest prime in an arithmetical progression.}
      \field{title}{Primzahlverteilung}
      \field{year}{1957}
      \true{nocite}
      \field{pages}{x+415 pp.}
      \range{pages}{-1}
      \verb{mrnumber}
      \verb 0087685
      \endverb
    \endentry
    \entry{MR0104638}{article}{}
      \name{author}{1}{}{%
        {{hash=4cabd3deed7eb4da21672a1718558bbe}{%
           family={Haselgrove},
           familyi={H\bibinitperiod},
           given={C.\bibnamedelimi B.},
           giveni={C\bibinitperiod\bibinitdelim B\bibinitperiod}}}%
      }
      \strng{namehash}{4cabd3deed7eb4da21672a1718558bbe}
      \strng{fullhash}{4cabd3deed7eb4da21672a1718558bbe}
      \strng{bibnamehash}{4cabd3deed7eb4da21672a1718558bbe}
      \strng{authorbibnamehash}{4cabd3deed7eb4da21672a1718558bbe}
      \strng{authornamehash}{4cabd3deed7eb4da21672a1718558bbe}
      \strng{authorfullhash}{4cabd3deed7eb4da21672a1718558bbe}
      \field{sortinit}{1}
      \field{sortinithash}{50c6687d7fc80f50136d75228e3c59ba}
      \field{labelnamesource}{author}
      \field{labeltitlesource}{title}
      \field{annotation}{Following Ingham's method \cite{MR0006202}, the author resolves Pólya's problem in the negative by showing that a truncated version of the explicit formula for $E^L(x)$, using zeros of $\zeta(s)$ up to height $1,000$, is positive at $x=e^{831.847}$; while this does not rigorously establish that $E^L(e^{831.847})$ is itself positive, the truncated value is a lower bound for $\limsup_{x\to\infty} E^L(x)$. Similarly, the author resolves Turán's problem in the negative by showing that the analogous truncation of the explicit formula for $E^{L_r}(x)$ is negative at $x=e^{853.853}$. The computations were done on an EDSAC~I and a Mark~I. \par This article cites~\cites{MR27305,MR0006202,1897.Mertens,1919.Polya}.}
      \field{issn}{0025-5793}
      \field{journaltitle}{Mathematika}
      \field{title}{A disproof of a conjecture of Pólya}
      \field{volume}{5}
      \field{year}{1958}
      \true{nocite}
      \field{pages}{141\bibrangedash 145}
      \range{pages}{5}
      \verb{doi}
      \verb 10.1112/S0025579300001480
      \endverb
      \verb{mrnumber}
      \verb 0104638
      \endverb
    \endentry
    \entry{MR0096622}{article}{}
      \name{author}{1}{}{%
        {{hash=d33db82fc1e83dc064a23d93e5a21891}{%
           family={Knapowski},
           familyi={K\bibinitperiod},
           given={S.},
           giveni={S\bibinitperiod}}}%
      }
      \strng{namehash}{d33db82fc1e83dc064a23d93e5a21891}
      \strng{fullhash}{d33db82fc1e83dc064a23d93e5a21891}
      \strng{bibnamehash}{d33db82fc1e83dc064a23d93e5a21891}
      \strng{authorbibnamehash}{d33db82fc1e83dc064a23d93e5a21891}
      \strng{authornamehash}{d33db82fc1e83dc064a23d93e5a21891}
      \strng{authorfullhash}{d33db82fc1e83dc064a23d93e5a21891}
      \field{extraname}{1}
      \field{sortinit}{1}
      \field{sortinithash}{50c6687d7fc80f50136d75228e3c59ba}
      \field{labelnamesource}{author}
      \field{labeltitlesource}{title}
      \field{annotation}{This article begins by recalling that $\Delta^\pi(x;q,1) \ll_\ep x^{\theta+\ep}$ implies $\Delta^\pi(x;q,a) \ll_\ep x^{\theta+\ep}$ for all $(a,q)=1$. The author establishes an explicit inequality relating the two, namely that \[ \max_{x\le T} |\Delta^\pi(x;q,a)| \le T^{\delta(T)} \exp\bigg( \frac{(1+q^{-1})\log T}{\sqrt{\log\log T}} \bigg) \big( \max_{x\le T} |\Delta^\pi(x;q,1)| + \phi(q)\sqrt T \big). \] Here $\delta(T) = \ep(\sqrt T)-\ep\big(\exp(\sqrt{\log\log T})\big)$, where $\ep(H)$ is the largest real part of the zeros of $\zeta(s)$ up to height~$H$; in particular, $\delta(T)$ is eventually identically~$0$ assuming~SA, and $\delta(T)\to0$ unconditionally. The author derives this result from the analogous result for $\Delta^\psi(x;q,a)$, which does not contain the term $\phi(q)\sqrt T$; he remarks upon a similar result for $\Delta^\Pi(x;q,a)$ assuming~SA. \par This article cites~\cites{1909.Landau,MR0060548}.}
      \field{issn}{0065-1036}
      \field{journaltitle}{Acta Arith.}
      \field{title}{On prime numbers in an arithmetical progression}
      \field{volume}{4}
      \field{year}{1958}
      \true{nocite}
      \field{pages}{57\bibrangedash 70}
      \range{pages}{14}
      \verb{doi}
      \verb 10.4064/aa-4-1-57-70
      \endverb
      \verb{mrnumber}
      \verb 0096622
      \endverb
      \verb{urlraw}
      \verb https://doi.org/10.4064/aa-4-1-57-7
      \endverb
      \verb{url}
      \verb https://doi.org/10.4064/aa-4-1-57-7
      \endverb
    \endentry
    \entry{MR0096630}{article}{}
      \name{author}{1}{}{%
        {{hash=d33db82fc1e83dc064a23d93e5a21891}{%
           family={Knapowski},
           familyi={K\bibinitperiod},
           given={S.},
           giveni={S\bibinitperiod}}}%
      }
      \strng{namehash}{d33db82fc1e83dc064a23d93e5a21891}
      \strng{fullhash}{d33db82fc1e83dc064a23d93e5a21891}
      \strng{bibnamehash}{d33db82fc1e83dc064a23d93e5a21891}
      \strng{authorbibnamehash}{d33db82fc1e83dc064a23d93e5a21891}
      \strng{authornamehash}{d33db82fc1e83dc064a23d93e5a21891}
      \strng{authorfullhash}{d33db82fc1e83dc064a23d93e5a21891}
      \field{extraname}{2}
      \field{sortinit}{1}
      \field{sortinithash}{50c6687d7fc80f50136d75228e3c59ba}
      \field{labelnamesource}{author}
      \field{labeltitlesource}{title}
      \field{annotation}{The author proves that the assumption $\int_1^T ({M(x)}/x)^2 \, dx \ll \log T$, which is stronger than the Mertens conjecture and is known to imply both RH and the simplicity of the zeros of~$\zeta(s)$, also implies that $M(x) = \Omega\big( x^{1/2}\exp(-{\log x}/{\sqrt{\log\log x}}) \big)$.\par This article cites~\cite{MR0060548}.}
      \field{issn}{0065-1036}
      \field{journaltitle}{Acta Arith.}
      \field{title}{On the Möbius function}
      \field{volume}{4}
      \field{year}{1958}
      \true{nocite}
      \field{pages}{209\bibrangedash 216}
      \range{pages}{8}
      \verb{doi}
      \verb 10.4064/aa-4-3-209-216
      \endverb
      \verb{mrnumber}
      \verb 0096630
      \endverb
    \endentry
    \entry{MR98723}{article}{}
      \name{author}{1}{}{%
        {{hash=86b5f5639455f1050b0d4e2301179477}{%
           family={Wintner},
           familyi={W\bibinitperiod},
           given={Aurel},
           giveni={A\bibinitperiod}}}%
      }
      \strng{namehash}{86b5f5639455f1050b0d4e2301179477}
      \strng{fullhash}{86b5f5639455f1050b0d4e2301179477}
      \strng{bibnamehash}{86b5f5639455f1050b0d4e2301179477}
      \strng{authorbibnamehash}{86b5f5639455f1050b0d4e2301179477}
      \strng{authornamehash}{86b5f5639455f1050b0d4e2301179477}
      \strng{authorfullhash}{86b5f5639455f1050b0d4e2301179477}
      \field{extraname}{5}
      \field{sortinit}{1}
      \field{sortinithash}{50c6687d7fc80f50136d75228e3c59ba}
      \field{labelnamesource}{author}
      \field{labeltitlesource}{title}
      \field{annotation}{The author considers an analogue of the weak Mertens conjecture for the summatory function of the Liouville function. The author shows that $L(x)\ll\sqrt{x}$ is equivalent to ${\sum_{n \leq x} {\lambda(n)}/{\sqrt{n}}=c\log x+O(1)}$, where $c=1/\zeta(1/2)<0$. The proof follows the similar lines as of the author's earlier work~\cite{MR29414}.\par This article cites~\cite{MR0006202,MR29414}.}
      \field{issn}{0002-9327}
      \field{journaltitle}{Amer. J. Math.}
      \field{title}{On the {$\lambda $}-variant of {M}ertens' {$\mu $}-hypothesis}
      \field{volume}{80}
      \field{year}{1958}
      \true{nocite}
      \field{pages}{639\bibrangedash 642}
      \range{pages}{4}
      \verb{mrnumber}
      \verb 98723
      \endverb
    \endentry
    \entry{MR0111722}{article}{}
      \name{author}{1}{}{%
        {{hash=d33db82fc1e83dc064a23d93e5a21891}{%
           family={Knapowski},
           familyi={K\bibinitperiod},
           given={S.},
           giveni={S\bibinitperiod}}}%
      }
      \list{language}{1}{%
        {English, with Russian summary}%
      }
      \strng{namehash}{d33db82fc1e83dc064a23d93e5a21891}
      \strng{fullhash}{d33db82fc1e83dc064a23d93e5a21891}
      \strng{bibnamehash}{d33db82fc1e83dc064a23d93e5a21891}
      \strng{authorbibnamehash}{d33db82fc1e83dc064a23d93e5a21891}
      \strng{authornamehash}{d33db82fc1e83dc064a23d93e5a21891}
      \strng{authorfullhash}{d33db82fc1e83dc064a23d93e5a21891}
      \field{extraname}{3}
      \field{sortinit}{1}
      \field{sortinithash}{50c6687d7fc80f50136d75228e3c59ba}
      \field{labelnamesource}{author}
      \field{labeltitlesource}{title}
      \field{annotation}{The author continues his work in~\cite{MR0096630} by considering lower bounds for the logarithmic averages $A^\psi_{|1|}(x)$ and $A^M_{|1|}(x)$. The author shows that if $\rho_0=\beta_0+i\gamma_0$ is a zero of $\zeta(s)$, then \[ A^\psi_{|1|}(x) > x^{\beta_0}\exp\bigg( -14\frac{\log{x}}{\sqrt{\log\log{x}}} \bigg) \] when $x$ is sufficiently large in terms of~$|\rho_0|$. He also shows that the conjectured upper bound $A^M_{|1|}(x) \ll \sqrt x$ actually implies the lower bound \[ A^M_{|1|}(x) > \sqrt x \exp\bigg( -\frac{\log x}{\sqrt{\log\log x}} \bigg) \] when $x$ is sufficiently large.\par This article cites~\cites{MR0096630, MR0068565, MR0077579, MR0043121, MR0060548}.}
      \field{issn}{0001-5954}
      \field{journaltitle}{Acta Math. Acad. Sci. Hungar.}
      \field{title}{On the mean values of certain functions in prime number theory}
      \field{volume}{10}
      \field{year}{1959}
      \true{nocite}
      \field{pages}{375\bibrangedash 390. (unbound insert)}
      \range{pages}{-1}
      \verb{mrnumber}
      \verb 0111722
      \endverb
    \endentry
    \entry{MR0108470}{article}{}
      \name{author}{1}{}{%
        {{hash=22b786017fa0c36179547815bccb7685}{%
           family={Shanks},
           familyi={S\bibinitperiod},
           given={D.},
           giveni={D\bibinitperiod}}}%
      }
      \strng{namehash}{22b786017fa0c36179547815bccb7685}
      \strng{fullhash}{22b786017fa0c36179547815bccb7685}
      \strng{bibnamehash}{22b786017fa0c36179547815bccb7685}
      \strng{authorbibnamehash}{22b786017fa0c36179547815bccb7685}
      \strng{authornamehash}{22b786017fa0c36179547815bccb7685}
      \strng{authorfullhash}{22b786017fa0c36179547815bccb7685}
      \field{sortinit}{1}
      \field{sortinithash}{50c6687d7fc80f50136d75228e3c59ba}
      \field{labelnamesource}{author}
      \field{labeltitlesource}{title}
      \field{annotation}{In this article, the author investigates Chebyshev's assertion that there are more primes of the form $4m-1$ than of the form $4m+1$. Define $\tau(n)=\pi(n;4,3,1)\sqrt{n}/\pi(n)$ (which is asymptotically equivalent to $\pi(n;4,3,1)\log n/\sqrt n$ but is easier to manipulate numerically). Upon computing $\pi(n;4,3,1)$ for values of $n$ up to 3 million, he analyzes the values $\tau(1000k)$ for $1\le k\le2000$, noting that their histogram is ``roughly normal with a mean of (nearly)~$1$''. The author conjectures that \[ \lim_{x\to\infty} \frac1x \sum_{n\le x} \tau(n) = \lim_{x\to\infty} \frac1x \sum_{n\le x} \frac{\pi(x;4,3,1)\sqrt{n}}{\pi(n)} = 1, \] and notes that weaker versions of the conjecture---namely, that the above limit holds under GRH for $L(s,\chi_{-4})$, or that the above limit either equals $1$ or fails to exist---are also open. (He can show, under GRH, that the mean value inside the limit is positive and bounded away from $0$ for sufficiently large~$x$.) \par Next, the author discusses the distribution of primes in the residue classes modulo $8$, $10$, and $12$. Both from examining the collected data and from combinatorial reasoning involving the multiplicative groups of those moduli, he concludes that the quadratic residues are the ones with a smaller number of primes (on average). For the specific modulus $4$, he outlines an argument, based on combinatorial reasoning with the quantities $\#\{ n\le x\colon \Omega(n)=a,\, n\equiv\pm1\mod4\}$, that shows that the mean value of $\tau(n)$ should be~$1$. Indeed, he remarks that the generalization of this mean value to the integers with $a$ prime factors (counted with multiplicity) predicts that it is the residue class $(-1)^a\mod4$ that should have more such integers; in other words, the bias switches according to the parity of the number of prime factors. A related remark is that there is a bias towards integers for which $(-1)^{\Omega(n)}\chi_{-4}(n)$ equals~$1$ over those for which it equals~$-1$. \par This article ends with a discussion of how similar arguments to those laid out in this article could be used to analyze the relationship between $\pi(x)$ and $\li(x)$. \par This article cites~\cites{1853.Chebyshev, 1891.Phragmen, MR1511360, MR1555148, MR1573970, MR0068565, MR0083001}.}
      \field{issn}{0891-6837}
      \field{journaltitle}{Math. Tables Aids Comput.}
      \field{title}{Quadratic residues and the distribution of primes}
      \field{volume}{13}
      \field{year}{1959}
      \true{nocite}
      \field{pages}{272\bibrangedash 284}
      \range{pages}{13}
      \verb{mrnumber}
      \verb 0108470
      \endverb
    \endentry
    \entry{MR115977}{article}{}
      \name{author}{1}{}{%
        {{hash=e68e101772b88cb1d48e688164bece92}{%
           family={Turán},
           familyi={T\bibinitperiod},
           given={P.},
           giveni={P\bibinitperiod}}}%
      }
      \strng{namehash}{e68e101772b88cb1d48e688164bece92}
      \strng{fullhash}{e68e101772b88cb1d48e688164bece92}
      \strng{bibnamehash}{e68e101772b88cb1d48e688164bece92}
      \strng{authorbibnamehash}{e68e101772b88cb1d48e688164bece92}
      \strng{authornamehash}{e68e101772b88cb1d48e688164bece92}
      \strng{authorfullhash}{e68e101772b88cb1d48e688164bece92}
      \field{extraname}{5}
      \field{sortinit}{1}
      \field{sortinithash}{50c6687d7fc80f50136d75228e3c59ba}
      \field{labelnamesource}{author}
      \field{labeltitlesource}{title}
      \field{annotation}{The main topic of this article is the connection between RH and zeros of partial sums of the series for $\zeta(s)$. Herein, however, the author mentions that the assertion that $L_r(x)$ is nonnegative (eventually or even always) has been called ``Turán's conjecture'', but that he never made such a claim even implicitly. The exact quote, in which ``${}^1$'' refers to~\cite{MR27305} and ``(3.2)'' refers to the inequality $L_r(x) \ge 0$, is: ``{\it Merkwürdigerweise bezeichnetendie an~${}^1$ anschlie{ß}enden Arbeiten schlechthin (3.2) als ,,Tur~insche Vermutung``, sogar für $n\ge1$ behauptet. (3.2) kommt in~${}^1$ nirgends vor, für $n\ge1$ nicht einmal implizite behauptet.}'' \par This article cites~\cites{MR27305,MR0104638,MR1511360,MR0046485}.}
      \field{issn}{0001-5954}
      \field{journaltitle}{Acta Math. Acad. Sci. Hungar.}
      \field{title}{Nachtrag zu meiner {A}bhandlung ``{O}n some approximative {D}irichlet polynomials in the theory of zeta-function of {R}iemann''}
      \field{volume}{10}
      \field{year}{1959}
      \true{nocite}
      \field{pages}{277\bibrangedash 298 (unbound insert)}
      \range{pages}{-1}
      \verb{doi}
      \verb 10.1007/BF02024493
      \endverb
      \verb{mrnumber}
      \verb 115977
      \endverb
      \verb{urlraw}
      \verb https://doi-org.ezproxy.library.ubc.ca/10.1007/BF02024493
      \endverb
      \verb{url}
      \verb https://doi-org.ezproxy.library.ubc.ca/10.1007/BF02024493
      \endverb
    \endentry
    \entry{MR0125822}{article}{}
      \name{author}{1}{}{%
        {{hash=d33db82fc1e83dc064a23d93e5a21891}{%
           family={Knapowski},
           familyi={K\bibinitperiod},
           given={S.},
           giveni={S\bibinitperiod}}}%
      }
      \strng{namehash}{d33db82fc1e83dc064a23d93e5a21891}
      \strng{fullhash}{d33db82fc1e83dc064a23d93e5a21891}
      \strng{bibnamehash}{d33db82fc1e83dc064a23d93e5a21891}
      \strng{authorbibnamehash}{d33db82fc1e83dc064a23d93e5a21891}
      \strng{authornamehash}{d33db82fc1e83dc064a23d93e5a21891}
      \strng{authorfullhash}{d33db82fc1e83dc064a23d93e5a21891}
      \field{extraname}{4}
      \field{sortinit}{1}
      \field{sortinithash}{50c6687d7fc80f50136d75228e3c59ba}
      \field{labelnamesource}{author}
      \field{labeltitlesource}{title}
      \field{annotation}{Assuming GRH$(q^7)$, the author shows that \[ \int_{T\exp(-(\log T)^{3/4})}^T \frac{|{\Delta^\psi(x;q,a)}|}{x} \,dx>{\phi(q)}T^{1/2}\exp\bigg(-\frac{2\log T}{\log\log T}\bigg) \] with $T$ is sufficiently large in terms of~$q$, which in particular implies an $\Omega$-result for $\Delta^\psi(x;q,a)$. Unconditionally the right-hand side can be replaced by $T^{1/4}$. \par This article cites~\cites{MR0096622, MR0087685,MR0060548}.}
      \field{issn}{0065-1036}
      \field{journaltitle}{Acta Arith.}
      \field{sortyear}{1960}
      \field{title}{Contributions to the theory of the distribution of prime numbers in arithmetical progressions. I}
      \field{volume}{6}
      \field{year}{1960/1961}
      \true{nocite}
      \field{pages}{415\bibrangedash 434}
      \range{pages}{20}
      \verb{doi}
      \verb 10.4064/aa-6-4-415-434
      \endverb
      \verb{mrnumber}
      \verb 0125822
      \endverb
    \endentry
    \entry{MR0120198}{article}{}
      \name{author}{1}{}{%
        {{hash=019b07c7177a5ca71ef4716a49411b64}{%
           family={Lehman},
           familyi={L\bibinitperiod},
           given={R.\bibnamedelimi Sherman},
           giveni={R\bibinitperiod\bibinitdelim S\bibinitperiod}}}%
      }
      \strng{namehash}{019b07c7177a5ca71ef4716a49411b64}
      \strng{fullhash}{019b07c7177a5ca71ef4716a49411b64}
      \strng{bibnamehash}{019b07c7177a5ca71ef4716a49411b64}
      \strng{authorbibnamehash}{019b07c7177a5ca71ef4716a49411b64}
      \strng{authornamehash}{019b07c7177a5ca71ef4716a49411b64}
      \strng{authorfullhash}{019b07c7177a5ca71ef4716a49411b64}
      \field{extraname}{1}
      \field{sortinit}{1}
      \field{sortinithash}{50c6687d7fc80f50136d75228e3c59ba}
      \field{labelnamesource}{author}
      \field{labeltitlesource}{title}
      \field{annotation}{The author describes computations leading to an explicit counterexample to Pólya's problem, showing that $L(906,180,359)=1$ is a small value for which the inequality fails. In order to speed up the computation for $L(x)$, which was performed on an IBM 704 at Berkeley, the author uses (a slightly more complicated version of) the recursive formula $$ L(x)=\sum_{m \leq x/m} \mu(m) \bigg \lfloor \sqrt{\frac{x}{m}}\bigg \rfloor -\sum_{k <v} \lambda(k)\bigg( \bigg \lfloor\frac{x}{km}\bigg \rfloor-\bigg \lfloor\frac{x}{mv}\bigg \rfloor \bigg) -\sum_{x/w<l \leq x/v} L \bigg(\frac{x}{l}\bigg) \sum_{\substack{m \mid l \\ m \leq x/w}} \mu(m). $$ The author uses a conditional truncated explicit formula for a weighted variant of $L(e^u)$ to identify promising candidates for positive values of $L(x)$.\par This article cites~\cite{1912.vonSterneck,1919.Polya,MR0006202,1932.Ingham,MR0104638,MR43841,MR51857}.}
      \field{issn}{0025-5718}
      \field{journaltitle}{Math. Comp.}
      \field{title}{On {L}iouville's function}
      \field{volume}{14}
      \field{year}{1960}
      \true{nocite}
      \field{pages}{311\bibrangedash 320}
      \range{pages}{10}
      \verb{doi}
      \verb 10.2307/2003890
      \endverb
      \verb{mrnumber}
      \verb 0120198
      \endverb
      \verb{urlraw}
      \verb https://doi.org/10.2307/2003890
      \endverb
      \verb{url}
      \verb https://doi.org/10.2307/2003890
      \endverb
    \endentry
    \entry{MR115965}{article}{}
      \name{author}{2}{}{%
        {{hash=63196bddebd8e713df6bac51c837c623}{%
           family={Lehmer},
           familyi={L\bibinitperiod},
           given={D.\bibnamedelimi H.},
           giveni={D\bibinitperiod\bibinitdelim H\bibinitperiod}}}%
        {{hash=74d133cbc34e02461a902602a69cfb30}{%
           family={Selberg},
           familyi={S\bibinitperiod},
           given={S.},
           giveni={S\bibinitperiod}}}%
      }
      \strng{namehash}{9aa6cea78280f1fd434f36a0a19ee629}
      \strng{fullhash}{9aa6cea78280f1fd434f36a0a19ee629}
      \strng{bibnamehash}{9aa6cea78280f1fd434f36a0a19ee629}
      \strng{authorbibnamehash}{9aa6cea78280f1fd434f36a0a19ee629}
      \strng{authornamehash}{9aa6cea78280f1fd434f36a0a19ee629}
      \strng{authorfullhash}{9aa6cea78280f1fd434f36a0a19ee629}
      \field{sortinit}{1}
      \field{sortinithash}{50c6687d7fc80f50136d75228e3c59ba}
      \field{labelnamesource}{author}
      \field{labeltitlesource}{title}
      \field{annotation}{Using Landau's Theorem, the authors show that $\fA_1^{M_r}(x)-K$ changes signs infinitely often for any constant~$K$. Numerical calculations show that the first $56$ sign changes of $\fA_1^{M_r}(x)-2$ are nearly in a geometric progression, and a heuristic explanation of this phenomenon is derived from the explicit formula $\fA_1^{M_r}(x)-2=-\sum {x^\rho}/{\rho(1-\rho) \zeta'(\rho)}$.}
      \field{issn}{0065-1036}
      \field{journaltitle}{Acta Arith.}
      \field{title}{A sum involving the function of {M}öbius}
      \field{volume}{6}
      \field{year}{1960}
      \true{nocite}
      \field{pages}{111\bibrangedash 114}
      \range{pages}{4}
      \verb{doi}
      \verb 10.4064/aa-6-1-111-114
      \endverb
      \verb{mrnumber}
      \verb 115965
      \endverb
      \verb{urlraw}
      \verb https://doi.org/10.4064/aa-6-1-111-114
      \endverb
      \verb{url}
      \verb https://doi.org/10.4064/aa-6-1-111-114
      \endverb
    \endentry
    \entry{MR0146153}{article}{}
      \name{author}{1}{}{%
        {{hash=36befda65f61ea63f097e5518d3ec577}{%
           family={Staś},
           familyi={S\bibinitperiod},
           given={W.},
           giveni={W\bibinitperiod}}}%
      }
      \list{language}{1}{%
        {German}%
      }
      \strng{namehash}{36befda65f61ea63f097e5518d3ec577}
      \strng{fullhash}{36befda65f61ea63f097e5518d3ec577}
      \strng{bibnamehash}{36befda65f61ea63f097e5518d3ec577}
      \strng{authorbibnamehash}{36befda65f61ea63f097e5518d3ec577}
      \strng{authornamehash}{36befda65f61ea63f097e5518d3ec577}
      \strng{authorfullhash}{36befda65f61ea63f097e5518d3ec577}
      \field{extraname}{1}
      \field{sortinit}{1}
      \field{sortinithash}{50c6687d7fc80f50136d75228e3c59ba}
      \field{labelnamesource}{author}
      \field{labeltitlesource}{title}
      \field{annotation}{Ingham \cite{1932.Ingham} showed that any precise information on the zero-free region of $\zeta(s)$ would lead to a correspondingly precise estimate on $\Delta^\psi(x)$. The author uses the power-sum method to obtain a partial converse of this theorem. As an illustration, the author shows that $\Delta^\psi(x) \ll x/ (\log x)^{1/10}$ would imply that $\zeta(s) \neq 0$ in the region $\sigma > 1 - \frac{1}{400}(\log t)t^{-20}$ when~$t$ is sufficiently large.\par This article cites~\cites{1932.Ingham,MR0060548}.}
      \field{issn}{0065-1036}
      \field{journaltitle}{Acta Arith.}
      \field{sortyear}{1960}
      \field{title}{Über die Umkehrung eines Satzes von Ingham}
      \field{volume}{6}
      \field{year}{1960/1961}
      \true{nocite}
      \field{pages}{435\bibrangedash 446}
      \range{pages}{12}
      \verb{doi}
      \verb 10.4064/aa-6-4-435-446
      \endverb
      \verb{mrnumber}
      \verb 0146153
      \endverb
    \endentry
    \entry{MR0142520}{article}{}
      \name{author}{1}{}{%
        {{hash=d33db82fc1e83dc064a23d93e5a21891}{%
           family={Knapowski},
           familyi={K\bibinitperiod},
           given={S.},
           giveni={S\bibinitperiod}}}%
      }
      \strng{namehash}{d33db82fc1e83dc064a23d93e5a21891}
      \strng{fullhash}{d33db82fc1e83dc064a23d93e5a21891}
      \strng{bibnamehash}{d33db82fc1e83dc064a23d93e5a21891}
      \strng{authorbibnamehash}{d33db82fc1e83dc064a23d93e5a21891}
      \strng{authornamehash}{d33db82fc1e83dc064a23d93e5a21891}
      \strng{authorfullhash}{d33db82fc1e83dc064a23d93e5a21891}
      \field{extraname}{5}
      \field{sortinit}{1}
      \field{sortinithash}{50c6687d7fc80f50136d75228e3c59ba}
      \field{labelnamesource}{author}
      \field{labeltitlesource}{title}
      \field{annotation}{Assuming GRH$(q^7)$, the author shows that \begin{align*} \int_{T\exp(-(\log T)^{3/4})}^T \frac{|\psi(x;q,a_1,a_2)|}{x} \,dx>T^{1/2}\exp\bigg(-\frac{2\log T}{\log\log T}\bigg) \end{align*} when $T$ is sufficiently large in terms of~$q$, and the analogous result with $\psi$ replaced by~$\Pi$. (Unconditionally the right-hand side can be replaced by $T^{1/4}$.) If both $a_1$ and $a_2$ are nonsquares\mod q, then the same result holds for~$\pi$ as well. \par This article cites~\cites{MR0125822,MR0087685}.}
      \field{issn}{0065-1036}
      \field{journaltitle}{Acta Arith}
      \field{sortyear}{1961}
      \field{title}{Contributions to the theory of the distributon of prime numbers in arithmetical progressions. II}
      \field{volume}{7}
      \field{year}{1961/1962}
      \true{nocite}
      \field{pages}{325\bibrangedash 335}
      \range{pages}{11}
      \verb{doi}
      \verb 10.4064/aa-7-4-325-335
      \endverb
      \verb{mrnumber}
      \verb 0142520
      \endverb
    \endentry
    \entry{MR0133287}{article}{}
      \name{author}{1}{}{%
        {{hash=d33db82fc1e83dc064a23d93e5a21891}{%
           family={Knapowski},
           familyi={K\bibinitperiod},
           given={S.},
           giveni={S\bibinitperiod}}}%
      }
      \strng{namehash}{d33db82fc1e83dc064a23d93e5a21891}
      \strng{fullhash}{d33db82fc1e83dc064a23d93e5a21891}
      \strng{bibnamehash}{d33db82fc1e83dc064a23d93e5a21891}
      \strng{authorbibnamehash}{d33db82fc1e83dc064a23d93e5a21891}
      \strng{authornamehash}{d33db82fc1e83dc064a23d93e5a21891}
      \strng{authorfullhash}{d33db82fc1e83dc064a23d93e5a21891}
      \field{extraname}{6}
      \field{sortinit}{1}
      \field{sortinithash}{50c6687d7fc80f50136d75228e3c59ba}
      \field{labelnamesource}{author}
      \field{labeltitlesource}{title}
      \field{annotation}{Supposing that $A^M_{|1|}(T)<aT^{1/2}$ for $T\ge1$, the author exhibits a constant $H(a)$ such that \[ \int_{T/H(a)}^T\frac{|M(x)|}{x}\,dx>\frac{T^{1/2}}{H(a)}. \] for $T\ge H(a)$, refining a theorem from~\cite{MR0111722} by following an idea due to Ingham in a letter to Turán. The author also states another theorem which is proved later in~\cite{MR0142500}.\par This article cites~\cites{MR0111722,MR0096630,MR0142500}.}
      \field{issn}{0065-1036}
      \field{journaltitle}{Acta Arith.}
      \field{title}{Mean-value estimations for the {M}öbius function. {I}}
      \field{volume}{7}
      \field{year}{1961}
      \true{nocite}
      \field{pages}{121\bibrangedash 130}
      \range{pages}{10}
      \verb{doi}
      \verb 10.4064/aa-7-2-121-130
      \endverb
      \verb{mrnumber}
      \verb 0133287
      \endverb
      \verb{urlraw}
      \verb https://doi.org/10.4064/aa-7-2-121-130
      \endverb
      \verb{url}
      \verb https://doi.org/10.4064/aa-7-2-121-130
      \endverb
    \endentry
    \entry{MR0142500}{article}{}
      \name{author}{1}{}{%
        {{hash=d33db82fc1e83dc064a23d93e5a21891}{%
           family={Knapowski},
           familyi={K\bibinitperiod},
           given={S.},
           giveni={S\bibinitperiod}}}%
      }
      \strng{namehash}{d33db82fc1e83dc064a23d93e5a21891}
      \strng{fullhash}{d33db82fc1e83dc064a23d93e5a21891}
      \strng{bibnamehash}{d33db82fc1e83dc064a23d93e5a21891}
      \strng{authorbibnamehash}{d33db82fc1e83dc064a23d93e5a21891}
      \strng{authornamehash}{d33db82fc1e83dc064a23d93e5a21891}
      \strng{authorfullhash}{d33db82fc1e83dc064a23d93e5a21891}
      \field{extraname}{7}
      \field{sortinit}{1}
      \field{sortinithash}{50c6687d7fc80f50136d75228e3c59ba}
      \field{labelnamesource}{author}
      \field{labeltitlesource}{title}
      \field{annotation}{This article presents a proof of the theorem which was announced in~\cite{MR0133287}: When $T$ is sufficiently large, RH($\omega$) implies that for $T \leq \exp(\omega^{10})$. \[ \int_{T\exp(-100\frac{\log T}{\log\log T}\log\log\log T)}^T\frac{|M(x)|}{x}\,dx>T^{1/2}\exp\biggl(-12\frac{\log T}{\log\log T}\log\log\log T\biggr). \] \par This article cites~\cites{MR0133287,MR0125822,MR0060548}.}
      \field{issn}{0065-1036}
      \field{journaltitle}{Acta Arith.}
      \field{title}{Mean-value estimations for the {M}öbius function. {II}}
      \field{volume}{7}
      \field{year}{1961}
      \true{nocite}
      \field{pages}{337\bibrangedash 343}
      \range{pages}{7}
      \verb{doi}
      \verb 10.4064/aa-7-4-337-343
      \endverb
      \verb{mrnumber}
      \verb 0142500
      \endverb
      \verb{urlraw}
      \verb https://doi.org/10.4064/aa-7-4-337-343
      \endverb
      \verb{url}
      \verb https://doi.org/10.4064/aa-7-4-337-343
      \endverb
    \endentry
    \entry{MR0133309}{article}{}
      \name{author}{1}{}{%
        {{hash=d33db82fc1e83dc064a23d93e5a21891}{%
           family={Knapowski},
           familyi={K\bibinitperiod},
           given={S.},
           giveni={S\bibinitperiod}}}%
      }
      \strng{namehash}{d33db82fc1e83dc064a23d93e5a21891}
      \strng{fullhash}{d33db82fc1e83dc064a23d93e5a21891}
      \strng{bibnamehash}{d33db82fc1e83dc064a23d93e5a21891}
      \strng{authorbibnamehash}{d33db82fc1e83dc064a23d93e5a21891}
      \strng{authornamehash}{d33db82fc1e83dc064a23d93e5a21891}
      \strng{authorfullhash}{d33db82fc1e83dc064a23d93e5a21891}
      \field{extraname}{8}
      \field{sortinit}{1}
      \field{sortinithash}{50c6687d7fc80f50136d75228e3c59ba}
      \field{labelnamesource}{author}
      \field{labeltitlesource}{title}
      \field{annotation}{The author shows that if $\rho_0 = \beta_0 + i \gamma_0$ is any zero of $\zeta(s)$, then for $T$ sufficiently large in terms of~$\gamma_0$, \begin{align*} \Delta^\psi(t) &= \Omega_\pm \bigg( t^{\beta_0} \exp\bigg( \frac{-15\log t}{\sqrt{\log \log t}} \bigg) \bigg), \end{align*} and the same for $\Delta^\Pi(t)$ (which implies an $\Omega_-$-result, though not an $\Omega_+$-result, for $\Delta^\pi(t)$). A slightly more precise version of this result implies that $\liminf_{T\to\infty} W^\psi(T) /\log\log T \ge 1/\log 2$.\par This article cites~\cites{1936.Ingham, MR0043121, MR0067145}.}
      \field{issn}{0024-6107}
      \field{journaltitle}{J. London Math. Soc.}
      \field{title}{On sign-changes in the remainder-term in the prime-number formula}
      \field{volume}{36}
      \field{year}{1961}
      \true{nocite}
      \field{pages}{451\bibrangedash 460}
      \range{pages}{10}
      \verb{doi}
      \verb 10.1112/jlms/s1-36.1.451
      \endverb
      \verb{mrnumber}
      \verb 0133309
      \endverb
      \verb{urlraw}
      \verb https://doi.org/10.1112/jlms/s1-36.1.451
      \endverb
      \verb{url}
      \verb https://doi.org/10.1112/jlms/s1-36.1.451
      \endverb
    \endentry
    \entry{MR0133308}{article}{}
      \name{author}{1}{}{%
        {{hash=d33db82fc1e83dc064a23d93e5a21891}{%
           family={Knapowski},
           familyi={K\bibinitperiod},
           given={S.},
           giveni={S\bibinitperiod}}}%
      }
      \strng{namehash}{d33db82fc1e83dc064a23d93e5a21891}
      \strng{fullhash}{d33db82fc1e83dc064a23d93e5a21891}
      \strng{bibnamehash}{d33db82fc1e83dc064a23d93e5a21891}
      \strng{authorbibnamehash}{d33db82fc1e83dc064a23d93e5a21891}
      \strng{authornamehash}{d33db82fc1e83dc064a23d93e5a21891}
      \strng{authorfullhash}{d33db82fc1e83dc064a23d93e5a21891}
      \field{extraname}{9}
      \field{sortinit}{1}
      \field{sortinithash}{50c6687d7fc80f50136d75228e3c59ba}
      \field{labelnamesource}{author}
      \field{labeltitlesource}{title}
      \field{annotation}{This article is concerned with explicit lower bounds for the number $W(T)$ of sign changes of the function $\li(x)-\pi(x)$. Previously, Ingham~\cite{1936.Ingham} had shown, assuming SA, that $ \liminf_{T \to \infty} {W(T)}/{\log T} >0, $ while the author~\cite{MR0133309} had proved unconditionally that $ \liminf_{T \to \infty} {W(T)}/{\log \log T} >0. $ Skewes~\cite{MR0067145} famously found that $W\big(e^{e^{e^{e^{7.705}}}}\big) \ge 1$. \par In this article, the author shows unconditionally that $W(T) \ge e^{-35} \log \log \log \log T$ for $T \ge e^{e^{e^{e^{35}}}}$ (the author did not try to optimize these constants). Similar to~\cite{MR0067145}, the proof is divided into two cases, first when RH is ``nearly true'' and then the contrary case. More specifically, setting $X = \sqrt[7]{\log \log T}$, the author defines a hypothesis (C) (the ``nearly true'' case) as follows: \emph{Every zero $\rho = \beta + i\gamma$ for which $|\gamma| \le X^3$ satisfies $\beta - \frac{1}{2} \le 2/(3X^3 \log X)$}. The author then proves the inequality above first assuming (C) and then again assuming its negation (NC).\par This article cites~\cites{1914.Littlewood, 1932.Ingham, MR0067145, MR0133309}.}
      \field{issn}{0065-1036}
      \field{journaltitle}{Acta Arith.}
      \field{sortyear}{1961}
      \field{title}{On sign-changes of the difference {$\pi (x)-{\rm li}\,x$}}
      \field{volume}{7}
      \field{year}{1961/1962}
      \true{nocite}
      \field{pages}{107\bibrangedash 119}
      \range{pages}{13}
      \verb{doi}
      \verb 10.4064/aa-7-2-107-119
      \endverb
      \verb{mrnumber}
      \verb 0133308
      \endverb
      \verb{urlraw}
      \verb https://doi.org/10.4064/aa-7-2-107-119
      \endverb
      \verb{url}
      \verb https://doi.org/10.4064/aa-7-2-107-119
      \endverb
    \endentry
    \entry{MR0131410}{article}{}
      \name{author}{2}{}{%
        {{hash=d33db82fc1e83dc064a23d93e5a21891}{%
           family={Knapowski},
           familyi={K\bibinitperiod},
           given={S.},
           giveni={S\bibinitperiod}}}%
        {{hash=36befda65f61ea63f097e5518d3ec577}{%
           family={Staś},
           familyi={S\bibinitperiod},
           given={W.},
           giveni={W\bibinitperiod}}}%
      }
      \strng{namehash}{3d6f57664844e9e6d022d0b7cd9c3631}
      \strng{fullhash}{3d6f57664844e9e6d022d0b7cd9c3631}
      \strng{bibnamehash}{3d6f57664844e9e6d022d0b7cd9c3631}
      \strng{authorbibnamehash}{3d6f57664844e9e6d022d0b7cd9c3631}
      \strng{authornamehash}{3d6f57664844e9e6d022d0b7cd9c3631}
      \strng{authorfullhash}{3d6f57664844e9e6d022d0b7cd9c3631}
      \field{sortinit}{1}
      \field{sortinithash}{50c6687d7fc80f50136d75228e3c59ba}
      \field{labelnamesource}{author}
      \field{labeltitlesource}{title}
      \field{annotation}{The authors examine the function $\Delta^{\psi_e}(x)=\sum_{n=1}^{\infty}(\Lambda(n)-1)e^{-n/x}$. They prove, unconditionally and effectively, that $$ \max_{1\le y\le x} |\Delta^{\psi_e}(y)| > x^{1/2} \exp\bigg( - \frac{4\log x\log \log \log x}{\log \log x} \bigg) $$ when $x$ is sufficiently large. The main tool of the proof is a version of the second main theorem of the power-sum method. \par This article cites~\cites{MR1555148,MR0125822,MR0060548}.}
      \field{issn}{0065-1036}
      \field{journaltitle}{Acta Arith.}
      \field{sortyear}{1961}
      \field{title}{A note on a theorem of Hardy and Littlewood}
      \field{volume}{7}
      \field{year}{1961/1962}
      \true{nocite}
      \field{pages}{161\bibrangedash 166}
      \range{pages}{6}
      \verb{doi}
      \verb 10.4064/aa-7-2-161-166
      \endverb
      \verb{mrnumber}
      \verb 0131410
      \endverb
    \endentry
    \entry{MR0132728}{article}{}
      \name{author}{1}{}{%
        {{hash=e68e101772b88cb1d48e688164bece92}{%
           family={Turán},
           familyi={T\bibinitperiod},
           given={P.},
           giveni={P\bibinitperiod}}}%
      }
      \list{language}{1}{%
        {English, with Russian summary}%
      }
      \strng{namehash}{e68e101772b88cb1d48e688164bece92}
      \strng{fullhash}{e68e101772b88cb1d48e688164bece92}
      \strng{bibnamehash}{e68e101772b88cb1d48e688164bece92}
      \strng{authorbibnamehash}{e68e101772b88cb1d48e688164bece92}
      \strng{authornamehash}{e68e101772b88cb1d48e688164bece92}
      \strng{authorfullhash}{e68e101772b88cb1d48e688164bece92}
      \field{extraname}{6}
      \field{sortinit}{1}
      \field{sortinithash}{50c6687d7fc80f50136d75228e3c59ba}
      \field{labelnamesource}{author}
      \field{labeltitlesource}{title}
      \field{annotation}{The author establishes a ``one-sided'' version of the second main theorem of the power-sum method with an argument restriction. Let $m$ be a nonnegative integer and choose $0 < \alpha \leq \frac{\pi}2 $, and suppose that $z_1,\ldots,z_n$ are complex numbers satisfying $1=|z_1| \geq |z_2| \geq \cdots \geq |z_n|$ and $|\arg(z_j)| > \alpha$, and that $b_1,\ldots,b_n$ are complex numbers with $\min_{1\le\mu\le n} \Re \sum_{j=1}^{\mu} b_j > 0$. Then there exists an integer $m+1 \le v \le m+n(3+\frac\pi\alpha)$ such that $$ \Re \sum_{j=1}^{n} b_jz_j^v \ge \frac1{2n+1}\bigg( \frac{n}{24e^3(m+n(3+\frac\pi\alpha))} \bigg)^{2n} \min_{1\le\mu\le n} \Re \sum_{j=1}^{\mu} b_j, $$ with the analogous result for $-\Re \sum_{j=1}^{n} b_jz_j^v$. \par This article cites~\cites{MR0060548}.}
      \field{issn}{0001-5954}
      \field{journaltitle}{Acta Math. Acad. Sci. Hungar.}
      \field{title}{On some further one-sided theorems of new type in the theory of Diophantine approximations}
      \field{volume}{12}
      \field{year}{1961}
      \true{nocite}
      \field{pages}{455\bibrangedash 468}
      \range{pages}{14}
      \verb{doi}
      \verb 10.1007/BF02023928
      \endverb
      \verb{mrnumber}
      \verb 0132728
      \endverb
    \endentry
    \entry{MR0142521}{article}{}
      \name{author}{1}{}{%
        {{hash=d33db82fc1e83dc064a23d93e5a21891}{%
           family={Knapowski},
           familyi={K\bibinitperiod},
           given={S.},
           giveni={S\bibinitperiod}}}%
      }
      \strng{namehash}{d33db82fc1e83dc064a23d93e5a21891}
      \strng{fullhash}{d33db82fc1e83dc064a23d93e5a21891}
      \strng{bibnamehash}{d33db82fc1e83dc064a23d93e5a21891}
      \strng{authorbibnamehash}{d33db82fc1e83dc064a23d93e5a21891}
      \strng{authornamehash}{d33db82fc1e83dc064a23d93e5a21891}
      \strng{authorfullhash}{d33db82fc1e83dc064a23d93e5a21891}
      \field{extraname}{10}
      \field{sortinit}{1}
      \field{sortinithash}{50c6687d7fc80f50136d75228e3c59ba}
      \field{labelnamesource}{author}
      \field{labeltitlesource}{title}
      \field{annotation}{Assuming GRH$(q^7)$, the author shows that \begin{align*} \int_{T\exp(-(\log T)^{3/4})}^T \frac{|\pi(x;q,a_1,a_2)|}{x} \,dx>T^{1/2}\exp\bigg(-\frac{7\log T}{\log\log T}\bigg); \end{align*} unlike in his earlier work, one no longer needs to assume that $a_1$ and $a_2$ are nonsquares. \par This article cites~\cites{MR0125822,MR0142520,MR0087685,MR0060548}.}
      \field{issn}{0065-1036}
      \field{journaltitle}{Acta Arith}
      \field{sortyear}{1962}
      \field{title}{Contributions to the theory of the distributon of prime numbers in arithmetical progressions. III}
      \field{volume}{8}
      \field{year}{1962/1963}
      \true{nocite}
      \field{pages}{97\bibrangedash 105}
      \range{pages}{9}
      \verb{doi}
      \verb 10.4064/aa-8-1-97-105
      \endverb
      \verb{mrnumber}
      \verb 0142521
      \endverb
    \endentry
    \entry{MR0155802}{article}{}
      \name{author}{1}{}{%
        {{hash=d33db82fc1e83dc064a23d93e5a21891}{%
           family={Knapowski},
           familyi={K\bibinitperiod},
           given={S.},
           giveni={S\bibinitperiod}}}%
      }
      \strng{namehash}{d33db82fc1e83dc064a23d93e5a21891}
      \strng{fullhash}{d33db82fc1e83dc064a23d93e5a21891}
      \strng{bibnamehash}{d33db82fc1e83dc064a23d93e5a21891}
      \strng{authorbibnamehash}{d33db82fc1e83dc064a23d93e5a21891}
      \strng{authornamehash}{d33db82fc1e83dc064a23d93e5a21891}
      \strng{authorfullhash}{d33db82fc1e83dc064a23d93e5a21891}
      \field{extraname}{11}
      \field{sortinit}{1}
      \field{sortinithash}{50c6687d7fc80f50136d75228e3c59ba}
      \field{labelnamesource}{author}
      \field{labeltitlesource}{title}
      \field{annotation}{Assuming RH$(H)$ (but with no assumptions on the simplicity of zeros of $\zeta(s)$), the author shows that \[ \max_{1\le x\le T} M(x) > \sqrt T \exp\bigg( -\frac{15\log T\log\log\log T}{\log\log T} \bigg) \] for $T\le e^{H^{10}}$, and similarly for $-M(x)$; in the same range he deduces that $W(M,T) \gg \log T$. \par This article cites~\cites{MR0096630,MR0111722,MR0133309,MR0146158,MR0132728,MR0133287}.}
      \field{issn}{0065-1036}
      \field{journaltitle}{Acta Arith.}
      \field{sortyear}{1962}
      \field{title}{On oscillations of certain means formed from the Möbius series. I}
      \field{volume}{8}
      \field{year}{1962/1963}
      \true{nocite}
      \field{pages}{311\bibrangedash 320}
      \range{pages}{10}
      \verb{doi}
      \verb 10.4064/aa-8-3-311-320
      \endverb
      \verb{mrnumber}
      \verb 0155802
      \endverb
    \endentry
    \entry{MR0146156}{article}{}
      \name{author}{2}{}{%
        {{hash=d33db82fc1e83dc064a23d93e5a21891}{%
           family={Knapowski},
           familyi={K\bibinitperiod},
           given={S.},
           giveni={S\bibinitperiod}}}%
        {{hash=e68e101772b88cb1d48e688164bece92}{%
           family={Turán},
           familyi={T\bibinitperiod},
           given={P.},
           giveni={P\bibinitperiod}}}%
      }
      \strng{namehash}{b6f197e0531d2bbfab6a59c048e3b28d}
      \strng{fullhash}{b6f197e0531d2bbfab6a59c048e3b28d}
      \strng{bibnamehash}{b6f197e0531d2bbfab6a59c048e3b28d}
      \strng{authorbibnamehash}{b6f197e0531d2bbfab6a59c048e3b28d}
      \strng{authornamehash}{b6f197e0531d2bbfab6a59c048e3b28d}
      \strng{authorfullhash}{b6f197e0531d2bbfab6a59c048e3b28d}
      \field{extraname}{1}
      \field{sortinit}{1}
      \field{sortinithash}{50c6687d7fc80f50136d75228e3c59ba}
      \field{labelnamesource}{author}
      \field{labeltitlesource}{title}
      \field{annotation}{The authors start by introducing ten problems of interest in ``comparative prime-number theory'' to the modulus~$k$, the first seven concerning the sign changes and extreme values of $\pi(x;k,\ell_1,\ell_2)$ and the natural density of the solutions to $\pi(x;k,\ell_1,\ell_2)>0$. The eighth problem, which the authors call the ``race-problem of Shanks--Rényi'' (which is perhaps the first time Rényi's name was linked to comparative prime number theory) is whether there are arbitrarily large solutions $x$ to $\pi(x;k,\ell_1) < \cdots < \pi(x;k,\ell_{\phi(k)})$; the last two problems concern the simultaneous inequalities $\pi(x;k,\ell_j) > \frac1{\phi(k)} \li(x)$.\par The authors allude to variants of these ten problems generated by replacing $\pi(x;k,\ell)$ by $\pi_e(x;k,\ell)$ (and, where needed, $\Li(x)$ by $\int_2^{\infty} \frac{e^{-t/x}}{\log{t}} \,dt$), and further vary these problems by replacing $\pi$ with $\psi$ or~$\Pi$. They are aware that such problems could be further varied (``\ldots{} the analogous problems concerning the distribution of primes in binary quadratic forms with fixed discriminant or of the prime ideals of a fixed field in various idealclasses'').\par In Section~4, the authors discuss how some of the problems involving $\psi(x;k,\ell)$ can be conditionally solved using Landau's theorem (and hence unconditionally for moduli dividing $24$). In Section~5, the authors discuss the results they have so far concerning the problems involving~$\pi(x;k,\ell)$, and in Section~8, they briefly discuss what is known about prime number races modulo~$4$. Throughout the rest of this article, the authors introduce the results that will be proved in the next seven articles of the series.\par This article cites~\cites{1853.Chebyshev, 1891.Phragmen, MR1555148, 1930.Polya, 1936.Ingham, MR0067145, MR0108470}.}
      \field{issn}{0001-5954}
      \field{journaltitle}{Acta Math. Acad. Sci. Hungar.}
      \field{title}{Comparative prime-number theory. {I}. {I}ntroduction}
      \field{volume}{13}
      \field{year}{1962}
      \true{nocite}
      \field{pages}{299\bibrangedash 314}
      \range{pages}{16}
      \verb{doi}
      \verb 10.1007/BF02020796
      \endverb
      \verb{mrnumber}
      \verb 0146156
      \endverb
      \verb{urlraw}
      \verb https://doi.org/10.1007/BF02020796
      \endverb
      \verb{url}
      \verb https://doi.org/10.1007/BF02020796
      \endverb
    \endentry
    \entry{MR0146157}{article}{}
      \name{author}{2}{}{%
        {{hash=d33db82fc1e83dc064a23d93e5a21891}{%
           family={Knapowski},
           familyi={K\bibinitperiod},
           given={S.},
           giveni={S\bibinitperiod}}}%
        {{hash=e68e101772b88cb1d48e688164bece92}{%
           family={Turán},
           familyi={T\bibinitperiod},
           given={P.},
           giveni={P\bibinitperiod}}}%
      }
      \strng{namehash}{b6f197e0531d2bbfab6a59c048e3b28d}
      \strng{fullhash}{b6f197e0531d2bbfab6a59c048e3b28d}
      \strng{bibnamehash}{b6f197e0531d2bbfab6a59c048e3b28d}
      \strng{authorbibnamehash}{b6f197e0531d2bbfab6a59c048e3b28d}
      \strng{authornamehash}{b6f197e0531d2bbfab6a59c048e3b28d}
      \strng{authorfullhash}{b6f197e0531d2bbfab6a59c048e3b28d}
      \field{extraname}{2}
      \field{sortinit}{1}
      \field{sortinithash}{50c6687d7fc80f50136d75228e3c59ba}
      \field{labelnamesource}{author}
      \field{labeltitlesource}{title}
      \field{annotation}{This article focuses on the race between the residue class $1\mod k$ and other residue classes $\ell\not\equiv 1 \mod k$, for a fixed modulus~$k$ for which HC holds.\par Fix a character $\chi\mod k$ such that $\chi(\ell)\ne1$, and let $\rho_0=\beta_0+i\gamma_0$ be a zero of $L(s,\chi)$. The authors prove that for $T$ large enough, \begin{align*} \max_{T^{1/3}\le x\le T}\psi(x;k,1,\ell) &> T^{\beta_0}\exp \left(-41\frac{\log T \log\log\log T}{\log \log T}\right)\\ \max_{T^{1/3}\le x\le T}\Pi(x;k,1,\ell) &> T^{\beta_0}\exp \left(-41\frac{\log T \log\log\log T}{\log \log T}\right) \end{align*} and symmetric results for the minimum. As one might expect, the oscillations obtained for $\pi(x;k,1,\ell)$ are worse, but the authors do prove \[ \max_{\exp (\log_3^{1/130}T)\le x\le T}\left(\frac{\log x}{\sqrt x}\right)\pi(x;k,1,\ell)> \frac{1}{100} \log\log\log\log\log T \] and the symmetric result for the minimum. Each of these results yields lower bounds on the corresponding number of sign changes, although for $\pi(x;k,1,\ell)$ the bound is very low---improving this bound is addressed directly in~\cite{MR0146158}. The authors' methods can also compare primes congruent to $1\mod k$ to the average number of primes in other residue classes\mod k: \[ \max_{\exp (\log_3^{1/130}T)\le x\le T}\left(\frac{\log x}{\sqrt x}\right)\bigg(\pi(x;k,1)-\frac1{\phi(k)-1}\sum_{\substack{(\ell,k)=1\\\ell\neq 1}}\pi(x;k,\ell)\bigg)> \frac{1}{100} \log_5 T \] and the symmetric result for the minimum.\par The proofs follow from the application of power-sum method for bounding exponential sums. Siegel's theorem on the existence of zeros in certain rectangles, coupled with the verification of HC for certain moduli up to~$24$, give for these moduli the first unconditional results about the size of the fluctuations of the above functions and the number of their sign changes.\par This article cites~\cites{1914.Littlewood, MR1555148, 1936.Ingham, MR0067145, MR0133309}.}
      \field{issn}{0001-5954}
      \field{journaltitle}{Acta Math. Acad. Sci. Hungar.}
      \field{title}{Comparative prime-number theory. {II}. {C}omparison of the progressions {$\equiv 1$} {${\rm mod}\ k$} and {$\equiv l$} {${\rm mod}\ k,\,l\not\equiv 1$} {${\rm mod}\ k$}}
      \field{volume}{13}
      \field{year}{1962}
      \true{nocite}
      \field{pages}{315\bibrangedash 342}
      \range{pages}{28}
      \verb{doi}
      \verb 10.1007/BF02020797
      \endverb
      \verb{mrnumber}
      \verb 0146157
      \endverb
      \verb{urlraw}
      \verb https://doi.org/10.1007/BF02020797
      \endverb
      \verb{url}
      \verb https://doi.org/10.1007/BF02020797
      \endverb
    \endentry
    \entry{MR0146158}{article}{}
      \name{author}{2}{}{%
        {{hash=d33db82fc1e83dc064a23d93e5a21891}{%
           family={Knapowski},
           familyi={K\bibinitperiod},
           given={S.},
           giveni={S\bibinitperiod}}}%
        {{hash=e68e101772b88cb1d48e688164bece92}{%
           family={Turán},
           familyi={T\bibinitperiod},
           given={P.},
           giveni={P\bibinitperiod}}}%
      }
      \strng{namehash}{b6f197e0531d2bbfab6a59c048e3b28d}
      \strng{fullhash}{b6f197e0531d2bbfab6a59c048e3b28d}
      \strng{bibnamehash}{b6f197e0531d2bbfab6a59c048e3b28d}
      \strng{authorbibnamehash}{b6f197e0531d2bbfab6a59c048e3b28d}
      \strng{authornamehash}{b6f197e0531d2bbfab6a59c048e3b28d}
      \strng{authorfullhash}{b6f197e0531d2bbfab6a59c048e3b28d}
      \field{extraname}{3}
      \field{sortinit}{1}
      \field{sortinithash}{50c6687d7fc80f50136d75228e3c59ba}
      \field{labelnamesource}{author}
      \field{labeltitlesource}{title}
      \field{annotation}{The authors continue their comparison of $\pi(x;k,\ell)$ and $\pi(x;k,1)$, where $\ell\not\equiv1\mod k$ and $(\ell,k)=1$, assuming HC for the modulus~$k$. They show that $W_{k;\ell,1}(T) > k^{-c}\log\log\log\log T$ for sufficiently large~$T$, where $c$ is an absolute effective constant. The proof technique involves Dirichlet's box principle and bounds obtained on $\Pi(x;k,\ell,1)$ in~\cite{MR0146157}.\par When $\ell$ is a quadratic residue\mod k, they also show that if $\rho_0=\beta_0+i\gamma_0$ is a zero of $L(s,\chi)$ for some character $\chi$ such that $\chi(\ell)\neq 1$, then for sufficiently large~$T$, \[ \max_{T^{1/3}\leq x \leq T} \pi(x;k,\ell,1) >T^{\beta_0} \exp\left(-42 \frac{\log T\log \log \log T}{\log \log T}\right),\] and a similar statement holds for the minimum; consequently, for such $k$ and $\ell$, the inequality $W_{k;\ell,1}(T) > \frac1{\log 3} \log\log T + O(1)$ holds for sufficiently large~$T$. The proof involves power-sum method for bounds on exponential sums.\par The authors remark that both theorems hold as well for $W\big(\Delta^\pi(x;k,1);T\big)$ and $W\big(\Deltadot^\pi(x;k,1);T\big)$.\par This article cites~\cites{MR1555148,1936.Ingham,MR0146156,MR0146157,MR1573970}.}
      \field{issn}{0001-5954}
      \field{journaltitle}{Acta Math. Acad. Sci. Hungar.}
      \field{title}{Comparative prime-number theory. {III}. {C}ontinuation of the study of comparison of the progressions {$\equiv 1$} {${\rm mod}\ k$} and {$\equiv l$} {${\rm mod}\ k$}}
      \field{volume}{13}
      \field{year}{1962}
      \true{nocite}
      \field{pages}{343\bibrangedash 364}
      \range{pages}{22}
      \verb{doi}
      \verb 10.1007/BF02020798
      \endverb
      \verb{mrnumber}
      \verb 0146158
      \endverb
      \verb{urlraw}
      \verb https://doi.org/10.1007/BF02020798
      \endverb
      \verb{url}
      \verb https://doi.org/10.1007/BF02020798
      \endverb
    \endentry
    \entry{MR0137689}{article}{}
      \name{author}{2}{}{%
        {{hash=14b76cd756023555a4b6632dbee72243}{%
           family={Rosser},
           familyi={R\bibinitperiod},
           given={J.\bibnamedelimi Barkley},
           giveni={J\bibinitperiod\bibinitdelim B\bibinitperiod}}}%
        {{hash=197eb6ee8cad0f814cbfe07e59001ef3}{%
           family={Schoenfeld},
           familyi={S\bibinitperiod},
           given={Lowell},
           giveni={L\bibinitperiod}}}%
      }
      \strng{namehash}{748c2386a8972928ef3f5dc56357fa0e}
      \strng{fullhash}{748c2386a8972928ef3f5dc56357fa0e}
      \strng{bibnamehash}{748c2386a8972928ef3f5dc56357fa0e}
      \strng{authorbibnamehash}{748c2386a8972928ef3f5dc56357fa0e}
      \strng{authornamehash}{748c2386a8972928ef3f5dc56357fa0e}
      \strng{authorfullhash}{748c2386a8972928ef3f5dc56357fa0e}
      \field{sortinit}{1}
      \field{sortinithash}{50c6687d7fc80f50136d75228e3c59ba}
      \field{labelnamesource}{author}
      \field{labeltitlesource}{title}
      \field{annotation}{Sprinkled among the many explicit inequalities for prime counting functions are a few comments directly related to comparative prime number theory, particularly on pages 72--73. The authors point out that it is unknown whether $n/\phi(n) \ge e^{C_0}\log\log n$ infinitely often, where $C_0$ is Euler's constant. They also speculate as to the existence of counterexamples to the inequalities \begin{align*} \log\log x + B &< \sum_{p\le x} \frac1p < \log\log x + B + \frac2{\sqrt x\log x} \\ \log x + E &< \sum_{p\le x} \frac{\log p}p < \log x + E + \frac{2.06123}{\sqrt x} \\ e^{C_0}\log x &< \prod_{p\le x} \frac p{p-1} < e^{C_0} \bigg( \log x + \frac2{\sqrt x} \bigg), \end{align*} where the constant $B$ is the unique real number with the property that the differences among the expressions on the first line tends to~$0$, and similarly with $E$ and the second line.}
      \field{issn}{0019-2082}
      \field{journaltitle}{Illinois J. Math.}
      \field{title}{Approximate formulas for some functions of prime numbers}
      \field{volume}{6}
      \field{year}{1962}
      \true{nocite}
      \field{pages}{64\bibrangedash 94}
      \range{pages}{31}
      \verb{mrnumber}
      \verb 0137689
      \endverb
    \endentry
    \entry{MR0146159}{article}{}
      \name{author}{2}{}{%
        {{hash=d33db82fc1e83dc064a23d93e5a21891}{%
           family={Knapowski},
           familyi={K\bibinitperiod},
           given={S.},
           giveni={S\bibinitperiod}}}%
        {{hash=e68e101772b88cb1d48e688164bece92}{%
           family={Turán},
           familyi={T\bibinitperiod},
           given={P.},
           giveni={P\bibinitperiod}}}%
      }
      \strng{namehash}{b6f197e0531d2bbfab6a59c048e3b28d}
      \strng{fullhash}{b6f197e0531d2bbfab6a59c048e3b28d}
      \strng{bibnamehash}{b6f197e0531d2bbfab6a59c048e3b28d}
      \strng{authorbibnamehash}{b6f197e0531d2bbfab6a59c048e3b28d}
      \strng{authornamehash}{b6f197e0531d2bbfab6a59c048e3b28d}
      \strng{authorfullhash}{b6f197e0531d2bbfab6a59c048e3b28d}
      \field{extraname}{4}
      \field{sortinit}{1}
      \field{sortinithash}{50c6687d7fc80f50136d75228e3c59ba}
      \field{labelnamesource}{author}
      \field{labeltitlesource}{title}
      \field{annotation}{The authors apply techniques from earlier in this series of articles to the modulus $k=8$, when $\ell_1,\ell_2 \in \{3,5,7\}$ are distinct quadratic nonresidues. They show that \[ \max_{T^{1/3}\leq x \leq T} \pi(x;8,\ell_1,\ell_2) > \sqrt T \exp\left(-23 \frac{\log T\log \log \log T}{\log \log T}\right),\] and similarly for $\Pi(x; 8, \ell_1,\ell_2)$ and $\psi(x; 8, \ell_1, \ell_2)$. Since $\ell_1$ and $\ell_2$ can be interchanged, this result implies the inequality $W_{8;\ell_1,\ell_2}(T) > \frac1{\log 3} \log\log T + O(1)$, and similarly for $W^\Pi_{8;\ell_1,\ell_2}(T)$ and $W^\psi_{8;\ell_1,\ell_2}(T)$.\par In the third section, the authors remark on the modulus $k=5$. The case $(\ell_1, \ell_2)=(\ell,1)$ has already been handled earlier in this series; the authors mention that the case $(\ell_1, \ell_2)=(2,3)$, where both are quadratic nonresidues, can be handled in a similar way to the $k=8$ cases treated in this article. The remaining cases have $\ell_1=4$, a quadratic residue not equal to $1\mod5$, and $\ell_2\in\{2,3\}$, and these cases yield ``an unpleasant (or pleasant?)~surprise'': the authors cannot establish sign changes for $\pi(x;5, 4, 2)$ and $\pi(x;5, 4, 3)$ even assuming GRH (although the methods do work for the $\Pi$ and $\psi$ versions, a situation the authors discuss further in the later articles of this series).\par This article cites~\cites{MR0146156,MR0146157, MR0146158}.}
      \field{issn}{0001-5954}
      \field{journaltitle}{Acta Math. Acad. Sci. Hungar.}
      \field{title}{Comparative prime-number theory. {IV}. {P}aradigma to the general case, {$k=8$} and {$5$}}
      \field{volume}{14}
      \field{year}{1963}
      \true{nocite}
      \field{pages}{31\bibrangedash 42}
      \range{pages}{12}
      \verb{doi}
      \verb 10.1007/BF01901928
      \endverb
      \verb{mrnumber}
      \verb 0146159
      \endverb
      \verb{urlraw}
      \verb https://doi.org/10.1007/BF01901928
      \endverb
      \verb{url}
      \verb https://doi.org/10.1007/BF01901928
      \endverb
    \endentry
    \entry{MR0146160}{article}{}
      \name{author}{2}{}{%
        {{hash=d33db82fc1e83dc064a23d93e5a21891}{%
           family={Knapowski},
           familyi={K\bibinitperiod},
           given={S.},
           giveni={S\bibinitperiod}}}%
        {{hash=e68e101772b88cb1d48e688164bece92}{%
           family={Turán},
           familyi={T\bibinitperiod},
           given={P.},
           giveni={P\bibinitperiod}}}%
      }
      \strng{namehash}{b6f197e0531d2bbfab6a59c048e3b28d}
      \strng{fullhash}{b6f197e0531d2bbfab6a59c048e3b28d}
      \strng{bibnamehash}{b6f197e0531d2bbfab6a59c048e3b28d}
      \strng{authorbibnamehash}{b6f197e0531d2bbfab6a59c048e3b28d}
      \strng{authornamehash}{b6f197e0531d2bbfab6a59c048e3b28d}
      \strng{authorfullhash}{b6f197e0531d2bbfab6a59c048e3b28d}
      \field{extraname}{5}
      \field{sortinit}{1}
      \field{sortinithash}{50c6687d7fc80f50136d75228e3c59ba}
      \field{labelnamesource}{author}
      \field{labeltitlesource}{title}
      \field{annotation}{Under the assumption of a ``finite Riemann--Piltz conjecture'' GRH$(H,E_k)$, the authors establish the following result for any distinct reduced residues $\ell_1,\ell_2\mod k$ and for $T$ sufficiently large (explicitly quantified in the article): \[ \max_{T^{1/3} \leq x \leq T } \Pi(x;k,\ell_1,\ell_2) > \sqrt{T}\exp \bigg( {-44\frac{\log{T}\log\log\log{T}}{\log\log{T}}} \bigg),\] and the same with $\Pi$ replaced by~$\psi$. A central element to their proof is the estimation of the integral \begin{align*} J(T) &=-\frac{1}{2\pi i}\int_2 \left(\frac{e^{y_1s}}{s}\right)^v\frac{(\omega_0L_1^{v_0})^s}{s^{v_0+1}}\cdot \frac{1}{\phi(k)} \bigg\{\sum_{\substack{\chi\mod k \\ \chi(\ell_1)\neq \chi(\ell_2)}}\left(\overline{\chi}(\ell_1)-\overline{\chi}(\ell_2)\right)\frac{L'}{L}(s,\chi)\bigg\}\ ds \\ &= \frac{1}{(v+v_0)!}\int_1^{Y_1} \Pi(x;k,\ell_1,\ell_2) \frac d{dx} \left( \bigg( \log\frac{Y_1}{x} \bigg)^{v+v_0} \log x \right) \, dx. \end{align*} Their approach involves an application of power-sum method similar to what appears in the previous articles of the series. The authors note that their main theorem gives similar bounds on $\Delta^\Pi(x;k,\ell)/\phi(k)$ and $\Delta^\psi(x;k,\ell)/\phi(k)$. In particular, this result implies the lower bounds $W^\Pi_{k;\ell_1,\ell_2}(T) > \frac1{\log 3}\log\log T + O(1)$ and $W^\psi_{k;\ell_1,\ell_2}(T) > \frac1{\log 3}\log\log T + O(1)$.\par This article cites~\cites{MR0125822,MR0146156,MR0146157,MR0146158,MR0146159}.}
      \field{issn}{0001-5954}
      \field{journaltitle}{Acta Math. Acad. Sci. Hungar.}
      \field{title}{Comparative prime-number theory. {V}. {S}ome theorems concerning the general case}
      \field{volume}{14}
      \field{year}{1963}
      \true{nocite}
      \field{pages}{43\bibrangedash 63}
      \range{pages}{21}
      \verb{doi}
      \verb 10.1007/BF01901929
      \endverb
      \verb{mrnumber}
      \verb 0146160
      \endverb
      \verb{urlraw}
      \verb https://doi.org/10.1007/BF01901929
      \endverb
      \verb{url}
      \verb https://doi.org/10.1007/BF01901929
      \endverb
    \endentry
    \entry{MR0146161}{article}{}
      \name{author}{2}{}{%
        {{hash=d33db82fc1e83dc064a23d93e5a21891}{%
           family={Knapowski},
           familyi={K\bibinitperiod},
           given={S.},
           giveni={S\bibinitperiod}}}%
        {{hash=e68e101772b88cb1d48e688164bece92}{%
           family={Turán},
           familyi={T\bibinitperiod},
           given={P.},
           giveni={P\bibinitperiod}}}%
      }
      \strng{namehash}{b6f197e0531d2bbfab6a59c048e3b28d}
      \strng{fullhash}{b6f197e0531d2bbfab6a59c048e3b28d}
      \strng{bibnamehash}{b6f197e0531d2bbfab6a59c048e3b28d}
      \strng{authorbibnamehash}{b6f197e0531d2bbfab6a59c048e3b28d}
      \strng{authornamehash}{b6f197e0531d2bbfab6a59c048e3b28d}
      \strng{authorfullhash}{b6f197e0531d2bbfab6a59c048e3b28d}
      \field{extraname}{6}
      \field{sortinit}{1}
      \field{sortinithash}{50c6687d7fc80f50136d75228e3c59ba}
      \field{labelnamesource}{author}
      \field{labeltitlesource}{title}
      \field{annotation}{Under the assumption of a ``finite Riemann--Piltz conjecture'' GRH$(H,E_k)$, the authors establish the following result in the case that $\ell_1$ and $\ell_2$ are either both quadratic residues or both quadratic nonresidues$\mod k$: when $T$ is sufficiently large in terms of $k$ (the authors give an explicit lower bound), the inequalities \begin{align*} \max_{T^{1/3}\leq x \leq T} \pi(x,k,\ell_1,\ell_2) &> \sqrt{T}\exp \bigg( {-44\frac{\log{T}\log\log\log{T}}{\log\log{T}}} \bigg) \\ \min_{T^{1/3}\leq x \leq T} \pi(x,k,\ell_1,\ell_2) &< -\sqrt{T}\exp \bigg( {-44\frac{\log{T}\log\log\log{T}}{\log\log{T}}} \bigg) \end{align*} hold. As usual, this result implies the lower bound $W_{k;\ell_1,\ell_2}(T) > \frac1{\log 3}\log\log T + O(1)$. The authors rely on multiple lemmas from their use of power-sum method in previous articles of this series. \par This article cites~\cites{MR0146157, MR0146158, MR0146160}.}
      \field{issn}{0001-5954}
      \field{journaltitle}{Acta Math. Acad. Sci. Hungar.}
      \field{title}{Comparative prime-number theory. {VI}. {C}ontinuation of the general case}
      \field{volume}{14}
      \field{year}{1963}
      \true{nocite}
      \field{pages}{65\bibrangedash 78}
      \range{pages}{14}
      \verb{doi}
      \verb 10.1007/BF01901930
      \endverb
      \verb{mrnumber}
      \verb 0146161
      \endverb
      \verb{urlraw}
      \verb https://doi.org/10.1007/BF01901930
      \endverb
      \verb{url}
      \verb https://doi.org/10.1007/BF01901930
      \endverb
    \endentry
    \entry{MR0156826}{article}{}
      \name{author}{2}{}{%
        {{hash=d33db82fc1e83dc064a23d93e5a21891}{%
           family={Knapowski},
           familyi={K\bibinitperiod},
           given={S.},
           giveni={S\bibinitperiod}}}%
        {{hash=e68e101772b88cb1d48e688164bece92}{%
           family={Turán},
           familyi={T\bibinitperiod},
           given={P.},
           giveni={P\bibinitperiod}}}%
      }
      \strng{namehash}{b6f197e0531d2bbfab6a59c048e3b28d}
      \strng{fullhash}{b6f197e0531d2bbfab6a59c048e3b28d}
      \strng{bibnamehash}{b6f197e0531d2bbfab6a59c048e3b28d}
      \strng{authorbibnamehash}{b6f197e0531d2bbfab6a59c048e3b28d}
      \strng{authornamehash}{b6f197e0531d2bbfab6a59c048e3b28d}
      \strng{authorfullhash}{b6f197e0531d2bbfab6a59c048e3b28d}
      \field{extraname}{7}
      \field{sortinit}{1}
      \field{sortinithash}{50c6687d7fc80f50136d75228e3c59ba}
      \field{labelnamesource}{author}
      \field{labeltitlesource}{title}
      \field{annotation}{This article gives a general conditional proof that $\psi(x; k, \ell_1,\ell_2)$ changes sign infinitely often. The authors show, assuming HC$(E_k)$ for the modulus $k$ for some constant $0<A_k\le1$, that there exists a positive constant $c$ such that $\psi(x; k, \ell_1,\ell_2)$ changes sign in every interval of the form $\omega \leq x \leq \exp(2\sqrt{\omega})$ as long as $$ \omega \geq \max \big\{ e^{k^c},e^{2/{E(k)^3}} \big\}. $$ This result immediately implies results for the first sign change of $\psi(x; k, \ell_1,\ell_2)$ and for its number of sign changes.\par This article cites~\cites{MR0146156, MR0146157, MR0146158, MR0146160, MR0146161}.}
      \field{issn}{0001-5954}
      \field{journaltitle}{Acta Math. Acad. Sci. Hungar}
      \field{title}{Comparative prime-number theory. {VII}. {T}he problem of sign-changes in the general case}
      \field{volume}{14}
      \field{year}{1963}
      \true{nocite}
      \field{pages}{241\bibrangedash 250}
      \range{pages}{10}
      \verb{doi}
      \verb 10.1007/BF01895712
      \endverb
      \verb{mrnumber}
      \verb 0156826
      \endverb
      \verb{urlraw}
      \verb https://doi.org/10.1007/BF01895712
      \endverb
      \verb{url}
      \verb https://doi.org/10.1007/BF01895712
      \endverb
    \endentry
    \entry{MR0156827}{article}{}
      \name{author}{2}{}{%
        {{hash=d33db82fc1e83dc064a23d93e5a21891}{%
           family={Knapowski},
           familyi={K\bibinitperiod},
           given={S.},
           giveni={S\bibinitperiod}}}%
        {{hash=e68e101772b88cb1d48e688164bece92}{%
           family={Turán},
           familyi={T\bibinitperiod},
           given={P.},
           giveni={P\bibinitperiod}}}%
      }
      \strng{namehash}{b6f197e0531d2bbfab6a59c048e3b28d}
      \strng{fullhash}{b6f197e0531d2bbfab6a59c048e3b28d}
      \strng{bibnamehash}{b6f197e0531d2bbfab6a59c048e3b28d}
      \strng{authorbibnamehash}{b6f197e0531d2bbfab6a59c048e3b28d}
      \strng{authornamehash}{b6f197e0531d2bbfab6a59c048e3b28d}
      \strng{authorfullhash}{b6f197e0531d2bbfab6a59c048e3b28d}
      \field{extraname}{8}
      \field{sortinit}{1}
      \field{sortinithash}{50c6687d7fc80f50136d75228e3c59ba}
      \field{labelnamesource}{author}
      \field{labeltitlesource}{title}
      \field{annotation}{Hardy--Littlewood and Landau had already shown that the assertion $\lim_{x\to\infty} \pi_e(x;4,1,3) = -\infty$ is equivalent to GRH for $L(s,\chi_{-4})$. In this article the authors obtain an analogous equivalence concerning the races between $1$ and a nonsquare\mod 8: slightly modifying the arguments for the\mod4 case, they show that the assertion $\lim_{x\to\infty} \theta_e(x;8,1,\ell) = -\infty$ for all $\ell\not\equiv1\mod8$ is equivalent to GRH for the three nonprincipal Dirichlet $L$-functions\mod8, and the same for the assertion $\lim_{x\to\infty} \pi_e(x;8,1,\ell) = -\infty$.\par They further show that the race between two nonsquares switches infinitely often---more precisely, for $\ell_1\not\equiv \ell_2\not\equiv 1\mod8$, they unconditionally show that when $T$ is large enough, \[ \max_{T^{1/3}\le x\le T}\theta_e(x;8,\ell_1,\ell_2)>\sqrt T \exp\left(-22\frac{\log T\log\log\log T}{\log\log T}\right). \] They indicate that this result is ``deeper'', and in particular that they cannot yet replace $\theta_e$ with $\pi_e$ in this result.\par The proofs rely on the power-sum method, as well as some explicit numerical data for the low-lying zeros of the $L$-functions\mod 8. \par This article cites~\cites{MR1555148, MR1544293A,MR1544293B}.}
      \field{issn}{0001-5954}
      \field{journaltitle}{Acta Math. Acad. Sci. Hungar}
      \field{title}{Comparative prime-number theory. {VIII}. {C}hebyshev's problem for {$k=8$}}
      \field{volume}{14}
      \field{year}{1963}
      \true{nocite}
      \field{pages}{251\bibrangedash 268}
      \range{pages}{18}
      \verb{doi}
      \verb 10.1007/BF01895713
      \endverb
      \verb{mrnumber}
      \verb 0156827
      \endverb
      \verb{urlraw}
      \verb https://doi.org/10.1007/BF01895713
      \endverb
      \verb{url}
      \verb https://doi.org/10.1007/BF01895713
      \endverb
    \endentry
    \entry{MR155787}{article}{}
      \name{author}{1}{}{%
        {{hash=36809eb1c6a9cf47ae1c9468614c32b9}{%
           family={Neubauer},
           familyi={N\bibinitperiod},
           given={Gerhard},
           giveni={G\bibinitperiod}}}%
      }
      \list{language}{1}{%
        {German}%
      }
      \strng{namehash}{36809eb1c6a9cf47ae1c9468614c32b9}
      \strng{fullhash}{36809eb1c6a9cf47ae1c9468614c32b9}
      \strng{bibnamehash}{36809eb1c6a9cf47ae1c9468614c32b9}
      \strng{authorbibnamehash}{36809eb1c6a9cf47ae1c9468614c32b9}
      \strng{authornamehash}{36809eb1c6a9cf47ae1c9468614c32b9}
      \strng{authorfullhash}{36809eb1c6a9cf47ae1c9468614c32b9}
      \field{sortinit}{1}
      \field{sortinithash}{50c6687d7fc80f50136d75228e3c59ba}
      \field{labelnamesource}{author}
      \field{labeltitlesource}{title}
      \field{annotation}{The author conducts an empirical study on the function $M(n)$, disproving two conjectures. The first conjecture, by von Sterneck~\cite{1913.vonSterneck} asserts that $|M(n)| \le \frac12\sqrt n$ for $n>200$, but the author finds that $n=7,760,000,000$ is a counterexample. The second conjecture, by Miller, asserts that $\sum_{n \leq x} M(n)<0$ for $x \geq 3$, but the author shows that it first fails at $x=21,067$. The author also points out some errors in the values of $M(n)$ listed in~\cite{1901.vonSterneck,1912.vonSterneck,1913.vonSterneck}. \par This article cites~\cite{1897.vonSterneck,1901.vonSterneck,1912.vonSterneck,1913.vonSterneck}.}
      \field{issn}{0029-599X}
      \field{journaltitle}{Numer. Math.}
      \field{title}{Eine empirische {U}ntersuchung zur {M}ertensschen {F}unktion}
      \field{volume}{5}
      \field{year}{1963}
      \true{nocite}
      \field{pages}{1\bibrangedash 13}
      \range{pages}{13}
      \verb{doi}
      \verb 10.1007/BF01385874
      \endverb
      \verb{mrnumber}
      \verb 155787
      \endverb
      \verb{urlraw}
      \verb https://doi-org.ezproxy.library.ubc.ca/10.1007/BF01385874
      \endverb
      \verb{url}
      \verb https://doi-org.ezproxy.library.ubc.ca/10.1007/BF01385874
      \endverb
    \endentry
    \entry{MR0184920}{book}{}
      \name{author}{1}{}{%
        {{hash=a9ad153a68cf8125996aecb67d0f1f8c}{%
           family={Ingham},
           familyi={I\bibinitperiod},
           given={A.\bibnamedelimi E.},
           giveni={A\bibinitperiod\bibinitdelim E\bibinitperiod}}}%
      }
      \list{publisher}{1}{%
        {Stechert-Hafner, Inc., New York}%
      }
      \strng{namehash}{a9ad153a68cf8125996aecb67d0f1f8c}
      \strng{fullhash}{a9ad153a68cf8125996aecb67d0f1f8c}
      \strng{bibnamehash}{a9ad153a68cf8125996aecb67d0f1f8c}
      \strng{authorbibnamehash}{a9ad153a68cf8125996aecb67d0f1f8c}
      \strng{authornamehash}{a9ad153a68cf8125996aecb67d0f1f8c}
      \strng{authorfullhash}{a9ad153a68cf8125996aecb67d0f1f8c}
      \field{extraname}{4}
      \field{sortinit}{1}
      \field{sortinithash}{50c6687d7fc80f50136d75228e3c59ba}
      \field{labelnamesource}{author}
      \field{labeltitlesource}{title}
      \field{annotation}{This is a reprint of \cite{1932.Ingham}.}
      \field{series}{Cambridge Tracts in Mathematics and Mathematical Physics, No. 30}
      \field{title}{The distribution of prime numbers}
      \field{year}{1964}
      \true{nocite}
      \field{pages}{v+114}
      \range{pages}{-1}
      \verb{mrnumber}
      \verb 0184920
      \endverb
    \endentry
    \entry{MR0176967}{article}{}
      \name{author}{1}{}{%
        {{hash=cc661844b1342c55df563c91418ea523}{%
           family={Kátai},
           familyi={K\bibinitperiod},
           given={I.},
           giveni={I\bibinitperiod}}}%
      }
      \list{language}{1}{%
        {German}%
      }
      \strng{namehash}{cc661844b1342c55df563c91418ea523}
      \strng{fullhash}{cc661844b1342c55df563c91418ea523}
      \strng{bibnamehash}{cc661844b1342c55df563c91418ea523}
      \strng{authorbibnamehash}{cc661844b1342c55df563c91418ea523}
      \strng{authornamehash}{cc661844b1342c55df563c91418ea523}
      \strng{authorfullhash}{cc661844b1342c55df563c91418ea523}
      \field{extraname}{1}
      \field{sortinit}{1}
      \field{sortinithash}{50c6687d7fc80f50136d75228e3c59ba}
      \field{labelnamesource}{author}
      \field{labeltitlesource}{title}
      \field{annotation}{Let $\ell_1$ and $\ell_2$ be distinct reduced residues\mod k. Assuming HC, the author proves that \[ \limsup_{x \to \infty} \frac{\psi(x,k,\ell_1,\ell_2)}{\sqrt{x}}>0\] (and hence the corresponding statement for $\liminf$). When $\ell_1$ and $\ell_2$ are either both quadratic residues or both quadratic nonresidues, it follows that \[ \limsup_{x \to \infty} \frac{\pi(x,k,\ell_1,\ell_2)}{{\sqrt{x}}/{\log x}} >0,\] (and the corresponding statement for $\liminf$). The proof uses an idea of Littlewood, namely to estimate the iterated integrals $\fA_n(x) = \int_2^x \fA_{n-1}(u) \ du$ where $\fA_0(x) = \psi(x;k,\ell_1,\ell_2) +O(\log x)$ is the explicit sum over zeros of Dirichlet $L$-functions\mod k.\par Assuming GRH, the author can make the above statements quantitative and localized to intervals of the form $(x_0,ax_0)$, thus obtaining the lower bounds $W^\psi_{k;\ell_1,\ell_2}(T) \gg \log x$ and (under the same assumption on $\ell_1$ and $\ell_2$) the same estimate for $W^\psi_{k;\ell_1,\ell_2}(T) \gg \log x$.\par This article cites~\cites{MR1511360, MR0146156,MR0146157,MR0146158,MR0146159,MR0146160,MR0146161,MR0156826,MR0156827, MR1574052}.}
      \field{issn}{0524-9007}
      \field{journaltitle}{Ann. Univ. Sci. Budapest. Eötvös Sect. Math.}
      \field{title}{Eine {B}emerkung zur ``{C}omparative prime-number theory {I}-{VIII}'' von {S}. {K}napowski und {P}. {T}urán}
      \field{volume}{7}
      \field{year}{1964}
      \true{nocite}
      \field{pages}{33\bibrangedash 40}
      \range{pages}{8}
      \verb{mrnumber}
      \verb 0176967
      \endverb
    \endentry
    \entry{MR0172856}{article}{}
      \name{author}{1}{}{%
        {{hash=d33db82fc1e83dc064a23d93e5a21891}{%
           family={Knapowski},
           familyi={K\bibinitperiod},
           given={S.},
           giveni={S\bibinitperiod}}}%
      }
      \strng{namehash}{d33db82fc1e83dc064a23d93e5a21891}
      \strng{fullhash}{d33db82fc1e83dc064a23d93e5a21891}
      \strng{bibnamehash}{d33db82fc1e83dc064a23d93e5a21891}
      \strng{authorbibnamehash}{d33db82fc1e83dc064a23d93e5a21891}
      \strng{authornamehash}{d33db82fc1e83dc064a23d93e5a21891}
      \strng{authorfullhash}{d33db82fc1e83dc064a23d93e5a21891}
      \field{extraname}{12}
      \field{sortinit}{1}
      \field{sortinithash}{50c6687d7fc80f50136d75228e3c59ba}
      \field{labelnamesource}{author}
      \field{labeltitlesource}{title}
      \field{annotation}{The author investigates oscillations of the Mertens sum $M(x)$ over intervals. Assuming RH, the author shows that for all~$T>1$, \[ \max_{I\subset[Te^{-6(\log T)^{5/6}},Te^{6(\log T)^{5/6}}]} M(I) > T^{1/2}e^{-(\log T)^{3/4}} \] and the symmetric result for the minimum; indeed, only the assumption RH$(H)$ is needed to establish these oscillations for $T\le e^{H^6}$. The author also establishes oscillations of the same size for the absolute logarithmic average $A^M_{|1|}\big( [Te^{-6(\log T)^{5/6}},Te^{6(\log T)^{5/6}}] \big)$. \par This article cites~\cites{MR0155802,MR0146158,MR0174538}.}
      \field{issn}{0065-1036}
      \field{journaltitle}{Acta Arith.}
      \field{title}{On oscillations of certain means formed from the {M}öbius series. {II}}
      \field{volume}{10}
      \field{year}{1964}
      \true{nocite}
      \field{pages}{377\bibrangedash 386}
      \range{pages}{10}
      \verb{doi}
      \verb 10.4064/aa-10-4-377-386
      \endverb
      \verb{mrnumber}
      \verb 0172856
      \endverb
      \verb{urlraw}
      \verb https://doi.org/10.4064/aa-10-4-377-386
      \endverb
      \verb{url}
      \verb https://doi.org/10.4064/aa-10-4-377-386
      \endverb
    \endentry
    \entry{MR0162771}{article}{}
      \name{author}{2}{}{%
        {{hash=d33db82fc1e83dc064a23d93e5a21891}{%
           family={Knapowski},
           familyi={K\bibinitperiod},
           given={S.},
           giveni={S\bibinitperiod}}}%
        {{hash=e68e101772b88cb1d48e688164bece92}{%
           family={Turán},
           familyi={T\bibinitperiod},
           given={P.},
           giveni={P\bibinitperiod}}}%
      }
      \strng{namehash}{b6f197e0531d2bbfab6a59c048e3b28d}
      \strng{fullhash}{b6f197e0531d2bbfab6a59c048e3b28d}
      \strng{bibnamehash}{b6f197e0531d2bbfab6a59c048e3b28d}
      \strng{authorbibnamehash}{b6f197e0531d2bbfab6a59c048e3b28d}
      \strng{authornamehash}{b6f197e0531d2bbfab6a59c048e3b28d}
      \strng{authorfullhash}{b6f197e0531d2bbfab6a59c048e3b28d}
      \field{extraname}{9}
      \field{sortinit}{1}
      \field{sortinithash}{50c6687d7fc80f50136d75228e3c59ba}
      \field{labelnamesource}{author}
      \field{labeltitlesource}{title}
      \field{annotation}{The first two sections offer a short summary of comparative prime number theory up to 1964. The authors classify the subject into 48 separate problems over 12 categories (and additional variants on these), which is of interest to those interested in the history of the field. They then move on to prove results about ``strongly localized accumulation problems''.\par Most generally, assuming HC for the modulus~$k$, they show that if $T$ is sufficiently large in terms of $k$, then for any $(\ell,k) = 1$ with $\ell\not\equiv1\mod k$, \[ \max_I \psi\big( I; k, \ell, 1 \big) > \sqrt{T}e^{-\log^{11/12} T} \quad\text{and}\quad \min_I \psi\big( I; k, \ell, 1 \big) < -\sqrt{T}e^{-\log^{11/12} T}, \] where the maximum and minimum are taken over all subintervals $I$ of $\big[ Te^{-\log^{11/12} T} , T \big]$. The central argument involves the evaluation of the integral \[ \frac 1{2\pi i} \int e^{As}\bigg(\frac{e^{Bs}-e^{-Bs}}{2Bs}\bigg)^r \bigg( \frac1{\phi(k)} \sum_{\chi\mod k} \overline{\chi}(\ell) \frac{L'}L(s,\chi) \bigg) \, ds \] for positive constants $A$ and~$B$.\par This article cites~\cites{1853.Chebyshev, MR1555148, MR1544293A, MR1544293B, MR0060548, MR0133309, MR0146156, MR0146157, MR0146158, MR0146159, MR0146160, MR0146161, MR0156826, MR0156827}.}
      \field{issn}{0065-1036}
      \field{journaltitle}{Acta Arith.}
      \field{title}{Further developments in the comparative prime-number theory. {I}}
      \field{volume}{9}
      \field{year}{1964}
      \true{nocite}
      \field{pages}{23\bibrangedash 40}
      \range{pages}{18}
      \verb{doi}
      \verb 10.4064/aa-9-1-23-40
      \endverb
      \verb{mrnumber}
      \verb 0162771
      \endverb
      \verb{urlraw}
      \verb https://doi.org/10.4064/aa-9-1-23-40
      \endverb
      \verb{url}
      \verb https://doi.org/10.4064/aa-9-1-23-40
      \endverb
    \endentry
    \entry{MR0174538}{article}{}
      \name{author}{2}{}{%
        {{hash=d33db82fc1e83dc064a23d93e5a21891}{%
           family={Knapowski},
           familyi={K\bibinitperiod},
           given={S.},
           giveni={S\bibinitperiod}}}%
        {{hash=e68e101772b88cb1d48e688164bece92}{%
           family={Turán},
           familyi={T\bibinitperiod},
           given={P.},
           giveni={P\bibinitperiod}}}%
      }
      \strng{namehash}{b6f197e0531d2bbfab6a59c048e3b28d}
      \strng{fullhash}{b6f197e0531d2bbfab6a59c048e3b28d}
      \strng{bibnamehash}{b6f197e0531d2bbfab6a59c048e3b28d}
      \strng{authorbibnamehash}{b6f197e0531d2bbfab6a59c048e3b28d}
      \strng{authornamehash}{b6f197e0531d2bbfab6a59c048e3b28d}
      \strng{authorfullhash}{b6f197e0531d2bbfab6a59c048e3b28d}
      \field{extraname}{10}
      \field{sortinit}{1}
      \field{sortinithash}{50c6687d7fc80f50136d75228e3c59ba}
      \field{labelnamesource}{author}
      \field{labeltitlesource}{title}
      \field{annotation}{The authors establish several theorems, all assuming HC$(E_k)$ with $E_k \ll \sqrt{\log k}/k$; most of their results concern the function $\theta_l(x,r;k,\ell_1,\ell_2)$ where $r=r(x,k)$ satisfies $\frac{\log k}{E_k} \ll r \le \log x$. Their most general result (Theorem VI) is that for any quadratic nonresidue $\ell_1\mod k$ and quadratic residue $\ell_2\mod k$, if $L(s,\chi)$ satisfies GRH for all characters $\chi\mod k$ such that $\chi(\ell_1) \ne \chi(\ell_2)$, then $\theta_l(x,r;k,\ell_1,\ell_2) \gg \sqrt{x}$ for $x$ sufficiently large. \par They also establish the following result. Let $\ell$ be a quadratic nonresidue\mod k, and suppose that there exists a character $\chi\mod k$ with $\chi(\ell)\ne1$ such that $L(s,\chi)$ has a zero $\rho_0=\beta_0+i\gamma_0$ with $\beta_0>\frac12$. Then for all $T$ (sufficiently large in terms of $k$ and $\rho_0$), there exist subintervals $I_\pm$ of $\big[ Te^{-5\log^{20/21}T} , Te^{5\log^{20/21}T} \big]$ such that $\pi(I_+;k,\ell,1) > T^{\beta_0} \exp\big( -(2+\gamma_0^2)(\log T)^{5/7} \big) $ and $ \pi(I_-;k,\ell,1) < - T^{\beta_0} \exp\big( -(2+\gamma_0^2)(\log T)^{5/7} \big)$. They deduce this theorem from an analogous theorem involving $\theta_l(x,r;k,\ell,1)$, which serves as a sort of inverse to Theorem VI. \par Together, these results imply, given a quadratic nonresidue $\ell\mod k$, that $\lim_{x\to\infty} \theta_l(x,r;k,\ell,1) = +\infty$ holds if and only if $L(s,\chi)$ satisfies GRH for all characters $\chi \mod k$ with $\chi(\ell)\ne1$. It follows that $\lim_{x\to\infty} \theta_l(x,r;k,\ell,1) = +\infty$ holds for all quadratic nonresidues\mod k if and only if $L(s,\chi)$ satisfies GRH for every nonprincipal character $\chi \mod k$. \par The authors also make some remarks about races between residue classes to different moduli, showing for example how the race between $\pi(x;3,1)$ and $\pi(x;4,1)$ reduces to a race between residue classes modulo~$12$, to which their results apply. \par This article cites~\cites{1853.Chebyshev, MR1555148, MR1544293A, MR0060548, MR0133309, MR0146157, MR0146159, MR0156827}.}
      \field{issn}{0065-1036}
      \field{journaltitle}{Acta Arith.}
      \field{title}{Further developments in the comparative prime-number theory. {II}. {A} modification of {C}hebyshev's assertion}
      \field{volume}{10}
      \field{year}{1964}
      \true{nocite}
      \field{pages}{293\bibrangedash 313}
      \range{pages}{21}
      \verb{doi}
      \verb 10.4064/aa-10-3-293-313
      \endverb
      \verb{mrnumber}
      \verb 0174538
      \endverb
      \verb{urlraw}
      \verb https://doi.org/10.4064/aa-10-3-293-313
      \endverb
      \verb{url}
      \verb https://doi.org/10.4064/aa-10-3-293-313
      \endverb
    \endentry
    \entry{MR0159791}{article}{}
      \name{author}{1}{}{%
        {{hash=044834fd4db0262d94552d21e42c5f4d}{%
           family={Makai},
           familyi={M\bibinitperiod},
           given={E.},
           giveni={E\bibinitperiod}}}%
      }
      \strng{namehash}{044834fd4db0262d94552d21e42c5f4d}
      \strng{fullhash}{044834fd4db0262d94552d21e42c5f4d}
      \strng{bibnamehash}{044834fd4db0262d94552d21e42c5f4d}
      \strng{authorbibnamehash}{044834fd4db0262d94552d21e42c5f4d}
      \strng{authornamehash}{044834fd4db0262d94552d21e42c5f4d}
      \strng{authorfullhash}{044834fd4db0262d94552d21e42c5f4d}
      \field{sortinit}{1}
      \field{sortinithash}{50c6687d7fc80f50136d75228e3c59ba}
      \field{labelnamesource}{author}
      \field{labeltitlesource}{title}
      \field{annotation}{The author shows that the constant~$B$ in the second main theorem of the power-sum method cannot be taken smaller than~$4e$. \par This article cites~\cites{MR0077579}.}
      \field{issn}{0001-5954}
      \field{journaltitle}{Acta Math. Acad. Sci. Hungar.}
      \field{title}{On a minimum problem. II}
      \field{volume}{15}
      \field{year}{1964}
      \true{nocite}
      \field{pages}{63\bibrangedash 66}
      \range{pages}{4}
      \verb{doi}
      \verb 10.1007/BF01897022
      \endverb
      \verb{mrnumber}
      \verb 0159791
      \endverb
    \endentry
    \entry{MR0177957}{article}{}
      \name{author}{1}{}{%
        {{hash=36befda65f61ea63f097e5518d3ec577}{%
           family={Staś},
           familyi={S\bibinitperiod},
           given={W.},
           giveni={W\bibinitperiod}}}%
      }
      \strng{namehash}{36befda65f61ea63f097e5518d3ec577}
      \strng{fullhash}{36befda65f61ea63f097e5518d3ec577}
      \strng{bibnamehash}{36befda65f61ea63f097e5518d3ec577}
      \strng{authorbibnamehash}{36befda65f61ea63f097e5518d3ec577}
      \strng{authornamehash}{36befda65f61ea63f097e5518d3ec577}
      \strng{authorfullhash}{36befda65f61ea63f097e5518d3ec577}
      \field{extraname}{2}
      \field{sortinit}{1}
      \field{sortinithash}{50c6687d7fc80f50136d75228e3c59ba}
      \field{labelnamesource}{author}
      \field{labeltitlesource}{title}
      \field{annotation}{Assuming RH, the author shows that \[ \max_{T^{1-\ep}\le x\le T} \bigg| \sum_{n=1}^\infty \frac{\mu(n)}n e^{-(x/n)^2} \bigg| \ge T^{-1/2-\ep} \] when~$T$ is sufficiently large in terms of~$\ep$. \par This article cites~\cite{MR0174538}.}
      \field{issn}{0065-1036}
      \field{journaltitle}{Acta Arith.}
      \field{sortyear}{1964}
      \field{title}{Some remarks on a series of Ramanujan}
      \field{volume}{10}
      \field{year}{1964/1965}
      \true{nocite}
      \field{pages}{359\bibrangedash 368}
      \range{pages}{10}
      \verb{doi}
      \verb 10.4064/aa-10-4-359-368
      \endverb
      \verb{mrnumber}
      \verb 0177957
      \endverb
    \endentry
    \entry{MR0229595}{inproceedings}{}
      \name{author}{1}{}{%
        {{hash=e68e101772b88cb1d48e688164bece92}{%
           family={Turán},
           familyi={T\bibinitperiod},
           given={P.},
           giveni={P\bibinitperiod}}}%
      }
      \list{publisher}{1}{%
        {Izdat. ``Nauka'', Leningrad}%
      }
      \strng{namehash}{e68e101772b88cb1d48e688164bece92}
      \strng{fullhash}{e68e101772b88cb1d48e688164bece92}
      \strng{bibnamehash}{e68e101772b88cb1d48e688164bece92}
      \strng{authorbibnamehash}{e68e101772b88cb1d48e688164bece92}
      \strng{authornamehash}{e68e101772b88cb1d48e688164bece92}
      \strng{authorfullhash}{e68e101772b88cb1d48e688164bece92}
      \field{extraname}{7}
      \field{sortinit}{1}
      \field{sortinithash}{50c6687d7fc80f50136d75228e3c59ba}
      \field{labelnamesource}{author}
      \field{labeltitlesource}{title}
      \field{annotation}{The author states some results whose detailed proofs later appeared in the series~\cite{MR0146156} with Knapowski. First, the interval $[T^{1/3},T]$ contains a sign change of $\psi(x;4,3,1)$ when $T$ is large, while $E^\pi(x;4,3,1) = \Omega_\pm(\log\log\log\log\log x)$. Analogous results hold for the moduli $q=3$ and $q=6$. For $q=8$, when $\ell=3,5,7$ the analogues hold for $\pi(x;8,\ell,1)$, while when $\ell_1,\ell_2\neq1$ are distinct then $\pi(x;8,\ell_1,\ell_2)=\Omega_\pm(\sqrt{x}e^{-{\log{x}}/{\sqrt{\log\log{x}}}})$ for values $x$ in all large intervals $[T^{1/3},T]$, and similarly for $\pi_e(x;8,\ell_1,\ell_2)$. When $\ell\in\{3,5,7\}$, he shows that $\lim_{x\to\infty} \pi_e(x;8,\ell,1) = \infty$ if and only GRH$(\chi)$ holds for all $\chi\mod8$ with $\chi(\ell)\neq1$. The analogues hold for $q=5$ (with the understanding that the distinction between $1$ and $3,5,7\mod8$ becomes the distinction between quadratic residues and nonresidues\mod5), with the exception that the $\Omega$-result for $\pi_e(x;5,4,1)$ requires that the complex $L(s,\chi)$ vanishes in the region $\sigma>1-(\log\log t)^{-1/3}$.}
      \field{booktitle}{Proc. {F}ourth {A}ll-{U}nion {M}ath. {C}ongr ({L}eningrad, 1961) ({R}ussian), {V}ol. {II}}
      \field{title}{On a comparative theory of primes}
      \field{year}{1964}
      \true{nocite}
      \field{pages}{137\bibrangedash 142}
      \range{pages}{6}
      \verb{mrnumber}
      \verb 0229595
      \endverb
    \endentry
    \entry{MR0198145}{article}{}
      \name{author}{1}{}{%
        {{hash=f2bf37798cd76913bab45e98f2a4d625}{%
           family={Grosswald},
           familyi={G\bibinitperiod},
           given={Emil},
           giveni={E\bibinitperiod}}}%
      }
      \strng{namehash}{f2bf37798cd76913bab45e98f2a4d625}
      \strng{fullhash}{f2bf37798cd76913bab45e98f2a4d625}
      \strng{bibnamehash}{f2bf37798cd76913bab45e98f2a4d625}
      \strng{authorbibnamehash}{f2bf37798cd76913bab45e98f2a4d625}
      \strng{authornamehash}{f2bf37798cd76913bab45e98f2a4d625}
      \strng{authorfullhash}{f2bf37798cd76913bab45e98f2a4d625}
      \field{extraname}{1}
      \field{sortinit}{1}
      \field{sortinithash}{50c6687d7fc80f50136d75228e3c59ba}
      \field{labelnamesource}{author}
      \field{labeltitlesource}{title}
      \field{annotation}{The author generalizes theorems of Landau~\cite{MR1511360} and Pólya~\cite{1930.Polya} on the relationship between sign changes of a real-valued function and singularities of the corresponding Mellin transform; the main results allow for logarithmic-type singularities as well as poles. The article asserts that its Theorem~6 follows from its Theorem~2, but it is remarked in~\cite{MR858395} that this implication is invalid.\par This article cites~\cites{1930.Polya,1914.Littlewood, MR1511360, MR1511206}.}
      \field{issn}{0021-2172}
      \field{journaltitle}{Israel J. Math.}
      \field{title}{On some generalizations of theorems by Landau and Pólya}
      \field{volume}{3}
      \field{year}{1965}
      \true{nocite}
      \field{pages}{211\bibrangedash 220}
      \range{pages}{10}
      \verb{doi}
      \verb 10.1007/BF03008399
      \endverb
      \verb{mrnumber}
      \verb 0198145
      \endverb
    \endentry
    \entry{MR0231801}{article}{}
      \name{author}{1}{}{%
        {{hash=fa8135c285260ff6fe48a19f384aa417}{%
           family={Kátai},
           familyi={K\bibinitperiod},
           given={Imre},
           giveni={I\bibinitperiod}}}%
      }
      \list{language}{1}{%
        {Hungarian}%
      }
      \strng{namehash}{fa8135c285260ff6fe48a19f384aa417}
      \strng{fullhash}{fa8135c285260ff6fe48a19f384aa417}
      \strng{bibnamehash}{fa8135c285260ff6fe48a19f384aa417}
      \strng{authorbibnamehash}{fa8135c285260ff6fe48a19f384aa417}
      \strng{authornamehash}{fa8135c285260ff6fe48a19f384aa417}
      \strng{authorfullhash}{fa8135c285260ff6fe48a19f384aa417}
      \field{extraname}{2}
      \field{sortinit}{1}
      \field{sortinithash}{50c6687d7fc80f50136d75228e3c59ba}
      \field{labelnamesource}{author}
      \field{labeltitlesource}{title}
      \field{annotation}{Assuming RH, the author proves that $M(x)=\Omega_{\pm}(Tx^{1/2}e^{-c(\log\log{x})^2})$ for some $c>0$. Analogously, assuming RH$(B)$ for any given $B>10^5$, the author proves that $M(x)=\Omega_{\pm}(x^{1/2-\delta}e^{-c(\log\log{x})^2})$, where $\delta\ll(\log\log B)/\log B$.\par This article cites~\cites{MR0046485,MR0096630,MR0133287}.}
      \field{journaltitle}{Magyar Tud. Akad. Mat. Fiz. Oszt. Közl.}
      \field{title}{The {$\Omega$}-estimation of the arithmetic mean of the Möbius function}
      \field{volume}{15}
      \field{year}{1965}
      \true{nocite}
      \field{pages}{15\bibrangedash 18}
      \range{pages}{4}
      \verb{mrnumber}
      \verb 0231801
      \endverb
    \endentry
    \entry{MR0180539}{article}{}
      \name{author}{2}{}{%
        {{hash=d33db82fc1e83dc064a23d93e5a21891}{%
           family={Knapowski},
           familyi={K\bibinitperiod},
           given={S.},
           giveni={S\bibinitperiod}}}%
        {{hash=e68e101772b88cb1d48e688164bece92}{%
           family={Turán},
           familyi={T\bibinitperiod},
           given={P.},
           giveni={P\bibinitperiod}}}%
      }
      \strng{namehash}{b6f197e0531d2bbfab6a59c048e3b28d}
      \strng{fullhash}{b6f197e0531d2bbfab6a59c048e3b28d}
      \strng{bibnamehash}{b6f197e0531d2bbfab6a59c048e3b28d}
      \strng{authorbibnamehash}{b6f197e0531d2bbfab6a59c048e3b28d}
      \strng{authornamehash}{b6f197e0531d2bbfab6a59c048e3b28d}
      \strng{authorfullhash}{b6f197e0531d2bbfab6a59c048e3b28d}
      \field{extraname}{11}
      \field{sortinit}{1}
      \field{sortinithash}{50c6687d7fc80f50136d75228e3c59ba}
      \field{labelnamesource}{author}
      \field{labeltitlesource}{title}
      \field{annotation}{This article concerns the weighted function $\theta_l(x,r;k,\ell,1)$, for a modulus $k$ satisfying HC and a quadratic residue $\ell\not\equiv1\mod k$. Let $\beta_0$ be the real part of any zero of an $L(s,\chi)$ where $\chi \mod k$ a character such that $\chi(\ell) \ne 1$. The authors exhibit extreme values of $\theta_l(x,r;k,\ell,1)$ where $x$ is near $T$ and $r$ is near $(\log T)^{2/3}$; more precisely, there exists a positive constant $c$ such that for $T$ sufficiently large, there exist $x_1,x_2 \in (Te^{-(\log T)^{5/6}}, Te^{(\log T)^{11/15}})$ such that for suitable $r_1,r_2 \in [(2\log T)^{2/3}, (2\log T)^{2/3} + (2\log T)^{2/5}]$, \[ \theta_l(x_1,r_1;k,\ell,1) > T^{\beta_0} e^{-c(\log T)^{5/6}} \quad\text{and}\quad \theta_l(x_2,r_2;k,\ell,1) < -T^{\beta_0} e^{-c(\log T)^{5/6}}. \] The authors then state that using the methods of their prior article~\cite{MR0174538}, it follows that for $T$ sufficiently large, there exist closed subintervals $I,J \subseteq [Te^{-(\log T)^{6/7}}, Te^{(\log T)^{6/7}}]$ such that one has the ``strongly localized accumulations'' \[ \pi(I;k,\ell,1) > \sqrt{T}e^{-c(\log T)^{5/6}} \quad\text{and}\quad \pi(J;k,\ell,1) < -\sqrt{T}e^{-c(\log T)^{5/6}}. \] \par This article cites~\cites{1853.Chebyshev, MR1555148, MR0156827, MR0174538, MR0146157, MR1544293A, MR1544293B, MR0060548}.}
      \field{issn}{0065-1036}
      \field{journaltitle}{Acta Arith.}
      \field{title}{Further developments in the comparative prime-number theory. {III}}
      \field{volume}{11}
      \field{year}{1965}
      \true{nocite}
      \field{pages}{115\bibrangedash 127}
      \range{pages}{13}
      \verb{doi}
      \verb 10.4064/aa-11-1-115-127
      \endverb
      \verb{mrnumber}
      \verb 0180539
      \endverb
      \verb{urlraw}
      \verb https://doi.org/10.4064/aa-11-1-115-127
      \endverb
      \verb{url}
      \verb https://doi.org/10.4064/aa-11-1-115-127
      \endverb
    \endentry
    \entry{MR0182616A}{article}{}
      \name{author}{2}{}{%
        {{hash=d33db82fc1e83dc064a23d93e5a21891}{%
           family={Knapowski},
           familyi={K\bibinitperiod},
           given={S.},
           giveni={S\bibinitperiod}}}%
        {{hash=e68e101772b88cb1d48e688164bece92}{%
           family={Turán},
           familyi={T\bibinitperiod},
           given={P.},
           giveni={P\bibinitperiod}}}%
      }
      \strng{namehash}{b6f197e0531d2bbfab6a59c048e3b28d}
      \strng{fullhash}{b6f197e0531d2bbfab6a59c048e3b28d}
      \strng{bibnamehash}{b6f197e0531d2bbfab6a59c048e3b28d}
      \strng{authorbibnamehash}{b6f197e0531d2bbfab6a59c048e3b28d}
      \strng{authornamehash}{b6f197e0531d2bbfab6a59c048e3b28d}
      \strng{authorfullhash}{b6f197e0531d2bbfab6a59c048e3b28d}
      \field{extraname}{12}
      \field{sortinit}{1}
      \field{sortinithash}{50c6687d7fc80f50136d75228e3c59ba}
      \field{labelnamesource}{author}
      \field{labeltitlesource}{title}
      \field{annotation}{Let $\ell_1$ and $\ell_2$ be quadratic nonresidues modulo a sufficiently large modulus~$k$. Let $\eta$ be sufficiently small in terms of~$k$, and suppose that the Dirichlet $L$-functions\mod k satisfy GRH$(2/\sqrt\eta,E_k)$ for suitable $E_k$. Then, when $T$ is sufficiently large in terms of $k$ and~$\eta$, there exist $x_+$ and $x_-$ in the interval $[T^{1-\sqrt{\eta}} , Te^{\log^{3/4}{T}}]$, and $\eta_1$ and $\eta_2$ in the interval $[2\eta \log{T} , 2\eta \log{T}+\sqrt{\log{T}}]$, such that $\theta_l(x_+,v_+;k,\ell_1,\ell_2) > T^{1/2 - 4\sqrt{\eta}}$ and $\theta_l(x_-,v_-;k,\ell_1,\ell_2) < - T^{1/2 - 4\sqrt{\eta}}$. Furthermore, under the same assumptions, there exist subintervals $I_+$, $I_-$ of $[T^{1-4\sqrt{\eta}},T^{1+4\sqrt{\eta}}]$ such that $\pi(I_+; k, \ell_1, \ell_2) >T^{1/2 - 5\sqrt{\eta}}$ and $\pi(I_-; k, \ell_1, \ell_2) < -T^{1/2 - 5\sqrt{\eta}}$. \par This article cites~\cites{MR0060548, MR0146158, MR0146160, MR0174538}.}
      \field{issn}{0065-1036}
      \field{journaltitle}{Acta Arith. 11 (1965), 147-161; ibid.}
      \field{title}{Further developments in the comparative prime-number theory. {IV}}
      \field{volume}{11}
      \field{year}{1965}
      \true{nocite}
      \field{pages}{147\bibrangedash 161}
      \range{pages}{15}
      \verb{doi}
      \verb 10.4064/aa-11-2-193-202
      \endverb
      \verb{mrnumber}
      \verb 0182616
      \endverb
      \verb{urlraw}
      \verb https://doi.org/10.4064/aa-11-2-193-202
      \endverb
      \verb{url}
      \verb https://doi.org/10.4064/aa-11-2-193-202
      \endverb
    \endentry
    \entry{MR0182616B}{article}{}
      \name{author}{2}{}{%
        {{hash=d33db82fc1e83dc064a23d93e5a21891}{%
           family={Knapowski},
           familyi={K\bibinitperiod},
           given={S.},
           giveni={S\bibinitperiod}}}%
        {{hash=e68e101772b88cb1d48e688164bece92}{%
           family={Turán},
           familyi={T\bibinitperiod},
           given={P.},
           giveni={P\bibinitperiod}}}%
      }
      \strng{namehash}{b6f197e0531d2bbfab6a59c048e3b28d}
      \strng{fullhash}{b6f197e0531d2bbfab6a59c048e3b28d}
      \strng{bibnamehash}{b6f197e0531d2bbfab6a59c048e3b28d}
      \strng{authorbibnamehash}{b6f197e0531d2bbfab6a59c048e3b28d}
      \strng{authornamehash}{b6f197e0531d2bbfab6a59c048e3b28d}
      \strng{authorfullhash}{b6f197e0531d2bbfab6a59c048e3b28d}
      \field{extraname}{13}
      \field{sortinit}{1}
      \field{sortinithash}{50c6687d7fc80f50136d75228e3c59ba}
      \field{labelnamesource}{author}
      \field{labeltitlesource}{title}
      \field{annotation}{This short article is distinct among the second series by Knapowski and Turán, in that, rather than making use of ``one-sided theorems'' of the power-sum method, it uses a different, ``two-sided'' theorem to obtain its results: If $m > 0$ and $z_1,\ldots,z_n \in \C$ nondecreasing in absolute value with $|z_1| = 1$, then for any $b_1,\ldots,b_n \in \C$, there exists an integer $\nu$ such that $m \le \nu \le m+n$ and \[ \bigg| \sum_{j=1}^n b_jz_j \bigg| \ge \frac 1{2n} \left( \frac n{8e(m+n)} \right)^n \min_{1\le k\le n} \bigg| \sum_{j=1}^k b_j \bigg|. \] In addition to the use of the ``two-sided'' theorem above, the authors use a modified idea attributed to Kreisel involving a sequence of integrals. \par The main result of this article is a single theorem, for residues $\ell\not\equiv1\mod k$ for sufficiently large moduli~$k$, under the assumption that there exists $0 < \delta < \frac 1{10}$ such that no function $L(s,\chi)$ with $\chi(\ell) \ne 1$ vanishes in the closed disk $|s-1| \le \frac 12 + 4\delta$. (This assumption is stronger than HC$(2\sqrt\delta)$ but weaker than HC$(\frac 12 + 4\delta)$.) For any sufficiently large $T$, the interval $I = [T,e^{(\log T)^2(\log\log T)^3}]$ contains $x_1,x_2$ such that \[ \psi(x_1;k,1,\ell) \ge x_1^{1/2-4\delta} \quad \text{and} \quad \psi(x_2;k,1,\ell) \le -x_1^{1/2-4\delta}. \] The authors compare their result to~\cite[Theorem 1.1]{MR0146157}, which uses more conventional methods and yields a more localized sign change. \par This article cites~\cites{MR0133308, MR0146157, MR0132728}.}
      \field{issn}{0065-1036}
      \field{journaltitle}{Acta Arith. 11 (1965), 147-161; ibid.}
      \field{title}{Further developments in the comparative prime-number theory. {V}}
      \field{volume}{11}
      \field{year}{1965}
      \true{nocite}
      \field{pages}{193\bibrangedash 202}
      \range{pages}{10}
      \verb{doi}
      \verb 10.4064/aa-11-2-193-202
      \endverb
      \verb{mrnumber}
      \verb 0182616
      \endverb
      \verb{urlraw}
      \verb https://doi.org/10.4064/aa-11-2-193-202
      \endverb
      \verb{url}
      \verb https://doi.org/10.4064/aa-11-2-193-202
      \endverb
    \endentry
    \entry{MR0177963}{article}{}
      \name{author}{2}{}{%
        {{hash=d33db82fc1e83dc064a23d93e5a21891}{%
           family={Knapowski},
           familyi={K\bibinitperiod},
           given={S.},
           giveni={S\bibinitperiod}}}%
        {{hash=e68e101772b88cb1d48e688164bece92}{%
           family={Turán},
           familyi={T\bibinitperiod},
           given={P.},
           giveni={P\bibinitperiod}}}%
      }
      \strng{namehash}{b6f197e0531d2bbfab6a59c048e3b28d}
      \strng{fullhash}{b6f197e0531d2bbfab6a59c048e3b28d}
      \strng{bibnamehash}{b6f197e0531d2bbfab6a59c048e3b28d}
      \strng{authorbibnamehash}{b6f197e0531d2bbfab6a59c048e3b28d}
      \strng{authornamehash}{b6f197e0531d2bbfab6a59c048e3b28d}
      \strng{authorfullhash}{b6f197e0531d2bbfab6a59c048e3b28d}
      \field{extraname}{14}
      \field{sortinit}{1}
      \field{sortinithash}{50c6687d7fc80f50136d75228e3c59ba}
      \field{labelnamesource}{author}
      \field{labeltitlesource}{title}
      \field{annotation}{The authors begin by remarking on some variants of the result of Hardy--Littlewood--Landau~\cite{MR1544293A, MR1544293B, MR1555148} that Chebyshev's assertion, namely that $\lim_{x\to\infty} \pi_e(x;4,1,3) = -\infty$, is equivalent to GRH for $L(s,\chi_{-4})$. The same methods would show the ``Abelian preponderance-relations'' that the limit $\lim_{x\to\infty} \pi_e(x;3,1,2) = -\infty$ if and only if GRH holds for $L(s,\chi_{-3})$, while $\lim_{x\to\infty} \pi_e(x;8,1,\ell) = -\infty$ for all $\ell\in\{3,5,7\}$ if and only if GRH holds for all nonprincipal Dirichlet $L$-functions\mod 8, and (``mutatis mutandis'') $\lim_{x\to\infty} \pi_e(x;12,1,\ell) = -\infty$ for all $\ell\in\{5,7,11\}$ if and only if GRH holds for all nonprincipal Dirichlet $L$-functions\mod{12}. All of these results, they point out, hold with $\pi_e$ replaced by~$\theta_e$. \par For the modulus $k=8$, in the case where $\ell_1,\ell_2\in\{3,5,7\}$ are distinct quadratic nonresidues, the authors had shown ~\cite{MR0156827} that \begin{align*} \max_{T^{1/3} \leq x \leq T}\theta_e(x; 8, \ell_1, \ell_2)> \sqrt{T} \exp\left(-22 \frac{\log T \log\log\log T}{\log \log T}\right); \end{align*} however, they point out that the method failed to yield the analogous result for the ``properly Čebyšev'' function~$\pi_e$. In this article, the authors do establish analogous large oscillations (without identifying the signs of those oscillations) in the form \begin{align*} \max_{T^{1/3} \leq x \leq T} \left|\pi_e(x; 8, \ell_1, \ell_2)\right| \geq \sqrt{T} \exp\left(-23 \frac{\log T \log\log\log T}{\log \log T}\right) \end{align*} for any distinct reduced residues $\ell_1,\ell_2\mod8$, as well as the analogous statement for $\pi_e(x;4,1,3)$. The additional technical tool is a result (then unpublished) of Szegö that derives estimates for $\sum_{j=1}^n b_j e^{-jy} \log j$ from estimates for $\sum_{j=1}^n b_j e^{-jy}$. \par This article cites~\cites{1853.Chebyshev, MR1544293A, MR1544293B, MR1555148, MR0156827}.}
      \field{issn}{0021-7670}
      \field{journaltitle}{J. Analyse Math.}
      \field{title}{On an assertion of Čebyšev}
      \field{volume}{14}
      \field{year}{1965}
      \true{nocite}
      \field{pages}{267\bibrangedash 274}
      \range{pages}{8}
      \verb{doi}
      \verb 10.1007/BF02806393
      \endverb
      \verb{mrnumber}
      \verb 0177963
      \endverb
      \verb{urlraw}
      \verb https://doi.org/10.1007/BF02806393
      \endverb
      \verb{url}
      \verb https://doi.org/10.1007/BF02806393
      \endverb
    \endentry
    \entry{MR0241374}{article}{}
      \name{author}{1}{}{%
        {{hash=fa8135c285260ff6fe48a19f384aa417}{%
           family={Kátai},
           familyi={K\bibinitperiod},
           given={Imre},
           giveni={I\bibinitperiod}}}%
      }
      \list{language}{1}{%
        {Hungarian, with English summary}%
      }
      \strng{namehash}{fa8135c285260ff6fe48a19f384aa417}
      \strng{fullhash}{fa8135c285260ff6fe48a19f384aa417}
      \strng{bibnamehash}{fa8135c285260ff6fe48a19f384aa417}
      \strng{authorbibnamehash}{fa8135c285260ff6fe48a19f384aa417}
      \strng{authornamehash}{fa8135c285260ff6fe48a19f384aa417}
      \strng{authorfullhash}{fa8135c285260ff6fe48a19f384aa417}
      \field{extraname}{3}
      \field{sortinit}{1}
      \field{sortinithash}{50c6687d7fc80f50136d75228e3c59ba}
      \field{labelnamesource}{author}
      \field{labeltitlesource}{title}
      \field{annotation}{The author summarizes many oscillation results from his earlier work, most unconditional and some assuming HC, for functions including $M(x)$, $M_r(x)$, $M_e(x)$, $\psi(x,k,\ell_1,\ell_2)$ and $\pi(x,k,\ell_1,\ell_2)$, $\psi_e(x,k,\ell_1,\ell_2)$, and $\sum_{n=1}^\infty \frac{\mu(n)}n e^{-(x/n)^2}$, as well as functions involving the distribution of $k$-free integers. An example that typifies the strengths of the results is $$ \max_{T<x<T^{7+4\sqrt3}} M(x) = \Omega_\pm(\sqrt x) \quad\text{and}\quad \max_{T<x<T^{1+\ep}} M(x) = \Omega_\pm(x^{\Theta-\ep}). $$\par This article cites~\cites{MR1555148,MR1511360,MR1544293B,1914.Littlewood,MR1574052,MR1575079,MR0146156,MR0146157,MR0146158,MR0146159,MR0146160,MR0146161,MR0156826,MR0156827,MR0096630,MR0142520,MR0155802,MR0172856,MR0087685,MR0177957,MR0046485,MR0067145}.}
      \field{journaltitle}{Magyar Tud. Akad. Mat. Fiz. Oszt. Közl.}
      \field{title}{Omega-type investigations in prime number theory}
      \field{volume}{16}
      \field{year}{1966}
      \true{nocite}
      \field{pages}{369\bibrangedash 396}
      \range{pages}{28}
      \verb{mrnumber}
      \verb 0241374
      \endverb
    \endentry
    \entry{MR0200250}{article}{}
      \name{author}{2}{}{%
        {{hash=d33db82fc1e83dc064a23d93e5a21891}{%
           family={Knapowski},
           familyi={K\bibinitperiod},
           given={S.},
           giveni={S\bibinitperiod}}}%
        {{hash=e68e101772b88cb1d48e688164bece92}{%
           family={Turán},
           familyi={T\bibinitperiod},
           given={P.},
           giveni={P\bibinitperiod}}}%
      }
      \strng{namehash}{b6f197e0531d2bbfab6a59c048e3b28d}
      \strng{fullhash}{b6f197e0531d2bbfab6a59c048e3b28d}
      \strng{bibnamehash}{b6f197e0531d2bbfab6a59c048e3b28d}
      \strng{authorbibnamehash}{b6f197e0531d2bbfab6a59c048e3b28d}
      \strng{authornamehash}{b6f197e0531d2bbfab6a59c048e3b28d}
      \strng{authorfullhash}{b6f197e0531d2bbfab6a59c048e3b28d}
      \field{extraname}{15}
      \field{sortinit}{1}
      \field{sortinithash}{50c6687d7fc80f50136d75228e3c59ba}
      \field{labelnamesource}{author}
      \field{labeltitlesource}{title}
      \field{annotation}{The authors consider a ``modified Abelian means'' race between two quadratic residues $\ell_1,\ell_2\mod k$, under the assumption $GRH(\frac3{\sqrt\eta},E_k)$ for suitable constants $\eta$ and~$E_k$. Their main result is that there exist $x\in [T^{1-\sqrt\eta},T\log T]$ and $\nu \sim 2\eta\log T$ such that \[ \theta_l(x,\nu;k,\ell_1,\ell_2) > T^{1/2-2\sqrt{\eta}} \] (and thus the symmetric result for a large negative value). A corollary is the existence of an interval $I \subset [T^{1-4\sqrt{\eta}},T^{1+4\sqrt{\eta}}]$ such that \[ \pi(I;k,\ell_1,\ell_2) > T^{1/2 - 3 \sqrt{\eta}}. \] The authors obtain similar bounds for two quadratic nonresidues in~\cite{MR0182616A}, but emphasize that they have not been able to extend the results to races where $\ell_1 \not\equiv 1 \mod{k}$ is a quadratic residue and $\ell_2$ is a quadratic nonresidue. They employ the power-sum method for exponential sums.\par This article cites~\cites{MR0146158,MR0174538,MR0180539,MR0182616A}.}
      \field{issn}{0065-1036}
      \field{journaltitle}{Acta Arith.}
      \field{title}{Further developments in the comparative prime-number theory. {VI}. {A}ccumulation theorems for residue-classes representing quadratic residues {${\rm mod}\,k$}}
      \field{volume}{12}
      \field{year}{1966}
      \true{nocite}
      \field{pages}{85\bibrangedash 96}
      \range{pages}{12}
      \verb{doi}
      \verb 10.4064/aa-12-1-85-96
      \endverb
      \verb{mrnumber}
      \verb 0200250
      \endverb
      \verb{urlraw}
      \verb https://doi.org/10.4064/aa-12-1-85-96
      \endverb
      \verb{url}
      \verb https://doi.org/10.4064/aa-12-1-85-96
      \endverb
    \endentry
    \entry{MR0202686}{article}{}
      \name{author}{1}{}{%
        {{hash=d67ac0f50fddf9c6d09a97d86037e761}{%
           family={Lehman},
           familyi={L\bibinitperiod},
           given={R.\bibnamedelimi S.},
           giveni={R\bibinitperiod\bibinitdelim S\bibinitperiod}}}%
      }
      \strng{namehash}{d67ac0f50fddf9c6d09a97d86037e761}
      \strng{fullhash}{d67ac0f50fddf9c6d09a97d86037e761}
      \strng{bibnamehash}{d67ac0f50fddf9c6d09a97d86037e761}
      \strng{authorbibnamehash}{d67ac0f50fddf9c6d09a97d86037e761}
      \strng{authornamehash}{d67ac0f50fddf9c6d09a97d86037e761}
      \strng{authorfullhash}{d67ac0f50fddf9c6d09a97d86037e761}
      \field{extraname}{2}
      \field{sortinit}{1}
      \field{sortinithash}{50c6687d7fc80f50136d75228e3c59ba}
      \field{labelnamesource}{author}
      \field{labeltitlesource}{title}
      \field{annotation}{The author shows that $\pi(x)>\li(x)$ for some $x<1.65\times10^{1,165}$; more precisely, he shows that there is an interval of length $10^{500}$ between $1.53\times10^{1,165}$ and $1.65\times10^{1,165}$ on which $\pi(x)>\li(x)$. The computations, involving the zeros of $\zeta(s)$ up to height $12,000$, were performed on an IBM 7090 at Berkeley. \par This article cites~\cites{MR0067145,MR0133308}.}
      \field{issn}{0065-1036}
      \field{journaltitle}{Acta Arith.}
      \field{title}{On the difference {$\pi (x)-\li (x)$}}
      \field{volume}{11}
      \field{year}{1966}
      \true{nocite}
      \field{pages}{397\bibrangedash 410}
      \range{pages}{14}
      \verb{doi}
      \verb 10.4064/aa-11-4-397-410
      \endverb
      \verb{mrnumber}
      \verb 0202686
      \endverb
    \endentry
    \entry{MR0199161}{article}{}
      \name{author}{1}{}{%
        {{hash=32597e08070c5aa88e8fb07dd295112a}{%
           family={Stark},
           familyi={S\bibinitperiod},
           given={H.\bibnamedelimi M.},
           giveni={H\bibinitperiod\bibinitdelim M\bibinitperiod}}}%
      }
      \strng{namehash}{32597e08070c5aa88e8fb07dd295112a}
      \strng{fullhash}{32597e08070c5aa88e8fb07dd295112a}
      \strng{bibnamehash}{32597e08070c5aa88e8fb07dd295112a}
      \strng{authorbibnamehash}{32597e08070c5aa88e8fb07dd295112a}
      \strng{authornamehash}{32597e08070c5aa88e8fb07dd295112a}
      \strng{authorfullhash}{32597e08070c5aa88e8fb07dd295112a}
      \field{extraname}{1}
      \field{sortinit}{1}
      \field{sortinithash}{50c6687d7fc80f50136d75228e3c59ba}
      \field{labelnamesource}{author}
      \field{labeltitlesource}{title}
      \field{annotation}{Evelyn and Linfoot~\cite{MR1502996} showed that $\Delta^{Q_k}(x)=\Omega_{\pm}(x^{1/2k})$. Following Ingham's method~\cites{MR0006202}, the author makes their result effective by showing that \[ \liminf_{x \to \infty} \Delta^{Q_k}(x) /x^{1/2k} < -C_k \quad \text{and} \quad \limsup_{x \to \infty} \Delta^{Q_k}(x) /x^{1/2k} > C_k; \] here $C_k=2(1-\frac{\gamma_1}{\gamma_2}) \bigl| \zeta(\frac{\rho_1}k)/\rho_1\zeta'(\rho_1) \bigr|$, where $\rho_j=\frac12+i\gamma_j$ ($j=1,2$) are the first two zeros of $\zeta(s)$ above the real axis.\par This article cites~\cites{MR0006202,MR1502996}.}
      \field{issn}{0002-9939,1088-6826}
      \field{journaltitle}{Proc. Amer. Math. Soc.}
      \field{title}{On the asymptotic density of the {$k$}-free integers}
      \field{volume}{17}
      \field{year}{1966}
      \true{nocite}
      \field{pages}{1211\bibrangedash 1214}
      \range{pages}{4}
      \verb{doi}
      \verb 10.2307/2036123
      \endverb
      \verb{mrnumber}
      \verb 199161
      \endverb
      \verb{urlraw}
      \verb https://doi.org/10.2307/2036123
      \endverb
      \verb{url}
      \verb https://doi.org/10.2307/2036123
      \endverb
    \endentry
    \entry{MR0202685}{article}{}
      \name{author}{1}{}{%
        {{hash=f2bf37798cd76913bab45e98f2a4d625}{%
           family={Grosswald},
           familyi={G\bibinitperiod},
           given={Emil},
           giveni={E\bibinitperiod}}}%
      }
      \strng{namehash}{f2bf37798cd76913bab45e98f2a4d625}
      \strng{fullhash}{f2bf37798cd76913bab45e98f2a4d625}
      \strng{bibnamehash}{f2bf37798cd76913bab45e98f2a4d625}
      \strng{authorbibnamehash}{f2bf37798cd76913bab45e98f2a4d625}
      \strng{authornamehash}{f2bf37798cd76913bab45e98f2a4d625}
      \strng{authorfullhash}{f2bf37798cd76913bab45e98f2a4d625}
      \field{extraname}{2}
      \field{sortinit}{1}
      \field{sortinithash}{50c6687d7fc80f50136d75228e3c59ba}
      \field{labelnamesource}{author}
      \field{labeltitlesource}{title}
      \field{annotation}{In this article, the author reproduces known oscillation theorems for $\Delta^\psi(x)$, $\Delta^\Pi(x)$, and $\Delta^\pi(x)$ with proofs based on extended versions of Landau's theorem~\cite{MR1511360}. The results on the number of sign changes of $\psi(q;a,b)$ can be compared with ones in~\cite{MR0146156, MR0146157, MR0146158, MR0146159, MR0146161}. Moreover, under a suitable assumption weaker than HC on the triples $(q,a,b)$ the author proves that the difference $\pi(x;q,a,b)$ changes sign infinitely often for $q=43, 47, 163$, noting that the cases when $q=4$ was proved by Hardy and Littlewood~\cite{MR1555148} and the cases when $q=3, 5, 8$ were proved by Knapowski and Turán~\cite{MR0146159}. It has been noted~\cite[page 2]{MR858395} that some of the implications used in this article to derive some of the theorems from the others may not be valid.\par This article cites~\cites{MR0198145,MR0104638,MR0146156,MR0146157,MR0146158,MR0146159,MR0146160,MR0146161,MR0156826,MR0156827,MR0162771,MR0177963,MR1511206,MR0060548,1936.Ingham,MR1511360,1914.Littlewood,1930.Polya,1853.Chebyshev}.}
      \field{issn}{0002-9947}
      \field{journaltitle}{Trans. Amer. Math. Soc.}
      \field{title}{Oscillation theorems of arithmetical functions}
      \field{volume}{126}
      \field{year}{1967}
      \true{nocite}
      \field{pages}{1\bibrangedash 28}
      \range{pages}{28}
      \verb{doi}
      \verb 10.2307/1994409
      \endverb
      \verb{mrnumber}
      \verb 0202685
      \endverb
    \endentry
    \entry{MR0207665}{article}{}
      \name{author}{1}{}{%
        {{hash=cc661844b1342c55df563c91418ea523}{%
           family={Kátai},
           familyi={K\bibinitperiod},
           given={I.},
           giveni={I\bibinitperiod}}}%
      }
      \list{language}{1}{%
        {Russian}%
      }
      \strng{namehash}{cc661844b1342c55df563c91418ea523}
      \strng{fullhash}{cc661844b1342c55df563c91418ea523}
      \strng{bibnamehash}{cc661844b1342c55df563c91418ea523}
      \strng{authorbibnamehash}{cc661844b1342c55df563c91418ea523}
      \strng{authornamehash}{cc661844b1342c55df563c91418ea523}
      \strng{authorfullhash}{cc661844b1342c55df563c91418ea523}
      \field{extraname}{4}
      \field{sortinit}{1}
      \field{sortinithash}{50c6687d7fc80f50136d75228e3c59ba}
      \field{labelnamesource}{author}
      \field{labeltitlesource}{title}
      \field{annotation}{The author establishes various theorems on bounds for various arithmetic functions using the ideas of Rodosskiĭ. These results include bounds related to $M(x)$ and $\psi(x;q,a,b)$ and $\psi_r(x;q,a)$, as well as estimates related to prime races modulo~$8$.\par This article cites~\cites{MR0087685,MR0146159,MR0146160,MR0146161,MR0156826,MR0156827,MR0177957}.}
      \field{issn}{0001-5954}
      \field{journaltitle}{Acta Math. Acad. Sci. Hungar}
      \field{title}{Comparative theory of prime numbers}
      \field{volume}{18}
      \field{year}{1967}
      \true{nocite}
      \field{pages}{133\bibrangedash 149}
      \range{pages}{17}
      \verb{doi}
      \verb 10.1007/BF02020967
      \endverb
      \verb{mrnumber}
      \verb 0207665
      \endverb
    \endentry
    \entry{MR0218318}{article}{}
      \name{author}{1}{}{%
        {{hash=cc661844b1342c55df563c91418ea523}{%
           family={Kátai},
           familyi={K\bibinitperiod},
           given={I.},
           giveni={I\bibinitperiod}}}%
      }
      \strng{namehash}{cc661844b1342c55df563c91418ea523}
      \strng{fullhash}{cc661844b1342c55df563c91418ea523}
      \strng{bibnamehash}{cc661844b1342c55df563c91418ea523}
      \strng{authorbibnamehash}{cc661844b1342c55df563c91418ea523}
      \strng{authornamehash}{cc661844b1342c55df563c91418ea523}
      \strng{authorfullhash}{cc661844b1342c55df563c91418ea523}
      \field{extraname}{5}
      \field{sortinit}{1}
      \field{sortinithash}{50c6687d7fc80f50136d75228e3c59ba}
      \field{labelnamesource}{author}
      \field{labeltitlesource}{title}
      \field{annotation}{The author establishes an oscillation theorem of Landau type for Dirichlet integrals with a nonreal pole of arbitrary multiplicity, with the additional feature that the oscillations can be localized to explicit intervals of the form $[T,T^K]$. From this, he deduces many unconditional number-theoretical results. For example, all of the following oscillations can be found in all sufficiently large intervals of the form $[T,T^{7+4\sqrt3}]$: \begin{itemize} \item $M(x) = \Omega_\pm(\sqrt x)$, and the same for $M(x,\chi_{-4})$ and $M(x;4,1)$ and $M(x;4,3)$ and $M_e(x)$ \item $M_r(x) = \Omega_\pm(1/\sqrt x)$, and the same for $\sum_{n=1}^\infty \frac{\mu(n)}n e^{-(x/n)^2}$ \item for any $k\ge2$, $\Delta^{Q_k}(x)=\Omega_\pm(x^{1/2k})$, and similarly for the sum of $e^{-n/x}$ over all $k$-free numbers \item $\psi(x;8,\ell_1,\ell_2) = \Omega_\pm(x^{1/2})$ and, if $\ell_1,\ell_2\not\equiv1\mod8$, then $\pi(x;8,\ell_1,\ell_2) = \Omega_\pm(x^{1/2})/\log x$ \end{itemize} Furthermore, by considering separately the cases where RH or GRH is true or false, the author finds all the following oscillations in intervals of the form $[T,T^{1+\ep}]$ for any fixed $\ep>0$: \begin{itemize} \item $M(x) = \Omega_\pm(x^{\Theta-\ep})$, and similarly for the other functions in the previous list \item under HC, $\psi(x;k,\ell_1,\ell_2) = \Omega_\pm(x^{\Theta(k)-\ep})$, and the same for $\psi_e(x;k,\ell_1,\ell_2)$ (indeed, a hypothesis slightly weaker than HC is required, in that real zeros of different $L$-functions could cancel each other out) \end{itemize} These last theorems imply qualitative improvements on the number of sign changes of their respective functions: for example, $W(M,T)/\log\log T\to\infty$ and, under HC, $W^\psi_{k,\ell_1,\ell_2}(T)/\log\log T\to\infty$. \par This article cites~\cites{MR1511360, MR1555148, MR0146156, MR0146157, MR0146158, MR0146159, MR0146160, MR0146161, MR0156826, MR0156827}.}
      \field{issn}{0001-5954}
      \field{journaltitle}{Acta Math. Acad. Sci. Hungar.}
      \field{title}{On investigations in the comparative prime number theory}
      \field{volume}{18}
      \field{year}{1967}
      \true{nocite}
      \field{pages}{379\bibrangedash 391}
      \range{pages}{13}
      \verb{doi}
      \verb 10.1007/BF02280297
      \endverb
      \verb{mrnumber}
      \verb 0218318
      \endverb
      \verb{urlraw}
      \verb https://doi.org/10.1007/BF02280297
      \endverb
      \verb{url}
      \verb https://doi.org/10.1007/BF02280297
      \endverb
    \endentry
    \entry{MR0219496}{article}{}
      \name{author}{1}{}{%
        {{hash=cc661844b1342c55df563c91418ea523}{%
           family={Kátai},
           familyi={K\bibinitperiod},
           given={I.},
           giveni={I\bibinitperiod}}}%
      }
      \strng{namehash}{cc661844b1342c55df563c91418ea523}
      \strng{fullhash}{cc661844b1342c55df563c91418ea523}
      \strng{bibnamehash}{cc661844b1342c55df563c91418ea523}
      \strng{authorbibnamehash}{cc661844b1342c55df563c91418ea523}
      \strng{authornamehash}{cc661844b1342c55df563c91418ea523}
      \strng{authorfullhash}{cc661844b1342c55df563c91418ea523}
      \field{extraname}{6}
      \field{sortinit}{1}
      \field{sortinithash}{50c6687d7fc80f50136d75228e3c59ba}
      \field{labelnamesource}{author}
      \field{labeltitlesource}{title}
      \field{annotation}{Assuming RH$(H)$, the author shows that $\max_{T^\kappa\le x\le T} M(x) \gg \sqrt x$ for~$T$ sufficiently large in terms of~$H$ and for an explicit $\kappa=\kappa(H)$ that increases to $1$ as $H\to\infty$ (and similarly with $M(x)$ replaced by $-M(x)$). All constants are effectively computable; for example, the value $\kappa=0.36$ follows from Rosser and Schoenfeld's calculations of zeros of $\zeta(s)$. The author proves analogous results for several weighted versions of $M(x)$, such as \[ S(x)=\sum_{n=1}^\infty \frac{\mu(n)}{n}e^{-(x/n)^2} = \sum_{k=1}^\infty \frac{(-x^2)^k}{k!\zeta(2k+1)}, \] and also establishes lower bounds of magnitude $T^{\kappa/2}$ for various averages of $M(x)$ and its variants. \par This article cites~\cites{MR0218318,MR0231801,MR0096630,MR0155802,MR0142520,MR1555148,MR0177957},}
      \field{issn}{0065-1036}
      \field{journaltitle}{Acta Arith.}
      \field{sortyear}{1967}
      \field{title}{On oscillations of number-theoretic functions}
      \field{volume}{13}
      \field{year}{1967/1968}
      \true{nocite}
      \field{pages}{107\bibrangedash 122}
      \range{pages}{16}
      \verb{doi}
      \verb 10.4064/aa-13-1-107-122
      \endverb
      \verb{mrnumber}
      \verb 0219496
      \endverb
    \endentry
    \entry{MR0207632}{article}{}
      \name{author}{1}{}{%
        {{hash=56c8a95f3c694f8332b16213f065326f}{%
           family={Ryan},
           familyi={R\bibinitperiod},
           given={J.\bibnamedelimi T.},
           giveni={J\bibinitperiod\bibinitdelim T\bibinitperiod}}}%
      }
      \strng{namehash}{56c8a95f3c694f8332b16213f065326f}
      \strng{fullhash}{56c8a95f3c694f8332b16213f065326f}
      \strng{bibnamehash}{56c8a95f3c694f8332b16213f065326f}
      \strng{authorbibnamehash}{56c8a95f3c694f8332b16213f065326f}
      \strng{authornamehash}{56c8a95f3c694f8332b16213f065326f}
      \strng{authorfullhash}{56c8a95f3c694f8332b16213f065326f}
      \field{sortinit}{1}
      \field{sortinithash}{50c6687d7fc80f50136d75228e3c59ba}
      \field{labelnamesource}{author}
      \field{labeltitlesource}{title}
      \field{annotation}{Let $\pi_a(x;m,b) = \#\{ n\le x\colon n\equiv b\mod m,\, \Omega(n)=a\}$. Via combinatorial reasoning, the author conjectures that $\pi_a(x;m,b)/\pi_a(x;1,1)$ is asymptotically \begin{align*} \frac 1{\phi(m)}-(-1)^{a+1} \bigg( 1-\frac1{c_m} \bigg) x^{-1/2} &\quad\text{if $b$ is a quadratic residue,} \\ \frac 1{\phi(m)}+\frac{(-1)^{a+1}}{c_m}x^{-1/2} &\quad\text{if $b$ is a quadratic nonresidue.} \end{align*} The author provides some numerical evidence supporting the conjecture.\par This article cites~\cite{1891.Phragmen, MR1511360, MR0068565, MR0108470}.}
      \field{issn}{0002-9890}
      \field{journaltitle}{Amer. Math. Monthly}
      \field{number}{1}
      \field{title}{One more ``many-more'' assertion}
      \field{volume}{74}
      \field{year}{1967}
      \true{nocite}
      \field{pages}{19\bibrangedash 24}
      \range{pages}{6}
      \verb{doi}
      \verb 10.2307/2314046
      \endverb
      \verb{mrnumber}
      \verb 0207632
      \endverb
    \endentry
    \entry{MR0233781}{article}{}
      \name{author}{2}{}{%
        {{hash=e0b33296fa8f31deb59123e86ff78d6e}{%
           family={Cohen},
           familyi={C\bibinitperiod},
           given={A.\bibnamedelimi M.},
           giveni={A\bibinitperiod\bibinitdelim M\bibinitperiod}}}%
        {{hash=ab9ca14356d28b56bf08d921be43f3b6}{%
           family={Mayhew},
           familyi={M\bibinitperiod},
           given={M.\bibnamedelimi J.\bibnamedelimi E.},
           giveni={M\bibinitperiod\bibinitdelim J\bibinitperiod\bibinitdelim E\bibinitperiod}}}%
      }
      \strng{namehash}{d09ba402a722f40a101c1ad5fa0eb9d0}
      \strng{fullhash}{d09ba402a722f40a101c1ad5fa0eb9d0}
      \strng{bibnamehash}{d09ba402a722f40a101c1ad5fa0eb9d0}
      \strng{authorbibnamehash}{d09ba402a722f40a101c1ad5fa0eb9d0}
      \strng{authornamehash}{d09ba402a722f40a101c1ad5fa0eb9d0}
      \strng{authorfullhash}{d09ba402a722f40a101c1ad5fa0eb9d0}
      \field{sortinit}{1}
      \field{sortinithash}{50c6687d7fc80f50136d75228e3c59ba}
      \field{labelnamesource}{author}
      \field{labeltitlesource}{title}
      \field{annotation}{Using an unpublished manuscript of Turing as a starting point, and using computations of zeros of $\zeta(s)$ by Haselgrove, the authors show that $\pi(x)-\li(x)>0$ for some $x\le10^{10^{529.7}}$.\par This article cites~\cites{MR0202686}.}
      \field{issn}{0024-6115}
      \field{journaltitle}{Proc. London Math. Soc. (3)}
      \field{title}{On the difference {$\pi (x)-\li (x)$}}
      \field{volume}{18}
      \field{year}{1968}
      \true{nocite}
      \field{pages}{691\bibrangedash 713}
      \range{pages}{23}
      \verb{doi}
      \verb 10.1112/plms/s3-18.4.691
      \endverb
      \verb{mrnumber}
      \verb 0233781
      \endverb
    \endentry
    \entry{MR240062}{article}{}
      \name{author}{2}{}{%
        {{hash=b86a87765e61c9d06430e684291fa9f9}{%
           family={Good},
           familyi={G\bibinitperiod},
           given={I.\bibnamedelimi J.},
           giveni={I\bibinitperiod\bibinitdelim J\bibinitperiod}}}%
        {{hash=6e79ce655dce4f2195319beea947f637}{%
           family={Churchhouse},
           familyi={C\bibinitperiod},
           given={R.\bibnamedelimi F.},
           giveni={R\bibinitperiod\bibinitdelim F\bibinitperiod}}}%
      }
      \strng{namehash}{8abf904253f009c4ff49d49aab43c597}
      \strng{fullhash}{8abf904253f009c4ff49d49aab43c597}
      \strng{bibnamehash}{8abf904253f009c4ff49d49aab43c597}
      \strng{authorbibnamehash}{8abf904253f009c4ff49d49aab43c597}
      \strng{authornamehash}{8abf904253f009c4ff49d49aab43c597}
      \strng{authorfullhash}{8abf904253f009c4ff49d49aab43c597}
      \field{sortinit}{1}
      \field{sortinithash}{50c6687d7fc80f50136d75228e3c59ba}
      \field{labelnamesource}{author}
      \field{labeltitlesource}{title}
      \field{annotation}{The authors describe a random model for partial sums of $\mu(n)$ over short intervals and conjecture that when $h$ is large, the limiting distribution of $M(x)-M(x-h)$ is normal with mean $0$ and variance $6h/\pi^2$. Based on the law of the iterated logarithm, they also conjecture that $\lim \sup_{x \to \infty} {M(x)}/{\sqrt{x \log \log x}}=\sqrt{12}/\pi$.\par This article cites~\cites{MR0104638,1932.Ingham,1914.Littlewood,MR155787,MR0046485}.}
      \field{issn}{0025-5718}
      \field{journaltitle}{Math. Comp.}
      \field{title}{The {R}iemann hypothesis and pseudorandom features of the {M}öbius sequence}
      \field{volume}{22}
      \field{year}{1968}
      \true{nocite}
      \field{pages}{857\bibrangedash 861}
      \range{pages}{5}
      \verb{doi}
      \verb 10.2307/2004584
      \endverb
      \verb{mrnumber}
      \verb 240062
      \endverb
      \verb{urlraw}
      \verb https://doi.org/10.2307/2004584
      \endverb
      \verb{url}
      \verb https://doi.org/10.2307/2004584
      \endverb
    \endentry
    \entry{MR0233782}{article}{}
      \name{author}{1}{}{%
        {{hash=cc661844b1342c55df563c91418ea523}{%
           family={Kátai},
           familyi={K\bibinitperiod},
           given={I.},
           giveni={I\bibinitperiod}}}%
      }
      \strng{namehash}{cc661844b1342c55df563c91418ea523}
      \strng{fullhash}{cc661844b1342c55df563c91418ea523}
      \strng{bibnamehash}{cc661844b1342c55df563c91418ea523}
      \strng{authorbibnamehash}{cc661844b1342c55df563c91418ea523}
      \strng{authornamehash}{cc661844b1342c55df563c91418ea523}
      \strng{authorfullhash}{cc661844b1342c55df563c91418ea523}
      \field{extraname}{7}
      \field{sortinit}{1}
      \field{sortinithash}{50c6687d7fc80f50136d75228e3c59ba}
      \field{labelnamesource}{author}
      \field{labeltitlesource}{title}
      \field{annotation}{For the moduli $q\in\{3,4,5,6,7,8,9,10,11,12,19,24\}$, and on every sufficiently large interval of the form $x\in[T,T^{7+4\sqrt3}]$, the author shows that $\Deltadot^\pi(x;q,a) = \Omega_\pm(\sqrt x/\log x)$ for all quadratic nonresidues~$a$, and also that $\pi(x;q,a,b) = \Omega_\pm(\sqrt x/\log x)$ when $a$ and $b$ are both quadratic residues or both quadratic nonresidues; the same results hold for the exponentially weighted functions $\Deltadot^\pi_e(x;q,a)$ and $\pi_e(x;q,a,b)$.\par This article cites~\cites{1914.Littlewood,MR0067145,MR0133309,MR0202686,MR0156827,MR0177963,MR0087685}.}
      \field{issn}{0001-6969}
      \field{journaltitle}{Acta Sci. Math. (Szeged)}
      \field{title}{On oscillation of the number of primes in an arithmetical progression.}
      \field{volume}{29}
      \field{year}{1968}
      \true{nocite}
      \field{pages}{271\bibrangedash 282}
      \range{pages}{12}
      \verb{mrnumber}
      \verb 0233782
      \endverb
    \endentry
    \entry{MR228456}{article}{}
      \name{author}{1}{}{%
        {{hash=e653ad19c187960606202bfb90c8f746}{%
           family={Spira},
           familyi={S\bibinitperiod},
           given={Robert},
           giveni={R\bibinitperiod}}}%
      }
      \strng{namehash}{e653ad19c187960606202bfb90c8f746}
      \strng{fullhash}{e653ad19c187960606202bfb90c8f746}
      \strng{bibnamehash}{e653ad19c187960606202bfb90c8f746}
      \strng{authorbibnamehash}{e653ad19c187960606202bfb90c8f746}
      \strng{authornamehash}{e653ad19c187960606202bfb90c8f746}
      \strng{authorfullhash}{e653ad19c187960606202bfb90c8f746}
      \field{sortinit}{1}
      \field{sortinithash}{50c6687d7fc80f50136d75228e3c59ba}
      \field{labelnamesource}{author}
      \field{labeltitlesource}{title}
      \field{annotation}{Haselgrove~\cite{MR0104638} resolved {P}ólya's problem in the negative; in this article, the author provides numerical evidence that confirms Haselgrove's computations. The author does a similar computation for the Mertens function, showing that \[ \liminf_{x \to \infty} E^M(x) < -0.6027 \quad \text{and} \quad \limsup_{x \to \infty} E^M(x) > 0.5355. \] The calculations were performed on an IBM 7040 at the University of Tennessee Computing Center. \par This article cites~\cites{MR27305,MR0104638,MR0137689,MR0120198}.}
      \field{issn}{0025-5718}
      \field{journaltitle}{Math. Comp.}
      \field{title}{Zeros of sections of the zeta function. {II}}
      \field{volume}{22}
      \field{year}{1968}
      \true{nocite}
      \field{pages}{163\bibrangedash 173}
      \range{pages}{11}
      \verb{doi}
      \verb 10.2307/2004774
      \endverb
      \verb{mrnumber}
      \verb 228456
      \endverb
      \verb{urlraw}
      \verb https://doi.org/10.2307/2004774
      \endverb
      \verb{url}
      \verb https://doi.org/10.2307/2004774
      \endverb
    \endentry
    \entry{MR0272729}{article}{}
      \name{author}{2}{}{%
        {{hash=d33db82fc1e83dc064a23d93e5a21891}{%
           family={Knapowski},
           familyi={K\bibinitperiod},
           given={S.},
           giveni={S\bibinitperiod}}}%
        {{hash=e68e101772b88cb1d48e688164bece92}{%
           family={Turán},
           familyi={T\bibinitperiod},
           given={P.},
           giveni={P\bibinitperiod}}}%
      }
      \list{language}{1}{%
        {German}%
      }
      \list{publisher}{1}{%
        {Plenum, New York}%
      }
      \strng{namehash}{b6f197e0531d2bbfab6a59c048e3b28d}
      \strng{fullhash}{b6f197e0531d2bbfab6a59c048e3b28d}
      \strng{bibnamehash}{b6f197e0531d2bbfab6a59c048e3b28d}
      \strng{authorbibnamehash}{b6f197e0531d2bbfab6a59c048e3b28d}
      \strng{authornamehash}{b6f197e0531d2bbfab6a59c048e3b28d}
      \strng{authorfullhash}{b6f197e0531d2bbfab6a59c048e3b28d}
      \field{extraname}{16}
      \field{sortinit}{1}
      \field{sortinithash}{50c6687d7fc80f50136d75228e3c59ba}
      \field{labelnamesource}{author}
      \field{labeltitlesource}{title}
      \field{annotation}{The authors prove that for all sufficiently large~$T$, there exist positive reals $x<y\le T$ such that $\theta\bigl((x,y];4,3,1\bigr)>\sqrt{y}$. The authors also prove an analogous bound for $\theta\bigl((x,y];4,1,3\bigr)$.\par This article cites~\cites{1853.Chebyshev,MR1555148,MR1544293A,MR1544293B,1914.Littlewood,MR0162771,MR1511360,MR0174538,MR0132728,MR0146157}}
      \field{booktitle}{Number Theory and Analysis (Papers in Honor of Edmund Landau)}
      \field{title}{Über einige Fragen der vergleichenden Primzahltheorie}
      \field{year}{1969}
      \true{nocite}
      \field{pages}{157\bibrangedash 171}
      \range{pages}{15}
      \verb{mrnumber}
      \verb 0272729
      \endverb
    \endentry
    \entry{MR0263757}{article}{}
      \name{author}{1}{}{%
        {{hash=85d5bb5658a4749b33270d8f0e8be228}{%
           family={Pólya},
           familyi={P\bibinitperiod},
           given={G.},
           giveni={G\bibinitperiod}}}%
      }
      \list{language}{1}{%
        {German}%
      }
      \list{publisher}{1}{%
        {Plenum, New York}%
      }
      \strng{namehash}{85d5bb5658a4749b33270d8f0e8be228}
      \strng{fullhash}{85d5bb5658a4749b33270d8f0e8be228}
      \strng{bibnamehash}{85d5bb5658a4749b33270d8f0e8be228}
      \strng{authorbibnamehash}{85d5bb5658a4749b33270d8f0e8be228}
      \strng{authornamehash}{85d5bb5658a4749b33270d8f0e8be228}
      \strng{authorfullhash}{85d5bb5658a4749b33270d8f0e8be228}
      \field{extraname}{4}
      \field{sortinit}{1}
      \field{sortinithash}{50c6687d7fc80f50136d75228e3c59ba}
      \field{labelnamesource}{author}
      \field{labeltitlesource}{title}
      \field{annotation}{This is a corrected version of the author's earlier article~\cite{1930.Polya}.}
      \field{booktitle}{Number Theory and Analysis (Papers in Honor of Edmund Landau)}
      \field{title}{Über das Vorzeichen des Restgliedes im Primzahlsatz}
      \field{year}{1969}
      \true{nocite}
      \field{pages}{233\bibrangedash 244}
      \range{pages}{12}
      \verb{mrnumber}
      \verb 0263757
      \endverb
    \endentry
    \entry{MR0257003}{article}{}
      \name{author}{1}{}{%
        {{hash=56a0c36192f225fe8031e07a1c09cdd1}{%
           family={Steinig},
           familyi={S\bibinitperiod},
           given={J.},
           giveni={J\bibinitperiod}}}%
      }
      \strng{namehash}{56a0c36192f225fe8031e07a1c09cdd1}
      \strng{fullhash}{56a0c36192f225fe8031e07a1c09cdd1}
      \strng{bibnamehash}{56a0c36192f225fe8031e07a1c09cdd1}
      \strng{authorbibnamehash}{56a0c36192f225fe8031e07a1c09cdd1}
      \strng{authornamehash}{56a0c36192f225fe8031e07a1c09cdd1}
      \strng{authorfullhash}{56a0c36192f225fe8031e07a1c09cdd1}
      \field{sortinit}{1}
      \field{sortinithash}{50c6687d7fc80f50136d75228e3c59ba}
      \field{labelnamesource}{author}
      \field{labeltitlesource}{title}
      \field{annotation}{The author fills a gap in the proof of a refinement of Landau's theorem by Pólya and applies it to summatory functions of arithmetic functions whose associated Dirichlet series have functional equations, obtaining lower bounds for the number of sign changes of the real and imaginary parts of the corresponding error terms. Particular applications include the divisor function $d(n)$, Ramanujan's function $\tau(n)$, and the number of representations of~$n$ as the sum of~$k$ squares (whose associated Dirichlet series is the Epstein zeta-function $\zeta_k(s)$).\par This article cites~\cites{MR1511360}.}
      \field{issn}{0010-2571}
      \field{journaltitle}{Comment. Math. Helv.}
      \field{title}{The changes of sign of certain arithmetical error-terms}
      \field{volume}{44}
      \field{year}{1969}
      \true{nocite}
      \field{pages}{385\bibrangedash 400}
      \range{pages}{16}
      \verb{doi}
      \verb 10.1007/BF02564539
      \endverb
      \verb{mrnumber}
      \verb 0257003
      \endverb
    \endentry
    \entry{MR280447}{article}{}
      \name{author}{1}{}{%
        {{hash=7bb09541026c884b0f336e16312810b6}{%
           family={Saffari},
           familyi={S\bibinitperiod},
           given={Bahman},
           giveni={B\bibinitperiod}}}%
      }
      \list{language}{1}{%
        {French}%
      }
      \strng{namehash}{7bb09541026c884b0f336e16312810b6}
      \strng{fullhash}{7bb09541026c884b0f336e16312810b6}
      \strng{bibnamehash}{7bb09541026c884b0f336e16312810b6}
      \strng{authorbibnamehash}{7bb09541026c884b0f336e16312810b6}
      \strng{authornamehash}{7bb09541026c884b0f336e16312810b6}
      \strng{authorfullhash}{7bb09541026c884b0f336e16312810b6}
      \field{sortinit}{1}
      \field{sortinithash}{50c6687d7fc80f50136d75228e3c59ba}
      \field{labelnamesource}{author}
      \field{labeltitlesource}{title}
      \field{annotation}{The author investigates the connection between Mertens conjecture and a finite version of LI, using a method similar to that of Ingham~\cite{MR0006202}. Let $\gamma_1<\gamma_2<\cdots$ denote the positive ordinates of the nontrivial zeros of $\zeta(s)$. Let $P(h)$ be the statement that there are no nontrivial linear relations $\sum_{k=1}^{h} a_{k} \gamma_{k} \ne 0$ with $\sum_{k=1}^h |a_h| \le h$, a finite and computationally verifiable assertion. The author shows that $P(28000)$ would imply $\lim \sup _{x \rightarrow \infty}|M(x)| / \sqrt{x}>1.179$, which disproves Mertens conjecture. \par This article cites~\cite{MR0330069,1897.Mertens,1897.vonSterneck,1901.vonSterneck,1913.vonSterneck,MR155787,MR0006202,1905.Stieltjes}.}
      \field{issn}{0151-0509}
      \field{journaltitle}{C. R. Acad. Sci. Paris Sér. A-B}
      \field{title}{Sur la fausseté de la conjecture de {M}ertens. ({W}ith discussion.)}
      \field{volume}{271}
      \field{year}{1970}
      \true{nocite}
      \field{pages}{A1097\bibrangedash A1101}
      \range{pages}{-1}
      \verb{mrnumber}
      \verb 280447
      \endverb
    \endentry
    \entry{MR0330069}{article}{}
      \name{author}{5}{}{%
        {{hash=7c2cfee9c4a819ab77cc793d9ce1bdab}{%
           family={Bateman},
           familyi={B\bibinitperiod},
           given={P.\bibnamedelimi T.},
           giveni={P\bibinitperiod\bibinitdelim T\bibinitperiod}}}%
        {{hash=6bdab84d1ccda38b53176b70d8f29951}{%
           family={Brown},
           familyi={B\bibinitperiod},
           given={J.\bibnamedelimi W.},
           giveni={J\bibinitperiod\bibinitdelim W\bibinitperiod}}}%
        {{hash=98ca3e62c2dc209695914011ea7ad03d}{%
           family={Hall},
           familyi={H\bibinitperiod},
           given={R.\bibnamedelimi S.},
           giveni={R\bibinitperiod\bibinitdelim S\bibinitperiod}}}%
        {{hash=82a56dc9f4059d331b665de076b6ccc8}{%
           family={Kloss},
           familyi={K\bibinitperiod},
           given={K.\bibnamedelimi E.},
           giveni={K\bibinitperiod\bibinitdelim E\bibinitperiod}}}%
        {{hash=af6770049aed53b1c6747cf202761a2c}{%
           family={Stemmler},
           familyi={S\bibinitperiod},
           given={Rosemarie\bibnamedelima M.},
           giveni={R\bibinitperiod\bibinitdelim M\bibinitperiod}}}%
      }
      \list{publisher}{1}{%
        {Academic Press, London}%
      }
      \strng{namehash}{818aad5926f1e301d22acb7b91a4fafa}
      \strng{fullhash}{e164d0e7c5218899980cab8c648e58c0}
      \strng{bibnamehash}{e164d0e7c5218899980cab8c648e58c0}
      \strng{authorbibnamehash}{e164d0e7c5218899980cab8c648e58c0}
      \strng{authornamehash}{818aad5926f1e301d22acb7b91a4fafa}
      \strng{authorfullhash}{e164d0e7c5218899980cab8c648e58c0}
      \field{sortinit}{1}
      \field{sortinithash}{50c6687d7fc80f50136d75228e3c59ba}
      \field{labelnamesource}{author}
      \field{labeltitlesource}{title}
      \field{annotation}{This article strengthens a result of Ingham~\cite{MR0006202} in a way that allows for computational exploration. Define a sum $\sum_{n=1}^N c_n\gamma_n$, where $\gamma_1, \gamma_2, \ldots$ are the positive ordinates of the nontrivial zeros of $\zeta(s)$, to be of type~(B) if each $c_n\in\{-2,-1,0,1,2\}$ with at most one $|c_n|$ equal to~$2$. The authors show that if at most finitely many sums of type (B) equal zero, then \begin{align*} \limsup_{x \to \infty} x^{-1/2}L(x)=\infty, &\quad \liminf_{x \to \infty} x^{-1/2}L(x)=-\infty, \\ \limsup_{x \to \infty} x^{-1/2}M(x)=\infty, &\quad \liminf_{x \to \infty} x^{-1/2}M(x)=-\infty. \end{align*} Tables of data are given for the smallest sums of type (B) for $N=1,\ldots,20$, as well as of type~(A) sums where only $c_n\in\{-1,0,1\}$ are allowed.\par This article cites~\cite{1932.Ingham, MR0006202}.}
      \field{booktitle}{Computers in number theory ({P}roc. {S}ci. {R}es. {C}ouncil {A}tlas {S}ympos. {N}o. 2, {O}xford, 1969)}
      \field{title}{Linear relations connecting the imaginary parts of the zeros of the zeta function}
      \field{year}{1971}
      \true{nocite}
      \field{pages}{11\bibrangedash 19}
      \range{pages}{9}
      \verb{mrnumber}
      \verb 0330069
      \endverb
    \endentry
    \entry{MR0300853}{article}{}
      \field{sortinit}{1}
      \field{sortinithash}{50c6687d7fc80f50136d75228e3c59ba}
      \field{labeltitlesource}{title}
      \field{annotation}{This article is a short biography of Stanisław Knapowski's life.}
      \field{issn}{0010-1354}
      \field{journaltitle}{Colloq. Math.}
      \field{title}{Stanisław {K}napowski (19. V. 1931--28. IX. 1967)}
      \field{volume}{23}
      \field{year}{1971}
      \true{nocite}
      \field{pages}{309\bibrangedash 310}
      \range{pages}{2}
      \verb{mrnumber}
      \verb 0300853
      \endverb
    \endentry
    \entry{MR0289452}{article}{}
      \name{author}{1}{}{%
        {{hash=32597e08070c5aa88e8fb07dd295112a}{%
           family={Stark},
           familyi={S\bibinitperiod},
           given={H.\bibnamedelimi M.},
           giveni={H\bibinitperiod\bibinitdelim M\bibinitperiod}}}%
      }
      \strng{namehash}{32597e08070c5aa88e8fb07dd295112a}
      \strng{fullhash}{32597e08070c5aa88e8fb07dd295112a}
      \strng{bibnamehash}{32597e08070c5aa88e8fb07dd295112a}
      \strng{authorbibnamehash}{32597e08070c5aa88e8fb07dd295112a}
      \strng{authornamehash}{32597e08070c5aa88e8fb07dd295112a}
      \strng{authorfullhash}{32597e08070c5aa88e8fb07dd295112a}
      \field{extraname}{2}
      \field{sortinit}{1}
      \field{sortinithash}{50c6687d7fc80f50136d75228e3c59ba}
      \field{labelnamesource}{author}
      \field{labeltitlesource}{title}
      \field{annotation}{This article gives a Tauberian theorem of Landau type (with a proof sketch and a reference to an unpublished manuscript of the author) and uses it to show that \begin{align*} \limsup_{x\to\infty} \bigl( & E^\pi(x;k,a) - E^\pi(x;q,b) \bigr) \ge c_q(b) - c_k(a) \\ &{} + \sum_{\substack{\chi\mod q \\ \chi\ne\chi_0}} \sum_{|\gamma|<T} \frac{\overline\chi(b)}\rho \bigg( 1 - \frac\gamma T \bigg) e^{(\rho-1/2)u} - \sum_{\substack{\chi\mod k \\ \chi\ne\chi_0}} \sum_{|\gamma|<T} \frac{\overline\chi(a)}\rho \bigg( 1 - \frac\gamma T \bigg) e^{(\rho-1/2)u} \end{align*} for any $T>0$ and any $u\in\R$, under the assumption of GRH$(0)$ for nonprincipal $L$-functions modulo~$k$ and~$q$. In particular, the $\limsup$ is positive if either $a$ is a nonsquare\mod k or $b$ is a square\mod q (these are the cases for which $c_q(b) - c_k(a) \ge 0$). Under the same assumption, the author further proves \begin{align*} \limsup_{x\to\infty} \bigl( E^\pi(x;k,a) & {}- E^\pi(x;q,b) \bigr) \ge c_q(b) - c_k(a) \\ &{} + \sum_{\substack{\chi\mod q \\ \chi\ne\chi_0}} \sum_{|\gamma|<T} \frac{\overline\chi(b)}\rho e^{(\rho-1/2)u} - \sum_{\substack{\chi\mod k \\ \chi\ne\chi_0}} \sum_{|\gamma|<T} \frac{\overline\chi(a)}\rho e^{(\rho-1/2)u} \end{align*} for any $T>0$ and any $u\ne0$. Finally, when $u<0$ the author obtains an exact formula for this right-hand side as $T\to\infty$, a version that implies that the $\limsup$ is infinite if $a\equiv1\mod k$ or $b\equiv1\mod q$ but not both. The author also uses this exact formula and some shrewd explicit computation to show for the first time that $\pi(x;5,4,2) = \Omega_+(\sqrt x/\log x)$.\par This article cites~\cites{MR0006202, MR0146156, MR0146157, MR0146158, MR0146159, MR0146160, MR0146161, MR0156826, MR0156827}.}
      \field{issn}{0065-1036}
      \field{journaltitle}{Acta Arith.}
      \field{title}{A problem in comparative prime number theory}
      \field{volume}{18}
      \field{year}{1971}
      \true{nocite}
      \field{pages}{311\bibrangedash 320}
      \range{pages}{10}
      \verb{doi}
      \verb 10.4064/aa-18-1-311-320
      \endverb
      \verb{mrnumber}
      \verb 0289452
      \endverb
      \verb{urlraw}
      \verb https://doi.org/10.4064/aa-18-1-311-320
      \endverb
      \verb{url}
      \verb https://doi.org/10.4064/aa-18-1-311-320
      \endverb
    \endentry
    \entry{MR0300854}{article}{}
      \name{author}{1}{}{%
        {{hash=f4ee93efcace35c40df58fa6a8eb26ba}{%
           family={Turán},
           familyi={T\bibinitperiod},
           given={Paul},
           giveni={P\bibinitperiod}}}%
      }
      \strng{namehash}{f4ee93efcace35c40df58fa6a8eb26ba}
      \strng{fullhash}{f4ee93efcace35c40df58fa6a8eb26ba}
      \strng{bibnamehash}{f4ee93efcace35c40df58fa6a8eb26ba}
      \strng{authorbibnamehash}{f4ee93efcace35c40df58fa6a8eb26ba}
      \strng{authornamehash}{f4ee93efcace35c40df58fa6a8eb26ba}
      \strng{authorfullhash}{f4ee93efcace35c40df58fa6a8eb26ba}
      \field{extraname}{8}
      \field{sortinit}{1}
      \field{sortinithash}{50c6687d7fc80f50136d75228e3c59ba}
      \field{labelnamesource}{author}
      \field{labeltitlesource}{title}
      \field{annotation}{The author recalls theorems from Knapowski's work after he died in a car accident at age~$36$.\par This article cites~\cites{MR0096622, MR0142520, MR0142521, MR0146156, MR0146157, MR0146158, MR0133287}.}
      \field{issn}{0010-1354}
      \field{journaltitle}{Colloq. Math.}
      \field{title}{Commemoration on Stanisław Knapowski}
      \field{volume}{23}
      \field{year}{1971}
      \true{nocite}
      \field{pages}{310\bibrangedash 318}
      \range{pages}{9}
      \verb{doi}
      \verb 10.4064/cm-23-2-309-321
      \endverb
      \verb{mrnumber}
      \verb 0300854
      \endverb
    \endentry
    \entry{MR0332684}{article}{}
      \name{author}{1}{}{%
        {{hash=5b3dabbd7db4c81dcf1e7cffedcb1b4d}{%
           family={Diamond},
           familyi={D\bibinitperiod},
           given={H.\bibnamedelimi G.},
           giveni={H\bibinitperiod\bibinitdelim G\bibinitperiod}}}%
      }
      \list{publisher}{1}{%
        {Springer, Berlin}%
      }
      \strng{namehash}{5b3dabbd7db4c81dcf1e7cffedcb1b4d}
      \strng{fullhash}{5b3dabbd7db4c81dcf1e7cffedcb1b4d}
      \strng{bibnamehash}{5b3dabbd7db4c81dcf1e7cffedcb1b4d}
      \strng{authorbibnamehash}{5b3dabbd7db4c81dcf1e7cffedcb1b4d}
      \strng{authornamehash}{5b3dabbd7db4c81dcf1e7cffedcb1b4d}
      \strng{authorfullhash}{5b3dabbd7db4c81dcf1e7cffedcb1b4d}
      \field{extraname}{1}
      \field{sortinit}{1}
      \field{sortinithash}{50c6687d7fc80f50136d75228e3c59ba}
      \field{labelnamesource}{author}
      \field{labeltitlesource}{title}
      \field{annotation}{The author presents two variants of oscillation theorems analogous to those of Ingham in~\cite{MR0006202}. Let $F(s)=\int_0^\infty e^{-su}f(u) \, du$ denote the Laplace transform of the measurable function $f\colon [0,\infty)\to\R$. We suppose that the integral defining $F(s)$ converges for $\Re(s)>0$, and that $F(s)$ can be continued as a meromorphic function to a neighborhood of the imaginary axis; suppose further that all the poles of $F(s)$ on the imaginary axis are simple. Let $T$ be the set of positive real numbers $t$ such that $it$ is a pole of~$F(s)$, and let $a_t$ be the residue of $F(s)$ at $s=it$; furthermore, let $a_0$ be the residue (possibly $0$) of $F(s)$ at $s=0$.\par The author defines a subset $W\subset T$ to be ``weakly independent of order $N$'' if the only way to find integers $|n_t|\leq N$ $(t\in W)$ such that $\sum_{t\in W} n_t t \in T$ is to choose one $n_t$ equal to~$1$ and the rest equal to~$0$. Given such a weakly independent subset $W\subset T$ of order $N$, the author proves that \begin{align*} \lim_{x\to\infty}\mathop{\rm ess\ sup}_{u\geq x} f(u) &\geq a_0+\frac {2N}{N+1}\sum_{j\in J}|a_j| \\ \lim_{x\to\infty}\mathop{\rm ess\ inf}_{u\geq x} f(u) &\leq a_0-\frac {2N}{N+1}\sum_{j\in J}|a_j| \end{align*} (where these essential supremum and infemum denote the supremum/infemum when we may ignore a set of inputs of measure~$0$); equivalently, if $ \frac {2N}{N+1}\sum_{j\in J}|a_j|>|a_0| $ then $f(x)$ has arbitrarily large sign changes. (The author gives a slight strengthening of this theorem as well.)\par This article cites~\cites{MR0006202, MR0104638, MR0198145, MR0202685, MR0330069,MR0120198,MR280447}.}
      \field{booktitle}{The theory of arithmetic functions ({P}roc. {C}onf., {W}estern {M}ichigan {U}niv., {K}alamazoo, {M}ich., 1971)}
      \field{title}{Two oscillation theorems}
      \field{year}{1972}
      \true{nocite}
      \field{pages}{113\bibrangedash 118. Lecture Notes in Math.\bibrangessep Vol. 251}
      \range{pages}{-1}
      \verb{mrnumber}
      \verb 0332684
      \endverb
    \endentry
    \entry{MR0332685}{incollection}{}
      \name{author}{1}{}{%
        {{hash=f2bf37798cd76913bab45e98f2a4d625}{%
           family={Grosswald},
           familyi={G\bibinitperiod},
           given={Emil},
           giveni={E\bibinitperiod}}}%
      }
      \list{publisher}{1}{%
        {Springer, Berlin}%
      }
      \strng{namehash}{f2bf37798cd76913bab45e98f2a4d625}
      \strng{fullhash}{f2bf37798cd76913bab45e98f2a4d625}
      \strng{bibnamehash}{f2bf37798cd76913bab45e98f2a4d625}
      \strng{authorbibnamehash}{f2bf37798cd76913bab45e98f2a4d625}
      \strng{authornamehash}{f2bf37798cd76913bab45e98f2a4d625}
      \strng{authorfullhash}{f2bf37798cd76913bab45e98f2a4d625}
      \field{extraname}{3}
      \field{sortinit}{1}
      \field{sortinithash}{50c6687d7fc80f50136d75228e3c59ba}
      \field{labelnamesource}{author}
      \field{labeltitlesource}{title}
      \field{annotation}{This article improves results due to Grosswald~\cites{MR0202685} that relied on an extended Landau's theorem and a theorem of Diamond~\cites{MR0332684}. Say that a finite set of real numbers is $k$-independent if there are no nontrivial vanishing linear combinations of that set with integer coefficients bounded by~$k$ in absolute value. The author shows that the verification of the $5$-independence of the (ordinates of the) first $30$ zeros of $\zeta$ is sufficient to show that $\pi(x)-\li(x)$ changes sign infinitely often. Similarly, the verification of the $13$-independence of the first $75$ zeros is sufficient to disprove the Mertens conjecture, and the verification of the $16$-independence of the first $13$ zeros is sufficient to prove that $L(x)$ and $L_r(x)$ change sign infinitely often. Moreover, for integers $q,q'$ with $\phi(q)=\phi(q')$, the author investigates the sign changes of $\pi(x;q,a)-\pi(x;q',a')$. In particular, the author shows that $\pi(x;p,a,a')$ changes sign infinitely often for all primes $p\le19$ and all appropriate $a,a'$.\par This article cites~\cites{MR0330069,MR0332684,MR0202685,MR1555148,MR0104638,MR0006202,MR0146156,MR0289452,1914.Littlewood,MR280447}.}
      \field{bookpagination}{none}
      \field{booktitle}{The theory of arithmetic functions ({P}roc. {C}onf., {W}estern {M}ichigan {U}niv., {K}alamazoo, {M}ich., 1971)}
      \field{note}{Lecture Notes in Math., Vol. 251}
      \field{title}{Oscillation theorems}
      \field{year}{1972}
      \true{nocite}
      \field{pages}{141\bibrangedash 168}
      \range{pages}{28}
      \verb{mrnumber}
      \verb 0332685
      \endverb
    \endentry
    \entry{MR0302585}{article}{}
      \name{author}{2}{}{%
        {{hash=d33db82fc1e83dc064a23d93e5a21891}{%
           family={Knapowski},
           familyi={K\bibinitperiod},
           given={S.},
           giveni={S\bibinitperiod}}}%
        {{hash=e68e101772b88cb1d48e688164bece92}{%
           family={Turán},
           familyi={T\bibinitperiod},
           given={P.},
           giveni={P\bibinitperiod}}}%
      }
      \strng{namehash}{b6f197e0531d2bbfab6a59c048e3b28d}
      \strng{fullhash}{b6f197e0531d2bbfab6a59c048e3b28d}
      \strng{bibnamehash}{b6f197e0531d2bbfab6a59c048e3b28d}
      \strng{authorbibnamehash}{b6f197e0531d2bbfab6a59c048e3b28d}
      \strng{authornamehash}{b6f197e0531d2bbfab6a59c048e3b28d}
      \strng{authorfullhash}{b6f197e0531d2bbfab6a59c048e3b28d}
      \field{extraname}{17}
      \field{sortinit}{1}
      \field{sortinithash}{50c6687d7fc80f50136d75228e3c59ba}
      \field{labelnamesource}{author}
      \field{labeltitlesource}{title}
      \field{annotation}{The authors show that for large enough $T$, there exist numbers $U_1, U_2, U_3, U_4$ with \begin{align*} \log \log \log T &\leq U_2 \exp\left(-\log^{15/16} U_2\right) \leq U_1 < U_2 \leq T\\ \log \log \log T &\leq U_4 \exp\left(-\log^{15/16} U_4\right) \leq U_3 < U_4 \leq T \end{align*} such that $\theta([U_1, U_2]; 4, 1, 3) > \sqrt{U_2}$ and $\theta([U_3, U_4]; 4, 1, 3)< - \sqrt{U_4}$. In particular, there exist consecutive primes $p_n$ and $p_{n+1}$, both congruent to $1\mod4$, satisfying $\log \log \log T \leq p_n < p_{n+1} \leq T$. \par This article cites~\cites{MR0146157, MR0174538, MR0060548}.}
      \field{issn}{0065-1036}
      \field{journaltitle}{Acta Arith.}
      \field{title}{Further developments in the comparative prime number theory. {VII}}
      \field{volume}{21}
      \field{year}{1972}
      \true{nocite}
      \field{pages}{193\bibrangedash 201}
      \range{pages}{9}
      \verb{doi}
      \verb 10.4064/aa-21-1-193-201
      \endverb
      \verb{mrnumber}
      \verb 0302585
      \endverb
      \verb{urlraw}
      \verb https://doi.org/10.4064/aa-21-1-193-201
      \endverb
      \verb{url}
      \verb https://doi.org/10.4064/aa-21-1-193-201
      \endverb
    \endentry
    \entry{MR0302590}{article}{}
      \name{author}{2}{}{%
        {{hash=41686ca8e3d6aab7a821e927fb151ebd}{%
           family={Shanks},
           familyi={S\bibinitperiod},
           given={Daniel},
           giveni={D\bibinitperiod}}}%
        {{hash=f6052b9f2b8df2f5c47462bdead6fa53}{%
           family={Lal},
           familyi={L\bibinitperiod},
           given={Mohan},
           giveni={M\bibinitperiod}}}%
      }
      \strng{namehash}{2b4db4c1ac7b9014f0cffbd9d5021787}
      \strng{fullhash}{2b4db4c1ac7b9014f0cffbd9d5021787}
      \strng{bibnamehash}{2b4db4c1ac7b9014f0cffbd9d5021787}
      \strng{authorbibnamehash}{2b4db4c1ac7b9014f0cffbd9d5021787}
      \strng{authornamehash}{2b4db4c1ac7b9014f0cffbd9d5021787}
      \strng{authorfullhash}{2b4db4c1ac7b9014f0cffbd9d5021787}
      \field{sortinit}{1}
      \field{sortinithash}{50c6687d7fc80f50136d75228e3c59ba}
      \field{labelnamesource}{author}
      \field{labeltitlesource}{title}
      \field{annotation}{Set $\alpha_a(p)=3$ if~$a$ is a cubic residue modulo~$p$ and $\alpha_a(p)=0$ otherwise, and define $k_a = \prod_{p\nmid a} (p-\alpha_a(p))/(p-1)$, a conditionally convergent constant relevant to the conjectured asymptotic formula for the number of prime values of $n^3+a$. The authors describe their computations of~$k_2$ and~$k_3$ accurate to six decimal places. Further heuristics invoke the race between those primes for which~$a$ is a cubic residue and those for which it is not, which is a Chebyshev density theorem race for the extension $\Q(\sqrt[3]a,e^{2\pi i/3})/\Q$.}
      \field{issn}{0025-5718}
      \field{journaltitle}{Math. Comp.}
      \field{title}{Bateman's constants reconsidered and the distribution of cubic residues}
      \field{volume}{26}
      \field{year}{1972}
      \true{nocite}
      \field{pages}{265\bibrangedash 285}
      \range{pages}{21}
      \verb{doi}
      \verb 10.2307/2004737
      \endverb
      \verb{mrnumber}
      \verb 0302590
      \endverb
      \verb{urlraw}
      \verb https://doi.org/10.2307/2004737
      \endverb
      \verb{url}
      \verb https://doi.org/10.2307/2004737
      \endverb
    \endentry
    \entry{MR0352012}{article}{}
      \name{author}{2}{}{%
        {{hash=d4045cfd1e1488ca61ed5c9c26bdd0a6}{%
           family={Dancs},
           familyi={D\bibinitperiod},
           given={S.},
           giveni={S\bibinitperiod}}}%
        {{hash=e68e101772b88cb1d48e688164bece92}{%
           family={Turán},
           familyi={T\bibinitperiod},
           given={P.},
           giveni={P\bibinitperiod}}}%
      }
      \strng{namehash}{e1ec24cc5bfe1ad55de291dec425eaea}
      \strng{fullhash}{e1ec24cc5bfe1ad55de291dec425eaea}
      \strng{bibnamehash}{e1ec24cc5bfe1ad55de291dec425eaea}
      \strng{authorbibnamehash}{e1ec24cc5bfe1ad55de291dec425eaea}
      \strng{authornamehash}{e1ec24cc5bfe1ad55de291dec425eaea}
      \strng{authorfullhash}{e1ec24cc5bfe1ad55de291dec425eaea}
      \field{sortinit}{1}
      \field{sortinithash}{50c6687d7fc80f50136d75228e3c59ba}
      \field{labelnamesource}{author}
      \field{labeltitlesource}{title}
      \field{annotation}{Spurred by unpublished work of Knapowski, the authors establish the following more flexible version of the second main theorem of the power-sum method. In addition to the usual condition $1=|z_1|\geq |z_2| \geq \cdots \geq |z_n|$, let $m \geq 0$ be an integer, choose $\delta_1$ and $\delta_2$ with $\frac{m}{m+n} \leq \delta_2 \leq \delta_1 \leq 1$, and choose indices $\ell_1$ and $\ell_2$ such that $|z_{\ell_1}| \geq \delta_1$ and $|z_{\ell_2}| \leq \delta_2$ (or $\ell_2=n$ if no such $z_{\ell_2}$ exists). Then \[ \max_{m+1 \leq v \leq m+n} \bigg| \sum_{j=1}^n b_j z_j^v \bigg| \geq 2\bigg( \frac{\delta_1-\delta_2}{8e} \bigg)^n \min_{\ell_1 \leq j \leq \ell_2} |b_1+ \cdots +b_j|. \]\par This article cites~\cites{MR0060548}.}
      \field{issn}{0524-9007}
      \field{journaltitle}{Ann. Univ. Sci. Budapest. Eötvös Sect. Math.}
      \field{title}{Investigations in the powersum theory. {I}}
      \field{volume}{16}
      \field{year}{1973}
      \true{nocite}
      \field{pages}{47\bibrangedash 52 (1974)}
      \range{pages}{-1}
      \verb{mrnumber}
      \verb 0352012
      \endverb
    \endentry
    \entry{MR0352026}{article}{}
      \name{author}{1}{}{%
        {{hash=b9201f1e3293ed411cec93a708d9b9b4}{%
           family={Jurkat},
           familyi={J\bibinitperiod},
           given={W.\bibnamedelimi B.},
           giveni={W\bibinitperiod\bibinitdelim B\bibinitperiod}}}%
      }
      \list{publisher}{1}{%
        {Amer. Math. Soc., Providence, R.I.}%
      }
      \strng{namehash}{b9201f1e3293ed411cec93a708d9b9b4}
      \strng{fullhash}{b9201f1e3293ed411cec93a708d9b9b4}
      \strng{bibnamehash}{b9201f1e3293ed411cec93a708d9b9b4}
      \strng{authorbibnamehash}{b9201f1e3293ed411cec93a708d9b9b4}
      \strng{authornamehash}{b9201f1e3293ed411cec93a708d9b9b4}
      \strng{authorfullhash}{b9201f1e3293ed411cec93a708d9b9b4}
      \field{sortinit}{1}
      \field{sortinithash}{50c6687d7fc80f50136d75228e3c59ba}
      \field{labelnamesource}{author}
      \field{labeltitlesource}{title}
      \field{annotation}{Define $C=2+\sum_{n=1}^\infty {(-1)^n(2\pi x)^{2n}}/{n(2n)!\zeta(2n+1)} \approx -.505$. The author proves unconditionally that $\limsup_{x\to\infty} E^M(x) > 1+C$ and $\liminf_{x\to\infty} E^M(x) < C$; in particular, the latter inequality disproves the conjecture of von Sterneck~\cite{1901.vonSterneck} that $|M(x)| \le \frac12\sqrt x$. The main tool is a general oscillation result for almost periodic functions in the distributional sense (combined with Landau's theorem to reduce to the RH case). The author's general theorems recover $\Omega$-results proved by Hardy and Littlewood for $\Delta^\psi(x)$ and $\Delta^D(x)$.\par This article cites~\cites{MR0330069,MR0332684,MR1576550,1932.Ingham,1936.Ingham,MR0006202,1914.Littlewood,MR155787,MR1573970,1901.vonSterneck,MR0046485}.}
      \field{booktitle}{Analytic number theory ({P}roc. {S}ympos. {P}ure {M}ath., {V}ol. {XXIV}, {S}t. {L}ouis {U}niv., {S}t. {L}ouis, {M}o., 1972)}
      \field{title}{On the {M}ertens conjecture and related general $\Omega$-theorems}
      \field{year}{1973}
      \true{nocite}
      \field{pages}{147\bibrangedash 158}
      \range{pages}{12}
      \verb{mrnumber}
      \verb 0352026
      \endverb
    \endentry
    \entry{MR0369287}{article}{}
      \name{author}{1}{}{%
        {{hash=a6d7fefac7d12a29485ec824a06cdbc6}{%
           family={Brent},
           familyi={B\bibinitperiod},
           given={Richard\bibnamedelima P.},
           giveni={R\bibinitperiod\bibinitdelim P\bibinitperiod}}}%
      }
      \strng{namehash}{a6d7fefac7d12a29485ec824a06cdbc6}
      \strng{fullhash}{a6d7fefac7d12a29485ec824a06cdbc6}
      \strng{bibnamehash}{a6d7fefac7d12a29485ec824a06cdbc6}
      \strng{authorbibnamehash}{a6d7fefac7d12a29485ec824a06cdbc6}
      \strng{authornamehash}{a6d7fefac7d12a29485ec824a06cdbc6}
      \strng{authorfullhash}{a6d7fefac7d12a29485ec824a06cdbc6}
      \field{sortinit}{1}
      \field{sortinithash}{50c6687d7fc80f50136d75228e3c59ba}
      \field{labelnamesource}{author}
      \field{labeltitlesource}{title}
      \field{annotation}{The author reports on extensive calculations of $\pi(x)$ and of the counting function for twin primes. He suggests from these computations that $E^\pi(x)$ has a limiting distribution with mean $-1$ and standard deviation approximately $0.21$ (which differs by less than $3\%$ from the true value under RH and LI). He verifies that $\pi(x)\le\li(x)$ for $x\le 8\times10^{10}$, and predicts from the observed distribution above that $\pi(x)\le\li(x)$ for~$x$ up to at least $10^{100}$. Shanks~\cite{MR0108470} conjectured that $\frac1x \int_0^x E^\pi(t)\, dt$ tends to~$-1$; the author suggests from his computations that this conjecture is false but that its logarithmic counterpart $\lim_{x\to\infty} \frac1{\log x} \int_1^x E^\pi(t)\, \frac{dt}t = -1$ should be true.\par The author also shows that the error term in the counting function of twin primes does not exceed $2.3x^{1/2}/\log x$ in absolute value for $x\leq 8\times 10^{10}$. He also gives the heuristic estimate $1.9021604\pm5\times10^{-7}$ for Brun's constant $\sum(\frac1q+\frac1{q+2})$ (where the sum is over twin primes~$q$), although no rigorous upper bound is given.\par This article cites~\cites{MR1555148,1932.Ingham, MR0202686,1914.Littlewood,MR0108470,MR0067145}.}
      \field{issn}{0025-5718}
      \field{journaltitle}{Math. Comp.}
      \field{note}{Collection of articles dedicated to Derrick Henry Lehmer on the occasion of his seventieth birthday}
      \field{title}{Irregularities in the distribution of primes and twin primes}
      \field{volume}{29}
      \field{year}{1975}
      \true{nocite}
      \field{pages}{43\bibrangedash 56}
      \range{pages}{14}
      \verb{doi}
      \verb 10.2307/2005460
      \endverb
      \verb{mrnumber}
      \verb 0369287
      \endverb
    \endentry
    \entry{MR0376566}{article}{}
      \name{author}{1}{}{%
        {{hash=4ceaa3cbf82f92b37edd652b07fd8881}{%
           family={Diamond},
           familyi={D\bibinitperiod},
           given={Harold\bibnamedelima G.},
           giveni={H\bibinitperiod\bibinitdelim G\bibinitperiod}}}%
      }
      \strng{namehash}{4ceaa3cbf82f92b37edd652b07fd8881}
      \strng{fullhash}{4ceaa3cbf82f92b37edd652b07fd8881}
      \strng{bibnamehash}{4ceaa3cbf82f92b37edd652b07fd8881}
      \strng{authorbibnamehash}{4ceaa3cbf82f92b37edd652b07fd8881}
      \strng{authornamehash}{4ceaa3cbf82f92b37edd652b07fd8881}
      \strng{authorfullhash}{4ceaa3cbf82f92b37edd652b07fd8881}
      \field{extraname}{2}
      \field{sortinit}{1}
      \field{sortinithash}{50c6687d7fc80f50136d75228e3c59ba}
      \field{labelnamesource}{author}
      \field{labeltitlesource}{title}
      \field{annotation}{The author gives a new proof that $\pi(x)-\li(x)$ changes sign infinitely often, without using an explicit formula for prime-counting functions but instead establishing a Tauberian theorem of Wiener--Ikehara type. The author believes that the arguments could be extended to gain the extra factor of $\log\log\log x$ in the oscillations found by Littlewood.\par This article cites~\cites{MR0137689, MR0006202}.}
      \field{issn}{0013-8584}
      \field{journaltitle}{Enseignement Math. (2)}
      \field{number}{1}
      \field{title}{Changes of sign of {$\pi (x)-\li (x)$}}
      \field{volume}{21}
      \field{year}{1975}
      \true{nocite}
      \field{pages}{1\bibrangedash 14}
      \range{pages}{14}
      \verb{mrnumber}
      \verb 0376566
      \endverb
    \endentry
    \entry{MR0417077}{book}{}
      \name{author}{1}{}{%
        {{hash=2f35569443e88f879c6356cd12f56c9c}{%
           family={Ellison},
           familyi={E\bibinitperiod},
           given={William\bibnamedelima John},
           giveni={W\bibinitperiod\bibinitdelim J\bibinitperiod}}}%
      }
      \list{language}{1}{%
        {French}%
      }
      \list{publisher}{1}{%
        {Hermann, Paris}%
      }
      \strng{namehash}{2f35569443e88f879c6356cd12f56c9c}
      \strng{fullhash}{2f35569443e88f879c6356cd12f56c9c}
      \strng{bibnamehash}{2f35569443e88f879c6356cd12f56c9c}
      \strng{authorbibnamehash}{2f35569443e88f879c6356cd12f56c9c}
      \strng{authornamehash}{2f35569443e88f879c6356cd12f56c9c}
      \strng{authorfullhash}{2f35569443e88f879c6356cd12f56c9c}
      \field{sortinit}{1}
      \field{sortinithash}{50c6687d7fc80f50136d75228e3c59ba}
      \field{labelnamesource}{author}
      \field{labeltitlesource}{title}
      \field{annotation}{Chapter~6 of the book is on irregularities in the distribution of prime counting functions. Complete proofs of classical oscillation results on $\Delta^{\psi}(x)$ and $\Delta^{\Pi}(x)$ are presented in Section~6.2. The Mertens conjecture and Pólya problem and their connections with Ingham's method~\cite{MR0006202} are discussed in Section~6.3. In Section~6.4, the authors discuss the disproof a conjecture of Shanks~\cite{MR0108470}, which is a precise version of Chebyshev's bias on the prime races between $1 \mod 4$ primes and $3 \mod 4$ primes. In the notes section of the same chapter, the authors discuss the difference between $\Delta^{\pi}(x)$ and its analogue in the setting of Gaussian primes, as well as sign changes of arithmetic functions.\par This chapter cites~\cites{1897.Mertens, MR155787,MR0006202,MR280447,MR0120198, MR0104638, MR0108470, 1853.Chebyshev, MR0083001, MR1555148, MR1511360, 1914.Littlewood, MR1573970, MR0067145, MR0137689, MR0233781, MR0202686, MR0257003,MR0202685,MR0133308,MR0146156,MR0146157,MR0146158,MR0146159,MR0146160,MR0146161}.}
      \field{note}{En collaboration avec Michel Mendès France; Publications de l'Institut de Mathématique de l'Université de Nancago, No. IX; Actualités Scientifiques et Industrielles, No. 1366}
      \field{title}{Les nombres premiers}
      \field{year}{1975}
      \true{nocite}
      \field{pages}{xiv+442}
      \range{pages}{-1}
      \verb{mrnumber}
      \verb 0417077
      \endverb
    \endentry
    \entry{MR429789}{article}{}
      \name{author}{2}{}{%
        {{hash=b603cd4e91f630615d9bcde63b8a7ffa}{%
           family={Jurkat},
           familyi={J\bibinitperiod},
           given={W.},
           giveni={W\bibinitperiod}}}%
        {{hash=c33ce895e089b8d5217e8f2e68f087ec}{%
           family={Peyerimhoff},
           familyi={P\bibinitperiod},
           given={A.},
           giveni={A\bibinitperiod}}}%
      }
      \strng{namehash}{a2de3531375e01cede375ae506c851d2}
      \strng{fullhash}{a2de3531375e01cede375ae506c851d2}
      \strng{bibnamehash}{a2de3531375e01cede375ae506c851d2}
      \strng{authorbibnamehash}{a2de3531375e01cede375ae506c851d2}
      \strng{authornamehash}{a2de3531375e01cede375ae506c851d2}
      \strng{authorfullhash}{a2de3531375e01cede375ae506c851d2}
      \field{sortinit}{1}
      \field{sortinithash}{50c6687d7fc80f50136d75228e3c59ba}
      \field{labelnamesource}{author}
      \field{labeltitlesource}{title}
      \field{annotation}{Similar to the first author's work~\cite{MR0352026}, the authors reduce the problem of improving the bounds on $\limsup E^M(x)$ and $\liminf E^M(x)$ to finding a reasonably good solution to an inhomogenous Diophantine approximation problem. Their new constructive algorithm leads to the following improved bounds related to the Mertens conjecture: \[ \liminf_{x \to \infty} E^M(x) < -0.638 \quad \text{and} \quad \limsup_{x \to \infty} E^M(x) > 0.779. \] The authors remark that the smallest counterexample to the Mertens conjecture is likely to be at least $\exp(4.16 \times 10^{14})$. \par This article cites~\cite{MR0330069,MR0332684,1897.Mertens,1897.vonSterneck,1901.vonSterneck,1912.vonSterneck,1913.vonSterneck,MR0046485,MR155787,MR0352026,MR0006202,MR228456}.}
      \field{issn}{0075-4102}
      \field{journaltitle}{J. Reine Angew. Math.}
      \field{title}{A constructive approach to {K}ronecker approximations and its application to the {M}ertens conjecture}
      \field{volume}{286(287)}
      \field{year}{1976}
      \true{nocite}
      \field{pages}{322\bibrangedash 340}
      \range{pages}{19}
      \verb{doi}
      \verb 10.1515/crll.1976.286-287.322
      \endverb
      \verb{mrnumber}
      \verb 429789
      \endverb
      \verb{urlraw}
      \verb https://doi.org/10.1515/crll.1976.286-287.322
      \endverb
      \verb{url}
      \verb https://doi.org/10.1515/crll.1976.286-287.322
      \endverb
    \endentry
    \entry{MR0439771}{inproceedings}{}
      \name{author}{2}{}{%
        {{hash=d33db82fc1e83dc064a23d93e5a21891}{%
           family={Knapowski},
           familyi={K\bibinitperiod},
           given={S.},
           giveni={S\bibinitperiod}}}%
        {{hash=e68e101772b88cb1d48e688164bece92}{%
           family={Turán},
           familyi={T\bibinitperiod},
           given={P.},
           giveni={P\bibinitperiod}}}%
      }
      \list{publisher}{1}{%
        {North-Holland, Amsterdam}%
      }
      \strng{namehash}{b6f197e0531d2bbfab6a59c048e3b28d}
      \strng{fullhash}{b6f197e0531d2bbfab6a59c048e3b28d}
      \strng{bibnamehash}{b6f197e0531d2bbfab6a59c048e3b28d}
      \strng{authorbibnamehash}{b6f197e0531d2bbfab6a59c048e3b28d}
      \strng{authornamehash}{b6f197e0531d2bbfab6a59c048e3b28d}
      \strng{authorfullhash}{b6f197e0531d2bbfab6a59c048e3b28d}
      \field{extraname}{18}
      \field{sortinit}{1}
      \field{sortinithash}{50c6687d7fc80f50136d75228e3c59ba}
      \field{labelnamesource}{author}
      \field{labeltitlesource}{title}
      \field{annotation}{The authors prove that $W^\pi(T) = \Omega((\log T)^{1/4}(\log\log T)^{-4})$ (a slightly weaker result is claimed). In light of Ingham's stronger result~\cite{1936.Ingham} assuming SA, it suffices to assume that RH is false. Given any zero $\rho=\beta+i\gamma$ of $\zeta(s)$ with $\beta>\frac12$, the authors show that for $x \in [Y,Y\exp((\log Y)^{3/4} (\log\log Y)^4)]$ one has $\Delta^\Pi(x) = \Omega_\pm(Y^\beta \exp(-\sqrt{\log Y}))$ when $Y$ is sufficiently large in terms of $\rho$. The authors make special mention of their application of a ``two-sided'' power-sum theorem to get a strong ``one-sided'' result.\par This article cites~\cite{MR1511206, 1914.Littlewood, MR1555148, MR1573970, MR1575079, MR0067145, MR0202686, MR0137689, 1930.Polya, 1936.Ingham, MR0133309, MR0182616B}.}
      \field{booktitle}{Topics in number theory ({P}roc. {C}olloq., {D}ebrecen, 1974)}
      \field{title}{On the sign changes of {$(\pi (x)-{\rm li}\ x)$}. {I}}
      \field{year}{1976}
      \true{nocite}
      \field{pages}{153\bibrangedash 169. Colloq. Math. Soc. János Bolyai\bibrangessep Vol. 13}
      \range{pages}{-1}
      \verb{mrnumber}
      \verb 0439771
      \endverb
    \endentry
    \entry{MR0439772}{article}{}
      \name{author}{2}{}{%
        {{hash=d33db82fc1e83dc064a23d93e5a21891}{%
           family={Knapowski},
           familyi={K\bibinitperiod},
           given={S.},
           giveni={S\bibinitperiod}}}%
        {{hash=e68e101772b88cb1d48e688164bece92}{%
           family={Turán},
           familyi={T\bibinitperiod},
           given={P.},
           giveni={P\bibinitperiod}}}%
      }
      \strng{namehash}{b6f197e0531d2bbfab6a59c048e3b28d}
      \strng{fullhash}{b6f197e0531d2bbfab6a59c048e3b28d}
      \strng{bibnamehash}{b6f197e0531d2bbfab6a59c048e3b28d}
      \strng{authorbibnamehash}{b6f197e0531d2bbfab6a59c048e3b28d}
      \strng{authornamehash}{b6f197e0531d2bbfab6a59c048e3b28d}
      \strng{authorfullhash}{b6f197e0531d2bbfab6a59c048e3b28d}
      \field{extraname}{19}
      \field{sortinit}{1}
      \field{sortinithash}{50c6687d7fc80f50136d75228e3c59ba}
      \field{labelnamesource}{author}
      \field{labeltitlesource}{title}
      \field{annotation}{Following ideas of Littlewood, Ingham, and Skewes, the authors show unconditionally that $W(Y) \gg \log\log\log Y$ for sufficiently large $Y$, where the implied constants are effective.\par The proof itself is divided into two cases. First, supposing the existence of an RH-violating zero $\beta+i\gamma$ of $\zeta(s)$ such that $\beta \ge \frac 12 + 2\log^{-1/5} Y$ and $0 < \gamma \le \log^{1/5} Y$, the authors establish the much stronger lower bound \[ V_1(Y) > \frac 12 \bigg(\frac{\log Y}{2\log^{5/6} Y} \bigg)^{1/5} > \frac 14 \log^{1/30} Y. \] The second case, where there is no such zero, is more technical and relies as usual upon Dirichlet's box principle.\par This article cites~\cites{1914.Littlewood, MR1555148, MR0067145, 1930.Polya, 1936.Ingham, MR0439771, MR0133308, 1932.Ingham}.}
      \field{journaltitle}{Monatsh. Math.}
      \field{number}{2}
      \field{title}{On the sign changes of {$(\pi (x)-{\rm li}\ x)$}. {II}}
      \field{volume}{82}
      \field{year}{1976}
      \true{nocite}
      \field{pages}{163\bibrangedash 175}
      \range{pages}{13}
      \verb{doi}
      \verb 10.1007/BF01305997
      \endverb
      \verb{mrnumber}
      \verb 0439772
      \endverb
      \verb{urlraw}
      \verb https://doi.org/10.1007/BF01305997
      \endverb
      \verb{url}
      \verb https://doi.org/10.1007/BF01305997
      \endverb
    \endentry
    \entry{MR0439774}{inproceedings}{}
      \name{author}{1}{}{%
        {{hash=66650e0cb2c39da0339e84e2bc736eab}{%
           family={Levinson},
           familyi={L\bibinitperiod},
           given={N.},
           giveni={N\bibinitperiod}}}%
      }
      \list{publisher}{1}{%
        {North-Holland, Amsterdam}%
      }
      \strng{namehash}{66650e0cb2c39da0339e84e2bc736eab}
      \strng{fullhash}{66650e0cb2c39da0339e84e2bc736eab}
      \strng{bibnamehash}{66650e0cb2c39da0339e84e2bc736eab}
      \strng{authorbibnamehash}{66650e0cb2c39da0339e84e2bc736eab}
      \strng{authornamehash}{66650e0cb2c39da0339e84e2bc736eab}
      \strng{authorfullhash}{66650e0cb2c39da0339e84e2bc736eab}
      \field{sortinit}{1}
      \field{sortinithash}{50c6687d7fc80f50136d75228e3c59ba}
      \field{labelnamesource}{author}
      \field{labeltitlesource}{title}
      \field{annotation}{Assuming that SA is false, the author shows that $\limsup_{T \to \infty} W^\pi(T)/\log T = \infty$. The author adapts a Landau-type argument of Pólya~\cite{1930.Polya,MR0263757} to treat functions with logarithmic singularities. \par This article cites~\cite{1936.Ingham,MR0133309,MR0133308,MR0439771,1930.Polya,MR0263757}.}
      \field{booktitle}{Topics in number theory ({P}roc. {C}olloq., {D}ebrecen, 1974)}
      \field{title}{On the number of sign changes of {$\pi (x)-\li x$}}
      \field{year}{1976}
      \true{nocite}
      \field{pages}{171\bibrangedash 177. Colloq. Math. Soc. János Bolyai\bibrangessep Vol. 13}
      \range{pages}{-1}
      \verb{mrnumber}
      \verb 0439774
      \endverb
    \endentry
    \entry{MR0439773}{article}{}
      \name{author}{1}{}{%
        {{hash=c6d965df16e600c701d4bcc2888ae34c}{%
           family={Pintz},
           familyi={P\bibinitperiod},
           given={J.},
           giveni={J\bibinitperiod}}}%
      }
      \list{language}{1}{%
        {German}%
      }
      \strng{namehash}{c6d965df16e600c701d4bcc2888ae34c}
      \strng{fullhash}{c6d965df16e600c701d4bcc2888ae34c}
      \strng{bibnamehash}{c6d965df16e600c701d4bcc2888ae34c}
      \strng{authorbibnamehash}{c6d965df16e600c701d4bcc2888ae34c}
      \strng{authornamehash}{c6d965df16e600c701d4bcc2888ae34c}
      \strng{authorfullhash}{c6d965df16e600c701d4bcc2888ae34c}
      \field{extraname}{1}
      \field{sortinit}{1}
      \field{sortinithash}{50c6687d7fc80f50136d75228e3c59ba}
      \field{labelnamesource}{author}
      \field{labeltitlesource}{title}
      \field{annotation}{The author shows that there exists $c>0$ such that $W(T) \gg (\log\log T)^{c}$ when $T$ is sufficiently large. Indeed, this is a special case of a more general result that establishes many large oscillations of $\Delta(x)$: let $D$ be sufficiently large and set $\mu=D/\log\log\log\log T$. Then there are at least $\exp\big( (\log\log\log T)^{1-\mu} \big)$ sign changes of $\Delta(x)$ up to $T$, with oscillations as large as \[ \Delta(x) > \bigg( \frac12-\frac{3\log D}{D} \bigg)\mu \cdot \frac{\sqrt x\log\log\log x}{\log x} \] (and the negative analogue), which, when $D$ is so large that $\mu\gg1$, provides oscillations as large as those established by Littlewood.\par This article cites~\cites{1936.Ingham, MR0439772, MR0067145, 1914.Littlewood}.}
      \field{journaltitle}{Monatsh. Math.}
      \field{number}{3}
      \field{title}{Bemerkungen zur {A}rbeit: ``{O}n the sign changes of {$(\pi (x)-{\rm li}\ x)$}. {II}'' ({M}onatsh. {M}ath. {\bf 82} (1976), no. 2, 163--175) von {S}. {K}napowski und {P}. {T}urán}
      \field{volume}{82}
      \field{year}{1976}
      \true{nocite}
      \field{pages}{199\bibrangedash 206}
      \range{pages}{8}
      \verb{doi}
      \verb 10.1007/BF01526326
      \endverb
      \verb{mrnumber}
      \verb 0439773
      \endverb
      \verb{urlraw}
      \verb https://doi.org/10.1007/BF01526326
      \endverb
      \verb{url}
      \verb https://doi.org/10.1007/BF01526326
      \endverb
    \endentry
    \entry{MR0429797}{article}{}
      \name{author}{2}{}{%
        {{hash=36befda65f61ea63f097e5518d3ec577}{%
           family={Staś},
           familyi={S\bibinitperiod},
           given={W.},
           giveni={W\bibinitperiod}}}%
        {{hash=b983eb20838b58a1908783452401fc09}{%
           family={Wiertelak},
           familyi={W\bibinitperiod},
           given={K.},
           giveni={K\bibinitperiod}}}%
      }
      \strng{namehash}{9af3bba0a834e4a5fc9ac880f2af2cce}
      \strng{fullhash}{9af3bba0a834e4a5fc9ac880f2af2cce}
      \strng{bibnamehash}{9af3bba0a834e4a5fc9ac880f2af2cce}
      \strng{authorbibnamehash}{9af3bba0a834e4a5fc9ac880f2af2cce}
      \strng{authornamehash}{9af3bba0a834e4a5fc9ac880f2af2cce}
      \strng{authorfullhash}{9af3bba0a834e4a5fc9ac880f2af2cce}
      \field{sortinit}{1}
      \field{sortinithash}{50c6687d7fc80f50136d75228e3c59ba}
      \field{labelnamesource}{author}
      \field{labeltitlesource}{title}
      \field{annotation}{Let $\mathfrak{K}_1,\mathfrak{K}_2$ be ideal classes\mod{\mathfrak{f}} in a number field. In this article, the authors use the second main theorem of the power-sum method to bound $\psi(x,\mathfrak{K}_1)-\psi(x,\mathfrak{K}_2)$ given a zero-free region for Hecke--Landau zeta functions of relevant Hecke characters~$\chi$, and establish a converse as well: Let~$\gamma_1$ be the supremum of numbers~$\gamma$ for which $\psi(x,\mathfrak{K}_1)-\psi(x,\mathfrak{K}_2) \ll xe^{-a(\log x)^\gamma}$ for some positive constant~$a$, and let~$\gamma_2$ be the infimum of numbers~$\gamma$ for which $\prod_{\chi(\mathfrak{K}_1)\ne\chi(\mathfrak{K}_2)}\zeta(s,\chi)$ does not vanish in the region $\sigma > 1-{b}/{(\log |t|)^\gamma}$ for some positive constant~$b$. The authors prove that $\gamma_1 = 1/(1+\gamma_2)$.\par This article cites~\cites{1932.Ingham,MR0087685,MR0146153,MR0060548}.}
      \field{issn}{0065-1036}
      \field{journaltitle}{Acta Arith.}
      \field{number}{2}
      \field{title}{Further applications of Turán's methods to the distribution of prime ideals in ideal classes\mod f}
      \field{volume}{31}
      \field{year}{1976}
      \true{nocite}
      \field{pages}{153\bibrangedash 165}
      \range{pages}{13}
      \verb{doi}
      \verb 10.4064/aa-31-2-153-165
      \endverb
      \verb{mrnumber}
      \verb 0429797
      \endverb
    \endentry
    \entry{MR0447090}{article}{}
      \name{author}{2}{}{%
        {{hash=e542667d9161d12cece57b09b9c6c693}{%
           family={Bays},
           familyi={B\bibinitperiod},
           given={C.},
           giveni={C\bibinitperiod}}}%
        {{hash=08e54e6381028d2260f342dd1c27f66e}{%
           family={Hudson},
           familyi={H\bibinitperiod},
           given={R.\bibnamedelimi H.},
           giveni={R\bibinitperiod\bibinitdelim H\bibinitperiod}}}%
      }
      \strng{namehash}{121290fcdac2ab7b7c2b8152b61c32dc}
      \strng{fullhash}{121290fcdac2ab7b7c2b8152b61c32dc}
      \strng{bibnamehash}{121290fcdac2ab7b7c2b8152b61c32dc}
      \strng{authorbibnamehash}{121290fcdac2ab7b7c2b8152b61c32dc}
      \strng{authornamehash}{121290fcdac2ab7b7c2b8152b61c32dc}
      \strng{authorfullhash}{121290fcdac2ab7b7c2b8152b61c32dc}
      \field{extraname}{1}
      \field{sortinit}{1}
      \field{sortinithash}{50c6687d7fc80f50136d75228e3c59ba}
      \field{labelnamesource}{author}
      \field{labeltitlesource}{title}
      \field{annotation}{The authors describe in detail a refinement of the segmented sieve of Eratosthenes, which they call the dual sieve, designed to lower the execution time. As an illustration, they record the number of primes in the eight reduced residue classes modulo $24$ (from which one can calculate the number of primes in residue classes modulo any divisor of~$24$) up to $10^{11},2 \times 10^{11},\ldots, 10^{12}$. From their table, one easily observes that $\pi(x;24,1)$ is consistently smaller than any other $\pi(x;24,a)$ by an amount that is very roughly $\frac12\pi(\sqrt x)$.}
      \field{journaltitle}{Nordisk Tidskr. Informationsbehandling (BIT)}
      \field{number}{2}
      \field{title}{The segmented sieve of {E}ratosthenes and primes in arithmetic progressions to {$10^{12}$}}
      \field{volume}{17}
      \field{year}{1977}
      \true{nocite}
      \field{pages}{121\bibrangedash 127}
      \range{pages}{7}
      \verb{doi}
      \verb 10.1007/bf01932283
      \endverb
      \verb{mrnumber}
      \verb 0447090
      \endverb
      \verb{urlraw}
      \verb https://doi.org/10.1007/bf01932283
      \endverb
      \verb{url}
      \verb https://doi.org/10.1007/bf01932283
      \endverb
    \endentry
    \entry{MR0506080}{article}{}
      \name{author}{1}{}{%
        {{hash=a8dc47e3fff14028c7ddc24efe554469}{%
           family={Hooley},
           familyi={H\bibinitperiod},
           given={C.},
           giveni={C\bibinitperiod}}}%
      }
      \strng{namehash}{a8dc47e3fff14028c7ddc24efe554469}
      \strng{fullhash}{a8dc47e3fff14028c7ddc24efe554469}
      \strng{bibnamehash}{a8dc47e3fff14028c7ddc24efe554469}
      \strng{authorbibnamehash}{a8dc47e3fff14028c7ddc24efe554469}
      \strng{authornamehash}{a8dc47e3fff14028c7ddc24efe554469}
      \strng{authorfullhash}{a8dc47e3fff14028c7ddc24efe554469}
      \field{sortinit}{1}
      \field{sortinithash}{50c6687d7fc80f50136d75228e3c59ba}
      \field{labelnamesource}{author}
      \field{labeltitlesource}{title}
      \field{annotation}{Assuming GRH and LI, the author shows that $E^\psi(x;q,a)/\sqrt{\phi(q)\log q}$ has a limiting logarithmic distribution function (depending only on~$q$ and not~$a$), as well as the central limit theorem that these functions tend to the standard normal distribution as $q\to\infty$. The proof proceeds by writing the characteristic function of this distribution as a product of Bessel functions. The same result holds for $E^\theta(x;q,a)/\sqrt{\phi(q)\log q}$, except that the mean of the limiting logarithmic distribution depends on whether~$a$ is a quadratic residue\mod q.\par This article cites~\cites{MR0004255}.}
      \field{issn}{0024-6107}
      \field{journaltitle}{J. London Math. Soc. (2)}
      \field{number}{1}
      \field{title}{On the {B}arban-{D}avenport-{H}alberstam theorem. {VII}}
      \field{volume}{16}
      \field{year}{1977}
      \true{nocite}
      \field{pages}{1\bibrangedash 8}
      \range{pages}{8}
      \verb{doi}
      \verb 10.1112/jlms/s2-16.1.1
      \endverb
      \verb{mrnumber}
      \verb 0506080
      \endverb
    \endentry
    \entry{MR0460261}{article}{}
      \name{author}{2}{}{%
        {{hash=9e0609c9931c1b620df97e8c58d07532}{%
           family={Hudson},
           familyi={H\bibinitperiod},
           given={Richard\bibnamedelima H.},
           giveni={R\bibinitperiod\bibinitdelim H\bibinitperiod}}}%
        {{hash=dcdb9f9c423bbcbfa2a60d519ccd6915}{%
           family={Bays},
           familyi={B\bibinitperiod},
           given={Carter},
           giveni={C\bibinitperiod}}}%
      }
      \strng{namehash}{7660af3606957b9dca1512f72f50250e}
      \strng{fullhash}{7660af3606957b9dca1512f72f50250e}
      \strng{bibnamehash}{7660af3606957b9dca1512f72f50250e}
      \strng{authorbibnamehash}{7660af3606957b9dca1512f72f50250e}
      \strng{authornamehash}{7660af3606957b9dca1512f72f50250e}
      \strng{authorfullhash}{7660af3606957b9dca1512f72f50250e}
      \field{sortinit}{1}
      \field{sortinithash}{50c6687d7fc80f50136d75228e3c59ba}
      \field{labelnamesource}{author}
      \field{labeltitlesource}{title}
      \field{annotation}{The goal of this article is to find an elementary explanation (not involving complex analysis) for biases in prime number races. They argue from the Meissel-type formula \begin{align*} \pi(x;q,a) = \Phi(x,x^{1/3};q,a) +& \pi(x^{1/3};q,a) - \sum_{\substack{x^{1/3}<p\le x^{1/2} \\ p\nmid q}} \pi\bigg( \frac xp;q,ap^{-1} \bigg) \\ +& \sum_{\substack{x^{1/3}<p\le x^{1/2} \\ p\nmid q}} \pi(p;q,ap^{-1}) - \sum_{\substack{b\mod q \\ b^2\equiv a\mod q}} \pi\big( (x^{1/3},x^{1/2}] ; q,b\big), \end{align*} where $\Phi(x,y;q,a)$ denotes the number of integers in $(1,x]$ congruent to $a\mod q$ that are free of prime factors up to~$y$. They identify the final term as a persistent bias against quadratic residues~$a$, and support their interpretation with numerical evidence.\par The authors refine the above formula to reduce the occurrences of $x^{1/3}$ to $x^{1/4}$, thus introducing contributions from products of three primes; while observable in the numerical data, these contributions seem to be less significant for large~$x$. From their analysis they make some heuristic predictions, in particular that \[ \fA_1^\pi(x;q,\cN) - (c_q-1) \fA_1^\pi(x;q,\cR) \sim \frac 12 \sum_{n=1}^x \pi(n^{1/2}). \]\par This article cites~\cite{MR0447090,MR0476615,MR1555148,MR0146156,MR0146157,MR0146158,MR0146159,MR0146160,MR0146161,MR0156826,MR0156827,MR0162771,MR0174538,MR0180539,MR0182616A,MR0182616B,MR0200250,MR0177963,MR0068565,MR1544293A,MR0083001,1853.Chebyshev,1891.Phragmen,MR0108470,MR0302590,MR0289452,MR0184920}.}
      \field{issn}{0075-4102}
      \field{journaltitle}{J. Reine Angew. Math.}
      \field{title}{The mean behavior of primes in arithmetic progressions}
      \field{volume}{296}
      \field{year}{1977}
      \true{nocite}
      \field{pages}{80\bibrangedash 99}
      \range{pages}{20}
      \verb{doi}
      \verb 10.1515/crll.1977.296.80
      \endverb
      \verb{mrnumber}
      \verb 0460261
      \endverb
    \endentry
    \entry{MR0466043}{incollection}{}
      \name{author}{2}{}{%
        {{hash=d33db82fc1e83dc064a23d93e5a21891}{%
           family={Knapowski},
           familyi={K\bibinitperiod},
           given={S.},
           giveni={S\bibinitperiod}}}%
        {{hash=e68e101772b88cb1d48e688164bece92}{%
           family={Turán},
           familyi={T\bibinitperiod},
           given={P.},
           giveni={P\bibinitperiod}}}%
      }
      \list{publisher}{1}{%
        {Academic Press, New York}%
      }
      \strng{namehash}{b6f197e0531d2bbfab6a59c048e3b28d}
      \strng{fullhash}{b6f197e0531d2bbfab6a59c048e3b28d}
      \strng{bibnamehash}{b6f197e0531d2bbfab6a59c048e3b28d}
      \strng{authorbibnamehash}{b6f197e0531d2bbfab6a59c048e3b28d}
      \strng{authornamehash}{b6f197e0531d2bbfab6a59c048e3b28d}
      \strng{authorfullhash}{b6f197e0531d2bbfab6a59c048e3b28d}
      \field{extraname}{20}
      \field{sortinit}{1}
      \field{sortinithash}{50c6687d7fc80f50136d75228e3c59ba}
      \field{labelnamesource}{author}
      \field{labeltitlesource}{title}
      \field{annotation}{The authors show, with $p_\nu$ denoting the $\nu$th prime, that when $T$ is sufficiently large, \[ \sum_{\substack{p_\nu\le T\\p_\nu\equiv p_{\nu+1}\equiv 1\mod 4}} 1 > \log^B T \] for some effective constant~$B$. The fact that the left-hand side tends to $\infty$ follows from Littlewood's oscillation theorem for primes\mod4, but no quantitative rate of growth had been established. The proof uses their result from~\cite{MR0302585}, as well as Ingham's application of Fejér kernels as in~\cite{1936.Ingham}. The authors note the open problem about the existence of infinitely many triples of consecutive primes congruent to~$1\mod4$. They also guess that the four possibilities for the pair of congruence classes $(p_\nu, p_{\nu+1}) \mod 4$ are not equally likely.\par This article cites \cite{MR0174538, 1932.Ingham, MR0302585}.}
      \field{booktitle}{Number theory and algebra}
      \field{title}{On prime numbers {$\equiv 1$} resp. {$3\mod4$}}
      \field{year}{1977}
      \true{nocite}
      \field{pages}{157\bibrangedash 165}
      \range{pages}{9}
      \verb{mrnumber}
      \verb 0466043
      \endverb
    \endentry
    \entry{MR607089}{article}{}
      \name{author}{1}{}{%
        {{hash=c6d965df16e600c701d4bcc2888ae34c}{%
           family={Pintz},
           familyi={P\bibinitperiod},
           given={J.},
           giveni={J\bibinitperiod}}}%
      }
      \strng{namehash}{c6d965df16e600c701d4bcc2888ae34c}
      \strng{fullhash}{c6d965df16e600c701d4bcc2888ae34c}
      \strng{bibnamehash}{c6d965df16e600c701d4bcc2888ae34c}
      \strng{authorbibnamehash}{c6d965df16e600c701d4bcc2888ae34c}
      \strng{authornamehash}{c6d965df16e600c701d4bcc2888ae34c}
      \strng{authorfullhash}{c6d965df16e600c701d4bcc2888ae34c}
      \field{extraname}{2}
      \field{sortinit}{1}
      \field{sortinithash}{50c6687d7fc80f50136d75228e3c59ba}
      \field{labelnamesource}{author}
      \field{labeltitlesource}{title}
      \field{annotation}{This article establishes new results on the number of sign changes of $\pi(x)-\li(x)$. In particular it proves the effective result $ W(T)\gg \sqrt{\log T}/\log\log T $ and corresponding results for $W^\Pi(T)$, $W^\theta(T)$, and $W^\psi(T)$ with the same lower bound. Moreover, it establishes that there is necessarily a sign change in the interval $[T,T\exp(63\sqrt{\log T}\log\log T)]$ for $T$ large enough in each of these cases, although the lower bound on such $T$ is effectively computable only in the $\Pi(x)$ and $\psi(x)$ versions.\par Under RH, Ingham's result from 1936 gives in fact a stronger localization theorem. The author here uses power-sum method, and in particular a result of Sós--Turán, to achieve a result under the assumption that RH fails. Ingham's idea to use Fejér kernels is also applied to prove the effective lower bound on $W(T)$ in the absence of an effective localization result.\par This article cites~\cites{1914.Littlewood, 1930.Polya, 1936.Ingham, MR0067145, MR0133309, MR0133308, MR0439771, MR0439772, MR0439773}.}
      \field{issn}{0081-6906}
      \field{journaltitle}{Studia Sci. Math. Hungar.}
      \field{number}{3-4}
      \field{title}{On the remainder term of the prime number formula. {III}. {S}ign changes of {$\pi (x)-{\rm li}x$}}
      \field{volume}{12}
      \field{year}{1977}
      \true{nocite}
      \field{pages}{345\bibrangedash 369 (1980)}
      \range{pages}{-1}
      \verb{mrnumber}
      \verb 607089
      \endverb
    \endentry
    \entry{MR0447151}{incollection}{}
      \name{author}{1}{}{%
        {{hash=c6d965df16e600c701d4bcc2888ae34c}{%
           family={Pintz},
           familyi={P\bibinitperiod},
           given={J.},
           giveni={J\bibinitperiod}}}%
      }
      \list{publisher}{1}{%
        {Soc. Math. France, Paris}%
      }
      \strng{namehash}{c6d965df16e600c701d4bcc2888ae34c}
      \strng{fullhash}{c6d965df16e600c701d4bcc2888ae34c}
      \strng{bibnamehash}{c6d965df16e600c701d4bcc2888ae34c}
      \strng{authorbibnamehash}{c6d965df16e600c701d4bcc2888ae34c}
      \strng{authornamehash}{c6d965df16e600c701d4bcc2888ae34c}
      \strng{authorfullhash}{c6d965df16e600c701d4bcc2888ae34c}
      \field{extraname}{3}
      \field{sortinit}{1}
      \field{sortinithash}{50c6687d7fc80f50136d75228e3c59ba}
      \field{labelnamesource}{author}
      \field{labeltitlesource}{title}
      \field{annotation}{The author begins with a thorough summary of results on sign changes of $\Delta^\pi(x)$ and related problems; he then announces, without proof, several new results of this type. He claims that for $T$ sufficiently large, unconditionally, $W(T) \gg \sqrt{\log T}/\log\log T$, and that there exists $c>0$ such that every interval of the form $[T^c,T]$ contains a sign change of $\Delta^\pi(x)$; ineffectively one can narrow these intervals to the form $[Te^{-\sqrt{\log T}\log\log T},T]$. Even if one restricts to ``big sign changes'', where $\Delta^\pi(x) = \Omega_\pm(\sqrt x\log\log\log x / \log x)$, the author asserts that the number of such sign changes up to $T$ is $\gg \sqrt{\log T} e^{-\sqrt{\log\log T}}$ effectively and $\gg \sqrt{\log T}/(\log\log T)^2$ ineffectively; these sign changes can be localized as well, and the latter inequality even holds for large sign changes of the average of $\Delta^\pi(x)$ over intervals of length $x/\log\log x$. The author further asserts that analogous theorems can be proved for the other prime counting functions, as well as for $\pi(x;4,1,3)$ and some other class of prime races.\par This article cites~\cites{1914.Littlewood, 1936.Ingham, MR0067145, MR0146156, MR0146157, MR0146158, MR0146159, MR0146160, MR0146161, MR0156826, MR0156827, MR0439771, MR0439772, MR0439773}.}
      \field{booktitle}{Journées {A}rithmétiques de {C}aen ({U}niv. {C}aen, {C}aen, 1976)}
      \field{title}{On the sign changes of {$\pi (x)-{\rm li}(x)$}}
      \field{year}{1977}
      \true{nocite}
      \field{pages}{255--265. Astérisque No. 41\bibrangedash 42}
      \range{pages}{-1}
      \verb{mrnumber}
      \verb 0447151
      \endverb
    \endentry
    \entry{MR0476616}{article}{}
      \name{author}{2}{}{%
        {{hash=e542667d9161d12cece57b09b9c6c693}{%
           family={Bays},
           familyi={B\bibinitperiod},
           given={C.},
           giveni={C\bibinitperiod}}}%
        {{hash=08e54e6381028d2260f342dd1c27f66e}{%
           family={Hudson},
           familyi={H\bibinitperiod},
           given={R.\bibnamedelimi H.},
           giveni={R\bibinitperiod\bibinitdelim H\bibinitperiod}}}%
      }
      \strng{namehash}{121290fcdac2ab7b7c2b8152b61c32dc}
      \strng{fullhash}{121290fcdac2ab7b7c2b8152b61c32dc}
      \strng{bibnamehash}{121290fcdac2ab7b7c2b8152b61c32dc}
      \strng{authorbibnamehash}{121290fcdac2ab7b7c2b8152b61c32dc}
      \strng{authornamehash}{121290fcdac2ab7b7c2b8152b61c32dc}
      \strng{authorfullhash}{121290fcdac2ab7b7c2b8152b61c32dc}
      \field{extraname}{2}
      \field{sortinit}{1}
      \field{sortinithash}{50c6687d7fc80f50136d75228e3c59ba}
      \field{labelnamesource}{author}
      \field{labeltitlesource}{title}
      \field{annotation}{The authors determine that $x = 608,981,813,029$ is the smallest $x$ such that $\pi(x;3,2,1) = -1$. A faster version of a previous program of theirs (which had run up to $2.5 \times 10^{11}$) was used to find this sign change. The authors provide graphs of $\pi(x;3,2,1)$ near this first sign change; they highlight that $\pi(x;3,2,1)$ becomes negative at two separate regions near the sign change, before taking on values shortly after that are much more positive. The authors observe that neither $\pi(x;3,2,1)$ nor $\pi(x;4,3,1)$ becomes very negative near the occurrence of its first negative values; in attempts to determine a smaller Skewes number, consequently, they recommend evaluation of $\Delta^\pi(x)$ in regular intervals in order to not miss a ``shallow'' sign change.\par This article cites~\cites{MR529694,1853.Chebyshev,1932.Ingham,MR0162771,MR0083001,MR0202686,1914.Littlewood,MR0108470,MR0476615}.}
      \field{issn}{0025-5718}
      \field{journaltitle}{Math. Comp.}
      \field{number}{142}
      \field{title}{Details of the first region of integers {$x$} with {$\pi _{3,2}(x)<\pi _{3,1}(x)$}}
      \field{volume}{32}
      \field{year}{1978}
      \true{nocite}
      \field{pages}{571\bibrangedash 576}
      \range{pages}{6}
      \verb{doi}
      \verb 10.2307/2006165
      \endverb
      \verb{mrnumber}
      \verb 0476616
      \endverb
      \verb{urlraw}
      \verb https://doi.org/10.2307/2006165
      \endverb
      \verb{url}
      \verb https://doi.org/10.2307/2006165
      \endverb
    \endentry
    \entry{MR0476615}{article}{}
      \name{author}{2}{}{%
        {{hash=dcdb9f9c423bbcbfa2a60d519ccd6915}{%
           family={Bays},
           familyi={B\bibinitperiod},
           given={Carter},
           giveni={C\bibinitperiod}}}%
        {{hash=9e0609c9931c1b620df97e8c58d07532}{%
           family={Hudson},
           familyi={H\bibinitperiod},
           given={Richard\bibnamedelima H.},
           giveni={R\bibinitperiod\bibinitdelim H\bibinitperiod}}}%
      }
      \strng{namehash}{0824391c51610edfe4bd3c404488dea5}
      \strng{fullhash}{0824391c51610edfe4bd3c404488dea5}
      \strng{bibnamehash}{0824391c51610edfe4bd3c404488dea5}
      \strng{authorbibnamehash}{0824391c51610edfe4bd3c404488dea5}
      \strng{authornamehash}{0824391c51610edfe4bd3c404488dea5}
      \strng{authorfullhash}{0824391c51610edfe4bd3c404488dea5}
      \field{extraname}{3}
      \field{sortinit}{1}
      \field{sortinithash}{50c6687d7fc80f50136d75228e3c59ba}
      \field{labelnamesource}{author}
      \field{labeltitlesource}{title}
      \field{annotation}{The authors describe the sixth ``axis crossing region'' (a term they define rigorously) for $\pi(x;4,3,1)$, in which there are $4.1 \times 10^8$ consecutive integers satisfying $\pi(n;4,3,1)<0$. The authors suggest from their findings of surprisingly large axis crossing regions that the number of integers number of integers $n\le x$ with $\pi(n;4,3,1)<0$ might be $\Omega(x/\log x)$, while still respecting the conjecture that this counting function is $o(x)$.\par This article cites~\cites{1853.Chebyshev,MR1555148,MR0460261,MR0184920,MR0146158,MR0156826,MR0162771,MR0174538,MR0068565,MR1544293A,MR0083001,1914.Littlewood,MR1576856,MR0108470}.}
      \field{issn}{0025-5718}
      \field{journaltitle}{Math. Comp.}
      \field{number}{141}
      \field{title}{On the fluctuations of {L}ittlewood for primes of the form {$4n \pm 1$}}
      \field{volume}{32}
      \field{year}{1978}
      \true{nocite}
      \field{pages}{281\bibrangedash 286}
      \range{pages}{6}
      \verb{doi}
      \verb 10.2307/2006277
      \endverb
      \verb{mrnumber}
      \verb 0476615
      \endverb
    \endentry
    \entry{MR0472726}{article}{}
      \name{author}{2}{}{%
        {{hash=dcdb9f9c423bbcbfa2a60d519ccd6915}{%
           family={Bays},
           familyi={B\bibinitperiod},
           given={Carter},
           giveni={C\bibinitperiod}}}%
        {{hash=9e0609c9931c1b620df97e8c58d07532}{%
           family={Hudson},
           familyi={H\bibinitperiod},
           given={Richard\bibnamedelima H.},
           giveni={R\bibinitperiod\bibinitdelim H\bibinitperiod}}}%
      }
      \strng{namehash}{0824391c51610edfe4bd3c404488dea5}
      \strng{fullhash}{0824391c51610edfe4bd3c404488dea5}
      \strng{bibnamehash}{0824391c51610edfe4bd3c404488dea5}
      \strng{authorbibnamehash}{0824391c51610edfe4bd3c404488dea5}
      \strng{authornamehash}{0824391c51610edfe4bd3c404488dea5}
      \strng{authorfullhash}{0824391c51610edfe4bd3c404488dea5}
      \field{extraname}{4}
      \field{sortinit}{1}
      \field{sortinithash}{50c6687d7fc80f50136d75228e3c59ba}
      \field{labelnamesource}{author}
      \field{labeltitlesource}{title}
      \field{annotation}{The authors describe their empirical observation that $\pi(x;24,13,1)<0$ for about a third of the integers between $0.978\times10^{12}$ and $1.094\times10^{12}$. They assert the expectation that the set of such real numbers~$x$ has density~$0$.\par This article cites~\cite{MR0476615, MR0447090, MR0476616, MR0460261, MR0146157, MR0162771, MR0174538, MR0068565, MR0108470}.}
      \field{issn}{0075-4102}
      \field{journaltitle}{J. Reine Angew. Math.}
      \field{title}{{The appearance of tens of billions of integers {$x$} with {$\pi_{24,13}(x)<\pi_{24,1}(x)$} in the vicinity of {$10^{12}$}}}
      \field{volume}{299/300}
      \field{year}{1978}
      \true{nocite}
      \field{pages}{234\bibrangedash 237}
      \range{pages}{4}
      \verb{doi}
      \verb 10.1515/crll.1978.299-300.234
      \endverb
      \verb{mrnumber}
      \verb 0472726
      \endverb
    \endentry
    \entry{MR630377}{article}{}
      \name{author}{1}{}{%
        {{hash=c6d965df16e600c701d4bcc2888ae34c}{%
           family={Pintz},
           familyi={P\bibinitperiod},
           given={J.},
           giveni={J\bibinitperiod}}}%
      }
      \strng{namehash}{c6d965df16e600c701d4bcc2888ae34c}
      \strng{fullhash}{c6d965df16e600c701d4bcc2888ae34c}
      \strng{bibnamehash}{c6d965df16e600c701d4bcc2888ae34c}
      \strng{authorbibnamehash}{c6d965df16e600c701d4bcc2888ae34c}
      \strng{authornamehash}{c6d965df16e600c701d4bcc2888ae34c}
      \strng{authorfullhash}{c6d965df16e600c701d4bcc2888ae34c}
      \field{extraname}{4}
      \field{sortinit}{1}
      \field{sortinithash}{50c6687d7fc80f50136d75228e3c59ba}
      \field{labelnamesource}{author}
      \field{labeltitlesource}{title}
      \field{annotation}{This article establishes lower bounds for the number of sign changes for the error terms of classical prime-counting functions. The main theorem of this article states that for $f=\pi$, $\Pi$, $\theta$, or $\psi$, there exists an absolute constant $Y(f)$ such that for $Y>Y(f)$, \[ W^f(Y)>\frac 1{10^{11}}\frac{\log{Y}}{(\log\log{Y})^3}. \] Interestingly, $Y(\Pi)$ and $Y(\psi)$ are effectively computable in the author's proof, whereas $Y(\pi)$ and $Y(\theta)$ are ineffective constants. \par This article cites~\cites{1914.Littlewood, 1936.Ingham, MR0067145, MR0133309, MR0133308, MR0439771, MR0439772, MR0439773, MR607089}.}
      \field{issn}{0081-6906}
      \field{journaltitle}{Studia Sci. Math. Hungar.}
      \field{number}{1-2}
      \field{title}{On the remainder term of the prime number formula. {IV}. {S}ign changes of {$\pi (x)-{\li}x$}}
      \field{volume}{13}
      \field{year}{1978}
      \true{nocite}
      \field{pages}{29\bibrangedash 42 (1981)}
      \range{pages}{-1}
      \verb{mrnumber}
      \verb 630377
      \endverb
    \endentry
    \entry{MR529694}{article}{}
      \name{author}{2}{}{%
        {{hash=e542667d9161d12cece57b09b9c6c693}{%
           family={Bays},
           familyi={B\bibinitperiod},
           given={C.},
           giveni={C\bibinitperiod}}}%
        {{hash=08e54e6381028d2260f342dd1c27f66e}{%
           family={Hudson},
           familyi={H\bibinitperiod},
           given={R.\bibnamedelimi H.},
           giveni={R\bibinitperiod\bibinitdelim H\bibinitperiod}}}%
      }
      \strng{namehash}{121290fcdac2ab7b7c2b8152b61c32dc}
      \strng{fullhash}{121290fcdac2ab7b7c2b8152b61c32dc}
      \strng{bibnamehash}{121290fcdac2ab7b7c2b8152b61c32dc}
      \strng{authorbibnamehash}{121290fcdac2ab7b7c2b8152b61c32dc}
      \strng{authornamehash}{121290fcdac2ab7b7c2b8152b61c32dc}
      \strng{authorfullhash}{121290fcdac2ab7b7c2b8152b61c32dc}
      \field{extraname}{5}
      \field{sortinit}{1}
      \field{sortinithash}{50c6687d7fc80f50136d75228e3c59ba}
      \field{labelnamesource}{author}
      \field{labeltitlesource}{title}
      \field{annotation}{The authors of this article determine by computation the locations where $\pi(x;4,3,1) < 0$, and where $\pi(x;8,a,1) < 0$ for any $a\in\{3,5,7\}$, for $x$ up to $10^{12}$. For $x < 10^9$, a check was made at every prime; for $10^9 \le x \le 10^{12}$, a check was made every $10^7$ integers, with additional checks in between if $\pi(x;q,a,b)$ was found to be near zero. They then organize these locations into ``axis-crossing regions'' (ACRs) $[m,n]$, where $\pi(m;q,a,1) = \pi(n;q,a,1) = -1$ and $\pi(x;q,a,1) \ge 0$ for all $x$ outside an ACR, with $m$ at least twice as large as the upper bound for the previous ACR. \par For $q=4$, they find six distinct ACRs under $10^{12}$. For $(q,a) = (8,5)$, they find two ACRs under $10^{12}$ and find no ACRs for $(q,a) = (8,3)$ or $(q,a) = (8,7)$. They compare their computations to earlier published results from Leech \cite{MR0083001}, Shanks \cite{MR0108470}, and an unpublished communication from Lehmer (dated October 29, 1975). While their results overlap with Leech and Shanks for $q=4$ for $x \le 3 \cdot 10^6$, they find that their new information contradicts a prior characterization of the ACRs as mostly consisting of sparse, tiny intervals. For example, one ACR below $x < 2 \cdot 10^{10}$ contains $5\cdot10^8$ integers where $\pi(x;4,3,1) < 0$; another ACR between $37\cdot10^9$ and $39\cdot10^9$ contains $1.2\cdot10^9$ integers with $\pi(x;8,5,1) < 0$. Consequently, they argue that for large $x$, such regions may be more typical than sign-changes being sparse, isolated points. \par This article cites~\cites{MR0476615,MR0447090,1853.Chebyshev,MR0146157,MR0146158,MR0083001,1914.Littlewood,MR0108470}.}
      \field{issn}{0161-1712}
      \field{journaltitle}{Internat. J. Math. Math. Sci.}
      \field{number}{1}
      \field{title}{Numerical and graphical description of all axis crossing regions for the moduli {$4$} and {$8$} which occur before {$10^{12}$}}
      \field{volume}{2}
      \field{year}{1979}
      \true{nocite}
      \field{pages}{111\bibrangedash 119}
      \range{pages}{9}
      \verb{doi}
      \verb 10.1155/S0161171279000119
      \endverb
      \verb{mrnumber}
      \verb 529694
      \endverb
      \verb{urlraw}
      \verb https://doi.org/10.1155/S0161171279000119
      \endverb
      \verb{url}
      \verb https://doi.org/10.1155/S0161171279000119
      \endverb
    \endentry
    \entry{MR534235}{article}{}
      \name{author}{1}{}{%
        {{hash=687d9a8bcde5764db23c9a6841909f8e}{%
           family={Besenfelder},
           familyi={B\bibinitperiod},
           given={H.-J.},
           giveni={H\bibinithyphendelim J\bibinitperiod}}}%
      }
      \list{language}{1}{%
        {German}%
      }
      \strng{namehash}{687d9a8bcde5764db23c9a6841909f8e}
      \strng{fullhash}{687d9a8bcde5764db23c9a6841909f8e}
      \strng{bibnamehash}{687d9a8bcde5764db23c9a6841909f8e}
      \strng{authorbibnamehash}{687d9a8bcde5764db23c9a6841909f8e}
      \strng{authornamehash}{687d9a8bcde5764db23c9a6841909f8e}
      \strng{authorfullhash}{687d9a8bcde5764db23c9a6841909f8e}
      \field{extraname}{1}
      \field{sortinit}{1}
      \field{sortinithash}{50c6687d7fc80f50136d75228e3c59ba}
      \field{labelnamesource}{author}
      \field{labeltitlesource}{title}
      \field{annotation}{Using an existing explicit formula for general Mellin-transform pairs, the author shows that \begin{equation*} \begin{split} 2 \sqrt{\pi y}\sum_{\substack{0<\sigma<1 \\ L(\sigma+i\gamma,\chi_{-4})=0}} e^{y(\sigma-1/2+i\gamma)^2 }&=\log \frac4\pi -2 \sum_{n=1}^\infty \frac{\Lambda(n)\chi_{-4}(n)}{\sqrt n}e^{-{(\log n)^2}/4y}\\ & \qquad{}-C_0+2\int_0^{\infty} \frac{e^{-x^2/4y + x/2}-1}{1-e^{2x}} \,dx. \end{split} \end{equation*} From this identity, he proves unconditionally that $$\lim_{x\to\infty} \sum_p \chi_{-4}(p) \frac{\log p }{\sqrt p} e^{-(\log^2{p})/x} = -\infty.$$ (Note: the author, Hans--Joachim Besenfelder, soon changed his last name to Bentz and began to publish under that name.)\par This article cites~\cites{MR1555148, MR0174538, MR1544293A, 1853.Chebyshev}.}
      \field{issn}{0075-4102}
      \field{journaltitle}{J. Reine Angew. Math.}
      \field{title}{Über eine {V}ermutung von {T}schebyschef. {I}}
      \field{volume}{307/308}
      \field{year}{1979}
      \true{nocite}
      \field{pages}{411\bibrangedash 417}
      \range{pages}{7}
      \verb{doi}
      \verb 10.1515/crll.1979.307-308.411
      \endverb
      \verb{mrnumber}
      \verb 534235
      \endverb
      \verb{urlraw}
      \verb https://doi.org/10.1515/crll.1979.307-308.411
      \endverb
      \verb{url}
      \verb https://doi.org/10.1515/crll.1979.307-308.411
      \endverb
    \endentry
    \entry{MR549977}{article}{}
      \name{author}{1}{}{%
        {{hash=f6fbcde367a0ed5df3da06b96b6d8eff}{%
           family={Riele},
           familyi={R\bibinitperiod},
           given={H.\bibnamedelimi J.\bibnamedelimi J.},
           giveni={H\bibinitperiod\bibinitdelim J\bibinitperiod\bibinitdelim J\bibinitperiod},
           prefix={te},
           prefixi={t\bibinitperiod}}}%
      }
      \strng{namehash}{f6fbcde367a0ed5df3da06b96b6d8eff}
      \strng{fullhash}{f6fbcde367a0ed5df3da06b96b6d8eff}
      \strng{bibnamehash}{f6fbcde367a0ed5df3da06b96b6d8eff}
      \strng{authorbibnamehash}{f6fbcde367a0ed5df3da06b96b6d8eff}
      \strng{authornamehash}{f6fbcde367a0ed5df3da06b96b6d8eff}
      \strng{authorfullhash}{f6fbcde367a0ed5df3da06b96b6d8eff}
      \field{extraname}{1}
      \field{sortinit}{1}
      \field{sortinithash}{50c6687d7fc80f50136d75228e3c59ba}
      \field{labelnamesource}{author}
      \field{labeltitlesource}{title}
      \field{annotation}{The author introduces some modifications of the method of Jurkat and Peyerimhoff~\cite{MR429789} and proves that \[ \liminf_{x \to \infty} E^M(x) < -0.843 \quad \text{and} \quad \limsup_{x \to \infty} E^M(x) > 0.860. \] The calculation took several hundred CPU-hours on a CDC Cyber 73/173 system. \par This article cites~\cite{1897.Mertens,MR0330069,MR228456,MR429789}.}
      \field{issn}{0075-4102}
      \field{journaltitle}{J. Reine Angew. Math.}
      \field{title}{Computations concerning the conjecture of {M}ertens}
      \field{volume}{311(312)}
      \field{year}{1979}
      \true{nocite}
      \field{pages}{356\bibrangedash 360}
      \range{pages}{5}
      \verb{doi}
      \verb 10.1515/crll.1979.311-312.356
      \endverb
      \verb{mrnumber}
      \verb 549977
      \endverb
      \verb{urlraw}
      \verb https://doi.org/10.1515/crll.1979.311-312.356
      \endverb
      \verb{url}
      \verb https://doi.org/10.1515/crll.1979.311-312.356
      \endverb
    \endentry
    \entry{MR595317}{article}{}
      \name{author}{2}{}{%
        {{hash=6bf42c5abf099333fae470886fff8616}{%
           family={Bentz},
           familyi={B\bibinitperiod},
           given={H.-J.},
           giveni={H\bibinithyphendelim J\bibinitperiod}}}%
        {{hash=c6d965df16e600c701d4bcc2888ae34c}{%
           family={Pintz},
           familyi={P\bibinitperiod},
           given={J.},
           giveni={J\bibinitperiod}}}%
      }
      \strng{namehash}{8ee030157e0efba486591e59a6b5401b}
      \strng{fullhash}{8ee030157e0efba486591e59a6b5401b}
      \strng{bibnamehash}{8ee030157e0efba486591e59a6b5401b}
      \strng{authorbibnamehash}{8ee030157e0efba486591e59a6b5401b}
      \strng{authornamehash}{8ee030157e0efba486591e59a6b5401b}
      \strng{authorfullhash}{8ee030157e0efba486591e59a6b5401b}
      \field{extraname}{1}
      \field{sortinit}{1}
      \field{sortinithash}{50c6687d7fc80f50136d75228e3c59ba}
      \field{labelnamesource}{author}
      \field{labeltitlesource}{title}
      \field{annotation}{The first section of this article offers a short, conventional history of prime number races, specifically citing Shanks's computational work and heuristics from~\cite{MR0108470} as motivation for its results. Let $\ell_1$ be a quadratic residue\mod q and $\ell_2$ a quadratic nonresidue\mod q. Suppose that Dirichlet $L$-functions\mod q satisfying the condition that all zeros $\beta+i\gamma$ satisfy the inequality $\beta^2-\gamma^2 \le \frac 14$ (a ``bowtie'' assumption). Then for $0 \le \alpha < 1/2$, \[ \sum_{p\equiv\ell_1\mod q} \frac{\log p}{p^\alpha} e^{-(\log p)^2/x} - \sum_{p\equiv\ell_2\mod q} \frac{\log p}{p^\alpha} e^{-(\log p)^2/x} \sim \frac{c_q}{\phi(q)} \sqrt{\pi x} \cdot e^{\frac x4 (\frac 12 - \alpha)^2}, \] and in particular tends to infinity. By computations of Spira, this result is unconditional for $q\le24$.\par This article cites~\cites{MR662791, MR534235, MR552462, 1853.Chebyshev, MR1555148, MR0146156, MR0146157, MR0146158, MR0146159, MR0146160, MR0146161, MR0156826, MR0156827, MR0162771, MR0174538, MR0180539, MR0182616A, MR0182616B, MR0200250, MR1544293A, MR1544293B, MR0083001, MR0108470}.}
      \field{issn}{0026-9255}
      \field{journaltitle}{Monatsh. Math.}
      \field{number}{2}
      \field{title}{Quadratic residues and the distribution of prime numbers}
      \field{volume}{90}
      \field{year}{1980}
      \true{nocite}
      \field{pages}{91\bibrangedash 100}
      \range{pages}{10}
      \verb{doi}
      \verb 10.1007/BF01303260
      \endverb
      \verb{mrnumber}
      \verb 595317
      \endverb
      \verb{urlraw}
      \verb https://doi.org/10.1007/BF01303260
      \endverb
      \verb{url}
      \verb https://doi.org/10.1007/BF01303260
      \endverb
    \endentry
    \entry{MR591612}{article}{}
      \name{author}{2}{}{%
        {{hash=9f64a229f3d656a4a04d7eda97b0e30a}{%
           family={Bentz},
           familyi={B\bibinitperiod},
           given={Hans-J.},
           giveni={H\bibinithyphendelim J\bibinitperiod}}}%
        {{hash=d192c3b49f427e2494b8aaea251e390b}{%
           family={Pintz},
           familyi={P\bibinitperiod},
           given={János},
           giveni={J\bibinitperiod}}}%
      }
      \list{language}{1}{%
        {German}%
      }
      \strng{namehash}{e6be682d0e7c51c0cf93fc1a6932936c}
      \strng{fullhash}{e6be682d0e7c51c0cf93fc1a6932936c}
      \strng{bibnamehash}{e6be682d0e7c51c0cf93fc1a6932936c}
      \strng{authorbibnamehash}{e6be682d0e7c51c0cf93fc1a6932936c}
      \strng{authornamehash}{e6be682d0e7c51c0cf93fc1a6932936c}
      \strng{authorfullhash}{e6be682d0e7c51c0cf93fc1a6932936c}
      \field{extraname}{2}
      \field{sortinit}{1}
      \field{sortinithash}{50c6687d7fc80f50136d75228e3c59ba}
      \field{labelnamesource}{author}
      \field{labeltitlesource}{title}
      \field{annotation}{Chebyshev conjectured that $\lim_{x \to \infty} \pi_e(x,\chi_{-4}) = -\infty$, which is equivalent~\cites{MR1555148,MR1544293A, MR1544293B} to GRH for $L(s,\chi_{-4})$. Knapowski and Turán showed that these assertions are further equivalent to $\lim_{x \to \infty} \theta_l(x,r,\chi_{-4}) = -\infty$. In this article, the authors establish a related implication: Let $\chi$ be a quadratic character\mod q. Suppose that all nonreal zeros $\rho= \beta + i\gamma$ of $L(s,\chi_1)$ satisfy $\beta^2-\gamma^2 < 1/4$, a property implied by GRH$(\frac{\sqrt3}2)$. Then for $0 \leq \alpha < 1/2$ they prove that \[ \lim_{x \to \infty} \sum_{p} \chi_1(p) \frac{\log p}{p^\alpha} \cdot \exp\biggl( - \frac{\log^2 p}{x} \biggr) = -\infty.\]\par This article cites~\cites{MR595317,MR552462,MR534235,MR1555148,MR0174538,MR0182616A,MR0087685,MR0108470,1853.Chebyshev}.}
      \field{issn}{0025-5874}
      \field{journaltitle}{Math. Z.}
      \field{number}{1}
      \field{title}{Über eine Verallgemeinerung des Tschebyschef-Problems}
      \field{volume}{174}
      \field{year}{1980}
      \true{nocite}
      \field{pages}{35\bibrangedash 41}
      \range{pages}{7}
      \verb{doi}
      \verb 10.1007/BF01215079
      \endverb
      \verb{mrnumber}
      \verb 591612
      \endverb
    \endentry
    \entry{MR552462}{article}{}
      \name{author}{1}{}{%
        {{hash=687d9a8bcde5764db23c9a6841909f8e}{%
           family={Besenfelder},
           familyi={B\bibinitperiod},
           given={H.-J.},
           giveni={H\bibinithyphendelim J\bibinitperiod}}}%
      }
      \list{language}{1}{%
        {German}%
      }
      \strng{namehash}{687d9a8bcde5764db23c9a6841909f8e}
      \strng{fullhash}{687d9a8bcde5764db23c9a6841909f8e}
      \strng{bibnamehash}{687d9a8bcde5764db23c9a6841909f8e}
      \strng{authorbibnamehash}{687d9a8bcde5764db23c9a6841909f8e}
      \strng{authornamehash}{687d9a8bcde5764db23c9a6841909f8e}
      \strng{authorfullhash}{687d9a8bcde5764db23c9a6841909f8e}
      \field{extraname}{2}
      \field{sortinit}{1}
      \field{sortinithash}{50c6687d7fc80f50136d75228e3c59ba}
      \field{labelnamesource}{author}
      \field{labeltitlesource}{title}
      \field{annotation}{The author proves unconditionally that for all $0\le\alpha\le\frac12$, $$\lim_{x\to\infty} \sum_{p} \frac{\chi_{-4} (p) \log p}{p^{\alpha}} e^{-(\log^2 p)/x} = -\infty,$$ using an explicit formula similar to the one in his prior work~\cite{MR534235}.\par This article cites~\cites{MR1555148, MR0174538, MR1544293A, MR0300853, MR0300854,1853.Chebyshev, MR534235, MR0156826, MR1544293B}.}
      \field{issn}{0075-4102}
      \field{journaltitle}{J. Reine Angew. Math.}
      \field{title}{Über eine {V}ermutung von {T}schebyschef. {II}}
      \field{volume}{313}
      \field{year}{1980}
      \true{nocite}
      \field{pages}{52\bibrangedash 58}
      \range{pages}{7}
      \verb{doi}
      \verb 10.1515/crll.1980.313.52
      \endverb
      \verb{mrnumber}
      \verb 552462
      \endverb
      \verb{urlraw}
      \verb https://doi.org/10.1515/crll.1980.313.52
      \endverb
      \verb{url}
      \verb https://doi.org/10.1515/crll.1980.313.52
      \endverb
    \endentry
    \entry{MR598886}{article}{}
      \name{author}{1}{}{%
        {{hash=e38330a7a234159955b75bdfae8cd3cd}{%
           family={Gallagher},
           familyi={G\bibinitperiod},
           given={P.\bibnamedelimi X.},
           giveni={P\bibinitperiod\bibinitdelim X\bibinitperiod}}}%
      }
      \strng{namehash}{e38330a7a234159955b75bdfae8cd3cd}
      \strng{fullhash}{e38330a7a234159955b75bdfae8cd3cd}
      \strng{bibnamehash}{e38330a7a234159955b75bdfae8cd3cd}
      \strng{authorbibnamehash}{e38330a7a234159955b75bdfae8cd3cd}
      \strng{authornamehash}{e38330a7a234159955b75bdfae8cd3cd}
      \strng{authorfullhash}{e38330a7a234159955b75bdfae8cd3cd}
      \field{sortinit}{1}
      \field{sortinithash}{50c6687d7fc80f50136d75228e3c59ba}
      \field{labelnamesource}{author}
      \field{labeltitlesource}{title}
      \field{annotation}{Under RH, the author proves that $E^\psi(x) \ll (\log \log x)^2$ except on a set of finite measure, and that for any function $f(x)$ tending to infinity, $E^\psi(x) \ll f(x)$ except on a set of density~$0$; the proofs are similar to showing that $E^\psi(x)$ has a limiting logarithmic distribution. The method also provides short proofs of Cramér's conditional estimates~\cite{MR1544509} $\int_1^X E^\psi(x)^2\, dx \ll X$ and $\int_1^X E^\psi(x)^2 \frac{dx}{x} \sim C \log X$ (for an explicit constant~$C$), as well as of Selberg's conditional result on the normal density of primes in short intervals.\par This article cites~\cites{MR1544509}.}
      \field{issn}{0065-1036}
      \field{journaltitle}{Acta Arith.}
      \field{title}{Some consequences of the {R}iemann hypothesis}
      \field{volume}{37}
      \field{year}{1980}
      \true{nocite}
      \field{pages}{339\bibrangedash 343}
      \range{pages}{5}
      \verb{doi}
      \verb 10.4064/aa-37-1-339-343
      \endverb
      \verb{mrnumber}
      \verb 598886
      \endverb
      \verb{urlraw}
      \verb https://doi.org/10.4064/aa-37-1-339-343
      \endverb
      \verb{url}
      \verb https://doi.org/10.4064/aa-37-1-339-343
      \endverb
    \endentry
    \entry{MR552467}{article}{}
      \name{author}{1}{}{%
        {{hash=9e0609c9931c1b620df97e8c58d07532}{%
           family={Hudson},
           familyi={H\bibinitperiod},
           given={Richard\bibnamedelima H.},
           giveni={R\bibinitperiod\bibinitdelim H\bibinitperiod}}}%
      }
      \strng{namehash}{9e0609c9931c1b620df97e8c58d07532}
      \strng{fullhash}{9e0609c9931c1b620df97e8c58d07532}
      \strng{bibnamehash}{9e0609c9931c1b620df97e8c58d07532}
      \strng{authorbibnamehash}{9e0609c9931c1b620df97e8c58d07532}
      \strng{authornamehash}{9e0609c9931c1b620df97e8c58d07532}
      \strng{authorfullhash}{9e0609c9931c1b620df97e8c58d07532}
      \field{extraname}{1}
      \field{sortinit}{1}
      \field{sortinithash}{50c6687d7fc80f50136d75228e3c59ba}
      \field{labelnamesource}{author}
      \field{labeltitlesource}{title}
      \field{annotation}{In this article, the author outlines a combinatorial principle that seeks to explain various effects and biases in comparative prime number theory. He highlights Riemann's original explicit formula $\pi(x)\sim \li(x)-\tfrac 12 \li (x^{1/2})-\tfrac 13 \li (x^{1/3})+\cdots$ and connects it to Chebyshev's observation, which can be seen as approximating $\pi(x;4,3,1)$ by half the number of prime squares. Arguing from a generalisation of an exact formula of Meissel, the author deduces, in the example of primes\mod 4, that an ``excess'' in the number of integers of the form $pq$, where $p$ and $q$ are prime, in the class $1\mod 4$ must result in a corresponding ``deficiency'' in the number of primes of exactly this magnitude, that is, half the number of prime squares. A combinatorial observation gives a reason for such an excess: in counting integers that are the product of two primes from a set, products of distinct primes are counted twice (as $pq$ and $qp$), while the prime squares are not. The author then provides similar arguments for why cubic and higher order effects should exist. Along with describing the combinatorial principle in generality, he provides details of some numerical investigations into these effects.\par This article cites~\cites{1853.Chebyshev, 1914.Littlewood, MR0108470, MR0162771, MR0174538, MR0289452}.}
      \field{issn}{0075-4102}
      \field{journaltitle}{J. Reine Angew. Math.}
      \field{title}{A common combinatorial principle underlies {R}iemann's formula, the {C}hebyshev phenomenon, and other subtle effects in comparative prime number theory. {I}}
      \field{volume}{313}
      \field{year}{1980}
      \true{nocite}
      \field{pages}{133\bibrangedash 150}
      \range{pages}{18}
      \verb{doi}
      \verb 10.1515/crll.1980.313.133
      \endverb
      \verb{mrnumber}
      \verb 552467
      \endverb
      \verb{urlraw}
      \verb https://doi.org/10.1515/crll.1980.313.133
      \endverb
      \verb{url}
      \verb https://doi.org/10.1515/crll.1980.313.133
      \endverb
    \endentry
    \entry{MR2631002}{book}{}
      \name{author}{1}{}{%
        {{hash=fbe9313e8e03d825339429599a21e678}{%
           family={Monach},
           familyi={M\bibinitperiod},
           given={William\bibnamedelima Reynolds},
           giveni={W\bibinitperiod\bibinitdelim R\bibinitperiod}}}%
      }
      \list{publisher}{1}{%
        {ProQuest LLC, Ann Arbor, MI}%
      }
      \strng{namehash}{fbe9313e8e03d825339429599a21e678}
      \strng{fullhash}{fbe9313e8e03d825339429599a21e678}
      \strng{bibnamehash}{fbe9313e8e03d825339429599a21e678}
      \strng{authorbibnamehash}{fbe9313e8e03d825339429599a21e678}
      \strng{authornamehash}{fbe9313e8e03d825339429599a21e678}
      \strng{authorfullhash}{fbe9313e8e03d825339429599a21e678}
      \field{sortinit}{1}
      \field{sortinithash}{50c6687d7fc80f50136d75228e3c59ba}
      \field{labelnamesource}{author}
      \field{labeltitlesource}{title}
      \field{annotation}{In Chapter 2 of this thesis, assuming RH and LI, the author shows that the limiting logarithmic distribution of $E^\psi$ is equal to the distribution function of the random variable $\sum_{\gamma>0} 2\sin(2\pi\theta_\gamma)/|\rho|$, where the~$\theta_\gamma$ are independent random variables uniformly distributed in $[0,1)$. The remainder of the chapter describes methods of computing bounds on this distribution function on the interval $[-1,1]$ (a table of values is included, as are the programs used to generate said table) and determining an asymptotic formula for its derivative.\par This article cites~\cite{MR0104638,MR0137689,MR1501802}.}
      \field{note}{Thesis (Ph.D.)--University of Michigan)}
      \field{pagetotal}{180}
      \field{title}{Numerical Investigation of Several Problems in Number Theory}
      \field{year}{1980}
      \true{nocite}
      \verb{mrnumber}
      \verb 2631002
      \endverb
      \verb{urlraw}
      \verb http://gateway.proquest.com/openurl?url_ver=Z39.88-2004&rft_val_fmt=info:ofi/fmt:kev:mtx:dissertation&res_dat=xri:pqdiss&rft_dat=xri:pqdiss:8106192
      \endverb
      \verb{url}
      \verb http://gateway.proquest.com/openurl?url_ver=Z39.88-2004&rft_val_fmt=info:ofi/fmt:kev:mtx:dissertation&res_dat=xri:pqdiss&rft_dat=xri:pqdiss:8106192
      \endverb
    \endentry
    \entry{MR634679}{inproceedings}{}
      \name{author}{1}{}{%
        {{hash=98296eedbd8e7fdb593797362d5ac2c7}{%
           family={Montgomery},
           familyi={M\bibinitperiod},
           given={H.\bibnamedelimi L.},
           giveni={H\bibinitperiod\bibinitdelim L\bibinitperiod}}}%
      }
      \list{publisher}{1}{%
        {Queen's Univ., Kingston, Ont.}%
      }
      \strng{namehash}{98296eedbd8e7fdb593797362d5ac2c7}
      \strng{fullhash}{98296eedbd8e7fdb593797362d5ac2c7}
      \strng{bibnamehash}{98296eedbd8e7fdb593797362d5ac2c7}
      \strng{authorbibnamehash}{98296eedbd8e7fdb593797362d5ac2c7}
      \strng{authornamehash}{98296eedbd8e7fdb593797362d5ac2c7}
      \strng{authorfullhash}{98296eedbd8e7fdb593797362d5ac2c7}
      \field{extraname}{1}
      \field{sortinit}{1}
      \field{sortinithash}{50c6687d7fc80f50136d75228e3c59ba}
      \field{labelnamesource}{author}
      \field{labeltitlesource}{title}
      \field{annotation}{Section~3 of this article examines random variables of the form $X=\sum_{k=1}^\infty r_k \sin(2\pi\theta_k)$ for $\{r_k\}$ a decreasing $\ell^2$ sequence, where the $\theta_k$ are independently uniformly distributed on~$\R/\Z$. The author establishes, for any integer $K\ge1$, the bounds \begin{align*} P\bigg(X\geq 2\sum_{k=1}^K r_k\bigg) &\leq \exp\left(-\frac 34 \bigg(\sum_{k=1}^K r_k\bigg)^2\bigg(\sum_{k=K+1}^\infty r_k^2\bigg)^{-1}\right) \\ P\bigg(X\geq \frac 12 \sum_{k=1}^K r_k\bigg) &\geq \frac 1{2^{40}}\exp\left(-100 \bigg(\sum_{k=1}^K r_k\bigg)^2\bigg(\sum_{k=K+1}^\infty r_k^2\bigg)^{-1}\right); \end{align*} in addition, if $\delta$ is sufficiently small and $\sum_{k\colon r_k>\delta}(r_k-\delta)\geq V$, then \[ P(X\geq V)\geq \frac 12 \exp\bigg(-\frac 12 \sum_{k\colon r_k>\delta} \log\frac{\pi^2 r_k}{2\delta}\bigg). \] These results can be applied to the limiting logarithmic distribution function of $E^\psi(x)$, which (assuming RH and LI) is the same as the distribution of the random variable $Y=\sum_{\gamma>0}\frac 2{|\rho|}\sin(2\pi\theta_\rho)$. In particular, the second result implies that there exist constants $0<c_1<c_2$ such that \[ \exp\big( -c_2\sqrt{v}e^{\sqrt{2\pi v}} \big) \le P(Y > v) \le \exp\big( -c_1\sqrt{v}e^{\sqrt{2\pi v}} \big), \] which suggests the conjecture $\displaystyle\limsup_{x\to\infty}\frac{E^\psi(x)}{(\log\log\log{x})^2}=\frac 1{2\pi}$ and $\displaystyle\liminf_{x\to\infty}\frac{E^\psi(x)}{(\log\log\log{x})^2}=-\frac 1{2\pi}$.}
      \field{booktitle}{Proceedings of the {Q}ueen's {N}umber {T}heory {C}onference, 1979 ({K}ingston, {O}nt., 1979)}
      \field{series}{Queen's Papers in Pure and Appl. Math.}
      \field{title}{The zeta function and prime numbers}
      \field{volume}{54}
      \field{year}{1980}
      \true{nocite}
      \field{pages}{1\bibrangedash 31}
      \range{pages}{31}
      \verb{mrnumber}
      \verb 634679
      \endverb
    \endentry
    \entry{MR585891}{article}{}
      \name{author}{1}{}{%
        {{hash=c6d965df16e600c701d4bcc2888ae34c}{%
           family={Pintz},
           familyi={P\bibinitperiod},
           given={J.},
           giveni={J\bibinitperiod}}}%
      }
      \strng{namehash}{c6d965df16e600c701d4bcc2888ae34c}
      \strng{fullhash}{c6d965df16e600c701d4bcc2888ae34c}
      \strng{bibnamehash}{c6d965df16e600c701d4bcc2888ae34c}
      \strng{authorbibnamehash}{c6d965df16e600c701d4bcc2888ae34c}
      \strng{authornamehash}{c6d965df16e600c701d4bcc2888ae34c}
      \strng{authorfullhash}{c6d965df16e600c701d4bcc2888ae34c}
      \field{extraname}{5}
      \field{sortinit}{1}
      \field{sortinithash}{50c6687d7fc80f50136d75228e3c59ba}
      \field{labelnamesource}{author}
      \field{labeltitlesource}{title}
      \field{annotation}{This article contains explicit oscillation results under the assumption that RH is false. Suppose that $\rho_0 = \beta_0+i\gamma_0$ is a nontrivial zero of~$\zeta(s)$. Let $0<\ep\le0.02$, and set $A=40000 \ep^{-2}\log\gamma_0$. Then for $H$ sufficiently large in terms of $\rho_0$, there exist $x_+$ and $x_-$ in the interval $[H,H^A]$ such that \[ \Delta^\pi(x_+) > (1-\ep)\frac{x_+^{\beta_0}}{|\rho_0|\log x_+} \quad\text{and}\quad \Delta^\pi(x_-) < -(1-\ep)\frac{x_-^{\beta_0}}{|\rho_0|\log x_-}, \] and the same for $\Delta^\Pi(x)$; similarly, the result holds without the factor $\log x$ in the denominator for $\Delta^\psi(x)$ and~$\Delta^\theta(x)$. In all these theorems, if in addition $\beta_0>\frac12+\ep$ and $\gamma_0$ is sufficiently large in terms of $\ep$, then by replacing the factor $(1-\ep)/|\rho_0|$ by the smaller $1/\gamma_0^{1+\ep}$, the localization can be improved to the interval $[H,H^{1+\ep}]$. Consequently, $W^f(T)/\log\log T$ tends to infinity for each of the four functions $f\in\{\pi,\Pi,\theta,\psi\}$ (the case where RH is true having been handled by Ingham~\cite{1936.Ingham}).\par This article cites~\cites{MR0146156, MR0146157, MR0146158, MR0146159, MR0146160, MR0146161, MR0156826, MR0156827, MR0162771, MR0174538, MR0180539, MR0182616A, MR0182616B, 1914.Littlewood, 1936.Ingham, MR0439771, MR634679}.}
      \field{issn}{0065-1036}
      \field{journaltitle}{Acta Arith.}
      \field{number}{4}
      \field{title}{On the remainder term of the prime number formula. {I}. {O}n a problem of {L}ittlewood}
      \field{volume}{36}
      \field{year}{1980}
      \true{nocite}
      \field{pages}{341\bibrangedash 365}
      \range{pages}{25}
      \verb{doi}
      \verb 10.4064/aa-36-4-341-365
      \endverb
      \verb{mrnumber}
      \verb 585891
      \endverb
      \verb{urlraw}
      \verb https://doi.org/10.4064/aa-36-4-341-365
      \endverb
      \verb{url}
      \verb https://doi.org/10.4064/aa-36-4-341-365
      \endverb
    \endentry
    \entry{MR598876}{article}{}
      \name{author}{1}{}{%
        {{hash=c6d965df16e600c701d4bcc2888ae34c}{%
           family={Pintz},
           familyi={P\bibinitperiod},
           given={J.},
           giveni={J\bibinitperiod}}}%
      }
      \strng{namehash}{c6d965df16e600c701d4bcc2888ae34c}
      \strng{fullhash}{c6d965df16e600c701d4bcc2888ae34c}
      \strng{bibnamehash}{c6d965df16e600c701d4bcc2888ae34c}
      \strng{authorbibnamehash}{c6d965df16e600c701d4bcc2888ae34c}
      \strng{authornamehash}{c6d965df16e600c701d4bcc2888ae34c}
      \strng{authorfullhash}{c6d965df16e600c701d4bcc2888ae34c}
      \field{extraname}{6}
      \field{sortinit}{1}
      \field{sortinithash}{50c6687d7fc80f50136d75228e3c59ba}
      \field{labelnamesource}{author}
      \field{labeltitlesource}{title}
      \field{annotation}{This article investigates the connection between the zero free region of $\zeta(s)$ and the size of the remainder term in the prime number theorem. Let $\eta\colon [1,\infty)\to(0,\frac12]$ be a continuous, decreasing function, and suppose that $\zeta(s)$ does not vanish when $\sigma>1-\eta(|t|)$. If we define $\omega(x) = \min_{t\ge1} (\eta(t)\log x + \log t)$, then for any $0<\ep<1$, \[ \Delta^\psi(x) \ll_\ep x/e^{(1-\ep)\omega(x)}, \] and the same is true for $\Delta^\theta(x)$ and $\Delta^\Pi(x)$ and~$\Delta^\pi(x)$. This is an improvement of a result of Ingham~\cite[Theorem 22]{1932.Ingham}, which had a factor of $\frac12$ in the exponent (and additional conditions upon~$\eta$). In particular, when combined with a 1960/61 theorem of Staś, this result provides a nearly lossless relationship between zero-free regions for $\zeta(s)$ and error terms in the prime number theorem. It follows that an $\Omega$-theorem for any of the four error terms given above actually implies $\Omega_\pm$-theorems, of the same order of magnitude up to an $\ep$ in the exponent, for all four error terms, an implication that seems extremely difficult to prove directly.\par This article cites~\cites{MR0146156, MR0146157, MR0146158, MR0146159, MR0146160, MR0146161, MR0156826, MR0156827, MR0162771, MR0174538, MR0180539, MR0182616A, MR0182616B, 1914.Littlewood, 1932.Ingham, MR0439771, MR634679}.}
      \field{issn}{0065-1036}
      \field{journaltitle}{Acta Arith.}
      \field{title}{On the remainder term of the prime number formula. {II}. {O}n a theorem of {I}ngham}
      \field{volume}{37}
      \field{year}{1980}
      \true{nocite}
      \field{pages}{209\bibrangedash 220}
      \range{pages}{12}
      \verb{doi}
      \verb 10.4064/aa-37-1-209-220
      \endverb
      \verb{mrnumber}
      \verb 598876
      \endverb
      \verb{urlraw}
      \verb https://doi.org/10.4064/aa-37-1-209-220
      \endverb
      \verb{url}
      \verb https://doi.org/10.4064/aa-37-1-209-220
      \endverb
    \endentry
    \entry{MR681441}{article}{}
      \name{author}{1}{}{%
        {{hash=c6d965df16e600c701d4bcc2888ae34c}{%
           family={Pintz},
           familyi={P\bibinitperiod},
           given={J.},
           giveni={J\bibinitperiod}}}%
      }
      \strng{namehash}{c6d965df16e600c701d4bcc2888ae34c}
      \strng{fullhash}{c6d965df16e600c701d4bcc2888ae34c}
      \strng{bibnamehash}{c6d965df16e600c701d4bcc2888ae34c}
      \strng{authorbibnamehash}{c6d965df16e600c701d4bcc2888ae34c}
      \strng{authornamehash}{c6d965df16e600c701d4bcc2888ae34c}
      \strng{authorfullhash}{c6d965df16e600c701d4bcc2888ae34c}
      \field{extraname}{7}
      \field{sortinit}{1}
      \field{sortinithash}{50c6687d7fc80f50136d75228e3c59ba}
      \field{labelnamesource}{author}
      \field{labeltitlesource}{title}
      \field{annotation}{For any of the functions $f\in\{\pi,\Pi,\theta,\psi\}$, the author establishes lower bounds for the integrated absolute error term $\fA_{|1|}^f(x)$. The main theorem of this article states that if $\beta_0+i\gamma_0$ is a zero of the Riemann zeta function, then $\fA_{|1|}^f(Y)/Y \ge Y^{\beta_0}e^{-2\sqrt{\log{Y}}(\log\log{Y})^2}$ when $Y$ is sufficiently large in terms of~$\gamma_0$. The author sketches a modification of the proof that yields the stronger lower bound $\fA_{|1|}^f(Y)/Y \ge Y^{\beta_0}e^{-18(\log{Y})^{1/3}(\log\log{Y})^{4/3}}$.\par This article cites~\cites{MR0125822, MR0043121, MR585891, MR598876, MR607089, MR630377, MR681442,MR1544509,MR1575079,MR0111722,MR0417077}.}
      \field{issn}{0081-6906}
      \field{journaltitle}{Studia Sci. Math. Hungar.}
      \field{number}{1-3}
      \field{title}{On the remainder term of the prime number formula. {V}. {E}ffective mean value theorems}
      \field{volume}{15}
      \field{year}{1980}
      \true{nocite}
      \field{pages}{215\bibrangedash 223}
      \range{pages}{9}
      \verb{mrnumber}
      \verb 681441
      \endverb
    \endentry
    \entry{MR681442}{article}{}
      \name{author}{1}{}{%
        {{hash=c6d965df16e600c701d4bcc2888ae34c}{%
           family={Pintz},
           familyi={P\bibinitperiod},
           given={J.},
           giveni={J\bibinitperiod}}}%
      }
      \strng{namehash}{c6d965df16e600c701d4bcc2888ae34c}
      \strng{fullhash}{c6d965df16e600c701d4bcc2888ae34c}
      \strng{bibnamehash}{c6d965df16e600c701d4bcc2888ae34c}
      \strng{authorbibnamehash}{c6d965df16e600c701d4bcc2888ae34c}
      \strng{authornamehash}{c6d965df16e600c701d4bcc2888ae34c}
      \strng{authorfullhash}{c6d965df16e600c701d4bcc2888ae34c}
      \field{extraname}{8}
      \field{sortinit}{1}
      \field{sortinithash}{50c6687d7fc80f50136d75228e3c59ba}
      \field{labelnamesource}{author}
      \field{labeltitlesource}{title}
      \field{annotation}{This article concerns the absolute averages $\fA_{|1|}$ of various standard error terms for prime counting functions. When $Y$ is sufficiently large (ineffectively), the author proves that \begin{align*} \fA_{|1|}^\pi(Y) &>0.62 \frac{Y^{3/2}}{\log Y}, & \fA_{|1|}^\Pi(Y) &>9 \cdot 10^{-5} \frac{Y^{3/2}}{\log Y}, \\ \fA_{|1|}^\theta(Y) &>0.62{Y^{3/2}}, & \fA_{|1|}^\psi(Y) &> 10^{-4} {Y^{3/2}}. \end{align*} Thanks to work of Cramér~\cite{MR1544509}, if RH is true then these bounds are best possible up to the leading constants (and even those constants are not too far off). Under RH, the author can improve some of these constants and also better localize the implied large values of the error terms; indeed, the lower bounds for the $\fA_{|1|}$ are derived from the existence of large oscillations of the error terms, rather than the other way around. \par This article cites~\cites{MR1544509,1936.Ingham,MR681441}.}
      \field{issn}{0081-6906}
      \field{journaltitle}{Studia Sci. Math. Hungar.}
      \field{number}{1-3}
      \field{title}{On the remainder term of the prime number formula. {VI}. {I}neffective mean value theorems}
      \field{volume}{15}
      \field{year}{1980}
      \true{nocite}
      \field{pages}{225\bibrangedash 230}
      \range{pages}{6}
      \verb{mrnumber}
      \verb 681442
      \endverb
    \endentry
    \entry{MR688630}{article}{}
      \name{author}{1}{}{%
        {{hash=c6d965df16e600c701d4bcc2888ae34c}{%
           family={Pintz},
           familyi={P\bibinitperiod},
           given={J.},
           giveni={J\bibinitperiod}}}%
      }
      \strng{namehash}{c6d965df16e600c701d4bcc2888ae34c}
      \strng{fullhash}{c6d965df16e600c701d4bcc2888ae34c}
      \strng{bibnamehash}{c6d965df16e600c701d4bcc2888ae34c}
      \strng{authorbibnamehash}{c6d965df16e600c701d4bcc2888ae34c}
      \strng{authornamehash}{c6d965df16e600c701d4bcc2888ae34c}
      \strng{authorfullhash}{c6d965df16e600c701d4bcc2888ae34c}
      \field{extraname}{9}
      \field{sortinit}{1}
      \field{sortinithash}{50c6687d7fc80f50136d75228e3c59ba}
      \field{labelnamesource}{author}
      \field{labeltitlesource}{title}
      \field{annotation}{The author shows that $M(x)$ changes sign in the interval $[Y \exp(-3 \log_2^{3/2} Y),Y]$ for all sufficiently large~$Y$, which improves the main result in~\cite{MR660714}. The proof uses upper and lower bounds for an integral of the form $\int_v^w |M(x)|x^{-1-\theta} \,dx$, where $\theta$ is the maximal real part of the zeros of $\zeta(s)$ for $|t| \lesssim \log Y$.\par This article cites~\cites{MR0207665,MR0219496,MR0155802,1909.Landau,MR630377,MR660714,MR678996,MR0087685}.}
      \field{issn}{0081-6906}
      \field{journaltitle}{Studia Sci. Math. Hungar.}
      \field{number}{4}
      \field{title}{Oscillatory properties of {$M(x)=\sum _{n\leq x}\mu (n)$}. {II}}
      \field{volume}{15}
      \field{year}{1980}
      \true{nocite}
      \field{pages}{491\bibrangedash 496}
      \range{pages}{6}
      \verb{mrnumber}
      \verb 688630
      \endverb
    \endentry
    \entry{MR584557}{article}{}
      \name{author}{1}{}{%
        {{hash=31d751abf6c11534c2c4933716b5a5c6}{%
           family={Tanaka},
           familyi={T\bibinitperiod},
           given={Minoru},
           giveni={M\bibinitperiod}}}%
      }
      \strng{namehash}{31d751abf6c11534c2c4933716b5a5c6}
      \strng{fullhash}{31d751abf6c11534c2c4933716b5a5c6}
      \strng{bibnamehash}{31d751abf6c11534c2c4933716b5a5c6}
      \strng{authorbibnamehash}{31d751abf6c11534c2c4933716b5a5c6}
      \strng{authornamehash}{31d751abf6c11534c2c4933716b5a5c6}
      \strng{authorfullhash}{31d751abf6c11534c2c4933716b5a5c6}
      \field{extraname}{1}
      \field{sortinit}{1}
      \field{sortinithash}{50c6687d7fc80f50136d75228e3c59ba}
      \field{labelnamesource}{author}
      \field{labeltitlesource}{title}
      \field{annotation}{This article reports on the sign changes of $L(x)$ for $x\leq 10^9$. In particular, the author shows that $906,150,256$ is the smallest integer $n \geq 2$ such that $L(n)>0$.\par This article cites~\cites{MR0006202,MR0120198}.}
      \field{issn}{0387-3870}
      \field{journaltitle}{Tokyo J. Math.}
      \field{number}{1}
      \field{title}{A numerical investigation on cumulative sum of the {L}iouville function}
      \field{volume}{3}
      \field{year}{1980}
      \true{nocite}
      \field{pages}{187\bibrangedash 189}
      \range{pages}{3}
      \verb{doi}
      \verb 10.3836/tjm/1270216093
      \endverb
      \verb{mrnumber}
      \verb 584557
      \endverb
      \verb{urlraw}
      \verb https://doi.org/10.3836/tjm/1270216093
      \endverb
      \verb{url}
      \verb https://doi.org/10.3836/tjm/1270216093
      \endverb
    \endentry
    \entry{MR0605090}{article}{}
      \name{author}{1}{}{%
        {{hash=31d751abf6c11534c2c4933716b5a5c6}{%
           family={Tanaka},
           familyi={T\bibinitperiod},
           given={Minoru},
           giveni={M\bibinitperiod}}}%
      }
      \strng{namehash}{31d751abf6c11534c2c4933716b5a5c6}
      \strng{fullhash}{31d751abf6c11534c2c4933716b5a5c6}
      \strng{bibnamehash}{31d751abf6c11534c2c4933716b5a5c6}
      \strng{authorbibnamehash}{31d751abf6c11534c2c4933716b5a5c6}
      \strng{authornamehash}{31d751abf6c11534c2c4933716b5a5c6}
      \strng{authorfullhash}{31d751abf6c11534c2c4933716b5a5c6}
      \field{extraname}{2}
      \field{sortinit}{1}
      \field{sortinithash}{50c6687d7fc80f50136d75228e3c59ba}
      \field{labelnamesource}{author}
      \field{labeltitlesource}{title}
      \field{annotation}{Define $B_k=\zeta\left(\frac{k}2\right) / \zeta\left(\frac{1}{2}\right) \zeta(k)$ if $k\ge3$ is odd and $B_k=1 / \zeta\left(\frac{1}{2}\right) \zeta\left(\frac{k}2\right)$ if $k\ge4$ is even. The author shows that $\sum_{n \leq x}\mu_k(x)-B_k \sqrt{x}=\Omega_{\pm}(\sqrt{x})$ for each $k\ge3$, generalizing classical oscillation results for $M(x)$ (which is the case $k=2$) and $L(x)$ (which is the case $k\to\infty$).\par This article cites~\cites{MR0068565,MR0006202}.}
      \field{issn}{0387-3870}
      \field{journaltitle}{Tokyo J. Math.}
      \field{number}{2}
      \field{title}{On the {M}öbius and allied functions}
      \field{volume}{3}
      \field{year}{1980}
      \true{nocite}
      \field{pages}{215\bibrangedash 218}
      \range{pages}{4}
      \verb{doi}
      \verb 10.3836/tjm/1270472994
      \endverb
      \verb{mrnumber}
      \verb 605090
      \endverb
      \verb{urlraw}
      \verb https://doi.org/10.3836/tjm/1270472994
      \endverb
      \verb{url}
      \verb https://doi.org/10.3836/tjm/1270472994
      \endverb
    \endentry
    \entry{MR654520}{incollection}{}
      \name{author}{2}{}{%
        {{hash=44f716ae9c52c6e9e0952988d11b7582}{%
           family={Anderson},
           familyi={A\bibinitperiod},
           given={R.\bibnamedelimi J.},
           giveni={R\bibinitperiod\bibinitdelim J\bibinitperiod}}}%
        {{hash=32597e08070c5aa88e8fb07dd295112a}{%
           family={Stark},
           familyi={S\bibinitperiod},
           given={H.\bibnamedelimi M.},
           giveni={H\bibinitperiod\bibinitdelim M\bibinitperiod}}}%
      }
      \list{publisher}{1}{%
        {Springer, Berlin-New York}%
      }
      \strng{namehash}{b69969b5d1cab3b7bcc8edd374f284cb}
      \strng{fullhash}{b69969b5d1cab3b7bcc8edd374f284cb}
      \strng{bibnamehash}{b69969b5d1cab3b7bcc8edd374f284cb}
      \strng{authorbibnamehash}{b69969b5d1cab3b7bcc8edd374f284cb}
      \strng{authornamehash}{b69969b5d1cab3b7bcc8edd374f284cb}
      \strng{authorfullhash}{b69969b5d1cab3b7bcc8edd374f284cb}
      \field{sortinit}{1}
      \field{sortinithash}{50c6687d7fc80f50136d75228e3c59ba}
      \field{labelnamesource}{author}
      \field{labeltitlesource}{title}
      \field{annotation}{The authors disprove several conjectures using a theorem of Landau and the locations of zeros of $\zeta(s)$. The authors also show that for conjectures of a certain type (such as von Sterneck's conjecture~\cites{1912.vonSterneck} $|M(x)|<\frac{1}{2}\sqrt{x}$ for $x>200$ and Pólya's problem~\cites{1919.Polya} $L(x)\leq 0$ for $x \geq 2$), a single counterexample implies infinitely many counterexamples. In particular, using the known counterexamples from Neubauer \cites{MR155787} that $E^M(x)>0.557$ for $x = 7.76 \times 10^9$, and from Lehman \cites{MR0120198} that $E^L(x)>0.023$ for $x=9.064 \times 10^8$, the authors conclude that $\limsup_{x \to \infty} E^M(x)>0.557$ and $\limsup_{x \to \infty} E^L(x)>0.023$.\par This article cites~\cites{MR0330069,MR0376566,MR0198145,MR0332685,MR0104638,MR0289452,MR0006202,MR0352026,MR0120198,MR549977,MR155787,MR280447,MR43841,MR51857}.}
      \field{bookpagination}{none}
      \field{booktitle}{Analytic number theory ({P}hiladelphia, {P}a., 1980)}
      \field{series}{Lecture Notes in Math.}
      \field{title}{Oscillation theorems}
      \field{volume}{899}
      \field{year}{1981}
      \true{nocite}
      \field{pages}{79\bibrangedash 106}
      \range{pages}{28}
      \verb{mrnumber}
      \verb 654520
      \endverb
    \endentry
    \entry{MR602236}{article}{}
      \name{author}{1}{}{%
        {{hash=97466ee083d3402c436ef85e5d255b89}{%
           family={Chen},
           familyi={C\bibinitperiod},
           given={W.\bibnamedelimi W.\bibnamedelimi L.},
           giveni={W\bibinitperiod\bibinitdelim W\bibinitperiod\bibinitdelim L\bibinitperiod}}}%
      }
      \strng{namehash}{97466ee083d3402c436ef85e5d255b89}
      \strng{fullhash}{97466ee083d3402c436ef85e5d255b89}
      \strng{bibnamehash}{97466ee083d3402c436ef85e5d255b89}
      \strng{authorbibnamehash}{97466ee083d3402c436ef85e5d255b89}
      \strng{authornamehash}{97466ee083d3402c436ef85e5d255b89}
      \strng{authorfullhash}{97466ee083d3402c436ef85e5d255b89}
      \field{sortinit}{1}
      \field{sortinithash}{50c6687d7fc80f50136d75228e3c59ba}
      \field{labelnamesource}{author}
      \field{labeltitlesource}{title}
      \field{annotation}{It was conjectured in~\cites{MR0108470} that $\sum_{n\le x} E^\pi(n) \sim x$ and $\sum_{n\le x} E^\pi(n;4,3,1) \sim x$; in this article the author disproves these conjectures. Defining the function $R(x)= \sum_{n \leq x} E^{\psi}(n)$ he shows that $R(x)=\Omega_\pm(x^{{1/2}+\Theta-\ep})$ for every $\ep>0$, which can be improved to $R(x)=\Omega_{\pm}(x^{1/2+\Theta})$ under SA. The same results hold for $P(x)-x$ in place of $R(x)$, where $P(x)= - \sum_{n \leq x} E^{\pi}(n)$. In these theorems, the author notes that $\psi(n)$ and $\pi(n)$ can also be replaced by $\psi(n;4,1,3)$ and $\pi(n;4,1,3)$.\par This article cites~\cites{MR0108470}.}
      \field{issn}{0024-6107}
      \field{journaltitle}{J. London Math. Soc. (2)}
      \field{number}{1}
      \field{title}{On the error term of the prime number theorem and the difference between the number of primes in the residue classes modulo $4$}
      \field{volume}{23}
      \field{year}{1981}
      \true{nocite}
      \field{pages}{24\bibrangedash 40}
      \range{pages}{17}
      \verb{doi}
      \verb 10.1112/jlms/s2-23.1.24
      \endverb
      \verb{mrnumber}
      \verb 602236
      \endverb
    \endentry
    \entry{MR660714}{article}{}
      \name{author}{1}{}{%
        {{hash=c6d965df16e600c701d4bcc2888ae34c}{%
           family={Pintz},
           familyi={P\bibinitperiod},
           given={J.},
           giveni={J\bibinitperiod}}}%
      }
      \strng{namehash}{c6d965df16e600c701d4bcc2888ae34c}
      \strng{fullhash}{c6d965df16e600c701d4bcc2888ae34c}
      \strng{bibnamehash}{c6d965df16e600c701d4bcc2888ae34c}
      \strng{authorbibnamehash}{c6d965df16e600c701d4bcc2888ae34c}
      \strng{authornamehash}{c6d965df16e600c701d4bcc2888ae34c}
      \strng{authorfullhash}{c6d965df16e600c701d4bcc2888ae34c}
      \field{extraname}{10}
      \field{sortinit}{1}
      \field{sortinithash}{50c6687d7fc80f50136d75228e3c59ba}
      \field{labelnamesource}{author}
      \field{labeltitlesource}{title}
      \field{annotation}{Assuming RH$(T)$, the author shows that $M(x)$ changes sign in the interval $[Y^{1-1/(T-2)}, Y^{1+1/(T-2)}]$ for every sufficiently large~$Y$. Consequently, using Brent's verification of $RH$ for the first $7 \cdot 10^7$ zeros of $\zeta(s)$, the author concludes that $M(x)$ changes sign in every interval of the form $[Y^{1-10^{-7}},Y]$ for all sufficiently large~$Y$, which improves the main result in ~\cite{MR0219496}.\par This article cites~\cites{MR0207665,MR0219496,MR0155802,MR0172856}.}
      \field{issn}{0174-4747}
      \field{journaltitle}{Analysis}
      \field{number}{3}
      \field{title}{On the sign changes of {$M(x)=\sum _{n\leq x}\mu (n)$}}
      \field{volume}{1}
      \field{year}{1981}
      \true{nocite}
      \field{pages}{191\bibrangedash 195}
      \range{pages}{5}
      \verb{doi}
      \verb 10.1524/anly.1981.1.3.191
      \endverb
      \verb{mrnumber}
      \verb 660714
      \endverb
      \verb{urlraw}
      \verb https://doi.org/10.1524/anly.1981.1.3.191
      \endverb
      \verb{url}
      \verb https://doi.org/10.1524/anly.1981.1.3.191
      \endverb
    \endentry
    \entry{MR675189}{article}{}
      \name{author}{1}{}{%
        {{hash=6bf42c5abf099333fae470886fff8616}{%
           family={Bentz},
           familyi={B\bibinitperiod},
           given={H.-J.},
           giveni={H\bibinithyphendelim J\bibinitperiod}}}%
      }
      \strng{namehash}{6bf42c5abf099333fae470886fff8616}
      \strng{fullhash}{6bf42c5abf099333fae470886fff8616}
      \strng{bibnamehash}{6bf42c5abf099333fae470886fff8616}
      \strng{authorbibnamehash}{6bf42c5abf099333fae470886fff8616}
      \strng{authornamehash}{6bf42c5abf099333fae470886fff8616}
      \strng{authorfullhash}{6bf42c5abf099333fae470886fff8616}
      \field{sortinit}{1}
      \field{sortinithash}{50c6687d7fc80f50136d75228e3c59ba}
      \field{labelnamesource}{author}
      \field{labeltitlesource}{title}
      \field{annotation}{For $0\le\alpha<\frac12$, the author shows unconditionally that \[ \sum_p \chi_{-4}(p) \frac{\log p}{p^\alpha} e^{-(\log x)^2/p} \sim -\frac{\sqrt{\pi x}}2 e^{x(1-2\alpha)^2/16}; \] when $\alpha=\frac12$, the right-hand side must be replaced by $\frac14\sqrt{\pi x}$. Both results remain valid if $\chi_{-4}$ is replaced by $\chi_{-3}$. These results can be interpreted as comparing (in a specific way) the residue class $1$ to the other reduced residue class modulo~$4$ or~$3$. Analogously, when $\alpha=\frac12$, the author establishes the same result when comparing $1\mod8$ to another reduced residue class\mod8; if two reduced residue classes\mod8 are compared, the resulting expression is bounded. The author asserts that the required hypotheses on zeros of relevant Dirichlet $L$-functions is that they do not vanish in the ``bowtie'' $\{s\colon \sigma>0,\, 0<|t|<|\sigma-\frac12|\}$. The author also presents some numerical data concerning the prime number race\mod 3.\par This article cites~\cites{MR595317, MR534235, MR552462, MR1555148, MR0460261, 1932.Ingham, MR0146156, MR0146157, MR0146158, MR0146159, MR0146160, MR0146161, MR0156826, MR0156827, MR0162771, MR0174538, MR0180539, MR0182616A, MR0182616B, MR0200250, MR0302585, MR1544293A, MR1544293B, MR0083001, MR0108470}.}
      \field{issn}{0022-314X}
      \field{journaltitle}{J. Number Theory}
      \field{number}{2}
      \field{title}{Discrepancies in the distribution of prime numbers}
      \field{volume}{15}
      \field{year}{1982}
      \true{nocite}
      \field{pages}{252\bibrangedash 274}
      \range{pages}{23}
      \verb{doi}
      \verb 10.1016/0022-314X(82)90030-0
      \endverb
      \verb{mrnumber}
      \verb 675189
      \endverb
      \verb{urlraw}
      \verb https://doi.org/10.1016/0022-314X(82)90030-0
      \endverb
      \verb{url}
      \verb https://doi.org/10.1016/0022-314X(82)90030-0
      \endverb
    \endentry
    \entry{MR662791}{article}{}
      \name{author}{2}{}{%
        {{hash=6bf42c5abf099333fae470886fff8616}{%
           family={Bentz},
           familyi={B\bibinitperiod},
           given={H.-J.},
           giveni={H\bibinithyphendelim J\bibinitperiod}}}%
        {{hash=c6d965df16e600c701d4bcc2888ae34c}{%
           family={Pintz},
           familyi={P\bibinitperiod},
           given={J.},
           giveni={J\bibinitperiod}}}%
      }
      \list{language}{1}{%
        {German}%
      }
      \strng{namehash}{8ee030157e0efba486591e59a6b5401b}
      \strng{fullhash}{8ee030157e0efba486591e59a6b5401b}
      \strng{bibnamehash}{8ee030157e0efba486591e59a6b5401b}
      \strng{authorbibnamehash}{8ee030157e0efba486591e59a6b5401b}
      \strng{authornamehash}{8ee030157e0efba486591e59a6b5401b}
      \strng{authorfullhash}{8ee030157e0efba486591e59a6b5401b}
      \field{extraname}{3}
      \field{sortinit}{1}
      \field{sortinithash}{50c6687d7fc80f50136d75228e3c59ba}
      \field{labelnamesource}{author}
      \field{labeltitlesource}{title}
      \field{annotation}{The authors give a shorter proof that $\displaystyle\lim_{x\to\infty}\sum_p \chi_{-4}(p)\log p\cdot\exp\left(-\frac{\log^2 p}{x}\right)=-\infty$.\par This article cites~\cites{MR534235,MR552462,MR1555148,MR0174538,MR1544293A,MR1544293B,MR0108470,1853.Chebyshev}.}
      \field{issn}{0378-6218}
      \field{journaltitle}{Resultate Math.}
      \field{number}{1}
      \field{title}{Über das {T}schebyschef-{P}roblem}
      \field{volume}{5}
      \field{year}{1982}
      \true{nocite}
      \field{pages}{1\bibrangedash 5}
      \range{pages}{5}
      \verb{doi}
      \verb 10.1007/bf03323296
      \endverb
      \verb{mrnumber}
      \verb 662791
      \endverb
      \verb{urlraw}
      \verb https://doi.org/10.1007/bf03323296
      \endverb
      \verb{url}
      \verb https://doi.org/10.1007/bf03323296
      \endverb
    \endentry
    \entry{MR678996}{article}{}
      \name{author}{1}{}{%
        {{hash=c6d965df16e600c701d4bcc2888ae34c}{%
           family={Pintz},
           familyi={P\bibinitperiod},
           given={J.},
           giveni={J\bibinitperiod}}}%
      }
      \strng{namehash}{c6d965df16e600c701d4bcc2888ae34c}
      \strng{fullhash}{c6d965df16e600c701d4bcc2888ae34c}
      \strng{bibnamehash}{c6d965df16e600c701d4bcc2888ae34c}
      \strng{authorbibnamehash}{c6d965df16e600c701d4bcc2888ae34c}
      \strng{authornamehash}{c6d965df16e600c701d4bcc2888ae34c}
      \strng{authorfullhash}{c6d965df16e600c701d4bcc2888ae34c}
      \field{extraname}{11}
      \field{sortinit}{1}
      \field{sortinithash}{50c6687d7fc80f50136d75228e3c59ba}
      \field{labelnamesource}{author}
      \field{labeltitlesource}{title}
      \field{annotation}{This article is concerned with oscillations in the Mertens sum. The natural difficulty of this problem comes from the fact that the explicit formula for $M(x)$ contains terms of the form $x^{\rho}/\rho\zeta'(\rho)$, which are more difficult to handle than the terms $x^{\rho}/\rho$ appearing in the explicit formula for $\Delta^\psi(x)$. The author proves that if $\rho_0=\beta_0+i\gamma_0$ is a zero of $\zeta(s)$, then for $Y> e^{|\gamma_0|+4}$, \begin{align*} \frac{1}{Y} \int_{Y/(100 \log Y)}^Y |M(x)|\, dx > \frac{1}{6|\rho_0|^3}Y^{\beta_0} \quad \text{and} \quad \max_{x \leq Y} |M(x)| \geq \frac{1}{6|\rho_0|^3}Y^{\beta_0}. \end{align*} Consequently, using the first zeta zero $\frac12 + i\gamma_1$ with $\gamma_1\approx14.13$, \begin{align*} \max_{Y/(100 \log Y) \le x \le Y} |M(x)| \geq \frac{1}{17000}\sqrt{Y} \end{align*} for $Y\ge2$; the constant $1/17000$ can be improved but not enough to disprove the Mertens conjecture.\par This article cites~\cites{MR0218318, MR0043121}.}
      \field{issn}{0065-1036}
      \field{journaltitle}{Acta Arith.}
      \field{number}{1}
      \field{title}{Oscillatory properties of {$M(x)=\sum _{n\leq x}\mu (n)$}. {I}}
      \field{volume}{42}
      \field{year}{1982}
      \true{nocite}
      \field{pages}{49\bibrangedash 55}
      \range{pages}{7}
      \verb{doi}
      \verb 10.4064/aa-42-1-49-55
      \endverb
      \verb{mrnumber}
      \verb 678996
      \endverb
      \verb{urlraw}
      \verb https://doi.org/10.4064/aa-42-1-49-55
      \endverb
      \verb{url}
      \verb https://doi.org/10.4064/aa-42-1-49-55
      \endverb
    \endentry
    \entry{MR817334}{article}{}
      \name{author}{1}{}{%
        {{hash=b39a346ae2b44071e01b12c1003e1246}{%
           family={Staś},
           familyi={S\bibinitperiod},
           given={Włodzimierz},
           giveni={W\bibinitperiod}}}%
      }
      \strng{namehash}{b39a346ae2b44071e01b12c1003e1246}
      \strng{fullhash}{b39a346ae2b44071e01b12c1003e1246}
      \strng{bibnamehash}{b39a346ae2b44071e01b12c1003e1246}
      \strng{authorbibnamehash}{b39a346ae2b44071e01b12c1003e1246}
      \strng{authornamehash}{b39a346ae2b44071e01b12c1003e1246}
      \strng{authorfullhash}{b39a346ae2b44071e01b12c1003e1246}
      \field{extraname}{3}
      \field{sortinit}{1}
      \field{sortinithash}{50c6687d7fc80f50136d75228e3c59ba}
      \field{labelnamesource}{author}
      \field{labeltitlesource}{title}
      \field{annotation}{In this article, the author proves an analogue of a result of Knapowski~\cite{MR0133309} for the error term of the prime ideal theorem. More precisely he shows that for a number field~$K$ whose Dedekind zeta function $\zeta_K(s)$ satisfies HC, and an arbitrary zero $\rho_0 = \beta_0 + i\gamma_0$ of $\zeta_K$, \[ \Delta^\psi(T,K) = \Omega_\pm\biggl( T^{\beta_0}\exp\biggl( -15\frac{\log T}{\sqrt{\log \log T}}\biggr)\biggr). \] The author also gives a lower bound on the number of sign changes of $\Delta^\psi(T,K)$.\par This article cites~\cites{MR0133309,MR0174538, MR585891,MR0060548,MR0132728}.}
      \field{issn}{0208-6573}
      \field{journaltitle}{Funct. Approx. Comment. Math.}
      \field{title}{On sign-changes in the remainder term of the prime ideal formula}
      \field{volume}{13}
      \field{year}{1982}
      \true{nocite}
      \field{pages}{159\bibrangedash 166}
      \range{pages}{8}
      \verb{mrnumber}
      \verb 817334
      \endverb
    \endentry
    \entry{MR820241}{incollection}{}
      \name{author}{2}{}{%
        {{hash=c5f23f7a8c75fe8b8720ab22fc314bda}{%
           family={Kolesnik},
           familyi={K\bibinitperiod},
           given={G.},
           giveni={G\bibinitperiod}}}%
        {{hash=b79ccf3603e8ce64e49a0d4684120097}{%
           family={Straus},
           familyi={S\bibinitperiod},
           given={E.\bibnamedelimi G.},
           giveni={E\bibinitperiod\bibinitdelim G\bibinitperiod}}}%
      }
      \list{publisher}{1}{%
        {Birkhäuser, Basel}%
      }
      \strng{namehash}{db6a61d464b7f706be1353a113c7a3d2}
      \strng{fullhash}{db6a61d464b7f706be1353a113c7a3d2}
      \strng{bibnamehash}{db6a61d464b7f706be1353a113c7a3d2}
      \strng{authorbibnamehash}{db6a61d464b7f706be1353a113c7a3d2}
      \strng{authornamehash}{db6a61d464b7f706be1353a113c7a3d2}
      \strng{authorfullhash}{db6a61d464b7f706be1353a113c7a3d2}
      \field{sortinit}{1}
      \field{sortinithash}{50c6687d7fc80f50136d75228e3c59ba}
      \field{labelnamesource}{author}
      \field{labeltitlesource}{title}
      \field{annotation}{The authors examine the second main theorem of the power-sum method, obtaining the constant $B=4e$ which is best possible by~\cites{MR0159791}. When $m\leq n$ they also give the larger lower bound \[ |s_v| \geq \frac{n!(2m+n)!}{2^n(2m+2n)! \sqrt{2m+2n+1}} \min_{1 \leq j \leq n}|b_1+ \cdots + b_j|. \]\par This article cites~\cites{MR0077579,MR0159791}.}
      \field{bookpagination}{none}
      \field{booktitle}{Studies in pure mathematics}
      \field{title}{On the sum of powers of complex numbers}
      \field{year}{1983}
      \true{nocite}
      \field{pages}{427\bibrangedash 442}
      \range{pages}{16}
      \verb{mrnumber}
      \verb 820241
      \endverb
    \endentry
    \entry{MR0724708}{article}{}
      \name{author}{1}{}{%
        {{hash=c6d965df16e600c701d4bcc2888ae34c}{%
           family={Pintz},
           familyi={P\bibinitperiod},
           given={J.},
           giveni={J\bibinitperiod}}}%
      }
      \strng{namehash}{c6d965df16e600c701d4bcc2888ae34c}
      \strng{fullhash}{c6d965df16e600c701d4bcc2888ae34c}
      \strng{bibnamehash}{c6d965df16e600c701d4bcc2888ae34c}
      \strng{authorbibnamehash}{c6d965df16e600c701d4bcc2888ae34c}
      \strng{authornamehash}{c6d965df16e600c701d4bcc2888ae34c}
      \strng{authorfullhash}{c6d965df16e600c701d4bcc2888ae34c}
      \field{extraname}{12}
      \field{sortinit}{1}
      \field{sortinithash}{50c6687d7fc80f50136d75228e3c59ba}
      \field{labelnamesource}{author}
      \field{labeltitlesource}{title}
      \field{annotation}{The author shows $\fA_{|1|}^{Q_k}(x) \gg Y^{1/2k}$ for each $k \geq 2$ using an unexpectedly simple argument, with effective estimates whose dependence on~$k$ is explicit. These results significantly strengthen work of Evelyn--Linfoot~\cites{MR1502996} and Kátai~\cite{MR0219496}.\par This article cites~\cites{MR1502996, MR0219496}.}
      \field{issn}{0024-6107,1469-7750}
      \field{journaltitle}{J. London Math. Soc. (2)}
      \field{number}{3}
      \field{title}{On the distribution of square-free numbers}
      \field{volume}{28}
      \field{year}{1983}
      \true{nocite}
      \field{pages}{401\bibrangedash 405}
      \range{pages}{5}
      \verb{doi}
      \verb 10.1112/jlms/s2-28.3.401
      \endverb
      \verb{mrnumber}
      \verb 724708
      \endverb
      \verb{urlraw}
      \verb https://doi.org/10.1112/jlms/s2-28.3.401
      \endverb
      \verb{url}
      \verb https://doi.org/10.1112/jlms/s2-28.3.401
      \endverb
    \endentry
    \entry{MR820251}{incollection}{}
      \name{author}{1}{}{%
        {{hash=c6d965df16e600c701d4bcc2888ae34c}{%
           family={Pintz},
           familyi={P\bibinitperiod},
           given={J.},
           giveni={J\bibinitperiod}}}%
      }
      \list{publisher}{1}{%
        {Birkhäuser, Basel}%
      }
      \strng{namehash}{c6d965df16e600c701d4bcc2888ae34c}
      \strng{fullhash}{c6d965df16e600c701d4bcc2888ae34c}
      \strng{bibnamehash}{c6d965df16e600c701d4bcc2888ae34c}
      \strng{authorbibnamehash}{c6d965df16e600c701d4bcc2888ae34c}
      \strng{authornamehash}{c6d965df16e600c701d4bcc2888ae34c}
      \strng{authorfullhash}{c6d965df16e600c701d4bcc2888ae34c}
      \field{extraname}{13}
      \field{sortinit}{1}
      \field{sortinithash}{50c6687d7fc80f50136d75228e3c59ba}
      \field{labelnamesource}{author}
      \field{labeltitlesource}{title}
      \field{annotation}{The author establishes two theorems that improve and simplify prior work of Turán and of Ingham. The first theorem states that if $\zeta(\beta_0 + i \gamma_0)=0$, then for $T$ sufficiently large in terms of $\gamma_0$, there exists an $x \in [T^{1/4}, T]$ for which $|\Delta^\psi(x)| \gg_{\gamma_0} x^{\beta_1}$ (with an explicit dependence on~$\gamma_0$). The second theorem assumes that $\zeta(s)\ne0$ in a region of the shape $\sigma \ge 1-\eta(t)$ where $\eta$ is continuous and decreasing, and, defining $\omega(x) = \min_{t\ge0} (\eta(t)\log x+\log t)$, concludes that $\Delta^\psi(x) = \Omega(x/e^{54\omega(x)})$. The main tool is the power-sum estimate of Sós and Turán. Near-optimal improvements (by the author) of these two results appeared slightly earlier~\cite{MR585891, MR598876}.\par This article cites~\cites{1891.Phragmen, MR585891, MR598876,1932.Ingham,MR0043121,MR0049219,MR0077579}.}
      \field{booktitle}{Studies in pure mathematics}
      \field{title}{Oscillatory properties of the remainder term of the prime number formula}
      \field{year}{1983}
      \true{nocite}
      \field{pages}{551\bibrangedash 560}
      \range{pages}{10}
      \verb{mrnumber}
      \verb 820251
      \endverb
    \endentry
    \entry{MR729173}{incollection}{}
      \name{author}{1}{}{%
        {{hash=4440c3774c4e033603367ee77d77249f}{%
           family={Robin},
           familyi={R\bibinitperiod},
           given={G.},
           giveni={G\bibinitperiod}}}%
      }
      \list{language}{1}{%
        {French}%
      }
      \list{publisher}{1}{%
        {Birkhäuser Boston, Boston, MA}%
      }
      \strng{namehash}{4440c3774c4e033603367ee77d77249f}
      \strng{fullhash}{4440c3774c4e033603367ee77d77249f}
      \strng{bibnamehash}{4440c3774c4e033603367ee77d77249f}
      \strng{authorbibnamehash}{4440c3774c4e033603367ee77d77249f}
      \strng{authornamehash}{4440c3774c4e033603367ee77d77249f}
      \strng{authorfullhash}{4440c3774c4e033603367ee77d77249f}
      \field{extraname}{1}
      \field{sortinit}{1}
      \field{sortinithash}{50c6687d7fc80f50136d75228e3c59ba}
      \field{labelnamesource}{author}
      \field{labeltitlesource}{title}
      \field{annotation}{The author investigates the function $\frac{\sigma(n)}n$ and proves that RH is true if and only if $\frac{\sigma(n)}n <e^{C_0}\log\log{n}$ when~$n$ is sufficiently large. If~$(C_k)$ is the sequence of colossally abundant numbers, the author further shows that $(\sigma(C_k)/C_k\log\log C_k)$ has infinitely many local extrema. Finally, the author shows that if RH is false, then both $\Delta^{\pi_r}(x)$ and $\prod_{p\leq x} (1-\frac1p)-1/e^{C_0}\log{x}$ change signs infinitely often.\par This article cites~\cites{MR0417077,MR0202685,MR0137689}}
      \field{bookpagination}{none}
      \field{booktitle}{Seminar on number theory, {P}aris 1981--82 ({P}aris, 1981/1982)}
      \field{series}{Progr. Math.}
      \field{title}{Sur l'ordre maximum de la fonction somme des diviseurs}
      \field{volume}{38}
      \field{year}{1983}
      \true{nocite}
      \field{pages}{233\bibrangedash 244}
      \range{pages}{12}
      \verb{mrnumber}
      \verb 729173
      \endverb
    \endentry
    \entry{MR777013}{article}{}
      \name{author}{1}{}{%
        {{hash=d1e4ad71dae945a6977ecf6f9ca65c30}{%
           family={Kaczorowski},
           familyi={K\bibinitperiod},
           given={J.},
           giveni={J\bibinitperiod}}}%
      }
      \strng{namehash}{d1e4ad71dae945a6977ecf6f9ca65c30}
      \strng{fullhash}{d1e4ad71dae945a6977ecf6f9ca65c30}
      \strng{bibnamehash}{d1e4ad71dae945a6977ecf6f9ca65c30}
      \strng{authorbibnamehash}{d1e4ad71dae945a6977ecf6f9ca65c30}
      \strng{authornamehash}{d1e4ad71dae945a6977ecf6f9ca65c30}
      \strng{authorfullhash}{d1e4ad71dae945a6977ecf6f9ca65c30}
      \field{extraname}{1}
      \field{sortinit}{1}
      \field{sortinithash}{50c6687d7fc80f50136d75228e3c59ba}
      \field{labelnamesource}{author}
      \field{labeltitlesource}{title}
      \field{annotation}{This article establishes a lower bound on the growth of the number of sign changes of $\Delta^\psi(x)$ and $\Delta^\Pi(x)$. Specifically, the author proves that $W^\psi(T) \ge \frac{\gamma_1}{4\pi} \log T$ (and the same for $W^\Pi(x)$) when $T$ is sufficiently large (effectively), where $\gamma_1 \approx 14.1347$. The key technique is to bound $W^\psi(T)$ below by $W\big( A_n^\psi;T \big)$, the number of sign changes of repeated logarithmic integrals of $\Delta^\psi(x)$; using the fact that the second-lowest nontrivial zero of $\zeta(s)$ has imaginary part exceeding $15$, the author derives an explicit formula for $W\big( A_n^\psi;T \big)$ when $n\asymp \log T$ is suitably chosen.\par This article cites~\cites{MR1555148, 1936.Ingham, 1932.Ingham, MR0133309, MR0133308, MR0439771, MR0439772, MR0202686, MR0439774, 1914.Littlewood, MR0439773, MR607089, MR630377, 1930.Polya, MR1511206, MR1573970, MR0067145}.}
      \field{issn}{0065-1036}
      \field{journaltitle}{Acta Arith.}
      \field{number}{4}
      \field{title}{On sign-changes in the remainder-term of the prime-number formula. {I}}
      \field{volume}{44}
      \field{year}{1984}
      \true{nocite}
      \field{pages}{365\bibrangedash 377}
      \range{pages}{13}
      \verb{doi}
      \verb 10.4064/aa-44-4-365-377
      \endverb
      \verb{mrnumber}
      \verb 777013
      \endverb
      \verb{urlraw}
      \verb https://doi.org/10.4064/aa-44-4-365-377
      \endverb
      \verb{url}
      \verb https://doi.org/10.4064/aa-44-4-365-377
      \endverb
    \endentry
    \entry{MR781183}{incollection}{}
      \name{author}{1}{}{%
        {{hash=c6d965df16e600c701d4bcc2888ae34c}{%
           family={Pintz},
           familyi={P\bibinitperiod},
           given={J.},
           giveni={J\bibinitperiod}}}%
      }
      \list{publisher}{1}{%
        {North-Holland, Amsterdam}%
      }
      \strng{namehash}{c6d965df16e600c701d4bcc2888ae34c}
      \strng{fullhash}{c6d965df16e600c701d4bcc2888ae34c}
      \strng{bibnamehash}{c6d965df16e600c701d4bcc2888ae34c}
      \strng{authorbibnamehash}{c6d965df16e600c701d4bcc2888ae34c}
      \strng{authornamehash}{c6d965df16e600c701d4bcc2888ae34c}
      \strng{authorfullhash}{c6d965df16e600c701d4bcc2888ae34c}
      \field{extraname}{14}
      \field{sortinit}{1}
      \field{sortinithash}{50c6687d7fc80f50136d75228e3c59ba}
      \field{labelnamesource}{author}
      \field{labeltitlesource}{title}
      \field{annotation}{The author investigates sign changes of $M(x)$ using lower bounds for $\fA^M_{|1|}(Y)$. He shows that if $\zeta(\rho_0)=0$, then $\fA^M_{|1|}(Y) > Y^{1+\beta_0}/6|\rho_0|^3$ for $Y>e^{|\gamma_0|+4}$, an effective improvement of an ineffective inequality of Kátai that provides the first effective disproof for the Mertens conjecture when RH is false. The author further concludes that $M(x)$ changes sign in every interval of the form $\bigl[ Y \exp(-3 \log _2 ^{3/2} Y), Y \bigr]$ when~$Y$ is sufficiently large.\par This article cites~\cites{MR0006202, MR0231801, MR0241374, MR0096630, MR0133287, MR0142500, 1909.Landau, 1897.Mertens, 1905.Stieltjes}.}
      \field{bookpagination}{none}
      \field{booktitle}{Topics in classical number theory, {V}ol. {I}, {II} ({B}udapest, 1981)}
      \field{series}{Colloq. Math. Soc. János Bolyai}
      \field{title}{On the partial sums of the {M}öbius function}
      \field{volume}{34}
      \field{year}{1984}
      \true{nocite}
      \field{pages}{1229\bibrangedash 1250}
      \range{pages}{22}
      \verb{mrnumber}
      \verb 781183
      \endverb
    \endentry
    \entry{MR756094}{incollection}{}
      \name{author}{1}{}{%
        {{hash=c6d965df16e600c701d4bcc2888ae34c}{%
           family={Pintz},
           familyi={P\bibinitperiod},
           given={J.},
           giveni={J\bibinitperiod}}}%
      }
      \list{publisher}{1}{%
        {Springer, Berlin}%
      }
      \strng{namehash}{c6d965df16e600c701d4bcc2888ae34c}
      \strng{fullhash}{c6d965df16e600c701d4bcc2888ae34c}
      \strng{bibnamehash}{c6d965df16e600c701d4bcc2888ae34c}
      \strng{authorbibnamehash}{c6d965df16e600c701d4bcc2888ae34c}
      \strng{authornamehash}{c6d965df16e600c701d4bcc2888ae34c}
      \strng{authorfullhash}{c6d965df16e600c701d4bcc2888ae34c}
      \field{extraname}{15}
      \field{sortinit}{1}
      \field{sortinithash}{50c6687d7fc80f50136d75228e3c59ba}
      \field{labelnamesource}{author}
      \field{labeltitlesource}{title}
      \field{annotation}{This article is primarily a summary of the results to be proved in the series \cite{MR585891,MR598876,MR688630,MR681441,MR681442, MR660714} by the author. The main functions of interest are $S(x) = \max_{0 \leq u \leq x} |\Delta^\psi(u)|$ and $\fA^\psi_{|1|}(x) = \int_0^x |\Delta^\psi(u)| \, du$. The following theorem is proved: Define $\omega(x) =\log \frac{x}{Z(x)}$, where $Z(x) = \max_{\rho} \frac{x^\beta}{|\gamma|}$. Then \[ \log \frac{x}{S(x)} \sim \log \frac{x^2}{\fA^\psi_{|1|}(x)} \sim \omega(x).\] In particular, this implies that $S(x)$ and $\frac1x\fA^\psi_{|1|}(x)$ are close in value, that is, the mean and maximum of $|\Delta^\psi(u)|$ are close. The proof uses a zero-density theorem of Carlson (for the upper bounds) and the power-sum method (for the lower bounds).\par This article cites~\cites{1932.Ingham,MR777013,MR0133309,MR0133308,MR0439771,MR0439772,1914.Littlewood,MR585891,MR598876,MR688630,MR681441,MR681442, MR660714,MR0067145,MR0111722,MR0146153}.}
      \field{booktitle}{Number theory, {N}oordwijkerhout 1983 ({N}oordwijkerhout, 1983)}
      \field{series}{Lecture Notes in Math.}
      \field{title}{On the remainder term of the prime number formula and the zeros of {R}iemann's zeta-function}
      \field{volume}{1068}
      \field{year}{1984}
      \true{nocite}
      \field{pages}{186\bibrangedash 197}
      \range{pages}{12}
      \verb{mrnumber}
      \verb 756094
      \endverb
    \endentry
    \entry{MR736725}{article}{}
      \name{author}{1}{}{%
        {{hash=c6d965df16e600c701d4bcc2888ae34c}{%
           family={Pintz},
           familyi={P\bibinitperiod},
           given={J.},
           giveni={J\bibinitperiod}}}%
      }
      \strng{namehash}{c6d965df16e600c701d4bcc2888ae34c}
      \strng{fullhash}{c6d965df16e600c701d4bcc2888ae34c}
      \strng{bibnamehash}{c6d965df16e600c701d4bcc2888ae34c}
      \strng{authorbibnamehash}{c6d965df16e600c701d4bcc2888ae34c}
      \strng{authornamehash}{c6d965df16e600c701d4bcc2888ae34c}
      \strng{authorfullhash}{c6d965df16e600c701d4bcc2888ae34c}
      \field{extraname}{16}
      \field{sortinit}{1}
      \field{sortinithash}{50c6687d7fc80f50136d75228e3c59ba}
      \field{labelnamesource}{author}
      \field{labeltitlesource}{title}
      \field{annotation}{By refining the proof method in his previous work~\cite{MR688630}, the author proves that if $\rho_0 = \beta_0 + i\gamma_0$ is a zero of $\zeta(s)$, then when $Y$ is sufficiently large in terms of~$\gamma_0$, \[ \max_{x\in [Y\exp( -5(\log\log Y)^{5/2} ),Y]} \frac{M(x)}{x^{\beta_0}} > \frac1{48|\rho_0|^3}, \quad \min_{x\in [Y\exp( -5(\log\log Y)^{5/2} ),Y]} \frac{M(x)}{x^{\beta_0}} < -\frac1{48|\rho_0|^3}. \] \par This article cites~\cites{MR0207665,MR630377,MR678996,MR688630,MR0155802,MR0231801,MR0219496}.}
      \field{issn}{0065-1036}
      \field{journaltitle}{Acta Arith.}
      \field{number}{2}
      \field{title}{Oscillatory properties of {$M(x)=\sum _{n\leq x}\mu (n)$}. {III}}
      \field{volume}{43}
      \field{year}{1984}
      \true{nocite}
      \field{pages}{105\bibrangedash 113}
      \range{pages}{9}
      \verb{doi}
      \verb 10.4064/aa-43-2-105-113
      \endverb
      \verb{mrnumber}
      \verb 736725
      \endverb
      \verb{urlraw}
      \verb https://doi.org/10.4064/aa-43-2-105-113
      \endverb
      \verb{url}
      \verb https://doi.org/10.4064/aa-43-2-105-113
      \endverb
    \endentry
    \entry{MR802311}{article}{}
      \name{author}{2}{}{%
        {{hash=c6d965df16e600c701d4bcc2888ae34c}{%
           family={Pintz},
           familyi={P\bibinitperiod},
           given={J.},
           giveni={J\bibinitperiod}}}%
        {{hash=c420715a04171fd2f8ed0fee9cd57cc0}{%
           family={Salerno},
           familyi={S\bibinitperiod},
           given={S.},
           giveni={S\bibinitperiod}}}%
      }
      \strng{namehash}{587c0ae365c9a78befc74bd5d1ffd220}
      \strng{fullhash}{587c0ae365c9a78befc74bd5d1ffd220}
      \strng{bibnamehash}{587c0ae365c9a78befc74bd5d1ffd220}
      \strng{authorbibnamehash}{587c0ae365c9a78befc74bd5d1ffd220}
      \strng{authornamehash}{587c0ae365c9a78befc74bd5d1ffd220}
      \strng{authorfullhash}{587c0ae365c9a78befc74bd5d1ffd220}
      \field{extraname}{1}
      \field{sortinit}{1}
      \field{sortinithash}{50c6687d7fc80f50136d75228e3c59ba}
      \field{labelnamesource}{author}
      \field{labeltitlesource}{title}
      \field{annotation}{The authors elaborate on their work in~\cite{MR756697} to handle prime number races where a bias is present. Again assuming a finite Riemann--Piltz conjecture, they show that when $Y$ is sufficiently large, \begin{align*} \frac{1}{Y}\int_{Y^{1-7/\lambda}}^{Y} |\pi(x;q,\ell_1,\ell_2)| \, dx \ge \sqrt Y \exp\left(-\frac{9 \log Y}{\lambda} - c_3 q \lambda (\log\log Y)^2\right), \end{align*} (and the same for $\theta$ in place of $\pi$) for any $\lambda$ satisfying \begin{align*} \frac{\sqrt{\log Y}}{\sqrt q \log \log Y} < \lambda < \frac{c_2 \log Y}{q \left(\log \log Y\right)^2}. \end{align*} They first deal with the case when both $\ell_1$ and $\ell_2$ are quadratic residues\mod q, using an explicit formula involving zeros of both $L(s,\chi)$ and $L(2s,\chi)$. In the remaining case when $\ell_1$ is a residue and $\ell_2$ is a nonresidue, there is an additional term corresponding to the pole of $L(2s, \chi_0)$ at $s=\frac12$.\par This article cites~\cites{MR756697,MR0142521,MR0142520,MR0125822}.}
      \field{issn}{0003-889X}
      \field{journaltitle}{Arch. Math. (Basel)}
      \field{number}{4}
      \field{title}{Irregularities in the distribution of primes in arithmetic progressions. {II}}
      \field{volume}{43}
      \field{year}{1984}
      \true{nocite}
      \field{pages}{351\bibrangedash 357}
      \range{pages}{7}
      \verb{doi}
      \verb 10.1007/BF01196659
      \endverb
      \verb{mrnumber}
      \verb 802311
      \endverb
      \verb{urlraw}
      \verb https://doi.org/10.1007/BF01196659
      \endverb
      \verb{url}
      \verb https://doi.org/10.1007/BF01196659
      \endverb
    \endentry
    \entry{MR764945}{article}{}
      \name{author}{2}{}{%
        {{hash=c6d965df16e600c701d4bcc2888ae34c}{%
           family={Pintz},
           familyi={P\bibinitperiod},
           given={J.},
           giveni={J\bibinitperiod}}}%
        {{hash=c420715a04171fd2f8ed0fee9cd57cc0}{%
           family={Salerno},
           familyi={S\bibinitperiod},
           given={S.},
           giveni={S\bibinitperiod}}}%
      }
      \strng{namehash}{587c0ae365c9a78befc74bd5d1ffd220}
      \strng{fullhash}{587c0ae365c9a78befc74bd5d1ffd220}
      \strng{bibnamehash}{587c0ae365c9a78befc74bd5d1ffd220}
      \strng{authorbibnamehash}{587c0ae365c9a78befc74bd5d1ffd220}
      \strng{authornamehash}{587c0ae365c9a78befc74bd5d1ffd220}
      \strng{authorfullhash}{587c0ae365c9a78befc74bd5d1ffd220}
      \field{extraname}{2}
      \field{sortinit}{1}
      \field{sortinithash}{50c6687d7fc80f50136d75228e3c59ba}
      \field{labelnamesource}{author}
      \field{labeltitlesource}{title}
      \field{annotation}{In this article, the authors obtain new estimates on $\psi(x;q,\ell_1,\ell_2)$ for arbitrary residues $\ell_1,\ell_2$. Assuming GRH$(cq^2 \log^6 q,E(q))$, the authors prove that for $Y$ sufficiently large, there exists \[ x \in \left[ Y\exp\left( - \frac{cq}{\sqrt{E(q)}}(\log Y)^{1/2} (\log \log Y)^{3/2} \right), Y \right]\] such that \[ \psi(x;q,\ell_1,\ell_2) > \sqrt{Y}\exp\left( - \frac{cq}{\sqrt{E(q)}}(\log Y)^{1/2} (\log \log Y)^{3/2} \right). \] This is an improvement over the work of Knapowski and Turán both in the localization and in the lower bound. The authors improve the power-sum bounds used by Knapowski and Turán to prove their results.\par This article cites~\cites{1909.Landau,MR0198145,MR0146156,MR0146157,MR0146158,MR0146159,MR0146160,MR0146161,MR0156826,MR0156827}.}
      \field{issn}{0391-173X}
      \field{journaltitle}{Ann. Scuola Norm. Sup. Pisa Cl. Sci. (4)}
      \field{number}{2}
      \field{title}{On the comparative theory of primes}
      \field{volume}{11}
      \field{year}{1984}
      \true{nocite}
      \field{pages}{245\bibrangedash 260}
      \range{pages}{16}
      \verb{mrnumber}
      \verb 764945
      \endverb
      \verb{urlraw}
      \verb http://www.numdam.org/item?id=ASNSP_1984_4_11_2_245_0
      \endverb
      \verb{url}
      \verb http://www.numdam.org/item?id=ASNSP_1984_4_11_2_245_0
      \endverb
    \endentry
    \entry{MR756697}{article}{}
      \name{author}{2}{}{%
        {{hash=d192c3b49f427e2494b8aaea251e390b}{%
           family={Pintz},
           familyi={P\bibinitperiod},
           given={János},
           giveni={J\bibinitperiod}}}%
        {{hash=d91f53901946e1db0f20eee95cc2a03e}{%
           family={Salerno},
           familyi={S\bibinitperiod},
           given={Saverio},
           giveni={S\bibinitperiod}}}%
      }
      \strng{namehash}{13a7448d38e929fce6e0619ccab09b13}
      \strng{fullhash}{13a7448d38e929fce6e0619ccab09b13}
      \strng{bibnamehash}{13a7448d38e929fce6e0619ccab09b13}
      \strng{authorbibnamehash}{13a7448d38e929fce6e0619ccab09b13}
      \strng{authornamehash}{13a7448d38e929fce6e0619ccab09b13}
      \strng{authorfullhash}{13a7448d38e929fce6e0619ccab09b13}
      \field{extraname}{3}
      \field{sortinit}{1}
      \field{sortinithash}{50c6687d7fc80f50136d75228e3c59ba}
      \field{labelnamesource}{author}
      \field{labeltitlesource}{title}
      \field{annotation}{Assuming a finite Riemann--Piltz conjecture, the authors show that when $Y$ is sufficiently large, \begin{align*} \int_{Y^{1-7/\lambda}}^{Y} \psi(x;q,\ell_1,\ell_2) \,\frac{dx}x \gg \sqrt Y \exp\left(-\frac{2 \log Y}{\lambda} - c_3 q \lambda \log^2 Y\right) \end{align*} (and the same for $\Pi$ in place of $\psi$) for any $\lambda$ satisfying \begin{align*} \frac{\sqrt{\log Y}}{\sqrt q \log \log Y} < \lambda < \frac{c_2 \log Y}{q (\log \log Y)^2}. \end{align*} Essentially any such choice of $\lambda$ improves upon analogous results of Knapowski~\cite{MR0125822, MR0142520, MR0142521}. The proof also works for $\psi$ replaced by $\theta$ or $\pi$, but only if $\ell_1$ and $\ell_2$ are both quadratic nonresidues\mod q.\par This article cites~\cites{MR0182616A,MR819064,MR0142521,MR0142520,MR0125822}.}
      \field{issn}{0003-889X}
      \field{journaltitle}{Arch. Math. (Basel)}
      \field{number}{5}
      \field{title}{Irregularities in the distribution of primes in arithmetic progressions. {I}}
      \field{volume}{42}
      \field{year}{1984}
      \true{nocite}
      \field{pages}{439\bibrangedash 447}
      \range{pages}{9}
      \verb{doi}
      \verb 10.1007/BF01190694
      \endverb
      \verb{mrnumber}
      \verb 756697
      \endverb
      \verb{urlraw}
      \verb https://doi.org/10.1007/BF01190694
      \endverb
      \verb{url}
      \verb https://doi.org/10.1007/BF01190694
      \endverb
    \endentry
    \entry{MR749389}{book}{}
      \name{author}{1}{}{%
        {{hash=f4ee93efcace35c40df58fa6a8eb26ba}{%
           family={Turán},
           familyi={T\bibinitperiod},
           given={Paul},
           giveni={P\bibinitperiod}}}%
      }
      \list{publisher}{1}{%
        {John Wiley \& Sons, Inc., New York}%
      }
      \strng{namehash}{f4ee93efcace35c40df58fa6a8eb26ba}
      \strng{fullhash}{f4ee93efcace35c40df58fa6a8eb26ba}
      \strng{bibnamehash}{f4ee93efcace35c40df58fa6a8eb26ba}
      \strng{authorbibnamehash}{f4ee93efcace35c40df58fa6a8eb26ba}
      \strng{authornamehash}{f4ee93efcace35c40df58fa6a8eb26ba}
      \strng{authorfullhash}{f4ee93efcace35c40df58fa6a8eb26ba}
      \field{extraname}{9}
      \field{sortinit}{1}
      \field{sortinithash}{50c6687d7fc80f50136d75228e3c59ba}
      \field{labelnamesource}{author}
      \field{labeltitlesource}{title}
      \field{annotation}{This is a comprehensive work on the power-sum method, completed with the assistance of G.~Halász and J.~Pintz and including a foreword by V.~T.~Sós. The first part of the book introduces the method, while the second and larger part discusses a wide range of applications. Chapters 48--56 are directly related to comparative prime number theory.}
      \field{isbn}{0-471-89255-6}
      \field{series}{Pure and Applied Mathematics (New York)}
      \field{title}{On a new method of analysis and its applications}
      \field{year}{1984}
      \true{nocite}
      \field{pages}{xvi+584}
      \range{pages}{-1}
      \verb{mrnumber}
      \verb 749389
      \endverb
    \endentry
    \entry{MR814687}{book}{}
      \name{author}{2}{}{%
        {{hash=5ddf9a712233e9bcfe9ad3ab3585dcf4}{%
           family={Ellison},
           familyi={E\bibinitperiod},
           given={William},
           giveni={W\bibinitperiod}}}%
        {{hash=641b110292dfa35db45ceb8c7790ba2d}{%
           family={Ellison},
           familyi={E\bibinitperiod},
           given={Fern},
           giveni={F\bibinitperiod}}}%
      }
      \list{publisher}{1}{%
        {John Wiley \& Sons, Inc., New York; Hermann, Paris}%
      }
      \strng{namehash}{2f130d3d62e0e56f43edeac5f8911bca}
      \strng{fullhash}{2f130d3d62e0e56f43edeac5f8911bca}
      \strng{bibnamehash}{2f130d3d62e0e56f43edeac5f8911bca}
      \strng{authorbibnamehash}{2f130d3d62e0e56f43edeac5f8911bca}
      \strng{authornamehash}{2f130d3d62e0e56f43edeac5f8911bca}
      \strng{authorfullhash}{2f130d3d62e0e56f43edeac5f8911bca}
      \field{sortinit}{1}
      \field{sortinithash}{50c6687d7fc80f50136d75228e3c59ba}
      \field{labelnamesource}{author}
      \field{labeltitlesource}{title}
      \field{annotation}{This book is an English translation of~\cite{MR0417077}.}
      \field{isbn}{0-471-82653-7}
      \field{series}{A Wiley-Interscience Publication}
      \field{title}{Prime numbers}
      \field{year}{1985}
      \true{nocite}
      \field{pages}{xii+417}
      \range{pages}{-1}
      \verb{mrnumber}
      \verb 814687
      \endverb
    \endentry
    \entry{MR777286}{article}{}
      \name{author}{1}{}{%
        {{hash=9e0609c9931c1b620df97e8c58d07532}{%
           family={Hudson},
           familyi={H\bibinitperiod},
           given={Richard\bibnamedelima H.},
           giveni={R\bibinitperiod\bibinitdelim H\bibinitperiod}}}%
      }
      \strng{namehash}{9e0609c9931c1b620df97e8c58d07532}
      \strng{fullhash}{9e0609c9931c1b620df97e8c58d07532}
      \strng{bibnamehash}{9e0609c9931c1b620df97e8c58d07532}
      \strng{authorbibnamehash}{9e0609c9931c1b620df97e8c58d07532}
      \strng{authornamehash}{9e0609c9931c1b620df97e8c58d07532}
      \strng{authorfullhash}{9e0609c9931c1b620df97e8c58d07532}
      \field{extraname}{2}
      \field{sortinit}{1}
      \field{sortinithash}{50c6687d7fc80f50136d75228e3c59ba}
      \field{labelnamesource}{author}
      \field{labeltitlesource}{title}
      \field{annotation}{The author presents data on $E^\pi(x;6,5,1)$ for $x\le2.5\cdot10^{11}$. He proves (acknowledging A.~Schinzel's help) that $\lim_{x\to\infty} E^\pi(x;6,5,1) \ne 1$, and gives a heuristic argument that $\lim_{x\to\infty} E^{A^\pi_m}(x;6,5,1) \ne 1$ for every $m\in\N$. On the other hand, he shows that $\lim_{m\to\infty} \limsup_{x\to\infty} E^{A^\pi_m}(x;6,5,1) = 1$ and similarly for $\liminf_{x\to\infty}$.\par This article cites~\cites{MR0162771, MR0174538, MR0177963, MR0180539, MR0182616A, MR0182616B,MR0417077}.}
      \field{issn}{0025-5718}
      \field{journaltitle}{Math. Comp.}
      \field{number}{170}
      \field{title}{Averaging effects on irregularities in the distribution of primes in arithmetic progressions}
      \field{volume}{44}
      \field{year}{1985}
      \true{nocite}
      \field{pages}{561\bibrangedash 571}
      \range{pages}{11}
      \verb{doi}
      \verb 10.2307/2007974
      \endverb
      \verb{mrnumber}
      \verb 777286
      \endverb
    \endentry
    \entry{MR791085}{article}{}
      \name{author}{1}{}{%
        {{hash=d1e4ad71dae945a6977ecf6f9ca65c30}{%
           family={Kaczorowski},
           familyi={K\bibinitperiod},
           given={J.},
           giveni={J\bibinitperiod}}}%
      }
      \strng{namehash}{d1e4ad71dae945a6977ecf6f9ca65c30}
      \strng{fullhash}{d1e4ad71dae945a6977ecf6f9ca65c30}
      \strng{bibnamehash}{d1e4ad71dae945a6977ecf6f9ca65c30}
      \strng{authorbibnamehash}{d1e4ad71dae945a6977ecf6f9ca65c30}
      \strng{authornamehash}{d1e4ad71dae945a6977ecf6f9ca65c30}
      \strng{authorfullhash}{d1e4ad71dae945a6977ecf6f9ca65c30}
      \field{extraname}{2}
      \field{sortinit}{1}
      \field{sortinithash}{50c6687d7fc80f50136d75228e3c59ba}
      \field{labelnamesource}{author}
      \field{labeltitlesource}{title}
      \field{annotation}{The author proves unconditionally that $W^\pi(T) \gg \log T$, though with an ineffective constant, and the same for $W^\theta(T)$. He also proves unconditionally that $\liminf_{T\to\infty} {W^\psi(T)}/{\log T} \ge {\gamma(\Theta)}/{\pi}$, where ${\gamma(\Theta)}$ is the smallest $\gamma>0$ such that $\zeta(\Theta+i\gamma)=0$ (or ${\gamma(\Theta)}=\infty$ if $\Theta$ is not attained); this improves a result of Pólya~\cite{1930.Polya}, which had $\limsup$ in place of $\liminf$. He remarks that if RH is false, then the proof of this latter result can be extended to $W^\theta(T)$ and (with a bit more difficulty) to $W^\Pi(T)$ and $W^\pi(T)$. As in a previous article, the proofs of both theorems make use of the iterated averages $A^f_n(x)$.\par This article cites~\cites{MR1555148, 1936.Ingham, MR777013, MR0133309, MR0133308, MR0439771, MR0439772, MR0439774, 1914.Littlewood, MR0439773, 1930.Polya}.}
      \field{issn}{0065-1036}
      \field{journaltitle}{Acta Arith.}
      \field{number}{1}
      \field{title}{On sign-changes in the remainder-term of the prime-number formula. {II}}
      \field{volume}{45}
      \field{year}{1985}
      \true{nocite}
      \field{pages}{65\bibrangedash 74}
      \range{pages}{10}
      \verb{doi}
      \verb 10.4064/aa-45-1-65-74
      \endverb
      \verb{mrnumber}
      \verb 791085
      \endverb
      \verb{urlraw}
      \verb https://doi.org/10.4064/aa-45-1-65-74
      \endverb
      \verb{url}
      \verb https://doi.org/10.4064/aa-45-1-65-74
      \endverb
    \endentry
    \entry{MR783538}{article}{}
      \name{author}{2}{}{%
        {{hash=5763e72244b318a02691f3fc37f706a2}{%
           family={Odlyzko},
           familyi={O\bibinitperiod},
           given={A.\bibnamedelimi M.},
           giveni={A\bibinitperiod\bibinitdelim M\bibinitperiod}}}%
        {{hash=f6fbcde367a0ed5df3da06b96b6d8eff}{%
           family={Riele},
           familyi={R\bibinitperiod},
           given={H.\bibnamedelimi J.\bibnamedelimi J.},
           giveni={H\bibinitperiod\bibinitdelim J\bibinitperiod\bibinitdelim J\bibinitperiod},
           prefix={te},
           prefixi={t\bibinitperiod}}}%
      }
      \strng{namehash}{d88121d28b2f42829b064ec49a086281}
      \strng{fullhash}{d88121d28b2f42829b064ec49a086281}
      \strng{bibnamehash}{d88121d28b2f42829b064ec49a086281}
      \strng{authorbibnamehash}{d88121d28b2f42829b064ec49a086281}
      \strng{authornamehash}{d88121d28b2f42829b064ec49a086281}
      \strng{authorfullhash}{d88121d28b2f42829b064ec49a086281}
      \field{sortinit}{1}
      \field{sortinithash}{50c6687d7fc80f50136d75228e3c59ba}
      \field{labelnamesource}{author}
      \field{labeltitlesource}{title}
      \field{annotation}{The authors disprove the Mertens conjecture by showing that \[ \liminf_{x \to \infty} E^M(x) <-1.009 \quad \text{and} \quad \limsup_{x \to \infty} E^M(x) > 1.06. \] The method, based upon work of Ingham, is to find values of $h_K(y) = \sum_{\rho}k(\gamma) \frac{e^{i\gamma y}}{\rho \zeta'(\rho)}$ that are large in absolute value, using the kernel $k(t) = g(t/T)$ with $T = 2515.286\ldots$ the height of the $2000$th zero of~$\zeta(s)$, where \[ g(t) = \begin{cases} (1-|t|)\cos(\pi t) + \pi^{-1}\sin(\pi |t|), & |t| \leq 1 \\ 0 , & |t| \geq 1 . \end{cases} \] The key development is the algorithm due to Lenstra, Lenstra, and Lovász for finding short vectors in lattices, which reduces the computation time needed to find an appropriate inhomogeneous Diophantine approximation. The authors begin with a summary of past work on the conjecture and conclude with remarks on future work towards finding explicit counterexamples to Mertens conjecture, discussing limitations of their method.\par This article cites~\cites{MR0046485,MR654520,MR0330069,MR0104638,MR0006202,MR0352026,MR736725,1897.vonSterneck,1901.vonSterneck,1912.vonSterneck,1913.vonSterneck,MR0120198,MR429789,MR549977,MR280447,MR240062}.}
      \field{issn}{0075-4102}
      \field{journaltitle}{J. Reine Angew. Math.}
      \field{title}{Disproof of the Mertens conjecture}
      \field{volume}{357}
      \field{year}{1985}
      \true{nocite}
      \field{pages}{138\bibrangedash 160}
      \range{pages}{23}
      \verb{doi}
      \verb 10.1515/crll.1985.357.138
      \endverb
      \verb{mrnumber}
      \verb 783538
      \endverb
    \endentry
    \entry{MR819064}{article}{}
      \name{author}{2}{}{%
        {{hash=c6d965df16e600c701d4bcc2888ae34c}{%
           family={Pintz},
           familyi={P\bibinitperiod},
           given={J.},
           giveni={J\bibinitperiod}}}%
        {{hash=c420715a04171fd2f8ed0fee9cd57cc0}{%
           family={Salerno},
           familyi={S\bibinitperiod},
           given={S.},
           giveni={S\bibinitperiod}}}%
      }
      \strng{namehash}{587c0ae365c9a78befc74bd5d1ffd220}
      \strng{fullhash}{587c0ae365c9a78befc74bd5d1ffd220}
      \strng{bibnamehash}{587c0ae365c9a78befc74bd5d1ffd220}
      \strng{authorbibnamehash}{587c0ae365c9a78befc74bd5d1ffd220}
      \strng{authornamehash}{587c0ae365c9a78befc74bd5d1ffd220}
      \strng{authorfullhash}{587c0ae365c9a78befc74bd5d1ffd220}
      \field{extraname}{4}
      \field{sortinit}{1}
      \field{sortinithash}{50c6687d7fc80f50136d75228e3c59ba}
      \field{labelnamesource}{author}
      \field{labeltitlesource}{title}
      \field{annotation}{Assuming HC$(E_q)$ and GRH$(D)$ with $D\geq D_0=c_0q^2\log^6q$ and $\lambda\geq 20D_0$, the authors show that when~$T$ is sufficiently large, there exists some~$k$ and some~$x$ near~$T$ such that \[ \psi_l(x,4k;q,a,b) > \sqrt{x}\exp\biggl( -\frac{LD^2}{\lambda^2}-\frac{cq^2}{E_q}\lambda\log^3L \biggr). \] The main tool in the proof is a ``one-sided'' version of the second main theorem of the power-sum method. When both~$a$ and~$b$ are quadratic residues or nonresidues, the same result holds with $\psi_l$ replaced by $\theta_l$.\par This article cites~\cites{1853.Chebyshev,1909.Landau,MR0162771,MR0174538,MR0180539,MR0182616A,MR0087685}.}
      \field{issn}{0236-5294}
      \field{journaltitle}{Acta Math. Hungar.}
      \field{number}{1-2}
      \field{title}{Accumulation theorems for primes in arithmetic progressions}
      \field{volume}{46}
      \field{year}{1985}
      \true{nocite}
      \field{pages}{151\bibrangedash 172}
      \range{pages}{22}
      \verb{doi}
      \verb 10.1007/BF01961016
      \endverb
      \verb{mrnumber}
      \verb 819064
      \endverb
    \endentry
    \entry{MR858395}{article}{}
      \name{author}{2}{}{%
        {{hash=d1e4ad71dae945a6977ecf6f9ca65c30}{%
           family={Kaczorowski},
           familyi={K\bibinitperiod},
           given={J.},
           giveni={J\bibinitperiod}}}%
        {{hash=c6d965df16e600c701d4bcc2888ae34c}{%
           family={Pintz},
           familyi={P\bibinitperiod},
           given={J.},
           giveni={J\bibinitperiod}}}%
      }
      \strng{namehash}{2529f5112eea5e5d50df838061479c33}
      \strng{fullhash}{2529f5112eea5e5d50df838061479c33}
      \strng{bibnamehash}{2529f5112eea5e5d50df838061479c33}
      \strng{authorbibnamehash}{2529f5112eea5e5d50df838061479c33}
      \strng{authornamehash}{2529f5112eea5e5d50df838061479c33}
      \strng{authorfullhash}{2529f5112eea5e5d50df838061479c33}
      \field{extraname}{1}
      \field{sortinit}{1}
      \field{sortinithash}{50c6687d7fc80f50136d75228e3c59ba}
      \field{labelnamesource}{author}
      \field{labeltitlesource}{title}
      \field{annotation}{In this article, the authors improve upon and extend results of Landau~\cite{MR1511360}, Pólya~\cite{1930.Polya}, and Grosswald~\cite{MR0198145}. Given a Dirichlet intergal $F(s) = \int_{x_0}^\infty f(x)x^{-s-1}\,dx$ converging on a right half-plane $\{\sigma>\theta\}$, with a continuation to a larger half-plane except for perhaps countably many poles or logarithmic singularities (in a precise sense), the authors show that $\liminf_{T\to\infty} \frac{W(f,T)}{\log T}\ge \frac\gamma\pi$ where $\gamma = \inf\{|t| \colon F(s)\text{ is not regular at }\theta+it\}$. This result implies $\gg \log T$ sign changes of the functions $M(x)$ and $\Delta^{Q_k}(x)$, as well as of $\psi(x;q,\ell_1,\ell_2)$ assuming~HC. For a slightly more restricted class of functions, the authors prove an effective version of a similar result, again guaranteeing $\gg \log T$ sign changes (essentially using a single singularity of $F(s)$) but now with effective constants.\par This article cites~\cites{MR1511360, 1930.Polya, MR0439774, MR0198145, MR0202685, MR0146160, MR0146161, MR688630, MR0218318, MR777013, MR791085}.}
      \field{issn}{0236-5294}
      \field{journaltitle}{Acta Math. Hungar.}
      \field{number}{1-2}
      \field{title}{Oscillatory properties of arithmetical functions. {I}}
      \field{volume}{48}
      \field{year}{1986}
      \true{nocite}
      \field{pages}{173\bibrangedash 185}
      \range{pages}{13}
      \verb{doi}
      \verb 10.1007/BF01949062
      \endverb
      \verb{mrnumber}
      \verb 858395
      \endverb
      \verb{urlraw}
      \verb https://doi.org/10.1007/BF01949062
      \endverb
      \verb{url}
      \verb https://doi.org/10.1007/BF01949062
      \endverb
    \endentry
    \entry{MR845900}{article}{}
      \name{author}{2}{}{%
        {{hash=c6d965df16e600c701d4bcc2888ae34c}{%
           family={Pintz},
           familyi={P\bibinitperiod},
           given={J.},
           giveni={J\bibinitperiod}}}%
        {{hash=c420715a04171fd2f8ed0fee9cd57cc0}{%
           family={Salerno},
           familyi={S\bibinitperiod},
           given={S.},
           giveni={S\bibinitperiod}}}%
      }
      \strng{namehash}{587c0ae365c9a78befc74bd5d1ffd220}
      \strng{fullhash}{587c0ae365c9a78befc74bd5d1ffd220}
      \strng{bibnamehash}{587c0ae365c9a78befc74bd5d1ffd220}
      \strng{authorbibnamehash}{587c0ae365c9a78befc74bd5d1ffd220}
      \strng{authornamehash}{587c0ae365c9a78befc74bd5d1ffd220}
      \strng{authorfullhash}{587c0ae365c9a78befc74bd5d1ffd220}
      \field{extraname}{5}
      \field{sortinit}{1}
      \field{sortinithash}{50c6687d7fc80f50136d75228e3c59ba}
      \field{labelnamesource}{author}
      \field{labeltitlesource}{title}
      \field{annotation}{In this article, the author establishes the effective results $\max_{Y^{7/8} \leq x \leq Y} E^\psi(x;q,a_1,a_2) > e^{-q^{c_1}}$ and $\max_{B_2(Y) \leq x \leq Y} E^\psi(x;q,a_1,a_2) \gg (\log Y)^{-q^{c_2}}$, where $B_2(Y) \gg Y(\log Y)^{-q^{c_3}}$, as well as the same statements where~$\psi$ is replaced by~$\theta$, $\Pi$, or~$\pi$. Furthermore, the author shows that $W^\psi(Y) \gg_q (\log Y)/\log\log Y$ with explicit constants, and similarly with~$\psi$ replaced by~$\Pi$ or, if~$a_1$ and~$a_2$ are both quadratic residues or both quadratic nonresidues modulo~$q$, with~$\theta$ or~$\pi$.\par This article cites~\cites{MR0146156, MR0146157, MR0146158, MR0146159, MR0146160,MR0146161, MR0156826, MR0156827, MR756697, MR764945}.}
      \field{issn}{0022-314X}
      \field{journaltitle}{J. Number Theory}
      \field{number}{2}
      \field{title}{Some consequences of the general Riemann hypothesis in the comparative theory of primes}
      \field{volume}{23}
      \field{year}{1986}
      \true{nocite}
      \field{pages}{183\bibrangedash 194}
      \range{pages}{12}
      \verb{doi}
      \verb 10.1016/0022-314X(86)90088-0
      \endverb
      \verb{mrnumber}
      \verb 845900
      \endverb
    \endentry
    \entry{MR928842}{article}{}
      \name{author}{1}{}{%
        {{hash=28b00c8013a39aa054af5aea7f8c5286}{%
           family={Robin},
           familyi={R\bibinitperiod},
           given={Guy},
           giveni={G\bibinitperiod}}}%
      }
      \list{language}{1}{%
        {French}%
      }
      \strng{namehash}{28b00c8013a39aa054af5aea7f8c5286}
      \strng{fullhash}{28b00c8013a39aa054af5aea7f8c5286}
      \strng{bibnamehash}{28b00c8013a39aa054af5aea7f8c5286}
      \strng{authorbibnamehash}{28b00c8013a39aa054af5aea7f8c5286}
      \strng{authornamehash}{28b00c8013a39aa054af5aea7f8c5286}
      \strng{authorfullhash}{28b00c8013a39aa054af5aea7f8c5286}
      \field{extraname}{2}
      \field{sortinit}{1}
      \field{sortinithash}{50c6687d7fc80f50136d75228e3c59ba}
      \field{labelnamesource}{author}
      \field{labeltitlesource}{title}
      \field{annotation}{This article examines, assuming HC, the asymptotic behaviour of the weighted average $ \cP(x)=\sum_{n\leq x} \mathring\Delta(n;k,\ell) n^{-\alpha}\log^\beta{n} $ where $\alpha$ and $\beta$ are fixed real numbers. If GRH is false, then $\cP(x) \ll 1+x^{1-\alpha+\Theta_k}\log^{\beta-1}{x}$ and $\cP(x)=\Omega_{\pm}(x^{1-\alpha+\Theta_k-\ep})$; under the additional assumption of SA, we have $\cP(x)=\Omega_\pm (x^{1-\alpha+\Theta_k}\log^{\beta-1}{x})$ (which is thus best possible for $\alpha < 1+\Theta_k$).\par If GRH is true, the behaviour depends more significantly upon~$\alpha$ and~$\beta$. When $\alpha>\frac32$, we have $\cP(x)\ll1$. When $\alpha=\frac32$, we have $\cP(x)\ll1$ if $\beta<0$, and $\cP(x)=(1-c_k(\ell))\log\log{x}+O(1)$ if $\beta=0$, and $\cP(x)=(1-c_k(\ell))(\log x)^\beta/\beta+O((\log{x})^{\beta-1}\log\log{x})$ if $\beta>0$. Finally, when $\alpha<\frac32$, we have $\cP(x) \ll x^{3/2-\alpha}\log^{\beta-1}{x}$. This theorem disproves Shanks's conjecture $\sum_{n\le x}\pi(n;4,3,1) n^{1/2}/{\pi(n)}\sim x$, as well as corresponding conjectures for other moduli. Moreover, it shows that Brent's conjecture $\sum_{n\le x}\pi(n;4,3,1) /n^{1/2}{\pi(n)}\sim \log x$ is equivalent to GRH.\par Again assuming GRH and $\alpha<\frac32$, the author asserts that for certain moduli including $3$, $4$, $5$, $6$, $7$, $8$, $9$, $10$, $12$, there exists a constant $\alpha_{k,\ell}$ such that for $\alpha>\alpha_{k,\ell}$, when $x$ is sufficiently large then $\cP(x)<0$ if $\ell$ is a quadratic residue and $\cP(x)>0$ if $\ell$ is a quadratic nonresidue. (It seems that this result actually holds for all moduli $k\ge3$.) On the other hand, for some moduli including $23, 43, 67, 163$, there exists a constant $\alpha'_{k,\ell}$ such that for $\alpha<\alpha'_{k,\ell}$, we have $\cP(x)=\Omega_\pm(x^{3/2-\alpha}\log^{\beta-1}{x})$.\par This article cites~\cites{MR1555148, MR1544293A, MR1544293B, MR0083001, MR0108470, MR0146156, MR0146157, MR0146158, MR0146159, MR0146160, MR0146161, MR0156826, MR0156827, MR0162771, MR0174538, MR0180539, MR0182616A, MR0182616B, MR0200250, MR0289452, MR0302585, MR0369287, MR0417077, MR595317, MR654520, MR602236, MR675189}.}
      \field{issn}{0240-2955}
      \field{journaltitle}{Ann. Fac. Sci. Toulouse Math. (5)}
      \field{number}{2}
      \field{title}{Irrégularités dans la distribution des nombres premiers dans les progressions arithmétiques}
      \field{volume}{8}
      \field{year}{1986}
      \true{nocite}
      \field{pages}{159\bibrangedash 173}
      \range{pages}{15}
      \verb{mrnumber}
      \verb 928842
      \endverb
      \verb{urlraw}
      \verb http://www.numdam.org/item?id=AFST_1986-1987_5_8_2_159_0
      \endverb
      \verb{url}
      \verb http://www.numdam.org/item?id=AFST_1986-1987_5_8_2_159_0
      \endverb
    \endentry
    \entry{MR0882550}{book}{}
      \name{author}{1}{}{%
        {{hash=baf219d417d1941af57ccbb309deb548}{%
           family={Titchmarsh},
           familyi={T\bibinitperiod},
           given={E.\bibnamedelimi C.},
           giveni={E\bibinitperiod\bibinitdelim C\bibinitperiod}}}%
      }
      \list{publisher}{1}{%
        {The Clarendon Press, Oxford University Press, New York}%
      }
      \strng{namehash}{baf219d417d1941af57ccbb309deb548}
      \strng{fullhash}{baf219d417d1941af57ccbb309deb548}
      \strng{bibnamehash}{baf219d417d1941af57ccbb309deb548}
      \strng{authorbibnamehash}{baf219d417d1941af57ccbb309deb548}
      \strng{authornamehash}{baf219d417d1941af57ccbb309deb548}
      \strng{authorfullhash}{baf219d417d1941af57ccbb309deb548}
      \field{extraname}{2}
      \field{sortinit}{1}
      \field{sortinithash}{50c6687d7fc80f50136d75228e3c59ba}
      \field{labelnamesource}{author}
      \field{labeltitlesource}{title}
      \field{annotation}{This book is the second edition of the classical and comprehensive treatise on the theory of the Riemann zeta-function, now with a large number of chapter-end notes by Heath--Brown. Chapter~XII contains the asymptotic formula for the error term in the divisor problem and relevant upper bounds and oscillation results. Chapter~XIV contains an explicit formula for $M(x)$ and its upper bounds and oscillations, as well as a connection between its mean-square average and the simplicity of zeros of $\zeta(s)$; in the chapter-end notes, it is noted that the Mertens conjecture has been disproved and Turán's problem resolved in the negative. \par This article cites~\cite{MR0006202, MR0104638, MR429789, MR783538}.}
      \field{edition}{Second}
      \field{isbn}{0-19-853369-1}
      \field{note}{Edited and with a preface by D. R. Heath-Brown}
      \field{title}{The theory of the {R}iemann zeta-function}
      \field{year}{1986}
      \true{nocite}
      \field{pages}{x+412}
      \range{pages}{-1}
      \verb{mrnumber}
      \verb 0882550
      \endverb
    \endentry
    \entry{MR927376}{article}{}
      \name{author}{1}{}{%
        {{hash=d1e4ad71dae945a6977ecf6f9ca65c30}{%
           family={Kaczorowski},
           familyi={K\bibinitperiod},
           given={J.},
           giveni={J\bibinitperiod}}}%
      }
      \strng{namehash}{d1e4ad71dae945a6977ecf6f9ca65c30}
      \strng{fullhash}{d1e4ad71dae945a6977ecf6f9ca65c30}
      \strng{bibnamehash}{d1e4ad71dae945a6977ecf6f9ca65c30}
      \strng{authorbibnamehash}{d1e4ad71dae945a6977ecf6f9ca65c30}
      \strng{authornamehash}{d1e4ad71dae945a6977ecf6f9ca65c30}
      \strng{authorfullhash}{d1e4ad71dae945a6977ecf6f9ca65c30}
      \field{extraname}{3}
      \field{sortinit}{1}
      \field{sortinithash}{50c6687d7fc80f50136d75228e3c59ba}
      \field{labelnamesource}{author}
      \field{labeltitlesource}{title}
      \field{annotation}{This article examines $W(\Delta^\psi_e;T)$, the number of sign changes of $\Delta^\psi_e(x) = \sum_{n=1}^\infty (\Lambda(n) - 1)e^{-n/x} $ in the interval $[0,T]$. The author first proves that SA (which he calls ``Ingham's condition'') is equivalent to the assertion that $W\big( \Delta^\psi_e;[T,eT] \big) \ll 1$ uniformly for $T>0$, which is further equivalent to each of $\limsup_{T\to\infty} W( \Delta^\psi_e;T ) < \infty$ and $\liminf_{T\to\infty} W( \Delta^\psi_e;T ) < \infty$. Assuming SA and LI$(\Theta)$, the author proves that $W(\Delta^\psi_e;T) \sim \kappa \log T$ as $T \to \infty$, for a constant $\kappa$ (given by an explicit integral) depending on the zeros of $\zeta(s)$ on the line $\sigma = \Theta$. Finally, assuming RH, the author proves that $W(\Delta^\psi_e;T) = \frac{\gamma_1}{\pi} \log T + O(1)$, where $\gamma_1 \approx 14.13$, and indeed that these sign changes are extremely regularly spaced and correspond to oscillations that are $\gg \sqrt x$. These results support the author's conjecture that $W(\Delta^\psi_e;T) \sim c\log T$ as $T \to \infty$. \par This article cites~\cites{MR1555148,1936.Ingham,MR777013,MR791085,MR0131410}.}
      \field{issn}{0065-1036}
      \field{journaltitle}{Acta Arith.}
      \field{number}{4}
      \field{title}{On sign-changes in the remainder-term of the prime-number formula. {III}}
      \field{volume}{48}
      \field{year}{1987}
      \true{nocite}
      \field{pages}{347\bibrangedash 371}
      \range{pages}{25}
      \verb{doi}
      \verb 10.4064/aa-48-4-347-371
      \endverb
      \verb{mrnumber}
      \verb 927376
      \endverb
      \verb{urlraw}
      \verb https://doi.org/10.4064/aa-48-4-347-371
      \endverb
      \verb{url}
      \verb https://doi.org/10.4064/aa-48-4-347-371
      \endverb
    \endentry
    \entry{MR891057}{article}{}
      \name{author}{2}{}{%
        {{hash=d1e4ad71dae945a6977ecf6f9ca65c30}{%
           family={Kaczorowski},
           familyi={K\bibinitperiod},
           given={J.},
           giveni={J\bibinitperiod}}}%
        {{hash=c6d965df16e600c701d4bcc2888ae34c}{%
           family={Pintz},
           familyi={P\bibinitperiod},
           given={J.},
           giveni={J\bibinitperiod}}}%
      }
      \strng{namehash}{2529f5112eea5e5d50df838061479c33}
      \strng{fullhash}{2529f5112eea5e5d50df838061479c33}
      \strng{bibnamehash}{2529f5112eea5e5d50df838061479c33}
      \strng{authorbibnamehash}{2529f5112eea5e5d50df838061479c33}
      \strng{authornamehash}{2529f5112eea5e5d50df838061479c33}
      \strng{authorfullhash}{2529f5112eea5e5d50df838061479c33}
      \field{extraname}{2}
      \field{sortinit}{1}
      \field{sortinithash}{50c6687d7fc80f50136d75228e3c59ba}
      \field{labelnamesource}{author}
      \field{labeltitlesource}{title}
      \field{annotation}{The authors extend their earlier results to obtain $\gg \log T$ sign changes for functions such as $\Delta^\Pi(x;q,a)$, $\Delta^\pi(x;q,a)$ where $a$ is a quadratic nonresidue, $\Delta^\Pi(x;q,a,b)$ where $a\not\equiv b\mod q$, and so on. They also similarly obtain sign changes (in relatively short intervals) for the error term in the asymptotic formula for the counting function of irreducible elements in the ring of integers $\mathcal{O}_K$ of a number field~$K$, assuming the Dedekind zeta function of the Hilbert class field of $K$ does not vanish on the interval $[1/2,1)$ and has at least one simple zero in the half-plane $\sigma>1/2$.\par This article cites~\cites{1936.Ingham, MR777013, MR791085, MR858395}.}
      \field{issn}{0236-5294}
      \field{journaltitle}{Acta Math. Hungar.}
      \field{number}{3-4}
      \field{title}{Oscillatory properties of arithmetical functions. {II}}
      \field{volume}{49}
      \field{year}{1987}
      \true{nocite}
      \field{pages}{441\bibrangedash 453}
      \range{pages}{13}
      \verb{doi}
      \verb 10.1007/BF01951008
      \endverb
      \verb{mrnumber}
      \verb 891057
      \endverb
      \verb{urlraw}
      \verb https://doi.org/10.1007/BF01951008
      \endverb
      \verb{url}
      \verb https://doi.org/10.1007/BF01951008
      \endverb
    \endentry
    \entry{MR891440}{article}{}
      \name{author}{1}{}{%
        {{hash=c6d965df16e600c701d4bcc2888ae34c}{%
           family={Pintz},
           familyi={P\bibinitperiod},
           given={J.},
           giveni={J\bibinitperiod}}}%
      }
      \strng{namehash}{c6d965df16e600c701d4bcc2888ae34c}
      \strng{fullhash}{c6d965df16e600c701d4bcc2888ae34c}
      \strng{bibnamehash}{c6d965df16e600c701d4bcc2888ae34c}
      \strng{authorbibnamehash}{c6d965df16e600c701d4bcc2888ae34c}
      \strng{authornamehash}{c6d965df16e600c701d4bcc2888ae34c}
      \strng{authorfullhash}{c6d965df16e600c701d4bcc2888ae34c}
      \field{extraname}{17}
      \field{sortinit}{1}
      \field{sortinithash}{50c6687d7fc80f50136d75228e3c59ba}
      \field{labelnamesource}{author}
      \field{labeltitlesource}{title}
      \field{annotation}{The disproof of the Mertens conjecture by Odlyzko and te Riele~\cites{MR783538} does not provide any effective counterexample. In this article, the author shows that there exists $x<\exp(3.21 \cdot 10^{64})$, such that $|M(x)| > \sqrt{x}$.\par This article cites~\cites{MR783538, 1897.Mertens, MR0006202,MR429789}.}
      \field{issn}{0303-1179}
      \field{journaltitle}{Astérisque}
      \field{number}{147-148}
      \field{title}{An effective disproof of the Mertens conjecture}
      \field{year}{1987}
      \true{nocite}
      \field{pages}{325\bibrangedash 333\bibrangessep 346}
      \range{pages}{10}
      \verb{mrnumber}
      \verb 891440
      \endverb
    \endentry
    \entry{MR866118}{article}{}
      \name{author}{1}{}{%
        {{hash=492b43ecdc1e86c9cf9402cb28b077f2}{%
           family={Riele},
           familyi={R\bibinitperiod},
           given={Herman\bibnamedelimb J.\bibnamedelimi J.},
           giveni={H\bibinitperiod\bibinitdelim J\bibinitperiod\bibinitdelim J\bibinitperiod},
           prefix={te},
           prefixi={t\bibinitperiod}}}%
      }
      \strng{namehash}{492b43ecdc1e86c9cf9402cb28b077f2}
      \strng{fullhash}{492b43ecdc1e86c9cf9402cb28b077f2}
      \strng{bibnamehash}{492b43ecdc1e86c9cf9402cb28b077f2}
      \strng{authorbibnamehash}{492b43ecdc1e86c9cf9402cb28b077f2}
      \strng{authornamehash}{492b43ecdc1e86c9cf9402cb28b077f2}
      \strng{authorfullhash}{492b43ecdc1e86c9cf9402cb28b077f2}
      \field{extraname}{2}
      \field{sortinit}{1}
      \field{sortinithash}{50c6687d7fc80f50136d75228e3c59ba}
      \field{labelnamesource}{author}
      \field{labeltitlesource}{title}
      \field{annotation}{The author shows that $\pi(x) > \li(x)$ for some $6.62 \times 10^{370} \leq x \leq 6.69 \times 10^{370}$, thereby improving the previous best estimate, $1.65 \times 10^{1165}$, for Skewes's number found by Lehman~\cite{MR0202686}. Using an explicit formula for $E^\pi(e^u)$ averaged by a Gaussian kernel, Lehman had found three candidates for $x$ near which $\pi(x) > \li(x)$, namely $e^{727.952}$, $e^{853.853}$, and $e^{2682.977}$. Lehman showed that $e^{2682.977}$ produced an actual example; using the zeros of $\zeta(s)$ up to height $5 \times 10^4$, found on a CYBER 205 supercomputer located at the Academic Computer Centre Amsterdam, the author shows that $e^{853.853}$ produces an actual example. The author speculates that zeros up to height $4 \times 10^5$ would be required to determine whether there is an actual example around $e^{727.952}$. \par This article cites~\cites{MR0202686,MR0067145,1914.Littlewood}.}
      \field{issn}{0025-5718}
      \field{journaltitle}{Math. Comp.}
      \field{number}{177}
      \field{title}{On the sign of the difference {$\pi(x)-\li(x)$}}
      \field{volume}{48}
      \field{year}{1987}
      \true{nocite}
      \field{pages}{323\bibrangedash 328}
      \range{pages}{6}
      \verb{doi}
      \verb 10.2307/2007893
      \endverb
      \verb{mrnumber}
      \verb 866118
      \endverb
      \verb{urlraw}
      \verb https://doi.org/10.2307/2007893
      \endverb
      \verb{url}
      \verb https://doi.org/10.2307/2007893
      \endverb
    \endentry
    \entry{MR974088}{article}{}
      \name{author}{1}{}{%
        {{hash=ed673dd9786336c3366d68e7f72549d4}{%
           family={Fujii},
           familyi={F\bibinitperiod},
           given={A.},
           giveni={A\bibinitperiod}}}%
      }
      \strng{namehash}{ed673dd9786336c3366d68e7f72549d4}
      \strng{fullhash}{ed673dd9786336c3366d68e7f72549d4}
      \strng{bibnamehash}{ed673dd9786336c3366d68e7f72549d4}
      \strng{authorbibnamehash}{ed673dd9786336c3366d68e7f72549d4}
      \strng{authornamehash}{ed673dd9786336c3366d68e7f72549d4}
      \strng{authorfullhash}{ed673dd9786336c3366d68e7f72549d4}
      \field{sortinit}{1}
      \field{sortinithash}{50c6687d7fc80f50136d75228e3c59ba}
      \field{labelnamesource}{author}
      \field{labeltitlesource}{title}
      \field{annotation}{The author shows if $0<\alpha\le4$, then the statement $\lim_{x\to\infty}\sum_p \chi_{-4}(p)e^{-(p/x)^\alpha}= - \infty$ is equivalent to GRH for $L(s,\chi_{-4})$, generalizing earlier results for $\alpha=1$. Further, the author defines $\xi(x,k)$ by $\Gamma(s)^k=\int_0^{\infty}x^{s-1}\xi(x,k)\, dx$ and shows the statement $\lim_{x\to\infty}\sum_p \chi_{-4}(p) \log{p}\cdot\xi(p/x,2)=-\infty$ is also equivalent to GRH for $L(s,\chi_{-4})$. The author conjectures that the first equivalence holds for any $\alpha>0$ and the second (without the factor $\log p$) for any integer $k\ge1$. He also remarks that the method can be used for other weight functions involving exponentials and Bessel functions, including Knapowski--Turán's function $e^{-\log^2(p/x)}$. \par This article cites~\cites{MR595317,MR1544293A,MR1544293B}.}
      \field{issn}{0386-2194}
      \field{journaltitle}{Proc. Japan Acad. Ser. A Math. Sci.}
      \field{number}{7}
      \field{title}{Some generalizations of {C}hebyshev's conjecture}
      \field{volume}{64}
      \field{year}{1988}
      \true{nocite}
      \field{pages}{260\bibrangedash 263}
      \range{pages}{4}
      \verb{mrnumber}
      \verb 974088
      \endverb
      \verb{urlraw}
      \verb http://projecteuclid.org/euclid.pja/1195513180
      \endverb
      \verb{url}
      \verb http://projecteuclid.org/euclid.pja/1195513180
      \endverb
    \endentry
    \entry{MR945273}{article}{}
      \name{author}{1}{}{%
        {{hash=d1e4ad71dae945a6977ecf6f9ca65c30}{%
           family={Kaczorowski},
           familyi={K\bibinitperiod},
           given={J.},
           giveni={J\bibinitperiod}}}%
      }
      \strng{namehash}{d1e4ad71dae945a6977ecf6f9ca65c30}
      \strng{fullhash}{d1e4ad71dae945a6977ecf6f9ca65c30}
      \strng{bibnamehash}{d1e4ad71dae945a6977ecf6f9ca65c30}
      \strng{authorbibnamehash}{d1e4ad71dae945a6977ecf6f9ca65c30}
      \strng{authornamehash}{d1e4ad71dae945a6977ecf6f9ca65c30}
      \strng{authorfullhash}{d1e4ad71dae945a6977ecf6f9ca65c30}
      \field{extraname}{4}
      \field{sortinit}{1}
      \field{sortinithash}{50c6687d7fc80f50136d75228e3c59ba}
      \field{labelnamesource}{author}
      \field{labeltitlesource}{title}
      \field{annotation}{The author proves, when $\Theta>\frac12$, that for any $\ep>0$ we have $\max_{T \le x \le (1+\ep)T} |\Delta_e^\psi(x)| \gg_\ep T^{\Theta-\ep}$. In light of the author's previous results~\cite{MR927376} that assumed RH, it follows that unconditionally (but ineffectively), $\max_{T \le x \le (1+\ep)T} |\Delta_e^\psi(x)| \gg_\ep \sqrt T$. The author also deduces that $W_e^\psi(T) = o(\log^2 T)$, and sketches a construction (of a ``barrier'') showing that this result cannot be improved without further information on the zeros of~$\zeta(s)$.\par This article cites~\cites{MR0125822,MR0131410,MR777013,MR791085,MR927376, MR0060548}.}
      \field{issn}{0065-1036}
      \field{journaltitle}{Acta Arith.}
      \field{number}{1}
      \field{title}{On sign-changes in the remainder-term of the prime-number formula. {IV}}
      \field{volume}{50}
      \field{year}{1988}
      \true{nocite}
      \field{pages}{15\bibrangedash 21}
      \range{pages}{7}
      \verb{doi}
      \verb 10.4064/aa-50-1-15-21
      \endverb
      \verb{mrnumber}
      \verb 945273
      \endverb
      \verb{urlraw}
      \verb https://doi.org/10.4064/aa-50-1-15-21
      \endverb
      \verb{url}
      \verb https://doi.org/10.4064/aa-50-1-15-21
      \endverb
    \endentry
    \entry{MR0980524}{article}{}
      \name{author}{2}{}{%
        {{hash=d1e4ad71dae945a6977ecf6f9ca65c30}{%
           family={Kaczorowski},
           familyi={K\bibinitperiod},
           given={J.},
           giveni={J\bibinitperiod}}}%
        {{hash=36befda65f61ea63f097e5518d3ec577}{%
           family={Staś},
           familyi={S\bibinitperiod},
           given={W.},
           giveni={W\bibinitperiod}}}%
      }
      \strng{namehash}{d5a8f4983f5607d22545d2b33d58100f}
      \strng{fullhash}{d5a8f4983f5607d22545d2b33d58100f}
      \strng{bibnamehash}{d5a8f4983f5607d22545d2b33d58100f}
      \strng{authorbibnamehash}{d5a8f4983f5607d22545d2b33d58100f}
      \strng{authornamehash}{d5a8f4983f5607d22545d2b33d58100f}
      \strng{authorfullhash}{d5a8f4983f5607d22545d2b33d58100f}
      \field{extraname}{1}
      \field{sortinit}{1}
      \field{sortinithash}{50c6687d7fc80f50136d75228e3c59ba}
      \field{labelnamesource}{author}
      \field{labeltitlesource}{title}
      \field{annotation}{The author shows that if $K$ is a number field such that HC holds for $\zeta_K(s)$, then for any $0<\ep <1$ one has $W^{\psi_K}(T) \geq (1-\ep) \frac{\gamma_K}{\pi} \log T$ when $T$ is sufficiently large in terms of $K$ and $\ep$. The assumption of no real zeros can be removed if the contribution from the real zeros is subtracted from $E^{\psi_K}(x)$. \par This article cites~\cite{1936.Ingham,MR777013,MR817334,MR0132728}.}
      \field{issn}{0010-1354}
      \field{journaltitle}{Colloq. Math.}
      \field{number}{1}
      \field{title}{On the number of sign changes in the remainder-term of the prime-ideal theorem}
      \field{volume}{56}
      \field{year}{1988}
      \true{nocite}
      \field{pages}{185\bibrangedash 197}
      \range{pages}{13}
      \verb{doi}
      \verb 10.4064/cm-56-1-185-197
      \endverb
      \verb{mrnumber}
      \verb 980524
      \endverb
      \verb{urlraw}
      \verb https://doi.org/10.4064/cm-56-1-185-197
      \endverb
      \verb{url}
      \verb https://doi.org/10.4064/cm-56-1-185-197
      \endverb
    \endentry
    \entry{MR1042465}{article}{}
      \name{author}{2}{}{%
        {{hash=6fa84c8cc1e455847a05b9946a1bd205}{%
           family={Kaczorowski},
           familyi={K\bibinitperiod},
           given={Jerzy},
           giveni={J\bibinitperiod}}}%
        {{hash=b39a346ae2b44071e01b12c1003e1246}{%
           family={Staś},
           familyi={S\bibinitperiod},
           given={Włodzimierz},
           giveni={W\bibinitperiod}}}%
      }
      \strng{namehash}{e170c0eabd81bb1165f3c0582dd5b3a2}
      \strng{fullhash}{e170c0eabd81bb1165f3c0582dd5b3a2}
      \strng{bibnamehash}{e170c0eabd81bb1165f3c0582dd5b3a2}
      \strng{authorbibnamehash}{e170c0eabd81bb1165f3c0582dd5b3a2}
      \strng{authornamehash}{e170c0eabd81bb1165f3c0582dd5b3a2}
      \strng{authorfullhash}{e170c0eabd81bb1165f3c0582dd5b3a2}
      \field{extraname}{2}
      \field{sortinit}{1}
      \field{sortinithash}{50c6687d7fc80f50136d75228e3c59ba}
      \field{labelnamesource}{author}
      \field{labeltitlesource}{title}
      \field{annotation}{This article is identical to~\cite{MR0980524}.}
      \field{issn}{0137-9747}
      \field{journaltitle}{Discuss. Math.}
      \field{title}{On the number of sign-changes in the remainder-term of the prime-ideal theorem}
      \field{volume}{9}
      \field{year}{1988}
      \true{nocite}
      \field{pages}{83\bibrangedash 102 (1989)}
      \range{pages}{-1}
      \verb{mrnumber}
      \verb 1042465
      \endverb
    \endentry
    \entry{MR0973808A}{article}{}
      \name{author}{1}{}{%
        {{hash=83e5472a007b4ceae548c8cea09a150a}{%
           family={Szydło},
           familyi={S\bibinitperiod},
           given={B.},
           giveni={B\bibinitperiod}}}%
      }
      \list{language}{1}{%
        {German}%
      }
      \strng{namehash}{83e5472a007b4ceae548c8cea09a150a}
      \strng{fullhash}{83e5472a007b4ceae548c8cea09a150a}
      \strng{bibnamehash}{83e5472a007b4ceae548c8cea09a150a}
      \strng{authorbibnamehash}{83e5472a007b4ceae548c8cea09a150a}
      \strng{authornamehash}{83e5472a007b4ceae548c8cea09a150a}
      \strng{authorfullhash}{83e5472a007b4ceae548c8cea09a150a}
      \field{extraname}{1}
      \field{sortinit}{1}
      \field{sortinithash}{50c6687d7fc80f50136d75228e3c59ba}
      \field{labelnamesource}{author}
      \field{labeltitlesource}{title}
      \field{annotation}{This article concerns effective results of the type $W(f;X) \ge c\log X$ for $X\ge X_0$, using a modification of Kaczorowski's method with a more general averaging operator $x^{-d-k}\int_0^{x}f(\xi) \xi^{d+k-1} \,d\xi$.\par This article cites~\cite{1914.Littlewood, MR777013, MR791085, MR858395, MR0973808B}.}
      \field{journaltitle}{Math. Ann.}
      \field{title}{{Über Vorzeichenwechsel einiger arithmetischer Funktionen. I}}
      \field{volume}{283}
      \field{year}{1989}
      \true{nocite}
      \field{pages}{139\bibrangedash 149}
      \range{pages}{11}
      \verb{mrnumber}
      \verb 0973808
      \endverb
    \endentry
    \entry{MR0973808B}{article}{}
      \name{author}{1}{}{%
        {{hash=83e5472a007b4ceae548c8cea09a150a}{%
           family={Szydło},
           familyi={S\bibinitperiod},
           given={B.},
           giveni={B\bibinitperiod}}}%
      }
      \list{language}{1}{%
        {German}%
      }
      \strng{namehash}{83e5472a007b4ceae548c8cea09a150a}
      \strng{fullhash}{83e5472a007b4ceae548c8cea09a150a}
      \strng{bibnamehash}{83e5472a007b4ceae548c8cea09a150a}
      \strng{authorbibnamehash}{83e5472a007b4ceae548c8cea09a150a}
      \strng{authornamehash}{83e5472a007b4ceae548c8cea09a150a}
      \strng{authorfullhash}{83e5472a007b4ceae548c8cea09a150a}
      \field{extraname}{2}
      \field{sortinit}{1}
      \field{sortinithash}{50c6687d7fc80f50136d75228e3c59ba}
      \field{labelnamesource}{author}
      \field{labeltitlesource}{title}
      \field{annotation}{The author uses his results from~\cite{MR0973808A} to give explicit lower bounds for the number of sign changes of $\Delta^\psi(X)$. He shows unconditionally that $W^\psi(X)\ge 0.013\log X$ when $X\ge 10^{2250}$, the proof requiring minimal computation, as well as the similar results $W^\psi(X)\ge 0.994 \frac{\gamma_1}{\pi} \log X$ for $X\ge \exp (198,594)$ and $W^\psi(X)\ge 0.99999997 \frac{\gamma_1}{\pi} \log X$ for $X\ge \exp (9 \times 10^{14})$. Assuming RH$(H)$ for some $H\ge501.5$, he further shows that $W^\psi(X)\ge (1-\frac{3}{H})\frac{\gamma_1}{\pi}\log X$ when $X\ge \exp(0.09\max\{4400,H\})$.\par This article cites~\cite{MR777013, MR791085, MR858395, MR0973808A}.}
      \field{journaltitle}{Math. Ann.}
      \field{title}{{Über Vorzeichenwechsel einiger arithmetischer Funktionen. II}}
      \field{volume}{283}
      \field{year}{1989}
      \true{nocite}
      \field{pages}{151\bibrangedash 163}
      \range{pages}{13}
      \verb{mrnumber}
      \verb 0973808
      \endverb
    \endentry
    \entry{MR1029966}{article}{}
      \name{author}{1}{}{%
        {{hash=83e5472a007b4ceae548c8cea09a150a}{%
           family={Szydło},
           familyi={S\bibinitperiod},
           given={B.},
           giveni={B\bibinitperiod}}}%
      }
      \list{language}{1}{%
        {German}%
      }
      \strng{namehash}{83e5472a007b4ceae548c8cea09a150a}
      \strng{fullhash}{83e5472a007b4ceae548c8cea09a150a}
      \strng{bibnamehash}{83e5472a007b4ceae548c8cea09a150a}
      \strng{authorbibnamehash}{83e5472a007b4ceae548c8cea09a150a}
      \strng{authornamehash}{83e5472a007b4ceae548c8cea09a150a}
      \strng{authorfullhash}{83e5472a007b4ceae548c8cea09a150a}
      \field{extraname}{3}
      \field{sortinit}{1}
      \field{sortinithash}{50c6687d7fc80f50136d75228e3c59ba}
      \field{labelnamesource}{author}
      \field{labeltitlesource}{title}
      \field{annotation}{The author generalizes his results~\cites{MR0973808A, MR0973808B} on the number of sign changes in $W^\psi(x)$ to algebraic number fields. Let $K$ be an algebraic number field of degree $n$ and $d$ the absolute value of its discriminant. In this article, he extends previous work~\cite{MR0980524} by giving effective lower bounds for the function $W^{\psi_K}(X)$ which counts the sign changes of $\sum_{{\mathfrak Np}^r \leq x}\log {\mathfrak Np}-x$. Assume GRH$(H,0)$ for the Dedekind zeta function~$\zeta_K$. When~$H$ and~$X$ are sufficiently large (in a precise way given with effective constants), the author shows that $W^{\psi_K}(X)\ge (1-\frac{10}{H})\frac{\gamma_K}{\pi}\log X$, where $\gamma_K$ is the positive imaginary part of the nontrivial zero of~$\zeta_K$ closest to the real axis.\par This article cites~\cite{MR777013, MR791085, MR858395, MR891057, MR0973808A, MR0973808B,MR0980524}.}
      \field{journaltitle}{Monatsh. Math.}
      \field{title}{{Über Vorzeichenwechsel einiger arithmetischer Funktionen. III}}
      \field{volume}{108}
      \field{year}{1989}
      \true{nocite}
      \field{pages}{325\bibrangedash 336}
      \range{pages}{12}
      \verb{mrnumber}
      \verb 1029966
      \endverb
    \endentry
    \entry{MR1083000}{article}{}
      \name{author}{1}{}{%
        {{hash=d1e4ad71dae945a6977ecf6f9ca65c30}{%
           family={Kaczorowski},
           familyi={K\bibinitperiod},
           given={J.},
           giveni={J\bibinitperiod}}}%
      }
      \strng{namehash}{d1e4ad71dae945a6977ecf6f9ca65c30}
      \strng{fullhash}{d1e4ad71dae945a6977ecf6f9ca65c30}
      \strng{bibnamehash}{d1e4ad71dae945a6977ecf6f9ca65c30}
      \strng{authorbibnamehash}{d1e4ad71dae945a6977ecf6f9ca65c30}
      \strng{authornamehash}{d1e4ad71dae945a6977ecf6f9ca65c30}
      \strng{authorfullhash}{d1e4ad71dae945a6977ecf6f9ca65c30}
      \field{extraname}{5}
      \field{sortinit}{1}
      \field{sortinithash}{50c6687d7fc80f50136d75228e3c59ba}
      \field{labelnamesource}{author}
      \field{labeltitlesource}{title}
      \field{annotation}{This is the first in a series of articles on the ``$k$-functions'' $k(z,\chi)$ and $K(z,\chi)$ and certain limiting values $F(x,\chi)$ of the latter (see Section~\ref{k-functions section} of this annotated bibliography for definitions). In Section~$3$, the author proves that $k(z,\chi)$ can be analytically continued to a meromorphic function on the Riemann surface~$\cM$ for $\log z$, and indeed that \begin{align*} k(z, \chi) - \frac{1}{2\pi i} \frac{e^z}{e^z-1} \log z \end{align*} is meromorphic and single-valued on $\mathbb C$. Indeed, the author finds all of the singularities of $k(z,\chi)$ on~$\cM$ (all simple poles) and their residues. He also establishes the functional equations \[ k(z,\chi) + e^z k(z^*, \overline{\chi}) = D(z,\chi), \qquad k(z,\chi) + \overline{k(z^c, \overline{\chi})} = e^zD(-z,\chi). \] In Section~$4$, the author establishes explicit formulas for $\psi_0(x,\chi)$ and $\psi_{0r}(x,\chi)$, stated in the forms \begin{align*} \qquad F(x, \chi) +\sum_{\substack{\beta>0 \\ L(\beta, \chi)=0}} \frac{e^{\beta x}}{\beta} = -\psi_0(e^x, \chi) - R_{\chi(-1)}(x) + B(\chi) + \begin{cases} e^x, &\text{if } \chi=\chi_0, \\ -x, &\text{if } \chi\ne\chi_0 \text{ and } \chi(-1)=1, \\ 0, &\text{if } \chi(-1) = -1 \end{cases} \end{align*} for $x>0$, and \begin{align*} \qquad \,\, F(x, \chi) + \sum_{\substack{\beta>0 \\ L(\beta, \chi)=0}} \frac{e^{\beta x}}{\beta} = \psi_{0r}(e^{|x|}, \chi) + R_{-\chi(-1)}(|x|) + C(\chi) + \begin{cases} e^x, &\text{if } \chi=\chi_0, \\ x, &\text{if } \chi\ne\chi_0 \text{ and } \chi(-1)=1, \\ 0, &\text{if } \chi(-1) = -1 \end{cases} \end{align*} for $x<0$. The author then shows that the left-hand side is equal to the series $\sum_{\rho} {e^{\rho x}}/{\rho}$ as in the classical explicit formulas for the right-hand sides.\par This article cites~\cites{1932.Ingham}.}
      \field{issn}{0065-1036}
      \field{journaltitle}{Acta Arith.}
      \field{number}{3}
      \field{title}{The {$k$}-functions in multiplicative number theory. {I}. {O}n complex explicit formulae}
      \field{volume}{56}
      \field{year}{1990}
      \true{nocite}
      \field{pages}{195\bibrangedash 211}
      \range{pages}{17}
      \verb{doi}
      \verb 10.4064/aa-56-3-195-211
      \endverb
      \verb{mrnumber}
      \verb 1083000
      \endverb
      \verb{urlraw}
      \verb https://doi.org/10.4064/aa-56-3-195-211
      \endverb
      \verb{url}
      \verb https://doi.org/10.4064/aa-56-3-195-211
      \endverb
    \endentry
    \entry{MR1083001}{article}{}
      \name{author}{1}{}{%
        {{hash=d1e4ad71dae945a6977ecf6f9ca65c30}{%
           family={Kaczorowski},
           familyi={K\bibinitperiod},
           given={J.},
           giveni={J\bibinitperiod}}}%
      }
      \strng{namehash}{d1e4ad71dae945a6977ecf6f9ca65c30}
      \strng{fullhash}{d1e4ad71dae945a6977ecf6f9ca65c30}
      \strng{bibnamehash}{d1e4ad71dae945a6977ecf6f9ca65c30}
      \strng{authorbibnamehash}{d1e4ad71dae945a6977ecf6f9ca65c30}
      \strng{authornamehash}{d1e4ad71dae945a6977ecf6f9ca65c30}
      \strng{authorfullhash}{d1e4ad71dae945a6977ecf6f9ca65c30}
      \field{extraname}{6}
      \field{sortinit}{1}
      \field{sortinithash}{50c6687d7fc80f50136d75228e3c59ba}
      \field{labelnamesource}{author}
      \field{labeltitlesource}{title}
      \field{annotation}{Let $0 < \gamma_1 \leq \gamma_2 \leq \cdots$ denote the imaginary parts of nontrivial zeros of $L(s, \chi)$ in the upper half plane. In this article Kaczorowski defines a positive Toeplitz matrix $\boldsymbol{A}=(a_{nk})$ by $a_{nk} = e^{-\gamma_k} \gamma_k^n \left(\sum_{h=1}^\infty e^{-\gamma_h}\gamma_h^n\right)^{-1}$ (for $n, k \geq 1$), and proves that for any nonzero real $x$, the sequence $(x\gamma_n)_{n=1}^\infty$ is $\boldsymbol{A}$-uniformly distributed\mod1; the known result that the $x\gamma_n$ are uniformly distributed\mod1 (in the sense of Weyl) follows as a corollary.\par This article cites~\cites{MR1083000}.}
      \field{issn}{0065-1036}
      \field{journaltitle}{Acta Arith.}
      \field{number}{3}
      \field{title}{The {$k$}-functions in multiplicative number theory. {II}. {U}niform distribution of zeta zeros}
      \field{volume}{56}
      \field{year}{1990}
      \true{nocite}
      \field{pages}{213\bibrangedash 224}
      \range{pages}{12}
      \verb{doi}
      \verb 10.4064/aa-56-3-213-224
      \endverb
      \verb{mrnumber}
      \verb 1083001
      \endverb
      \verb{urlraw}
      \verb https://doi.org/10.4064/aa-56-3-213-224
      \endverb
      \verb{url}
      \verb https://doi.org/10.4064/aa-56-3-213-224
      \endverb
    \endentry
    \entry{MR1109990}{article}{}
      \name{author}{1}{}{%
        {{hash=3a6f6cc642ae49cd63edd92858935cb9}{%
           family={Bartz},
           familyi={B\bibinitperiod},
           given={K.\bibnamedelimi M.},
           giveni={K\bibinitperiod\bibinitdelim M\bibinitperiod}}}%
      }
      \strng{namehash}{3a6f6cc642ae49cd63edd92858935cb9}
      \strng{fullhash}{3a6f6cc642ae49cd63edd92858935cb9}
      \strng{bibnamehash}{3a6f6cc642ae49cd63edd92858935cb9}
      \strng{authorbibnamehash}{3a6f6cc642ae49cd63edd92858935cb9}
      \strng{authornamehash}{3a6f6cc642ae49cd63edd92858935cb9}
      \strng{authorfullhash}{3a6f6cc642ae49cd63edd92858935cb9}
      \field{sortinit}{1}
      \field{sortinithash}{50c6687d7fc80f50136d75228e3c59ba}
      \field{labelnamesource}{author}
      \field{labeltitlesource}{title}
      \field{annotation}{Assuming RH and the simplicity of the zeros of $\zeta(s)$, Titchmarsh showed that \begin{align*} M_0(x)=\sum_\rho \frac{x^{\rho}}{\rho\zeta'(\rho)}-2-\sum_{n=1}^{\infty}\frac{(-1)^n(2\pi/x)^{2n}}{(2n)!n\zeta(2n+1)}. \end{align*} In Part I, the author investigates the function $m(z)=\sum_{\Im\rho>0}{e^{\rho z}}/{\zeta'(\rho)}$ (still assuming the simplicity of the zeros). He shows that $m(z)$ is a holomorphic function for $\Im z>0$ that has an analytic continuation to a meromorphic function on $\mathbb{C}$ satisfying the functional equation \begin{align*} m(z)+\overline{m(\overline{z})}=-2\sum_{n=1}^{\infty}\frac{\mu(n)}{n}\cos\bigg(\frac{2\pi}{n}e^{-z}\bigg). \end{align*} Using analytic properties of $m(z)$, he establishes Titchmarsh's formula for $M_0(x)$ without assuming RH. \par This article cites~\cites{1932.Ingham,MR0046485,MR783538,MR1083000}.}
      \field{issn}{0065-1036}
      \field{journaltitle}{Acta Arith.}
      \field{number}{4}
      \field{title}{On some complex explicit formulae connected with the {M}öbius function. {I}, {II}}
      \field{volume}{57}
      \field{year}{1991}
      \true{nocite}
      \field{pages}{283\bibrangedash 293\bibrangessep 295\bibrangedash 305}
      \range{pages}{22}
      \verb{doi}
      \verb 10.4064/aa-57-4-283-293
      \endverb
      \verb{mrnumber}
      \verb 1109990
      \endverb
      \verb{urlraw}
      \verb https://doi.org/10.4064/aa-57-4-283-293
      \endverb
      \verb{url}
      \verb https://doi.org/10.4064/aa-57-4-283-293
      \endverb
    \endentry
    \entry{MR1105605}{article}{}
      \name{author}{1}{}{%
        {{hash=d1e4ad71dae945a6977ecf6f9ca65c30}{%
           family={Kaczorowski},
           familyi={K\bibinitperiod},
           given={J.},
           giveni={J\bibinitperiod}}}%
      }
      \strng{namehash}{d1e4ad71dae945a6977ecf6f9ca65c30}
      \strng{fullhash}{d1e4ad71dae945a6977ecf6f9ca65c30}
      \strng{bibnamehash}{d1e4ad71dae945a6977ecf6f9ca65c30}
      \strng{authorbibnamehash}{d1e4ad71dae945a6977ecf6f9ca65c30}
      \strng{authornamehash}{d1e4ad71dae945a6977ecf6f9ca65c30}
      \strng{authorfullhash}{d1e4ad71dae945a6977ecf6f9ca65c30}
      \field{extraname}{7}
      \field{sortinit}{1}
      \field{sortinithash}{50c6687d7fc80f50136d75228e3c59ba}
      \field{labelnamesource}{author}
      \field{labeltitlesource}{title}
      \field{annotation}{Continuing the previous article in this series, the author defines an ``$\boldsymbol{A}$-discrepancy'' \begin{align*} D_n^*(x) = \sup_{0\leq t \leq 1} \bigg| \bigg( \sum_{\substack{k\ge1\\ \{x\gamma_k\} < t}} e^{-\gamma_k} \gamma_k^n \bigg) \bigg/ \bigg( \sum_{k=1}^\infty e^{-\gamma_k} \gamma_k^n \bigg) - t\bigg|. \end{align*} He shows that $D_n^*(x) \ll (\log \log n/\log n)^{2/3}$ for every real number $x \ne 0$. Under a certain $\boldsymbol{A}$-variant of zero-density theorems for Dirichlet $L$-functions, he proves that $D_n^*(x) \ll 1/\log n$ and conjectures that $D_n^*(x) \sim \alpha(x)/\log n$ for some constant~$\alpha(x)$.\par This article cites~\cites{MR1083001}.}
      \field{issn}{0065-1036}
      \field{journaltitle}{Acta Arith.}
      \field{number}{3}
      \field{title}{The {$k$}-functions in multiplicative number theory. {III}. {U}niform distribution of zeta zeros; discrepancy}
      \field{volume}{57}
      \field{year}{1991}
      \true{nocite}
      \field{pages}{199\bibrangedash 210}
      \range{pages}{12}
      \verb{doi}
      \verb 10.4064/aa-57-3-199-210
      \endverb
      \verb{mrnumber}
      \verb 1105605
      \endverb
      \verb{urlraw}
      \verb https://doi.org/10.4064/aa-57-3-199-210
      \endverb
      \verb{url}
      \verb https://doi.org/10.4064/aa-57-3-199-210
      \endverb
    \endentry
    \entry{MR1105608}{article}{}
      \name{author}{1}{}{%
        {{hash=d1e4ad71dae945a6977ecf6f9ca65c30}{%
           family={Kaczorowski},
           familyi={K\bibinitperiod},
           given={J.},
           giveni={J\bibinitperiod}}}%
      }
      \strng{namehash}{d1e4ad71dae945a6977ecf6f9ca65c30}
      \strng{fullhash}{d1e4ad71dae945a6977ecf6f9ca65c30}
      \strng{bibnamehash}{d1e4ad71dae945a6977ecf6f9ca65c30}
      \strng{authorbibnamehash}{d1e4ad71dae945a6977ecf6f9ca65c30}
      \strng{authornamehash}{d1e4ad71dae945a6977ecf6f9ca65c30}
      \strng{authorfullhash}{d1e4ad71dae945a6977ecf6f9ca65c30}
      \field{extraname}{8}
      \field{sortinit}{1}
      \field{sortinithash}{50c6687d7fc80f50136d75228e3c59ba}
      \field{labelnamesource}{author}
      \field{labeltitlesource}{title}
      \field{annotation}{Using $k$-functions, the author proves an oscillation theorem for a specific class of almost periodic functions, generalizing results of Ingham. He then applies the theorem to prime counting functions: for instance, he shows that for every positive constant $A>0$, there exists a number $c>1$ such that for all $x>1$, we have $\sup E^\psi(x) \ge A$ and $\inf E^\psi(x) \le -A$ on the interval $(x, cx)$. The author observes that Ingham's refutation of the weak Mertens conjecture under RH and LI can be adapted to give the analogous results for $E^\pi(x;q,a)$ and $E^\pi(x;q,a,b)$; he reproves these results, assuming GRH but replacing LI by a condition on the intersection of the appropriate infinite subtorus and the diagonal subtorus. Finally, the author establishes an interesting connection between the range of $E^\psi(x;q,a)$ and the range of $x^{1/2}\Delta_r^\psi(x;q,a^{-1})$.\par This article cites~\cite{MR654520, MR0330069, MR0332684, MR0104638, 1936.Ingham, MR0006202, MR0352026, MR1083000, MR783538, MR0289452}.}
      \field{issn}{0065-1036}
      \field{journaltitle}{Acta Arith.}
      \field{number}{3}
      \field{title}{The {$k$}-functions in multiplicative number theory. {IV}. {O}n a method of {A}. {E}. {I}ngham}
      \field{volume}{57}
      \field{year}{1991}
      \true{nocite}
      \field{pages}{231\bibrangedash 244}
      \range{pages}{14}
      \verb{doi}
      \verb 10.4064/aa-57-3-231-244
      \endverb
      \verb{mrnumber}
      \verb 1105608
      \endverb
      \verb{urlraw}
      \verb https://doi.org/10.4064/aa-57-3-231-244
      \endverb
      \verb{url}
      \verb https://doi.org/10.4064/aa-57-3-231-244
      \endverb
    \endentry
    \entry{MR1133236}{article}{}
      \name{author}{1}{}{%
        {{hash=d1e4ad71dae945a6977ecf6f9ca65c30}{%
           family={Kaczorowski},
           familyi={K\bibinitperiod},
           given={J.},
           giveni={J\bibinitperiod}}}%
      }
      \strng{namehash}{d1e4ad71dae945a6977ecf6f9ca65c30}
      \strng{fullhash}{d1e4ad71dae945a6977ecf6f9ca65c30}
      \strng{bibnamehash}{d1e4ad71dae945a6977ecf6f9ca65c30}
      \strng{authorbibnamehash}{d1e4ad71dae945a6977ecf6f9ca65c30}
      \strng{authornamehash}{d1e4ad71dae945a6977ecf6f9ca65c30}
      \strng{authorfullhash}{d1e4ad71dae945a6977ecf6f9ca65c30}
      \field{extraname}{9}
      \field{sortinit}{1}
      \field{sortinithash}{50c6687d7fc80f50136d75228e3c59ba}
      \field{labelnamesource}{author}
      \field{labeltitlesource}{title}
      \field{annotation}{Assuming GRH and HC, the author shows that \begin{align*} \liminf_{T \to \infty} \frac{W_{q, a}^{\psi} (T)}{\log T} \geq \frac{\gamma_0}{\pi}+ \kappa, \end{align*} where $\gamma$ is the lowest ``uncancelled'' zero of Dirichlet $L$-functions modulo~$q$---more precisely, the minimal $\gamma>0$ such that $\frac12+i \gamma$ is a pole of $\sum_{\chi \mod q} \overline\chi(a) \frac{L'}L(s, \chi)$---and where $\kappa=\kappa(q,a)$ is a nonnegative number, defined as the density of zeros of a certain linear combination of $K$-functions. When $q=a=1$, this result implies that \begin{align*} \liminf_{T \to \infty} \frac{W^{\psi} (T)}{\log T} \geq \frac{\gamma_1}{\pi}+10^{-250}, \end{align*} with $\gamma_1=14.135\ldots$; this result breaks the barrier $\frac{\gamma_1}{\pi}$ that had been achieved in prior results.\par This article cites~\cites{MR1083000, MR1105608, MR777013, MR791085, MR927376, MR858395, MR891057}.}
      \field{issn}{0065-1036}
      \field{journaltitle}{Acta Arith.}
      \field{number}{1}
      \field{title}{The {$k$}-functions in multiplicative number theory. {V}. {C}hanges of sign of some arithmetical error terms}
      \field{volume}{59}
      \field{year}{1991}
      \true{nocite}
      \field{pages}{37\bibrangedash 58}
      \range{pages}{22}
      \verb{doi}
      \verb 10.4064/aa-59-1-37-58
      \endverb
      \verb{mrnumber}
      \verb 1133236
      \endverb
      \verb{urlraw}
      \verb https://doi.org/10.4064/aa-59-1-37-58
      \endverb
      \verb{url}
      \verb https://doi.org/10.4064/aa-59-1-37-58
      \endverb
    \endentry
    \entry{MR1159354}{article}{}
      \name{author}{1}{}{%
        {{hash=5895b52bfd9795ff19a630ff28e5fe17}{%
           family={Heath-Brown},
           familyi={H\bibinithyphendelim B\bibinitperiod},
           given={D.\bibnamedelimi R.},
           giveni={D\bibinitperiod\bibinitdelim R\bibinitperiod}}}%
      }
      \strng{namehash}{5895b52bfd9795ff19a630ff28e5fe17}
      \strng{fullhash}{5895b52bfd9795ff19a630ff28e5fe17}
      \strng{bibnamehash}{5895b52bfd9795ff19a630ff28e5fe17}
      \strng{authorbibnamehash}{5895b52bfd9795ff19a630ff28e5fe17}
      \strng{authornamehash}{5895b52bfd9795ff19a630ff28e5fe17}
      \strng{authorfullhash}{5895b52bfd9795ff19a630ff28e5fe17}
      \field{sortinit}{1}
      \field{sortinithash}{50c6687d7fc80f50136d75228e3c59ba}
      \field{labelnamesource}{author}
      \field{labeltitlesource}{title}
      \field{annotation}{Let $a_1(x),a_2(x),\ldots$ be continuous real-valued functions of period $1$, and let $\gamma_1, \gamma_2,\ldots$ be nonzero constants such that $\lim_{N\to\infty}\limsup_{T\to\infty} \frac{1}{T}\int_0^T \min\{1,|F(x)-\sum_{n\leq N}a_n(\gamma_nx)|\} \,dx=0$; the author shows that the mean value $\frac1T \int_0^T p(F(t))\,dt$ converges to a limit as $T\to\infty$ for every continuous, integrable, piecewise differentiable function $p$ for which $\hat p$ is also integrable. Under further hypotheses, including the linear independence of the $\gamma_j$ over $\Q$, the author shows that $F(x)$ has a limiting distribution function. When such a limiting distribution function exists, the author establishes a necessary condition for certain normalized moments of $F(x)$ to converge in the limit. The author applies these theorems when $F(x)$ is the error term in: the classical divisor problem, the circle problem, the Piltz divisor problem for $\tau_3$, or the second moment for $\zeta(s)$ on the critical line. An important lemma is the evaluation, for continuous functions $b_1(t),\ldots,b_k(t)$ of period $1$ from $\R$ to $\C$, of $\lim_{T\to\infty} \frac{1}{T}\int_0^T b_1(\gamma_1 t)\ldots b_k(\gamma_k t)\,dt$.\par This article cites~\cites{MR1544509,MR0882550}.}
      \field{issn}{0065-1036}
      \field{journaltitle}{Acta Arith.}
      \field{number}{4}
      \field{title}{The distribution and moments of the error term in the {D}irichlet divisor problem}
      \field{volume}{60}
      \field{year}{1992}
      \true{nocite}
      \field{pages}{389\bibrangedash 415}
      \range{pages}{27}
      \verb{doi}
      \verb 10.4064/aa-60-4-389-415
      \endverb
      \verb{mrnumber}
      \verb 1159354
      \endverb
      \verb{urlraw}
      \verb https://doi.org/10.4064/aa-60-4-389-415
      \endverb
      \verb{url}
      \verb https://doi.org/10.4064/aa-60-4-389-415
      \endverb
    \endentry
    \entry{MR1251926}{article}{}
      \name{author}{1}{}{%
        {{hash=d1e4ad71dae945a6977ecf6f9ca65c30}{%
           family={Kaczorowski},
           familyi={K\bibinitperiod},
           given={J.},
           giveni={J\bibinitperiod}}}%
      }
      \strng{namehash}{d1e4ad71dae945a6977ecf6f9ca65c30}
      \strng{fullhash}{d1e4ad71dae945a6977ecf6f9ca65c30}
      \strng{bibnamehash}{d1e4ad71dae945a6977ecf6f9ca65c30}
      \strng{authorbibnamehash}{d1e4ad71dae945a6977ecf6f9ca65c30}
      \strng{authornamehash}{d1e4ad71dae945a6977ecf6f9ca65c30}
      \strng{authorfullhash}{d1e4ad71dae945a6977ecf6f9ca65c30}
      \field{extraname}{10}
      \field{sortinit}{1}
      \field{sortinithash}{50c6687d7fc80f50136d75228e3c59ba}
      \field{labelnamesource}{author}
      \field{labeltitlesource}{title}
      \field{annotation}{The author shows, assuming GRH, that there are infinitely many integers $m$ such that $\pi(m;q,1)$ is larger than any other $\pi(m;q,a)$, as well as infinitely many integers such that $\pi(m;q,1)$ is the smallest. Indeed, the sets of such integers have positive lower density, even if a stronger inequality is demanded: the author proves that for any positive real number~$u$, there exist constants $c>1$ and $b>0$ such that for every $T\geq 1$, \begin{align*} \#\big\{T\leq m\leq c T\colon E^\pi(m;q,1) - \max_{a\not\equiv 1\mod q} E^\pi(m;q,a) \ge u \big\}\geq b T, \end{align*} and similarly for $E^\psi$; and the same result holds with${}\ge u$ replaced by${}\le-u$. The author uses his $k$-functions that appeared in prior work, as well as an examination of the boundary values of Dirichlet series.\par The author formulates the following ``Strong Race Hypothesis'': for every permutation of the set $\{a_1,a_2,\ldots,a_{\phi(q)}\}$ of reduced residue classes modulo~$q$, the set of integers $m$ such that \[ \pi(m;q,a_1)<\pi(m;q,a_2)<\cdots<\pi(m;q,a_{\phi(q)}) \] has positive lower density. \par This article cites~\cites{MR0108470, MR0146156, MR1083000, MR1105608}.}
      \field{issn}{0033-5606}
      \field{journaltitle}{Quart. J. Math. Oxford Ser. (2)}
      \field{number}{176}
      \field{title}{A contribution to the {S}hanks-{R}ényi race problem}
      \field{volume}{44}
      \field{year}{1993}
      \true{nocite}
      \field{pages}{451\bibrangedash 458}
      \range{pages}{8}
      \verb{doi}
      \verb 10.1093/qmath/44.4.451
      \endverb
      \verb{mrnumber}
      \verb 1251926
      \endverb
      \verb{urlraw}
      \verb https://doi.org/10.1093/qmath/44.4.451
      \endverb
      \verb{url}
      \verb https://doi.org/10.1093/qmath/44.4.451
      \endverb
    \endentry
    \entry{MR1193094}{article}{}
      \name{author}{1}{}{%
        {{hash=c5aa7061ae0ee03e7795baf500605b3a}{%
           family={Sankaranarayanan},
           familyi={S\bibinitperiod},
           given={A.},
           giveni={A\bibinitperiod}}}%
      }
      \strng{namehash}{c5aa7061ae0ee03e7795baf500605b3a}
      \strng{fullhash}{c5aa7061ae0ee03e7795baf500605b3a}
      \strng{bibnamehash}{c5aa7061ae0ee03e7795baf500605b3a}
      \strng{authorbibnamehash}{c5aa7061ae0ee03e7795baf500605b3a}
      \strng{authornamehash}{c5aa7061ae0ee03e7795baf500605b3a}
      \strng{authorfullhash}{c5aa7061ae0ee03e7795baf500605b3a}
      \field{sortinit}{1}
      \field{sortinithash}{50c6687d7fc80f50136d75228e3c59ba}
      \field{labelnamesource}{author}
      \field{labeltitlesource}{title}
      \field{annotation}{Using a method of Kaczorowski~\cite{MR777013}, the author shows that the number of sign changes of $\Delta^{Q_2}(x)$ for $2\le x\le T$ is $\gg \log T$ with an effective constant. The author claims the analogous result for $\Delta^{Q_k}$ (with the implicit constant depending on~$k$) for all $2 \le k\le 10^8$. \par This article cites~\cites{MR945273,MR891057,MR1029966,MR0973808A,MR0973808B, MR777013}.}
      \field{issn}{0003-889X}
      \field{journaltitle}{Arch. Math. (Basel)}
      \field{number}{1}
      \field{title}{On the sign changes in the remainder term of an asymptotic formula for the number of squarefree numbers}
      \field{volume}{60}
      \field{year}{1993}
      \true{nocite}
      \field{pages}{51\bibrangedash 57}
      \range{pages}{7}
      \verb{doi}
      \verb 10.1007/BF01194239
      \endverb
      \verb{mrnumber}
      \verb 1193094
      \endverb
    \endentry
    \entry{MR1256149}{article}{}
      \name{author}{1}{}{%
        {{hash=6fa84c8cc1e455847a05b9946a1bd205}{%
           family={Kaczorowski},
           familyi={K\bibinitperiod},
           given={Jerzy},
           giveni={J\bibinitperiod}}}%
      }
      \strng{namehash}{6fa84c8cc1e455847a05b9946a1bd205}
      \strng{fullhash}{6fa84c8cc1e455847a05b9946a1bd205}
      \strng{bibnamehash}{6fa84c8cc1e455847a05b9946a1bd205}
      \strng{authorbibnamehash}{6fa84c8cc1e455847a05b9946a1bd205}
      \strng{authornamehash}{6fa84c8cc1e455847a05b9946a1bd205}
      \strng{authorfullhash}{6fa84c8cc1e455847a05b9946a1bd205}
      \field{extraname}{11}
      \field{sortinit}{1}
      \field{sortinithash}{50c6687d7fc80f50136d75228e3c59ba}
      \field{labelnamesource}{author}
      \field{labeltitlesource}{title}
      \field{annotation}{In this article, the author uses $k$-functions to prove several results. Assuming RH, he shows that the set $\{ m\in\N\colon E^\pi(x) > u \}$ has positive lower density for any $u\in\R$, and similarly with $E^\pi(x)$ replaced by $-E^\pi(x)$, $E^\psi(x)$, $-E^\psi(x)$; assuming GRH, $E^\pi(x)$ may be further replaced by $E^\pi(x;q,1,\max)$ and $-E^\pi(x;q,1,\min)$ and their counterparts with~$\psi$. Unconditionally he gives analogous weaker density statements such as $\# \{ m\le M\colon E^\pi(x) > u \} = \Omega(M^{1-\ep})$. All these results are ineffective, but the author provides effective counterparts such as $\liminf_{N\to\infty} \sum_{k=1}^{N} 2^{-k} \# \{2^k \leq m < 2^{k+1}\colon E^\pi(x) > u \} > 0$.\par This article cites~\cite{MR0476616, MR0476615, 1936.Ingham, MR1083000, MR1105608, MR1251926, MR1310738, MR1256149, MR0146156, MR1511360, 1914.Littlewood, MR0087685, MR0108470}.}
      \field{issn}{0075-4102}
      \field{journaltitle}{J. Reine Angew. Math.}
      \field{title}{Results on the distribution of primes}
      \field{volume}{446}
      \field{year}{1994}
      \true{nocite}
      \field{pages}{89\bibrangedash 113}
      \range{pages}{25}
      \verb{doi}
      \verb 10.1515/crll.1994.446.89
      \endverb
      \verb{mrnumber}
      \verb 1256149
      \endverb
    \endentry
    \entry{MR1297543}{book}{}
      \name{author}{1}{}{%
        {{hash=d6fc16fe2b971a8009f6a09376b27581}{%
           family={Montgomery},
           familyi={M\bibinitperiod},
           given={Hugh\bibnamedelima L.},
           giveni={H\bibinitperiod\bibinitdelim L\bibinitperiod}}}%
      }
      \list{publisher}{1}{%
        {Published for the Conference Board of the Mathematical Sciences, Washington, DC; by the American Mathematical Society, Providence, RI}%
      }
      \strng{namehash}{d6fc16fe2b971a8009f6a09376b27581}
      \strng{fullhash}{d6fc16fe2b971a8009f6a09376b27581}
      \strng{bibnamehash}{d6fc16fe2b971a8009f6a09376b27581}
      \strng{authorbibnamehash}{d6fc16fe2b971a8009f6a09376b27581}
      \strng{authornamehash}{d6fc16fe2b971a8009f6a09376b27581}
      \strng{authorfullhash}{d6fc16fe2b971a8009f6a09376b27581}
      \field{extraname}{2}
      \field{sortinit}{1}
      \field{sortinithash}{50c6687d7fc80f50136d75228e3c59ba}
      \field{labelnamesource}{author}
      \field{labeltitlesource}{title}
      \field{annotation}{In Chapter~$5$ the author gives proofs of the first and second main theorems of the power-sum method using duality, and improves the first main theorem for longer ranges of~$v$ to $\max_{1\leq v\leq N^2} |s_v|\gg {|s_0|}/{N}$. The author also deduces Fabry's gap theorem from the first main theorem, and gives improved lower bounds for the specific cases where all the coefficients~$b_j$ are all nonnegative or all equal to~$1$.\par This chapter cites~\cites{MR0133309,MR0133308,MR1575079,MR820241,MR749389}.}
      \field{isbn}{0-8218-0737-4}
      \field{series}{CBMS Regional Conference Series in Mathematics}
      \field{title}{Ten lectures on the interface between analytic number theory and harmonic analysis}
      \field{volume}{84}
      \field{year}{1994}
      \true{nocite}
      \field{pages}{xiv+220}
      \range{pages}{-1}
      \verb{doi}
      \verb 10.1090/cbms/084
      \endverb
      \verb{mrnumber}
      \verb 1297543
      \endverb
    \endentry
    \entry{MR1329368}{article}{}
      \name{author}{2}{}{%
        {{hash=9073d8f33230938a75c73529bcd5e1ef}{%
           family={Rubinstein},
           familyi={R\bibinitperiod},
           given={M.},
           giveni={M\bibinitperiod}}}%
        {{hash=ea78b04077f0596f70f50f0559b2a1f2}{%
           family={Sarnak},
           familyi={S\bibinitperiod},
           given={P.},
           giveni={P\bibinitperiod}}}%
      }
      \strng{namehash}{78bb8618e127f18952100dffae455403}
      \strng{fullhash}{78bb8618e127f18952100dffae455403}
      \strng{bibnamehash}{78bb8618e127f18952100dffae455403}
      \strng{authorbibnamehash}{78bb8618e127f18952100dffae455403}
      \strng{authornamehash}{78bb8618e127f18952100dffae455403}
      \strng{authorfullhash}{78bb8618e127f18952100dffae455403}
      \field{sortinit}{1}
      \field{sortinithash}{50c6687d7fc80f50136d75228e3c59ba}
      \field{labelnamesource}{author}
      \field{labeltitlesource}{title}
      \field{annotation}{This is the article that really placed in central roles the logarithmic limiting distributions and logarithmic densities of prime number races.\par Assuming GRH: the authors show that \begin{align*} E_{q;a_1,\ldots,a_r}(x) = \frac{\log x}{\sqrt x} \big( \phi(q) \pi(x;q,a_1) - \pi(x), \ldots, \phi(q) \pi(x;q,a_r) - \pi(x) \big) \end{align*} has a limiting logarithmic distribution $\mu_{q;a_1,\ldots,a_r}$ on $\mathbb R^r$. They give an exponential upper bound for the ``tail'' of $\mu_{q;a_1,\ldots,a_r}$ (that is, the mass assigned to the exterior of a large ball), as well as a doubly exponential lower bound for the portion of that tail lying in certain specific orthants. They note the analogous results for the race between $\pi(x)$ and $\Li(x)$, as well as for the race between $\pi(x;q,N)$ and $\pi(x;q,R)$; in these two-way races, it follows that $\underline\delta(\pi,\Li)$, $\underline\delta(\Li,\pi)$, $\underline\delta_{q;R,N}$, and $\underline\delta_{q;N,R}$ are strictly positive.\par Assuming GRH and LI: they give the formula for the Fourier transform of $\mu_{q;a_1,\ldots,a_r}$. From it they deduce that the densities $\delta_{q;a_1,\ldots,a_r}$ exist and are strictly positive. They characterize the races (all with $r\le 3$) for which $\mu_{q;a_1,\ldots ,a_r}$ is symmetric under all permutations of the coordinates. They show that $\delta_{q;a_1,\ldots ,a_r}$ tends to $1/r!$ as $q$ tends to infinity, and establish a central limit theorem for $E_{q; N,R}(x)/\sqrt{\log q}$. They also compute $\delta(\Li,\pi)$, and $\delta_{q;N,R}$ for $q\in\{3,4,5,7,11,13\}$, to several decimal places.\par This article cites~\cites{1914.Littlewood,MR1573970,MR0004255,MR0083001,MR0108470,MR0146156,MR0476616,MR866118,1938.Wintner,1932.Ingham,MR0506080,MR1159354}.}
      \field{issn}{1058-6458}
      \field{journaltitle}{Experiment. Math.}
      \field{number}{3}
      \field{title}{Chebyshev's bias}
      \field{volume}{3}
      \field{year}{1994}
      \true{nocite}
      \field{pages}{173\bibrangedash 197}
      \range{pages}{25}
      \verb{mrnumber}
      \verb 1329368
      \endverb
      \verb{urlraw}
      \verb http://projecteuclid.org/euclid.em/1048515870
      \endverb
      \verb{url}
      \verb http://projecteuclid.org/euclid.em/1048515870
      \endverb
    \endentry
    \entry{MR1344249}{article}{}
      \name{author}{1}{}{%
        {{hash=6fa84c8cc1e455847a05b9946a1bd205}{%
           family={Kaczorowski},
           familyi={K\bibinitperiod},
           given={Jerzy},
           giveni={J\bibinitperiod}}}%
      }
      \strng{namehash}{6fa84c8cc1e455847a05b9946a1bd205}
      \strng{fullhash}{6fa84c8cc1e455847a05b9946a1bd205}
      \strng{bibnamehash}{6fa84c8cc1e455847a05b9946a1bd205}
      \strng{authorbibnamehash}{6fa84c8cc1e455847a05b9946a1bd205}
      \strng{authornamehash}{6fa84c8cc1e455847a05b9946a1bd205}
      \strng{authorfullhash}{6fa84c8cc1e455847a05b9946a1bd205}
      \field{extraname}{12}
      \field{sortinit}{1}
      \field{sortinithash}{50c6687d7fc80f50136d75228e3c59ba}
      \field{labelnamesource}{author}
      \field{labeltitlesource}{title}
      \field{annotation}{The author examines upper and lower natural densities in the prime number race modulo~$4$. Assuming GRH for $L(s, \chi_{-4})$, he proves that $\mathfrak d^*(4;1,3) \geq 0.04054045$ and $\mathfrak d^*(4;3,1) \geq 0.99998936$, and that $\mathfrak d_*(4; 1,3) <0.0000106$ and $\mathfrak d_*(4; 3,1) < 0.9594595$.\par This article cites~\cite{MR0476615, 1936.Ingham, MR1083000, MR1105608, MR1251926, MR1310738, MR1256149, MR0146156, MR1511360, 1914.Littlewood, MR0087685, MR0108470}.}
      \field{issn}{0174-4747}
      \field{journaltitle}{Analysis}
      \field{number}{2}
      \field{title}{On the distribution of primes (mod $4$)}
      \field{volume}{15}
      \field{year}{1995}
      \true{nocite}
      \field{pages}{159\bibrangedash 171}
      \range{pages}{13}
      \verb{doi}
      \verb 10.1524/anly.1995.15.2.159
      \endverb
      \verb{mrnumber}
      \verb 1344249
      \endverb
    \endentry
    \entry{MR1310738}{article}{}
      \name{author}{1}{}{%
        {{hash=6fa84c8cc1e455847a05b9946a1bd205}{%
           family={Kaczorowski},
           familyi={K\bibinitperiod},
           given={Jerzy},
           giveni={J\bibinitperiod}}}%
      }
      \strng{namehash}{6fa84c8cc1e455847a05b9946a1bd205}
      \strng{fullhash}{6fa84c8cc1e455847a05b9946a1bd205}
      \strng{bibnamehash}{6fa84c8cc1e455847a05b9946a1bd205}
      \strng{authorbibnamehash}{6fa84c8cc1e455847a05b9946a1bd205}
      \strng{authornamehash}{6fa84c8cc1e455847a05b9946a1bd205}
      \strng{authorfullhash}{6fa84c8cc1e455847a05b9946a1bd205}
      \field{extraname}{13}
      \field{sortinit}{1}
      \field{sortinithash}{50c6687d7fc80f50136d75228e3c59ba}
      \field{labelnamesource}{author}
      \field{labeltitlesource}{title}
      \field{annotation}{Assuming GRH, the author shows that for any permutation $(a_1,a_2,a_3,a_4)$ of $(1,2,3,4)$, there exist constants $b>0$ and $c_0>1$ such that \begin{align*} \#\big\{T\le x\le c_0T \colon \psi(x;5,a_1)>\psi(x;5,a_2)>&\psi(x;5,a_3)>\psi(x;5,a_4), \\&\min_{1\le j\le 3}\big( \psi(x;5,a_j,a_{j+1}) \big) \geq b\sqrt x \big\} \gg T. \end{align*} The proof is another application of the author's theory of $k$-functions and involves explicit calculations using the Dirichlet $L$-functions\mod 5 and exponential sums corresponding to each permutation.\par This article cites~\cites{MR1083000, MR1251926, MR0146156}.}
      \field{issn}{0022-314X}
      \field{journaltitle}{J. Number Theory}
      \field{number}{1}
      \field{title}{On the {S}hanks-{R}ényi race problem mod {$5$}}
      \field{volume}{50}
      \field{year}{1995}
      \true{nocite}
      \field{pages}{106\bibrangedash 118}
      \range{pages}{13}
      \verb{doi}
      \verb 10.1006/jnth.1995.1006
      \endverb
      \verb{mrnumber}
      \verb 1310738
      \endverb
      \verb{urlraw}
      \verb https://doi.org/10.1006/jnth.1995.1006
      \endverb
      \verb{url}
      \verb https://doi.org/10.1006/jnth.1995.1006
      \endverb
    \endentry
    \entry{MR1367576}{article}{}
      \name{author}{1}{}{%
        {{hash=d1e4ad71dae945a6977ecf6f9ca65c30}{%
           family={Kaczorowski},
           familyi={K\bibinitperiod},
           given={J.},
           giveni={J\bibinitperiod}}}%
      }
      \strng{namehash}{d1e4ad71dae945a6977ecf6f9ca65c30}
      \strng{fullhash}{d1e4ad71dae945a6977ecf6f9ca65c30}
      \strng{bibnamehash}{d1e4ad71dae945a6977ecf6f9ca65c30}
      \strng{authorbibnamehash}{d1e4ad71dae945a6977ecf6f9ca65c30}
      \strng{authornamehash}{d1e4ad71dae945a6977ecf6f9ca65c30}
      \strng{authorfullhash}{d1e4ad71dae945a6977ecf6f9ca65c30}
      \field{extraname}{14}
      \field{sortinit}{1}
      \field{sortinithash}{50c6687d7fc80f50136d75228e3c59ba}
      \field{labelnamesource}{author}
      \field{labeltitlesource}{title}
      \field{annotation}{The author's earlier result~\cite{MR1251926} implies that there exists a permutation $\{\sigma_j\}$ of the reduced residue classes\mod q that begins with $1$ such that the set of $x$ for which the $\pi(x;q,\sigma_j)$ are in the given order has positive lower density. In this article, he provides a method for computing explicit permutations with this property. As an application, he gives permutations for each prime modulus $\le 29$ that satisfy these conditions and therefore provably occur with positive lower density; for example, modulo $13$ he provides the permutation $(1,7,8,9,2,6,12,10,11,5,3,4)$ with this property. Similar results apply to permutations with $1$ in last place, and to the functions $\psi(x;q,\sigma_j)$.\par This article cites~\cites{MR0146156, MR1083000, MR1105608,MR1251926, MR1256149, MR1310738}.}
      \field{issn}{0065-1036}
      \field{journaltitle}{Acta Arith.}
      \field{number}{1}
      \field{title}{On the {S}hanks-{R}ényi race problem}
      \field{volume}{74}
      \field{year}{1996}
      \true{nocite}
      \field{pages}{31\bibrangedash 46}
      \range{pages}{16}
      \verb{doi}
      \verb 10.4064/aa-74-1-31-46
      \endverb
      \verb{mrnumber}
      \verb 1367576
      \endverb
      \verb{urlraw}
      \verb https://doi.org/10.4064/aa-74-1-31-46
      \endverb
      \verb{url}
      \verb https://doi.org/10.4064/aa-74-1-31-46
      \endverb
    \endentry
    \entry{MR1440155}{article}{}
      \name{author}{2}{}{%
        {{hash=a5c0de227f384c5812b85df50f0ac6e8}{%
           family={Stechkin},
           familyi={S\bibinitperiod},
           given={S.\bibnamedelimi B.},
           giveni={S\bibinitperiod\bibinitdelim B\bibinitperiod}}}%
        {{hash=548c78bd85b3dd0d14bb46a59f54e09e}{%
           family={Popov},
           familyi={P\bibinitperiod},
           given={A.\bibnamedelimi Yu.},
           giveni={A\bibinitperiod\bibinitdelim Y\bibinitperiod}}}%
      }
      \strng{namehash}{78455df493ce58458371e5d01ed8136b}
      \strng{fullhash}{78455df493ce58458371e5d01ed8136b}
      \strng{bibnamehash}{78455df493ce58458371e5d01ed8136b}
      \strng{authorbibnamehash}{78455df493ce58458371e5d01ed8136b}
      \strng{authornamehash}{78455df493ce58458371e5d01ed8136b}
      \strng{authorfullhash}{78455df493ce58458371e5d01ed8136b}
      \field{sortinit}{1}
      \field{sortinithash}{50c6687d7fc80f50136d75228e3c59ba}
      \field{labelnamesource}{author}
      \field{labeltitlesource}{title}
      \field{annotation}{The authors prove that when~$x$ sufficiently large, \[ \int_x^{2x} |\Delta^\psi(u)| \,du \ge \frac{x^{3/2}}{200} \quad\text{and}\quad \int_x^{2x} |\Delta^\pi(u)| \,du \ge \frac{x^{3/2}}{\log x}, \] and also that there exists a constant $A>1$ such that \[ \int_x^{Ax} \max\{\Delta^\psi(u)\} \,du \ge x^{3/2} \quad\text{and}\quad \int_x^{Ax} \min\{\Delta^\psi(u)\} \,du \le -x^{3/2}. \] Assuming RH, these latter inequalities hold with $A=212$. If RH is false, however, then the constant~$A$ (as well as the constant implied by ``sufficiently large'') is ineffective.\par This article cites~\cites{MR1511206, MR1544509}.}
      \field{issn}{0042-1316}
      \field{journaltitle}{Uspekhi Mat. Nauk}
      \field{number}{6(312)}
      \field{title}{Asymptotic distribution of prime numbers in the mean}
      \field{volume}{51}
      \field{year}{1996}
      \true{nocite}
      \field{pages}{21\bibrangedash 88}
      \range{pages}{68}
      \verb{doi}
      \verb 10.1070/RM1996v051n06ABEH003000
      \endverb
      \verb{mrnumber}
      \verb 1440155
      \endverb
      \verb{urlraw}
      \verb https://doi.org/10.1070/RM1996v051n06ABEH003000
      \endverb
      \verb{url}
      \verb https://doi.org/10.1070/RM1996v051n06ABEH003000
      \endverb
    \endentry
    \entry{MR1645811}{incollection}{}
      \name{author}{1}{}{%
        {{hash=6fa84c8cc1e455847a05b9946a1bd205}{%
           family={Kaczorowski},
           familyi={K\bibinitperiod},
           given={Jerzy},
           giveni={J\bibinitperiod}}}%
      }
      \list{publisher}{1}{%
        {Birkhäuser, Basel}%
      }
      \strng{namehash}{6fa84c8cc1e455847a05b9946a1bd205}
      \strng{fullhash}{6fa84c8cc1e455847a05b9946a1bd205}
      \strng{bibnamehash}{6fa84c8cc1e455847a05b9946a1bd205}
      \strng{authorbibnamehash}{6fa84c8cc1e455847a05b9946a1bd205}
      \strng{authornamehash}{6fa84c8cc1e455847a05b9946a1bd205}
      \strng{authorfullhash}{6fa84c8cc1e455847a05b9946a1bd205}
      \field{extraname}{15}
      \field{sortinit}{1}
      \field{sortinithash}{50c6687d7fc80f50136d75228e3c59ba}
      \field{labelnamesource}{author}
      \field{labeltitlesource}{title}
      \field{annotation}{From the Math Review by M.~Jutila: ``Various aspects and problems as well as the history of comparative prime number theory are surveyed, including recent important work by the author [using $K$-functions]. The central topic is the comparative study of the frequency of primes in different arithmetic progressions. The article ends with a list of open problems and an extensive bibliography with 30 references.''\par This article cites~\cites{MR0476615,MR1159354,MR0506080,MR777013,MR791085,MR927376,MR945273,MR1083000,MR1105608,MR1133236,MR1251926,MR1310738,MR1256149,MR1344249,MR858395,MR891057,MR0146156,MR0162771,MR1544293A,MR1544293B,MR783538,MR634679,MR1329368,MR0137689,MR0004255}.}
      \field{bookpagination}{none}
      \field{booktitle}{European {C}ongress of {M}athematics, {V}ol. {I} ({B}udapest, 1996)}
      \field{series}{Progr. Math.}
      \field{title}{Boundary values of {D}irichlet series and the distribution of primes}
      \field{volume}{168}
      \field{year}{1998}
      \true{nocite}
      \field{pages}{237\bibrangedash 254}
      \range{pages}{18}
      \verb{mrnumber}
      \verb 1645811
      \endverb
    \endentry
    \entry{MR1651741}{article}{}
      \name{author}{2}{}{%
        {{hash=e542667d9161d12cece57b09b9c6c693}{%
           family={Bays},
           familyi={B\bibinitperiod},
           given={C.},
           giveni={C\bibinitperiod}}}%
        {{hash=08e54e6381028d2260f342dd1c27f66e}{%
           family={Hudson},
           familyi={H\bibinitperiod},
           given={R.\bibnamedelimi H.},
           giveni={R\bibinitperiod\bibinitdelim H\bibinitperiod}}}%
      }
      \strng{namehash}{121290fcdac2ab7b7c2b8152b61c32dc}
      \strng{fullhash}{121290fcdac2ab7b7c2b8152b61c32dc}
      \strng{bibnamehash}{121290fcdac2ab7b7c2b8152b61c32dc}
      \strng{authorbibnamehash}{121290fcdac2ab7b7c2b8152b61c32dc}
      \strng{authornamehash}{121290fcdac2ab7b7c2b8152b61c32dc}
      \strng{authorfullhash}{121290fcdac2ab7b7c2b8152b61c32dc}
      \field{extraname}{6}
      \field{sortinit}{2}
      \field{sortinithash}{ed39bb39cf854d5250e95b1c1f94f4ed}
      \field{labelnamesource}{author}
      \field{labeltitlesource}{title}
      \field{annotation}{The authors describe computations of both $\pi(x;4,3,1)$ for values of $x$ past $10^{12}$, as well as its estimate using a truncated explicit formula with 12,000 zeros of $L(s,\chi{-4})$ for values of $x$ up to $10^{1000}$. The estimate duplicates the true distribution with satisfying accuracy, rediscovers all known axis-crossing regions, and finds probable new axis-crossing regions. The method extends to other differences such as $\Delta^\pi(x)$ and $\pi(x;q,\cN,\cR)$.\par This article cites~\cites{MR0476616, MR529694, MR1651741, MR0447090, MR1310738, MR0083001, MR1329368, MR0108470}.}
      \field{issn}{0025-5718}
      \field{journaltitle}{Math. Comp.}
      \field{number}{230}
      \field{title}{Zeroes of {D}irichlet {$L$}-functions and irregularities in the distribution of primes}
      \field{volume}{69}
      \field{year}{2000}
      \true{nocite}
      \field{pages}{861\bibrangedash 866}
      \range{pages}{6}
      \verb{doi}
      \verb 10.1090/S0025-5718-99-01105-9
      \endverb
      \verb{mrnumber}
      \verb 1651741
      \endverb
      \verb{urlraw}
      \verb https://doi.org/10.1090/S0025-5718-99-01105-9
      \endverb
      \verb{url}
      \verb https://doi.org/10.1090/S0025-5718-99-01105-9
      \endverb
    \endentry
    \entry{MR1752093}{article}{}
      \name{author}{2}{}{%
        {{hash=dcdb9f9c423bbcbfa2a60d519ccd6915}{%
           family={Bays},
           familyi={B\bibinitperiod},
           given={Carter},
           giveni={C\bibinitperiod}}}%
        {{hash=9e0609c9931c1b620df97e8c58d07532}{%
           family={Hudson},
           familyi={H\bibinitperiod},
           given={Richard\bibnamedelima H.},
           giveni={R\bibinitperiod\bibinitdelim H\bibinitperiod}}}%
      }
      \strng{namehash}{0824391c51610edfe4bd3c404488dea5}
      \strng{fullhash}{0824391c51610edfe4bd3c404488dea5}
      \strng{bibnamehash}{0824391c51610edfe4bd3c404488dea5}
      \strng{authorbibnamehash}{0824391c51610edfe4bd3c404488dea5}
      \strng{authornamehash}{0824391c51610edfe4bd3c404488dea5}
      \strng{authorfullhash}{0824391c51610edfe4bd3c404488dea5}
      \field{extraname}{7}
      \field{sortinit}{2}
      \field{sortinithash}{ed39bb39cf854d5250e95b1c1f94f4ed}
      \field{labelnamesource}{author}
      \field{labeltitlesource}{title}
      \field{annotation}{Using Lehman's theorem~\cite{MR0202686} together with the first $10^6$ zeros of $\zeta(s)$ (supplied by Odlyzko), the authors show that there exist many integers~$x$ in the range $[\exp(727.95209-.002),\exp(727.95209+.002)]$ for which $\pi(x)>\li(x)$. (This corresponds to the interval $[1.3954272\times10^{316}, 1.4010201\times10^{316})$, although the authors claim that the interval is $[1.398201\times10^{316}, 1.398244\times10^{316}]$.) They also report on computations approximating $\pi(x)-\li(x)$ for $x$ up to about $4\times10^{12370}$ (the significance being that the 20th region with many solutions to $\pi(x)>\li(x)$ is probably around this number) that support the conditional result of Rubinstein and Sarnak that $\delta(\li,\pi) \approx 0.99999974$.\par This article cites~\cites{MR1651741,MR1816036,1932.Ingham,MR0202686,1914.Littlewood,MR866118,MR1329368,MR1573970,MR0067145}.}
      \field{issn}{0025-5718}
      \field{journaltitle}{Math. Comp.}
      \field{number}{231}
      \field{title}{A new bound for the smallest {$x$} with {$\pi (x)>\li (x)$}}
      \field{volume}{69}
      \field{year}{2000}
      \true{nocite}
      \field{pages}{1285\bibrangedash 1296}
      \range{pages}{12}
      \verb{mrnumber}
      \verb 1752093
      \endverb
    \endentry
    \entry{MR1806291}{article}{}
      \name{author}{2}{}{%
        {{hash=0644385a54ab456bd85da1f460b152fd}{%
           family={Feuerverger},
           familyi={F\bibinitperiod},
           given={A.},
           giveni={A\bibinitperiod}}}%
        {{hash=feda7d3cba2a0ad4e772307fca655937}{%
           family={Martin},
           familyi={M\bibinitperiod},
           given={G.},
           giveni={G\bibinitperiod}}}%
      }
      \strng{namehash}{ad494ec99168c86b41ad9bf70473e1a4}
      \strng{fullhash}{ad494ec99168c86b41ad9bf70473e1a4}
      \strng{bibnamehash}{ad494ec99168c86b41ad9bf70473e1a4}
      \strng{authorbibnamehash}{ad494ec99168c86b41ad9bf70473e1a4}
      \strng{authornamehash}{ad494ec99168c86b41ad9bf70473e1a4}
      \strng{authorfullhash}{ad494ec99168c86b41ad9bf70473e1a4}
      \field{sortinit}{2}
      \field{sortinithash}{ed39bb39cf854d5250e95b1c1f94f4ed}
      \field{labelnamesource}{author}
      \field{labeltitlesource}{title}
      \field{annotation}{Assuming GRH and LI, the authors establish a general formula for $\delta_{q;a_1,\ldots,a_r}$ as a linear combination of principal values of integers of dimensions up to $r-1$. As an example, they numerically compute all possible race densities when $q\in\{5,8,12\}$, all densities for two- and three-way races when $q\in\{7,9\}$, and all two-way race densities when $q=11$. They also list some families of symmetries among these densities, as well as some families of provable inequalities for three-way race densities. The authors conjecture that ``two-ties'' configurations are possible in the following sense: for any $1\le i,j,k,\ell\le r$, there exist arbitrarily large $x$ such that $\pi(x;q,a_i) = \pi(x;q,a_j)$ and $\pi(x;q,a_k) = \pi(x;q,a_\ell)$ (even if the relative sizes of all $r$ contestants are prescribed); they further conjecture that analogous ``three-ties'' configurations occur only finitely often. They also make some speculative conjectures about ``bias factors'' and ``order equivalences'' to introduce the idea of comparing densities to one another; they also describe the question of the possible asymptotic sizes of $\delta_{q;a_1,\ldots,a_r}$ if $r$ is a function of~$q$. Finally, they remark that their techniques also resolve a question of Knapowski and Turán~\cite{MR0146156} about simultaneous solutions to $\Delta^\pi(x;q,a_j) > 0$.\par This article cites~\cites{MR1367576,MR0146156,1914.Littlewood,MR1329368,MR0108470}.}
      \field{issn}{1058-6458}
      \field{journaltitle}{Experiment. Math.}
      \field{number}{4}
      \field{title}{Biases in the {S}hanks-{R}ényi prime number race}
      \field{volume}{9}
      \field{year}{2000}
      \true{nocite}
      \field{pages}{535\bibrangedash 570}
      \range{pages}{36}
      \verb{mrnumber}
      \verb 1806291
      \endverb
      \verb{urlraw}
      \verb http://projecteuclid.org/euclid.em/1045759521
      \endverb
      \verb{url}
      \verb http://projecteuclid.org/euclid.em/1045759521
      \endverb
    \endentry
    \entry{MR1756780}{book}{}
      \name{author}{1}{}{%
        {{hash=8c77965dc9a09815d7992a73cfd267ae}{%
           family={Narkiewicz},
           familyi={N\bibinitperiod},
           given={Władysław},
           giveni={W\bibinitperiod}}}%
      }
      \list{publisher}{1}{%
        {Springer-Verlag, Berlin}%
      }
      \strng{namehash}{8c77965dc9a09815d7992a73cfd267ae}
      \strng{fullhash}{8c77965dc9a09815d7992a73cfd267ae}
      \strng{bibnamehash}{8c77965dc9a09815d7992a73cfd267ae}
      \strng{authorbibnamehash}{8c77965dc9a09815d7992a73cfd267ae}
      \strng{authornamehash}{8c77965dc9a09815d7992a73cfd267ae}
      \strng{authorfullhash}{8c77965dc9a09815d7992a73cfd267ae}
      \field{sortinit}{2}
      \field{sortinithash}{ed39bb39cf854d5250e95b1c1f94f4ed}
      \field{labelnamesource}{author}
      \field{labeltitlesource}{title}
      \field{annotation}{Section~6.6 is devoted to the sign of the difference $\pi(x)-\li(x)$. The author sketches the proof of Littlewood's theorem and gives a detailed survey of results about $W(T)$ and $W^{\psi}(T)$.\par This article cites~\cites{1914.Littlewood, 1930.Polya, 1936.Ingham,MR1511206, MR1573970, MR0067145, MR0133309, MR0133308, MR0202686, MR0233782, MR0439771, MR0439772, MR0439774, MR607089, MR0447151, MR866118}.}
      \field{isbn}{3-540-66289-8}
      \field{series}{Springer Monographs in Mathematics}
      \field{title}{The development of prime number theory}
      \field{year}{2000}
      \true{nocite}
      \field{pages}{xii+448}
      \range{pages}{-1}
      \verb{doi}
      \verb 10.1007/978-3-662-13157-2
      \endverb
      \verb{mrnumber}
      \verb 1756780
      \endverb
      \verb{urlraw}
      \verb https://doi.org/10.1007/978-3-662-13157-2
      \endverb
      \verb{url}
      \verb https://doi.org/10.1007/978-3-662-13157-2
      \endverb
    \endentry
    \entry{2000.Ng}{online}{}
      \name{author}{1}{}{%
        {{hash=585b691d1c27d2773434ed1ac9664075}{%
           family={Ng},
           familyi={N\bibinitperiod},
           given={N.},
           giveni={N\bibinitperiod}}}%
      }
      \strng{namehash}{585b691d1c27d2773434ed1ac9664075}
      \strng{fullhash}{585b691d1c27d2773434ed1ac9664075}
      \strng{bibnamehash}{585b691d1c27d2773434ed1ac9664075}
      \strng{authorbibnamehash}{585b691d1c27d2773434ed1ac9664075}
      \strng{authornamehash}{585b691d1c27d2773434ed1ac9664075}
      \strng{authorfullhash}{585b691d1c27d2773434ed1ac9664075}
      \field{extraname}{1}
      \field{sortinit}{2}
      \field{sortinithash}{ed39bb39cf854d5250e95b1c1f94f4ed}
      \field{labelnamesource}{author}
      \field{labeltitlesource}{title}
      \field{annotation}{In Section~2 of this thesis, the author gives examples of Artin $L$-functions for which LI is false. In Section~5.1, the author derives (under GRH) an explicit formula for $\frac{\log{x}}{\sqrt{x}}\left(\frac{|G|}{|C|}\pi_C(x)-\pi_K(x)\right)$ as $x\to\infty$, where $L/K$ is a normal extension of number fields, $G=\operatorname{Gal}(L/K)$ and $C$ is a conjugacy class of~$G$, and $\pi_K(x)=\#\{\mathfrak p\subset\mathcal{O}_K \colon \text{N}\mathfrak p\leq x\}$ and $\pi_C(x)=\#\{ \mathfrak p\subset\mathcal{O}_K \colon \text{N}\mathfrak p\leq x,\, \sigma_{\mathfrak p}=C\}$ where $\sigma_{\mathfrak p}$ is the Frobenius. Using the explicit formula, the author obtains analogues of Chebychev's bias in Galois groups. In Section~6.1, an explicit formula is given for $\frac{\log{x}}{\sqrt{x}}\pi(x,\mathfrak{a_1},\mathfrak{a_2})$, where $\mathfrak{a_1}$ and $\mathfrak{a_2}$ are ideal classes of a number field $K$. This formula is used to prove analogues of Chebychev's bias for primes in ideal classes. In Chapter 7, the author conditionally establishes the existence of a limiting distribution for $E^M(x)$ and conjectures that its maximal order is $(\log\log\log x)^{5/4}$.\par This article cites~ \cites{MR0330069, MR1816036, MR1544509, MR1806291, MR598886, MR0104638, MR1159354, MR0006202, MR634679, MR783538, MR1329368, MR0882550}}
      \field{note}{Thesis (Ph.D.)--University of British Columbia}
      \field{pagetotal}{161}
      \field{title}{{Limiting processes and Zeros of Artin L-Functions}}
      \field{year}{2000}
      \true{nocite}
      \verb{urlraw}
      \verb http://www.cs.uleth.ca/~nathanng/RESEARCH/phd.thesis.pdf
      \endverb
      \verb{url}
      \verb http://www.cs.uleth.ca/~nathanng/RESEARCH/phd.thesis.pdf
      \endverb
    \endentry
    \entry{MR1761276}{article}{}
      \name{author}{1}{}{%
        {{hash=2e461074b5178ec74d42e5ab1268d913}{%
           family={Puchta},
           familyi={P\bibinitperiod},
           given={J.-C.},
           giveni={J\bibinithyphendelim C\bibinitperiod}}}%
      }
      \strng{namehash}{2e461074b5178ec74d42e5ab1268d913}
      \strng{fullhash}{2e461074b5178ec74d42e5ab1268d913}
      \strng{bibnamehash}{2e461074b5178ec74d42e5ab1268d913}
      \strng{authorbibnamehash}{2e461074b5178ec74d42e5ab1268d913}
      \strng{authornamehash}{2e461074b5178ec74d42e5ab1268d913}
      \strng{authorfullhash}{2e461074b5178ec74d42e5ab1268d913}
      \field{sortinit}{2}
      \field{sortinithash}{ed39bb39cf854d5250e95b1c1f94f4ed}
      \field{labelnamesource}{author}
      \field{labeltitlesource}{title}
      \field{annotation}{Assuming GRH, the author shows that $E^\psi(x;q,1,\min) = \Omega_+(\log\log\log x)$ and $E^\psi(x;q,1,\max) = \Omega_-(\log\log\log x)$ where all constants are effective. In particular, the $q=1$ case gives an effective version of the classical result of Littlewood~\cites{1914.Littlewood}. One tool is an evaluation of the moments of $E^{\psi_\chi}(u)$ for Dirichlet characters~$\chi$: the author shows that $\frac{1}{y}\int_0^y E^{\psi_\chi}(u)^k \, \frac{du}u \to (-1)^k \sum_{\gamma_1 + \cdots + \gamma_k=0} {1/\rho_1 \cdots \rho_k}$, where the sum runs over all nontrivial zeros of $L(s,\chi)$.\par This article cites \cites{1914.Littlewood, MR0004255,MR1083000,MR1083001,MR1105605,MR1105608, MR1310738,MR1251926}.}
      \field{issn}{0236-5294}
      \field{journaltitle}{Acta Math. Hungar.}
      \field{number}{3}
      \field{title}{On large oscillations of the remainder of the prime number theorems}
      \field{volume}{87}
      \field{year}{2000}
      \true{nocite}
      \field{pages}{213\bibrangedash 227}
      \range{pages}{15}
      \verb{mrnumber}
      \verb 1761276
      \endverb
    \endentry
    \entry{MR1816036}{article}{}
      \name{author}{4}{}{%
        {{hash=e542667d9161d12cece57b09b9c6c693}{%
           family={Bays},
           familyi={B\bibinitperiod},
           given={C.},
           giveni={C\bibinitperiod}}}%
        {{hash=e7e31df1803902b56a78f32b1c733598}{%
           family={Ford},
           familyi={F\bibinitperiod},
           given={K.},
           giveni={K\bibinitperiod}}}%
        {{hash=08e54e6381028d2260f342dd1c27f66e}{%
           family={Hudson},
           familyi={H\bibinitperiod},
           given={R.\bibnamedelimi H.},
           giveni={R\bibinitperiod\bibinitdelim H\bibinitperiod}}}%
        {{hash=9073d8f33230938a75c73529bcd5e1ef}{%
           family={Rubinstein},
           familyi={R\bibinitperiod},
           given={M.},
           giveni={M\bibinitperiod}}}%
      }
      \strng{namehash}{0543a8fb7e634d40caee3a99730ca685}
      \strng{fullhash}{4815739125bfd0a086513390034cdd23}
      \strng{bibnamehash}{4815739125bfd0a086513390034cdd23}
      \strng{authorbibnamehash}{4815739125bfd0a086513390034cdd23}
      \strng{authornamehash}{0543a8fb7e634d40caee3a99730ca685}
      \strng{authorfullhash}{4815739125bfd0a086513390034cdd23}
      \field{sortinit}{2}
      \field{sortinithash}{ed39bb39cf854d5250e95b1c1f94f4ed}
      \field{labelnamesource}{author}
      \field{labeltitlesource}{title}
      \field{annotation}{This article is concerned with estimating the densities $\delta_{q,\cN,\cR}$, in particular by using the fact that the truncated explicit formula is an almost-periodic function whose largest ``quasi-period'' is dictated by its lowest zero. The authors present data supporting the fact that these densities depends strongly on the location of the first few zeros of $L(s,\chi_{\pm q})$ and the size of the first zero. Plots of $E(x;q,\cN,\cR)$ for all primes $q$ with $h(-q)=1$ are provided, and the difficulty of accurately estimating the densities from such data is pointed out. The authors note that the Chowla--Selberg formula implies that if $Q(\sqrt{-q} )$ is an imaginary quadratic field with class number~$1$ then $L(s,\chi_{-q})$ has a relatively low-lying zero; they also analyze the connection between low-lying zeros and class numbers $3$ and~$5$.\par This article cites~\cites{MR0476616, MR529694, MR1651741, 1853.Chebyshev, MR552467, MR1251926, MR1310738, MR0146156, MR0083001, 1914.Littlewood, MR1329368, MR0108470}.}
      \field{issn}{0022-314X}
      \field{journaltitle}{J. Number Theory}
      \field{number}{1}
      \field{title}{Zeros of {D}irichlet {$L$}-functions near the real axis and {C}hebyshev's bias}
      \field{volume}{87}
      \field{year}{2001}
      \true{nocite}
      \field{pages}{54\bibrangedash 76}
      \range{pages}{23}
      \verb{doi}
      \verb 10.1006/jnth.2000.2601
      \endverb
      \verb{mrnumber}
      \verb 1816036
      \endverb
      \verb{urlraw}
      \verb https://doi.org/10.1006/jnth.2000.2601
      \endverb
      \verb{url}
      \verb https://doi.org/10.1006/jnth.2000.2601
      \endverb
    \endentry
    \entry{MR1862054}{article}{}
      \name{author}{2}{}{%
        {{hash=bb758f27d76cefc7bd5244190e91843b}{%
           family={Ford},
           familyi={F\bibinitperiod},
           given={Kevin},
           giveni={K\bibinitperiod}}}%
        {{hash=9e0609c9931c1b620df97e8c58d07532}{%
           family={Hudson},
           familyi={H\bibinitperiod},
           given={Richard\bibnamedelima H.},
           giveni={R\bibinitperiod\bibinitdelim H\bibinitperiod}}}%
      }
      \strng{namehash}{3bfafe737308f5813c05132ae0dd9deb}
      \strng{fullhash}{3bfafe737308f5813c05132ae0dd9deb}
      \strng{bibnamehash}{3bfafe737308f5813c05132ae0dd9deb}
      \strng{authorbibnamehash}{3bfafe737308f5813c05132ae0dd9deb}
      \strng{authornamehash}{3bfafe737308f5813c05132ae0dd9deb}
      \strng{authorfullhash}{3bfafe737308f5813c05132ae0dd9deb}
      \field{sortinit}{2}
      \field{sortinithash}{ed39bb39cf854d5250e95b1c1f94f4ed}
      \field{labelnamesource}{author}
      \field{labeltitlesource}{title}
      \field{annotation}{The authors generalize Lehman's method~\cite{MR0202686} to arithmetic progressions, enabling them to give new results on the location of negative values of $\pi(x;q,b,1)$ for moduli $q \mid 24$. They show that each $\pi(x;8,b,1)$ is negative for some $x < 5 \cdot 10^{19}$, each $\pi(x;12,b,1)$ is negative for some $x <10^{84}$, and each $\pi(x;24,b,1)$ is negative for some $x<10^{253}$. Further, if GRH$(630,000)$ holds for $L(s,\chi_{-4})$, then $\pi(x;4,1,3) > \sqrt{x}/\log x$ for some~$x \approx 8\cdot10^{34,171}$. \par This article cites \cites{MR0202686, MR0447090, MR0476616, MR0472726, MR1752093,MR1816036,MR1251926,MR1256149,MR1367576,MR0146159,MR0068565,MR0083001,MR0137689,MR1329368,MR0108470,MR1573970, MR0067145}.}
      \field{issn}{0065-1036}
      \field{journaltitle}{Acta Arith.}
      \field{number}{4}
      \field{title}{Sign changes in {$\pi _{q,a}(x)-\pi _{q,b}(x)$}}
      \field{volume}{100}
      \field{year}{2001}
      \true{nocite}
      \field{pages}{297\bibrangedash 314}
      \range{pages}{18}
      \verb{doi}
      \verb 10.4064/aa100-4-1
      \endverb
      \verb{mrnumber}
      \verb 1862054
      \endverb
    \endentry
    \entry{MR1908877}{article}{}
      \name{author}{1}{}{%
        {{hash=b8f6e7b5ec0a55f3a2fbfd17aab7dfa9}{%
           family={Ruzsa},
           familyi={R\bibinitperiod},
           given={Imre\bibnamedelima Z.},
           giveni={I\bibinitperiod\bibinitdelim Z\bibinitperiod}}}%
      }
      \strng{namehash}{b8f6e7b5ec0a55f3a2fbfd17aab7dfa9}
      \strng{fullhash}{b8f6e7b5ec0a55f3a2fbfd17aab7dfa9}
      \strng{bibnamehash}{b8f6e7b5ec0a55f3a2fbfd17aab7dfa9}
      \strng{authorbibnamehash}{b8f6e7b5ec0a55f3a2fbfd17aab7dfa9}
      \strng{authornamehash}{b8f6e7b5ec0a55f3a2fbfd17aab7dfa9}
      \strng{authorfullhash}{b8f6e7b5ec0a55f3a2fbfd17aab7dfa9}
      \field{sortinit}{2}
      \field{sortinithash}{ed39bb39cf854d5250e95b1c1f94f4ed}
      \field{labelnamesource}{author}
      \field{labeltitlesource}{title}
      \field{annotation}{The author proves that the number of pairs of consecutive primes up to~$x$ that are both congruent to $1\mod4$ is $\gg x\log\log x/\log^2x$, improving a result of Shiu. A generalization holds where the single residue class $1\mod4$ is replaced by an arbitrary set of reduced residue classes modulo~$q$ of size $\phi(q)/2$. The proof uses Maier's method. \par This article cites \cites{MR0466043,MR1297543}.}
      \field{issn}{0019-3577}
      \field{journaltitle}{Indag. Math. (N.S.)}
      \field{number}{4}
      \field{title}{Consecutive primes modulo 4}
      \field{volume}{12}
      \field{year}{2001}
      \true{nocite}
      \field{pages}{489\bibrangedash 503}
      \range{pages}{15}
      \verb{doi}
      \verb 10.1016/S0019-3577(01)80038-0
      \endverb
      \verb{mrnumber}
      \verb 1908877
      \endverb
    \endentry
    \entry{MR1985941}{incollection}{}
      \name{author}{2}{}{%
        {{hash=e7e31df1803902b56a78f32b1c733598}{%
           family={Ford},
           familyi={F\bibinitperiod},
           given={K.},
           giveni={K\bibinitperiod}}}%
        {{hash=c48086b9b194fded979a790687062dbb}{%
           family={Konyagin},
           familyi={K\bibinitperiod},
           given={S.},
           giveni={S\bibinitperiod}}}%
      }
      \list{publisher}{1}{%
        {Mosk. Gos. Univ. im. Lomonosova, Mekh.-Mat. Fak., Moscow}%
      }
      \strng{namehash}{3cc0b23765677fa7230ff0cbba627702}
      \strng{fullhash}{3cc0b23765677fa7230ff0cbba627702}
      \strng{bibnamehash}{3cc0b23765677fa7230ff0cbba627702}
      \strng{authorbibnamehash}{3cc0b23765677fa7230ff0cbba627702}
      \strng{authornamehash}{3cc0b23765677fa7230ff0cbba627702}
      \strng{authorfullhash}{3cc0b23765677fa7230ff0cbba627702}
      \field{extraname}{1}
      \field{sortinit}{2}
      \field{sortinithash}{ed39bb39cf854d5250e95b1c1f94f4ed}
      \field{labelnamesource}{author}
      \field{labeltitlesource}{title}
      \field{annotation}{This article describes nine families of problems that are central to the study of comparative prime number theory. The first eight are taken from or inspired by the problems listed by Knapowski and Turán in~\cite{MR0146156}; the ninth problem, entitled ``Union-problems'', examines the distribution of $\frac1{\#\cA} \pi(x;q,\cA) - \frac1{\#\cB} \pi(x;q,\cB)$. Throughout the rest of the article, the authors provide an overview of what is already known about the first seven problems.\par This article cites~\cites{1853.Chebyshev, MR1511360, 1914.Littlewood, MR1555148, MR1544293A, MR1544293B, MR1573970, 1938.Wintner, MR0004255, MR0006202, MR0067145, MR0083001, MR0108470, MR0146156, MR0146157, MR0146158, MR0146159, MR0146160, MR0146161, MR0156826, MR0156827, MR0176967, MR0162771, MR0174538, MR0180539, MR0182616A, MR0182616B, MR0202686, MR0200250, MR0202685, MR0330069, MR0332684, MR0332685, MR0289452, MR0302585, MR0447090, MR0506080, MR0472726, MR0476615, MR529694, MR552467, MR634679, MR866118, MR1083000, MR1105608, MR1251926, MR1329368, MR1310738, MR1367576, MR1651741, MR1806291, MR1816036, MR1862054, MR1909220, MR1956261, MR1957110}.}
      \field{booktitle}{I{V} {I}nternational {C}onference ``{M}odern {P}roblems of {N}umber {T}heory and its {A}pplications'': {C}urrent {P}roblems, {P}art {II} ({R}ussian) ({T}ula, 2001)}
      \field{title}{Chebyshev's conjecture and the prime number race}
      \field{year}{2002}
      \true{nocite}
      \field{pages}{67\bibrangedash 91}
      \range{pages}{25}
      \verb{mrnumber}
      \verb 1985941
      \endverb
    \endentry
    \entry{MR1909220}{article}{}
      \name{author}{2}{}{%
        {{hash=e7e31df1803902b56a78f32b1c733598}{%
           family={Ford},
           familyi={F\bibinitperiod},
           given={K.},
           giveni={K\bibinitperiod}}}%
        {{hash=c48086b9b194fded979a790687062dbb}{%
           family={Konyagin},
           familyi={K\bibinitperiod},
           given={S.},
           giveni={S\bibinitperiod}}}%
      }
      \strng{namehash}{3cc0b23765677fa7230ff0cbba627702}
      \strng{fullhash}{3cc0b23765677fa7230ff0cbba627702}
      \strng{bibnamehash}{3cc0b23765677fa7230ff0cbba627702}
      \strng{authorbibnamehash}{3cc0b23765677fa7230ff0cbba627702}
      \strng{authornamehash}{3cc0b23765677fa7230ff0cbba627702}
      \strng{authorfullhash}{3cc0b23765677fa7230ff0cbba627702}
      \field{extraname}{2}
      \field{sortinit}{2}
      \field{sortinithash}{ed39bb39cf854d5250e95b1c1f94f4ed}
      \field{labelnamesource}{author}
      \field{labeltitlesource}{title}
      \field{annotation}{This article introduces the term ``barrier'' for a hypothetical configuration of zeros of Dirichlet $L$-functions that causes prime number races not to be exhaustive. The authors show (through several complementary constructions) that for every three-way prime number races, there is a finite barrier that prevents at least one of the six possible orderings from occurring for large~$x$; moreover, the zeros in these barriers can be arbitrarily close to the critical line and arbitrarily far from the real axis. While most of their constructions of barriers involve zeros with linearly dependent ordinates, the authors construct in the final section a barrier with linearly independent ordinates.\par This article cites~\cites{MR0476616,1853.Chebyshev,MR1251926,MR1310738,MR1367576,MR0146156,MR0146157,MR0146158,MR0146159,MR0146160,MR0146161,MR0156826,MR0156827,MR0162771,MR0174538,MR0180539,MR0182616A,MR0182616B,MR0200250,1914.Littlewood,MR1329368,MR0108470}.}
      \field{issn}{0012-7094}
      \field{journaltitle}{Duke Math. J.}
      \field{number}{2}
      \field{title}{The prime number race and zeros of {$L$}-functions off the critical line}
      \field{volume}{113}
      \field{year}{2002}
      \true{nocite}
      \field{pages}{313\bibrangedash 330}
      \range{pages}{18}
      \verb{doi}
      \verb 10.1215/S0012-7094-02-11324-6
      \endverb
      \verb{mrnumber}
      \verb 1909220
      \endverb
      \verb{urlraw}
      \verb https://doi.org/10.1215/S0012-7094-02-11324-6
      \endverb
      \verb{url}
      \verb https://doi.org/10.1215/S0012-7094-02-11324-6
      \endverb
    \endentry
    \entry{MR1916279}{article}{}
      \name{author}{1}{}{%
        {{hash=5cf7064e6cd6ca7d20eeb2d3e5e04811}{%
           family={Lau},
           familyi={L\bibinitperiod},
           given={Yuk-Kam},
           giveni={Y\bibinithyphendelim K\bibinitperiod}}}%
      }
      \strng{namehash}{5cf7064e6cd6ca7d20eeb2d3e5e04811}
      \strng{fullhash}{5cf7064e6cd6ca7d20eeb2d3e5e04811}
      \strng{bibnamehash}{5cf7064e6cd6ca7d20eeb2d3e5e04811}
      \strng{authorbibnamehash}{5cf7064e6cd6ca7d20eeb2d3e5e04811}
      \strng{authornamehash}{5cf7064e6cd6ca7d20eeb2d3e5e04811}
      \strng{authorfullhash}{5cf7064e6cd6ca7d20eeb2d3e5e04811}
      \field{sortinit}{2}
      \field{sortinithash}{ed39bb39cf854d5250e95b1c1f94f4ed}
      \field{labelnamesource}{author}
      \field{labeltitlesource}{title}
      \field{annotation}{This article investigates the limiting distribution of almost periodic functions. The author proves the existence of the limiting distribution of a class of functions which are bounded and can be approximated by periodic functions in $L_1$-norm. By using the quantitative version of the continuity theorem, the author is able to investigate the rate of convergence of some cases. Compared to Heath-Brown's work~\cites{MR1159354} on distribution of the error term in the Dirichlet divisor problem, the result here is more general and requires weaker hypotheses.\par This article cites ~\cites{MR1159354,MR1329368}.}
      \field{issn}{0022-314X}
      \field{journaltitle}{J. Number Theory}
      \field{number}{2}
      \field{title}{On the existence of limiting distributions of some number-theoretic error terms}
      \field{volume}{94}
      \field{year}{2002}
      \true{nocite}
      \field{pages}{359\bibrangedash 374}
      \range{pages}{16}
      \verb{doi}
      \verb 10.1006/jnth.2001.2734
      \endverb
      \verb{mrnumber}
      \verb 1916279
      \endverb
    \endentry
    \entry{MR1956261}{inproceedings}{}
      \name{author}{1}{}{%
        {{hash=feda7d3cba2a0ad4e772307fca655937}{%
           family={Martin},
           familyi={M\bibinitperiod},
           given={G.},
           giveni={G\bibinitperiod}}}%
      }
      \list{publisher}{1}{%
        {A K Peters, Natick, MA}%
      }
      \strng{namehash}{feda7d3cba2a0ad4e772307fca655937}
      \strng{fullhash}{feda7d3cba2a0ad4e772307fca655937}
      \strng{bibnamehash}{feda7d3cba2a0ad4e772307fca655937}
      \strng{authorbibnamehash}{feda7d3cba2a0ad4e772307fca655937}
      \strng{authornamehash}{feda7d3cba2a0ad4e772307fca655937}
      \strng{authorfullhash}{feda7d3cba2a0ad4e772307fca655937}
      \field{sortinit}{2}
      \field{sortinithash}{ed39bb39cf854d5250e95b1c1f94f4ed}
      \field{labelnamesource}{author}
      \field{labeltitlesource}{title}
      \field{annotation}{The author begins with the known values (assuming GRH and LI), for $q=8$ and $q=12$, of $\delta_{q;a,1}$ where $a\not\equiv1\mod q$ and of $\delta_{q;a,b,c}$ where $\{a,b,c\}$ is a permutation of the three nonidentity elements of $(\Z/q\Z)^\times$. He investigates how one could have predicted the relative sizes of these densities using the values and conductors of the nonprincipal Dirichlet characters\mod q, by arguing by analogy with independent random variables with variances of different sizes.\par The author also comments on the equality of a family of variant definitions of the logarithmic density of a set of positive real numbers, as well as some conjectures on the rarity of ties $\pi(x;q,a)=\pi(x;q,b)$.\par This article cites~\cites{MR0083001,MR0476616,MR529694,MR552467,MR1329368,MR1806291}.}
      \field{booktitle}{Number theory for the millennium, {II} ({U}rbana, {IL}, 2000)}
      \field{title}{Asymmetries in the {S}hanks-{R}ényi prime number race}
      \field{year}{2002}
      \true{nocite}
      \field{pages}{403\bibrangedash 415}
      \range{pages}{13}
      \verb{mrnumber}
      \verb 1956261
      \endverb
    \endentry
    \entry{MR2013259}{article}{}
      \name{author}{2}{}{%
        {{hash=f8a84ad9ba564b4272e977c2c3bc0e70}{%
           family={Balanzario},
           familyi={B\bibinitperiod},
           given={E.\bibnamedelimi P.},
           giveni={E\bibinitperiod\bibinitdelim P\bibinitperiod}}}%
        {{hash=92c6e72ba89fceafe31367461d88d360}{%
           family={Hernández},
           familyi={H\bibinitperiod},
           given={S.},
           giveni={S\bibinitperiod}}}%
      }
      \strng{namehash}{63612cf56d8d86ce4a5d5b3a7102abf9}
      \strng{fullhash}{63612cf56d8d86ce4a5d5b3a7102abf9}
      \strng{bibnamehash}{63612cf56d8d86ce4a5d5b3a7102abf9}
      \strng{authorbibnamehash}{63612cf56d8d86ce4a5d5b3a7102abf9}
      \strng{authornamehash}{63612cf56d8d86ce4a5d5b3a7102abf9}
      \strng{authorfullhash}{63612cf56d8d86ce4a5d5b3a7102abf9}
      \field{sortinit}{2}
      \field{sortinithash}{ed39bb39cf854d5250e95b1c1f94f4ed}
      \field{labelnamesource}{author}
      \field{labeltitlesource}{title}
      \field{annotation}{Fix distinct integers $a\ge0$ and $b\ge1$ and a real number $0<\theta<\frac1{2b}$, and define \[ h(x)=\sum_{n=1}^\infty \ell(n)e^{-n/x}-\frac{\Gamma(1/a)}{a\zeta(b/a)} x^{1/a}, \quad\text{where } \frac{\zeta(as)}{\zeta(bs)}=\sum_{n=1}^\infty\frac{\ell(n)}{n^s}. \] (When $a=0$, the second term in the definition of~$h$ is omitted.) Special cases include $\ell(n)=\mu(n)$ when $(a,b)=(0,1)$ and $\ell(n)=\lambda(n)$ when $(a,b)=(2,1)$. Assuming that the zeros of $\zeta(s)$, up to a height depending on $\theta$ and $b$, are simple and lie on the critical line, the authors show that $W(h;T;t^\theta) \gg \log T$. For example, in the special cases mentioned above, the authors verify the assumption up to height $10^5$ to deduce oscillations of size $t^{0.278}$. \par This article cites~\cites{MR777013,MR1193094}.}
      \field{issn}{0003-889X}
      \field{journaltitle}{Arch. Math. (Basel)}
      \field{number}{3}
      \field{title}{On the number of large oscillations of some arithmetical power series}
      \field{volume}{81}
      \field{year}{2003}
      \true{nocite}
      \field{pages}{285\bibrangedash 290}
      \range{pages}{6}
      \verb{doi}
      \verb 10.1007/s00013-003-4704-2
      \endverb
      \verb{mrnumber}
      \verb 2013259
      \endverb
    \endentry
    \entry{MR2075622}{article}{}
      \name{author}{2}{}{%
        {{hash=e7e31df1803902b56a78f32b1c733598}{%
           family={Ford},
           familyi={F\bibinitperiod},
           given={K.},
           giveni={K\bibinitperiod}}}%
        {{hash=c48086b9b194fded979a790687062dbb}{%
           family={Konyagin},
           familyi={K\bibinitperiod},
           given={S.},
           giveni={S\bibinitperiod}}}%
      }
      \list{publisher}{1}{%
        {Univ. Bonn, Bonn}%
      }
      \strng{namehash}{3cc0b23765677fa7230ff0cbba627702}
      \strng{fullhash}{3cc0b23765677fa7230ff0cbba627702}
      \strng{bibnamehash}{3cc0b23765677fa7230ff0cbba627702}
      \strng{authorbibnamehash}{3cc0b23765677fa7230ff0cbba627702}
      \strng{authornamehash}{3cc0b23765677fa7230ff0cbba627702}
      \strng{authorfullhash}{3cc0b23765677fa7230ff0cbba627702}
      \field{extraname}{3}
      \field{sortinit}{2}
      \field{sortinithash}{ed39bb39cf854d5250e95b1c1f94f4ed}
      \field{labelnamesource}{author}
      \field{labeltitlesource}{title}
      \field{annotation}{The authors continue to explore ``barriers'' for various statements in comparative prime number theory. In this article they make the distinction between the barrier itself, which is a multiset of complex numbers, and the consequences to orderings of prime-counting functions (or functions that are sufficiently close to such) that would follow from the $L(s,\chi)$ having their rightmost zeros precisely at the elements of the barrier. After establishing some lemmas on values of trigonometric polynomials, the authors prove several results concerning whether or not bounded (or finite) barriers exist for various races of the form $\pi(x;q,1)$ against many other $\pi(x;q,a)$. They show that if every pair from a set of $r$ functions is exhaustive, then the total number of $r$-way orderings is at least $r(r-1)/2+1$; they investigate whether barriers can exist that limit the number of orderings to this minimal value (``extremal barriers''), and construct barriers that do force at most $r(r-1)$ orderings.\par This article cites~\cites{MR1909220,MR1985941,MR1083000,MR1251926,MR0146156,MR0146157,MR0146158,MR0146159,MR0146160,MR0146161,MR0156826,MR0156827,MR0162771,MR0174538,MR0180539,MR0182616A,MR0182616B,MR0200250,1914.Littlewood}.}
      \field{booktitle}{Proceedings of the {S}ession in {A}nalytic {N}umber {T}heory and {D}iophantine {E}quations}
      \field{series}{Bonner Math. Schriften}
      \field{title}{The prime number race and zeros of {$L$}-functions off the critical line. {II}}
      \field{volume}{360}
      \field{year}{2003}
      \true{nocite}
      \field{pages}{40}
      \range{pages}{1}
      \verb{mrnumber}
      \verb 2075622
      \endverb
    \endentry
    \entry{MR1957110}{article}{}
      \name{author}{2}{}{%
        {{hash=6fa84c8cc1e455847a05b9946a1bd205}{%
           family={Kaczorowski},
           familyi={K\bibinitperiod},
           given={Jerzy},
           giveni={J\bibinitperiod}}}%
        {{hash=d6ae4b9844776294c92ff23952e91b53}{%
           family={Ramaré},
           familyi={R\bibinitperiod},
           given={Olivier},
           giveni={O\bibinitperiod}}}%
      }
      \strng{namehash}{ea4952aa08d573beb06051bb8ebb2f86}
      \strng{fullhash}{ea4952aa08d573beb06051bb8ebb2f86}
      \strng{bibnamehash}{ea4952aa08d573beb06051bb8ebb2f86}
      \strng{authorbibnamehash}{ea4952aa08d573beb06051bb8ebb2f86}
      \strng{authornamehash}{ea4952aa08d573beb06051bb8ebb2f86}
      \strng{authorfullhash}{ea4952aa08d573beb06051bb8ebb2f86}
      \field{sortinit}{2}
      \field{sortinithash}{ed39bb39cf854d5250e95b1c1f94f4ed}
      \field{labelnamesource}{author}
      \field{labeltitlesource}{title}
      \field{annotation}{This article investigates, assuming GRH, the distribution of values of a large class of functions of arithmetic significance (related to the Selberg class) using boundary values of $k$-functions. They establish an explicit formula for the appropriate functions as well as almost periodicity in the $L^2$ sense of Stepanov and the existence of a limiting logarithmic density.\par This article cites~\cites{MR1083000,MR1105608,MR1344249,MR1645811,MR1329368}.}
      \field{issn}{0065-1036}
      \field{journaltitle}{Acta Arith.}
      \field{number}{3}
      \field{title}{Almost periodicity of some error terms in prime number theory}
      \field{volume}{106}
      \field{year}{2003}
      \true{nocite}
      \field{pages}{277\bibrangedash 297}
      \range{pages}{21}
      \verb{doi}
      \verb 10.4064/aa106-3-6
      \endverb
      \verb{mrnumber}
      \verb 1957110
      \endverb
    \endentry
    \entry{MR2037293}{article}{}
      \name{author}{1}{}{%
        {{hash=13724bf2435ff7e0fb4fe7a4020ba08c}{%
           family={Leboeuf},
           familyi={L\bibinitperiod},
           given={P.},
           giveni={P\bibinitperiod}}}%
      }
      \strng{namehash}{13724bf2435ff7e0fb4fe7a4020ba08c}
      \strng{fullhash}{13724bf2435ff7e0fb4fe7a4020ba08c}
      \strng{bibnamehash}{13724bf2435ff7e0fb4fe7a4020ba08c}
      \strng{authorbibnamehash}{13724bf2435ff7e0fb4fe7a4020ba08c}
      \strng{authornamehash}{13724bf2435ff7e0fb4fe7a4020ba08c}
      \strng{authorfullhash}{13724bf2435ff7e0fb4fe7a4020ba08c}
      \field{sortinit}{2}
      \field{sortinithash}{ed39bb39cf854d5250e95b1c1f94f4ed}
      \field{labelnamesource}{author}
      \field{labeltitlesource}{title}
      \field{annotation}{In Section~$4$, the author heuristically computes higher moments of $E^\pi(x)$ from the explicit formula, determining that the third moment should be asymptotic to $-\frac3{\sqrt x} \sum_{\gamma>0} (\frac14+\gamma^2)^2$; while this vanishes in the limit as $x\to\infty$, it demonstrates an asymmetry for finite~$x$. The author also calculates the limiting fourth moment and notes that it differs from that of the Gaussian approximation. In addition, the author shows that there are persistent correlations between values of $E^\pi(x)$ at arguments separated by up to~$\frac x{\gamma_1}$, where $\gamma_1\approx 14.135$ is the smallest ordinate of a nontrivial of $\zeta(s)$. These observations are supported by numerical calculations. \par This article cites~\cites{MR1329368,MR1651741,MR0202686}.}
      \field{issn}{1424-0637}
      \field{journaltitle}{Ann. Henri Poincaré}
      \field{number}{suppl. 2}
      \field{title}{Prime correlations and fluctuations}
      \field{volume}{4}
      \field{year}{2003}
      \true{nocite}
      \field{pages}{S727\bibrangedash S752}
      \range{pages}{-1}
      \verb{doi}
      \verb 10.1007/s00023-003-0958-2
      \endverb
      \verb{mrnumber}
      \verb 2037293
      \endverb
    \endentry
    \entry{MR2047102}{article}{}
      \name{author}{3}{}{%
        {{hash=6f53fb227b01894d40689a82f6e1c775}{%
           family={Deléglise},
           familyi={D\bibinitperiod},
           given={Marc},
           giveni={M\bibinitperiod}}}%
        {{hash=34a11ef93294af142e84c4243d02208b}{%
           family={Dusart},
           familyi={D\bibinitperiod},
           given={Pierre},
           giveni={P\bibinitperiod}}}%
        {{hash=75094d519df9a079290d93e149ccf600}{%
           family={Roblot},
           familyi={R\bibinitperiod},
           given={Xavier-François},
           giveni={X\bibinithyphendelim F\bibinitperiod}}}%
      }
      \strng{namehash}{1237cd8bf0b8b34be0815571948c8580}
      \strng{fullhash}{1237cd8bf0b8b34be0815571948c8580}
      \strng{bibnamehash}{1237cd8bf0b8b34be0815571948c8580}
      \strng{authorbibnamehash}{1237cd8bf0b8b34be0815571948c8580}
      \strng{authornamehash}{1237cd8bf0b8b34be0815571948c8580}
      \strng{authorfullhash}{1237cd8bf0b8b34be0815571948c8580}
      \field{sortinit}{2}
      \field{sortinithash}{ed39bb39cf854d5250e95b1c1f94f4ed}
      \field{labelnamesource}{author}
      \field{labeltitlesource}{title}
      \field{annotation}{Extending a well-known method of Meissel for computing $\pi(x)$, the authors give an algorithm for computing $\pi(x;q,a)$ in time $O(x^{2/3}/\log^2 x)$. They use this algorithm is used to compute $\pi(x;4,1)$ and $\pi(x;4,3)$ for $x = 10^j,2 \cdot 10^j,\ldots,9 \cdot 10^j$ for each $10 \leq j \leq 19$. In particular, a new region where $\pi(x;4,3)<\pi(x;4,1)$ is discovered near $x = 10^{18}$. \par This article cites~\cites{MR0476615,MR1816036,MR1806291,MR0184920,MR0083001,1914.Littlewood,MR1329368}.}
      \field{issn}{0025-5718}
      \field{journaltitle}{Math. Comp.}
      \field{number}{247}
      \field{title}{Counting primes in residue classes}
      \field{volume}{73}
      \field{year}{2004}
      \true{nocite}
      \field{pages}{1565\bibrangedash 1575}
      \range{pages}{11}
      \verb{doi}
      \verb 10.1090/S0025-5718-04-01649-7
      \endverb
      \verb{mrnumber}
      \verb 2047102
      \endverb
    \endentry
    \entry{MR2118272}{article}{}
      \name{author}{2}{}{%
        {{hash=9491eb14d70047b89ad523bb740e7125}{%
           family={Kotnik},
           familyi={K\bibinitperiod},
           given={Tadej},
           giveni={T\bibinitperiod}}}%
        {{hash=1ed18e6c700f1143da76930e91f0e194}{%
           family={Lune},
           familyi={L\bibinitperiod},
           given={Jan},
           giveni={J\bibinitperiod},
           prefix={van\bibnamedelima de},
           prefixi={v\bibinitperiod\bibinitdelim d\bibinitperiod}}}%
      }
      \strng{namehash}{01c804ddeb08d741c890deb05472a70b}
      \strng{fullhash}{01c804ddeb08d741c890deb05472a70b}
      \strng{bibnamehash}{01c804ddeb08d741c890deb05472a70b}
      \strng{authorbibnamehash}{01c804ddeb08d741c890deb05472a70b}
      \strng{authornamehash}{01c804ddeb08d741c890deb05472a70b}
      \strng{authorfullhash}{01c804ddeb08d741c890deb05472a70b}
      \field{sortinit}{2}
      \field{sortinithash}{ed39bb39cf854d5250e95b1c1f94f4ed}
      \field{labelnamesource}{author}
      \field{labeltitlesource}{title}
      \field{annotation}{The authors use the first million zeros of $\zeta(s)$ to approximate $E^M(x)$ for $x \leq 10^{10^{10}}$. Based on the data, the authors conjecture that $E^M(x) = \Omega_{\pm} (\sqrt{\log \log \log x} )$, with the data suggesting approximately the constant $\frac12$ for each sign. In particular, the authors propose that earlier conjectures~\cite{MR280447,MR240062,MR2078705} on the extreme values of $E^M(x)$ do not hold. \par This article cites~\cites{1914.Littlewood,1897.Mertens,MR2078705,MR783538,MR891440,MR280447,MR240062}.}
      \field{issn}{1058-6458}
      \field{journaltitle}{Experiment. Math.}
      \field{number}{4}
      \field{title}{On the order of the {M}ertens function}
      \field{volume}{13}
      \field{year}{2004}
      \true{nocite}
      \field{pages}{473\bibrangedash 481}
      \range{pages}{9}
      \verb{mrnumber}
      \verb 2118272
      \endverb
    \endentry
    \entry{MR2034131}{article}{}
      \name{author}{1}{}{%
        {{hash=40a927b0e919f747a0218929b31aef12}{%
           family={Moree},
           familyi={M\bibinitperiod},
           given={Pieter},
           giveni={P\bibinitperiod}}}%
      }
      \strng{namehash}{40a927b0e919f747a0218929b31aef12}
      \strng{fullhash}{40a927b0e919f747a0218929b31aef12}
      \strng{bibnamehash}{40a927b0e919f747a0218929b31aef12}
      \strng{authorbibnamehash}{40a927b0e919f747a0218929b31aef12}
      \strng{authornamehash}{40a927b0e919f747a0218929b31aef12}
      \strng{authorfullhash}{40a927b0e919f747a0218929b31aef12}
      \field{sortinit}{2}
      \field{sortinithash}{ed39bb39cf854d5250e95b1c1f94f4ed}
      \field{labelnamesource}{author}
      \field{labeltitlesource}{title}
      \field{annotation}{Define $N(x;q,a)$ to be the number of integers up to $x$ all of whose prime factors are congruent to $a\mod q$. The author shows that $\min \{N(x;3,2), N(x;4,3)\} \geq \max\{N(x;3,1), N(x;4,1)\}$ for all $x\ge1$, using the fact that $N(x;q,a)$ is the summatory function of a multiplicative function whose values on primes alone has summatory function $\pi(x;q,a)$; consequently, biases in the distribution of small primes in residue classes modulo~$q$ can be magnified into complete biases among the $N(x;q,a)$. Indeed, Wirsing's method is already enough to establish such inequalities for sufficiently large~$x$; the author develops effective versions of Wirsing's method to conclude the same for all $x\ge1$. \par This article cites~\cites{MR0476616,1853.Chebyshev,MR1862054,MR1645811,MR0162771,1914.Littlewood,MR1756780,MR0137689,MR1329368}.}
      \field{issn}{0025-5718}
      \field{journaltitle}{Math. Comp.}
      \field{number}{245}
      \field{title}{{C}hebyshev's bias for composite numbers with restricted prime divisors}
      \field{volume}{73}
      \field{year}{2004}
      \true{nocite}
      \field{pages}{425\bibrangedash 449}
      \range{pages}{25}
      \verb{doi}
      \verb 10.1090/S0025-5718-03-01536-9
      \endverb
      \verb{mrnumber}
      \verb 2034131
      \endverb
    \endentry
    \entry{MR2078705}{article}{}
      \name{author}{1}{}{%
        {{hash=585b691d1c27d2773434ed1ac9664075}{%
           family={Ng},
           familyi={N\bibinitperiod},
           given={N.},
           giveni={N\bibinitperiod}}}%
      }
      \strng{namehash}{585b691d1c27d2773434ed1ac9664075}
      \strng{fullhash}{585b691d1c27d2773434ed1ac9664075}
      \strng{bibnamehash}{585b691d1c27d2773434ed1ac9664075}
      \strng{authorbibnamehash}{585b691d1c27d2773434ed1ac9664075}
      \strng{authornamehash}{585b691d1c27d2773434ed1ac9664075}
      \strng{authorfullhash}{585b691d1c27d2773434ed1ac9664075}
      \field{extraname}{2}
      \field{sortinit}{2}
      \field{sortinithash}{ed39bb39cf854d5250e95b1c1f94f4ed}
      \field{labelnamesource}{author}
      \field{labeltitlesource}{title}
      \field{annotation}{This article contains several results on the Mertens sum $M(x)$, all conditional on both RH and the estimate $\sum_{0<\gamma\le T} 1/{|\zeta'(\rho)|^2}\ll T$. The author shows that $M(x)\ll x^{1/2}(\log{x})^{3/2}$, and furthermore that $M(x)\ll x^{1/2}(\log\log{x})^{3/2}$ except on a set of finite logarithmic measure. He also proves that \[ \int_1^X \bigg( \frac{M(x)}x \bigg)^2\,dx \sim \log X \cdot \sum_{\gamma>0}\frac 2{|\rho\zeta'(\rho)|^2}, \] which in particular implies the ``weak Mertens conjecture'' $\int_2^Y ({M(x)}/x)^2 \,dx\ll \log{Y}$. Under the additional assumption of LI, he shows that $E^M(x)$ has a limiting logarithmic distribution whose Fourier transform can be written down explicitly. Partly building on unpublished work by Gonek, the author conjectures that $E^M(x)=\Omega_{\pm} \big((\log \log \log x)^{5/4}\big).$ \par This article cites~\cites{MR1544509, 1932.Ingham, MR0006202, MR0104638, MR0330069, MR634679, MR783538, MR1159354, MR1329368, 2000.Ng}.}
      \field{issn}{0024-6115}
      \field{journaltitle}{Proc. London Math. Soc. (3)}
      \field{number}{2}
      \field{title}{The distribution of the summatory function of the {M}öbius function}
      \field{volume}{89}
      \field{year}{2004}
      \true{nocite}
      \field{pages}{361\bibrangedash 389}
      \range{pages}{29}
      \verb{doi}
      \verb 10.1112/S0024611504014741
      \endverb
      \verb{mrnumber}
      \verb 2078705
      \endverb
      \verb{urlraw}
      \verb https://doi.org/10.1112/S0024611504014741
      \endverb
      \verb{url}
      \verb https://doi.org/10.1112/S0024611504014741
      \endverb
    \endentry
    \entry{MR2040112}{article}{}
      \name{author}{1}{}{%
        {{hash=b772e92036ff60f0205b68346a3f61db}{%
           family={Schlage--Puchta},
           familyi={S\bibinithyphendelim \bibinithyphendelim P\bibinitperiod},
           given={J.-C.},
           giveni={J\bibinithyphendelim C\bibinitperiod}}}%
      }
      \strng{namehash}{b772e92036ff60f0205b68346a3f61db}
      \strng{fullhash}{b772e92036ff60f0205b68346a3f61db}
      \strng{bibnamehash}{b772e92036ff60f0205b68346a3f61db}
      \strng{authorbibnamehash}{b772e92036ff60f0205b68346a3f61db}
      \strng{authornamehash}{b772e92036ff60f0205b68346a3f61db}
      \strng{authorfullhash}{b772e92036ff60f0205b68346a3f61db}
      \field{extraname}{1}
      \field{sortinit}{2}
      \field{sortinithash}{ed39bb39cf854d5250e95b1c1f94f4ed}
      \field{labelnamesource}{author}
      \field{labeltitlesource}{title}
      \field{annotation}{This article examines the race between $\pi(x;q,1)$ and $\pi(x;q,{\ne}1) = \max_{a\not\equiv1\mod q} \pi(x;q,a)$. Set $C = \exp\big(\max\{q,e^{1260}\}^{170}+e^{18c_q}\big)$. Assuming GRH, the author proves that $W(\pi(x;q,1), \pi(x;q,{\ne}1); T) > (\log T)/C-1$, and in particular that there exists $x<e^C$ such that $\pi(x;q,1) > \pi(x;q,{\ne}1)$. The analogous results are proved for the race between $\pi(x)$ and $\li(x)$, with $C=\exp(e^{16.7})$.\par This article cites~\cites{1936.Ingham, MR1251926, MR1310738, MR0202686, MR607089, MR630377, MR1761276, MR866118, MR1573970}.}
      \field{issn}{0236-5294}
      \field{journaltitle}{Acta Math. Hungar.}
      \field{number}{4}
      \field{title}{Sign changes of {$\pi(x,q,1)-\pi(x,q,a)$}}
      \field{volume}{102}
      \field{year}{2004}
      \true{nocite}
      \field{pages}{305\bibrangedash 320}
      \range{pages}{16}
      \verb{doi}
      \verb 10.1023/B:AMHU.0000024681.23784.d1
      \endverb
      \verb{mrnumber}
      \verb 2040112
      \endverb
      \verb{urlraw}
      \verb https://doi.org/10.1023/B:AMHU.0000024681.23784.d1
      \endverb
      \verb{url}
      \verb https://doi.org/10.1023/B:AMHU.0000024681.23784.d1
      \endverb
    \endentry
    \entry{MR2256805}{article}{}
      \name{author}{1}{}{%
        {{hash=1ad9cfd8df25881b7891382dee75005c}{%
           family={Karatsuba},
           familyi={K\bibinitperiod},
           given={A.\bibnamedelimi A.},
           giveni={A\bibinitperiod\bibinitdelim A\bibinitperiod}}}%
      }
      \list{language}{1}{%
        {Russian}%
      }
      \strng{namehash}{1ad9cfd8df25881b7891382dee75005c}
      \strng{fullhash}{1ad9cfd8df25881b7891382dee75005c}
      \strng{bibnamehash}{1ad9cfd8df25881b7891382dee75005c}
      \strng{authorbibnamehash}{1ad9cfd8df25881b7891382dee75005c}
      \strng{authornamehash}{1ad9cfd8df25881b7891382dee75005c}
      \strng{authorfullhash}{1ad9cfd8df25881b7891382dee75005c}
      \field{extraname}{1}
      \field{sortinit}{2}
      \field{sortinithash}{ed39bb39cf854d5250e95b1c1f94f4ed}
      \field{labelnamesource}{author}
      \field{labeltitlesource}{title}
      \field{annotation}{Let $R_1(x)=\pi(x)-\li(x)+\frac{1}{2}\li(\sqrt{x})$. The author follows Kaczorowski's method in~\cite{MR777013, MR791085} to show that for $W(R_1,T) \gg \log T$ and $W(\int_0^x R_1(t)\,dt,T) \gg \log T$. \par This article cites~\cites{MR2169423,MR777013, MR791085}.}
      \field{issn}{0869-5652}
      \field{journaltitle}{Dokl. Akad. Nauk}
      \field{number}{4}
      \field{title}{Behavior of the function {$R_1(x)$} and of its mean value}
      \field{volume}{404}
      \field{year}{2005}
      \true{nocite}
      \field{pages}{439\bibrangedash 442}
      \range{pages}{4}
      \verb{mrnumber}
      \verb 2256805
      \endverb
    \endentry
    \entry{MR2169423}{article}{}
      \name{author}{1}{}{%
        {{hash=1ad9cfd8df25881b7891382dee75005c}{%
           family={Karatsuba},
           familyi={K\bibinitperiod},
           given={A.\bibnamedelimi A.},
           giveni={A\bibinitperiod\bibinitdelim A\bibinitperiod}}}%
      }
      \list{language}{1}{%
        {Russian}%
      }
      \strng{namehash}{1ad9cfd8df25881b7891382dee75005c}
      \strng{fullhash}{1ad9cfd8df25881b7891382dee75005c}
      \strng{bibnamehash}{1ad9cfd8df25881b7891382dee75005c}
      \strng{authorbibnamehash}{1ad9cfd8df25881b7891382dee75005c}
      \strng{authornamehash}{1ad9cfd8df25881b7891382dee75005c}
      \strng{authorfullhash}{1ad9cfd8df25881b7891382dee75005c}
      \field{extraname}{2}
      \field{sortinit}{2}
      \field{sortinithash}{ed39bb39cf854d5250e95b1c1f94f4ed}
      \field{labelnamesource}{author}
      \field{labeltitlesource}{title}
      \field{annotation}{Let $R_1(x)=\pi(x)-\li(x)+\frac{1}{2}\li(\sqrt{x})$, and recursively define $$ R_{j+1}(x) = \int_2^{x} R_j(t)\,dt = -\frac{x^{j+\frac{1}{2}}}{\log x} \sum_{\gamma} \frac{x^{\frac{1}{2}+i \gamma}}{(\frac{1}{2}+i\gamma) \cdots (j+\frac{1}{2}+i \gamma)} +O\bigg(\frac{x^{j+\frac{1}{2}}}{\log^2 x}\bigg) $$ for each $j \geq 1$. Ingham~\cite{1932.Ingham} conjectured that~$\pi(x)$ is better approximated by $\li(x)-\frac{1}{2}\li(\sqrt{x})$ than by $\li(x)$. Under RH, the author validates this conjecture by showing that each of the functions $R_j(x)$ changes signs infinitely often, in contrast to $\fA^\pi_1(x) < -\frac35 x^{3/2}/\log x + O(x^{3/2}/\log^2 x)$. \par This article cites~\cites{1932.Ingham}.}
      \field{issn}{2226-8383}
      \field{journaltitle}{Chebyshevskii Sb.}
      \field{number}{4(12)}
      \field{title}{On the approximation of {$\pi(x)$}}
      \field{volume}{5}
      \field{year}{2005}
      \true{nocite}
      \field{pages}{5\bibrangedash 20}
      \range{pages}{16}
      \verb{mrnumber}
      \verb 2169423
      \endverb
    \endentry
    \entry{MR2262605}{article}{}
      \name{author}{1}{}{%
        {{hash=1ad9cfd8df25881b7891382dee75005c}{%
           family={Karatsuba},
           familyi={K\bibinitperiod},
           given={A.\bibnamedelimi A.},
           giveni={A\bibinitperiod\bibinitdelim A\bibinitperiod}}}%
      }
      \list{language}{1}{%
        {Russian}%
      }
      \strng{namehash}{1ad9cfd8df25881b7891382dee75005c}
      \strng{fullhash}{1ad9cfd8df25881b7891382dee75005c}
      \strng{bibnamehash}{1ad9cfd8df25881b7891382dee75005c}
      \strng{authorbibnamehash}{1ad9cfd8df25881b7891382dee75005c}
      \strng{authornamehash}{1ad9cfd8df25881b7891382dee75005c}
      \strng{authorfullhash}{1ad9cfd8df25881b7891382dee75005c}
      \field{extraname}{3}
      \field{sortinit}{2}
      \field{sortinithash}{ed39bb39cf854d5250e95b1c1f94f4ed}
      \field{labelnamesource}{author}
      \field{labeltitlesource}{title}
      \field{annotation}{Let $R_1(x)=\pi(x)-\li(x)+\frac{1}{2}\li(\sqrt{x})$ and $R_{j+1}(x)=\int_2^{x} R_j(t)\,dt$ for $j\ge1$. The author generalizes his earlier article~\cite{MR2256805} to show that for $W(R_j,T) \gg \log T$ for all positive integers $j$. \par This article cites~\cites{MR0973808A, MR0973808B, MR777013, MR791085, MR2169423}.}
      \field{isbn}{5-87954-388-9}
      \field{issn}{2226-8383}
      \field{journaltitle}{Chebyshevskii Sb.}
      \field{number}{2(14)}
      \field{title}{On the number of sign changes of the function {$R_1(x)$} and its mean values}
      \field{volume}{6}
      \field{year}{2005}
      \true{nocite}
      \field{pages}{163\bibrangedash 183}
      \range{pages}{21}
      \verb{mrnumber}
      \verb 2262605
      \endverb
    \endentry
    \entry{MR2110393}{article}{}
      \name{author}{1}{}{%
        {{hash=2a44d6571c930fba6c9dbf36ab288ad2}{%
           family={Radziejewski},
           familyi={R\bibinitperiod},
           given={Maciej},
           giveni={M\bibinitperiod}}}%
      }
      \strng{namehash}{2a44d6571c930fba6c9dbf36ab288ad2}
      \strng{fullhash}{2a44d6571c930fba6c9dbf36ab288ad2}
      \strng{bibnamehash}{2a44d6571c930fba6c9dbf36ab288ad2}
      \strng{authorbibnamehash}{2a44d6571c930fba6c9dbf36ab288ad2}
      \strng{authornamehash}{2a44d6571c930fba6c9dbf36ab288ad2}
      \strng{authorfullhash}{2a44d6571c930fba6c9dbf36ab288ad2}
      \field{extraname}{1}
      \field{sortinit}{2}
      \field{sortinithash}{ed39bb39cf854d5250e95b1c1f94f4ed}
      \field{labelnamesource}{author}
      \field{labeltitlesource}{title}
      \field{annotation}{Given any number field~$K$ and any subgroup~$\Gamma$ of the narrow class group of~$K$, there is a notion of irreducibility of ideals in~$\Gamma$; the corresponding factorization into irreducibles can be nonunique, and indeed the lengths of such factorizations can also be nonunique. The author examines the counting functions of ideals in~$\Gamma$ whose set of factorization lengths has cardinality lying in a prescribed interval, or (alternatively) contains one of a prescribed set of positive integers. The author obtains oscillation results for the corresponding error terms of the form $\Omega(x^{1/2-\ep})$, as well as lower bounds of size $\log X$ for the number of sign changes up to~$X$. \par This article cites~\cites{MR891057,MR2123959}.}
      \field{issn}{0065-1036}
      \field{journaltitle}{Acta Arith.}
      \field{number}{2}
      \field{title}{On the distribution of algebraic numbers with prescribed factorization properties}
      \field{volume}{116}
      \field{year}{2005}
      \true{nocite}
      \field{pages}{153\bibrangedash 171}
      \range{pages}{19}
      \verb{doi}
      \verb 10.4064/aa116-2-4
      \endverb
      \verb{mrnumber}
      \verb 2110393
      \endverb
    \endentry
    \entry{MR2123959}{article}{}
      \name{author}{1}{}{%
        {{hash=2a44d6571c930fba6c9dbf36ab288ad2}{%
           family={Radziejewski},
           familyi={R\bibinitperiod},
           given={Maciej},
           giveni={M\bibinitperiod}}}%
      }
      \strng{namehash}{2a44d6571c930fba6c9dbf36ab288ad2}
      \strng{fullhash}{2a44d6571c930fba6c9dbf36ab288ad2}
      \strng{bibnamehash}{2a44d6571c930fba6c9dbf36ab288ad2}
      \strng{authorbibnamehash}{2a44d6571c930fba6c9dbf36ab288ad2}
      \strng{authornamehash}{2a44d6571c930fba6c9dbf36ab288ad2}
      \strng{authorfullhash}{2a44d6571c930fba6c9dbf36ab288ad2}
      \field{extraname}{2}
      \field{sortinit}{2}
      \field{sortinithash}{ed39bb39cf854d5250e95b1c1f94f4ed}
      \field{labelnamesource}{author}
      \field{labeltitlesource}{title}
      \field{annotation}{Kaczorowski and Pintz~\cite{MR891057} considered functions of a real variable whose Mellin transforms had singularities of the form $(s-\rho)^w P(\log(s-\rho))$ for some $w\in\C$. The author extends this class of functions to those with singularities that are linear combinations of this type, showing that such functions (suitably normalized) have oscillations of the form $\Omega(x^{1/2-\ep})$, as well as lower bounds of size $\log X$ for the number of sign changes up to~$X$. The motivation was to address error terms of counting functions corresponding to certain ideal factorization problems in number fields~\cite{MR2110393}, and the author provides an application to counting the number of ideals in a subgroup of the narrow class group all of whose restricted factorizations have the same length.\par This article cites~\cites{MR0198145,MR1555148,1936.Ingham,MR1256149,MR0219496,MR1511360,MR0439774,1930.Polya,MR891057,MR858395}.}
      \field{issn}{0026-9255}
      \field{journaltitle}{Monatsh. Math.}
      \field{number}{2}
      \field{title}{Oscillations of error terms associated with certain arithmetical functions}
      \field{volume}{144}
      \field{year}{2005}
      \true{nocite}
      \field{pages}{113\bibrangedash 130}
      \range{pages}{18}
      \verb{doi}
      \verb 10.1007/s00605-003-0147-x
      \endverb
      \verb{mrnumber}
      \verb 2123959
      \endverb
    \endentry
    \entry{MR2202918}{article}{}
      \name{author}{2}{}{%
        {{hash=7c95bde388d46e453494c7b2fbd1b088}{%
           family={Granville},
           familyi={G\bibinitperiod},
           given={A.},
           giveni={A\bibinitperiod}}}%
        {{hash=feda7d3cba2a0ad4e772307fca655937}{%
           family={Martin},
           familyi={M\bibinitperiod},
           given={G.},
           giveni={G\bibinitperiod}}}%
      }
      \strng{namehash}{c4ece346b079ec803e063600350bbf92}
      \strng{fullhash}{c4ece346b079ec803e063600350bbf92}
      \strng{bibnamehash}{c4ece346b079ec803e063600350bbf92}
      \strng{authorbibnamehash}{c4ece346b079ec803e063600350bbf92}
      \strng{authornamehash}{c4ece346b079ec803e063600350bbf92}
      \strng{authorfullhash}{c4ece346b079ec803e063600350bbf92}
      \field{sortinit}{2}
      \field{sortinithash}{ed39bb39cf854d5250e95b1c1f94f4ed}
      \field{labelnamesource}{author}
      \field{labeltitlesource}{title}
      \field{annotation}{The authors present an accessible survey of prime number races, explicit formulas for $\pi(x)$ and $\pi(x;q,a)$, the biases caused by squares of primes, and limiting distributions and densities. The final section describes unpublished research of G.~Davidoff (in connection with an REU group), including the theorem that $W_{q;\cN,\cR}(T)$ is unbounded for any prime modulus~$q$, assuming only that $L(s,\chi_{\pm q})$ has no real zero in the interval $[\Theta(\chi_{\pm q}),1]$.\par This article cites~\cites{1914.Littlewood,MR0004255,MR0108470,MR0146156,MR0146157,MR0146158,MR0146159,MR0146160,MR0146161,MR0156826,MR0156827,MR0289452,MR0476616,MR552467,MR1329368,MR1344249,MR1367576,2007.Sarnak,2000.Ng,MR1752093,MR1816036,MR1806291,MR1956261,MR1909220, MR2075622,MR0882550}.}
      \field{issn}{0002-9890}
      \field{journaltitle}{Amer. Math. Monthly}
      \field{number}{1}
      \field{title}{Prime number races}
      \field{volume}{113}
      \field{year}{2006}
      \true{nocite}
      \field{pages}{1\bibrangedash 33}
      \range{pages}{33}
      \verb{doi}
      \verb 10.2307/27641834
      \endverb
      \verb{mrnumber}
      \verb 2202918
      \endverb
      \verb{urlraw}
      \verb https://doi.org/10.2307/27641834
      \endverb
      \verb{url}
      \verb https://doi.org/10.2307/27641834
      \endverb
    \endentry
    \entry{MR2282922}{incollection}{}
      \name{author}{2}{}{%
        {{hash=9491eb14d70047b89ad523bb740e7125}{%
           family={Kotnik},
           familyi={K\bibinitperiod},
           given={Tadej},
           giveni={T\bibinitperiod}}}%
        {{hash=c83d046dcdb536f5163e46eb8fe07827}{%
           family={Riele},
           familyi={R\bibinitperiod},
           given={Herman},
           giveni={H\bibinitperiod},
           prefix={te},
           prefixi={t\bibinitperiod}}}%
      }
      \list{publisher}{1}{%
        {Springer, Berlin}%
      }
      \strng{namehash}{ed0be9de3d43cca36b74eb71c7163f7a}
      \strng{fullhash}{ed0be9de3d43cca36b74eb71c7163f7a}
      \strng{bibnamehash}{ed0be9de3d43cca36b74eb71c7163f7a}
      \strng{authorbibnamehash}{ed0be9de3d43cca36b74eb71c7163f7a}
      \strng{authornamehash}{ed0be9de3d43cca36b74eb71c7163f7a}
      \strng{authorfullhash}{ed0be9de3d43cca36b74eb71c7163f7a}
      \field{sortinit}{2}
      \field{sortinithash}{ed39bb39cf854d5250e95b1c1f94f4ed}
      \field{labelnamesource}{author}
      \field{labeltitlesource}{title}
      \field{annotation}{Based on the work of Odlyzko and te Riele~\cite{MR783538}, a refined version of the LLL algorithm, and improved precision for zeros of $\zeta$ function, the authors show that \[ \liminf_{x \to \infty} E^M(x) < -1.229 \quad \text{and} \quad \limsup_{x \to \infty} E^M(x) > 1.218. \] Using a similar idea, the authors show that a counterexample to the Mertens conjecture occurs before $\exp(1.59 \times 10^{40})$, improving Pintz's result~\cite{MR891440}. The authors also provide new numerical evidence to support the conjecture $M(x) = \Omega_{\pm} (\sqrt{x\log \log \log x} )$ made in~\cite{MR2118272}. \par This article cites~\cite{MR0046485,MR891440,MR2078705,MR783538,1897.Mertens,MR2118272,MR429789,MR0330069,MR0352026,MR0006202,MR549977,MR280447,MR240062}.}
      \field{booktitle}{Algorithmic number theory}
      \field{series}{Lecture Notes in Comput. Sci.}
      \field{title}{The {M}ertens conjecture revisited}
      \field{volume}{4076}
      \field{year}{2006}
      \true{nocite}
      \field{pages}{156\bibrangedash 167}
      \range{pages}{12}
      \verb{mrnumber}
      \verb 2282922
      \endverb
    \endentry
    \entry{MR2271616}{article}{}
      \name{author}{2}{}{%
        {{hash=d6fc16fe2b971a8009f6a09376b27581}{%
           family={Montgomery},
           familyi={M\bibinitperiod},
           given={Hugh\bibnamedelima L.},
           giveni={H\bibinitperiod\bibinitdelim L\bibinitperiod}}}%
        {{hash=855d3da694fbeeb0998eb3736ff8e5d1}{%
           family={Vorhauer},
           familyi={V\bibinitperiod},
           given={Ulrike\bibnamedelimb M.\bibnamedelimi A.},
           giveni={U\bibinitperiod\bibinitdelim M\bibinitperiod\bibinitdelim A\bibinitperiod}}}%
      }
      \strng{namehash}{a14ec9a880532023786b3f0b7cf9fee9}
      \strng{fullhash}{a14ec9a880532023786b3f0b7cf9fee9}
      \strng{bibnamehash}{a14ec9a880532023786b3f0b7cf9fee9}
      \strng{authorbibnamehash}{a14ec9a880532023786b3f0b7cf9fee9}
      \strng{authornamehash}{a14ec9a880532023786b3f0b7cf9fee9}
      \strng{authorfullhash}{a14ec9a880532023786b3f0b7cf9fee9}
      \field{sortinit}{2}
      \field{sortinithash}{ed39bb39cf854d5250e95b1c1f94f4ed}
      \field{labelnamesource}{author}
      \field{labeltitlesource}{title}
      \field{annotation}{The authors show that $\Delta^\psi(x)$ takes values of both signs in every interval $[x,19x]$ for $x\geq 1$ (where the~$19$ is best possible), and takes values of both signs in every interval $[x,2.02x]$ when $x$ is sufficiently large.\par This article cites~\cites{1936.Ingham, MR777013, MR1133236, 1914.Littlewood, MR1574052,1930.Polya}.}
      \field{issn}{0208-6573}
      \field{journaltitle}{Funct. Approx. Comment. Math.}
      \field{title}{Changes of sign of the error term in the prime number theorem}
      \field{volume}{35}
      \field{year}{2006}
      \true{nocite}
      \field{pages}{235\bibrangedash 247}
      \range{pages}{13}
      \verb{doi}
      \verb 10.7169/facm/1229442626
      \endverb
      \verb{mrnumber}
      \verb 2271616
      \endverb
    \endentry
    \entry{MR2340242}{article}{}
      \name{author}{1}{}{%
        {{hash=6fa84c8cc1e455847a05b9946a1bd205}{%
           family={Kaczorowski},
           familyi={K\bibinitperiod},
           given={Jerzy},
           giveni={J\bibinitperiod}}}%
      }
      \strng{namehash}{6fa84c8cc1e455847a05b9946a1bd205}
      \strng{fullhash}{6fa84c8cc1e455847a05b9946a1bd205}
      \strng{bibnamehash}{6fa84c8cc1e455847a05b9946a1bd205}
      \strng{authorbibnamehash}{6fa84c8cc1e455847a05b9946a1bd205}
      \strng{authornamehash}{6fa84c8cc1e455847a05b9946a1bd205}
      \strng{authorfullhash}{6fa84c8cc1e455847a05b9946a1bd205}
      \field{extraname}{16}
      \field{sortinit}{2}
      \field{sortinithash}{ed39bb39cf854d5250e95b1c1f94f4ed}
      \field{labelnamesource}{author}
      \field{labeltitlesource}{title}
      \field{annotation}{The author shows that \begin{align*} \sum_{n=1}^\infty \mu(n) \left(\cos \left(\frac{x}{n}\right) - 1\right) = \Omega_{\pm} (x^{1/2} \log \log \log x), \end{align*} which comes tantalizingly close to the assertion that $E^M(x) = \Omega_\pm(\log\log\log x)$. The proof uses an explicit formula for a relative of $k(z)$, namely \begin{align*} m(z) = \frac{1}{2\pi i} \int_{\mathcal C} \frac{e^{zs}}{\zeta(s)} \, ds \end{align*} where $\mathcal C$ is the boundary of a half strip enclosing the part of the critical strip lying in the upper half plane. \par This article cites~\cite{MR1109990, MR0006202, MR1105608, MR2329690, MR2118272, MR2078705, MR783538, MR891440, MR0882550}.}
      \field{issn}{0024-6107}
      \field{journaltitle}{J. Lond. Math. Soc. (2)}
      \field{number}{2}
      \field{title}{Results on the {M}öbius function}
      \field{volume}{75}
      \field{year}{2007}
      \true{nocite}
      \field{pages}{509\bibrangedash 521}
      \range{pages}{13}
      \verb{doi}
      \verb 10.1112/jlms/jdm006
      \endverb
      \verb{mrnumber}
      \verb 2340242
      \endverb
      \verb{urlraw}
      \verb https://doi.org/10.1112/jlms/jdm006
      \endverb
      \verb{url}
      \verb https://doi.org/10.1112/jlms/jdm006
      \endverb
    \endentry
    \entry{MR2329690}{article}{}
      \name{author}{2}{}{%
        {{hash=6fa84c8cc1e455847a05b9946a1bd205}{%
           family={Kaczorowski},
           familyi={K\bibinitperiod},
           given={Jerzy},
           giveni={J\bibinitperiod}}}%
        {{hash=87807ce1c4509fc3bd17938ccd4fe3e6}{%
           family={Wiertelak},
           familyi={W\bibinitperiod},
           given={Kazimierz},
           giveni={K\bibinitperiod}}}%
      }
      \strng{namehash}{5e2b3cb8731c2c2fbd035cbb54743005}
      \strng{fullhash}{5e2b3cb8731c2c2fbd035cbb54743005}
      \strng{bibnamehash}{5e2b3cb8731c2c2fbd035cbb54743005}
      \strng{authorbibnamehash}{5e2b3cb8731c2c2fbd035cbb54743005}
      \strng{authornamehash}{5e2b3cb8731c2c2fbd035cbb54743005}
      \strng{authorfullhash}{5e2b3cb8731c2c2fbd035cbb54743005}
      \field{extraname}{1}
      \field{sortinit}{2}
      \field{sortinithash}{ed39bb39cf854d5250e95b1c1f94f4ed}
      \field{labelnamesource}{author}
      \field{labeltitlesource}{title}
      \field{annotation}{The authors consider a general class of functions of the form $F(z)=\sum_{n=1}^\infty a_ne^{i\omega_nz}$ on the upper-half plane, whose boundary values $P(x) = \lim_{y\to x^+} \Re F(x+iy)$ exist for sufficiently large $x$. If certain growth conditions are satisfied, they establish oscillation results for $P(x)$. Applying their result to $K(z)e^{-z/2}$ recovers Littlewood's theorem $\Delta^\psi(x)=\Omega_\pm(x^\frac{1}{2}\log\log\log{x})$, and similar consequences hold for $\Delta^\psi(x;q,1)$ and $\psi(x;q,1,a)$. As another application, the authors establish that the error term in the asymptotic formula for $\sum_{n\leq x}2^{\omega(n)}$ is $\Omega_\pm\big(x^{1/4} (\log\log{x})^{1/2} / (\log\log\log{x})^{3/2} \big)$. (As an aside, in this variant of the divisor problem, obtaining even $\ll x^{1/2-\delta}$ for the error term requires~RH.)\par This article cites~\cites{MR1083000,MR1105608,MR1645811,1914.Littlewood}.}
      \field{issn}{0305-0041}
      \field{journaltitle}{Math. Proc. Cambridge Philos. Soc.}
      \field{number}{3}
      \field{title}{{$\Omega$}-estimates for a class of arithmetic error terms}
      \field{volume}{142}
      \field{year}{2007}
      \true{nocite}
      \field{pages}{385\bibrangedash 394}
      \range{pages}{10}
      \verb{doi}
      \verb 10.1017/S0305004107000035
      \endverb
      \verb{mrnumber}
      \verb 2329690
      \endverb
    \endentry
    \entry{2007.Sarnak}{online}{}
      \name{author}{1}{}{%
        {{hash=ea78b04077f0596f70f50f0559b2a1f2}{%
           family={Sarnak},
           familyi={S\bibinitperiod},
           given={P.},
           giveni={P\bibinitperiod}}}%
      }
      \strng{namehash}{ea78b04077f0596f70f50f0559b2a1f2}
      \strng{fullhash}{ea78b04077f0596f70f50f0559b2a1f2}
      \strng{bibnamehash}{ea78b04077f0596f70f50f0559b2a1f2}
      \strng{authorbibnamehash}{ea78b04077f0596f70f50f0559b2a1f2}
      \strng{authornamehash}{ea78b04077f0596f70f50f0559b2a1f2}
      \strng{authorfullhash}{ea78b04077f0596f70f50f0559b2a1f2}
      \field{sortinit}{2}
      \field{sortinithash}{ed39bb39cf854d5250e95b1c1f94f4ed}
      \field{labelnamesource}{author}
      \field{labeltitlesource}{title}
      \field{annotation}{Let $\lambda(p)=\tau(p)/p^{11/2}$ denote the normalized $p$th Fourier coefficient of the Ramanujan function $\Delta(z)$, a holomorphic cusp form of weight~$12$. Under various assumptions (functional equation, RH, LI) for the symmetric power $L$-functions associated to $\Delta(z)$, the author proves the existence of a limiting logarithmic distribution for the cumulative sum $x^{-1/2}\log x \cdot \sum_{p\le x}\lambda(p)$. This distribution has mean~$1$ but infinite variance, so that even though there is a bias towards being positive, the set of such $x$ has logarithmic density $\frac12$. The author provides a similar analysis for $x^{-1/2}\log x \cdot \sum_{p\le x}a_E(p)/\sqrt p$, where $a_E(p)$ are the coefficients of the weight-$2$ cusp form attached to an elliptic curve $E/\mathbb{Q}$; here the mean is $1-2r(E)$ where $r(E)$ is the rank of~$E$. Here the variance is (conjecturally) finite, and so the logarithmic density of the set of $x$ for which $\sum_{p\le x}a_E(p)/\sqrt p > 0$ is strictly between~$\frac12$ and~$1$ when $r(E)=0$, but strictly between~$0$ and~$\frac12$ when $r(E)\ge1$. The author makes analogous remarks about the symmetric powers of these elliptic curve $L$-functions, where the finiteness of the variance corresponds to whether $E$ has complex multiplication.\par This letter cites~\cites{MR1329368}.}
      \field{title}{Letter to {B}arry {M}azur on `{C}hebyshev's bias' for {$\tau (p)$}}
      \field{year}{2007}
      \true{nocite}
      \verb{urlraw}
      \verb http://web.math.princeton.edu/sarnak/MazurLtrMay08.PDF
      \endverb
      \verb{url}
      \verb http://web.math.princeton.edu/sarnak/MazurLtrMay08.PDF
      \endverb
    \endentry
    \entry{MR2398787}{article}{}
      \name{author}{3}{}{%
        {{hash=86e7b84d8adeabc3f69abb52df304f6a}{%
           family={Borwein},
           familyi={B\bibinitperiod},
           given={Peter},
           giveni={P\bibinitperiod}}}%
        {{hash=a15eabe7e88a5dca9dbe96e2754e374a}{%
           family={Ferguson},
           familyi={F\bibinitperiod},
           given={Ron},
           giveni={R\bibinitperiod}}}%
        {{hash=454d3dedc781ec4a49112de10313b411}{%
           family={Mossinghoff},
           familyi={M\bibinitperiod},
           given={Michael\bibnamedelima J.},
           giveni={M\bibinitperiod\bibinitdelim J\bibinitperiod}}}%
      }
      \strng{namehash}{96894e6dfcaffbe33c893760d2cc6cc2}
      \strng{fullhash}{96894e6dfcaffbe33c893760d2cc6cc2}
      \strng{bibnamehash}{96894e6dfcaffbe33c893760d2cc6cc2}
      \strng{authorbibnamehash}{96894e6dfcaffbe33c893760d2cc6cc2}
      \strng{authornamehash}{96894e6dfcaffbe33c893760d2cc6cc2}
      \strng{authorfullhash}{96894e6dfcaffbe33c893760d2cc6cc2}
      \field{sortinit}{2}
      \field{sortinithash}{ed39bb39cf854d5250e95b1c1f94f4ed}
      \field{labelnamesource}{author}
      \field{labeltitlesource}{title}
      \field{annotation}{The authors develop refined algorithms for computing the sums $L(x)$ and $L_r(x)$. They show that the smallest positive integer~$n$ for which $L_r(n)<0$ (that is, the first witness to the negative resolution of Turán's problem) is $72,185,376,951,205$; they also describe the many sign changes of $L_r(x)$ among the subsequent integers, and suggest that its next sign change occurs near $1.16\cdot10^{19}$. Similar results are obtained for $L(n)$, including a new region of sign changes starting at $351,100,332,278,253$; their calculation allows them to conclude that $L(n) > 0.061867\sqrt n$ infinitely often.\par This article \cites{MR27305,MR654520,MR0120198,MR0006202,MR0104638,MR0003644,MR584557}.}
      \field{issn}{0025-5718}
      \field{journaltitle}{Math. Comp.}
      \field{number}{263}
      \field{title}{Sign changes in sums of the {L}iouville function}
      \field{volume}{77}
      \field{year}{2008}
      \true{nocite}
      \field{pages}{1681\bibrangedash 1694}
      \range{pages}{14}
      \verb{doi}
      \verb 10.1090/S0025-5718-08-02036-X
      \endverb
      \verb{mrnumber}
      \verb 2398787
      \endverb
    \endentry
    \entry{MR2474313}{article}{}
      \name{author}{1}{}{%
        {{hash=7871a490afe3f3963421804f312b0ad7}{%
           family={Cha},
           familyi={C\bibinitperiod},
           given={Byungchul},
           giveni={B\bibinitperiod}}}%
      }
      \strng{namehash}{7871a490afe3f3963421804f312b0ad7}
      \strng{fullhash}{7871a490afe3f3963421804f312b0ad7}
      \strng{bibnamehash}{7871a490afe3f3963421804f312b0ad7}
      \strng{authorbibnamehash}{7871a490afe3f3963421804f312b0ad7}
      \strng{authornamehash}{7871a490afe3f3963421804f312b0ad7}
      \strng{authorfullhash}{7871a490afe3f3963421804f312b0ad7}
      \field{extraname}{1}
      \field{sortinit}{2}
      \field{sortinithash}{ed39bb39cf854d5250e95b1c1f94f4ed}
      \field{labelnamesource}{author}
      \field{labeltitlesource}{title}
      \field{annotation}{The author adapts the arguments of Rubinstein and Sarnak~\cite{MR1329368} to the function field setting, establishing (under the appropriate generalization of LI) the existence of limiting distributions and their symmetries, a bias towards quadratic nonresidues in two-way races, and central limit theorems. The author also provides some examples where LI is violated and the expected biases are overridden, as well as an example where LI can actually be verified.\par This article cites~\cites{1853.Chebyshev,MR2202918,MR1329368}.}
      \field{issn}{0010-437X}
      \field{journaltitle}{Compos. Math.}
      \field{number}{6}
      \field{title}{{C}hebyshev's bias in function fields}
      \field{volume}{144}
      \field{year}{2008}
      \true{nocite}
      \field{pages}{1351\bibrangedash 1374}
      \range{pages}{24}
      \verb{doi}
      \verb 10.1112/S0010437X08003631
      \endverb
      \verb{mrnumber}
      \verb 2474313
      \endverb
    \endentry
    \entry{MR2420864}{article}{}
      \name{author}{1}{}{%
        {{hash=9491eb14d70047b89ad523bb740e7125}{%
           family={Kotnik},
           familyi={K\bibinitperiod},
           given={Tadej},
           giveni={T\bibinitperiod}}}%
      }
      \strng{namehash}{9491eb14d70047b89ad523bb740e7125}
      \strng{fullhash}{9491eb14d70047b89ad523bb740e7125}
      \strng{bibnamehash}{9491eb14d70047b89ad523bb740e7125}
      \strng{authorbibnamehash}{9491eb14d70047b89ad523bb740e7125}
      \strng{authornamehash}{9491eb14d70047b89ad523bb740e7125}
      \strng{authorfullhash}{9491eb14d70047b89ad523bb740e7125}
      \field{sortinit}{2}
      \field{sortinithash}{ed39bb39cf854d5250e95b1c1f94f4ed}
      \field{labelnamesource}{author}
      \field{labeltitlesource}{title}
      \field{annotation}{The author shows that $\pi(x) <\li(x)$ for $2 \leq x \leq 10^{14}$, improving unpublished work of Odlyzko showing that $\pi(x) <\li(x)$ for $x \leq 1.59 \times 10^{13}$. The rigorous computation used an Eratosthenes sieve for $\pi(x)$ and Ramanujan's formula $\li(x)=\sqrt{x} \sum_{n=1}^{\infty} a_n \log ^n x + \log \log x + C_0$, truncated at $n=75$ and linearly interpolated when~$x>10^{10}$. Based on the data set, the author conjectures that $|\Delta^\pi(x)|<\sqrt{x}$ and ${-\frac{2}{5}x^{3/2} <\fA^\pi_1(x)<0}$ for $x>2$. The program for computation and storage of the data was written in Delphi~6.0 and run on a PC with a 2.4~GHz Intel Pentium~4 processor and 512~MB of RAM.\par This article cites~\cites{MR1752093,MR0369287,1932.Ingham,MR0202686,1914.Littlewood,MR0137689,MR0067145,MR866118}.}
      \field{issn}{1019-7168}
      \field{journaltitle}{Adv. Comput. Math.}
      \field{number}{1}
      \field{title}{The prime-counting function and its analytic approximations: {$\pi (x)$} and its approximations}
      \field{volume}{29}
      \field{year}{2008}
      \true{nocite}
      \field{pages}{55\bibrangedash 70}
      \range{pages}{16}
      \verb{doi}
      \verb 10.1007/s10444-007-9039-2
      \endverb
      \verb{mrnumber}
      \verb 2420864
      \endverb
    \endentry
    \entry{MR2605532}{article}{}
      \name{author}{2}{}{%
        {{hash=5b3dabbd7db4c81dcf1e7cffedcb1b4d}{%
           family={Diamond},
           familyi={D\bibinitperiod},
           given={H.\bibnamedelimi G.},
           giveni={H\bibinitperiod\bibinitdelim G\bibinitperiod}}}%
        {{hash=c6d965df16e600c701d4bcc2888ae34c}{%
           family={Pintz},
           familyi={P\bibinitperiod},
           given={J.},
           giveni={J\bibinitperiod}}}%
      }
      \strng{namehash}{d3c019510c17f40c3aa1bd8a2407a489}
      \strng{fullhash}{d3c019510c17f40c3aa1bd8a2407a489}
      \strng{bibnamehash}{d3c019510c17f40c3aa1bd8a2407a489}
      \strng{authorbibnamehash}{d3c019510c17f40c3aa1bd8a2407a489}
      \strng{authornamehash}{d3c019510c17f40c3aa1bd8a2407a489}
      \strng{authorfullhash}{d3c019510c17f40c3aa1bd8a2407a489}
      \field{sortinit}{2}
      \field{sortinithash}{ed39bb39cf854d5250e95b1c1f94f4ed}
      \field{labelnamesource}{author}
      \field{labeltitlesource}{title}
      \field{annotation}{In this article, the authors show that $-\sum_{p\leq x}\log(1-\frac{1}{p})-\log\log x-C_0=\Omega_\pm(1/(\sqrt{x}\log x))$, which implies that $\prod_{p\leq x}(1-\frac{1}{p})^{-1} - e^{{C_0}}\log x$ changes sign infinitely often. If RH is false, the result is a consequence of Landau's theorem. If RH is true, the authors first show that $$ \int_1^x\frac{d\Pi(t)}{t}-\int_1^x\frac{1-t^{-1}}{t\log t}\,dt=\Omega_\pm\biggl( \frac{\log\log\log x}{\sqrt{x}\log x} \biggr) $$ and then use Littlewood's~\cites{1914.Littlewood} result on the sign changes of $\Delta^\pi(x)$ and Cramér's bound $\Delta^{E^{\psi}}_1(x)\ll x$. The authors include a second proof in the RH case using a variant of the Wiener--Ikehara method due to Ingham~\cites{MR0006202}.\par This article cites~\cites{MR0376566,MR0006202,1914.Littlewood,MR0137689,MR43841,MR51857}.}
      \field{issn}{1246-7405}
      \field{journaltitle}{J. Théor. Nombres Bordeaux}
      \field{number}{3}
      \field{title}{Oscillation of {M}ertens' product formula}
      \field{volume}{21}
      \field{year}{2009}
      \true{nocite}
      \field{pages}{523\bibrangedash 533}
      \range{pages}{11}
      \verb{mrnumber}
      \verb 2605532
      \endverb
      \verb{urlraw}
      \verb http://jtnb.cedram.org/item?id=JTNB_2009__21_3_523_0
      \endverb
      \verb{url}
      \verb http://jtnb.cedram.org/item?id=JTNB_2009__21_3_523_0
      \endverb
    \endentry
    \entry{MR2470105}{article}{}
      \name{author}{1}{}{%
        {{hash=6fa84c8cc1e455847a05b9946a1bd205}{%
           family={Kaczorowski},
           familyi={K\bibinitperiod},
           given={Jerzy},
           giveni={J\bibinitperiod}}}%
      }
      \strng{namehash}{6fa84c8cc1e455847a05b9946a1bd205}
      \strng{fullhash}{6fa84c8cc1e455847a05b9946a1bd205}
      \strng{bibnamehash}{6fa84c8cc1e455847a05b9946a1bd205}
      \strng{authorbibnamehash}{6fa84c8cc1e455847a05b9946a1bd205}
      \strng{authornamehash}{6fa84c8cc1e455847a05b9946a1bd205}
      \strng{authorfullhash}{6fa84c8cc1e455847a05b9946a1bd205}
      \field{extraname}{17}
      \field{sortinit}{2}
      \field{sortinithash}{ed39bb39cf854d5250e95b1c1f94f4ed}
      \field{labelnamesource}{author}
      \field{labeltitlesource}{title}
      \field{annotation}{Given a number field~$K$, the author studies the oscillations of the error term $E_K(x)$ for the counting function of irreducible elements of~$K$ (up to units). Let $\zeta_{K_H}(s)$ denote the Dedekind zeta function of the Hilbert class field of~$K$, and let $\rho$ denote a nontrivial zero of $\zeta_{K_H}$. The author introduces a complicated integer-valued quantity $m^*(\rho,K)$ related to the multiplicity of $\rho$; under the assumption that some $m^*(\rho,K)$ is nonzero, he proves that $E_K(x) = \Omega_\pm\big({\sqrt{x}(\log\log x)^{D(K)-1}/\log{x}}\big)$, where $D(K)$ is the Davenport constant of the class group of~$K$. As a result, when $K$ with class number~$2$, we unconditionally have $E_k(x) = \Omega_\pm\big({\sqrt{x}\log\log{x}/\log{x}}\big)$; when the class number is an odd prime~$p$, we have $E_k(x)=\Omega_\pm({\sqrt{x}(\log\log{x})^{p-1}/\log{x}})$ provided that $\zeta_{K_H}(s)$ has at least one nontrivial zero of multiplicity not divisible by~$p$. All of these oscillations have logarithmic frequency. Assuming RH for $\zeta_{K_H}(s)$, the author sketches a proof of $H_K(x)\ll \sqrt x(\log{x})^{D(K)+1}$.\par This article cites~\cites{MR791085,MR891057,MR1957110,MR2329690,MR2506435,1914.Littlewood,MR0882550}.}
      \field{issn}{0026-9255}
      \field{journaltitle}{Monatsh. Math.}
      \field{number}{1}
      \field{title}{On the distribution of irreducible algebraic integers}
      \field{volume}{156}
      \field{year}{2009}
      \true{nocite}
      \field{pages}{47\bibrangedash 71}
      \range{pages}{25}
      \verb{doi}
      \verb 10.1007/s00605-008-0559-8
      \endverb
      \verb{mrnumber}
      \verb 2470105
      \endverb
    \endentry
    \entry{MR2506435}{article}{}
      \name{author}{2}{}{%
        {{hash=6fa84c8cc1e455847a05b9946a1bd205}{%
           family={Kaczorowski},
           familyi={K\bibinitperiod},
           given={Jerzy},
           giveni={J\bibinitperiod}}}%
        {{hash=87807ce1c4509fc3bd17938ccd4fe3e6}{%
           family={Wiertelak},
           familyi={W\bibinitperiod},
           given={Kazimierz},
           giveni={K\bibinitperiod}}}%
      }
      \strng{namehash}{5e2b3cb8731c2c2fbd035cbb54743005}
      \strng{fullhash}{5e2b3cb8731c2c2fbd035cbb54743005}
      \strng{bibnamehash}{5e2b3cb8731c2c2fbd035cbb54743005}
      \strng{authorbibnamehash}{5e2b3cb8731c2c2fbd035cbb54743005}
      \strng{authornamehash}{5e2b3cb8731c2c2fbd035cbb54743005}
      \strng{authorfullhash}{5e2b3cb8731c2c2fbd035cbb54743005}
      \field{extraname}{2}
      \field{sortinit}{2}
      \field{sortinithash}{ed39bb39cf854d5250e95b1c1f94f4ed}
      \field{labelnamesource}{author}
      \field{labeltitlesource}{title}
      \field{annotation}{The authors use a combination of methods from~\cites{MR777013,MR791085,MR858395,MR2329690} to prove that $$ W\bigl( E^\psi(T);T;\log\log H(T) \bigr)\gg (\log{T})/H(T) $$ for sufficiently large~$T$, where $H(T)$ is a function with $1\ll H(T)<\log T$. Letting $H(T)=\exp((\log\log{T})^c)$ with a small $c>0$, this theorem recovers the classical result of Littlewood that $\Delta^\psi(x)=\Omega_\pm(\sqrt{x}\log\log\log{x})$. They show the analogous bound for $E^\pi(T;q,a,1)$ assuming HC for Dirichlet $L$-functions. The authors also investigate the summatory function $D_2(x)=\sum_{n\leq x}2^{\omega(n)}$. Assuming that all the simple zeros $\rho=1/2+i\gamma$ of $\zeta(s)$ on the critical line satisfy ${\zeta'(\rho)}\gg|\gamma|^{-O(1)}$, they show that $$ W\bigl( D_2(T);T;T^{1/4}(\log H(T))^{1/2}(\log\log H(T))^{-3/2} \bigr) \gg (\log T)/H(T) $$ for any function $1\ll H(T)\leq \log T$.\par This article cites~\cites{MR777013,MR791085,MR927376,MR945273,MR858395,MR1083000,MR1105608,MR2329690,MR2271616,1914.Littlewood,MR0882550}.}
      \field{issn}{0002-9947}
      \field{journaltitle}{Trans. Amer. Math. Soc.}
      \field{number}{9}
      \field{title}{Oscillations of a given size of some arithmetic error terms}
      \field{volume}{361}
      \field{year}{2009}
      \true{nocite}
      \field{pages}{5023\bibrangedash 5039}
      \range{pages}{17}
      \verb{doi}
      \verb 10.1090/S0002-9947-09-04803-X
      \endverb
      \verb{mrnumber}
      \verb 2506435
      \endverb
    \endentry
    \entry{MR2753165}{book}{}
      \name{author}{1}{}{%
        {{hash=9670e498e8fb9912d411c47aed724ede}{%
           family={Sneed},
           familyi={S\bibinitperiod},
           given={J.\bibnamedelimi P.},
           giveni={J\bibinitperiod\bibinitdelim P\bibinitperiod}}}%
      }
      \list{publisher}{1}{%
        {ProQuest LLC, Ann Arbor, MI}%
      }
      \strng{namehash}{9670e498e8fb9912d411c47aed724ede}
      \strng{fullhash}{9670e498e8fb9912d411c47aed724ede}
      \strng{bibnamehash}{9670e498e8fb9912d411c47aed724ede}
      \strng{authorbibnamehash}{9670e498e8fb9912d411c47aed724ede}
      \strng{authornamehash}{9670e498e8fb9912d411c47aed724ede}
      \strng{authorfullhash}{9670e498e8fb9912d411c47aed724ede}
      \field{sortinit}{2}
      \field{sortinithash}{ed39bb39cf854d5250e95b1c1f94f4ed}
      \field{labelnamesource}{author}
      \field{labeltitlesource}{title}
      \field{annotation}{In Chapter~2, the author verifies HC for moduli $q\le100$ and shows that all (two-way) prime number races for such moduli have infinitely many lead changes. In Chapter~3, the author explores semiprime races, which concern the functions $\pi_2(x;q,\ell)=\#\{p_1p_2\leq x\colon p_1p_2\equiv\ell\mod q\}$ and $E_2(x;q,\ell_1,\ell_2)=\big( \pi_2(x;q,\ell_1)-\pi_2(x;q,\ell_2) \big)\log x/\sqrt x\log\log x$. Assuming GRH and LI, he proves that $E_2(x;4,3,1)$ has mean $-\frac12$ (negative, unlike $E(x;4,3,1)$) and takes negative values with logarithmic density $\approx 0.894280$. \par This article cites~\cites{MR0447090,1853.Chebyshev,MR0332684,MR1862054,MR1985941,MR0202685,MR2202918,MR0006202,MR0176967,MR1511360,MR0146157,MR0083001,1914.Littlewood,MR1329368,MR0289452,MR0004255,MR429789}.}
      \field{isbn}{978-1124-04248-0}
      \field{note}{Thesis (Ph.D.)--University of Illinois at Urbana-Champaign}
      \field{pagetotal}{83}
      \field{title}{Prime and quasi-prime number races}
      \field{year}{2009}
      \true{nocite}
      \verb{mrnumber}
      \verb 2753165
      \endverb
      \verb{urlraw}
      \verb http://gateway.proquest.com/openurl?url_ver=Z39.88-2004&rft_val_fmt=info:ofi/fmt:kev:mtx:dissertation&res_dat=xri:pqdiss&rft_dat=xri:pqdiss:3411454
      \endverb
      \verb{url}
      \verb http://gateway.proquest.com/openurl?url_ver=Z39.88-2004&rft_val_fmt=info:ofi/fmt:kev:mtx:dissertation&res_dat=xri:pqdiss&rft_dat=xri:pqdiss:3411454
      \endverb
    \endentry
    \entry{MR2600420}{article}{}
      \name{author}{2}{}{%
        {{hash=7871a490afe3f3963421804f312b0ad7}{%
           family={Cha},
           familyi={C\bibinitperiod},
           given={Byungchul},
           giveni={B\bibinitperiod}}}%
        {{hash=1470ddce431ad26e652280417656b586}{%
           family={Kim},
           familyi={K\bibinitperiod},
           given={Seick},
           giveni={S\bibinitperiod}}}%
      }
      \strng{namehash}{6505b468496e6694d7b4710ddb7b703d}
      \strng{fullhash}{6505b468496e6694d7b4710ddb7b703d}
      \strng{bibnamehash}{6505b468496e6694d7b4710ddb7b703d}
      \strng{authorbibnamehash}{6505b468496e6694d7b4710ddb7b703d}
      \strng{authornamehash}{6505b468496e6694d7b4710ddb7b703d}
      \strng{authorfullhash}{6505b468496e6694d7b4710ddb7b703d}
      \field{sortinit}{2}
      \field{sortinithash}{ed39bb39cf854d5250e95b1c1f94f4ed}
      \field{labelnamesource}{author}
      \field{labeltitlesource}{title}
      \field{annotation}{The authors derive a formula, in the function field setting, for the logarithmic densities of races among the counting functions of irreducible polynomials of degree up to~$X$ in residue classes. The proofs follow~\cite{MR1329368} closely, although the authors introduce a smoothing function on the Fourier side to overcome the problem that the Fourier transforms decay slowly due to the finiteness of the number of zeros of the relevant $L$-functions.\par This article cites ~\cites{MR1329368,MR2474313,MR1806291,MR2202918}.}
      \field{issn}{0022-314X}
      \field{journaltitle}{J. Number Theory}
      \field{number}{4}
      \field{title}{Biases in the prime number race of function fields}
      \field{volume}{130}
      \field{year}{2010}
      \true{nocite}
      \field{pages}{1048\bibrangedash 1055}
      \range{pages}{8}
      \verb{doi}
      \verb 10.1016/j.jnt.2009.09.015
      \endverb
      \verb{mrnumber}
      \verb 2600420
      \endverb
    \endentry
    \entry{MR2652902}{article}{}
      \name{author}{2}{}{%
        {{hash=69f328db05bb9b15874d29cfcf79dfdb}{%
           family={Chao},
           familyi={C\bibinitperiod},
           given={Kuok\bibnamedelima Fai},
           giveni={K\bibinitperiod\bibinitdelim F\bibinitperiod}}}%
        {{hash=81a915ac987d508458c06ee0891cc3c2}{%
           family={Plymen},
           familyi={P\bibinitperiod},
           given={Roger},
           giveni={R\bibinitperiod}}}%
      }
      \strng{namehash}{0737652978438d918ff3cb160495ab49}
      \strng{fullhash}{0737652978438d918ff3cb160495ab49}
      \strng{bibnamehash}{0737652978438d918ff3cb160495ab49}
      \strng{authorbibnamehash}{0737652978438d918ff3cb160495ab49}
      \strng{authornamehash}{0737652978438d918ff3cb160495ab49}
      \strng{authorfullhash}{0737652978438d918ff3cb160495ab49}
      \field{sortinit}{2}
      \field{sortinithash}{ed39bb39cf854d5250e95b1c1f94f4ed}
      \field{labelnamesource}{author}
      \field{labeltitlesource}{title}
      \field{annotation}{The authors modify Lehman's theorem~\cite{MR0202686} by improving the bound on $\theta_1(x) = \frac{2\log x}3 \big( \frac{\pi(x)}{x\log x}-1)$ using an inequality of Panaitopol. Together with $2\times10^6$ zeros of $\zeta(s)$ (supplied by Odlyzko), they show that there exist at least $10^{154}$ consecutive integers $x$ in the range $[1.3978965\times 10^{316}, 1.398344\times 10^{316}]$ for which $\pi(x)>\li(x)$.\par This article cites~\cites{MR1752093,MR2420864,MR0202686,1914.Littlewood,MR866118,MR0137689,MR1573970,MR0067145}.}
      \field{issn}{1793-0421}
      \field{journaltitle}{Int. J. Number Theory}
      \field{number}{3}
      \field{title}{A new bound for the smallest {$x$} with {$\pi (x)>\li (x)$}}
      \field{volume}{6}
      \field{year}{2010}
      \true{nocite}
      \field{pages}{681\bibrangedash 690}
      \range{pages}{10}
      \verb{doi}
      \verb 10.1142/S1793042110003125
      \endverb
      \verb{mrnumber}
      \verb 2652902
      \endverb
    \endentry
    \entry{MR2778652}{article}{}
      \name{author}{2}{}{%
        {{hash=e7e31df1803902b56a78f32b1c733598}{%
           family={Ford},
           familyi={F\bibinitperiod},
           given={K.},
           giveni={K\bibinitperiod}}}%
        {{hash=ee6277cd0f78c9744eaedab7ee5add23}{%
           family={Sneed},
           familyi={S\bibinitperiod},
           given={J.},
           giveni={J\bibinitperiod}}}%
      }
      \strng{namehash}{06ab6a04f799f3b33364b611b68406f5}
      \strng{fullhash}{06ab6a04f799f3b33364b611b68406f5}
      \strng{bibnamehash}{06ab6a04f799f3b33364b611b68406f5}
      \strng{authorbibnamehash}{06ab6a04f799f3b33364b611b68406f5}
      \strng{authornamehash}{06ab6a04f799f3b33364b611b68406f5}
      \strng{authorfullhash}{06ab6a04f799f3b33364b611b68406f5}
      \field{sortinit}{2}
      \field{sortinithash}{ed39bb39cf854d5250e95b1c1f94f4ed}
      \field{labelnamesource}{author}
      \field{labeltitlesource}{title}
      \field{annotation}{The authors examine Chebyshev's bias for integers which are the product of two primes. Let $\pi_2(x;q,a)$ denote the number of integers $n$ up to $x$ such that $n\equiv a\mod q$ and $\Omega(n)=2$.\par Assume GRH for Dirichlet $L$-functions modulo $q$ and also that the zeros of $L(s,\chi)$ are simple for each nonprincipal character $\chi\mod q$. If $f(x_1,\ldots,x_r)$ is the logarithmic density function of $\big( E(x;q,a_1,b_1),\ldots,E(x;q,a_r,b_r) \big),$ the authors show that the logarithmic density function of $\big( E_2(x;q,a_1,b_1),\ldots,E_2(x;q,a_r,b_r) \big)$ is \[f\left(\frac{c_q(b_1)-c_q(a_1)}2-x_1,\ldots,\frac{c_q(b_r)-c_q(a_r)}2-x_r\right).\] Consequently, assuming both GRH and LI for modulus $q$, the authors show that $\delta_2\big( \pi_2(x;q,a),\pi_2(x;q,b) \big)$ exists, and equals $\frac12$ if $a$ and $b$ are both quadratic residues or both quadratic nonresidues\mod q. Otherwise, if $a$ is a quadratic nonresidue and $b$ is a quadratic residue, then \[1-\delta_{q;a_1,a_2}<\delta_2\big( \pi_2(x;q,a),\pi_2(x;q,b) \big)<\frac 12.\]\par This article cites~\cites{1853.Chebyshev, 1914.Littlewood, MR0083001, MR0146156, MR1329368, MR1985941, MR2202918, MR2753165, MR3028758}.}
      \field{issn}{1058-6458}
      \field{journaltitle}{Experiment. Math.}
      \field{number}{4}
      \field{title}{Chebyshev's bias for products of two primes}
      \field{volume}{19}
      \field{year}{2010}
      \true{nocite}
      \field{pages}{385\bibrangedash 398}
      \range{pages}{14}
      \verb{doi}
      \verb 10.1080/10586458.2010.10390630
      \endverb
      \verb{mrnumber}
      \verb 2778652
      \endverb
      \verb{urlraw}
      \verb https://doi.org/10.1080/10586458.2010.10390630
      \endverb
      \verb{url}
      \verb https://doi.org/10.1080/10586458.2010.10390630
      \endverb
    \endentry
    \entry{MR2730494}{article}{}
      \name{author}{1}{}{%
        {{hash=6fa84c8cc1e455847a05b9946a1bd205}{%
           family={Kaczorowski},
           familyi={K\bibinitperiod},
           given={Jerzy},
           giveni={J\bibinitperiod}}}%
      }
      \strng{namehash}{6fa84c8cc1e455847a05b9946a1bd205}
      \strng{fullhash}{6fa84c8cc1e455847a05b9946a1bd205}
      \strng{bibnamehash}{6fa84c8cc1e455847a05b9946a1bd205}
      \strng{authorbibnamehash}{6fa84c8cc1e455847a05b9946a1bd205}
      \strng{authornamehash}{6fa84c8cc1e455847a05b9946a1bd205}
      \strng{authorfullhash}{6fa84c8cc1e455847a05b9946a1bd205}
      \field{extraname}{18}
      \field{sortinit}{2}
      \field{sortinithash}{ed39bb39cf854d5250e95b1c1f94f4ed}
      \field{labelnamesource}{author}
      \field{labeltitlesource}{title}
      \field{annotation}{The author continues his examination of~$E_K(x)$ from~\cite{MR2470105}, under the assumption that no entire function from the Selberg class vanishes at $s=1$ (which follows from Selberg's orthonormality conjecture). He shows that this assumption implies that at least one of the quantities $m^*(\rho,K)$ from his previous work is nonzero, and therefore that this nonvanishing assumption implies that $E_K(x) = \Omega_\pm\big({\sqrt{x}(\log\log x)^{D(K)-1}/\log{x}}\big)$, where $D(K)$ is the Davenport constant of the class group of~$K$. He also shows unconditionally that $E_K(x)=\Omega_\pm({\sqrt{x}/(\log{x})^{O_K(1)}})$. All these oscillations have logarithmic frequency.\par This article cites~\cites{MR777013,MR791085,MR2470105,MR891057,MR2506435}.}
      \field{issn}{0025-584X}
      \field{journaltitle}{Math. Nachr.}
      \field{number}{9}
      \field{title}{{$\Omega$}-estimates related to irreducible algebraic integers}
      \field{volume}{283}
      \field{year}{2010}
      \true{nocite}
      \field{pages}{1291\bibrangedash 1303}
      \range{pages}{13}
      \verb{doi}
      \verb 10.1002/mana.200710158
      \endverb
      \verb{mrnumber}
      \verb 2730494
      \endverb
    \endentry
    \entry{MR2684490}{article}{}
      \name{author}{2}{}{%
        {{hash=6fa84c8cc1e455847a05b9946a1bd205}{%
           family={Kaczorowski},
           familyi={K\bibinitperiod},
           given={Jerzy},
           giveni={J\bibinitperiod}}}%
        {{hash=87807ce1c4509fc3bd17938ccd4fe3e6}{%
           family={Wiertelak},
           familyi={W\bibinitperiod},
           given={Kazimierz},
           giveni={K\bibinitperiod}}}%
      }
      \strng{namehash}{5e2b3cb8731c2c2fbd035cbb54743005}
      \strng{fullhash}{5e2b3cb8731c2c2fbd035cbb54743005}
      \strng{bibnamehash}{5e2b3cb8731c2c2fbd035cbb54743005}
      \strng{authorbibnamehash}{5e2b3cb8731c2c2fbd035cbb54743005}
      \strng{authornamehash}{5e2b3cb8731c2c2fbd035cbb54743005}
      \strng{authorfullhash}{5e2b3cb8731c2c2fbd035cbb54743005}
      \field{extraname}{3}
      \field{sortinit}{2}
      \field{sortinithash}{ed39bb39cf854d5250e95b1c1f94f4ed}
      \field{labelnamesource}{author}
      \field{labeltitlesource}{title}
      \field{annotation}{In this article, the authors study the oscillations of $f(x)=\sum_{n\leq x}\phi(n)-\frac{3}{\pi^2}x^2$, the best known results being Walfisz's estimate $f(x)\ll x(\log{x})^{2/3}(\log\log{x})^{4/3}$ and Montgomery's result $f(x)=\Omega_\pm(x\sqrt{\log\log{x}})$. The authors show that $f(x) = f^{AR}(x)+f^{AN}(x)$, where the arithmetic part $f^{AR}(x) = \sum_{n=1}^\infty\frac{\mu(n)}{n}\{\frac{x}{n}\}$ and the analytic part $f^{AN}(x) = -\int_0^x f^{AR}(t)\frac{dt}{t} = \frac{1}{2}\sum_{n=1}^\infty\mu(n)\{\frac{x}{n}\}^2+\frac{1}{2}$. They establish Montgomery-type oscillations for $f^{AR}(x)$ and similar functions, including examples arising from coefficients of newforms. The analytic part $f^{AN}(x)$ has an explicit formula (thus is $o(x)$ unconditionally in particular) and is interesting in its own right, and the authors establish Littlewood-type oscillations for it. \par This article cites~\cites{MR777013,MR791085,MR2470105,MR2730494,MR858395,MR2329690,MR2506435,MR2759800}.}
      \field{issn}{0022-314X}
      \field{journaltitle}{J. Number Theory}
      \field{number}{12}
      \field{title}{Oscillations of the remainder term related to the {E}uler totient function}
      \field{volume}{130}
      \field{year}{2010}
      \true{nocite}
      \field{pages}{2683\bibrangedash 2700}
      \range{pages}{18}
      \verb{doi}
      \verb 10.1016/j.jnt.2010.06.010
      \endverb
      \verb{mrnumber}
      \verb 2684490
      \endverb
    \endentry
    \entry{MR2759800}{article}{}
      \name{author}{2}{}{%
        {{hash=6fa84c8cc1e455847a05b9946a1bd205}{%
           family={Kaczorowski},
           familyi={K\bibinitperiod},
           given={Jerzy},
           giveni={J\bibinitperiod}}}%
        {{hash=87807ce1c4509fc3bd17938ccd4fe3e6}{%
           family={Wiertelak},
           familyi={W\bibinitperiod},
           given={Kazimierz},
           giveni={K\bibinitperiod}}}%
      }
      \strng{namehash}{5e2b3cb8731c2c2fbd035cbb54743005}
      \strng{fullhash}{5e2b3cb8731c2c2fbd035cbb54743005}
      \strng{bibnamehash}{5e2b3cb8731c2c2fbd035cbb54743005}
      \strng{authorbibnamehash}{5e2b3cb8731c2c2fbd035cbb54743005}
      \strng{authornamehash}{5e2b3cb8731c2c2fbd035cbb54743005}
      \strng{authorfullhash}{5e2b3cb8731c2c2fbd035cbb54743005}
      \field{extraname}{4}
      \field{sortinit}{2}
      \field{sortinithash}{ed39bb39cf854d5250e95b1c1f94f4ed}
      \field{labelnamesource}{author}
      \field{labeltitlesource}{title}
      \field{annotation}{The authors observe that the order of magnitude of arithmetic error terms is often uninfluenced by smoothing; for example, $\Delta^\psi(x) = \Omega_\pm (x\log\log\log x)$ unconditionally and $\Delta^\psi(x) \ll x(\log x)^2$ on RH, while $A_1^\psi(x)=\Omega_\pm(\sqrt{x})$ unconditionally and $A_1^\psi(x)=O(\sqrt{x})$ on RH. If we define $f(x)=\sum_{n\leq x}\phi(n)-\frac{3}{\pi^2}x^2$, then Walfisz showed that $f(x)\ll x(\log{x})^{2/3}(\log\log{x})^{4/3}$ and Montgomery showed that $f(x)=\Omega_\pm(x\sqrt{\log\log{x}})$. In contrast, the authors show that $A_1^f(x)=\Omega_\pm(\sqrt{x}\log\log\log{x})$ unconditionally and $A_1^f(x) \le \sqrt{x}\exp\big( O(\log{x}/\log\log{x}) \big)$ on RH.\par This article cites~\cites{MR777013, MR791085,MR2470105,MR2730494,MR858395,MR1083000,MR1105608,MR1105608,MR2329690,MR2506435,1914.Littlewood}.}
      \field{issn}{0025-584X}
      \field{journaltitle}{Math. Nachr.}
      \field{number}{11}
      \field{title}{Smoothing arithmetic error terms: the case of the {E}uler {$\phi$} function}
      \field{volume}{283}
      \field{year}{2010}
      \true{nocite}
      \field{pages}{1637\bibrangedash 1645}
      \range{pages}{9}
      \verb{doi}
      \verb 10.1002/mana.200810048
      \endverb
      \verb{mrnumber}
      \verb 2759800
      \endverb
    \endentry
    \entry{MR2684372}{article}{}
      \name{author}{2}{}{%
        {{hash=99298fc0cd61b25871f9b3ad564fa78b}{%
           family={Saouter},
           familyi={S\bibinitperiod},
           given={Yannick},
           giveni={Y\bibinitperiod}}}%
        {{hash=447f22d7de07f4cbd918f6b2abb97288}{%
           family={Demichel},
           familyi={D\bibinitperiod},
           given={Patrick},
           giveni={P\bibinitperiod}}}%
      }
      \strng{namehash}{74715dbf8807c6d6381a8af29072c8b3}
      \strng{fullhash}{74715dbf8807c6d6381a8af29072c8b3}
      \strng{bibnamehash}{74715dbf8807c6d6381a8af29072c8b3}
      \strng{authorbibnamehash}{74715dbf8807c6d6381a8af29072c8b3}
      \strng{authornamehash}{74715dbf8807c6d6381a8af29072c8b3}
      \strng{authorfullhash}{74715dbf8807c6d6381a8af29072c8b3}
      \field{sortinit}{2}
      \field{sortinithash}{ed39bb39cf854d5250e95b1c1f94f4ed}
      \field{labelnamesource}{author}
      \field{labeltitlesource}{title}
      \field{annotation}{The authors improve the error term for the function $I(\omega, \eta)$ in Lehman's theorem~\cite{MR0202686}, which allows them to conclude that there are more than $6.09 \times 10^{150}$ successive integers in the vicinity of $\exp(727.951335792) \approx 1.397166707819\times10^{316}$ such that $\pi(x)>\li(x)$. They show that $\pi(x)-\li(x)>9.1472 \times 10^{149}$ for some $x \in [1.39715131\times10^{316}, 1.39718211\times10^{316}]$; assuming RH, they show that $\pi(x)-\li(x)>1.7503\times 10^{148}$ for some~$x$ in the smaller interval $[1.3971619476\times10^{316},1.3971714624\times10^{316}]$.\par This article cites \cite{MR0137689, MR0202686, MR866118, MR552467, MR1752093, MR2652902}.}
      \field{issn}{0025-5718}
      \field{journaltitle}{Math. Comp.}
      \field{number}{272}
      \field{title}{A sharp region where {$\pi (x)-\li (x)$} is positive}
      \field{volume}{79}
      \field{year}{2010}
      \true{nocite}
      \field{pages}{2395\bibrangedash 2405}
      \range{pages}{11}
      \verb{doi}
      \verb 10.1090/S0025-5718-10-02351-3
      \endverb
      \verb{mrnumber}
      \verb 2684372
      \endverb
    \endentry
    \entry{MR2811555}{article}{}
      \name{author}{2}{}{%
        {{hash=3417f8b5237d274ddd67adbd43553547}{%
           family={Cha},
           familyi={C\bibinitperiod},
           given={B},
           giveni={B\bibinitperiod}}}%
        {{hash=5fe0c1b9c49f664236b5bbd12a052411}{%
           family={Im},
           familyi={I\bibinitperiod},
           given={B.-H.},
           giveni={B\bibinithyphendelim H\bibinitperiod}}}%
      }
      \strng{namehash}{ea08ea1554b13d2b20e9dca09ca0fd3e}
      \strng{fullhash}{ea08ea1554b13d2b20e9dca09ca0fd3e}
      \strng{bibnamehash}{ea08ea1554b13d2b20e9dca09ca0fd3e}
      \strng{authorbibnamehash}{ea08ea1554b13d2b20e9dca09ca0fd3e}
      \strng{authornamehash}{ea08ea1554b13d2b20e9dca09ca0fd3e}
      \strng{authorfullhash}{ea08ea1554b13d2b20e9dca09ca0fd3e}
      \field{sortinit}{2}
      \field{sortinithash}{ed39bb39cf854d5250e95b1c1f94f4ed}
      \field{labelnamesource}{author}
      \field{labeltitlesource}{title}
      \field{annotation}{The authors study Chebyshev's bias in a finite Galois extension~$F$ of a global function field~$K$. Let $\pi_C(N)$ be the number of prime elements in~$K$ of degree~$N$ whose Frobenius lies in~$C$. Following the strategy of Rubinstein and Sarnak~\cite{MR1329368}, the authors study the limiting distribution of $\frac{X}{q^{X/2}}\sum_{N=1}^X \frac{|G|}{|C|} \big( \pi_c(N)-\pi(N) \big)$. The authors show that when $F/K$ is geometric and satisfies LI, there is a bias towards conjugacy class of the Galois group containing fewer square elements, as in the number field case. The authors give a complete description of the prime element race in the case of scalar field extensions of~$K$.\par This article cites~\cite{MR1329368,MR2202918,MR1806291,MR2600420,MR2474313}.}
      \field{issn}{0022-314X}
      \field{journaltitle}{J. Number Theory}
      \field{number}{10}
      \field{title}{{C}hebyshev's bias in {G}alois extensions of global function fields}
      \field{volume}{131}
      \field{year}{2011}
      \true{nocite}
      \field{pages}{1875\bibrangedash 1886}
      \range{pages}{12}
      \verb{doi}
      \verb 10.1016/j.jnt.2011.03.011
      \endverb
      \verb{mrnumber}
      \verb 2811555
      \endverb
    \endentry
    \entry{MR3103752}{book}{}
      \name{author}{1}{}{%
        {{hash=889ecd9d3a804c5a7e9a5d8287766657}{%
           family={Fiorilli},
           familyi={F\bibinitperiod},
           given={D.},
           giveni={D\bibinitperiod}}}%
      }
      \list{language}{1}{%
        {French}%
      }
      \list{publisher}{1}{%
        {ProQuest LLC, Ann Arbor, MI}%
      }
      \strng{namehash}{889ecd9d3a804c5a7e9a5d8287766657}
      \strng{fullhash}{889ecd9d3a804c5a7e9a5d8287766657}
      \strng{bibnamehash}{889ecd9d3a804c5a7e9a5d8287766657}
      \strng{authorbibnamehash}{889ecd9d3a804c5a7e9a5d8287766657}
      \strng{authornamehash}{889ecd9d3a804c5a7e9a5d8287766657}
      \strng{authorfullhash}{889ecd9d3a804c5a7e9a5d8287766657}
      \field{extraname}{1}
      \field{sortinit}{2}
      \field{sortinithash}{ed39bb39cf854d5250e95b1c1f94f4ed}
      \field{labelnamesource}{author}
      \field{labeltitlesource}{title}
      \field{annotation}{Chapter 2 of this thesis is the article~\cite{MR3028758}. A result in the introduction explains why the method of~\cite{MR777013} is not capable of producing more sign changes in the $\pi$-vs.-$\li$ race: the author shows that if \[ \psi_k(x) = \frac1{k!} \sum_{n\le x} \Lambda(n) \big( \log \tfrac xn \big)^k = x - \sum_\rho \frac{x^\rho}{\rho^{k+1}} + \sum_{j=0}^k \frac{a_{k-j}}{j!} (\log x)^j \] is the $k$-fold logarithmic average of $\psi(x)$, then $W(\psi_k;T) = \frac{\gamma_1}{2\pi} \log T + O_k(1)$ for $k\ge5$, where $\gamma_1 \approx 14.135$ is the smallest ordinate of a nontrivial zero of $\zeta(s)$. The phenomenon is that the repeated averaging washes out all of the small-scale sign changes expected for $\psi(x)$ itself.}
      \field{isbn}{978-0494-89015-8}
      \field{note}{Thesis (Ph.D.)--Université de Montréal}
      \field{pagetotal}{261}
      \field{title}{Irrégularités dans la distribution des nombres premiers et des suites plus générales dans les progressions arithmétiques}
      \field{year}{2011}
      \true{nocite}
      \verb{mrnumber}
      \verb 3103752
      \endverb
    \endentry
    \entry{MR2813366}{article}{}
      \name{author}{2}{}{%
        {{hash=ba700e1ffe0ab3bd31564865c0880f36}{%
           family={Stoll},
           familyi={S\bibinitperiod},
           given={Douglas\bibnamedelima A.},
           giveni={D\bibinitperiod\bibinitdelim A\bibinitperiod}}}%
        {{hash=447f22d7de07f4cbd918f6b2abb97288}{%
           family={Demichel},
           familyi={D\bibinitperiod},
           given={Patrick},
           giveni={P\bibinitperiod}}}%
      }
      \strng{namehash}{307314f703dd0ce2792f1f0923e08ca4}
      \strng{fullhash}{307314f703dd0ce2792f1f0923e08ca4}
      \strng{bibnamehash}{307314f703dd0ce2792f1f0923e08ca4}
      \strng{authorbibnamehash}{307314f703dd0ce2792f1f0923e08ca4}
      \strng{authornamehash}{307314f703dd0ce2792f1f0923e08ca4}
      \strng{authorfullhash}{307314f703dd0ce2792f1f0923e08ca4}
      \field{sortinit}{2}
      \field{sortinithash}{ed39bb39cf854d5250e95b1c1f94f4ed}
      \field{labelnamesource}{author}
      \field{labeltitlesource}{title}
      \field{annotation}{The authors analyze $E^\pi(x)$ for $x<10^{10^{13}}$ using the first $2\times10^{11}$ nontrivial zeros of $\zeta(s)$. Based on numerical computation, the authors suggest that there may exist an~$x$ near $1.397162914 \times 10^{316}$ such that $\pi(x)>\li(x)$. The authors also conjecture that $|E^\pi(x)|<\frac1e(\log \log \log x+e+1)$.\par This article cites~\cites{MR1752093,MR0369287,MR2652902,1932.Ingham,MR2420864,MR0202686,1914.Littlewood,MR0137689,MR1329368,MR2684372,MR1573970,MR0067145}.}
      \field{issn}{0025-5718}
      \field{journaltitle}{Math. Comp.}
      \field{number}{276}
      \field{title}{The impact of {$\zeta(s)$} complex zeros on {$\pi(x)$} for {$x<10^{10^{13}}$}}
      \field{volume}{80}
      \field{year}{2011}
      \true{nocite}
      \field{pages}{2381\bibrangedash 2394}
      \range{pages}{14}
      \verb{doi}
      \verb 10.1090/S0025-5718-2011-02477-4
      \endverb
      \verb{mrnumber}
      \verb 2813366
      \endverb
    \endentry
    \entry{MR2987108}{article}{}
      \name{author}{2}{}{%
        {{hash=3028d3b5860b77006bb46a3a93b688b4}{%
           family={Kunik},
           familyi={K\bibinitperiod},
           given={Matthias},
           giveni={M\bibinitperiod}}}%
        {{hash=cdaf4b1735b29361604e44f0ef594303}{%
           family={Lucht},
           familyi={L\bibinitperiod},
           given={Lutz\bibnamedelima G.},
           giveni={L\bibinitperiod\bibinitdelim G\bibinitperiod}}}%
      }
      \strng{namehash}{0ff5403c60f18b96de7464513572a2ac}
      \strng{fullhash}{0ff5403c60f18b96de7464513572a2ac}
      \strng{bibnamehash}{0ff5403c60f18b96de7464513572a2ac}
      \strng{authorbibnamehash}{0ff5403c60f18b96de7464513572a2ac}
      \strng{authornamehash}{0ff5403c60f18b96de7464513572a2ac}
      \strng{authorfullhash}{0ff5403c60f18b96de7464513572a2ac}
      \field{sortinit}{2}
      \field{sortinithash}{ed39bb39cf854d5250e95b1c1f94f4ed}
      \field{labelnamesource}{author}
      \field{labeltitlesource}{title}
      \field{annotation}{In this article, the authors investigate the series $F(w)= \sum_{n=1}^{\infty} \frac{\Lambda(n)}{n}w^n$ for $|w|\le1$. The authors use a more convenient direct method compared to that of Hardy and Littlewood in~\cite{MR1555148} to derive the explicit formulas for the series $\sum \Lambda(n)e^{-nz} = 1/z + (\cosh(z)-1)\log z + T(z) - \sum_ \rho \Gamma(\rho)z^{-\rho}$ and for $F(e^{-z})$; in particular the authors obtain a closed form for the entire function $T(z)$. The main result reveals logarithmic singularities of $F(e^{2\pi it})$ at the reduced rational numbers $t=\frac{a}{q}$ with squarefree denominator~$q \in \N$.\par This article cites~\cites{MR1555148}.}
      \field{isbn}{978-83-232-2442-6}
      \field{issn}{0208-6573}
      \field{journaltitle}{Funct. Approx. Comment. Math.}
      \field{number}{part 1}
      \field{part}{part 1}
      \field{title}{Power series with the von {M}angoldt function}
      \field{volume}{47}
      \field{year}{2012}
      \true{nocite}
      \field{pages}{15\bibrangedash 33}
      \range{pages}{19}
      \verb{doi}
      \verb 10.7169/facm/2012.47.1.2
      \endverb
      \verb{mrnumber}
      \verb 2987108
      \endverb
    \endentry
    \entry{MR2943671}{article}{}
      \name{author}{1}{}{%
        {{hash=a01ad5bbada0f5c69f81f1149c07aefc}{%
           family={Lamzouri},
           familyi={L\bibinitperiod},
           given={Y.},
           giveni={Y\bibinitperiod}}}%
      }
      \strng{namehash}{a01ad5bbada0f5c69f81f1149c07aefc}
      \strng{fullhash}{a01ad5bbada0f5c69f81f1149c07aefc}
      \strng{bibnamehash}{a01ad5bbada0f5c69f81f1149c07aefc}
      \strng{authorbibnamehash}{a01ad5bbada0f5c69f81f1149c07aefc}
      \strng{authornamehash}{a01ad5bbada0f5c69f81f1149c07aefc}
      \strng{authorfullhash}{a01ad5bbada0f5c69f81f1149c07aefc}
      \field{extraname}{1}
      \field{sortinit}{2}
      \field{sortinithash}{ed39bb39cf854d5250e95b1c1f94f4ed}
      \field{labelnamesource}{author}
      \field{labeltitlesource}{title}
      \field{annotation}{The author refines tail estimates for the distribution $\mu_{q;a_1,\ldots,a_r}$ given by Rubinstein and Sarnak~\cite{MR1329368}. From the Math Review by D.~R.~Heath-Brown, with slight changes in notation: ``Define \[ \sigma(q) = \biggl( 2 \sum_{\substack{\chi\mod q \\ \chi\ne\chi_0}} \sum_{\gamma>0} \frac1{1/4+\gamma^2} \biggr)^{1/2}, \] with $\gamma$ running over ordinates of zeros of $L(s,\chi)$. \ldots{} It is proved here that if $0<\lambda\le\sqrt{\log\log q}$ then \[ \mu_{q;a_1,\ldots,a_r} \bigl( \|t\|^2 > \lambda\sigma(q) \bigr) = (2\pi)^{-r/2} \int_{\|t\|_2>\lambda} \exp\bigl( -\tfrac12\|t\|_2 \bigr) \,dt + O_r\bigl( (\log q)^{-2} \bigr). \] \ldots{} For large~$V$ it is shown that \[ \exp\biggl( -c_1 \frac{V^2}{\phi(q)\log q} \biggr) \ll_q \mu_{q;a_1,\ldots,a_r} \bigl( \|t\|^2 > V \bigr) \ll_q \exp\biggl( -c_2 \frac{V^2}{\phi(q)\log q} \biggr) \] for $(\phi(q)\log q)^{1/2} \ll V \le A\phi(q)\log q$, with constants $c_1,c_2$ depending only on~$r$ and~$A$. For even larger~$V$ the paper gives further results, showing that there are transitions in the behaviour, occuring around $V=\phi(q)\log q$ and $V=\phi(q)\log^2 q$.''\par This article cites~\cites{MR3028758,MR0146156,MR0146157,MR0146158,MR0146157,MR0146160,MR0146161,MR0156826,MR0156827,MR1251926,MR1367576,MR2631002,MR634679,MR1329368}.}
      \field{issn}{0305-0041}
      \field{journaltitle}{Math. Proc. Cambridge Philos. Soc.}
      \field{number}{1}
      \field{title}{Large deviations of the limiting distribution in the {S}hanks--{R}ényi prime number race}
      \field{volume}{153}
      \field{year}{2012}
      \true{nocite}
      \field{pages}{147\bibrangedash 166}
      \range{pages}{20}
      \verb{doi}
      \verb 10.1017/S030500411200014X
      \endverb
      \verb{mrnumber}
      \verb 2943671
      \endverb
      \verb{urlraw}
      \verb https://doi.org/10.1017/S030500411200014X
      \endverb
      \verb{url}
      \verb https://doi.org/10.1017/S030500411200014X
      \endverb
    \endentry
    \entry{MR2998146}{article}{}
      \name{author}{1}{}{%
        {{hash=a01ad5bbada0f5c69f81f1149c07aefc}{%
           family={Lamzouri},
           familyi={L\bibinitperiod},
           given={Y.},
           giveni={Y\bibinitperiod}}}%
      }
      \strng{namehash}{a01ad5bbada0f5c69f81f1149c07aefc}
      \strng{fullhash}{a01ad5bbada0f5c69f81f1149c07aefc}
      \strng{bibnamehash}{a01ad5bbada0f5c69f81f1149c07aefc}
      \strng{authorbibnamehash}{a01ad5bbada0f5c69f81f1149c07aefc}
      \strng{authornamehash}{a01ad5bbada0f5c69f81f1149c07aefc}
      \strng{authorfullhash}{a01ad5bbada0f5c69f81f1149c07aefc}
      \field{extraname}{2}
      \field{sortinit}{2}
      \field{sortinithash}{ed39bb39cf854d5250e95b1c1f94f4ed}
      \field{labelnamesource}{author}
      \field{labeltitlesource}{title}
      \field{annotation}{In this article, the author examines the order of magnitude of the density $\delta_{q;a_1,\ldots,a_r}$ when the number of contestants~$r$ tends to infinity as $q\to\infty$. Assuming GRH and LI, he shows that $\delta_{q;a_1,\ldots,a_r}=\frac{1}{r!}(1+O(\frac{r^2}{\log{q}}))$ when $r\leq \sqrt{\log{q}}$, while $\delta_{q;a_1,\ldots,a_r}=\exp(-r\log{r}+r+O(\log{r}+\frac{r^2}{\log{q}}))$ when $\sqrt{\log{q}}\ll r\leq (1-\ep)\log{q}/\log\log{q}$. For larger~$r$, he provides the upper bound $\delta_{q;a_1,\ldots,a_r} \ll_{\ep} q^{-1+2\ep}$ for $(1-\ep)\log{q}/\log\log{q}\leq r\leq\phi(q)$.\par This article cites~\cites{MR0460261,MR1806291,MR3028758,MR1909220,MR2075622,MR1251926,MR1256149,MR1367576,MR0146156,MR0146157,MR0146158,MR0146157,MR0146160,MR0146161,MR0156826,MR0156827,MR2943671,1914.Littlewood,MR1956261,MR1329368}.}
      \field{issn}{1073-2780}
      \field{journaltitle}{Math. Res. Lett.}
      \field{number}{3}
      \field{title}{The {S}hanks-{R}ényi prime number race with many contestants}
      \field{volume}{19}
      \field{year}{2012}
      \true{nocite}
      \field{pages}{649\bibrangedash 666}
      \range{pages}{18}
      \verb{doi}
      \verb 10.4310/MRL.2012.v19.n3.a11
      \endverb
      \verb{mrnumber}
      \verb 2998146
      \endverb
      \verb{urlraw}
      \verb https://doi.org/10.4310/MRL.2012.v19.n3.a11
      \endverb
      \verb{url}
      \verb https://doi.org/10.4310/MRL.2012.v19.n3.a11
      \endverb
    \endentry
    \entry{MR3062002}{article}{}
      \name{author}{2}{}{%
        {{hash=454d3dedc781ec4a49112de10313b411}{%
           family={Mossinghoff},
           familyi={M\bibinitperiod},
           given={Michael\bibnamedelima J.},
           giveni={M\bibinitperiod\bibinitdelim J\bibinitperiod}}}%
        {{hash=66490333fa52c597ed64f7abf27c7710}{%
           family={Trudgian},
           familyi={T\bibinitperiod},
           given={Timothy\bibnamedelima S.},
           giveni={T\bibinitperiod\bibinitdelim S\bibinitperiod}}}%
      }
      \strng{namehash}{901cb3f7453290119b926d8dd538a562}
      \strng{fullhash}{901cb3f7453290119b926d8dd538a562}
      \strng{bibnamehash}{901cb3f7453290119b926d8dd538a562}
      \strng{authorbibnamehash}{901cb3f7453290119b926d8dd538a562}
      \strng{authornamehash}{901cb3f7453290119b926d8dd538a562}
      \strng{authorfullhash}{901cb3f7453290119b926d8dd538a562}
      \field{extraname}{1}
      \field{sortinit}{2}
      \field{sortinithash}{ed39bb39cf854d5250e95b1c1f94f4ed}
      \field{labelnamesource}{author}
      \field{labeltitlesource}{title}
      \field{annotation}{For $\alpha\in[0,1]$, define $L_\alpha(x) = \sum_{n\le x} (-1)^{\Omega(n)}/n^\alpha$. The authors generalize classical results of Pólya and Turán by proving that RH is equivalent to the estimate $L_\alpha(x) \ll_{\alpha,\ep} x^{1/2-\alpha+\ep}$ for all $\ep>0$, the equivalence being for any (hence all) $\alpha\in[0,1]$. Modifying the function slightly by defining \[ {\mathcal L}_\alpha(x) = L_\alpha(x) - \begin{cases} 0, &\text{if }0\le\alpha<\frac12, \\ \log x/2\zeta(\frac12), &\text{if }\alpha=\frac12, \\ \zeta(2\alpha)/\zeta(\alpha), &\text{if }\frac12<\alpha\le1, \end{cases} \] they also show that the assertion, for any $c\in\R$, that ${\mathcal L}_\alpha(x) - cx^{1/2-\alpha}$ has constant sign implies that all the zeros of $\zeta(s)$ are simple but is incompatible with~LI, even with finitely many exceptions. (The assertion that $L_{1/2}(x) - c$ itself has constant sign is shown to imply that all zeros of $\zeta(s)$ have multiplicity at most~$2$, but is not known to be inconsistent with~LI.) The authors state the problem of showing that ${\mathcal L}_\alpha(x)$ has infinitely many sign changes, which is known only for $\alpha=0$ and $\alpha=1$~\cite{MR0104638}, and also of showing that $L_{1/2}(x)\le0$ for~$x$ sufficiently large, perhaps $x\ge17$.\par This article cites~\cites{MR27305,MR0006202,MR0104638,MR1329368,MR2398787,MR0120198,MR0882550,MR584557}.}
      \field{issn}{1446-7887}
      \field{journaltitle}{J. Aust. Math. Soc.}
      \field{number}{1--2}
      \field{title}{Between the problems of {P}ólya and {T}urán}
      \field{volume}{93}
      \field{year}{2012}
      \true{nocite}
      \field{pages}{157\bibrangedash 171}
      \range{pages}{15}
      \verb{doi}
      \verb 10.1017/S1446788712000201
      \endverb
      \verb{mrnumber}
      \verb 3062002
      \endverb
    \endentry
    \entry{MR3028758}{article}{}
      \name{author}{2}{}{%
        {{hash=889ecd9d3a804c5a7e9a5d8287766657}{%
           family={Fiorilli},
           familyi={F\bibinitperiod},
           given={D.},
           giveni={D\bibinitperiod}}}%
        {{hash=feda7d3cba2a0ad4e772307fca655937}{%
           family={Martin},
           familyi={M\bibinitperiod},
           given={G.},
           giveni={G\bibinitperiod}}}%
      }
      \strng{namehash}{8f8b9fef368ad32f0c071fbab617f456}
      \strng{fullhash}{8f8b9fef368ad32f0c071fbab617f456}
      \strng{bibnamehash}{8f8b9fef368ad32f0c071fbab617f456}
      \strng{authorbibnamehash}{8f8b9fef368ad32f0c071fbab617f456}
      \strng{authornamehash}{8f8b9fef368ad32f0c071fbab617f456}
      \strng{authorfullhash}{8f8b9fef368ad32f0c071fbab617f456}
      \field{extraname}{1}
      \field{sortinit}{2}
      \field{sortinithash}{ed39bb39cf854d5250e95b1c1f94f4ed}
      \field{labelnamesource}{author}
      \field{labeltitlesource}{title}
      \field{annotation}{Assuming GRH and LI, the authors establish an asymptotic series for $\delta_{q;a,b}$ whose error term can be taken to be any negative power of~$q$. The first asymptotic formula given by this series, in the case where~$a$ is a quadratic nonresidue and~$b$ is a quadratic residue\mod q, is \[ \delta_{q;a,b} = \frac12 + \frac{\rho(q)}{\sqrt{2\pi V(q;a,b)}} + O_\ep(q^{-3/2+\ep}), \] where $V(q;a,b) = \sum_{\chi\mod q} |\chi(b)-\chi(a)|^2\sum_\rho 1/(\frac14+\gamma^2)$ and $\rho(q)$ is the number of square roots of~$1\mod q$. The authors give a closed form for $V(q;a,b)$ in terms of arithmetic properties of its arguments, which allow them to compare various values of $\delta_{q;a,b}$; for example, for fixed $a\ne-1$, they show (always assuming GRH and LI) that when~$q$ is sufficiently large, $\delta_{q;-1,1} < \delta_{q;a,1}$ whenever~$-1$ and~$a$ are both quadratic nonresidues\mod q. Explicit bounds for the error terms in this formula allow the authors to calculate the list of all~$117$ two-way prime number races (up to symmetries) for which $\delta_{q;a,b} > 0.9$, the largest of which is $\delta_{24,5,1} \approx 0.999988$.\par This article cites~\cites{MR0108470,MR0137689,MR0506080,MR1329368,MR1806291,MR1816036}.}
      \field{issn}{0075-4102}
      \field{journaltitle}{J. Reine Angew. Math.}
      \field{title}{Inequities in the {S}hanks-{R}ényi prime number race: an asymptotic formula for the densities}
      \field{volume}{676}
      \field{year}{2013}
      \true{nocite}
      \field{pages}{121\bibrangedash 212}
      \range{pages}{92}
      \verb{mrnumber}
      \verb 3028758
      \endverb
    \endentry
    \entry{MR3151605}{article}{}
      \name{author}{3}{}{%
        {{hash=e7e31df1803902b56a78f32b1c733598}{%
           family={Ford},
           familyi={F\bibinitperiod},
           given={K.},
           giveni={K\bibinitperiod}}}%
        {{hash=a01ad5bbada0f5c69f81f1149c07aefc}{%
           family={Lamzouri},
           familyi={L\bibinitperiod},
           given={Y.},
           giveni={Y\bibinitperiod}}}%
        {{hash=c48086b9b194fded979a790687062dbb}{%
           family={Konyagin},
           familyi={K\bibinitperiod},
           given={S.},
           giveni={S\bibinitperiod}}}%
      }
      \strng{namehash}{63af0df95eb4cedba38dcfa76456f7cf}
      \strng{fullhash}{63af0df95eb4cedba38dcfa76456f7cf}
      \strng{bibnamehash}{63af0df95eb4cedba38dcfa76456f7cf}
      \strng{authorbibnamehash}{63af0df95eb4cedba38dcfa76456f7cf}
      \strng{authornamehash}{63af0df95eb4cedba38dcfa76456f7cf}
      \strng{authorfullhash}{63af0df95eb4cedba38dcfa76456f7cf}
      \field{sortinit}{2}
      \field{sortinithash}{ed39bb39cf854d5250e95b1c1f94f4ed}
      \field{labelnamesource}{author}
      \field{labeltitlesource}{title}
      \field{annotation}{The authors construct ``barriers'' for two-way prime number races, including $\pi(x) - \li(x)$. They show that a certain hypothetical configuration of zeros of $L(s,\chi)$ would imply that $\{x \geq 2 \colon \pi(x;q,b)>\pi(x;q,a)\}$ has asymptotic density~$0$. Furthermore, they show that a certain hypothetical configuration of the zeros of $\zeta(s)$ would imply that $\{x \geq 2 \colon \pi(x) > \li(x)\}$ has asymptotic density~$0$, while a related configuration would result in the same set having asymptotic density~$1$.\par This article cites~\cites{MR1909220,MR1985941,MR2075622,MR2202918,MR1251926,MR1256149,MR0146156,MR0146157,MR0146158,MR0146159,MR0146160,MR0146161,MR0156826,MR0156827,MR0162771,MR0174538,MR0180539,MR0182616A,MR0182616B,MR0200250,MR0302585,1914.Littlewood,MR1329368,MR2753165}.}
      \field{issn}{0033-5606}
      \field{journaltitle}{Q. J. Math.}
      \field{number}{4}
      \field{title}{The prime number race and zeros of {D}irichlet {$L$}-functions off the critical line: {P}art {III}}
      \field{volume}{64}
      \field{year}{2013}
      \true{nocite}
      \field{pages}{1091\bibrangedash 1098}
      \range{pages}{8}
      \verb{doi}
      \verb 10.1093/qmath/has021
      \endverb
      \verb{mrnumber}
      \verb 3151605
      \endverb
      \verb{urlraw}
      \verb https://doi.org/10.1093/qmath/has021
      \endverb
      \verb{url}
      \verb https://doi.org/10.1093/qmath/has021
      \endverb
    \endentry
    \entry{MR2994374}{article}{}
      \name{author}{1}{}{%
        {{hash=48017d687b31d05c591036588e04dd66}{%
           family={Humphries},
           familyi={H\bibinitperiod},
           given={Peter},
           giveni={P\bibinitperiod}}}%
      }
      \strng{namehash}{48017d687b31d05c591036588e04dd66}
      \strng{fullhash}{48017d687b31d05c591036588e04dd66}
      \strng{bibnamehash}{48017d687b31d05c591036588e04dd66}
      \strng{authorbibnamehash}{48017d687b31d05c591036588e04dd66}
      \strng{authornamehash}{48017d687b31d05c591036588e04dd66}
      \strng{authorfullhash}{48017d687b31d05c591036588e04dd66}
      \field{extraname}{1}
      \field{sortinit}{2}
      \field{sortinithash}{ed39bb39cf854d5250e95b1c1f94f4ed}
      \field{labelnamesource}{author}
      \field{labeltitlesource}{title}
      \field{annotation}{The author gives an excellent review of the history and known results concerning the weighted Liouville summatory function $L_\alpha(x) = \sum_{n\le x} \lambda(n)/n^\alpha$. Define $\delta_\alpha$ to be the logarithmic density of the set of positive real numbers~$x$ such that $L_\alpha(x) \le 0$. Assuming RH and LI and the estimate $\sum_{0<\gamma<T} |\zeta'(\rho)|^{-2} \ll T$, the author shows when $0\le\alpha<\frac12$ that $\frac12\le\delta_\alpha<1$, so that in particular, $L_\alpha(x)$ changes sign infinitely often. The author additionally shows under the same hypotheses that $\lim_{\alpha\to1/2-} \delta_\alpha = 1$. Moreover, the author proves (without needing the hypothesis LI) that $\delta_{1/2} = 1$, lending some credence to the conjecture that $L_{1/2}(x) \le 0$ for $x\ge17$.\par This article cites~\cites{1919.Polya,2000.Ng,MR4082248,MR0006202,MR0104638,MR0330069,MR634679,MR654520,MR1329368,MR1806291,MR2078705,MR2398787,MR3062002,MR3063909,MR98723,MR584557,MR43841}.}
      \field{issn}{0022-314X}
      \field{journaltitle}{J. Number Theory}
      \field{number}{2}
      \field{title}{The distribution of weighted sums of the {L}iouville function and {P}ólya's conjecture}
      \field{volume}{133}
      \field{year}{2013}
      \true{nocite}
      \field{pages}{545\bibrangedash 582}
      \range{pages}{38}
      \verb{doi}
      \verb 10.1016/j.jnt.2012.08.011
      \endverb
      \verb{mrnumber}
      \verb 2994374
      \endverb
    \endentry
    \entry{MR3063909}{article}{}
      \name{author}{1}{}{%
        {{hash=a01ad5bbada0f5c69f81f1149c07aefc}{%
           family={Lamzouri},
           familyi={L\bibinitperiod},
           given={Y.},
           giveni={Y\bibinitperiod}}}%
      }
      \strng{namehash}{a01ad5bbada0f5c69f81f1149c07aefc}
      \strng{fullhash}{a01ad5bbada0f5c69f81f1149c07aefc}
      \strng{bibnamehash}{a01ad5bbada0f5c69f81f1149c07aefc}
      \strng{authorbibnamehash}{a01ad5bbada0f5c69f81f1149c07aefc}
      \strng{authornamehash}{a01ad5bbada0f5c69f81f1149c07aefc}
      \strng{authorfullhash}{a01ad5bbada0f5c69f81f1149c07aefc}
      \field{extraname}{3}
      \field{sortinit}{2}
      \field{sortinithash}{ed39bb39cf854d5250e95b1c1f94f4ed}
      \field{labelnamesource}{author}
      \field{labeltitlesource}{title}
      \field{annotation}{Assuming GRH and LI, the author describes phenomena that occur for multi-way prime number races that are not present in two-way races. When $r\ge3$, the author shows that $\bigl| \delta_{q;a_1,\ldots,a_r} - \frac1{r!} \bigr| = \Omega(1/\log q)$, in contrast to the $r=2$ case~\cite{MR3028758} where $\bigl| \delta_{q;a_1,a_2} - \frac12 \bigr| \ll q^{-1/2+o(1)}$. The author also shows that when~$q$ is sufficiently large in terms of~$r$, there are always~$r$ squares $a_1,\ldots,a_r$ modulo~$q$ for which $\delta_{q;a_1,\ldots,a_r} \ne \frac1{r!}$ (and similarly~$r$ nonsquares). The method of proof is to study the approximation of the characteristic function $\hat{\mu}_{q;a_1,\ldots,a_r}$ by that of a multivariate Gaussian, paying particular attention to secondary main terms that arise.\par This article cites~\cites{MR0447090,MR1806291,MR3028758,MR1909220,MR2202918,MR0146156,MR0146157,MR0146158,MR0146159,MR0146160,MR0146161,MR0156826,MR0156827,MR1251926,MR1256149,MR1367576,MR2998146,MR1329368}.}
      \field{issn}{0025-5831}
      \field{journaltitle}{Math. Ann.}
      \field{number}{3}
      \field{title}{Prime number races with three or more competitors}
      \field{volume}{356}
      \field{year}{2013}
      \true{nocite}
      \field{pages}{1117\bibrangedash 1162}
      \range{pages}{46}
      \verb{doi}
      \verb 10.1007/s00208-012-0874-1
      \endverb
      \verb{mrnumber}
      \verb 3063909
      \endverb
      \verb{urlraw}
      \verb https://doi.org/10.1007/s00208-012-0874-1
      \endverb
      \verb{url}
      \verb https://doi.org/10.1007/s00208-012-0874-1
      \endverb
    \endentry
    \entry{arXiv.Myerscough}{online}{}
      \name{author}{1}{}{%
        {{hash=3a7e707950739f23c95b5c3b3aee3641}{%
           family={Myerscough},
           familyi={M\bibinitperiod},
           given={C.},
           giveni={C\bibinitperiod}}}%
      }
      \strng{namehash}{3a7e707950739f23c95b5c3b3aee3641}
      \strng{fullhash}{3a7e707950739f23c95b5c3b3aee3641}
      \strng{bibnamehash}{3a7e707950739f23c95b5c3b3aee3641}
      \strng{authorbibnamehash}{3a7e707950739f23c95b5c3b3aee3641}
      \strng{authornamehash}{3a7e707950739f23c95b5c3b3aee3641}
      \strng{authorfullhash}{3a7e707950739f23c95b5c3b3aee3641}
      \field{sortinit}{2}
      \field{sortinithash}{ed39bb39cf854d5250e95b1c1f94f4ed}
      \field{labelnamesource}{author}
      \field{labeltitlesource}{title}
      \field{annotation}{The author studies the family $P_\mu(t)$ of density functions of the sum of random variables \begin{align*} \sum_\chi|\chi(a)-\chi(b)|\sum_{\gamma>\mu}\frac{2\Re Z_\gamma}{\sqrt{1/4+\gamma^2}} \end{align*} where $Z_\gamma$ are independent random variables, uniformly distributed on the unit circle, indexed by the ordinates of zeros of $L(s,\chi)$; the motivation is that $P_0(t)$ has the same distribution as the prime number race measure $\mu_{q;a,b}$. The author discovers that when $t$ is small, the behaviour of $P_\mu(t)$ approaches the normal distribution over a wider and wider range of $t/\sigma_\mu$ as $\mu$ increases, while when $t$ is large, $P_\mu(t)$ decays significantly faster than a normal distribution.\par To explicitly evaluate $P_\mu(t)$, the author compares the methods of the steepest descent (to third order), numerical convolution methods, and the Rubinstein--Sarnak method, and notes that there is an agreement to within $0.001\%$ for $1\leq t\leq 1.2$ (and similar observations for other ranges of~$t$). The author points out that Rubinstein--Sarnak is the best way to obtain results which are much larger than the absolute accuracy of computation, while the steepest descent method is desirable for extreme deviations, and the convolution method is valuable for some intermediate ranges.\par This article cites~\cites{MR1651741, MR1752093, MR1816036, MR1806291, MR3272280, MR3028758, 1932.Ingham, MR2943671, MR2037293, MR0202686, 1914.Littlewood, MR1956261, MR2631002, MR1329368}}
      \field{title}{Application of an accurate remainder term in the calculation of residue class distributions}
      \field{year}{2013}
      \true{nocite}
      \verb{urlraw}
      \verb https://arxiv.org/abs/1301.1434
      \endverb
      \verb{url}
      \verb https://arxiv.org/abs/1301.1434
      \endverb
    \endentry
    \entry{MR3206041}{article}{}
      \name{author}{1}{}{%
        {{hash=9f2b575b642afbcf502733e2015df915}{%
           family={Petrushov},
           familyi={P\bibinitperiod},
           given={O.\bibnamedelimi A.},
           giveni={O\bibinitperiod\bibinitdelim A\bibinitperiod}}}%
      }
      \strng{namehash}{9f2b575b642afbcf502733e2015df915}
      \strng{fullhash}{9f2b575b642afbcf502733e2015df915}
      \strng{bibnamehash}{9f2b575b642afbcf502733e2015df915}
      \strng{authorbibnamehash}{9f2b575b642afbcf502733e2015df915}
      \strng{authornamehash}{9f2b575b642afbcf502733e2015df915}
      \strng{authorfullhash}{9f2b575b642afbcf502733e2015df915}
      \field{sortinit}{2}
      \field{sortinithash}{ed39bb39cf854d5250e95b1c1f94f4ed}
      \field{labelnamesource}{author}
      \field{labeltitlesource}{title}
      \field{annotation}{The author establishes results on the limit supremum and infimum of a real continuous function $\nu(t)$ under assumptions on the behaviour of the Laplace transform $\int_0^\infty e^{-st} \,d\nu(t)$ near its singular point. As a corollary, the author obtains results on the limiting behaviour of the function $M_a(x) = \sum_{n < x} \mu(n) n^{-a}$ for $0\leq a \le \frac12$ in terms of functions related to $1/\zeta(s)$.\par This article cites~\cite{MR783538}.}
      \field{issn}{0001-4346}
      \field{journaltitle}{Math. Notes}
      \field{number}{5--6}
      \field{title}{Asymptotic estimates of functions based on the behavior of their {L}aplace transforms near singular points}
      \field{volume}{93}
      \field{year}{2013}
      \true{nocite}
      \field{pages}{906\bibrangedash 916}
      \range{pages}{11}
      \verb{doi}
      \verb 10.1134/S0001434613050283
      \endverb
      \verb{mrnumber}
      \verb 3206041
      \endverb
    \endentry
    \entry{MR3261965}{article}{}
      \name{author}{3}{}{%
        {{hash=3d7986c1f1aad61da9a4782349cfc695}{%
           family={Akbary},
           familyi={A\bibinitperiod},
           given={A.},
           giveni={A\bibinitperiod}}}%
        {{hash=585b691d1c27d2773434ed1ac9664075}{%
           family={Ng},
           familyi={N\bibinitperiod},
           given={N.},
           giveni={N\bibinitperiod}}}%
        {{hash=77b8d807a8642e05d0c1744c3120354e}{%
           family={Shahabi},
           familyi={S\bibinitperiod},
           given={M.},
           giveni={M\bibinitperiod}}}%
      }
      \strng{namehash}{87ed2f0893dc8a9968837b23ecc3e66b}
      \strng{fullhash}{87ed2f0893dc8a9968837b23ecc3e66b}
      \strng{bibnamehash}{87ed2f0893dc8a9968837b23ecc3e66b}
      \strng{authorbibnamehash}{87ed2f0893dc8a9968837b23ecc3e66b}
      \strng{authornamehash}{87ed2f0893dc8a9968837b23ecc3e66b}
      \strng{authorfullhash}{87ed2f0893dc8a9968837b23ecc3e66b}
      \field{sortinit}{2}
      \field{sortinithash}{ed39bb39cf854d5250e95b1c1f94f4ed}
      \field{labelnamesource}{author}
      \field{labeltitlesource}{title}
      \field{annotation}{The authors establish a sufficient condition for when general functions of the form \[\Re \sum_{\lambda_n\le X} r_n e^{i\lambda_n y} + E(y,X)\] (or vector-valued analogues) are $B^2$-almost-periodic and therefore possess limiting distributions. They apply this theorem, assuming the appropriate RH, to (logarithmically scaled versions of) functions such as: $\sum_{n\le x} \Lambda(n) a_n(\pi)$ for automorphic $L$-functions $L(s,\chi)$; the Mertens and Liouville functions $M(x)$ and $L(x)$, analogues of these for arithmetic progressions, and interpolations between these and $M_r(x)$ and $L_r(x)$; and (on the further assumption that Artin $L$-functions are entire) for prime counting functions in Chebotarev conjugacy classes. They provide further conclusions if the linear independence of the $\lambda_n$ is also assumed.\par This article cites~\cites{1938.Wintner,2000.Ng,MR0104638,MR598886,MR783538,MR1159354,MR1329368,MR1501802,MR1507933,MR1544509,MR1806291,MR2398787,MR2994374,MR3028758,MR3063909,MR3264673,MR3272280,MR2078705,MR584557}.}
      \field{issn}{0033-5606}
      \field{journaltitle}{Q. J. Math.}
      \field{number}{3}
      \field{title}{Limiting distributions of the classical error terms of prime number theory}
      \field{volume}{65}
      \field{year}{2014}
      \true{nocite}
      \field{pages}{743\bibrangedash 780}
      \range{pages}{38}
      \verb{doi}
      \verb 10.1093/qmath/hat059
      \endverb
      \verb{mrnumber}
      \verb 3261965
      \endverb
      \verb{urlraw}
      \verb https://doi.org/10.1093/qmath/hat059
      \endverb
      \verb{url}
      \verb https://doi.org/10.1093/qmath/hat059
      \endverb
    \endentry
    \entry{MR3306734}{article}{}
      \name{author}{3}{}{%
        {{hash=550b5a0b5fd42a70ac3723323eef8887}{%
           family={Chaubey},
           familyi={C\bibinitperiod},
           given={Sneha},
           giveni={S\bibinitperiod}}}%
        {{hash=259966712b379a8c7cdc8ece86803924}{%
           family={Lanius},
           familyi={L\bibinitperiod},
           given={Melinda},
           giveni={M\bibinitperiod}}}%
        {{hash=db9b605a4c37f759795e5a1699044f16}{%
           family={Zaharescu},
           familyi={Z\bibinitperiod},
           given={Alexandru},
           giveni={A\bibinitperiod}}}%
      }
      \strng{namehash}{f307e9735e20f51a6fae94d9f56d696c}
      \strng{fullhash}{f307e9735e20f51a6fae94d9f56d696c}
      \strng{bibnamehash}{f307e9735e20f51a6fae94d9f56d696c}
      \strng{authorbibnamehash}{f307e9735e20f51a6fae94d9f56d696c}
      \strng{authornamehash}{f307e9735e20f51a6fae94d9f56d696c}
      \strng{authorfullhash}{f307e9735e20f51a6fae94d9f56d696c}
      \field{sortinit}{2}
      \field{sortinithash}{ed39bb39cf854d5250e95b1c1f94f4ed}
      \field{labelnamesource}{author}
      \field{labeltitlesource}{title}
      \field{annotation}{Atanassov defined the ``irrational factor'' of the number $n=p_1^{r_1}\cdots p_k^{r_k}$ to be $I(n) = p_1^{1/r_1}\cdots p_k^{1/r_k}$. The authors establish an asymptotic formula for the summatory function of $I(n)$ over a reduced residue class. They then use a Landau-type argument to show that the race between the summatory functions of $I(n)$ over integers that are $1\mod 3$ and integers that are $2\mod 3$ has infinitely many lead changes; they remark that a theorem from~\cite{MR2506435} shows that the number of sign changes in $[1,T]$ is $\gg\log T$ when~$T$ is sufficiently large.\par This article cites~\cites{1914.Littlewood,MR0146156,MR858395,MR1511360,MR2506435}.}
      \field{issn}{0253-4142}
      \field{journaltitle}{Proc. Indian Acad. Sci. Math. Sci.}
      \field{number}{4}
      \field{title}{Irrational factor races}
      \field{volume}{124}
      \field{year}{2014}
      \true{nocite}
      \field{pages}{471\bibrangedash 479}
      \range{pages}{9}
      \verb{doi}
      \verb 10.1007/s12044-014-0198-z
      \endverb
      \verb{mrnumber}
      \verb 3306734
      \endverb
    \endentry
    \entry{MR3264673}{article}{}
      \name{author}{1}{}{%
        {{hash=889ecd9d3a804c5a7e9a5d8287766657}{%
           family={Fiorilli},
           familyi={F\bibinitperiod},
           given={D.},
           giveni={D\bibinitperiod}}}%
      }
      \strng{namehash}{889ecd9d3a804c5a7e9a5d8287766657}
      \strng{fullhash}{889ecd9d3a804c5a7e9a5d8287766657}
      \strng{bibnamehash}{889ecd9d3a804c5a7e9a5d8287766657}
      \strng{authorbibnamehash}{889ecd9d3a804c5a7e9a5d8287766657}
      \strng{authornamehash}{889ecd9d3a804c5a7e9a5d8287766657}
      \strng{authorfullhash}{889ecd9d3a804c5a7e9a5d8287766657}
      \field{extraname}{2}
      \field{sortinit}{2}
      \field{sortinithash}{ed39bb39cf854d5250e95b1c1f94f4ed}
      \field{labelnamesource}{author}
      \field{labeltitlesource}{title}
      \field{annotation}{In this article, the author extends the results from~\cite{2007.Sarnak} on ``elliptic curve prime number races'', the races between primes for which the number of points on a given elliptic curve over $\mathbb{F}_p$ is greater or less than $p+1$. Assuming GRH and LI for $L$-functions of elliptic curves, the author establishes an equivalence between two phenomena. The first phenomenon is unbounded analytic ranks for elliptic curves---more precisely, the conjecture that $\limsup_{N_E\to\infty}r_{\text{an}}(E)/\sqrt{\log(N_E)}=\infty$ where $N_E$ is the conductor of~$E$. The second phenomenon is the existence of arbitrarily biased elliptic curve prime races, that is, races for which the logarithmic density of the set $\{t:S(t)=-\sum_{p\le t}a_pp^{-1/2} \geq 0\}$ is arbitrarily close to~$1$. When showing that the first phenomenon implies the second, LI can be weakened to the hypothesis that the nonreal zeros of $L$-functions of any elliptic curve have multiplicities bounded by a universal constant. When showing that the second phenomenon implies the first, only LI is required, and in fact the author shows that the existence of arbitrarily biased races under LI would imply GRH for an infinite family of $L(s,E)$.\par This article cites \cite{2007.Sarnak, MR3261965, MR3028758, MR2202918, 1914.Littlewood, 2000.Ng, MR1761276, MR1329368}.}
      \field{issn}{1073-7928}
      \field{journaltitle}{Int. Math. Res. Not. IMRN}
      \field{number}{18}
      \field{title}{Elliptic curves of unbounded rank and {C}hebyshev's bias}
      \field{year}{2014}
      \true{nocite}
      \field{pages}{4997\bibrangedash 5024}
      \range{pages}{28}
      \verb{doi}
      \verb 10.1093/imrn/rnt103
      \endverb
      \verb{mrnumber}
      \verb 3264673
      \endverb
      \verb{urlraw}
      \verb https://doi.org/10.1093/imrn/rnt103
      \endverb
      \verb{url}
      \verb https://doi.org/10.1093/imrn/rnt103
      \endverb
    \endentry
    \entry{MR3272280}{article}{}
      \name{author}{1}{}{%
        {{hash=889ecd9d3a804c5a7e9a5d8287766657}{%
           family={Fiorilli},
           familyi={F\bibinitperiod},
           given={D.},
           giveni={D\bibinitperiod}}}%
      }
      \strng{namehash}{889ecd9d3a804c5a7e9a5d8287766657}
      \strng{fullhash}{889ecd9d3a804c5a7e9a5d8287766657}
      \strng{bibnamehash}{889ecd9d3a804c5a7e9a5d8287766657}
      \strng{authorbibnamehash}{889ecd9d3a804c5a7e9a5d8287766657}
      \strng{authornamehash}{889ecd9d3a804c5a7e9a5d8287766657}
      \strng{authorfullhash}{889ecd9d3a804c5a7e9a5d8287766657}
      \field{extraname}{3}
      \field{sortinit}{2}
      \field{sortinithash}{ed39bb39cf854d5250e95b1c1f94f4ed}
      \field{labelnamesource}{author}
      \field{labeltitlesource}{title}
      \field{annotation}{It is known, assuming GRH and LI, that $\delta(p;N,R)$ tends to $\frac12$ as the prime modulus $p$ tends to infinity. In contrast, the author shows that the analogous density $\delta(q;N,R)$ takes a set of values that is dense in $(\frac12,1)$, under the same two assumptions (although LI can be weakened to a mild bound on the multiplicities of zeros of $L(s,\chi)$), with large biases corresponding to highly composite moduli in a quantitative sense (including a conjecture for the asymptotic size of the analogues of Skewes's number for these races). The author proves similar results for fairly general linear combinations of reduced residue classes.\par This article cites~\cites{MR3261965, MR1752093, MR1816036, MR1806291, MR3028758, MR2202918, MR1555148, MR1344249, MR2943671, MR1544293A, MR1544293B, MR2631002, MR634679, arXiv.Myerscough, 2000.Ng, MR2078705, MR1329368, MR1761276, MR1573970}.}
      \field{issn}{1937-0652}
      \field{journaltitle}{Algebra Number Theory}
      \field{number}{7}
      \field{title}{Highly biased prime number races}
      \field{volume}{8}
      \field{year}{2014}
      \true{nocite}
      \field{pages}{1733\bibrangedash 1767}
      \range{pages}{35}
      \verb{doi}
      \verb 10.2140/ant.2014.8.1733
      \endverb
      \verb{mrnumber}
      \verb 3272280
      \endverb
      \verb{urlraw}
      \verb https://doi.org/10.2140/ant.2014.8.1733
      \endverb
      \verb{url}
      \verb https://doi.org/10.2140/ant.2014.8.1733
      \endverb
    \endentry
    \entry{MR3163001}{article}{}
      \name{author}{1}{}{%
        {{hash=48017d687b31d05c591036588e04dd66}{%
           family={Humphries},
           familyi={H\bibinitperiod},
           given={Peter},
           giveni={P\bibinitperiod}}}%
      }
      \strng{namehash}{48017d687b31d05c591036588e04dd66}
      \strng{fullhash}{48017d687b31d05c591036588e04dd66}
      \strng{bibnamehash}{48017d687b31d05c591036588e04dd66}
      \strng{authorbibnamehash}{48017d687b31d05c591036588e04dd66}
      \strng{authornamehash}{48017d687b31d05c591036588e04dd66}
      \strng{authorfullhash}{48017d687b31d05c591036588e04dd66}
      \field{extraname}{2}
      \field{sortinit}{2}
      \field{sortinithash}{ed39bb39cf854d5250e95b1c1f94f4ed}
      \field{labelnamesource}{author}
      \field{labeltitlesource}{title}
      \field{annotation}{The author gives a necessary and sufficient condition for an elliptic curve over a finite field to satisfy the function field analogue of the Mertens conjecture $\limsup_{X \to \infty} |M_{E/\mathbb{F}_q}(X)|/q^{X/2}\leq 1$ (as studied by the author in~\cite{MR3189983}), in terms of the order of the finite field and the trace of the Frobenius endomorphism acting on the curve. Moreover, the author shows if the Mertens conjecture holds for a given elliptic curve, then in fact $\limsup_{X \to \infty} |M_{E/\mathbb{F}_q}(X)|/q^{X/2} = 1$.\par This article cites~\cites{MR0330069,MR3335903,MR3684399,MR0006202,MR2282922,1897.Mertens,MR783538}.}
      \field{issn}{0004-9727}
      \field{journaltitle}{Bull. Aust. Math. Soc.}
      \field{number}{1}
      \field{title}{On the {M}ertens conjecture for elliptic curves over finite fields}
      \field{volume}{89}
      \field{year}{2014}
      \true{nocite}
      \field{pages}{19\bibrangedash 32}
      \range{pages}{14}
      \verb{doi}
      \verb 10.1017/S0004972712001116
      \endverb
      \verb{mrnumber}
      \verb 3163001
      \endverb
      \verb{urlraw}
      \verb https://doi.org/10.1017/S0004972712001116
      \endverb
      \verb{url}
      \verb https://doi.org/10.1017/S0004972712001116
      \endverb
    \endentry
    \entry{MR3189983}{article}{}
      \name{author}{1}{}{%
        {{hash=48017d687b31d05c591036588e04dd66}{%
           family={Humphries},
           familyi={H\bibinitperiod},
           given={Peter},
           giveni={P\bibinitperiod}}}%
      }
      \strng{namehash}{48017d687b31d05c591036588e04dd66}
      \strng{fullhash}{48017d687b31d05c591036588e04dd66}
      \strng{bibnamehash}{48017d687b31d05c591036588e04dd66}
      \strng{authorbibnamehash}{48017d687b31d05c591036588e04dd66}
      \strng{authornamehash}{48017d687b31d05c591036588e04dd66}
      \strng{authorfullhash}{48017d687b31d05c591036588e04dd66}
      \field{extraname}{3}
      \field{sortinit}{2}
      \field{sortinithash}{ed39bb39cf854d5250e95b1c1f94f4ed}
      \field{labelnamesource}{author}
      \field{labeltitlesource}{title}
      \field{annotation}{Cha~\cite{MR3684399} introduced the Mertens function $M_{C/\mathbb{F}_q}(X)$ of a smooth projective curve defined over a finite field, and showed under LI that $E^{M_{C/\mathbb{F}_q}}(X)$ is bounded where $E^{M_{C/\mathbb{F}_q}}(X)=M_{C/\mathbb{F}_q}(X)/q^{X/2}$. Thus, a natural analogue of the Mertens conjecture for function fields would be $$\limsup_{X \to \infty} |E^{M_{C/\mathbb{F}_q}}(X)| \leq 1.$$ For fixed~$q$ and~$g$, let $\mathcal{H}_{2g+1,q^n}$ denote the set of hyperelliptic curves $y^2=f(x)$ over $\mathbb{F}_{q^n}$ arising from squarefree monic polynomials in $\mathbb{F}_{q^n}[x]$ of degree $2g+1$. The author shows that as $n \to \infty$, almost all curves in $\mathcal{H}_{2g+1,q^n}$ satisfy $\limsup_{X \to \infty} |E^{M_{C/\mathbb{F}_{q^n}}}(X)| > 1$, while for any $\beta>1$, a positive proportion of curves in $\mathcal{H}_{2g+1,q^n}$ satisfy $\limsup_{X \to \infty} |E^{M_{C/\mathbb{F}_{q^n}}}(X)| \leq \beta$.\par This article cites~\cites{MR0330069,MR3163001,MR3335903,MR3684399,MR0006202,MR2282922,1897.Mertens,MR2078705,MR783538,1913.vonSterneck}.}
      \field{issn}{1793-0421}
      \field{journaltitle}{Int. J. Number Theory}
      \field{number}{2}
      \field{title}{On the {M}ertens conjecture for function fields}
      \field{volume}{10}
      \field{year}{2014}
      \true{nocite}
      \field{pages}{341\bibrangedash 361}
      \range{pages}{21}
      \verb{doi}
      \verb 10.1142/S1793042113500978
      \endverb
      \verb{mrnumber}
      \verb 3189983
      \endverb
      \verb{urlraw}
      \verb https://doi.org/10.1142/S1793042113500978
      \endverb
      \verb{url}
      \verb https://doi.org/10.1142/S1793042113500978
      \endverb
    \endentry
    \entry{MR3179660}{article}{}
      \name{author}{1}{}{%
        {{hash=2a44d6571c930fba6c9dbf36ab288ad2}{%
           family={Radziejewski},
           familyi={R\bibinitperiod},
           given={Maciej},
           giveni={M\bibinitperiod}}}%
      }
      \strng{namehash}{2a44d6571c930fba6c9dbf36ab288ad2}
      \strng{fullhash}{2a44d6571c930fba6c9dbf36ab288ad2}
      \strng{bibnamehash}{2a44d6571c930fba6c9dbf36ab288ad2}
      \strng{authorbibnamehash}{2a44d6571c930fba6c9dbf36ab288ad2}
      \strng{authornamehash}{2a44d6571c930fba6c9dbf36ab288ad2}
      \strng{authorfullhash}{2a44d6571c930fba6c9dbf36ab288ad2}
      \field{extraname}{3}
      \field{sortinit}{2}
      \field{sortinithash}{ed39bb39cf854d5250e95b1c1f94f4ed}
      \field{labelnamesource}{author}
      \field{labeltitlesource}{title}
      \field{annotation}{The author summarizes many existing Landau-type oscillation theorems, and establishes a new oscillation theorem for functions whose Mellin transforms satisfy certain growth conditions (weaker than usual) and have certain singularities (more general than usual). As an application, he shows that the counting function of numbers that can be written as the sum of two squares has oscillations of size $x^{1/2}(\log x)^{-3/2-\ep}$; he generalizes this result to $\sum_{n\le x} (\frac14r(n))^z$, where $r(n)$ is the number of representations of~$n$ as the sum of two squares, for generic~$z$ (but not $z=1$).\par This article cites~\cites{1891.Phragmen,MR0146156,MR0146159,MR0156827,MR858395,MR891057,MR1511360,MR2110393,MR2123959,MR2470105,MR2730494}.}
      \field{issn}{0033-5606}
      \field{journaltitle}{Q. J. Math.}
      \field{number}{1}
      \field{title}{Oscillatory properties of real functions with weakly bounded Mellin transform}
      \field{volume}{65}
      \field{year}{2014}
      \true{nocite}
      \field{pages}{249\bibrangedash 266}
      \range{pages}{18}
      \verb{doi}
      \verb 10.1093/qmath/has036
      \endverb
      \verb{mrnumber}
      \verb 3179660
      \endverb
    \endentry
    \entry{MR3120597}{article}{}
      \name{author}{2}{}{%
        {{hash=99298fc0cd61b25871f9b3ad564fa78b}{%
           family={Saouter},
           familyi={S\bibinitperiod},
           given={Yannick},
           giveni={Y\bibinitperiod}}}%
        {{hash=c83d046dcdb536f5163e46eb8fe07827}{%
           family={Riele},
           familyi={R\bibinitperiod},
           given={Herman},
           giveni={H\bibinitperiod},
           prefix={te},
           prefixi={t\bibinitperiod}}}%
      }
      \strng{namehash}{c854407c034da79afefc129f9cbec9ec}
      \strng{fullhash}{c854407c034da79afefc129f9cbec9ec}
      \strng{bibnamehash}{c854407c034da79afefc129f9cbec9ec}
      \strng{authorbibnamehash}{c854407c034da79afefc129f9cbec9ec}
      \strng{authornamehash}{c854407c034da79afefc129f9cbec9ec}
      \strng{authorfullhash}{c854407c034da79afefc129f9cbec9ec}
      \field{sortinit}{2}
      \field{sortinithash}{ed39bb39cf854d5250e95b1c1f94f4ed}
      \field{labelnamesource}{author}
      \field{labeltitlesource}{title}
      \field{annotation}{The authors refine Pintz's effective disproof~\cite{MR891440} of the Mertens conjecture and show that $|M(x)|>1.0088 \sqrt{x}$ for some $x<\exp(1.004 \times 10^{33})$. The authors also discuss possibilities for obtaining smaller counterexamples. \par This article cites~\cite{1897.Mertens,MR2282922,MR783538,MR891440,MR2684372,MR2340242}.}
      \field{issn}{0025-5718}
      \field{journaltitle}{Math. Comp.}
      \field{number}{285}
      \field{title}{Improved results on the {M}ertens conjecture}
      \field{volume}{83}
      \field{year}{2014}
      \true{nocite}
      \field{pages}{421\bibrangedash 433}
      \range{pages}{13}
      \verb{doi}
      \verb 10.1090/S0025-5718-2013-02716-0
      \endverb
      \verb{mrnumber}
      \verb 3120597
      \endverb
      \verb{urlraw}
      \verb https://doi.org/10.1090/S0025-5718-2013-02716-0
      \endverb
      \verb{url}
      \verb https://doi.org/10.1090/S0025-5718-2013-02716-0
      \endverb
    \endentry
    \entry{MR3335903}{article}{}
      \name{author}{2}{}{%
        {{hash=251a94d49687997eac18f56ef1823028}{%
           family={Best},
           familyi={B\bibinitperiod},
           given={D.\bibnamedelimi G.},
           giveni={D\bibinitperiod\bibinitdelim G\bibinitperiod}}}%
        {{hash=f8faf70b3f8ac0edff4b9b5337e687bf}{%
           family={Trudgian},
           familyi={T\bibinitperiod},
           given={T.\bibnamedelimi S.},
           giveni={T\bibinitperiod\bibinitdelim S\bibinitperiod}}}%
      }
      \strng{namehash}{4ff3faf470a39e73d27ec82924abaeb8}
      \strng{fullhash}{4ff3faf470a39e73d27ec82924abaeb8}
      \strng{bibnamehash}{4ff3faf470a39e73d27ec82924abaeb8}
      \strng{authorbibnamehash}{4ff3faf470a39e73d27ec82924abaeb8}
      \strng{authornamehash}{4ff3faf470a39e73d27ec82924abaeb8}
      \strng{authorfullhash}{4ff3faf470a39e73d27ec82924abaeb8}
      \field{sortinit}{2}
      \field{sortinithash}{ed39bb39cf854d5250e95b1c1f94f4ed}
      \field{labelnamesource}{author}
      \field{labeltitlesource}{title}
      \field{annotation}{Using the method of Grosswald~\cite{MR0332685}, the authors prove that \[ \liminf_{x \to \infty} E^M(x) < -1.6383 \quad \text{and} \quad \limsup_{x \to \infty} E^M(x) > 1.6383. \] The proof uses a refined variant of the LLL-algorithm and a kernel function used by Odlyzko and te Riele~\cite{MR783538}, as well as several techniques for speeding up the computation.\par In the appendix of the article, the authors list new lower bounds on $m$ for the ``$m$-dependence'' of the first $n$ positive ordinates of nontrivial zeros of $\zeta(s)$, for $n \leq 500$. This numerical evidence, combined with Grosswald's theorem, provides new proofs that $E^\pi(x)$, $L(x)$, and $L_r(x)$ change signs infinitely often.\par This article cites~\cite{MR429789,MR2282922,MR783538,MR0006202,MR0104638,MR0332685,MR0332684,MR0330069,MR654520}.}
      \field{issn}{0025-5718}
      \field{journaltitle}{Math. Comp.}
      \field{number}{294}
      \field{title}{Linear relations of zeroes of the zeta-function}
      \field{volume}{84}
      \field{year}{2015}
      \true{nocite}
      \field{pages}{2047\bibrangedash 2058}
      \range{pages}{12}
      \verb{doi}
      \verb 10.1090/S0025-5718-2014-02916-5
      \endverb
      \verb{mrnumber}
      \verb 3335903
      \endverb
      \verb{urlraw}
      \verb https://doi.org/10.1090/S0025-5718-2014-02916-5
      \endverb
      \verb{url}
      \verb https://doi.org/10.1090/S0025-5718-2014-02916-5
      \endverb
    \endentry
    \entry{MR3414306}{article}{}
      \name{author}{1}{}{%
        {{hash=b58171fe4e01fa0b55ff6f0dcd610186}{%
           family={Büthe},
           familyi={B\bibinitperiod},
           given={J.},
           giveni={J\bibinitperiod}}}%
      }
      \strng{namehash}{b58171fe4e01fa0b55ff6f0dcd610186}
      \strng{fullhash}{b58171fe4e01fa0b55ff6f0dcd610186}
      \strng{bibnamehash}{b58171fe4e01fa0b55ff6f0dcd610186}
      \strng{authorbibnamehash}{b58171fe4e01fa0b55ff6f0dcd610186}
      \strng{authornamehash}{b58171fe4e01fa0b55ff6f0dcd610186}
      \strng{authorfullhash}{b58171fe4e01fa0b55ff6f0dcd610186}
      \field{sortinit}{2}
      \field{sortinithash}{ed39bb39cf854d5250e95b1c1f94f4ed}
      \field{labelnamesource}{author}
      \field{labeltitlesource}{title}
      \field{annotation}{Rosser and Schoenfeld \cite{MR0137689} observed that $\Delta^{\pi_r}(x) \ge 0$ for $1\le x\le 10^8$. The author proves that there exists an $x_0 \approx 1.91\times 10^{215}$ such that $\Delta^{\pi_r}(x) < 0$ for $x \in [x_0-6\times 10^{103},x_0]$. The proof is an adaptation of a method of Lehman~\cite{MR0202686} that improved upon Skewes's number, though with a different choice of kernel function.\par This article cites~\cites{MR0202686, MR729173, MR3356033, MR1573970, MR0137689}.}
      \field{issn}{0065-1036}
      \field{journaltitle}{Acta Arith.}
      \field{number}{2}
      \field{title}{On the first sign change in {M}ertens' theorem}
      \field{volume}{171}
      \field{year}{2015}
      \true{nocite}
      \field{pages}{183\bibrangedash 195}
      \range{pages}{13}
      \verb{doi}
      \verb 10.4064/aa171-2-5
      \endverb
      \verb{mrnumber}
      \verb 3414306
      \endverb
      \verb{urlraw}
      \verb https://doi.org/10.4064/aa171-2-5
      \endverb
      \verb{url}
      \verb https://doi.org/10.4064/aa171-2-5
      \endverb
    \endentry
    \entry{MR3356760}{article}{}
      \name{author}{1}{}{%
        {{hash=889ecd9d3a804c5a7e9a5d8287766657}{%
           family={Fiorilli},
           familyi={F\bibinitperiod},
           given={D.},
           giveni={D\bibinitperiod}}}%
      }
      \strng{namehash}{889ecd9d3a804c5a7e9a5d8287766657}
      \strng{fullhash}{889ecd9d3a804c5a7e9a5d8287766657}
      \strng{bibnamehash}{889ecd9d3a804c5a7e9a5d8287766657}
      \strng{authorbibnamehash}{889ecd9d3a804c5a7e9a5d8287766657}
      \strng{authornamehash}{889ecd9d3a804c5a7e9a5d8287766657}
      \strng{authorfullhash}{889ecd9d3a804c5a7e9a5d8287766657}
      \field{extraname}{4}
      \field{sortinit}{2}
      \field{sortinithash}{ed39bb39cf854d5250e95b1c1f94f4ed}
      \field{labelnamesource}{author}
      \field{labeltitlesource}{title}
      \field{annotation}{The author studies the ``variance'' of primes in arithmetic progressions \[ V(x;q) = \sum_{(a,q)=1} \bigg| \psi(x;q,a) - \frac{\psi(x;\chi_0)}{\phi(q)} \bigg|^2 \] that was first examined by Hooley and is known to be asymptotic to $x\log q$ in certain ranges of~$q$. The author conjectures that $V(x;q) \sim x\log q$ uniformly for $(\log\log x)^{1+\delta} \le q \le x$ for any fixed $\delta>0$, a much wider range than previously anticipated. The conjecture arises from the computation, assuming GRH and LI, of the limiting logarithmic distribution of $V(x;q)$ for a given~$q$ and studying its large deviations.\par This article cites~\cites{2000.Ng,MR0004255,MR0506080,MR634679,MR1329368,MR1806291,MR1956261,MR2078705,MR2631002,MR2943671,MR2998146,MR3028758,MR3063909,MR3264673,MR3272280}.}
      \field{issn}{1073-7928}
      \field{journaltitle}{Int. Math. Res. Not. IMRN}
      \field{number}{12}
      \field{title}{The distribution of the variance of primes in arithmetic progressions}
      \field{year}{2015}
      \true{nocite}
      \field{pages}{4421\bibrangedash 4448}
      \range{pages}{28}
      \verb{doi}
      \verb 10.1093/imrn/rnu074
      \endverb
      \verb{mrnumber}
      \verb 3356760
      \endverb
    \endentry
    \entry{MR3414302}{article}{}
      \name{author}{2}{}{%
        {{hash=6b393826a3dd69e5d16ac8319b2612dd}{%
           family={Kisilevsky},
           familyi={K\bibinitperiod},
           given={H.},
           giveni={H\bibinitperiod}}}%
        {{hash=8aabc8fd166d5aa179c76a42ffba2cc5}{%
           family={Rubinstein},
           familyi={R\bibinitperiod},
           given={M.\bibnamedelimi O.},
           giveni={M\bibinitperiod\bibinitdelim O\bibinitperiod}}}%
      }
      \strng{namehash}{f6f50421343a63c3dbf092af65a8d0d5}
      \strng{fullhash}{f6f50421343a63c3dbf092af65a8d0d5}
      \strng{bibnamehash}{f6f50421343a63c3dbf092af65a8d0d5}
      \strng{authorbibnamehash}{f6f50421343a63c3dbf092af65a8d0d5}
      \strng{authornamehash}{f6f50421343a63c3dbf092af65a8d0d5}
      \strng{authorfullhash}{f6f50421343a63c3dbf092af65a8d0d5}
      \field{sortinit}{2}
      \field{sortinithash}{ed39bb39cf854d5250e95b1c1f94f4ed}
      \field{labelnamesource}{author}
      \field{labeltitlesource}{title}
      \field{annotation}{Let $P_{\rm odd}= \{ 2, 5, 11, 17, 23, 31, \ldots \}$ be the set consisting of every other prime (the odd-index primes). In this article, the authors show that $P_{\rm odd}$ cannot be written as a finite union of sets of the form $\{p\colon p\equiv a\mod q\}$, even up to finitely many exceptions. More generally, call a set of $P$ prime ideals of a number field $K$ a Chebotarev set if there are finitely many finite Galois extensions $L_i/K$ and conjugacy classes $C_i$ such that the symmetric difference of the sets $P$ and $\bigcup_{i=1} \pi_{C_i}(L_i/K)$ is finite. Using explicit formulas, the authors show that if~$P$ is a Chebotarev set of density $\beta\in\mathbb{Q}$ with $0<\beta<1$, then $\Delta_{P}(x)=P(x)-\beta\pi_K(x)=\Omega(x^{1/2}/\log x)$. In particular, $P_{\rm odd}$ is not a Chebatorev set (hence not a finite union of sets of primes in residue classes) since the counting function of $P_{\rm odd}$ is $\frac12\pi(x) + O(1)$.\par This article cites~\cites{1932.Ingham,MR1329368}.}
      \field{issn}{0065-1036}
      \field{journaltitle}{Acta Arith.}
      \field{number}{2}
      \field{title}{Chebotarev sets}
      \field{volume}{171}
      \field{year}{2015}
      \true{nocite}
      \field{pages}{97\bibrangedash 124}
      \range{pages}{28}
      \verb{doi}
      \verb 10.4064/aa171-2-1
      \endverb
      \verb{mrnumber}
      \verb 3414302
      \endverb
      \verb{urlraw}
      \verb https://doi.org/10.4064/aa171-2-1
      \endverb
      \verb{url}
      \verb https://doi.org/10.4064/aa171-2-1
      \endverb
    \endentry
    \entry{arXiv.Lay}{online}{}
      \name{author}{1}{}{%
        {{hash=00f8cf5ecf192c1d41eced997c115cb8}{%
           family={Lay},
           familyi={L\bibinitperiod},
           given={J.},
           giveni={J\bibinitperiod}}}%
      }
      \strng{namehash}{00f8cf5ecf192c1d41eced997c115cb8}
      \strng{fullhash}{00f8cf5ecf192c1d41eced997c115cb8}
      \strng{bibnamehash}{00f8cf5ecf192c1d41eced997c115cb8}
      \strng{authorbibnamehash}{00f8cf5ecf192c1d41eced997c115cb8}
      \strng{authornamehash}{00f8cf5ecf192c1d41eced997c115cb8}
      \strng{authorfullhash}{00f8cf5ecf192c1d41eced997c115cb8}
      \field{sortinit}{2}
      \field{sortinithash}{ed39bb39cf854d5250e95b1c1f94f4ed}
      \field{labelnamesource}{author}
      \field{labeltitlesource}{title}
      \field{annotation}{The author shows that the functions $E^\theta_r(x)$ and $E^\pi_r(x)$ change sign infinitely often. Under RH, following a similar proof used by Diamond and Pintz~\cite{MR2605532}, the author shows $E^\theta_r(x) = \Omega_{\pm} \big( \log \log \log x \big)$ and the same for $E^\pi_r(x)$. Similar to the work of Lamzouri~\cite{MR3455269}, assuming RH and LI, the author proves that both the logarithmic densities of $\{x>1\colon E^\theta_r(x)>0\}$ and $\{x>1\colon E^\pi_r(x)>0\}$ equal $\delta(\li,\pi) \approx 0.99999974$.\par This article cites~\cites{MR3261965,MR2605532,1932.Ingham,MR3455269,MR0137689,MR1329368,MR0004255,MR3455269}.}
      \field{title}{Sign changes in Mertens' first and second theorems}
      \field{year}{2015}
      \true{nocite}
      \verb{urlraw}
      \verb https://arxiv.org/abs/1505.03589
      \endverb
      \verb{url}
      \verb https://arxiv.org/abs/1505.03589
      \endverb
    \endentry
    \entry{MR3356033}{article}{}
      \name{author}{3}{}{%
        {{hash=99298fc0cd61b25871f9b3ad564fa78b}{%
           family={Saouter},
           familyi={S\bibinitperiod},
           given={Yannick},
           giveni={Y\bibinitperiod}}}%
        {{hash=427a3ba3d2c76d68bc8c1bfc231edce0}{%
           family={Trudgian},
           familyi={T\bibinitperiod},
           given={Timothy},
           giveni={T\bibinitperiod}}}%
        {{hash=447f22d7de07f4cbd918f6b2abb97288}{%
           family={Demichel},
           familyi={D\bibinitperiod},
           given={Patrick},
           giveni={P\bibinitperiod}}}%
      }
      \strng{namehash}{fcb20b42ca212ff9c12bc49fe9926ddf}
      \strng{fullhash}{fcb20b42ca212ff9c12bc49fe9926ddf}
      \strng{bibnamehash}{fcb20b42ca212ff9c12bc49fe9926ddf}
      \strng{authorbibnamehash}{fcb20b42ca212ff9c12bc49fe9926ddf}
      \strng{authornamehash}{fcb20b42ca212ff9c12bc49fe9926ddf}
      \strng{authorfullhash}{fcb20b42ca212ff9c12bc49fe9926ddf}
      \field{sortinit}{2}
      \field{sortinithash}{ed39bb39cf854d5250e95b1c1f94f4ed}
      \field{labelnamesource}{author}
      \field{labeltitlesource}{title}
      \field{annotation}{The authors incorporate both theoretical and computational improvements to show that there are more than $7.17\times10^{152}$ consecutive integers in the interval $[1.397165243588\times10^{316}, 1.397167149324\times10^{316}]$ for which $\pi(x)>\li(x)$.\par This article cites~\cites{MR2420864,MR0202686,1914.Littlewood,MR2684372}.}
      \field{issn}{0025-5718}
      \field{journaltitle}{Math. Comp.}
      \field{number}{295}
      \field{title}{A still sharper region where {$\pi (x)-\li (x)$} is positive}
      \field{volume}{84}
      \field{year}{2015}
      \true{nocite}
      \field{pages}{2433\bibrangedash 2446}
      \range{pages}{14}
      \verb{doi}
      \verb 10.1090/S0025-5718-2015-02930-5
      \endverb
      \verb{mrnumber}
      \verb 3356033
      \endverb
    \endentry
    \entry{MR3509711}{article}{}
      \name{author}{3}{}{%
        {{hash=922e48f9d159c924a55b336dfeb65261}{%
           family={Bhowmik},
           familyi={B\bibinitperiod},
           given={Gautami},
           giveni={G\bibinitperiod}}}%
        {{hash=d6ae4b9844776294c92ff23952e91b53}{%
           family={Ramaré},
           familyi={R\bibinitperiod},
           given={Olivier},
           giveni={O\bibinitperiod}}}%
        {{hash=dd239e78c668d9270b067bfae5530a1f}{%
           family={Schlage--Puchta},
           familyi={S\bibinithyphendelim \bibinithyphendelim P\bibinitperiod},
           given={Jan-Christoph},
           giveni={J\bibinithyphendelim C\bibinitperiod}}}%
      }
      \list{language}{2}{%
        {English, with English}%
        {French summaries}%
      }
      \strng{namehash}{83e0fb0de8538ce919bef0bbac9f375f}
      \strng{fullhash}{83e0fb0de8538ce919bef0bbac9f375f}
      \strng{bibnamehash}{83e0fb0de8538ce919bef0bbac9f375f}
      \strng{authorbibnamehash}{83e0fb0de8538ce919bef0bbac9f375f}
      \strng{authornamehash}{83e0fb0de8538ce919bef0bbac9f375f}
      \strng{authorfullhash}{83e0fb0de8538ce919bef0bbac9f375f}
      \field{sortinit}{2}
      \field{sortinithash}{ed39bb39cf854d5250e95b1c1f94f4ed}
      \field{labelnamesource}{author}
      \field{labeltitlesource}{title}
      \field{annotation}{The authors establish an effective version of Landau's theorem for Dirichlet integrals, which they apply to the function $G_k(n) = \sum_{n_1 + \cdots n_k = n} \Lambda(n_1) \cdots \Lambda(n_k)$ counting the weighted number of representations of $n$ as the sum of $k$ primes. Further defining \[ \Delta_k(x) = \sum_{n\le x} G_k(n) - \frac{x^k}{k!} + k \sum_\rho \frac{x^{\rho+k-1}}{\rho(\rho+1)\cdots(\rho+k-1)}, \] they show under RH that $\Delta_2(n) = \Omega_\pm(x)$ and, for particular constants $c_k$, that $\Delta_k(n) - c_kx^{k-1} = \Omega_\pm(x^{k-1})$ for $k\ge3$.\par This article cites~\cites{MR0146156,MR2040112}.}
      \field{issn}{1246-7405}
      \field{journaltitle}{J. Théor. Nombres Bordeaux}
      \field{number}{2}
      \field{title}{Tauberian oscillation theorems and the distribution of {G}oldbach numbers}
      \field{volume}{28}
      \field{year}{2016}
      \true{nocite}
      \field{pages}{291\bibrangedash 299}
      \range{pages}{9}
      \verb{mrnumber}
      \verb 3509711
      \endverb
    \endentry
    \entry{MR3581815}{article}{}
      \name{author}{3}{}{%
        {{hash=3417f8b5237d274ddd67adbd43553547}{%
           family={Cha},
           familyi={C\bibinitperiod},
           given={B.},
           giveni={B\bibinitperiod}}}%
        {{hash=889ecd9d3a804c5a7e9a5d8287766657}{%
           family={Fiorilli},
           familyi={F\bibinitperiod},
           given={D.},
           giveni={D\bibinitperiod}}}%
        {{hash=dfca5f0fbbbf59635f02988f421953ae}{%
           family={Jouve},
           familyi={J\bibinitperiod},
           given={F.},
           giveni={F\bibinitperiod}}}%
      }
      \strng{namehash}{5cfc0c94602fe7b3e73bd312fbce4b01}
      \strng{fullhash}{5cfc0c94602fe7b3e73bd312fbce4b01}
      \strng{bibnamehash}{5cfc0c94602fe7b3e73bd312fbce4b01}
      \strng{authorbibnamehash}{5cfc0c94602fe7b3e73bd312fbce4b01}
      \strng{authornamehash}{5cfc0c94602fe7b3e73bd312fbce4b01}
      \strng{authorfullhash}{5cfc0c94602fe7b3e73bd312fbce4b01}
      \field{sortinit}{2}
      \field{sortinithash}{ed39bb39cf854d5250e95b1c1f94f4ed}
      \field{labelnamesource}{author}
      \field{labeltitlesource}{title}
      \field{annotation}{In this article, the authors describe the prime number race for elliptic curves over the function field of a proper, smooth, and geometrically connected curve over a finite field. Let ${E/K}$ be a family of elliptic curves of unbounded conductor for which $L(E/K,T)$ satisfies LI (which can be confirmed in some cases) and such that $\text{rank}(E/K)=o(\sqrt{N_{E/K}})$ as $N_{E/K}\to\infty$. Then, the random variable ${\frac{\sqrt q-1}{\sqrt q}}X_E/\sqrt{N_{E/K}}$ converges in distribution to the standard Gaussian as $N_{E/K}\to\infty$. As a consequence, $\delta(E)$ (the proportion of those~$x$ for which more primes up to~$x$ have positive trace of Frobenius $a_p(E)$ than negative) tends to $\frac 12$.\par Moreover, the authors also study the behaviour of the function $T_E(X)= -{X}{q^{-X/2}}\sum_{\deg(v)\leq X} 2\cos{\theta_v}$ associated to the elliptic curves of Ulmer's family of elliptic curves over $\Bbb F_q[t]$. They discuss the cases of extreme bias and moderate bias for Ulmer's family. Moreover, through proving a central limit theorem, the authors shows that $\delta(E_f)-\frac 12 = \Omega(\frac1{\sqrt{d}})$ for quadratic twists~$E_f$.\par This article cites~\cites{MR3261965,MR2474313,MR3264673,MR3272280,MR3028758,MR1329368,2007.Sarnak}.}
      \field{issn}{0012-9593}
      \field{journaltitle}{Ann. Sci. Éc. Norm. Supér. (4)}
      \field{number}{5}
      \field{title}{Prime number races for elliptic curves over function fields}
      \field{volume}{49}
      \field{year}{2016}
      \true{nocite}
      \field{pages}{1239\bibrangedash 1277}
      \range{pages}{39}
      \verb{doi}
      \verb 10.24033/asens.2308
      \endverb
      \verb{mrnumber}
      \verb 3581815
      \endverb
      \verb{urlraw}
      \verb https://doi.org/10.24033/asens.2308
      \endverb
      \verb{url}
      \verb https://doi.org/10.24033/asens.2308
      \endverb
    \endentry
    \entry{MR3521338}{article}{}
      \name{author}{3}{}{%
        {{hash=d7aee8a672f37558eb9119074764e7f5}{%
           family={Dummit},
           familyi={D\bibinitperiod},
           given={D.},
           giveni={D\bibinitperiod}}}%
        {{hash=7c95bde388d46e453494c7b2fbd1b088}{%
           family={Granville},
           familyi={G\bibinitperiod},
           given={A.},
           giveni={A\bibinitperiod}}}%
        {{hash=a699d3ec8d37bdcc6b5f75e108b4b9d3}{%
           family={Kisilevsky},
           familyi={K\bibinitperiod},
           given={B.},
           giveni={B\bibinitperiod}}}%
      }
      \strng{namehash}{e1122a9ecacc1c1ca72d1f02ad518714}
      \strng{fullhash}{e1122a9ecacc1c1ca72d1f02ad518714}
      \strng{bibnamehash}{e1122a9ecacc1c1ca72d1f02ad518714}
      \strng{authorbibnamehash}{e1122a9ecacc1c1ca72d1f02ad518714}
      \strng{authornamehash}{e1122a9ecacc1c1ca72d1f02ad518714}
      \strng{authorfullhash}{e1122a9ecacc1c1ca72d1f02ad518714}
      \field{sortinit}{2}
      \field{sortinithash}{ed39bb39cf854d5250e95b1c1f94f4ed}
      \field{labelnamesource}{author}
      \field{labeltitlesource}{title}
      \field{annotation}{This article establishes permanent, relatively large biases in races involving products of two primes. The authors show that if $\chi$ is a quadratic character\mod d, then \[ \frac{\#\{ pq\le x\colon \chi(p)=\chi(q)=-1 \}}{\#\{ pq\le x\colon (pq,d)=1 \}} = 1 - \bigg( \sum_p \frac{\chi(p)}p + o(1) \bigg) \frac1{\log\log x}, \] and the same statement with both minus signs changed to plus signs. In particular, the race between integers $pq$ with $\chi(p)=\chi(q)=-1$ and those with $\chi(p)=\chi(q)=1$ has a bias in favor of the sign of (and proportional to) $\sum_p {\chi(p)}/p$.\par For example, $\sum_p {\chi_5(p)}/p \approx -1.008$, and correspondingly integers $pq$ with both $p$ and $q$ quadratic nonresidues\mod 5 are 41.6\% more numerous up to $10^7$ than random chance would suggest. The authors conjecture that there exists $d\le x$ such that the right-hand side is as large as $1+\log\log\log x/{\log\log x}$. They also note that the same proof gives a bias for the ratio \[ \sum_{\substack{p\le x\\\chi(p)=1}}\frac1p\; \bigg/\sum_{\substack{p\le x\\\chi(p)=-1}}\frac1p = 1 + \bigg( 2\sum_p \frac{\chi(p)}p + o(1) \bigg) \frac1{\log\log x}, \] giving a permanent bias involving only primes (in contrast to the unweighted race between $\pi(x;d,R)$ and $\pi(x;d,N)$); the proof also generalizes to products of $k$ primes with prescribed quadratic character values (for fixed $k$). They remark on the possibility of counting $pq\le x$ where $p$ and $q$ are restricted to prescribed but arbitrary residue classes.\par This article cites~\cites{MR2778652}.}
      \field{issn}{0025-5793}
      \field{journaltitle}{Mathematika}
      \field{number}{2}
      \field{title}{Big biases amongst products of two primes}
      \field{volume}{62}
      \field{year}{2016}
      \true{nocite}
      \field{pages}{502\bibrangedash 507}
      \range{pages}{6}
      \verb{doi}
      \verb 10.1112/S0025579315000339
      \endverb
      \verb{mrnumber}
      \verb 3521338
      \endverb
      \verb{urlraw}
      \verb https://doi.org/10.1112/S0025579315000339
      \endverb
      \verb{url}
      \verb https://doi.org/10.1112/S0025579315000339
      \endverb
    \endentry
    \entry{MR3455269}{article}{}
      \name{author}{1}{}{%
        {{hash=a01ad5bbada0f5c69f81f1149c07aefc}{%
           family={Lamzouri},
           familyi={L\bibinitperiod},
           given={Y.},
           giveni={Y\bibinitperiod}}}%
      }
      \strng{namehash}{a01ad5bbada0f5c69f81f1149c07aefc}
      \strng{fullhash}{a01ad5bbada0f5c69f81f1149c07aefc}
      \strng{bibnamehash}{a01ad5bbada0f5c69f81f1149c07aefc}
      \strng{authorbibnamehash}{a01ad5bbada0f5c69f81f1149c07aefc}
      \strng{authornamehash}{a01ad5bbada0f5c69f81f1149c07aefc}
      \strng{authorfullhash}{a01ad5bbada0f5c69f81f1149c07aefc}
      \field{extraname}{4}
      \field{sortinit}{2}
      \field{sortinithash}{ed39bb39cf854d5250e95b1c1f94f4ed}
      \field{labelnamesource}{author}
      \field{labeltitlesource}{title}
      \field{annotation}{Let $\mathcal{M}$ be the set of real numbers $x>1$ such that $\prod_{p \leq x } ( 1- 1/{p})^{-1} > e^{C_0} \log x$. Assuming RH, the author proves that both $\mathcal{M}$ and its complement have positive lower logarithmic density. Further assuming LI, the author shows that the logarithmic density of $\mathcal{M}$ equals $\delta(\li,\pi) \approx 0.99999974$. The author also conjectures that \[ \limsup_{x\to\infty} \bigg( \bigg( \prod_{p\le x} \bigg( 1-\frac1p \bigg)^{-1} - e^{C_0}\log x \bigg) \bigg/ \frac{(\log\log\log x)^2}{\sqrt x} \bigg) = \frac{e^{C_0}}{2\pi} \] and the symmetric result for the $\liminf$. \par This article cites~\cites{MR2605532,MR3151605,MR2202918,1914.Littlewood,MR2631002,MR634679,MR0137689,MR1329368}, as well as an earlier draft of this annotated bibliography.}
      \field{issn}{1793-0421}
      \field{journaltitle}{Int. J. Number Theory}
      \field{number}{1}
      \field{title}{A bias in {M}ertens' product formula}
      \field{volume}{12}
      \field{year}{2016}
      \true{nocite}
      \field{pages}{97\bibrangedash 109}
      \range{pages}{13}
      \verb{doi}
      \verb 10.1142/S1793042116500068
      \endverb
      \verb{mrnumber}
      \verb 3455269
      \endverb
      \verb{urlraw}
      \verb https://doi.org/10.1142/S1793042116500068
      \endverb
      \verb{url}
      \verb https://doi.org/10.1142/S1793042116500068
      \endverb
    \endentry
    \entry{MR3624386}{article}{}
      \name{author}{2}{}{%
        {{hash=5cb6f78b503191fd0bfbfcba676ac134}{%
           family={Lemke\bibnamedelima Oliver},
           familyi={L\bibinitperiod\bibinitdelim O\bibinitperiod},
           given={R.\bibnamedelimi J.},
           giveni={R\bibinitperiod\bibinitdelim J\bibinitperiod}}}%
        {{hash=549bdccc730fda930fd587b6ae3caaba}{%
           family={Soundararajan},
           familyi={S\bibinitperiod},
           given={K.},
           giveni={K\bibinitperiod}}}%
      }
      \strng{namehash}{ae4a066eca6b7243e3cef24fd0e898c7}
      \strng{fullhash}{ae4a066eca6b7243e3cef24fd0e898c7}
      \strng{bibnamehash}{ae4a066eca6b7243e3cef24fd0e898c7}
      \strng{authorbibnamehash}{ae4a066eca6b7243e3cef24fd0e898c7}
      \strng{authornamehash}{ae4a066eca6b7243e3cef24fd0e898c7}
      \strng{authorfullhash}{ae4a066eca6b7243e3cef24fd0e898c7}
      \field{extraname}{1}
      \field{sortinit}{2}
      \field{sortinithash}{ed39bb39cf854d5250e95b1c1f94f4ed}
      \field{labelnamesource}{author}
      \field{labeltitlesource}{title}
      \field{annotation}{Given a tuple $\textbf{a} = (a_1,\ldots,a_r)$ of reduced residues modulo~$q$, let \[ \pi(x;q,\textbf{a}) = \# \{ p_n \leq x \colon p_{n+i-1} \equiv a_i \mod{q} \text{ for each } 1 \leq i \leq r\}\] count the occurrences of the pattern of residues defined by~$\textbf{a}$. The authors observe (among other things) that repeated residues appear less frequently than changing ones; for example, for $x_0 = p_{10^8}$, we have $\pi(x_0;10,(1,3)) = 7,429,438$ but $\pi(x_0;10,(1,1)) = 4,623,042$. The authors conjecture that all patterns occur equally often in the limit, but that lower order terms create predictable biases; in contrast to prime number races and their infinity of sign changes, some of these inequalities should always hold, such as $\pi(x;3,(1,-1)) > \pi(x;3,(1,1))$ for $x \geq 5$. The authors provide numerical evidence for their conjectures, as well as heuristic justification related to the Hardy--Littlewood prime $k$-tuples conjecture.\par This article cites~\cites{MR1329368,MR0466043,MR2202918,MR3521338}.}
      \field{issn}{1091-6490}
      \field{journaltitle}{Proc. Natl. Acad. Sci. USA}
      \field{number}{31}
      \field{title}{Unexpected biases in the distribution of consecutive primes}
      \field{volume}{113}
      \field{year}{2016}
      \true{nocite}
      \field{pages}{E4446\bibrangedash E4454}
      \range{pages}{-1}
      \verb{doi}
      \verb 10.1073/pnas.1605366113
      \endverb
      \verb{mrnumber}
      \verb 3624386
      \endverb
      \verb{urlraw}
      \verb https://doi.org/10.1073/pnas.1605366113
      \endverb
      \verb{url}
      \verb https://doi.org/10.1073/pnas.1605366113
      \endverb
    \endentry
    \entry{MR3454375}{article}{}
      \name{author}{2}{}{%
        {{hash=b7137f17a7ecf01647eb7f9483185c94}{%
           family={Platt},
           familyi={P\bibinitperiod},
           given={D.\bibnamedelimi J.},
           giveni={D\bibinitperiod\bibinitdelim J\bibinitperiod}}}%
        {{hash=f8faf70b3f8ac0edff4b9b5337e687bf}{%
           family={Trudgian},
           familyi={T\bibinitperiod},
           given={T.\bibnamedelimi S.},
           giveni={T\bibinitperiod\bibinitdelim S\bibinitperiod}}}%
      }
      \strng{namehash}{c14c29110c40a8b34a8d4a8b0fedcbce}
      \strng{fullhash}{c14c29110c40a8b34a8d4a8b0fedcbce}
      \strng{bibnamehash}{c14c29110c40a8b34a8d4a8b0fedcbce}
      \strng{authorbibnamehash}{c14c29110c40a8b34a8d4a8b0fedcbce}
      \strng{authornamehash}{c14c29110c40a8b34a8d4a8b0fedcbce}
      \strng{authorfullhash}{c14c29110c40a8b34a8d4a8b0fedcbce}
      \field{extraname}{1}
      \field{sortinit}{2}
      \field{sortinithash}{ed39bb39cf854d5250e95b1c1f94f4ed}
      \field{labelnamesource}{author}
      \field{labeltitlesource}{title}
      \field{annotation}{The authors compute that $\Delta^\theta(x) < 0$ for $0 \le x \le 1.39 \times 10^{17}$. By partial summation, this implies that $\Delta^\pi(x) < 0$ for $2 < x \le 1.39 \times 10^{17}$. The authors also prove that there exists $x \approx 1.3971623 \times 10^{316}$ for which $\Delta^\theta(x) > 0$.\par This article cites~\cites{MR1752093, 1932.Ingham, MR2420864, MR0202686, 1914.Littlewood, MR2271616, MR3356033, MR0067145}.}
      \field{issn}{0025-5718}
      \field{journaltitle}{Math. Comp.}
      \field{number}{299}
      \field{title}{On the first sign change of {$\theta(x)-x$}}
      \field{volume}{85}
      \field{year}{2016}
      \true{nocite}
      \field{pages}{1539\bibrangedash 1547}
      \range{pages}{9}
      \verb{doi}
      \verb 10.1090/mcom/3021
      \endverb
      \verb{mrnumber}
      \verb 3454375
      \endverb
      \verb{urlraw}
      \verb https://doi.org/10.1090/mcom/3021
      \endverb
      \verb{url}
      \verb https://doi.org/10.1090/mcom/3021
      \endverb
    \endentry
    \entry{MR3684399}{article}{}
      \name{author}{1}{}{%
        {{hash=7871a490afe3f3963421804f312b0ad7}{%
           family={Cha},
           familyi={C\bibinitperiod},
           given={Byungchul},
           giveni={B\bibinitperiod}}}%
      }
      \strng{namehash}{7871a490afe3f3963421804f312b0ad7}
      \strng{fullhash}{7871a490afe3f3963421804f312b0ad7}
      \strng{bibnamehash}{7871a490afe3f3963421804f312b0ad7}
      \strng{authorbibnamehash}{7871a490afe3f3963421804f312b0ad7}
      \strng{authornamehash}{7871a490afe3f3963421804f312b0ad7}
      \strng{authorfullhash}{7871a490afe3f3963421804f312b0ad7}
      \field{extraname}{2}
      \field{sortinit}{2}
      \field{sortinithash}{ed39bb39cf854d5250e95b1c1f94f4ed}
      \field{labelnamesource}{author}
      \field{labeltitlesource}{title}
      \field{annotation}{In this article, the author investigates an analogue of the Mertens sum for function fields. Let~$C$ be a nonsingular projective curve defined over a finite field $\mathbb{F}_q$ of characteristic $p>2$. The author defines the {M}öbius function $\mu_{C/\mathbb{F}_q}(D)$ of $C/\mathbb{F}_q$ for all effective divisors $D$ of $C$ to be $$ \mu_{C/\mathbb{F}_q}(D) = \begin{cases} 1, & \text{if $D=0$},\\ 0, & \text{if a prime divisor divides $D$ with order at least 2},\\ (-1)^t, & \text{if $D$ is a sum of $t$ distinct prime divisors}. \end{cases}$$ Let $c_{\mu}(N)= \sum_{\deg(D)=N} \mu_{C/\mathbb{F}_q}(D)$ and $M_{C/\mathbb{F}_q}(X)= \sum_{\deg D \le X} \mu_{C/\mathbb{F}_q}(D)$. The author justifies these new definitions by showing that the Dirichlet series $Z_\mu(u)$ associated with $\mu_{C/\mathbb{F}_q}(D)$ for a divisor~$D$ is closely related to the the zeta function $Z_{C/\mathbb{F}_q}(u)$ associated with the curve~$C$: $$ \frac{1}{Z_{C/\mathbb{F}_q}(u)}=Z_\mu(u)= \sum_{D \geq 0}\frac{\mu_{C/\mathbb{F}_q}(D)}{\mathcal{N}D^s} = \sum_{N=0}^\infty c_\mu(N)u^N, $$ where $u=q^{-s}$ and $\mathcal{N}D$ is the absolute norm of~$D$, so that $M_{C/\mathbb{F}_q}(X)=\sum_{N=0}^\infty c_\mu(N)$. By working with the coefficients, the author shows that $M_{C/\mathbb{F}_q}(X) \ll X^{2g-1}q^{X/2}$.\par Let $\gamma_1 = \sqrt{q}e^{i \theta_1} , \ldots, \gamma_{2g} = \sqrt{q}e^{i \theta_{2g}}$ be the inverse zeros of $Z_{C/\mathbb{F}_q}(u)$, where~$g$ is the genus of~$C$. The curve~$C$ is said to satisfy LI if the set $\{\theta_j \colon 0 \le \theta_j \le \pi \} \cup \{\pi\}$ is linearly independent over~$\mathbb{Q}$. The author shows that LI implies $$ D(C/\mathbb{F}_q) = \limsup_{X \to \infty} \frac{|M_{C/\mathbb{F}_q}(X)|}{q^{X/2}}=\sum_{j=1}^{2g} \left| \frac{\gamma_j}{Z_{C/\mathbb{F}_q}'(\gamma_j^{-1})}\frac{\gamma_j}{\gamma_j-1} \right|<\infty, $$ while $D(C/\mathbb{F}_q)=\infty$ if there is a zero with multiple order. The author also shows that a family of hyperelliptic curves satisfy LI and computes the geometric average of $D(C/\mathbb{F}_q)$ over this family using an equidistribution theorem due to Deligne.\par This article cites~\cites{MR2474313,MR2811555,MR3189983,MR0006202,1897.Mertens,MR2078705,MR783538,MR1329368}.}
      \field{issn}{0065-1036}
      \field{journaltitle}{Acta Arith.}
      \field{number}{4}
      \field{title}{The summatory function of the {M}öbius function in function fields}
      \field{volume}{179}
      \field{year}{2017}
      \true{nocite}
      \field{pages}{375\bibrangedash 395}
      \range{pages}{21}
      \verb{doi}
      \verb 10.4064/aa8590-1-2017
      \endverb
      \verb{mrnumber}
      \verb 3684399
      \endverb
      \verb{urlraw}
      \verb https://doi.org/10.4064/aa8590-1-2017
      \endverb
      \verb{url}
      \verb https://doi.org/10.4064/aa8590-1-2017
      \endverb
    \endentry
    \entry{MR3692499}{article}{}
      \name{author}{1}{}{%
        {{hash=16aef6e7e9117011b23f5df0b0f8e46b}{%
           family={Hough},
           familyi={H\bibinitperiod},
           given={Patrick},
           giveni={P\bibinitperiod}}}%
      }
      \strng{namehash}{16aef6e7e9117011b23f5df0b0f8e46b}
      \strng{fullhash}{16aef6e7e9117011b23f5df0b0f8e46b}
      \strng{bibnamehash}{16aef6e7e9117011b23f5df0b0f8e46b}
      \strng{authorbibnamehash}{16aef6e7e9117011b23f5df0b0f8e46b}
      \strng{authornamehash}{16aef6e7e9117011b23f5df0b0f8e46b}
      \strng{authorfullhash}{16aef6e7e9117011b23f5df0b0f8e46b}
      \field{sortinit}{2}
      \field{sortinithash}{ed39bb39cf854d5250e95b1c1f94f4ed}
      \field{labelnamesource}{author}
      \field{labeltitlesource}{title}
      \field{annotation}{The author establishes a stronger version of a conjecture from~\cite{MR3521338} on the biases between products of two primes, by showing that sufficiently large $x$, there are at least $\exp\big((C-\ep)\sqrt{\log x}\big)$ integers $d \leq \exp(C\sqrt{\log x})$ such that \[ \frac{\#\{pq \leq x \colon \chi_d(p) = \chi_d(q) = -1 \} } {\frac{1}{4}\# \{pq \leq x \colon (pq,d) = 1\} } \geq 1 + \frac{\log \log \log x + O(1)}{\log \log x}, \] and similarly with ``${}\ge1+\cdots$'' replaced by ``${}\le1-\cdots$'' and/or with $-1$ replaced by $1$ on the left-hand side. The author also shows that GRH implies that these oscillations are best possible. The proof uses a result of Granville and Soundararajan on extremal values of $L(1,\chi)$.\par This article cites~\cite{MR3028758,MR2202918,MR3521338}.}
      \field{issn}{2363-9555}
      \field{journaltitle}{Res. Number Theory}
      \field{title}{A lower bound for biases amongst products of two primes}
      \field{volume}{3}
      \field{year}{2017}
      \true{nocite}
      \field{pages}{Art. 19\bibrangessep 11}
      \range{pages}{-1}
      \verb{doi}
      \verb 10.1007/s40993-017-0083-9
      \endverb
      \verb{mrnumber}
      \verb 3692499
      \endverb
    \endentry
    \entry{MR3604916}{article}{}
      \name{author}{1}{}{%
        {{hash=5d7cfdf14768e19ed402154ca9e7f128}{%
           family={Meng},
           familyi={M\bibinitperiod},
           given={X.},
           giveni={X\bibinitperiod}}}%
      }
      \strng{namehash}{5d7cfdf14768e19ed402154ca9e7f128}
      \strng{fullhash}{5d7cfdf14768e19ed402154ca9e7f128}
      \strng{bibnamehash}{5d7cfdf14768e19ed402154ca9e7f128}
      \strng{authorbibnamehash}{5d7cfdf14768e19ed402154ca9e7f128}
      \strng{authornamehash}{5d7cfdf14768e19ed402154ca9e7f128}
      \strng{authorfullhash}{5d7cfdf14768e19ed402154ca9e7f128}
      \field{extraname}{1}
      \field{sortinit}{2}
      \field{sortinithash}{ed39bb39cf854d5250e95b1c1f94f4ed}
      \field{labelnamesource}{author}
      \field{labeltitlesource}{title}
      \field{annotation}{Assuming RH, this article connects the normalized error term $E^{Q_k}(x) = (Q_k(x) - x/\zeta(k))/x^{1/2k}$ for the distribution of $k$-free numbers with the sum $J_{-1}(T) = \sum _{0<\gamma \leq T} |\zeta'(\rho)|^{-2}$ over nontrivial zeros of $\zeta(s)$. The author first shows that $J_{-1}(T) \ll_{\ep} T^{1+\ep}$ holding for all $\ep>0$ is equivalent to $\int^{X}_1 E^{Q_k}(x)^2 \,\frac{dx}{x} \ll_{k}\log X$ holding for all $k \geq 2$. If in fact $J_{-1}(T)\ll _{\ep}T^{1+\ep}$ for all $\ep>0$, the author proves that for each $k \geq 2$, \[ \lim_{X\to\infty} \frac1{\log X} \int^{X} _1 E^{Q_k}(x)^2 \frac{dx}{x} = \sum_{\gamma>0}\frac{2|\zeta(\rho/k)|^2}{|\rho \zeta'(\rho)|^2} \] and $E^{Q_k}(x)$ has a limiting logarithmic distribution. The author establishes analogous results for $M(x)$. Assuming (RH still and) the assumption $J_{-1}(T)\ll T^{2-\delta}$ for any fixed $\delta>0$, he establishes a weaker version of Mertens conjecture, namely $\int_2^X (M(x)/x)^2 \,dx \ll \log X$, and proves that $M(x) \ll x^{1/2}(\log \log x)^{3/2}$ except on a set of finite logarithmic measure. He also shows that \[ \lim_{X\to\infty} \frac1{\log X} \int^{X}_1 E^M(x)^2 \,\frac{dx}{x} = \sum_{\gamma>0}\frac{2}{|\rho \zeta'(\rho)|^2} \] and $E^M(x)$ has a limiting logarithmic distribution.\par This article cites \cites{MR3261965,MR0882550,MR2078705,MR3261965,MR1502996}.}
      \field{issn}{0305-0041}
      \field{journaltitle}{Math. Proc. Cambridge Philos. Soc.}
      \field{number}{2}
      \field{title}{The distribution of {$k$}-free numbers and the derivative of the {R}iemann zeta-function}
      \field{volume}{162}
      \field{year}{2017}
      \true{nocite}
      \field{pages}{293\bibrangedash 317}
      \range{pages}{25}
      \verb{doi}
      \verb 10.1017/S0305004116000554
      \endverb
      \verb{mrnumber}
      \verb 3604916
      \endverb
      \verb{urlraw}
      \verb https://doi.org/10.1017/S0305004116000554
      \endverb
      \verb{url}
      \verb https://doi.org/10.1017/S0305004116000554
      \endverb
    \endentry
    \entry{MR3700043}{incollection}{}
      \name{author}{2}{}{%
        {{hash=454d3dedc781ec4a49112de10313b411}{%
           family={Mossinghoff},
           familyi={M\bibinitperiod},
           given={Michael\bibnamedelima J.},
           giveni={M\bibinitperiod\bibinitdelim J\bibinitperiod}}}%
        {{hash=66490333fa52c597ed64f7abf27c7710}{%
           family={Trudgian},
           familyi={T\bibinitperiod},
           given={Timothy\bibnamedelima S.},
           giveni={T\bibinitperiod\bibinitdelim S\bibinitperiod}}}%
      }
      \list{publisher}{1}{%
        {Springer, Cham}%
      }
      \strng{namehash}{901cb3f7453290119b926d8dd538a562}
      \strng{fullhash}{901cb3f7453290119b926d8dd538a562}
      \strng{bibnamehash}{901cb3f7453290119b926d8dd538a562}
      \strng{authorbibnamehash}{901cb3f7453290119b926d8dd538a562}
      \strng{authornamehash}{901cb3f7453290119b926d8dd538a562}
      \strng{authorfullhash}{901cb3f7453290119b926d8dd538a562}
      \field{extraname}{2}
      \field{sortinit}{2}
      \field{sortinithash}{ed39bb39cf854d5250e95b1c1f94f4ed}
      \field{labelnamesource}{author}
      \field{labeltitlesource}{title}
      \field{annotation}{The authors show that \begin{align*} \liminf_{x \to \infty} E^L(x) < -2.3723 \quad &\text{and} \quad \limsup_{x \to \infty} E^L(x) > 1.0028 \\ \liminf_{x \to \infty} E^{L_r}(x) < -1.0028 \quad &\text{and} \quad \limsup_{x \to \infty} E^{L_r}(x) > 2.3723. \end{align*} They further define $L_\alpha(x)=\sum_{n \leq x} \lambda(n)n^{-\alpha}$ for all $\alpha \in (0,1)$, and show that $L_\alpha(x)$ changes signs infinitely often for each $\alpha \in (0, 0.29714\ldots)$, while $L_\alpha(x)-\zeta(2\alpha)/\zeta(\alpha)$ changes signs infinitely often for each $\alpha \in (0.70285,\ldots,1)$. The theoretical and computational reasoning employed is similar to~\cite{MR3335903}. \par This article cites~\cite{MR27305,MR584557,1905.Stieltjes,MR2398787,MR783538,MR0120198,MR429789,MR0006202,MR2994374,MR0104638,MR0332685,MR0202685,MR0332684,MR3335903,MR0330069,MR654520,MR3261965}.}
      \field{booktitle}{Exploring the {R}iemann zeta function}
      \field{title}{The {L}iouville function and the {R}iemann hypothesis}
      \field{year}{2017}
      \true{nocite}
      \field{pages}{201\bibrangedash 221}
      \range{pages}{21}
      \verb{mrnumber}
      \verb 3700043
      \endverb
    \endentry
    \entry{MR3773805}{article}{}
      \name{author}{2}{}{%
        {{hash=d6e30aa748c8fb70505ca61ee684d871}{%
           family={Harper},
           familyi={H\bibinitperiod},
           given={Adam\bibnamedelima J.},
           giveni={A\bibinitperiod\bibinitdelim J\bibinitperiod}}}%
        {{hash=64363cb2afef9e33f967130adc38684b}{%
           family={Lamzouri},
           familyi={L\bibinitperiod},
           given={Youness},
           giveni={Y\bibinitperiod}}}%
      }
      \strng{namehash}{d2f0ae4433a2a9de0fd6046d871e9ead}
      \strng{fullhash}{d2f0ae4433a2a9de0fd6046d871e9ead}
      \strng{bibnamehash}{d2f0ae4433a2a9de0fd6046d871e9ead}
      \strng{authorbibnamehash}{d2f0ae4433a2a9de0fd6046d871e9ead}
      \strng{authornamehash}{d2f0ae4433a2a9de0fd6046d871e9ead}
      \strng{authorfullhash}{d2f0ae4433a2a9de0fd6046d871e9ead}
      \field{sortinit}{2}
      \field{sortinithash}{ed39bb39cf854d5250e95b1c1f94f4ed}
      \field{labelnamesource}{author}
      \field{labeltitlesource}{title}
      \field{annotation}{The authors investigate, assuming GRH and LI, the asymptotic behavior of $\delta_{q;a_1,\ldots,a_n}$ when the number of competitors~$n$ grows as a function of the modulus~$q$. They show that if $n\le \log q/(\log\log q)^4$ then $\delta_{q;a_1,\ldots,a_n} \sim \frac{1}{n!}$, resolving an unpublished conjecture by Ford and Lamzouri and strengthening results of Rubinstein and Sarnak~\cite{MR1329368} and Lamzouri~\cite{MR3063909}. They prove that this is not necessarily true for larger~$n$, as predicted by Feuerverger and Martin~\cite{MR1806291}: when $\phi(q)^{\ep} \leq n \leq \phi(q)$ one has $\liminf_{q\to\infty} n!\delta_{q;a_1,\ldots,a_n} < 1$.\par They further discuss the first~$k$ leaders in a prime number race: for each integer $1\leq k\leq n$, they define $\delta_k(q;a_1, \ldots, a_n)$ to be the logarithmic density of the set of real numbers $x\geq 2$ such that $$ \pi(x;q, a_1)>\pi(x;q, a_2) > \cdots> \pi(x;q, a_k)> \max_{k+1\leq j\leq n}\pi(x;q, a_j). $$ They show that if $2 \leq n \leq \phi(q)^{1 / 32}$ then $\delta_1(q;a_1, \ldots, a_n) \sim \frac1n$, and that the analogous result $\delta_k(q;a_1, \ldots, a_n) \sim \frac{(n-k)!}{n!}$ holds for $k(\log k)^{10}\le(\log q)/\log n$ but not for $n\ge \phi(q)^\ep$.\par In addition to using the circle method to control the average size of correlations in prime number races, the authors develop sophisticated probabilistic tools including an exchangeable pairs version of Stein's method and variants of ``normal comparison'' lemmas of Slepian.\par This article cites~\cites{MR1806291, MR3272280, MR3028758, MR1909220, MR3151605, MR2202918, MR1251926, MR1256149, MR1367576, MR0146156, MR0146157, MR0146158, MR0146159, MR0146160, MR0156826, MR0156827, MR2998146, MR3063909, 1914.Littlewood, MR4082248, MR1297543, MR1329368}, as well as an earlier draft of this annotated bibliography.}
      \field{issn}{0178-8051}
      \field{journaltitle}{Probab. Theory Related Fields}
      \field{number}{3-4}
      \field{title}{Orderings of weakly correlated random variables, and prime number races with many contestants}
      \field{volume}{170}
      \field{year}{2018}
      \true{nocite}
      \field{pages}{961\bibrangedash 1010}
      \range{pages}{50}
      \verb{doi}
      \verb 10.1007/s00440-017-0800-2
      \endverb
      \verb{mrnumber}
      \verb 3773805
      \endverb
      \verb{urlraw}
      \verb https://doi.org/10.1007/s00440-017-0800-2
      \endverb
      \verb{url}
      \verb https://doi.org/10.1007/s00440-017-0800-2
      \endverb
    \endentry
    \entry{MR3739227}{article}{}
      \name{author}{1}{}{%
        {{hash=28c386582a9310ce9c19f006559d9e7a}{%
           family={Hurst},
           familyi={H\bibinitperiod},
           given={Greg},
           giveni={G\bibinitperiod}}}%
      }
      \strng{namehash}{28c386582a9310ce9c19f006559d9e7a}
      \strng{fullhash}{28c386582a9310ce9c19f006559d9e7a}
      \strng{bibnamehash}{28c386582a9310ce9c19f006559d9e7a}
      \strng{authorbibnamehash}{28c386582a9310ce9c19f006559d9e7a}
      \strng{authornamehash}{28c386582a9310ce9c19f006559d9e7a}
      \strng{authorfullhash}{28c386582a9310ce9c19f006559d9e7a}
      \field{sortinit}{2}
      \field{sortinithash}{ed39bb39cf854d5250e95b1c1f94f4ed}
      \field{labelnamesource}{author}
      \field{labeltitlesource}{title}
      \field{annotation}{The authors extend Odlyzko and te Riele's disproof of Mertens conjecture~\cite{MR783538} by obtaining stronger bounds using more modern algorithms and computers. In particular the authors prove that \[ \liminf_{x \to \infty} E^M(x) < -1.837625 \quad \text{and} \quad \limsup_{x \to \infty} E^M(x) > 1.826054. \] The authors compute $M(x)$ for $x\le10^{16}$ and for $x=2^{54},\ldots,2^{73}$. Their algorithm operates in logspace, lowering the amount of storage required during computations. The authors note that cache misses were a significant source of increased computing time for larger values of $M(x)$; they suggest that a future algorithm could further reduce cache size by storing Möbius values in two bits, rather than bytes.\par This article cites~\cite{MR783538,MR0006202,MR0202685,MR2078705,MR2118272,MR0046485,MR3335903}.}
      \field{issn}{0025-5718}
      \field{journaltitle}{Math. Comp.}
      \field{number}{310}
      \field{title}{Computations of the Mertens function and improved bounds on the Mertens conjecture}
      \field{volume}{87}
      \field{year}{2018}
      \true{nocite}
      \field{pages}{1013\bibrangedash 1028}
      \range{pages}{16}
      \verb{doi}
      \verb 10.1090/mcom/3275
      \endverb
      \verb{mrnumber}
      \verb 3739227
      \endverb
    \endentry
    \entry{MR3803705}{article}{}
      \name{author}{1}{}{%
        {{hash=5d7cfdf14768e19ed402154ca9e7f128}{%
           family={Meng},
           familyi={M\bibinitperiod},
           given={X.},
           giveni={X\bibinitperiod}}}%
      }
      \strng{namehash}{5d7cfdf14768e19ed402154ca9e7f128}
      \strng{fullhash}{5d7cfdf14768e19ed402154ca9e7f128}
      \strng{bibnamehash}{5d7cfdf14768e19ed402154ca9e7f128}
      \strng{authorbibnamehash}{5d7cfdf14768e19ed402154ca9e7f128}
      \strng{authornamehash}{5d7cfdf14768e19ed402154ca9e7f128}
      \strng{authorfullhash}{5d7cfdf14768e19ed402154ca9e7f128}
      \field{extraname}{2}
      \field{sortinit}{2}
      \field{sortinithash}{ed39bb39cf854d5250e95b1c1f94f4ed}
      \field{labelnamesource}{author}
      \field{labeltitlesource}{title}
      \field{annotation}{Define $\pi_k(x;q,a) = \#\{n \leq x \colon n \equiv a \mod{q},\, \omega(n) = k \}$, and define $N_k(x;q,a)$ to be the \mbox{analogous} function with $\omega(n)$ replaced by $\Omega(n)$; define $\pi_k(x;q,a,b)=\pi_k(x;q,a)-\pi_k(x;q,b)$ and similarly for~$N_k(x;q,a,b)$. Assuming GRH and the simplicity of the zeros of $L(s,\chi)$, the author establishes asymptotic formulas relating $\pi_k(x;q,a,b)$ and $N_k(x;q,a,b)$ to $N_1(x;q,a,b) = \pi(x;q,a,b)$. Further assuming LI, the author establishes the existence of the densities $\delta(\pi_k(x;q,a,b))$ and $\delta(N_k(x;q,a,b))$ (including computations of several such), characterizes when these densities are less than or greater than $\frac12$, and demonstrates how the distances from such densities to $\frac12$ decrease as functions of~$k$. \par This article cites~\cites{1853.Chebyshev,MR3028758,MR1985941,MR2778652,MR2202918,MR552467,MR0146156,MR3063909,MR0068565,MR0083001,1914.Littlewood,MR1329368,MR0882550}.}
      \field{issn}{1937-0652}
      \field{journaltitle}{Algebra Number Theory}
      \field{number}{2}
      \field{title}{Chebyshev's bias for products of $k$ primes}
      \field{volume}{12}
      \field{year}{2018}
      \true{nocite}
      \field{pages}{305\bibrangedash 341}
      \range{pages}{37}
      \verb{doi}
      \verb 10.2140/ant.2018.12.305
      \endverb
      \verb{mrnumber}
      \verb 3803705
      \endverb
    \endentry
    \entry{MR3778223}{article}{}
      \name{author}{1}{}{%
        {{hash=5d7cfdf14768e19ed402154ca9e7f128}{%
           family={Meng},
           familyi={M\bibinitperiod},
           given={X.},
           giveni={X\bibinitperiod}}}%
      }
      \strng{namehash}{5d7cfdf14768e19ed402154ca9e7f128}
      \strng{fullhash}{5d7cfdf14768e19ed402154ca9e7f128}
      \strng{bibnamehash}{5d7cfdf14768e19ed402154ca9e7f128}
      \strng{authorbibnamehash}{5d7cfdf14768e19ed402154ca9e7f128}
      \strng{authornamehash}{5d7cfdf14768e19ed402154ca9e7f128}
      \strng{authorfullhash}{5d7cfdf14768e19ed402154ca9e7f128}
      \field{extraname}{3}
      \field{sortinit}{2}
      \field{sortinithash}{ed39bb39cf854d5250e95b1c1f94f4ed}
      \field{labelnamesource}{author}
      \field{labeltitlesource}{title}
      \field{annotation}{Given $k \geq 2$ and a multiset $\{a_1,\ldots,a_k\}$ of reduced residues modulo~$q$, the author determines an asymptotic formula for the number of squarefree integers up to~$x$ that can be written as $p_1p_2\cdots p_k$ with each $p_j \equiv a_j \mod q$, with an explicit second-order term of relative size $1/\log\log x$ compared to the main term. In particular, certain such multisets are biased over others, and those biases are eventually permanent. The constants in this second-order term are closely related to the constants in Mertens sum for primes in arithmetic progressions, which are heavily influenced by the least prime appearing. The author's results are given both for fixed $k$ and with $k$ growing like a constant multiple of $\log\log x$. \par This article cites~\cites{MR3521338,MR2778652,MR3803705,MR2034131}.}
      \field{issn}{0025-5793}
      \field{journaltitle}{Mathematika}
      \field{number}{1}
      \field{title}{Large bias for integers with prime factors in arithmetic progressions}
      \field{volume}{64}
      \field{year}{2018}
      \true{nocite}
      \field{pages}{237\bibrangedash 252}
      \range{pages}{16}
      \verb{mrnumber}
      \verb 3778223
      \endverb
    \endentry
    \entry{MR3871592}{article}{}
      \name{author}{1}{}{%
        {{hash=b772e92036ff60f0205b68346a3f61db}{%
           family={Schlage--Puchta},
           familyi={S\bibinithyphendelim \bibinithyphendelim P\bibinitperiod},
           given={J.-C.},
           giveni={J\bibinithyphendelim C\bibinitperiod}}}%
      }
      \strng{namehash}{b772e92036ff60f0205b68346a3f61db}
      \strng{fullhash}{b772e92036ff60f0205b68346a3f61db}
      \strng{bibnamehash}{b772e92036ff60f0205b68346a3f61db}
      \strng{authorbibnamehash}{b772e92036ff60f0205b68346a3f61db}
      \strng{authornamehash}{b772e92036ff60f0205b68346a3f61db}
      \strng{authorfullhash}{b772e92036ff60f0205b68346a3f61db}
      \field{extraname}{2}
      \field{sortinit}{2}
      \field{sortinithash}{ed39bb39cf854d5250e95b1c1f94f4ed}
      \field{labelnamesource}{author}
      \field{labeltitlesource}{title}
      \field{annotation}{Using the power-sum method and Pintz's technique of kernel functions, the author improves Pintz's result ~\cites{MR820251} on localized oscillations of $\Delta^\psi$. The main result is that if $0<\ep<\frac1e$, and $\rho_0=\sigma_0+i \gamma_0$ is a zero of $\zeta(s)$ with $\sigma_0 \geq \frac{1}{2}+\ep$ and $\gamma_0 > 8.3 ^{1/\ep}$, then for each~$T$ sufficiently large (depending explicitly on~$\gamma_0$ and~$\ep$), there exists $x \in [T, T^{1+\ep}]$ such that $|\Delta^\psi (x)| \gg {x^{\sigma_0}}/{\gamma_0^{1+\ep}}$.\par This article cites~\cites{MR820241,MR1575079, MR820251,MR749389,MR0043121}.}
      \field{issn}{0236-5294}
      \field{journaltitle}{Acta Math. Hungar.}
      \field{number}{2}
      \field{title}{Oscillations of the error term in the prime number theorem}
      \field{volume}{156}
      \field{year}{2018}
      \true{nocite}
      \field{pages}{303\bibrangedash 308}
      \range{pages}{6}
      \verb{doi}
      \verb 10.1007/s10474-018-0884-x
      \endverb
      \verb{mrnumber}
      \verb 3871592
      \endverb
      \verb{urlraw}
      \verb https://doi.org/10.1007/s10474-018-0884-x
      \endverb
      \verb{url}
      \verb https://doi.org/10.1007/s10474-018-0884-x
      \endverb
    \endentry
    \entry{MR3961320}{article}{}
      \name{author}{3}{}{%
        {{hash=bb758f27d76cefc7bd5244190e91843b}{%
           family={Ford},
           familyi={F\bibinitperiod},
           given={Kevin},
           giveni={K\bibinitperiod}}}%
        {{hash=d6e30aa748c8fb70505ca61ee684d871}{%
           family={Harper},
           familyi={H\bibinitperiod},
           given={Adam\bibnamedelima J.},
           giveni={A\bibinitperiod\bibinitdelim J\bibinitperiod}}}%
        {{hash=64363cb2afef9e33f967130adc38684b}{%
           family={Lamzouri},
           familyi={L\bibinitperiod},
           given={Youness},
           giveni={Y\bibinitperiod}}}%
      }
      \strng{namehash}{edca729ace4d98a40acfda652484ce64}
      \strng{fullhash}{edca729ace4d98a40acfda652484ce64}
      \strng{bibnamehash}{edca729ace4d98a40acfda652484ce64}
      \strng{authorbibnamehash}{edca729ace4d98a40acfda652484ce64}
      \strng{authornamehash}{edca729ace4d98a40acfda652484ce64}
      \strng{authorfullhash}{edca729ace4d98a40acfda652484ce64}
      \field{sortinit}{2}
      \field{sortinithash}{ed39bb39cf854d5250e95b1c1f94f4ed}
      \field{labelnamesource}{author}
      \field{labeltitlesource}{title}
      \field{annotation}{The authors show that prime race densities $\delta_{q;a_1,\ldots,a_r}$ can have size significantly different from $\frac1{r!}$ when~$r$ is sufficiently large in terms of~$q$. More precisely, they show that there exists $\eta>0$ and $a_1,\ldots,a_r\mod q$ such that \[ \delta_{q;a_1,\ldots,a_r} \le \exp\bigg( - \frac{\eta\min\{n,\phi(q)^{1/50}\}}{\log q} \bigg) \frac1{r!}, \] and an analogous lower bound for large values without the negative sign on the right-hand side. They derive this result from a similar result concerning the density of real numbers~$x$ such that $\pi(x;q,a_1) > \cdots > \pi(x;q,a_k)$ are the~$k$ largest of the~$r$ prime counting functions while simultaneously $\pi(x;q,a_{k+1}) < \cdots < \pi(x;q,a_{2k})$ are the~$k$ smallest. The authors exploit the fact that certain $E(x;q,a_j)$ are known to have large correlations of size $\Omega_-(1/\log q)$; they also develop a comparison theorem to multivariate normal distributions with a relative rather than absolute error.\par This article cites~\cite{MR1985941,MR0506080,MR1329368,MR1806291,MR2075622,MR2202918,MR2998146,MR3028758,MR3063909,MR3272280,MR3773805}, as well as an earlier draft of this annotated bibliography.}
      \field{issn}{0025-5831}
      \field{journaltitle}{Math. Ann.}
      \field{number}{1-2}
      \field{title}{Extreme biases in prime number races with many contestants}
      \field{volume}{374}
      \field{year}{2019}
      \true{nocite}
      \field{pages}{517\bibrangedash 551}
      \range{pages}{35}
      \verb{doi}
      \verb 10.1007/s00208-019-01810-x
      \endverb
      \verb{mrnumber}
      \verb 3961320
      \endverb
      \verb{urlraw}
      \verb https://doi.org/10.1007/s00208-019-01810-x
      \endverb
      \verb{url}
      \verb https://doi.org/10.1007/s00208-019-01810-x
      \endverb
    \endentry
    \entry{MR3996180}{article}{}
      \name{author}{3}{}{%
        {{hash=48017d687b31d05c591036588e04dd66}{%
           family={Humphries},
           familyi={H\bibinitperiod},
           given={Peter},
           giveni={P\bibinitperiod}}}%
        {{hash=d346d7afb6639c668a44a4e3ce8a8b98}{%
           family={Shekatkar},
           familyi={S\bibinitperiod},
           given={Snehal\bibnamedelima M.},
           giveni={S\bibinitperiod\bibinitdelim M\bibinitperiod}}}%
        {{hash=1a3f18d66237c8283c6ec23cd9a8c12d}{%
           family={Wong},
           familyi={W\bibinitperiod},
           given={Tian\bibnamedelima An},
           giveni={T\bibinitperiod\bibinitdelim A\bibinitperiod}}}%
      }
      \strng{namehash}{0212718ca31c4c1791497f3ee79e82ac}
      \strng{fullhash}{0212718ca31c4c1791497f3ee79e82ac}
      \strng{bibnamehash}{0212718ca31c4c1791497f3ee79e82ac}
      \strng{authorbibnamehash}{0212718ca31c4c1791497f3ee79e82ac}
      \strng{authornamehash}{0212718ca31c4c1791497f3ee79e82ac}
      \strng{authorfullhash}{0212718ca31c4c1791497f3ee79e82ac}
      \field{sortinit}{2}
      \field{sortinithash}{ed39bb39cf854d5250e95b1c1f94f4ed}
      \field{labelnamesource}{author}
      \field{labeltitlesource}{title}
      \field{annotation}{Given a set~$S$ of integers, define $\Omega_S(n) = \#\{p\mid n\colon p\in S\}$. The authors study $\Sigma_S(x) = \sum_{n\le x} (-1)^{\Omega_S(n)}$ when~$S$ is a union of residue classes modulo~$q$. If~$S$ is the complete set of reduced residue classes\mod q then $\Sigma_S(x)$ is similar to $L(x)$; if~$q$ is prime and~$S$ is the set of quadratic nonresidues\mod q, then $\Sigma_S(x)$ is similar to $\sum_{n\le x} (\frac nq)$ except that $(-1)^{\Omega_S(n)}$ takes nonzero values at multiples of~$q$ as well.\par The authors show that the general behaviour of $\Sigma_S(x)$ is more predictable than these special cases: if~$S$ is the union of~$r$ reduced residue classes\mod q with $r\notin \{ \phi(q), \frac{\phi(q)}2 \}$, then $\Sigma_S(x) \sim b_0 {x}/{(\log x)^{2-2r/\varphi(q)}}$ for some explicit constant~$b_0$ that has the same sign as $\frac{\phi(q)}2-r$. If $r=\frac{\phi(q)}2$, then still $\Sigma_S(x)=o(x)$ provided that~$S$ is not a set of the form $\{n\colon \chi(n)=-1\}$ for some quadratic Dirichlet character $\chi\mod q$. Finally, suppose that $S = \{n\colon \chi(n)\ne-1\}$ for some quadratic $\chi\mod q$. Assuming GRH, LI, and $\sum_{0<\gamma\leq T} {|L'(\rho,\chi_q )|^-2} \ll T^{\theta}$ for some $1<\theta<3-\sqrt{3}$, the authors show that the logarithmic density~$\delta$ of $\{x>0\colon \Sigma_S(x)\ge 0 \}$ exists and satisfies $1/2\leq \delta<1$.\par This article cites~\cites{1919.Polya, MR3261965, MR3335903, MR2398787, MR0104638, MR2994374, MR0006202, MR3604916, MR3778223}.}
      \field{issn}{1246-7405}
      \field{journaltitle}{J. Théor. Nombres Bordeaux}
      \field{number}{1}
      \field{title}{Biases in prime factorizations and {L}iouville functions for arithmetic progressions}
      \field{volume}{31}
      \field{year}{2019}
      \true{nocite}
      \field{pages}{1\bibrangedash 25}
      \range{pages}{25}
      \verb{mrnumber}
      \verb 3996180
      \endverb
      \verb{urlraw}
      \verb http://jtnb.cedram.org/item?id=JTNB_2019__31_1_1_0
      \endverb
      \verb{url}
      \verb http://jtnb.cedram.org/item?id=JTNB_2019__31_1_1_0
      \endverb
    \endentry
    \entry{MR4012359}{article}{}
      \name{author}{2}{}{%
        {{hash=64363cb2afef9e33f967130adc38684b}{%
           family={Lamzouri},
           familyi={L\bibinitperiod},
           given={Youness},
           giveni={Y\bibinitperiod}}}%
        {{hash=5509e2d8aa4558a4fd44867b61f0c363}{%
           family={Martin},
           familyi={M\bibinitperiod},
           given={Bruno},
           giveni={B\bibinitperiod}}}%
      }
      \strng{namehash}{52cf16fd2d29d30832e4c5dcecfe7ec4}
      \strng{fullhash}{52cf16fd2d29d30832e4c5dcecfe7ec4}
      \strng{bibnamehash}{52cf16fd2d29d30832e4c5dcecfe7ec4}
      \strng{authorbibnamehash}{52cf16fd2d29d30832e4c5dcecfe7ec4}
      \strng{authornamehash}{52cf16fd2d29d30832e4c5dcecfe7ec4}
      \strng{authorfullhash}{52cf16fd2d29d30832e4c5dcecfe7ec4}
      \field{sortinit}{2}
      \field{sortinithash}{ed39bb39cf854d5250e95b1c1f94f4ed}
      \field{labelnamesource}{author}
      \field{labeltitlesource}{title}
      \field{annotation}{Let~$A$ and~$B$ be disjoint sets of reduced residue classes modulo~$q$ with $\#A=\#B$. If~$A$ and~$B$ contain the same number of quadratic residues, the authors prove that~HC implies $\pi(x;q,A,B)=\Omega_{\pm} (\sqrt{x}/\log x)$. Under GRH and LI, the authors show that $\delta_{q;A,B}=\frac12$ if~$A$ and~$B$ contain the same number of quadratic residues while $\delta_{q;A,B}<\frac12$ if~$A$ contains more quadratic residues than~$B$.\par Using these results, the authors study the oscillation and the limiting logarithmic distribution of the function $S_k(x)=\sum_{p \leq x} (-1)^{s_k(p)}$, where $s_k(p)$ is the sum of the last $k$ binary digits of $p$ (note $S_2(x)=\pi(x;4,3,1)-1$ recovers Chebyshev's bias \cite{1853.Chebyshev}). The first result above implies that $S_k(x)$ changes sign infinitely often for all $k\ge2$, while the second result above implies that $\delta(\{x\colon S_k(x)>0\}) = \frac12$ for $k\ge4$. (In both cases the assumptions refer to Dirichlet $L$-functions modulo~$2^k$.) The authors further compute that $\delta(\{x\colon S_3(x)>0\})=\delta(\{x\colon \pi(x;8,3)+\pi(x;8,5)>\pi(x;8,1)+\pi(x;8,7)\})\approx 0.9822$ using the approach of Rubinstein--Sarnak~\cite{MR1329368}.\par This article cites \cite{MR1985941,MR3151605,MR1555148,MR1251926,MR0233782,MR0146157,MR1329368}.}
      \field{issn}{0208-6573}
      \field{journaltitle}{Funct. Approx. Comment. Math.}
      \field{number}{1}
      \field{title}{On the race between primes with an odd versus an even sum of the last {$k$} binary digits}
      \field{volume}{61}
      \field{year}{2019}
      \true{nocite}
      \field{pages}{7\bibrangedash 25}
      \range{pages}{19}
      \verb{doi}
      \verb 10.7169/facm/1687
      \endverb
      \verb{mrnumber}
      \verb 4012359
      \endverb
      \verb{urlraw}
      \verb https://doi.org/10.7169/facm/1687
      \endverb
      \verb{url}
      \verb https://doi.org/10.7169/facm/1687
      \endverb
    \endentry
    \entry{MR3993767}{article}{}
      \name{author}{3}{}{%
        {{hash=d6c5b18d6dc761cc36407ac5d5441ece}{%
           family={Lichtman},
           familyi={L\bibinitperiod},
           given={J.\bibnamedelimi D.},
           giveni={J\bibinitperiod\bibinitdelim D\bibinitperiod}}}%
        {{hash=feda7d3cba2a0ad4e772307fca655937}{%
           family={Martin},
           familyi={M\bibinitperiod},
           given={G.},
           giveni={G\bibinitperiod}}}%
        {{hash=a658c262e68541245018e2e8dd4064a7}{%
           family={Pomerance},
           familyi={P\bibinitperiod},
           given={C.},
           giveni={C\bibinitperiod}}}%
      }
      \strng{namehash}{c132fc8c2248d719686bb8b2ea4c74aa}
      \strng{fullhash}{c132fc8c2248d719686bb8b2ea4c74aa}
      \strng{bibnamehash}{c132fc8c2248d719686bb8b2ea4c74aa}
      \strng{authorbibnamehash}{c132fc8c2248d719686bb8b2ea4c74aa}
      \strng{authornamehash}{c132fc8c2248d719686bb8b2ea4c74aa}
      \strng{authorfullhash}{c132fc8c2248d719686bb8b2ea4c74aa}
      \field{sortinit}{2}
      \field{sortinithash}{ed39bb39cf854d5250e95b1c1f94f4ed}
      \field{labelnamesource}{author}
      \field{labeltitlesource}{title}
      \field{annotation}{Assuming GRH and LI, the authors show that the set of primes~$q$ for which $\pi(q)>\li(q)$ has a well-defined logarithmic relative density in the set of all primes, whose value equals $\delta(\pi,\li)$ (approximately $2.6\cdot10^{-7}$). The same is true of the set of primes~$q$ such that $e^{C_0} \log q > \prod_{p\le q} (1-1/p)^{-1}$, and the set of primes~$q$ such that $1/\log q > \sum_{p\ge q} 1/p\log p$. The methods apply to discrete subsets of~$\R$ that are well-distributed in almost all short intervals (such as the integers themselves). \par This article cites~\cite{MR3261965,MR1752093,MR4120786,MR3028758,MR3455269,1914.Littlewood,MR1329368}.}
      \field{journaltitle}{Proc. Amer. Math. Soc.}
      \field{number}{9}
      \field{title}{Primes in prime number races}
      \field{volume}{147}
      \field{year}{2019}
      \true{nocite}
      \field{pages}{3743\bibrangedash 3757}
      \range{pages}{15}
    \endentry
    \entry{MR3925757}{article}{}
      \name{author}{2}{}{%
        {{hash=7876864d71873e82aad2393b8cf18fe4}{%
           family={Platt},
           familyi={P\bibinitperiod},
           given={Dave},
           giveni={D\bibinitperiod}}}%
        {{hash=af8f95636cfb8ee32db62432b50fcea2}{%
           family={Trudgian},
           familyi={T\bibinitperiod},
           given={Tim},
           giveni={T\bibinitperiod}}}%
      }
      \strng{namehash}{d0b290fda16e446719a77447f29d0e95}
      \strng{fullhash}{d0b290fda16e446719a77447f29d0e95}
      \strng{bibnamehash}{d0b290fda16e446719a77447f29d0e95}
      \strng{authorbibnamehash}{d0b290fda16e446719a77447f29d0e95}
      \strng{authornamehash}{d0b290fda16e446719a77447f29d0e95}
      \strng{authorfullhash}{d0b290fda16e446719a77447f29d0e95}
      \field{extraname}{2}
      \field{sortinit}{2}
      \field{sortinithash}{ed39bb39cf854d5250e95b1c1f94f4ed}
      \field{labelnamesource}{author}
      \field{labeltitlesource}{title}
      \field{annotation}{The authors show that GRH for $L(s,\chi_{-4})$ is equivalent, for any $0<\alpha<20.40442$, to \[ \lim_{x\to\infty} \sum_p \chi_{-4}(p) e^{-(x/n)^\alpha} = -\infty; \] they show also that the given method cannot be extended to $\alpha<20.40443$.\par This article cites~\cite{MR1985941,MR974088,MR1555148,MR1544293A,MR0202686,MR1329368}.}
      \field{issn}{1793-0421}
      \field{journaltitle}{Int. J. Number Theory}
      \field{number}{3}
      \field{title}{Fujii's development on {C}hebyshev's conjecture}
      \field{volume}{15}
      \field{year}{2019}
      \true{nocite}
      \field{pages}{639\bibrangedash 644}
      \range{pages}{6}
      \verb{doi}
      \verb 10.1142/S1793042119500337
      \endverb
      \verb{mrnumber}
      \verb 3925757
      \endverb
      \verb{urlraw}
      \verb https://doi.org/10.1142/S1793042119500337
      \endverb
      \verb{url}
      \verb https://doi.org/10.1142/S1793042119500337
      \endverb
    \endentry
    \entry{MR4120786}{article}{}
      \name{author}{1}{}{%
        {{hash=c9734df31e44a7c0ee5c006cb13a5e86}{%
           family={Devin},
           familyi={D\bibinitperiod},
           given={Lucile},
           giveni={L\bibinitperiod}}}%
      }
      \strng{namehash}{c9734df31e44a7c0ee5c006cb13a5e86}
      \strng{fullhash}{c9734df31e44a7c0ee5c006cb13a5e86}
      \strng{bibnamehash}{c9734df31e44a7c0ee5c006cb13a5e86}
      \strng{authorbibnamehash}{c9734df31e44a7c0ee5c006cb13a5e86}
      \strng{authornamehash}{c9734df31e44a7c0ee5c006cb13a5e86}
      \strng{authorfullhash}{c9734df31e44a7c0ee5c006cb13a5e86}
      \field{extraname}{1}
      \field{sortinit}{2}
      \field{sortinithash}{ed39bb39cf854d5250e95b1c1f94f4ed}
      \field{labelnamesource}{author}
      \field{labeltitlesource}{title}
      \field{annotation}{In this article, the author extends the study of prime number races to counting functions with explicit formulas associated to general $L$-functions from the Selberg class, including Dirichlet $L$-functions and Hasse--Weil $L$-functions. She unconditionally proves the existence of a limiting logarithmic distribution for these counting functions, and establishes properties of the limiting distribution (such as absolute continuity) under much weaker conditions than GRH and LI.\par This article cites~\cites{MR1329368,MR2202918,MR3261965, MR1909220}.}
      \field{issn}{0305-0041}
      \field{journaltitle}{Math. Proc. Cambridge Philos. Soc.}
      \field{number}{1}
      \field{title}{Chebyshev's bias for analytic {L}-functions}
      \field{volume}{169}
      \field{year}{2020}
      \true{nocite}
      \field{pages}{103\bibrangedash 140}
      \range{pages}{38}
      \verb{doi}
      \verb 10.1017/s0305004119000100
      \endverb
      \verb{mrnumber}
      \verb 4120786
      \endverb
      \verb{urlraw}
      \verb https://doi.org/10.1017/s0305004119000100
      \endverb
      \verb{url}
      \verb https://doi.org/10.1017/s0305004119000100
      \endverb
    \endentry
    \entry{MR4176773}{article}{}
      \name{author}{1}{}{%
        {{hash=c9734df31e44a7c0ee5c006cb13a5e86}{%
           family={Devin},
           familyi={D\bibinitperiod},
           given={Lucile},
           giveni={L\bibinitperiod}}}%
      }
      \strng{namehash}{c9734df31e44a7c0ee5c006cb13a5e86}
      \strng{fullhash}{c9734df31e44a7c0ee5c006cb13a5e86}
      \strng{bibnamehash}{c9734df31e44a7c0ee5c006cb13a5e86}
      \strng{authorbibnamehash}{c9734df31e44a7c0ee5c006cb13a5e86}
      \strng{authornamehash}{c9734df31e44a7c0ee5c006cb13a5e86}
      \strng{authorfullhash}{c9734df31e44a7c0ee5c006cb13a5e86}
      \field{extraname}{2}
      \field{sortinit}{2}
      \field{sortinithash}{ed39bb39cf854d5250e95b1c1f94f4ed}
      \field{labelnamesource}{author}
      \field{labeltitlesource}{title}
      \field{annotation}{In this article, the author extends her earlier the weakening of the LI assumption to the setting of prime number races with many contestants. Assuming GRH but only weakened versions of the LI assumption, she shows the existence of logarithmic densities associated to $r$-way prime number races. Her methods apply equally well to the number field and function field settings (in the latter case, GRH is no longer a hypothesis). An erratum appeared in the same journal in 2021 (MR\href{http://www.ams.org/mathscinet-getitem?mr=4352666}{\nolinkurl{4352666}}).\par This article cites~\cites{MR1329368,MR3993767,MR4082248,MR1957110,MR2202918,MR1985941,MR3803705,MR4120786,MR2600420,MR2811555,MR3581815,MR2474313,MR3261965}.}
      \field{issn}{0008-4395}
      \field{journaltitle}{Canad. Math. Bull.}
      \field{number}{4}
      \field{title}{Limiting properties of the distribution of primes in an arbitrarily large number of residue classes}
      \field{volume}{63}
      \field{year}{2020}
      \true{nocite}
      \field{pages}{837\bibrangedash 849}
      \range{pages}{13}
      \verb{doi}
      \verb 10.4153/s0008439520000089
      \endverb
      \verb{mrnumber}
      \verb 4176773
      \endverb
      \verb{urlraw}
      \verb https://doi.org/10.4153/s0008439520000089
      \endverb
      \verb{url}
      \verb https://doi.org/10.4153/s0008439520000089
      \endverb
    \endentry
    \entry{MR4043824}{article}{}
      \name{author}{2}{}{%
        {{hash=f5fefda26be806a5e5f1affbb4e80fe9}{%
           family={Lemke\bibnamedelima Oliver},
           familyi={L\bibinitperiod\bibinitdelim O\bibinitperiod},
           given={Robert\bibnamedelima J.},
           giveni={R\bibinitperiod\bibinitdelim J\bibinitperiod}}}%
        {{hash=db8c7134ab46023e3103bbf8edcf2ca5}{%
           family={Soundararajan},
           familyi={S\bibinitperiod},
           given={Kannan},
           giveni={K\bibinitperiod}}}%
      }
      \strng{namehash}{b4ea39c0bfa1f7e30e1685d948198229}
      \strng{fullhash}{b4ea39c0bfa1f7e30e1685d948198229}
      \strng{bibnamehash}{b4ea39c0bfa1f7e30e1685d948198229}
      \strng{authorbibnamehash}{b4ea39c0bfa1f7e30e1685d948198229}
      \strng{authornamehash}{b4ea39c0bfa1f7e30e1685d948198229}
      \strng{authorfullhash}{b4ea39c0bfa1f7e30e1685d948198229}
      \field{extraname}{2}
      \field{sortinit}{2}
      \field{sortinithash}{ed39bb39cf854d5250e95b1c1f94f4ed}
      \field{labelnamesource}{author}
      \field{labeltitlesource}{title}
      \field{annotation}{The authors continue their study of biases in the patterns of consecutive primes\mod q. Their conjectures in~\cite{MR3624386} contained a tertiary main term with complicated constants $c_2(q,{\bf a})$ whose distribution was not easily understood. In this article, they connect these constants with both the discrete Fourier transform of Dedekind sums and the error term $R(x) = \sum_{n\le x} \phi(n) - \frac{3}{\pi^2}x^2$. They show that $q^{-1}c_2(q, (a,b))$ has a continuous limiting distribution as $q\to\infty$ (connected to that of the related Dedekind sums), and that $x^{-1}R(x)$ has a limiting distribution with doubly exponential decay. \par This article cites~\cites{MR3624386, MR1297543}.}
      \field{issn}{0305-0041}
      \field{journaltitle}{Math. Proc. Cambridge Philos. Soc.}
      \field{number}{1}
      \field{title}{The distribution of consecutive prime biases and sums of sawtooth random variables}
      \field{volume}{168}
      \field{year}{2020}
      \true{nocite}
      \field{pages}{149\bibrangedash 169}
      \range{pages}{21}
      \verb{doi}
      \verb 10.1017/s0305004118000592
      \endverb
      \verb{mrnumber}
      \verb 4043824
      \endverb
      \verb{urlraw}
      \verb https://doi-org.ezproxy.library.ubc.ca/10.1017/s0305004118000592
      \endverb
      \verb{url}
      \verb https://doi-org.ezproxy.library.ubc.ca/10.1017/s0305004118000592
      \endverb
    \endentry
    \entry{MR4082248}{article}{}
      \name{author}{2}{}{%
        {{hash=c06cf923cb7824f2455f1750eec3b292}{%
           family={Martin},
           familyi={M\bibinitperiod},
           given={Greg},
           giveni={G\bibinitperiod}}}%
        {{hash=3d2daa810172929af247bf36116d8aaa}{%
           family={Ng},
           familyi={N\bibinitperiod},
           given={Nathan},
           giveni={N\bibinitperiod}}}%
      }
      \strng{namehash}{11060c47f44aeca22f1c306f4f674b92}
      \strng{fullhash}{11060c47f44aeca22f1c306f4f674b92}
      \strng{bibnamehash}{11060c47f44aeca22f1c306f4f674b92}
      \strng{authorbibnamehash}{11060c47f44aeca22f1c306f4f674b92}
      \strng{authornamehash}{11060c47f44aeca22f1c306f4f674b92}
      \strng{authorfullhash}{11060c47f44aeca22f1c306f4f674b92}
      \field{sortinit}{2}
      \field{sortinithash}{ed39bb39cf854d5250e95b1c1f94f4ed}
      \field{labelnamesource}{author}
      \field{labeltitlesource}{title}
      \field{annotation}{The authors investigate, under GRH, how to weaken the LI assumption in earlier results. Say that a zero $\frac12+i\gamma$ of some Dirichlet $L$-function\mod q is ``self-sufficient'' if it cannot be written as a rational linear combination of other such zeros (so that LI is true if and only if all such zeros are self-sufficient). The authors show that a certain finite number of self-sufficient zeros is enough for the logarithmic limiting distributions of $E^\pi(x;q,a_1,\ldots,a_r)$ to exist. If the sum of reciprocals of ordinates of self-sufficient zeros is sufficiently large in terms of~$q$, the authors show that all logarithmic densities $\delta_{q;a_1,\ldots,a_r}$ exist and are positive; if that reciprocal sum diverges, then the distribution of $E^\pi(x;q,a_1,\ldots,a_r)$ assigns mass to every open set in $\R^r$.\par This article cites~\cites{MR3261965,MR4120786,MR0332684,MR3028758,MR3151605,MR2075622,MR2202918,MR0332685,MR2994374,MR0006202,MR1105608,MR1251926,MR0176967,MR1329368,MR2753165,MR0289452,1938.Wintner}.}
      \field{issn}{0002-9947}
      \field{journaltitle}{Trans. Amer. Math. Soc.}
      \field{number}{5}
      \field{title}{Inclusive prime number races}
      \field{volume}{373}
      \field{year}{2020}
      \true{nocite}
      \field{pages}{3561\bibrangedash 3607}
      \range{pages}{47}
      \verb{doi}
      \verb 10.1090/tran/7996
      \endverb
      \verb{mrnumber}
      \verb 4082248
      \endverb
      \verb{urlraw}
      \verb https://doi.org/10.1090/tran/7996
      \endverb
      \verb{url}
      \verb https://doi.org/10.1090/tran/7996
      \endverb
    \endentry
    \entry{MR4077187}{article}{}
      \name{author}{1}{}{%
        {{hash=0e36a3a1f7c6c0678ce294e777ceb2f9}{%
           family={Meng},
           familyi={M\bibinitperiod},
           given={Xianchang},
           giveni={X\bibinitperiod}}}%
      }
      \strng{namehash}{0e36a3a1f7c6c0678ce294e777ceb2f9}
      \strng{fullhash}{0e36a3a1f7c6c0678ce294e777ceb2f9}
      \strng{bibnamehash}{0e36a3a1f7c6c0678ce294e777ceb2f9}
      \strng{authorbibnamehash}{0e36a3a1f7c6c0678ce294e777ceb2f9}
      \strng{authornamehash}{0e36a3a1f7c6c0678ce294e777ceb2f9}
      \strng{authorfullhash}{0e36a3a1f7c6c0678ce294e777ceb2f9}
      \field{extraname}{4}
      \field{sortinit}{2}
      \field{sortinithash}{ed39bb39cf854d5250e95b1c1f94f4ed}
      \field{labelnamesource}{author}
      \field{labeltitlesource}{title}
      \field{annotation}{Assuming GRH and LI, the author proves that both $\sum_{n\le x} \omega(n)\chi_{-4}(n) < 0$ and $\sum_{n\le x} \Omega(n)\chi_{-4}(n) > 0$ on sets of logarithmic density~$1$. This result gives a conditional confirmation of a conjecture of G.~Martin.\par This article cites~\cites{MR3261965,1853.Chebyshev, MR3604916, MR3803705, MR3778223, MR1329368}.}
      \field{issn}{0033-5606}
      \field{journaltitle}{Q. J. Math.}
      \field{number}{1}
      \field{title}{Number of prime factors over arithmetic progressions}
      \field{volume}{71}
      \field{year}{2020}
      \true{nocite}
      \field{pages}{97\bibrangedash 121}
      \range{pages}{25}
      \verb{doi}
      \verb 10.1093/qmathj/haz040
      \endverb
      \verb{mrnumber}
      \verb 4077187
      \endverb
      \verb{urlraw}
      \verb https://doi.org/10.1093/qmathj/haz040
      \endverb
      \verb{url}
      \verb https://doi.org/10.1093/qmathj/haz040
      \endverb
    \endentry
    \entry{MR4132130}{incollection}{}
      \name{author}{2}{}{%
        {{hash=454d3dedc781ec4a49112de10313b411}{%
           family={Mossinghoff},
           familyi={M\bibinitperiod},
           given={Michael\bibnamedelima J.},
           giveni={M\bibinitperiod\bibinitdelim J\bibinitperiod}}}%
        {{hash=66490333fa52c597ed64f7abf27c7710}{%
           family={Trudgian},
           familyi={T\bibinitperiod},
           given={Timothy\bibnamedelima S.},
           giveni={T\bibinitperiod\bibinitdelim S\bibinitperiod}}}%
      }
      \list{publisher}{1}{%
        {Amer. Math. Soc., [Providence], RI}%
      }
      \strng{namehash}{901cb3f7453290119b926d8dd538a562}
      \strng{fullhash}{901cb3f7453290119b926d8dd538a562}
      \strng{bibnamehash}{901cb3f7453290119b926d8dd538a562}
      \strng{authorbibnamehash}{901cb3f7453290119b926d8dd538a562}
      \strng{authornamehash}{901cb3f7453290119b926d8dd538a562}
      \strng{authorfullhash}{901cb3f7453290119b926d8dd538a562}
      \field{extraname}{3}
      \field{sortinit}{2}
      \field{sortinithash}{ed39bb39cf854d5250e95b1c1f94f4ed}
      \field{labelnamesource}{author}
      \field{labeltitlesource}{title}
      \field{annotation}{Let $H(x)=\sum_{n\leq x}(-1)^{\omega(n)}$. The authors prove that $H(x)>1.7\sqrt{x}$ and $H(x)<-1.7\sqrt{x}$ infinitely often, complementing existing results on oscillations for $L(x)$. The proof is very similar to that of the follow-up paper~\cite{MR4194167}.\par This article cites~\cites{MR654520,MR0330069, MR3335903,MR2398787,MR0332684,MR0202685,MR0332685,MR0104638,MR2994374,MR3739227,MR0006202,MR429789,MR0120198,1897.Mertens,MR3062002,MR3700043,MR2078705,MR783538,1919.Polya,MR1329368,MR0882550}.}
      \field{booktitle}{75 years of mathematics of computation}
      \field{isbn}{978-1-4704-5163-9}
      \field{series}{Contemp. Math.}
      \field{title}{A tale of two omegas}
      \field{volume}{754}
      \field{year}{2020}
      \true{nocite}
      \field{pages}{343\bibrangedash 364}
      \range{pages}{22}
      \verb{mrnumber}
      \verb 4132130
      \endverb
    \endentry
    \entry{MR4249594}{book}{}
      \name{author}{1}{}{%
        {{hash=81a915ac987d508458c06ee0891cc3c2}{%
           family={Plymen},
           familyi={P\bibinitperiod},
           given={Roger},
           giveni={R\bibinitperiod}}}%
      }
      \list{publisher}{1}{%
        {American Mathematical Society, Providence, RI}%
      }
      \strng{namehash}{81a915ac987d508458c06ee0891cc3c2}
      \strng{fullhash}{81a915ac987d508458c06ee0891cc3c2}
      \strng{bibnamehash}{81a915ac987d508458c06ee0891cc3c2}
      \strng{authorbibnamehash}{81a915ac987d508458c06ee0891cc3c2}
      \strng{authornamehash}{81a915ac987d508458c06ee0891cc3c2}
      \strng{authorfullhash}{81a915ac987d508458c06ee0891cc3c2}
      \field{sortinit}{2}
      \field{sortinithash}{ed39bb39cf854d5250e95b1c1f94f4ed}
      \field{labelnamesource}{author}
      \field{labeltitlesource}{title}
      \field{annotation}{Chapter~6 provides a comprehensive overview of the methods used to study the oscillations of $\pi(x)-\li(x)$. Chapter 7 outlines known results on the logarithmic density of the race between $\pi(x)$ and $\li(x)$, as well as upper bounds for the Skewes number. \par This book cites~\cites{MR1752093, MR2652902,1932.Ingham, MR0202686, MR3993767, 1914.Littlewood, MR866118, MR1329368, MR2684372, MR1573970, MR0067145,MR2813366,MR0882550}.}
      \field{isbn}{978-1-4704-6257-4}
      \field{series}{Student Mathematical Library}
      \field{title}{The {G}reat {P}rime {N}umber {R}ace}
      \field{volume}{92}
      \field{year}{2020}
      \true{nocite}
      \field{pages}{138}
      \range{pages}{1}
      \verb{mrnumber}
      \verb 4249594
      \endverb
    \endentry
    \entry{MR4265352}{article}{}
      \name{author}{1}{}{%
        {{hash=49275b2a1e8fa5af0a0e7ef3941655a7}{%
           family={Bailleul},
           familyi={B\bibinitperiod},
           given={Alexandre},
           giveni={A\bibinitperiod}}}%
      }
      \strng{namehash}{49275b2a1e8fa5af0a0e7ef3941655a7}
      \strng{fullhash}{49275b2a1e8fa5af0a0e7ef3941655a7}
      \strng{bibnamehash}{49275b2a1e8fa5af0a0e7ef3941655a7}
      \strng{authorbibnamehash}{49275b2a1e8fa5af0a0e7ef3941655a7}
      \strng{authornamehash}{49275b2a1e8fa5af0a0e7ef3941655a7}
      \strng{authorfullhash}{49275b2a1e8fa5af0a0e7ef3941655a7}
      \field{extraname}{1}
      \field{sortinit}{2}
      \field{sortinithash}{ed39bb39cf854d5250e95b1c1f94f4ed}
      \field{labelnamesource}{author}
      \field{labeltitlesource}{title}
      \field{annotation}{This article studies a generalization of Chebyshev's bias to the Chebotarev density theorem, which was first investigated by Ng~\cite{2000.Ng}. Namely, if $L/K$ is a Galois extension of number fields and $C$ is a conjugacy class of $\operatorname{Gal}(L/K)$, then $\pi_C(x)$ denotes the number of unramified prime ideals $\mathfrak{p}$ of $L$ with norm ${N(\mathfrak{p}) \leq x}$ and whose Artin symbol $(\frac{\mathfrak{p}}{L/K})$ equals $C$. For any two conjugacy classes $C_1$ and $C_2$ of $\operatorname{Gal}(L/K)$, a ``Chebotarev race'' measures the logarithmic density of values $x$ such that ${\pi_{C_1}(x)}/{\#C_1} > {\pi_{C_2}(x)}/{\#C_2}$. Fiorilli and Jouve~\cite{arXiv.FiorilliJouve} have shown that some Chebotarev races are extremely biased (in that the logarithmic density is close to 1) and the author continues this study in a new context.\par Using class field theory constructions and results on root numbers of quaternion extensions, the author analyzes Chebyshev's bias in certain families of number field extensions with Galois groups of $2$-power order, namely dihedral or generalized quaternion. Conditional on GRH and variants of LI, he performs an analysis in two different aspects: the ``horizontal aspect'' (fixed Galois group) and the ``vertical aspect'' (high-degree towers of $2$-power order). Each direction reveals an interesting connection between the central zeros of Artin $L$-functions and the bias of the Chebotarev race; this relationship is a key novelty of the article.\par This article cites~\cites{MR2474313,MR2811555,1853.Chebyshev,MR4120786,MR3803705,arXiv.FiorilliJouve,MR4400872,2000.Ng,MR3028758,MR3961320,MR3773805,MR1251926,MR1344249,MR0146156,MR0162771,MR2998146,MR3063909,MR4082248,MR1329368,MR0004255}.}
      \field{issn}{1937-0652}
      \field{journaltitle}{Algebra Number Theory}
      \field{number}{4}
      \field{title}{Chebyshev's bias in dihedral and generalized quaternion {G}alois groups}
      \field{volume}{15}
      \field{year}{2021}
      \true{nocite}
      \field{pages}{999\bibrangedash 1041}
      \range{pages}{43}
      \verb{doi}
      \verb 10.2140/ant.2021.15.999
      \endverb
      \verb{mrnumber}
      \verb 4265352
      \endverb
      \verb{urlraw}
      \verb https://doi.org/10.2140/ant.2021.15.999
      \endverb
      \verb{url}
      \verb https://doi.org/10.2140/ant.2021.15.999
      \endverb
    \endentry
    \entry{arXiv.Devin}{online}{}
      \name{author}{1}{}{%
        {{hash=c9734df31e44a7c0ee5c006cb13a5e86}{%
           family={Devin},
           familyi={D\bibinitperiod},
           given={Lucile},
           giveni={L\bibinitperiod}}}%
      }
      \strng{namehash}{c9734df31e44a7c0ee5c006cb13a5e86}
      \strng{fullhash}{c9734df31e44a7c0ee5c006cb13a5e86}
      \strng{bibnamehash}{c9734df31e44a7c0ee5c006cb13a5e86}
      \strng{authorbibnamehash}{c9734df31e44a7c0ee5c006cb13a5e86}
      \strng{authornamehash}{c9734df31e44a7c0ee5c006cb13a5e86}
      \strng{authorfullhash}{c9734df31e44a7c0ee5c006cb13a5e86}
      \field{extraname}{3}
      \field{sortinit}{2}
      \field{sortinithash}{ed39bb39cf854d5250e95b1c1f94f4ed}
      \field{labelnamesource}{author}
      \field{labeltitlesource}{title}
      \field{annotation}{Every $p\equiv1\mod4$ can be written uniquely as $p=a^2+4b^2$ with positive integers $a,b$. The author formulates two conjectures concerning counting functions related to the distribution of these numbers~$a$ and $2b$. First, she conjectures that the function $\sum_{p\le x,\, p\equiv1\mod4} \mathop{\rm sign}(a-2b)$ is negative for a set of logarithmic density strictly between~$\frac12$ and~$1$, and thus there is a bias towards the even square being larger than the odd square in such representations. Second, she conjectures that $\sum_{p\le x,\, p\equiv1\mod4} \chi_{-4}(a)$ is eventually always positive, so that there is a complete bias towards the positive odd number in such representations being $1\mod4$ rather than $3\mod4$. Both conjectures rely on her extension of existing results to cover counting functions that are governed by a sum of infinitely many $L$-functions (such as Hecke $L$-functions).\par This article cites~\cites{MR3261965,MR3581815,MR4120786,MR3264673,MR3272280,MR1957110,MR1329368,2007.Sarnak,MR0004255}, as well as an earlier draft of this annotated bibliography.}
      \field{title}{Discrepancies in the distribution of Gaussian primes}
      \field{year}{2021}
      \true{nocite}
      \verb{urlraw}
      \verb https://arxiv.org/abs/2105.02492
      \endverb
      \verb{url}
      \verb https://arxiv.org/abs/2105.02492
      \endverb
    \endentry
    \entry{MR4316675}{article}{}
      \name{author}{2}{}{%
        {{hash=c9734df31e44a7c0ee5c006cb13a5e86}{%
           family={Devin},
           familyi={D\bibinitperiod},
           given={Lucile},
           giveni={L\bibinitperiod}}}%
        {{hash=0e36a3a1f7c6c0678ce294e777ceb2f9}{%
           family={Meng},
           familyi={M\bibinitperiod},
           given={Xianchang},
           giveni={X\bibinitperiod}}}%
      }
      \strng{namehash}{452055ee754af6a220856a0aadedfe82}
      \strng{fullhash}{452055ee754af6a220856a0aadedfe82}
      \strng{bibnamehash}{452055ee754af6a220856a0aadedfe82}
      \strng{authorbibnamehash}{452055ee754af6a220856a0aadedfe82}
      \strng{authornamehash}{452055ee754af6a220856a0aadedfe82}
      \strng{authorfullhash}{452055ee754af6a220856a0aadedfe82}
      \field{sortinit}{2}
      \field{sortinithash}{ed39bb39cf854d5250e95b1c1f94f4ed}
      \field{labelnamesource}{author}
      \field{labeltitlesource}{title}
      \field{annotation}{The authors study the counting function of polynomials in $\mathbb F_q[t]$ in sets of invertible residue classes modulo a fixed polynomial that have $k$ irreducible factors (these can be counted either with or without multiplicity). They prove asymptotic formulas for the normalized differences of two such counting functions in terms of the zeros of the relevant $L$-functions; these asymptotic formulas are valid for $k$ almost as large as $\sqrt{\log X}$ where $X$ bounds the degree of the polynomials being counted. In the important special case where the sets of residue classes are the quadratic residues and nonresidues, the formulas simplify and the sign of the bias analyzed: assuming LI, the bias is towards quadratic residues when we count irreducible factors without multiplicity, while the bias depends on the parity of~$k$ when we count them with multiplicity. Finally, the authors use the fact that some $L$-functions vanish at the critical point $q^{-1/2}$ to show that there exist completely biased and unbiased races, contrary to what is expected in the number field setting.\par This article cites~\cites{MR3261965,MR3581815, MR2474313, MR3272280, MR1985941, MR3028758, MR2778652, MR3803705, MR3778223, MR4082248, MR2034131, MR1329368,MR4176773}.}
      \field{issn}{0001-8708}
      \field{journaltitle}{Adv. Math.}
      \field{title}{Chebyshev's bias for products of irreducible polynomials}
      \field{volume}{392}
      \field{year}{2021}
      \true{nocite}
      \field{pages}{Paper No. 108040\bibrangessep 45}
      \range{pages}{-1}
      \verb{doi}
      \verb 10.1016/j.aim.2021.108040
      \endverb
      \verb{mrnumber}
      \verb 4316675
      \endverb
      \verb{urlraw}
      \verb https://doi.org/10.1016/j.aim.2021.108040
      \endverb
      \verb{url}
      \verb https://doi.org/10.1016/j.aim.2021.108040
      \endverb
    \endentry
    \entry{MR4194167}{article}{}
      \name{author}{3}{}{%
        {{hash=454d3dedc781ec4a49112de10313b411}{%
           family={Mossinghoff},
           familyi={M\bibinitperiod},
           given={Michael\bibnamedelima J.},
           giveni={M\bibinitperiod\bibinitdelim J\bibinitperiod}}}%
        {{hash=ca1fa3eb66152a77dcfaa9c5daf01ef9}{%
           family={Oliveira\bibnamedelimb e\bibnamedelima Silva},
           familyi={O\bibinitperiod\bibinitdelim e\bibinitperiod\bibinitdelim S\bibinitperiod},
           given={Tomás},
           giveni={T\bibinitperiod}}}%
        {{hash=66490333fa52c597ed64f7abf27c7710}{%
           family={Trudgian},
           familyi={T\bibinitperiod},
           given={Timothy\bibnamedelima S.},
           giveni={T\bibinitperiod\bibinitdelim S\bibinitperiod}}}%
      }
      \strng{namehash}{7efc60c309fdd97a9edaf0af2c43c3d9}
      \strng{fullhash}{7efc60c309fdd97a9edaf0af2c43c3d9}
      \strng{bibnamehash}{7efc60c309fdd97a9edaf0af2c43c3d9}
      \strng{authorbibnamehash}{7efc60c309fdd97a9edaf0af2c43c3d9}
      \strng{authornamehash}{7efc60c309fdd97a9edaf0af2c43c3d9}
      \strng{authorfullhash}{7efc60c309fdd97a9edaf0af2c43c3d9}
      \field{sortinit}{2}
      \field{sortinithash}{ed39bb39cf854d5250e95b1c1f94f4ed}
      \field{labelnamesource}{author}
      \field{labeltitlesource}{title}
      \field{annotation}{For each $k \geq 2$, set $$ C_k=\min\biggl\{ \limsup_{x \to \infty} \frac{\Delta^{Q_k}(x)}{x^{1/2k}}, \biggl| \liminf_{x \to \infty} \frac{\Delta^{Q_k}(x)}{x^{1/2k}} \biggr| \biggr\}. $$ The authors show that $C_k>3$ for $2 \leq k \leq 5$ and $C_k>2$ for $6\leq k \leq 10$, and that $C_k>0.74969$ for all sufficiently large $k$, improving previously known lower bounds on $C_k$ significantly. The proof is based on a variant of Ingham's method developed by Anderson and Stark~\cites{MR654520}. To achieve that, the authors establish weak linear dependence relations among a subset of ordinates of $\zeta$ via the LLL algorithm. The massive computation required approximately~$5$ core-years.\par This article cites~\cites{MR654520, MR3335903,MR0006202,MR429789,MR3604916,MR3062002,MR3700043,MR4132130,MR2078705,MR0199161,MR0882550,MR1502996,MR0724708}.}
      \field{issn}{0025-5718,1088-6842}
      \field{journaltitle}{Math. Comp.}
      \field{number}{328}
      \field{title}{The distribution of {$k$}-free numbers}
      \field{volume}{90}
      \field{year}{2021}
      \true{nocite}
      \field{pages}{907\bibrangedash 929}
      \range{pages}{23}
      \verb{doi}
      \verb 10.1090/mcom/3581
      \endverb
      \verb{mrnumber}
      \verb 4194167
      \endverb
      \verb{urlraw}
      \verb https://doi.org/10.1090/mcom/3581
      \endverb
      \verb{url}
      \verb https://doi.org/10.1090/mcom/3581
      \endverb
    \endentry
    \entry{web.Shchebetov}{online}{}
      \name{author}{1}{}{%
        {{hash=a61d2f8619a2d8b3611da9d61d4a659a}{%
           family={Shchebetov},
           familyi={S\bibinitperiod},
           given={A.},
           giveni={A\bibinitperiod}}}%
      }
      \strng{namehash}{a61d2f8619a2d8b3611da9d61d4a659a}
      \strng{fullhash}{a61d2f8619a2d8b3611da9d61d4a659a}
      \strng{bibnamehash}{a61d2f8619a2d8b3611da9d61d4a659a}
      \strng{authorbibnamehash}{a61d2f8619a2d8b3611da9d61d4a659a}
      \strng{authornamehash}{a61d2f8619a2d8b3611da9d61d4a659a}
      \strng{authorfullhash}{a61d2f8619a2d8b3611da9d61d4a659a}
      \field{sortinit}{2}
      \field{sortinithash}{ed39bb39cf854d5250e95b1c1f94f4ed}
      \field{labelnamesource}{author}
      \field{labeltitlesource}{title}
      \field{annotation}{This website contains the downloadable application ``Chebyshev's bias visualizer'', which permits graphical exploration of prime number races with many parameters customizable by the user. The author has explored several races for larger values of~$x$ than previously and, in particular, discovered the first region where $\pi(x;12,5,1)<0$ (starting at $25,726,067,172,577$) and the first region where $\pi(x;12,7,1)<0$ (starting at $27,489,101,529,529$).}
      \field{title}{Chebyshev's bias visualizer}
      \field{year}{2021}
      \true{nocite}
      \verb{urlraw}
      \verb http://math101.guru/en/downloads-2/repository/
      \endverb
      \verb{url}
      \verb http://math101.guru/en/downloads-2/repository/
      \endverb
    \endentry
    \entry{MR4369787}{article}{}
      \name{author}{1}{}{%
        {{hash=53e2b48286c0b37330d03fd888d20e26}{%
           family={Aymone},
           familyi={A\bibinitperiod},
           given={Marco},
           giveni={M\bibinitperiod}}}%
      }
      \strng{namehash}{53e2b48286c0b37330d03fd888d20e26}
      \strng{fullhash}{53e2b48286c0b37330d03fd888d20e26}
      \strng{bibnamehash}{53e2b48286c0b37330d03fd888d20e26}
      \strng{authorbibnamehash}{53e2b48286c0b37330d03fd888d20e26}
      \strng{authornamehash}{53e2b48286c0b37330d03fd888d20e26}
      \strng{authorfullhash}{53e2b48286c0b37330d03fd888d20e26}
      \field{sortinit}{2}
      \field{sortinithash}{ed39bb39cf854d5250e95b1c1f94f4ed}
      \field{labelnamesource}{author}
      \field{labeltitlesource}{title}
      \field{annotation}{The author observes that GRH for $L(s, \chi_{-4})$ would imply that $\sum_{p \leq x} \chi_{-4}(p)p^{-\sigma}$ would be eventually negative when $\sigma \to 1/2^+$. He investigates the connection between the sign changes of the weighted prime number races associated to a real nonprincipal Dirichlet character~$\chi$ and the zero-free region of the corresponding $L$-function $L(s, \chi)$, proving that there exists $0\leq\sigma<1$ such that $\sum_{p\leq x}\chi(p)p^{-\sigma}$ has only finitely many sign changes if and only if there exists $\ep>0$ such that $L(s,\chi)\neq 0$ in the half plane $\sigma>1-\ep$.\par This article cites~\cites{MR0882550, MR3773805, MR2202918}.}
      \field{issn}{1793-0421}
      \field{journaltitle}{Int. J. Number Theory}
      \field{number}{1}
      \field{title}{A note on prime number races and zero free regions for {$L$} functions}
      \field{volume}{18}
      \field{year}{2022}
      \true{nocite}
      \field{pages}{1\bibrangedash 8}
      \range{pages}{8}
      \verb{doi}
      \verb 10.1142/S1793042122500014
      \endverb
      \verb{mrnumber}
      \verb 4369787
      \endverb
      \verb{urlraw}
      \verb https://doi.org/10.1142/S1793042122500014
      \endverb
      \verb{url}
      \verb https://doi.org/10.1142/S1793042122500014
      \endverb
    \endentry
    \entry{MR4447414}{article}{}
      \name{author}{1}{}{%
        {{hash=49275b2a1e8fa5af0a0e7ef3941655a7}{%
           family={Bailleul},
           familyi={B\bibinitperiod},
           given={Alexandre},
           giveni={A\bibinitperiod}}}%
      }
      \strng{namehash}{49275b2a1e8fa5af0a0e7ef3941655a7}
      \strng{fullhash}{49275b2a1e8fa5af0a0e7ef3941655a7}
      \strng{bibnamehash}{49275b2a1e8fa5af0a0e7ef3941655a7}
      \strng{authorbibnamehash}{49275b2a1e8fa5af0a0e7ef3941655a7}
      \strng{authornamehash}{49275b2a1e8fa5af0a0e7ef3941655a7}
      \strng{authorfullhash}{49275b2a1e8fa5af0a0e7ef3941655a7}
      \field{extraname}{2}
      \field{sortinit}{2}
      \field{sortinithash}{ed39bb39cf854d5250e95b1c1f94f4ed}
      \field{labelnamesource}{author}
      \field{labeltitlesource}{title}
      \field{annotation}{The author proves general versions of Kronecker--Weyl theorems, both in the discrete and continuous settings and with no linear independence assumption, where the sets on which the equidistribution is guaranteed are explicitly constructed. The article consists of many applications of these explicit theorems, including sufficient conditions for lower bounds for the lower densities in certain Chebotarev races. The applications are given in a more general setting for random variables, where the densities are bounded by probabilities of certain explicitly given events. As a concrete application, in the last part of the article, the author studies densities for prime divisor races in geometric Galois extensions of function fields of one variable over finite fields, including the cases for which the LI hypothesis fails to hold.\par This article cites~\cites{MR3261965, MR2474313, MR3581815, MR2811555, MR4120786, MR4176773, MR4316675, arXiv.FiorilliJouve, MR1957110, MR4082248, MR4082248, 2000.Ng, MR1329368}.}
      \field{issn}{2522-0160}
      \field{journaltitle}{Res. Number Theory}
      \field{number}{3}
      \field{title}{Explicit {K}ronecker--{W}eyl theorems and applications to prime number races}
      \field{volume}{8}
      \field{year}{2022}
      \true{nocite}
      \field{pages}{Paper No. 43\bibrangessep 34}
      \range{pages}{-1}
      \verb{doi}
      \verb 10.1007/s40993-022-00349-2
      \endverb
      \verb{mrnumber}
      \verb 4447414
      \endverb
      \verb{urlraw}
      \verb https://doi.org/10.1007/s40993-022-00349-2
      \endverb
      \verb{url}
      \verb https://doi.org/10.1007/s40993-022-00349-2
      \endverb
    \endentry
    \entry{MR4400872}{article}{}
      \name{author}{2}{}{%
        {{hash=c511fbe6d97a5da2025e6199f40ab277}{%
           family={Fiorilli},
           familyi={F\bibinitperiod},
           given={Daniel},
           giveni={D\bibinitperiod}}}%
        {{hash=8e65f0b511fecc96916507c9bdad4520}{%
           family={Jouve},
           familyi={J\bibinitperiod},
           given={Florent},
           giveni={F\bibinitperiod}}}%
      }
      \strng{namehash}{927d699f5bb8bac56bfe890ba7e7af7f}
      \strng{fullhash}{927d699f5bb8bac56bfe890ba7e7af7f}
      \strng{bibnamehash}{927d699f5bb8bac56bfe890ba7e7af7f}
      \strng{authorbibnamehash}{927d699f5bb8bac56bfe890ba7e7af7f}
      \strng{authornamehash}{927d699f5bb8bac56bfe890ba7e7af7f}
      \strng{authorfullhash}{927d699f5bb8bac56bfe890ba7e7af7f}
      \field{extraname}{1}
      \field{sortinit}{2}
      \field{sortinithash}{ed39bb39cf854d5250e95b1c1f94f4ed}
      \field{labelnamesource}{author}
      \field{labeltitlesource}{title}
      \field{annotation}{If $L/K$ is a Galois extension of number fields and~$C$ is a conjugacy class in $\operatorname{Gal}(L/K)$, then $\pi(x; C, L/K)$ denotes the number of unramified prime ideals $\mathfrak{p}$ of $L$ with norm ${N(\mathfrak{p}) \leq x}$ whose Frobenius conjugacy class $\text{Frob}_{\mathfrak{p}}$ equals~$C$. For two conjugacy classes $C_1, \, C_2$, let $\delta_{L/K; C_1, C_2}$ be the logarithmic density of the set $\bigl\{x\geq 1\colon \frac1{\#C_1}\pi(x; C_1, L/K) > \frac1{\#C_2}\pi(x; C_2, L/K) \bigr\}$. The authors unconditionally show that there are infinitely many Galois extensions $L/K$ and conjugacy classes $C_1, \, C_2$ such that $\delta_{L/K; C_1, C_2}=1$.\par This article cites~\cites{1853.Chebyshev, MR1344249, 2000.Ng,arXiv.FiorilliJouve,MR1329368}.}
      \field{issn}{2429-7100}
      \field{journaltitle}{J. Éc. polytech. Math.}
      \field{title}{Unconditional {C}hebyshev biases in number fields}
      \field{volume}{9}
      \field{year}{2022}
      \true{nocite}
      \field{pages}{671\bibrangedash 679}
      \range{pages}{9}
      \verb{doi}
      \verb 10.5802/jep.19
      \endverb
      \verb{mrnumber}
      \verb 4400872
      \endverb
      \verb{urlraw}
      \verb https://doi.org/10.5802/jep.19
      \endverb
      \verb{url}
      \verb https://doi.org/10.5802/jep.19
      \endverb
    \endentry
    \entry{MR4379402}{article}{}
      \name{author}{2}{}{%
        {{hash=ab91abc046ab17798c39d00cb11edfba}{%
           family={Lin},
           familyi={L\bibinitperiod},
           given={Jiawei},
           giveni={J\bibinitperiod}}}%
        {{hash=c06cf923cb7824f2455f1750eec3b292}{%
           family={Martin},
           familyi={M\bibinitperiod},
           given={Greg},
           giveni={G\bibinitperiod}}}%
      }
      \strng{namehash}{2a2204ec026337a0144e11ff56499eee}
      \strng{fullhash}{2a2204ec026337a0144e11ff56499eee}
      \strng{bibnamehash}{2a2204ec026337a0144e11ff56499eee}
      \strng{authorbibnamehash}{2a2204ec026337a0144e11ff56499eee}
      \strng{authornamehash}{2a2204ec026337a0144e11ff56499eee}
      \strng{authorfullhash}{2a2204ec026337a0144e11ff56499eee}
      \field{sortinit}{2}
      \field{sortinithash}{ed39bb39cf854d5250e95b1c1f94f4ed}
      \field{labelnamesource}{author}
      \field{labeltitlesource}{title}
      \field{annotation}{Assuming GRH and LI, the authors establish an asymptotic formula for $\delta_{q;a_1,a_2,a_3}$ in the case when $a_1^2\equiv a_2^2\equiv a_3^2\mod q$, with a power savings in $q$ in the error term; the main term, which is the arctangent of a specified algebraic function of the quantities $b(\chi) = \sum_{L(\rho,\chi)=0} 1/|\rho|^2$, arises as an orthant probability for a multivariate normal random variable. The congruence hypothesis allows for an atypical normalization of the $E^\pi(x;q,a_j)$ for which any given $\chi\mod q$ appears in at most one of the three such expressions, allowing one to model the error terms by independent random variables. The authors propose that to minimize their error terms, asymptotic formulas for $\delta_{q;a_1,\ldots,a_r}$ always be phrased in terms of orthant probabilities.\par This article cites~\cites{MR3028758,MR3773805,MR2998146,MR3063909,MR1956261,MR1329368}.}
      \field{issn}{0008-414X}
      \field{journaltitle}{Canad. J. Math.}
      \field{number}{1}
      \field{title}{Densities in certain three-way prime number races}
      \field{volume}{74}
      \field{year}{2022}
      \true{nocite}
      \field{pages}{232\bibrangedash 265}
      \range{pages}{34}
      \verb{doi}
      \verb 10.4153/S0008414X20000747
      \endverb
      \verb{mrnumber}
      \verb 4379402
      \endverb
      \verb{urlraw}
      \verb https://doi.org/10.4153/S0008414X20000747
      \endverb
      \verb{url}
      \verb https://doi.org/10.4153/S0008414X20000747
      \endverb
    \endentry
    \entry{MR4449835}{article}{}
      \name{author}{1}{}{%
        {{hash=b1b4d664f6bf440515ac846f7b16d530}{%
           family={Sedrati},
           familyi={S\bibinitperiod},
           given={Youssef},
           giveni={Y\bibinitperiod}}}%
      }
      \strng{namehash}{b1b4d664f6bf440515ac846f7b16d530}
      \strng{fullhash}{b1b4d664f6bf440515ac846f7b16d530}
      \strng{bibnamehash}{b1b4d664f6bf440515ac846f7b16d530}
      \strng{authorbibnamehash}{b1b4d664f6bf440515ac846f7b16d530}
      \strng{authornamehash}{b1b4d664f6bf440515ac846f7b16d530}
      \strng{authorfullhash}{b1b4d664f6bf440515ac846f7b16d530}
      \field{sortinit}{2}
      \field{sortinithash}{ed39bb39cf854d5250e95b1c1f94f4ed}
      \field{labelnamesource}{author}
      \field{labeltitlesource}{title}
      \field{annotation}{The author extends results for classical prime number races~\cites{MR1806291,MR3063909} to the function field setting. Given a prime power~$q$ and a monic polynomial $m\in \mathbb{F}_q[T]$, the function $\pi_q(a,m, N)$ counts the number of irreducible polynomials $P\in\mathbb{F}_q[T]$ of degree~$N$ such that $P \equiv a \mod m$. Let $\delta_{m;a_1,\ldots, a_r}$ be the density of integers $X$ such that $\sum_{N=1}^X \pi_q(a_1,m,N)>\cdots > \sum_{N=1}^X \pi_q(a_r,m,N)$. Assuming LI, the author establishes an asymptotic formula for $\delta_{m;a_1,\ldots, a_r}$ as the degree of~$m$ tends to infinity. The author further shows that races with three or more competitors behave differently than two-way races in the sense that $\delta_{m;a_1,a_2} - \frac12 \ll q^{(-1/2+o(1))\deg m}$ while $|\delta_{m;a_1,\ldots, a_r} - \frac1{r!}| = \Omega(1/\deg m)$ when $r\ge3$. When $a_1, \ldots, a_r$ have bounded degree, the author provides a simple criteria for $\delta_{m;a_1,\ldots, a_r}$ to exhibit an extreme bias, and proves (still assuming LI) that when $r\ge3$, there are always unbiased races involving only quadratic residues or only quadratic residues when $\deg m$ is large enough. The author concludes by giving a few examples of races where LI is actually false.\par This article cites~\cites{MR4265352,MR2474313,MR3581815,MR2811555,MR2600420,MR4176773,MR3803705,MR1985941,MR1806291,MR3028758,MR2202918,MR3063909,1914.Littlewood,MR1329368}.}
      \field{issn}{0025-5793}
      \field{journaltitle}{Mathematika}
      \field{number}{3}
      \field{title}{Inequities in the {S}hanks-{R}enyi prime number race over function fields}
      \field{volume}{68}
      \field{year}{2022}
      \true{nocite}
      \field{pages}{840\bibrangedash 895}
      \range{pages}{56}
      \verb{doi}
      \verb 10.1112/mtk.12150
      \endverb
      \verb{mrnumber}
      \verb 4449835
      \endverb
      \verb{urlraw}
      \verb https://doi.org/10.1112/mtk.12150
      \endverb
      \verb{url}
      \verb https://doi.org/10.1112/mtk.12150
      \endverb
    \endentry
    \entry{arXiv.FiorilliJouve}{online}{}
      \name{author}{2}{}{%
        {{hash=889ecd9d3a804c5a7e9a5d8287766657}{%
           family={Fiorilli},
           familyi={F\bibinitperiod},
           given={D.},
           giveni={D\bibinitperiod}}}%
        {{hash=dfca5f0fbbbf59635f02988f421953ae}{%
           family={Jouve},
           familyi={J\bibinitperiod},
           given={F.},
           giveni={F\bibinitperiod}}}%
      }
      \strng{namehash}{bdf4ae134537f1f2a294e43fbbf433aa}
      \strng{fullhash}{bdf4ae134537f1f2a294e43fbbf433aa}
      \strng{bibnamehash}{bdf4ae134537f1f2a294e43fbbf433aa}
      \strng{authorbibnamehash}{bdf4ae134537f1f2a294e43fbbf433aa}
      \strng{authornamehash}{bdf4ae134537f1f2a294e43fbbf433aa}
      \strng{authorfullhash}{bdf4ae134537f1f2a294e43fbbf433aa}
      \field{extraname}{2}
      \field{sortinit}{2}
      \field{sortinithash}{ed39bb39cf854d5250e95b1c1f94f4ed}
      \field{labelnamesource}{author}
      \field{labeltitlesource}{title}
      \field{annotation}{For a class function $t$ on $G=\mathop{\rm Gal}(L/K)$, the authors consider the Frobenius counting function \[ \pi(x;L/K,t) = \sum_{\substack{\mathfrak{p}\subset \mathcal{O}_K \\ {\rm N}\mathfrak{p}\leq x}} t(\text{Frob}_\mathfrak{p}) \] and study the ``normalized Chebotarev error term'' \[ E(x;L/K,t) = x^{-\Theta}\log x\cdot (\pi(x;L/K,t) - \widehat{t}(1)\text{Li}(x)), \] where $\Theta$ is the supremum of real parts of zeros of certain Artin $L$-functions of $L/\mathbb{Q}$ relevant to~$t$. They show that this error admits a limiting logarithmic distribution, and they compute its mean, along with unconditional and conditional bounds on its variance. They also prove improved bounds on the error term of the Chebotarev density theorem assuming GRH and Artin's conjecture. They further consider the logarithmic densities $\delta(L/K;t)$ of the set of values of~$x$ for which $E(x;L/K,t)$ is positive, giving criteria for when the density is close to 1 or $\frac{1}{2}$, indicating a bias or lack thereof. Finally, these ideas are applied to the class function $t_{C_1,C_2} = \frac{|G|}{|C_1|}1_{C_1}-\frac{|G|}{|C_2|}1_{C_2}$ for special families of extensions ($S_n$, dihedral, radical, abelian, and Hilbert class fields of quadratic extensions), giving asymptotic bounds for the densities in each case.\par This article cites~\cites{MR3261965,MR2474313,MR2811555,MR3581815,MR4120786,MR4316675,MR3521338,MR3272280,MR3264673,MR3028758,MR2778652,MR2202918,MR1344249,MR1310738,MR0146156,MR0146157,MR0146158,MR2943671,MR3624386,1914.Littlewood,MR4082248,MR3803705,MR2631002,2000.Ng,MR1761276,MR1329368,2007.Sarnak,MR0004255,MR4265352,MR4194167} and this bibliography.}
      \field{note}{J. Inst. Math. Jussieu, to appear}
      \field{title}{Distributions of Frobenius elements in families of Galois extensions}
      \field{year}{2023}
      \true{nocite}
      \verb{urlraw}
      \verb https://arxiv.org/abs/2001.05428
      \endverb
      \verb{url}
      \verb https://arxiv.org/abs/2001.05428
      \endverb
    \endentry
    \entry{arXiv.FM}{online}{}
      \name{author}{2}{}{%
        {{hash=c511fbe6d97a5da2025e6199f40ab277}{%
           family={{Fiorilli}},
           familyi={F\bibinitperiod},
           given={Daniel},
           giveni={D\bibinitperiod}}}%
        {{hash=c06cf923cb7824f2455f1750eec3b292}{%
           family={{Martin}},
           familyi={M\bibinitperiod},
           given={Greg},
           giveni={G\bibinitperiod}}}%
      }
      \strng{namehash}{887572e6be3fbbeb636b6dec60af9d05}
      \strng{fullhash}{887572e6be3fbbeb636b6dec60af9d05}
      \strng{bibnamehash}{887572e6be3fbbeb636b6dec60af9d05}
      \strng{authorbibnamehash}{887572e6be3fbbeb636b6dec60af9d05}
      \strng{authornamehash}{887572e6be3fbbeb636b6dec60af9d05}
      \strng{authorfullhash}{887572e6be3fbbeb636b6dec60af9d05}
      \field{extraname}{2}
      \field{sortinit}{2}
      \field{sortinithash}{ed39bb39cf854d5250e95b1c1f94f4ed}
      \field{labelnamesource}{author}
      \field{labeltitlesource}{title}
      \field{annotation}{Define $G(x;q) = \sum_{(a,q)=1} \bigl( \Delta^\theta(x;q,a)/\phi(q) \bigr)^2$. When~$q$ is fixed, results of the form $E(x,\chi) = \Omega(\log\log\log x)$ imply that $G(x;q) \not\ll x\log q$; such oscillation results were proved by Littlewood~\cite{1914.Littlewood} for certain characters, and extended by G.~Davidoff (unpublished) to all real nonprincipal characters. Hooley conjectured that $G(x;q) \ll x\log q$ as soon as~$q$ tends to infinity with~$x$. In this article, the authors disprove this conjecture. While the basic approach using Diophantine approximation and the explicit formula is the same as in earlier results, the need to have estimates that are uniform in~$q$ requires significant technical attention.\par This article cites~\cites{MR3356760, MR3028758, MR0506080, MR858395, 1914.Littlewood,MR820251, MR1329368, MR3871592}.}
      \field{note}{J. Eur. Math. Soc. (JEMS), to appear}
      \field{title}{A disproof of Hooley's conjecture}
      \field{year}{2023}
      \true{nocite}
      \verb{urlraw}
      \verb https://arxiv.org/abs/2008.05837
      \endverb
      \verb{url}
      \verb https://arxiv.org/abs/2008.05837
      \endverb
    \endentry
    \entry{MR4591762}{article}{}
      \name{author}{3}{}{%
        {{hash=c06cf923cb7824f2455f1750eec3b292}{%
           family={Martin},
           familyi={M\bibinitperiod},
           given={Greg},
           giveni={G\bibinitperiod}}}%
        {{hash=95a7b4ff9d58f1b7a8d6dc0d7ba8604f}{%
           family={Mossinghoff},
           familyi={M\bibinitperiod},
           given={Michael},
           giveni={M\bibinitperiod}}}%
        {{hash=427a3ba3d2c76d68bc8c1bfc231edce0}{%
           family={Trudgian},
           familyi={T\bibinitperiod},
           given={Timothy},
           giveni={T\bibinitperiod}}}%
      }
      \strng{namehash}{8e4514da6e580a9113eaccb5513e18fc}
      \strng{fullhash}{8e4514da6e580a9113eaccb5513e18fc}
      \strng{bibnamehash}{8e4514da6e580a9113eaccb5513e18fc}
      \strng{authorbibnamehash}{8e4514da6e580a9113eaccb5513e18fc}
      \strng{authornamehash}{8e4514da6e580a9113eaccb5513e18fc}
      \strng{authorfullhash}{8e4514da6e580a9113eaccb5513e18fc}
      \field{sortinit}{2}
      \field{sortinithash}{ed39bb39cf854d5250e95b1c1f94f4ed}
      \field{labelnamesource}{author}
      \field{labeltitlesource}{title}
      \field{annotation}{The authors investigate comparative number theoretic results for a family of arithmetic functions called ``fake~$\mu$'s'', which are multiplicative functions $\digamma(n)$ taking values in $\{-1,0,1\}$; fake~$\mu$'s include many commonly studied functions such as $\mu(n)$, $\mu^2(n)$, $(-1)^{\omega(n)}$, and $\mu_k(n)$. Generalizing a technique of Tanaka~\cite{MR0605090}, the authors show that if $f(n)$ is a fake~$\mu$ with $f(p)=-1$ and $f(p^2)=1$ for all primes~$p$, then its summatory function $F(x)$ satisfes $$ \liminf _{x \rightarrow \infty} \frac{F(x)}{\sqrt{x}}<b \quad \text { and } \quad \limsup _{x \rightarrow \infty} \frac{F(x)}{\sqrt{x}}>b, $$ where~$b$ is twice the residue at~$\frac12$ of the Dirichlet series corresponding to~$f(n)$. The authors also determine the minimum and maximum of the above constant $b$ and show that both extreme values can be achieved by particular fake~$\mu$'s.\par This article cites~\cites{MR2398787, MR1806291, MR0104638, MR2994374, MR3996180,MR3739227, MR0006202,MR0120198,MR3604916, 1897.Mertens, MR3062002,MR3700043,MR2078705,MR783538,1919.Polya,MR1329368,MR584557,MR0882550,MR0605090, MR4194167,MR4132130}.}
      \field{issn}{0002-9939}
      \field{journaltitle}{Proc. Amer. Math. Soc.}
      \field{number}{8}
      \field{title}{Fake mu's}
      \field{volume}{151}
      \field{year}{2023}
      \true{nocite}
      \field{pages}{3229\bibrangedash 3244}
      \range{pages}{16}
      \verb{doi}
      \verb 10.1090/proc/16186
      \endverb
      \verb{mrnumber}
      \verb 4591762
      \endverb
      \verb{urlraw}
      \verb https://doi.org/10.1090/proc/16186
      \endverb
      \verb{url}
      \verb https://doi.org/10.1090/proc/16186
      \endverb
    \endentry
  \enddatalist
\endrefsection
\endinput

